% Copyright (c) 2026 Mark Hunter. SPDX-License-Identifier: GPL-3.0-or-later
% !TeX program = xelatex
% Compile this master directly; latexmk -pdf also runs makeindex and reruns.
\documentclass[10pt,a4paper,oneside,headings=small]{scrreprt}
\usepackage[margin=22mm,headheight=15pt]{geometry}
\usepackage{iftex}
\ifPDFTeX
  \usepackage[T1]{fontenc}
  \usepackage[utf8]{inputenc}
  \usepackage{DejaVuSans}
  \usepackage[scaled=0.87]{DejaVuSansMono}
  \usepackage{textcomp,amssymb,accsupp}
  % Preserve the source's Unicode symbols with conventional pdfLaTeX fonts.
  \DeclareTextCommandDefault{\textleftrightarrow}{\ensuremath{\leftrightarrow}}
  \DeclareTextCommandDefault{\textneq}{\BeginAccSupp{method=hex,unicode,ActualText=2260}\ensuremath{\neq}\EndAccSupp{}}
  \DeclareTextCommandDefault{\textDOWNarrow}{\ensuremath{\blacktriangledown}}
  \DeclareUnicodeCharacter{2192}{\textrightarrow}
  \DeclareUnicodeCharacter{2193}{\textdownarrow}
  \DeclareUnicodeCharacter{2194}{\textleftrightarrow}
  \DeclareUnicodeCharacter{2260}{\textneq}
  \DeclareUnicodeCharacter{25BC}{\textDOWNarrow}
\else
\usepackage{fontspec}
\setmainfont{DejaVu Sans}
\setsansfont{DejaVu Sans}
\setmonofont{DejaVu Sans Mono}[Scale=0.87]
\fi
\renewcommand{\familydefault}{\sfdefault}
\usepackage[table]{xcolor}
\definecolor{Ink}{HTML}{19384C}
\definecolor{TableHead}{HTML}{E6EDF2}
\definecolor{Link}{HTML}{155B83}
\usepackage{array,longtable,tabularx,seqsplit,needspace,enumitem,fvextra}
\ifPDFTeX
  % seqsplit's byte scanner splits UTF-8 characters under pdfTeX. Map complete
  % text characters instead, preserving Unicode in editable code/identifier text.
  \newcommand{\HarmoniaSequenceHasUnicode}[1]{\pdfmatch{[^ -\string~]}{\detokenize{#1}}}
  \ExplSyntaxOn
  \cs_new_eq:NN \__harmonia_ascii_seqsplit:n \seqsplit
  \bool_new:N \l__harmonia_split_started_bool
  \RenewDocumentCommand{\seqsplit}{m}{
    \int_compare:nNnTF {\HarmoniaSequenceHasUnicode{#1}} = {0}
    {\__harmonia_ascii_seqsplit:n {#1}}
    {
    \group_begin:
    \bool_set_false:N \l__harmonia_split_started_bool
    \text_map_inline:nn {#1}{
      \bool_if:NT \l__harmonia_split_started_bool {\seqinsert}
      ##1
      \bool_set_true:N \l__harmonia_split_started_bool
    }
    \group_end:
    }
  }
  \ExplSyntaxOff
\fi
\usepackage{makeidx}
\makeindex
\usepackage{scrlayer-scrpage}
\clearpairofpagestyles
\ihead{Harmonia Architecture Continuum RAD}
\ohead{RAD 1.1.16}
\ifoot{Architecture Continuum · RAD 1.1.16}
\ofoot{\pagemark}
\usepackage[unicode,colorlinks,linkcolor=Link,urlcolor=Link,hyperindex=true,
  pdftitle={Harmonia Architecture Continuum RAD 1.1.16},
  pdfauthor={Harmonia},bookmarksnumbered]{hyperref}
\usepackage{bookmark}
\setlength{\parindent}{0pt}
\setlength{\parskip}{0.55em plus 0.15em minus 0.1em}
\setlength{\tabcolsep}{4pt}
\setlength{\LTpre}{0.35em}
\setlength{\LTpost}{0.5em}
\renewcommand{\arraystretch}{1.15}
\newlength{\TableWidth}
\setlist{nosep,leftmargin=1.4em,topsep=0.3em}
\setlistdepth{8}
\renewlist{itemize}{itemize}{8}
\renewlist{enumerate}{enumerate}{8}
\setlist[itemize]{label=\textbullet}
\setlist[enumerate]{label=\arabic*.}
% Canonical navigation coordinates are explicit at every depth. LaTeX never
% supplies a competing number scheme. The outline includes every entry; the
% printed contents select chapters/sections for a readable document spine.
\setcounter{tocdepth}{2}
\hypersetup{bookmarksdepth=7}
\bookmarksetup{open,openlevel=0}
\RedeclareSectionCommands[runin=false,beforeskip=1.4ex,afterskip=.6ex]{paragraph,subparagraph}
\setcounter{secnumdepth}{-2}
\newcommand{\cataloguefield}[2]{%
  \begingroup
  \renewcommand{\addcontentsline}[3]{}%
  \csname #1\endcsname{#2}%
  \endgroup
}
\DeclareTOCStyleEntries[pagenumberwidth=4em,dynnumwidth]{tocline}{chapter,section,subsection,subsubsection,paragraph,subparagraph}
\DeclareTOCStyleEntries[dynindent]{tocline}{section,subsection,subsubsection,paragraph,subparagraph}
\DeclareTOCStyleEntry[numwidth=2.6em]{tocline}{chapter}
\DeclareTOCStyleEntry[numwidth=3.5em]{tocline}{section}
\DeclareTOCStyleEntry[numwidth=4.5em]{tocline}{subsection}
\DeclareTOCStyleEntry[numwidth=5.5em]{tocline}{subsubsection}
\DeclareTOCStyleEntry[numwidth=6.5em]{tocline}{paragraph}
\DeclareTOCStyleEntry[numwidth=7.5em]{tocline}{subparagraph}
\DeclareTOCStyleEntry[pagenumberwidth=4em]{tocline}{part}
% Preserve the deepest existing model folders and element-reference levels.
% These are typesetting levels, not a publication taxonomy.
\DeclareNewSectionCommand[style=section,level=6,indent=0pt,
  beforeskip=1.4ex,afterskip=.6ex,font=\normalsize\bfseries,
  counterwithin=subparagraph,tocstyle=tocline,tocindent=8em,tocnumwidth=9em]{subsubparagraph}
\DeclareNewSectionCommand[style=section,level=7,indent=0pt,
  beforeskip=1.4ex,afterskip=.6ex,font=\normalsize\bfseries,
  counterwithin=subsubparagraph,tocstyle=tocline,tocindent=9em,tocnumwidth=10em]{subsubsubparagraph}
\addtokomafont{disposition}{\color{Ink}}
\tolerance=2000
\emergencystretch=5em
\raggedbottom
\usepackage{tikz}
\usetikzlibrary{arrows.meta}
\usepackage{xurl}
\begin{document}
\hypersetup{pageanchor=false}
\begin{titlepage}
\vspace*{24mm}
{\Huge\bfseries\color{Ink} Harmonia\par}
\vspace{8mm}
{\LARGE Enterprise Architecture\par Architecture Continuum RAD\par}
\vspace{12mm}
{\Large Release 1.1.16\par Previous release 1.1.15\par}
\vfill
Narrative publication of the canonical Architecture Continuum\par
\vspace{8mm}
The Continuum supplies architectural authority. This document selects, organises
and explains its existing content; publication confers no additional acceptance.
Individual standing and unresolved matters remain as recorded in each source.
\end{titlepage}
\pagenumbering{roman}
\hypersetup{pageanchor=true}
\tableofcontents
\chapter*{Reading the Architecture Continuum}
\addcontentsline{toc}{chapter}{Reading the Architecture Continuum}
Architecture Governance introduces the existing authority, review, evidence and
repository arrangements. Parts I--VI then follow Motivation, Strategy, Business,
Application, Technology, and Implementation \& Migration. Chapters and nested
subjects follow the canonical Continuum taxonomy. Unpopulated domains explicitly
record absence; no missing architecture is supplied for publication.

Each architectural node carries its Continuum navigation coordinate: G for
Governance and I--VI for architectural Parts; R denotes cross-cutting reference
navigation outside architectural decomposition. Coordinates continue through every
published descendant. They are locations that can change when architecture moves.
Canonical IDs remain immutable identities, and names retain their architectural
meaning. Native ArchiMate relationships preserve semantic assertions; navigation
nesting asserts no additional containment or ownership.

Strategy preserves the governed Enterprise, Business and Business Enabling
capability decompositions. Governed predecessor identities appear in the historical
reference index and link to their successors; they are not peer current subjects.
Business behaviour repeats the complete recorded ancestry needed to contextualise
represented behaviour. Each contextual occurrence identifies its canonical Strategy
Capability, Strategy coordinate and immutable ID. This is contextual navigation,
not another Capability instance or a semantic relationship. Application separates logical Topics
from scoped information structures and organises software by its existing
application responsibilities. A native DataObject is never promoted to a Topic
or Class for publication convenience. Stakeholders appear contextually against
Motivation material only through existing canonical associations. All Stakeholders
remain discoverable in R's Stakeholder Index; index placement asserts no association
and creates no Motivation decomposition subject.

The Architecture Continuum was accepted as the canonical baseline in RAD-1.1.7.
Each element and assertion retains its independent standing. Currentness supplies
no acceptance. Candidate, Proposed, historical and superseded material remains
qualified within its architectural context. Unclassified branches preserve genuine
classification gaps without manufacturing decisions. Historical definitions and
full source contributions remain available in the linked canonical records.

Use the full PDF outline for navigation down to every entry, and the Architectural
Name Index for names and their coordinates. Each narrative entry links to canonical
Markdown. Source-cited reading notes and selected existing relationship diagrams
explain represented architecture; their layout adds no semantic assertions.
The \href{structure.json}{publication correspondence} records each source, its
coordinate and preserved qualifications. The
\href{../../../continuum-manifest/published/1.1.16/continuum.html}{Continuum Manifest browser}
exposes the same numbered hierarchy and the complete qualified catalogue.
The \href{../../../../governance/architectural-axioms.md}{Architectural Axioms}
remain solely defined by their linked canonical Continuum Principles.
This document is a derived publication and supplies no independent authority.
Neither generation nor successful checks establish architectural acceptance.

\clearpage
\pagenumbering{arabic}
\chapter*{G Architecture Governance}
\label{navigation:G}
\addcontentsline{toc}{chapter}{G Architecture Governance}
The existing governance contracts establish authority, decision rights, review and evidence. These links expose that repository through the architectural navigation; they create no parallel governance records.
\section*{G.1 Governance Model}
\phantomsection
\label{navigation:G.1}
\addcontentsline{toc}{section}{G.1 Governance Model}

\subsection*{G.1.1 Architecture Authority}
\phantomsection
\label{navigation:G.1.1}
\addcontentsline{toc}{subsection}{G.1.1 Architecture Authority}

The Architecture Continuum supplies canonical architectural authority. Its canonical Principle elements solely define the Architectural Axioms. Governance contracts and accepted decisions retain their defined roles; publications and projections supply no independent architectural authority.\par
\href{../../../../governance/architecture-governance.md#authority-and-adoption}{Governed source: Architecture Authority}.\par
\subsection*{G.1.2 Decision Rights}
\phantomsection
\label{navigation:G.1.2}
\addcontentsline{toc}{subsection}{G.1.2 Decision Rights}

Mark Hunter, Solution Architect, holds final architectural decision and acceptance authority. The Architecture Review Board is the collaborative review and adjudication forum. Agents may analyse and implement authorised work; they cannot approve architecture or confirm their own architectural closure.\par
\href{../../../../governance/architecture-governance.md#architectural-decision-authority}{Governed source: Decision Rights}.\par
\subsection*{G.1.3 Roles and Responsibilities}
\phantomsection
\label{navigation:G.1.3}
\addcontentsline{toc}{subsection}{G.1.3 Roles and Responsibilities}

The Solution Architect adjudicates and confirms the incorporated result. Architecture peers and the ARB may challenge and recommend. Executing agents preserve operative instructions, actual application, evidence and uncertainty within the authorised scope.\par
\href{../../../../governance/architecture-governance.md#architectural-decision-authority}{Governed source: Roles and Responsibilities}.\par
\section*{G.2 Architecture Evidence}
\phantomsection
\label{navigation:G.2}
\addcontentsline{toc}{section}{G.2 Architecture Evidence}

\subsection*{G.2.1 Evidence before change}
\phantomsection
\label{navigation:G.2.1}
\addcontentsline{toc}{subsection}{G.2.1 Evidence before change}

The existing iteration-scoped before/ locations hold activity inputs and input provenance.\par
\href{../../../../design/evidence}{Governed source: Evidence before change}.\par
\subsection*{G.2.2 Evidence after change}
\phantomsection
\label{navigation:G.2.2}
\addcontentsline{toc}{subsection}{G.2.2 Evidence after change}

The existing iteration-scoped after/ locations hold outputs and output provenance.\par
\href{../../../../design/evidence}{Governed source: Evidence after change}.\par
\section*{G.3 Architecture Review}
\phantomsection
\label{navigation:G.3}
\addcontentsline{toc}{section}{G.3 Architecture Review}

\subsection*{G.3.1 Review records}
\phantomsection
\label{navigation:G.3.1}
\addcontentsline{toc}{subsection}{G.3.1 Review records}

Pending HCRLs preserve material proposals, attributable dispositions, conditions and actual application facts for the governed transition. The review repository supplies no acceptance merely through location.\par
\href{../../../../design/review}{Governed source: Review records}.\par
\subsection*{G.3.2 Pending architectural decisions/dispositions}
\phantomsection
\label{navigation:G.3.2}
\addcontentsline{toc}{subsection}{G.3.2 Pending architectural decisions/dispositions}

Unresolved architectural decisions and dispositions remain explicit in review records and qualified canonical elements. The accepted 1.1.10 transition defers 76 contextual classification questions; their meanings and standing remain unchanged. The completed 1.1.10 register records the accepted transition and its limits.\par
\href{../../../../design/review}{Governed source: Pending architectural decisions/dispositions}.\par
\section*{G.4 Architecture Register}
\phantomsection
\label{navigation:G.4}
\addcontentsline{toc}{section}{G.4 Architecture Register}

\subsection*{G.4.1 Accepted/completed governance records}
\phantomsection
\label{navigation:G.4.1}
\addcontentsline{toc}{subsection}{G.4.1 Accepted/completed governance records}

The existing register retains completed HCRLs after Solution Architect confirmation of the authorised change and actual incorporated result. Completed records are immutable under common governance.\par
\href{../../../../design/register}{Governed source: Accepted/completed governance records}.\par
\section*{G.5 Architecture Governance Processes}
\phantomsection
\label{navigation:G.5}
\addcontentsline{toc}{section}{G.5 Architecture Governance Processes}

\subsection*{G.5.1 Governed architecture Praxes}
\phantomsection
\label{navigation:G.5.1}
\addcontentsline{toc}{subsection}{G.5.1 Governed architecture Praxes}

\subsubsection*{G.5.1.1 Architecture Review Praxis}
\phantomsection
\label{navigation:G.5.1.1}
\addcontentsline{toc}{subsubsection}{G.5.1.1 Architecture Review Praxis}

\href{../../../../governance/architecture-review-praxis.md}{Governed source: Architecture Review Praxis}.\par
\subsubsection*{G.5.1.2 Architecture Authoring Praxis}
\phantomsection
\label{navigation:G.5.1.2}
\addcontentsline{toc}{subsubsection}{G.5.1.2 Architecture Authoring Praxis}

\href{../../../../governance/architecture-authoring-praxis.md}{Governed source: Architecture Authoring Praxis}.\par
\subsubsection*{G.5.1.3 Architectural Element Incorporation Praxis}
\phantomsection
\label{navigation:G.5.1.3}
\addcontentsline{toc}{subsubsection}{G.5.1.3 Architectural Element Incorporation Praxis}

\href{../../../../governance/continuum-incorporation-praxis.md}{Governed source: Architectural Element Incorporation Praxis}.\par
\subsubsection*{G.5.1.4 ArchiMate Round-Trip Engineering Praxis}
\phantomsection
\label{navigation:G.5.1.4}
\addcontentsline{toc}{subsubsection}{G.5.1.4 ArchiMate Round-Trip Engineering Praxis}

\href{../../../../governance/archimate-round-trip-engineering-praxis.md}{Governed source: ArchiMate Round-Trip Engineering Praxis}.\par
\section*{G.6 Architecture Principles and Compliance}
\phantomsection
\label{navigation:G.6}
\addcontentsline{toc}{section}{G.6 Architecture Principles and Compliance}

\subsection*{G.6.1 Application of canonical Architecture Principles}
\phantomsection
\label{navigation:G.6.1}
\addcontentsline{toc}{subsection}{G.6.1 Application of canonical Architecture Principles}

Canonical Continuum Principle elements are the sole normative axiom definitions; this source supplies application context and links.\par
\href{../../../../governance/architectural-axioms.md}{Governed source: Application of canonical Architecture Principles}.\par
\subsection*{G.6.2 Compliance/conformance}
\phantomsection
\label{navigation:G.6.2}
\addcontentsline{toc}{subsection}{G.6.2 Compliance/conformance}

Current validation and reproducibility mechanisms provide bounded evidence; they do not establish architectural acceptance.\par
\href{../../../../tooling/README.md}{Governed source: Compliance/conformance}.\par
\section*{G.7 Architecture Repository and Authority}
\phantomsection
\label{navigation:G.7}
\addcontentsline{toc}{section}{G.7 Architecture Repository and Authority}

\subsection*{G.7.1 Architecture Continuum}
\phantomsection
\label{navigation:G.7.1}
\addcontentsline{toc}{subsection}{G.7.1 Architecture Continuum}

Individually identified canonical Markdown elements own architectural meaning. The ordered canonical navigation describes architectural organisation; native types and relationships retain semantic meaning independently of that organisation.\par
\href{../../../../governance/continuum-representation.md}{Governed source: Architecture Continuum}.\par
\subsection*{G.7.2 Architecture publications and projections}
\phantomsection
\label{navigation:G.7.2}
\addcontentsline{toc}{subsection}{G.7.2 Architecture publications and projections}

The Continuum Manifest is the structural catalogue and the Continuum RAD is the human-readable narrative. Archi/Open Exchange is the projection, workbench and interchange representation. All derive from canonical Continuum architecture and navigation.\par
\href{../../../README.md}{Governed source: Architecture publications and projections}.\par
\subsection*{G.7.3 Evidence}
\phantomsection
\label{navigation:G.7.3}
\addcontentsline{toc}{subsection}{G.7.3 Evidence}

The existing iteration-scoped evidence repository retains required Tasks, Reports and purpose-selected supporting material under the Architecture Change Provenance contract.\par
\href{../../../../design/evidence}{Governed source: Evidence}.\par
\subsection*{G.7.4 Review}
\phantomsection
\label{navigation:G.7.4}
\addcontentsline{toc}{subsection}{G.7.4 Review}

The existing review repository retains pending architectural proposals, adjudication and result-confirmation material with its recorded standing.\par
\href{../../../../design/review}{Governed source: Review}.\par
\subsection*{G.7.5 Register}
\phantomsection
\label{navigation:G.7.5}
\addcontentsline{toc}{subsection}{G.7.5 Register}

The existing register retains accepted/completed governance records; directory placement never manufactures the required Solution Architect decision.\par
\href{../../../../design/register}{Governed source: Register}.\par

\part*{Part I — Motivation}
\label{navigation:I}
\addcontentsline{toc}{part}{Part I — Motivation}
Stakeholders are cross-cutting Motivation perspectives. Existing associations are retained; no new stakeholder-to-concern associations are inferred.\par
Read the axioms as constraints on the rest of the architecture. Harmonia owns its semantics while using underlying engines for machinery; native engine behaviour does not establish information authority. Semantic distinctions and explicit uncertainty must survive downstream representation. These points explain why the following chapters distinguish currentness, acceptance, responsibility, persistence and execution rather than treating them as interchangeable.

{\small Reading sources: \hyperref[element:54508cda80484a6caf95b901780eaf7d]{AX-04 Harmonia Owns the Semantics; Engines Provide the Machinery}; \hyperref[element:d3fc6025dc7646d1b2020b51f7a44bfa]{AX-14 Semantic Distinctions Are Preserved}; \hyperref[element:8882573df7ce43c6bc0f66da0a676139]{AX-17 Architectural Authority and Explicit Uncertainty}.}\par
\begin{center}\begin{tikzpicture}[>=Stealth,every node/.style={draw=Ink!40,rounded corners,align=left,font=\small,text width=5.2cm,inner sep=6pt}]
\node (source) at (0,0) {\hyperref[element:f4a3060eba464f4ab89b50aec4628e5b]{[RequirementGroup] Non-Functional-Requirements}};
\node (target0) at (8,1.7) {\hyperref[element:abf16578801a466b8f08d80ae292e644]{[RequirementGroup] Availability}};
\draw[->] (source.east) -- (3,0) |- (target0.west);
\node[draw=none,text width=1.8cm,font=\tiny,align=center,inner sep=0pt] at (4.15,2.08) {Composition\newline{}source-qualified / unclassified proposal};
\node (target1) at (8,0.0) {\hyperref[element:cb090e968a25409fb8ceb0fcc9a11ee7]{[RequirementGroup] Disaster Recovery Requirements}};
\draw[->] (source.east) -- (3,0) |- (target1.west);
\node[draw=none,text width=1.8cm,font=\tiny,align=center,inner sep=0pt] at (4.15,0.38) {Composition\newline{}source-qualified / unclassified proposal};
\node (target2) at (8,-1.7) {\hyperref[element:dc3f2a927e284da88734df4f7e0f9311]{[RequirementGroup] Performance Requirements}};
\draw[->] (source.east) -- (3,0) |- (target2.west);
\node[draw=none,text width=1.8cm,font=\tiny,align=center,inner sep=0pt] at (4.15,-1.32) {Composition\newline{}source-qualified / unclassified proposal};
\end{tikzpicture}\end{center}
{\small Selected outgoing Composition assertions; arrows retain source-to-target direction. Labels retain assertion standing and record state; element standing accompanies each description. \href{structure.json}{Complete relationship qualifications and provenance} remain source-defined. No omitted edge or semantic effect is inferred.}\par
\chapter*{I.1 Drivers}
\phantomsection
\label{navigation:I.1}
\addcontentsline{toc}{chapter}{I.1 Drivers}

\section*{I.1.1 24/7 Availability}
\phantomsection
\label{navigation:I.1.1}
\addcontentsline{toc}{section}{I.1.1 24/7 Availability}

\label{element:782f53e2190f4a8a8d0dc0a33cd6fbb3}
\index{24/7 Availability (I.1.1)}
{\small Driver · Historical — current applicability unresolved · historical\_reference}\par
\href{../../../../continuum/motivation/driver/782f53e2190f4a8a8d0dc0a33cd6fbb3.md}{Historical canonical definition and provenance}.\par
\subsection*{I.1.1.1 Stakeholder references — Reference navigation}
\phantomsection
\label{navigation:I.1.1.1}
\addcontentsline{toc}{subsection}{I.1.1.1 Stakeholder references — Reference navigation}

Context from existing canonical Association relationships only; these references do not allocate Stakeholders.\par
\subsubsection*{I.1.1.1.1 Healthcare Provider — Stakeholder reference [historical\_reference]}
\phantomsection
\label{navigation:I.1.1.1.1}
\addcontentsline{toc}{subsubsection}{I.1.1.1.1 Healthcare Provider — Stakeholder reference [historical\_reference]}

{\small Canonical Stakeholder: R.1.6 \hyperref[element:2018200a493d458096b72031cc2ae779]{Healthcare Provider}; Canonical ID: \texttt{\seqsplit{2018200a493d458096b72031cc2ae779}}. Historical — current applicability unresolved · historical\_reference.}\par
Contextual reference only; no additional architectural instance or semantic relationship.\par
Existing canonical associations: id-78a34270ced656e7901d367bfd405999\par
\section*{I.1.2 Audit \& PHI Privacy}
\phantomsection
\label{navigation:I.1.2}
\addcontentsline{toc}{section}{I.1.2 Audit \& PHI Privacy}

\label{element:499313a0a5d84520a3061120ae9b6dac}
\index{Audit \& PHI Privacy (I.1.2)}
{\small Driver · Historical — current applicability unresolved · historical\_reference}\par
\href{../../../../continuum/motivation/driver/499313a0a5d84520a3061120ae9b6dac.md}{Historical canonical definition and provenance}.\par
\subsection*{I.1.2.1 Stakeholder references — Reference navigation}
\phantomsection
\label{navigation:I.1.2.1}
\addcontentsline{toc}{subsection}{I.1.2.1 Stakeholder references — Reference navigation}

Context from existing canonical Association relationships only; these references do not allocate Stakeholders.\par
\subsubsection*{I.1.2.1.1 Patient — Stakeholder reference [historical\_reference]}
\phantomsection
\label{navigation:I.1.2.1.1}
\addcontentsline{toc}{subsubsection}{I.1.2.1.1 Patient — Stakeholder reference [historical\_reference]}

{\small Canonical Stakeholder: R.1.7 \hyperref[element:5f6d333e94634420b8a8bbf8b9cb6e2a]{Patient}; Canonical ID: \texttt{\seqsplit{5f6d333e94634420b8a8bbf8b9cb6e2a}}. Historical — current applicability unresolved · historical\_reference.}\par
Contextual reference only; no additional architectural instance or semantic relationship.\par
Existing canonical associations: id-daee593684c55e55bf1ce04c223cf585\par
\section*{I.1.3 Clinical Interoperability}
\phantomsection
\label{navigation:I.1.3}
\addcontentsline{toc}{section}{I.1.3 Clinical Interoperability}

\label{element:babb9f5d63d8435890c1db3579ec78d9}
\index{Clinical Interoperability (I.1.3)}
{\small Driver · Historical — current applicability unresolved · historical\_reference}\par
\href{../../../../continuum/motivation/driver/babb9f5d63d8435890c1db3579ec78d9.md}{Historical canonical definition and provenance}.\par
\subsection*{I.1.3.1 Stakeholder references — Reference navigation}
\phantomsection
\label{navigation:I.1.3.1}
\addcontentsline{toc}{subsection}{I.1.3.1 Stakeholder references — Reference navigation}

Context from existing canonical Association relationships only; these references do not allocate Stakeholders.\par
\subsubsection*{I.1.3.1.1 Health Network — Stakeholder reference [historical\_reference]}
\phantomsection
\label{navigation:I.1.3.1.1}
\addcontentsline{toc}{subsubsection}{I.1.3.1.1 Health Network — Stakeholder reference [historical\_reference]}

{\small Canonical Stakeholder: R.1.3 \hyperref[element:55e57c01dfe146a2a225228b31196862]{Health Network}; Canonical ID: \texttt{\seqsplit{55e57c01dfe146a2a225228b31196862}}. Historical — current applicability unresolved · historical\_reference.}\par
Contextual reference only; no additional architectural instance or semantic relationship.\par
Existing canonical associations: id-d5fa8c64975852f88132df579c806836\par
\subsubsection*{I.1.3.1.2 Care Team — Stakeholder reference [historical\_reference]}
\phantomsection
\label{navigation:I.1.3.1.2}
\addcontentsline{toc}{subsubsection}{I.1.3.1.2 Care Team — Stakeholder reference [historical\_reference]}

{\small Canonical Stakeholder: R.1.1 \hyperref[element:6f5a203d81e24a81b384a80781987daa]{Care Team}; Canonical ID: \texttt{\seqsplit{6f5a203d81e24a81b384a80781987daa}}. Historical — current applicability unresolved · historical\_reference.}\par
Contextual reference only; no additional architectural instance or semantic relationship.\par
Existing canonical associations: id-457a3dd70bb358ff8989bc3c9c725c70\par
\section*{I.1.4 Clinical Interoperability Across Heterogeneous Standards}
\phantomsection
\label{navigation:I.1.4}
\addcontentsline{toc}{section}{I.1.4 Clinical Interoperability Across Heterogeneous Standards}

\label{element:921d7cf84d2b4abca0deb0d300fa75ea}
\index{Clinical Interoperability Across Heterogeneous Standards (I.1.4)}
{\small Driver · Agreed · current}\par

\par\medskip\textbf{Context \& Pressure}\par


Regional healthcare environments are characterized by decades of technical layering. Acute care hospitals, outpatient clinics, pathology laboratories, and diagnostic imaging centers operate software solutions spanning multiple technological eras:
- Legacy hospital information systems communicating via HL7 v2.x pipe-and-hat (\texttt{\seqsplit{ER7}}) messaging over Minimal Lower Layer Protocol (MLLP).
- Modern clinical repositories exposing HL7 FHIR RESTful APIs (R4/R5) using JSON and XML encodings.
- Specialized imaging platforms exchanging DICOM objects.
- Proprietary file drops, CSV exports, and ad-hoc database integrations.



\par\medskip\textbf{Architectural Imperative}\par


Harmonia must bridge these heterogeneous standards without forcing disruptive, multi-million-dollar system replacements across delivery organizations. The platform must accept standards-based messages natively at its integration membrane while maintaining canonical semantic coherence internally, preventing message corruption or loss of clinical intent during translation.



\begin{itemize}


\item \textbf{Primary Stakeholders}: \href{\detokenize{../../../../../docs/markdown/01-motivation/stakeholders/enterprise-stakeholders.md}}{Regional Health Network Operator}, \href{\detokenize{../../../../../docs/markdown/01-motivation/stakeholders/enterprise-stakeholders.md}}{Healthcare Delivery Organizations}

\item \textbf{External Constraint}: \href{\detokenize{../../../../../docs/markdown/01-motivation/requirements-constraints/external-constraints.md}}{Mandated External Interoperability Contracts}

\item \textbf{Associated Goal}: \href{\detokenize{../../../../../docs/markdown/01-motivation/goals-outcomes/strategic-goals.md}}{Vendor-Independent Longitudinal Clinical Information Coherence}

\item \textbf{Governing Axioms}: \hyperref[element:eee8c3783d964d18a7f9dcea85fcd4cf]{AX-02 (Standards at Boundary)}, \hyperref[element:e99a4e2bf99f476d876394300d6b9c3d]{AX-03 (Native Standards Representations)}

\end{itemize}

{\footnotesize \href{../../../../continuum/motivation/driver/921d7cf84d2b4abca0deb0d300fa75ea.md}{Canonical record: 921d7cf84d2b4abca0deb0d300fa75ea}; full qualifications and provenance.}\par
\section*{I.1.5 Continuous Clinical Service Availability}
\phantomsection
\label{navigation:I.1.5}
\addcontentsline{toc}{section}{I.1.5 Continuous Clinical Service Availability}

\label{element:80e2d5f142da4dae8b7bb70e9ddcab06}
\index{Continuous Clinical Service Availability (I.1.5)}
{\small Driver · Agreed · current}\par

\par\medskip\textbf{Context \& Pressure}\par


Healthcare delivery is continuous, mission-critical, and time-sensitive (24/7/365). In emergency departments, intensive care units, and operating theaters, clinical decisions cannot be paused for integration system downtime, database maintenance, or middleware restarts.



A failure in the integration environment directly impacts patient care: an emergency physician cannot view critical drug allergy records; a surgeon cannot inspect pre-operative pathology results; an inpatient unit cannot receive bed transfer notifications.



\par\medskip\textbf{Architectural Imperative}\par


Harmonia must provide continuous clinical service availability. The platform must withstand hardware failures, network partitions, and downstream destination outages without dropping incoming clinical events or blocking point-of-care clinical queries.



\begin{itemize}


\item \textbf{Primary Stakeholders}: \href{\detokenize{../../../../../docs/markdown/01-motivation/stakeholders/enterprise-stakeholders.md}}{Clinicians \& Care Teams}, \href{\detokenize{../../../../../docs/markdown/01-motivation/stakeholders/enterprise-stakeholders.md}}{Platform Operations \& Integration Engineers}, \href{\detokenize{../../../../../docs/markdown/01-motivation/stakeholders/enterprise-stakeholders.md}}{Healthcare Delivery Organizations}

\item \textbf{Key Assessments}: \href{\detokenize{../../../../../docs/markdown/01-motivation/drivers-assessments/assessments.md}}{Silent Data Loss via False Acceptance}, \href{\detokenize{../../../../../docs/markdown/01-motivation/drivers-assessments/assessments.md}}{Centralised Synchronous Persistence Can Constrain Concurrency}

\item \textbf{Associated Goals}: \href{\detokenize{../../../../../docs/markdown/01-motivation/goals-outcomes/strategic-goals.md}}{Durable Acceptance \& Preservation of Clinical Events}, \href{\detokenize{../../../../../docs/markdown/01-motivation/goals-outcomes/strategic-goals.md}}{Responsive Access to Managed Information}

\item \textbf{Governing Axioms}: \hyperref[element:0d9673833eb840eab03804aa076c1b0c]{AX-10 (Failure as Normal Condition)}, \hyperref[element:8a1e995f16954484a78ae1c5c114d546]{AX-11 (Responsive \& Resilient Access)}, \hyperref[element:0412a2084f0145babc4f022c71fea65f]{AX-05 (State Separation)}

\end{itemize}

{\footnotesize \href{../../../../continuum/motivation/driver/80e2d5f142da4dae8b7bb70e9ddcab06.md}{Canonical record: 80e2d5f142da4dae8b7bb70e9ddcab06}; full qualifications and provenance.}\par
\section*{I.1.6 Coordinated Operational Activity and Entity State}
\phantomsection
\label{navigation:I.1.6}
\addcontentsline{toc}{section}{I.1.6 Coordinated Operational Activity and Entity State}

\label{element:b3fb294d95024e4382b2c251c389fa66}
\index{Coordinated Operational Activity and Entity State (I.1.6)}
{\small Driver · Agreed · current}\par

\par\medskip\textbf{Context \& Pressure}\par


Healthcare integration is rarely a simple point-to-point message transmission. A single real-world clinical event—such as a patient admission or an emergency department transfer—triggers a cascade of multi-stage operational activities across multiple downstream systems (e.g., notifying pharmacy, updating bed management, triggering pre-admission lab orders, broadcasting HL7 ADT messages to billing and clinical portals).



When integration middleware treats these complex interactions simply as a stream of disconnected, fire-and-forget messages, failures in downstream destinations cause silent state divergence. The hospital's billing system may record the patient as discharged while the pharmacy continues dispensing medication.



\par\medskip\textbf{Architectural Imperative}\par


Harmonia must coordinate multi-stage operational activity in lockstep with the governed state of the real-world clinical entity. Activities must progress through explicit, observable state transitions, ensuring that fan-out status, acknowledgements, and retry progressions are fully trackable and recoverable.



\begin{itemize}


\item \textbf{Primary Stakeholders}: \href{\detokenize{../../../../../docs/markdown/01-motivation/stakeholders/enterprise-stakeholders.md}}{Healthcare Delivery Organizations}, \href{\detokenize{../../../../../docs/markdown/01-motivation/stakeholders/enterprise-stakeholders.md}}{Platform Operations \& Integration Engineers}

\item \textbf{Key Assessment}: \href{\detokenize{../../../../../docs/markdown/01-motivation/drivers-assessments/assessments.md}}{Unmonitored Destination Failure / Fan-Out Divergence}

\item \textbf{Associated Goal}: \href{\detokenize{../../../../../docs/markdown/01-motivation/goals-outcomes/strategic-goals.md}}{Coordinated Progression of Operational Activities \& Associated Entity State}

\item \textbf{Governing Axiom}: \hyperref[element:de7e90f46127455a9ec9ad3b9b48ba54]{AX-16 (Operational Activity \& Entity State Progress Together)}

\end{itemize}

{\footnotesize \href{../../../../continuum/motivation/driver/b3fb294d95024e4382b2c251c389fa66.md}{Canonical record: b3fb294d95024e4382b2c251c389fa66}; full qualifications and provenance.}\par
\section*{I.1.7 Durable Clinical Information Independence}
\phantomsection
\label{navigation:I.1.7}
\addcontentsline{toc}{section}{I.1.7 Durable Clinical Information Independence}

\label{element:3d81912dd8c94603ba0e10c98b7a62d8}
\index{Durable Clinical Information Independence (I.1.7)}
{\small Driver · Agreed · current}\par

\par\medskip\textbf{Context \& Pressure}\par


Clinical information must remain intelligible, authoritative, and actionable across a patient's entire lifetime (often 80+ years). In contrast, commercial Electronic Medical Record (EMR) software products, proprietary database schemas, and integration broker technologies have commercial lifecycles of only 5 to 10 years.



When healthcare networks store longitudinal patient data exclusively inside proprietary EMR schemas, they suffer severe vendor lock-in, exorbitant data extraction fees, and the perpetual risk of data obsolescence during vendor transitions.



\par\medskip\textbf{Architectural Imperative}\par


Harmonia must establish durable clinical information independence. Patient health data managed by Harmonia must exist as a vendor-neutral, sovereign regional asset. The clinical representation must outlive the commercial life of any individual software vendor, database engine, or middleware framework.



\begin{itemize}


\item \textbf{Primary Stakeholders}: \href{\detokenize{../../../../../docs/markdown/01-motivation/stakeholders/enterprise-stakeholders.md}}{Clinicians \& Care Teams}, \href{\detokenize{../../../../../docs/markdown/01-motivation/stakeholders/enterprise-stakeholders.md}}{Healthcare Delivery Organizations}

\item \textbf{Associated Goal}: \href{\detokenize{../../../../../docs/markdown/01-motivation/goals-outcomes/strategic-goals.md}}{Vendor-Independent Longitudinal Clinical Information Coherence}

\item \textbf{Governing Axioms}: \hyperref[element:e9bd61acaaea4577bc4086a5639e52c3]{AX-01 (Health-Information Centricity)}, \hyperref[element:952c3b6148ce4fce94b69f99a0c2fec7]{AX-06 (Explicit Information Authority)}

\end{itemize}

{\footnotesize \href{../../../../continuum/motivation/driver/3d81912dd8c94603ba0e10c98b7a62d8.md}{Canonical record: 3d81912dd8c94603ba0e10c98b7a62d8}; full qualifications and provenance.}\par
\section*{I.1.8 National Healthcare Directory / Endpoint Governance}
\phantomsection
\label{navigation:I.1.8}
\addcontentsline{toc}{section}{I.1.8 National Healthcare Directory / Endpoint Governance}

\label{element:5ea7e57dcde04ef88c19a88ade9af6ca}
\index{National Healthcare Directory / Endpoint Governance (I.1.8)}
{\small Driver · Agreed · current}\par

\par\medskip\textbf{Context \& Pressure}\par


Healthcare delivery requires seamless coordination between hundreds of clinics, hospitals, diagnostic practices, and thousands of accredited practitioners. Clinical correspondence, diagnostic reports, and e-referrals frequently go astray due to outdated contact details, retired practitioners, or misconfigured electronic messaging endpoints.



National regulatory frameworks (such as the Australian HI Service and AHPRA) publish authoritative registers of practitioners and healthcare organizations. Electronic messaging networks require verified service endpoints to ensure confidential messages are delivered only to legitimate recipients.



\par\medskip\textbf{Architectural Imperative}\par


Harmonia must govern practitioner credentials and electronic service endpoints through authoritative directory mechanisms. Integration workflows must validate that practitioner identities, roles, and destination addresses are authenticated and currently authorized.



\begin{itemize}


\item \textbf{Primary Stakeholders}: \href{\detokenize{../../../../../docs/markdown/01-motivation/stakeholders/enterprise-stakeholders.md}}{Healthcare Directory Stewards \& Registrars}, \href{\detokenize{../../../../../docs/markdown/01-motivation/stakeholders/enterprise-stakeholders.md}}{Clinicians \& Care Teams}

\item \textbf{External Authorities}: \href{\detokenize{../../../../../docs/markdown/01-motivation/stakeholders/external-authorities.md}}{HI Service}, \href{\detokenize{../../../../../docs/markdown/01-motivation/stakeholders/external-authorities.md}}{AHPRA}

\item \textbf{Associated Goal}: \href{\detokenize{../../../../../docs/markdown/01-motivation/goals-outcomes/strategic-goals.md}}{Reliable Subject Identity \& Referential Integrity}

\item \textbf{Governing Axiom}: \hyperref[element:d3fc6025dc7646d1b2020b51f7a44bfa]{AX-14 (Preservation of Semantic Distinctions)}

\end{itemize}

{\footnotesize \href{../../../../continuum/motivation/driver/5ea7e57dcde04ef88c19a88ade9af6ca.md}{Canonical record: 5ea7e57dcde04ef88c19a88ade9af6ca}; full qualifications and provenance.}\par
\section*{I.1.9 Patient Safety \& Identity Integrity}
\phantomsection
\label{navigation:I.1.9}
\addcontentsline{toc}{section}{I.1.9 Patient Safety \& Identity Integrity}

\label{element:b835c7b1055e4dbc9e1cd3a0c545c1c8}
\index{Patient Safety \& Identity Integrity (I.1.9)}
{\small Driver · Agreed · current}\par

\par\medskip\textbf{Context \& Pressure}\par


Patient safety is a foundational operational priority in healthcare. Accurately associating a clinical observation, pathology result, or medication order with the correct individual human subject is essential.



Mismatched patient records lead directly to severe clinical harm: administering contraindicated blood transfusions, giving incorrect medication doses to patients with unrecorded allergies, or performing wrong-patient interventions. In distributed regional networks, patients move between multiple facilities where different local medical record numbers (MRNs) are assigned.



\par\medskip\textbf{Architectural Imperative}\par


Harmonia must rigorously preserve and validate subject referential integrity across all managed clinical events. It must prevent ambiguous, invalid, or contradictory subject associations from being silently treated as valid.



\emph{Note on Architectural Boundary}: Harmonia does not act as an enterprise Master Patient Index (MPI) or deterministic golden-record matcher. Rather, it enforces strict referential validation and transparently captures identifier provenance without lossy unification.



\begin{itemize}


\item \textbf{Primary Stakeholders}: \href{\detokenize{../../../../../docs/markdown/01-motivation/stakeholders/enterprise-stakeholders.md}}{Patients \& Care Recipients}, \href{\detokenize{../../../../../docs/markdown/01-motivation/stakeholders/enterprise-stakeholders.md}}{Clinicians \& Care Teams}

\item \textbf{External Constraint}: \href{\detokenize{../../../../../docs/markdown/01-motivation/requirements-constraints/external-constraints.md}}{Applicable National / Jurisdictional Healthcare Identifier Obligations}

\item \textbf{Associated Goal}: \href{\detokenize{../../../../../docs/markdown/01-motivation/goals-outcomes/strategic-goals.md}}{Reliable Subject Identity \& Referential Integrity}

\item \textbf{Governing Axioms}: \hyperref[element:952c3b6148ce4fce94b69f99a0c2fec7]{AX-06 (Explicit Information Authority)}, \hyperref[element:d3fc6025dc7646d1b2020b51f7a44bfa]{AX-14 (Preservation of Semantic Distinctions)}

\end{itemize}

{\footnotesize \href{../../../../continuum/motivation/driver/b835c7b1055e4dbc9e1cd3a0c545c1c8.md}{Canonical record: b835c7b1055e4dbc9e1cd3a0c545c1c8}; full qualifications and provenance.}\par
\section*{I.1.10 Statutory Health Information Privacy \& Protection}
\phantomsection
\label{navigation:I.1.10}
\addcontentsline{toc}{section}{I.1.10 Statutory Health Information Privacy \& Protection}

\label{element:1dbf34b452b6497f914c704e037aec24}
\index{Statutory Health Information Privacy \& Protection (I.1.10)}
{\small Driver · Agreed · current}\par

\par\medskip\textbf{Context \& Pressure}\par


Health records contain the most sensitive category of personal data recognized under law (medical histories, psychiatric evaluations, sexual health diagnoses, genetic profiles). Healthcare networks are subject to rigorous statutory privacy regimes (e.g., the Australian \emph{Privacy Act 1988}, My Health Records Act, HIPAA, GDPR).



Statutory authorities enforce severe civil penalties, mandatory breach reporting, and reputational sanctions for unauthorized disclosure, eavesdropping, or accidental data leakage.



\par\medskip\textbf{Architectural Imperative}\par


Harmonia must implement intrinsic, default-deny security and privacy controls across all operations. Sensitive Protected Health Information (PHI) must never be emitted into non-clinical logging streams, telemetry systems, or debug channels. Access must require explicit cryptographic authentication and fine-grained authorization.



\begin{itemize}


\item \textbf{Primary Stakeholders}: \href{\detokenize{../../../../../docs/markdown/01-motivation/stakeholders/enterprise-stakeholders.md}}{Patients \& Care Recipients}, \href{\detokenize{../../../../../docs/markdown/01-motivation/stakeholders/enterprise-stakeholders.md}}{Platform Operations \& Integration Engineers}

\item \textbf{External Constraint}: \href{\detokenize{../../../../../docs/markdown/01-motivation/requirements-constraints/external-constraints.md}}{Applicable Health-Information Privacy and Data-Protection Obligations}

\item \textbf{Key Assessments}: \href{\detokenize{../../../../../docs/markdown/01-motivation/drivers-assessments/assessments.md}}{Implicit Perimeter Trust}, \href{\detokenize{../../../../../docs/markdown/01-motivation/drivers-assessments/assessments.md}}{PHI Leakage through Operational Logging}

\item \textbf{Governing Axioms}: \hyperref[element:5de2ae2c3b794c2d9f3a14efd51d564f]{AX-07 (Security Is Intrinsic to Managed Operations)}, \hyperref[element:21d6864657964ba5b03e3f6b5eb29cd5]{AX-08 (Evidence Records Meaning)}

\end{itemize}

{\footnotesize \href{../../../../continuum/motivation/driver/1dbf34b452b6497f914c704e037aec24.md}{Canonical record: 1dbf34b452b6497f914c704e037aec24}; full qualifications and provenance.}\par
\chapter*{I.2 Strategic Goals}
\phantomsection
\label{navigation:I.2}
\addcontentsline{toc}{chapter}{I.2 Strategic Goals}

\section*{I.2.1 Coordinated Progression of Operational Activities and Associated Entity State}
\phantomsection
\label{navigation:I.2.1}
\addcontentsline{toc}{section}{I.2.1 Coordinated Progression of Operational Activities and Associated Entity State}

\label{element:221eaf405d454c749112ece28ac595f1}
\index{Coordinated Progression of Operational Activities and Associated Entity State (I.2.1)}
{\small Goal · Agreed · current}\par

\par\medskip\textbf{Goal Statement}\par


Harmonia shall coordinate multi-stage operational activities through explicit, observable state transitions that remain strictly synchronized with the governed state of the associated real-world clinical entity.



\par\medskip\textbf{Rationale \& Context}\par


Healthcare integration activities (such as ADT distribution, lab order dispatch, and multi-destination fan-out) involve multi-stage distributed operations. Treating integration simply as a stream of fire-and-forget message transfers causes destination divergence and unmonitored failures (as diagnosed in \href{\detokenize{../../../../../docs/markdown/01-motivation/drivers-assessments/assessments.md}}{Unmonitored Destination Failure}).



\par\medskip\textbf{Architectural Boundary \& Intent}\par


\begin{itemize}


\item \textbf{Beyond Conventional Workflow Engines}: Avoid language implying that Harmonia is merely a generic workflow engine. Harmonia does not simply execute arbitrary BPMN scripts; it coordinates distributed integration activities in lockstep with the real-world lifecycle of clinical entities (e.g., patient admission state, lab specimen lifecycle, encounter status).

\item Requires explicit, observable progression state across multi-stage activities (\href{\detokenize{../../../../../docs/markdown/01-motivation/requirements-constraints/foundational-requirements.md}}{\texttt{\seqsplit{REQ-FND-002}}}).

\item Governed by \hyperref[element:de7e90f46127455a9ec9ad3b9b48ba54]{AX-16 (Operational Activity \& Entity State Progress Together)}.

\end{itemize}

{\footnotesize \href{../../../../continuum/motivation/goal/221eaf405d454c749112ece28ac595f1.md}{Canonical record: 221eaf405d454c749112ece28ac595f1}; full qualifications and provenance.}\par
\section*{I.2.2 Durable Acceptance and Preservation of Clinical Events}
\phantomsection
\label{navigation:I.2.2}
\addcontentsline{toc}{section}{I.2.2 Durable Acceptance and Preservation of Clinical Events}

\label{element:44daa77317aa44b7967b62b295d7288d}
\index{Durable Acceptance and Preservation of Clinical Events (I.2.2)}
{\small Goal · Agreed · current}\par

\par\medskip\textbf{Goal Statement}\par


Harmonia shall ensure that once an incoming clinical event is accepted by the platform, that event is guaranteed to be durably preserved and cannot be silently lost due to subsequent crashes, power interruptions, hardware faults, or transient network partitions.



\par\medskip\textbf{Rationale \& Context}\par


Clinical events—such as patient admissions, pathology test orders, critical laboratory alerts, and clinical discharge summaries—carry direct patient safety implications. The platform must eliminate the critical vulnerability identified in \href{\detokenize{../../../../../docs/markdown/01-motivation/drivers-assessments/assessments.md}}{Silent Data Loss via False Acceptance}, ensuring that acceptance at the ingress membrane is backed by authoritative durability before responsibility is released.



\par\medskip\textbf{Architectural Boundary}\par


\begin{itemize}


\item This goal defines the requirement for durable boundary crossing.

\item It does \textbf{not} mandate a specific storage technology (such as a relational database, append-only log, or distributed file system).

\item Downstream realization is governed by \href{\detokenize{../../../../../docs/markdown/01-motivation/requirements-constraints/foundational-requirements.md}}{\texttt{\seqsplit{REQ-FND-001}} Durable Ingress Acceptance Boundary}.

\end{itemize}

{\footnotesize \href{../../../../continuum/motivation/goal/44daa77317aa44b7967b62b295d7288d.md}{Canonical record: 44daa77317aa44b7967b62b295d7288d}; full qualifications and provenance.}\par
\section*{I.2.3 Real-Time Ingestion}
\phantomsection
\label{navigation:I.2.3}
\addcontentsline{toc}{section}{I.2.3 Real-Time Ingestion}

\label{element:d43c345839e548589e2b6aa677022f8b}
\index{Real-Time Ingestion (I.2.3)}
{\small Goal · Historical — current applicability unresolved · historical\_reference}\par
\href{../../../../continuum/motivation/goal/d43c345839e548589e2b6aa677022f8b.md}{Historical canonical definition and provenance}.\par
\section*{I.2.4 Reliable Subject Identity and Referential Integrity}
\phantomsection
\label{navigation:I.2.4}
\addcontentsline{toc}{section}{I.2.4 Reliable Subject Identity and Referential Integrity}

\label{element:e2f8b62e2cbd486282d68b6009c06de6}
\index{Reliable Subject Identity and Referential Integrity (I.2.4)}
{\small Goal · Agreed · current}\par

\par\medskip\textbf{Goal Statement}\par


Harmonia shall ensure that all managed clinical events, observations, and diagnostic records maintain verified, internally consistent, and referentially sound associations with the correct human patient subject.



\par\medskip\textbf{Rationale \& Context}\par


Accurate patient identification is fundamental to patient safety. Cross-associating clinical records between different individuals leads directly to contraindicated treatments and severe diagnostic errors. In regional health networks where patients possess multiple local Medical Record Numbers (MRNs) across different hospitals, the platform must preserve and validate identity references across every managed event.



\par\medskip\textbf{Architectural Boundaries \& Explicit Anti-Goals}\par


Harmonia explicitly defines what this goal \textbf{does NOT entail}:
- \textbf{Harmonia is NOT a Master Patient Index (MPI)}: Harmonia does not maintain an enterprise patient demographic master record.
- \textbf{Harmonia is NOT a Golden-Record Authority}: The platform does not synthesize or assert a single "golden" demographic record.
- \textbf{Harmonia is NOT a Deterministic Identity Matcher}: Harmonia does not execute probabilistic or deterministic patient matching algorithms to merge patient identities.
- \emph{Downstream Handoff}: Questions of enterprise identity matching, demographic curation, and golden-record master indices are intentionally deferred to \textbf{Domain 02 Strategy} and \textbf{Domain 04 Information Architecture}. In Domain 01, this goal enforces strict referential integrity, identifier validation, and non-destructive provenance capture (\href{\detokenize{../../../../../docs/markdown/01-motivation/requirements-constraints/foundational-requirements.md}}{\texttt{\seqsplit{REQ-FND-003}}}).


{\footnotesize \href{../../../../continuum/motivation/goal/e2f8b62e2cbd486282d68b6009c06de6.md}{Canonical record: e2f8b62e2cbd486282d68b6009c06de6}; full qualifications and provenance.}\par
\section*{I.2.5 Responsive Access to Managed Information}
\phantomsection
\label{navigation:I.2.5}
\addcontentsline{toc}{section}{I.2.5 Responsive Access to Managed Information}

\label{element:c7b8639a922f4d1cb2fcc68933cf7a90}
\index{Responsive Access to Managed Information (I.2.5)}
{\small Goal · Agreed · current}\par

\par\medskip\textbf{Goal Statement}\par


Harmonia shall provide consumers and internal services with timely, highly available, and responsive access to managed information and active coordination state, decoupled from the latency and lock contention of central synchronous persistence.



\par\medskip\textbf{Rationale \& Context}\par


In a regional network processing thousands of concurrent clinical interactions, forcing every active query and coordination check through central database locks introduces severe contention and availability bottlenecks (as diagnosed in \href{\detokenize{../../../../../docs/markdown/01-motivation/drivers-assessments/assessments.md}}{Centralised Synchronous Persistence Can Constrain Concurrent Processing}). Consumers require responsive, concurrent read paths that remain available even during background batch processing or heavy ingestion spikes.



\par\medskip\textbf{Architectural Boundary \& Technology Independence}\par


\begin{itemize}


\item \textbf{Strictly Technology-Independent}: This goal does \textbf{not} mandate caching, in-memory grids, Infinispan, Redis, or non-blocking reactive frameworks.

\item It defines the architectural intent: separating active, responsive query and coordination paths from heavy transactional persistence locks.

\item Realized downstream under \hyperref[element:8a1e995f16954484a78ae1c5c114d546]{AX-11 (Responsive \& Resilient Access)} and \hyperref[element:0412a2084f0145babc4f022c71fea65f]{AX-05 (State Separation)}.

\end{itemize}

{\footnotesize \href{../../../../continuum/motivation/goal/c7b8639a922f4d1cb2fcc68933cf7a90.md}{Canonical record: c7b8639a922f4d1cb2fcc68933cf7a90}; full qualifications and provenance.}\par
\section*{I.2.6 Sub-ms Cache Access}
\phantomsection
\label{navigation:I.2.6}
\addcontentsline{toc}{section}{I.2.6 Sub-ms Cache Access}

\label{element:f699096b8d0f4f5690eb57b25001b5e5}
\index{Sub-ms Cache Access (I.2.6)}
{\small Goal · Historical — current applicability unresolved · historical\_reference}\par
\href{../../../../continuum/motivation/goal/f699096b8d0f4f5690eb57b25001b5e5.md}{Historical canonical definition and provenance}.\par
\section*{I.2.7 Timely Availability of Clinical Information}
\phantomsection
\label{navigation:I.2.7}
\addcontentsline{toc}{section}{I.2.7 Timely Availability of Clinical Information}

\label{element:26e84fb0acfc49f1ac501b64a6e4bbf7}
\index{Timely Availability of Clinical Information (I.2.7)}
{\small Goal · Agreed · current}\par

\par\medskip\textbf{Goal Statement}\par


Harmonia shall make clinical information, event notifications, and diagnostic updates available to care teams and consuming hospital systems within operational latency thresholds appropriate to patient care.



\par\medskip\textbf{Rationale \& Context}\par


Healthcare delivery operates under tight clinical timelines. Point-of-care decisions in emergency departments and intensive care units require rapid turnaround of diagnostic findings. While high throughput is vital, low latency for critical alerts (e.g., life-threatening pathology results) is paramount to prevent diagnostic delays.



\par\medskip\textbf{Architectural Boundary}\par


\begin{itemize}


\item Governs the latency and throughput profile of integration pipelines.

\item Differentiates between urgent real-time clinical feeds and bulk historical synchronizations.

\item Hands over directly to \href{\detokenize{../../../../../docs/markdown/01-motivation/goals-outcomes/business-outcomes.md}}{Continuous and Resilient Regional Health Information Exchange (O2)}.

\end{itemize}

{\footnotesize \href{../../../../continuum/motivation/goal/26e84fb0acfc49f1ac501b64a6e4bbf7.md}{Canonical record: 26e84fb0acfc49f1ac501b64a6e4bbf7}; full qualifications and provenance.}\par
\section*{I.2.8 Vendor-Independent Longitudinal Clinical Information Coherence}
\phantomsection
\label{navigation:I.2.8}
\addcontentsline{toc}{section}{I.2.8 Vendor-Independent Longitudinal Clinical Information Coherence}

\label{element:905aca03c85748d79c3c065f6f16fc0a}
\index{Vendor-Independent Longitudinal Clinical Information Coherence (I.2.8)}
{\small Goal · Agreed · current}\par

\par\medskip\textbf{Goal Statement}\par


Harmonia shall establish and maintain a vendor-neutral, semantically coherent representation of a patient's longitudinal health record that outlives individual EMR software lifecycles and remains independent of proprietary vendor schemas.



\par\medskip\textbf{Rationale \& Context}\par


Patient care spans decades across multiple disparate healthcare institutions. Commercial EMR platforms have short market lifecycles and store data in proprietary, closed schemas that create vendor lock-in. To ensure continuity of care, the regional health network requires an enduring clinical asset that preserves the semantic meaning, units of measure, clinical nuance, and temporal sequence of health events across disparate provider systems.



\par\medskip\textbf{Architectural Boundary}\par


\begin{itemize}


\item Governed by \hyperref[element:e9bd61acaaea4577bc4086a5639e52c3]{AX-01 (Health-Information Centricity)}, \hyperref[element:eee8c3783d964d18a7f9dcea85fcd4cf]{AX-02 (Standards at Boundary)}, and \hyperref[element:952c3b6148ce4fce94b69f99a0c2fec7]{AX-06 (Explicit Information Authority)}.

\item Hands over directly to \href{\detokenize{../../../../../docs/markdown/01-motivation/goals-outcomes/business-outcomes.md}}{Reduced Clinical Risk (O1)} and \href{\detokenize{../../../../../docs/markdown/01-motivation/goals-outcomes/business-outcomes.md}}{Continuous Regional Exchange (O2)}.

\end{itemize}

{\footnotesize \href{../../../../continuum/motivation/goal/905aca03c85748d79c3c065f6f16fc0a.md}{Canonical record: 905aca03c85748d79c3c065f6f16fc0a}; full qualifications and provenance.}\par
\section*{I.2.9 Zero Message Loss}
\phantomsection
\label{navigation:I.2.9}
\addcontentsline{toc}{section}{I.2.9 Zero Message Loss}

\label{element:8f70dbb9c28f4dd9afe91c21243e5914}
\index{Zero Message Loss (I.2.9)}
{\small Goal · Historical — current applicability unresolved · historical\_reference}\par
\href{../../../../continuum/motivation/goal/8f70dbb9c28f4dd9afe91c21243e5914.md}{Historical canonical definition and provenance}.\par
\chapter*{I.3 Business Outcomes}
\phantomsection
\label{navigation:I.3}
\addcontentsline{toc}{chapter}{I.3 Business Outcomes}

\section*{I.3.1 Continuous and Resilient Regional Health Information Exchange (O2)}
\phantomsection
\label{navigation:I.3.1}
\addcontentsline{toc}{section}{I.3.1 Continuous and Resilient Regional Health Information Exchange (O2)}

\label{element:aed90b4a88f24b78adfeadbf0bca8980}
\index{Continuous and Resilient Regional Health Information Exchange (O2) (I.3.1)}
{\small Outcome · Agreed · current}\par

\par\medskip\textbf{Outcome Definition}\par


Regional healthcare facilities operate within an uninterrupted, highly available, and fault-tolerant digital integration fabric that reliably exchanges clinical events across heterogeneous systems without data loss or performance degradation.



\par\medskip\textbf{Clinical \& Operational Value}\par


\begin{itemize}


\item \textbf{Zero Silent Data Loss}: Admission notifications, emergency referrals, and diagnostic alerts are guaranteed to cross durable boundaries before positive transport acknowledgement, ensuring that integration failures never result in dropped clinical events.

\item \textbf{High Concurrency \& Responsiveness}: Point-of-care clinical queries and vital notification streams remain highly responsive, unblocked by background batch synchronization or database transaction locks.

\item \textbf{Coordinated Inter-Facility Workflows}: Multi-stage clinical activities (e.g., patient transfers, regional lab dispatch, pharmacy updates) progress deterministically across multiple destination systems, preventing divergence of clinical state between facilities.

\item \textbf{Seamless Standards Interoperability}: Seamless bridging of legacy HL7 v2 systems and modern FHIR RESTful repositories without expensive, disruptive rip-and-replace software upgrades.

\end{itemize}


\par\medskip\textbf{Direct Contributing Goals}\par


\begin{itemize}


\item \href{\detokenize{../../../../../docs/markdown/01-motivation/goals-outcomes/strategic-goals.md}}{Vendor-Independent Longitudinal Clinical Information Coherence} (Thread 1)

\item \href{\detokenize{../../../../../docs/markdown/01-motivation/goals-outcomes/strategic-goals.md}}{Durable Acceptance and Preservation of Clinical Events} (Thread 3)

\item \href{\detokenize{../../../../../docs/markdown/01-motivation/goals-outcomes/strategic-goals.md}}{Responsive Access to Managed Information} (Thread 4)

\item \href{\detokenize{../../../../../docs/markdown/01-motivation/goals-outcomes/strategic-goals.md}}{Coordinated Progression of Operational Activities and Associated Entity State} (Thread 7)

\end{itemize}

{\footnotesize \href{../../../../continuum/motivation/outcome/aed90b4a88f24b78adfeadbf0bca8980.md}{Canonical record: aed90b4a88f24b78adfeadbf0bca8980}; full qualifications and provenance.}\par
\section*{I.3.2 Demonstrable Protection and Accountable Handling of Health Information (O3)}
\phantomsection
\label{navigation:I.3.2}
\addcontentsline{toc}{section}{I.3.2 Demonstrable Protection and Accountable Handling of Health Information (O3)}

\label{element:75519ed8f7de4a619cc9065f525d703a}
\index{Demonstrable Protection and Accountable Handling of Health Information (O3) (I.3.2)}
{\small Outcome · Agreed · current}\par

\par\medskip\textbf{Outcome Definition}\par


Sensitive health information is demonstrably safeguarded across all integration operations, backed by default-deny authorization, zero-PHI diagnostic logging, and tamper-evident audit provenance that satisfies rigorous statutory data protection regimes.



\par\medskip\textbf{Clinical \& Operational Value}\par


\begin{itemize}


\item \textbf{Patient Trust \& Confidentiality}: Patients are assured that their intimate health data is protected from unauthorized access, eavesdropping, and accidental leakage into operational logs.

\item \textbf{Statutory Compliance \& Legal Assurance}: The regional health network demonstrably complies with national privacy legislation (e.g., the Australian \emph{Privacy Act 1988}, My Health Records Act, and Notifiable Data Breaches scheme), eliminating legal liability and multi-million-dollar breach penalties.

\item \textbf{Forensic Auditability}: Every state-changing clinical transaction maintains unambiguous, tamper-evident cryptographic provenance (who authorized it, what information was conveyed, when it occurred), enabling rigorous clinical incident retrospectives.

\end{itemize}


\par\medskip\textbf{Relationship to Motivational Principles}\par


Unlike Outcomes O1 and O2 (which receive direct handovers from specific strategic goals), Outcome O3 is an \textbf{intrinsic platform property}:
- It is realized through the foundational architectural principles \hyperref[element:5de2ae2c3b794c2d9f3a14efd51d564f]{AX-07 (Security Is Intrinsic to Managed Operations)} and \hyperref[element:21d6864657964ba5b03e3f6b5eb29cd5]{AX-08 (Evidence Records Meaning)}.
- Security and privacy are enforced intrinsically across all governed ingress, processing, and egress paths rather than being isolated to a single integration workflow.
- As established in the \href{\detokenize{../../../../../docs/markdown/01-motivation/orientation-view.md}}{Orientation View}, AX-07 and AX-08 do not maintain causal arrows to outcomes; O3 acts as a common consequence sink for the entire platform.


{\footnotesize \href{../../../../continuum/motivation/outcome/75519ed8f7de4a619cc9065f525d703a.md}{Canonical record: 75519ed8f7de4a619cc9065f525d703a}; full qualifications and provenance.}\par
\section*{I.3.3 Reduced Clinical Risk from Unavailable, Fragmented or Incorrectly Associated Information (O1)}
\phantomsection
\label{navigation:I.3.3}
\addcontentsline{toc}{section}{I.3.3 Reduced Clinical Risk from Unavailable, Fragmented or Incorrectly Associated Information (O1)}

\label{element:0582153fb1f1402b83e5358b61fd60af}
\index{Reduced Clinical Risk from Unavailable, Fragmented or Incorrectly Associated Information (O1) (I.3.3)}
{\small Outcome · Agreed · current}\par

\par\medskip\textbf{Outcome Definition}\par


Patients experience significantly safer care and reduced preventable harm across the regional healthcare network, achieved through the elimination of clinical data blind spots, the prevention of erroneous subject associations, and the provision of coherent longitudinal patient records.



\par\medskip\textbf{Clinical \& Operational Value}\par


\begin{itemize}


\item \textbf{Prevention of Adverse Medical Events}: Clinicians have point-of-care access to complete allergy lists, current medications, and past medical histories, preventing severe adverse drug events and contraindicated treatments.

\item \textbf{Elimination of Diagnostic Blind Spots}: Hospital emergency departments can immediately view recent laboratory results and diagnostic imaging reports performed at external facilities, avoiding redundant radiation exposure and delays in critical therapy.

\item \textbf{Accurate Subject Attribution}: Clinical records are referentially validated, ensuring that test results and treatment orders are never erroneously attributed to the wrong patient.

\end{itemize}


\par\medskip\textbf{Direct Contributing Goals}\par


\begin{itemize}


\item \href{\detokenize{../../../../../docs/markdown/01-motivation/goals-outcomes/strategic-goals.md}}{Vendor-Independent Longitudinal Clinical Information Coherence} (Thread 2)

\item \href{\detokenize{../../../../../docs/markdown/01-motivation/goals-outcomes/strategic-goals.md}}{Reliable Subject Identity and Referential Integrity} (Thread 5)

\end{itemize}

{\footnotesize \href{../../../../continuum/motivation/outcome/0582153fb1f1402b83e5358b61fd60af.md}{Canonical record: 0582153fb1f1402b83e5358b61fd60af}; full qualifications and provenance.}\par
\chapter*{I.4 Assessments}
\phantomsection
\label{navigation:I.4}
\addcontentsline{toc}{chapter}{I.4 Assessments}

\section*{I.4.1 Centralised Synchronous Persistence Can Constrain Concurrent Processing}
\phantomsection
\label{navigation:I.4.1}
\addcontentsline{toc}{section}{I.4.1 Centralised Synchronous Persistence Can Constrain Concurrent Processing}

\label{element:4f7b18cecb08415bad2b4247068a511e}
\index{Centralised Synchronous Persistence Can Constrain Concurrent Processing (I.4.1)}
{\small Assessment · Agreed · current}\par

\par\medskip\textbf{Assessment Summary}\par


Requiring central synchronous persistence interaction for every unit of active processing can introduce contention, latency, unnecessary I/O and coordination bottlenecks that limit concurrent throughput and scalability.



\par\medskip\textbf{Operational Hazard}\par


In a high-volume regional integration hub, thousands of clinical messages arrive concurrently from emergency rooms, pathology analyzers, and patient portals. If the architecture forces every active query, routing step, and transient coordination check to execute synchronous disk I/O and transactional row locking against a single central database:
- Transaction lock contention and wait states escalate as concurrency increases.
- Increased I/O wait times can cause ingress queues to fill and upstream connections to time out.
- High-priority emergency point-of-care clinical queries become stalled behind long-running background batch imports.
- Scalability is constrained to the physical I/O limits of the central persistence appliance.



\emph{Note on Technology Neutrality}: This assessment diagnoses the fundamental hazards of synchronous persistence contention, latency, and coordination bottlenecks. It is an architectural evaluation of concurrency dynamics; it does not mandate any specific product, cache, database, or queueing technology.



\par\medskip\textbf{Architectural Implication \& Response}\par


\begin{itemize}


\item The architecture must decouple active operational coordination from authoritative durable persistence.

\item High-frequency operational checks and client-facing queries must be served through responsive, concurrent access paths without imposing transactional lock contention on authoritative storage.

\item \textbf{Related Driver}: \href{\detokenize{../../../../../docs/markdown/01-motivation/drivers-assessments/drivers.md}}{Continuous Clinical Service Availability}

\item \textbf{Target Goal}: \href{\detokenize{../../../../../docs/markdown/01-motivation/goals-outcomes/strategic-goals.md}}{Responsive Access to Managed Information}

\item \textbf{Governing Axioms}: \hyperref[element:8a1e995f16954484a78ae1c5c114d546]{AX-11 (Responsive \& Resilient Information Access)}, \hyperref[element:0412a2084f0145babc4f022c71fea65f]{AX-05 (Active Coordination ≠ Authoritative State)}, \hyperref[element:9a5cb9e8876844e58e660f909f7df9fc]{AX-09 (Default Ephemerality of Operational State)}

\end{itemize}

{\footnotesize \href{../../../../continuum/motivation/assessment/4f7b18cecb08415bad2b4247068a511e.md}{Canonical record: 4f7b18cecb08415bad2b4247068a511e}; full qualifications and provenance.}\par
\section*{I.4.2 Implicit Perimeter Trust}
\phantomsection
\label{navigation:I.4.2}
\addcontentsline{toc}{section}{I.4.2 Implicit Perimeter Trust}

\label{element:f1b3b73329ae46b9b5105fb8ff236be5}
\index{Implicit Perimeter Trust (I.4.2)}
{\small Assessment · Agreed · current}\par

\par\medskip\textbf{Assessment Summary}\par


Assuming that internal network segments, container fabrics, or service-to-service calls are trustworthy by default leaves the platform vulnerable to lateral movement, credential abuse, and unauthorized clinical data manipulation.



\par\medskip\textbf{Operational Hazard}\par


Traditional enterprise architectures frequently rely on perimeter firewalls: once a request crosses the external network boundary, all internal service-to-service communication is unauthenticated, unencrypted, and fully trusted.



In modern distributed and containerized integration environments, perimeter-only trust is hazardous:
- Compromise of a single peripheral container or gateway allows an attacker to make unhindered calls to core clinical repositories.
- Internal integration routes can bypass clinical consent rules and authorization checks because the caller is presumed to be "internal."
- Audit logs cannot prove who authorized a specific data retrieval or state mutation because individual operations lack verifiable caller security context.



\par\medskip\textbf{Architectural Implication \& Response}\par


\begin{itemize}


\item Security must be intrinsic to every governed operation rather than dependent upon voluntary caller behaviour or perimeter assumptions.

\item Every internal request, service invocation, and state mutation must execute within an established, verifiable security context subject to platform-enforced security policy.

\item Downstream security architecture (Domain 08) derives specific enforcement patterns, such as default-deny authorization and authenticated service communication, from governing axiom AX-07.

\item \textbf{Related Driver}: \href{\detokenize{../../../../../docs/markdown/01-motivation/drivers-assessments/drivers.md}}{Statutory Health Information Privacy \& Protection}

\item \textbf{Governing Axiom}: \hyperref[element:5de2ae2c3b794c2d9f3a14efd51d564f]{AX-07 (Security Is Intrinsic to Managed Operations)}

\item \textbf{Downstream Requirements}: \texttt{\seqsplit{SEC-001 / SEC-002}} (Trusted Authentication \& Default-Deny Authorization).

\end{itemize}

{\footnotesize \href{../../../../continuum/motivation/assessment/f1b3b73329ae46b9b5105fb8ff236be5.md}{Canonical record: f1b3b73329ae46b9b5105fb8ff236be5}; full qualifications and provenance.}\par
\section*{I.4.3 PHI Leakage through Operational Logging}
\phantomsection
\label{navigation:I.4.3}
\addcontentsline{toc}{section}{I.4.3 PHI Leakage through Operational Logging}

\label{element:6d8848b700df4f6eacc0347993ac541a}
\index{PHI Leakage through Operational Logging (I.4.3)}
{\small Assessment · Agreed · current}\par

\par\medskip\textbf{Assessment Summary}\par


Emitting unmasked Protected Health Information (PHI) into operational diagnostic logs, monitoring systems, and tracing tools creates severe data breach exposures and violates statutory privacy obligations.



\par\medskip\textbf{Operational Hazard}\par


Software developers and integration engineers require diagnostic logging to troubleshoot routing errors, transform bugs, and network timeouts. A common developer practice is logging raw message payloads (e.g., logging an entire raw HL7 message string or complete FHIR resource JSON) at \texttt{\seqsplit{DEBUG}} or \texttt{\seqsplit{INFO}} levels.



In a healthcare environment, these logs are ingested into centralized search clusters, log forwarders, monitoring dashboards, and cloud observability tools:
- Patient names, home addresses, phone numbers, Medicare numbers, and intimate clinical diagnoses become stored in unencrypted, broadly accessible log repositories.
- Operational support personnel and external vendors gain unauthorized access to sensitive clinical data without clinical need-to-know authorization.
- Routine log retention and aggregation triggers substantial statutory data breach notifications under national privacy legislation (e.g., the Australian \emph{Privacy Act 1988} and Notifiable Data Breaches scheme).



\par\medskip\textbf{Architectural Implication \& Response}\par


\begin{itemize}


\item Operational log streams must be structurally segregated from clinical data payloads.

\item Diagnostic logs must emit only non-sensitive operational metadata (transaction UUIDs, message control IDs, timestamps, masked identifiers) and strictly prohibit unmasked clinical observation values, demographic text, or free-text clinical notes.

\item \textbf{Related Driver}: \href{\detokenize{../../../../../docs/markdown/01-motivation/drivers-assessments/drivers.md}}{Statutory Health Information Privacy \& Protection}

\item \textbf{Governing Axiom}: \hyperref[element:5de2ae2c3b794c2d9f3a14efd51d564f]{AX-07 (Security Is Intrinsic to Managed Operations)}

\item \textbf{Derived Guardrail \& Requirement}: Classifies downstream as a mandatory Security Architecture requirement (\texttt{\seqsplit{SEC-010}} PHI-Safe Operational Logging), with formal governance guardrails owned under Domain 13.

\end{itemize}

{\footnotesize \href{../../../../continuum/motivation/assessment/6d8848b700df4f6eacc0347993ac541a.md}{Canonical record: 6d8848b700df4f6eacc0347993ac541a}; full qualifications and provenance.}\par
\section*{I.4.4 Silent Data Loss via False Acceptance}
\phantomsection
\label{navigation:I.4.4}
\addcontentsline{toc}{section}{I.4.4 Silent Data Loss via False Acceptance}

\label{element:9f185cd26a0f4118a0b28bb08134f472}
\index{Silent Data Loss via False Acceptance (I.4.4)}
{\small Assessment · Agreed · current}\par

\par\medskip\textbf{Assessment Summary}\par


Prematurely acknowledging receipt of an incoming clinical event before responsibility for that event has crossed an authoritative durable boundary creates a critical risk of unrecoverable data loss.



\par\medskip\textbf{Operational Hazard}\par


When an ingress component positively acknowledges an incoming operational or clinical event before responsibility for that event has crossed an authoritative durable boundary, upstream senders legitimately consider their delivery obligation satisfied and release local responsibility.



If an unrecoverable failure occurs (such as process termination, node failure, network partition, or power loss) after positive acknowledgement but before durable persistence or guaranteed downstream delivery:
- The upstream sender marks the event as delivered and will not retransmit.
- The receiving platform loses the event representation without having persisted or processed it.
- The clinical event is silently dropped without any error indication or notification.



In clinical practice, this leads directly to missing critical diagnostic results, unrecorded medication orders, or undetected admission events, creating severe risks for patient safety.



\par\medskip\textbf{Downstream Integration Architecture Example}\par


In traditional HL7 v2 and MLLP transport middleware, inbound gateways often return an immediate positive acknowledgement (\texttt{\seqsplit{MSA-1 = AA}}, Application Accept) upon successful memory parsing or volatile buffering. If the gateway process crashes or hardware fails before the message is committed to durable storage, silent data loss occurs because the sender received an \texttt{\seqsplit{AA}} and will never retry. Downstream integration adapters (Domain 06) must enforce durable acceptance boundaries and emit negative acknowledgements (\texttt{\seqsplit{AE}}) when durability is not established.



\par\medskip\textbf{Architectural Implication \& Response}\par


\begin{itemize}


\item Positive acceptance must be strictly coupled to crossing an authoritative durable boundary.

\item Ingress boundaries must guarantee that incoming event responsibility is durably accepted before emitting a positive acceptance response to the sender. If durable acceptance cannot be established, an explicit negative or indeterminate response must be emitted so the sender retains retry responsibility.

\item \textbf{Related Driver}: \href{\detokenize{../../../../../docs/markdown/01-motivation/drivers-assessments/drivers.md}}{Continuous Clinical Service Availability}

\item \textbf{Target Goal}: \href{\detokenize{../../../../../docs/markdown/01-motivation/goals-outcomes/strategic-goals.md}}{Durable Acceptance \& Preservation of Clinical Events}

\item \textbf{Governing Axioms}: \hyperref[element:0d9673833eb840eab03804aa076c1b0c]{AX-10 (Failure as Normal Condition)}, \hyperref[element:0412a2084f0145babc4f022c71fea65f]{AX-05 (State Separation)}, \hyperref[element:ef8abbaec37e4890ad78b9f263e67906]{AX-15 (Explicit Representation of Uncertainty)}

\item \textbf{Foundational Requirement}: \href{\detokenize{../../../../../docs/markdown/01-motivation/requirements-constraints/foundational-requirements.md}}{\texttt{\seqsplit{REQ-FND-001}} Durable Ingress Acceptance Boundary}

\end{itemize}

{\footnotesize \href{../../../../continuum/motivation/assessment/9f185cd26a0f4118a0b28bb08134f472.md}{Canonical record: 9f185cd26a0f4118a0b28bb08134f472}; full qualifications and provenance.}\par
\section*{I.4.5 Unmonitored Destination Failure / Fan-Out Divergence}
\phantomsection
\label{navigation:I.4.5}
\addcontentsline{toc}{section}{I.4.5 Unmonitored Destination Failure / Fan-Out Divergence}

\label{element:3357fcd2f7ce457386e33a2f06ae2182}
\index{Unmonitored Destination Failure / Fan-Out Divergence (I.4.5)}
{\small Assessment · Agreed · current}\par

\par\medskip\textbf{Assessment Summary}\par


Broadcasting clinical events across multiple downstream destination systems without tracking granular, per-destination delivery status leads to silent divergence between destination systems and fractured clinical state.



\par\medskip\textbf{Operational Hazard}\par


A single clinical event often requires distribution to multiple disparate destinations (fan-out). For example, an HL7 \texttt{\seqsplit{ADT\textasciicircum{}A08}} (Update Patient Information) or \texttt{\seqsplit{ADT\textasciicircum{}A02}} (Transfer Patient) may need to update:
1. The regional clinical data repository.
2. The local pharmacy dispensing system.
3. The radiology information system (RIS).
4. The hospital billing and claims engine.
5. An external primary health network notifications feed.



If the integration platform considers the parent activity "completed" merely because the message was posted to an internal dispatch mechanism, failures at individual destination endpoints go unmonitored:
- The clinical repository is updated, but the pharmacy feed times out.
- The parent workflow reports success.
- The pharmacy continues operating on obsolete patient medication allergy data.
- Systemic clinical state diverges silently across the healthcare network.



\par\medskip\textbf{Architectural Implication \& Response}\par


\begin{itemize}


\item Operational workflows must track explicit progression state per destination.

\item Parent activities must maintain fine-grained sub-status across the entire fan-out lifecycle until every individual destination delivery is deterministically confirmed or explicitly flagged as failed/indeterminate.

\item \textbf{Related Driver}: \href{\detokenize{../../../../../docs/markdown/01-motivation/drivers-assessments/drivers.md}}{Coordinated Operational Activity and Entity State}

\item \textbf{Target Goal}: \href{\detokenize{../../../../../docs/markdown/01-motivation/goals-outcomes/strategic-goals.md}}{Coordinated Progression of Operational Activities \& Associated Entity State}

\item \textbf{Governing Axiom}: \hyperref[element:de7e90f46127455a9ec9ad3b9b48ba54]{AX-16 (Operational Activity \& Entity State Progress Together)}

\item \textbf{Foundational Requirement}: \href{\detokenize{../../../../../docs/markdown/01-motivation/requirements-constraints/foundational-requirements.md}}{\texttt{\seqsplit{REQ-FND-002}} Operational Activity Progression State}

\end{itemize}

{\footnotesize \href{../../../../continuum/motivation/assessment/3357fcd2f7ce457386e33a2f06ae2182.md}{Canonical record: 3357fcd2f7ce457386e33a2f06ae2182}; full qualifications and provenance.}\par
\chapter*{I.5 Requirements}
\phantomsection
\label{navigation:I.5}
\addcontentsline{toc}{chapter}{I.5 Requirements}

\section*{I.5.1 AI Assurance Architecture Discovery Before Implementation}
\phantomsection
\label{navigation:I.5.1}
\addcontentsline{toc}{section}{I.5.1 AI Assurance Architecture Discovery Before Implementation}

\label{element:24a609117a4f40839f8b669ebf720cde}
\index{AI Assurance Architecture Discovery Before Implementation (I.5.1)}
{\small Requirement · Adjudicated Domain05 direction — U02/U06/U19; holistic result acceptance pending · current}\par

Domain05 must undertake explicit Information Systems Layer Architecture discovery before implementation models are finalised, including: AI model/service selection and abstraction; context/input construction; information disclosure/security; model/service provenance and versioning; confidence/uncertainty/result representation; deterministic versus non-deterministic behaviour; reproducibility; result validation; human review/escalation; assurance evidence representation; workflow integration; failure/timeout/unavailability behaviour; and reporting/observability consistent with Dokimasia's terminal-boundary rules.



This assertion is normalised from the governed DesignNote under U02/U06/U19 and the ARB release Praxis. Original candidate/release qualifications and separately deferred implementation details remain in force.



Architectural subject: \hyperref[element:fb3f37a6ac614415ac486a6de9d35994]{Energeia}.


{\footnotesize \href{../../../../continuum/motivation/requirement/24a609117a4f40839f8b669ebf720cde.md}{Canonical record: 24a609117a4f40839f8b669ebf720cde}; full qualifications and provenance.}\par
\section*{I.5.2 Artemis HA Failover}
\phantomsection
\label{navigation:I.5.2}
\addcontentsline{toc}{section}{I.5.2 Artemis HA Failover}

\label{element:8c0f08d0d12f48b9ae6a001c19ede377}
\index{Artemis HA Failover (I.5.2)}
{\small Requirement · Historical — current applicability unresolved · historical\_reference}\par
\href{../../../../continuum/motivation/requirement/8c0f08d0d12f48b9ae6a001c19ede377.md}{Historical canonical definition and provenance}.\par
\section*{I.5.3 Audit Corrections Create New Linked Events}
\phantomsection
\label{navigation:I.5.3}
\addcontentsline{toc}{section}{I.5.3 Audit Corrections Create New Linked Events}

\label{element:ddaa9bff2b684af7a6ad18fe4f2da8ad}
\index{Audit Corrections Create New Linked Events (I.5.3)}
{\small Requirement · Adjudicated Domain05 direction — U07/U15; holistic result acceptance pending · current}\par

Where an earlier AuditEvent is subsequently determined to be incorrect, incomplete, misleading, or otherwise requires qualification, Kleio SHALL create a \textbf{new AuditEvent} that refers back to the earlier AuditEvent and records the subsequent assertion.



This assertion is normalised from the governed DesignNote under U07/U15 and the ARB release Praxis. Original candidate/release qualifications and separately deferred implementation details remain in force.



Architectural subject: \hyperref[element:fefda110e2274884ad58906733606efd]{Kleio}.


{\footnotesize \href{../../../../continuum/motivation/requirement/ddaa9bff2b684af7a6ad18fe4f2da8ad.md}{Canonical record: ddaa9bff2b684af7a6ad18fe4f2da8ad}; full qualifications and provenance.}\par
\section*{I.5.4 CLIN-001 — Vendor-Neutral Clinical Information Layer}
\phantomsection
\label{navigation:I.5.4}
\addcontentsline{toc}{section}{I.5.4 CLIN-001 — Vendor-Neutral Clinical Information Layer}

\label{element:44eec8cf17f44593824a1f723d949352}
\index{CLIN-001 — Vendor-Neutral Clinical Information Layer (I.5.4)}
{\small Requirement · Historical — current owning-domain standing not inferred · historical\_reference}\par
\href{../../../../continuum/motivation/requirement/44eec8cf17f44593824a1f723d949352.md}{Historical canonical definition and provenance}.\par
\section*{I.5.5 CLIN-002 — FHIR-Based Longitudinal Clinical Representation}
\phantomsection
\label{navigation:I.5.5}
\addcontentsline{toc}{section}{I.5.5 CLIN-002 — FHIR-Based Longitudinal Clinical Representation}

\label{element:9ad88c6b7f934923800dd46c8522b9a1}
\index{CLIN-002 — FHIR-Based Longitudinal Clinical Representation (I.5.5)}
{\small Requirement · Historical — current owning-domain standing not inferred · historical\_reference}\par
\href{../../../../continuum/motivation/requirement/9ad88c6b7f934923800dd46c8522b9a1.md}{Historical canonical definition and provenance}.\par
\section*{I.5.6 CLIN-003 — Medication as a First-Class Clinical Domain}
\phantomsection
\label{navigation:I.5.6}
\addcontentsline{toc}{section}{I.5.6 CLIN-003 — Medication as a First-Class Clinical Domain}

\label{element:f5a6dd8f23f047bb9f4c327bbd3e85f2}
\index{CLIN-003 — Medication as a First-Class Clinical Domain (I.5.6)}
{\small Requirement · Historical — current owning-domain standing not inferred · historical\_reference}\par
\href{../../../../continuum/motivation/requirement/f5a6dd8f23f047bb9f4c327bbd3e85f2.md}{Historical canonical definition and provenance}.\par
\section*{I.5.7 CLIN-004 — Source-System Independence}
\phantomsection
\label{navigation:I.5.7}
\addcontentsline{toc}{section}{I.5.7 CLIN-004 — Source-System Independence}

\label{element:a5a3d2ddb02d4ac4b3eabdfb2ecd22bb}
\index{CLIN-004 — Source-System Independence (I.5.7)}
{\small Requirement · Historical — current owning-domain standing not inferred · historical\_reference}\par
\href{../../../../continuum/motivation/requirement/a5a3d2ddb02d4ac4b3eabdfb2ecd22bb.md}{Historical canonical definition and provenance}.\par
\section*{I.5.8 CLIN-005 — Source Authority Preservation}
\phantomsection
\label{navigation:I.5.8}
\addcontentsline{toc}{section}{I.5.8 CLIN-005 — Source Authority Preservation}

\label{element:5f399a338688438c922df969e724a24a}
\index{CLIN-005 — Source Authority Preservation (I.5.8)}
{\small Requirement · Historical — current owning-domain standing not inferred · historical\_reference}\par
\href{../../../../continuum/motivation/requirement/5f399a338688438c922df969e724a24a.md}{Historical canonical definition and provenance}.\par
\section*{I.5.9 CLIN-006 — Authority-Aware Clinical Reconciliation}
\phantomsection
\label{navigation:I.5.9}
\addcontentsline{toc}{section}{I.5.9 CLIN-006 — Authority-Aware Clinical Reconciliation}

\label{element:3d7c709a435145e3821184b2c9b980fa}
\index{CLIN-006 — Authority-Aware Clinical Reconciliation (I.5.9)}
{\small Requirement · Historical — current owning-domain standing not inferred · historical\_reference}\par
\href{../../../../continuum/motivation/requirement/3d7c709a435145e3821184b2c9b980fa.md}{Historical canonical definition and provenance}.\par
\section*{I.5.10 CLIN-007 — Preservation of Conflicting Assertions}
\phantomsection
\label{navigation:I.5.10}
\addcontentsline{toc}{section}{I.5.10 CLIN-007 — Preservation of Conflicting Assertions}

\label{element:cd669acb9a76499db70b5d092917a56b}
\index{CLIN-007 — Preservation of Conflicting Assertions (I.5.10)}
{\small Requirement · Historical — current owning-domain standing not inferred · historical\_reference}\par
\href{../../../../continuum/motivation/requirement/cd669acb9a76499db70b5d092917a56b.md}{Historical canonical definition and provenance}.\par
\section*{I.5.11 CLIN-008 — Clinical Source Context}
\phantomsection
\label{navigation:I.5.11}
\addcontentsline{toc}{section}{I.5.11 CLIN-008 — Clinical Source Context}

\label{element:591d9b8007504d2382ae7977a6589742}
\index{CLIN-008 — Clinical Source Context (I.5.11)}
{\small Requirement · Historical — current owning-domain standing not inferred · historical\_reference}\par
\href{../../../../continuum/motivation/requirement/591d9b8007504d2382ae7977a6589742.md}{Historical canonical definition and provenance}.\par
\section*{I.5.12 CLIN-009 — Governed Patient Identity}
\phantomsection
\label{navigation:I.5.12}
\addcontentsline{toc}{section}{I.5.12 CLIN-009 — Governed Patient Identity}

\label{element:4dded2374b9346e4b8d0e326e7d9c97c}
\index{CLIN-009 — Governed Patient Identity (I.5.12)}
{\small Requirement · Historical — current owning-domain standing not inferred · historical\_reference}\par
\href{../../../../continuum/motivation/requirement/4dded2374b9346e4b8d0e326e7d9c97c.md}{Historical canonical definition and provenance}.\par
\section*{I.5.13 CLIN-010 — Read-Oriented Initial LHR}
\phantomsection
\label{navigation:I.5.13}
\addcontentsline{toc}{section}{I.5.13 CLIN-010 — Read-Oriented Initial LHR}

\label{element:96dad3bf0c764ca0849227eae64848aa}
\index{CLIN-010 — Read-Oriented Initial LHR (I.5.13)}
{\small Requirement · Historical — current owning-domain standing not inferred · historical\_reference}\par
\href{../../../../continuum/motivation/requirement/96dad3bf0c764ca0849227eae64848aa.md}{Historical canonical definition and provenance}.\par
\section*{I.5.14 CLIN-011 — Selective Clinical Authoring}
\phantomsection
\label{navigation:I.5.14}
\addcontentsline{toc}{section}{I.5.14 CLIN-011 — Selective Clinical Authoring}

\label{element:bd593ed8f3b3465ca499151f18767585}
\index{CLIN-011 — Selective Clinical Authoring (I.5.14)}
{\small Requirement · Historical — current owning-domain standing not inferred · historical\_reference}\par
\href{../../../../continuum/motivation/requirement/bd593ed8f3b3465ca499151f18767585.md}{Historical canonical definition and provenance}.\par
\section*{I.5.15 CORE-001 — Information Authority Classification}
\phantomsection
\label{navigation:I.5.15}
\addcontentsline{toc}{section}{I.5.15 CORE-001 — Information Authority Classification}

\label{element:6e76bb1db8f2464496355d47d2f981ef}
\index{CORE-001 — Information Authority Classification (I.5.15)}
{\small Requirement · Historical — current owning-domain standing not inferred · historical\_reference}\par
\href{../../../../continuum/motivation/requirement/6e76bb1db8f2464496355d47d2f981ef.md}{Historical canonical definition and provenance}.\par
\section*{I.5.16 CORE-002 — Separation of Architectural Concerns}
\phantomsection
\label{navigation:I.5.16}
\addcontentsline{toc}{section}{I.5.16 CORE-002 — Separation of Architectural Concerns}

\label{element:a5c09af151984e989590b9198e2def70}
\index{CORE-002 — Separation of Architectural Concerns (I.5.16)}
{\small Requirement · Historical — current owning-domain standing not inferred · historical\_reference}\par
\href{../../../../continuum/motivation/requirement/a5c09af151984e989590b9198e2def70.md}{Historical canonical definition and provenance}.\par
\section*{I.5.17 CORE-003 — Governed Canonical Semantics}
\phantomsection
\label{navigation:I.5.17}
\addcontentsline{toc}{section}{I.5.17 CORE-003 — Governed Canonical Semantics}

\label{element:006383a5cce64780960fda32b6682d8d}
\index{CORE-003 — Governed Canonical Semantics (I.5.17)}
{\small Requirement · Historical — current owning-domain standing not inferred · historical\_reference}\par
\href{../../../../continuum/motivation/requirement/006383a5cce64780960fda32b6682d8d.md}{Historical canonical definition and provenance}.\par
\section*{I.5.18 CORE-004 — Correlation and Causation}
\phantomsection
\label{navigation:I.5.18}
\addcontentsline{toc}{section}{I.5.18 CORE-004 — Correlation and Causation}

\label{element:f2e504e7d6374e41a81303265adda6de}
\index{CORE-004 — Correlation and Causation (I.5.18)}
{\small Requirement · Historical — current owning-domain standing not inferred · historical\_reference}\par
\href{../../../../continuum/motivation/requirement/f2e504e7d6374e41a81303265adda6de.md}{Historical canonical definition and provenance}.\par
\section*{I.5.19 DATA-001 — Durable Governed State}
\phantomsection
\label{navigation:I.5.19}
\addcontentsline{toc}{section}{I.5.19 DATA-001 — Durable Governed State}

\label{element:15c71ccd5d714503891aece2a23445e1}
\index{DATA-001 — Durable Governed State (I.5.19)}
{\small Requirement · Historical — current owning-domain standing not inferred · historical\_reference}\par
\href{../../../../continuum/motivation/requirement/15c71ccd5d714503891aece2a23445e1.md}{Historical canonical definition and provenance}.\par
\section*{I.5.20 DATA-002 — Authoritative Read Semantics}
\phantomsection
\label{navigation:I.5.20}
\addcontentsline{toc}{section}{I.5.20 DATA-002 — Authoritative Read Semantics}

\label{element:bc55838493c24617823410798c8d7598}
\index{DATA-002 — Authoritative Read Semantics (I.5.20)}
{\small Requirement · Historical — current owning-domain standing not inferred · historical\_reference}\par
\href{../../../../continuum/motivation/requirement/bc55838493c24617823410798c8d7598.md}{Historical canonical definition and provenance}.\par
\section*{I.5.21 DATA-003 — Cache Transparency}
\phantomsection
\label{navigation:I.5.21}
\addcontentsline{toc}{section}{I.5.21 DATA-003 — Cache Transparency}

\label{element:ff04c3dab8a04359b9c4d62939839be7}
\index{DATA-003 — Cache Transparency (I.5.21)}
{\small Requirement · Historical — current owning-domain standing not inferred · historical\_reference}\par
\href{../../../../continuum/motivation/requirement/ff04c3dab8a04359b9c4d62939839be7.md}{Historical canonical definition and provenance}.\par
\section*{I.5.22 DATA-004 — Search Scalability}
\phantomsection
\label{navigation:I.5.22}
\addcontentsline{toc}{section}{I.5.22 DATA-004 — Search Scalability}

\label{element:f27e74b6cb13454aa4339932436971e3}
\index{DATA-004 — Search Scalability (I.5.22)}
{\small Requirement · Historical — current owning-domain standing not inferred · historical\_reference}\par
\href{../../../../continuum/motivation/requirement/f27e74b6cb13454aa4339932436971e3.md}{Historical canonical definition and provenance}.\par
\section*{I.5.23 DATA-005 — Concurrency and Version Integrity}
\phantomsection
\label{navigation:I.5.23}
\addcontentsline{toc}{section}{I.5.23 DATA-005 — Concurrency and Version Integrity}

\label{element:7f40b2f4947845d287564b43f812b6c8}
\index{DATA-005 — Concurrency and Version Integrity (I.5.23)}
{\small Requirement · Historical — current owning-domain standing not inferred · historical\_reference}\par
\href{../../../../continuum/motivation/requirement/7f40b2f4947845d287564b43f812b6c8.md}{Historical canonical definition and provenance}.\par
\section*{I.5.24 DATA-006 — Referential Integrity}
\phantomsection
\label{navigation:I.5.24}
\addcontentsline{toc}{section}{I.5.24 DATA-006 — Referential Integrity}

\label{element:4a8a79f520394cfe863f473ec2a31b34}
\index{DATA-006 — Referential Integrity (I.5.24)}
{\small Requirement · Historical — current owning-domain standing not inferred · historical\_reference}\par
\href{../../../../continuum/motivation/requirement/4a8a79f520394cfe863f473ec2a31b34.md}{Historical canonical definition and provenance}.\par
\section*{I.5.25 DATA-007 — Persistence Failure Safety}
\phantomsection
\label{navigation:I.5.25}
\addcontentsline{toc}{section}{I.5.25 DATA-007 — Persistence Failure Safety}

\label{element:40ed2565d0d1467cbbfcf5706924de94}
\index{DATA-007 — Persistence Failure Safety (I.5.25)}
{\small Requirement · Historical — current owning-domain standing not inferred · historical\_reference}\par
\href{../../../../continuum/motivation/requirement/40ed2565d0d1467cbbfcf5706924de94.md}{Historical canonical definition and provenance}.\par
\section*{I.5.26 DATA-008 — Recovery Fidelity}
\phantomsection
\label{navigation:I.5.26}
\addcontentsline{toc}{section}{I.5.26 DATA-008 — Recovery Fidelity}

\label{element:7c7a33be31d042459adb83c71ae044fa}
\index{DATA-008 — Recovery Fidelity (I.5.26)}
{\small Requirement · Historical — current owning-domain standing not inferred · historical\_reference}\par
\href{../../../../continuum/motivation/requirement/7c7a33be31d042459adb83c71ae044fa.md}{Historical canonical definition and provenance}.\par
\section*{I.5.27 Development Representations of Runtime Governed Datasets}
\phantomsection
\label{navigation:I.5.27}
\addcontentsline{toc}{section}{I.5.27 Development Representations of Runtime Governed Datasets}

\label{element:0a4a5c73048c4e49a82d40e349e70a95}
\index{Development Representations of Runtime Governed Datasets (I.5.27)}
{\small Requirement · Adjudicated Domain05 direction — U13/U15; holistic result acceptance pending · current}\par

\begin{quote}


\textbf{Calliope SHALL provide development-time representations of known
runtime-managed DataSets, CodeSets and ValueSets sufficient to support
efficient human, IDE and AI-assisted software development, while
preserving runtime resolution against authoritative Calliope-DataSets
state.}



\end{quote}


This assertion is normalised from the governed DesignNote under U13/U15 and the ARB release Praxis. Original candidate/release qualifications and separately deferred implementation details remain in force.



Architectural subject: \hyperref[element:466ad8f75a954c0581b5bba090443849]{Calliope}.


{\footnotesize \href{../../../../continuum/motivation/requirement/0a4a5c73048c4e49a82d40e349e70a95.md}{Canonical record: 0a4a5c73048c4e49a82d40e349e70a95}; full qualifications and provenance.}\par
\section*{I.5.28 DIR-001 — Standards-Based Provider Directory}
\phantomsection
\label{navigation:I.5.28}
\addcontentsline{toc}{section}{I.5.28 DIR-001 — Standards-Based Provider Directory}

\label{element:c1b1d82f8a9b4c649a3216c5e09a22c6}
\index{DIR-001 — Standards-Based Provider Directory (I.5.28)}
{\small Requirement · Historical — current owning-domain standing not inferred · historical\_reference}\par
\href{../../../../continuum/motivation/requirement/c1b1d82f8a9b4c649a3216c5e09a22c6.md}{Historical canonical definition and provenance}.\par
\section*{I.5.29 DIR-002 — Provider Directory Discovery}
\phantomsection
\label{navigation:I.5.29}
\addcontentsline{toc}{section}{I.5.29 DIR-002 — Provider Directory Discovery}

\label{element:27998c2a8f2e497a814b32e7a0bc5e03}
\index{DIR-002 — Provider Directory Discovery (I.5.29)}
{\small Requirement · Historical — current owning-domain standing not inferred · historical\_reference}\par
\href{../../../../continuum/motivation/requirement/27998c2a8f2e497a814b32e7a0bc5e03.md}{Historical canonical definition and provenance}.\par
\section*{I.5.30 DIR-003 — Governed Directory Change}
\phantomsection
\label{navigation:I.5.30}
\addcontentsline{toc}{section}{I.5.30 DIR-003 — Governed Directory Change}

\label{element:934462c11133409cb3aa59bf2502fa81}
\index{DIR-003 — Governed Directory Change (I.5.30)}
{\small Requirement · Historical — current owning-domain standing not inferred · historical\_reference}\par
\href{../../../../continuum/motivation/requirement/934462c11133409cb3aa59bf2502fa81.md}{Historical canonical definition and provenance}.\par
\section*{I.5.31 DIR-004 — Directory Referential Integrity}
\phantomsection
\label{navigation:I.5.31}
\addcontentsline{toc}{section}{I.5.31 DIR-004 — Directory Referential Integrity}

\label{element:f866ae5adb08499794331bc600db827c}
\index{DIR-004 — Directory Referential Integrity (I.5.31)}
{\small Requirement · Historical — current owning-domain standing not inferred · historical\_reference}\par
\href{../../../../continuum/motivation/requirement/f866ae5adb08499794331bc600db827c.md}{Historical canonical definition and provenance}.\par
\section*{I.5.32 DIR-005 — Healthcare Identifier Governance}
\phantomsection
\label{navigation:I.5.32}
\addcontentsline{toc}{section}{I.5.32 DIR-005 — Healthcare Identifier Governance}

\label{element:097368b4b0b4468dbf1c57c11d9ed86c}
\index{DIR-005 — Healthcare Identifier Governance (I.5.32)}
{\small Requirement · Historical — current owning-domain standing not inferred · historical\_reference}\par
\href{../../../../continuum/motivation/requirement/097368b4b0b4468dbf1c57c11d9ed86c.md}{Historical canonical definition and provenance}.\par
\section*{I.5.33 DIR-006 — Provider Self-Service Governance}
\phantomsection
\label{navigation:I.5.33}
\addcontentsline{toc}{section}{I.5.33 DIR-006 — Provider Self-Service Governance}

\label{element:4230f93461a4470383fb99fddc35f8bf}
\index{DIR-006 — Provider Self-Service Governance (I.5.33)}
{\small Requirement · Historical — current owning-domain standing not inferred · historical\_reference}\par
\href{../../../../continuum/motivation/requirement/4230f93461a4470383fb99fddc35f8bf.md}{Historical canonical definition and provenance}.\par
\section*{I.5.34 Efficient Materialisation of Anticipated Context Resources}
\phantomsection
\label{navigation:I.5.34}
\addcontentsline{toc}{section}{I.5.34 Efficient Materialisation of Anticipated Context Resources}

\label{element:1f3431ca25284783ac69f8f6affcd103}
\index{Efficient Materialisation of Anticipated Context Resources (I.5.34)}
{\small Requirement · Adjudicated Domain05 direction — U10/U18; holistic result acceptance pending · current}\par

Mneme SHALL support efficient materialisation of the anticipated working resource set declared by the Ergo definition.



This assertion is normalised from the governed DesignNote under U10/U18 and the ARB release Praxis. Original candidate/release qualifications and separately deferred implementation details remain in force.



Architectural subject: \hyperref[element:1af8a56358f7429fabfb57be66542a99]{Mneme}.


{\footnotesize \href{../../../../continuum/motivation/requirement/1f3431ca25284783ac69f8f6affcd103.md}{Canonical record: 1f3431ca25284783ac69f8f6affcd103}; full qualifications and provenance.}\par
\section*{I.5.35 Establish Effective Security Attributes through Governed Rules}
\phantomsection
\label{navigation:I.5.35}
\addcontentsline{toc}{section}{I.5.35 Establish Effective Security Attributes through Governed Rules}

\label{element:d14061c03b9042948ba33cae91e02efb}
\index{Establish Effective Security Attributes through Governed Rules (I.5.35)}
{\small Requirement · Adjudicated Domain05 direction — U12/U15/U16; holistic result acceptance pending · current}\par

Where required Information.Compartment and/or Information.Classification attributes are absent, the applicable enforcement mechanism SHALL establish effective values through governed, deterministic rules.



This assertion is normalised from the governed DesignNote under U12/U15/U16 and the ARB release Praxis. Original candidate/release qualifications and separately deferred implementation details remain in force.



Architectural subject: \hyperref[element:899d34ec491948b68aad09e507dea6d1]{Themis}.


{\footnotesize \href{../../../../continuum/motivation/requirement/d14061c03b9042948ba33cae91e02efb.md}{Canonical record: d14061c03b9042948ba33cae91e02efb}; full qualifications and provenance.}\par
\section*{I.5.36 Explain Effective PEP Criteria Derivation}
\phantomsection
\label{navigation:I.5.36}
\addcontentsline{toc}{section}{I.5.36 Explain Effective PEP Criteria Derivation}

\label{element:f027de050273400eaac0de083b820238}
\index{Explain Effective PEP Criteria Derivation (I.5.36)}
{\small Requirement · Adjudicated Domain05 direction — U12/U15/U16; holistic result acceptance pending · current}\par

Themis SHALL retain sufficient provenance to explain the derivation of effective PEP-Criteria, including the authoritative Criteria, Criteria collections, AuthorisationRole assignments, direct assignments and applicable policy state from which authorisation was derived.



This assertion is normalised from the governed DesignNote under U12/U15/U16 and the ARB release Praxis. Original candidate/release qualifications and separately deferred implementation details remain in force.



Architectural subject: \hyperref[element:899d34ec491948b68aad09e507dea6d1]{Themis}.


{\footnotesize \href{../../../../continuum/motivation/requirement/f027de050273400eaac0de083b820238.md}{Canonical record: f027de050273400eaac0de083b820238}; full qualifications and provenance.}\par
\section*{I.5.37 FHIR R5 Schemas}
\phantomsection
\label{navigation:I.5.37}
\addcontentsline{toc}{section}{I.5.37 FHIR R5 Schemas}

\label{element:8d8c62b8cb7b40b0be2ce23f7f41f154}
\index{FHIR R5 Schemas (I.5.37)}
{\small Requirement · Historical / Reclassified — Domain04/06 · historical\_reference}\par
\href{../../../../continuum/motivation/requirement/8d8c62b8cb7b40b0be2ce23f7f41f154.md}{Historical canonical definition and provenance}.\par
\section*{I.5.38 INT-001 — Standards-Based Healthcare Interoperability}
\phantomsection
\label{navigation:I.5.38}
\addcontentsline{toc}{section}{I.5.38 INT-001 — Standards-Based Healthcare Interoperability}

\label{element:44dba056ada04f6583b77859aac9ac68}
\index{INT-001 — Standards-Based Healthcare Interoperability (I.5.38)}
{\small Requirement · Historical — current owning-domain standing not inferred · historical\_reference}\par
\href{../../../../continuum/motivation/requirement/44dba056ada04f6583b77859aac9ac68.md}{Historical canonical definition and provenance}.\par
\section*{I.5.39 INT-002 — Reliable Messaging}
\phantomsection
\label{navigation:I.5.39}
\addcontentsline{toc}{section}{I.5.39 INT-002 — Reliable Messaging}

\label{element:aac9b66911e849cb9792c3a2b7eacb88}
\index{INT-002 — Reliable Messaging (I.5.39)}
{\small Requirement · Historical — current owning-domain standing not inferred · historical\_reference}\par
\href{../../../../continuum/motivation/requirement/aac9b66911e849cb9792c3a2b7eacb88.md}{Historical canonical definition and provenance}.\par
\section*{I.5.40 INT-003 — Messaging Recovery}
\phantomsection
\label{navigation:I.5.40}
\addcontentsline{toc}{section}{I.5.40 INT-003 — Messaging Recovery}

\label{element:6e72cf09e7cc4acf9466dd39e1a15ae6}
\index{INT-003 — Messaging Recovery (I.5.40)}
{\small Requirement · Historical — current owning-domain standing not inferred · historical\_reference}\par
\href{../../../../continuum/motivation/requirement/6e72cf09e7cc4acf9466dd39e1a15ae6.md}{Historical canonical definition and provenance}.\par
\section*{I.5.41 INT-004 — FHIR R5 Boundary}
\phantomsection
\label{navigation:I.5.41}
\addcontentsline{toc}{section}{I.5.41 INT-004 — FHIR R5 Boundary}

\label{element:927b0cf4710e424391120b124067051e}
\index{INT-004 — FHIR R5 Boundary (I.5.41)}
{\small Requirement · Historical — current owning-domain standing not inferred · historical\_reference}\par
\href{../../../../continuum/motivation/requirement/927b0cf4710e424391120b124067051e.md}{Historical canonical definition and provenance}.\par
\section*{I.5.42 INT-005 — Opaque Resource Identity}
\phantomsection
\label{navigation:I.5.42}
\addcontentsline{toc}{section}{I.5.42 INT-005 — Opaque Resource Identity}

\label{element:91bd7a8732b3468b84be2493103433e7}
\index{INT-005 — Opaque Resource Identity (I.5.42)}
{\small Requirement · Historical — current owning-domain standing not inferred · historical\_reference}\par
\href{../../../../continuum/motivation/requirement/91bd7a8732b3468b84be2493103433e7.md}{Historical canonical definition and provenance}.\par
\section*{I.5.43 INT-006 — Governed Validation}
\phantomsection
\label{navigation:I.5.43}
\addcontentsline{toc}{section}{I.5.43 INT-006 — Governed Validation}

\label{element:6c14251b5d4e4eff96d2fd812834d124}
\index{INT-006 — Governed Validation (I.5.43)}
{\small Requirement · Historical — current owning-domain standing not inferred · historical\_reference}\par
\href{../../../../continuum/motivation/requirement/6c14251b5d4e4eff96d2fd812834d124.md}{Historical canonical definition and provenance}.\par
\section*{I.5.44 INT-007 — Terminology Abstraction}
\phantomsection
\label{navigation:I.5.44}
\addcontentsline{toc}{section}{I.5.44 INT-007 — Terminology Abstraction}

\label{element:6fe86be142044e5aa7720e2e6493a72a}
\index{INT-007 — Terminology Abstraction (I.5.44)}
{\small Requirement · Historical — current owning-domain standing not inferred · historical\_reference}\par
\href{../../../../continuum/motivation/requirement/6fe86be142044e5aa7720e2e6493a72a.md}{Historical canonical definition and provenance}.\par
\section*{I.5.45 Normal Reporting Does Not Require Raw PubSub Topology Interpretation}
\phantomsection
\label{navigation:I.5.45}
\addcontentsline{toc}{section}{I.5.45 Normal Reporting Does Not Require Raw PubSub Topology Interpretation}

\label{element:aa569bb1448648f685bb73cf9c6563e3}
\index{Normal Reporting Does Not Require Raw PubSub Topology Interpretation (I.5.45)}
{\small Requirement · Adjudicated Domain05 direction — U04/U06/U14/U18; holistic result acceptance pending · current}\par

Harmonia SHALL NOT require humans to reason directly from raw Pub/Sub topology and telemetry for normal reporting and assurance purposes.



This assertion is normalised from the governed DesignNote under U04/U06/U14/U18 and the ARB release Praxis. Original candidate/release qualifications and separately deferred implementation details remain in force.



Architectural subject: \hyperref[element:fb3f37a6ac614415ac486a6de9d35994]{Energeia}.


{\footnotesize \href{../../../../continuum/motivation/requirement/aa569bb1448648f685bb73cf9c6563e3.md}{Canonical record: aa569bb1448648f685bb73cf9c6563e3}; full qualifications and provenance.}\par
\section*{I.5.46 Operations and Assurance Supports Semantic Preserving Partitioning}
\phantomsection
\label{navigation:I.5.46}
\addcontentsline{toc}{section}{I.5.46 Operations and Assurance Supports Semantic Preserving Partitioning}

\label{element:3238d37e27a7421e905b705c5f30fc85}
\index{Operations and Assurance Supports Semantic Preserving Partitioning (I.5.46)}
{\small Requirement · Adjudicated Domain05 direction — U11; holistic result acceptance pending · current}\par

\textbf{Harmonia Operations \& Assurance (O\&A)} SHALL permit horizontal partitioning without changing semantic identity, provenance, relationships or governed meaning.



This assertion is normalised from the governed DesignNote under U11 and the ARB release Praxis. Original candidate/release qualifications and separately deferred implementation details remain in force.



Architectural subject: \hyperref[element:7985e647bf9e4ca3ba70faf9e3e4ae02]{Mnemosyne}.


{\footnotesize \href{../../../../continuum/motivation/requirement/3238d37e27a7421e905b705c5f30fc85.md}{Canonical record: 3238d37e27a7421e905b705c5f30fc85}; full qualifications and provenance.}\par
\section*{I.5.47 OPS-001 — Reproducible Deployment}
\phantomsection
\label{navigation:I.5.47}
\addcontentsline{toc}{section}{I.5.47 OPS-001 — Reproducible Deployment}

\label{element:c4373b1fd8444f10b502a5c3216152d7}
\index{OPS-001 — Reproducible Deployment (I.5.47)}
{\small Requirement · Historical — current owning-domain standing not inferred · historical\_reference}\par
\href{../../../../continuum/motivation/requirement/c4373b1fd8444f10b502a5c3216152d7.md}{Historical canonical definition and provenance}.\par
\section*{I.5.48 OPS-002 — Durable Infrastructure State}
\phantomsection
\label{navigation:I.5.48}
\addcontentsline{toc}{section}{I.5.48 OPS-002 — Durable Infrastructure State}

\label{element:d5d6a9ea450047b1bd62b13177e8fca6}
\index{OPS-002 — Durable Infrastructure State (I.5.48)}
{\small Requirement · Historical — current owning-domain standing not inferred · historical\_reference}\par
\href{../../../../continuum/motivation/requirement/d5d6a9ea450047b1bd62b13177e8fca6.md}{Historical canonical definition and provenance}.\par
\section*{I.5.49 OPS-003 — Meaningful Health and Readiness}
\phantomsection
\label{navigation:I.5.49}
\addcontentsline{toc}{section}{I.5.49 OPS-003 — Meaningful Health and Readiness}

\label{element:4c02fec3af264b878ee3e0cfaecbf55f}
\index{OPS-003 — Meaningful Health and Readiness (I.5.49)}
{\small Requirement · Historical — current owning-domain standing not inferred · historical\_reference}\par
\href{../../../../continuum/motivation/requirement/4c02fec3af264b878ee3e0cfaecbf55f.md}{Historical canonical definition and provenance}.\par
\section*{I.5.50 OPS-004 — Operational Observability}
\phantomsection
\label{navigation:I.5.50}
\addcontentsline{toc}{section}{I.5.50 OPS-004 — Operational Observability}

\label{element:b5d3447a05b249a8af154136fc36c4d9}
\index{OPS-004 — Operational Observability (I.5.50)}
{\small Requirement · Historical — current owning-domain standing not inferred · historical\_reference}\par
\href{../../../../continuum/motivation/requirement/b5d3447a05b249a8af154136fc36c4d9.md}{Historical canonical definition and provenance}.\par
\section*{I.5.51 OPS-005 — Production Resilience}
\phantomsection
\label{navigation:I.5.51}
\addcontentsline{toc}{section}{I.5.51 OPS-005 — Production Resilience}

\label{element:01e7d6168d554048a424ef75a6556040}
\index{OPS-005 — Production Resilience (I.5.51)}
{\small Requirement · Historical — current owning-domain standing not inferred · historical\_reference}\par
\href{../../../../continuum/motivation/requirement/01e7d6168d554048a424ef75a6556040.md}{Historical canonical definition and provenance}.\par
\section*{I.5.52 OPS-006 — Architecture Conformance}
\phantomsection
\label{navigation:I.5.52}
\addcontentsline{toc}{section}{I.5.52 OPS-006 — Architecture Conformance}

\label{element:dee1792d83bf413f974fc24802084872}
\index{OPS-006 — Architecture Conformance (I.5.52)}
{\small Requirement · Historical — current owning-domain standing not inferred · historical\_reference}\par
\href{../../../../continuum/motivation/requirement/dee1792d83bf413f974fc24802084872.md}{Historical canonical definition and provenance}.\par
\section*{I.5.53 Preserve Necessary Evidence When Historical Resolution Is Unavailable}
\phantomsection
\label{navigation:I.5.53}
\addcontentsline{toc}{section}{I.5.53 Preserve Necessary Evidence When Historical Resolution Is Unavailable}

\label{element:81ce6a25725e41bb8e92c945f53c8b4c}
\index{Preserve Necessary Evidence When Historical Resolution Is Unavailable (I.5.53)}
{\small Requirement · Adjudicated Domain05 direction — U07/U15; holistic result acceptance pending · current}\par

Where the authoritative source cannot guarantee longitudinal retrieval of the information as it existed at the relevant audited point in time, Kleio SHALL preserve the information necessary to maintain the integrity of the audit evidence.



This assertion is normalised from the governed DesignNote under U07/U15 and the ARB release Praxis. Original candidate/release qualifications and separately deferred implementation details remain in force.



Architectural subject: \hyperref[element:fefda110e2274884ad58906733606efd]{Kleio}.


{\footnotesize \href{../../../../continuum/motivation/requirement/81ce6a25725e41bb8e92c945f53c8b4c.md}{Canonical record: 81ce6a25725e41bb8e92c945f53c8b4c}; full qualifications and provenance.}\par
\section*{I.5.54 PROV-001 — End-to-End Provenance}
\phantomsection
\label{navigation:I.5.54}
\addcontentsline{toc}{section}{I.5.54 PROV-001 — End-to-End Provenance}

\label{element:3e45e60affb44bf384d54de26d3add6a}
\index{PROV-001 — End-to-End Provenance (I.5.54)}
{\small Requirement · Historical — current owning-domain standing not inferred · historical\_reference}\par
\href{../../../../continuum/motivation/requirement/3e45e60affb44bf384d54de26d3add6a.md}{Historical canonical definition and provenance}.\par
\section*{I.5.55 PROV-002 — Source Context}
\phantomsection
\label{navigation:I.5.55}
\addcontentsline{toc}{section}{I.5.55 PROV-002 — Source Context}

\label{element:929ed42779a843dd84b683e3dcc35eb3}
\index{PROV-002 — Source Context (I.5.55)}
{\small Requirement · Historical — current owning-domain standing not inferred · historical\_reference}\par
\href{../../../../continuum/motivation/requirement/929ed42779a843dd84b683e3dcc35eb3.md}{Historical canonical definition and provenance}.\par
\section*{I.5.56 PROV-003 — Dual-Actor Provenance}
\phantomsection
\label{navigation:I.5.56}
\addcontentsline{toc}{section}{I.5.56 PROV-003 — Dual-Actor Provenance}

\label{element:c016286faa3d4dfbbad133f73e855f1c}
\index{PROV-003 — Dual-Actor Provenance (I.5.56)}
{\small Requirement · Historical — current owning-domain standing not inferred · historical\_reference}\par
\href{../../../../continuum/motivation/requirement/c016286faa3d4dfbbad133f73e855f1c.md}{Historical canonical definition and provenance}.\par
\section*{I.5.57 PROV-004 — Provenance and Authority Independence}
\phantomsection
\label{navigation:I.5.57}
\addcontentsline{toc}{section}{I.5.57 PROV-004 — Provenance and Authority Independence}

\label{element:3fb9a89d345a4cb6a56d62823316664e}
\index{PROV-004 — Provenance and Authority Independence (I.5.57)}
{\small Requirement · Historical — current owning-domain standing not inferred · historical\_reference}\par
\href{../../../../continuum/motivation/requirement/3fb9a89d345a4cb6a56d62823316664e.md}{Historical canonical definition and provenance}.\par
\section*{I.5.58 PROV-005 — Immutable Audit}
\phantomsection
\label{navigation:I.5.58}
\addcontentsline{toc}{section}{I.5.58 PROV-005 — Immutable Audit}

\label{element:0f7b23cfd94e450c8a2772879bc9bb1e}
\index{PROV-005 — Immutable Audit (I.5.58)}
{\small Requirement · Historical — current owning-domain standing not inferred · historical\_reference}\par
\href{../../../../continuum/motivation/requirement/0f7b23cfd94e450c8a2772879bc9bb1e.md}{Historical canonical definition and provenance}.\par
\section*{I.5.59 PROV-006 — Audit Discoverability}
\phantomsection
\label{navigation:I.5.59}
\addcontentsline{toc}{section}{I.5.59 PROV-006 — Audit Discoverability}

\label{element:927c65d2d665451396782b4c1d7c8c86}
\index{PROV-006 — Audit Discoverability (I.5.59)}
{\small Requirement · Historical — current owning-domain standing not inferred · historical\_reference}\par
\href{../../../../continuum/motivation/requirement/927c65d2d665451396782b4c1d7c8c86.md}{Historical canonical definition and provenance}.\par
\section*{I.5.60 REQ-FND-001: Durable Ingress Acceptance Boundary}
\phantomsection
\label{navigation:I.5.60}
\addcontentsline{toc}{section}{I.5.60 REQ-FND-001: Durable Ingress Acceptance Boundary}

\label{element:01c834537a1e41dfb523d48f287d93e9}
\index{REQ-FND-001: Durable Ingress Acceptance Boundary (I.5.60)}
{\small Requirement · Accepted · current}\par

\par\medskip\textbf{Normative Requirement Statement}\par


\begin{quote}


\textbf{Harmonia shall not positively acknowledge acceptance of an incoming event until responsibility for that event has crossed the required durable acceptance boundary.}



\end{quote}


\par\medskip\textbf{Architectural Intent \& Rationale}\par


This requirement eliminates the critical vulnerability identified in \href{\detokenize{../../../../../docs/markdown/01-motivation/drivers-assessments/assessments.md}}{Silent Data Loss via False Acceptance}. When an upstream external system transmits a clinical event (such as an emergency admission, critical lab result, or medication order), it remains the responsibility of that upstream sender until Harmonia has durably secured the event.



Releasing an upstream sender from responsibility (by returning a positive transport acknowledgement) before the event is committed to durable storage creates a vulnerable window where gateway crashes, memory exhaustions, or hardware reboots cause silent, unrecoverable data loss.



\par\medskip\textbf{Protocol-Neutral Delineation \& Downstream Realisations}\par


\begin{itemize}


\item \textbf{Normative Rule}: The requirement mandates that the durable acceptance boundary must be crossed before positive acknowledgement. It is completely protocol-neutral.

\item \textbf{Downstream Realisations}:

\item In HL7 v2 over MLLP integration (owned by Domain 06 Integration Architecture), this rule is realised by forbidding an \texttt{\seqsplit{MSA-1 = AA}} (Application Accept) acknowledgement until downstream persistence or queue publication succeeds. If persistence fails, the gateway MUST return \texttt{\seqsplit{MSA-1 = AE}} (Application Error) to trigger upstream sender retry.

\item In FHIR RESTful APIs, this rule is realised by returning \texttt{\seqsplit{201 Created}} or \texttt{\seqsplit{200 OK}} only after transactional database commitment, returning \texttt{\seqsplit{202 Accepted}} only when durable asynchronous queueing is verified, and returning \texttt{\seqsplit{500 Internal Server Error}} or \texttt{\seqsplit{503 Service Unavailable}} on durable failure.

\item 

In Ingress Dual-Write safety (\texttt{\seqsplit{REC-001}}), this rule enforces that in-memory active caching alone is insufficient to satisfy positive acceptance.




\item 

\textbf{Primary Motivation}: \href{\detokenize{../../../../../docs/markdown/01-motivation/drivers-assessments/assessments.md}}{Silent Data Loss via False Acceptance}, \href{\detokenize{../../../../../docs/markdown/01-motivation/drivers-assessments/drivers.md}}{Continuous Clinical Service Availability}




\item \textbf{Governing Principles}: \hyperref[element:0412a2084f0145babc4f022c71fea65f]{AX-05 (State Separation)}, \hyperref[element:0d9673833eb840eab03804aa076c1b0c]{AX-10 (Failure as Normal Condition)}, \hyperref[element:ef8abbaec37e4890ad78b9f263e67906]{AX-15 (Explicit Representation of Uncertainty)}

\item \textbf{Historical Identifier Reference}: \texttt{\seqsplit{REQ-INGRESS-001}}

\end{itemize}

{\footnotesize \href{../../../../continuum/motivation/requirement/01c834537a1e41dfb523d48f287d93e9.md}{Canonical record: 01c834537a1e41dfb523d48f287d93e9}; full qualifications and provenance.}\par
\section*{I.5.61 REQ-FND-002: Operational Activity Progression State}
\phantomsection
\label{navigation:I.5.61}
\addcontentsline{toc}{section}{I.5.61 REQ-FND-002: Operational Activity Progression State}

\label{element:eaabcdb96aad4f32a050c94f9ecec049}
\index{REQ-FND-002: Operational Activity Progression State (I.5.61)}
{\small Requirement · Accepted · current}\par

\par\medskip\textbf{Normative Requirement Statement}\par


\begin{quote}


\textbf{Harmonia shall maintain sufficient explicit and observable state to determine the progression and outcome of managed multi-stage operational activities.}



\end{quote}


\par\medskip\textbf{Architectural Intent \& Rationale}\par


Complex health integration involves distributed, multi-stage activities—such as routing ADT events, dispatching laboratory orders, and executing multi-destination fan-out deliveries. If integration middleware treats these operations as disconnected fire-and-forget messages, failures at individual destination endpoints cause unmonitored divergence between hospital systems (as diagnosed in \href{\detokenize{../../../../../docs/markdown/01-motivation/drivers-assessments/assessments.md}}{Unmonitored Destination Failure}).



Harmonia mandates that every multi-stage operational activity must maintain explicit, queryable, and observable progression state across its entire lifecycle. Operational state must capture checkpoints, destination delivery statuses, error codes, and completion outcomes, enabling operational engineers to monitor, audit, and reconcile distributed activities.



\par\medskip\textbf{Downstream Scope}\par


\begin{itemize}


\item Realised in Domain 05 Application Architecture and Domain 06 Integration Architecture through Pragma envelopes, FHIR \texttt{\seqsplit{Task}} resources, and task progression tracking.

\item 

Does not mandate a specific workflow engine; it mandates observable progression state synchronized with real-world entity state.




\item 

\textbf{Primary Motivation}: \href{\detokenize{../../../../../docs/markdown/01-motivation/drivers-assessments/assessments.md}}{Unmonitored Destination Failure / Fan-Out Divergence}, \href{\detokenize{../../../../../docs/markdown/01-motivation/drivers-assessments/drivers.md}}{Coordinated Operational Activity and Entity State}




\item \textbf{Governing Principle}: \hyperref[element:de7e90f46127455a9ec9ad3b9b48ba54]{AX-16 (Operational Activity \& Entity State Progress Together)}

\item \textbf{Historical Identifier Reference}: \texttt{\seqsplit{REQ-COORD-001}}

\end{itemize}

{\footnotesize \href{../../../../continuum/motivation/requirement/eaabcdb96aad4f32a050c94f9ecec049.md}{Canonical record: eaabcdb96aad4f32a050c94f9ecec049}; full qualifications and provenance.}\par
\section*{I.5.62 REQ-FND-003: Subject Referential Integrity}
\phantomsection
\label{navigation:I.5.62}
\addcontentsline{toc}{section}{I.5.62 REQ-FND-003: Subject Referential Integrity}

\label{element:43ecbfe8f9b14fcbbc59ee0c03c8774b}
\index{REQ-FND-003: Subject Referential Integrity (I.5.62)}
{\small Requirement · Accepted · current}\par

\par\medskip\textbf{Normative Requirement Statement}\par


\begin{quote}


\textbf{Harmonia shall preserve and validate the integrity of subject references and associations used within managed information, preventing known-invalid or internally inconsistent associations from being silently treated as valid.}



\end{quote}


\par\medskip\textbf{Architectural Intent \& Rationale}\par


Clinical records derive their entire medical validity from their association with the correct human patient subject. In regional health networks, patient identifiers arrive from disparate sources (national IHIs, regional MRNs, local facility patient numbers). Mismatched or corrupted subject references lead directly to patient harm.



Harmonia mandates that any subject reference attached to a managed clinical event must be structurally validated, referentially verified, and checked for internal consistency. Known-invalid identifiers (e.g., failed check-digits, malformed namespaces, or conflicting patient identifiers within a single clinical bundle) must be rejected or explicitly flagged rather than silently propagated as valid.



\par\medskip\textbf{Explicit Boundary: Non-MPI Demarcation}\par


\begin{itemize}


\item \textbf{Not a Master Patient Index (MPI)}: This requirement does \textbf{not} establish Harmonia as an MPI or enterprise demographic matching engine.

\item \textbf{Not a Golden-Record Determiner}: Harmonia does not execute probabilistic matching or merge patient demographics into a synthetic "golden record."

\item 

\textbf{Intent}: Preserving the referential integrity, authenticity, and provenance of subject identifiers without performing lossy or speculative demographic unification.




\item 

\textbf{Primary Motivation}: \href{\detokenize{../../../../../docs/markdown/01-motivation/drivers-assessments/drivers.md}}{Patient Safety \& Identity Integrity}




\item \textbf{External Constraint}: \href{\detokenize{../../../../../docs/markdown/01-motivation/requirements-constraints/external-constraints.md}}{Applicable National / Jurisdictional Healthcare Identifier Obligations}

\item \textbf{Governing Principles}: \hyperref[element:952c3b6148ce4fce94b69f99a0c2fec7]{AX-06 (Explicit Information Authority)}, \hyperref[element:d3fc6025dc7646d1b2020b51f7a44bfa]{AX-14 (Preservation of Semantic Distinctions)}

\item \textbf{Historical Identifier Reference}: \texttt{\seqsplit{REQ-IDENT-001}}

\end{itemize}

{\footnotesize \href{../../../../continuum/motivation/requirement/43ecbfe8f9b14fcbbc59ee0c03c8774b.md}{Canonical record: 43ecbfe8f9b14fcbbc59ee0c03c8774b}; full qualifications and provenance.}\par
\section*{I.5.63 REQ-FND-004: Explicit Indeterminate Outcome}
\phantomsection
\label{navigation:I.5.63}
\addcontentsline{toc}{section}{I.5.63 REQ-FND-004: Explicit Indeterminate Outcome}

\label{element:fc732b64fc564127b73174e28d9ccce8}
\index{REQ-FND-004: Explicit Indeterminate Outcome (I.5.63)}
{\small Requirement · Accepted · current}\par

\par\medskip\textbf{Normative Requirement Statement}\par


\begin{quote}


\textbf{Where Harmonia cannot establish the outcome of a managed operation with sufficient certainty, the outcome shall be represented explicitly as indeterminate and shall not be interpreted as either success or failure without subsequent reconciliation or another authoritative determination.}



\end{quote}


\par\medskip\textbf{Architectural Intent \& Rationale}\par


In distributed networks, communication timeouts, socket disconnects, and unacknowledged dispatches leave the status of remote operations uncertain. In the presence of a network partition or destination crash, a client cannot know whether an update was processed and the ACK was lost, or whether the update was never received.



A frequent architectural failure mode is "binary collapsing": assuming an unconfirmed operation failed (triggering duplicate side effects on blind retry) or assuming it succeeded (ignoring unperformed clinical orders). Harmonia mandates that uncertainty must be preserved and explicitly represented as \texttt{\seqsplit{INDETERMINATE}}.



\par\medskip\textbf{Reconciliation \& Retry Semantics}\par


\begin{itemize}


\item This requirement does \textbf{not} prohibit automated retries or resilient recovery.

\item 

It mandates that retry mechanisms must be safe, idempotent, and reconciliation-aware. An indeterminate outcome must be resolved through explicit reconciliation (such as querying destination state or verifying idempotency tokens) rather than blind, speculative assumption.




\item 

\textbf{Primary Motivation}: \href{\detokenize{../../../../../docs/markdown/01-motivation/drivers-assessments/assessments.md}}{Silent Data Loss via False Acceptance}, \href{\detokenize{../../../../../docs/markdown/01-motivation/drivers-assessments/drivers.md}}{Continuous Clinical Service Availability}




\item \textbf{Governing Principle}: \hyperref[element:ef8abbaec37e4890ad78b9f263e67906]{AX-15 (Uncertainty Is Preserved Until Resolved)}

\item \textbf{Historical Identifier Reference}: \texttt{\seqsplit{REQ-UNCERT-001}}

\end{itemize}

{\footnotesize \href{../../../../continuum/motivation/requirement/fc732b64fc564127b73174e28d9ccce8.md}{Canonical record: fc732b64fc564127b73174e28d9ccce8}; full qualifications and provenance.}\par
\section*{I.5.64 REQ-FND-005: Independent Assurance of Governed Activity}
\phantomsection
\label{navigation:I.5.64}
\addcontentsline{toc}{section}{I.5.64 REQ-FND-005: Independent Assurance of Governed Activity}

\label{element:bbb7438e78604308abba93afc14f4e2e}
\index{REQ-FND-005: Independent Assurance of Governed Activity (I.5.64)}
{\small Requirement · APPROVED · current}\par

\textbf{Classification:} Foundational Platform Requirement. \textbf{Status:} APPROVED — human architectural approval recorded on 2026-10-08; normative wording and relationships unchanged. The authorised Independent Assurance reconciliation is APPROVED / CLOSED; Domain01 Motivation Architecture is CLOSED / FROZEN.



\par\medskip\textbf{Normative Requirement Statement}\par


\begin{quote}


\textbf{Harmonia shall support independent governed assurance of activities, information and outcomes for which assurance is required, such that assurance conclusions are established from sufficient and trustworthy evidence against applicable obligations, expectations or controls.}



\textbf{The activity or mechanism responsible for performing or managing the subject being assured shall not be solely responsible for determining, suppressing, manufacturing or retrospectively altering its assurance outcome. Harmonia shall establish sufficient separation of assurance responsibility and authority from responsibility for performance or control of the subject so that assurance is not merely self-attestation by the subject being assured.}



\textbf{Assurance may identify findings, exceptions and the need for operational response, but shall not thereby assume responsibility for management, assignment, delegation, escalation or remediation of the activity being assured.}



\textbf{Assurance may manage its own assurance activity without thereby assuming management of the subject activity.}



\textbf{Where available evidence is insufficient to establish an assurance conclusion with the required degree of confidence, that insufficiency shall remain explicit and shall not be interpreted as either assurance or non-assurance.}



\end{quote}


\par\medskip\textbf{Architectural Intent \& Rationale}\par


Those accountable for health-information handling and operational activity need credible grounds for demonstrating that applicable obligations, expectations or controls have been met. The \hyperref[element:e76216c494fa4500bb80eda6f83513c1]{Regional Health Network Operator's compliance and forensic-review concerns}, \hyperref[element:c20aa17ce4aa4a58be5eb3cff2ecdd1d]{engineers' incident review and deterministic verification concerns}, and \hyperref[element:75519ed8f7de4a619cc9065f525d703a]{O3's demonstrable protection and accountable handling} already establish the value of substantive review and trustworthy evidence.



Those existing concerns do not themselves establish performer independence. The approved reconciliation adds that obligation explicitly: a conclusion controlled solely by the activity or mechanism being assured can conceal findings or manufacture confidence and cannot provide the required independent assurance. Evidence availability and operational completion alone cannot establish whether the applicable obligations, expectations or controls have been satisfied.



\par\medskip\textbf{Applicability \& Responsibility Boundary}\par


The requirement applies where assurance is required. It does not require assurance of every activity, information item or outcome. The criteria establishing applicability, applicable obligations and the required degree of confidence remain for subsequent governed derivation; this requirement does not define an applicability model.



Independence concerns responsibility and authority. It does not prescribe physical, technical or organisational separation, nor require external assurance. It does not allocate an assessor or prohibit the subject's performer from contributing evidence. Assurance manages its own work and may inform an operational response; responsibility for managing that response remains distinct. The explanatory shorthand is: \textbf{the assurance capability is the QA service, not the Job Foreman.}



\par\medskip\textbf{Semantic \& Uncertainty Boundaries}\par


The following distinctions constrain this requirement; they do not define new information types or an assurance outcome taxonomy:



\begin{quote}\small\ttfamily
Provenance ≠ Assurance\\
Audit information ≠ Assurance\\
Evidence ≠ Assurance conclusion\\
Operational progression ≠ Assurance\\
Workflow completion ≠ Assurance\\
Business Outcome ≠ Assurance Outcome\\
Operational uncertainty ≠ Insufficient assurance evidence\\
Independent assurance ≠ External assurance\\
Assurance finding ≠ Operational response\\
Assurance responsibility ≠ Operational management responsibility\\

\end{quote}


\hyperref[element:fc732b64fc564127b73174e28d9ccce8]{\texttt{\seqsplit{REQ-FND-004}}} and \hyperref[element:ef8abbaec37e4890ad78b9f263e67906]{AX-15} continue to govern uncertainty about an operation's state, effect or outcome, including an assurance activity's own execution where applicable. Separately, an assurance activity may execute correctly yet lack sufficient evidence to establish a conclusion about its subject. This requirement preserves that insufficiency without treating it as an uncertain execution outcome, execution failure, assurance or non-assurance. Neither existing operational rule is rewritten or extended into an assurance-evidence sufficiency rule.



\par\medskip\textbf{Motivation Relationships}\par


These are branching relationships, not a sequential motivation pipeline. The independence, management boundary and evidence-insufficiency obligations are newly approved intent for reconciliation, not meanings retrospectively attributed to the reused elements.



\begin{longtable}{>{\raggedright\arraybackslash}p{\dimexpr\linewidth/3-2\tabcolsep\relax}>{\raggedright\arraybackslash}p{\dimexpr\linewidth/3-2\tabcolsep\relax}>{\raggedright\arraybackslash}p{\dimexpr\linewidth/3-2\tabcolsep\relax}}
Existing element & Relationship to \texttt{\seqsplit{REQ-FND-005}} & Semantic limit \\
\hyperref[element:e76216c494fa4500bb80eda6f83513c1]{Regional Health Network Operator} & Compliance demonstration and forensic-review concerns motivate the requirement directly. & No new Operator-to-Privacy-Driver relationship or statutory independence mandate is inferred. \\
\hyperref[element:c20aa17ce4aa4a58be5eb3cff2ecdd1d]{Platform Operations \& Integration Engineers} & Incident review and deterministic verification concerns motivate the requirement directly. & Operational diagnosis and monitoring remain distinct from independent assurance. \\
\hyperref[element:1dbf34b452b6497f914c704e037aec24]{Statutory Health Information Privacy \& Protection Driver} & Supports the requirement's need to evaluate handling against applicable obligations and controls. & Privacy supplies one valid motivation, not the entire assurance scope or independence obligation. \\
\hyperref[element:75519ed8f7de4a619cc9065f525d703a]{O3} & The requirement supports O3's existing demonstrability and accountability purpose. & Textual contribution only; no new causal handover, intermediate Goal or assurance-outcome equivalence. \\
\hyperref[element:0ed6fb9fcc9547659ce76b917881274e]{CST-EXT-001} & Bounds assurance concerning health-information privacy and protection by applicable obligations. & No universal external requirement for independent assurance is claimed. \\
\hyperref[element:952c3b6148ce4fce94b69f99a0c2fec7]{AX-06} & Co-governs interpretation of evidence authority and credibility independently of technical state. & Information authority does not establish assessor independence. \\
\hyperref[element:5de2ae2c3b794c2d9f3a14efd51d564f]{AX-07} & Co-governs security context and policy for assurance activity. & Security enforcement does not establish independent conclusion authority. \\
\hyperref[element:21d6864657964ba5b03e3f6b5eb29cd5]{AX-08} & Co-governs meaningful provenance and audit evidence used in assurance. & Preserving evidence does not perform assurance or establish a conclusion. \\
\hyperref[element:d3fc6025dc7646d1b2020b51f7a44bfa]{AX-14} & Co-governs preservation of the requirement's semantic distinctions. & Distinct meanings do not prescribe separate technologies or information types. \\
\end{longtable}


\hyperref[element:9a5cb9e8876844e58e660f909f7df9fc]{AX-09} remains an evidence-selection and retention boundary: observations and progression state are not automatically provenance or audit evidence. \hyperref[element:eaabcdb96aad4f32a050c94f9ecec049]{\texttt{\seqsplit{REQ-FND-002}}} can supply observable operational facts; it is a complementary evidence relationship, not an independence parent or assurance conclusion. \texttt{\seqsplit{REQ-FND-004}} / AX-15 have the complementary uncertainty relationship described above. \hyperref[element:8882573df7ce43c6bc0f66da0a676139]{AX-17} governs this explicit reconciliation and preservation of unresolved downstream derivation; it is not a source of the real-world assurance need.



\par\medskip\textbf{Downstream Derivation Boundary}\par


The requirement supplies upstream motivation from which Strategy responsibility, Business assurance behaviour and Information Architecture may subsequently be derived, followed by Application Architecture and an eventual solution. Those relationships remain unestablished here. No capability allocation, workflow specification, solution construct, outcome values, evidence model, execution mechanism or deployment boundary is defined by this requirement.


{\footnotesize \href{../../../../continuum/motivation/requirement/bbb7438e78604308abba93afc14f4e2e.md}{Canonical record: bbb7438e78604308abba93afc14f4e2e}; full qualifications and provenance.}\par
\section*{I.5.65 SEC-001 — Trusted Authentication}
\phantomsection
\label{navigation:I.5.65}
\addcontentsline{toc}{section}{I.5.65 SEC-001 — Trusted Authentication}

\label{element:e4998b0a132c41afbe610b11869c3297}
\index{SEC-001 — Trusted Authentication (I.5.65)}
{\small Requirement · Historical — current owning-domain standing not inferred · historical\_reference}\par
\href{../../../../continuum/motivation/requirement/e4998b0a132c41afbe610b11869c3297.md}{Historical canonical definition and provenance}.\par
\section*{I.5.66 SEC-002 — Default-Deny Authorisation}
\phantomsection
\label{navigation:I.5.66}
\addcontentsline{toc}{section}{I.5.66 SEC-002 — Default-Deny Authorisation}

\label{element:3bdecb66e3f34110aa0871ba34091bfd}
\index{SEC-002 — Default-Deny Authorisation (I.5.66)}
{\small Requirement · Historical — current owning-domain standing not inferred · historical\_reference}\par
\href{../../../../continuum/motivation/requirement/3bdecb66e3f34110aa0871ba34091bfd.md}{Historical canonical definition and provenance}.\par
\section*{I.5.67 SEC-003 — Authentication and Authorisation Separation}
\phantomsection
\label{navigation:I.5.67}
\addcontentsline{toc}{section}{I.5.67 SEC-003 — Authentication and Authorisation Separation}

\label{element:88f2ace9a07d4ed7b5a54a4af93af3c9}
\index{SEC-003 — Authentication and Authorisation Separation (I.5.67)}
{\small Requirement · Historical — current owning-domain standing not inferred · historical\_reference}\par
\href{../../../../continuum/motivation/requirement/88f2ace9a07d4ed7b5a54a4af93af3c9.md}{Historical canonical definition and provenance}.\par
\section*{I.5.68 SEC-004 — Contextual Resource Authorisation}
\phantomsection
\label{navigation:I.5.68}
\addcontentsline{toc}{section}{I.5.68 SEC-004 — Contextual Resource Authorisation}

\label{element:1f6822aef4bc44548c6708fcc018e4ba}
\index{SEC-004 — Contextual Resource Authorisation (I.5.68)}
{\small Requirement · Historical — current owning-domain standing not inferred · historical\_reference}\par
\href{../../../../continuum/motivation/requirement/1f6822aef4bc44548c6708fcc018e4ba.md}{Historical canonical definition and provenance}.\par
\section*{I.5.69 SEC-005 — End-to-End Security Context}
\phantomsection
\label{navigation:I.5.69}
\addcontentsline{toc}{section}{I.5.69 SEC-005 — End-to-End Security Context}

\label{element:d70e0a51440640b1bd17562b6ea4ff98}
\index{SEC-005 — End-to-End Security Context (I.5.69)}
{\small Requirement · Historical — current owning-domain standing not inferred · historical\_reference}\par
\href{../../../../continuum/motivation/requirement/d70e0a51440640b1bd17562b6ea4ff98.md}{Historical canonical definition and provenance}.\par
\section*{I.5.70 SEC-006 — Initiating and Executing Identity}
\phantomsection
\label{navigation:I.5.70}
\addcontentsline{toc}{section}{I.5.70 SEC-006 — Initiating and Executing Identity}

\label{element:5eaac47ebef94bbb818e4ba751601ab1}
\index{SEC-006 — Initiating and Executing Identity (I.5.70)}
{\small Requirement · Historical — current owning-domain standing not inferred · historical\_reference}\par
\href{../../../../continuum/motivation/requirement/5eaac47ebef94bbb818e4ba751601ab1.md}{Historical canonical definition and provenance}.\par
\section*{I.5.71 SEC-007 — Anti-Spoofing}
\phantomsection
\label{navigation:I.5.71}
\addcontentsline{toc}{section}{I.5.71 SEC-007 — Anti-Spoofing}

\label{element:8996bdd1790e4efeb78bedb85269c916}
\index{SEC-007 — Anti-Spoofing (I.5.71)}
{\small Requirement · Historical — current owning-domain standing not inferred · historical\_reference}\par
\href{../../../../continuum/motivation/requirement/8996bdd1790e4efeb78bedb85269c916.md}{Historical canonical definition and provenance}.\par
\section*{I.5.72 SEC-008 — Security Metadata Preservation}
\phantomsection
\label{navigation:I.5.72}
\addcontentsline{toc}{section}{I.5.72 SEC-008 — Security Metadata Preservation}

\label{element:c037846f015645728b4afe4d011ae0cd}
\index{SEC-008 — Security Metadata Preservation (I.5.72)}
{\small Requirement · Historical — current owning-domain standing not inferred · historical\_reference}\par
\href{../../../../continuum/motivation/requirement/c037846f015645728b4afe4d011ae0cd.md}{Historical canonical definition and provenance}.\par
\section*{I.5.73 SEC-009 — Least Privilege}
\phantomsection
\label{navigation:I.5.73}
\addcontentsline{toc}{section}{I.5.73 SEC-009 — Least Privilege}

\label{element:f760e0cc08354fc88c8da049958b45ed}
\index{SEC-009 — Least Privilege (I.5.73)}
{\small Requirement · Historical — current owning-domain standing not inferred · historical\_reference}\par
\href{../../../../continuum/motivation/requirement/f760e0cc08354fc88c8da049958b45ed.md}{Historical canonical definition and provenance}.\par
\section*{I.5.74 SEC-010 — PHI-Safe Operational Logging}
\phantomsection
\label{navigation:I.5.74}
\addcontentsline{toc}{section}{I.5.74 SEC-010 — PHI-Safe Operational Logging}

\label{element:4a078bd44e874a48abc650f87d64f16f}
\index{SEC-010 — PHI-Safe Operational Logging (I.5.74)}
{\small Requirement · Historical — current owning-domain standing not inferred · historical\_reference}\par
\href{../../../../continuum/motivation/requirement/4a078bd44e874a48abc650f87d64f16f.md}{Historical canonical definition and provenance}.\par
\section*{I.5.75 UX-001 — Consistent Iris Experience}
\phantomsection
\label{navigation:I.5.75}
\addcontentsline{toc}{section}{I.5.75 UX-001 — Consistent Iris Experience}

\label{element:6bea2694012442a39034a6cbfcc79e90}
\index{UX-001 — Consistent Iris Experience (I.5.75)}
{\small Requirement · Historical — current owning-domain standing not inferred · historical\_reference}\par
\href{../../../../continuum/motivation/requirement/6bea2694012442a39034a6cbfcc79e90.md}{Historical canonical definition and provenance}.\par
\section*{I.5.76 UX-002 — Information-Dense Operational Design}
\phantomsection
\label{navigation:I.5.76}
\addcontentsline{toc}{section}{I.5.76 UX-002 — Information-Dense Operational Design}

\label{element:73354d29efa74895afebec8437c875e8}
\index{UX-002 — Information-Dense Operational Design (I.5.76)}
{\small Requirement · Historical — current owning-domain standing not inferred · historical\_reference}\par
\href{../../../../continuum/motivation/requirement/73354d29efa74895afebec8437c875e8.md}{Historical canonical definition and provenance}.\par
\section*{I.5.77 UX-003 — Honest System State}
\phantomsection
\label{navigation:I.5.77}
\addcontentsline{toc}{section}{I.5.77 UX-003 — Honest System State}

\label{element:752452c7367d4d4dbc1cb9ffe0bfd78f}
\index{UX-003 — Honest System State (I.5.77)}
{\small Requirement · Historical — current owning-domain standing not inferred · historical\_reference}\par
\href{../../../../continuum/motivation/requirement/752452c7367d4d4dbc1cb9ffe0bfd78f.md}{Historical canonical definition and provenance}.\par
\section*{I.5.78 UX-004 — Explicit Availability States}
\phantomsection
\label{navigation:I.5.78}
\addcontentsline{toc}{section}{I.5.78 UX-004 — Explicit Availability States}

\label{element:0e4b5219260e419d987195c9832a0bbe}
\index{UX-004 — Explicit Availability States (I.5.78)}
{\small Requirement · Historical — current owning-domain standing not inferred · historical\_reference}\par
\href{../../../../continuum/motivation/requirement/0e4b5219260e419d987195c9832a0bbe.md}{Historical canonical definition and provenance}.\par
\section*{I.5.79 UX-005 — Patient Search and Workspace}
\phantomsection
\label{navigation:I.5.79}
\addcontentsline{toc}{section}{I.5.79 UX-005 — Patient Search and Workspace}

\label{element:32c55c722ae147329784976793abd458}
\index{UX-005 — Patient Search and Workspace (I.5.79)}
{\small Requirement · Historical — current owning-domain standing not inferred · historical\_reference}\par
\href{../../../../continuum/motivation/requirement/32c55c722ae147329784976793abd458.md}{Historical canonical definition and provenance}.\par
\section*{I.5.80 UX-006 — Longitudinal Clinical Navigation}
\phantomsection
\label{navigation:I.5.80}
\addcontentsline{toc}{section}{I.5.80 UX-006 — Longitudinal Clinical Navigation}

\label{element:09318a29f8b747c8ac6b0e276e2febb4}
\index{UX-006 — Longitudinal Clinical Navigation (I.5.80)}
{\small Requirement · Historical — current owning-domain standing not inferred · historical\_reference}\par
\href{../../../../continuum/motivation/requirement/09318a29f8b747c8ac6b0e276e2febb4.md}{Historical canonical definition and provenance}.\par
\section*{I.5.81 UX-007 — Persistent Patient Context}
\phantomsection
\label{navigation:I.5.81}
\addcontentsline{toc}{section}{I.5.81 UX-007 — Persistent Patient Context}

\label{element:873b8b2c9a514ab387833beb5552feda}
\index{UX-007 — Persistent Patient Context (I.5.81)}
{\small Requirement · Historical — current owning-domain standing not inferred · historical\_reference}\par
\href{../../../../continuum/motivation/requirement/873b8b2c9a514ab387833beb5552feda.md}{Historical canonical definition and provenance}.\par
\section*{I.5.82 UX-008 — Provenance-Aware Clinical Presentation}
\phantomsection
\label{navigation:I.5.82}
\addcontentsline{toc}{section}{I.5.82 UX-008 — Provenance-Aware Clinical Presentation}

\label{element:bb298bafe5894d9b9c057b896f669c06}
\index{UX-008 — Provenance-Aware Clinical Presentation (I.5.82)}
{\small Requirement · Historical — current owning-domain standing not inferred · historical\_reference}\par
\href{../../../../continuum/motivation/requirement/bb298bafe5894d9b9c057b896f669c06.md}{Historical canonical definition and provenance}.\par
\section*{I.5.83 UX-009 — Initial Clinical Scope}
\phantomsection
\label{navigation:I.5.83}
\addcontentsline{toc}{section}{I.5.83 UX-009 — Initial Clinical Scope}

\label{element:807861850fa54d6c849e3e26d7172b9c}
\index{UX-009 — Initial Clinical Scope (I.5.83)}
{\small Requirement · Historical — current owning-domain standing not inferred · historical\_reference}\par
\href{../../../../continuum/motivation/requirement/807861850fa54d6c849e3e26d7172b9c.md}{Historical canonical definition and provenance}.\par
\chapter*{I.6 Constraints}
\phantomsection
\label{navigation:I.6}
\addcontentsline{toc}{chapter}{I.6 Constraints}

\section*{I.6.1 AuthorisationRole Criteria Are Additive}
\phantomsection
\label{navigation:I.6.1}
\addcontentsline{toc}{section}{I.6.1 AuthorisationRole Criteria Are Additive}

\label{element:d5cfaf6f8c224d1b9769b1edd7c27591}
\index{AuthorisationRole Criteria Are Additive (I.6.1)}
{\small Constraint · Adjudicated Domain05 direction — U12/U15/U16; holistic result acceptance pending · current}\par

AuthorisationRole containment SHALL NOT be interpreted as requiring all contained Criteria to be simultaneously satisfied.



This assertion is normalised from the governed DesignNote under U12/U15/U16 and the ARB release Praxis. Original candidate/release qualifications and separately deferred implementation details remain in force.



Architectural subject: \hyperref[element:899d34ec491948b68aad09e507dea6d1]{Themis}.


{\footnotesize \href{../../../../continuum/motivation/constraint/d5cfaf6f8c224d1b9769b1edd7c27591.md}{Canonical record: d5cfaf6f8c224d1b9769b1edd7c27591}; full qualifications and provenance.}\par
\section*{I.6.2 Authoritative AuditEvents Are Immutable}
\phantomsection
\label{navigation:I.6.2}
\addcontentsline{toc}{section}{I.6.2 Authoritative AuditEvents Are Immutable}

\label{element:21eb20cdc600411cb71fbfe2abce4dbe}
\index{Authoritative AuditEvents Are Immutable (I.6.2)}
{\small Constraint · Adjudicated Domain05 direction — U07/U15; holistic result acceptance pending · current}\par

\begin{quote}


\textbf{An existing authoritative AuditEvent SHALL NOT be corrected, amended, replaced or overwritten.}



\end{quote}


This assertion is normalised from the governed DesignNote under U07/U15 and the ARB release Praxis. Original candidate/release qualifications and separately deferred implementation details remain in force.



Architectural subject: \hyperref[element:fefda110e2274884ad58906733606efd]{Kleio}.


{\footnotesize \href{../../../../continuum/motivation/constraint/21eb20cdc600411cb71fbfe2abce4dbe.md}{Canonical record: 21eb20cdc600411cb71fbfe2abce4dbe}; full qualifications and provenance.}\par
\section*{I.6.3 Authority Forwarding Requires FORWARD Authorisation}
\phantomsection
\label{navigation:I.6.3}
\addcontentsline{toc}{section}{I.6.3 Authority Forwarding Requires FORWARD Authorisation}

\label{element:5aed12f6540445ed85c1866adbef93af}
\index{Authority Forwarding Requires FORWARD Authorisation (I.6.3)}
{\small Constraint · Adjudicated Domain05 direction — U12/U15/U16; holistic result acceptance pending · current}\par

Forwarding SHALL itself be subject to applicable FORWARD authorisation.



This assertion is normalised from the governed DesignNote under U12/U15/U16 and the ARB release Praxis. Original candidate/release qualifications and separately deferred implementation details remain in force.



Architectural subject: \hyperref[element:899d34ec491948b68aad09e507dea6d1]{Themis}.


{\footnotesize \href{../../../../continuum/motivation/constraint/5aed12f6540445ed85c1866adbef93af.md}{Canonical record: 5aed12f6540445ed85c1866adbef93af}; full qualifications and provenance.}\par
\section*{I.6.4 Configuration and Governance Has Coherent Authority}
\phantomsection
\label{navigation:I.6.4}
\addcontentsline{toc}{section}{I.6.4 Configuration and Governance Has Coherent Authority}

\label{element:7cdbc752ca2e4ebf9ee5acff84fd7400}
\index{Configuration and Governance Has Coherent Authority (I.6.4)}
{\small Constraint · Adjudicated Domain05 direction — U11; holistic result acceptance pending · current}\par

\textbf{Harmonia Configuration \& Governance (C\&G)} SHALL remain a single coherent governed information space and SHALL NOT be partitioned in a manner that creates independent or divergent governance authority.



This assertion is normalised from the governed DesignNote under U11 and the ARB release Praxis. Original candidate/release qualifications and separately deferred implementation details remain in force.



Architectural subject: \hyperref[element:7985e647bf9e4ca3ba70faf9e3e4ae02]{Mnemosyne}.


{\footnotesize \href{../../../../continuum/motivation/constraint/7cdbc752ca2e4ebf9ee5acff84fd7400.md}{Canonical record: 7cdbc752ca2e4ebf9ee5acff84fd7400}; full qualifications and provenance.}\par
\section*{I.6.5 CST-EXT-001: Applicable Health-Information Privacy and Data-Protection Obligations}
\phantomsection
\label{navigation:I.6.5}
\addcontentsline{toc}{section}{I.6.5 CST-EXT-001: Applicable Health-Information Privacy and Data-Protection Obligations}

\label{element:0ed6fb9fcc9547659ce76b917881274e}
\index{CST-EXT-001: Applicable Health-Information Privacy and Data-Protection Obligations (I.6.5)}
{\small Constraint · Accepted · current}\par

\par\medskip\textbf{Categorical Definition}\par


\begin{quote}


\textbf{Harmonia must operate within the health-information privacy, confidentiality, disclosure, retention, and data-protection obligations applicable to the jurisdiction and deployment context.}



\end{quote}


\par\medskip\textbf{Context \& Architectural Boundary}\par


Health data is categorized as hypersensitive personal data across all legal jurisdictions. Wherever Harmonia is deployed, it is bound by the statutory data protection regime governing that region. The platform must be architecturally designed to satisfy these obligations without requiring bespoke rewrites of core domain logic.



\par\medskip\textbf{Representative Instances \& Examples}\par


\begin{itemize}


\item \textbf{Australia}: The \emph{Privacy Act 1988 (Cth)}, the Australian Privacy Principles (APPs), the \emph{My Health Records Act 2012}, and state-based health records legislation (e.g., \emph{Health Records Act 2001 (Vic)}).

\item \textbf{United States}: The Health Insurance Portability and Accountability Act (HIPAA) Security and Privacy Rules.

\item \textbf{European Union}: The General Data Protection Regulation (GDPR) and European Health Data Space (EHDS) regulations.

\end{itemize}


\par\medskip\textbf{Direct Architectural Impacts}\par


\begin{itemize}


\item Requires downstream security architecture to govern access to protected health information and enforce applicable access controls (\href{\detokenize{../../../../../docs/markdown/01-motivation/requirements-constraints/master-requirements-catalogue.md}}{\texttt{\seqsplit{SEC-001 / SEC-002}}}).

\item Requires downstream operational logging policies and guardrails to protect health information from unauthorized exposure in diagnostic telemetry (\href{\detokenize{../../../../../docs/markdown/01-motivation/requirements-constraints/master-requirements-catalogue.md}}{\texttt{\seqsplit{SEC-010}}}).

\item Requires platform communications and stored data representations to apply appropriate protection mechanisms according to deployment risk and statutory requirements.

\item Informs audit and provenance requirements for demonstrating authorized handling of health information (\href{\detokenize{../../../../../docs/markdown/01-motivation/requirements-constraints/master-requirements-catalogue.md}}{\texttt{\seqsplit{PROV-001}}}).

\item Bounds Platform Driver: \href{\detokenize{../../../../../docs/markdown/01-motivation/drivers-assessments/drivers.md}}{Statutory Health Information Privacy \& Protection}.

\end{itemize}

{\footnotesize \href{../../../../continuum/motivation/constraint/0ed6fb9fcc9547659ce76b917881274e.md}{Canonical record: 0ed6fb9fcc9547659ce76b917881274e}; full qualifications and provenance.}\par
\section*{I.6.6 CST-EXT-002: Applicable National / Jurisdictional Healthcare Identifier Obligations}
\phantomsection
\label{navigation:I.6.6}
\addcontentsline{toc}{section}{I.6.6 CST-EXT-002: Applicable National / Jurisdictional Healthcare Identifier Obligations}

\label{element:edb1db6b18d84487ab48b9416f26e980}
\index{CST-EXT-002: Applicable National / Jurisdictional Healthcare Identifier Obligations (I.6.6)}
{\small Constraint · Accepted · current}\par

\par\medskip\textbf{Categorical Definition}\par


\begin{quote}


\textbf{Where a deployment participates in a national or jurisdictional healthcare identifier ecosystem, Harmonia must conform to applicable identifier representation, validation, authority, and usage obligations.}



\end{quote}


\par\medskip\textbf{Context \& Architectural Boundary}\par


Many modern nations operate statutory healthcare identifier services to uniquely identify patients, clinicians, and health organizations. Participation in these ecosystems is legally regulated: the law dictates how identifiers may be retrieved, verified, stored, cross-referenced, and disclosed. Harmonia must conform strictly to these statutory rules while maintaining internal referential integrity.



\par\medskip\textbf{Representative Instances \& Examples}\par


\begin{itemize}


\item \textbf{Australia}: The \emph{Healthcare Identifiers Act 2010 (Cth)} and HI Service Operating Rules:

\item \textbf{IHI} (Individual Healthcare Identifier): Strict prohibitions against using the IHI as a general-purpose identity token or publishing it outside authorized clinical communication channels.

\item \textbf{HPI-I} (Healthcare Provider Identifier - Individual): Validation against national registers and credentialing rules.

\item \textbf{HPI-O} (Healthcare Provider Identifier - Organisation): Association with authorized electronic service delivery endpoints.

\item \textbf{United Kingdom}: NHS Number validation and Spine integration rules.

\item \textbf{United States}: National Provider Identifier (NPI) specifications.

\end{itemize}


\par\medskip\textbf{Direct Architectural Impacts}\par


\begin{itemize}


\item Requires strict validation of identifier check digits, namespaces, and authority types before accepting clinical records (\href{\detokenize{../../../../../docs/markdown/01-motivation/requirements-constraints/foundational-requirements.md}}{\texttt{\seqsplit{REQ-FND-003}}}).

\item Dictates that external national identifiers must not be conflated with internal synthetic database keys or process identifiers.

\item Requires integration boundaries interacting with national directory or identifier services to adhere to governing access, query, and communication rules.

\item Bounds Platform Driver: \href{\detokenize{../../../../../docs/markdown/01-motivation/drivers-assessments/drivers.md}}{Patient Safety \& Identity Integrity}.

\end{itemize}

{\footnotesize \href{../../../../continuum/motivation/constraint/edb1db6b18d84487ab48b9416f26e980.md}{Canonical record: edb1db6b18d84487ab48b9416f26e980}; full qualifications and provenance.}\par
\section*{I.6.7 CST-EXT-003: Mandated External Interoperability Contracts}
\phantomsection
\label{navigation:I.6.7}
\addcontentsline{toc}{section}{I.6.7 CST-EXT-003: Mandated External Interoperability Contracts}

\label{element:cad6d41364ef4bb1ae0a2a519db11da4}
\index{CST-EXT-003: Mandated External Interoperability Contracts (I.6.7)}
{\small Constraint · Accepted · current}\par

\par\medskip\textbf{Categorical Definition}\par


\begin{quote}


\textbf{Where an external system or ecosystem mandates a protocol, information standard, transport contract, or serialization format, Harmonia must conform to that contract at the applicable interoperability boundary.}



\end{quote}


\par\medskip\textbf{Context \& Architectural Boundary}\par


Harmonia operates as an integration hub connecting hundreds of disparate health systems. The platform cannot dictate the internal wire protocols of external hospitals, laboratory analyzers, or national repositories. At its ingress and egress boundaries, Harmonia must conform strictly to the wire formats, message structures, and transport handshakes mandated by external systems.



However, in accordance with \hyperref[element:eee8c3783d964d18a7f9dcea85fcd4cf]{AX-02 (Standards at Boundary; Sovereignty Within)}, external contracts govern only the external membrane. Once an event crosses into Harmonia, the platform enforces its own canonical domain semantics.



\par\medskip\textbf{Representative Instances \& Examples}\par


\begin{itemize}


\item \textbf{HL7 v2.x over MLLP}: Minimal Lower Layer Protocol (MLLP) framing over raw TCP, utilizing pipe-and-hat (\texttt{\seqsplit{ER7}}) segment syntax for admissions (ADT), orders (ORM), and observations (ORU).

\item \textbf{HL7 FHIR RESTful APIs}: HTTP/JSON or HTTP/XML operations conforming to specific regional profiles (e.g., Australian Core AU-Base / AU-Core, US Core).

\item \textbf{Secure Messaging Delivery (SMD)}: Encrypted Web Services (WS-*) protocols mandated for inter-clinic electronic correspondence.

\item \textbf{DICOM}: Digital Imaging and Communications in Medicine protocol for radiological studies.

\end{itemize}


\par\medskip\textbf{Direct Architectural Impacts}\par


\begin{itemize}


\item Demarcates the boundary between Pylai (the protocol membrane gateways) and core internal platform engines.

\item Governs wire-level error signaling (e.g., mapping internal durable persistence failures to protocol-appropriate error responses).

\item Prevents external wire formats from polluting internal domain models.

\item Imposes on Platform Driver: \href{\detokenize{../../../../../docs/markdown/01-motivation/drivers-assessments/drivers.md}}{Clinical Interoperability Across Heterogeneous Standards}.

\end{itemize}

{\footnotesize \href{../../../../continuum/motivation/constraint/cad6d41364ef4bb1ae0a2a519db11da4.md}{Canonical record: cad6d41364ef4bb1ae0a2a519db11da4}; full qualifications and provenance.}\par
\section*{I.6.8 Dataset Access Has No Runtime Bootstrap Cycle}
\phantomsection
\label{navigation:I.6.8}
\addcontentsline{toc}{section}{I.6.8 Dataset Access Has No Runtime Bootstrap Cycle}

\label{element:c8c479dc9aad450ea58d07dec9aea1a6}
\index{Dataset Access Has No Runtime Bootstrap Cycle (I.6.8)}
{\small Constraint · Adjudicated Domain05 direction — U13/U15; holistic result acceptance pending · current}\par

\begin{quote}


\textbf{The structural contracts necessary to access and manage
Calliope-DataSets SHALL be established through Calliope-Models and the
software build/deployment lifecycle. Runtime Calliope-DataSets SHALL
provide the authoritative mutable content governed by those contracts
and SHALL NOT be required to establish the structural contracts
necessary for its own access.}



\end{quote}


This assertion is normalised from the governed DesignNote under U13/U15 and the ARB release Praxis. Original candidate/release qualifications and separately deferred implementation details remain in force.



Architectural subject: \hyperref[element:466ad8f75a954c0581b5bba090443849]{Calliope}.


{\footnotesize \href{../../../../continuum/motivation/constraint/c8c479dc9aad450ea58d07dec9aea1a6.md}{Canonical record: c8c479dc9aad450ea58d07dec9aea1a6}; full qualifications and provenance.}\par
\section*{I.6.9 Dokimasia Evidence Replaces a Separate Self AuditEvent}
\phantomsection
\label{navigation:I.6.9}
\addcontentsline{toc}{section}{I.6.9 Dokimasia Evidence Replaces a Separate Self AuditEvent}

\label{element:e1ce978457e543d5894de24e9be9ac2c}
\index{Dokimasia Evidence Replaces a Separate Self AuditEvent (I.6.9)}
{\small Constraint · Adjudicated Domain05 direction — U02/U06/U19; holistic result acceptance pending · current}\par

A Dokimasia execution SHALL NOT generate a separate AuditEvent for its own execution. The governed Dokimasia assurance record/result provides the authoritative evidence of that assurance activity.



This assertion is normalised from the governed DesignNote under U02/U06/U19 and the ARB release Praxis. Original candidate/release qualifications and separately deferred implementation details remain in force.



Architectural subject: \hyperref[element:fb3f37a6ac614415ac486a6de9d35994]{Energeia}.


{\footnotesize \href{../../../../continuum/motivation/constraint/e1ce978457e543d5894de24e9be9ac2c.md}{Canonical record: e1ce978457e543d5894de24e9be9ac2c}; full qualifications and provenance.}\par
\section*{I.6.10 Dokimasia Has a Finite Assurance Boundary}
\phantomsection
\label{navigation:I.6.10}
\addcontentsline{toc}{section}{I.6.10 Dokimasia Has a Finite Assurance Boundary}

\label{element:3ba73277abf146728ca8c1f5846227d7}
\index{Dokimasia Has a Finite Assurance Boundary (I.6.10)}
{\small Constraint · Adjudicated Domain05 direction — U02/U06/U19; holistic result acceptance pending · current}\par

A Dokimasia execution SHALL NOT itself be subject to monitoring, assurance or assessment by another Dokimasia execution. This establishes a finite assurance boundary.



This assertion is normalised from the governed DesignNote under U02/U06/U19 and the ARB release Praxis. Original candidate/release qualifications and separately deferred implementation details remain in force.



Architectural subject: \hyperref[element:fb3f37a6ac614415ac486a6de9d35994]{Energeia}.


{\footnotesize \href{../../../../continuum/motivation/constraint/3ba73277abf146728ca8c1f5846227d7.md}{Canonical record: 3ba73277abf146728ca8c1f5846227d7}; full qualifications and provenance.}\par
\section*{I.6.11 Every Petasos Queue Has a Defined Functional Purpose}
\phantomsection
\label{navigation:I.6.11}
\addcontentsline{toc}{section}{I.6.11 Every Petasos Queue Has a Defined Functional Purpose}

\label{element:6fa0f390bfa140aa8bb6acc56f1a5ec6}
\index{Every Petasos Queue Has a Defined Functional Purpose (I.6.11)}
{\small Constraint · Adjudicated Domain05 direction — U09; holistic result acceptance pending · current}\par

A Petasos queue SHALL exist for a defined Harmonia functional purpose.



This assertion is normalised from the governed DesignNote under U09 and the ARB release Praxis. Original candidate/release qualifications and separately deferred implementation details remain in force.



Architectural subject: \hyperref[element:195e5f10db8e419391354197b5d7a8cc]{Petasos}.


{\footnotesize \href{../../../../continuum/motivation/constraint/6fa0f390bfa140aa8bb6acc56f1a5ec6.md}{Canonical record: 6fa0f390bfa140aa8bb6acc56f1a5ec6}; full qualifications and provenance.}\par
\section*{I.6.12 Harmonisation Does Not Resume the Originating Ergon}
\phantomsection
\label{navigation:I.6.12}
\addcontentsline{toc}{section}{I.6.12 Harmonisation Does Not Resume the Originating Ergon}

\label{element:1739ad181a754d01ae07f94de49cd65e}
\index{Harmonisation Does Not Resume the Originating Ergon (I.6.12)}
{\small Constraint · Adjudicated Domain05 direction — U10/U18; holistic result acceptance pending · current}\par

Harmonisation SHALL NOT be implemented by resuming inside the originating Ergon.



This assertion is normalised from the governed DesignNote under U10/U18 and the ARB release Praxis. Original candidate/release qualifications and separately deferred implementation details remain in force.



Architectural subject: \hyperref[element:1af8a56358f7429fabfb57be66542a99]{Mneme}.


{\footnotesize \href{../../../../continuum/motivation/constraint/1739ad181a754d01ae07f94de49cd65e.md}{Canonical record: 1739ad181a754d01ae07f94de49cd65e}; full qualifications and provenance.}\par
\section*{I.6.13 Harmonisation Reuses the Common Execution Framework}
\phantomsection
\label{navigation:I.6.13}
\addcontentsline{toc}{section}{I.6.13 Harmonisation Reuses the Common Execution Framework}

\label{element:59e1a236500240609a78b6529fe8e358}
\index{Harmonisation Reuses the Common Execution Framework (I.6.13)}
{\small Constraint · Adjudicated Domain05 direction — U10/U18; holistic result acceptance pending · current}\par

The RHM SHALL NOT duplicate or embed a second execution framework within Mneme.



This assertion is normalised from the governed DesignNote under U10/U18 and the ARB release Praxis. Original candidate/release qualifications and separately deferred implementation details remain in force.



Architectural subject: \hyperref[element:1af8a56358f7429fabfb57be66542a99]{Mneme}.


{\footnotesize \href{../../../../continuum/motivation/constraint/59e1a236500240609a78b6529fe8e358.md}{Canonical record: 59e1a236500240609a78b6529fe8e358}; full qualifications and provenance.}\par
\section*{I.6.14 Host Components Do Not Reinterpret Themis Policy}
\phantomsection
\label{navigation:I.6.14}
\addcontentsline{toc}{section}{I.6.14 Host Components Do Not Reinterpret Themis Policy}

\label{element:80e5d019948b4c16acb510547d7a3ac6}
\index{Host Components Do Not Reinterpret Themis Policy (I.6.14)}
{\small Constraint · Adjudicated Domain05 direction — U08/U12/U15; holistic result acceptance pending · current}\par

A Harmonia component SHALL NOT implement an independent interpretation of Themis policy in place of the Themis PEP.



This assertion is normalised from the governed DesignNote under U08/U12/U15 and the ARB release Praxis. Original candidate/release qualifications and separately deferred implementation details remain in force.



Architectural subject: \hyperref[element:899d34ec491948b68aad09e507dea6d1]{Themis}.


{\footnotesize \href{../../../../continuum/motivation/constraint/80e5d019948b4c16acb510547d7a3ac6.md}{Canonical record: 80e5d019948b4c16acb510547d7a3ac6}; full qualifications and provenance.}\par
\section*{I.6.15 Independent Audit Targets Are Not Assumed Replicas}
\phantomsection
\label{navigation:I.6.15}
\addcontentsline{toc}{section}{I.6.15 Independent Audit Targets Are Not Assumed Replicas}

\label{element:e84fc7e020474eda99a0e32def0a6371}
\index{Independent Audit Targets Are Not Assumed Replicas (I.6.15)}
{\small Constraint · Adjudicated Domain05 direction — U07/U15; holistic result acceptance pending · current}\par

These targets SHALL NOT be described as replicas unless and until convergence/replication semantics are actually established. For R1.x/R2.x they are \textbf{independent authoritative audit persistence targets}.



This assertion is normalised from the governed DesignNote under U07/U15 and the ARB release Praxis. Original candidate/release qualifications and separately deferred implementation details remain in force.



Architectural subject: \hyperref[element:fefda110e2274884ad58906733606efd]{Kleio}.


{\footnotesize \href{../../../../continuum/motivation/constraint/e84fc7e020474eda99a0e32def0a6371.md}{Canonical record: e84fc7e020474eda99a0e32def0a6371}; full qualifications and provenance.}\par
\section*{I.6.16 Information Topic and Messaging Destination Remain Distinct}
\phantomsection
\label{navigation:I.6.16}
\addcontentsline{toc}{section}{I.6.16 Information Topic and Messaging Destination Remain Distinct}

\label{element:ba909d70784349cd94b8adb4d6a9ab69}
\index{Information Topic and Messaging Destination Remain Distinct (I.6.16)}
{\small Constraint · Adjudicated Domain05 direction — U05/U09/U18; holistic result acceptance pending · current}\par

Care must be taken not to confuse Petasos queue/destination classification with InformationUnit semantic Topic.



This assertion is normalised from the governed DesignNote under U05/U09/U18 and the ARB release Praxis. Original candidate/release qualifications and separately deferred implementation details remain in force.



Architectural subject: \hyperref[element:1af8a56358f7429fabfb57be66542a99]{Mneme}.


{\footnotesize \href{../../../../continuum/motivation/constraint/ba909d70784349cd94b8adb4d6a9ab69.md}{Canonical record: ba909d70784349cd94b8adb4d6a9ab69}; full qualifications and provenance.}\par
\section*{I.6.17 Ingress Context Loading Does Not Interpret Raw Payload Semantics}
\phantomsection
\label{navigation:I.6.17}
\addcontentsline{toc}{section}{I.6.17 Ingress Context Loading Does Not Interpret Raw Payload Semantics}

\label{element:f5af90e09e9a422d9b1c9a7f3c44090f}
\index{Ingress Context Loading Does Not Interpret Raw Payload Semantics (I.6.17)}
{\small Constraint · Adjudicated Domain05 direction — U14/U15; holistic result acceptance pending · current}\par

It must \textbf{not} perform interface-specific semantic interpretation or transformation of the raw payload. That belongs to the paired developer-written \texttt{\seqsplit{Ergo.BusinessLogic}}.



This assertion is normalised from the governed DesignNote under U14/U15 and the ARB release Praxis. Original candidate/release qualifications and separately deferred implementation details remain in force.



Architectural subject: \hyperref[element:e3e4147fea3640198dc251ed41ee854d]{Pylai}.


{\footnotesize \href{../../../../continuum/motivation/constraint/f5af90e09e9a422d9b1c9a7f3c44090f.md}{Canonical record: f5af90e09e9a422d9b1c9a7f3c44090f}; full qualifications and provenance.}\par
\section*{I.6.18 Internal Mneme Responsibilities Are Not Automatically Exposed Services}
\phantomsection
\label{navigation:I.6.18}
\addcontentsline{toc}{section}{I.6.18 Internal Mneme Responsibilities Are Not Automatically Exposed Services}

\label{element:f0a85fcf515946e2b7c88543061ebac1}
\index{Internal Mneme Responsibilities Are Not Automatically Exposed Services (I.6.18)}
{\small Constraint · Adjudicated Domain05 direction — U10/U18; holistic result acceptance pending · current}\par

Other Mneme responsibilities described in this note are supporting/internal responsibilities and SHALL NOT automatically be interpreted as additional exposed services.



This assertion is normalised from the governed DesignNote under U10/U18 and the ARB release Praxis. Original candidate/release qualifications and separately deferred implementation details remain in force.



Architectural subject: \hyperref[element:1af8a56358f7429fabfb57be66542a99]{Mneme}.


{\footnotesize \href{../../../../continuum/motivation/constraint/f0a85fcf515946e2b7c88543061ebac1.md}{Canonical record: f0a85fcf515946e2b7c88543061ebac1}; full qualifications and provenance.}\par
\section*{I.6.19 Journey Observations Remain Distinct from Assurance Inference}
\phantomsection
\label{navigation:I.6.19}
\addcontentsline{toc}{section}{I.6.19 Journey Observations Remain Distinct from Assurance Inference}

\label{element:ae2fa71039034290b019eb1cc6eb1996}
\index{Journey Observations Remain Distinct from Assurance Inference (I.6.19)}
{\small Constraint · Adjudicated Domain05 direction — U06/U09/U19; holistic result acceptance pending · current}\par

Journey SHALL remain a representation of what Harmonia observed. Assurance interpretation, including AI-assisted interpretation, SHALL NOT rewrite Journey into a representation of what an assurance process inferred happened.



This assertion is normalised from the governed DesignNote under U06/U09/U19 and the ARB release Praxis. Original candidate/release qualifications and separately deferred implementation details remain in force.



Architectural subject: \hyperref[element:fb3f37a6ac614415ac486a6de9d35994]{Energeia}.


{\footnotesize \href{../../../../continuum/motivation/constraint/ae2fa71039034290b019eb1cc6eb1996.md}{Canonical record: ae2fa71039034290b019eb1cc6eb1996}; full qualifications and provenance.}\par
\section*{I.6.20 Kleio Constructs AuditEvent Representations for Clients}
\phantomsection
\label{navigation:I.6.20}
\addcontentsline{toc}{section}{I.6.20 Kleio Constructs AuditEvent Representations for Clients}

\label{element:1aa42a9aa5a644bc9072522e0a123352}
\index{Kleio Constructs AuditEvent Representations for Clients (I.6.20)}
{\small Constraint · Adjudicated Domain05 direction — U07/U15; holistic result acceptance pending · current}\par

\begin{quote}


\textbf{Clients SHALL NOT be required to construct FHIR AuditEvent resources. Clients submit auditable occurrences through the Kleio Audit Capture contract. Kleio is responsible for constructing and managing the authoritative FHIR AuditEvent representation.}



\end{quote}


This assertion is normalised from the governed DesignNote under U07/U15 and the ARB release Praxis. Original candidate/release qualifications and separately deferred implementation details remain in force.



Architectural subject: \hyperref[element:fefda110e2274884ad58906733606efd]{Kleio}.


{\footnotesize \href{../../../../continuum/motivation/constraint/1aa42a9aa5a644bc9072522e0a123352.md}{Canonical record: 1aa42a9aa5a644bc9072522e0a123352}; full qualifications and provenance.}\par
\section*{I.6.21 Local PEP Enforcement Is Constrained and Deterministic}
\phantomsection
\label{navigation:I.6.21}
\addcontentsline{toc}{section}{I.6.21 Local PEP Enforcement Is Constrained and Deterministic}

\label{element:8c9ef077cb5d4cc9aacb55d743b1af92}
\index{Local PEP Enforcement Is Constrained and Deterministic (I.6.21)}
{\small Constraint · Adjudicated Domain05 direction — U12/U15/U16; holistic result acceptance pending · current}\par

The PEP is intentionally constrained and deterministic. It SHALL NOT become a general-purpose policy-language interpreter.



This assertion is normalised from the governed DesignNote under U12/U15/U16 and the ARB release Praxis. Original candidate/release qualifications and separately deferred implementation details remain in force.



Architectural subject: \hyperref[element:899d34ec491948b68aad09e507dea6d1]{Themis}.


{\footnotesize \href{../../../../continuum/motivation/constraint/8c9ef077cb5d4cc9aacb55d743b1af92.md}{Canonical record: 8c9ef077cb5d4cc9aacb55d743b1af92}; full qualifications and provenance.}\par
\section*{I.6.22 Mneme Is Insulated from Ponos Implementation Internals}
\phantomsection
\label{navigation:I.6.22}
\addcontentsline{toc}{section}{I.6.22 Mneme Is Insulated from Ponos Implementation Internals}

\label{element:c30ae90fc19b465982b28d61ac1edca5}
\index{Mneme Is Insulated from Ponos Implementation Internals (I.6.22)}
{\small Constraint · Adjudicated Domain05 direction — U10/U18; holistic result acceptance pending · current}\par

The cross-component dependency SHALL remain narrow and architectural rather than coupling Mneme to Ponos implementation internals.



This assertion is normalised from the governed DesignNote under U10/U18 and the ARB release Praxis. Original candidate/release qualifications and separately deferred implementation details remain in force.



Architectural subject: \hyperref[element:1af8a56358f7429fabfb57be66542a99]{Mneme}.


{\footnotesize \href{../../../../continuum/motivation/constraint/c30ae90fc19b465982b28d61ac1edca5.md}{Canonical record: c30ae90fc19b465982b28d61ac1edca5}; full qualifications and provenance.}\par
\section*{I.6.23 Negative Criteria Require a Demonstrated Business Requirement}
\phantomsection
\label{navigation:I.6.23}
\addcontentsline{toc}{section}{I.6.23 Negative Criteria Require a Demonstrated Business Requirement}

\label{element:820e647c43614052aa25e385d6fe1dac}
\index{Negative Criteria Require a Demonstrated Business Requirement (I.6.23)}
{\small Constraint · Adjudicated Domain05 direction — U12/U15/U16; holistic result acceptance pending · current}\par

The Candidate Architecture does not introduce explicit negative/DENY Criteria. Such a mechanism SHALL NOT be introduced unless a demonstrated business requirement justifies the additional policy semantics.



This assertion is normalised from the governed DesignNote under U12/U15/U16 and the ARB release Praxis. Original candidate/release qualifications and separately deferred implementation details remain in force.



Architectural subject: \hyperref[element:899d34ec491948b68aad09e507dea6d1]{Themis}.


{\footnotesize \href{../../../../continuum/motivation/constraint/820e647c43614052aa25e385d6fe1dac.md}{Canonical record: 820e647c43614052aa25e385d6fe1dac}; full qualifications and provenance.}\par
\section*{I.6.24 New Policy Context Requires Authoritative PDP Processing}
\phantomsection
\label{navigation:I.6.24}
\addcontentsline{toc}{section}{I.6.24 New Policy Context Requires Authoritative PDP Processing}

\label{element:80629d4c229e4b7095d011bf77d704aa}
\index{New Policy Context Requires Authoritative PDP Processing (I.6.24)}
{\small Constraint · Adjudicated Domain05 direction — U08/U12/U15; holistic result acceptance pending · current}\par

If the PDP cannot establish the required new PEP-Context, the activity SHALL NOT proceed beyond the point requiring that context and must wait until authoritative policy context can be established.



This assertion is normalised from the governed DesignNote under U08/U12/U15 and the ARB release Praxis. Original candidate/release qualifications and separately deferred implementation details remain in force.



Architectural subject: \hyperref[element:899d34ec491948b68aad09e507dea6d1]{Themis}.


{\footnotesize \href{../../../../continuum/motivation/constraint/80629d4c229e4b7095d011bf77d704aa.md}{Canonical record: 80629d4c229e4b7095d011bf77d704aa}; full qualifications and provenance.}\par
\section*{I.6.25 PDP Unavailability Is Not Policy Denial}
\phantomsection
\label{navigation:I.6.25}
\addcontentsline{toc}{section}{I.6.25 PDP Unavailability Is Not Policy Denial}

\label{element:8d0e87312d1f4d6f9298843ce840e654}
\index{PDP Unavailability Is Not Policy Denial (I.6.25)}
{\small Constraint · Adjudicated Domain05 direction — U08/U12/U15; holistic result acceptance pending · current}\par

\begin{enumerate}


\item \textbf{PDP unavailability SHALL NOT be represented as a policy denial.}

\end{enumerate}


This assertion is normalised from the governed DesignNote under U08/U12/U15 and the ARB release Praxis. Original candidate/release qualifications and separately deferred implementation details remain in force.



Architectural subject: \hyperref[element:899d34ec491948b68aad09e507dea6d1]{Themis}.


{\footnotesize \href{../../../../continuum/motivation/constraint/8d0e87312d1f4d6f9298843ce840e654.md}{Canonical record: 8d0e87312d1f4d6f9298843ce840e654}; full qualifications and provenance.}\par
\section*{I.6.26 Persistence Preserves Applicable Compartments}
\phantomsection
\label{navigation:I.6.26}
\addcontentsline{toc}{section}{I.6.26 Persistence Preserves Applicable Compartments}

\label{element:6e963678cf6f4c3488c798d6afc3b910}
\index{Persistence Preserves Applicable Compartments (I.6.26)}
{\small Constraint · Adjudicated Domain05 direction — U11; holistic result acceptance pending · current}\par

Persistence SHALL NOT accidentally strip, weaken or lose the applicable compartment classification.



This assertion is normalised from the governed DesignNote under U11 and the ARB release Praxis. Original candidate/release qualifications and separately deferred implementation details remain in force.



Architectural subject: \hyperref[element:7985e647bf9e4ca3ba70faf9e3e4ae02]{Mnemosyne}.


{\footnotesize \href{../../../../continuum/motivation/constraint/6e963678cf6f4c3488c798d6afc3b910.md}{Canonical record: 6e963678cf6f4c3488c798d6afc3b910}; full qualifications and provenance.}\par
\section*{I.6.27 PetasosParcel Is Transported by Value}
\phantomsection
\label{navigation:I.6.27}
\addcontentsline{toc}{section}{I.6.27 PetasosParcel Is Transported by Value}

\label{element:b599ebffaf2f4e9e8deaec2ed98f2cd7}
\index{PetasosParcel Is Transported by Value (I.6.27)}
{\small Constraint · Adjudicated Domain05 direction — U09; holistic result acceptance pending · current}\par

Petasos SHALL NOT be described as transporting only references. It transports PetasosParcel instances by value.



This assertion is normalised from the governed DesignNote under U09 and the ARB release Praxis. Original candidate/release qualifications and separately deferred implementation details remain in force.



Architectural subject: \hyperref[element:195e5f10db8e419391354197b5d7a8cc]{Petasos}.


{\footnotesize \href{../../../../continuum/motivation/constraint/b599ebffaf2f4e9e8deaec2ed98f2cd7.md}{Canonical record: b599ebffaf2f4e9e8deaec2ed98f2cd7}; full qualifications and provenance.}\par
\section*{I.6.28 R1 and R2 Observation Does Not Authorise Autonomous Adaptation}
\phantomsection
\label{navigation:I.6.28}
\addcontentsline{toc}{section}{I.6.28 R1 and R2 Observation Does Not Authorise Autonomous Adaptation}

\label{element:7ca8481aa04a4dfdb34766511e002a5b}
\index{R1 and R2 Observation Does Not Authorise Autonomous Adaptation (I.6.28)}
{\small Constraint · Adjudicated Domain05 direction — U06/U09/U19; holistic result acceptance pending · current}\par

R1.x/R2.x SHALL NOT introduce autonomous adaptation merely because Energeia can observe Harmonia operation.



This assertion is normalised from the governed DesignNote under U06/U09/U19 and the ARB release Praxis. Original candidate/release qualifications and separately deferred implementation details remain in force.



Architectural subject: \hyperref[element:fb3f37a6ac614415ac486a6de9d35994]{Energeia}.


{\footnotesize \href{../../../../continuum/motivation/constraint/7ca8481aa04a4dfdb34766511e002a5b.md}{Canonical record: 7ca8481aa04a4dfdb34766511e002a5b}; full qualifications and provenance.}\par
\section*{I.6.29 Revocation Is Distinct from Temporal Expiry}
\phantomsection
\label{navigation:I.6.29}
\addcontentsline{toc}{section}{I.6.29 Revocation Is Distinct from Temporal Expiry}

\label{element:462d94a9a14d4eb8a74710fff4b12402}
\index{Revocation Is Distinct from Temporal Expiry (I.6.29)}
{\small Constraint · Adjudicated Domain05 direction — U12/U15/U16; holistic result acceptance pending · current}\par

Revocation SHALL NOT be treated merely as temporal expiry.



This assertion is normalised from the governed DesignNote under U12/U15/U16 and the ARB release Praxis. Original candidate/release qualifications and separately deferred implementation details remain in force.



Architectural subject: \hyperref[element:899d34ec491948b68aad09e507dea6d1]{Themis}.


{\footnotesize \href{../../../../continuum/motivation/constraint/462d94a9a14d4eb8a74710fff4b12402.md}{Canonical record: 462d94a9a14d4eb8a74710fff4b12402}; full qualifications and provenance.}\par
\section*{I.6.30 Runtime Vocabulary Membership Is Authoritative}
\phantomsection
\label{navigation:I.6.30}
\addcontentsline{toc}{section}{I.6.30 Runtime Vocabulary Membership Is Authoritative}

\label{element:872a5f1aa5a44fd4ba68d793c022adf2}
\index{Runtime Vocabulary Membership Is Authoritative (I.6.30)}
{\small Constraint · Adjudicated Domain05 direction — U13/U15; holistic result acceptance pending · current}\par

Software built against an earlier representation SHALL NOT automatically
treat a runtime value as invalid merely because that value was unknown
at build time.



This assertion is normalised from the governed DesignNote under U13/U15 and the ARB release Praxis. Original candidate/release qualifications and separately deferred implementation details remain in force.



Architectural subject: \hyperref[element:466ad8f75a954c0581b5bba090443849]{Calliope}.


{\footnotesize \href{../../../../continuum/motivation/constraint/872a5f1aa5a44fd4ba68d793c022adf2.md}{Canonical record: 872a5f1aa5a44fd4ba68d793c022adf2}; full qualifications and provenance.}\par
\section*{I.6.31 Semantic Persistence Domains Do Not Prescribe Physical Topology}
\phantomsection
\label{navigation:I.6.31}
\addcontentsline{toc}{section}{I.6.31 Semantic Persistence Domains Do Not Prescribe Physical Topology}

\label{element:eca9a99c2ad44a1ea7833d2b0c5a3e93}
\index{Semantic Persistence Domains Do Not Prescribe Physical Topology (I.6.31)}
{\small Constraint · Adjudicated Domain05 direction — U11; holistic result acceptance pending · current}\par

They SHALL NOT automatically be interpreted as three databases, schemas, services, workloads or physical deployments.



This assertion is normalised from the governed DesignNote under U11 and the ARB release Praxis. Original candidate/release qualifications and separately deferred implementation details remain in force.



Architectural subject: \hyperref[element:7985e647bf9e4ca3ba70faf9e3e4ae02]{Mnemosyne}.


{\footnotesize \href{../../../../continuum/motivation/constraint/eca9a99c2ad44a1ea7833d2b0c5a3e93.md}{Canonical record: eca9a99c2ad44a1ea7833d2b0c5a3e93}; full qualifications and provenance.}\par
\section*{I.6.32 Separate Models and DataSets Lifecycles}
\phantomsection
\label{navigation:I.6.32}
\addcontentsline{toc}{section}{I.6.32 Separate Models and DataSets Lifecycles}

\label{element:5c15fa0761da496b8194f3a6f742ecb5}
\index{Separate Models and DataSets Lifecycles (I.6.32)}
{\small Constraint · Adjudicated Domain05 direction — U13/U15; holistic result acceptance pending · current}\par

The two modules collaborate but SHALL NOT be treated as having a single
lifecycle.



This assertion is normalised from the governed DesignNote under U13/U15 and the ARB release Praxis. Original candidate/release qualifications and separately deferred implementation details remain in force.



Architectural subject: \hyperref[element:466ad8f75a954c0581b5bba090443849]{Calliope}.


{\footnotesize \href{../../../../continuum/motivation/constraint/5c15fa0761da496b8194f3a6f742ecb5.md}{Canonical record: 5c15fa0761da496b8194f3a6f742ecb5}; full qualifications and provenance.}\par
\section*{I.6.33 Specialised Execution Reuses Common Platform Capabilities}
\phantomsection
\label{navigation:I.6.33}
\addcontentsline{toc}{section}{I.6.33 Specialised Execution Reuses Common Platform Capabilities}

\label{element:968e208d336c48b781a4eabd76e6337b}
\index{Specialised Execution Reuses Common Platform Capabilities (I.6.33)}
{\small Constraint · Adjudicated Domain05 direction — U02/U06/U19; holistic result acceptance pending · current}\par

Specialised Harmonia execution constructs SHALL reuse common platform management, assurance, reporting and integration capabilities inherited through their base architectural constructs. Specialisation introduces only behaviour and information semantics necessary for the specialised role.



This assertion is normalised from the governed DesignNote under U02/U06/U19 and the ARB release Praxis. Original candidate/release qualifications and separately deferred implementation details remain in force.



Architectural subject: \hyperref[element:fb3f37a6ac614415ac486a6de9d35994]{Energeia}.


{\footnotesize \href{../../../../continuum/motivation/constraint/968e208d336c48b781a4eabd76e6337b.md}{Canonical record: 968e208d336c48b781a4eabd76e6337b}; full qualifications and provenance.}\par
\section*{I.6.34 Unavailable Authoritative Audit Capture Fails Closed}
\phantomsection
\label{navigation:I.6.34}
\addcontentsline{toc}{section}{I.6.34 Unavailable Authoritative Audit Capture Fails Closed}

\label{element:6c5d5f1d872642f2b7770248d4e1e1f4}
\index{Unavailable Authoritative Audit Capture Fails Closed (I.6.34)}
{\small Constraint · Adjudicated Domain05 direction — U07/U15; holistic result acceptance pending · current}\par

The inability to establish required audit evidence is therefore a platform-level loss of essential Harmonia service and SHALL result in fail-closed behaviour.



This assertion is normalised from the governed DesignNote under U07/U15 and the ARB release Praxis. Original candidate/release qualifications and separately deferred implementation details remain in force.



Architectural subject: \hyperref[element:fefda110e2274884ad58906733606efd]{Kleio}.


{\footnotesize \href{../../../../continuum/motivation/constraint/6c5d5f1d872642f2b7770248d4e1e1f4.md}{Canonical record: 6c5d5f1d872642f2b7770248d4e1e1f4}; full qualifications and provenance.}\par
\section*{I.6.35 Unknown Required Security Attributes Fail Closed}
\phantomsection
\label{navigation:I.6.35}
\addcontentsline{toc}{section}{I.6.35 Unknown Required Security Attributes Fail Closed}

\label{element:7ea0d91ed25848f4add672ddf08f42a2}
\index{Unknown Required Security Attributes Fail Closed (I.6.35)}
{\small Constraint · Adjudicated Domain05 direction — U12/U15/U16; holistic result acceptance pending · current}\par

If required effective security attributes cannot be established, enforcement SHALL fail closed.



This assertion is normalised from the governed DesignNote under U12/U15/U16 and the ARB release Praxis. Original candidate/release qualifications and separately deferred implementation details remain in force.



Architectural subject: \hyperref[element:899d34ec491948b68aad09e507dea6d1]{Themis}.


{\footnotesize \href{../../../../continuum/motivation/constraint/7ea0d91ed25848f4add672ddf08f42a2.md}{Canonical record: 7ea0d91ed25848f4add672ddf08f42a2}; full qualifications and provenance.}\par
\section*{I.6.36 Valid Context Permits Enforcement During PDP Unavailability}
\phantomsection
\label{navigation:I.6.36}
\addcontentsline{toc}{section}{I.6.36 Valid Context Permits Enforcement During PDP Unavailability}

\label{element:5a6285eb1a1646e1badd693b8c175f01}
\index{Valid Context Permits Enforcement During PDP Unavailability (I.6.36)}
{\small Constraint · Adjudicated Domain05 direction — U08/U12/U15; holistic result acceptance pending · current}\par

Themis PDP availability SHALL NOT, by itself, interrupt an in-flight activity that already possesses a valid PEP-Context sufficient for the policy enforcement decisions required by that activity.



This assertion is normalised from the governed DesignNote under U08/U12/U15 and the ARB release Praxis. Original candidate/release qualifications and separately deferred implementation details remain in force.



Architectural subject: \hyperref[element:899d34ec491948b68aad09e507dea6d1]{Themis}.


{\footnotesize \href{../../../../continuum/motivation/constraint/5a6285eb1a1646e1badd693b8c175f01.md}{Canonical record: 5a6285eb1a1646e1badd693b8c175f01}; full qualifications and provenance.}\par
\chapter*{I.7 Architecture Axioms}
\phantomsection
\label{navigation:I.7}
\addcontentsline{toc}{chapter}{I.7 Architecture Axioms}

\section*{I.7.1 AX-01 Harmonia Is Health-Information Centric}
\phantomsection
\label{navigation:I.7.1}
\addcontentsline{toc}{section}{I.7.1 AX-01 Harmonia Is Health-Information Centric}

\label{element:e9bd61acaaea4577bc4086a5639e52c3}
\index{AX-01 Harmonia Is Health-Information Centric (I.7.1)}
{\small Principle · Canonical · current}\par

\par\medskip\textbf{Axiom}\par


Harmonia is a distributed health-information management and
interoperability framework.



Its primary purpose is the ingestion, validation, transformation,
governance, processing, persistence, distribution and exposure of
health-related information.



Harmonia SHALL support applicable health-information standards and
models while retaining the ability to manage information according to
Harmonia's own internal operational requirements.



\par\medskip\textbf{Rationale}\par


Harmonia is not intended to become a generic distributed application
framework to which healthcare happens to be attached.



Its architecture exists to solve health-information problems:
interoperability, information authority, provenance, security,
terminology, workflow, persistence, distribution, resilience and
longitudinal information management.



\par\medskip\textbf{Consequences}\par


Health-information semantics take precedence over generic infrastructure
convenience.



Standards such as HL7 FHIR and HL7 v2 are first-class architectural
concerns.



Generic infrastructure capabilities SHOULD be used where useful, but
SHALL NOT dictate Harmonia's health-information semantics.



\textbf{This does not mean}



\begin{itemize}


\item Harmonia is a FHIR server.

\item Harmonia is restricted to FHIR information.

\item Every internal Harmonia object must correspond to a healthcare interoperability standard.

\end{itemize}

{\footnotesize \href{../../../../continuum/motivation/principle/e9bd61acaaea4577bc4086a5639e52c3.md}{Canonical record: e9bd61acaaea4577bc4086a5639e52c3}; full qualifications and provenance.}\par
\section*{I.7.2 AX-02 Standards at the Boundary; Harmonia Within the Boundary}
\phantomsection
\label{navigation:I.7.2}
\addcontentsline{toc}{section}{I.7.2 AX-02 Standards at the Boundary; Harmonia Within the Boundary}

\label{element:eee8c3783d964d18a7f9dcea85fcd4cf}
\index{AX-02 Standards at the Boundary; Harmonia Within the Boundary (I.7.2)}
{\small Principle · Canonical · current}\par

\par\medskip\textbf{Axiom}\par


Harmonia SHALL ingest and expose health information using the standards,
information models and interoperability protocols applicable to an
external interface contract.



Within Harmonia, information MAY be augmented, represented, indexed,
cached, distributed, governed, versioned, correlated and persisted using
Harmonia-specific semantics and metadata necessary to satisfy Harmonia's
operational requirements.



Harmonia's internal operational semantics are private to Harmonia.
External systems SHALL NOT be required to understand, preserve,
reproduce or participate in them.



\par\medskip\textbf{Rationale}\par


Interoperability requires agreement at system boundaries. It does not
require participating systems to share an internal architecture.



Allowing external contracts to dictate Harmonia's internal
implementation would couple Harmonia to the assumptions and limitations
of external systems.



\par\medskip\textbf{Consequences}\par


Pylai acts as an interoperability membrane, not merely a protocol
adapter.



External representations SHALL conform to their applicable interface
contracts.



Harmonia-private operational metadata SHALL NOT leak into an external
interoperability contract.



\par\medskip\textbf{Maxim}\par


Standards at the boundary. Harmonia within the boundary.


{\footnotesize \href{../../../../continuum/motivation/principle/eee8c3783d964d18a7f9dcea85fcd4cf.md}{Canonical record: eee8c3783d964d18a7f9dcea85fcd4cf}; full qualifications and provenance.}\par
\section*{I.7.3 AX-03 Native Standards Models Remain Native}
\phantomsection
\label{navigation:I.7.3}
\addcontentsline{toc}{section}{I.7.3 AX-03 Native Standards Models Remain Native}

\label{element:e99a4e2bf99f476d876394300d6b9c3d}
\index{AX-03 Native Standards Models Remain Native (I.7.3)}
{\small Principle · Canonical · current}\par

\par\medskip\textbf{Axiom}\par


Standards-defined information objects SHALL retain their
standards-defined representation within Harmonia wherever practicable.



Harmonia SHALL NOT introduce parallel or derivative representations
solely to accommodate internal management concerns where the standard
representation provides an appropriate extensibility mechanism.



\par\medskip\textbf{Rationale}\par


Parallel representations introduce mapping complexity, semantic drift
and unnecessary ownership questions.



For FHIR information, the FHIR resource is already an extensible
information model.



\par\medskip\textbf{Consequences}\par


FHIR R5 and FHIR R4 AU resources SHOULD remain native FHIR resources
represented using the applicable HAPI FHIR models.



Where appropriate, Harmonia-specific durable metadata MAY be represented
using private FHIR extensions.



Jurisdictional and implementation-guide extensions, such as extensions
required for Australian FHIR conformance, remain part of the
standards-based representation and SHALL NOT be confused with
Harmonia-private operational extensions.



Non-FHIR managed objects MAY use representation-appropriate Harmonia
metadata mechanisms.



\textbf{This does not mean:}



\begin{itemize}


\item All Harmonia metadata must be represented as FHIR extensions.

\item Cache-specific state must be embedded into FHIR resources.

\item Non-FHIR Harmonia objects must be converted into FHIR.

\end{itemize}

{\footnotesize \href{../../../../continuum/motivation/principle/e99a4e2bf99f476d876394300d6b9c3d.md}{Canonical record: e99a4e2bf99f476d876394300d6b9c3d}; full qualifications and provenance.}\par
\section*{I.7.4 AX-04 Harmonia Owns the Semantics; Engines Provide the Machinery}
\phantomsection
\label{navigation:I.7.4}
\addcontentsline{toc}{section}{I.7.4 AX-04 Harmonia Owns the Semantics; Engines Provide the Machinery}

\label{element:54508cda80484a6caf95b901780eaf7d}
\index{AX-04 Harmonia Owns the Semantics; Engines Provide the Machinery (I.7.4)}
{\small Principle · Canonical · current}\par

\par\medskip\textbf{Axiom}\par


Harmonia owns the architectural semantics governing information
authority, security, lifecycle, concurrency, provenance, auditability,
resilience and operational integrity.



Harmonia SHALL preferentially use capabilities supplied by underlying
technology engines where those capabilities satisfy Harmonia's
architectural invariants rather than reproduce equivalent functionality.



\par\medskip\textbf{Rationale}\par


Harmonia should own its meaning without unnecessarily owning commodity
infrastructure machinery.



\par\medskip\textbf{Consequences}\par


HAPI FHIR MAY provide FHIR persistence, indexing, versioning and search.



Infinispan MAY provide distributed caching, entry metadata and atomic
coordination primitives.



ActiveMQ Artemis MAY provide durable messaging.



PostgreSQL MAY provide transactional relational persistence.



These technologies implement Harmonia capabilities. They do not define
Harmonia's architectural semantics.



\par\medskip\textbf{Maxim}\par


Use the machinery. Own the semantics.



\textbf{This does not mean:}



\begin{itemize}


\item Harmonia should wrap every engine capability behind a bespoke implementation.

\item An engine's native behaviour automatically satisfies a Harmonia invariant.

\item Using HAPI FHIR makes Harmonia a HAPI FHIR application.

\end{itemize}

{\footnotesize \href{../../../../continuum/motivation/principle/54508cda80484a6caf95b901780eaf7d.md}{Canonical record: 54508cda80484a6caf95b901780eaf7d}; full qualifications and provenance.}\par
\section*{I.7.5 AX-05 Active State and Authoritative Durable State Are Distinct}
\phantomsection
\label{navigation:I.7.5}
\addcontentsline{toc}{section}{I.7.5 AX-05 Active State and Authoritative Durable State Are Distinct}

\label{element:0412a2084f0145babc4f022c71fea65f}
\index{AX-05 Active State and Authoritative Durable State Are Distinct (I.7.5)}
{\small Principle · Canonical · current}\par

\par\medskip\textbf{Axiom}\par


Mneme owns Harmonia's application-facing access to managed information
and the distributed active-state representation, observation and
coordination required to use that information safely.



Mnemosyne owns Harmonia's authoritative durable representation of
managed information. It atomically establishes authoritative state and
authoritative version progression and persists the durable management
metadata required to interpret that state.



Mneme manages active use; Mnemosyne establishes durable truth.



\par\medskip\textbf{Rationale}\par


Distributed availability and authoritative persistence are different
responsibilities.



Conflating them makes cache state authoritative, persistence state
application-facing, or both.



\par\medskip\textbf{Consequences}\par


Mneme MAY reject or coordinate a proposed state progression before
persistence.



Only Mnemosyne can establish a new authoritative durable state.



Following authoritative commit, Mneme SHALL converge its active
representation toward the authoritative state.



Loss of Mneme SHALL NOT cause process-local or cached state to become
authoritative.



Loss of Mnemosyne SHALL NOT cause Mneme to promote its active
representation into authoritative durable state.



Mneme active state SHALL be reconstructable from authoritative state
where applicable.



\textbf{This does not mean:}



\begin{itemize}


\item Every Mneme operation requires a database operation.

\item Mnemosyne is exposed as an application database API.

\item Mneme and Mnemosyne must use the same version or concurrency mechanism.

\end{itemize}


Mneme active-state generation SHALL represent successfully established active state. A valid active-state token SHALL identify an observed generation of that state.



Failed or degraded convergence following authoritative state progression SHALL NOT establish or advance a valid active-state generation. Where the active representation cannot be successfully converged, its coordination state SHALL be treated as untrusted until reconciled with authoritative state.



Mneme active-state generation and Mnemosyne authoritative version SHALL remain distinct concurrency domains. Neither SHALL be inferred from, substituted for, or treated as an alias of the other.


{\footnotesize \href{../../../../continuum/motivation/principle/0412a2084f0145babc4f022c71fea65f.md}{Canonical record: 0412a2084f0145babc4f022c71fea65f}; full qualifications and provenance.}\par
\section*{I.7.6 AX-06 Information Authority Is Explicit}
\phantomsection
\label{navigation:I.7.6}
\addcontentsline{toc}{section}{I.7.6 AX-06 Information Authority Is Explicit}

\label{element:952c3b6148ce4fce94b69f99a0c2fec7}
\index{AX-06 Information Authority Is Explicit (I.7.6)}
{\small Principle · Canonical · current}\par

\par\medskip\textbf{Axiom}\par


Harmonia SHALL represent the authority and credibility of managed
information independently of its technical transport, cache state,
persistence state or concurrency state.



Harmonia SHALL support governed information-authority concepts including:



\begin{itemize}


\item Authoritative,

\item Informational, and

\item Anecdotal.

\end{itemize}


\par\medskip\textbf{Rationale}\par


Technical correctness and information credibility answer different
questions.



A successfully persisted assertion is not necessarily authoritative. A
stale assertion is not necessarily false. A technically reliable source
is not necessarily the authoritative source for every fact it supplies.



\par\medskip\textbf{Consequences}\par


Concurrency determines whether information has changed.



Semantic analysis determines what changed.



Information authority and assurance describe the standing of competing
assertions.



Policy determines what Harmonia does about them.



Calliope governs the semantics of information authority. Appropriate
Harmonia processes and policies apply those semantics.



\textbf{This does not mean:}



\begin{itemize}


\item Authoritative information automatically wins every conflict.

\item Anecdotal information is discarded.

\item Persistence success establishes information authority.

\end{itemize}

{\footnotesize \href{../../../../continuum/motivation/principle/952c3b6148ce4fce94b69f99a0c2fec7.md}{Canonical record: 952c3b6148ce4fce94b69f99a0c2fec7}; full qualifications and provenance.}\par
\section*{I.7.7 AX-07 Security Is Intrinsic to Managed Operations}
\phantomsection
\label{navigation:I.7.7}
\addcontentsline{toc}{section}{I.7.7 AX-07 Security Is Intrinsic to Managed Operations}

\label{element:5de2ae2c3b794c2d9f3a14efd51d564f}
\index{AX-07 Security Is Intrinsic to Managed Operations (I.7.7)}
{\small Principle · Canonical · current}\par

\par\medskip\textbf{Axiom}\par


Every governed operation SHALL execute within an established security
context and SHALL be subject to Harmonia security policy.



Security SHALL be enforced by framework and platform boundaries where
practicable and SHALL NOT depend solely upon developer knowledge, coding
convention or voluntary caller behaviour.



\par\medskip\textbf{Rationale}\par


Security is part of the semantics of performing an operation, not an
optional service invoked by well-behaved code.



\par\medskip\textbf{Consequences}\par


Themis owns Harmonia security policy evaluation.



Trusted identity and security context SHALL be established before
governed actions are performed.



Security context is operational context and SHALL NOT automatically
become persistent information content.



Security-significant events MAY become durable audit evidence according
to policy.







\par\medskip\textbf{Kleio mandatory authoritative audit capture}\par


Kleio SHALL perform mandatory authoritative audit capture under intrinsic
platform authority associated with its established audit responsibility.
That capture SHALL execute within an established security context and remain
subject to Harmonia security policy; Themis retains policy evaluation ownership.
This is governed platform authority, not a security bypass or policy exception.



The Identity associated with the audited activity SHALL remain audit evidence.
It is not the authority under which the mandatory AuditEvent is persisted.
Kleio SHALL NOT borrow the originating Identity's PEP-Criteria merely to persist
evidence of that Identity's activity. Identity-derived authorisation resolved
and enforced through Themis remains distinct from Kleio's capture authority.



Where authoritative audit capture is required for a governed operation and
Kleio cannot exercise that platform authority sufficiently to capture the
required evidence, the governed operation SHALL fail closed.



This authority is established specifically for Kleio mandatory authoritative
audit capture. It SHALL NOT be inferred for other Kleio activities or generalised
to other Harmonia components. Any proposed equivalent authority elsewhere is
a separate architecture question under AX-17. Detailed authority representation
and enforcement mechanisms remain deferred.



\textbf{This does not mean:}



\begin{itemize}


\item Every security decision belongs in a FHIR resource.

\item Every authorization evaluation belongs in the audit repository.

\item Internal security context should be exposed to external systems.

\end{itemize}

{\footnotesize \href{../../../../continuum/motivation/principle/5de2ae2c3b794c2d9f3a14efd51d564f.md}{Canonical record: 5de2ae2c3b794c2d9f3a14efd51d564f}; full qualifications and provenance.}\par
\section*{I.7.8 AX-08 Evidence Records Meaning, Not Machinery}
\phantomsection
\label{navigation:I.7.8}
\addcontentsline{toc}{section}{I.7.8 AX-08 Evidence Records Meaning, Not Machinery}

\label{element:21d6864657964ba5b03e3f6b5eb29cd5}
\index{AX-08 Evidence Records Meaning, Not Machinery (I.7.8)}
{\small Principle · Canonical · current}\par

\par\medskip\textbf{Axiom}\par


Provenance and audit evidence SHALL describe information-significant,
security-significant and business-significant events, assertions and
decisions.



Internal implementation mechanics SHALL NOT ordinarily become provenance
or audit evidence merely because they occurred.



\par\medskip\textbf{Rationale}\par


A durable evidence repository should explain what happened to
information and why, rather than become a permanent distributed debug
log.



\par\medskip\textbf{Consequences}\par


Authoritative state acceptance, significant information transformations,
stewardship decisions and security-significant actions MAY require
durable evidence.



Cache misses, cache rebalancing, CAS retries, connection retries,
reconstruction and ordinary infrastructure mechanics normally belong in
logging, metrics or transient diagnostics.



Kleio preserves durable evidence.



\par\medskip\textbf{Maxim}\par


Preserve the meaningful fact, not every mechanism that produced it.


{\footnotesize \href{../../../../continuum/motivation/principle/21d6864657964ba5b03e3f6b5eb29cd5.md}{Canonical record: 21d6864657964ba5b03e3f6b5eb29cd5}; full qualifications and provenance.}\par
\section*{I.7.9 AX-09 Transient Operational State Is Ephemeral by Default}
\phantomsection
\label{navigation:I.7.9}
\addcontentsline{toc}{section}{I.7.9 AX-09 Transient Operational State Is Ephemeral by Default}

\label{element:9a5cb9e8876844e58e660f909f7df9fc}
\index{AX-09 Transient Operational State Is Ephemeral by Default (I.7.9)}
{\small Principle · Canonical · current}\par

\par\medskip\textbf{Axiom}\par


Harmonia MAY maintain transient operational state associated with
managed information where required for caching, persistence,
concurrency, security, routing, resilience, diagnostics and processing
integrity.



Such state SHALL be retained only for as long as required by its
operational purpose and SHALL normally disappear when that purpose has
been satisfied.



Transient operational state SHALL NOT automatically become provenance or
audit evidence.



\par\medskip\textbf{Diagnostic observability}\par


Harmonia SHALL provide controlled mechanisms by which designated
transient operational state MAY be observed or reported when diagnostics
are explicitly enabled at framework, subsystem or component scope.



Where transient state directly contributes to an information-significant
decision, assertion or evidentiary fact, the meaningful information
SHALL be promoted into the appropriate durable provenance or audit
representation.



\par\medskip\textbf{Rationale}\par


Distributed systems require substantial operational state. Persisting
all of it indefinitely would confuse implementation mechanics with
information history.



Conversely, making operational state completely invisible would make the
framework extremely difficult to diagnose.



\textbf{This does not mean:}



\begin{itemize}


\item Transient state is unimportant.

\item Transient state cannot be logged or inspected.

\item Diagnostic traces should automatically become permanent evidence.

\end{itemize}

{\footnotesize \href{../../../../continuum/motivation/principle/9a5cb9e8876844e58e660f909f7df9fc.md}{Canonical record: 9a5cb9e8876844e58e660f909f7df9fc}; full qualifications and provenance.}\par
\section*{I.7.10 AX-10 Distribution, Load and Failure Are Normal Operating Conditions}
\phantomsection
\label{navigation:I.7.10}
\addcontentsline{toc}{section}{I.7.10 AX-10 Distribution, Load and Failure Are Normal Operating Conditions}

\label{element:0d9673833eb840eab03804aa076c1b0c}
\index{AX-10 Distribution, Load and Failure Are Normal Operating Conditions (I.7.10)}
{\small Principle · Canonical · current}\par

\par\medskip\textbf{Axiom}\par


Harmonia SHALL be designed on the assumption of distributed deployment,
sustained processing load and failure or degradation of individual
runtime components.



Loss or degradation of implementation machinery SHALL NOT silently alter
the authoritative, security, governance or information-authority
semantics of managed information.



\par\medskip\textbf{Rationale}\par


Distribution and failure are not exceptional architectural cases for
Harmonia. They are part of its normal operating environment.



\par\medskip\textbf{Consequences}\par


Work that has crossed an explicit durable processing boundary SHALL be
recoverable according to that boundary's contract.



Processing pressure SHOULD preferentially be represented as bounded,
durable backlog rather than unbounded growth of JVM heap, threads,
database connections or volatile work.



Failure SHALL be visible rather than silently replaced by semantically
weaker local behaviour.



\textbf{This does not mean:}



\begin{itemize}


\item Every component must be independently distributed.

\item Every operation must be asynchronous.

\item Every internal Java boundary should become a network boundary.

\end{itemize}

{\footnotesize \href{../../../../continuum/motivation/principle/0d9673833eb840eab03804aa076c1b0c.md}{Canonical record: 0d9673833eb840eab03804aa076c1b0c}; full qualifications and provenance.}\par
\section*{I.7.11 AX-11 Managed Information Access Is Highly Available and Responsive}
\phantomsection
\label{navigation:I.7.11}
\addcontentsline{toc}{section}{I.7.11 AX-11 Managed Information Access Is Highly Available and Responsive}

\label{element:8a1e995f16954484a78ae1c5c114d546}
\index{AX-11 Managed Information Access Is Highly Available and Responsive (I.7.11)}
{\small Principle · Canonical · current}\par

\par\medskip\textbf{Axiom}\par


Mneme SHALL provide highly available, responsive and load-tolerant
access to Harmonia-managed active information.



Its architecture SHALL favour bounded resource consumption, distributed
workload, reconstructable active state and graceful expression of
processing pressure.



Availability SHALL NOT be achieved by weakening authoritative-state,
concurrency, security, governance or information-authority guarantees.



\par\medskip\textbf{Rationale}\par


Mneme exists partly to prevent every application interaction from
becoming an authoritative persistence interaction while still preserving
Harmonia's correctness guarantees.



\par\medskip\textbf{Consequences}\par


Different categories of Mneme state MAY use different distribution and
consistency strategies.



Small coordination state need not have the same topology as large
resource caches.



Cache topology SHALL be selected according to the semantics and
operational requirements of the managed state rather than applied
universally.



\textbf{This does not mean:}



\begin{itemize}


\item Every cache should use REPL\_SYNC.

\item Every managed resource must always reside in memory.

\item Cache availability is equivalent to authoritative information availability.

\end{itemize}

{\footnotesize \href{../../../../continuum/motivation/principle/8a1e995f16954484a78ae1c5c114d546.md}{Canonical record: 8a1e995f16954484a78ae1c5c114d546}; full qualifications and provenance.}\par
\section*{I.7.12 AX-12 Hide Plumbing, Not Information}
\phantomsection
\label{navigation:I.7.12}
\addcontentsline{toc}{section}{I.7.12 AX-12 Hide Plumbing, Not Information}

\label{element:ce85889bd9b34f79a53d4c907ceb19d4}
\index{AX-12 Hide Plumbing, Not Information (I.7.12)}
{\small Principle · Canonical · current}\par

\par\medskip\textbf{Axiom}\par


Developers implementing Harmonia work units SHALL have direct and fluent
access to applicable standards-defined health-information models and
Harmonia's managed-information capabilities.



Framework plumbing SHALL provide security, governance, authority,
persistence, concurrency, provenance, messaging and operational services
without unnecessarily obscuring or replacing the underlying information
models.



\par\medskip\textbf{Rationale}\par


Ergon and other Harmonia developers should
implement health-information behaviour, not repeatedly reconstruct the
framework's infrastructure contracts.



At the same time, hiding the actual information model behind generic
wrappers makes health semantics harder rather than easier to understand.



\par\medskip\textbf{Consequences}\par


A developer working with a FHIR Practitioner, Observation,
DiagnosticReport or other resource SHOULD be able to work with the real
standards model.



Harmonia SHALL make the safe and governed path the normal and convenient
path.



Raw infrastructure mechanisms SHOULD NOT become the ordinary application
programming model.



\par\medskip\textbf{Maxim}\par


Hide plumbing, not information.



\textbf{This does not mean:}



\begin{itemize}


\item Every capability must be exposed through one universal API.

\item Standards models should contain Harmonia infrastructure logic.

\item Developers should directly manipulate RemoteCache, EntityManager, JMS or other infrastructure merely because Harmonia uses those technologies internally.

\end{itemize}

{\footnotesize \href{../../../../continuum/motivation/principle/ce85889bd9b34f79a53d4c907ceb19d4.md}{Canonical record: ce85889bd9b34f79a53d4c907ceb19d4}; full qualifications and provenance.}\par
\section*{I.7.13 AX-13 Harmonia Management Has an Explicit Boundary}
\phantomsection
\label{navigation:I.7.13}
\addcontentsline{toc}{section}{I.7.13 AX-13 Harmonia Management Has an Explicit Boundary}

\label{element:b6a55decfeec4f1bb3d7f18534a2ffea}
\index{AX-13 Harmonia Management Has an Explicit Boundary (I.7.13)}
{\small Principle · Canonical · current}\par

\par\medskip\textbf{Axiom}\par


Ingress establishes Harmonia management. Internal processing preserves
Harmonia management. Egress terminates Harmonia management of the
emitted representation.



Harmonia MAY retain evidence describing the information from which an
external representation was derived, the transformation applied, the
transmission performed and the externally observable outcome.



Harmonia SHALL NOT attribute its internal authority, governance,
concurrency, security, availability or operational guarantees to an
emitted representation after that representation crosses an external
egress boundary.



\par\medskip\textbf{Rationale}\par


Harmonia can govern information while it is under Harmonia's control. It
cannot govern a copy after that copy has been transferred to another
system.



Provenance can describe what Harmonia emitted; it cannot extend
Harmonia's operational control into another system.



\par\medskip\textbf{Consequences}\par


An HL7 v2 message ceases to be a Harmonia-managed representation when
emitted through its external egress boundary.



A FHIR resource returned through REST similarly becomes an externally
held representation after transmission.



A subsequent FHIR PUT, HL7 v2 message or other submission represents a
new ingress interaction and is assessed accordingly.



An HTTP ETag, FHIR meta.versionId, HL7 acknowledgement or equivalent
interoperability mechanism MAY communicate externally meaningful
transaction or version semantics. It SHALL NOT be interpreted as
extending Harmonia's internal operational control beyond its boundary.



\par\medskip\textbf{Maxim}\par


Provenance may describe or cross the boundary; operational control does
not.



\textbf{This does not mean:}



\begin{itemize}


\item The underlying information inside Harmonia ceases to be managed when a representation is published.

\item Harmonia forgets what it transmitted.

\item Acknowledgements and transaction outcomes are irrelevant to audit or provenance.

\end{itemize}

{\footnotesize \href{../../../../continuum/motivation/principle/b6a55decfeec4f1bb3d7f18534a2ffea.md}{Canonical record: b6a55decfeec4f1bb3d7f18534a2ffea}; full qualifications and provenance.}\par
\section*{I.7.14 AX-14 Semantic Distinctions Are Preserved}
\phantomsection
\label{navigation:I.7.14}
\addcontentsline{toc}{section}{I.7.14 AX-14 Semantic Distinctions Are Preserved}

\label{element:d3fc6025dc7646d1b2020b51f7a44bfa}
\index{AX-14 Semantic Distinctions Are Preserved (I.7.14)}
{\small Principle · Canonical · current}\par

\par\medskip\textbf{Axiom}\par


Harmonia SHALL preserve meaningful distinctions between states, outcomes and concepts throughout its internal processing and across subsystem boundaries.



A distinction that is significant to information meaning, authority, security, concurrency, persistence, processing outcome or operational correctness SHALL NOT be collapsed merely because an underlying technology, transport, API or implementation abstraction does not represent that distinction directly.



Where an external interface contract intentionally presents a simpler or different semantic model, Harmonia MAY project its internal semantics into that contract explicitly at the applicable boundary.



\par\medskip\textbf{Rationale}\par


Loss of semantic distinction causes information about what Harmonia knows, what occurred and what may safely happen next to be lost.



Two conditions may produce similar implementation behaviour while having materially different meanings. Resource absence, access denial, persistence failure and indeterminate outcome may all prevent information from being returned, but they do not describe the same state.



Likewise, persistence version, active-state generation, interoperability version and information authority describe different properties even where an implementation happens to correlate them.



Implementation convenience SHALL NOT redefine those meanings.



\par\medskip\textbf{Consequences}\par


Semantically distinct outcomes SHALL remain distinguishable through Harmonia's internal contracts.



Security denial SHALL NOT be represented as resource absence merely because both prevent access to information.



A failure to establish authoritative state SHALL NOT be represented as authoritative absence.



An indeterminate outcome SHALL NOT be represented as either success or failure unless subsequent processing establishes that outcome.



ActiveStateToken, AuthoritativeVersion, externally meaningful version identifiers and standards-defined version metadata SHALL remain distinct where they represent different semantic domains.



Where an underlying engine or protocol lacks a distinction required by Harmonia, Harmonia SHALL introduce an appropriate semantic representation rather than infer meaning from incidental implementation details such as exception classes, message text or correlated version values.



Boundary projection MAY deliberately reduce or transform semantic distinctions where required by an external contract or policy, but such projection SHALL be explicit and SHALL NOT redefine Harmonia's internal semantics.



\par\medskip\textbf{Maxim}\par


Preserve semantics internally; project deliberately at boundaries.



This does not mean: Harmonia must expose every internal distinction externally. Similar outcomes can never share implementation machinery. Every semantic distinction requires a separate Java type. External standards or interface contracts may be ignored in favour of Harmonia's internal model.


{\footnotesize \href{../../../../continuum/motivation/principle/d3fc6025dc7646d1b2020b51f7a44bfa.md}{Canonical record: d3fc6025dc7646d1b2020b51f7a44bfa}; full qualifications and provenance.}\par
\section*{I.7.15 AX-15 Uncertainty Is Preserved Until Resolved}
\phantomsection
\label{navigation:I.7.15}
\addcontentsline{toc}{section}{I.7.15 AX-15 Uncertainty Is Preserved Until Resolved}

\label{element:ef8abbaec37e4890ad78b9f263e67906}
\index{AX-15 Uncertainty Is Preserved Until Resolved (I.7.15)}
{\small Principle · Canonical · current}\par

\par\medskip\textbf{Axiom}\par


Where Harmonia cannot establish the state, outcome or effect of an operation with the certainty required by its governing contract, that uncertainty SHALL be represented explicitly.



Harmonia SHALL NOT infer, manufacture or assume a more certain outcome merely to simplify processing, recovery or application behaviour.



Uncertainty MAY be resolved through authoritative observation, reconciliation or other governed evidence capable of establishing the required state.



\par\medskip\textbf{Rationale}\par


Distributed operations can fail at points where the observable failure does not establish whether the requested operation occurred.



For example, loss of a connection after transmitting a request does not establish whether the receiving subsystem processed or committed that request. Treating such an outcome as either success or failure introduces information that Harmonia does not possess.



Preserving uncertainty allows subsequent processing to determine what actually occurred without compounding an unknown outcome through unsafe assumptions or repeated actions.



\par\medskip\textbf{Consequences}\par


An indeterminate operation outcome SHALL remain distinguishable from both confirmed success and confirmed non-occurrence.



Harmonia SHALL NOT automatically retry a state-changing operation where the previous attempt may have succeeded unless the operation's contract makes that retry demonstrably safe.



Recovery from an uncertain outcome SHOULD preferentially establish the current authoritative state and reconcile against it before further state-changing action is attempted.



Timeouts, connection failures and similar technical events SHALL NOT, by themselves, determine semantic outcome. Their interpretation SHALL depend upon what can be established about the operation's progress and the governing boundary contract.



Where subsequent authoritative observation or other governed evidence resolves an uncertainty, Harmonia MAY continue processing from the newly established state.



Uncertainty SHALL NOT cause transient, cached or process-local state to be promoted into authoritative state.



\par\medskip\textbf{Maxim}\par


\textbf{Do not turn “unknown” into “yes” or “no”.}



\textbf{This does not mean:} Harmonia must retain uncertainty indefinitely. Every technical failure produces an uncertain outcome. Operations known not to have crossed the relevant state-changing boundary cannot be safely retried. Reconciliation must always require human intervention.


{\footnotesize \href{../../../../continuum/motivation/principle/ef8abbaec37e4890ad78b9f263e67906.md}{Canonical record: ef8abbaec37e4890ad78b9f263e67906}; full qualifications and provenance.}\par
\section*{I.7.16 AX-16 Operational Activity \& Entity State Progress Together}
\phantomsection
\label{navigation:I.7.16}
\addcontentsline{toc}{section}{I.7.16 AX-16 Operational Activity \& Entity State Progress Together}

\label{element:de7e90f46127455a9ec9ad3b9b48ba54}
\index{AX-16 Operational Activity \& Entity State Progress Together (I.7.16)}
{\small Principle · Canonical · current}\par

This axiom is maintained in this register. Current Harmonia realisation describes downstream machinery and is non-normative.



\par\medskip\textbf{Axiom}\par


\begin{quote}


\textbf{Operational activity associated with a real-world entity progresses through explicit and observable state transitions coordinated with the governed state of that entity. Operational activity is not treated merely as a sequence of disconnected message transfers.}



\end{quote}


\par\medskip\textbf{Architectural Implications}\par


\begin{itemize}


\item Health integration involves multi-stage workflows (e.g., ADT distribution, lab specimen processing).

\item Operational tasks must be tracked in lockstep with the real-world state of the clinical entity.

\item Prevents silent destination divergence during fan-out processing: every destination hop is an observable sub-state of the parent operational activity.

\end{itemize}


\par\medskip\textbf{Current Harmonia Realisation (Non-Normative)}\par


\begin{itemize}


\item Implemented in \texttt{\seqsplit{energeia-erga}} (e.g., \texttt{\seqsplit{AdtDistributionErgon}} tracking destination fan-out checkpoints).

\item Multi-destination statuses (\texttt{\seqsplit{FANOUT\_DISPATCH\_INITIATED}}, \texttt{\seqsplit{QUEUED}}, \texttt{\seqsplit{DELIVERED}}) are recorded in Pragma envelopes and FHIR \texttt{\seqsplit{Task.output}} extensions.

\end{itemize}


\par\medskip\textbf{Approved interpretation and related authority}\par


The 2026-10-09 migration adjudication preserves coordinated, explicit and
observable progression. The "lockstep" implication above does not mandate
synchronous execution, atomic entity/activity updates or a literal lockstep
implementation. AX-16 does not require universal Digital Twin routing, a Twin
for population/search activity, or clinical-work authority.



The \href{\detokenize{../../../../../docs/markdown/02-strategy/strategic-views/logical-component-responsibilities.md#entity-specific-operational-coordination}}{Strategy entity-specific Digital Twin coordination rule}
remains the controlling Strategy-level specialisation of AX-16.
\hyperref[element:b8cafe3cf5904d2bb9a9767260fddf78]{AX-18} constrains inference of domain authority from coordination.


{\footnotesize \href{../../../../continuum/motivation/principle/de7e90f46127455a9ec9ad3b9b48ba54.md}{Canonical record: de7e90f46127455a9ec9ad3b9b48ba54}; full qualifications and provenance.}\par
\section*{I.7.17 AX-17 Architectural Authority and Explicit Uncertainty}
\phantomsection
\label{navigation:I.7.17}
\addcontentsline{toc}{section}{I.7.17 AX-17 Architectural Authority and Explicit Uncertainty}

\label{element:8882573df7ce43c6bc0f66da0a676139}
\index{AX-17 Architectural Authority and Explicit Uncertainty (I.7.17)}
{\small Principle · Canonical · current}\par

\par\medskip\textbf{Axiom}\par


Documented architectural decisions, definitions, relationships, boundaries
and constraints are authoritative and SHALL NOT be silently reinterpreted,
replaced, bypassed or contradicted by downstream architecture or
implementation.



Where the architecture does not establish a fact, relationship,
responsibility or decision, that absence SHALL be preserved explicitly
rather than completed through inference, convention, structural
convenience, lexical similarity or anticipated implementation.



Unresolved architecture MAY be accompanied by clearly identified
candidates, suggestions or recommendations, but these SHALL remain
explicitly non-authoritative until explicitly accepted into the
architecture.



Where evidence indicates that authoritative architecture may be incorrect,
incomplete or internally inconsistent, the conflict SHALL be raised
explicitly for architectural review. The existing architecture remains
authoritative until an approved architectural change is made. That
approved change becomes the new authoritative architecture.



\par\medskip\textbf{Rationale}\par


Downstream derivation must distinguish established architectural knowledge
from gaps and proposed resolutions. A plausible relationship supplies no
authority merely because it makes a model or implementation appear
complete. Absence of an established relationship is architecturally
meaningful information.



Repeated architectural review during Domain04 Package2 G1 confirmed this
need across different traceability relationships:



\begin{itemize}


\item K9: unresolved Capability hierarchy must not be completed through
  invented ancestry.

\item K10: Behaviour / Process traceability must not be completed through
  invented stage equivalence.

\item K11: Function / Service traceability must not be completed through
  unsupported Feature association.

\item K12: cross-capability dependency traceability must not be completed
  through invented Service ownership or consumption.

\item K13: downstream representation traceability must not be completed through
  anticipated FHIR, application or technology mappings.

\end{itemize}


These examples justify an architecture-wide rule; its application does
not depend on Domain04 or G1 terminology.



\par\medskip\textbf{Consequences}\par


Human architects, human developers and AI development agents SHALL apply
this axiom to architecture derivation, documentation generation and
implementation.



Architectural traceability SHALL represent established relationships and
explicitly preserve unresolved or unestablished relationships.
Completeness SHALL NOT be manufactured through:



\begin{itemize}


\item inferred hierarchy;

\item structural convenience;

\item lexical similarity;

\item conventional modelling patterns;

\item assumed equivalence;

\item anticipated implementation;

\item familiar standards;

\item probable application design.

\end{itemize}


Candidates, suggestions and recommendations SHALL be distinguishable from
documented, approved architecture. A proposed resolution or apparently
better alternative SHALL NOT acquire authority through downstream use.



AI development agents SHOULD challenge architecture when evidence
warrants architectural review. They SHALL NOT silently bypass it because
they recognise a familiar destination.



This axiom applies the authority and explicit-review obligations in
\href{\detokenize{../../../../governance/architectural-axioms.md#1-purpose}}{Purpose}, \href{\detokenize{../../../../governance/architectural-axioms.md#6-relationship-to-architecture-decisions}}{Relationship to Architecture Decisions}
and \href{\detokenize{../../../../governance/architectural-axioms.md#7-architectural-review-rule}}{Architectural Review Rule} to downstream
derivation and implementation.



\hyperref[element:d3fc6025dc7646d1b2020b51f7a44bfa]{AX-14} preserves semantic
distinctions; this axiom preserves the distinction between established,
unresolved and proposed architecture.
\hyperref[element:ef8abbaec37e4890ad78b9f263e67906]{AX-15} governs uncertainty
about operational state, outcome or effect; this axiom governs uncertainty
about the architecture itself. Neither form of uncertainty may be silently
resolved by assumption.



\hyperref[element:eee8c3783d964d18a7f9dcea85fcd4cf]{AX-02}
and \hyperref[element:54508cda80484a6caf95b901780eaf7d]{AX-04}
preserve Harmonia's internal semantics from external standards and engine
assumptions. Familiar standards or technology patterns do not establish
otherwise undocumented architectural relationships.
\hyperref[element:952c3b6148ce4fce94b69f99a0c2fec7]{AX-06} governs the authority
of managed information; architectural authority remains a distinct concern.



\par\medskip\textbf{Maxim}\par


\textbf{Traceability must be truthful, not artificially complete.}


{\footnotesize \href{../../../../continuum/motivation/principle/8882573df7ce43c6bc0f66da0a676139.md}{Canonical record: 8882573df7ce43c6bc0f66da0a676139}; full qualifications and provenance.}\par
\section*{I.7.18 AX-18 Integration and Coordination Do Not Confer Domain Authority}
\phantomsection
\label{navigation:I.7.18}
\addcontentsline{toc}{section}{I.7.18 AX-18 Integration and Coordination Do Not Confer Domain Authority}

\label{element:b8cafe3cf5904d2bb9a9767260fddf78}
\index{AX-18 Integration and Coordination Do Not Confer Domain Authority (I.7.18)}
{\small Principle · Canonical · current}\par

\par\medskip\textbf{Axiom}\par


\begin{quote}


Harmonia MAY integrate information concerning, and coordinate operational activity associated with, an externally governed domain activity without thereby acquiring authority to determine or manage that domain activity.



Authority to determine or manage a domain activity SHALL remain with the responsible actor, system or governed capability unless that authority is explicitly assigned to Harmonia by the architecture.



Harmonia's observation, processing, persistence, distribution, presentation or use of information concerning a domain activity SHALL NOT, by itself, establish authority over that activity.



\end{quote}


\par\medskip\textbf{Rationale}\par


Participation in a domain activity and authority over that activity are
different architectural concerns. Information possession, observation,
integration, coordination and processing do not, by themselves, confer
domain authority. Authority requires explicit architectural assignment.



This new human-approved axiom generalises an already-established class of
architectural boundary. It is established at the 2026-10-09 migration
boundary; it is not a claim that AX-01 through AX-17 always contained this
axiom.



\par\medskip\textbf{Consequences and qualifications}\par


AX-18 SHALL NOT be interpreted as preventing Harmonia from performing
operational activity for which architecture has assigned Harmonia
responsibility. Explicit architectural assignment may establish authority
to determine or manage a domain activity.



In particular, AX-18 does not mean:



\begin{itemize}


\item Harmonia cannot execute Praxis/Ergo activity;

\item Harmonia cannot coordinate entity-specific operational activity through a Digital Twin;

\item Harmonia cannot process or transform clinical information;

\item Harmonia cannot perform workflow/activity coordination;

\item Harmonia cannot manage responsibilities explicitly assigned to Harmonia;

\item Harmonia cannot enforce policy, security, information governance or operational controls assigned to it.

\end{itemize}


Participation, information, integration and coordination do not, by
themselves, establish authority to determine or manage the underlying domain
activity. Assigned Harmonia operational responsibility remains applicable.



\par\medskip\textbf{Examples and established specialisations (non-normative)}\par


These examples explain the boundary; they do not expand the axiom or
establish Application Components, Information Concepts, Services, Processes,
AI allocation or implementation mechanisms.



\begin{itemize}


\item \textbf{Clinical work:} Harmonia may integrate and coordinate around clinical work
  without thereby acquiring authority to determine clinical work, manage
  clinical handover, own clinical worklists, allocate clinical work, make
  clinical decisions or become the EMR workflow authority. The detailed
  \href{\detokenize{../../../../../docs/markdown/03-business-architecture/metamodel/business-architecture-metamodel.md#37-clinical-work-integration-boundary}}{Domain03 clinical-work integration boundary}
  remains the Business specialisation.

\item \textbf{Digital Twins:} A Twin may coordinate operational activity associated
  with its real-world entity without thereby acquiring authority over
  originating clinical facts or domain decisions. The
  \href{\detokenize{../../../../../docs/markdown/02-strategy/strategic-views/logical-component-responsibilities.md#entity-specific-operational-coordination}}{Strategy entity-specific coordination rule}
  remains controlling for that coordination.

\item \textbf{Capacity:} Higher-order Service Capacity Management may consume
  context-specific capacity without acquiring the originating responsibility
  for Clinic, Ward, Theatre, ED, Outpatient or Mobile Service Capacity
  Management. The detailed
  \href{\detokenize{../../../../../docs/markdown/03-business-architecture/behaviours/04-health-service-operations.md#capacity-management-responsibility-boundary}}{Business capacity responsibility boundary}
  remains in Business Architecture.

\item \textbf{Assurance:} Observation, access to evidence or establishment of an
  assurance conclusion does not confer management authority over the activity
  being assured. Independent Assurance remains governed by
  \hyperref[element:bbb7438e78604308abba93afc14f4e2e]{REQ-FND-005}
  and its established
  \href{\detokenize{../../../../../docs/markdown/02-strategy/capability-maps/health-service-assurance-derivation.md}}{Strategy}
  and \href{\detokenize{../../../../../docs/markdown/03-business-architecture/behaviours/health-service-assurance.md}}{Business}
  derivation.

\item \textbf{Presentation / UI:} Presentation or collection of information does not,
  by itself, confer originating information or domain-decision authority on
  the presenting component. The existing
  \href{\detokenize{../../../../../docs/markdown/02-strategy/strategic-views/logical-component-responsibilities.md#component-6-iris-contextual-human-interaction}}{Iris contextual human interaction responsibility}
  remains distinct.

\item \textbf{AI:} Assistance with interpretation, pattern recognition, criteria
  evaluation or decision support does not, by itself, confer the authority of
  the governed actor, role or capability using that assistance. This example
  assigns no AI responsibility or capability.

\end{itemize}


\par\medskip\textbf{Related architectural authority}\par


\begin{itemize}


\item \hyperref[element:952c3b6148ce4fce94b69f99a0c2fec7]{AX-06} governs information authority and credibility.

\item \hyperref[element:b6a55decfeec4f1bb3d7f18534a2ffea]{AX-13} governs the Harmonia management boundary at ingress/egress.

\item \hyperref[element:d3fc6025dc7646d1b2020b51f7a44bfa]{AX-14} preserves meaningful semantic distinctions.

\item \hyperref[element:de7e90f46127455a9ec9ad3b9b48ba54]{AX-16} governs coordinated operational activity/entity progression.

\item \hyperref[element:8882573df7ce43c6bc0f66da0a676139]{AX-17} governs architectural authority and explicit uncertainty.

\item AX-18 governs inference of domain authority from participation, information,
  integration or coordination.

\end{itemize}

{\footnotesize \href{../../../../continuum/motivation/principle/b8cafe3cf5904d2bb9a9767260fddf78.md}{Canonical record: b8cafe3cf5904d2bb9a9767260fddf78}; full qualifications and provenance.}\par
\chapter*{I.8 Unclassified Motivation}
\phantomsection
\label{navigation:I.8}
\addcontentsline{toc}{chapter}{I.8 Unclassified Motivation}

\section*{I.8.1 At-Least-Once}
\phantomsection
\label{navigation:I.8.1}
\addcontentsline{toc}{section}{I.8.1 At-Least-Once}

\label{element:bd4c8726125648e68d5746bab3c765cd}
\index{At-Least-Once (I.8.1)}
{\small Principle · Historical / Reclassified — Domain06 · historical\_reference}\par
\href{../../../../continuum/motivation/principle/bd4c8726125648e68d5746bab3c765cd.md}{Historical canonical definition and provenance}.\par
\section*{I.8.2 Canonical Modeling with Standardized FHIR R5 Boundaries}
\phantomsection
\label{navigation:I.8.2}
\addcontentsline{toc}{section}{I.8.2 Canonical Modeling with Standardized FHIR R5 Boundaries}

\label{element:5d97974406c546078b1d64e5ff85bb13}
\index{Canonical Modeling with Standardized FHIR R5 Boundaries (I.8.2)}
{\small Principle · Historical / Reclassified · historical\_reference}\par
\href{../../../../continuum/motivation/principle/5d97974406c546078b1d64e5ff85bb13.md}{Historical canonical definition and provenance}.\par
\section*{I.8.3 Governed Asynchronous Change Management for Directory Entities}
\phantomsection
\label{navigation:I.8.3}
\addcontentsline{toc}{section}{I.8.3 Governed Asynchronous Change Management for Directory Entities}

\label{element:11001a6ea3df45689a27216df641dac2}
\index{Governed Asynchronous Change Management for Directory Entities (I.8.3)}
{\small Principle · Historical — current applicability unresolved · historical\_reference}\par
\href{../../../../continuum/motivation/principle/11001a6ea3df45689a27216df641dac2.md}{Historical canonical definition and provenance}.\par
\section*{I.8.4 Governed PHI Logging \& Dual-Logger Privacy Controls}
\phantomsection
\label{navigation:I.8.4}
\addcontentsline{toc}{section}{I.8.4 Governed PHI Logging \& Dual-Logger Privacy Controls}

\label{element:96ed278624c34ba4bfade0df08d45d74}
\index{Governed PHI Logging \& Dual-Logger Privacy Controls (I.8.4)}
{\small Principle · Historical / Reclassified · historical\_reference}\par
\href{../../../../continuum/motivation/principle/96ed278624c34ba4bfade0df08d45d74.md}{Historical canonical definition and provenance}.\par
\section*{I.8.5 Provenance and Information Authority}
\phantomsection
\label{navigation:I.8.5}
\addcontentsline{toc}{section}{I.8.5 Provenance and Information Authority}

\label{element:638c501ba4df4c67aee92d1f328aad84}
\index{Provenance and Information Authority (I.8.5)}
{\small Principle · Historical — current applicability unresolved · historical\_reference}\par
\href{../../../../continuum/motivation/principle/638c501ba4df4c67aee92d1f328aad84.md}{Historical canonical definition and provenance}.\par
\section*{I.8.6 Separation of Transport from Business Logic}
\phantomsection
\label{navigation:I.8.6}
\addcontentsline{toc}{section}{I.8.6 Separation of Transport from Business Logic}

\label{element:5ae28648334e4f6184044af0462eb9a8}
\index{Separation of Transport from Business Logic (I.8.6)}
{\small Principle · Historical / Reclassified · historical\_reference}\par
\href{../../../../continuum/motivation/principle/5ae28648334e4f6184044af0462eb9a8.md}{Historical canonical definition and provenance}.\par
\section*{I.8.7 Write-Behind}
\phantomsection
\label{navigation:I.8.7}
\addcontentsline{toc}{section}{I.8.7 Write-Behind}

\label{element:a5dffc0325d44915971eea600bba48f1}
\index{Write-Behind (I.8.7)}
{\small Principle · Retired / Superseded · historical\_reference}\par
\href{../../../../continuum/motivation/principle/a5dffc0325d44915971eea600bba48f1.md}{Historical canonical definition and provenance}.\par
\chapter*{I.9 Requirement Groups}
\phantomsection
\label{navigation:I.9}
\addcontentsline{toc}{chapter}{I.9 Requirement Groups}

\section*{I.9.1 [RequirementGroup] Functional Requirements}
\phantomsection
\label{navigation:I.9.1}
\addcontentsline{toc}{section}{I.9.1 [RequirementGroup] Functional Requirements}

\label{element:6acde68c43b14b6d8c7707779abdfd53}
\index{[RequirementGroup] Functional Requirements (I.9.1)}
{\small Grouping / RequirementGroup · Accepted bounded architecture — HCRL 1.1.15-001 DR-1 · current}\par

Functional Requirements grouping hierarchy root accepted under HCRL 1.1.15-001 DR-1.


{\small Recorded unresolved matters: canonicalName, friendlyName; exact scope remains in the canonical record.}\par
{\footnotesize \href{../../../../continuum/grouping/requirement-group/6acde68c43b14b6d8c7707779abdfd53.md}{Canonical record: 6acde68c43b14b6d8c7707779abdfd53}; full qualifications and provenance.}\par
\subsection*{I.9.1.1 [RequirementGroup] Information Management}
\phantomsection
\label{navigation:I.9.1.1}
\addcontentsline{toc}{subsection}{I.9.1.1 [RequirementGroup] Information Management}

\label{element:c34c366824ff452eb575345c16db720b}
\index{[RequirementGroup] Information Management (I.9.1.1)}
{\small Grouping / RequirementGroup · Accepted bounded architecture — HCRL 1.1.15-001 DR-1 · current}\par

Information management requirements grouping nested under Functional Requirements accepted under HCRL 1.1.15-001 DR-1.


{\small Recorded unresolved matters: canonicalName, friendlyName; exact scope remains in the canonical record.}\par
{\footnotesize \href{../../../../continuum/grouping/requirement-group/c34c366824ff452eb575345c16db720b.md}{Canonical record: c34c366824ff452eb575345c16db720b}; full qualifications and provenance.}\par
\subsubsection*{I.9.1.1.1 [Requirement] Information Durability During Platform Failure}
\phantomsection
\label{navigation:I.9.1.1.1}
\addcontentsline{toc}{subsubsection}{I.9.1.1.1 [Requirement] Information Durability During Platform Failure}

\label{element:6a2a3de03beb47b884293797b92a91eb}
\index{[Requirement] Information Durability During Platform Failure (I.9.1.1.1)}
{\small Requirement · Accepted bounded architecture — HCRL 1.1.15-001 DR-1 · current}\par

Intentionally incomplete requirement placeholder / thinking point accepted under HCRL 1.1.15-001 DR-1. Contained beneath [RequirementGroup] Information Management.


{\small Recorded unresolved matters: canonicalName, friendlyName; exact scope remains in the canonical record.}\par
{\footnotesize \href{../../../../continuum/motivation/requirement/6a2a3de03beb47b884293797b92a91eb.md}{Canonical record: 6a2a3de03beb47b884293797b92a91eb}; full qualifications and provenance.}\par
\section*{I.9.2 [RequirementGroup] Non-Functional-Requirements}
\phantomsection
\label{navigation:I.9.2}
\addcontentsline{toc}{section}{I.9.2 [RequirementGroup] Non-Functional-Requirements}

\label{element:f4a3060eba464f4ab89b50aec4628e5b}
\index{[RequirementGroup] Non-Functional-Requirements (I.9.2)}
{\small Grouping / RequirementGroup · Accepted bounded architecture — HCRL 1.1.15-001 DR-1 · current}\par

Non-Functional Requirements grouping hierarchy root accepted under HCRL 1.1.15-001 DR-1.


{\small Recorded unresolved matters: canonicalName, friendlyName; exact scope remains in the canonical record.}\par
{\footnotesize \href{../../../../continuum/grouping/requirement-group/f4a3060eba464f4ab89b50aec4628e5b.md}{Canonical record: f4a3060eba464f4ab89b50aec4628e5b}; full qualifications and provenance.}\par
\subsection*{I.9.2.1 [RequirementGroup] Performance Requirements}
\phantomsection
\label{navigation:I.9.2.1}
\addcontentsline{toc}{subsection}{I.9.2.1 [RequirementGroup] Performance Requirements}

\label{element:dc3f2a927e284da88734df4f7e0f9311}
\index{[RequirementGroup] Performance Requirements (I.9.2.1)}
{\small Grouping / RequirementGroup · Accepted bounded architecture — HCRL 1.1.15-001 DR-1 · current}\par

Performance requirements grouping nested under Non-Functional-Requirements accepted under HCRL 1.1.15-001 DR-1.


{\small Recorded unresolved matters: canonicalName, friendlyName; exact scope remains in the canonical record.}\par
{\footnotesize \href{../../../../continuum/grouping/requirement-group/dc3f2a927e284da88734df4f7e0f9311.md}{Canonical record: dc3f2a927e284da88734df4f7e0f9311}; full qualifications and provenance.}\par
\subsection*{I.9.2.2 [RequirementGroup] Disaster Recovery Requirements}
\phantomsection
\label{navigation:I.9.2.2}
\addcontentsline{toc}{subsection}{I.9.2.2 [RequirementGroup] Disaster Recovery Requirements}

\label{element:cb090e968a25409fb8ceb0fcc9a11ee7}
\index{[RequirementGroup] Disaster Recovery Requirements (I.9.2.2)}
{\small Grouping / RequirementGroup · Accepted bounded architecture — HCRL 1.1.15-001 DR-1 · current}\par

Disaster recovery requirements grouping nested under Non-Functional-Requirements accepted under HCRL 1.1.15-001 DR-1.


{\small Recorded unresolved matters: canonicalName, friendlyName; exact scope remains in the canonical record.}\par
{\footnotesize \href{../../../../continuum/grouping/requirement-group/cb090e968a25409fb8ceb0fcc9a11ee7.md}{Canonical record: cb090e968a25409fb8ceb0fcc9a11ee7}; full qualifications and provenance.}\par
\subsection*{I.9.2.3 [RequirementGroup] Availability}
\phantomsection
\label{navigation:I.9.2.3}
\addcontentsline{toc}{subsection}{I.9.2.3 [RequirementGroup] Availability}

\label{element:abf16578801a466b8f08d80ae292e644}
\index{[RequirementGroup] Availability (I.9.2.3)}
{\small Grouping / RequirementGroup · Accepted bounded architecture — HCRL 1.1.15-001 DR-1 · current}\par

Availability requirements grouping nested under Non-Functional-Requirements accepted under HCRL 1.1.15-001 DR-1.


{\small Recorded unresolved matters: canonicalName, friendlyName; exact scope remains in the canonical record.}\par
{\footnotesize \href{../../../../continuum/grouping/requirement-group/abf16578801a466b8f08d80ae292e644.md}{Canonical record: abf16578801a466b8f08d80ae292e644}; full qualifications and provenance.}\par
\subsubsection*{I.9.2.3.1 [RequirementGroup] Information Availability}
\phantomsection
\label{navigation:I.9.2.3.1}
\addcontentsline{toc}{subsubsection}{I.9.2.3.1 [RequirementGroup] Information Availability}

\label{element:1862d183cf7d45c4972e6aa3411d1193}
\index{[RequirementGroup] Information Availability (I.9.2.3.1)}
{\small Grouping / RequirementGroup · Accepted bounded architecture — HCRL 1.1.15-001 DR-1 · current}\par

Information availability requirements grouping nested under Availability accepted under HCRL 1.1.15-001 DR-1.


{\small Recorded unresolved matters: canonicalName, friendlyName; exact scope remains in the canonical record.}\par
{\footnotesize \href{../../../../continuum/grouping/requirement-group/1862d183cf7d45c4972e6aa3411d1193.md}{Canonical record: 1862d183cf7d45c4972e6aa3411d1193}; full qualifications and provenance.}\par

\part*{Part II — Strategy}
\label{navigation:II}
\addcontentsline{toc}{part}{Part II — Strategy}
The model separates Business Capabilities, Business Enabling Capabilities and Enterprise Capabilities within Strategy. Follow those existing branches to move from health-service responsibilities to their enabling concerns and platform capabilities. The folder called Application Layer here remains inside Strategy: it contains Enterprise Capabilities, and its navigation label does not reclassify those capabilities as application components.

{\small Reading sources: \hyperref[element:ed342cb8656b4f80b55ad519ecbd2d3a]{[CT1] Care Access \& Coordination}; \hyperref[element:b5a92874f8ac4509a30af5253590edc1]{[CT1] Entity Information Enablement}; \hyperref[element:15e57021490c4e738aee9009ff84a7c3]{[CT1] Governed Information Operations}.}\par
Health Service Governance illustrates a boundary that matters when reading capability descriptions: it establishes obligations, authority and compliance expectations, while operational management progresses means within that envelope and assurance evaluates evidence. Follow the capability's stated boundary and its qualified links rather than reading all neighbouring responsibilities as one owner. Proposed placement, relevance or enablement remains independently qualified even where the responsibility is canonical.

{\small Reading sources: \hyperref[element:209aada781e844e28202a2535a264c21]{[CT1] Health Service Governance}.}\par
\begin{center}\begin{tikzpicture}[>=Stealth,every node/.style={draw=Ink!40,rounded corners,align=left,font=\small,text width=5.2cm,inner sep=6pt}]
\node (source) at (0,0) {\hyperref[element:b13ee2c5a4ef4fd6be8a44795696dfc2]{[CT1] Health Service Operations Enablement}};
\node (target0) at (8,4.25) {\hyperref[element:f117ccec0bfa4f738e6ddb646b3534a2]{[CT2] Bed \& Care-Place Management}};
\draw[->] (source.east) -- (3,0) |- (target0.west);
\node[draw=none,text width=1.8cm,font=\tiny,align=center,inner sep=0pt] at (4.15,4.63) {Composition\newline{}Proposed / current};
\node (target1) at (8,2.55) {\hyperref[element:3dffe63e03124efab255204ee4e9a5d9]{[CT2] Clinic \& Practice Operations}};
\draw[->] (source.east) -- (3,0) |- (target1.west);
\node[draw=none,text width=1.8cm,font=\tiny,align=center,inner sep=0pt] at (4.15,2.93) {Composition\newline{}Proposed / current};
\node (target2) at (8,0.85) {\hyperref[element:8a99f4a6e5054d6c997bacec04805c01]{[CT2] Clinical Competency Administration}};
\draw[->] (source.east) -- (3,0) |- (target2.west);
\node[draw=none,text width=1.8cm,font=\tiny,align=center,inner sep=0pt] at (4.15,1.23) {Composition\newline{}Proposed / current};
\node (target3) at (8,-0.85) {\hyperref[element:0c27be11b4a44cb19e992309b8146a17]{[CT2] Clinical Logistics Coordination}};
\draw[->] (source.east) -- (3,0) |- (target3.west);
\node[draw=none,text width=1.8cm,font=\tiny,align=center,inner sep=0pt] at (4.15,-0.47) {Composition\newline{}Proposed / current};
\node (target4) at (8,-2.55) {\hyperref[element:0511a705d74a4f48a0ce92590ae764fa]{[CT2] Clinical Qualification Administration}};
\draw[->] (source.east) -- (3,0) |- (target4.west);
\node[draw=none,text width=1.8cm,font=\tiny,align=center,inner sep=0pt] at (4.15,-2.17) {Composition\newline{}Proposed / current};
\node (target5) at (8,-4.25) {\hyperref[element:192ca88a289544dcbf6adaf643f39969]{[CT2] Clinical Rostering}};
\draw[->] (source.east) -- (3,0) |- (target5.west);
\node[draw=none,text width=1.8cm,font=\tiny,align=center,inner sep=0pt] at (4.15,-3.87) {Composition\newline{}Proposed / current};
\end{tikzpicture}\end{center}
{\small Selected outgoing Composition assertions; arrows retain source-to-target direction. Labels retain assertion standing and record state; element standing accompanies each description. \href{structure.json}{Complete relationship qualifications and provenance} remain source-defined. No omitted edge or semantic effect is inferred.}\par
\chapter*{II.1 Capabilities}
\phantomsection
\label{navigation:II.1}
\addcontentsline{toc}{chapter}{II.1 Capabilities}

\section*{II.1.1 Enterprise Capabilities}
\phantomsection
\label{navigation:II.1.1}
\addcontentsline{toc}{section}{II.1.1 Enterprise Capabilities}

\subsection*{II.1.1.1 [CT1] Coordinated Activity \& Interaction}
\phantomsection
\label{navigation:II.1.1.1}
\addcontentsline{toc}{subsection}{II.1.1.1 [CT1] Coordinated Activity \& Interaction}

\label{element:5187c3db390e444db3eb809c02a2f1c1}
\index{[CT1] Coordinated Activity \& Interaction (II.1.1.1)}
{\small Capability · Proposed · current}\par

Enable governed interaction, event participation, boundary exchange and accountable activity progression.



\textbf{Overview}



Requests, events, work and human interaction are related ways to cause or observe change, with distinct contracts and progression responsibilities.



\textbf{Intent}



Make requests, event participation, work and human interaction progress according to their different contracts.



\textbf{Boundary}



Exchange is not execution, notification is not outcome, and interaction is not clinical authority.


{\footnotesize \href{../../../../continuum/strategy/capability/5187c3db390e444db3eb809c02a2f1c1.md}{Canonical record: 5187c3db390e444db3eb809c02a2f1c1}; full qualifications and provenance.}\par
\subsubsection*{II.1.1.1.1 [CT2] Activity \& Execution}
\phantomsection
\label{navigation:II.1.1.1.1}
\addcontentsline{toc}{subsubsection}{II.1.1.1.1 [CT2] Activity \& Execution}

\label{element:e5ff80e655084524a8dcee8e4ded788a}
\index{[CT2] Activity \& Execution (II.1.1.1.1)}
{\small Capability · Canonical · current}\par

\begin{itemize}


\item \textbf{Architectural Scope}: Define, instantiate, match, assign, dispatch, supervise, and progress coordinated operational and clinical activity units.

\item \textbf{Functional Responsibilities}:

\item Instantiating operational work items (Work Orders, To Dos, synthetic platform Tasks) in response to triggers.

\item Matching and assigning work units to eligible human or automated performers and roles based on policy rules.

\item Dispatching work units to their assigned performers with acknowledgement according to the applicable interaction contract.

\item Supervising execution progress, tracking timeouts, managing escalations, and handling activity interruptions.

\item \textbf{Decoupling Note}: Free of worker threads, thread pools, execution daemons, or specific workflow runtime engines.

\end{itemize}


Activity \& Execution provides reusable instantiation, matching, dispatch and supervision for governed activity units. Human Work Orders, human decision To Dos and automated Synthetic Tasks retain different performer and authority semantics while sharing a need for accountable progression. This reuse supports operational coordination without making a runtime engine the definition of the work.



\textbf{Boundary}



Coordination does not confer clinical work allocation or judgement. The capability specifies neither worker threads nor a particular execution engine.


{\footnotesize \href{../../../../continuum/strategy/capability/e5ff80e655084524a8dcee8e4ded788a.md}{Canonical record: e5ff80e655084524a8dcee8e4ded788a}; full qualifications and provenance.}\par
\paragraph*{II.1.1.1.1.1 [CT3] Entity-Relative Operational Progression}
\phantomsection
\label{navigation:II.1.1.1.1.1}
\addcontentsline{toc}{paragraph}{II.1.1.1.1.1 [CT3] Entity-Relative Operational Progression}

\label{element:69c82474fb534152b151c92f706197d0}
\index{[CT3] Entity-Relative Operational Progression (II.1.1.1.1.1)}
{\small Capability · Proposed · current}\par

Coordinate serial, entity-relative operational inputs and their interested behaviours.



\textbf{Overview}



Patient, practitioner and bed concerns recur across different business contexts; each can need a coherent progression frame without one central entity service.



\textbf{Intent}



Maintain a coherent operational reference frame as one entity reacts to successive inputs.



\textbf{Boundary}



Interest fan-out is distinct from intra-entity concurrency; entity progression does not acquire originating facts.


{\footnotesize \href{../../../../continuum/strategy/capability/69c82474fb534152b151c92f706197d0.md}{Canonical record: 69c82474fb534152b151c92f706197d0}; full qualifications and provenance.}\par
\paragraph*{II.1.1.1.1.2 [CT3] Governed Activity Execution Control}
\phantomsection
\label{navigation:II.1.1.1.1.2}
\addcontentsline{toc}{paragraph}{II.1.1.1.1.2 [CT3] Governed Activity Execution Control}

\label{element:1e351cdfdc1045e384c3bd768eed8e80}
\index{[CT3] Governed Activity Execution Control (II.1.1.1.1.2)}
{\small Capability · Proposed · current}\par

Instantiate, dispatch, supervise and progress authorised operational undertakings with explicit outcome standing.



\textbf{Overview}



Reusable execution control respects work-unit archetype, participation and uncertain effects rather than equating dispatch with completion.



\textbf{Intent}



Progress accountable operational work without equating dispatch, execution and successful outcome.



\textbf{Boundary}



Does not own underlying clinical-work meaning, clinical decisions or originator business authority.


{\footnotesize \href{../../../../continuum/strategy/capability/1e351cdfdc1045e384c3bd768eed8e80.md}{Canonical record: 1e351cdfdc1045e384c3bd768eed8e80}; full qualifications and provenance.}\par
\paragraph*{II.1.1.1.1.3 [CT3] Operational Fulfilment Definition}
\phantomsection
\label{navigation:II.1.1.1.1.3}
\addcontentsline{toc}{paragraph}{II.1.1.1.1.3 [CT3] Operational Fulfilment Definition}

\label{element:2b52c0de17b146da962f303e99f90a4a}
\index{[CT3] Operational Fulfilment Definition (II.1.1.1.1.3)}
{\small Capability · Proposed · current}\par

Define the context requirements, expected outcomes and housekeeping of an operational undertaking.



\textbf{Overview}



The semantic execution envelope applies to one-step and multi-step activity and stays distinct from dispatch and supervision.



\textbf{Intent}



Keep the semantic promise of an undertaking separate from the machinery that executes it.



\textbf{Boundary}



A fulfilment definition does not execute, orchestrate or choreograph its activity.


{\footnotesize \href{../../../../continuum/strategy/capability/2b52c0de17b146da962f303e99f90a4a.md}{Canonical record: 2b52c0de17b146da962f303e99f90a4a}; full qualifications and provenance.}\par
\subsubsection*{II.1.1.1.2 [CT2] Event \& Subscription}
\phantomsection
\label{navigation:II.1.1.1.2}
\addcontentsline{toc}{subsubsection}{II.1.1.1.2 [CT2] Event \& Subscription}

\label{element:b768b129971e4bf5bccc03d33d2fcb4e}
\index{[CT2] Event \& Subscription (II.1.1.1.2)}
{\small Capability · Canonical · current}\par

\begin{itemize}


\item \textbf{Architectural Scope}: Publish, ingest, correlate, filter, and syndicate events and subscription-driven changes across the enterprise landscape.

\item \textbf{Functional Responsibilities}:

\item Ingesting operational and clinical event triggers from external systems and internal state changes.

\item Managing topic- and criteria-based event subscriptions for internal and external consumers.

\item Evaluating dynamic subscription filters to match event payloads with interested subscriber criteria.

\item Delivering asynchronous event notifications reliably to registered subscribers.

\item \textbf{Decoupling Note}: Defined strictly in terms of event publishing, criteria matching, and notification delivery; entirely independent of broker topics, partitions, or queues.

\end{itemize}


Event \& Subscription connects significant changes with participants whose registered interests match them. It makes event ingestion, correlation, interest criteria and notification reusable across clinical and operational contexts. The architectural meaning is the relationship between a change and an interested consumer, rather than the broker constructs used to carry it.



\textbf{Boundary}



Receiving a notification does not determine the consumer's business response; topics and queues do not define the capability.


{\footnotesize \href{../../../../continuum/strategy/capability/b768b129971e4bf5bccc03d33d2fcb4e.md}{Canonical record: b768b129971e4bf5bccc03d33d2fcb4e}; full qualifications and provenance.}\par
\subsubsection*{II.1.1.1.3 [CT2] Interaction \& Experience}
\phantomsection
\label{navigation:II.1.1.1.3}
\addcontentsline{toc}{subsubsection}{II.1.1.1.3 [CT2] Interaction \& Experience}

\label{element:4ec9757b1c6947b4b70b0f1159b760d3}
\index{[CT2] Interaction \& Experience (II.1.1.1.3)}
{\small Capability · Canonical · current}\par

\begin{itemize}


\item \textbf{Architectural Scope}: Provide contextual interaction between human actors, Harmonia, and enterprise presentation services, including visualization, alerting, and collaboration integration.

\item \textbf{Functional Responsibilities}:

\item Mediating presentation interactions via dedicated backend-for-frontend gateways decoupled from internal databases.

\item Rendering clinical timelines, provider directory management interfaces, and operational console dashboards.

\item Providing contextual notifications and alerts to human users across web and mobile channels.

\item Projecting longitudinal record context into secure multidisciplinary collaboration spaces.

\item \textbf{Boundary Note}: Does not absorb the healthcare semantics of Clinical Collaboration or medical documentation. Consumes established identity and security context; does not own authentication, consent, or clinical truth.

\end{itemize}


Interaction \& Experience makes contextual information and governed actions available to human participants through presentation and collaboration integration. Clinical timelines, directory administration and operational dashboards have different purposes, but share a need to retain established identity, security and information context. The recurring responsibility is usable interaction within those boundaries.



\textbf{Boundary}



Presentation does not own consent, authentication or clinical truth, and collaboration projection does not absorb clinical documentation semantics.


{\footnotesize \href{../../../../continuum/strategy/capability/4ec9757b1c6947b4b70b0f1159b760d3.md}{Canonical record: 4ec9757b1c6947b4b70b0f1159b760d3}; full qualifications and provenance.}\par
\subsubsection*{II.1.1.1.4 [CT2] Interoperability \& Exchange}
\phantomsection
\label{navigation:II.1.1.1.4}
\addcontentsline{toc}{subsubsection}{II.1.1.1.4 [CT2] Interoperability \& Exchange}

\label{element:0f2e5222874949d8b0cce234941ecb20}
\index{[CT2] Interoperability \& Exchange (II.1.1.1.4)}
{\small Capability · Canonical · current}\par

\begin{itemize}


\item \textbf{Architectural Scope}: Resolve destinations, transform schemas, route, transport, deliver, acknowledge, and coordinate standards-based information exchange across enterprise boundaries.

\item \textbf{Functional Responsibilities}:

\item Providing multi-protocol boundary adaptation for healthcare standards (HL7 v2, FHIR REST, DICOM, secure messaging).

\item Executing bidirectional structural and semantic transformations between legacy formats and canonical representations.

\item Ensuring reliable addressed delivery and fan-out distribution to external clinical endpoints.

\item Processing and distinguishing technical delivery acknowledgements from business acknowledgements.

\item \textbf{Boundary Rule}: Egress terminates Harmonia management of the emitted representation (Axiom AX-13).

\end{itemize}


Interoperability \& Exchange moves information across participating boundaries through destination resolution, transformation, delivery and acknowledgements. It preserves the difference between technical transmission and the recipient's business response so that an exchange can be understood under its actual contract. Standards govern the external representation while internal management semantics remain private.



\textbf{Boundary}



Egress ends Harmonia management of the emitted representation. Retained delivery evidence does not extend control over an external recipient or transfer payload authority.


{\footnotesize \href{../../../../continuum/strategy/capability/0f2e5222874949d8b0cce234941ecb20.md}{Canonical record: 0f2e5222874949d8b0cce234941ecb20}; full qualifications and provenance.}\par
\subsection*{II.1.1.2 [CT1] Governed Information Operations}
\phantomsection
\label{navigation:II.1.1.2}
\addcontentsline{toc}{subsection}{II.1.1.2 [CT1] Governed Information Operations}

\label{element:15e57021490c4e738aee9009ff84a7c3}
\index{[CT1] Governed Information Operations (II.1.1.2)}
{\small Capability · Proposed · current}\par

Maintain identifiable, meaningful, context-qualified and retrievable information through valid state progression.



\textbf{Overview}



Entity integrity, context, lifecycle, information maintenance, discovery and semantic conformance are constituent reusable abilities; none is an implementation subsystem.



\textbf{Intent}



Keep information identity, meaning, context, state and retrieval compatible across diverse healthcare uses.



\textbf{Boundary}



Information management and semantic governance do not grant originating domain authority or erase active/durable state separation.


{\footnotesize \href{../../../../continuum/strategy/capability/15e57021490c4e738aee9009ff84a7c3.md}{Canonical record: 15e57021490c4e738aee9009ff84a7c3}; full qualifications and provenance.}\par
\subsubsection*{II.1.1.2.1 [CT2] Context Management}
\phantomsection
\label{navigation:II.1.1.2.1}
\addcontentsline{toc}{subsubsection}{II.1.1.2.1 [CT2] Context Management}

\label{element:3cfd3e56c23a4943bd1b476ec248ac49}
\index{[CT2] Context Management (II.1.1.2.1)}
{\small Capability · Canonical · current}\par

\begin{itemize}


\item \textbf{Architectural Scope}: Establish, validate, propagate, and preserve the execution, clinical, security, and administrative context within which information and activities have meaning.

\item \textbf{Functional Responsibilities}:

\item Binding the subject of care, encounter, ordering provider, and healthcare service context to in-flight messages and transactions.

\item Propagating tamper-evident, attributable security context across asynchronous, multi-hop operational boundaries.

\item Ensuring context preservation during cross-protocol transformations across disparate healthcare exchange formats and protocols.

\item Detecting and rejecting requests where mandatory operational or clinical context is missing or inconsistent.

\item \textbf{Collaborative Nature}: Bounded reusable context responsibility may be realised through participation across gateways, activity execution and information-access boundaries. Participation does not transfer semantic ownership.

\end{itemize}


Context Management preserves the circumstances in which information and activity have meaning as they cross processing boundaries. Subject, encounter, provider, service and security context must remain attributable and consistent rather than be reconstructed from convenient local assumptions. Its reuse is especially significant in multi-hop activity where a syntactically valid message can otherwise lose the context required for legitimate interpretation.



\textbf{Boundary}



Context participation does not transfer semantic ownership. Generic context management and assurance-specific context semantics contribute together without allocating a new central Context Service.


{\footnotesize \href{../../../../continuum/strategy/capability/3cfd3e56c23a4943bd1b476ec248ac49.md}{Canonical record: 3cfd3e56c23a4943bd1b476ec248ac49}; full qualifications and provenance.}\par
\subsubsection*{II.1.1.2.2 [CT2] Information Management}
\phantomsection
\label{navigation:II.1.1.2.2}
\addcontentsline{toc}{subsubsection}{II.1.1.2.2 [CT2] Information Management}

\label{element:43051efd406b4040bff8e794b7a631e0}
\index{[CT2] Information Management (II.1.1.2.2)}
{\small Capability · Canonical · current}\par

\begin{itemize}


\item \textbf{Architectural Scope}: Validate, represent, persist, retrieve, version, and manage governed clinical, administrative, and operational information assets.

\item \textbf{Functional Responsibilities}:

\item Validating inbound information structures against canonical information constraints.

\item Providing authoritative durable persistence for historical health records and administrative registries.

\item Managing active distributed working memory representations for low-latency retrieval.

\item Governing document versioning, supersession, amendments, and retractions without data loss.

\item \textbf{Architectural Principle}: Enforces the foundational separation between \emph{active distributed state access} and \emph{authoritative durable state preservation} (Axiom AX-05).

\end{itemize}


Information Management provides the recurring ability to represent and maintain governed information through validation, preservation, retrieval and versioned change. It accommodates both responsive active use and authoritative durable preservation while keeping those responsibilities separate. This distinction allows availability and working state to support care without turning a cache into durable truth.



\textbf{Boundary}



Engine capabilities provide machinery under Harmonia semantics. Persistence success does not establish information authority or make the platform the originator of the represented facts.


{\footnotesize \href{../../../../continuum/strategy/capability/43051efd406b4040bff8e794b7a631e0.md}{Canonical record: 43051efd406b4040bff8e794b7a631e0}; full qualifications and provenance.}\par
\paragraph*{II.1.1.2.2.1 [CT3] Active Information Availability}
\phantomsection
\label{navigation:II.1.1.2.2.1}
\addcontentsline{toc}{paragraph}{II.1.1.2.2.1 [CT3] Active Information Availability}

\label{element:bb711f4d8fb74af2917869b9b05f7f84}
\index{[CT3] Active Information Availability (II.1.1.2.2.1)}
{\small Capability · Proposed · current}\par

Maintain responsive, reconstructable active information and coordinate its safe use.



\textbf{Overview}



Active representation, observation and convergence support availability while preserving an explicit trust basis for use.



\textbf{Intent}



Provide responsive managed use without promoting active representation into durable truth.



\textbf{Boundary}



An active generation is not an authoritative version; failure to converge does not establish valid active state.


{\footnotesize \href{../../../../continuum/strategy/capability/bb711f4d8fb74af2917869b9b05f7f84.md}{Canonical record: bb711f4d8fb74af2917869b9b05f7f84}; full qualifications and provenance.}\par
\paragraph*{II.1.1.2.2.2 [CT3] Authoritative Durable State Establishment}
\phantomsection
\label{navigation:II.1.1.2.2.2}
\addcontentsline{toc}{paragraph}{II.1.1.2.2.2 [CT3] Authoritative Durable State Establishment}

\label{element:0d5bd68ff50f410aaceed05377e1a4df}
\index{[CT3] Authoritative Durable State Establishment (II.1.1.2.2.2)}
{\small Capability · Proposed · current}\par

Atomically establish durable managed state and its authoritative version progression.



\textbf{Overview}



Durable acceptance gives managed information an explicit committed state from which applicable active representations can be reconstructed.



\textbf{Intent}



Give managed change an explicit committed state and authoritative progression basis.



\textbf{Boundary}



Commit is distinct from semantic source authority and from successful active convergence; applications do not gain a raw persistence path.


{\footnotesize \href{../../../../continuum/strategy/capability/0d5bd68ff50f410aaceed05377e1a4df.md}{Canonical record: 0d5bd68ff50f410aaceed05377e1a4df}; full qualifications and provenance.}\par
\paragraph*{II.1.1.2.2.3 [CT3] Binary Content Association \& Resolution}
\phantomsection
\label{navigation:II.1.1.2.2.3}
\addcontentsline{toc}{paragraph}{II.1.1.2.2.3 [CT3] Binary Content Association \& Resolution}

\label{element:5314959509dc4b558711ec169ddee5a0}
\index{[CT3] Binary Content Association \& Resolution (II.1.1.2.2.3)}
{\small Capability · Proposed · current}\par

Maintain content associations and custody knowledge sufficient for governed payload resolution, retrieval and assembly.



\textbf{Overview}



Document attachments, diagnostic media and operational artefacts share a need to locate usable content while distinguishing logical meaning from physical possession.



\textbf{Intent}



Make logical information and its usable payloads consistently associable while retaining custody boundaries.



\textbf{Boundary}



No universal binary custody, storage topology or persistence authority is implied.


{\footnotesize \href{../../../../continuum/strategy/capability/5314959509dc4b558711ec169ddee5a0.md}{Canonical record: 5314959509dc4b558711ec169ddee5a0}; full qualifications and provenance.}\par
\subsubsection*{II.1.1.2.3 [CT2] Managed Entity \& Relationship}
\phantomsection
\label{navigation:II.1.1.2.3}
\addcontentsline{toc}{subsubsection}{II.1.1.2.3 [CT2] Managed Entity \& Relationship}

\label{element:910634fb1621416c818a022ab0b5d3a1}
\index{[CT2] Managed Entity \& Relationship (II.1.1.2.3)}
{\small Capability · Canonical · current}\par

\begin{itemize}


\item \textbf{Architectural Scope}: Identify, classify, relate, and maintain the structural integrity of managed healthcare and administrative entities throughout the enterprise landscape.

\item \textbf{Functional Responsibilities}:

\item Uniquely identifying core domain entities (Persons, Practitioners, Roles, Organisations, Locations, Services, Devices).

\item Resolving and maintaining multi-authority identifier systems and aliases without destructive overwriting.

\item Representing and traversing multi-dimensional relationships (affiliations, organizational hierarchies, physical location trees, care team memberships).

\item Maintaining referential integrity across interconnected entity graphs.

\item \textbf{Anti-Responsibilities}: Does not own physical database relational tables or foreign key constraints (Domain 04/07).

\end{itemize}


This reusable capability keeps managed healthcare entities identifiable and their relationships structurally intelligible. It supplies common identity, alias and relationship behaviour recurring in persons, practitioners, organisations, locations and devices. Reuse concerns those architectural responsibilities, allowing the healthcare-specific meaning of each entity to remain with its governing context.



\textbf{Boundary}



Relationship integrity is not a physical database design. Correlating entities does not establish originating clinical identity authority or a new standalone Correlation capability.


{\footnotesize \href{../../../../continuum/strategy/capability/910634fb1621416c818a022ab0b5d3a1.md}{Canonical record: 910634fb1621416c818a022ab0b5d3a1}; full qualifications and provenance.}\par
\subsubsection*{II.1.1.2.4 [CT2] Managed State \& Lifecycle}
\phantomsection
\label{navigation:II.1.1.2.4}
\addcontentsline{toc}{subsubsection}{II.1.1.2.4 [CT2] Managed State \& Lifecycle}

\label{element:6296abeee5074b9e80f21aea3a297e42}
\index{[CT2] Managed State \& Lifecycle (II.1.1.2.4)}
{\small Capability · Canonical · current}\par

\begin{itemize}


\item \textbf{Architectural Scope}: Govern the current and historical state, effective temporal periods, state transition validations, and lifecycle progression of managed integration artifacts and entities.

\item \textbf{Functional Responsibilities}:

\item Defining and enforcing valid state transition state machines (e.g., \texttt{\seqsplit{REQUESTED}} \$\textbackslash{}to\$ \texttt{\seqsplit{VALIDATED}} \$\textbackslash{}to\$ \texttt{\seqsplit{IN\_PROGRESS}} \$\textbackslash{}to\$ \texttt{\seqsplit{COMPLETED}}).

\item Maintaining temporal validity intervals (effective start and end dates) for relationships, credentials, and states.

\item Supporting non-destructive state progression, preserving historical state versions for longitudinal review.

\item Rejecting illegal or out-of-order state transitions while accommodating idempotent replay.

\item \textbf{Anti-Responsibilities}: Does not define thread-level mutexes or database transaction isolation levels.

\end{itemize}


Managed State \& Lifecycle provides reusable responsibility for valid progression, temporal applicability and non-destructive change. A directory relationship, clinical document and operational work item have different business states, yet all need an intelligible distinction between current standing, historical state and invalid transitions. The capability captures that recurrence without replacing each subject's lifecycle meaning.



\textbf{Boundary}



Business lifecycle integrity does not prescribe locks, database isolation or worker implementation, and illustrative state names do not define every subject's transitions.


{\footnotesize \href{../../../../continuum/strategy/capability/6296abeee5074b9e80f21aea3a297e42.md}{Canonical record: 6296abeee5074b9e80f21aea3a297e42}; full qualifications and provenance.}\par
\subsubsection*{II.1.1.2.5 [CT2] Search \& Discovery}
\phantomsection
\label{navigation:II.1.1.2.5}
\addcontentsline{toc}{subsubsection}{II.1.1.2.5 [CT2] Search \& Discovery}

\label{element:6a45701634f04b43919d9152a61e5f2e}
\index{[CT2] Search \& Discovery (II.1.1.2.5)}
{\small Capability · Canonical · current}\par

\begin{itemize}


\item \textbf{Architectural Scope}: Index, search, resolve, filter, and discover information assets, entities, directory endpoints, and healthcare services across the enterprise.

\item \textbf{Functional Responsibilities}:

\item Providing multi-criteria indexed search across complex entity graphs (e.g., find active cardiologists at Location X accepting new referrals).

\item Resolving service and practitioner communication endpoints for addressed delivery.

\item Supporting bounded, paginated clinical queries across longitudinal patient records.

\item Applying dynamic security and consent filters to search results prior to egress.

\item \textbf{Anti-Responsibilities}: Does not dictate specific search engine indexes, inverted index implementations, or SQL query syntaxes.

\end{itemize}


Search \& Discovery makes information, entities and service endpoints findable through governed criteria and bounded queries. It addresses the recurring need to locate relevant material or a destination before it can be used. Security-qualified results preserve the distinction between knowing that something exists and being entitled to obtain or disclose it.



\textbf{Boundary}



Discovery does not decide policy or prescribe a search engine, index format or query language.


{\footnotesize \href{../../../../continuum/strategy/capability/6a45701634f04b43919d9152a61e5f2e.md}{Canonical record: 6a45701634f04b43919d9152a61e5f2e}; full qualifications and provenance.}\par
\subsubsection*{II.1.1.2.6 [CT2] Semantic Governance \& Conformance}
\phantomsection
\label{navigation:II.1.1.2.6}
\addcontentsline{toc}{subsubsection}{II.1.1.2.6 [CT2] Semantic Governance \& Conformance}

\label{element:9ee41437bd4c4bd18389f05c12108f97}
\index{[CT2] Semantic Governance \& Conformance (II.1.1.2.6)}
{\small Capability · Canonical · current}\par

\begin{itemize}


\item \textbf{Architectural Scope}: Govern canonical models, schemas, terminologies, value sets, mapping tables, constraint rules, and semantic conformance validation across the enterprise.

\item \textbf{Functional Responsibilities}:

\item Defining and publishing canonical healthcare data models and profile structures.

\item Governing clinical terminology mappings, value sets, and concept lookups (SNOMED CT, LOINC).

\item Validating inbound and outbound message representations against published semantic rules.

\item Governing semantic versioning, deprecation, and backward compatibility across schema revisions.

\item \textbf{Component Relationship}: Informs the Calliope strategic logical component, but Calliope is the responsibility centre, not the capability definition.

\end{itemize}


This capability makes canonical information models, terminology bindings, mappings and conformance rules governable across enterprise use. It addresses both what representations mean and whether they respect the published rules, keeping version and compatibility concerns explicit. A conformant representation can therefore be distinguished from one that merely uses a supported transport or parser.



\textbf{Boundary}



Conformance does not establish clinical truth, information authority or disclosure permission. Calliope is an established responsibility centre, not the definition of the capability.


{\footnotesize \href{../../../../continuum/strategy/capability/9ee41437bd4c4bd18389f05c12108f97.md}{Canonical record: 9ee41437bd4c4bd18389f05c12108f97}; full qualifications and provenance.}\par
\subsection*{II.1.1.3 [CT1] Trust, Accountability \& Assurance}
\phantomsection
\label{navigation:II.1.1.3}
\addcontentsline{toc}{subsection}{II.1.1.3 [CT1] Trust, Accountability \& Assurance}

\label{element:321de197612f49089d8b4a4cc1a2e003}
\index{[CT1] Trust, Accountability \& Assurance (II.1.1.3)}
{\small Capability · Proposed · current}\par

Make governed activity permissible, attributable, operationally dependable and independently assessable.



\textbf{Overview}



Policy, evidence, resilience and independent assurance contribute to confidence through different responsibilities. Keeping them as peers prevents technical success or recorded evidence from becoming a business assurance conclusion.



\textbf{Intent}



Keep permission, attribution, operational integrity and independent confidence distinguishable and usable together.



\textbf{Boundary}



Cross-cutting participation does not distribute semantic ownership, mandate centralisation or authorise the assured subject to determine its own conclusion.


{\footnotesize \href{../../../../continuum/strategy/capability/321de197612f49089d8b4a4cc1a2e003.md}{Canonical record: 321de197612f49089d8b4a4cc1a2e003}; full qualifications and provenance.}\par
\subsubsection*{II.1.1.3.1 [CT2] Information Authority Qualification}
\phantomsection
\label{navigation:II.1.1.3.1}
\addcontentsline{toc}{subsubsection}{II.1.1.3.1 [CT2] Information Authority Qualification}

\label{element:cbba9984dd404deabed41fdcc96088b1}
\index{[CT2] Information Authority Qualification (II.1.1.3.1)}
{\small Capability · Proposed · current}\par

Represent and qualify the source-authority and credibility standing of information independently of transport, persistence and concurrency.



\textbf{Overview}



Identity correlation, clinical-history synthesis, reference-content use and assurance evidence all need to distinguish the standing of an assertion from technical acceptance. This recurring ability is insufficiently explicit in the existing fourteen EC definitions and is proposed as a new peer responsibility.



\textbf{Intent}



Keep source credibility explicit even when the technical handling of competing assertions is equally valid.



\textbf{Boundary}



This is not independent service assurance, a universal conflict arbiter or clinical truth determination. Existing EC-13 governs authority nomenclature; this capability applies and preserves the qualified standing.


{\footnotesize \href{../../../../continuum/strategy/capability/cbba9984dd404deabed41fdcc96088b1.md}{Canonical record: cbba9984dd404deabed41fdcc96088b1}; full qualifications and provenance.}\par
\subsubsection*{II.1.1.3.2 [CT2] Operational Assurance}
\phantomsection
\label{navigation:II.1.1.3.2}
\addcontentsline{toc}{subsubsection}{II.1.1.3.2 [CT2] Operational Assurance}

\label{element:938cd44380ab431fa8dae13c702d83c1}
\index{[CT2] Operational Assurance (II.1.1.3.2)}
{\small Capability · Canonical · current}\par

\begin{itemize}


\item \textbf{Architectural Scope}: Ensure concurrency integrity, system resilience, duplicate detection, exception handling, automated recovery, and operational telemetry across platform operations.

\item \textbf{Functional Responsibilities}:

\item Enforcing sliding-window duplicate message detection and idempotent suppression.

\item Guaranteeing ingress dual-write safety (durable acceptance before emitting positive application ACK).

\item Managing structured failure recovery, exception escalation, and dead-letter isolation.

\item Emitting platform telemetry, operational health metrics, and queue processing metrics.

\item \textbf{Collaborative Nature}: Realized collaboratively across boundary gateways, messaging infrastructure, execution workers, and storage layers.

\end{itemize}


\par\medskip\textbf{EC-12 / EC-14 Boundary}\par


EC-12 retains the established scope and responsibilities above. It may contribute directly to generic operational/processing assurance, limited to that contribution in the \href{\detokenize{../../../../../docs/markdown/02-strategy/capability-maps/health-service-assurance-derivation.md#collaborative-enterprise-capability-contribution-matrix}}{approved assurance matrix}. It does not own the general semantics of assurance disposition, assurance-specific criteria applicability, evidentiary association, assessment against assurance criteria, adjudication, or independent assurance findings and conclusions. Those distinguishing responsibilities are introduced through EC-14 Service Guardian. Operational health and business-level processing assurance remain distinct; EC-12 is not expanded to absorb EC-14.



Operational Assurance preserves processing integrity under concurrency, duplicate delivery and failure. It brings together resilience, durable ingress acceptance, recovery and operational telemetry so that technical degradation does not silently weaken the managed operation's semantics. Its contribution is collaborative across operational boundaries while remaining a bounded architectural responsibility.



\textbf{Boundary}



Operational health is distinct from independent assessment of a governed subject. Assurance-specific criteria applicability, adjudication and conclusions belong to the established EC-14 contribution.


{\footnotesize \href{../../../../continuum/strategy/capability/938cd44380ab431fa8dae13c702d83c1.md}{Canonical record: 938cd44380ab431fa8dae13c702d83c1}; full qualifications and provenance.}\par
\subsubsection*{II.1.1.3.3 [CT2] Policy \& Control}
\phantomsection
\label{navigation:II.1.1.3.3}
\addcontentsline{toc}{subsubsection}{II.1.1.3.3 [CT2] Policy \& Control}

\label{element:0c5d986cc948421a8bebcab3b74d108f}
\index{[CT2] Policy \& Control (II.1.1.3.3)}
{\small Capability · Canonical · current}\par

\begin{itemize}


\item \textbf{Architectural Scope}: Evaluate and enforce security, consent, privacy, and regulatory policies including authority, access control, and information handling rules.

\item \textbf{Functional Responsibilities}:

\item Executing impartial, default-deny policy evaluations for all integration transactions.

\item Enforcing Role-Based and Attribute-Based Access Control (RBAC/ABAC) policies.

\item Evaluating patient consent directives and privacy restrictions at boundary and query checkpoints.

\item Enforcing statutory information disclosure boundaries (e.g., mental health, VIP records).

\item \textbf{Collaborative Nature}: Evaluated across boundary gateways, internal bus dispatch, storage interfaces, and presentation layers.

\end{itemize}


Policy \& Control supplies the reusable evaluation and enforcement of authority, consent, privacy and handling constraints. Its purpose is consistent governed action across ingress, information access, execution and presentation boundaries. Default-deny evaluation preserves the need for established authority rather than treating a caller's intent or a technically valid request as permission.



\textbf{Boundary}



Participation at multiple enforcement points does not distribute policy ownership or turn operational security context into persistent clinical content.


{\footnotesize \href{../../../../continuum/strategy/capability/0c5d986cc948421a8bebcab3b74d108f.md}{Canonical record: 0c5d986cc948421a8bebcab3b74d108f}; full qualifications and provenance.}\par
\paragraph*{II.1.1.3.3.1 [CT3] Operation Policy Resolution}
\phantomsection
\label{navigation:II.1.1.3.3.1}
\addcontentsline{toc}{paragraph}{II.1.1.3.3.1 [CT3] Operation Policy Resolution}

\label{element:7eca9fc4d2244b0facc83376e57ebbc4}
\index{[CT3] Operation Policy Resolution (II.1.1.3.3.1)}
{\small Capability · Proposed · current}\par

Resolve effective, temporally applicable permission criteria for an established identity and governed operation.



\textbf{Overview}



Permission depends on controlled criteria, role assignments, action and information scope; revocation is distinct from natural expiry.



\textbf{Intent}



Determine the effective permitted action under the applicable identity, information and temporal criteria.



\textbf{Boundary}



Does not authenticate identity or grant permission from caller-controlled tags or mere role names.


{\footnotesize \href{../../../../continuum/strategy/capability/7eca9fc4d2244b0facc83376e57ebbc4.md}{Canonical record: 7eca9fc4d2244b0facc83376e57ebbc4}; full qualifications and provenance.}\par
\paragraph*{II.1.1.3.3.2 [CT3] Operation-Bound Policy Enforcement}
\phantomsection
\label{navigation:II.1.1.3.3.2}
\addcontentsline{toc}{paragraph}{II.1.1.3.3.2 [CT3] Operation-Bound Policy Enforcement}

\label{element:cb32b32c657445b3b9b27bebc069ace9}
\index{[CT3] Operation-Bound Policy Enforcement (II.1.1.3.3.2)}
{\small Capability · Proposed · current}\par

Enforce resolved policy constraints at the boundary where a governed action is attempted.



\textbf{Overview}



Distributed participants apply the applicable effective authorisation and fail closed when authority is absent or untrustworthy.



\textbf{Intent}



Ensure an operation cannot proceed merely because its caller omitted or bypassed a voluntary policy check.



\textbf{Boundary}



Enforcement participation does not distribute policy ownership or make security context persistent content.


{\footnotesize \href{../../../../continuum/strategy/capability/cb32b32c657445b3b9b27bebc069ace9.md}{Canonical record: cb32b32c657445b3b9b27bebc069ace9}; full qualifications and provenance.}\par
\subsubsection*{II.1.1.3.4 [CT2] Provenance \& Traceability}
\phantomsection
\label{navigation:II.1.1.3.4}
\addcontentsline{toc}{subsubsection}{II.1.1.3.4 [CT2] Provenance \& Traceability}

\label{element:68a0755aa4e14a509ac5e6c0bbcff3e7}
\index{[CT2] Provenance \& Traceability (II.1.1.3.4)}
{\small Capability · Canonical · current}\par

\begin{itemize}


\item \textbf{Architectural Scope}: Preserve required origin attribution, meaningful transformation history, significant execution checkpoints, tamper-evident audit evidence and end-to-end traceability. Harmonia SHALL preserve provenance, audit and operational evidence where required by governing obligations, architectural significance, accountability, reconstruction, security or assurance needs. Routine technical activity SHALL NOT automatically require durable evidentiary recording.

\item \textbf{Functional Responsibilities}:

\item Recording attributable, verifiable origin metadata for information sources, authors, and timestamps.

\item Tracking information-significant transformation history and required mapping checkpoints as information crosses system seams.

\item Capturing non-PHI operational and security audit trails for regulatory compliance (HIPAA, GDPR).

\item Supporting end-to-end transaction tracing across asynchronous processing pipelines.

\item \textbf{Collaborative Nature}: Participating elements preserve required provenance and audit context according to policy. Telemetry, operational facts, provenance, audit and assurance evidence remain distinct. Being evidence is contextual, not an intrinsic property of every recorded fact; preservation does not establish an assurance conclusion.

\end{itemize}


Provenance \& Traceability preserves the meaningful attribution and history needed to explain information and significant activity. It connects source, transformation and required checkpoints without making every internal mechanism permanent evidence. Its recurring contribution supports accountability, reconstruction and assurance use while retaining the context that makes a recorded fact evidentially relevant.



\textbf{Boundary}



Telemetry, provenance, audit and assurance evidence remain distinct. Recording evidence neither establishes an assurance conclusion nor automatically requires retention of routine technical activity.


{\footnotesize \href{../../../../continuum/strategy/capability/68a0755aa4e14a509ac5e6c0bbcff3e7.md}{Canonical record: 68a0755aa4e14a509ac5e6c0bbcff3e7}; full qualifications and provenance.}\par
\paragraph*{II.1.1.3.4.1 [CT3] Governed Activity Audit Evidence}
\phantomsection
\label{navigation:II.1.1.3.4.1}
\addcontentsline{toc}{paragraph}{II.1.1.3.4.1 [CT3] Governed Activity Audit Evidence}

\label{element:c8d93ddc5aa84960964b78345fff1b29}
\index{[CT3] Governed Activity Audit Evidence (II.1.1.3.4.1)}
{\small Capability · Proposed · current}\par

Capture required significant activity and decision evidence under its applicable governed capture authority.



\textbf{Overview}



Denied actions and successful mutations may both require trustworthy evidence; the audited identity and capture authority remain distinct.



\textbf{Intent}



Preserve required significant action evidence under its proper capture authority.



\textbf{Boundary}



Mandatory capture authority is specific to the established audit responsibility, not a platform-wide permission shortcut.


{\footnotesize \href{../../../../continuum/strategy/capability/c8d93ddc5aa84960964b78345fff1b29.md}{Canonical record: c8d93ddc5aa84960964b78345fff1b29}; full qualifications and provenance.}\par
\paragraph*{II.1.1.3.4.2 [CT3] Information Attribution \& Lineage}
\phantomsection
\label{navigation:II.1.1.3.4.2}
\addcontentsline{toc}{paragraph}{II.1.1.3.4.2 [CT3] Information Attribution \& Lineage}

\label{element:8bc66c9bc1024a68a1bfd010b4dd1b79}
\index{[CT3] Information Attribution \& Lineage (II.1.1.3.4.2)}
{\small Capability · Proposed · current}\par

Preserve source attribution and significant transformation relationships for managed information.



\textbf{Overview}



Lineage explains how an assertion reached its present form without turning every infrastructure event into provenance.



\textbf{Intent}



Explain the information-significant origin and transformation of an assertion.



\textbf{Boundary}



Does not infer clinical credibility from a successful transport or establish independent assurance findings.


{\footnotesize \href{../../../../continuum/strategy/capability/8bc66c9bc1024a68a1bfd010b4dd1b79.md}{Canonical record: 8bc66c9bc1024a68a1bfd010b4dd1b79}; full qualifications and provenance.}\par
\subsubsection*{II.1.1.3.5 [CT2] Service Guardian}
\phantomsection
\label{navigation:II.1.1.3.5}
\addcontentsline{toc}{subsubsection}{II.1.1.3.5 [CT2] Service Guardian}

\label{element:d767f12f9b1c495883216ace464ee1c2}
\index{[CT2] Service Guardian (II.1.1.3.5)}
{\small Capability · Canonical · current}\par

\begin{itemize}


\item \textbf{Architectural Scope}: Service Guardian provides the reusable capability to independently establish whether governed service behaviour, state or outcomes remain within defined assurance expectations by associating applicable criteria and trustworthy evidence with a governed subject, assessing that evidence, adjudicating the assessment, and establishing reportable assurance findings and conclusions.

\item \textbf{Functional Responsibilities}:

\item Supplying assurance-specific semantics for assurance disposition and assurance criteria applicability.

\item Associating a governed subject/context with applicable criteria and sufficient trustworthy evidence for a particular assurance purpose.

\item Assessing evidence against applicable criteria and adjudicating assessed detail to establish reportable findings, exceptions and conclusions.

\item Preserving assurance independence, temporal/version-specific evidence and criteria applicability, and explicit inability to establish a conclusion where evidence is insufficient.

\item \textbf{Collaborative Nature}: Supplies the assurance-specific semantics required to bind existing reusable capabilities into governed assurance. It contributes directly to Assurance Design, Assurance Criteria Management and Governed Assurance; direct contribution is not sole realisation. The \href{\detokenize{../../../../../docs/markdown/02-strategy/capability-maps/health-service-assurance-derivation.md#collaborative-enterprise-capability-contribution-matrix}}{approved matrix} preserves the contributions of EC-01 through EC-13.

\item \textbf{Anti-Responsibilities}: Does not own the underlying governed subject, authoritative governance requirements, or source information used as evidence. Does not perform, manage, remediate, assign, delegate or operationally escalate the subject activity. Does not replace System Operations or EC-12 Operational Assurance, and does not provide Clinical Services Delivery Assurance.

\item \textbf{Allocation Boundary}: No strategic logical component, application component, runtime service or implementation construct is allocated EC-14. The catalogue defines a capability, not a software component or deployment boundary.

\end{itemize}


\par\medskip\textbf{Conceptual Behavioural Sequence}\par


\begin{quote}\small\ttfamily
Define Good Behaviour\\
    ↓\\
Gather Evidence\\
    ↓\\
Assess Evidence against Good Behaviour\\
    ↓\\
Adjudicate Assessed Detail\\
    ↓\\
Report\\

\end{quote}


This sequence explains the capability; it is not an implementation workflow specification or a derived Business Process. “Define Good Behaviour” identifies assurance expectations and criteria; it does not transfer ownership of governing requirements to Service Guardian. Source information ownership remains with its responsible capability; assurance owns its contextual evidentiary association.



The \href{\detokenize{../../../../../docs/markdown/02-strategy/capabilities/business-enabling-capabilities.md#health-service-assurance-approved-business-enabling-capabilities}}{Business Enabling assurance principles} govern disposition, generic processing assurance, evidence/criteria temporality, assessment/adjudication, independence, failure/recovery and explicit evidence insufficiency. These semantics are not universally reducible to PASS / FAIL / UNKNOWN and do not imply an assurance information model or persistence mechanism.



Service Guardian supplies the reusable assurance semantics that bind a governed subject to applicable criteria and trustworthy evidence, support assessment and adjudication, and establish reportable findings. Its contribution distinguishes explicit evidence insufficiency from a conclusion about the subject. Existing Enterprise Capabilities contribute collaboratively; EC-14 does not duplicate all their information, policy or execution responsibilities.



Independence concerns control of evaluation and conclusion. A subject may supply evidence but cannot solely determine, suppress or retrospectively change its own assurance outcome.



\textbf{Boundary}



The capability does not manage or remediate the subject, replace operational assurance, or assure clinical adequacy. No application component, runtime service or deployment boundary is allocated, and the explanatory sequence is not a derived Business Process.


{\footnotesize \href{../../../../continuum/strategy/capability/d767f12f9b1c495883216ace464ee1c2.md}{Canonical record: d767f12f9b1c495883216ace464ee1c2}; full qualifications and provenance.}\par
\section*{II.1.2 Business Capabilities}
\phantomsection
\label{navigation:II.1.2}
\addcontentsline{toc}{section}{II.1.2 Business Capabilities}

\subsection*{II.1.2.1 [CT1] Care Access \& Coordination}
\phantomsection
\label{navigation:II.1.2.1}
\addcontentsline{toc}{subsection}{II.1.2.1 [CT1] Care Access \& Coordination}

{\small Historical predecessor: R.2.1 \hyperref[element:0dec07b6ddbe48bda93a2295dd3e2d7d]{[CT1] Care Access \& Coordination}; Canonical ID: \texttt{\seqsplit{0dec07b6ddbe48bda93a2295dd3e2d7d}}. Superseded historical identity — direct correspondence · superseded.}\par
\label{element:ed342cb8656b4f80b55ad519ecbd2d3a}
\index{[CT1] Care Access \& Coordination (II.1.2.1)}
{\small Capability · Canonical · current}\par

Managing patient intake, referral triage, appointment scheduling, multidisciplinary care team coordination, cross-facility patient transfers, and care continuity across organisational boundaries.



This capability makes care reachable and operationally coherent across participating services. It encompasses entering care, obtaining scheduled opportunities and usable capacity, coordinating participation and responsibility transitions, and arranging bounded support logistics. Its decomposition separates access decisions, resource use, continuity and supporting work so that movement or message delivery cannot stand in for acceptance of care responsibility.



\textbf{Intent}



Help a person's care requirements progress across organisational and operational boundaries with the relevant access and continuity context intact.



\textbf{Boundary}



Coordination does not determine clinical work or confer clinical authority. Corporate resource governance and general hospital logistics are outside the bounded care-support responsibility.



\textbf{Harmonia Architectural Relevance}



Referrals, encounters, appointments, operational state and communication provide material information for continuity. Harmonia's integration and coordination support does not make it the owner of clinical handover or the originating workforce and resource facts.


{\footnotesize \href{../../../../continuum/strategy/capability/ed342cb8656b4f80b55ad519ecbd2d3a.md}{Canonical record: ed342cb8656b4f80b55ad519ecbd2d3a}; full qualifications and provenance.}\par
\subsubsection*{II.1.2.1.1 [CT2] Care Access Management}
\phantomsection
\label{navigation:II.1.2.1.1}
\addcontentsline{toc}{subsubsection}{II.1.2.1.1 [CT2] Care Access Management}

\label{element:acbcad8330e64b59be5e3ec4a100a544}
\index{[CT2] Care Access Management (II.1.2.1.1)}
{\small Capability · Canonical · current}\par

Establish and manage access to applicable care.



Care Access Management establishes the circumstances in which an individual seeks or enters care. Intake creates the access context, referral requests another participant's involvement, and eligibility determines whether applicable access conditions are met. These are separate contributions to access: a person can be eligible without having a booking, and a received referral can remain undecided.



\textbf{Intent}



Make entry requirements and referral disposition intelligible without implying that care has already been accepted or delivered.



\textbf{Boundary}



Eligibility does not allocate care, and referral receipt does not establish an assumption of clinical responsibility.


{\small Recorded unresolved matters: source-qualified matters; exact scope remains in the canonical record.}\par
{\footnotesize \href{../../../../continuum/strategy/capability/acbcad8330e64b59be5e3ec4a100a544.md}{Canonical record: acbcad8330e64b59be5e3ec4a100a544}; full qualifications and provenance.}\par
\paragraph*{II.1.2.1.1.1 [CT3] Care Eligibility Determination}
\phantomsection
\label{navigation:II.1.2.1.1.1}
\addcontentsline{toc}{paragraph}{II.1.2.1.1.1 [CT3] Care Eligibility Determination}

\label{element:0922a1141f82416ca9de12054c28a69a}
\index{[CT3] Care Eligibility Determination (II.1.2.1.1.1)}
{\small Capability · Canonical · current}\par

Determine whether applicable eligibility conditions for care access are satisfied.



Care Eligibility Determination establishes whether the conditions governing access to applicable care are met. It contributes a qualified access decision alongside intake and referral, so that an identifiable request can be distinguished from one that satisfies the relevant criteria. Its purpose is eligibility in the applicable context rather than clinical diagnosis or service allocation.



\textbf{Boundary}



Eligibility alone does not reserve capacity, allocate a care place or deliver care.


{\small Recorded unresolved matters: source-qualified matters; exact scope remains in the canonical record.}\par
{\footnotesize \href{../../../../continuum/strategy/capability/0922a1141f82416ca9de12054c28a69a.md}{Canonical record: 0922a1141f82416ca9de12054c28a69a}; full qualifications and provenance.}\par
\paragraph*{II.1.2.1.1.2 [CT3] Care Intake Management}
\phantomsection
\label{navigation:II.1.2.1.1.2}
\addcontentsline{toc}{paragraph}{II.1.2.1.1.2 [CT3] Care Intake Management}

\label{element:1582467fff3b4dafa348aa100e9ca2de}
\index{[CT3] Care Intake Management (II.1.2.1.1.2)}
{\small Capability · Canonical · current}\par

Establish and manage the intake context through which an individual seeks or enters care.



Care Intake Management establishes the access context through which a person seeks or enters care. It makes the entry circumstance manageable without assuming that every intake follows a referral or already satisfies eligibility conditions. Within Care Access Management, intake is the responsibility for entry context rather than a decision that all subsequent care arrangements are complete.



\textbf{Boundary}



An intake record does not establish acceptance of a referral, care responsibility or delivery of care.


{\small Recorded unresolved matters: source-qualified matters; exact scope remains in the canonical record.}\par
{\footnotesize \href{../../../../continuum/strategy/capability/1582467fff3b4dafa348aa100e9ca2de.md}{Canonical record: 1582467fff3b4dafa348aa100e9ca2de}; full qualifications and provenance.}\par
\paragraph*{II.1.2.1.1.3 [CT3] Care Referral Management}
\phantomsection
\label{navigation:II.1.2.1.1.3}
\addcontentsline{toc}{paragraph}{II.1.2.1.1.3 [CT3] Care Referral Management}

\label{element:ebdaa781bba542f8b744ef34a836a269}
\index{[CT3] Care Referral Management (II.1.2.1.1.3)}
{\small Capability · Canonical · current}\par

Manage a request for another participant/service to assess, manage or assume applicable care responsibility.



Care Referral Management concerns a request for another participant or service to assess, manage or assume applicable care responsibility. It preserves the distinction between the requester's intent and the recipient's disposition, allowing care coordination to recognise an unresolved referral rather than assume a successful handoff. Referral is therefore broader than sending a document.



\textbf{Boundary}



Submission or receipt does not establish acceptance, an appointment, delivered care or responsibility actually assumed.


{\small Recorded unresolved matters: source-qualified matters; exact scope remains in the canonical record.}\par
{\footnotesize \href{../../../../continuum/strategy/capability/ebdaa781bba542f8b744ef34a836a269.md}{Canonical record: ebdaa781bba542f8b744ef34a836a269}; full qualifications and provenance.}\par
\subsubsection*{II.1.2.1.2 [CT2] Care Coordination \& Transition}
\phantomsection
\label{navigation:II.1.2.1.2}
\addcontentsline{toc}{subsubsection}{II.1.2.1.2 [CT2] Care Coordination \& Transition}

\label{element:174a9b93d19c44828f20360e9d543925}
\index{[CT2] Care Coordination \& Transition (II.1.2.1.2)}
{\small Capability · Canonical · current}\par

Coordinate care participation, actual responsibility transition and discharge.



Care Coordination \& Transition supports continuity by keeping participation, actual responsibility assumption and discharge progression distinct. Care teams may collaborate while responsibility is shared or changes, but physical transit and information exchange do not themselves establish who has accepted care. The capability therefore addresses continuity of accountable participation as well as the practical circumstances of departure.



\textbf{Intent}



Avoid gaps in care continuity caused by treating a movement or communication as an accepted responsibility transition.



\textbf{Boundary}



Harmonia may assemble handover information without managing clinical handover. Arrival, discharge readiness and physical departure are not equivalent to responsibility acceptance.


{\small Recorded unresolved matters: source-qualified matters; exact scope remains in the canonical record.}\par
{\footnotesize \href{../../../../continuum/strategy/capability/174a9b93d19c44828f20360e9d543925.md}{Canonical record: 174a9b93d19c44828f20360e9d543925}; full qualifications and provenance.}\par
\paragraph*{II.1.2.1.2.1 [CT3] Care Discharge Coordination}
\phantomsection
\label{navigation:II.1.2.1.2.1}
\addcontentsline{toc}{paragraph}{II.1.2.1.2.1 [CT3] Care Discharge Coordination}

\label{element:62e436db8ce44a98ba855de68ac6b876}
\index{[CT3] Care Discharge Coordination (II.1.2.1.2.1)}
{\small Capability · Canonical · current}\par

Coordinate applicable discharge requirements and progression.



Care Discharge Coordination brings applicable discharge requirements and progression into a coherent operational context. It supports the transition out of a care circumstance while keeping readiness, physical exit and continuing responsibility distinguishable. The capability concerns coordinating those requirements rather than deciding that every clinical condition for discharge has been met.



\textbf{Boundary}



Discharge readiness and departure are not themselves evidence that another participant has accepted care responsibility.


{\small Recorded unresolved matters: source-qualified matters; exact scope remains in the canonical record.}\par
{\footnotesize \href{../../../../continuum/strategy/capability/62e436db8ce44a98ba855de68ac6b876.md}{Canonical record: 62e436db8ce44a98ba855de68ac6b876}; full qualifications and provenance.}\par
\paragraph*{II.1.2.1.2.2 [CT3] Care Responsibility Transition}
\phantomsection
\label{navigation:II.1.2.1.2.2}
\addcontentsline{toc}{paragraph}{II.1.2.1.2.2 [CT3] Care Responsibility Transition}

\label{element:f8c9f7e2dcf14ceb9f228daaf9fed45e}
\index{[CT3] Care Responsibility Transition (II.1.2.1.2.2)}
{\small Capability · Canonical · current}\par

Establish and coordinate who has actually assumed applicable care responsibility.



Care Responsibility Transition establishes who has actually assumed applicable care responsibility and coordinates that transition. Its purpose is to prevent physical movement or exchanged information from being mistaken for accountable acceptance. It gives care continuity a responsibility meaning that a transport checkpoint or delivered referral cannot establish alone.



\textbf{Boundary}



Patient arrival, message receipt and departure do not by themselves demonstrate accepted responsibility.


{\small Recorded unresolved matters: source-qualified matters; exact scope remains in the canonical record.}\par
{\footnotesize \href{../../../../continuum/strategy/capability/f8c9f7e2dcf14ceb9f228daaf9fed45e.md}{Canonical record: f8c9f7e2dcf14ceb9f228daaf9fed45e}; full qualifications and provenance.}\par
\paragraph*{II.1.2.1.2.3 [CT3] Care-Team Coordination}
\phantomsection
\label{navigation:II.1.2.1.2.3}
\addcontentsline{toc}{paragraph}{II.1.2.1.2.3 [CT3] Care-Team Coordination}

\label{element:bbd138aab70d4b06b43662fb3f5373db}
\index{[CT3] Care-Team Coordination (II.1.2.1.2.3)}
{\small Capability · Canonical · current}\par

Coordinate participating care-team context and collaboration around care.



Care-Team Coordination maintains the participating team context and collaboration around care. It supports continuity by making participation intelligible across people and organisational boundaries, without assuming that membership itself assigns every clinical responsibility. Within Care Coordination \& Transition it concerns collaboration around care rather than proof of a responsibility handoff.



\textbf{Boundary}



Harmonia's support for team information and collaboration does not confer management of clinical handover or clinical worklists.


{\small Recorded unresolved matters: source-qualified matters; exact scope remains in the canonical record.}\par
{\footnotesize \href{../../../../continuum/strategy/capability/bbd138aab70d4b06b43662fb3f5373db.md}{Canonical record: bbd138aab70d4b06b43662fb3f5373db}; full qualifications and provenance.}\par
\subsubsection*{II.1.2.1.3 [CT2] Care Scheduling \& Capacity Management}
\phantomsection
\label{navigation:II.1.2.1.3}
\addcontentsline{toc}{subsubsection}{II.1.2.1.3 [CT2] Care Scheduling \& Capacity Management}

\label{element:61dd4b247df14546bf012587164bfe91}
\index{[CT2] Care Scheduling \& Capacity Management (II.1.2.1.3)}
{\small Capability · Canonical · current}\par

Manage scheduled opportunities, care places, usable operational capacity and care-demand resourcing.



This capability connects care demand with scheduled opportunities, suitable care places and legitimately usable operational resources. Its children distinguish a booking, a place allocation, the readiness of that place, resource deployment, flow coordination and service-level usable capacity. An empty bed or an available practitioner is relevant evidence, but neither alone establishes that the service can accommodate care.



\textbf{Intent}



Support care requirements with an operationally credible combination of time, place, resources and capacity.



\textbf{Boundary}



Deployment consumes roster, qualification and privilege context without changing their authoritative records. Operational capacity differs from strategic capacity planning and from the unresolved collective service-readiness concern.


{\small Recorded unresolved matters: source-qualified matters; exact scope remains in the canonical record.}\par
{\footnotesize \href{../../../../continuum/strategy/capability/61dd4b247df14546bf012587164bfe91.md}{Canonical record: 61dd4b247df14546bf012587164bfe91}; full qualifications and provenance.}\par
\paragraph*{II.1.2.1.3.1 [CT3] Care Appointment Management}
\phantomsection
\label{navigation:II.1.2.1.3.1}
\addcontentsline{toc}{paragraph}{II.1.2.1.3.1 [CT3] Care Appointment Management}

\label{element:14dc21b28938425498ff7b04e49efd30}
\index{[CT3] Care Appointment Management (II.1.2.1.3.1)}
{\small Capability · Canonical · current}\par

Establish and manage a scheduled opportunity for care.



Care Appointment Management establishes a scheduled opportunity for care and maintains that opportunity as arrangements change. Within scheduling and capacity it supplies the temporal commitment around which participants can organise, while leaving the actual availability of places and resources distinguishable. A scheduled opportunity is useful coordination information even before attendance or care occurs.



\textbf{Boundary}



A booking does not establish attendance, intervention performance or satisfaction of diagnostic-service-specific scheduling constraints.


{\small Recorded unresolved matters: source-qualified matters; exact scope remains in the canonical record.}\par
{\footnotesize \href{../../../../continuum/strategy/capability/14dc21b28938425498ff7b04e49efd30.md}{Canonical record: 14dc21b28938425498ff7b04e49efd30}; full qualifications and provenance.}\par
\paragraph*{II.1.2.1.3.2 [CT3] Care Resource Deployment}
\phantomsection
\label{navigation:II.1.2.1.3.2}
\addcontentsline{toc}{paragraph}{II.1.2.1.3.2 [CT3] Care Resource Deployment}

\label{element:7854224c3b064116b0d48f01ce5537d8}
\index{[CT3] Care Resource Deployment (II.1.2.1.3.2)}
{\small Capability · Canonical · current}\par

Apply legitimately available and appropriately authorised resources to a current or planned operational care requirement, particularly using roster/on-call and service-coverage context.



Care Resource Deployment applies legitimately available resources to a current or planned care requirement. It considers coverage, role, qualification, engagement and privilege context to distinguish availability in principle from deployment to a particular need. Its contribution is the operational application of resources, without changing the authoritative facts that make their use legitimate.



\textbf{Boundary}



Deployment is neither roster authoring nor clinical work assignment and does not confer clinical authority or ownership of enterprise assets.


{\small Recorded unresolved matters: source-qualified matters; exact scope remains in the canonical record.}\par
{\footnotesize \href{../../../../continuum/strategy/capability/7854224c3b064116b0d48f01ce5537d8.md}{Canonical record: 7854224c3b064116b0d48f01ce5537d8}; full qualifications and provenance.}\par
\paragraph*{II.1.2.1.3.3 [CT3] Care-Place Allocation}
\phantomsection
\label{navigation:II.1.2.1.3.3}
\addcontentsline{toc}{paragraph}{II.1.2.1.3.3 [CT3] Care-Place Allocation}

\label{element:f8a4df47bea742a0869802472136d55f}
\index{[CT3] Care-Place Allocation (II.1.2.1.3.3)}
{\small Capability · Canonical · current}\par

Allocate an applicable care place to a care requirement.



Care-Place Allocation associates a care requirement with an applicable care place. Its outcome is the allocation itself, distinct from defining the place or establishing its readiness. This separation supports reasoning about reserved, allocated and usable places without assuming that an apparently empty location is fit for the requirement.



\textbf{Boundary}



Allocation does not establish clinical suitability beyond the applicable context or make an unready place ready.


{\small Recorded unresolved matters: source-qualified matters; exact scope remains in the canonical record.}\par
{\footnotesize \href{../../../../continuum/strategy/capability/f8a4df47bea742a0869802472136d55f.md}{Canonical record: f8a4df47bea742a0869802472136d55f}; full qualifications and provenance.}\par
\paragraph*{II.1.2.1.3.4 [CT3] Care-Place Readiness Management}
\phantomsection
\label{navigation:II.1.2.1.3.4}
\addcontentsline{toc}{paragraph}{II.1.2.1.3.4 [CT3] Care-Place Readiness Management}

\label{element:4dd0f1b9f9334c5c80066a8974fdfb7a}
\index{[CT3] Care-Place Readiness Management (II.1.2.1.3.4)}
{\small Capability · Canonical · current}\par

Establish and manage whether a particular care place is ready for allocation or use.



Care-Place Readiness Management establishes whether a particular place can be allocated or used in its operational context. It connects the place's condition with its intended use so that vacancy can be distinguished from readiness. This is a place-specific responsibility within scheduling and capacity, supporting credible allocation rather than merely counting physical spaces.



\textbf{Boundary}



Readiness of one place does not establish collective readiness or usable capacity of the whole service.


{\small Recorded unresolved matters: source-qualified matters; exact scope remains in the canonical record.}\par
{\footnotesize \href{../../../../continuum/strategy/capability/4dd0f1b9f9334c5c80066a8974fdfb7a.md}{Canonical record: 4dd0f1b9f9334c5c80066a8974fdfb7a}; full qualifications and provenance.}\par
\paragraph*{II.1.2.1.3.5 [CT3] Operational Care Capacity Management}
\phantomsection
\label{navigation:II.1.2.1.3.5}
\addcontentsline{toc}{paragraph}{II.1.2.1.3.5 [CT3] Operational Care Capacity Management}

\label{element:3040519e1fec43699c14a3cafb2d945d}
\index{[CT3] Operational Care Capacity Management (II.1.2.1.3.5)}
{\small Capability · Canonical · current}\par

Understand and manage the service’s usable capacity to accommodate care in its current operational context.



Operational Care Capacity Management establishes and manages the service's usable ability to accommodate care under current conditions. It considers capacity as a service-level concern rather than a simple count of staff, devices or vacant beds. This supports decisions grounded in what can actually be accommodated, while keeping current capacity distinct from plans for future provision.



\textbf{Boundary}



Resource availability, strategic capacity requirements and collective service readiness are distinct; the latter's placement remains unresolved.


{\small Recorded unresolved matters: source-qualified matters; exact scope remains in the canonical record.}\par
{\footnotesize \href{../../../../continuum/strategy/capability/3040519e1fec43699c14a3cafb2d945d.md}{Canonical record: 3040519e1fec43699c14a3cafb2d945d}; full qualifications and provenance.}\par
\paragraph*{II.1.2.1.3.6 [CT3] Operational Care Flow Coordination}
\phantomsection
\label{navigation:II.1.2.1.3.6}
\addcontentsline{toc}{paragraph}{II.1.2.1.3.6 [CT3] Operational Care Flow Coordination}

\label{element:d1dedb10ea694af5b88dbe1b6192f3ac}
\index{[CT3] Operational Care Flow Coordination (II.1.2.1.3.6)}
{\small Capability · Canonical · current}\par

Coordinate operational care flow and its progression within the care-delivery environment.



Operational Care Flow Coordination addresses progression of care flow within the delivery environment. It provides the coordinating responsibility around operational movement and impediments, rather than treating each scheduled opportunity or place allocation as an isolated fact. Its outcome concerns coherent flow under current conditions, not the clinical determination of what care should occur.



\textbf{Boundary}



Coordinating flow does not transfer clinical work, workforce governance or general hospital-logistics authority.


{\small Recorded unresolved matters: source-qualified matters; exact scope remains in the canonical record.}\par
{\footnotesize \href{../../../../continuum/strategy/capability/d1dedb10ea694af5b88dbe1b6192f3ac.md}{Canonical record: d1dedb10ea694af5b88dbe1b6192f3ac}; full qualifications and provenance.}\par
\subsubsection*{II.1.2.1.4 [CT2] Care Support Logistics Management}
\phantomsection
\label{navigation:II.1.2.1.4}
\addcontentsline{toc}{subsubsection}{II.1.2.1.4 [CT2] Care Support Logistics Management}

\label{element:92be036674254f719cf4bb613120128a}
\index{[CT2] Care Support Logistics Management (II.1.2.1.4)}
{\small Capability · Canonical · current}\par

Coordinate bounded operational support needed for care access and fulfilment.



Care Support Logistics Management coordinates the bounded physical and operational support needed to fulfil care. Patient transit, clinical-material movement and authorised support work have different subjects and completion evidence, even when the same operational environment supports them. This capability provides their common care-support frame without absorbing every logistical activity of a healthcare enterprise.



\textbf{Intent}



Ensure necessary supporting activity can progress with the relevant patient, material and operational constraints visible.



\textbf{Boundary}



Assignment does not establish a worker's qualification or clinical authority. Enterprise supply chains, facilities and workforce governance retain their own responsibilities.


{\small Recorded unresolved matters: source-qualified matters; exact scope remains in the canonical record.}\par
{\footnotesize \href{../../../../continuum/strategy/capability/92be036674254f719cf4bb613120128a.md}{Canonical record: 92be036674254f719cf4bb613120128a}; full qualifications and provenance.}\par
\paragraph*{II.1.2.1.4.1 [CT3] Care Support Work Coordination}
\phantomsection
\label{navigation:II.1.2.1.4.1}
\addcontentsline{toc}{paragraph}{II.1.2.1.4.1 [CT3] Care Support Work Coordination}

\label{element:f9f5c69b0f0b472287cd4ac56ccff6f3}
\index{[CT3] Care Support Work Coordination (II.1.2.1.4.1)}
{\small Capability · Canonical · current}\par

Coordinate authorised operational support work: what needs doing, who performs it and how progression, completion or escalation is coordinated.



Care Support Work Coordination organises authorised operational support work around care: what needs doing, who will perform it and how progress, completion or escalation is coordinated. It gives supporting activity an explicit progression responsibility, rather than assuming that issuing a request means the work occurred. The work remains bounded by applicable authority and resource constraints.



\textbf{Boundary}



Dispatch does not establish worker availability, qualification or privilege, and support coordination does not determine clinical work.


{\small Recorded unresolved matters: source-qualified matters; exact scope remains in the canonical record.}\par
{\footnotesize \href{../../../../continuum/strategy/capability/f9f5c69b0f0b472287cd4ac56ccff6f3.md}{Canonical record: f9f5c69b0f0b472287cd4ac56ccff6f3}; full qualifications and provenance.}\par
\paragraph*{II.1.2.1.4.2 [CT3] Clinical Material Logistics Coordination}
\phantomsection
\label{navigation:II.1.2.1.4.2}
\addcontentsline{toc}{paragraph}{II.1.2.1.4.2 [CT3] Clinical Material Logistics Coordination}

\label{element:4b66444de58a48a9a9433568bf3a60d3}
\index{[CT3] Clinical Material Logistics Coordination (II.1.2.1.4.2)}
{\small Capability · Canonical · current}\par

Coordinate clinically relevant material movement, custody, availability or disposition needed for care or diagnostic fulfilment.



This capability coordinates clinically relevant material movement and the custody, availability or disposition needed for care or diagnostic fulfilment. It makes the material's operational journey intelligible while retaining why the material matters to the clinical activity. The responsibility differs from moving a patient and from managing the diagnostic meaning of information derived from a specimen.



\textbf{Boundary}



Bounded clinical-material logistics does not establish a general procurement, supply-chain or enterprise inventory architecture.


{\small Recorded unresolved matters: source-qualified matters; exact scope remains in the canonical record.}\par
{\footnotesize \href{../../../../continuum/strategy/capability/4b66444de58a48a9a9433568bf3a60d3.md}{Canonical record: 4b66444de58a48a9a9433568bf3a60d3}; full qualifications and provenance.}\par
\paragraph*{II.1.2.1.4.3 [CT3] Patient Movement Coordination}
\phantomsection
\label{navigation:II.1.2.1.4.3}
\addcontentsline{toc}{paragraph}{II.1.2.1.4.3 [CT3] Patient Movement Coordination}

\label{element:f404b6647fe04a3390e57022bdef0843}
\index{[CT3] Patient Movement Coordination (II.1.2.1.4.3)}
{\small Capability · Canonical · current}\par

Coordinate physical patient transit and applicable mobility constraints.



Patient Movement Coordination concerns the person's physical transit and the mobility constraints relevant to it. Within care-support logistics it supports movement appropriate to the patient's circumstances without treating transport completion as a clinical transition. Its responsibility is useful wherever care requires a change of physical location.



\textbf{Boundary}



Transport progression does not establish clinical handover, destination acceptance of care or authority over the clinical plan.


{\small Recorded unresolved matters: source-qualified matters; exact scope remains in the canonical record.}\par
{\footnotesize \href{../../../../continuum/strategy/capability/f404b6647fe04a3390e57022bdef0843.md}{Canonical record: f404b6647fe04a3390e57022bdef0843}; full qualifications and provenance.}\par
\subsubsection*{II.1.2.1.5 [CT2] Service Delivery Preparedness Management}
\phantomsection
\label{navigation:II.1.2.1.5}
\addcontentsline{toc}{subsubsection}{II.1.2.1.5 [CT2] Service Delivery Preparedness Management}

\label{element:9665d894c4c9438999d8cc10afae695d}
\index{[CT2] Service Delivery Preparedness Management (II.1.2.1.5)}
{\small Capability · Proposed · current}\par

Establish whether service offerings can be delivered in their intended operating contexts through the combined standing of their dependencies.



\textbf{Overview}



Preparedness asks whether the service as a whole can operate, distinct from scheduling capacity or readiness of a single place. This is a proposed missing business frame for the already identified readiness concern.



\textbf{Intent}



Recognise whether the whole service can actually operate in its intended context.



\textbf{Boundary}



Does not own the dependencies, their management or clinical execution.


{\footnotesize \href{../../../../continuum/strategy/capability/9665d894c4c9438999d8cc10afae695d.md}{Canonical record: 9665d894c4c9438999d8cc10afae695d}; full qualifications and provenance.}\par
\paragraph*{II.1.2.1.5.1 [CT3] Health Service Operational Readiness}
\phantomsection
\label{navigation:II.1.2.1.5.1}
\addcontentsline{toc}{paragraph}{II.1.2.1.5.1 [CT3] Health Service Operational Readiness}

\label{element:cc91eb4017034939ab82b152bddf7b1c}
\index{[CT3] Health Service Operational Readiness (II.1.2.1.5.1)}
{\small Capability · Candidate · current}\par

Definition unresolved in the baseline. No definition has been synthesised from the name.



\textbf{Proposed overview}



Collective readiness is an access prerequisite and has a coherent proposed preparedness frame; it is not reduced to capacity, resource state or service alignment.



\textbf{Proposed definition}



Establish collective operational readiness of a service offering in a site or session context from its applicable dependencies.



\textbf{Proposed boundary}



The element remains Candidate. The definition/placement elaboration is Proposed and does not adjudicate the unresolved RAD question or transfer source-domain authority.



\medskip\hrule\medskip

{\small Recorded unresolved matters: source-qualified matters; exact scope remains in the canonical record.}\par
{\footnotesize \href{../../../../continuum/strategy/capability/cc91eb4017034939ab82b152bddf7b1c.md}{Canonical record: cc91eb4017034939ab82b152bddf7b1c}; full qualifications and provenance.}\par
\subsection*{II.1.2.2 [CT1] Clinical Quality, Safety \& Improvement}
\phantomsection
\label{navigation:II.1.2.2}
\addcontentsline{toc}{subsection}{II.1.2.2 [CT1] Clinical Quality, Safety \& Improvement}

{\small Historical predecessor: R.2.11 \hyperref[element:b967c383dbcd4c24abe6a83ea5e79e00]{[CT1] Clinical Quality, Safety \& Improvement}; Canonical ID: \texttt{\seqsplit{b967c383dbcd4c24abe6a83ea5e79e00}}. Superseded historical identity — direct correspondence · superseded.}\par
\label{element:63837bf848774572a06b935af95581a2}
\index{[CT1] Clinical Quality, Safety \& Improvement (II.1.2.2)}
{\small Capability · Canonical · current}\par

Operational clinical quality, safety and improvement: monitoring clinical outcomes; identifying quality and safety concerns; adverse event and incident management, reporting and response; infection control; operational quality and safety monitoring; continuous quality improvement; and coordination of operational responses to identified concerns.



This capability addresses operational clinical quality and safety, including concerns, incidents, response and improvement. Its established safety branch distinguishes risk identification, prospective mitigation and active response. Its quality branch distinguishes criteria, evidence, assessment and reporting. Together these make it possible to reason about clinical quality without equating an observation, a judgement and the action taken in response.



\textbf{Intent}



Support recognition and management of clinical safety concerns and informed improvement of clinical quality.



\textbf{Boundary}



This is neither a generic enterprise incident container nor a replacement for macro-level Health Service Assurance. Its bounded decomposition does not exhaust clinical quality and safety practice.



\textbf{Harmonia Architectural Relevance}



Information about outcomes, risks, incidents and operational responses materially supports these responsibilities. Integrating that evidence does not assign clinical quality judgement or originating clinical authority to the platform.


{\footnotesize \href{../../../../continuum/strategy/capability/63837bf848774572a06b935af95581a2.md}{Canonical record: 63837bf848774572a06b935af95581a2}; full qualifications and provenance.}\par
\subsubsection*{II.1.2.2.1 [CT2] Clinical Quality Management}
\phantomsection
\label{navigation:II.1.2.2.1}
\addcontentsline{toc}{subsubsection}{II.1.2.2.1 [CT2] Clinical Quality Management}

{\small Historical predecessor: R.2.6 \hyperref[element:dc1bece79eb7452e951d03a2ee6dfaa6]{[CT2] Clinical Quality Management}; Canonical ID: \texttt{\seqsplit{dc1bece79eb7452e951d03a2ee6dfaa6}}. Superseded historical identity — direct correspondence · superseded.}\par
\label{element:89761f2d38f846c39df1df0975964288}
\index{[CT2] Clinical Quality Management (II.1.2.2.1)}
{\small Capability · Canonical · current}\par

Manage the clinical-quality distinctions materially required for R1/R2.



Clinical Quality Management supplies the criteria and evidence basis for evaluating clinical quality and communicating established findings. Its four children preserve the difference between the expected standard, the material considered, the assessment reached and the report of that assessment. This supports scrutiny of why a quality conclusion is justified rather than treating the presence of measurements as a judgement.



\textbf{Intent}



Make clinical-quality conclusions understandable in relation to their criteria, evidence and assessment authority.



\textbf{Boundary}



Reporting does not originate the conclusion or perform subsequent improvement. Neither a separate monitoring provider nor generic Continuous Improvement placement is established by this decomposition.


{\small Recorded unresolved matters: source-qualified matters; exact scope remains in the canonical record.}\par
{\footnotesize \href{../../../../continuum/strategy/capability/89761f2d38f846c39df1df0975964288.md}{Canonical record: 89761f2d38f846c39df1df0975964288}; full qualifications and provenance.}\par
\paragraph*{II.1.2.2.1.1 [CT3] Clinical Quality Criteria Management}
\phantomsection
\label{navigation:II.1.2.2.1.1}
\addcontentsline{toc}{paragraph}{II.1.2.2.1.1 [CT3] Clinical Quality Criteria Management}

{\small Historical predecessor: R.2.2 \hyperref[element:7a51321b00464016920976b269981edd]{[CT3] Clinical Quality Criteria Management}; Canonical ID: \texttt{\seqsplit{7a51321b00464016920976b269981edd}}. Superseded historical identity — direct correspondence · superseded.}\par
\label{element:1f48527a9a7c4382a7f8a660ccb9f29f}
\index{[CT3] Clinical Quality Criteria Management (II.1.2.2.1.1)}
{\small Capability · Canonical · current}\par

Establish and manage measures, indicators, expectations and criteria against which clinical quality is considered.



Clinical Quality Criteria Management establishes the measures, indicators and expectations used to consider clinical quality. It makes the basis of a later judgement explicit so that evidence can be assessed against an identifiable expectation rather than an unstated standard. Criteria are distinct from the measurements or findings produced when they are applied.



\textbf{Boundary}



Managing criteria does not originate the underlying evidence or establish a quality conclusion.


{\small Recorded unresolved matters: source-qualified matters; exact scope remains in the canonical record.}\par
{\footnotesize \href{../../../../continuum/strategy/capability/1f48527a9a7c4382a7f8a660ccb9f29f.md}{Canonical record: 1f48527a9a7c4382a7f8a660ccb9f29f}; full qualifications and provenance.}\par
\paragraph*{II.1.2.2.1.2 [CT3] Clinical Quality Evidence Management}
\phantomsection
\label{navigation:II.1.2.2.1.2}
\addcontentsline{toc}{paragraph}{II.1.2.2.1.2 [CT3] Clinical Quality Evidence Management}

{\small Historical predecessor: R.2.3 \hyperref[element:f31902d474894670aeef3e0abbedccf0]{[CT3] Clinical Quality Evidence Management}; Canonical ID: \texttt{\seqsplit{f31902d474894670aeef3e0abbedccf0}}. Superseded historical identity — direct correspondence · superseded.}\par
\label{element:1b0a1c6945ed4fadba045ff4753c104c}
\index{[CT3] Clinical Quality Evidence Management (II.1.2.2.1.2)}
{\small Capability · Canonical · current}\par

Manage the body of applicable evidence used for clinical-quality consideration without owning its originating activity.



Clinical Quality Evidence Management maintains the body of material applicable to quality consideration. Its purpose is to retain evidence in a form and context that can support review without treating availability as proof of quality. The originating activity and its authority remain distinguishable from the later quality purpose for which its information is used.



\textbf{Boundary}



Evidentiary use transfers neither source ownership nor authority to determine the activity being reviewed.


{\small Recorded unresolved matters: source-qualified matters; exact scope remains in the canonical record.}\par
{\footnotesize \href{../../../../continuum/strategy/capability/1b0a1c6945ed4fadba045ff4753c104c.md}{Canonical record: 1b0a1c6945ed4fadba045ff4753c104c}; full qualifications and provenance.}\par
\paragraph*{II.1.2.2.1.3 [CT3] Clinical Quality Reporting}
\phantomsection
\label{navigation:II.1.2.2.1.3}
\addcontentsline{toc}{paragraph}{II.1.2.2.1.3 [CT3] Clinical Quality Reporting}

{\small Historical predecessor: R.2.4 \hyperref[element:e5812fcdf3f5465b9b1d8ba536075c5f]{[CT3] Clinical Quality Reporting}; Canonical ID: \texttt{\seqsplit{e5812fcdf3f5465b9b1d8ba536075c5f}}. Superseded historical identity — direct correspondence · superseded.}\par
\label{element:70f34839c23042d59469e6ef96944057}
\index{[CT3] Clinical Quality Reporting (II.1.2.2.1.3)}
{\small Capability · Canonical · current}\par

Communicate established clinical-quality measures, assessments, findings and conclusions to authorised consumers.



Clinical Quality Reporting communicates established quality measures, assessments and findings to authorised consumers. It makes the outcome of quality consideration usable beyond the assessment context while retaining its basis and standing. Reporting therefore represents a communication responsibility, not another opportunity to originate or silently change the reported judgement.



\textbf{Boundary}



Communicating a quality concern does not perform the operational or clinical improvement response.


{\small Recorded unresolved matters: source-qualified matters; exact scope remains in the canonical record.}\par
{\footnotesize \href{../../../../continuum/strategy/capability/70f34839c23042d59469e6ef96944057.md}{Canonical record: 70f34839c23042d59469e6ef96944057}; full qualifications and provenance.}\par
\paragraph*{II.1.2.2.1.4 [CT3] Clinical Quality Review \& Assessment Management}
\phantomsection
\label{navigation:II.1.2.2.1.4}
\addcontentsline{toc}{paragraph}{II.1.2.2.1.4 [CT3] Clinical Quality Review \& Assessment Management}

{\small Historical predecessor: R.2.5 \hyperref[element:4b17ec92f0af4c4b9556912dbdd0f29f]{[CT3] Clinical Quality Review \& Assessment Management}; Canonical ID: \texttt{\seqsplit{4b17ec92f0af4c4b9556912dbdd0f29f}}. Superseded historical identity — direct correspondence · superseded.}\par
\label{element:8318c35b302b4dcaa0e795fbfdc93e3b}
\index{[CT3] Clinical Quality Review \& Assessment Management (II.1.2.2.1.4)}
{\small Capability · Canonical · current}\par

Review evidence against applicable criteria and establish a clinical-quality assessment, finding or conclusion.



This capability evaluates quality evidence against applicable criteria and establishes a finding, assessment or conclusion. It is the judgement-bearing responsibility within Clinical Quality Management, distinct from maintaining the material and from communicating the result. Making that distinction supports scrutiny of the criterion and evidentiary basis for a quality assessment.



\textbf{Boundary}



An assessment is not a report of itself or the subsequent improvement action it may inform.


{\small Recorded unresolved matters: source-qualified matters; exact scope remains in the canonical record.}\par
{\footnotesize \href{../../../../continuum/strategy/capability/8318c35b302b4dcaa0e795fbfdc93e3b.md}{Canonical record: 8318c35b302b4dcaa0e795fbfdc93e3b}; full qualifications and provenance.}\par
\subsubsection*{II.1.2.2.2 [CT2] Clinical Safety Management}
\phantomsection
\label{navigation:II.1.2.2.2}
\addcontentsline{toc}{subsubsection}{II.1.2.2.2 [CT2] Clinical Safety Management}

{\small Historical predecessor: R.2.10 \hyperref[element:650ccec030fa46cc8c9d6b508a841778]{[CT2] Clinical Safety Management}; Canonical ID: \texttt{\seqsplit{650ccec030fa46cc8c9d6b508a841778}}. Superseded historical identity — direct correspondence · superseded.}\par
\label{element:28d1bf3e514e4c29837ebedaf9e4917b}
\index{[CT2] Clinical Safety Management (II.1.2.2.2)}
{\small Capability · Canonical · current}\par

Manage the clinical-risk distinctions materially required for R1/R2.



Clinical Safety Management distinguishes recognising a safety risk from reducing future exposure and actively dealing with a developing or materialised concern. These responsibilities may use overlapping evidence, but their purposes differ: characterisation makes the risk intelligible, mitigation changes prospective risk, and response addresses the situation requiring action.



\textbf{Intent}



Support proportionate safety action grounded in the nature of the identified clinical risk.



\textbf{Boundary}



Monitoring evidence is not automatically risk identification. Other capabilities can perform required actions without transferring their responsibilities; the three children are not a mandatory sequence or exhaustive safety taxonomy.


{\small Recorded unresolved matters: source-qualified matters; exact scope remains in the canonical record.}\par
{\footnotesize \href{../../../../continuum/strategy/capability/28d1bf3e514e4c29837ebedaf9e4917b.md}{Canonical record: 28d1bf3e514e4c29837ebedaf9e4917b}; full qualifications and provenance.}\par
\paragraph*{II.1.2.2.2.1 [CT3] Clinical Risk Identification}
\phantomsection
\label{navigation:II.1.2.2.2.1}
\addcontentsline{toc}{paragraph}{II.1.2.2.2.1 [CT3] Clinical Risk Identification}

{\small Historical predecessor: R.2.7 \hyperref[element:690da55f1f5b4c44ba58cf67a717a564]{[CT3] Clinical Risk Identification}; Canonical ID: \texttt{\seqsplit{690da55f1f5b4c44ba58cf67a717a564}}. Superseded historical identity — direct correspondence · superseded.}\par
\label{element:d54c19c6efaf42af82ad71a4142149d4}
\index{[CT3] Clinical Risk Identification (II.1.2.2.2.1)}
{\small Capability · Canonical · current}\par

Recognise and characterise a clinical safety risk from applicable evidence.



Clinical Risk Identification recognises and characterises a safety risk from relevant evidence. It establishes the concern to which mitigation or active response may be directed, rather than assuming that every monitored deviation is already a clinical risk. Its value is the meaningful characterisation that makes proportionate safety consideration possible.



\textbf{Boundary}



Evidence observation alone does not identify risk or establish a mitigation or response decision.


{\small Recorded unresolved matters: source-qualified matters; exact scope remains in the canonical record.}\par
{\footnotesize \href{../../../../continuum/strategy/capability/d54c19c6efaf42af82ad71a4142149d4.md}{Canonical record: d54c19c6efaf42af82ad71a4142149d4}; full qualifications and provenance.}\par
\paragraph*{II.1.2.2.2.2 [CT3] Clinical Risk Mitigation Management}
\phantomsection
\label{navigation:II.1.2.2.2.2}
\addcontentsline{toc}{paragraph}{II.1.2.2.2.2 [CT3] Clinical Risk Mitigation Management}

{\small Historical predecessor: R.2.8 \hyperref[element:0bd3311d89af45bb935fde20e719369c]{[CT3] Clinical Risk Mitigation Management}; Canonical ID: \texttt{\seqsplit{0bd3311d89af45bb935fde20e719369c}}. Superseded historical identity — direct correspondence · superseded.}\par
\label{element:2c232ade21ec45619c3ee4e928bd97ee}
\index{[CT3] Clinical Risk Mitigation Management (II.1.2.2.2.2)}
{\small Capability · Canonical · current}\par

Manage prospective/preventative measures reducing identified risk likelihood, exposure or consequence.



Clinical Risk Mitigation Management addresses prospective measures intended to reduce an identified risk's likelihood, exposure or consequences. It gives prevention a distinct responsibility alongside risk characterisation and active response. Required measures can involve other capabilities without transferring the ownership of their actions or information to safety management.



\textbf{Boundary}



Mitigation differs from containing or remediating a concern that has materialised or requires active intervention.


{\small Recorded unresolved matters: source-qualified matters; exact scope remains in the canonical record.}\par
{\footnotesize \href{../../../../continuum/strategy/capability/2c232ade21ec45619c3ee4e928bd97ee.md}{Canonical record: 2c232ade21ec45619c3ee4e928bd97ee}; full qualifications and provenance.}\par
\paragraph*{II.1.2.2.2.3 [CT3] Clinical Risk Response Management}
\phantomsection
\label{navigation:II.1.2.2.2.3}
\addcontentsline{toc}{paragraph}{II.1.2.2.2.3 [CT3] Clinical Risk Response Management}

{\small Historical predecessor: R.2.9 \hyperref[element:20f719ab4b6741b393ef570b4d1b05df]{[CT3] Clinical Risk Response Management}; Canonical ID: \texttt{\seqsplit{20f719ab4b6741b393ef570b4d1b05df}}. Superseded historical identity — direct correspondence · superseded.}\par
\label{element:5e36cb137e4e4782b2ddbf78cdf0da7b}
\index{[CT3] Clinical Risk Response Management (II.1.2.2.2.3)}
{\small Capability · Canonical · current}\par

Manage active response, containment and remediation where risk has materialised, is emerging or requires active intervention.



Clinical Risk Response Management manages active intervention, containment or remediation where a safety risk is emerging, materialised or otherwise demands action. Its purpose is to address the developing situation rather than only reduce future exposure. The responsibility remains clinical-risk-specific while the actual actions may belong to other established abilities.



\textbf{Boundary}



It is not a general enterprise incident-management container and does not acquire the clinical authority of responding participants.


{\small Recorded unresolved matters: source-qualified matters; exact scope remains in the canonical record.}\par
{\footnotesize \href{../../../../continuum/strategy/capability/5e36cb137e4e4782b2ddbf78cdf0da7b.md}{Canonical record: 5e36cb137e4e4782b2ddbf78cdf0da7b}; full qualifications and provenance.}\par
\subsection*{II.1.2.3 [CT1] Community Health \& Wellbeing}
\phantomsection
\label{navigation:II.1.2.3}
\addcontentsline{toc}{subsection}{II.1.2.3 [CT1] Community Health \& Wellbeing}

{\small Historical predecessor: R.2.19 \hyperref[element:175174c5cce549d2b154a4f6704721dc]{[CT1] Community Health \& Wellbeing}; Canonical ID: \texttt{\seqsplit{175174c5cce549d2b154a4f6704721dc}}. Superseded historical identity — direct correspondence · superseded.}\par
\label{element:b326b4aaf02a41a986ed22d184cbc152}
\index{[CT1] Community Health \& Wellbeing (II.1.2.3)}
{\small Capability · Canonical · current}\par

Support community health and wellbeing, including the bounded preventative-health abilities required for R1/R2. Wider environmental health, outreach and social determinants remain enterprise context without inferred additional decomposition.



Community Health \& Wellbeing frames organised activity supporting health beyond a single person's episode of care. The governed decomposition focuses on preventative health: identifying needs, defining programmes and cohorts, educating, observing programme operation and reporting. Wider outreach, environmental and social-determinant concerns remain enterprise context rather than additional modelled children.



\textbf{Intent}



Support organised prevention with a clear relationship between the health problem, intended population and programme activity.



\textbf{Boundary}



Programme management does not acquire the clinical or appointment responsibilities needed to deliver a programme.



\textbf{Harmonia Architectural Relevance}



Programme and cohort distinctions materially shape the information, notifications and care enablement required. This supports the Relevant classification without assigning wider public-health operating authority to Harmonia.


{\footnotesize \href{../../../../continuum/strategy/capability/b326b4aaf02a41a986ed22d184cbc152.md}{Canonical record: b326b4aaf02a41a986ed22d184cbc152}; full qualifications and provenance.}\par
\subsubsection*{II.1.2.3.1 [CT2] Preventative Health}
\phantomsection
\label{navigation:II.1.2.3.1}
\addcontentsline{toc}{subsubsection}{II.1.2.3.1 [CT2] Preventative Health}

{\small Historical predecessor: R.2.18 \hyperref[element:80affa7bcd284ecbbcbc5092f8497ddf]{[CT2] Preventative Health}; Canonical ID: \texttt{\seqsplit{80affa7bcd284ecbbcbc5092f8497ddf}}. Superseded historical identity — direct correspondence · superseded.}\par
\label{element:81be04a2d58548a4aea94f9aae90feb6}
\index{[CT2] Preventative Health (II.1.2.3.1)}
{\small Capability · Canonical · current}\par

Address preventative-health needs through organised activity with sufficient R1/R2 distinctions.



Preventative Health organises a response to a health problem before adverse outcomes arise or worsen. The established children distinguish why attention is warranted, how an organised programme is managed, who it targets, the education provided, what programme operation shows and what is reported. This allows a screening or immunisation programme to be explained without creating a new capability for every programme type.



\textbf{Intent}



Connect a preventative purpose with an identifiable target population and observable organised activity.



\textbf{Boundary}



Programme performance differs from population-health assessment and individual clinical monitoring. Unmodelled specialist responsibilities are not silently assigned to the existing children.


{\small Recorded unresolved matters: source-qualified matters; exact scope remains in the canonical record.}\par
{\footnotesize \href{../../../../continuum/strategy/capability/81be04a2d58548a4aea94f9aae90feb6.md}{Canonical record: 81be04a2d58548a4aea94f9aae90feb6}; full qualifications and provenance.}\par
\paragraph*{II.1.2.3.1.1 [CT3] Preventative Health Cohort Management}
\phantomsection
\label{navigation:II.1.2.3.1.1}
\addcontentsline{toc}{paragraph}{II.1.2.3.1.1 [CT3] Preventative Health Cohort Management}

{\small Historical predecessor: R.2.12 \hyperref[element:fef479f0fd5048b19b2aec770f0233fe]{[CT3] Preventative Health Cohort Management}; Canonical ID: \texttt{\seqsplit{fef479f0fd5048b19b2aec770f0233fe}}. Superseded historical identity — direct correspondence · superseded.}\par
\label{element:97625bb28e5745d5994e1bffd4037188}
\index{[CT3] Preventative Health Cohort Management (II.1.2.3.1.1)}
{\small Capability · Canonical · current}\par

Establish and manage the target population, inclusion/exclusion criteria, membership, stratification and applicable entry/exit context.



Preventative Health Cohort Management establishes who belongs to the population targeted for a preventative purpose. Inclusion, exclusion, stratification and entry or exit context make membership explainable rather than a list detached from its rationale. It provides a governed population context through which programme activity and reporting can be understood.



\textbf{Boundary}



A population at risk is not automatically a defined cohort, and cohort membership is not an individual clinical determination.


{\small Recorded unresolved matters: source-qualified matters; exact scope remains in the canonical record.}\par
{\footnotesize \href{../../../../continuum/strategy/capability/97625bb28e5745d5994e1bffd4037188.md}{Canonical record: 97625bb28e5745d5994e1bffd4037188}; full qualifications and provenance.}\par
\paragraph*{II.1.2.3.1.2 [CT3] Preventative Health Education Services}
\phantomsection
\label{navigation:II.1.2.3.1.2}
\addcontentsline{toc}{paragraph}{II.1.2.3.1.2 [CT3] Preventative Health Education Services}

{\small Historical predecessor: R.2.13 \hyperref[element:dbe10959237f405caeadef2a5748303a]{[CT3] Preventative Health Education Services}; Canonical ID: \texttt{\seqsplit{dbe10959237f405caeadef2a5748303a}}. Superseded historical identity — direct correspondence · superseded.}\par
\label{element:7541d9b3b65d47d286dabdc093cb0eed}
\index{[CT3] Preventative Health Education Services (II.1.2.3.1.2)}
{\small Capability · Canonical · current}\par

Provide education intended to prevent adverse health outcomes and influence preventative knowledge or behaviour.



Preventative Health Education Services provides education intended to influence knowledge or behaviour so that adverse health outcomes can be prevented. Its purpose distinguishes it from general organisational communication: the educational contribution is directed to a preventative-health need. It can support organised programme activity without being equivalent to programme management or clinical intervention.



\textbf{Boundary}



Communicating educational material does not establish individual clinical assessment, treatment or programme effectiveness.


{\small Recorded unresolved matters: source-qualified matters; exact scope remains in the canonical record.}\par
{\footnotesize \href{../../../../continuum/strategy/capability/7541d9b3b65d47d286dabdc093cb0eed.md}{Canonical record: 7541d9b3b65d47d286dabdc093cb0eed}; full qualifications and provenance.}\par
\paragraph*{II.1.2.3.1.3 [CT3] Preventative Health Problem Identification}
\phantomsection
\label{navigation:II.1.2.3.1.3}
\addcontentsline{toc}{paragraph}{II.1.2.3.1.3 [CT3] Preventative Health Problem Identification}

{\small Historical predecessor: R.2.14 \hyperref[element:b35d3521141c419fb676a79b2ede40bd]{[CT3] Preventative Health Problem Identification}; Canonical ID: \texttt{\seqsplit{b35d3521141c419fb676a79b2ede40bd}}. Superseded historical identity — direct correspondence · superseded.}\par
\label{element:c1567022671f433fa1eeb37a9d156d92}
\index{[CT3] Preventative Health Problem Identification (II.1.2.3.1.3)}
{\small Capability · Canonical · current}\par

Identify and characterise a preventative-health problem warranting organised attention.



This capability establishes the preventative-health problem warranting organised attention. It gives an intended programme a reason grounded in the concern being addressed rather than assuming that a population statistic already defines the response. The identified problem can frame subsequent cohort and programme considerations without becoming an individual's clinical diagnosis.



\textbf{Boundary}



Population evidence informs problem identification but does not by itself establish a programme, cohort or individual care need.


{\small Recorded unresolved matters: source-qualified matters; exact scope remains in the canonical record.}\par
{\footnotesize \href{../../../../continuum/strategy/capability/c1567022671f433fa1eeb37a9d156d92.md}{Canonical record: c1567022671f433fa1eeb37a9d156d92}; full qualifications and provenance.}\par
\paragraph*{II.1.2.3.1.4 [CT3] Preventative Health Programme Management}
\phantomsection
\label{navigation:II.1.2.3.1.4}
\addcontentsline{toc}{paragraph}{II.1.2.3.1.4 [CT3] Preventative Health Programme Management}

{\small Historical predecessor: R.2.15 \hyperref[element:d71def9e76a544ff9761c672bb24f3be]{[CT3] Preventative Health Programme Management}; Canonical ID: \texttt{\seqsplit{d71def9e76a544ff9761c672bb24f3be}}. Superseded historical identity — direct correspondence · superseded.}\par
\label{element:563d36b179d34cb6913253d06e8d77f6}
\index{[CT3] Preventative Health Programme Management (II.1.2.3.1.4)}
{\small Capability · Canonical · current}\par

Establish and manage an organised response to a preventative-health need/objective without absorbing its delivery capabilities.



Preventative Health Programme Management establishes and manages an organised response to a preventative need or objective. It keeps the programme purpose and coordinated activity intelligible while allowing appointments, interventions and other delivery abilities to retain their own responsibilities. Screening and immunisation illustrate programmes rather than creating additional capability tiers.



\textbf{Boundary}



Managing the programme does not acquire the clinical or operational authority of the activities delivering it.


{\small Recorded unresolved matters: source-qualified matters; exact scope remains in the canonical record.}\par
{\footnotesize \href{../../../../continuum/strategy/capability/563d36b179d34cb6913253d06e8d77f6.md}{Canonical record: 563d36b179d34cb6913253d06e8d77f6}; full qualifications and provenance.}\par
\paragraph*{II.1.2.3.1.5 [CT3] Preventative Health Programme Monitoring}
\phantomsection
\label{navigation:II.1.2.3.1.5}
\addcontentsline{toc}{paragraph}{II.1.2.3.1.5 [CT3] Preventative Health Programme Monitoring}

{\small Historical predecessor: R.2.16 \hyperref[element:e05f6310f6d544efaca5568604f27dd8]{[CT3] Preventative Health Programme Monitoring}; Canonical ID: \texttt{\seqsplit{e05f6310f6d544efaca5568604f27dd8}}. Superseded historical identity — direct correspondence · superseded.}\par
\label{element:945b40ced8004e1cbf7c98c731d17daf}
\index{[CT3] Preventative Health Programme Monitoring (II.1.2.3.1.5)}
{\small Capability · Canonical · current}\par

Observe programme operation, reach, participation, coverage, progress and applicable performance.



This capability observes how a preventative programme is operating, including reach, participation, coverage and progress. It supplies an account of programme performance in the programme's own scope, rather than substituting population-health state for evidence of programme activity. That account can inform reporting or consideration of the programme without itself being an assurance conclusion.



\textbf{Boundary}



Programme monitoring is distinct from individual health monitoring, population assessment and independent assurance.


{\small Recorded unresolved matters: source-qualified matters; exact scope remains in the canonical record.}\par
{\footnotesize \href{../../../../continuum/strategy/capability/945b40ced8004e1cbf7c98c731d17daf.md}{Canonical record: 945b40ced8004e1cbf7c98c731d17daf}; full qualifications and provenance.}\par
\paragraph*{II.1.2.3.1.6 [CT3] Preventative Health Reporting}
\phantomsection
\label{navigation:II.1.2.3.1.6}
\addcontentsline{toc}{paragraph}{II.1.2.3.1.6 [CT3] Preventative Health Reporting}

{\small Historical predecessor: R.2.17 \hyperref[element:4faef45b4ebe4fa69e57c3d81bfea05f]{[CT3] Preventative Health Reporting}; Canonical ID: \texttt{\seqsplit{4faef45b4ebe4fa69e57c3d81bfea05f}}. Superseded historical identity — direct correspondence · superseded.}\par
\label{element:d2ef01c4cab348df822e802ab084faae}
\index{[CT3] Preventative Health Reporting (II.1.2.3.1.6)}
{\small Capability · Canonical · current}\par

Report established preventative-health problems, programmes, cohorts, activity, performance and outcomes.



Preventative Health Reporting communicates established information about preventative problems, programmes, cohorts, activity and outcomes. It brings these different subjects together for a reporting purpose while preserving their distinct meanings. A programme's reach and an observed population outcome can both be relevant without implying that one proves the other.



\textbf{Boundary}



Reporting does not originate the underlying problem, cohort or evaluation conclusion.


{\small Recorded unresolved matters: source-qualified matters; exact scope remains in the canonical record.}\par
{\footnotesize \href{../../../../continuum/strategy/capability/d2ef01c4cab348df822e802ab084faae.md}{Canonical record: d2ef01c4cab348df822e802ab084faae}; full qualifications and provenance.}\par
\subsection*{II.1.2.4 [CT1] Connected Health Services}
\phantomsection
\label{navigation:II.1.2.4}
\addcontentsline{toc}{subsection}{II.1.2.4 [CT1] Connected Health Services}

{\small Historical predecessor: R.2.20 \hyperref[element:48f01ef8b0e6431c85cf4f40cf8f94c3]{[CT1] Connected Health Services}; Canonical ID: \texttt{\seqsplit{48f01ef8b0e6431c85cf4f40cf8f94c3}}. Superseded historical identity — direct correspondence · superseded.}\par
\label{element:3ef397463e584737b367e93f0bf1d85d}
\index{[CT1] Connected Health Services (II.1.2.4)}
{\small Capability · Canonical · current}\par

Exchange health information appropriately and provide digitally mediated communication, notifications and collaborative interaction between health-service participants.



Connected Health Services represents the organisation's ability to exchange health information and communicate digitally between health-service participants. Its two children distinguish exchange of health-information artefacts from participant communication, notification and collaboration. Both support connected care, but a user viewing information through a screen is not thereby participating in a communication channel.



\textbf{Intent}



Enable appropriate information movement and interaction between otherwise separate participants and services.



\textbf{Boundary}



This organisational capability is not defined by an interface engine or messaging product. General presentation remains a separate enabling concern.



\textbf{Harmonia Architectural Relevance}



Information exchange and digitally mediated interaction are direct platform responsibilities within their assigned contracts, supporting the Core classification without transferring payload or clinical authority.


{\footnotesize \href{../../../../continuum/strategy/capability/3ef397463e584737b367e93f0bf1d85d.md}{Canonical record: 3ef397463e584737b367e93f0bf1d85d}; full qualifications and provenance.}\par
\subsubsection*{II.1.2.4.1 [CT2] Digital Communication Channel}
\phantomsection
\label{navigation:II.1.2.4.1}
\addcontentsline{toc}{subsubsection}{II.1.2.4.1 [CT2] Digital Communication Channel}

\label{element:f5defe141228490598115fe41db89bd7}
\index{[CT2] Digital Communication Channel (II.1.2.4.1)}
{\small Capability · Canonical · current}\par

Provide digitally mediated communication, notification and collaborative interaction between patients, practitioners and other health-service participants.



Digital Communication Channel supports communication between patients, practitioners and other participants through digitally mediated messages, notifications and collaborative interaction. Its organisational purpose is the exchange between participants: for example, a notification can connect someone to a care event without itself delivering care or changing the event's originating facts.



\textbf{Boundary}



A presentation interface used to exercise another capability is not automatically this communication ability. No CT3 set or general presentation BEC placement is established here.


{\small Recorded unresolved matters: source-qualified matters; exact scope remains in the canonical record.}\par
{\footnotesize \href{../../../../continuum/strategy/capability/f5defe141228490598115fe41db89bd7.md}{Canonical record: f5defe141228490598115fe41db89bd7}; full qualifications and provenance.}\par
\subsubsection*{II.1.2.4.2 [CT2] Health Information Exchange}
\phantomsection
\label{navigation:II.1.2.4.2}
\addcontentsline{toc}{subsubsection}{II.1.2.4.2 [CT2] Health Information Exchange}

\label{element:baad6b6117084311acbf5fc4044dfabf}
\index{[CT2] Health Information Exchange (II.1.2.4.2)}
{\small Capability · Canonical · current}\par

Exchange health information appropriately between participating parties, systems and services.



As a Business Capability, Health Information Exchange is the organisational ability to move health information appropriately between participating parties, systems and services. It explains the enterprise purpose served by exchange without defining that purpose as a protocol, server or software product. Its placement under Connected Health Services distinguishes information exchange from the broader interaction and notification purposes of Digital Communication Channel.



\textbf{Boundary}



No CT3 decomposition is established. The identically named BEC is a distinct enabling responsibility; the naming conflict remains unresolved, and this entry does not merge the two or assign the sharing-programme concern a parent.


{\small Recorded unresolved matters: source-qualified matters; exact scope remains in the canonical record.}\par
{\footnotesize \href{../../../../continuum/strategy/capability/baad6b6117084311acbf5fc4044dfabf.md}{Canonical record: baad6b6117084311acbf5fc4044dfabf}; full qualifications and provenance.}\par
\paragraph*{II.1.2.4.2.1 [CT3] Health Information Sharing Programme Management}
\phantomsection
\label{navigation:II.1.2.4.2.1}
\addcontentsline{toc}{paragraph}{II.1.2.4.2.1 [CT3] Health Information Sharing Programme Management}

\label{element:6c5c2cb625414b149ac7f2ccb4b8a897}
\index{[CT3] Health Information Sharing Programme Management (II.1.2.4.2.1)}
{\small Capability · Candidate · current}\par

Definition unresolved in the baseline. No definition has been synthesised from the name.



\textbf{Proposed overview}



Programme-level coordination supplies a plausible frame for the preserved name without replacing the established exchange abilities.



\textbf{Proposed definition}



Coordinate the governed scope, participants and change context of a health-information-sharing programme.



\textbf{Proposed boundary}



The element remains Candidate. The definition/placement elaboration is Proposed and does not adjudicate the unresolved RAD question or transfer source-domain authority.



\medskip\hrule\medskip

{\small Recorded unresolved matters: source-qualified matters; exact scope remains in the canonical record.}\par
{\footnotesize \href{../../../../continuum/strategy/capability/6c5c2cb625414b149ac7f2ccb4b8a897.md}{Canonical record: 6c5c2cb625414b149ac7f2ccb4b8a897}; full qualifications and provenance.}\par
\subsection*{II.1.2.5 [CT1] Corporate Resources \& Enterprise Services}
\phantomsection
\label{navigation:II.1.2.5}
\addcontentsline{toc}{subsection}{II.1.2.5 [CT1] Corporate Resources \& Enterprise Services}

{\small Historical predecessor: R.2.21 \hyperref[element:a2dbd74847a44f4c879983229f5c9d08]{[CT1] Corporate Resources \& Enterprise Services}; Canonical ID: \texttt{\seqsplit{a2dbd74847a44f4c879983229f5c9d08}}. Superseded historical identity — direct correspondence · superseded.}\par
\label{element:b79628f156e14a8cbb82a0dea2fc57ff}
\index{[CT1] Corporate Resources \& Enterprise Services (II.1.2.5)}
{\small Capability · Canonical · current}\par

Provide the bounded clinical-activity charging information and enterprise support context required for R1/R2, without assigning enterprise finance, HR, facilities or ICT operating authority to Harmonia.



This capability retains the financial and corporate-support context material to R1/R2. The financial branch concentrates on information about clinical activity relevant to charging or cost; corporate support covers administrative, physical, technical, people and communication conditions affecting service activity. The limited decomposition makes these dependencies visible without assigning a complete corporate enterprise architecture to Harmonia.



\textbf{Intent}



Support health-service activity with the required corporate context and legitimate activity information.



\textbf{Boundary}



Evidence about cost-relevant activity does not perform billing or accounting. Consuming corporate resource state does not acquire the authority to operate those resources.



\textbf{Harmonia Architectural Relevance}



Clinical activity evidence provides material financial support. Corporate constraints may also affect enablement, but contextual consumption does not classify every corporate descendant as Relevant or Core.


{\footnotesize \href{../../../../continuum/strategy/capability/b79628f156e14a8cbb82a0dea2fc57ff.md}{Canonical record: b79628f156e14a8cbb82a0dea2fc57ff}; full qualifications and provenance.}\par
\subsubsection*{II.1.2.5.1 [CT2] Corporate Support Services}
\phantomsection
\label{navigation:II.1.2.5.1}
\addcontentsline{toc}{subsubsection}{II.1.2.5.1 [CT2] Corporate Support Services}

\label{element:c8777bccbe304e36b7d9fbd955d67579}
\index{[CT2] Corporate Support Services (II.1.2.5.1)}
{\small Capability · Canonical · current}\par

Provide enterprise support whose resources, availability, state or constraints may materially affect service activity.



Corporate Support Services frames enterprise support whose resources, availability or constraints affect health-service activity. Its children retain recognisable administrative, facilities, infrastructure, communications, people and operational-oversight responsibilities. This explains why their context matters to service delivery while keeping the people who operate those concerns accountable for them.



\textbf{Intent}



Make material support conditions visible to health-service activity without confusing consumption with ownership.



\textbf{Boundary}



No further R1/R2 decomposition is established. Context from these abilities does not settle the collective service-readiness placement question.


{\small Recorded unresolved matters: source-qualified matters; exact scope remains in the canonical record.}\par
{\footnotesize \href{../../../../continuum/strategy/capability/c8777bccbe304e36b7d9fbd955d67579.md}{Canonical record: c8777bccbe304e36b7d9fbd955d67579}; full qualifications and provenance.}\par
\paragraph*{II.1.2.5.1.1 [CT3] Back-Office Services}
\phantomsection
\label{navigation:II.1.2.5.1.1}
\addcontentsline{toc}{paragraph}{II.1.2.5.1.1 [CT3] Back-Office Services}

\label{element:9f4d3705403647f4a1e25bc3dece035f}
\index{[CT3] Back-Office Services (II.1.2.5.1.1)}
{\small Capability · Canonical · current}\par

Provide administrative/support services on which health-service activity depends.



Back-Office Services provides the administrative and support activity on which health-service operations depend. Within Corporate Support Services it retains an enterprise context for dependencies that may affect whether service activity can proceed. The model stops at this broad responsibility rather than allocating each administrative function to Harmonia.



\textbf{Boundary}



Consumption of support status does not transfer administrative operating authority or establish a more detailed decomposition.


{\small Recorded unresolved matters: source-qualified matters; exact scope remains in the canonical record.}\par
{\footnotesize \href{../../../../continuum/strategy/capability/9f4d3705403647f4a1e25bc3dece035f.md}{Canonical record: 9f4d3705403647f4a1e25bc3dece035f}; full qualifications and provenance.}\par
\paragraph*{II.1.2.5.1.2 [CT3] Communications \& Publicity}
\phantomsection
\label{navigation:II.1.2.5.1.2}
\addcontentsline{toc}{paragraph}{II.1.2.5.1.2 [CT3] Communications \& Publicity}

\label{element:1fec53fb12444abcac41094dace1d3c2}
\index{[CT3] Communications \& Publicity (II.1.2.5.1.2)}
{\small Capability · Canonical · current}\par

Manage organisational communication/publication arrangements relevant to information availability/distribution.



Communications \& Publicity manages organisational publication and communication arrangements relevant to information availability and distribution. It provides enterprise communication context without being identical to clinical record exchange or digitally mediated participant interaction. The governed decomposition retains that broad support role without adding specific publicity or communications functions.



\textbf{Boundary}



Organisational communication does not originate the clinical facts or governed conclusions it may convey.


{\small Recorded unresolved matters: source-qualified matters; exact scope remains in the canonical record.}\par
{\footnotesize \href{../../../../continuum/strategy/capability/1fec53fb12444abcac41094dace1d3c2.md}{Canonical record: 1fec53fb12444abcac41094dace1d3c2}; full qualifications and provenance.}\par
\paragraph*{II.1.2.5.1.3 [CT3] Facilities Management}
\phantomsection
\label{navigation:II.1.2.5.1.3}
\addcontentsline{toc}{paragraph}{II.1.2.5.1.3 [CT3] Facilities Management}

\label{element:9c82e623bd174409b90d24af6b58ebf0}
\index{[CT3] Facilities Management (II.1.2.5.1.3)}
{\small Capability · Canonical · current}\par

Manage physical facilities/services and applicable resource, location, availability and facility-state context.



Facilities Management manages physical facilities and services, including the state and availability context relevant to health activity. It contributes the real-world conditions that location definitions alone cannot establish. This distinguishes a place recognised in a directory from a facility being operated and maintained for use.



\textbf{Boundary}



Harmonia's consumption of facility constraints does not make it the facilities operator or establish collective service readiness.


{\small Recorded unresolved matters: source-qualified matters; exact scope remains in the canonical record.}\par
{\footnotesize \href{../../../../continuum/strategy/capability/9c82e623bd174409b90d24af6b58ebf0.md}{Canonical record: 9c82e623bd174409b90d24af6b58ebf0}; full qualifications and provenance.}\par
\paragraph*{II.1.2.5.1.4 [CT3] IT Infrastructure Management}
\phantomsection
\label{navigation:II.1.2.5.1.4}
\addcontentsline{toc}{paragraph}{II.1.2.5.1.4 [CT3] IT Infrastructure Management}

\label{element:9a37413582314e0dbac006f82900858b}
\index{[CT3] IT Infrastructure Management (II.1.2.5.1.4)}
{\small Capability · Canonical · current}\par

Operate/manage underlying technology infrastructure and its availability.



IT Infrastructure Management operates the underlying technology infrastructure and its availability. It supplies an enterprise support context for health-service activity that relies on functioning technology. Its business outcome is operational infrastructure, distinct from governance of health-information meaning or the application responsibilities realised on that infrastructure.



\textbf{Boundary}



Infrastructure machinery does not define Harmonia's information semantics or acquire clinical authority.


{\small Recorded unresolved matters: source-qualified matters; exact scope remains in the canonical record.}\par
{\footnotesize \href{../../../../continuum/strategy/capability/9a37413582314e0dbac006f82900858b.md}{Canonical record: 9a37413582314e0dbac006f82900858b}; full qualifications and provenance.}\par
\paragraph*{II.1.2.5.1.5 [CT3] Operational Oversight}
\phantomsection
\label{navigation:II.1.2.5.1.5}
\addcontentsline{toc}{paragraph}{II.1.2.5.1.5 [CT3] Operational Oversight}

\label{element:d6f39d94cd2e4bf2ab91ff58ba604ab0}
\index{[CT3] Operational Oversight (II.1.2.5.1.5)}
{\small Capability · Canonical · current}\par

Observe broader operational conditions affecting the organisation’s ability to work or deliver services.



Operational Oversight observes broader operating conditions affecting the organisation's ability to work and deliver services. It provides context beyond an individual work item or care place, helping make material constraints visible. Within corporate support it remains a broad enterprise responsibility rather than an inferred owner for every operational-readiness question.



\textbf{Boundary}



Oversight does not resolve collective service-readiness placement or constitute independent assurance merely through observation.


{\small Recorded unresolved matters: source-qualified matters; exact scope remains in the canonical record.}\par
{\footnotesize \href{../../../../continuum/strategy/capability/d6f39d94cd2e4bf2ab91ff58ba604ab0.md}{Canonical record: d6f39d94cd2e4bf2ab91ff58ba604ab0}; full qualifications and provenance.}\par
\paragraph*{II.1.2.5.1.6 [CT3] People Management}
\phantomsection
\label{navigation:II.1.2.5.1.6}
\addcontentsline{toc}{paragraph}{II.1.2.5.1.6 [CT3] People Management}

\label{element:df8b9d309b354efcb15991b3d75ed852}
\index{[CT3] People Management (II.1.2.5.1.6)}
{\small Capability · Canonical · current}\par

Manage people across the organisation, distinct from practitioner-specific Health Workforce Management.



People Management addresses people across the organisation as an enterprise support responsibility. It retains a scope broader than the practitioner-specific standing, privilege and coverage concerns of Health Workforce Management. This distinction allows health-service activity to consume relevant people context without redefining all workforce administration as clinical governance.



\textbf{Boundary}



General people administration does not itself establish practitioner clinical qualification or privilege.


{\small Recorded unresolved matters: source-qualified matters; exact scope remains in the canonical record.}\par
{\footnotesize \href{../../../../continuum/strategy/capability/df8b9d309b354efcb15991b3d75ed852.md}{Canonical record: df8b9d309b354efcb15991b3d75ed852}; full qualifications and provenance.}\par
\subsubsection*{II.1.2.5.2 [CT2] Financial \& Commercial Services}
\phantomsection
\label{navigation:II.1.2.5.2}
\addcontentsline{toc}{subsubsection}{II.1.2.5.2 [CT2] Financial \& Commercial Services}

\label{element:290b574a02634f559b35a28daad0dea8}
\index{[CT2] Financial \& Commercial Services (II.1.2.5.2)}
{\small Capability · Canonical · current}\par

Provide the financial/commercial business context materially required for R1/R2.



Financial \& Commercial Services provides the financial business frame needed for the bounded R1/R2 responsibility. Its established child identifies and supplies information about chargeable or cost-relevant clinical activity. The purpose is to make applicable activity evidence available to the responsible financial context without expanding the model into pricing, claims or revenue management.



\textbf{Intent}



Support legitimate funding, financial and charging purposes with relevant clinical-activity information.



\textbf{Boundary}



Not every exchanged request or activity is chargeable. Provision of evidence does not establish a charging decision or an enterprise financial system.


{\small Recorded unresolved matters: source-qualified matters; exact scope remains in the canonical record.}\par
{\footnotesize \href{../../../../continuum/strategy/capability/290b574a02634f559b35a28daad0dea8.md}{Canonical record: 290b574a02634f559b35a28daad0dea8}; full qualifications and provenance.}\par
\paragraph*{II.1.2.5.2.1 [CT3] Clinical Activity Charging Information Provision}
\phantomsection
\label{navigation:II.1.2.5.2.1}
\addcontentsline{toc}{paragraph}{II.1.2.5.2.1 [CT3] Clinical Activity Charging Information Provision}

\label{element:ce9655c00f044ca996c9d4a1222785f5}
\index{[CT3] Clinical Activity Charging Information Provision (II.1.2.5.2.1)}
{\small Capability · Canonical · current}\par

Identify and provide information about chargeable or cost-relevant clinical/health-service activity where required for applicable funding, financial or charging purposes.



This capability identifies and supplies information about clinical or health-service activity relevant to legitimate charging, funding or cost purposes. Requests, service types, quantities and utilisation context may contribute evidence, but chargeability cannot be inferred from the fact that Harmonia mediated an interaction. Its outcome is usable activity information for the responsible financial context.



\textbf{Boundary}



Evidence provision does not perform pricing, billing, claims, accounting or revenue management.


{\small Recorded unresolved matters: source-qualified matters; exact scope remains in the canonical record.}\par
{\footnotesize \href{../../../../continuum/strategy/capability/ce9655c00f044ca996c9d4a1222785f5.md}{Canonical record: ce9655c00f044ca996c9d4a1222785f5}; full qualifications and provenance.}\par
\subsection*{II.1.2.6 [CT1] Diagnostic, Therapeutic \& Clinical Support Services}
\phantomsection
\label{navigation:II.1.2.6}
\addcontentsline{toc}{subsection}{II.1.2.6 [CT1] Diagnostic, Therapeutic \& Clinical Support Services}

{\small Historical predecessor: R.2.22 \hyperref[element:e003b9cd0730459a84775d2837d1d173]{[CT1] Diagnostic, Therapeutic \& Clinical Support Services}; Canonical ID: \texttt{\seqsplit{e003b9cd0730459a84775d2837d1d173}}. Superseded historical identity — direct correspondence · superseded.}\par
\label{element:4d98c3503a9f4890a677f09b0c3b6f82}
\index{[CT1] Diagnostic, Therapeutic \& Clinical Support Services (II.1.2.6)}
{\small Capability · Canonical · current}\par

The specialised clinical services that support direct patient care, including pathology and laboratory testing, diagnostic imaging (radiology), pharmacy and medication dispensing, and allied health therapies.



Specialised services contribute evidence, professional interpretation, medication provision and therapeutic support to individual care. The established decomposition distinguishes diagnostic provision from medication and pharmacy services, while retaining wider therapeutic and allied-health scope as enterprise context. The catalogue deliberately explains the modelled branches without presenting them as an exhaustive inventory of clinical support.



\textbf{Intent}



Provide the specialised contributions that treating participants need to assess and care for individuals.



\textbf{Boundary}



Producing diagnostic evidence, authorising a report, dispensing medication and administering treatment establish different outcomes and authorities.



\textbf{Harmonia Architectural Relevance}



Requests, specimen associations, reports and medication events require reliable correlation and exchange. Supporting those information flows does not make Harmonia the diagnostic practitioner, pharmacy or therapy provider.


{\footnotesize \href{../../../../continuum/strategy/capability/4d98c3503a9f4890a677f09b0c3b6f82.md}{Canonical record: 4d98c3503a9f4890a677f09b0c3b6f82}; full qualifications and provenance.}\par
\subsubsection*{II.1.2.6.1 [CT2] Diagnostic Support Provision}
\phantomsection
\label{navigation:II.1.2.6.1}
\addcontentsline{toc}{subsubsection}{II.1.2.6.1 [CT2] Diagnostic Support Provision}

\label{element:827910cf74ae4061afbda52dde3da106}
\index{[CT2] Diagnostic Support Provision (II.1.2.6.1)}
{\small Capability · Canonical · current}\par

Provide diagnostic services and manage their evidence, interpretation and governed reporting.



Diagnostic Support Provision covers the distinct responsibilities through which a diagnostic service obtains evidence and makes its findings available. Diagnostic scheduling, service performance, specimen stewardship, data management, interpretation, report governance and report distribution each establish a different outcome. Their separation allows preliminary, amended or final information to retain its professional and evidentiary standing without treating delivery as authorisation.



\textbf{Intent}



Make diagnostic evidence and professionally governed findings usable by the appropriate care participants.



\textbf{Boundary}



Diagnostic interpretation is limited to the service's professional scope; patient-level Clinical Determination considers the wider clinical context. The children do not prescribe a universal diagnostic workflow or establish a Diagnostic Order Management capability.


{\small Recorded unresolved matters: source-qualified matters; exact scope remains in the canonical record.}\par
{\footnotesize \href{../../../../continuum/strategy/capability/827910cf74ae4061afbda52dde3da106.md}{Canonical record: 827910cf74ae4061afbda52dde3da106}; full qualifications and provenance.}\par
\paragraph*{II.1.2.6.1.1 [CT3] Diagnostic Data Management}
\phantomsection
\label{navigation:II.1.2.6.1.1}
\addcontentsline{toc}{paragraph}{II.1.2.6.1.1 [CT3] Diagnostic Data Management}

\label{element:aeb8e6d2225a4df38e61ae116850b32c}
\index{[CT3] Diagnostic Data Management (II.1.2.6.1.1)}
{\small Capability · Canonical · current}\par

Manage acquired or produced diagnostic evidence as information, including identity, association, integrity, organisation, availability and lifecycle.



Diagnostic Data Management maintains acquired or produced diagnostic evidence as usable information. Its responsibility includes identity, association, integrity, organisation and lifecycle so that evidence can be found and understood in the right diagnostic context. It addresses the evidence record independently of the professional conclusion subsequently drawn from it.



\textbf{Boundary}



Managing diagnostic data neither interprets that data nor governs the authoritative expression of findings in a report.


{\small Recorded unresolved matters: source-qualified matters; exact scope remains in the canonical record.}\par
{\footnotesize \href{../../../../continuum/strategy/capability/aeb8e6d2225a4df38e61ae116850b32c.md}{Canonical record: aeb8e6d2225a4df38e61ae116850b32c}; full qualifications and provenance.}\par
\paragraph*{II.1.2.6.1.2 [CT3] Diagnostic Interpretation}
\phantomsection
\label{navigation:II.1.2.6.1.2}
\addcontentsline{toc}{paragraph}{II.1.2.6.1.2 [CT3] Diagnostic Interpretation}

\label{element:82ad867109f1438eb1b0916ed9cd820f}
\index{[CT3] Diagnostic Interpretation (II.1.2.6.1.2)}
{\small Capability · Canonical · current}\par

Establish what diagnostic evidence demonstrates within the diagnostic service’s professional scope.



Diagnostic Interpretation establishes what diagnostic evidence demonstrates within the service's professional scope. It turns evidence into diagnostic findings without claiming the wider patient-level judgement that considers all clinical circumstances. Separating interpretation from acquisition and reporting makes the professional contribution clear even when one service performs all three.



\textbf{Boundary}



Patient-level Clinical Determination remains distinct, as do report authorisation and distribution.


{\small Recorded unresolved matters: source-qualified matters; exact scope remains in the canonical record.}\par
{\footnotesize \href{../../../../continuum/strategy/capability/82ad867109f1438eb1b0916ed9cd820f.md}{Canonical record: 82ad867109f1438eb1b0916ed9cd820f}; full qualifications and provenance.}\par
\paragraph*{II.1.2.6.1.3 [CT3] Diagnostic Report Distribution}
\phantomsection
\label{navigation:II.1.2.6.1.3}
\addcontentsline{toc}{paragraph}{II.1.2.6.1.3 [CT3] Diagnostic Report Distribution}

\label{element:c8fdba3d18e143acb2bd77ae828afd4b}
\index{[CT3] Diagnostic Report Distribution (II.1.2.6.1.3)}
{\small Capability · Canonical · current}\par

Provide an established report to applicable authorised recipients or consuming contexts under applicable access, privacy, security and communication controls.



Diagnostic Report Distribution provides an already established report to applicable authorised recipients or consuming contexts. Its purpose is to make the governed findings available where they are needed, while respecting access, privacy, security and communication controls. Distribution retains the report's standing rather than originating or silently reinterpreting its findings.



\textbf{Boundary}



Delivery does not establish clinical review, report authorisation or a patient-level clinical decision.


{\small Recorded unresolved matters: source-qualified matters; exact scope remains in the canonical record.}\par
{\footnotesize \href{../../../../continuum/strategy/capability/c8fdba3d18e143acb2bd77ae828afd4b.md}{Canonical record: c8fdba3d18e143acb2bd77ae828afd4b}; full qualifications and provenance.}\par
\paragraph*{II.1.2.6.1.4 [CT3] Diagnostic Report Management}
\phantomsection
\label{navigation:II.1.2.6.1.4}
\addcontentsline{toc}{paragraph}{II.1.2.6.1.4 [CT3] Diagnostic Report Management}

\label{element:977ed170c34445ddb5095945f54f6907}
\index{[CT3] Diagnostic Report Management (II.1.2.6.1.4)}
{\small Capability · Canonical · current}\par

Establish and govern the authoritative expression of diagnostic findings, including applicable creation, review, amendment, verification, finalisation and authorisation/signing.



Diagnostic Report Management governs the authoritative expression of diagnostic findings. Creation, review, verification, amendment, finalisation and signing matter because they establish the standing of the report rather than merely its existence as data. The responsibility preserves a meaningful distinction between evidence, interpretation and the report that a service is prepared to authorise.



\textbf{Boundary}



Signing is behaviour within this capability, not another CT3. Sending a report does not establish its authorisation.


{\small Recorded unresolved matters: source-qualified matters; exact scope remains in the canonical record.}\par
{\footnotesize \href{../../../../continuum/strategy/capability/977ed170c34445ddb5095945f54f6907.md}{Canonical record: 977ed170c34445ddb5095945f54f6907}; full qualifications and provenance.}\par
\paragraph*{II.1.2.6.1.5 [CT3] Diagnostic Service Appointment Management}
\phantomsection
\label{navigation:II.1.2.6.1.5}
\addcontentsline{toc}{paragraph}{II.1.2.6.1.5 [CT3] Diagnostic Service Appointment Management}

\label{element:5410a0275dfa4460ad49100fa5a2f2ec}
\index{[CT3] Diagnostic Service Appointment Management (II.1.2.6.1.5)}
{\small Capability · Canonical · current}\par

Manage the diagnostic-service-specific scheduled opportunity, including preparation, modality, urgency, duration and other service-specific constraints.



Diagnostic Service Appointment Management makes a scheduled diagnostic opportunity credible in the service's own context. Preparation, modality, urgency and duration can materially affect whether an appointment is usable, so the responsibility is richer than recording a generic time slot. It supports service access while remaining separate from performance of the diagnostic activity.



\textbf{Boundary}



A diagnostic booking establishes neither evidence acquisition nor a report and is distinct from generic Care Appointment Management.


{\small Recorded unresolved matters: source-qualified matters; exact scope remains in the canonical record.}\par
{\footnotesize \href{../../../../continuum/strategy/capability/5410a0275dfa4460ad49100fa5a2f2ec.md}{Canonical record: 5410a0275dfa4460ad49100fa5a2f2ec}; full qualifications and provenance.}\par
\paragraph*{II.1.2.6.1.6 [CT3] Diagnostic Service Delivery}
\phantomsection
\label{navigation:II.1.2.6.1.6}
\addcontentsline{toc}{paragraph}{II.1.2.6.1.6 [CT3] Diagnostic Service Delivery}

\label{element:7a1ae953164d456c93637f8e35b4940e}
\index{[CT3] Diagnostic Service Delivery (II.1.2.6.1.6)}
{\small Capability · Canonical · current}\par

Perform the diagnostic activity that obtains or produces diagnostic evidence.



Diagnostic Service Delivery performs the activity that obtains or produces diagnostic evidence. Within Diagnostic Support Provision it is the point at which the diagnostic procedure actually occurs, rather than being scheduled or requested. The evidence obtained can then be managed and interpreted under the distinct responsibilities that give it professional meaning.



\textbf{Boundary}



Performance does not itself establish interpretation, report authorisation or report delivery.


{\small Recorded unresolved matters: source-qualified matters; exact scope remains in the canonical record.}\par
{\footnotesize \href{../../../../continuum/strategy/capability/7a1ae953164d456c93637f8e35b4940e.md}{Canonical record: 7a1ae953164d456c93637f8e35b4940e}; full qualifications and provenance.}\par
\paragraph*{II.1.2.6.1.7 [CT3] Diagnostic Specimen Management}
\phantomsection
\label{navigation:II.1.2.6.1.7}
\addcontentsline{toc}{paragraph}{II.1.2.6.1.7 [CT3] Diagnostic Specimen Management}

\label{element:93ab99a055314129984a18ea85db5cc2}
\index{[CT3] Diagnostic Specimen Management (II.1.2.6.1.7)}
{\small Capability · Canonical · current}\par

Maintain a specimen’s identity, integrity, suitability, custody and applicable lifecycle context.



Diagnostic Specimen Management preserves the specimen context needed to trust its use: which specimen it is, whether it remains suitable, and how integrity and custody are maintained. It concerns the physical diagnostic subject across its applicable lifecycle rather than only the results generated from it. This distinction protects the relationship between evidence and the material from which it arose.



\textbf{Boundary}



Specimen stewardship is separate from diagnostic-data management and interpretation of results.


{\small Recorded unresolved matters: source-qualified matters; exact scope remains in the canonical record.}\par
{\footnotesize \href{../../../../continuum/strategy/capability/93ab99a055314129984a18ea85db5cc2.md}{Canonical record: 93ab99a055314129984a18ea85db5cc2}; full qualifications and provenance.}\par
\subsubsection*{II.1.2.6.2 [CT2] Medication \& Pharmacy Services}
\phantomsection
\label{navigation:II.1.2.6.2}
\addcontentsline{toc}{subsubsection}{II.1.2.6.2 [CT2] Medication \& Pharmacy Services}

\label{element:ab2a4e6a7acc4868bc4e7b6c0d292cb4}
\index{[CT2] Medication \& Pharmacy Services (II.1.2.6.2)}
{\small Capability · Canonical · current}\par

Provide professional medication/pharmacy services under applicable controls.



This capability frames professional medication provision and medication-specific advice. Prescription dispensing, non-prescription provision, advice and dose administration aids answer different needs, even where exercised by one pharmacy. The decomposition keeps the authorised regimen, provision of medicine and actual administration separate so that a supply event is not mistaken for treatment performed.



\textbf{Intent}



Support safe, professionally controlled access to and use of medication within the established pharmacy responsibility.



\textbf{Boundary}



Advice does not confer prescribing authority, and dose aids do not establish a regimen. Whether specialised compounding/preparation warrants a separate capability remains unresolved under OI-BC-001.


{\small Recorded unresolved matters: source-qualified matters; exact scope remains in the canonical record.}\par
{\footnotesize \href{../../../../continuum/strategy/capability/ab2a4e6a7acc4868bc4e7b6c0d292cb4.md}{Canonical record: ab2a4e6a7acc4868bc4e7b6c0d292cb4}; full qualifications and provenance.}\par
\paragraph*{II.1.2.6.2.1 [CT3] Medication Advice}
\phantomsection
\label{navigation:II.1.2.6.2.1}
\addcontentsline{toc}{paragraph}{II.1.2.6.2.1 [CT3] Medication Advice}

\label{element:b867a2fde009437196a5fed682abf073}
\index{[CT3] Medication Advice (II.1.2.6.2.1)}
{\small Capability · Canonical · current}\par

Provide professional medication-specific advice about appropriate use, precautions, interactions, administration, adherence or safe handling.



Medication Advice provides professional guidance about medication use and its precautions, interactions, administration, adherence or handling. It supports understanding of how medication can be used safely without turning every advisory interaction into a prescribing decision. Within pharmacy services, the advice outcome is distinct from provision of the medicine itself.



\textbf{Boundary}



Advice does not confer authority to diagnose, prescribe or alter treatment; detailed scope remains subject to refinement.


{\small Recorded unresolved matters: source-qualified matters; exact scope remains in the canonical record.}\par
{\footnotesize \href{../../../../continuum/strategy/capability/b867a2fde009437196a5fed682abf073.md}{Canonical record: b867a2fde009437196a5fed682abf073}; full qualifications and provenance.}\par
\paragraph*{II.1.2.6.2.2 [CT3] Medication Compounding \& Preparation}
\phantomsection
\label{navigation:II.1.2.6.2.2}
\addcontentsline{toc}{paragraph}{II.1.2.6.2.2 [CT3] Medication Compounding \& Preparation}

\label{element:07380933da66427c865ca8e64d414e59}
\index{[CT3] Medication Compounding \& Preparation (II.1.2.6.2.2)}
{\small Capability · Proposed · current}\par

Prepare or compound a clinically required medication formulation under the applicable professional and product-governance mandate.



\textbf{Overview}



Prepare or compound a clinically required medication formulation under the applicable professional and product-governance mandate. The proposed decomposition separates this outcome from the other responsibilities within Medication \& Pharmacy Services.



\textbf{Intent}



Provide the required medication formulation when ordinary supply alone does not satisfy the authorised preparation need.



\textbf{Boundary}



Preparation differs from dispensing and administration; this proposal assigns no compounding activity to Harmonia.


{\footnotesize \href{../../../../continuum/strategy/capability/07380933da66427c865ca8e64d414e59.md}{Canonical record: 07380933da66427c865ca8e64d414e59}; full qualifications and provenance.}\par
\paragraph*{II.1.2.6.2.3 [CT3] Medication Dose Administration Aid Provision}
\phantomsection
\label{navigation:II.1.2.6.2.3}
\addcontentsline{toc}{paragraph}{II.1.2.6.2.3 [CT3] Medication Dose Administration Aid Provision}

\label{element:bd597608766b4d709703977c02cc5f07}
\index{[CT3] Medication Dose Administration Aid Provision (II.1.2.6.2.3)}
{\small Capability · Canonical · current}\par

Prepare and provide organised medication doses supporting administration according to an already established regimen.



This capability prepares and supplies organised doses to support administration under an already established medication regimen. Its purpose is to make that regimen easier to follow in practice while preserving the difference between organising doses and choosing therapy. The business responsibility is not defined by any particular commercial packaging arrangement.



\textbf{Boundary}



Preparing a dose aid does not establish the regimen or demonstrate that a dose was administered.


{\small Recorded unresolved matters: source-qualified matters; exact scope remains in the canonical record.}\par
{\footnotesize \href{../../../../continuum/strategy/capability/bd597608766b4d709703977c02cc5f07.md}{Canonical record: bd597608766b4d709703977c02cc5f07}; full qualifications and provenance.}\par
\paragraph*{II.1.2.6.2.4 [CT3] Non-Prescription Medication Provision}
\phantomsection
\label{navigation:II.1.2.6.2.4}
\addcontentsline{toc}{paragraph}{II.1.2.6.2.4 [CT3] Non-Prescription Medication Provision}

\label{element:7bbe4810fb974a8ab6d47b30d6475df3}
\index{[CT3] Non-Prescription Medication Provision (II.1.2.6.2.4)}
{\small Capability · Canonical · current}\par

Provide medication not requiring an applicable prescription, subject to professional, regulatory and safety controls.



Non-Prescription Medication Provision makes medication available where the applicable conditions do not require a prescription. It remains a professional, regulated and safety-constrained responsibility rather than an unrestricted supply transaction. Its distinction from prescription dispensing preserves the different authority basis under which provision occurs.



\textbf{Boundary}



Provision neither establishes diagnosis nor proves use or administration of the medication.


{\small Recorded unresolved matters: source-qualified matters; exact scope remains in the canonical record.}\par
{\footnotesize \href{../../../../continuum/strategy/capability/7bbe4810fb974a8ab6d47b30d6475df3.md}{Canonical record: 7bbe4810fb974a8ab6d47b30d6475df3}; full qualifications and provenance.}\par
\paragraph*{II.1.2.6.2.5 [CT3] Prescription Medication Dispensing}
\phantomsection
\label{navigation:II.1.2.6.2.5}
\addcontentsline{toc}{paragraph}{II.1.2.6.2.5 [CT3] Prescription Medication Dispensing}

\label{element:fef5b79eea77438d9d453a7daef744fc}
\index{[CT3] Prescription Medication Dispensing (II.1.2.6.2.5)}
{\small Capability · Canonical · current}\par

Professionally prepare and provide medication according to an applicable authorised prescription/order and dispensing requirements.



Prescription Medication Dispensing is the professional preparation and provision of medication against an applicable authorised prescription or order. It fulfils a pharmacy responsibility grounded in the prescribed intent and dispensing requirements. The resulting dispensing event provides different evidence from the prescribing decision and from later administration to the individual.



\textbf{Boundary}



Dispensing does not establish that medication was administered; unresolved specialised compounding must not be silently assigned here.


{\small Recorded unresolved matters: source-qualified matters; exact scope remains in the canonical record.}\par
{\footnotesize \href{../../../../continuum/strategy/capability/fef5b79eea77438d9d453a7daef744fc.md}{Canonical record: fef5b79eea77438d9d453a7daef744fc}; full qualifications and provenance.}\par
\subsection*{II.1.2.7 [CT1] Health Information \& Knowledge Management}
\phantomsection
\label{navigation:II.1.2.7}
\addcontentsline{toc}{subsection}{II.1.2.7 [CT1] Health Information \& Knowledge Management}

{\small Historical predecessor: R.2.23 \hyperref[element:a1c557bd263047d0b3ff566cf50b98c1]{[CT1] Health Information \& Knowledge Management}; Canonical ID: \texttt{\seqsplit{a1c557bd263047d0b3ff566cf50b98c1}}. Superseded historical identity — direct correspondence · superseded.}\par
\label{element:cca5d5abec204b339eb5b9bd924ed64e}
\index{[CT1] Health Information \& Knowledge Management (II.1.2.7)}
{\small Capability · Canonical · current}\par

Manage governed health-information content, structures, lifecycle, catalogues, compartments, correlation and retrieval, alongside reusable clinical knowledge and its governed standing.



This capability governs both health information and reusable clinical knowledge while preserving their different meanings. Health Information Management concerns content, structures, lifecycle, discovery context, grouping, correlation and retrieval. Clinical Knowledge Management concerns identifying and formulating reusable knowledge and establishing its governed standing. A source clinical fact can contribute to knowledge without becoming a general clinical conclusion merely through reuse.



\textbf{Intent}



Keep information interpretable and usable across time and context while making the standing and limits of reusable knowledge explicit.



\textbf{Boundary}



Managing or preserving information does not originate its clinical facts. Governing structures and binding concepts does not acquire the authority of the terminology they reference.



\textbf{Harmonia Architectural Relevance}



These are substantive information-management responsibilities within Harmonia's assigned boundary. The Core classification remains bounded by information authority, active/durable-state separation and clinical decision authority.


{\footnotesize \href{../../../../continuum/strategy/capability/cca5d5abec204b339eb5b9bd924ed64e.md}{Canonical record: cca5d5abec204b339eb5b9bd924ed64e}; full qualifications and provenance.}\par
\subsubsection*{II.1.2.7.1 [CT2] Clinical Knowledge Management}
\phantomsection
\label{navigation:II.1.2.7.1}
\addcontentsline{toc}{subsubsection}{II.1.2.7.1 [CT2] Clinical Knowledge Management}

\label{element:76b375248aee49aeb5a5044176a688d8}
\index{[CT2] Clinical Knowledge Management (II.1.2.7.1)}
{\small Capability · Canonical · current}\par

Identify, extract, formulate, associate, govern and share reusable clinical knowledge.



Clinical Knowledge Management turns potentially useful source material into explicit reusable knowledge with an intelligible basis and governed standing. Context identification and extraction preserve where material came from; formulation expresses applicability and rationale; concept association supplies semantic context; review and publication govern its intended use. These responsibilities distinguish source information, formulated knowledge and approved knowledge without prescribing a compulsory workflow.



\textbf{Intent}



Make reusable clinical knowledge available with the qualifications and approval standing needed to judge its use.



\textbf{Boundary}



Publication does not establish applicability to an individual or a clinical decision. Concept association does not redefine terminology, and formulation does not constitute approval.


{\small Recorded unresolved matters: source-qualified matters; exact scope remains in the canonical record.}\par
{\footnotesize \href{../../../../continuum/strategy/capability/76b375248aee49aeb5a5044176a688d8.md}{Canonical record: 76b375248aee49aeb5a5044176a688d8}; full qualifications and provenance.}\par
\paragraph*{II.1.2.7.1.1 [CT3] Clinical Knowledge Concept Association}
\phantomsection
\label{navigation:II.1.2.7.1.1}
\addcontentsline{toc}{paragraph}{II.1.2.7.1.1 [CT3] Clinical Knowledge Concept Association}

\label{element:3666bf40ec894f5aac3714dc76df392a}
\index{[CT3] Clinical Knowledge Concept Association (II.1.2.7.1.1)}
{\small Capability · Canonical · current}\par

Associate knowledge with applicable governed clinical concepts without acquiring concept authority.



Clinical Knowledge Concept Association relates reusable knowledge to the governed clinical concepts relevant to its meaning. It makes the semantic context of the knowledge intelligible without allowing local association to redefine the referenced concepts. The association assists interpretation while preserving the distinction between knowledge authority and concept authority.



\textbf{Boundary}



Associating knowledge with a terminology concept does not acquire terminology governance or establish approval of the knowledge.


{\small Recorded unresolved matters: source-qualified matters; exact scope remains in the canonical record.}\par
{\footnotesize \href{../../../../continuum/strategy/capability/3666bf40ec894f5aac3714dc76df392a.md}{Canonical record: 3666bf40ec894f5aac3714dc76df392a}; full qualifications and provenance.}\par
\paragraph*{II.1.2.7.1.2 [CT3] Clinical Knowledge Context Extraction}
\phantomsection
\label{navigation:II.1.2.7.1.2}
\addcontentsline{toc}{paragraph}{II.1.2.7.1.2 [CT3] Clinical Knowledge Context Extraction}

\label{element:0911997eb4584172afd10b1d37946061}
\index{[CT3] Clinical Knowledge Context Extraction (II.1.2.7.1.2)}
{\small Capability · Canonical · current}\par

Extract relevant material while preserving its source/context relationship.



Clinical Knowledge Context Extraction selects relevant material while retaining its relationship to the source and context. Its purpose is to make potential knowledge material available without stripping away the circumstances needed to understand it. Extracted material remains distinguishable from the reusable knowledge subsequently formulated from it.



\textbf{Boundary}



Extraction does not establish approval or turn source information into an individual clinical conclusion.


{\small Recorded unresolved matters: source-qualified matters; exact scope remains in the canonical record.}\par
{\footnotesize \href{../../../../continuum/strategy/capability/0911997eb4584172afd10b1d37946061.md}{Canonical record: 0911997eb4584172afd10b1d37946061}; full qualifications and provenance.}\par
\paragraph*{II.1.2.7.1.3 [CT3] Clinical Knowledge Context Identification}
\phantomsection
\label{navigation:II.1.2.7.1.3}
\addcontentsline{toc}{paragraph}{II.1.2.7.1.3 [CT3] Clinical Knowledge Context Identification}

\label{element:a0facf9015514223a928e39fcbbb6bee}
\index{[CT3] Clinical Knowledge Context Identification (II.1.2.7.1.3)}
{\small Capability · Canonical · current}\par

Identify sources/contexts potentially containing reusable clinical knowledge.



This capability identifies sources or contexts that may contain reusable clinical knowledge. It establishes where potentially useful material can be sought without assuming that all source information is already formulated knowledge. Within knowledge management it provides a contextual starting point rather than an approval or publication outcome.



\textbf{Boundary}



Identifying a source does not extract its material, establish reusable knowledge or determine its clinical applicability.


{\small Recorded unresolved matters: source-qualified matters; exact scope remains in the canonical record.}\par
{\footnotesize \href{../../../../continuum/strategy/capability/a0facf9015514223a928e39fcbbb6bee.md}{Canonical record: a0facf9015514223a928e39fcbbb6bee}; full qualifications and provenance.}\par
\paragraph*{II.1.2.7.1.4 [CT3] Clinical Knowledge Information Formulation}
\phantomsection
\label{navigation:II.1.2.7.1.4}
\addcontentsline{toc}{paragraph}{II.1.2.7.1.4 [CT3] Clinical Knowledge Information Formulation}

\label{element:05421be130ff4020b119c314aca46573}
\index{[CT3] Clinical Knowledge Information Formulation (II.1.2.7.1.4)}
{\small Capability · Canonical · current}\par

Formulate explicit reusable knowledge, including applicability, qualifications, rationale and evidence relationships.



This capability expresses reusable clinical knowledge with its applicability, qualifications, rationale and evidence relationships. It moves beyond extracted material by making the knowledge claim and its limits explicit. That formulation supports later review and use while retaining the distinction between a proposed knowledge expression and one approved for a governed purpose.



\textbf{Boundary}



Formulation does not approve knowledge, govern terminology or define health-information structures.


{\small Recorded unresolved matters: source-qualified matters; exact scope remains in the canonical record.}\par
{\footnotesize \href{../../../../continuum/strategy/capability/05421be130ff4020b119c314aca46573.md}{Canonical record: 05421be130ff4020b119c314aca46573}; full qualifications and provenance.}\par
\paragraph*{II.1.2.7.1.5 [CT3] Clinical Knowledge Publication \& Sharing}
\phantomsection
\label{navigation:II.1.2.7.1.5}
\addcontentsline{toc}{paragraph}{II.1.2.7.1.5 [CT3] Clinical Knowledge Publication \& Sharing}

\label{element:962a5e3048424ff7b2413e21a8e51f53}
\index{[CT3] Clinical Knowledge Publication \& Sharing (II.1.2.7.1.5)}
{\small Capability · Canonical · current}\par

Make approved knowledge available to applicable authorised consumers and use contexts.



Clinical Knowledge Publication \& Sharing makes approved knowledge available to applicable authorised consumers and contexts. It preserves the standing and limits under which knowledge was approved rather than reducing publication to unrestricted dissemination. Its contribution is availability for appropriate use, distinct from both approval and the individual clinical application that may follow.



\textbf{Boundary}



Published knowledge is not a patient-specific conclusion, decision or instruction merely because a clinician can access it.


{\small Recorded unresolved matters: source-qualified matters; exact scope remains in the canonical record.}\par
{\footnotesize \href{../../../../continuum/strategy/capability/962a5e3048424ff7b2413e21a8e51f53.md}{Canonical record: 962a5e3048424ff7b2413e21a8e51f53}; full qualifications and provenance.}\par
\paragraph*{II.1.2.7.1.6 [CT3] Clinical Knowledge Review \& Approval Management}
\phantomsection
\label{navigation:II.1.2.7.1.6}
\addcontentsline{toc}{paragraph}{II.1.2.7.1.6 [CT3] Clinical Knowledge Review \& Approval Management}

\label{element:ee04d00077d548a7a0a2dffd424c3d78}
\index{[CT3] Clinical Knowledge Review \& Approval Management (II.1.2.7.1.6)}
{\small Capability · Canonical · current}\par

Manage review, endorsement, approval/rejection, conditions and applicable history for intended governed use.



This capability manages the review and governed standing of reusable knowledge, including endorsement, rejection, conditions and history. It establishes how the knowledge may legitimately be treated for its intended use rather than assuming that formulation makes it trustworthy or approved. Temporal standing matters where a later review changes the conditions for use.



\textbf{Boundary}



Approval does not establish applicability to a particular individual or replace the responsible clinical judgement.


{\small Recorded unresolved matters: source-qualified matters; exact scope remains in the canonical record.}\par
{\footnotesize \href{../../../../continuum/strategy/capability/ee04d00077d548a7a0a2dffd424c3d78.md}{Canonical record: ee04d00077d548a7a0a2dffd424c3d78}; full qualifications and provenance.}\par
\subsubsection*{II.1.2.7.2 [CT2] Health Information Management}
\phantomsection
\label{navigation:II.1.2.7.2}
\addcontentsline{toc}{subsubsection}{II.1.2.7.2 [CT2] Health Information Management}

\label{element:eb2984f0ca79420e83012f9e06a005fb}
\index{[CT2] Health Information Management (II.1.2.7.2)}
{\small Capability · Canonical · current}\par

Manage health information as governed, semantically meaningful content.



Health Information Management maintains health information as meaningful governed content rather than isolated persisted values. Its children distinguish the content itself from the structures that explain it, lifecycle and history, discovery aids, contextual groupings, relationships and retrieval. These distinctions let an architect ask whether information exists, what it means, which state applies and whether it may be obtained without conflating those questions.



\textbf{Intent}



Preserve usable health information whose meaning, relationships and temporal standing can be understood.



\textbf{Boundary}



Discoverability is not permission, compartment membership is not an access decision, and correlation is not identity or clinical interpretation. Current state and immutable historical operational truth remain distinct.


{\small Recorded unresolved matters: source-qualified matters; exact scope remains in the canonical record.}\par
{\footnotesize \href{../../../../continuum/strategy/capability/eb2984f0ca79420e83012f9e06a005fb.md}{Canonical record: eb2984f0ca79420e83012f9e06a005fb}; full qualifications and provenance.}\par
\paragraph*{II.1.2.7.2.1 [CT3] Health Information Catalogue and Index Management}
\phantomsection
\label{navigation:II.1.2.7.2.1}
\addcontentsline{toc}{paragraph}{II.1.2.7.2.1 [CT3] Health Information Catalogue and Index Management}

\label{element:4c834101307546d39a0ac36b07358faf}
\index{[CT3] Health Information Catalogue and Index Management (II.1.2.7.2.1)}
{\small Capability · Canonical · current}\par

Manage descriptive catalogues/indexes through which information existence and location can be established.



This capability maintains the descriptive aids through which information existence and location can be established. It makes information discoverable without equating the catalogue entry with the content it describes. Within information management it answers where relevant information may be found before the separate question of obtaining it under an authorised context.



\textbf{Boundary}



Discoverability does not authorise access, retrieve the content or communicate it to a recipient.


{\small Recorded unresolved matters: source-qualified matters; exact scope remains in the canonical record.}\par
{\footnotesize \href{../../../../continuum/strategy/capability/4c834101307546d39a0ac36b07358faf.md}{Canonical record: 4c834101307546d39a0ac36b07358faf}; full qualifications and provenance.}\par
\paragraph*{II.1.2.7.2.2 [CT3] Health Information Compartmentalisation}
\phantomsection
\label{navigation:II.1.2.7.2.2}
\addcontentsline{toc}{paragraph}{II.1.2.7.2.2 [CT3] Health Information Compartmentalisation}

\label{element:660d78d6a34041939d1c290266d47f03}
\index{[CT3] Health Information Compartmentalisation (II.1.2.7.2.2)}
{\small Capability · Canonical · current}\par

Manage contextual association, partitioning and bounded information groupings.



Health Information Compartmentalisation establishes contextual associations and bounded groupings of information. It makes the relationship of information to a business context explicit, allowing groupings to carry meaning beyond physical storage location. That context can inform governed use while remaining distinct from sensitivity classification and a particular access decision.



\textbf{Boundary}



Compartment membership neither establishes permission nor replaces classification or terminology governance.


{\small Recorded unresolved matters: source-qualified matters; exact scope remains in the canonical record.}\par
{\footnotesize \href{../../../../continuum/strategy/capability/660d78d6a34041939d1c290266d47f03.md}{Canonical record: 660d78d6a34041939d1c290266d47f03}; full qualifications and provenance.}\par
\paragraph*{II.1.2.7.2.3 [CT3] Health Information Data Correlation}
\phantomsection
\label{navigation:II.1.2.7.2.3}
\addcontentsline{toc}{paragraph}{II.1.2.7.2.3 [CT3] Health Information Data Correlation}

\label{element:6c41e90a8cf24ed689d9d6a363c8d3fe}
\index{[CT3] Health Information Data Correlation (II.1.2.7.2.3)}
{\small Capability · Canonical · current}\par

Establish meaningful relationships between independently managed information items/sets.



Health Information Data Correlation establishes meaningful relationships between independently managed information items or sets. It helps information be understood together without requiring a merger of their identities, histories or originating authority. The responsibility is the association itself, rather than a clinical conclusion inferred from the associated content.



\textbf{Boundary}



Correlation does not establish person identity, clinical interpretation or ownership transfer between the related items.


{\small Recorded unresolved matters: source-qualified matters; exact scope remains in the canonical record.}\par
{\footnotesize \href{../../../../continuum/strategy/capability/6c41e90a8cf24ed689d9d6a363c8d3fe.md}{Canonical record: 6c41e90a8cf24ed689d9d6a363c8d3fe}; full qualifications and provenance.}\par
\paragraph*{II.1.2.7.2.4 [CT3] Health Information Data Management}
\phantomsection
\label{navigation:II.1.2.7.2.4}
\addcontentsline{toc}{paragraph}{II.1.2.7.2.4 [CT3] Health Information Data Management}

\label{element:9565887a92164432933ea97e95af2368}
\index{[CT3] Health Information Data Management (II.1.2.7.2.4)}
{\small Capability · Canonical · current}\par

Manage actual health-information content within applicable governed structures.



Health Information Data Management concerns the actual content maintained within applicable governed structures. It keeps the represented information distinct from definitions describing how that information should be organised and interpreted. Within Health Information Management, its responsibility is meaningful content rather than mere successful persistence of values.



\textbf{Boundary}



Managing content does not originate its clinical facts or acquire the authority of its semantic references.


{\small Recorded unresolved matters: source-qualified matters; exact scope remains in the canonical record.}\par
{\footnotesize \href{../../../../continuum/strategy/capability/9565887a92164432933ea97e95af2368.md}{Canonical record: 9565887a92164432933ea97e95af2368}; full qualifications and provenance.}\par
\paragraph*{II.1.2.7.2.5 [CT3] Health Information Data Retrieval}
\phantomsection
\label{navigation:II.1.2.7.2.5}
\addcontentsline{toc}{paragraph}{II.1.2.7.2.5 [CT3] Health Information Data Retrieval}

\label{element:082760fc4804489a81eaed3e4fd9daa9}
\index{[CT3] Health Information Data Retrieval (II.1.2.7.2.5)}
{\small Capability · Canonical · current}\par

Obtain identified information under an applicable authorised request/context.



Health Information Data Retrieval obtains identified information under an applicable authorised request and context. It supplies the content whose existence may already be known through a catalogue or index. The distinction matters because locating an item, obtaining it and presenting it to a person each establish different outcomes.



\textbf{Boundary}



Retrieval does not itself grant permission, provide participant communication or confer authority over the retrieved facts.


{\small Recorded unresolved matters: source-qualified matters; exact scope remains in the canonical record.}\par
{\footnotesize \href{../../../../continuum/strategy/capability/082760fc4804489a81eaed3e4fd9daa9.md}{Canonical record: 082760fc4804489a81eaed3e4fd9daa9}; full qualifications and provenance.}\par
\paragraph*{II.1.2.7.2.6 [CT3] Health Information Lifecycle Management}
\phantomsection
\label{navigation:II.1.2.7.2.6}
\addcontentsline{toc}{paragraph}{II.1.2.7.2.6 [CT3] Health Information Lifecycle Management}

\label{element:874cd3b294184f928dbb7f656abd97f0}
\index{[CT3] Health Information Lifecycle Management (II.1.2.7.2.6)}
{\small Capability · Canonical · current}\par

Govern state, amendment, correction, supersession, currency, retention and disposition while preserving provenance and historical meaning.



Health Information Lifecycle Management governs how information's state and temporal standing change through amendment, correction, supersession, retention and disposition. It allows current knowledge to change without retrospectively rewriting what was previously received, known or managed. That distinction supports interpretation of both the present representation and its historical operational meaning.



\textbf{Boundary}



Current authoritative state differs from historical truth; preserving either does not confer originating authority for the represented facts.


{\small Recorded unresolved matters: source-qualified matters; exact scope remains in the canonical record.}\par
{\footnotesize \href{../../../../continuum/strategy/capability/874cd3b294184f928dbb7f656abd97f0.md}{Canonical record: 874cd3b294184f928dbb7f656abd97f0}; full qualifications and provenance.}\par
\paragraph*{II.1.2.7.2.7 [CT3] Health Information Structure Management}
\phantomsection
\label{navigation:II.1.2.7.2.7}
\addcontentsline{toc}{paragraph}{II.1.2.7.2.7 [CT3] Health Information Structure Management}

\label{element:c3b3e39dd141448e96b7dd7b4c4516f2}
\index{[CT3] Health Information Structure Management (II.1.2.7.2.7)}
{\small Capability · Canonical · current}\par

Manage structures, element/reference definitions, constraints and semantic bindings giving content unambiguous meaning.



This capability manages the structures, element definitions, references and bindings that make content interpretable. It addresses how a represented item can have an unambiguous meaning without treating the structure as the clinical fact itself. Semantic bindings identify applicable governed references while preserving their external authority.



\textbf{Boundary}



Structure management is not terminology governance and does not prescribe FHIR resources, physical schemas or storage technology.


{\small Recorded unresolved matters: source-qualified matters; exact scope remains in the canonical record.}\par
{\footnotesize \href{../../../../continuum/strategy/capability/c3b3e39dd141448e96b7dd7b4c4516f2.md}{Canonical record: c3b3e39dd141448e96b7dd7b4c4516f2}; full qualifications and provenance.}\par
\subsection*{II.1.2.8 [CT1] Health Products \& Clinical Technology}
\phantomsection
\label{navigation:II.1.2.8}
\addcontentsline{toc}{subsection}{II.1.2.8 [CT1] Health Products \& Clinical Technology}

{\small Historical predecessor: R.2.24 \hyperref[element:a631a8600b8b49df9ef452c48c91349a]{[CT1] Health Products \& Clinical Technology}; Canonical ID: \texttt{\seqsplit{a631a8600b8b49df9ef452c48c91349a}}. Superseded historical identity — direct correspondence · superseded.}\par
\label{element:54388e3202c2403089b5c476c78c4567}
\index{[CT1] Health Products \& Clinical Technology (II.1.2.8)}
{\small Capability · Canonical · current}\par

Manage health products and clinical technology in the care environment, with distinct device inventory, functional/release, deployment, utilisation, integration and lifecycle responsibilities. Product procurement and supply context is retained without inferring further Health Product Management decomposition.



This capability manages health products and clinical technology as business concerns in the care environment. Its established clinical-technology branch distinguishes knowledge of the device estate, functional state, deployment, user enablement, integration and lifecycle disposition. Product procurement and supply remain relevant context without a newly inferred Health Product Management decomposition.



\textbf{Intent}



Keep the technology context needed for safe service activity and information exchange coherent and governed.



\textbf{Boundary}



Device-specific enablement, clinical privilege and use in care are different concerns. The seven established device children do not resolve the separate maintenance-membership discrepancy.



\textbf{Harmonia Architectural Relevance}



Integration provision and device-data integration materially affect exchanged information and its care context. Supply-chain, maintenance and clinical-use authority are not transferred to Harmonia.


{\footnotesize \href{../../../../continuum/strategy/capability/54388e3202c2403089b5c476c78c4567.md}{Canonical record: 54388e3202c2403089b5c476c78c4567}; full qualifications and provenance.}\par
\subsubsection*{II.1.2.8.1 [CT2] Clinical Technology Management}
\phantomsection
\label{navigation:II.1.2.8.1}
\addcontentsline{toc}{subsubsection}{II.1.2.8.1 [CT2] Clinical Technology Management}

\label{element:6c91476a45eb4432b1f90b47a4c2ad35}
\index{[CT2] Clinical Technology Management (II.1.2.8.1)}
{\small Capability · Canonical · current}\par

Manage clinical technology as a coherent business concern rather than only the device information Harmonia happens to use.



Clinical Technology Management treats a device as more than a data endpoint. Its responsibilities distinguish what device exists, what functionality and release it has, where and to whom it is deployed, how authorised users are supported, whether exchange is provisioned, what data is integrated and whether it should remain in service. These distinctions preserve the business meaning of the device estate beyond the subset of information used by Harmonia.



\textbf{Intent}



Maintain intelligible device capability, deployment and lifecycle context for clinical service use.



\textbf{Boundary}



Provisioning integration does not demonstrate that data has flowed. Monitoring, restoring device condition and deciding lifecycle disposition remain distinct; maintenance's tier membership requires adjudication.


{\small Recorded unresolved matters: source-qualified matters; exact scope remains in the canonical record.}\par
{\footnotesize \href{../../../../continuum/strategy/capability/6c91476a45eb4432b1f90b47a4c2ad35.md}{Canonical record: 6c91476a45eb4432b1f90b47a4c2ad35}; full qualifications and provenance.}\par
\paragraph*{II.1.2.8.1.1 [CT3] Clinical Device Data Integration}
\phantomsection
\label{navigation:II.1.2.8.1.1}
\addcontentsline{toc}{paragraph}{II.1.2.8.1.1 [CT3] Clinical Device Data Integration}

\label{element:382a6fe3546f491f940cba1c1e5df5ef}
\index{[CT3] Clinical Device Data Integration (II.1.2.8.1.1)}
{\small Capability · Canonical · current}\par

Acquire, exchange or convey device data through established integration capability.



Clinical Device Data Integration acquires, exchanges or conveys device data using an established integration capability. It represents actual information participation rather than the prior configuration enabling it. The distinction makes it possible to preserve the device and care context of received data while retaining separate responsibility for its subsequent interpretation.



\textbf{Boundary}



Data movement does not confer clinical meaning, originating authority or device-maintenance responsibility.


{\small Recorded unresolved matters: source-qualified matters; exact scope remains in the canonical record.}\par
{\footnotesize \href{../../../../continuum/strategy/capability/382a6fe3546f491f940cba1c1e5df5ef.md}{Canonical record: 382a6fe3546f491f940cba1c1e5df5ef}; full qualifications and provenance.}\par
\paragraph*{II.1.2.8.1.2 [CT3] Clinical Device Feature \& Release Management}
\phantomsection
\label{navigation:II.1.2.8.1.2}
\addcontentsline{toc}{paragraph}{II.1.2.8.1.2 [CT3] Clinical Device Feature \& Release Management}

\label{element:23d87ead4cdc40b785b006ca80f6d7b1}
\index{[CT3] Clinical Device Feature \& Release Management (II.1.2.8.1.2)}
{\small Capability · Canonical · current}\par

Manage functional/release state, including applicable software, firmware, configuration and enabled device characteristics.



This capability governs the functional state of a device, including applicable software, firmware, configuration and enabled characteristics. Two instances of the same model may therefore have materially different capabilities in service. The responsibility makes those differences intelligible without reducing device governance to physical inventory or location.



\textbf{Boundary}



Release state is distinct from maintenance condition, deployment and professional authority to use the device.


{\small Recorded unresolved matters: source-qualified matters; exact scope remains in the canonical record.}\par
{\footnotesize \href{../../../../continuum/strategy/capability/23d87ead4cdc40b785b006ca80f6d7b1.md}{Canonical record: 23d87ead4cdc40b785b006ca80f6d7b1}; full qualifications and provenance.}\par
\paragraph*{II.1.2.8.1.3 [CT3] Clinical Device Integration Provision}
\phantomsection
\label{navigation:II.1.2.8.1.3}
\addcontentsline{toc}{paragraph}{II.1.2.8.1.3 [CT3] Clinical Device Integration Provision}

\label{element:7abab1d7d68f48cea565df752fad34c4}
\index{[CT3] Clinical Device Integration Provision (II.1.2.8.1.3)}
{\small Capability · Canonical · current}\par

Establish and maintain the device’s capability and configuration to participate in applicable information exchanges.



Clinical Device Integration Provision establishes the device capability and configuration required to participate in applicable information exchanges. Its outcome is an integration-ready arrangement, which must remain distinguishable from data actually acquired or exchanged. This allows an architect to discuss whether connectivity has been enabled without assuming that a particular clinical feed has operated successfully.



\textbf{Boundary}



Provisioning does not demonstrate data integration, clinical interpretation or downstream authority for the data.


{\small Recorded unresolved matters: source-qualified matters; exact scope remains in the canonical record.}\par
{\footnotesize \href{../../../../continuum/strategy/capability/7abab1d7d68f48cea565df752fad34c4.md}{Canonical record: 7abab1d7d68f48cea565df752fad34c4}; full qualifications and provenance.}\par
\paragraph*{II.1.2.8.1.4 [CT3] Clinical Device Inventory Management}
\phantomsection
\label{navigation:II.1.2.8.1.4}
\addcontentsline{toc}{paragraph}{II.1.2.8.1.4 [CT3] Clinical Device Inventory Management}

\label{element:f8205d4cf8144f34b2da1139ff45c868}
\index{[CT3] Clinical Device Inventory Management (II.1.2.8.1.4)}
{\small Capability · Canonical · current}\par

Establish what devices exist and their managed asset identity, specification, ownership/custody and inventory context.



Clinical Device Inventory Management establishes the managed identity and inventory context of the device estate. Specification, ownership and custody make it possible to distinguish which devices exist from their current release, deployment or use. That stable asset context supports subsequent technology-management responsibilities without claiming their different outcomes.



\textbf{Boundary}



Inventory presence does not establish readiness, functional release state or current assignment to a care requirement.


{\small Recorded unresolved matters: source-qualified matters; exact scope remains in the canonical record.}\par
{\footnotesize \href{../../../../continuum/strategy/capability/f8205d4cf8144f34b2da1139ff45c868.md}{Canonical record: f8205d4cf8144f34b2da1139ff45c868}; full qualifications and provenance.}\par
\paragraph*{II.1.2.8.1.5 [CT3] Clinical Device Lifecycle Management}
\phantomsection
\label{navigation:II.1.2.8.1.5}
\addcontentsline{toc}{paragraph}{II.1.2.8.1.5 [CT3] Clinical Device Lifecycle Management}

\label{element:186667f5a3844d10a236af73bc91a0de}
\index{[CT3] Clinical Device Lifecycle Management (II.1.2.8.1.5)}
{\small Capability · Canonical · current}\par

Govern introduction, continued service, review, replacement, withdrawal, decommissioning and retirement within the overall device estate.



Clinical Device Lifecycle Management governs whether devices enter, remain in or leave the managed estate. Introduction, review, replacement, withdrawal and retirement concern the device's overall service disposition. This gives technology management a longer-lived responsibility than the device's current loan, release or data-exchange session.



\textbf{Boundary}



Lifecycle disposition differs from monitoring and restoring device condition; the maintenance-membership discrepancy remains unresolved.


{\small Recorded unresolved matters: source-qualified matters; exact scope remains in the canonical record.}\par
{\footnotesize \href{../../../../continuum/strategy/capability/186667f5a3844d10a236af73bc91a0de.md}{Canonical record: 186667f5a3844d10a236af73bc91a0de}; full qualifications and provenance.}\par
\paragraph*{II.1.2.8.1.6 [CT3] Clinical Device Loan \& Deployment Management}
\phantomsection
\label{navigation:II.1.2.8.1.6}
\addcontentsline{toc}{paragraph}{II.1.2.8.1.6 [CT3] Clinical Device Loan \& Deployment Management}

\label{element:5ac802cf322d408f82c4a28cd24cb43b}
\index{[CT3] Clinical Device Loan \& Deployment Management (II.1.2.8.1.6)}
{\small Capability · Canonical · current}\par

Manage device loan, issue, assignment, placement, return and deployment.



Clinical Device Loan \& Deployment Management manages where devices are placed and the issue, assignment, return or loan context under which they are available. It connects an identifiable device with its operational deployment rather than assuming ownership or inventory presence explains its current use. This supports an intelligible account of the device estate as placements change.



\textbf{Boundary}



Deployment is not device-specific user enablement, clinical use or evidence of clinical privilege.


{\small Recorded unresolved matters: source-qualified matters; exact scope remains in the canonical record.}\par
{\footnotesize \href{../../../../continuum/strategy/capability/5ac802cf322d408f82c4a28cd24cb43b.md}{Canonical record: 5ac802cf322d408f82c4a28cd24cb43b}; full qualifications and provenance.}\par
\paragraph*{II.1.2.8.1.7 [CT3] Clinical Device Maintenance Management}
\phantomsection
\label{navigation:II.1.2.8.1.7}
\addcontentsline{toc}{paragraph}{II.1.2.8.1.7 [CT3] Clinical Device Maintenance Management}

\label{element:731ea36d36d8466d84196ccda2e61cf9}
\index{[CT3] Clinical Device Maintenance Management (II.1.2.8.1.7)}
{\small Capability · Candidate · current}\par

Definition unresolved in the baseline. No definition has been synthesised from the name.



\textbf{Proposed overview}



Maintenance/restoration is materially distinct from overall lifecycle disposition and monitoring; include the Candidate concern as a proposed eighth CT3.



\textbf{Proposed definition}



Maintain or restore acceptable device condition through testing, calibration, inspection, maintenance and repair.



\textbf{Proposed boundary}



The element remains Candidate. The definition/placement elaboration is Proposed and does not adjudicate the unresolved RAD question or transfer source-domain authority.



\medskip\hrule\medskip

{\small Recorded unresolved matters: source-qualified matters; exact scope remains in the canonical record.}\par
{\footnotesize \href{../../../../continuum/strategy/capability/731ea36d36d8466d84196ccda2e61cf9.md}{Canonical record: 731ea36d36d8466d84196ccda2e61cf9}; full qualifications and provenance.}\par
\paragraph*{II.1.2.8.1.8 [CT3] Clinical Device Utilisation Support}
\phantomsection
\label{navigation:II.1.2.8.1.8}
\addcontentsline{toc}{paragraph}{II.1.2.8.1.8 [CT3] Clinical Device Utilisation Support}

\label{element:0100cdb80a0c41c5930b1ca66248adf9}
\index{[CT3] Clinical Device Utilisation Support (II.1.2.8.1.8)}
{\small Capability · Canonical · current}\par

Enable and support otherwise appropriately authorised users to use the particular device safely and effectively.



Clinical Device Utilisation Support helps otherwise appropriately authorised users use a particular device safely and effectively. It represents the device-specific enablement needed beyond general professional qualifications. Within technology management it addresses the practical support for use, without making that support the clinical action performed with the device.



\textbf{Boundary}



Professional qualification or privilege alone does not establish device-specific enablement; enablement does not confer clinical authority.


{\small Recorded unresolved matters: source-qualified matters; exact scope remains in the canonical record.}\par
{\footnotesize \href{../../../../continuum/strategy/capability/0100cdb80a0c41c5930b1ca66248adf9.md}{Canonical record: 0100cdb80a0c41c5930b1ca66248adf9}; full qualifications and provenance.}\par
\subsection*{II.1.2.9 [CT1] Health Research \& Innovation}
\phantomsection
\label{navigation:II.1.2.9}
\addcontentsline{toc}{subsection}{II.1.2.9 [CT1] Health Research \& Innovation}

{\small Historical predecessor: R.2.25 \hyperref[element:2cc31384c8864c73989025cbe96dcfa7]{[CT1] Health Research \& Innovation}; Canonical ID: \texttt{\seqsplit{2cc31384c8864c73989025cbe96dcfa7}}. Superseded historical identity — direct correspondence · superseded.}\par
\label{element:582bf62145fd42118e5d57bb714841e8}
\index{[CT1] Health Research \& Innovation (II.1.2.9)}
{\small Capability · Canonical · current}\par

Provide existing health information for authorised research under applicable constraints and evaluate health innovation; deliberately under-decomposed for R1/R2.



This capability frames the deliberately bounded research and innovation responsibilities material to Harmonia's R1/R2 architecture. Research Information Support concerns provision of existing information for authorised research; Health Innovation concerns evaluation of proposed or emerging innovation. These establish useful enterprise context without implying a complete research organisation or innovation lifecycle.



\textbf{Intent}



Support legitimate research information needs and informed consideration of changes to health-service or clinical practice.



\textbf{Boundary}



Information provision does not own trials, ethics approval or specialist research analysis. Evaluating innovation does not authorise adoption or implementation.



\textbf{Harmonia Architectural Relevance}



Governed provision of existing information materially enables authorised research. Innovation evaluation has no separate platform assignment; the CT1 classification must not be inherited by each descendant.


{\footnotesize \href{../../../../continuum/strategy/capability/582bf62145fd42118e5d57bb714841e8.md}{Canonical record: 582bf62145fd42118e5d57bb714841e8}; full qualifications and provenance.}\par
\subsubsection*{II.1.2.9.1 [CT2] Health Innovation}
\phantomsection
\label{navigation:II.1.2.9.1}
\addcontentsline{toc}{subsubsection}{II.1.2.9.1 [CT2] Health Innovation}

\label{element:f263989e25c9469cac1bbad6491432f0}
\index{[CT2] Health Innovation (II.1.2.9.1)}
{\small Capability · Canonical · current}\par

Consider health-service and clinical innovation within the R1/R2 business context.



Health Innovation provides the business frame for considering changes in health-service and clinical practice. Its established evaluation child examines suitability, value and implications so that an innovation can be understood in the context in which it might be used. The model stops at that evaluative responsibility rather than presenting a complete pipeline from idea to operational adoption.



\textbf{Intent}



Make consideration of innovation informed by its prospective benefits, consequences and suitability.



\textbf{Boundary}



Evaluation does not establish an adoption decision, professional authority or an implementation programme.


{\small Recorded unresolved matters: source-qualified matters; exact scope remains in the canonical record.}\par
{\footnotesize \href{../../../../continuum/strategy/capability/f263989e25c9469cac1bbad6491432f0.md}{Canonical record: f263989e25c9469cac1bbad6491432f0}; full qualifications and provenance.}\par
\paragraph*{II.1.2.9.1.1 [CT3] Health Innovation Evaluation}
\phantomsection
\label{navigation:II.1.2.9.1.1}
\addcontentsline{toc}{paragraph}{II.1.2.9.1.1 [CT3] Health Innovation Evaluation}

\label{element:adc555bcb54e4de290869c7deebd128c}
\index{[CT3] Health Innovation Evaluation (II.1.2.9.1.1)}
{\small Capability · Canonical · current}\par

Evaluate proposed/emerging innovation for suitability, value, implications and potential adoption.



Health Innovation Evaluation is the organisation's ability to form a considered view of a proposed or emerging innovation before treating it as suitable for use. Within Health Innovation it brings suitability, value and implications into the same evaluative context, so that potential benefit can be considered alongside consequences for the health-service or clinical setting. The outcome supports informed consideration of adoption rather than making adoption inevitable.



\textbf{Intent}



Give decision makers an intelligible basis for considering whether an innovation merits further adoption consideration.



\textbf{Boundary}



Evaluation does not establish adoption, authorise a clinical change or create an implementation responsibility. The single established child does not imply a complete innovation lifecycle.


{\small Recorded unresolved matters: source-qualified matters; exact scope remains in the canonical record.}\par
{\footnotesize \href{../../../../continuum/strategy/capability/adc555bcb54e4de290869c7deebd128c.md}{Canonical record: adc555bcb54e4de290869c7deebd128c}; full qualifications and provenance.}\par
\subsubsection*{II.1.2.9.2 [CT2] Research Information Support}
\phantomsection
\label{navigation:II.1.2.9.2}
\addcontentsline{toc}{subsubsection}{II.1.2.9.2 [CT2] Research Information Support}

\label{element:ef18f999eef3488ba584d263a551732b}
\index{[CT2] Research Information Support (II.1.2.9.2)}
{\small Capability · Canonical · current}\par

Provide the existing health information needed for authorised research within the deliberately bounded R1/R2 responsibility.



Research Information Support represents the ability to make existing health information available for a legitimate research purpose. Its single established child concentrates on governed provision, preserving the constraints under which information can be used rather than treating technical availability as research permission. The narrow decomposition reflects R1/R2 sufficiency rather than a claim to encompass every form of research support.



\textbf{Intent}



Support authorised research with information supplied under the applicable ethical, privacy and governance conditions.



\textbf{Boundary}



Research governance, trial oversight and specialist analysis remain unmodelled and are not absorbed into provision.


{\small Recorded unresolved matters: source-qualified matters; exact scope remains in the canonical record.}\par
{\footnotesize \href{../../../../continuum/strategy/capability/ef18f999eef3488ba584d263a551732b.md}{Canonical record: ef18f999eef3488ba584d263a551732b}; full qualifications and provenance.}\par
\paragraph*{II.1.2.9.2.1 [CT3] Governed Research Information Provision}
\phantomsection
\label{navigation:II.1.2.9.2.1}
\addcontentsline{toc}{paragraph}{II.1.2.9.2.1 [CT3] Governed Research Information Provision}

\label{element:6907401e6786470083c1d34a65aa225b}
\index{[CT3] Governed Research Information Provision (II.1.2.9.2.1)}
{\small Capability · Canonical · current}\par

Provide existing health information for an authorised research purpose under applicable ethics, governance, privacy, consent, security and information-use constraints.



Governed Research Information Provision makes existing health information available for an authorised research purpose under its applicable constraints. It represents a substantive information contribution while keeping ethics, governance, consent, privacy, security and permitted use material to provision. The responsibility concerns supplying existing information legitimately, rather than running the research that uses it.



\textbf{Boundary}



Availability is not research authorisation. Trials, genomic analysis and research-governance responsibilities are not subsumed by this deliberately narrow capability.


{\small Recorded unresolved matters: source-qualified matters; exact scope remains in the canonical record.}\par
{\footnotesize \href{../../../../continuum/strategy/capability/6907401e6786470083c1d34a65aa225b.md}{Canonical record: 6907401e6786470083c1d34a65aa225b}; full qualifications and provenance.}\par
\subsection*{II.1.2.10 [CT1] Health Rights, Advocacy \& Participation}
\phantomsection
\label{navigation:II.1.2.10}
\addcontentsline{toc}{subsection}{II.1.2.10 [CT1] Health Rights, Advocacy \& Participation}

{\small Historical predecessor: R.2.26 \hyperref[element:58376e5c53bf4af0a54f99f21ac52666]{[CT1] Health Rights, Advocacy \& Participation}; Canonical ID: \texttt{\seqsplit{58376e5c53bf4af0a54f99f21ac52666}}. Superseded historical identity — direct correspondence · superseded.}\par
\label{element:7ab5c82a500b45f39267883a6e78c482}
\index{[CT1] Health Rights, Advocacy \& Participation (II.1.2.10)}
{\small Capability · Canonical · current}\par

Protecting patient rights, facilitating informed consent, managing advance care directives, supporting patient advocacy, and enabling consumer participation in health decisions.



This capability represents the organisation's responsibility to protect patients' rights and enable their voice in health decisions. Consent, advance care directives, advocacy and consumer participation concern the legitimacy and personal consequences of care, rather than merely the availability of clinical information. The governed model retains this broad enterprise responsibility without an established subordinate capability set.



\textbf{Intent}



Support informed participation and respect for the individual's rights in decisions affecting their health.



\textbf{Boundary}



Recording a directive or restricting access can support rights without performing advocacy or managing consumer participation.



\textbf{Harmonia Architectural Relevance}



The established Adjacent classification reflects the use and supply of consent and record-access evidence. It does not assign the underlying rights, advocacy or participation activity to Harmonia.


{\footnotesize \href{../../../../continuum/strategy/capability/7ab5c82a500b45f39267883a6e78c482.md}{Canonical record: 7ab5c82a500b45f39267883a6e78c482}; full qualifications and provenance.}\par
\subsection*{II.1.2.11 [CT1] Health Service Assurance}
\phantomsection
\label{navigation:II.1.2.11}
\addcontentsline{toc}{subsection}{II.1.2.11 [CT1] Health Service Assurance}

{\small Historical predecessor: R.2.27 \hyperref[element:6a82dc1423374e809dbf542eeaa69482]{[CT1] Health Service Assurance}; Canonical ID: \texttt{\seqsplit{6a82dc1423374e809dbf542eeaa69482}}. Superseded historical identity — direct correspondence · superseded.}\par
\label{element:14860102964c441681d41ec3da423fe0}
\index{[CT1] Health Service Assurance (II.1.2.11)}
{\small Capability · Canonical · current}\par

Establish and maintain justified confidence that health services can be, and are being, delivered in accordance with applicable expectations, obligations, constraints and required levels of performance.



Health Service Assurance establishes justified confidence in the ability to deliver health services appropriately. Specification defines how confidence will be established; evaluation determines what the available evidence supports. Its subject is macro-level service delivery, with workforce, clinical, information, security and other concerns contributing evidence or criteria without becoming separate assurance trees.



Where independent governed assurance is required, the subject cannot control its own evaluation or conclusion. Independence concerns control of the assurance judgement rather than requiring a separate organisation or technology installation.



\textbf{Intent}



Distinguish justified confidence, qualified conclusions and insufficient evidence so that assurance communicates what is actually supported.



\textbf{Boundary}



Assurance manages its own evaluation and does not assign, remediate or progress the subject activity. Harmonia does not thereby assume responsibility for clinical adequacy or professional practice.



\textbf{Harmonia Architectural Relevance}



Approved independent governed assurance materially supports service confidence. The three established assurance BECs provide bounded support; they are not remapped to the newer BC descendants and are not assigned implementation components.


{\footnotesize \href{../../../../continuum/strategy/capability/14860102964c441681d41ec3da423fe0.md}{Canonical record: 14860102964c441681d41ec3da423fe0}; full qualifications and provenance.}\par
\subsubsection*{II.1.2.11.1 [CT2] Assurance Evaluation}
\phantomsection
\label{navigation:II.1.2.11.1}
\addcontentsline{toc}{subsubsection}{II.1.2.11.1 [CT2] Assurance Evaluation}

\label{element:d1b20e0fb91e4bd1950ad57b84e2cafd}
\index{[CT2] Assurance Evaluation (II.1.2.11.1)}
{\small Capability · Canonical · current}\par

Determine what justified confidence the applicable evidence supports.



Assurance Evaluation determines what confidence the evidence legitimately supports under the established approach and criteria. Evidence Assessment considers the material against that basis; Conclusion Determination establishes the resulting position, including qualifications and unresolved concerns. The distinction prevents accumulated evidence or execution completion from being mistaken for a justified assurance conclusion.



\textbf{Intent}



State the supported level and limits of confidence without concealing insufficient evidence.



\textbf{Boundary}



Insufficient evidence establishes neither assurance nor non-assurance. Operational response remains with the activity's responsible capability, and the existing assurance BEC ownership is unchanged.


{\small Recorded unresolved matters: source-qualified matters; exact scope remains in the canonical record.}\par
{\footnotesize \href{../../../../continuum/strategy/capability/d1b20e0fb91e4bd1950ad57b84e2cafd.md}{Canonical record: d1b20e0fb91e4bd1950ad57b84e2cafd}; full qualifications and provenance.}\par
\paragraph*{II.1.2.11.1.1 [CT3] Assurance Conclusion Determination}
\phantomsection
\label{navigation:II.1.2.11.1.1}
\addcontentsline{toc}{paragraph}{II.1.2.11.1.1 [CT3] Assurance Conclusion Determination}

\label{element:10f538feae4042829fafeae472882128}
\index{[CT3] Assurance Conclusion Determination (II.1.2.11.1.1)}
{\small Capability · Canonical · current}\par

Establish the legitimately supported conclusion, including insufficiency, qualification and unresolved concerns where applicable.



Assurance Conclusion Determination establishes what confidence is legitimately supported by the assessment. It preserves qualifications and unresolved concerns rather than forcing incomplete evidence into a positive or negative result. Within Assurance Evaluation it is the responsibility that gives the conclusion its standing, distinct from gathering or assessing the supporting material.



\textbf{Boundary}



Insufficiency proves neither assurance nor non-assurance, and the conclusion does not manage the subject's operational response.


{\small Recorded unresolved matters: source-qualified matters; exact scope remains in the canonical record.}\par
{\footnotesize \href{../../../../continuum/strategy/capability/10f538feae4042829fafeae472882128.md}{Canonical record: 10f538feae4042829fafeae472882128}; full qualifications and provenance.}\par
\paragraph*{II.1.2.11.1.2 [CT3] Assurance Evidence Assessment}
\phantomsection
\label{navigation:II.1.2.11.1.2}
\addcontentsline{toc}{paragraph}{II.1.2.11.1.2 [CT3] Assurance Evidence Assessment}

\label{element:1b295763b37f4622bd1bb80c6ea3741c}
\index{[CT3] Assurance Evidence Assessment (II.1.2.11.1.2)}
{\small Capability · Canonical · current}\par

Evaluate available evidence against the established assurance approach and criteria.



Assurance Evidence Assessment evaluates available evidence against the established approach and criteria. It produces the assessed basis from which a legitimate conclusion may be determined, retaining qualifications and evidentiary limits. The responsibility is evaluative rather than the operational production of the source information or remediation of what it reveals.



\textbf{Boundary}



Assessment is distinct from conclusion determination; unavailable or insufficient evidence must remain explicit.


{\small Recorded unresolved matters: source-qualified matters; exact scope remains in the canonical record.}\par
{\footnotesize \href{../../../../continuum/strategy/capability/1b295763b37f4622bd1bb80c6ea3741c.md}{Canonical record: 1b295763b37f4622bd1bb80c6ea3741c}; full qualifications and provenance.}\par
\subsubsection*{II.1.2.11.2 [CT2] Health Service Assurance Specification}
\phantomsection
\label{navigation:II.1.2.11.2}
\addcontentsline{toc}{subsubsection}{II.1.2.11.2 [CT2] Health Service Assurance Specification}

\label{element:ec5c4d1219464433b2cee5c5a988fef7}
\index{[CT2] Health Service Assurance Specification (II.1.2.11.2)}
{\small Capability · Canonical · current}\par

Specify how macro-level service-delivery assurance will be undertaken and what sufficiently supports confidence.



This capability defines the approach and evaluation criteria through which service-delivery assurance can be justified. Approach Definition addresses scope, method, responsibility, evidence and applicable independence; Criteria Governance establishes the basis for judging sufficiency. Their separation distinguishes the design of an assurance activity from the normative or evaluative material it applies.



\textbf{Intent}



Make the intended basis of confidence explicit before a conclusion is treated as assurance.



\textbf{Boundary}



An approach cannot waive a higher-level independence requirement. Contributing specialist concerns do not create separate Business Assurance hierarchies.


{\small Recorded unresolved matters: source-qualified matters; exact scope remains in the canonical record.}\par
{\footnotesize \href{../../../../continuum/strategy/capability/ec5c4d1219464433b2cee5c5a988fef7.md}{Canonical record: ec5c4d1219464433b2cee5c5a988fef7}; full qualifications and provenance.}\par
\paragraph*{II.1.2.11.2.1 [CT3] Assurance Approach Definition}
\phantomsection
\label{navigation:II.1.2.11.2.1}
\addcontentsline{toc}{paragraph}{II.1.2.11.2.1 [CT3] Assurance Approach Definition}

\label{element:d672f8831df54f2099927f820d0025f5}
\index{[CT3] Assurance Approach Definition (II.1.2.11.2.1)}
{\small Capability · Canonical · current}\par

Define scope, evidence expectations, responsibilities, methods, applicable degree of independence and timing/frequency for service-delivery assurance.



Assurance Approach Definition establishes how service-delivery assurance will be undertaken, including scope, evidence, methods, responsibilities, timing and required independence. It gives assurance activity a defensible approach rather than assuming any available observation can support confidence. The approach must respect applicable higher-level requirements when deciding what constitutes sufficient evaluation.



\textbf{Boundary}



Self-attestation cannot replace required independent governed assurance, and an approach does not itself establish a conclusion.


{\small Recorded unresolved matters: source-qualified matters; exact scope remains in the canonical record.}\par
{\footnotesize \href{../../../../continuum/strategy/capability/d672f8831df54f2099927f820d0025f5.md}{Canonical record: d672f8831df54f2099927f820d0025f5}; full qualifications and provenance.}\par
\paragraph*{II.1.2.11.2.2 [CT3] Assurance Evaluation Criteria Governance}
\phantomsection
\label{navigation:II.1.2.11.2.2}
\addcontentsline{toc}{paragraph}{II.1.2.11.2.2 [CT3] Assurance Evaluation Criteria Governance}

\label{element:f3e32f2a2c9d47a89eaf124647161522}
\index{[CT3] Assurance Evaluation Criteria Governance (II.1.2.11.2.2)}
{\small Capability · Canonical · current}\par

Establish and maintain criteria against which sufficient assurance can be determined.



This capability establishes and maintains the criteria against which assurance sufficiency is judged. It makes the evaluative basis of confidence explicit across the different concerns that can affect service delivery. Criteria can draw on clinical, workforce, technology, information or other contexts without turning those subjects into independent assurance capability trees.



\textbf{Boundary}



Criteria governance is distinct from applying criteria to evidence and does not remap the existing assurance BEC responsibilities.


{\small Recorded unresolved matters: source-qualified matters; exact scope remains in the canonical record.}\par
{\footnotesize \href{../../../../continuum/strategy/capability/f3e32f2a2c9d47a89eaf124647161522.md}{Canonical record: f3e32f2a2c9d47a89eaf124647161522}; full qualifications and provenance.}\par
\subsection*{II.1.2.12 [CT1] Health Service Direction \& Stewardship}
\phantomsection
\label{navigation:II.1.2.12}
\addcontentsline{toc}{subsection}{II.1.2.12 [CT1] Health Service Direction \& Stewardship}

{\small Historical predecessor: R.2.28 \hyperref[element:360bf50f92484448a01cc57f2dbce23c]{[CT1] Health Service Direction \& Stewardship}; Canonical ID: \texttt{\seqsplit{360bf50f92484448a01cc57f2dbce23c}}. Superseded historical identity — direct correspondence · superseded.}\par
\label{element:d51e7125bf9c4d2c9c59b60346b4d86f}
\index{[CT1] Health Service Direction \& Stewardship (II.1.2.12)}
{\small Capability · Canonical · current}\par

Model needed health services and their viability requirements, support establishment/improvement and align services with the wider ecosystem of dependencies, resources, financial arrangements and organisational change.



This capability considers what health services should exist, what makes them viable and how they fit within the wider service ecosystem. Modelling frames the need and required abilities; alignment addresses dependencies, resources, financial arrangements and surrounding change. The distinction keeps the rationale for a service separate from its identifiable offering and its current operational delivery state.



\textbf{Intent}



Support viable, coherent service development and improvement within organisational and ecosystem constraints.



\textbf{Boundary}



Identifying dependencies and resources does not acquire their operating authority. Service modelling and alignment do not establish the authoritative service catalogue.



\textbf{Harmonia Architectural Relevance}



Service need, viability and alignment facts materially shape service information and derivation requirements. That relevance does not assign executive authority or establish a future BEC mapping.


{\footnotesize \href{../../../../continuum/strategy/capability/d51e7125bf9c4d2c9c59b60346b4d86f.md}{Canonical record: d51e7125bf9c4d2c9c59b60346b4d86f}; full qualifications and provenance.}\par
\subsubsection*{II.1.2.12.1 [CT2] Health Service Alignment}
\phantomsection
\label{navigation:II.1.2.12.1}
\addcontentsline{toc}{subsubsection}{II.1.2.12.1 [CT2] Health Service Alignment}

\label{element:d4fc5367352741a99bd334f7e8bf5856}
\index{[CT2] Health Service Alignment (II.1.2.12.1)}
{\small Capability · Canonical · current}\par

Position the service coherently in its wider health-service ecosystem.



Health Service Alignment positions a service among the other services, resources and arrangements on which it depends. It considers reciprocal dependencies, required resources, financial fit and the organisational changes needed for introduction or material change. Where modelling asks what service is needed, alignment asks how it can fit coherently into its environment.



\textbf{Intent}



Expose the dependencies and surrounding commitments that affect service viability.



\textbf{Boundary}



Alignment does not establish the service offering, own the required resources or settle the unresolved ancestry of collective operational readiness.


{\small Recorded unresolved matters: source-qualified matters; exact scope remains in the canonical record.}\par
{\footnotesize \href{../../../../continuum/strategy/capability/d4fc5367352741a99bd334f7e8bf5856.md}{Canonical record: d4fc5367352741a99bd334f7e8bf5856}; full qualifications and provenance.}\par
\paragraph*{II.1.2.12.1.1 [CT3] Health Service Co-Dependency Management}
\phantomsection
\label{navigation:II.1.2.12.1.1}
\addcontentsline{toc}{paragraph}{II.1.2.12.1.1 [CT3] Health Service Co-Dependency Management}

\label{element:f031ee3ddc5e47c48d4308d4077b3967}
\index{[CT3] Health Service Co-Dependency Management (II.1.2.12.1.1)}
{\small Capability · Canonical · current}\par

Identify and manage what the service needs from others and what others need from it.



Health Service Co-Dependency Management identifies and manages reciprocal needs between a service and its wider ecosystem. It makes explicit both what the service requires from others and what others rely on it to provide. This supports alignment by exposing dependencies that a service description considered in isolation may conceal.



\textbf{Boundary}



Recognising a dependency does not acquire the provider's responsibility or establish a new technical Service dependency.


{\small Recorded unresolved matters: source-qualified matters; exact scope remains in the canonical record.}\par
{\footnotesize \href{../../../../continuum/strategy/capability/f031ee3ddc5e47c48d4308d4077b3967.md}{Canonical record: f031ee3ddc5e47c48d4308d4077b3967}; full qualifications and provenance.}\par
\paragraph*{II.1.2.12.1.2 [CT3] Health Service Financial Management}
\phantomsection
\label{navigation:II.1.2.12.1.2}
\addcontentsline{toc}{paragraph}{II.1.2.12.1.2 [CT3] Health Service Financial Management}

\label{element:2fbf31e974574f288f087179e4023dc2}
\index{[CT3] Health Service Financial Management (II.1.2.12.1.2)}
{\small Capability · Canonical · current}\par

Establish the service’s fit within funding, financial and charging arrangements without absorbing enterprise finance/accounting.



Health Service Financial Management establishes how the service fits applicable funding, financial and charging arrangements. It connects service viability with its surrounding financial context without attempting to become the enterprise finance function. The responsibility is service-specific alignment rather than the production of all financial transactions associated with care.



\textbf{Boundary}



Financial fit does not confer pricing, accounting, billing or enterprise revenue-management authority.


{\small Recorded unresolved matters: source-qualified matters; exact scope remains in the canonical record.}\par
{\footnotesize \href{../../../../continuum/strategy/capability/2fbf31e974574f288f087179e4023dc2.md}{Canonical record: 2fbf31e974574f288f087179e4023dc2}; full qualifications and provenance.}\par
\paragraph*{II.1.2.12.1.3 [CT3] Health Service Implementation Planning}
\phantomsection
\label{navigation:II.1.2.12.1.3}
\addcontentsline{toc}{paragraph}{II.1.2.12.1.3 [CT3] Health Service Implementation Planning}

\label{element:67bd795d91ce4dd49fdd481a3c451366}
\index{[CT3] Health Service Implementation Planning (II.1.2.12.1.3)}
{\small Capability · Canonical · current}\par

Establish and coordinate surrounding organisational/service changes needed to introduce or materially change the service.



Health Service Implementation Planning establishes and coordinates the organisational and surrounding service changes needed for introduction or material change. It considers how the service fits its environment, complementing modelling's establishment support with an alignment view of interdependent change. Its purpose is coherence of the planned transition rather than ownership of every change activity.



\textbf{Boundary}



Planning does not perform all supporting changes or define the authoritative service offering.


{\small Recorded unresolved matters: source-qualified matters; exact scope remains in the canonical record.}\par
{\footnotesize \href{../../../../continuum/strategy/capability/67bd795d91ce4dd49fdd481a3c451366.md}{Canonical record: 67bd795d91ce4dd49fdd481a3c451366}; full qualifications and provenance.}\par
\paragraph*{II.1.2.12.1.4 [CT3] Health Service Resource Identification}
\phantomsection
\label{navigation:II.1.2.12.1.4}
\addcontentsline{toc}{paragraph}{II.1.2.12.1.4 [CT3] Health Service Resource Identification}

\label{element:7d46575486c847cd8fdbb665b3ed1090}
\index{[CT3] Health Service Resource Identification (II.1.2.12.1.4)}
{\small Capability · Canonical · current}\par

Identify required workforce, facilities, equipment, clinical technology, ICT, administration and support resources without owning them.



This capability establishes the resources a service requires across workforce, facilities, equipment, clinical technology, ICT and support. It makes resource needs visible in the service's alignment context without treating identification as procurement, ownership or allocation. Required resources can therefore be discussed distinctly from resources presently available or deployed.



\textbf{Boundary}



Identifying resources does not operate them, assign them to care or establish collective service readiness.


{\small Recorded unresolved matters: source-qualified matters; exact scope remains in the canonical record.}\par
{\footnotesize \href{../../../../continuum/strategy/capability/7d46575486c847cd8fdbb665b3ed1090.md}{Canonical record: 7d46575486c847cd8fdbb665b3ed1090}; full qualifications and provenance.}\par
\subsubsection*{II.1.2.12.2 [CT2] Health Service Modelling}
\phantomsection
\label{navigation:II.1.2.12.2}
\addcontentsline{toc}{subsubsection}{II.1.2.12.2 [CT2] Health Service Modelling}

\label{element:9a3e90db97b94d1888a2a5c5effb8999}
\index{[CT2] Health Service Modelling (II.1.2.12.2)}
{\small Capability · Canonical · current}\par

Establish what service should exist and what it needs to become viable.



Health Service Modelling establishes the service needed and the ordinary business abilities required to make it viable. Its children distinguish recognising unmet need, identifying required abilities, supporting establishment and evaluating an existing service for improvement. The model can therefore revisit service assumptions without treating the original definition as permanently sufficient.



\textbf{Intent}



Connect a service's purpose with the abilities and preparation required for it to exist usefully.



\textbf{Boundary}



Capability Identification uses ordinary business meaning; it does not allocate Harmonia capability tiers. Implementation support does not take over recruitment, facilities, procurement or ICT responsibilities.


{\small Recorded unresolved matters: source-qualified matters; exact scope remains in the canonical record.}\par
{\footnotesize \href{../../../../continuum/strategy/capability/9a3e90db97b94d1888a2a5c5effb8999.md}{Canonical record: 9a3e90db97b94d1888a2a5c5effb8999}; full qualifications and provenance.}\par
\paragraph*{II.1.2.12.2.1 [CT3] Health Service Capability Identification}
\phantomsection
\label{navigation:II.1.2.12.2.1}
\addcontentsline{toc}{paragraph}{II.1.2.12.2.1 [CT3] Health Service Capability Identification}

\label{element:53f64a402905438aa87f9d16eff876da}
\index{[CT3] Health Service Capability Identification (II.1.2.12.2.1)}
{\small Capability · Canonical · current}\par

Identify the ordinary business abilities needed to satisfy the service need across clinical, workforce, facility, technology, information and support concerns.



This capability identifies the ordinary business abilities a proposed service needs across clinical, workforce, facility, information, technology and support concerns. It broadens service modelling beyond a clinical description to the conditions needed for viability. Its contribution is an understanding of required abilities, leaving their fulfilment and governing authority with the appropriate responsibilities.



\textbf{Boundary}



The word capability here does not allocate Harmonia CT tiers or establish new catalogue elements.


{\small Recorded unresolved matters: source-qualified matters; exact scope remains in the canonical record.}\par
{\footnotesize \href{../../../../continuum/strategy/capability/53f64a402905438aa87f9d16eff876da.md}{Canonical record: 53f64a402905438aa87f9d16eff876da}; full qualifications and provenance.}\par
\paragraph*{II.1.2.12.2.2 [CT3] Health Service Continuous Improvement}
\phantomsection
\label{navigation:II.1.2.12.2.2}
\addcontentsline{toc}{paragraph}{II.1.2.12.2.2 [CT3] Health Service Continuous Improvement}

\label{element:b2b2b91b1a094ef1b13ce6870ff9716f}
\index{[CT3] Health Service Continuous Improvement (II.1.2.12.2.2)}
{\small Capability · Canonical · current}\par

Evaluate the established service and identify correction/improvement needs that may feed back into modelling and implementation support.



This capability evaluates an established service and identifies correction or improvement needs across its service-wide purpose. It allows experience of actual operation to inform reconsideration of the service model and implementation support. The scope is broader than reporting a clinical-quality finding, while remaining anchored to the service being improved.



\textbf{Boundary}



It does not establish generic Continuous Improvement placement or transfer the authority of specialist activities affected by improvement.


{\small Recorded unresolved matters: source-qualified matters; exact scope remains in the canonical record.}\par
{\footnotesize \href{../../../../continuum/strategy/capability/b2b2b91b1a094ef1b13ce6870ff9716f.md}{Canonical record: b2b2b91b1a094ef1b13ce6870ff9716f}; full qualifications and provenance.}\par
\paragraph*{II.1.2.12.2.3 [CT3] Health Service Implementation Support}
\phantomsection
\label{navigation:II.1.2.12.2.3}
\addcontentsline{toc}{paragraph}{II.1.2.12.2.3 [CT3] Health Service Implementation Support}

\label{element:f74f5944573b408d94fe8c88e33e128a}
\index{[CT3] Health Service Implementation Support (II.1.2.12.2.3)}
{\small Capability · Canonical · current}\par

Support establishment of the framed service through applicable organisational preparation and programmes of work.



Health Service Implementation Support assists establishment of the framed service through organisational preparation and programmes of work. It connects the service model to practical readiness for introduction while recognising that recruitment, training, equipment and technology changes retain their own owners. Support therefore contributes to establishment without becoming an all-purpose implementation authority.



\textbf{Boundary}



It does not absorb HR, facilities, procurement or ICT responsibilities or establish operational-readiness ancestry.


{\small Recorded unresolved matters: source-qualified matters; exact scope remains in the canonical record.}\par
{\footnotesize \href{../../../../continuum/strategy/capability/f74f5944573b408d94fe8c88e33e128a.md}{Canonical record: f74f5944573b408d94fe8c88e33e128a}; full qualifications and provenance.}\par
\paragraph*{II.1.2.12.2.4 [CT3] Health Service Need Identification}
\phantomsection
\label{navigation:II.1.2.12.2.4}
\addcontentsline{toc}{paragraph}{II.1.2.12.2.4 [CT3] Health Service Need Identification}

\label{element:976c90e05d0e44b1b24c89bf61eac005}
\index{[CT3] Health Service Need Identification (II.1.2.12.2.4)}
{\small Capability · Canonical · current}\par

Establish an unmet health-service need requiring response.



Health Service Need Identification establishes the unmet service need for which an organisational response is being considered. It gives service modelling a reason grounded in what is missing rather than beginning with a predetermined solution. That identified need can then frame the abilities and resources required without assuming that a new offering has already been established.



\textbf{Boundary}



Identifying need does not define an offering, approve funding or determine the surrounding implementation arrangements.


{\small Recorded unresolved matters: source-qualified matters; exact scope remains in the canonical record.}\par
{\footnotesize \href{../../../../continuum/strategy/capability/976c90e05d0e44b1b24c89bf61eac005.md}{Canonical record: 976c90e05d0e44b1b24c89bf61eac005}; full qualifications and provenance.}\par
\subsection*{II.1.2.13 [CT1] Health Service Governance}
\phantomsection
\label{navigation:II.1.2.13}
\addcontentsline{toc}{subsection}{II.1.2.13 [CT1] Health Service Governance}

{\small Historical predecessor: R.2.29 \hyperref[element:24f3dd30223142efa35a5826552d53ab]{[CT1] Health Service Governance}; Canonical ID: \texttt{\seqsplit{24f3dd30223142efa35a5826552d53ab}}. Superseded historical identity — direct correspondence · superseded.}\par
\label{element:209aada781e844e28202a2535a264c21}
\index{[CT1] Health Service Governance (II.1.2.13)}
{\small Capability · Canonical · current}\par

Establish and maintain the obligations, authorities, constraints, policies and governance framework under which health-service activity must operate, and define how adherence will be assessed, checked and reported.



Health Service Governance establishes the normative envelope within which health-service activity operates: what obligations apply, who may decide, what constraints bind decisions and how adherence will be assessed and reported. Its framework and compliance branches distinguish the authoritative requirements from the arrangements used to judge satisfaction. Governance therefore shapes legitimate activity without becoming the manager of every governed operation.



\textbf{Intent}



Make obligations, decision authority and compliance expectations explicit and traceable.



\textbf{Boundary}



Operational management chooses and progresses means within the governing envelope; assurance evaluates evidence. Neither responsibility is absorbed merely because governance defines its constraints.



\textbf{Harmonia Architectural Relevance}



Formal relevance remains Unresolved. The requirements plainly constrain managed activity, but the RAD has not assigned enough platform participation or enabling derivation to adjudicate its relevance classification.


{\small Recorded unresolved matters: source-qualified matters; exact scope remains in the canonical record.}\par
{\footnotesize \href{../../../../continuum/strategy/capability/209aada781e844e28202a2535a264c21.md}{Canonical record: 209aada781e844e28202a2535a264c21}; full qualifications and provenance.}\par
\subsubsection*{II.1.2.13.1 [CT2] Compliance Governance}
\phantomsection
\label{navigation:II.1.2.13.1}
\addcontentsline{toc}{subsubsection}{II.1.2.13.1 [CT2] Compliance Governance}

\label{element:02a52157967f4b498807cac70e3410e9}
\index{[CT2] Compliance Governance (II.1.2.13.1)}
{\small Capability · Canonical · current}\par

Define how adherence to the authoritative WHAT will be determined, checked and reported.



Compliance Governance defines how adherence to established obligations, authorities and constraints will be determined. Assessment Definition specifies the evidence and basis for evaluation; Reporting \& Checkpoint Definition specifies where adherence is checked and what must be communicated. This makes the compliance test explicit without treating the design of a test as its successful execution.



\textbf{Intent}



Provide a defensible basis for checking and reporting whether the governing envelope has been respected.



\textbf{Boundary}



Defining assessment arrangements does not perform assurance or manage the governed activity. The obligations being assessed retain their authoritative source.


{\small Recorded unresolved matters: source-qualified matters; exact scope remains in the canonical record.}\par
{\footnotesize \href{../../../../continuum/strategy/capability/02a52157967f4b498807cac70e3410e9.md}{Canonical record: 02a52157967f4b498807cac70e3410e9}; full qualifications and provenance.}\par
\paragraph*{II.1.2.13.1.1 [CT3] Compliance Assessment Definition}
\phantomsection
\label{navigation:II.1.2.13.1.1}
\addcontentsline{toc}{paragraph}{II.1.2.13.1.1 [CT3] Compliance Assessment Definition}

\label{element:8a6db4bd0ed845dda125896c4e3ad3d0}
\index{[CT3] Compliance Assessment Definition (II.1.2.13.1.1)}
{\small Capability · Canonical · current}\par

Define the evidence, measures, criteria, controls and assessment arrangements sufficient to evaluate compliance with established obligations, authorities and constraints.



Compliance Assessment Definition specifies the evidence, measures, criteria and arrangements sufficient to evaluate adherence. It makes the assessment basis explicit after obligations, authorities and constraints have been established. This preserves the distinction between defining what compliance means and determining whether the available evidence demonstrates it in a particular case.



\textbf{Boundary}



Assessment definition does not replace the authoritative obligation, perform assurance or establish a compliance finding.


{\small Recorded unresolved matters: source-qualified matters; exact scope remains in the canonical record.}\par
{\footnotesize \href{../../../../continuum/strategy/capability/8a6db4bd0ed845dda125896c4e3ad3d0.md}{Canonical record: 8a6db4bd0ed845dda125896c4e3ad3d0}; full qualifications and provenance.}\par
\paragraph*{II.1.2.13.1.2 [CT3] Governance Reporting \& Checkpoint Definition}
\phantomsection
\label{navigation:II.1.2.13.1.2}
\addcontentsline{toc}{paragraph}{II.1.2.13.1.2 [CT3] Governance Reporting \& Checkpoint Definition}

\label{element:0ff9f332acdc435e98f4f527c74f60c9}
\index{[CT3] Governance Reporting \& Checkpoint Definition (II.1.2.13.1.2)}
{\small Capability · Canonical · current}\par

Establish where/how adherence is checked and what governance/compliance information must be reported.



This capability specifies where adherence must be checked and what governance or compliance information must be reported. It gives the governance framework observable expectations without treating reporting requirements as completed reporting. Within Compliance Governance it complements the assessment basis by making checking and communication obligations explicit.



\textbf{Boundary}



Defining a checkpoint does not perform the check or take over progression of the governed operation.


{\small Recorded unresolved matters: source-qualified matters; exact scope remains in the canonical record.}\par
{\footnotesize \href{../../../../continuum/strategy/capability/0ff9f332acdc435e98f4f527c74f60c9.md}{Canonical record: 0ff9f332acdc435e98f4f527c74f60c9}; full qualifications and provenance.}\par
\subsubsection*{II.1.2.13.2 [CT2] Governance Framework Management}
\phantomsection
\label{navigation:II.1.2.13.2}
\addcontentsline{toc}{subsubsection}{II.1.2.13.2 [CT2] Governance Framework Management}

\label{element:b39f53b3ffb441f388a89b45cc517deb}
\index{[CT2] Governance Framework Management (II.1.2.13.2)}
{\small Capability · Canonical · current}\par

Establish the authoritative WHAT and legitimate working parameters.



Governance Framework Management establishes what must be satisfied and the legitimate authority and bounds for organisational decisions. Its two children separate obligations attributable to source authority from the allocation and limitation of decision authority. Together they supply the authoritative reference against which operational arrangements and compliance assessments can be considered.



\textbf{Intent}



Keep required outcomes and legitimate decision parameters distinguishable and traceable to their authority.



\textbf{Boundary}



Interpreting an external obligation or engaging with a regulator does not transfer that body's authority. Operational means remain with management.


{\small Recorded unresolved matters: source-qualified matters; exact scope remains in the canonical record.}\par
{\footnotesize \href{../../../../continuum/strategy/capability/b39f53b3ffb441f388a89b45cc517deb.md}{Canonical record: b39f53b3ffb441f388a89b45cc517deb}; full qualifications and provenance.}\par
\paragraph*{II.1.2.13.2.1 [CT3] Authority \& Constraint Governance}
\phantomsection
\label{navigation:II.1.2.13.2.1}
\addcontentsline{toc}{paragraph}{II.1.2.13.2.1 [CT3] Authority \& Constraint Governance}

\label{element:ad7f41dc48a046b185f044f0a24517ca}
\index{[CT3] Authority \& Constraint Governance (II.1.2.13.2.1)}
{\small Capability · Canonical · current}\par

Establish who may make, approve or govern organisational decisions/artefacts and the bounds within which authority may be exercised.



This capability establishes who may make or approve organisational decisions and the bounds within which that authority applies. It complements obligation governance by distinguishing required outcomes from legitimate decision powers. A participant can therefore be responsible for acting within constraints without being authorised to redefine the governing requirement.



\textbf{Boundary}



Authority allocation does not itself perform an operational decision or a runtime access-control evaluation.


{\small Recorded unresolved matters: source-qualified matters; exact scope remains in the canonical record.}\par
{\footnotesize \href{../../../../continuum/strategy/capability/ad7f41dc48a046b185f044f0a24517ca.md}{Canonical record: ad7f41dc48a046b185f044f0a24517ca}; full qualifications and provenance.}\par
\paragraph*{II.1.2.13.2.2 [CT3] Policy / Obligation Governance}
\phantomsection
\label{navigation:II.1.2.13.2.2}
\addcontentsline{toc}{paragraph}{II.1.2.13.2.2 [CT3] Policy / Obligation Governance}

\label{element:73c6e5b36bdf4bf0ad2b8b50c9b9a00b}
\index{[CT3] Policy / Obligation Governance (II.1.2.13.2.2)}
{\small Capability · Canonical · current}\par

Identify, interpret, establish and maintain authoritative legislative, regulatory, policy, standards, guideline, contractual and other obligations, attributable and traceable to source authority.



Policy / Obligation Governance establishes and maintains the authoritative requirements under which the organisation operates. It interprets applicable legislative, regulatory, contractual, policy and other sources while preserving traceability to their authority. The purpose is a legitimate account of what must be satisfied, rather than an operational choice about how to satisfy it.



\textbf{Boundary}



Interpreting obligations does not acquire the regulator's authority or determine operational means.


{\small Recorded unresolved matters: source-qualified matters; exact scope remains in the canonical record.}\par
{\footnotesize \href{../../../../continuum/strategy/capability/73c6e5b36bdf4bf0ad2b8b50c9b9a00b.md}{Canonical record: 73c6e5b36bdf4bf0ad2b8b50c9b9a00b}; full qualifications and provenance.}\par
\subsection*{II.1.2.14 [CT1] Individual Care Delivery}
\phantomsection
\label{navigation:II.1.2.14}
\addcontentsline{toc}{subsection}{II.1.2.14 [CT1] Individual Care Delivery}

{\small Historical predecessor: R.2.48 \hyperref[element:162deb7e292042f78f677a14476a3a1d]{[CT1] Individual Care Delivery}; Canonical ID: \texttt{\seqsplit{162deb7e292042f78f677a14476a3a1d}}. Superseded historical identity — direct correspondence · superseded.}\par
\label{element:b964a11cde594318aa14ff935414158a}
\index{[CT1] Individual Care Delivery (II.1.2.14)}
{\small Capability · Canonical · current}\par

Assess, clinically determine, plan, intervene, monitor and review individual care independently of setting, mode, time pattern or organisational arrangement.



Individual Care Delivery is the organisation's ability to understand an individual's health and act on that understanding through purposeful care. Its six constituent abilities distinguish establishing findings, interpreting their clinical meaning, choosing a response, performing care, observing change and evaluating care against its objectives. These distinctions describe responsibilities and outcomes rather than a compulsory sequence: a consultation or continuing episode may revisit several of them.



The same ability can be exercised in a hospital, a general practice, a home or a virtual consultation. Setting and time pattern affect the circumstances of care without changing the responsibility represented by this capability.



\textbf{Intent}



Support clinically meaningful care whose intent, performance, observed effects and continuing appropriateness can be distinguished.



\textbf{Boundary}



Care access, appointments and operational coordination enable care but do not establish that it has been performed. Clinical judgement and the authority to deliver care remain with the responsible participants; possession or coordination of their information does not transfer that authority to Harmonia.



\textbf{Harmonia Architectural Relevance}



Harmonia supplies and connects information used in individual care. The approved Health Information Access support trace is bounded: records can inform care, but direct assessment can establish findings independently, and access to a proposed intervention does not authorise its performance.


{\footnotesize \href{../../../../continuum/strategy/capability/b964a11cde594318aa14ff935414158a.md}{Canonical record: b964a11cde594318aa14ff935414158a}; full qualifications and provenance.}\par
\subsubsection*{II.1.2.14.1 [CT2] Clinical Determination}
\phantomsection
\label{navigation:II.1.2.14.1}
\addcontentsline{toc}{subsubsection}{II.1.2.14.1 [CT2] Clinical Determination}

{\small Historical predecessor: R.2.33 \hyperref[element:018ccd818ec142fba1ffb086d900acc6]{[CT2] Clinical Determination}; Canonical ID: \texttt{\seqsplit{018ccd818ec142fba1ffb086d900acc6}}. Superseded historical identity — direct correspondence · superseded.}\par
\label{element:39219464b6304b59bd7c3de4e237066e}
\index{[CT2] Clinical Determination (II.1.2.14.1)}
{\small Capability · Canonical · current}\par

The capability to interpret assessed health information and establish clinically meaningful conclusions about an individual's health state, condition, risk or prognosis.



Clinical Determination turns assessed findings into clinically interpreted understanding of a person's condition, risk or prognosis. It represents the professional responsibility for deciding what the findings mean in the individual's clinical context, including the significance of changes identified through monitoring. This is the interpretive responsibility within Individual Care Delivery, rather than another method of collecting evidence.



\textbf{Intent}



Establish clinical conclusions on which an intended response can be based.



\textbf{Boundary}



Determining clinical meaning does not itself select the care response. No CT3 taxonomy has been adjudicated; diagnosis and prognostic judgement illustrate this responsibility without becoming additional catalogue elements.


{\small Recorded unresolved matters: source-qualified matters; exact scope remains in the canonical record.}\par
{\footnotesize \href{../../../../continuum/strategy/capability/39219464b6304b59bd7c3de4e237066e.md}{Canonical record: 39219464b6304b59bd7c3de4e237066e}; full qualifications and provenance.}\par
\paragraph*{II.1.2.14.1.1 [CT3] Clinical Finding Interpretation}
\phantomsection
\label{navigation:II.1.2.14.1.1}
\addcontentsline{toc}{paragraph}{II.1.2.14.1.1 [CT3] Clinical Finding Interpretation}

{\small Historical predecessor: R.2.30 \hyperref[element:cf8ddfca7a8e4ed893ec39b795200d77]{[CT3] Clinical Finding Interpretation}; Canonical ID: \texttt{\seqsplit{cf8ddfca7a8e4ed893ec39b795200d77}}. Superseded historical identity — direct correspondence · superseded.}\par
\label{element:deb78ba3b1514074a8ca5bc5da0673f5}
\index{[CT3] Clinical Finding Interpretation (II.1.2.14.1.1)}
{\small Capability · Proposed · current}\par

Interpret observed clinical findings and their significance for the individual care concern.



\textbf{Overview}



Interpret observed clinical findings and their significance for the individual care concern. The proposed decomposition separates this outcome from the other responsibilities within Clinical Determination.



\textbf{Intent}



Make the clinical significance of observations intelligible before diagnostic or response decisions are taken.



\textbf{Boundary}



Does not itself establish a diagnosis, prescribe intervention or confer this clinical responsibility on Harmonia.


{\footnotesize \href{../../../../continuum/strategy/capability/deb78ba3b1514074a8ca5bc5da0673f5.md}{Canonical record: deb78ba3b1514074a8ca5bc5da0673f5}; full qualifications and provenance.}\par
\paragraph*{II.1.2.14.1.2 [CT3] Diagnosis \& Problem Determination}
\phantomsection
\label{navigation:II.1.2.14.1.2}
\addcontentsline{toc}{paragraph}{II.1.2.14.1.2 [CT3] Diagnosis \& Problem Determination}

{\small Historical predecessor: R.2.31 \hyperref[element:e575c583783546ef85cf898523d18797]{[CT3] Diagnosis \& Problem Determination}; Canonical ID: \texttt{\seqsplit{e575c583783546ef85cf898523d18797}}. Superseded historical identity — direct correspondence · superseded.}\par
\label{element:d965032637e64275b7a61c0413834590}
\index{[CT3] Diagnosis \& Problem Determination (II.1.2.14.1.2)}
{\small Capability · Proposed · current}\par

Establish and qualify the clinical problems or diagnoses applicable to an individual.



\textbf{Overview}



Establish and qualify the clinical problems or diagnoses applicable to an individual. The proposed decomposition separates this outcome from the other responsibilities within Clinical Determination.



\textbf{Intent}



Establish the clinically accountable problem statement that subsequent care decisions address.



\textbf{Boundary}



Clinical determination remains with the mandated professional; a coded record is not evidence that Harmonia made the diagnosis.


{\footnotesize \href{../../../../continuum/strategy/capability/d965032637e64275b7a61c0413834590.md}{Canonical record: d965032637e64275b7a61c0413834590}; full qualifications and provenance.}\par
\paragraph*{II.1.2.14.1.3 [CT3] Health Risk \& Prognostic Determination}
\phantomsection
\label{navigation:II.1.2.14.1.3}
\addcontentsline{toc}{paragraph}{II.1.2.14.1.3 [CT3] Health Risk \& Prognostic Determination}

{\small Historical predecessor: R.2.32 \hyperref[element:3215f0b0e1ce45e8ab16dae272ea6fe7]{[CT3] Health Risk \& Prognostic Determination}; Canonical ID: \texttt{\seqsplit{3215f0b0e1ce45e8ab16dae272ea6fe7}}. Superseded historical identity — direct correspondence · superseded.}\par
\label{element:5bba0ea969814b5fa7ad11b67ab40a0c}
\index{[CT3] Health Risk \& Prognostic Determination (II.1.2.14.1.3)}
{\small Capability · Proposed · current}\par

Establish clinically meaningful risk and expected-course judgements for an individual.



\textbf{Overview}



Establish clinically meaningful risk and expected-course judgements for an individual. The proposed decomposition separates this outcome from the other responsibilities within Clinical Determination.



\textbf{Intent}



Make anticipated clinical risk and course explicit for the individual care concern.



\textbf{Boundary}



Risk/prognostic judgement is distinct from operational urgency, source-information authority and platform assurance.


{\footnotesize \href{../../../../continuum/strategy/capability/5bba0ea969814b5fa7ad11b67ab40a0c.md}{Canonical record: 5bba0ea969814b5fa7ad11b67ab40a0c}; full qualifications and provenance.}\par
\subsubsection*{II.1.2.14.2 [CT2] Clinical Intervention}
\phantomsection
\label{navigation:II.1.2.14.2}
\addcontentsline{toc}{subsubsection}{II.1.2.14.2 [CT2] Clinical Intervention}

{\small Historical predecessor: R.2.34 \hyperref[element:958345ee65f446d3b1259b8895ab100f]{[CT2] Clinical Intervention}; Canonical ID: \texttt{\seqsplit{958345ee65f446d3b1259b8895ab100f}}. Superseded historical identity — direct correspondence · superseded.}\par
\label{element:5359d1b0ff814a0e871384469f5ebcd7}
\index{[CT2] Clinical Intervention (II.1.2.14.2)}
{\small Capability · Canonical · current}\par

The capability to perform clinical actions intended to alter, maintain or otherwise manage an individual's health state in accordance with an established care intent.



Clinical Intervention is the ability to perform a clinical action under an established care intent. It marks the transition from intended care to care actually performed, whether the action seeks to alter, maintain or otherwise manage health. Its exercise can establish that treatment occurred even when the desired clinical result was not achieved.



\textbf{Intent}



Put the intended clinical response into practice through appropriately authorised care.



\textbf{Boundary}



Orders and appointments express or arrange future activity; they are not evidence that the intervention occurred. Therapy, medication and procedural examples do not constitute an approved three-child taxonomy: that proposed decomposition remains provisional.


{\small Recorded unresolved matters: source-qualified matters; exact scope remains in the canonical record.}\par
{\footnotesize \href{../../../../continuum/strategy/capability/5359d1b0ff814a0e871384469f5ebcd7.md}{Canonical record: 5359d1b0ff814a0e871384469f5ebcd7}; full qualifications and provenance.}\par
\subsubsection*{II.1.2.14.3 [CT2] Individual Care Planning}
\phantomsection
\label{navigation:II.1.2.14.3}
\addcontentsline{toc}{subsubsection}{II.1.2.14.3 [CT2] Individual Care Planning}

{\small Historical predecessor: R.2.37 \hyperref[element:574e61e230e643f78c58649adb4db313]{[CT2] Individual Care Planning}; Canonical ID: \texttt{\seqsplit{574e61e230e643f78c58649adb4db313}}. Superseded historical identity — direct correspondence · superseded.}\par
\label{element:f036a9866459429a8fb5f5e10727b5e8}
\index{[CT2] Individual Care Planning (II.1.2.14.3)}
{\small Capability · Canonical · current}\par

The capability to establish the intended clinical response to an individual's assessed and clinically interpreted health state, including the objectives, interventions and monitoring required to progress their care.



Individual Care Planning establishes the intended course of care in response to assessed and clinically interpreted health information. It connects clinical understanding to objectives, intended interventions and monitoring needs. Its outcome is an expression of what should happen, allowing intended care to remain distinguishable from requests, operational arrangements and actions subsequently performed.



\textbf{Intent}



Make the clinical response and its objectives explicit enough to guide care and later evaluate its effectiveness.



\textbf{Boundary}



Prescribing can express clinical intent without establishing dispensing or administration. Scheduling or allocating resources does not perform the planned intervention. No CT3 decomposition is established.


{\small Recorded unresolved matters: source-qualified matters; exact scope remains in the canonical record.}\par
{\footnotesize \href{../../../../continuum/strategy/capability/f036a9866459429a8fb5f5e10727b5e8.md}{Canonical record: f036a9866459429a8fb5f5e10727b5e8}; full qualifications and provenance.}\par
\paragraph*{II.1.2.14.3.1 [CT3] Care Goal Definition}
\phantomsection
\label{navigation:II.1.2.14.3.1}
\addcontentsline{toc}{paragraph}{II.1.2.14.3.1 [CT3] Care Goal Definition}

{\small Historical predecessor: R.2.35 \hyperref[element:8bfbcc0a455d423c90ae2f1cfb732ca8]{[CT3] Care Goal Definition}; Canonical ID: \texttt{\seqsplit{8bfbcc0a455d423c90ae2f1cfb732ca8}}. Superseded historical identity — direct correspondence · superseded.}\par
\label{element:609e1bef516b44bb976b5dbebe7f684b}
\index{[CT3] Care Goal Definition (II.1.2.14.3.1)}
{\small Capability · Proposed · current}\par

Establish the intended outcomes and priorities against which an individual care response is planned.



\textbf{Overview}



Establish the intended outcomes and priorities against which an individual care response is planned. The proposed decomposition separates this outcome from the other responsibilities within Individual Care Planning.



\textbf{Intent}



Establish the intended clinical benefit and priorities against which care responses can be considered.



\textbf{Boundary}



Goal definition remains distinct from selecting an intervention and from recording an externally agreed care plan.


{\footnotesize \href{../../../../continuum/strategy/capability/609e1bef516b44bb976b5dbebe7f684b.md}{Canonical record: 609e1bef516b44bb976b5dbebe7f684b}; full qualifications and provenance.}\par
\paragraph*{II.1.2.14.3.2 [CT3] Care Response Planning}
\phantomsection
\label{navigation:II.1.2.14.3.2}
\addcontentsline{toc}{paragraph}{II.1.2.14.3.2 [CT3] Care Response Planning}

{\small Historical predecessor: R.2.36 \hyperref[element:d64569acaf2c44f3b2a5332f74fff7c7]{[CT3] Care Response Planning}; Canonical ID: \texttt{\seqsplit{d64569acaf2c44f3b2a5332f74fff7c7}}. Superseded historical identity — direct correspondence · superseded.}\par
\label{element:2636e43903e44a97ad5b8db32692ce8b}
\index{[CT3] Care Response Planning (II.1.2.14.3.2)}
{\small Capability · Proposed · current}\par

Determine the intended clinical response and its sequencing in relation to agreed care goals.



\textbf{Overview}



Determine the intended clinical response and its sequencing in relation to agreed care goals. The proposed decomposition separates this outcome from the other responsibilities within Individual Care Planning.



\textbf{Intent}



Make the intended clinical response coherent in relation to the agreed care goals.



\textbf{Boundary}



Planning is distinct from intervention delivery and from the operational definition of a synthetic activity.


{\footnotesize \href{../../../../continuum/strategy/capability/2636e43903e44a97ad5b8db32692ce8b.md}{Canonical record: 2636e43903e44a97ad5b8db32692ce8b}; full qualifications and provenance.}\par
\subsubsection*{II.1.2.14.4 [CT2] Individual Care Review}
\phantomsection
\label{navigation:II.1.2.14.4}
\addcontentsline{toc}{subsubsection}{II.1.2.14.4 [CT2] Individual Care Review}

{\small Historical predecessor: R.2.40 \hyperref[element:e1048f6c635d4e16b5384cc2cb45a71b]{[CT2] Individual Care Review}; Canonical ID: \texttt{\seqsplit{e1048f6c635d4e16b5384cc2cb45a71b}}. Superseded historical identity — direct correspondence · superseded.}\par
\label{element:625a93f132f242f5a8e2ae1dc19979fb}
\index{[CT2] Individual Care Review (II.1.2.14.4)}
{\small Capability · Canonical · current}\par

The capability to evaluate the progress and effectiveness of an individual's care against its intended objectives and establish the continuing appropriateness of the current care intent.



Individual Care Review asks whether care is achieving its intended objectives and whether the existing care intent remains suitable. Its children separate effectiveness from appropriateness: an intervention may have an effect yet no longer be suitable to continue, or remain a reasonable choice despite an inadequate response. Review therefore evaluates the relationship between care, outcomes and intent rather than simply repeating assessment.



\textbf{Intent}



Establish whether continuing the current approach is justified and whether reconsideration is needed.



\textbf{Boundary}



Recognising a need to change care does not establish the replacement plan. Monitoring supplies relevant observations; planning establishes revised intent.


{\small Recorded unresolved matters: source-qualified matters; exact scope remains in the canonical record.}\par
{\footnotesize \href{../../../../continuum/strategy/capability/625a93f132f242f5a8e2ae1dc19979fb.md}{Canonical record: 625a93f132f242f5a8e2ae1dc19979fb}; full qualifications and provenance.}\par
\paragraph*{II.1.2.14.4.1 [CT3] Care Appropriateness Review}
\phantomsection
\label{navigation:II.1.2.14.4.1}
\addcontentsline{toc}{paragraph}{II.1.2.14.4.1 [CT3] Care Appropriateness Review}

{\small Historical predecessor: R.2.39 \hyperref[element:905eb2de1ac341ac96331e3d7f76079b]{[CT3] Care Intent Appropriateness Review}; Canonical ID: \texttt{\seqsplit{905eb2de1ac341ac96331e3d7f76079b}}. Superseded historical identity — evolved correspondence · superseded.}\par
\label{element:5eac5fb73e31416288d8c992805b33cb}
\index{[CT3] Care Appropriateness Review (II.1.2.14.4.1)}
{\small Capability · Canonical · current}\par

Evaluate whether the current intended care remains suitable in the individual’s current context.



Care Appropriateness Review evaluates whether the current care intent remains suitable for the individual in their present context. Its question concerns continuing suitability rather than whether treatment has produced a measurable effect. That distinction permits a review to recognise that effective care may no longer be appropriate to continue.



\textbf{Boundary}



An appropriateness judgement can identify a need to reconsider care; establishing revised intent remains Individual Care Planning.


{\small Recorded unresolved matters: source-qualified matters; exact scope remains in the canonical record.}\par
{\footnotesize \href{../../../../continuum/strategy/capability/5eac5fb73e31416288d8c992805b33cb.md}{Canonical record: 5eac5fb73e31416288d8c992805b33cb}; full qualifications and provenance.}\par
\paragraph*{II.1.2.14.4.2 [CT3] Care Effectiveness Review}
\phantomsection
\label{navigation:II.1.2.14.4.2}
\addcontentsline{toc}{paragraph}{II.1.2.14.4.2 [CT3] Care Effectiveness Review}

{\small Historical predecessor: R.2.38 \hyperref[element:31f49cdea9f44f188459503580c2b86e]{[CT3] Care Effectiveness Review}; Canonical ID: \texttt{\seqsplit{31f49cdea9f44f188459503580c2b86e}}. Superseded historical identity — direct correspondence · superseded.}\par
\label{element:79cfeaefeb174a54ba43343d37f035c1}
\index{[CT3] Care Effectiveness Review (II.1.2.14.4.2)}
{\small Capability · Canonical · current}\par

Evaluate whether delivered care is achieving or has achieved its intended goals and outcomes.



Care Effectiveness Review evaluates delivered care against its intended goals and outcomes. It connects observed response to the reason the care was undertaken, making a judgement that monitoring alone does not supply. An effectiveness finding can inform review of appropriateness while preserving the possibility that a reasonable care choice has not achieved its objective.



\textbf{Boundary}



Observed improvement is not automatically proof that the intended objectives have been achieved or that continuing care is appropriate.


{\small Recorded unresolved matters: source-qualified matters; exact scope remains in the canonical record.}\par
{\footnotesize \href{../../../../continuum/strategy/capability/79cfeaefeb174a54ba43343d37f035c1.md}{Canonical record: 79cfeaefeb174a54ba43343d37f035c1}; full qualifications and provenance.}\par
\subsubsection*{II.1.2.14.5 [CT2] Individual Health Assessment}
\phantomsection
\label{navigation:II.1.2.14.5}
\addcontentsline{toc}{subsubsection}{II.1.2.14.5 [CT2] Individual Health Assessment}

{\small Historical predecessor: R.2.44 \hyperref[element:9111446aa2d045bd9acd98542b0eef1f]{[CT2] Individual Health Assessment}; Canonical ID: \texttt{\seqsplit{9111446aa2d045bd9acd98542b0eef1f}}. Superseded historical identity — direct correspondence · superseded.}\par
\label{element:045f819bb3694d3c9db7f20886e2993f}
\index{[CT2] Individual Health Assessment (II.1.2.14.5)}
{\small Capability · Canonical · current}\par

The capability to establish an understanding of an individual's current health state through the collection, consideration and evaluation of relevant clinical, health and contextual information.



Assessment establishes what is currently known about an individual's health through history, examination and measurement. Within Individual Care Delivery it supplies assessed findings, including findings obtained directly from the person rather than retrieved from an existing record. Its three children distinguish ways of establishing knowledge, rather than professions, clinical specialties or locations.



\textbf{Intent}



Give subsequent clinical interpretation a defensible account of the person's current state and relevant history.



\textbf{Boundary}



Finding an abnormal measurement is distinguishable from deciding what it means clinically. Diagnosis, prognosis, selection of care and performance of care belong to adjacent abilities.


{\small Recorded unresolved matters: source-qualified matters; exact scope remains in the canonical record.}\par
{\footnotesize \href{../../../../continuum/strategy/capability/045f819bb3694d3c9db7f20886e2993f.md}{Canonical record: 045f819bb3694d3c9db7f20886e2993f}; full qualifications and provenance.}\par
\paragraph*{II.1.2.14.5.1 [CT3] Clinical Examination}
\phantomsection
\label{navigation:II.1.2.14.5.1}
\addcontentsline{toc}{paragraph}{II.1.2.14.5.1 [CT3] Clinical Examination}

{\small Historical predecessor: R.2.41 \hyperref[element:8bc15e2b8f114deebcbbe3f2013a8147]{[CT3] Clinical Examination}; Canonical ID: \texttt{\seqsplit{8bc15e2b8f114deebcbbe3f2013a8147}}. Superseded historical identity — direct correspondence · superseded.}\par
\label{element:8855d4d9e2954e0ab66ab65a556fc61a}
\index{[CT3] Clinical Examination (II.1.2.14.5.1)}
{\small Capability · Canonical · current}\par

Establish health findings through clinical examination.



Clinical Examination establishes findings through examination of the individual. It distinguishes directly acquired clinical knowledge from the historical account and the measurement-focused contribution of its assessment siblings. Its outcome is a set of examination findings that can inform clinical interpretation, irrespective of the setting or profession exercising the ability.



\textbf{Boundary}



An examination finding is not itself a diagnosis, care plan or performed treatment.


{\small Recorded unresolved matters: source-qualified matters; exact scope remains in the canonical record.}\par
{\footnotesize \href{../../../../continuum/strategy/capability/8855d4d9e2954e0ab66ab65a556fc61a.md}{Canonical record: 8855d4d9e2954e0ab66ab65a556fc61a}; full qualifications and provenance.}\par
\paragraph*{II.1.2.14.5.2 [CT3] Health History Assessment}
\phantomsection
\label{navigation:II.1.2.14.5.2}
\addcontentsline{toc}{paragraph}{II.1.2.14.5.2 [CT3] Health History Assessment}

{\small Historical predecessor: R.2.42 \hyperref[element:bcaacf446f45443eafd8e21a1fd74ee1]{[CT3] Health History Assessment}; Canonical ID: \texttt{\seqsplit{bcaacf446f45443eafd8e21a1fd74ee1}}. Superseded historical identity — direct correspondence · superseded.}\par
\label{element:7ea7e3f74e2940a7934b588a1a038834}
\index{[CT3] Health History Assessment (II.1.2.14.5.2)}
{\small Capability · Canonical · current}\par

Establish relevant health knowledge from the individual’s history and applicable historical/contextual information.



Health History Assessment establishes clinically relevant understanding from a person's history and its surrounding context. Within assessment it contributes knowledge that a current examination or measurement alone may not reveal, while requiring consideration of the relevance of historical information to the present situation. Access to previous records can assist, but the responsibility is the assessment of history rather than retrieval itself.



\textbf{Boundary}



Historical information does not become a current clinical conclusion merely because it is available; interpretation of the assessed findings remains Clinical Determination.


{\small Recorded unresolved matters: source-qualified matters; exact scope remains in the canonical record.}\par
{\footnotesize \href{../../../../continuum/strategy/capability/7ea7e3f74e2940a7934b588a1a038834.md}{Canonical record: 7ea7e3f74e2940a7934b588a1a038834}; full qualifications and provenance.}\par
\paragraph*{II.1.2.14.5.3 [CT3] Health Status Measurement}
\phantomsection
\label{navigation:II.1.2.14.5.3}
\addcontentsline{toc}{paragraph}{II.1.2.14.5.3 [CT3] Health Status Measurement}

{\small Historical predecessor: R.2.43 \hyperref[element:3717906f7f8840d3a8a3e3877543dd18]{[CT3] Health Status Measurement}; Canonical ID: \texttt{\seqsplit{3717906f7f8840d3a8a3e3877543dd18}}. Superseded historical identity — direct correspondence · superseded.}\par
\label{element:a796bfb3335d44f6aeb9970628ed8554}
\index{[CT3] Health Status Measurement (II.1.2.14.5.3)}
{\small Capability · Canonical · current}\par

Establish individual health observations or measurements.



Health Status Measurement establishes observations or measurements of an individual's health. Its contribution to assessment is evidence about the measured state, with a purpose distinct from developing a view over time or deciding the clinical meaning of a result. Measurements can inform both present assessment and later monitoring without collapsing those abilities.



\textbf{Boundary}



Obtaining a value does not establish a temporal trend or the clinical significance of that value.


{\small Recorded unresolved matters: source-qualified matters; exact scope remains in the canonical record.}\par
{\footnotesize \href{../../../../continuum/strategy/capability/a796bfb3335d44f6aeb9970628ed8554.md}{Canonical record: a796bfb3335d44f6aeb9970628ed8554}; full qualifications and provenance.}\par
\subsubsection*{II.1.2.14.6 [CT2] Individual Health Monitoring}
\phantomsection
\label{navigation:II.1.2.14.6}
\addcontentsline{toc}{subsubsection}{II.1.2.14.6 [CT2] Individual Health Monitoring}

{\small Historical predecessor: R.2.47 \hyperref[element:31e7c27556314cafb74d040441d2dac6]{[CT2] Individual Health Monitoring}; Canonical ID: \texttt{\seqsplit{31e7c27556314cafb74d040441d2dac6}}. Superseded historical identity — direct correspondence · superseded.}\par
\label{element:c7d0e498e1ad44138df940af2f91cd2d}
\index{[CT2] Individual Health Monitoring (II.1.2.14.6)}
{\small Capability · Canonical · current}\par

The capability to observe and evaluate an individual's health state over time in order to identify clinically meaningful change, stability or deviation from an expected state or trajectory.



Monitoring develops a view of an individual's health over time so that change, stability or deviation can be recognised. Health State Monitoring establishes the aggregate temporally situated view; Care Response Monitoring considers that view in relation to a particular condition and treatment. Together they distinguish observation over time from a single measurement and from the clinical judgement prompted by a trend.



\textbf{Intent}



Make relevant trajectories visible while preserving what was observed and what still requires interpretation.



\textbf{Boundary}



Monitoring does not establish the clinical significance of a change or judge whether care achieved its objectives. Those responsibilities remain Clinical Determination and Individual Care Review respectively.


{\small Recorded unresolved matters: source-qualified matters; exact scope remains in the canonical record.}\par
{\footnotesize \href{../../../../continuum/strategy/capability/c7d0e498e1ad44138df940af2f91cd2d.md}{Canonical record: c7d0e498e1ad44138df940af2f91cd2d}; full qualifications and provenance.}\par
\paragraph*{II.1.2.14.6.1 [CT3] Care Response Monitoring}
\phantomsection
\label{navigation:II.1.2.14.6.1}
\addcontentsline{toc}{paragraph}{II.1.2.14.6.1 [CT3] Care Response Monitoring}

{\small Historical predecessor: R.2.45 \hyperref[element:615f93d05fbb4a8481f07778bd4f4803]{[CT3] Care Response Monitoring}; Canonical ID: \texttt{\seqsplit{615f93d05fbb4a8481f07778bd4f4803}}. Superseded historical identity — direct correspondence · superseded.}\par
\label{element:d6951c9357fa40a4943cde581d95d5b1}
\index{[CT3] Care Response Monitoring (II.1.2.14.6.1)}
{\small Capability · Canonical · current}\par

Scope that health-state view to a particular problem/condition and the treatment applied to it.



Care Response Monitoring considers the monitored health-state view in relation to a particular problem or condition and the treatment applied. It makes a response trajectory visible in its treatment context, helping distinguish a general change in health from a change relevant to the care being followed. It consumes the aggregate view established through Health State Monitoring.



\textbf{Boundary}



Identifying the response does not judge treatment effectiveness against objectives or determine a replacement treatment.


{\small Recorded unresolved matters: source-qualified matters; exact scope remains in the canonical record.}\par
{\footnotesize \href{../../../../continuum/strategy/capability/d6951c9357fa40a4943cde581d95d5b1.md}{Canonical record: d6951c9357fa40a4943cde581d95d5b1}; full qualifications and provenance.}\par
\paragraph*{II.1.2.14.6.2 [CT3] Health State Monitoring}
\phantomsection
\label{navigation:II.1.2.14.6.2}
\addcontentsline{toc}{paragraph}{II.1.2.14.6.2 [CT3] Health State Monitoring}

{\small Historical predecessor: R.2.46 \hyperref[element:346f024f74b447a8b6056b9d521837ff]{[CT3] Health State Monitoring}; Canonical ID: \texttt{\seqsplit{346f024f74b447a8b6056b9d521837ff}}. Superseded historical identity — direct correspondence · superseded.}\par
\label{element:d77f22909f8b47dba4cbd95982e576df}
\index{[CT3] Health State Monitoring (II.1.2.14.6.2)}
{\small Capability · Canonical · current}\par

Establish an aggregate, temporally situated view of the individual’s health from applicable monitored data.



Health State Monitoring brings applicable monitored data into an aggregate view situated in time. Its purpose is to make the individual's changing or stable health state intelligible across observations, rather than leaving each observation isolated. Within monitoring it provides the broader view that can also be considered in relation to a particular treatment or condition.



\textbf{Boundary}



Aggregation is not clinical interpretation. Care Response Monitoring uses the established view rather than independently recreating it.


{\small Recorded unresolved matters: source-qualified matters; exact scope remains in the canonical record.}\par
{\footnotesize \href{../../../../continuum/strategy/capability/d77f22909f8b47dba4cbd95982e576df.md}{Canonical record: d77f22909f8b47dba4cbd95982e576df}; full qualifications and provenance.}\par
\subsection*{II.1.2.15 [CT1] Population Health \& Health-System Planning}
\phantomsection
\label{navigation:II.1.2.15}
\addcontentsline{toc}{subsection}{II.1.2.15 [CT1] Population Health \& Health-System Planning}

{\small Historical predecessor: R.2.49 \hyperref[element:a5003436a7f348b1bff60ba593adfa27]{[CT1] Population Health \& Health-System Planning}; Canonical ID: \texttt{\seqsplit{a5003436a7f348b1bff60ba593adfa27}}. Superseded historical identity — direct correspondence · superseded.}\par
\label{element:c63b38b37e9b4e9b97dc87ed54604def}
\index{[CT1] Population Health \& Health-System Planning (II.1.2.15)}
{\small Capability · Canonical · current}\par

Understand population health through assessment/surveillance and applicable event/notification management; assess health-service demand and plan required capacity and its distribution.



This capability connects an understanding of population health with consideration of future health-service needs. Surveillance addresses population state, significant events and notification; planning assesses service demand and required capacity and distribution. Population evidence can inform planning without predetermining its conclusions or becoming the operational allocation of resources.



\textbf{Intent}



Support informed public-health awareness and a reasoned view of the services a population will require.



\textbf{Boundary}



Population conclusions differ from individual clinical findings, and planned capacity differs from capacity currently usable by a service.



\textbf{Harmonia Architectural Relevance}



Governed population information and formal notifications materially support health activity. Consuming planning context does not make the platform the health-system planner or owner of its resources.


{\footnotesize \href{../../../../continuum/strategy/capability/c63b38b37e9b4e9b97dc87ed54604def.md}{Canonical record: c63b38b37e9b4e9b97dc87ed54604def}; full qualifications and provenance.}\par
\subsubsection*{II.1.2.15.1 [CT2] Health-System Planning}
\phantomsection
\label{navigation:II.1.2.15.1}
\addcontentsline{toc}{subsubsection}{II.1.2.15.1 [CT2] Health-System Planning}

\label{element:9ca4df07eceb4150a0190cc488270ce2}
\index{[CT2] Health-System Planning (II.1.2.15.1)}
{\small Capability · Canonical · current}\par

Assess health-service demand and plan the capacity and footprint required to meet it.



Health-System Planning considers the capacity and service footprint required to meet assessed demand. Demand Assessment establishes the basis and scale of need, Capacity Planning determines what is required, and Distribution Planning addresses where that capacity should sit across the service network. The decomposition preserves the difference between understanding demand and deciding a response to it.



\textbf{Intent}



Support a defensible relationship between expected demand and the planned scale and distribution of health services.



\textbf{Boundary}



Plans do not establish available resources or operational allocations. Detailed commissioning, capital, funding and workforce planning remain deliberately unmodelled.


{\small Recorded unresolved matters: source-qualified matters; exact scope remains in the canonical record.}\par
{\footnotesize \href{../../../../continuum/strategy/capability/9ca4df07eceb4150a0190cc488270ce2.md}{Canonical record: 9ca4df07eceb4150a0190cc488270ce2}; full qualifications and provenance.}\par
\paragraph*{II.1.2.15.1.1 [CT3] Health Service Capacity Planning}
\phantomsection
\label{navigation:II.1.2.15.1.1}
\addcontentsline{toc}{paragraph}{II.1.2.15.1.1 [CT3] Health Service Capacity Planning}

\label{element:c1d1db0e2c854077a41721615e86119d}
\index{[CT3] Health Service Capacity Planning (II.1.2.15.1.1)}
{\small Capability · Canonical · current}\par

Determine required strategic/tactical capacity to address assessed demand.



Health Service Capacity Planning determines the strategic or tactical capacity required to address assessed demand. It considers what should be available, preserving a distinction between requirement and the resources currently in service. Within Health-System Planning it gives demand a capacity response without performing operational allocation.



\textbf{Boundary}



Planned capacity is neither established resource availability nor the service's currently usable operational capacity.


{\small Recorded unresolved matters: source-qualified matters; exact scope remains in the canonical record.}\par
{\footnotesize \href{../../../../continuum/strategy/capability/c1d1db0e2c854077a41721615e86119d.md}{Canonical record: c1d1db0e2c854077a41721615e86119d}; full qualifications and provenance.}\par
\paragraph*{II.1.2.15.1.2 [CT3] Health Service Demand Assessment}
\phantomsection
\label{navigation:II.1.2.15.1.2}
\addcontentsline{toc}{paragraph}{II.1.2.15.1.2 [CT3] Health Service Demand Assessment}

\label{element:e56d354eafd94ff7a495fc2ca2b8a593}
\index{[CT3] Health Service Demand Assessment (II.1.2.15.1.2)}
{\small Capability · Canonical · current}\par

Assess current/expected health-service demand across a population, geography, service domain or health-system context.



Health Service Demand Assessment establishes a view of current or expected demand within a population, geography or service context. It translates relevant evidence into a demand judgement rather than assuming population-health state directly determines service requirements. Its assumptions and scope matter because capacity planning depends on what demand has actually been assessed.



\textbf{Boundary}



System-level demand is distinct from a particular person's care requirement and does not itself determine planned capacity.


{\small Recorded unresolved matters: source-qualified matters; exact scope remains in the canonical record.}\par
{\footnotesize \href{../../../../continuum/strategy/capability/e56d354eafd94ff7a495fc2ca2b8a593.md}{Canonical record: e56d354eafd94ff7a495fc2ca2b8a593}; full qualifications and provenance.}\par
\paragraph*{II.1.2.15.1.3 [CT3] Health Service Distribution Planning}
\phantomsection
\label{navigation:II.1.2.15.1.3}
\addcontentsline{toc}{paragraph}{II.1.2.15.1.3 [CT3] Health Service Distribution Planning}

\label{element:4fc525fccbe142a58ee43dcc58b11c92}
\index{[CT3] Health Service Distribution Planning (II.1.2.15.1.3)}
{\small Capability · Canonical · current}\par

Determine the distribution of planned capacity across a geographical, organisational, site or service-network footprint.



Health Service Distribution Planning determines how planned capacity should be distributed across a geographical, organisational or service-network footprint. It addresses the placement of the planned response rather than only its aggregate scale. This allows the planning context to distinguish enough total capacity from capacity positioned where it is needed.



\textbf{Boundary}



A planned footprint does not establish recognised location records, identifiable offerings or actual resource deployment.


{\small Recorded unresolved matters: source-qualified matters; exact scope remains in the canonical record.}\par
{\footnotesize \href{../../../../continuum/strategy/capability/4fc525fccbe142a58ee43dcc58b11c92.md}{Canonical record: 4fc525fccbe142a58ee43dcc58b11c92}; full qualifications and provenance.}\par
\subsubsection*{II.1.2.15.2 [CT2] Population Health Surveillance}
\phantomsection
\label{navigation:II.1.2.15.2}
\addcontentsline{toc}{subsubsection}{II.1.2.15.2 [CT2] Population Health Surveillance}

\label{element:83e568fa8af14d7791a7253765ea390f}
\index{[CT2] Population Health Surveillance (II.1.2.15.2)}
{\small Capability · Canonical · current}\par

Understand population-health state and manage organised surveillance, significant events and applicable notification.



Population Health Surveillance establishes patterns and change across a defined population and manages the surveillance arrangements through which they are considered. Its children separate longitudinal assessment, surveillance management, event identification, formal notification and reporting. These distinctions matter because an observed pattern can be significant without being notifiable, and identifying a notifiable event does not demonstrate that the notification obligation has been discharged.



\textbf{Intent}



Make population-health concerns and applicable communication obligations visible in their population and temporal context.



\textbf{Boundary}



Notification management does not define the obligation or own the operational response. Surveillance evidence can inform other abilities without establishing their findings.


{\small Recorded unresolved matters: source-qualified matters; exact scope remains in the canonical record.}\par
{\footnotesize \href{../../../../continuum/strategy/capability/83e568fa8af14d7791a7253765ea390f.md}{Canonical record: 83e568fa8af14d7791a7253765ea390f}; full qualifications and provenance.}\par
\paragraph*{II.1.2.15.2.1 [CT3] Notifiable Health Event Management}
\phantomsection
\label{navigation:II.1.2.15.2.1}
\addcontentsline{toc}{paragraph}{II.1.2.15.2.1 [CT3] Notifiable Health Event Management}

\label{element:e605436254ed4dad819f452448e054a1}
\index{[CT3] Notifiable Health Event Management (II.1.2.15.2.1)}
{\small Capability · Canonical · current}\par

Manage formal notification applicability, required information, recipient, status, submission, correction/supplementation and applicable follow-up.



Notifiable Health Event Management addresses the formal communication responsibilities applicable to a health event. It manages applicability, required information, recipient and submission status, including correction or supplementation where relevant. Its purpose is to make the notification obligation and its discharge distinguishable from both event identification and ordinary reporting.



\textbf{Boundary}



Notification management does not define the governing obligation or own the operational response to the event.


{\small Recorded unresolved matters: source-qualified matters; exact scope remains in the canonical record.}\par
{\footnotesize \href{../../../../continuum/strategy/capability/e605436254ed4dad819f452448e054a1.md}{Canonical record: e605436254ed4dad819f452448e054a1}; full qualifications and provenance.}\par
\paragraph*{II.1.2.15.2.2 [CT3] Population Health Event Identification}
\phantomsection
\label{navigation:II.1.2.15.2.2}
\addcontentsline{toc}{paragraph}{II.1.2.15.2.2 [CT3] Population Health Event Identification}

\label{element:e084ba370bb545d0a4c28c37e2e8e850}
\index{[CT3] Population Health Event Identification (II.1.2.15.2.2)}
{\small Capability · Canonical · current}\par

Identify and characterise significant population-health events/patterns from applicable evidence.



Population Health Event Identification recognises and characterises a significant event or pattern from relevant evidence. It establishes the population-level concern rather than merely noting that observations changed. The resulting event context can inform notification or response decisions while retaining a separate responsibility for those decisions.



\textbf{Boundary}



An identified event is not automatically notifiable, and notifiability does not demonstrate that an obligation was discharged.


{\small Recorded unresolved matters: source-qualified matters; exact scope remains in the canonical record.}\par
{\footnotesize \href{../../../../continuum/strategy/capability/e084ba370bb545d0a4c28c37e2e8e850.md}{Canonical record: e084ba370bb545d0a4c28c37e2e8e850}; full qualifications and provenance.}\par
\paragraph*{II.1.2.15.2.3 [CT3] Population Health Longitudinal Assessment}
\phantomsection
\label{navigation:II.1.2.15.2.3}
\addcontentsline{toc}{paragraph}{II.1.2.15.2.3 [CT3] Population Health Longitudinal Assessment}

\label{element:4d52b8acda504329a595be27a736e63c}
\index{[CT3] Population Health Longitudinal Assessment (II.1.2.15.2.3)}
{\small Capability · Canonical · current}\par

Establish population-health state, patterns and change over time from applicable evidence.



This capability establishes population-health state, patterns and change over time from relevant evidence. It interprets information at a population scope rather than treating a collection of individual observations as a population conclusion by itself. Its longitudinal purpose makes the temporal and population basis important to the meaning of the assessment.



\textbf{Boundary}



Population assessment is not preventative-programme performance, an individual health conclusion or a service-demand decision.


{\small Recorded unresolved matters: source-qualified matters; exact scope remains in the canonical record.}\par
{\footnotesize \href{../../../../continuum/strategy/capability/4d52b8acda504329a595be27a736e63c.md}{Canonical record: 4d52b8acda504329a595be27a736e63c}; full qualifications and provenance.}\par
\paragraph*{II.1.2.15.2.4 [CT3] Population Health Reporting}
\phantomsection
\label{navigation:II.1.2.15.2.4}
\addcontentsline{toc}{paragraph}{II.1.2.15.2.4 [CT3] Population Health Reporting}

\label{element:50b8b0884f8444c4b11d9e0f7a7754a3}
\index{[CT3] Population Health Reporting (II.1.2.15.2.4)}
{\small Capability · Canonical · current}\par

Communicate established population-health assessments, surveillance, events, trends and findings.



Population Health Reporting communicates established assessments, surveillance findings, events and trends. Its role is to make population understanding available with its scope intact, rather than replace the assessment or event-identification work underlying it. This supports use of the findings by other abilities without assuming what those abilities must conclude.



\textbf{Boundary}



General reporting is distinct from formal notification management and does not establish a response or planning decision.


{\small Recorded unresolved matters: source-qualified matters; exact scope remains in the canonical record.}\par
{\footnotesize \href{../../../../continuum/strategy/capability/50b8b0884f8444c4b11d9e0f7a7754a3.md}{Canonical record: 50b8b0884f8444c4b11d9e0f7a7754a3}; full qualifications and provenance.}\par
\paragraph*{II.1.2.15.2.5 [CT3] Population Health Surveillance Management}
\phantomsection
\label{navigation:II.1.2.15.2.5}
\addcontentsline{toc}{paragraph}{II.1.2.15.2.5 [CT3] Population Health Surveillance Management}

\label{element:f0791768ef494e619006ee0d89d0d9d8}
\index{[CT3] Population Health Surveillance Management (II.1.2.15.2.5)}
{\small Capability · Canonical · current}\par

Establish and manage surveillance purposes, populations, concerns and regimes.



Population Health Surveillance Management establishes the purposes, populations and regimes under which surveillance operates. It concerns the organised basis for watching health concerns, rather than any particular event detected through that activity. Keeping the regime distinct from its findings makes it possible to understand why evidence is being collected and considered.



\textbf{Boundary}



Managing surveillance neither identifies every significant event nor determines the operational response to one.


{\small Recorded unresolved matters: source-qualified matters; exact scope remains in the canonical record.}\par
{\footnotesize \href{../../../../continuum/strategy/capability/f0791768ef494e619006ee0d89d0d9d8.md}{Canonical record: f0791768ef494e619006ee0d89d0d9d8}; full qualifications and provenance.}\par
\subsection*{II.1.2.16 [CT1] Security, Privacy \& Digital Trust}
\phantomsection
\label{navigation:II.1.2.16}
\addcontentsline{toc}{subsection}{II.1.2.16 [CT1] Security, Privacy \& Digital Trust}

{\small Historical predecessor: R.2.50 \hyperref[element:b01441b2f5584fc2b9c3988643a7c93b]{[CT1] Security, Privacy \& Digital Trust}; Canonical ID: \texttt{\seqsplit{b01441b2f5584fc2b9c3988643a7c93b}}. Superseded historical identity — direct correspondence · superseded.}\par
\label{element:71f0431fdfcd4f779e16d9e5ea609833}
\index{[CT1] Security, Privacy \& Digital Trust (II.1.2.16)}
{\small Capability · Canonical · current}\par

Establish adequate identity confidence, coherent access controls, privacy/protection safeguards, governed activity attribution and accountability for security-relevant activity and information use.



This capability establishes confidence in identity, coherent control of access, protection of information and accountability for its use. Its three branches distinguish identity/access governance, protection and privacy, and attribution/accountability. The scope includes physical and procedural arrangements as well as digital controls, preventing a software authorisation check from standing in for the organisation's whole trust responsibility.



\textbf{Intent}



Support information use and governed activity within legitimate identity, authority, privacy and protection constraints.



\textbf{Boundary}



Identifying an actor, permitting an action, recording participation and judging appropriateness are different outcomes. Evidence production is not itself independent assurance.



\textbf{Harmonia Architectural Relevance}



Security and governance are intrinsic to Harmonia-managed operations. The Core classification concerns assigned platform controls and evidence; it does not acquire every organisational security or privacy responsibility.


{\footnotesize \href{../../../../continuum/strategy/capability/71f0431fdfcd4f779e16d9e5ea609833.md}{Canonical record: 71f0431fdfcd4f779e16d9e5ea609833}; full qualifications and provenance.}\par
\subsubsection*{II.1.2.16.1 [CT2] Digital Accountability \& Trust}
\phantomsection
\label{navigation:II.1.2.16.1}
\addcontentsline{toc}{subsubsection}{II.1.2.16.1 [CT2] Digital Accountability \& Trust}

\label{element:9a600233ae1f4f3289c09a0b08ce029d}
\index{[CT2] Digital Accountability \& Trust (II.1.2.16.1)}
{\small Capability · Canonical · current}\par

Establish governed attribution and accountability for activity and information use.



Digital Accountability \& Trust makes activity attributable and evaluates whether security-relevant activity and information use were appropriate. Attribution records participation together with acting capacity and relevant permission or constraint context; accountability uses that context to consider conduct against policy and obligations. A record of who acted is therefore necessary context, rather than proof that the action was permitted.



\textbf{Intent}



Support meaningful responsibility for information use and security-relevant conduct.



\textbf{Boundary}



Authentication, attribution and authorisation must remain distinguishable. Expected downstream evidentiary use does not establish new BC-to-BEC relationships.


{\small Recorded unresolved matters: source-qualified matters; exact scope remains in the canonical record.}\par
{\footnotesize \href{../../../../continuum/strategy/capability/9a600233ae1f4f3289c09a0b08ce029d.md}{Canonical record: 9a600233ae1f4f3289c09a0b08ce029d}; full qualifications and provenance.}\par
\paragraph*{II.1.2.16.1.1 [CT3] Governed Activity Attribution}
\phantomsection
\label{navigation:II.1.2.16.1.1}
\addcontentsline{toc}{paragraph}{II.1.2.16.1.1 [CT3] Governed Activity Attribution}

\label{element:646b3fd80a8b4f269691db7942e0f809}
\index{[CT3] Governed Activity Attribution (II.1.2.16.1.1)}
{\small Capability · Canonical · current}\par

Preserve who participated, what was done, the attributable person/role/unit/organisation capacity and applicable permission/authority/constraint context.



Governed Activity Attribution preserves who participated and in what acting capacity, together with the relevant permission, authority and constraint context. It makes conduct explainable beyond a bare actor identifier and timestamp. Within Digital Accountability \& Trust, this supplies the contextual record needed to distinguish participation from entitlement.



\textbf{Boundary}



Attributing an action to a person or role does not establish that the action was permitted or appropriate.


{\small Recorded unresolved matters: source-qualified matters; exact scope remains in the canonical record.}\par
{\footnotesize \href{../../../../continuum/strategy/capability/646b3fd80a8b4f269691db7942e0f809.md}{Canonical record: 646b3fd80a8b4f269691db7942e0f809}; full qualifications and provenance.}\par
\paragraph*{II.1.2.16.1.2 [CT3] Security \& Information-Use Accountability}
\phantomsection
\label{navigation:II.1.2.16.1.2}
\addcontentsline{toc}{paragraph}{II.1.2.16.1.2 [CT3] Security \& Information-Use Accountability}

\label{element:a33d0cc14c384bd48012def4ded54a67}
\index{[CT3] Security \& Information-Use Accountability (II.1.2.16.1.2)}
{\small Capability · Canonical · current}\par

Evaluate whether attributed security-relevant activity/information use was appropriate within policy, obligation and authority constraints, and establish accountability.



This capability evaluates attributed activity and information use against applicable policy, obligations and authority constraints. It establishes accountability rather than simply recording participation. Its purpose depends on preserving the circumstances of action so that appropriateness is considered within the relevant mandate instead of inferred from identity or technical success.



\textbf{Boundary}



An authenticated actor and a successful operation are not proof of legitimate information use.


{\small Recorded unresolved matters: source-qualified matters; exact scope remains in the canonical record.}\par
{\footnotesize \href{../../../../continuum/strategy/capability/a33d0cc14c384bd48012def4ded54a67.md}{Canonical record: a33d0cc14c384bd48012def4ded54a67}; full qualifications and provenance.}\par
\subsubsection*{II.1.2.16.2 [CT2] Digital Identity \& Access Governance}
\phantomsection
\label{navigation:II.1.2.16.2}
\addcontentsline{toc}{subsubsection}{II.1.2.16.2 [CT2] Digital Identity \& Access Governance}

\label{element:b9cb3666f9134ba29eb896e67eba7250}
\index{[CT2] Digital Identity \& Access Governance (II.1.2.16.2)}
{\small Capability · Canonical · current}\par

Govern sufficient identity confidence and coherent access controls across the organisation.



This capability governs who may act and how controlled concerns are identified and protected across the organisation. Its children separate confidence in a claimed identity, identity models, contextual identity quality, authentication, access-control context, assurance, collaboration and enforcement. That separation prevents successful authentication from being treated as sufficient authority for every use.



\textbf{Intent}



Make identity confidence and control requirements appropriate to the circumstances in which access or activity occurs.



\textbf{Boundary}



Federation or Single Sign-On does not guarantee identity quality. Enforcement spans physical, procedural and digital controls; it is broader than a runtime policy engine.


{\small Recorded unresolved matters: source-qualified matters; exact scope remains in the canonical record.}\par
{\footnotesize \href{../../../../continuum/strategy/capability/b9cb3666f9134ba29eb896e67eba7250.md}{Canonical record: b9cb3666f9134ba29eb896e67eba7250}; full qualifications and provenance.}\par
\paragraph*{II.1.2.16.2.1 [CT3] Access Control Assurance \& Reporting}
\phantomsection
\label{navigation:II.1.2.16.2.1}
\addcontentsline{toc}{paragraph}{II.1.2.16.2.1 [CT3] Access Control Assurance \& Reporting}

\label{element:ba0bf23c2b5b4de1bdfbc5df90f57a2f}
\index{[CT3] Access Control Assurance \& Reporting (II.1.2.16.2.1)}
{\small Capability · Canonical · current}\par

Establish organisational assurance/evidence of whether access-control arrangements are appropriately established and operating.



This capability establishes organisational evidence and assurance about whether access-control arrangements are appropriately established and operating. It makes control effectiveness a distinct concern from defining requirements or performing an individual enforcement action. Reporting communicates that established position without making an isolated successful check proof of end-to-end control.



\textbf{Boundary}



Subject-specific access-control assurance does not create another macro-level Health Service Assurance hierarchy.


{\small Recorded unresolved matters: source-qualified matters; exact scope remains in the canonical record.}\par
{\footnotesize \href{../../../../continuum/strategy/capability/ba0bf23c2b5b4de1bdfbc5df90f57a2f.md}{Canonical record: ba0bf23c2b5b4de1bdfbc5df90f57a2f}; full qualifications and provenance.}\par
\paragraph*{II.1.2.16.2.2 [CT3] Access Control Collaboration}
\phantomsection
\label{navigation:II.1.2.16.2.2}
\addcontentsline{toc}{paragraph}{II.1.2.16.2.2 [CT3] Access Control Collaboration}

\label{element:c1a39b83d27a4a148f997c0a709e82ac}
\index{[CT3] Access Control Collaboration (II.1.2.16.2.2)}
{\small Capability · Canonical · current}\par

Coordinate arrangements across independently controlled domains/parties while retaining their control ownership.



Access Control Collaboration coordinates arrangements across parties or domains that retain independent control ownership. Its purpose is coherent cross-boundary control where a strong local arrangement alone cannot establish the whole result. Collaboration therefore concerns agreed control participation rather than silently centralising every participant's authority.



\textbf{Boundary}



Coordination does not transfer the parties' control ownership or guarantee identity quality across domains.


{\small Recorded unresolved matters: source-qualified matters; exact scope remains in the canonical record.}\par
{\footnotesize \href{../../../../continuum/strategy/capability/c1a39b83d27a4a148f997c0a709e82ac.md}{Canonical record: c1a39b83d27a4a148f997c0a709e82ac}; full qualifications and provenance.}\par
\paragraph*{II.1.2.16.2.3 [CT3] Access Control Context Management}
\phantomsection
\label{navigation:II.1.2.16.2.3}
\addcontentsline{toc}{paragraph}{II.1.2.16.2.3 [CT3] Access Control Context Management}

\label{element:a51f7a8045d24d0f982320025ba3a6ae}
\index{[CT3] Access Control Context Management (II.1.2.16.2.3)}
{\small Capability · Canonical · current}\par

Determine what information, functions, activities or other concerns require control, under what circumstances and how requirements are modelled.



Access Control Context Management establishes what requires control, the circumstances in which control applies and how those requirements are modelled. It gives access governance an explicit subject and context rather than reducing it to assembling runtime attributes. The responsibility can encompass information, functions and activities under physical, procedural or digital arrangements.



\textbf{Boundary}



Defining control context does not itself enforce the controls or establish that they operate effectively.


{\small Recorded unresolved matters: source-qualified matters; exact scope remains in the canonical record.}\par
{\footnotesize \href{../../../../continuum/strategy/capability/a51f7a8045d24d0f982320025ba3a6ae.md}{Canonical record: a51f7a8045d24d0f982320025ba3a6ae}; full qualifications and provenance.}\par
\paragraph*{II.1.2.16.2.4 [CT3] Access Control Enforcement}
\phantomsection
\label{navigation:II.1.2.16.2.4}
\addcontentsline{toc}{paragraph}{II.1.2.16.2.4 [CT3] Access Control Enforcement}

\label{element:b72a7b3fb7384870a9790b04d3c9e94a}
\index{[CT3] Access Control Enforcement (II.1.2.16.2.4)}
{\small Capability · Canonical · current}\par

Ensure physical, procedural and digital control mechanisms coherently impose intended controls.



Access Control Enforcement ensures that intended controls are coherently imposed through the applicable physical, procedural and digital mechanisms. It represents the organisational ability to make controls effective, rather than equating enforcement with one software decision. Its outcome depends on the controlled circumstances and the arrangements intended to apply there.



\textbf{Boundary}



A local software authorisation result alone does not establish coherent end-to-end organisational enforcement.


{\small Recorded unresolved matters: source-qualified matters; exact scope remains in the canonical record.}\par
{\footnotesize \href{../../../../continuum/strategy/capability/b72a7b3fb7384870a9790b04d3c9e94a.md}{Canonical record: b72a7b3fb7384870a9790b04d3c9e94a}; full qualifications and provenance.}\par
\paragraph*{II.1.2.16.2.5 [CT3] Identity Assurance}
\phantomsection
\label{navigation:II.1.2.16.2.5}
\addcontentsline{toc}{paragraph}{II.1.2.16.2.5 [CT3] Identity Assurance}

\label{element:56ccd6d3bf5b4ec980a333dfadf8d9fa}
\index{[CT3] Identity Assurance (II.1.2.16.2.5)}
{\small Capability · Canonical · current}\par

Establish sufficient confidence that a person is who they claim to be through appropriate evidence and processes.



Identity Assurance establishes confidence that a person is who they claim to be using appropriate evidence and processes. It concerns the basis and strength of the identity claim, rather than only successful use of an authentication mechanism. Within identity governance it makes confidence a question that can be considered separately from contextual sufficiency and access authority.



\textbf{Boundary}



An assured identity is not automatically sufficient for every context or permitted to perform every action.


{\small Recorded unresolved matters: source-qualified matters; exact scope remains in the canonical record.}\par
{\footnotesize \href{../../../../continuum/strategy/capability/56ccd6d3bf5b4ec980a333dfadf8d9fa.md}{Canonical record: 56ccd6d3bf5b4ec980a333dfadf8d9fa}; full qualifications and provenance.}\par
\paragraph*{II.1.2.16.2.6 [CT3] Identity Identification \& Authentication}
\phantomsection
\label{navigation:II.1.2.16.2.6}
\addcontentsline{toc}{paragraph}{II.1.2.16.2.6 [CT3] Identity Identification \& Authentication}

\label{element:ba4ad62548984d63ad75da3561566431}
\index{[CT3] Identity Identification \& Authentication (II.1.2.16.2.6)}
{\small Capability · Canonical · current}\par

Identify an actor and authenticate the asserted identity.



This capability identifies the acting participant and authenticates the asserted identity. It establishes who is acting under the applicable identity arrangements so that later control and accountability can be grounded in an identified actor. The outcome remains distinct from confidence established through identity evidence and from authority to access a controlled concern.



\textbf{Boundary}



Authentication neither establishes every required identity-quality condition nor authorises the requested activity.


{\small Recorded unresolved matters: source-qualified matters; exact scope remains in the canonical record.}\par
{\footnotesize \href{../../../../continuum/strategy/capability/ba4ad62548984d63ad75da3561566431.md}{Canonical record: ba4ad62548984d63ad75da3561566431}; full qualifications and provenance.}\par
\paragraph*{II.1.2.16.2.7 [CT3] Identity Modelling \& Interoperability}
\phantomsection
\label{navigation:II.1.2.16.2.7}
\addcontentsline{toc}{paragraph}{II.1.2.16.2.7 [CT3] Identity Modelling \& Interoperability}

\label{element:c40a00b0fe56405d99b6bb6ede72e18e}
\index{[CT3] Identity Modelling \& Interoperability (II.1.2.16.2.7)}
{\small Capability · Canonical · current}\par

Establish identity models and interoperability arrangements appropriate to organisational identity use.



This capability establishes identity models and arrangements through which identities can be understood across organisational use. It addresses the coherence needed when parties exchange or rely on identity information without assuming that their identity assertions have identical quality. Its purpose is meaningful interoperability, rather than a particular federation or Single Sign-On technology.



\textbf{Boundary}



Interoperability does not guarantee assurance quality or transfer each party's control ownership.


{\small Recorded unresolved matters: source-qualified matters; exact scope remains in the canonical record.}\par
{\footnotesize \href{../../../../continuum/strategy/capability/c40a00b0fe56405d99b6bb6ede72e18e.md}{Canonical record: c40a00b0fe56405d99b6bb6ede72e18e}; full qualifications and provenance.}\par
\paragraph*{II.1.2.16.2.8 [CT3] Identity Qualification Control}
\phantomsection
\label{navigation:II.1.2.16.2.8}
\addcontentsline{toc}{paragraph}{II.1.2.16.2.8 [CT3] Identity Qualification Control}

\label{element:3d53064f2fb7475b8f49fb58ec8ad4a5}
\index{[CT3] Identity Qualification Control (II.1.2.16.2.8)}
{\small Capability · Canonical · current}\par

Determine whether presented identity characteristics, provenance and assurance quality are sufficient for its intended context.



Identity Qualification Control evaluates whether presented identity characteristics, provenance and assurance are sufficient for the intended context. It adds a contextual judgement beyond whether the actor can be authenticated. This permits an identity adequate for one purpose to remain insufficient for another without claiming that the identity is necessarily false.



\textbf{Boundary}



Identity qualification does not establish practitioner clinical qualification or a general permission to act.


{\small Recorded unresolved matters: source-qualified matters; exact scope remains in the canonical record.}\par
{\footnotesize \href{../../../../continuum/strategy/capability/3d53064f2fb7475b8f49fb58ec8ad4a5.md}{Canonical record: 3d53064f2fb7475b8f49fb58ec8ad4a5}; full qualifications and provenance.}\par
\subsubsection*{II.1.2.16.3 [CT2] Health Information Protection}
\phantomsection
\label{navigation:II.1.2.16.3}
\addcontentsline{toc}{subsubsection}{II.1.2.16.3 [CT2] Health Information Protection}

\label{element:4ae0de8fa16840d09ab2ef3f5545a5e7}
\index{[CT2] Health Information Protection (II.1.2.16.3)}
{\small Capability · Canonical · current}\par

Govern privacy consequences and adequate protection of the information estate and its environment.



Health Information Protection governs the consequences of information exposure and the safeguards protecting information and its operating environment. Privacy safeguards, estate protection, access-practice monitoring, compartment assurance and classification assurance address different questions. Together they distinguish appropriate arrangements from evidence of how information is actually handled, including physical and printed handling.



\textbf{Intent}



Protect information and the people affected by its use, with an intelligible basis for evaluating protective arrangements.



\textbf{Boundary}



Compartment assurance does not decide access. Subject-specific protection assurance does not create additional macro-level Health Service Assurance trees.


{\small Recorded unresolved matters: source-qualified matters; exact scope remains in the canonical record.}\par
{\footnotesize \href{../../../../continuum/strategy/capability/4ae0de8fa16840d09ab2ef3f5545a5e7.md}{Canonical record: 4ae0de8fa16840d09ab2ef3f5545a5e7}; full qualifications and provenance.}\par
\paragraph*{II.1.2.16.3.1 [CT3] Health Information Access Monitoring}
\phantomsection
\label{navigation:II.1.2.16.3.1}
\addcontentsline{toc}{paragraph}{II.1.2.16.3.1 [CT3] Health Information Access Monitoring}

\label{element:01854c7311404bbf9f698fce4656b54d}
\index{[CT3] Health Information Access Monitoring (II.1.2.16.3.1)}
{\small Capability · Canonical · current}\par

Evaluate actual organisational access/handling practice, including controls, accessibility, compartments, classification and physical/printed handling.



Health Information Access Monitoring evaluates actual access and handling practice across organisational circumstances. Digital audit information can contribute, but so can physical access, printed handling, compartments and classification practice. Its purpose is to understand what is happening in use rather than merely observe software events or assume configured controls describe reality.



\textbf{Boundary}



Evidence supply does not transfer business accountability, and observation is not automatically independent assurance.


{\small Recorded unresolved matters: source-qualified matters; exact scope remains in the canonical record.}\par
{\footnotesize \href{../../../../continuum/strategy/capability/01854c7311404bbf9f698fce4656b54d.md}{Canonical record: 01854c7311404bbf9f698fce4656b54d}; full qualifications and provenance.}\par
\paragraph*{II.1.2.16.3.2 [CT3] Health Information Classification Assurance}
\phantomsection
\label{navigation:II.1.2.16.3.2}
\addcontentsline{toc}{paragraph}{II.1.2.16.3.2 [CT3] Health Information Classification Assurance}

\label{element:4a4a871081bf45829eab2a93c9d1c21f}
\index{[CT3] Health Information Classification Assurance (II.1.2.16.3.2)}
{\small Capability · Canonical · current}\par

Assure appropriate classification according to sensitivity and potential consequences of exposure/mishandling.



This capability evaluates whether classification is appropriate to information sensitivity and the consequences of exposure or mishandling. It examines the application of classification rather than merely the presence of a label. The responsibility supports protection decisions while retaining the distinction between the classification basis and the control applied to a particular request.



\textbf{Boundary}



Classification assurance does not define compartment membership or grant permission to access classified information.


{\small Recorded unresolved matters: source-qualified matters; exact scope remains in the canonical record.}\par
{\footnotesize \href{../../../../continuum/strategy/capability/4a4a871081bf45829eab2a93c9d1c21f.md}{Canonical record: 4a4a871081bf45829eab2a93c9d1c21f}; full qualifications and provenance.}\par
\paragraph*{II.1.2.16.3.3 [CT3] Health Information Compartment Assurance}
\phantomsection
\label{navigation:II.1.2.16.3.3}
\addcontentsline{toc}{paragraph}{II.1.2.16.3.3 [CT3] Health Information Compartment Assurance}

\label{element:be47d201de1a45a88f48234d69deb2eb}
\index{[CT3] Health Information Compartment Assurance (II.1.2.16.3.3)}
{\small Capability · Canonical · current}\par

Assure appropriate establishment and application of compartmentalisation without establishing access decisions.



This capability evaluates whether contextual information compartments are appropriately established and applied. It treats the correctness of the grouping as a concern distinct from permission to disclose a particular item. Such assurance can make compartment use more trustworthy without changing the business relationships that define membership.



\textbf{Boundary}



Assuring compartmentalisation does not establish access decisions or sensitivity classification.


{\small Recorded unresolved matters: source-qualified matters; exact scope remains in the canonical record.}\par
{\footnotesize \href{../../../../continuum/strategy/capability/be47d201de1a45a88f48234d69deb2eb.md}{Canonical record: be47d201de1a45a88f48234d69deb2eb}; full qualifications and provenance.}\par
\paragraph*{II.1.2.16.3.4 [CT3] Health Information Safeguarding}
\phantomsection
\label{navigation:II.1.2.16.3.4}
\addcontentsline{toc}{paragraph}{II.1.2.16.3.4 [CT3] Health Information Safeguarding}

\label{element:6e89db75de5b4c4c9274f78c8d9e665a}
\index{[CT3] Health Information Safeguarding (II.1.2.16.3.4)}
{\small Capability · Canonical · current}\par

Ensure information and its operational environment are protected and resilient against loss, compromise, inappropriate access or operational failure.



Health Information Safeguarding protects information and its operational environment against loss, compromise, inappropriate access and operational failure. It frames resilience and protection as organisational outcomes, allowing physical controls, recovery arrangements and testing to contribute without defining the capability as a product. The concern is broader than any individual access transaction.



\textbf{Boundary}



Estate protection remains distinct from evaluating privacy consequences and from assurance conclusions about service delivery.


{\small Recorded unresolved matters: source-qualified matters; exact scope remains in the canonical record.}\par
{\footnotesize \href{../../../../continuum/strategy/capability/6e89db75de5b4c4c9274f78c8d9e665a.md}{Canonical record: 6e89db75de5b4c4c9274f78c8d9e665a}; full qualifications and provenance.}\par
\paragraph*{II.1.2.16.3.5 [CT3] Health Privacy Safeguard Management}
\phantomsection
\label{navigation:II.1.2.16.3.5}
\addcontentsline{toc}{paragraph}{II.1.2.16.3.5 [CT3] Health Privacy Safeguard Management}

\label{element:089cd8744ed6410388daeff463507fc1}
\index{[CT3] Health Privacy Safeguard Management (II.1.2.16.3.5)}
{\small Capability · Canonical · current}\par

Establish organisational safeguards and assurance processes concerning privacy consequences of information supply, exposure or use.



This capability establishes safeguards and assurance arrangements concerning the privacy consequences of supplying, exposing or using information. It focuses on the effect on people and legitimate information use, distinguishing that concern from protection of the information estate itself. Technical access restrictions may contribute without exhausting the organisational privacy responsibility.



\textbf{Boundary}



Privacy safeguards do not replace broader estate resilience, access decisions or originating consent authority.


{\small Recorded unresolved matters: source-qualified matters; exact scope remains in the canonical record.}\par
{\footnotesize \href{../../../../continuum/strategy/capability/089cd8744ed6410388daeff463507fc1.md}{Canonical record: 089cd8744ed6410388daeff463507fc1}; full qualifications and provenance.}\par
\subsection*{II.1.2.17 [CT1] Standards, Semantics \& Reference Governance}
\phantomsection
\label{navigation:II.1.2.17}
\addcontentsline{toc}{subsection}{II.1.2.17 [CT1] Standards, Semantics \& Reference Governance}

{\small Historical predecessor: R.2.51 \hyperref[element:3b65a3d007f54aa9b875e655475e1786]{[CT1] Standards, Semantics \& Reference Governance}; Canonical ID: \texttt{\seqsplit{3b65a3d007f54aa9b875e655475e1786}}. Superseded historical identity — direct correspondence · superseded.}\par
\label{element:5ffeb7401d314f06b05d79c3aff98242}
\index{[CT1] Standards, Semantics \& Reference Governance (II.1.2.17)}
{\small Capability · Canonical · current}\par

Govern organisational adoption, oversight and appropriate use of health-information models, standards, terminology, classifications and semantic mappings.



This capability establishes the organisational standing and appropriate use of models, standards and semantic artefacts. Standards governance addresses discovery, adoption, education, compliance and impacts; clinical-semantics governance addresses terminology, classifications and mappings. It provides the governed meaning within which information can be represented and exchanged, rather than merely maintaining technical schemas.



\textbf{Intent}



Support consistent, appropriate interpretation and use of health-information standards and semantic references.



\textbf{Boundary}



Using a terminology binding, mapping or classification does not transfer its governing authority. The established peer set must not be expanded simply to create a symmetrical taxonomy.



\textbf{Harmonia Architectural Relevance}



Semantic and conformance governance is a substantive platform responsibility. Its organisational meaning constrains technical representations; implementation products do not define that meaning.


{\footnotesize \href{../../../../continuum/strategy/capability/5ffeb7401d314f06b05d79c3aff98242.md}{Canonical record: 5ffeb7401d314f06b05d79c3aff98242}; full qualifications and provenance.}\par
\subsubsection*{II.1.2.17.1 [CT2] Clinical Semantics Governance}
\phantomsection
\label{navigation:II.1.2.17.1}
\addcontentsline{toc}{subsubsection}{II.1.2.17.1 [CT2] Clinical Semantics Governance}

\label{element:5cc8319e1dde4c818cc248a47233b67b}
\index{[CT2] Clinical Semantics Governance (II.1.2.17.1)}
{\small Capability · Canonical · current}\par

Govern clinical meaning and appropriate organisational use of semantic artefacts.



Clinical Semantics Governance maintains the standing and appropriate use of terminology, classifications and semantic mappings. Terminology provides consistent meaning; classification concerns sensitivity and consequences of exposure or mishandling. The established children distinguish governing artefacts from overseeing actual use and permit controlled terminology localisation where adopted material does not meet legitimate needs.



\textbf{Intent}



Keep meaning and classification intelligible across organisational use without concealing semantic drift or local variation.



\textbf{Boundary}



Compartmentalisation concerns contextual relationships rather than classification or terminology. Classification Localisation is not an established peer; mappings must preserve the separate meanings and authorities they connect.


{\small Recorded unresolved matters: source-qualified matters; exact scope remains in the canonical record.}\par
{\footnotesize \href{../../../../continuum/strategy/capability/5cc8319e1dde4c818cc248a47233b67b.md}{Canonical record: 5cc8319e1dde4c818cc248a47233b67b}; full qualifications and provenance.}\par
\paragraph*{II.1.2.17.1.1 [CT3] Classification Governance}
\phantomsection
\label{navigation:II.1.2.17.1.1}
\addcontentsline{toc}{paragraph}{II.1.2.17.1.1 [CT3] Classification Governance}

\label{element:4520ce37454e4de68580d75fdb08efe6}
\index{[CT3] Classification Governance (II.1.2.17.1.1)}
{\small Capability · Canonical · current}\par

Govern applicable classification definitions and their organisational standing.



Classification Governance establishes the definitions and organisational standing of applicable classifications. In this governed context, classification concerns sensitivity and the consequences of exposure or mishandling, rather than the clinical meaning supplied by terminology. Its outcome supplies a basis for consistent application and protection considerations.



\textbf{Boundary}



Classification is orthogonal to contextual compartment membership and does not by itself determine an access decision.


{\small Recorded unresolved matters: source-qualified matters; exact scope remains in the canonical record.}\par
{\footnotesize \href{../../../../continuum/strategy/capability/4520ce37454e4de68580d75fdb08efe6.md}{Canonical record: 4520ce37454e4de68580d75fdb08efe6}; full qualifications and provenance.}\par
\paragraph*{II.1.2.17.1.2 [CT3] Classification Utilisation Oversight}
\phantomsection
\label{navigation:II.1.2.17.1.2}
\addcontentsline{toc}{paragraph}{II.1.2.17.1.2 [CT3] Classification Utilisation Oversight}

\label{element:1d7c50aa80464613864dd6efb3fb56cc}
\index{[CT3] Classification Utilisation Oversight (II.1.2.17.1.2)}
{\small Capability · Canonical · current}\par

Oversee appropriate organisational application of classifications.



Classification Utilisation Oversight evaluates appropriate organisational application of classifications. It addresses whether information is treated consistently with the governed sensitivity meaning, rather than assuming that an established classification scheme guarantees correct use. This retains a distinct responsibility for practice alongside definition governance.



\textbf{Boundary}



Oversight neither creates classification definitions nor establishes a new localisation peer.


{\small Recorded unresolved matters: source-qualified matters; exact scope remains in the canonical record.}\par
{\footnotesize \href{../../../../continuum/strategy/capability/1d7c50aa80464613864dd6efb3fb56cc.md}{Canonical record: 1d7c50aa80464613864dd6efb3fb56cc}; full qualifications and provenance.}\par
\paragraph*{II.1.2.17.1.3 [CT3] Semantic Mapping Governance}
\phantomsection
\label{navigation:II.1.2.17.1.3}
\addcontentsline{toc}{paragraph}{II.1.2.17.1.3 [CT3] Semantic Mapping Governance}

\label{element:17dc6ab9f4394049a5327415d19628f2}
\index{[CT3] Semantic Mapping Governance (II.1.2.17.1.3)}
{\small Capability · Canonical · current}\par

Govern mappings between applicable semantic artefacts without collapsing their source meanings or authorities.



Semantic Mapping Governance maintains mappings between applicable semantic artefacts while preserving the meaning and authority of each side. Its purpose is intelligible correspondence rather than assumed equivalence. A governed mapping can support exchange or interpretation without suggesting that the mapped artefacts have become the same concept or share one authority.



\textbf{Boundary}



A mapping does not erase source distinctions or give a consuming system authority to redefine either source.


{\small Recorded unresolved matters: source-qualified matters; exact scope remains in the canonical record.}\par
{\footnotesize \href{../../../../continuum/strategy/capability/17dc6ab9f4394049a5327415d19628f2.md}{Canonical record: 17dc6ab9f4394049a5327415d19628f2}; full qualifications and provenance.}\par
\paragraph*{II.1.2.17.1.4 [CT3] Terminology Governance}
\phantomsection
\label{navigation:II.1.2.17.1.4}
\addcontentsline{toc}{paragraph}{II.1.2.17.1.4 [CT3] Terminology Governance}

\label{element:e926a47881fa482088859f59fa6fe7e4}
\index{[CT3] Terminology Governance (II.1.2.17.1.4)}
{\small Capability · Canonical · current}\par

Manage governed terminology/reference-set lifecycle and standing, including acceptance, release, update, supersession/retirement and authority/provider coordination.



Terminology Governance maintains the lifecycle and standing of terminology and reference sets used by the organisation. Acceptance, release, updates and retirement require coordination with the appropriate authority or provider so that a code's meaning is not treated as timeless or locally self-defined. This supplies a governed semantic basis for information bindings and knowledge associations.



\textbf{Boundary}



Governing terminology differs from binding a structure to it or overseeing its actual use.


{\small Recorded unresolved matters: source-qualified matters; exact scope remains in the canonical record.}\par
{\footnotesize \href{../../../../continuum/strategy/capability/e926a47881fa482088859f59fa6fe7e4.md}{Canonical record: e926a47881fa482088859f59fa6fe7e4}; full qualifications and provenance.}\par
\paragraph*{II.1.2.17.1.5 [CT3] Terminology Localisation}
\phantomsection
\label{navigation:II.1.2.17.1.5}
\addcontentsline{toc}{paragraph}{II.1.2.17.1.5 [CT3] Terminology Localisation}

\label{element:3aa80e1d84ea4e009b149c4f9a23f0f8}
\index{[CT3] Terminology Localisation (II.1.2.17.1.5)}
{\small Capability · Canonical · current}\par

Manage controlled local responses where adopted terminology does not meet legitimate organisational semantic requirements.



Terminology Localisation manages controlled local responses where adopted terminology does not satisfy legitimate organisational semantic needs. It makes the local variation explicit rather than allowing ad hoc usage to conceal a gap or change in meaning. Within clinical-semantics governance, this preserves a distinction between governed local adaptation and uncontrolled drift.



\textbf{Boundary}



The established localisation responsibility concerns terminology; no corresponding Classification Localisation capability is inferred.


{\small Recorded unresolved matters: source-qualified matters; exact scope remains in the canonical record.}\par
{\footnotesize \href{../../../../continuum/strategy/capability/3aa80e1d84ea4e009b149c4f9a23f0f8.md}{Canonical record: 3aa80e1d84ea4e009b149c4f9a23f0f8}; full qualifications and provenance.}\par
\paragraph*{II.1.2.17.1.6 [CT3] Terminology Utilisation Oversight}
\phantomsection
\label{navigation:II.1.2.17.1.6}
\addcontentsline{toc}{paragraph}{II.1.2.17.1.6 [CT3] Terminology Utilisation Oversight}

\label{element:4ca436cd90b745daa290b067c57a0450}
\index{[CT3] Terminology Utilisation Oversight (II.1.2.17.1.6)}
{\small Capability · Canonical · current}\par

Oversee actual terminology/reference-set use and identify inappropriate, inconsistent or drifting usage.



Terminology Utilisation Oversight considers how terminology and reference sets are actually used. It identifies inappropriate application, inconsistency or semantic drift that the existence of an approved terminology does not prevent. Its purpose is to preserve coherent use without making each consuming capability the owner of terminology meaning.



\textbf{Boundary}



Use oversight does not redefine terminology or automatically authorise a local semantic variation.


{\small Recorded unresolved matters: source-qualified matters; exact scope remains in the canonical record.}\par
{\footnotesize \href{../../../../continuum/strategy/capability/4ca436cd90b745daa290b067c57a0450.md}{Canonical record: 4ca436cd90b745daa290b067c57a0450}; full qualifications and provenance.}\par
\subsubsection*{II.1.2.17.2 [CT2] Health Information Standards Governance}
\phantomsection
\label{navigation:II.1.2.17.2}
\addcontentsline{toc}{subsubsection}{II.1.2.17.2 [CT2] Health Information Standards Governance}

\label{element:b228039d8e764f9b98cb810b68af329d}
\index{[CT2] Health Information Standards Governance (II.1.2.17.2)}
{\small Capability · Canonical · current}\par

Govern organisational adoption and use of health-information models and applicable standards.



This capability governs which information models and standards the organisation adopts and how their use is supported and evaluated. The model distinguishes discovering candidates, adopting models, educating users, addressing conformity and overseeing governance. Standards impact identification has a separate established place because its implications may extend into clinical practice, policy and obligations beyond information-system implementation.



\textbf{Intent}



Make organisational standards use deliberate and its consequences understandable.



\textbf{Boundary}



Oversight evaluates governance without performing every peer activity. Adoption includes its own model-impact consideration; the listed abilities do not establish a sequential adoption Process.


{\small Recorded unresolved matters: source-qualified matters; exact scope remains in the canonical record.}\par
{\footnotesize \href{../../../../continuum/strategy/capability/b228039d8e764f9b98cb810b68af329d.md}{Canonical record: b228039d8e764f9b98cb810b68af329d}; full qualifications and provenance.}\par
\paragraph*{II.1.2.17.2.1 [CT3] Information Model Adoption}
\phantomsection
\label{navigation:II.1.2.17.2.1}
\addcontentsline{toc}{paragraph}{II.1.2.17.2.1 [CT3] Information Model Adoption}

\label{element:0c3ce0778d6342b6b100e1e30e777fbe}
\index{[CT3] Information Model Adoption (II.1.2.17.2.1)}
{\small Capability · Canonical · current}\par

Adopt applicable information models, including intrinsic assessment of adoption impacts.



Information Model Adoption establishes the organisation's use of an applicable information model, including consideration of the consequences of adopting it. It gives a discovered model an organisational standing distinct from mere availability. The intrinsic impact consideration helps adoption remain a business governance responsibility rather than a mechanical import of definitions.



\textbf{Boundary}



Adoption does not perform every education, compliance or governance-oversight responsibility associated with subsequent use.


{\small Recorded unresolved matters: source-qualified matters; exact scope remains in the canonical record.}\par
{\footnotesize \href{../../../../continuum/strategy/capability/0c3ce0778d6342b6b100e1e30e777fbe.md}{Canonical record: 0c3ce0778d6342b6b100e1e30e777fbe}; full qualifications and provenance.}\par
\paragraph*{II.1.2.17.2.2 [CT3] Information Model Compliance}
\phantomsection
\label{navigation:II.1.2.17.2.2}
\addcontentsline{toc}{paragraph}{II.1.2.17.2.2 [CT3] Information Model Compliance}

\label{element:160db14dbe5149db91223a93f6ee9102}
\index{[CT3] Information Model Compliance (II.1.2.17.2.2)}
{\small Capability · Canonical · current}\par

Address organisational conformity to applicable information-model requirements.



Information Model Compliance addresses conformity to applicable model requirements. It makes actual organisational use accountable to the adopted semantic frame rather than treating model availability or user education as proof of conformity. Within standards governance it concerns the requirements for use, not the separate task of discovering or adopting models.



\textbf{Boundary}



Conformity does not redefine the model or absorb governance oversight of the wider adoption arrangements.


{\small Recorded unresolved matters: source-qualified matters; exact scope remains in the canonical record.}\par
{\footnotesize \href{../../../../continuum/strategy/capability/160db14dbe5149db91223a93f6ee9102.md}{Canonical record: 160db14dbe5149db91223a93f6ee9102}; full qualifications and provenance.}\par
\paragraph*{II.1.2.17.2.3 [CT3] Information Model Discovery}
\phantomsection
\label{navigation:II.1.2.17.2.3}
\addcontentsline{toc}{paragraph}{II.1.2.17.2.3 [CT3] Information Model Discovery}

\label{element:f33dc5b8ba9643329bb7b9431e1f9a6f}
\index{[CT3] Information Model Discovery (II.1.2.17.2.3)}
{\small Capability · Canonical · current}\par

Identify information models relevant to the organisation’s health-information needs.



Information Model Discovery identifies candidate models relevant to the organisation's health-information needs. It supports informed standards governance by making applicable possibilities visible before adoption. Its business purpose is finding suitable semantic frames rather than selecting a software representation or assuming every discovered model should be used.



\textbf{Boundary}



Discovery does not adopt a model, establish compliance or grant it organisational authority.


{\small Recorded unresolved matters: source-qualified matters; exact scope remains in the canonical record.}\par
{\footnotesize \href{../../../../continuum/strategy/capability/f33dc5b8ba9643329bb7b9431e1f9a6f.md}{Canonical record: f33dc5b8ba9643329bb7b9431e1f9a6f}; full qualifications and provenance.}\par
\paragraph*{II.1.2.17.2.4 [CT3] Information Model Education}
\phantomsection
\label{navigation:II.1.2.17.2.4}
\addcontentsline{toc}{paragraph}{II.1.2.17.2.4 [CT3] Information Model Education}

\label{element:7dcb6427f9134fb7a4a63a5ec1d5c85b}
\index{[CT3] Information Model Education (II.1.2.17.2.4)}
{\small Capability · Canonical · current}\par

Support organisational understanding and appropriate use of adopted information models.



Information Model Education supports understanding and appropriate use of adopted models. It addresses the people and practice dimension of semantic consistency: a published model is not necessarily a model understood in the contexts where it is applied. Its responsibility therefore complements adoption and compliance without becoming either.



\textbf{Boundary}



Educational support neither changes the authoritative model nor establishes that actual use complies with it.


{\small Recorded unresolved matters: source-qualified matters; exact scope remains in the canonical record.}\par
{\footnotesize \href{../../../../continuum/strategy/capability/7dcb6427f9134fb7a4a63a5ec1d5c85b.md}{Canonical record: 7dcb6427f9134fb7a4a63a5ec1d5c85b}; full qualifications and provenance.}\par
\paragraph*{II.1.2.17.2.5 [CT3] Information Model Governance Oversight}
\phantomsection
\label{navigation:II.1.2.17.2.5}
\addcontentsline{toc}{paragraph}{II.1.2.17.2.5 [CT3] Information Model Governance Oversight}

\label{element:1e71b707beab450f9cbfaf0bf88b6ec4}
\index{[CT3] Information Model Governance Oversight (II.1.2.17.2.5)}
{\small Capability · Canonical · current}\par

Oversee and evaluate effectiveness, consistency and appropriate operation of information-model governance.



This capability evaluates the effectiveness and consistency of information-model governance itself. It asks whether discovery, adoption and use are being governed appropriately, rather than repeating each operational peer. Its purpose is coherent oversight of the arrangements that give models standing and guide their application.



\textbf{Boundary}



Oversight does not itself perform discovery, adoption, education or compliance activities.


{\small Recorded unresolved matters: source-qualified matters; exact scope remains in the canonical record.}\par
{\footnotesize \href{../../../../continuum/strategy/capability/1e71b707beab450f9cbfaf0bf88b6ec4.md}{Canonical record: 1e71b707beab450f9cbfaf0bf88b6ec4}; full qualifications and provenance.}\par
\paragraph*{II.1.2.17.2.6 [CT3] Standards Adoption}
\phantomsection
\label{navigation:II.1.2.17.2.6}
\addcontentsline{toc}{paragraph}{II.1.2.17.2.6 [CT3] Standards Adoption}

\label{element:0d59f956f5714d3985398f09aeb7c541}
\index{[CT3] Standards Adoption (II.1.2.17.2.6)}
{\small Capability · Canonical · current}\par

Adopt applicable standards within the organisation.



Standards Adoption establishes which applicable standards the organisation adopts. It gives standards an intentional place in organisational practice rather than allowing familiar products or conventions to define requirements by default. The business responsibility concerns adoption in context, including standards whose effects extend beyond information representation.



\textbf{Boundary}



Adoption does not establish a complete standards lifecycle taxonomy or perform the distinct identification of broader standards impacts.


{\small Recorded unresolved matters: source-qualified matters; exact scope remains in the canonical record.}\par
{\footnotesize \href{../../../../continuum/strategy/capability/0d59f956f5714d3985398f09aeb7c541.md}{Canonical record: 0d59f956f5714d3985398f09aeb7c541}; full qualifications and provenance.}\par
\paragraph*{II.1.2.17.2.7 [CT3] Standards Impact Identification}
\phantomsection
\label{navigation:II.1.2.17.2.7}
\addcontentsline{toc}{paragraph}{II.1.2.17.2.7 [CT3] Standards Impact Identification}

\label{element:6b9c96c7d9174e7fbfd11267673d42f3}
\index{[CT3] Standards Impact Identification (II.1.2.17.2.7)}
{\small Capability · Canonical · current}\par

Identify standards impacts across clinical/work practice, permitted medications, procedures, policy, safety, obligations and reporting as applicable.



Standards Impact Identification recognises the consequences of standards for relevant organisational activity. Impacts may concern clinical practice, medications, procedures, policy, safety, obligations or reporting, so the responsibility cannot be reduced to technical compatibility. It makes those consequences visible for the responsible capabilities to consider within their own authority.



\textbf{Boundary}



Identifying an impact neither adopts the standard nor acquires the operational authority needed to respond to it.


{\small Recorded unresolved matters: source-qualified matters; exact scope remains in the canonical record.}\par
{\footnotesize \href{../../../../continuum/strategy/capability/6b9c96c7d9174e7fbfd11267673d42f3.md}{Canonical record: 6b9c96c7d9174e7fbfd11267673d42f3}; full qualifications and provenance.}\par
\subsection*{II.1.2.18 [CT1] Workforce \& Organisational Capability}
\phantomsection
\label{navigation:II.1.2.18}
\addcontentsline{toc}{subsection}{II.1.2.18 [CT1] Workforce \& Organisational Capability}

{\small Historical predecessor: R.2.52 \hyperref[element:b1161fe1a315420fbc7d048a1ea58126]{[CT1] Workforce \& Organisational Capability}; Canonical ID: \texttt{\seqsplit{b1161fe1a315420fbc7d048a1ea58126}}. Superseded historical identity — direct correspondence · superseded.}\par
\label{element:631b43664bfb4c50ac8cf4958b2c347d}
\index{[CT1] Workforce \& Organisational Capability (II.1.2.18)}
{\small Capability · Canonical · current}\par

Govern practitioner qualification, engagement/privilege and enterprise service coverage, together with organisation identity/structure, recognised locations and identifiable healthcare service offerings.



This capability governs the workforce and organisational conditions under which health services can operate. Its workforce branch addresses professional standing, participation, privilege and enterprise coverage; its structural branch addresses organisations, recognised locations and identifiable service offerings. Directories can represent these facts, but their representation is separate from business authority to establish or change them.



\textbf{Intent}



Make the organisational and workforce context of care explicit enough to constrain legitimate participation and service availability.



\textbf{Boundary}



Qualification, engagement, privilege, service coverage and deployment are distinct. Organisation or location association does not substitute for material Healthcare Service context.



\textbf{Harmonia Architectural Relevance}



These facts materially constrain care enablement, routing and operational context. Directory support preserves the originating workforce and organisational governance rather than making it platform-owned.


{\footnotesize \href{../../../../continuum/strategy/capability/631b43664bfb4c50ac8cf4958b2c347d.md}{Canonical record: 631b43664bfb4c50ac8cf4958b2c347d}; full qualifications and provenance.}\par
\subsubsection*{II.1.2.18.1 [CT2] Health Workforce Management}
\phantomsection
\label{navigation:II.1.2.18.1}
\addcontentsline{toc}{subsubsection}{II.1.2.18.1 [CT2] Health Workforce Management}

\label{element:7ec329cdc070472eb68966345d82031a}
\index{[CT2] Health Workforce Management (II.1.2.18.1)}
{\small Capability · Canonical · current}\par

Govern practitioners’ professional standing, engagement/privilege and whole-enterprise service coverage.



Health Workforce Management governs practitioners' professional standing, authorised participation and the ability to cover required enterprise work. Its three children separate qualification from engagement and privilege, and both from whole-enterprise coverage. This prevents an available or rostered person from being treated as automatically qualified and privileged for a particular activity.



\textbf{Intent}



Support legitimate practitioner participation and an informed understanding of service-coverage gaps and risks.



\textbf{Boundary}



Coverage is broader than rostering. Applying an available resource to a care requirement and progressing a particular support work item remain separate operational responsibilities.


{\small Recorded unresolved matters: source-qualified matters; exact scope remains in the canonical record.}\par
{\footnotesize \href{../../../../continuum/strategy/capability/7ec329cdc070472eb68966345d82031a.md}{Canonical record: 7ec329cdc070472eb68966345d82031a}; full qualifications and provenance.}\par
\paragraph*{II.1.2.18.1.1 [CT3] Practitioner Engagement \& Privilege Management}
\phantomsection
\label{navigation:II.1.2.18.1.1}
\addcontentsline{toc}{paragraph}{II.1.2.18.1.1 [CT3] Practitioner Engagement \& Privilege Management}

\label{element:36d5bf4c19d34172b72f4b77099ff235}
\index{[CT3] Practitioner Engagement \& Privilege Management (II.1.2.18.1.1)}
{\small Capability · Canonical · current}\par

Manage practitioners as service participants, their roles/capacities, applicable privileges and qualification context above ordinary HR employment administration.



This capability manages how practitioners participate in services and the privileges applicable to that participation. It connects professional standing with roles and permitted capacities in the service context, beyond ordinary employment administration. Separating engagement and privilege from qualification makes it possible to recognise a qualified practitioner who is not authorised for a particular service activity.



\textbf{Boundary}



Representation or consumption of privilege information does not transfer the underlying workforce decision authority to Harmonia.


{\small Recorded unresolved matters: source-qualified matters; exact scope remains in the canonical record.}\par
{\footnotesize \href{../../../../continuum/strategy/capability/36d5bf4c19d34172b72f4b77099ff235.md}{Canonical record: 36d5bf4c19d34172b72f4b77099ff235}; full qualifications and provenance.}\par
\paragraph*{II.1.2.18.1.2 [CT3] Practitioner Qualification Management}
\phantomsection
\label{navigation:II.1.2.18.1.2}
\addcontentsline{toc}{paragraph}{II.1.2.18.1.2 [CT3] Practitioner Qualification Management}

\label{element:279bab5de1b74c17a5ee2db5cae8d385}
\index{[CT3] Practitioner Qualification Management (II.1.2.18.1.2)}
{\small Capability · Canonical · current}\par

Govern professional/clinical standing and qualifications, including applicable review, escalation, complaints, penalties/issues, suspension and revocation.



Practitioner Qualification Management governs professional and clinical standing, including review and changes such as suspension or revocation. It makes qualification a continuing governed concern rather than a static directory attribute. That standing can constrain engagement and privilege while remaining distinguishable from either an employment relationship or permission for a particular activity.



\textbf{Boundary}



Qualification alone does not establish engagement, clinical privilege, device-specific enablement or operational assignment.


{\small Recorded unresolved matters: source-qualified matters; exact scope remains in the canonical record.}\par
{\footnotesize \href{../../../../continuum/strategy/capability/279bab5de1b74c17a5ee2db5cae8d385.md}{Canonical record: 279bab5de1b74c17a5ee2db5cae8d385}; full qualifications and provenance.}\par
\paragraph*{II.1.2.18.1.3 [CT3] Workforce Service Coverage Management}
\phantomsection
\label{navigation:II.1.2.18.1.3}
\addcontentsline{toc}{paragraph}{II.1.2.18.1.3 [CT3] Workforce Service Coverage Management}

\label{element:cfb04e13895745619171a3e8f15d349c}
\index{[CT3] Workforce Service Coverage Management (II.1.2.18.1.3)}
{\small Capability · Canonical · current}\par

Observe whether required enterprise work/service obligations can be fulfilled, identify coverage gaps/risks and adjust coverage in response.



Workforce Service Coverage Management considers whether required enterprise work and service obligations can be fulfilled and how coverage should change when gaps or risks arise. Its concern is the adequacy of coverage, broader than the existence of a roster. This supplies a service-wide workforce view distinct from deploying a person to a particular care requirement.



\textbf{Boundary}



Coverage management is not individual support-work progression or automatic authority to assign clinical work.


{\small Recorded unresolved matters: source-qualified matters; exact scope remains in the canonical record.}\par
{\footnotesize \href{../../../../continuum/strategy/capability/cfb04e13895745619171a3e8f15d349c.md}{Canonical record: cfb04e13895745619171a3e8f15d349c}; full qualifications and provenance.}\par
\subsubsection*{II.1.2.18.2 [CT2] Healthcare Organisation \& Service Structure Management}
\phantomsection
\label{navigation:II.1.2.18.2}
\addcontentsline{toc}{subsubsection}{II.1.2.18.2 [CT2] Healthcare Organisation \& Service Structure Management}

\label{element:de8c1b408c774e59899662495f4d6ced}
\index{[CT2] Healthcare Organisation \& Service Structure Management (II.1.2.18.2)}
{\small Capability · Canonical · current}\par

Manage organisational identity/structure, recognised locations and identifiable organisational service offerings.



This capability establishes the organisation, location and service-offering context through which health services are recognised. Its children distinguish the identity and structure of participants, the definition of places, and what identifiable services are offered under temporal and conditional availability constraints. The distinction allows a service to exist without assuming it is offered everywhere or currently available.



\textbf{Intent}



Provide an intelligible account of who offers what service, where and under which conditions.



\textbf{Boundary}



Location definition is not facilities operation. Offering availability does not establish that all resources and dependencies are collectively ready to deliver the service.


{\small Recorded unresolved matters: source-qualified matters; exact scope remains in the canonical record.}\par
{\footnotesize \href{../../../../continuum/strategy/capability/de8c1b408c774e59899662495f4d6ced.md}{Canonical record: de8c1b408c774e59899662495f4d6ced}; full qualifications and provenance.}\par
\paragraph*{II.1.2.18.2.1 [CT3] Healthcare Location Definition Management}
\phantomsection
\label{navigation:II.1.2.18.2.1}
\addcontentsline{toc}{paragraph}{II.1.2.18.2.1 [CT3] Healthcare Location Definition Management}

\label{element:d77a5c8f1818465c98125acc4bc2ac81}
\index{[CT3] Healthcare Location Definition Management (II.1.2.18.2.1)}
{\small Capability · Canonical · current}\par

Establish and manage recognised locations and their applicable business characteristics and relationships.



Healthcare Location Definition Management establishes recognised places and their business characteristics and relationships. It gives health-service activity an intelligible location context while remaining separate from maintaining the physical facilities. A defined location can be referenced in offerings and operational activity without implying that it is currently usable or ready for care.



\textbf{Boundary}



Location definition is not facilities operation, care-place readiness or service availability.


{\small Recorded unresolved matters: source-qualified matters; exact scope remains in the canonical record.}\par
{\footnotesize \href{../../../../continuum/strategy/capability/d77a5c8f1818465c98125acc4bc2ac81.md}{Canonical record: d77a5c8f1818465c98125acc4bc2ac81}; full qualifications and provenance.}\par
\paragraph*{II.1.2.18.2.2 [CT3] Healthcare Organisation Identity \& Structure Management}
\phantomsection
\label{navigation:II.1.2.18.2.2}
\addcontentsline{toc}{paragraph}{II.1.2.18.2.2 [CT3] Healthcare Organisation Identity \& Structure Management}

\label{element:4684348b4d144e61aa7dd203facfebdd}
\index{[CT3] Healthcare Organisation Identity \& Structure Management (II.1.2.18.2.2)}
{\small Capability · Canonical · current}\par

Establish and manage identity and structural relationships of participating organisations and organisational units.



This capability establishes the identity and structural relationships of participating organisations and organisational units. It provides the business account of who the participants are and how they relate, which directory representations can then convey. The distinction matters because a represented hierarchy should not create organisational authority merely through document or interface nesting.



\textbf{Boundary}



Organisation identity does not substitute for location, practitioner or material Healthcare Service context.


{\small Recorded unresolved matters: source-qualified matters; exact scope remains in the canonical record.}\par
{\footnotesize \href{../../../../continuum/strategy/capability/4684348b4d144e61aa7dd203facfebdd.md}{Canonical record: 4684348b4d144e61aa7dd203facfebdd}; full qualifications and provenance.}\par
\paragraph*{II.1.2.18.2.3 [CT3] Healthcare Service Offering Management}
\phantomsection
\label{navigation:II.1.2.18.2.3}
\addcontentsline{toc}{paragraph}{II.1.2.18.2.3 [CT3] Healthcare Service Offering Management}

\label{element:9fe177e4422b4853acb1dd0d0b74c4ba}
\index{[CT3] Healthcare Service Offering Management (II.1.2.18.2.3)}
{\small Capability · Canonical · current}\par

Establish/manage what identifiable service is offered, by whom, where, when and under which availability conditions/constraints.



This capability establishes the identifiable service offered, including provider, place, timing and availability conditions. It makes a service discoverable as an offering rather than merely an abstract service idea. Temporal and conditional availability are material: an offering may exist but be unavailable in the present context, and another pathway may cover different hours.



\textbf{Boundary}



An available offering does not establish collective readiness to deliver it; organisation and location alone do not supply the service context.


{\small Recorded unresolved matters: source-qualified matters; exact scope remains in the canonical record.}\par
{\footnotesize \href{../../../../continuum/strategy/capability/9fe177e4422b4853acb1dd0d0b74c4ba.md}{Canonical record: 9fe177e4422b4853acb1dd0d0b74c4ba}; full qualifications and provenance.}\par
\section*{II.1.3 Business Enabling Capabilities}
\phantomsection
\label{navigation:II.1.3}
\addcontentsline{toc}{section}{II.1.3 Business Enabling Capabilities}

\subsection*{II.1.3.1 [CT1] Care Administration Enablement}
\phantomsection
\label{navigation:II.1.3.1}
\addcontentsline{toc}{subsection}{II.1.3.1 [CT1] Care Administration Enablement}

{\small Historical predecessor: R.2.71 \hyperref[element:b611a1fcfadb4b30af9ba74a4c993d8a]{[CT1] Care Administration Enablement}; Canonical ID: \texttt{\seqsplit{b611a1fcfadb4b30af9ba74a4c993d8a}}. Superseded historical identity — direct correspondence · superseded.}\par
\label{element:28df27d55c514a37995145ab46945f78}
\index{[CT1] Care Administration Enablement (II.1.3.1)}
{\small Capability · Proposed · current}\par

Administer the information and progression of care access, requests, encounters, records and plans across responsible care systems.



\textbf{Overview}



Administer the information and progression of care access, requests, encounters, records and plans across responsible care systems.



\textbf{Intent}



Administer the information and progression of care access, requests, encounters, records and plans across responsible care systems.



\textbf{Boundary}



Information administration does not perform the clinical act, own host booking engines or infer Transfer of Care from referral acceptance.


{\footnotesize \href{../../../../continuum/strategy/capability/28df27d55c514a37995145ab46945f78.md}{Canonical record: 28df27d55c514a37995145ab46945f78}; full qualifications and provenance.}\par
\subsubsection*{II.1.3.1.1 [CT2] Care Access Transaction Enablement [superseded]}
\phantomsection
\label{navigation:II.1.3.1.1}
\addcontentsline{toc}{subsubsection}{II.1.3.1.1 [CT2] Care Access Transaction Enablement [superseded]}

\label{element:b65973b576154d3a96cf458bf33df53a}
\index{[CT2] Care Access Transaction Enablement (II.1.3.1.1)}
{\small Capability · Proposed · superseded}\par
\href{../../../../continuum/strategy/capability/b65973b576154d3a96cf458bf33df53a.md}{Historical canonical definition and provenance}.\par
\subsubsection*{II.1.3.1.2 [CT2] Care Documentation \& Communication Enablement [superseded]}
\phantomsection
\label{navigation:II.1.3.1.2}
\addcontentsline{toc}{subsubsection}{II.1.3.1.2 [CT2] Care Documentation \& Communication Enablement [superseded]}

\label{element:5ad2e381ec484d1d9062a65eee842fa7}
\index{[CT2] Care Documentation \& Communication Enablement (II.1.3.1.2)}
{\small Capability · Proposed · superseded}\par
\href{../../../../continuum/strategy/capability/5ad2e381ec484d1d9062a65eee842fa7.md}{Historical canonical definition and provenance}.\par
\subsubsection*{II.1.3.1.3 [CT2] Care Financial Information Enablement [superseded]}
\phantomsection
\label{navigation:II.1.3.1.3}
\addcontentsline{toc}{subsubsection}{II.1.3.1.3 [CT2] Care Financial Information Enablement [superseded]}

\label{element:59a756b89c604e3e97ee5dede2d6ed89}
\index{[CT2] Care Financial Information Enablement (II.1.3.1.3)}
{\small Capability · Proposed · superseded}\par
\href{../../../../continuum/strategy/capability/59a756b89c604e3e97ee5dede2d6ed89.md}{Historical canonical definition and provenance}.\par
\subsubsection*{II.1.3.1.4 [CT2] Care Fulfilment Administration Enablement [superseded]}
\phantomsection
\label{navigation:II.1.3.1.4}
\addcontentsline{toc}{subsubsection}{II.1.3.1.4 [CT2] Care Fulfilment Administration Enablement [superseded]}

\label{element:9f0f4379c1954524ba5e691d0c7d95ab}
\index{[CT2] Care Fulfilment Administration Enablement (II.1.3.1.4)}
{\small Capability · Proposed · superseded}\par
\href{../../../../continuum/strategy/capability/9f0f4379c1954524ba5e691d0c7d95ab.md}{Historical canonical definition and provenance}.\par
\subsubsection*{II.1.3.1.5 [CT2] Care Plan Administration}
\phantomsection
\label{navigation:II.1.3.1.5}
\addcontentsline{toc}{subsubsection}{II.1.3.1.5 [CT2] Care Plan Administration}

{\small Historical predecessor: R.2.54 \hyperref[element:e31407cea7624632a8ed2dc53b09d726]{[CT2] Care Plan Administration}; Canonical ID: \texttt{\seqsplit{e31407cea7624632a8ed2dc53b09d726}}. Superseded historical identity — direct correspondence · superseded.}\par
\label{element:cf06eec5699f406fa30a00d2f76e3098}
\index{[CT2] Care Plan Administration (II.1.3.1.5)}
{\small Capability · Canonical · current}\par

Care Plan Administration receives care plans and coordinates the distribution of updates to authorised team members. It supports shared understanding of established clinical intent across participating organisations. The responsibility concerns the plan as managed information, distinct from deciding the intended course of care or performing the actions expressed by it.



\textbf{Overview}



Care Plan Administration receives care plans and coordinates the distribution of updates to authorised team members. It supports shared understanding of established clinical intent across participating organisations. The responsibility concerns the plan as managed information, distinct from deciding the intended course of care or performing the actions expressed by it.



\textbf{Intent}



Care Plan Administration receives care plans and coordinates the distribution of updates to authorised team members. It supports shared understanding of established clinical intent across participating organisations. The responsibility concerns the plan as managed information, distinct from deciding the intended course of care or performing the actions expressed by it.



\textbf{Boundary}



Information administration does not confer authority over the represented person, organisation, service, place, device or clinical activity.



\medskip\hrule\medskip

{\footnotesize \href{../../../../continuum/strategy/capability/cf06eec5699f406fa30a00d2f76e3098.md}{Canonical record: cf06eec5699f406fa30a00d2f76e3098}; full qualifications and provenance.}\par
\paragraph*{II.1.3.1.5.1 [CT3] Care Plan Lifecycle \& Coordination}
\phantomsection
\label{navigation:II.1.3.1.5.1}
\addcontentsline{toc}{paragraph}{II.1.3.1.5.1 [CT3] Care Plan Lifecycle \& Coordination}

{\small Historical predecessor: R.2.53 \hyperref[element:689bb3f8a24249b79f787e1f48eb268d]{[CT3] Care Plan Lifecycle \& Coordination}; Canonical ID: \texttt{\seqsplit{689bb3f8a24249b79f787e1f48eb268d}}. Superseded historical identity — direct correspondence · superseded.}\par
\label{element:4cb7db65392f4fab801d7794af4d4b29}
\index{[CT3] Care Plan Lifecycle \& Coordination (II.1.3.1.5.1)}
{\small Capability · Proposed · current}\par

Administer received care plans, goals, revisions and cross-participant coordination while retaining clinical planning authority at its source.



\textbf{Overview}



Administer received care plans, goals, revisions and cross-participant coordination while retaining clinical planning authority at its source. The reviewed hierarchy assigns this enduring ability to its stated business-enabling context.



\textbf{Intent}



Administer received care plans, goals, revisions and cross-participant coordination while retaining clinical planning authority at its source.



\textbf{Boundary}



Information administration and coordination preserve the authority of the responsible source, clinical actor and governed domain; they confer no additional real-world management authority.


{\footnotesize \href{../../../../continuum/strategy/capability/4cb7db65392f4fab801d7794af4d4b29.md}{Canonical record: 4cb7db65392f4fab801d7794af4d4b29}; full qualifications and provenance.}\par
\subparagraph*{II.1.3.1.5.1.1 [Feature] Care Goal \& Plan Progression Representation}
\phantomsection
\label{navigation:II.1.3.1.5.1.1}
\addcontentsline{toc}{subparagraph}{II.1.3.1.5.1.1 [Feature] Care Goal \& Plan Progression Representation}

\label{element:7332de92e6c14bc8b071cf18f78f9525}
\index{[Feature] Care Goal \& Plan Progression Representation (II.1.3.1.5.1.1)}
{\small Capability / Feature · Proposed · current}\par

Represent externally agreed care goals, participants and plan changes with their effective standing.



\textbf{Overview}



Plan ingestion/distribution is complemented by the ability to distinguish the applicable goal and plan context as a shared plan progresses.



\textbf{Boundary}



Representing a goal does not choose treatment, determine clinical completion or grant plan approval authority.


{\footnotesize \href{../../../../continuum/strategy/capability/feature/7332de92e6c14bc8b071cf18f78f9525.md}{Canonical record: 7332de92e6c14bc8b071cf18f78f9525}; full qualifications and provenance.}\par
\subparagraph*{II.1.3.1.5.1.2 [Feature] Care Plan Ingestion \& Distribution}
\phantomsection
\label{navigation:II.1.3.1.5.1.2}
\addcontentsline{toc}{subparagraph}{II.1.3.1.5.1.2 [Feature] Care Plan Ingestion \& Distribution}

\label{element:6d4b374a6bca4b58a86c912a94bbd6b1}
\index{[Feature] Care Plan Ingestion \& Distribution (II.1.3.1.5.1.2)}
{\small Capability / Feature · Canonical · current}\par

Ingest multidisciplinary care plans and syndicate updates to authorized care team members.



Care-plan exchange keeps authorised team members aware of established multidisciplinary intent and its updates. It provides information continuity while preserving clinical planning authority and the distinction between a distributed plan and care actually performed.


{\footnotesize \href{../../../../continuum/strategy/capability/feature/6d4b374a6bca4b58a86c912a94bbd6b1.md}{Canonical record: 6d4b374a6bca4b58a86c912a94bbd6b1}; full qualifications and provenance.}\par
\paragraph*{II.1.3.1.5.2 [CT3] Community Care Coordination}
\phantomsection
\label{navigation:II.1.3.1.5.2}
\addcontentsline{toc}{paragraph}{II.1.3.1.5.2 [CT3] Community Care Coordination}

\label{element:8c8c6266d9ed43e09dbe44fcdd669132}
\index{[CT3] Community Care Coordination (II.1.3.1.5.2)}
{\small Capability · Canonical · current}\par

Definition unresolved in the baseline. No definition has been synthesised from the name.



\medskip\hrule\medskip

{\footnotesize \href{../../../../continuum/strategy/capability/8c8c6266d9ed43e09dbe44fcdd669132.md}{Canonical record: 8c8c6266d9ed43e09dbe44fcdd669132}; full qualifications and provenance.}\par
\subparagraph*{II.1.3.1.5.2.1 [Feature] Community Care Plan Synchronization}
\phantomsection
\label{navigation:II.1.3.1.5.2.1}
\addcontentsline{toc}{subparagraph}{II.1.3.1.5.2.1 [Feature] Community Care Plan Synchronization}

\label{element:30a0a2a6f1044d9ca9337b655f303e24}
\index{[Feature] Community Care Plan Synchronization (II.1.3.1.5.2.1)}
{\small Capability / Feature · Canonical · current}\par

Synchronize community care plan updates between hospital outreach teams and community health organizations.



Community plan synchronisation keeps hospital outreach and community participants informed of established plan changes. It supports shared clinical intent without allowing synchronisation to become authority to choose that intent or allocate clinical work.


{\footnotesize \href{../../../../continuum/strategy/capability/feature/30a0a2a6f1044d9ca9337b655f303e24.md}{Canonical record: 30a0a2a6f1044d9ca9337b655f303e24}; full qualifications and provenance.}\par
\subparagraph*{II.1.3.1.5.2.2 [Feature] Community Encounter Ingestion}
\phantomsection
\label{navigation:II.1.3.1.5.2.2}
\addcontentsline{toc}{subparagraph}{II.1.3.1.5.2.2 [Feature] Community Encounter Ingestion}

\label{element:7afe8ae69ed7491192d884576ca0dc15}
\index{[Feature] Community Encounter Ingestion (II.1.3.1.5.2.2)}
{\small Capability / Feature · Canonical · current}\par

Ingest home visit notes, community nursing observations, and social support assessments.



Community encounter ingestion brings home-visit and community assessments into the longitudinal information context. It supports continuity across organisational boundaries while preserving source attribution and the difference between the documentation received and care performed by the platform.


{\footnotesize \href{../../../../continuum/strategy/capability/feature/7afe8ae69ed7491192d884576ca0dc15.md}{Canonical record: 7afe8ae69ed7491192d884576ca0dc15}; full qualifications and provenance.}\par
\paragraph*{II.1.3.1.5.3 [CT3] Outreach Care Continuity}
\phantomsection
\label{navigation:II.1.3.1.5.3}
\addcontentsline{toc}{paragraph}{II.1.3.1.5.3 [CT3] Outreach Care Continuity}

\label{element:8526d971b07f40de84ef05b0d192089c}
\index{[CT3] Outreach Care Continuity (II.1.3.1.5.3)}
{\small Capability · Canonical · current}\par

Definition unresolved in the baseline. No definition has been synthesised from the name.



\medskip\hrule\medskip

{\footnotesize \href{../../../../continuum/strategy/capability/8526d971b07f40de84ef05b0d192089c.md}{Canonical record: 8526d971b07f40de84ef05b0d192089c}; full qualifications and provenance.}\par
\subparagraph*{II.1.3.1.5.3.1 [Feature] Outreach Encounter Outcome Ingestion}
\phantomsection
\label{navigation:II.1.3.1.5.3.1}
\addcontentsline{toc}{subparagraph}{II.1.3.1.5.3.1 [Feature] Outreach Encounter Outcome Ingestion}

\label{element:464ee3ad8d3648d992c6e51e3632d593}
\index{[Feature] Outreach Encounter Outcome Ingestion (II.1.3.1.5.3.1)}
{\small Capability / Feature · Canonical · current}\par

Ingest and reconcile clinical documentation gathered during rural and remote outreach visits.



Outreach outcome ingestion brings documentation collected during remote visits back into the governed record. It reconnects previously separated activity with the regional context without erasing when, where and by whom the source information was established.


{\footnotesize \href{../../../../continuum/strategy/capability/feature/464ee3ad8d3648d992c6e51e3632d593.md}{Canonical record: 464ee3ad8d3648d992c6e51e3632d593}; full qualifications and provenance.}\par
\subparagraph*{II.1.3.1.5.3.2 [Feature] Outreach Visit Context Provision}
\phantomsection
\label{navigation:II.1.3.1.5.3.2}
\addcontentsline{toc}{subparagraph}{II.1.3.1.5.3.2 [Feature] Outreach Visit Context Provision}

\label{element:44049d029a41437685272dfe94db4a2a}
\index{[Feature] Outreach Visit Context Provision (II.1.3.1.5.3.2)}
{\small Capability / Feature · Canonical · current}\par

Package and deliver portable clinical summaries and offline-capable care summaries for remote outreach clinics.



Outreach context provision prepares portable summaries for care beyond the usual connected environment. Their usefulness depends on retaining clinical and temporal context; an offline-capable copy is not automatically current authoritative state or a new origin of clinical facts.


{\footnotesize \href{../../../../continuum/strategy/capability/feature/44049d029a41437685272dfe94db4a2a.md}{Canonical record: 44049d029a41437685272dfe94db4a2a}; full qualifications and provenance.}\par
\subsubsection*{II.1.3.1.6 [CT2] Client-Centred Administration}
\phantomsection
\label{navigation:II.1.3.1.6}
\addcontentsline{toc}{subsubsection}{II.1.3.1.6 [CT2] Client-Centred Administration}

{\small Historical predecessor: R.2.55 \hyperref[element:232b3fb2ff204c9b9ff7a60ff6a7a32d]{[CT2] Client-Centred Administration}; Canonical ID: \texttt{\seqsplit{232b3fb2ff204c9b9ff7a60ff6a7a32d}}. Superseded historical identity — direct correspondence · superseded.}\par
\label{element:8752b7a937754b2faa296300c6a48911}
\index{[CT2] Client-Centred Administration (II.1.3.1.6)}
{\small Capability · Canonical · current}\par

Client-Centred Administration receives demographic and registration changes and makes verified updates available to participating clinical systems. It connects administrative events from external PAS environments with the client context used across care. The enabling responsibility concerns integration of changes, rather than independently deciding the authoritative demographic truth whenever a source sends an update.



\textbf{Overview}



Client-Centred Administration receives demographic and registration changes and makes verified updates available to participating clinical systems. It connects administrative events from external PAS environments with the client context used across care. The enabling responsibility concerns integration of changes, rather than independently deciding the authoritative demographic truth whenever a source sends an update.



\textbf{Intent}



Client-Centred Administration receives demographic and registration changes and makes verified updates available to participating clinical systems. It connects administrative events from external PAS environments with the client context used across care. The enabling responsibility concerns integration of changes, rather than independently deciding the authoritative demographic truth whenever a source sends an update.



\textbf{Boundary}



Information administration does not confer authority over the represented person, organisation, service, place, device or clinical activity.



\medskip\hrule\medskip

{\footnotesize \href{../../../../continuum/strategy/capability/8752b7a937754b2faa296300c6a48911.md}{Canonical record: 8752b7a937754b2faa296300c6a48911}; full qualifications and provenance.}\par
\paragraph*{II.1.3.1.6.1 [CT3] Client Intake \& Demographic Update}
\phantomsection
\label{navigation:II.1.3.1.6.1}
\addcontentsline{toc}{paragraph}{II.1.3.1.6.1 [CT3] Client Intake \& Demographic Update}

\label{element:f51f8d6b7556424987c76d9a7465c67c}
\index{[CT3] Client Intake \& Demographic Update (II.1.3.1.6.1)}
{\small Capability · Proposed · current}\par

Administer the established client intake \& demographic update information and coordination responsibility while preserving its source authority.



\textbf{Overview}



Administer the established client intake \& demographic update information and coordination responsibility while preserving its source authority. The reviewed hierarchy assigns this enduring ability to its stated business-enabling context.



\textbf{Intent}



Administer the established client intake \& demographic update information and coordination responsibility while preserving its source authority.



\textbf{Boundary}



Information administration and coordination preserve the authority of the responsible source, clinical actor and governed domain; they confer no additional real-world management authority.


{\footnotesize \href{../../../../continuum/strategy/capability/f51f8d6b7556424987c76d9a7465c67c.md}{Canonical record: f51f8d6b7556424987c76d9a7465c67c}; full qualifications and provenance.}\par
\subparagraph*{II.1.3.1.6.1.1 [Feature] Demographic Change Distribution}
\phantomsection
\label{navigation:II.1.3.1.6.1.1}
\addcontentsline{toc}{subparagraph}{II.1.3.1.6.1.1 [Feature] Demographic Change Distribution}

\label{element:bfd2ce654a79404dbef2694d30fb39fc}
\index{[Feature] Demographic Change Distribution (II.1.3.1.6.1.1)}
{\small Capability / Feature · Canonical · current}\par

Fan out verified client demographic updates to downstream clinical subsystems.



Change distribution propagates verified demographic updates to participating clinical systems. It supports consistent subject context across the network while keeping delivery to each destination distinguishable from that destination's eventual incorporation or business response.


{\footnotesize \href{../../../../continuum/strategy/capability/feature/bfd2ce654a79404dbef2694d30fb39fc.md}{Canonical record: bfd2ce654a79404dbef2694d30fb39fc}; full qualifications and provenance.}\par
\subparagraph*{II.1.3.1.6.1.2 [Feature] Demographic Update Ingestion}
\phantomsection
\label{navigation:II.1.3.1.6.1.2}
\addcontentsline{toc}{subparagraph}{II.1.3.1.6.1.2 [Feature] Demographic Update Ingestion}

\label{element:594ccbd95f5841e4b874fd1bade34a5d}
\index{[Feature] Demographic Update Ingestion (II.1.3.1.6.1.2)}
{\small Capability / Feature · Canonical · current}\par

Ingest and validate patient admission and registration updates from external PAS systems.



Demographic ingestion receives administrative changes from upstream admission and registration systems. It makes externally established updates available for governed processing without treating each incoming assertion as an independent master-person adjudication by Harmonia.


{\footnotesize \href{../../../../continuum/strategy/capability/feature/594ccbd95f5841e4b874fd1bade34a5d.md}{Canonical record: 594ccbd95f5841e4b874fd1bade34a5d}; full qualifications and provenance.}\par
\paragraph*{II.1.3.1.6.2 [CT3] Screening and Immunisation Enablement}
\phantomsection
\label{navigation:II.1.3.1.6.2}
\addcontentsline{toc}{paragraph}{II.1.3.1.6.2 [CT3] Screening and Immunisation Enablement}

\label{element:3e9016eaebae4ed0b67a5cacba2652f5}
\index{[CT3] Screening and Immunisation Enablement (II.1.3.1.6.2)}
{\small Capability · Canonical · current}\par

Definition unresolved in the baseline. No definition has been synthesised from the name.



\medskip\hrule\medskip

{\footnotesize \href{../../../../continuum/strategy/capability/3e9016eaebae4ed0b67a5cacba2652f5.md}{Canonical record: 3e9016eaebae4ed0b67a5cacba2652f5}; full qualifications and provenance.}\par
\subparagraph*{II.1.3.1.6.2.1 [Feature] Immunisation Event Submission}
\phantomsection
\label{navigation:II.1.3.1.6.2.1}
\addcontentsline{toc}{subparagraph}{II.1.3.1.6.2.1 [Feature] Immunisation Event Submission}

\label{element:47bab2681bf943899e83fab8c9941676}
\index{[Feature] Immunisation Event Submission (II.1.3.1.6.2.1)}
{\small Capability / Feature · Canonical · current}\par

Ingest point-of-care immunisation administration records and syndicate to national immunization registers.



Immunisation submission connects a source-recorded administration event with the national register context. It supports continuity and reporting without performing the immunisation or making successful transmission equivalent to every downstream registry outcome.


{\footnotesize \href{../../../../continuum/strategy/capability/feature/47bab2681bf943899e83fab8c9941676.md}{Canonical record: 47bab2681bf943899e83fab8c9941676}; full qualifications and provenance.}\par
\subparagraph*{II.1.3.1.6.2.2 [Feature] Screening Recall Notification Distribution}
\phantomsection
\label{navigation:II.1.3.1.6.2.2}
\addcontentsline{toc}{subparagraph}{II.1.3.1.6.2.2 [Feature] Screening Recall Notification Distribution}

\label{element:7ec85487625349a7b3292585fadb20c7}
\index{[Feature] Screening Recall Notification Distribution (II.1.3.1.6.2.2)}
{\small Capability / Feature · Canonical · current}\par

Distribute clinical screening reminders and recall notices to designated healthcare providers.



Screening recall distribution conveys established reminders to the designated care providers. Its notification responsibility is expressible through approved event and workflow composition; the catalogue does not invent a specialised Function or transfer the clinical decision to recall the patient.


{\footnotesize \href{../../../../continuum/strategy/capability/feature/7ec85487625349a7b3292585fadb20c7.md}{Canonical record: 7ec85487625349a7b3292585fadb20c7}; full qualifications and provenance.}\par
\subsubsection*{II.1.3.1.7 [CT2] Clinical Charging \& Claims Administration}
\phantomsection
\label{navigation:II.1.3.1.7}
\addcontentsline{toc}{subsubsection}{II.1.3.1.7 [CT2] Clinical Charging \& Claims Administration}

{\small Historical predecessor: R.2.56 \hyperref[element:44433108cc8f4f5f9d7581f9e5b29eb5]{[CT2] Clinical Charging \& Claims Administration}; Canonical ID: \texttt{\seqsplit{44433108cc8f4f5f9d7581f9e5b29eb5}}. Superseded historical identity — direct correspondence · superseded.}\par
\label{element:af78ab7c89454a608cd4d34ec7703331}
\index{[CT2] Clinical Charging \& Claims Administration (II.1.3.1.7)}
{\small Capability · Canonical · current}\par

This capability supplies clinical activity summaries and charging records to billing contexts. It preserves a route from activity information to the systems responsible for financial processing without treating the export as a completed bill or claim. No Feature decomposition or comprehensive revenue-cycle responsibility is established.



\textbf{Overview}



This capability supplies clinical activity summaries and charging records to billing contexts. It preserves a route from activity information to the systems responsible for financial processing without treating the export as a completed bill or claim. No Feature decomposition or comprehensive revenue-cycle responsibility is established.



\textbf{Intent}



This capability supplies clinical activity summaries and charging records to billing contexts. It preserves a route from activity information to the systems responsible for financial processing without treating the export as a completed bill or claim. No Feature decomposition or comprehensive revenue-cycle responsibility is established.



\textbf{Boundary}



Information administration does not confer authority over the represented person, organisation, service, place, device or clinical activity.



\medskip\hrule\medskip

{\footnotesize \href{../../../../continuum/strategy/capability/af78ab7c89454a608cd4d34ec7703331.md}{Canonical record: af78ab7c89454a608cd4d34ec7703331}; full qualifications and provenance.}\par
\subsubsection*{II.1.3.1.8 [CT2] Clinical Communication Administration}
\phantomsection
\label{navigation:II.1.3.1.8}
\addcontentsline{toc}{subsubsection}{II.1.3.1.8 [CT2] Clinical Communication Administration}

{\small Historical predecessor: R.2.57 \hyperref[element:bd4d6f4a7f03402ba7857383c7be7a4c]{[CT2] Clinical Communication Administration}; Canonical ID: \texttt{\seqsplit{bd4d6f4a7f03402ba7857383c7be7a4c}}. Superseded historical identity — direct correspondence · superseded.}\par
\label{element:b38f2c2811724734b5887b6dbaaabbfc}
\index{[CT2] Clinical Communication Administration (II.1.3.1.8)}
{\small Capability · Canonical · current}\par

This capability manages secure clinical communication transactions, delivery and receipt evidence, and the required enduring audit trail. Verified provider and team endpoints support appropriate addressing while transaction correlation makes the communication history intelligible. The responsibility preserves separate meanings for transport delivery, recipient acknowledgement and clinical or business disposition.



\textbf{Overview}



This capability manages secure clinical communication transactions, delivery and receipt evidence, and the required enduring audit trail. Verified provider and team endpoints support appropriate addressing while transaction correlation makes the communication history intelligible. The responsibility preserves separate meanings for transport delivery, recipient acknowledgement and clinical or business disposition.



\textbf{Intent}



This capability manages secure clinical communication transactions, delivery and receipt evidence, and the required enduring audit trail. Verified provider and team endpoints support appropriate addressing while transaction correlation makes the communication history intelligible. The responsibility preserves separate meanings for transport delivery, recipient acknowledgement and clinical or business disposition.



\textbf{Boundary}



Information administration does not confer authority over the represented person, organisation, service, place, device or clinical activity.



\medskip\hrule\medskip

{\footnotesize \href{../../../../continuum/strategy/capability/b38f2c2811724734b5887b6dbaaabbfc.md}{Canonical record: b38f2c2811724734b5887b6dbaaabbfc}; full qualifications and provenance.}\par
\paragraph*{II.1.3.1.8.1 [CT3] Secure Clinical Communication}
\phantomsection
\label{navigation:II.1.3.1.8.1}
\addcontentsline{toc}{paragraph}{II.1.3.1.8.1 [CT3] Secure Clinical Communication}

\label{element:eddb86c7922145439b9dcf0259f4aed5}
\index{[CT3] Secure Clinical Communication (II.1.3.1.8.1)}
{\small Capability · Proposed · current}\par

Administer the established secure clinical communication information and coordination responsibility while preserving its source authority.



\textbf{Overview}



Administer the established secure clinical communication information and coordination responsibility while preserving its source authority. The reviewed hierarchy assigns this enduring ability to its stated business-enabling context.



\textbf{Intent}



Administer the established secure clinical communication information and coordination responsibility while preserving its source authority.



\textbf{Boundary}



Information administration and coordination preserve the authority of the responsible source, clinical actor and governed domain; they confer no additional real-world management authority.


{\footnotesize \href{../../../../continuum/strategy/capability/eddb86c7922145439b9dcf0259f4aed5.md}{Canonical record: eddb86c7922145439b9dcf0259f4aed5}; full qualifications and provenance.}\par
\subparagraph*{II.1.3.1.8.1.1 [Feature] Clinical Communication Distribution}
\phantomsection
\label{navigation:II.1.3.1.8.1.1}
\addcontentsline{toc}{subparagraph}{II.1.3.1.8.1.1 [Feature] Clinical Communication Distribution}

\label{element:8527f36d5d5a460b92403f9d84f36490}
\index{[Feature] Clinical Communication Distribution (II.1.3.1.8.1.1)}
{\small Capability / Feature · Canonical · current}\par

Distribute secure clinical communications between verified healthcare providers and designated care team endpoints.



Clinical communication distribution connects verified providers and designated team endpoints through governed addressed delivery. The responsibility concerns reaching the intended participant, rather than assuming that delivered information has been read, incorporated or acted upon clinically.


{\footnotesize \href{../../../../continuum/strategy/capability/feature/8527f36d5d5a460b92403f9d84f36490.md}{Canonical record: 8527f36d5d5a460b92403f9d84f36490}; full qualifications and provenance.}\par
\subparagraph*{II.1.3.1.8.1.2 [Feature] Communication Audit Logging}
\phantomsection
\label{navigation:II.1.3.1.8.1.2}
\addcontentsline{toc}{subparagraph}{II.1.3.1.8.1.2 [Feature] Communication Audit Logging}

\label{element:b9df49b2f221486392f5cc393369d888}
\index{[Feature] Communication Audit Logging (II.1.3.1.8.1.2)}
{\small Capability / Feature · Canonical · current}\par

Maintain an enduring, non-PHI, tamper-evident log of all clinical communication transactions.



Communication audit evidence preserves the attributable transaction history required for scrutiny of clinical communication. It records meaningful communication facts under the zero-PHI constraints rather than turning the audit trail into a copy of message content or a permanent diagnostic trace.


{\footnotesize \href{../../../../continuum/strategy/capability/feature/b9df49b2f221486392f5cc393369d888.md}{Canonical record: b9df49b2f221486392f5cc393369d888}; full qualifications and provenance.}\par
\subparagraph*{II.1.3.1.8.1.3 [Feature] Delivery Acknowledgement Tracking}
\phantomsection
\label{navigation:II.1.3.1.8.1.3}
\addcontentsline{toc}{subparagraph}{II.1.3.1.8.1.3 [Feature] Delivery Acknowledgement Tracking}

\label{element:6884eb06bb794c5b8513b9b63d69af7e}
\index{[Feature] Delivery Acknowledgement Tracking (II.1.3.1.8.1.3)}
{\small Capability / Feature · Canonical · current}\par

Correlate recipient delivery and read receipts with outbound clinical communications, distinguishing transport delivery from recipient acknowledgement. Receipt/read evidence does not by itself establish business acceptance, clinical review, incorporation, approval, finality or originating authority.



Acknowledgement tracking associates delivery and read evidence with the communication transaction. Each receipt retains its contractual meaning; neither delivery nor reading alone proves business acceptance, clinical review, approval, finality or originating authority.


{\footnotesize \href{../../../../continuum/strategy/capability/feature/6884eb06bb794c5b8513b9b63d69af7e.md}{Canonical record: 6884eb06bb794c5b8513b9b63d69af7e}; full qualifications and provenance.}\par
\subsubsection*{II.1.3.1.9 [CT2] Clinical Record Administration}
\phantomsection
\label{navigation:II.1.3.1.9}
\addcontentsline{toc}{subsubsection}{II.1.3.1.9 [CT2] Clinical Record Administration}

{\small Historical predecessor: R.2.59 \hyperref[element:be3f2e695eda403fbce8118daf83d386]{[CT2] Clinical Record Administration}; Canonical ID: \texttt{\seqsplit{be3f2e695eda403fbce8118daf83d386}}. Superseded historical identity — direct correspondence · superseded.}\par
\label{element:ce9e7e7b22054cefb707a650d5be7219}
\index{[CT2] Clinical Record Administration (II.1.3.1.9)}
{\small Capability · Canonical · current}\par

Clinical Record Administration manages clinical documents as governed information with authorship, classification, versions and supersession history. It distinguishes a new document from an amendment or withdrawal and maintains metadata for discovery. Its purpose is to preserve an intelligible documentary record across change, rather than merely accumulate files.



\textbf{Overview}



Clinical Record Administration manages clinical documents as governed information with authorship, classification, versions and supersession history. It distinguishes a new document from an amendment or withdrawal and maintains metadata for discovery. Its purpose is to preserve an intelligible documentary record across change, rather than merely accumulate files.



\textbf{Intent}



Clinical Record Administration manages clinical documents as governed information with authorship, classification, versions and supersession history. It distinguishes a new document from an amendment or withdrawal and maintains metadata for discovery. Its purpose is to preserve an intelligible documentary record across change, rather than merely accumulate files.



\textbf{Boundary}



Information administration does not confer authority over the represented person, organisation, service, place, device or clinical activity.



\medskip\hrule\medskip

{\footnotesize \href{../../../../continuum/strategy/capability/ce9e7e7b22054cefb707a650d5be7219.md}{Canonical record: ce9e7e7b22054cefb707a650d5be7219}; full qualifications and provenance.}\par
\paragraph*{II.1.3.1.9.1 [CT3] Clinical Documentation Lifecycle Administration}
\phantomsection
\label{navigation:II.1.3.1.9.1}
\addcontentsline{toc}{paragraph}{II.1.3.1.9.1 [CT3] Clinical Documentation Lifecycle Administration}

{\small Historical predecessor: R.2.58 \hyperref[element:5b1636373a8d4e788b3ddbc1659a47ef]{[CT3] Clinical Documentation Lifecycle Administration}; Canonical ID: \texttt{\seqsplit{5b1636373a8d4e788b3ddbc1659a47ef}}. Superseded historical identity — direct correspondence · superseded.}\par
\label{element:df683100067b47e785037c27cfd17a3c}
\index{[CT3] Clinical Documentation Lifecycle Administration (II.1.3.1.9.1)}
{\small Capability · Proposed · current}\par

Administer ingestion, indexing, versioning, supersession and governed availability of clinical documentation.



\textbf{Overview}



Administer ingestion, indexing, versioning, supersession and governed availability of clinical documentation. The reviewed hierarchy assigns this enduring ability to its stated business-enabling context.



\textbf{Intent}



Administer ingestion, indexing, versioning, supersession and governed availability of clinical documentation.



\textbf{Boundary}



Information administration and coordination preserve the authority of the responsible source, clinical actor and governed domain; they confer no additional real-world management authority.


{\footnotesize \href{../../../../continuum/strategy/capability/df683100067b47e785037c27cfd17a3c.md}{Canonical record: df683100067b47e785037c27cfd17a3c}; full qualifications and provenance.}\par
\subparagraph*{II.1.3.1.9.1.1 [Feature] Clinical Document Ingestion}
\phantomsection
\label{navigation:II.1.3.1.9.1.1}
\addcontentsline{toc}{subparagraph}{II.1.3.1.9.1.1 [Feature] Clinical Document Ingestion}

\label{element:54c79edadc7a478fb594ff4eb47f8368}
\index{[Feature] Clinical Document Ingestion (II.1.3.1.9.1.1)}
{\small Capability / Feature · Canonical · current}\par

Ingest structured and unstructured clinical documents (discharge summaries, consult notes, letters).



Clinical document ingestion receives authored documents into governed information management. It supports both structured and unstructured contributions while preserving that receipt establishes a managed contribution, not independent clinical verification of every assertion in the document.


{\footnotesize \href{../../../../continuum/strategy/capability/feature/54c79edadc7a478fb594ff4eb47f8368.md}{Canonical record: 54c79edadc7a478fb594ff4eb47f8368}; full qualifications and provenance.}\par
\subparagraph*{II.1.3.1.9.1.2 [Feature] Document Metadata Indexing}
\phantomsection
\label{navigation:II.1.3.1.9.1.2}
\addcontentsline{toc}{subparagraph}{II.1.3.1.9.1.2 [Feature] Document Metadata Indexing}

\label{element:0c85a269e90a4329b83d934f9061cd49}
\index{[Feature] Document Metadata Indexing (II.1.3.1.9.1.2)}
{\small Capability / Feature · Canonical · current}\par

Index clinical document metadata (author, specialty, encounter, date, confidentiality) for rapid discovery.



Metadata indexing makes documents discoverable by clinically and operationally relevant attributes. It describes content for search without making the index the document itself or granting access merely because a matching record can be located.


{\footnotesize \href{../../../../continuum/strategy/capability/feature/0c85a269e90a4329b83d934f9061cd49.md}{Canonical record: 0c85a269e90a4329b83d934f9061cd49}; full qualifications and provenance.}\par
\subparagraph*{II.1.3.1.9.1.3 [Feature] Document Versioning \& Supersession}
\phantomsection
\label{navigation:II.1.3.1.9.1.3}
\addcontentsline{toc}{subparagraph}{II.1.3.1.9.1.3 [Feature] Document Versioning \& Supersession}

\label{element:c8e2e623dd4e4c3b8de6c3bfb8c2a3ee}
\index{[Feature] Document Versioning \& Supersession (II.1.3.1.9.1.3)}
{\small Capability / Feature · Canonical · current}\par

Manage document classification, author attribution, amendment attribution, versioning, and supersession or withdrawal without destructive data loss.



Versioning and supersession make document change understandable without destructive loss of the previous account. Author and amendment attribution distinguish who established each contribution; withdrawal or correction changes current standing without retrospectively erasing what Harmonia previously received or managed.


{\footnotesize \href{../../../../continuum/strategy/capability/feature/c8e2e623dd4e4c3b8de6c3bfb8c2a3ee.md}{Canonical record: c8e2e623dd4e4c3b8de6c3bfb8c2a3ee}; full qualifications and provenance.}\par
\subsubsection*{II.1.3.1.10 [CT2] Diagnostic Administration}
\phantomsection
\label{navigation:II.1.3.1.10}
\addcontentsline{toc}{subsubsection}{II.1.3.1.10 [CT2] Diagnostic Administration}

{\small Historical predecessor: R.2.60 \hyperref[element:2b1a27d5d9944e05b4cf0fcaad4ca433]{[CT2] Diagnostic Administration}; Canonical ID: \texttt{\seqsplit{2b1a27d5d9944e05b4cf0fcaad4ca433}}. Superseded historical identity — direct correspondence · superseded.}\par
\label{element:3cc720ba2e8b4d7a94bb349391d61895}
\index{[CT2] Diagnostic Administration (II.1.3.1.10)}
{\small Capability · Canonical · current}\par

Diagnostic Administration correlates requests with specimens and encounters and integrates diagnostic reports into their request context. It connects the laboratory or imaging contribution with the participants expecting an outcome, preserving report identity and standing as updates occur. Its enabling role differs from performing the diagnostic procedure or professionally interpreting evidence.



\textbf{Overview}



Diagnostic Administration correlates requests with specimens and encounters and integrates diagnostic reports into their request context. It connects the laboratory or imaging contribution with the participants expecting an outcome, preserving report identity and standing as updates occur. Its enabling role differs from performing the diagnostic procedure or professionally interpreting evidence.



\textbf{Intent}



Diagnostic Administration correlates requests with specimens and encounters and integrates diagnostic reports into their request context. It connects the laboratory or imaging contribution with the participants expecting an outcome, preserving report identity and standing as updates occur. Its enabling role differs from performing the diagnostic procedure or professionally interpreting evidence.



\textbf{Boundary}



Information administration does not confer authority over the represented person, organisation, service, place, device or clinical activity.



\medskip\hrule\medskip

{\footnotesize \href{../../../../continuum/strategy/capability/3cc720ba2e8b4d7a94bb349391d61895.md}{Canonical record: 3cc720ba2e8b4d7a94bb349391d61895}; full qualifications and provenance.}\par
\paragraph*{II.1.3.1.10.1 [CT3] Diagnostic Service Integration}
\phantomsection
\label{navigation:II.1.3.1.10.1}
\addcontentsline{toc}{paragraph}{II.1.3.1.10.1 [CT3] Diagnostic Service Integration}

\label{element:9ded2f3a28d24017bb921a0cdcf65a2b}
\index{[CT3] Diagnostic Service Integration (II.1.3.1.10.1)}
{\small Capability · Canonical · current}\par

Definition unresolved in the baseline. No definition has been synthesised from the name.



\medskip\hrule\medskip

{\footnotesize \href{../../../../continuum/strategy/capability/9ded2f3a28d24017bb921a0cdcf65a2b.md}{Canonical record: 9ded2f3a28d24017bb921a0cdcf65a2b}; full qualifications and provenance.}\par
\subparagraph*{II.1.3.1.10.1.1 [Feature] Diagnostic History Consolidation}
\phantomsection
\label{navigation:II.1.3.1.10.1.1}
\addcontentsline{toc}{subparagraph}{II.1.3.1.10.1.1 [Feature] Diagnostic History Consolidation}

\label{element:0b3fa0b7374842ff8ce199191a3479bd}
\index{[Feature] Diagnostic History Consolidation (II.1.3.1.10.1.1)}
{\small Capability / Feature · Canonical · current}\par

Consolidate historical diagnostic trends (e.g., cumulative laboratory results) across multiple regional providers into an integrated view.



Diagnostic history consolidation makes prior evidence visible across regional providers. A combined historical view can support comparison over time without itself establishing a diagnostic interpretation, reconciling conflicting clinical assertions or replacing each report's source standing.


{\footnotesize \href{../../../../continuum/strategy/capability/feature/0b3fa0b7374842ff8ce199191a3479bd.md}{Canonical record: 0b3fa0b7374842ff8ce199191a3479bd}; full qualifications and provenance.}\par
\subparagraph*{II.1.3.1.10.1.2 [Feature] Laboratory \& Imaging Result Distribution}
\phantomsection
\label{navigation:II.1.3.1.10.1.2}
\addcontentsline{toc}{subparagraph}{II.1.3.1.10.1.2 [Feature] Laboratory \& Imaging Result Distribution}

\label{element:93d42becda8f4850b73ff2c93bb7e8f3}
\index{[Feature] Laboratory \& Imaging Result Distribution (II.1.3.1.10.1.2)}
{\small Capability / Feature · Canonical · current}\par

Reliably distribute structured diagnostic reports (pathology, radiology) to ordering clinicians and designated secondary recipients.



Result distribution connects established diagnostic reports with ordering and designated secondary recipients. It supports continuity of diagnostic information across care settings while retaining the distinction between delivery, recipient review and the clinical decision based on the result.


{\footnotesize \href{../../../../continuum/strategy/capability/feature/93d42becda8f4850b73ff2c93bb7e8f3.md}{Canonical record: 93d42becda8f4850b73ff2c93bb7e8f3}; full qualifications and provenance.}\par
\paragraph*{II.1.3.1.10.2 [CT3] Diagnostic Workflow Management}
\phantomsection
\label{navigation:II.1.3.1.10.2}
\addcontentsline{toc}{paragraph}{II.1.3.1.10.2 [CT3] Diagnostic Workflow Management}

\label{element:270ed0e32b294458aa52cf57df7e55db}
\index{[CT3] Diagnostic Workflow Management (II.1.3.1.10.2)}
{\small Capability · Proposed · current}\par

Administer the established diagnostic workflow management information and coordination responsibility while preserving its source authority.



\textbf{Overview}



Administer the established diagnostic workflow management information and coordination responsibility while preserving its source authority. The reviewed hierarchy assigns this enduring ability to its stated business-enabling context.



\textbf{Intent}



Administer the established diagnostic workflow management information and coordination responsibility while preserving its source authority.



\textbf{Boundary}



Information administration and coordination preserve the authority of the responsible source, clinical actor and governed domain; they confer no additional real-world management authority.


{\footnotesize \href{../../../../continuum/strategy/capability/270ed0e32b294458aa52cf57df7e55db.md}{Canonical record: 270ed0e32b294458aa52cf57df7e55db}; full qualifications and provenance.}\par
\subparagraph*{II.1.3.1.10.2.1 [Feature] Diagnostic Preliminary \& Corrected Report Handling}
\phantomsection
\label{navigation:II.1.3.1.10.2.1}
\addcontentsline{toc}{subparagraph}{II.1.3.1.10.2.1 [Feature] Diagnostic Preliminary \& Corrected Report Handling}

\label{element:7449ee4dada7424ea6d39785258b4981}
\index{[Feature] Diagnostic Preliminary \& Corrected Report Handling (II.1.3.1.10.2.1)}
{\small Capability / Feature · Proposed · current}\par

Receive, qualify and distribute preliminary, corrected and addendum diagnostic information with its relationship to prior reports.



\textbf{Overview}



The finalised-report Feature remains unchanged. The broader documented diagnostic behaviour needs a separate frame for information whose preliminary or correction standing is material.



\textbf{Boundary}



Does not originate diagnostic interpretation or silently overwrite a previously received report.


{\footnotesize \href{../../../../continuum/strategy/capability/feature/7449ee4dada7424ea6d39785258b4981.md}{Canonical record: 7449ee4dada7424ea6d39785258b4981}; full qualifications and provenance.}\par
\subparagraph*{II.1.3.1.10.2.2 [Feature] Diagnostic Report Distribution}
\phantomsection
\label{navigation:II.1.3.1.10.2.2}
\addcontentsline{toc}{subparagraph}{II.1.3.1.10.2.2 [Feature] Diagnostic Report Distribution}

\label{element:3813b3b5175e461ba74805d8fa2aa648}
\index{[Feature] Diagnostic Report Distribution (II.1.3.1.10.2.2)}
{\small Capability / Feature · Canonical · current}\par

Distribute finalized laboratory, pathology, and imaging reports to ordering clinicians and designated care teams.



This Feature distributes finalised diagnostic reports to the ordering and designated care participants. Its narrower finalised-report scope is preserved even though broader Diagnostic Administration behaviour also handles preliminary, corrected and addendum information; those behaviours do not silently redefine the Feature.


{\footnotesize \href{../../../../continuum/strategy/capability/feature/3813b3b5175e461ba74805d8fa2aa648.md}{Canonical record: 3813b3b5175e461ba74805d8fa2aa648}; full qualifications and provenance.}\par
\subparagraph*{II.1.3.1.10.2.3 [Feature] Diagnostic Report Ingestion \& Binding}
\phantomsection
\label{navigation:II.1.3.1.10.2.3}
\addcontentsline{toc}{subparagraph}{II.1.3.1.10.2.3 [Feature] Diagnostic Report Ingestion \& Binding}

\label{element:a00823445a8a44c1a00f3ec93a7e0e6f}
\index{[Feature] Diagnostic Report Ingestion \& Binding (II.1.3.1.10.2.3)}
{\small Capability / Feature · Canonical · current}\par

Ingest finalized laboratory and radiology reports, verifying content integrity and binding them to the parent order.



Report ingestion binds a finalised diagnostic contribution to the order it answers and verifies content integrity. The report's professional standing comes from its diagnostic source; successful ingestion does not independently authorise its findings or establish a patient-level clinical conclusion.


{\footnotesize \href{../../../../continuum/strategy/capability/feature/a00823445a8a44c1a00f3ec93a7e0e6f.md}{Canonical record: a00823445a8a44c1a00f3ec93a7e0e6f}; full qualifications and provenance.}\par
\subparagraph*{II.1.3.1.10.2.4 [Feature] Diagnostic Request Correlation}
\phantomsection
\label{navigation:II.1.3.1.10.2.4}
\addcontentsline{toc}{subparagraph}{II.1.3.1.10.2.4 [Feature] Diagnostic Request Correlation}

\label{element:30cd73b8d7d941648affbb379196f532}
\index{[Feature] Diagnostic Request Correlation (II.1.3.1.10.2.4)}
{\small Capability / Feature · Canonical · current}\par

Correlate diagnostic test orders with collected clinical specimens and patient encounters.



Diagnostic request correlation connects the order with its specimen and encounter context. It helps distinguish evidence belonging to the intended fulfilment from superficially similar results, without treating association as proof of specimen suitability or diagnostic interpretation.


{\footnotesize \href{../../../../continuum/strategy/capability/feature/30cd73b8d7d941648affbb379196f532.md}{Canonical record: 30cd73b8d7d941648affbb379196f532}; full qualifications and provenance.}\par
\subsubsection*{II.1.3.1.11 [CT2] Episode \& Encounter Administration}
\phantomsection
\label{navigation:II.1.3.1.11}
\addcontentsline{toc}{subsubsection}{II.1.3.1.11 [CT2] Episode \& Encounter Administration}

{\small Historical predecessor: R.2.61 \hyperref[element:2778f988e8d44e3091749542c39d38a4]{[CT2] Episode \& Encounter Administration}; Canonical ID: \texttt{\seqsplit{2778f988e8d44e3091749542c39d38a4}}. Superseded historical identity — direct correspondence · superseded.}\par
\label{element:61d7067aee0540259b3daca4287d66d2}
\index{[CT2] Episode \& Encounter Administration (II.1.3.1.11)}
{\small Capability · Canonical · current}\par

This capability maintains identifiable encounter contexts and their governed progression, including association with care teams and locations. It enables information from an admission, visit or community contact to be understood within the relevant episode of activity. Encounter progression is an administrative and operational account, separate from the clinical acts that occur during it.



\textbf{Overview}



This capability maintains identifiable encounter contexts and their governed progression, including association with care teams and locations. It enables information from an admission, visit or community contact to be understood within the relevant episode of activity. Encounter progression is an administrative and operational account, separate from the clinical acts that occur during it.



\textbf{Intent}



This capability maintains identifiable encounter contexts and their governed progression, including association with care teams and locations. It enables information from an admission, visit or community contact to be understood within the relevant episode of activity. Encounter progression is an administrative and operational account, separate from the clinical acts that occur during it.



\textbf{Boundary}



Information administration does not confer authority over the represented person, organisation, service, place, device or clinical activity.



\medskip\hrule\medskip

{\footnotesize \href{../../../../continuum/strategy/capability/61d7067aee0540259b3daca4287d66d2.md}{Canonical record: 61d7067aee0540259b3daca4287d66d2}; full qualifications and provenance.}\par
\paragraph*{II.1.3.1.11.1 [CT3] Acute Care Monitoring \& Event Coordination}
\phantomsection
\label{navigation:II.1.3.1.11.1}
\addcontentsline{toc}{paragraph}{II.1.3.1.11.1 [CT3] Acute Care Monitoring \& Event Coordination}

\label{element:ee17b26b8c874cc3af50d07a787ce488}
\index{[CT3] Acute Care Monitoring \& Event Coordination (II.1.3.1.11.1)}
{\small Capability · Canonical · current}\par

Definition unresolved in the baseline. No definition has been synthesised from the name.



\medskip\hrule\medskip

{\footnotesize \href{../../../../continuum/strategy/capability/ee17b26b8c874cc3af50d07a787ce488.md}{Canonical record: ee17b26b8c874cc3af50d07a787ce488}; full qualifications and provenance.}\par
\subparagraph*{II.1.3.1.11.1.1 [Feature] Acute Care Clinical Context Provision}
\phantomsection
\label{navigation:II.1.3.1.11.1.1}
\addcontentsline{toc}{subparagraph}{II.1.3.1.11.1.1 [Feature] Acute Care Clinical Context Provision}

\label{element:90de621f7ac1410eb6fc770cb07a9e03}
\index{[Feature] Acute Care Clinical Context Provision (II.1.3.1.11.1.1)}
{\small Capability / Feature · Canonical · current}\par

Provide consolidated longitudinal history and active medication lists to acute care clinicians at point of triage.



Acute context provision brings relevant history and active medications together for clinicians facing a current presentation. It supports rapid understanding without substituting the consolidated view for clinical assessment, triage or a decision about treatment.


{\footnotesize \href{../../../../continuum/strategy/capability/feature/90de621f7ac1410eb6fc770cb07a9e03.md}{Canonical record: 90de621f7ac1410eb6fc770cb07a9e03}; full qualifications and provenance.}\par
\subparagraph*{II.1.3.1.11.1.2 [Feature] Acute Clinical State Event Capture}
\phantomsection
\label{navigation:II.1.3.1.11.1.2}
\addcontentsline{toc}{subparagraph}{II.1.3.1.11.1.2 [Feature] Acute Clinical State Event Capture}

\label{element:566f6934c9b04627b0876f91bedd8f30}
\index{[Feature] Acute Clinical State Event Capture (II.1.3.1.11.1.2)}
{\small Capability / Feature · Canonical · current}\par

Ingest clinical deterioration notifications and acute clinical status updates.



Acute-state capture records deterioration notifications and current clinical status as externally established information. Approved reusable clinical-information composition is sufficient for this responsibility; routing a critical-result alert does not become a specialised capture Function merely because both concern urgent care.


{\footnotesize \href{../../../../continuum/strategy/capability/feature/566f6934c9b04627b0876f91bedd8f30.md}{Canonical record: 566f6934c9b04627b0876f91bedd8f30}; full qualifications and provenance.}\par
\paragraph*{II.1.3.1.11.2 [CT3] Encounter Lifecycle Management}
\phantomsection
\label{navigation:II.1.3.1.11.2}
\addcontentsline{toc}{paragraph}{II.1.3.1.11.2 [CT3] Encounter Lifecycle Management}

\label{element:2ebf1bb20ab34755b41567c156cc5928}
\index{[CT3] Encounter Lifecycle Management (II.1.3.1.11.2)}
{\small Capability · Proposed · current}\par

Administer the established encounter lifecycle management information and coordination responsibility while preserving its source authority.



\textbf{Overview}



Administer the established encounter lifecycle management information and coordination responsibility while preserving its source authority. The reviewed hierarchy assigns this enduring ability to its stated business-enabling context.



\textbf{Intent}



Administer the established encounter lifecycle management information and coordination responsibility while preserving its source authority.



\textbf{Boundary}



Information administration and coordination preserve the authority of the responsible source, clinical actor and governed domain; they confer no additional real-world management authority.


{\footnotesize \href{../../../../continuum/strategy/capability/2ebf1bb20ab34755b41567c156cc5928.md}{Canonical record: 2ebf1bb20ab34755b41567c156cc5928}; full qualifications and provenance.}\par
\subparagraph*{II.1.3.1.11.2.1 [Feature] Encounter Bed/Location Association}
\phantomsection
\label{navigation:II.1.3.1.11.2.1}
\addcontentsline{toc}{subparagraph}{II.1.3.1.11.2.1 [Feature] Encounter Bed/Location Association}

\label{element:47f0fca36aa4454eae95f4b151c7076f}
\index{[Feature] Encounter Bed/Location Association (II.1.3.1.11.2.1)}
{\small Capability / Feature · Canonical · current}\par

Dynamically maintain the association between an encounter, care team, and physical bed or ward location.



Encounter association keeps the relevant place and care-team context aligned with the episode as the patient moves. Location changes are meaningful operational facts, but they do not themselves establish a transfer of clinical responsibility.


{\footnotesize \href{../../../../continuum/strategy/capability/feature/47f0fca36aa4454eae95f4b151c7076f.md}{Canonical record: 47f0fca36aa4454eae95f4b151c7076f}; full qualifications and provenance.}\par
\subparagraph*{II.1.3.1.11.2.2 [Feature] Encounter Context Creation}
\phantomsection
\label{navigation:II.1.3.1.11.2.2}
\addcontentsline{toc}{subparagraph}{II.1.3.1.11.2.2 [Feature] Encounter Context Creation}

\label{element:817f1046d7e84699a2e2050b84285f95}
\index{[Feature] Encounter Context Creation (II.1.3.1.11.2.2)}
{\small Capability / Feature · Canonical · current}\par

Create and maintain authoritative encounter tracking contexts for hospital admissions, outpatient visits, and community contacts.



Encounter context creation supplies an identifiable frame for information concerning a particular admission, visit or contact. Its governed tracking authority concerns the encounter representation within Harmonia; it does not make Harmonia the originator of every clinical fact associated with that encounter.


{\footnotesize \href{../../../../continuum/strategy/capability/feature/817f1046d7e84699a2e2050b84285f95.md}{Canonical record: 817f1046d7e84699a2e2050b84285f95}; full qualifications and provenance.}\par
\subparagraph*{II.1.3.1.11.2.3 [Feature] Encounter History Retrieval}
\phantomsection
\label{navigation:II.1.3.1.11.2.3}
\addcontentsline{toc}{subparagraph}{II.1.3.1.11.2.3 [Feature] Encounter History Retrieval}

\label{element:4088e1408f554143a0f1469efd1b43a7}
\index{[Feature] Encounter History Retrieval (II.1.3.1.11.2.3)}
{\small Capability / Feature · Proposed · current}\par

Expose qualified episode, movement and encounter-history context.



\textbf{Overview}



Consumers need to interpret prior care-place and episode relationships, not just the latest encounter state; historical retrieval remains bounded by source attribution and access policy.



\textbf{Boundary}



A historical view does not retrospectively alter clinical events or acquire longitudinal-record ownership.


{\footnotesize \href{../../../../continuum/strategy/capability/feature/4088e1408f554143a0f1469efd1b43a7.md}{Canonical record: 4088e1408f554143a0f1469efd1b43a7}; full qualifications and provenance.}\par
\subparagraph*{II.1.3.1.11.2.4 [Feature] Encounter State Progression}
\phantomsection
\label{navigation:II.1.3.1.11.2.4}
\addcontentsline{toc}{subparagraph}{II.1.3.1.11.2.4 [Feature] Encounter State Progression}

\label{element:c049039bf57e4e4aa97aab5967196383}
\index{[Feature] Encounter State Progression (II.1.3.1.11.2.4)}
{\small Capability / Feature · Canonical · current}\par

Track encounter status transitions (PLANNED → ARRIVED → IN\_PROGRESS → ON\_LEAVE → DISCHARGED).



Encounter progression preserves the difference between planned, arrived, active, leave and discharged circumstances. These states make administrative activity interpretable without implying a clinical outcome. The source's ON\_LEAVE transition detail remains incomplete in the corresponding Process and is not repaired here.


{\footnotesize \href{../../../../continuum/strategy/capability/feature/c049039bf57e4e4aa97aab5967196383.md}{Canonical record: c049039bf57e4e4aa97aab5967196383}; full qualifications and provenance.}\par
\subsubsection*{II.1.3.1.12 [CT2] Funding \& Coverage Administration}
\phantomsection
\label{navigation:II.1.3.1.12}
\addcontentsline{toc}{subsubsection}{II.1.3.1.12 [CT2] Funding \& Coverage Administration}

{\small Historical predecessor: R.2.62 \hyperref[element:58beafcbcd084331be710eb00d27f6cf]{[CT2] Funding \& Coverage Administration}; Canonical ID: \texttt{\seqsplit{58beafcbcd084331be710eb00d27f6cf}}. Superseded historical identity — direct correspondence · superseded.}\par
\label{element:640432d413cf47928ea2b3d221f92313}
\index{[CT2] Funding \& Coverage Administration (II.1.3.1.12)}
{\small Capability · Canonical · current}\par

Funding \& Coverage Administration represents externally supplied insurance coverage and eligibility flags needed as care context. It allows those administrative conditions to be available without implying that Harmonia determines entitlements. The established Adjacent standing and absence of Features bound the current responsibility to this information contribution.



\textbf{Overview}



Funding \& Coverage Administration represents externally supplied insurance coverage and eligibility flags needed as care context. It allows those administrative conditions to be available without implying that Harmonia determines entitlements. The established Adjacent standing and absence of Features bound the current responsibility to this information contribution.



\textbf{Intent}



Funding \& Coverage Administration represents externally supplied insurance coverage and eligibility flags needed as care context. It allows those administrative conditions to be available without implying that Harmonia determines entitlements. The established Adjacent standing and absence of Features bound the current responsibility to this information contribution.



\textbf{Boundary}



Information administration does not confer authority over the represented person, organisation, service, place, device or clinical activity.



\medskip\hrule\medskip

{\footnotesize \href{../../../../continuum/strategy/capability/640432d413cf47928ea2b3d221f92313.md}{Canonical record: 640432d413cf47928ea2b3d221f92313}; full qualifications and provenance.}\par
\subsubsection*{II.1.3.1.13 [CT2] Medication Administration}
\phantomsection
\label{navigation:II.1.3.1.13}
\addcontentsline{toc}{subsubsection}{II.1.3.1.13 [CT2] Medication Administration}

{\small Historical predecessor: R.2.63 \hyperref[element:5a38bf2527a8404bac714121f2dc76a9]{[CT2] Medication Administration}; Canonical ID: \texttt{\seqsplit{5a38bf2527a8404bac714121f2dc76a9}}. Superseded historical identity — direct correspondence · superseded.}\par
\label{element:c907e2a95629435b9951941ad4e2ebe9}
\index{[CT2] Medication Administration (II.1.3.1.13)}
{\small Capability · Canonical · current}\par

In this enabling catalogue, Medication Administration integrates medication orders, dispensing evidence and point-of-care administration records. It preserves the distinct events through which intended therapy becomes supplied and, where reported, administered. This administrative integration supports a coherent medication account across systems without itself dispensing or administering a medicine.



\textbf{Overview}



In this enabling catalogue, Medication Administration integrates medication orders, dispensing evidence and point-of-care administration records. It preserves the distinct events through which intended therapy becomes supplied and, where reported, administered. This administrative integration supports a coherent medication account across systems without itself dispensing or administering a medicine.



\textbf{Intent}



In this enabling catalogue, Medication Administration integrates medication orders, dispensing evidence and point-of-care administration records. It preserves the distinct events through which intended therapy becomes supplied and, where reported, administered. This administrative integration supports a coherent medication account across systems without itself dispensing or administering a medicine.



\textbf{Boundary}



Information administration does not confer authority over the represented person, organisation, service, place, device or clinical activity.



\medskip\hrule\medskip

{\footnotesize \href{../../../../continuum/strategy/capability/c907e2a95629435b9951941ad4e2ebe9.md}{Canonical record: c907e2a95629435b9951941ad4e2ebe9}; full qualifications and provenance.}\par
\paragraph*{II.1.3.1.13.1 [CT3] Medication Order \& Dispense Management}
\phantomsection
\label{navigation:II.1.3.1.13.1}
\addcontentsline{toc}{paragraph}{II.1.3.1.13.1 [CT3] Medication Order \& Dispense Management}

\label{element:9874098228764b62b1f640ea749b742d}
\index{[CT3] Medication Order \& Dispense Management (II.1.3.1.13.1)}
{\small Capability · Proposed · current}\par

Administer the established medication order \& dispense management information and coordination responsibility while preserving its source authority.



\textbf{Overview}



Administer the established medication order \& dispense management information and coordination responsibility while preserving its source authority. The reviewed hierarchy assigns this enduring ability to its stated business-enabling context.



\textbf{Intent}



Administer the established medication order \& dispense management information and coordination responsibility while preserving its source authority.



\textbf{Boundary}



Information administration and coordination preserve the authority of the responsible source, clinical actor and governed domain; they confer no additional real-world management authority.


{\footnotesize \href{../../../../continuum/strategy/capability/9874098228764b62b1f640ea749b742d.md}{Canonical record: 9874098228764b62b1f640ea749b742d}; full qualifications and provenance.}\par
\subparagraph*{II.1.3.1.13.1.1 [Feature] Dispense Event Tracking}
\phantomsection
\label{navigation:II.1.3.1.13.1.1}
\addcontentsline{toc}{subparagraph}{II.1.3.1.13.1.1 [Feature] Dispense Event Tracking}

\label{element:430bf990d4194255b7ae09c2fe27367e}
\index{[Feature] Dispense Event Tracking (II.1.3.1.13.1.1)}
{\small Capability / Feature · Canonical · current}\par

Ingest medication dispense verification records from pharmacy management systems.



Dispense tracking receives the pharmacy's evidence that medication was supplied under the applicable context. It connects supply with the medication account without inferring that the patient received a dose or that Harmonia performed the professional dispensing activity.


{\footnotesize \href{../../../../continuum/strategy/capability/feature/430bf990d4194255b7ae09c2fe27367e.md}{Canonical record: 430bf990d4194255b7ae09c2fe27367e}; full qualifications and provenance.}\par
\subparagraph*{II.1.3.1.13.1.2 [Feature] Medication Administration Record (MAR) Ingestion}
\phantomsection
\label{navigation:II.1.3.1.13.1.2}
\addcontentsline{toc}{subparagraph}{II.1.3.1.13.1.2 [Feature] Medication Administration Record (MAR) Ingestion}

\label{element:779d9cb66bc84585a8dd1f456b0ee37c}
\index{[Feature] Medication Administration Record (MAR) Ingestion (II.1.3.1.13.1.2)}
{\small Capability / Feature · Canonical · current}\par

Ingest point-of-care medication administration events.



MAR ingestion captures source-reported point-of-care administration events. It preserves the crucial distinction between intended or supplied medication and an administration recorded as performed, without originating that bedside event or its clinical authority.


{\footnotesize \href{../../../../continuum/strategy/capability/feature/779d9cb66bc84585a8dd1f456b0ee37c.md}{Canonical record: 779d9cb66bc84585a8dd1f456b0ee37c}; full qualifications and provenance.}\par
\subparagraph*{II.1.3.1.13.1.3 [Feature] Medication Order Ingestion}
\phantomsection
\label{navigation:II.1.3.1.13.1.3}
\addcontentsline{toc}{subparagraph}{II.1.3.1.13.1.3 [Feature] Medication Order Ingestion}

\label{element:5247b47f181243bb975ae2cd70db65ba}
\index{[Feature] Medication Order Ingestion (II.1.3.1.13.1.3)}
{\small Capability / Feature · Canonical · current}\par

Ingest inpatient and discharge prescription orders.



Medication order ingestion preserves the intended therapy expressed by inpatient or discharge prescriptions. The received order is evidence of clinical intent, distinct from pharmacy dispensing and the patient's subsequent administration record.


{\footnotesize \href{../../../../continuum/strategy/capability/feature/5247b47f181243bb975ae2cd70db65ba.md}{Canonical record: 5247b47f181243bb975ae2cd70db65ba}; full qualifications and provenance.}\par
\paragraph*{II.1.3.1.13.2 [CT3] Medication Therapy Provision and Monitoring}
\phantomsection
\label{navigation:II.1.3.1.13.2}
\addcontentsline{toc}{paragraph}{II.1.3.1.13.2 [CT3] Medication Therapy Provision and Monitoring}

\label{element:d1821b4f6d5746f9a697dfe3b792a186}
\index{[CT3] Medication Therapy Provision and Monitoring (II.1.3.1.13.2)}
{\small Capability · Canonical · current}\par

Definition unresolved in the baseline. No definition has been synthesised from the name.



\medskip\hrule\medskip

{\footnotesize \href{../../../../continuum/strategy/capability/d1821b4f6d5746f9a697dfe3b792a186.md}{Canonical record: d1821b4f6d5746f9a697dfe3b792a186}; full qualifications and provenance.}\par
\subparagraph*{II.1.3.1.13.2.1 [Feature] Medication Administration Tracking}
\phantomsection
\label{navigation:II.1.3.1.13.2.1}
\addcontentsline{toc}{subparagraph}{II.1.3.1.13.2.1 [Feature] Medication Administration Tracking}

\label{element:6e01c72ef9c149e2ac231463810c2265}
\index{[Feature] Medication Administration Tracking (II.1.3.1.13.2.1)}
{\small Capability / Feature · Canonical · current}\par

Distribute medication therapy orders, dispensing updates, and administration records across clinical care boundaries.



Medication tracking conveys orders, dispensing changes and administration information across care boundaries while keeping their meanings separate. This supports a coherent therapy account without treating a supplied medicine as an administered dose or a distributed order as treatment performed.


{\footnotesize \href{../../../../continuum/strategy/capability/feature/6e01c72ef9c149e2ac231463810c2265.md}{Canonical record: 6e01c72ef9c149e2ac231463810c2265}; full qualifications and provenance.}\par
\subparagraph*{II.1.3.1.13.2.2 [Feature] Medication History Reconciliation Ingestion}
\phantomsection
\label{navigation:II.1.3.1.13.2.2}
\addcontentsline{toc}{subparagraph}{II.1.3.1.13.2.2 [Feature] Medication History Reconciliation Ingestion}

\label{element:561e1ef130e04f38b854665bb8fb9c77}
\index{[Feature] Medication History Reconciliation Ingestion (II.1.3.1.13.2.2)}
{\small Capability / Feature · Canonical · current}\par

Ingest reconciled medication histories compiled during clinical transitions of care.



This Feature receives medication histories already reconciled by clinical participants during care transitions. The platform preserves and connects that account; ingestion is distinct from performing the clinical reconciliation or independently deciding which regimen is correct.


{\footnotesize \href{../../../../continuum/strategy/capability/feature/561e1ef130e04f38b854665bb8fb9c77.md}{Canonical record: 561e1ef130e04f38b854665bb8fb9c77}; full qualifications and provenance.}\par
\subsubsection*{II.1.3.1.14 [CT2] Order Administration}
\phantomsection
\label{navigation:II.1.3.1.14}
\addcontentsline{toc}{subsubsection}{II.1.3.1.14 [CT2] Order Administration}

{\small Historical predecessor: R.2.65 \hyperref[element:ebb735a2915744819d4001d9b8425cc6]{[CT2] Order Administration}; Canonical ID: \texttt{\seqsplit{ebb735a2915744819d4001d9b8425cc6}}. Superseded historical identity — direct correspondence · superseded.}\par
\label{element:c3c5abd7987f4433a46cbfed8167914c}
\index{[CT2] Order Administration (II.1.3.1.14)}
{\small Capability · Canonical · current}\par

Order Administration manages the closed loop around a request for defined fulfilment. It receives the request, resolves and routes to a destination, distinguishes delivery from business disposition, tracks progress and associates the eventual outcome. This makes the difference between an order sent and an order fulfilled visible throughout the requester's interaction.



\textbf{Overview}



Order Administration manages the closed loop around a request for defined fulfilment. It receives the request, resolves and routes to a destination, distinguishes delivery from business disposition, tracks progress and associates the eventual outcome. This makes the difference between an order sent and an order fulfilled visible throughout the requester's interaction.



\textbf{Intent}



Order Administration manages the closed loop around a request for defined fulfilment. It receives the request, resolves and routes to a destination, distinguishes delivery from business disposition, tracks progress and associates the eventual outcome. This makes the difference between an order sent and an order fulfilled visible throughout the requester's interaction.



\textbf{Boundary}



Information administration does not confer authority over the represented person, organisation, service, place, device or clinical activity.



\medskip\hrule\medskip

{\footnotesize \href{../../../../continuum/strategy/capability/c3c5abd7987f4433a46cbfed8167914c.md}{Canonical record: c3c5abd7987f4433a46cbfed8167914c}; full qualifications and provenance.}\par
\paragraph*{II.1.3.1.14.1 [CT3] Closed-Loop Order Progression}
\phantomsection
\label{navigation:II.1.3.1.14.1}
\addcontentsline{toc}{paragraph}{II.1.3.1.14.1 [CT3] Closed-Loop Order Progression}

{\small Historical predecessor: R.2.64 \hyperref[element:2ad29cebc1e94053a300d869e7797232]{[CT3] Closed-Loop Order Progression}; Canonical ID: \texttt{\seqsplit{2ad29cebc1e94053a300d869e7797232}}. Superseded historical identity — direct correspondence · superseded.}\par
\label{element:2f79ebd7392c46ef92def6f25bc9c7ce}
\index{[CT3] Closed-Loop Order Progression (II.1.3.1.14.1)}
{\small Capability · Proposed · current}\par

Administer the traceable progression of an order from request through destination, delivery, acknowledgement, fulfilment and outcome association.



\textbf{Overview}



Administer the traceable progression of an order from request through destination, delivery, acknowledgement, fulfilment and outcome association. The reviewed hierarchy assigns this enduring ability to its stated business-enabling context.



\textbf{Intent}



Administer the traceable progression of an order from request through destination, delivery, acknowledgement, fulfilment and outcome association.



\textbf{Boundary}



Information administration and coordination preserve the authority of the responsible source, clinical actor and governed domain; they confer no additional real-world management authority.


{\footnotesize \href{../../../../continuum/strategy/capability/2f79ebd7392c46ef92def6f25bc9c7ce.md}{Canonical record: 2f79ebd7392c46ef92def6f25bc9c7ce}; full qualifications and provenance.}\par
\subparagraph*{II.1.3.1.14.1.1 [Feature] Order Cancellation \& Modification Coordination}
\phantomsection
\label{navigation:II.1.3.1.14.1.1}
\addcontentsline{toc}{subparagraph}{II.1.3.1.14.1.1 [Feature] Order Cancellation \& Modification Coordination}

\label{element:6d12d9ddf4f8425ca60b289f3356f683}
\index{[Feature] Order Cancellation \& Modification Coordination (II.1.3.1.14.1.1)}
{\small Capability / Feature · Canonical · current}\par

Coordinate order cancellation and amendment requests, ensuring synchronization with performing systems.



Cancellation and amendment coordination carries a requested change into the performing-system context. It preserves the difference between asking to change an order and evidence that the change took effect, rather than silently rewriting the original request as though no prior intent existed.


{\footnotesize \href{../../../../continuum/strategy/capability/feature/6d12d9ddf4f8425ca60b289f3356f683.md}{Canonical record: 6d12d9ddf4f8425ca60b289f3356f683}; full qualifications and provenance.}\par
\subparagraph*{II.1.3.1.14.1.2 [Feature] Order Closed-Loop Progression Tracking}
\phantomsection
\label{navigation:II.1.3.1.14.1.2}
\addcontentsline{toc}{subparagraph}{II.1.3.1.14.1.2 [Feature] Order Closed-Loop Progression Tracking}

\label{element:9e63463942b94a6ba3bc045a04d3580b}
\index{[Feature] Order Closed-Loop Progression Tracking (II.1.3.1.14.1.2)}
{\small Capability / Feature · Canonical · current}\par

Track order progression across destination determination, routing, delivery, receipt acknowledgement (distinguishing technical delivery acknowledgement from business acknowledgement), execution progress, completion, and outcome association.



Closed-loop tracking keeps delivery, technical acknowledgement, business disposition, execution and returned outcome distinct. Its purpose is to expose what has actually been established at each stage, preventing a successful handoff from being reported as eventual fulfilment.


{\footnotesize \href{../../../../continuum/strategy/capability/feature/9e63463942b94a6ba3bc045a04d3580b.md}{Canonical record: 9e63463942b94a6ba3bc045a04d3580b}; full qualifications and provenance.}\par
\subparagraph*{II.1.3.1.14.1.3 [Feature] Order Destination Resolution \& Routing}
\phantomsection
\label{navigation:II.1.3.1.14.1.3}
\addcontentsline{toc}{subparagraph}{II.1.3.1.14.1.3 [Feature] Order Destination Resolution \& Routing}

\label{element:9f8f1542efc64e89ad226689291f2635}
\index{[Feature] Order Destination Resolution \& Routing (II.1.3.1.14.1.3)}
{\small Capability / Feature · Canonical · current}\par

Resolve fulfillment endpoints and reliably route orders to performing service systems.



Destination resolution connects an order with the service endpoint intended to fulfil it. Routing advances the request towards that destination without proving execution or changing the original clinical intent simply because a transport route was chosen.


{\footnotesize \href{../../../../continuum/strategy/capability/feature/9f8f1542efc64e89ad226689291f2635.md}{Canonical record: 9f8f1542efc64e89ad226689291f2635}; full qualifications and provenance.}\par
\subparagraph*{II.1.3.1.14.1.4 [Feature] Order Outcome Association}
\phantomsection
\label{navigation:II.1.3.1.14.1.4}
\addcontentsline{toc}{subparagraph}{II.1.3.1.14.1.4 [Feature] Order Outcome Association}

\label{element:4ad2e11162f34775adfea01b949d2803}
\index{[Feature] Order Outcome Association (II.1.3.1.14.1.4)}
{\small Capability / Feature · Proposed · current}\par

Associate an externally established outcome with the applicable bounded order and its progression.



\textbf{Overview}



Execution completion and returned outcome are different checkpoints. This Feature makes the result link explicit without equating an order acknowledgement with an outcome.



\textbf{Boundary}



Does not interpret diagnostic findings or declare fulfilment from transport delivery alone.


{\footnotesize \href{../../../../continuum/strategy/capability/feature/4ad2e11162f34775adfea01b949d2803.md}{Canonical record: 4ad2e11162f34775adfea01b949d2803}; full qualifications and provenance.}\par
\subparagraph*{II.1.3.1.14.1.5 [Feature] Order Request Ingestion}
\phantomsection
\label{navigation:II.1.3.1.14.1.5}
\addcontentsline{toc}{subparagraph}{II.1.3.1.14.1.5 [Feature] Order Request Ingestion}

\label{element:fc1f8168fd6949beb810e4ff57a2368e}
\index{[Feature] Order Request Ingestion (II.1.3.1.14.1.5)}
{\small Capability / Feature · Canonical · current}\par

Ingest clinical diagnostic, medication, and procedure order requests with validated ordering clinician context.



Order ingestion establishes the received bounded fulfilment request with its ordering-clinician context. It begins the information account needed for a closed loop while leaving the performing service's business acceptance and execution distinguishable from technical ingress acceptance.


{\footnotesize \href{../../../../continuum/strategy/capability/feature/fc1f8168fd6949beb810e4ff57a2368e.md}{Canonical record: fc1f8168fd6949beb810e4ff57a2368e}; full qualifications and provenance.}\par
\subsubsection*{II.1.3.1.15 [CT2] Procedure Administration}
\phantomsection
\label{navigation:II.1.3.1.15}
\addcontentsline{toc}{subsubsection}{II.1.3.1.15 [CT2] Procedure Administration}

{\small Historical predecessor: R.2.66 \hyperref[element:466a4057001846bb8ae0481453d3e410]{[CT2] Procedure Administration}; Canonical ID: \texttt{\seqsplit{466a4057001846bb8ae0481453d3e410}}. Superseded historical identity — direct correspondence · superseded.}\par
\label{element:649dd7cc26094e539386a050086dfb91}
\index{[CT2] Procedure Administration (II.1.3.1.15)}
{\small Capability · Canonical · current}\par

Procedure Administration integrates procedural bookings and the documentation of procedures performed. It keeps the scheduled case distinct from operative, anaesthetic and post-procedure information received later. That separation supports continuity across scheduling and clinical documentation without inferring that a booking establishes performance.



\textbf{Overview}



Procedure Administration integrates procedural bookings and the documentation of procedures performed. It keeps the scheduled case distinct from operative, anaesthetic and post-procedure information received later. That separation supports continuity across scheduling and clinical documentation without inferring that a booking establishes performance.



\textbf{Intent}



Procedure Administration integrates procedural bookings and the documentation of procedures performed. It keeps the scheduled case distinct from operative, anaesthetic and post-procedure information received later. That separation supports continuity across scheduling and clinical documentation without inferring that a booking establishes performance.



\textbf{Boundary}



Information administration does not confer authority over the represented person, organisation, service, place, device or clinical activity.



\medskip\hrule\medskip

{\footnotesize \href{../../../../continuum/strategy/capability/649dd7cc26094e539386a050086dfb91.md}{Canonical record: 649dd7cc26094e539386a050086dfb91}; full qualifications and provenance.}\par
\paragraph*{II.1.3.1.15.1 [CT3] Procedure \& Surgical Administration}
\phantomsection
\label{navigation:II.1.3.1.15.1}
\addcontentsline{toc}{paragraph}{II.1.3.1.15.1 [CT3] Procedure \& Surgical Administration}

\label{element:3c4a53dd8c8f468f83ca5d92397dd7e8}
\index{[CT3] Procedure \& Surgical Administration (II.1.3.1.15.1)}
{\small Capability · Proposed · current}\par

Administer the established procedure \& surgical administration information and coordination responsibility while preserving its source authority.



\textbf{Overview}



Administer the established procedure \& surgical administration information and coordination responsibility while preserving its source authority. The reviewed hierarchy assigns this enduring ability to its stated business-enabling context.



\textbf{Intent}



Administer the established procedure \& surgical administration information and coordination responsibility while preserving its source authority.



\textbf{Boundary}



Information administration and coordination preserve the authority of the responsible source, clinical actor and governed domain; they confer no additional real-world management authority.


{\footnotesize \href{../../../../continuum/strategy/capability/3c4a53dd8c8f468f83ca5d92397dd7e8.md}{Canonical record: 3c4a53dd8c8f468f83ca5d92397dd7e8}; full qualifications and provenance.}\par
\subparagraph*{II.1.3.1.15.1.1 [Feature] Procedural Documentation Ingestion}
\phantomsection
\label{navigation:II.1.3.1.15.1.1}
\addcontentsline{toc}{subparagraph}{II.1.3.1.15.1.1 [Feature] Procedural Documentation Ingestion}

\label{element:324aefc4578f4cae934101ae5ad74ff8}
\index{[Feature] Procedural Documentation Ingestion (II.1.3.1.15.1.1)}
{\small Capability / Feature · Canonical · current}\par

Ingest operative reports, anesthesia records, and post-procedure summaries.



Procedural documentation ingestion preserves the operative, anaesthetic and post-procedure account supplied by the clinical service. It connects performed-care documentation with the wider record without giving the integration process authority to originate or judge the clinical findings.


{\footnotesize \href{../../../../continuum/strategy/capability/feature/324aefc4578f4cae934101ae5ad74ff8.md}{Canonical record: 324aefc4578f4cae934101ae5ad74ff8}; full qualifications and provenance.}\par
\subparagraph*{II.1.3.1.15.1.2 [Feature] Procedure Booking Ingestion}
\phantomsection
\label{navigation:II.1.3.1.15.1.2}
\addcontentsline{toc}{subparagraph}{II.1.3.1.15.1.2 [Feature] Procedure Booking Ingestion}

\label{element:12e0f64cff2043339be17fb38708247a}
\index{[Feature] Procedure Booking Ingestion (II.1.3.1.15.1.2)}
{\small Capability / Feature · Canonical · current}\par

Ingest scheduled surgical case bookings and procedural requests.



Procedure booking ingestion makes scheduled cases and procedural requests available in the wider care context. It conveys intended or arranged activity while retaining the distinction between a booked procedure and one actually performed.


{\footnotesize \href{../../../../continuum/strategy/capability/feature/12e0f64cff2043339be17fb38708247a.md}{Canonical record: 12e0f64cff2043339be17fb38708247a}; full qualifications and provenance.}\par
\subsubsection*{II.1.3.1.16 [CT2] Referral Administration}
\phantomsection
\label{navigation:II.1.3.1.16}
\addcontentsline{toc}{subsubsection}{II.1.3.1.16 [CT2] Referral Administration}

{\small Historical predecessor: R.2.68 \hyperref[element:5b0ccf4b30054145952c298b348ed266]{[CT2] Referral Administration}; Canonical ID: \texttt{\seqsplit{5b0ccf4b30054145952c298b348ed266}}. Superseded historical identity — direct correspondence · superseded.}\par
\label{element:ae0de25600ee416aa24d2a167b64f230}
\index{[CT2] Referral Administration (II.1.3.1.16)}
{\small Capability · Canonical · current}\par

Referral Administration receives referral material, associates the supporting clinical context and tracks the recipient's explicit disposition and outcome. Its purpose is to keep the request and its progression intelligible across care boundaries. It distinguishes a request for another participant's involvement from an order for bounded fulfilment and from actual acceptance of care responsibility.



\textbf{Overview}



Referral Administration receives referral material, associates the supporting clinical context and tracks the recipient's explicit disposition and outcome. Its purpose is to keep the request and its progression intelligible across care boundaries. It distinguishes a request for another participant's involvement from an order for bounded fulfilment and from actual acceptance of care responsibility.



\textbf{Intent}



Referral Administration receives referral material, associates the supporting clinical context and tracks the recipient's explicit disposition and outcome. Its purpose is to keep the request and its progression intelligible across care boundaries. It distinguishes a request for another participant's involvement from an order for bounded fulfilment and from actual acceptance of care responsibility.



\textbf{Boundary}



Information administration does not confer authority over the represented person, organisation, service, place, device or clinical activity.



\medskip\hrule\medskip

{\footnotesize \href{../../../../continuum/strategy/capability/ae0de25600ee416aa24d2a167b64f230.md}{Canonical record: ae0de25600ee416aa24d2a167b64f230}; full qualifications and provenance.}\par
\paragraph*{II.1.3.1.16.1 [CT3] Referral Intake \& Progression}
\phantomsection
\label{navigation:II.1.3.1.16.1}
\addcontentsline{toc}{paragraph}{II.1.3.1.16.1 [CT3] Referral Intake \& Progression}

{\small Historical predecessor: R.2.67 \hyperref[element:7180f6265d5c4687a001220c6126e301]{[CT3] Referral Intake \& Progression}; Canonical ID: \texttt{\seqsplit{7180f6265d5c4687a001220c6126e301}}. Superseded historical identity — direct correspondence · superseded.}\par
\label{element:d11bdb9464b2476894aa372ab9b929a7}
\index{[CT3] Referral Intake \& Progression (II.1.3.1.16.1)}
{\small Capability · Proposed · current}\par

Receive and correlate referral requests, supporting information, destinations, disposition and progression through responsible care systems.



\textbf{Overview}



Receive and correlate referral requests, supporting information, destinations, disposition and progression through responsible care systems. The reviewed hierarchy assigns this enduring ability to its stated business-enabling context.



\textbf{Intent}



Receive and correlate referral requests, supporting information, destinations, disposition and progression through responsible care systems.



\textbf{Boundary}



Information administration and coordination preserve the authority of the responsible source, clinical actor and governed domain; they confer no additional real-world management authority.


{\footnotesize \href{../../../../continuum/strategy/capability/d11bdb9464b2476894aa372ab9b929a7.md}{Canonical record: d11bdb9464b2476894aa372ab9b929a7}; full qualifications and provenance.}\par
\subparagraph*{II.1.3.1.16.1.1 [Feature] Referral Destination Coordination}
\phantomsection
\label{navigation:II.1.3.1.16.1.1}
\addcontentsline{toc}{subparagraph}{II.1.3.1.16.1.1 [Feature] Referral Destination Coordination}

\label{element:7d7ddffcda9d4854a02be5c3265ad775}
\index{[Feature] Referral Destination Coordination (II.1.3.1.16.1.1)}
{\small Capability / Feature · Proposed · current}\par

Coordinate qualified referral destination resolution and convey the request under its applicable disposition contract.



\textbf{Overview}



Referral intake and status capture leave a distinct routing responsibility: finding the intended service and preserving redirection without turning technical delivery into assumed care responsibility.



\textbf{Boundary}



Does not make clinical triage decisions or infer Transfer of Care from acceptance.


{\footnotesize \href{../../../../continuum/strategy/capability/feature/7d7ddffcda9d4854a02be5c3265ad775.md}{Canonical record: 7d7ddffcda9d4854a02be5c3265ad775}; full qualifications and provenance.}\par
\subparagraph*{II.1.3.1.16.1.2 [Feature] Referral Document Ingestion}
\phantomsection
\label{navigation:II.1.3.1.16.1.2}
\addcontentsline{toc}{subparagraph}{II.1.3.1.16.1.2 [Feature] Referral Document Ingestion}

\label{element:1d1c5f8e8ca1498fb68dba601a7079dc}
\index{[Feature] Referral Document Ingestion (II.1.3.1.16.1.2)}
{\small Capability / Feature · Canonical · current}\par

Ingest structured clinical referral messages across integration boundaries.



Referral ingestion establishes the received referral material within the integration boundary. It preserves the request for another participant's involvement without treating successful receipt as clinical acceptance, responsibility assumption or a promise that care will be delivered.


{\footnotesize \href{../../../../continuum/strategy/capability/feature/1d1c5f8e8ca1498fb68dba601a7079dc.md}{Canonical record: 1d1c5f8e8ca1498fb68dba601a7079dc}; full qualifications and provenance.}\par
\subparagraph*{II.1.3.1.16.1.3 [Feature] Referral Status \& Outcome Tracking}
\phantomsection
\label{navigation:II.1.3.1.16.1.3}
\addcontentsline{toc}{subparagraph}{II.1.3.1.16.1.3 [Feature] Referral Status \& Outcome Tracking}

\label{element:836fa704f7b24bfc8f7eeb72a36bb897}
\index{[Feature] Referral Status \& Outcome Tracking (II.1.3.1.16.1.3)}
{\small Capability / Feature · Canonical · current}\par

Track referral operational progression and communicate referral outcomes (acceptance, decline, or redirection) to referring providers.



Status and outcome tracking makes explicit referral dispositions visible to the referring participant. Acceptance, decline and redirection remain meaningful alternatives; a transport acknowledgement does not fill in a missing business decision or establish a universal referral progression.


{\footnotesize \href{../../../../continuum/strategy/capability/feature/836fa704f7b24bfc8f7eeb72a36bb897.md}{Canonical record: 836fa704f7b24bfc8f7eeb72a36bb897}; full qualifications and provenance.}\par
\subparagraph*{II.1.3.1.16.1.4 [Feature] Referral Supporting Information}
\phantomsection
\label{navigation:II.1.3.1.16.1.4}
\addcontentsline{toc}{subparagraph}{II.1.3.1.16.1.4 [Feature] Referral Supporting Information}

\label{element:c8a50c0e932847d7b7d0289101aacd98}
\index{[Feature] Referral Supporting Information (II.1.3.1.16.1.4)}
{\small Capability / Feature · Canonical · current}\par

Collate and associate supporting clinical documentation, diagnostic history, and patient context with referral requests.



Supporting information gives the referral its relevant history and clinical context. Association makes the material available for the recipient's consideration while preserving source provenance and the distinction between evidence supplied and the recipient's own clinical decision.


{\footnotesize \href{../../../../continuum/strategy/capability/feature/c8a50c0e932847d7b7d0289101aacd98.md}{Canonical record: c8a50c0e932847d7b7d0289101aacd98}; full qualifications and provenance.}\par
\subsubsection*{II.1.3.1.17 [CT2] Scheduling Administration}
\phantomsection
\label{navigation:II.1.3.1.17}
\addcontentsline{toc}{subsubsection}{II.1.3.1.17 [CT2] Scheduling Administration}

{\small Historical predecessor: R.2.70 \hyperref[element:1a6630940ffb44fcad4543b93a367b67]{[CT2] Scheduling Administration}; Canonical ID: \texttt{\seqsplit{1a6630940ffb44fcad4543b93a367b67}}. Superseded historical identity — direct correspondence · superseded.}\par
\label{element:35fb3744c900470f9895056246a13fb1}
\index{[CT2] Scheduling Administration (II.1.3.1.17)}
{\small Capability · Canonical · current}\par

Scheduling Administration integrates appointment information established by departmental scheduling systems. It receives scheduled events and propagates cancellation, rescheduling and attendance changes so that care-coordination contexts do not diverge. Its concern is the coherence of scheduling information across participants rather than originating every booking decision.



\textbf{Overview}



Scheduling Administration integrates appointment information established by departmental scheduling systems. It receives scheduled events and propagates cancellation, rescheduling and attendance changes so that care-coordination contexts do not diverge. Its concern is the coherence of scheduling information across participants rather than originating every booking decision.



\textbf{Intent}



Scheduling Administration integrates appointment information established by departmental scheduling systems. It receives scheduled events and propagates cancellation, rescheduling and attendance changes so that care-coordination contexts do not diverge. Its concern is the coherence of scheduling information across participants rather than originating every booking decision.



\textbf{Boundary}



Information administration does not confer authority over the represented person, organisation, service, place, device or clinical activity.



\medskip\hrule\medskip

{\footnotesize \href{../../../../continuum/strategy/capability/35fb3744c900470f9895056246a13fb1.md}{Canonical record: 35fb3744c900470f9895056246a13fb1}; full qualifications and provenance.}\par
\paragraph*{II.1.3.1.17.1 [CT3] Appointment \& Booking Administration}
\phantomsection
\label{navigation:II.1.3.1.17.1}
\addcontentsline{toc}{paragraph}{II.1.3.1.17.1 [CT3] Appointment \& Booking Administration}

{\small Historical predecessor: R.2.69 \hyperref[element:cdc16dd0f84049f8af3f1ec76f1e1428]{[CT3] Appointment \& Booking Administration}; Canonical ID: \texttt{\seqsplit{cdc16dd0f84049f8af3f1ec76f1e1428}}. Superseded historical identity — direct correspondence · superseded.}\par
\label{element:7b3ba9bcac564dd4a7a08631581db2ca}
\index{[CT3] Appointment \& Booking Administration (II.1.3.1.17.1)}
{\small Capability · Proposed · current}\par

Administer externally supplied appointment and booking information and communicate changes in their scheduling status.



\textbf{Overview}



Administer externally supplied appointment and booking information and communicate changes in their scheduling status. The reviewed hierarchy assigns this enduring ability to its stated business-enabling context.



\textbf{Intent}



Administer externally supplied appointment and booking information and communicate changes in their scheduling status.



\textbf{Boundary}



Information administration and coordination preserve the authority of the responsible source, clinical actor and governed domain; they confer no additional real-world management authority.


{\footnotesize \href{../../../../continuum/strategy/capability/7b3ba9bcac564dd4a7a08631581db2ca.md}{Canonical record: 7b3ba9bcac564dd4a7a08631581db2ca}; full qualifications and provenance.}\par
\subparagraph*{II.1.3.1.17.1.1 [Feature] Appointment Notification Ingestion}
\phantomsection
\label{navigation:II.1.3.1.17.1.1}
\addcontentsline{toc}{subparagraph}{II.1.3.1.17.1.1 [Feature] Appointment Notification Ingestion}

\label{element:1d5bc2b1742544728357fbea42a55ec4}
\index{[Feature] Appointment Notification Ingestion (II.1.3.1.17.1.1)}
{\small Capability / Feature · Canonical · current}\par

Ingest scheduled appointment events from departmental scheduling systems.



Appointment ingestion receives scheduling facts established by departmental systems. It supports shared awareness of planned care opportunities without claiming authority over the underlying booking decision or proving that attendance and clinical fulfilment will occur.


{\footnotesize \href{../../../../continuum/strategy/capability/feature/1d5bc2b1742544728357fbea42a55ec4.md}{Canonical record: 1d5bc2b1742544728357fbea42a55ec4}; full qualifications and provenance.}\par
\subparagraph*{II.1.3.1.17.1.2 [Feature] Appointment Status Synchronization}
\phantomsection
\label{navigation:II.1.3.1.17.1.2}
\addcontentsline{toc}{subparagraph}{II.1.3.1.17.1.2 [Feature] Appointment Status Synchronization}

\label{element:a7a9932600d84434a1a37005f726747d}
\index{[Feature] Appointment Status Synchronization (II.1.3.1.17.1.2)}
{\small Capability / Feature · Canonical · current}\par

Propagate booking cancellations, reschedules, and attendances across care coordination systems.



Status synchronisation keeps cancellation, rescheduling and attendance information coherent across coordination contexts. The purpose is propagation of the externally established booking change rather than independent scheduling or inference of a clinical outcome from appointment state.


{\footnotesize \href{../../../../continuum/strategy/capability/feature/a7a9932600d84434a1a37005f726747d.md}{Canonical record: a7a9932600d84434a1a37005f726747d}; full qualifications and provenance.}\par
\subsection*{II.1.3.2 [CT1] Care Context Integration Enablement [superseded]}
\phantomsection
\label{navigation:II.1.3.2}
\addcontentsline{toc}{subsection}{II.1.3.2 [CT1] Care Context Integration Enablement [superseded]}

\label{element:a84a7e96e7874a4b87d18c53b1b5cf29}
\index{[CT1] Care Context Integration Enablement (II.1.3.2)}
{\small Capability · Proposed · superseded}\par
\href{../../../../continuum/strategy/capability/a84a7e96e7874a4b87d18c53b1b5cf29.md}{Historical canonical definition and provenance}.\par
\subsubsection*{II.1.3.2.1 [CT2] Facility \& Episodic Care Integration [superseded]}
\phantomsection
\label{navigation:II.1.3.2.1}
\addcontentsline{toc}{subsubsection}{II.1.3.2.1 [CT2] Facility \& Episodic Care Integration [superseded]}

\label{element:c1eb24ed7e36417cb426288ff57512c0}
\index{[CT2] Facility \& Episodic Care Integration (II.1.3.2.1)}
{\small Capability · Proposed · superseded}\par
\href{../../../../continuum/strategy/capability/c1eb24ed7e36417cb426288ff57512c0.md}{Historical canonical definition and provenance}.\par
\subsubsection*{II.1.3.2.2 [CT2] Longitudinal \& Community Care Integration [superseded]}
\phantomsection
\label{navigation:II.1.3.2.2}
\addcontentsline{toc}{subsubsection}{II.1.3.2.2 [CT2] Longitudinal \& Community Care Integration [superseded]}

\label{element:8430d50adc504ecc959f15b7cbbcd095}
\index{[CT2] Longitudinal \& Community Care Integration (II.1.3.2.2)}
{\small Capability · Proposed · superseded}\par
\href{../../../../continuum/strategy/capability/8430d50adc504ecc959f15b7cbbcd095.md}{Historical canonical definition and provenance}.\par
\subsubsection*{II.1.3.2.3 [CT2] Specialist Clinical Context Integration [superseded]}
\phantomsection
\label{navigation:II.1.3.2.3}
\addcontentsline{toc}{subsubsection}{II.1.3.2.3 [CT2] Specialist Clinical Context Integration [superseded]}

\label{element:062802a425d14530b0ce9434c5980d85}
\index{[CT2] Specialist Clinical Context Integration (II.1.3.2.3)}
{\small Capability · Proposed · superseded}\par
\href{../../../../continuum/strategy/capability/062802a425d14530b0ce9434c5980d85.md}{Historical canonical definition and provenance}.\par
\subsection*{II.1.3.3 [CT1] Clinical Service Dependency Management}
\phantomsection
\label{navigation:II.1.3.3}
\addcontentsline{toc}{subsection}{II.1.3.3 [CT1] Clinical Service Dependency Management}

{\small Historical predecessor: R.2.88 \hyperref[element:984ee17123eb4f92b3e55e55d24459b8]{[CT1] Clinical Service Dependency Management}; Canonical ID: \texttt{\seqsplit{984ee17123eb4f92b3e55e55d24459b8}}. Superseded historical identity — direct correspondence · superseded.}\par
\label{element:232db9b8b47c450cb0aa9acaf539c09d}
\index{[CT1] Clinical Service Dependency Management (II.1.3.3)}
{\small Capability · Proposed · current}\par

Identify and track operational and infrastructure dependencies necessary for clinical services to function, including care-place and environmental readiness, devices, materiel, forms, workstations, connectivity, infrastructure and basic working materials.



\textbf{Overview}



Identify and track operational and infrastructure dependencies necessary for clinical services to function, including care-place and environmental readiness, devices, materiel, forms, workstations, connectivity, infrastructure and basic working materials. The reviewed hierarchy assigns this enduring ability to its stated business-enabling context.



\textbf{Intent}



Identify and track operational and infrastructure dependencies necessary for clinical services to function, including care-place and environmental readiness, devices, materiel, forms, workstations, connectivity, infrastructure and basic working materials.



\textbf{Boundary}



Information administration and coordination preserve the authority of the responsible source, clinical actor and governed domain; they confer no additional real-world management authority.


{\footnotesize \href{../../../../continuum/strategy/capability/232db9b8b47c450cb0aa9acaf539c09d.md}{Canonical record: 232db9b8b47c450cb0aa9acaf539c09d}; full qualifications and provenance.}\par
\subsubsection*{II.1.3.3.1 [CT2] Acute Care}
\phantomsection
\label{navigation:II.1.3.3.1}
\addcontentsline{toc}{subsubsection}{II.1.3.3.1 [CT2] Acute Care}

{\small Historical predecessor: R.2.72 \hyperref[element:7b591957f7424424abc80f699e9f1960]{[CT2] Acute Care}; Canonical ID: \texttt{\seqsplit{7b591957f7424424abc80f699e9f1960}}. Superseded historical identity — direct correspondence · superseded.}\par
\label{element:9bb988c526404522b83aeed5069e0eb5}
\index{[CT2] Acute Care (II.1.3.3.1)}
{\small Capability · Proposed · current}\par

Identify and track the operational and infrastructure dependencies necessary for acute care services to function, retaining each dependency provider's responsibility.



\textbf{Overview}



Identify and track the operational and infrastructure dependencies necessary for acute care services to function, retaining each dependency provider's responsibility. The reviewed hierarchy assigns this enduring ability to its stated business-enabling context.



\textbf{Intent}



Identify and track the operational and infrastructure dependencies necessary for acute care services to function, retaining each dependency provider's responsibility.



\textbf{Boundary}



Information administration and coordination preserve the authority of the responsible source, clinical actor and governed domain; they confer no additional real-world management authority.


{\footnotesize \href{../../../../continuum/strategy/capability/9bb988c526404522b83aeed5069e0eb5.md}{Canonical record: 9bb988c526404522b83aeed5069e0eb5}; full qualifications and provenance.}\par
\subsubsection*{II.1.3.3.2 [CT2] Allied Health}
\phantomsection
\label{navigation:II.1.3.3.2}
\addcontentsline{toc}{subsubsection}{II.1.3.3.2 [CT2] Allied Health}

{\small Historical predecessor: R.2.73 \hyperref[element:6cd45c27d0c344b9807174dbd812d572]{[CT2] Allied Health}; Canonical ID: \texttt{\seqsplit{6cd45c27d0c344b9807174dbd812d572}}. Superseded historical identity — direct correspondence · superseded.}\par
\label{element:63cd714a3e2a4354aa5383b189fea566}
\index{[CT2] Allied Health (II.1.3.3.2)}
{\small Capability · Canonical · current}\par

Identify and track the operational and infrastructure dependencies necessary for allied health services to function, retaining each dependency provider's responsibility.



\textbf{Overview}



Identify and track the operational and infrastructure dependencies necessary for allied health services to function, retaining each dependency provider's responsibility.



\textbf{Intent}



Identify and track the operational and infrastructure dependencies necessary for allied health services to function, retaining each dependency provider's responsibility.



\textbf{Boundary}



Information administration does not confer authority over the represented person, organisation, service, place, device or clinical activity.



\medskip\hrule\medskip

{\footnotesize \href{../../../../continuum/strategy/capability/63cd714a3e2a4354aa5383b189fea566.md}{Canonical record: 63cd714a3e2a4354aa5383b189fea566}; full qualifications and provenance.}\par
\subsubsection*{II.1.3.3.3 [CT2] Ambulatory Care}
\phantomsection
\label{navigation:II.1.3.3.3}
\addcontentsline{toc}{subsubsection}{II.1.3.3.3 [CT2] Ambulatory Care}

{\small Historical predecessor: R.2.74 \hyperref[element:e996b61b6e7b4619913f427aed2b55b1]{[CT2] Ambulatory Care}; Canonical ID: \texttt{\seqsplit{e996b61b6e7b4619913f427aed2b55b1}}. Superseded historical identity — direct correspondence · superseded.}\par
\label{element:f74d7c3a39c540c4b462576e8f09fc58}
\index{[CT2] Ambulatory Care (II.1.3.3.3)}
{\small Capability · Canonical · current}\par

Identify and track the operational and infrastructure dependencies necessary for ambulatory care services to function, retaining each dependency provider's responsibility.



\textbf{Overview}



Identify and track the operational and infrastructure dependencies necessary for ambulatory care services to function, retaining each dependency provider's responsibility.



\textbf{Intent}



Identify and track the operational and infrastructure dependencies necessary for ambulatory care services to function, retaining each dependency provider's responsibility.



\textbf{Boundary}



Information administration does not confer authority over the represented person, organisation, service, place, device or clinical activity.



\medskip\hrule\medskip

{\footnotesize \href{../../../../continuum/strategy/capability/f74d7c3a39c540c4b462576e8f09fc58.md}{Canonical record: f74d7c3a39c540c4b462576e8f09fc58}; full qualifications and provenance.}\par
\subsubsection*{II.1.3.3.4 [CT2] Clinical Assessment}
\phantomsection
\label{navigation:II.1.3.3.4}
\addcontentsline{toc}{subsubsection}{II.1.3.3.4 [CT2] Clinical Assessment}

{\small Historical predecessor: R.2.75 \hyperref[element:1c3b553c87704e27ba211fbbd02d95cb]{[CT2] Clinical Assessment}; Canonical ID: \texttt{\seqsplit{1c3b553c87704e27ba211fbbd02d95cb}}. Superseded historical identity — direct correspondence · superseded.}\par
\label{element:a7d84c95e48c453797419813ca80a66c}
\index{[CT2] Clinical Assessment (II.1.3.3.4)}
{\small Capability · Canonical · current}\par

Identify and track the operational and infrastructure dependencies necessary for clinical assessment services to function, retaining each dependency provider's responsibility.



\textbf{Overview}



Identify and track the operational and infrastructure dependencies necessary for clinical assessment services to function, retaining each dependency provider's responsibility.



\textbf{Intent}



Identify and track the operational and infrastructure dependencies necessary for clinical assessment services to function, retaining each dependency provider's responsibility.



\textbf{Boundary}



Information administration does not confer authority over the represented person, organisation, service, place, device or clinical activity.



\medskip\hrule\medskip

{\footnotesize \href{../../../../continuum/strategy/capability/a7d84c95e48c453797419813ca80a66c.md}{Canonical record: a7d84c95e48c453797419813ca80a66c}; full qualifications and provenance.}\par
\subsubsection*{II.1.3.3.5 [CT2] Community Care}
\phantomsection
\label{navigation:II.1.3.3.5}
\addcontentsline{toc}{subsubsection}{II.1.3.3.5 [CT2] Community Care}

{\small Historical predecessor: R.2.76 \hyperref[element:4b7177f3f40449518fba381acae6ba1e]{[CT2] Community Care}; Canonical ID: \texttt{\seqsplit{4b7177f3f40449518fba381acae6ba1e}}. Superseded historical identity — direct correspondence · superseded.}\par
\label{element:942aa78cad004ef99de6b3aae8ddd014}
\index{[CT2] Community Care (II.1.3.3.5)}
{\small Capability · Proposed · current}\par

Identify and track the operational and infrastructure dependencies necessary for community care services to function, retaining each dependency provider's responsibility.



\textbf{Overview}



Identify and track the operational and infrastructure dependencies necessary for community care services to function, retaining each dependency provider's responsibility. The reviewed hierarchy assigns this enduring ability to its stated business-enabling context.



\textbf{Intent}



Identify and track the operational and infrastructure dependencies necessary for community care services to function, retaining each dependency provider's responsibility.



\textbf{Boundary}



Information administration and coordination preserve the authority of the responsible source, clinical actor and governed domain; they confer no additional real-world management authority.


{\footnotesize \href{../../../../continuum/strategy/capability/942aa78cad004ef99de6b3aae8ddd014.md}{Canonical record: 942aa78cad004ef99de6b3aae8ddd014}; full qualifications and provenance.}\par
\subsubsection*{II.1.3.3.6 [CT2] Diagnostic Services}
\phantomsection
\label{navigation:II.1.3.3.6}
\addcontentsline{toc}{subsubsection}{II.1.3.3.6 [CT2] Diagnostic Services}

{\small Historical predecessor: R.2.77 \hyperref[element:03a57e0c9ea34f7fb7b78e5f5c85379f]{[CT2] Diagnostic Services}; Canonical ID: \texttt{\seqsplit{03a57e0c9ea34f7fb7b78e5f5c85379f}}. Superseded historical identity — direct correspondence · superseded.}\par
\label{element:b825ebab42f444d8b738b4c1f489f785}
\index{[CT2] Diagnostic Services (II.1.3.3.6)}
{\small Capability · Proposed · current}\par

Identify and track the operational and infrastructure dependencies necessary for diagnostic services services to function, retaining each dependency provider's responsibility.



\textbf{Overview}



Identify and track the operational and infrastructure dependencies necessary for diagnostic services services to function, retaining each dependency provider's responsibility. The reviewed hierarchy assigns this enduring ability to its stated business-enabling context.



\textbf{Intent}



Identify and track the operational and infrastructure dependencies necessary for diagnostic services services to function, retaining each dependency provider's responsibility.



\textbf{Boundary}



Information administration and coordination preserve the authority of the responsible source, clinical actor and governed domain; they confer no additional real-world management authority.


{\footnotesize \href{../../../../continuum/strategy/capability/b825ebab42f444d8b738b4c1f489f785.md}{Canonical record: b825ebab42f444d8b738b4c1f489f785}; full qualifications and provenance.}\par
\subsubsection*{II.1.3.3.7 [CT2] Emergency Care}
\phantomsection
\label{navigation:II.1.3.3.7}
\addcontentsline{toc}{subsubsection}{II.1.3.3.7 [CT2] Emergency Care}

{\small Historical predecessor: R.2.78 \hyperref[element:e5edbb54258c490aaba0816d980dc624]{[CT2] Emergency Care}; Canonical ID: \texttt{\seqsplit{e5edbb54258c490aaba0816d980dc624}}. Superseded historical identity — direct correspondence · superseded.}\par
\label{element:aa088f0647ad4b78966fe28fa80af674}
\index{[CT2] Emergency Care (II.1.3.3.7)}
{\small Capability · Proposed · current}\par

Identify and track the operational and infrastructure dependencies necessary for emergency care services to function, retaining each dependency provider's responsibility.



\textbf{Overview}



Identify and track the operational and infrastructure dependencies necessary for emergency care services to function, retaining each dependency provider's responsibility. The reviewed hierarchy assigns this enduring ability to its stated business-enabling context.



\textbf{Intent}



Identify and track the operational and infrastructure dependencies necessary for emergency care services to function, retaining each dependency provider's responsibility.



\textbf{Boundary}



Information administration and coordination preserve the authority of the responsible source, clinical actor and governed domain; they confer no additional real-world management authority.


{\footnotesize \href{../../../../continuum/strategy/capability/aa088f0647ad4b78966fe28fa80af674.md}{Canonical record: aa088f0647ad4b78966fe28fa80af674}; full qualifications and provenance.}\par
\subsubsection*{II.1.3.3.8 [CT2] Inpatient Care}
\phantomsection
\label{navigation:II.1.3.3.8}
\addcontentsline{toc}{subsubsection}{II.1.3.3.8 [CT2] Inpatient Care}

{\small Historical predecessor: R.2.79 \hyperref[element:bcc40a14465e48aca42c9505a6853081]{[CT2] Inpatient Care}; Canonical ID: \texttt{\seqsplit{bcc40a14465e48aca42c9505a6853081}}. Superseded historical identity — direct correspondence · superseded.}\par
\label{element:b39fb49c01b545a5b51d3313531efb1b}
\index{[CT2] Inpatient Care (II.1.3.3.8)}
{\small Capability · Proposed · current}\par

Identify and track the operational and infrastructure dependencies necessary for inpatient care services to function, retaining each dependency provider's responsibility.



\textbf{Overview}



Identify and track the operational and infrastructure dependencies necessary for inpatient care services to function, retaining each dependency provider's responsibility. The reviewed hierarchy assigns this enduring ability to its stated business-enabling context.



\textbf{Intent}



Identify and track the operational and infrastructure dependencies necessary for inpatient care services to function, retaining each dependency provider's responsibility.



\textbf{Boundary}



Information administration and coordination preserve the authority of the responsible source, clinical actor and governed domain; they confer no additional real-world management authority.


{\footnotesize \href{../../../../continuum/strategy/capability/b39fb49c01b545a5b51d3313531efb1b.md}{Canonical record: b39fb49c01b545a5b51d3313531efb1b}; full qualifications and provenance.}\par
\subsubsection*{II.1.3.3.9 [CT2] Medication Therapy}
\phantomsection
\label{navigation:II.1.3.3.9}
\addcontentsline{toc}{subsubsection}{II.1.3.3.9 [CT2] Medication Therapy}

{\small Historical predecessor: R.2.80 \hyperref[element:db064e0cd4794479be1374e97a3d7e63]{[CT2] Medication Therapy}; Canonical ID: \texttt{\seqsplit{db064e0cd4794479be1374e97a3d7e63}}. Superseded historical identity — direct correspondence · superseded.}\par
\label{element:bb04996164c94e0c827a020432f4e498}
\index{[CT2] Medication Therapy (II.1.3.3.9)}
{\small Capability · Proposed · current}\par

Identify and track the operational and infrastructure dependencies necessary for medication therapy services to function, retaining each dependency provider's responsibility.



\textbf{Overview}



Identify and track the operational and infrastructure dependencies necessary for medication therapy services to function, retaining each dependency provider's responsibility. The reviewed hierarchy assigns this enduring ability to its stated business-enabling context.



\textbf{Intent}



Identify and track the operational and infrastructure dependencies necessary for medication therapy services to function, retaining each dependency provider's responsibility.



\textbf{Boundary}



Information administration and coordination preserve the authority of the responsible source, clinical actor and governed domain; they confer no additional real-world management authority.


{\footnotesize \href{../../../../continuum/strategy/capability/bb04996164c94e0c827a020432f4e498.md}{Canonical record: bb04996164c94e0c827a020432f4e498}; full qualifications and provenance.}\par
\subsubsection*{II.1.3.3.10 [CT2] Outreach Care}
\phantomsection
\label{navigation:II.1.3.3.10}
\addcontentsline{toc}{subsubsection}{II.1.3.3.10 [CT2] Outreach Care}

{\small Historical predecessor: R.2.81 \hyperref[element:75f8b32a3d5048ccbdfb30d61ae73dad]{[CT2] Outreach Care}; Canonical ID: \texttt{\seqsplit{75f8b32a3d5048ccbdfb30d61ae73dad}}. Superseded historical identity — direct correspondence · superseded.}\par
\label{element:b9ddbe2dd6a7485ba9cd3482215ea09e}
\index{[CT2] Outreach Care (II.1.3.3.10)}
{\small Capability · Proposed · current}\par

Identify and track the operational and infrastructure dependencies necessary for outreach care services to function, retaining each dependency provider's responsibility.



\textbf{Overview}



Identify and track the operational and infrastructure dependencies necessary for outreach care services to function, retaining each dependency provider's responsibility. The reviewed hierarchy assigns this enduring ability to its stated business-enabling context.



\textbf{Intent}



Identify and track the operational and infrastructure dependencies necessary for outreach care services to function, retaining each dependency provider's responsibility.



\textbf{Boundary}



Information administration and coordination preserve the authority of the responsible source, clinical actor and governed domain; they confer no additional real-world management authority.


{\footnotesize \href{../../../../continuum/strategy/capability/b9ddbe2dd6a7485ba9cd3482215ea09e.md}{Canonical record: b9ddbe2dd6a7485ba9cd3482215ea09e}; full qualifications and provenance.}\par
\subsubsection*{II.1.3.3.11 [CT2] Preventive Care}
\phantomsection
\label{navigation:II.1.3.3.11}
\addcontentsline{toc}{subsubsection}{II.1.3.3.11 [CT2] Preventive Care}

{\small Historical predecessor: R.2.82 \hyperref[element:97f0d1acc37b4ca994c6c1ca006f2daa]{[CT2] Preventive Care}; Canonical ID: \texttt{\seqsplit{97f0d1acc37b4ca994c6c1ca006f2daa}}. Superseded historical identity — direct correspondence · superseded.}\par
\label{element:dfdc4f0273ba4910848b3b65c9ef6ab2}
\index{[CT2] Preventive Care (II.1.3.3.11)}
{\small Capability · Proposed · current}\par

Identify and track the operational and infrastructure dependencies necessary for preventive care services to function, retaining each dependency provider's responsibility.



\textbf{Overview}



Identify and track the operational and infrastructure dependencies necessary for preventive care services to function, retaining each dependency provider's responsibility. The reviewed hierarchy assigns this enduring ability to its stated business-enabling context.



\textbf{Intent}



Identify and track the operational and infrastructure dependencies necessary for preventive care services to function, retaining each dependency provider's responsibility.



\textbf{Boundary}



Information administration and coordination preserve the authority of the responsible source, clinical actor and governed domain; they confer no additional real-world management authority.


{\footnotesize \href{../../../../continuum/strategy/capability/dfdc4f0273ba4910848b3b65c9ef6ab2.md}{Canonical record: dfdc4f0273ba4910848b3b65c9ef6ab2}; full qualifications and provenance.}\par
\subsubsection*{II.1.3.3.12 [CT2] Primary Care}
\phantomsection
\label{navigation:II.1.3.3.12}
\addcontentsline{toc}{subsubsection}{II.1.3.3.12 [CT2] Primary Care}

{\small Historical predecessor: R.2.83 \hyperref[element:fe9d58b5cfe14e8cb7a83a1ac6686d65]{[CT2] Primary Care}; Canonical ID: \texttt{\seqsplit{fe9d58b5cfe14e8cb7a83a1ac6686d65}}. Superseded historical identity — direct correspondence · superseded.}\par
\label{element:d118673d77d74c8cb8cec24f03adddd8}
\index{[CT2] Primary Care (II.1.3.3.12)}
{\small Capability · Proposed · current}\par

Identify and track the operational and infrastructure dependencies necessary for primary care services to function, retaining each dependency provider's responsibility.



\textbf{Overview}



Identify and track the operational and infrastructure dependencies necessary for primary care services to function, retaining each dependency provider's responsibility. The reviewed hierarchy assigns this enduring ability to its stated business-enabling context.



\textbf{Intent}



Identify and track the operational and infrastructure dependencies necessary for primary care services to function, retaining each dependency provider's responsibility.



\textbf{Boundary}



Information administration and coordination preserve the authority of the responsible source, clinical actor and governed domain; they confer no additional real-world management authority.


{\footnotesize \href{../../../../continuum/strategy/capability/d118673d77d74c8cb8cec24f03adddd8.md}{Canonical record: d118673d77d74c8cb8cec24f03adddd8}; full qualifications and provenance.}\par
\subsubsection*{II.1.3.3.13 [CT2] Rehabilitation}
\phantomsection
\label{navigation:II.1.3.3.13}
\addcontentsline{toc}{subsubsection}{II.1.3.3.13 [CT2] Rehabilitation}

{\small Historical predecessor: R.2.84 \hyperref[element:28ec21d5918a49aaac12740c69f42d77]{[CT2] Rehabilitation}; Canonical ID: \texttt{\seqsplit{28ec21d5918a49aaac12740c69f42d77}}. Superseded historical identity — direct correspondence · superseded.}\par
\label{element:99ed51abcc5e4e8ba3b3d0161b17544a}
\index{[CT2] Rehabilitation (II.1.3.3.13)}
{\small Capability · Canonical · current}\par

Identify and track the operational and infrastructure dependencies necessary for rehabilitation services to function, retaining each dependency provider's responsibility.



\textbf{Overview}



Identify and track the operational and infrastructure dependencies necessary for rehabilitation services to function, retaining each dependency provider's responsibility.



\textbf{Intent}



Identify and track the operational and infrastructure dependencies necessary for rehabilitation services to function, retaining each dependency provider's responsibility.



\textbf{Boundary}



Information administration does not confer authority over the represented person, organisation, service, place, device or clinical activity.



\medskip\hrule\medskip

{\footnotesize \href{../../../../continuum/strategy/capability/99ed51abcc5e4e8ba3b3d0161b17544a.md}{Canonical record: 99ed51abcc5e4e8ba3b3d0161b17544a}; full qualifications and provenance.}\par
\subsubsection*{II.1.3.3.14 [CT2] Remote \& Virtual Care}
\phantomsection
\label{navigation:II.1.3.3.14}
\addcontentsline{toc}{subsubsection}{II.1.3.3.14 [CT2] Remote \& Virtual Care}

{\small Historical predecessor: R.2.85 \hyperref[element:75e2b447659c4fdf94107d94758337ed]{[CT2] Remote \& Virtual Care}; Canonical ID: \texttt{\seqsplit{75e2b447659c4fdf94107d94758337ed}}. Superseded historical identity — direct correspondence · superseded.}\par
\label{element:30dff22d40b64da29ede7c41673db085}
\index{[CT2] Remote \& Virtual Care (II.1.3.3.14)}
{\small Capability · Proposed · current}\par

Identify and track the operational and infrastructure dependencies necessary for remote \& virtual care services to function, retaining each dependency provider's responsibility.



\textbf{Overview}



Identify and track the operational and infrastructure dependencies necessary for remote \& virtual care services to function, retaining each dependency provider's responsibility. The reviewed hierarchy assigns this enduring ability to its stated business-enabling context.



\textbf{Intent}



Identify and track the operational and infrastructure dependencies necessary for remote \& virtual care services to function, retaining each dependency provider's responsibility.



\textbf{Boundary}



Information administration and coordination preserve the authority of the responsible source, clinical actor and governed domain; they confer no additional real-world management authority.


{\footnotesize \href{../../../../continuum/strategy/capability/30dff22d40b64da29ede7c41673db085.md}{Canonical record: 30dff22d40b64da29ede7c41673db085}; full qualifications and provenance.}\par
\subsubsection*{II.1.3.3.15 [CT2] Surgical \& Procedural Care}
\phantomsection
\label{navigation:II.1.3.3.15}
\addcontentsline{toc}{subsubsection}{II.1.3.3.15 [CT2] Surgical \& Procedural Care}

{\small Historical predecessor: R.2.86 \hyperref[element:8bc2a106165a4702b5f72fe46a206b69]{[CT2] Surgical \& Procedural Care}; Canonical ID: \texttt{\seqsplit{8bc2a106165a4702b5f72fe46a206b69}}. Superseded historical identity — direct correspondence · superseded.}\par
\label{element:c9a26d1aa1014545b9ae759b663c807b}
\index{[CT2] Surgical \& Procedural Care (II.1.3.3.15)}
{\small Capability · Canonical · current}\par

Identify and track the operational and infrastructure dependencies necessary for surgical \& procedural care services to function, retaining each dependency provider's responsibility.



\textbf{Overview}



Identify and track the operational and infrastructure dependencies necessary for surgical \& procedural care services to function, retaining each dependency provider's responsibility.



\textbf{Intent}



Identify and track the operational and infrastructure dependencies necessary for surgical \& procedural care services to function, retaining each dependency provider's responsibility.



\textbf{Boundary}



Information administration does not confer authority over the represented person, organisation, service, place, device or clinical activity.



\medskip\hrule\medskip

{\footnotesize \href{../../../../continuum/strategy/capability/c9a26d1aa1014545b9ae759b663c807b.md}{Canonical record: c9a26d1aa1014545b9ae759b663c807b}; full qualifications and provenance.}\par
\subsubsection*{II.1.3.3.16 [CT2] Therapeutic Services}
\phantomsection
\label{navigation:II.1.3.3.16}
\addcontentsline{toc}{subsubsection}{II.1.3.3.16 [CT2] Therapeutic Services}

{\small Historical predecessor: R.2.87 \hyperref[element:54032b52ee444bfe8b611fb9e8e2ba10]{[CT2] Therapeutic Services}; Canonical ID: \texttt{\seqsplit{54032b52ee444bfe8b611fb9e8e2ba10}}. Superseded historical identity — direct correspondence · superseded.}\par
\label{element:b512b6d326004cb7b91a627fe6dbd8be}
\index{[CT2] Therapeutic Services (II.1.3.3.16)}
{\small Capability · Canonical · current}\par

Identify and track the operational and infrastructure dependencies necessary for therapeutic services services to function, retaining each dependency provider's responsibility.



\textbf{Overview}



Identify and track the operational and infrastructure dependencies necessary for therapeutic services services to function, retaining each dependency provider's responsibility.



\textbf{Intent}



Identify and track the operational and infrastructure dependencies necessary for therapeutic services services to function, retaining each dependency provider's responsibility.



\textbf{Boundary}



Information administration does not confer authority over the represented person, organisation, service, place, device or clinical activity.



\medskip\hrule\medskip

{\footnotesize \href{../../../../continuum/strategy/capability/b512b6d326004cb7b91a627fe6dbd8be.md}{Canonical record: b512b6d326004cb7b91a627fe6dbd8be}; full qualifications and provenance.}\par
\subsection*{II.1.3.4 [CT1] Entity Information Enablement}
\phantomsection
\label{navigation:II.1.3.4}
\addcontentsline{toc}{subsection}{II.1.3.4 [CT1] Entity Information Enablement}

{\small Historical predecessor: R.2.114 \hyperref[element:e4d959d6c71d498c9ce00d2be4c048e3]{[CT1] Entity Information Enablement}; Canonical ID: \texttt{\seqsplit{e4d959d6c71d498c9ce00d2be4c048e3}}. Superseded historical identity — direct correspondence · superseded.}\par
\label{element:b5a92874f8ac4509a30af5253590edc1}
\index{[CT1] Entity Information Enablement (II.1.3.4)}
{\small Capability · Proposed · current}\par

Administer identifying, descriptive and qualified relationship information concerning healthcare participants, organisations, locations, services, materiel and devices.



\textbf{Overview}



Administer identifying, descriptive and qualified relationship information concerning healthcare participants, organisations, locations, services, materiel and devices.



\textbf{Intent}



Administer identifying, descriptive and qualified relationship information concerning healthcare participants, organisations, locations, services, materiel and devices.



\textbf{Boundary}



Registry information does not establish clinical privilege, service readiness, occupancy or originating national identity authority.


{\footnotesize \href{../../../../continuum/strategy/capability/b5a92874f8ac4509a30af5253590edc1.md}{Canonical record: b5a92874f8ac4509a30af5253590edc1}; full qualifications and provenance.}\par
\subsubsection*{II.1.3.4.1 [CT2] Client Administration}
\phantomsection
\label{navigation:II.1.3.4.1}
\addcontentsline{toc}{subsubsection}{II.1.3.4.1 [CT2] Client Administration}

{\small Historical predecessor: R.2.93 \hyperref[element:70d2a44a0442441f8b306b82b297bb8a]{[CT2] Client Administration}; Canonical ID: \texttt{\seqsplit{70d2a44a0442441f8b306b82b297bb8a}}. Superseded historical identity — direct correspondence · superseded.}\par
\label{element:f0885472cced4bccbcce1568bf0360dc}
\index{[CT2] Client Administration (II.1.3.4.1)}
{\small Capability · Canonical · current}\par

Client Administration provides the representation and demographic context needed to recognise individuals receiving care. It brings identity, healthcare-subject context, relationships and privacy into the same administrative frame while preserving their different responsibilities. The established partial chain to Person Identity is retained; the catalogue's other contextual descriptions do not supply a complete capability hierarchy.



\textbf{Overview}



Client Administration provides the representation and demographic context needed to recognise individuals receiving care. It brings identity, healthcare-subject context, relationships and privacy into the same administrative frame while preserving their different responsibilities. The established partial chain to Person Identity is retained; the catalogue's other contextual descriptions do not supply a complete capability hierarchy.



\textbf{Intent}



Client Administration provides the representation and demographic context needed to recognise individuals receiving care. It brings identity, healthcare-subject context, relationships and privacy into the same administrative frame while preserving their different responsibilities. The established partial chain to Person Identity is retained; the catalogue's other contextual descriptions do not supply a complete capability hierarchy.



\textbf{Boundary}



Information administration does not confer authority over the represented person, organisation, service, place, device or clinical activity.



\medskip\hrule\medskip

{\small Recorded unresolved matters: source-qualified matters; exact scope remains in the canonical record.}\par
{\footnotesize \href{../../../../continuum/strategy/capability/f0885472cced4bccbcce1568bf0360dc.md}{Canonical record: f0885472cced4bccbcce1568bf0360dc}; full qualifications and provenance.}\par
\paragraph*{II.1.3.4.1.1 [CT3] Client Privacy Guardrail Management}
\phantomsection
\label{navigation:II.1.3.4.1.1}
\addcontentsline{toc}{paragraph}{II.1.3.4.1.1 [CT3] Client Privacy Guardrail Management}

{\small Historical predecessor: R.2.89 \hyperref[element:d2e6078b5db041cb904e0cadcebce6eb]{[CT3] Client Privacy Guardrail Management}; Canonical ID: \texttt{\seqsplit{d2e6078b5db041cb904e0cadcebce6eb}}. Superseded historical identity — direct correspondence · superseded.}\par
\label{element:7577a6e1815243648dd5356c21cca151}
\index{[CT3] Client Privacy Guardrail Management (II.1.3.4.1.1)}
{\small Capability · Canonical · current}\par

Enforcing client-level privacy preferences, consent directives, and confidentiality markers.



Definition — reviewed candidate



Manage client-associated privacy constraints, preferences, permissions, restrictions and consent directives, and evaluate their applicability to information use.



\textbf{Overview}



Manage client-associated privacy constraints, preferences, permissions, restrictions and consent directives, and evaluate their applicability to information use.



\textbf{Intent}



Manage client-associated privacy constraints, preferences, permissions, restrictions and consent directives, and evaluate their applicability to information use.



\textbf{Boundary}



Owns client-associated constraints and consent directives; general enterprise privacy policy remains outside this capability.



\medskip\hrule\medskip

{\footnotesize \href{../../../../continuum/strategy/capability/7577a6e1815243648dd5356c21cca151.md}{Canonical record: 7577a6e1815243648dd5356c21cca151}; full qualifications and provenance.}\par
\subparagraph*{II.1.3.4.1.1.1 [Feature] Confidentiality Flag Enforcement}
\phantomsection
\label{navigation:II.1.3.4.1.1.1}
\addcontentsline{toc}{subparagraph}{II.1.3.4.1.1.1 [Feature] Confidentiality Flag Enforcement}

\label{element:1c18e7c2ccf640c9afc2525503b505a6}
\index{[Feature] Confidentiality Flag Enforcement (II.1.3.4.1.1.1)}
{\small Capability / Feature · Canonical · current}\par

Enforce restricted access flags (e.g., VIP, sensitive diagnosis, restricted record) across query and presentation interfaces.



Confidentiality enforcement makes a restricted-record condition effective across discovery and presentation paths. Its purpose is consistent governed treatment of sensitive information, rather than a flag displayed in one interface while other paths disclose the same content without restriction.


{\footnotesize \href{../../../../continuum/strategy/capability/feature/1c18e7c2ccf640c9afc2525503b505a6.md}{Canonical record: 1c18e7c2ccf640c9afc2525503b505a6}; full qualifications and provenance.}\par
\subparagraph*{II.1.3.4.1.1.2 [Feature] Consent Directive Evaluation}
\phantomsection
\label{navigation:II.1.3.4.1.1.2}
\addcontentsline{toc}{subparagraph}{II.1.3.4.1.1.2 [Feature] Consent Directive Evaluation}

\label{element:baf6dc6ef38a4256a9a67a65fd3f31ba}
\index{[Feature] Consent Directive Evaluation (II.1.3.4.1.1.2)}
{\small Capability / Feature · Canonical · current}\par

Evaluate client consent directives against requested information disclosure and exchange actions.



Consent evaluation brings the client's applicable directives into consideration of a requested disclosure or exchange. It distinguishes the recorded directive from the particular action being evaluated, preserving the need for established policy and authority rather than treating consent data as a universal permission token.


{\footnotesize \href{../../../../continuum/strategy/capability/feature/baf6dc6ef38a4256a9a67a65fd3f31ba.md}{Canonical record: baf6dc6ef38a4256a9a67a65fd3f31ba}; full qualifications and provenance.}\par
\paragraph*{II.1.3.4.1.2 [CT3] Client Relationship Administration}
\phantomsection
\label{navigation:II.1.3.4.1.2}
\addcontentsline{toc}{paragraph}{II.1.3.4.1.2 [CT3] Client Relationship Administration}

{\small Historical predecessor: R.2.90 \hyperref[element:2c799ec0c4c74c97a75c12e949887b95]{[CT3] Client Relationship Administration}; Canonical ID: \texttt{\seqsplit{2c799ec0c4c74c97a75c12e949887b95}}. Superseded historical identity — direct correspondence · superseded.}\par
\label{element:d682c11d3eb04770b1f20be478fd475e}
\index{[CT3] Client Relationship Administration (II.1.3.4.1.2)}
{\small Capability · Canonical · current}\par

Representing relationships between the client and carers, guardians, or contacts.



Definition — reviewed candidate



Administer qualified client relationships with carers, guardians, contacts and other nominated care participants.



\textbf{Overview}



Administer qualified client relationships with carers, guardians, contacts and other nominated care participants.



\textbf{Intent}



Administer qualified client relationships with carers, guardians, contacts and other nominated care participants.



\textbf{Boundary}



Information administration does not confer authority over the represented person, organisation, service, place, device or clinical activity.



\medskip\hrule\medskip

{\footnotesize \href{../../../../continuum/strategy/capability/d682c11d3eb04770b1f20be478fd475e.md}{Canonical record: d682c11d3eb04770b1f20be478fd475e}; full qualifications and provenance.}\par
\subparagraph*{II.1.3.4.1.2.1 [Feature] Next-of-Kin \& Guardian Association}
\phantomsection
\label{navigation:II.1.3.4.1.2.1}
\addcontentsline{toc}{subparagraph}{II.1.3.4.1.2.1 [Feature] Next-of-Kin \& Guardian Association}

\label{element:05de74a8a7dc4334ba60543a228e73cb}
\index{[Feature] Next-of-Kin \& Guardian Association (II.1.3.4.1.2.1)}
{\small Capability / Feature · Canonical · current}\par

Associate nominated emergency contacts, legal guardians, and care agents with the client record.



This Feature records the nominated people and capacities through which a client can be contacted or represented. Association preserves who the contact or guardian is in relation to the client; it does not make every nominated person an unrestricted recipient of clinical information.


{\footnotesize \href{../../../../continuum/strategy/capability/feature/05de74a8a7dc4334ba60543a228e73cb.md}{Canonical record: 05de74a8a7dc4334ba60543a228e73cb}; full qualifications and provenance.}\par
\subparagraph*{II.1.3.4.1.2.2 [Feature] Relationship Validity Tracking}
\phantomsection
\label{navigation:II.1.3.4.1.2.2}
\addcontentsline{toc}{subparagraph}{II.1.3.4.1.2.2 [Feature] Relationship Validity Tracking}

\label{element:309603a8330545abb5db5f9889aa60d9}
\index{[Feature] Relationship Validity Tracking (II.1.3.4.1.2.2)}
{\small Capability / Feature · Canonical · current}\par

Maintain effective date intervals and authority scopes for nominated client relationships.



Relationship validity makes effective periods and authority scope part of the relationship's meaning. A contact association that was legitimate in one period or capacity need not be treated as currently valid for every purpose.


{\footnotesize \href{../../../../continuum/strategy/capability/feature/309603a8330545abb5db5f9889aa60d9.md}{Canonical record: 309603a8330545abb5db5f9889aa60d9}; full qualifications and provenance.}\par
\paragraph*{II.1.3.4.1.3 [CT3] Healthcare Subject Profile Administration}
\phantomsection
\label{navigation:II.1.3.4.1.3}
\addcontentsline{toc}{paragraph}{II.1.3.4.1.3 [CT3] Healthcare Subject Profile Administration}

{\small Historical predecessor: R.2.91 \hyperref[element:d653c3247ddd49479f985f6b435933a0]{[CT3] Healthcare Subject Profile Administration}; Canonical ID: \texttt{\seqsplit{d653c3247ddd49479f985f6b435933a0}}. Superseded historical identity — direct correspondence · superseded.}\par
\label{element:df214b5499694b768eb4a7a175a9d471}
\index{[CT3] Healthcare Subject Profile Administration (II.1.3.4.1.3)}
{\small Capability · Canonical · current}\par

Managing the core healthcare subject context bound to clinical records and operational activities.



Definition — reviewed candidate



Administer demographic and subject-of-care profile information bound to clinical records and operational activities.



\textbf{Overview}



Administer demographic and subject-of-care profile information bound to clinical records and operational activities.



\textbf{Intent}



Administer demographic and subject-of-care profile information bound to clinical records and operational activities.



\textbf{Boundary}



Information administration does not confer authority over the represented person, organisation, service, place, device or clinical activity.



\medskip\hrule\medskip

{\footnotesize \href{../../../../continuum/strategy/capability/df214b5499694b768eb4a7a175a9d471.md}{Canonical record: df214b5499694b768eb4a7a175a9d471}; full qualifications and provenance.}\par
\subparagraph*{II.1.3.4.1.3.1 [Feature] Demographic Context Governance}
\phantomsection
\label{navigation:II.1.3.4.1.3.1}
\addcontentsline{toc}{subparagraph}{II.1.3.4.1.3.1 [Feature] Demographic Context Governance}

\label{element:3d0377451a1b44d0bb41404207cdd867}
\index{[Feature] Demographic Context Governance (II.1.3.4.1.3.1)}
{\small Capability / Feature · Canonical · current}\par

Manage essential demographic attributes (date of birth, sex, vital status) with full provenance and timestamping.



Demographic governance retains the provenance and temporal context of attributes used to interpret the subject of care. A changed date, status or demographic assertion can therefore be understood as sourced information rather than an unexplained overwrite of previously known context.


{\footnotesize \href{../../../../continuum/strategy/capability/feature/3d0377451a1b44d0bb41404207cdd867.md}{Canonical record: 3d0377451a1b44d0bb41404207cdd867}; full qualifications and provenance.}\par
\subparagraph*{II.1.3.4.1.3.2 [Feature] Subject Context Binding}
\phantomsection
\label{navigation:II.1.3.4.1.3.2}
\addcontentsline{toc}{subparagraph}{II.1.3.4.1.3.2 [Feature] Subject Context Binding}

\label{element:5d956689bc714e2d8be4fff7766c0b2d}
\index{[Feature] Subject Context Binding (II.1.3.4.1.3.2)}
{\small Capability / Feature · Canonical · current}\par

Bind authenticated person identity to clinical encounters, orders, and observations as the enduring subject of care with verifiable association.



Subject binding connects resolved identity with the records and activities concerning that person. It makes the association verifiable across encounters, orders and observations rather than assuming that proximity in a message or screen establishes the enduring subject of care.


{\footnotesize \href{../../../../continuum/strategy/capability/feature/5d956689bc714e2d8be4fff7766c0b2d.md}{Canonical record: 5d956689bc714e2d8be4fff7766c0b2d}; full qualifications and provenance.}\par
\paragraph*{II.1.3.4.1.4 [CT3] Person Identity Administration}
\phantomsection
\label{navigation:II.1.3.4.1.4}
\addcontentsline{toc}{paragraph}{II.1.3.4.1.4 [CT3] Person Identity Administration}

{\small Historical predecessor: R.2.92 \hyperref[element:7550fd5b17e643118fa76d2616e1ba20]{[CT3] Person Identity Administration}; Canonical ID: \texttt{\seqsplit{7550fd5b17e643118fa76d2616e1ba20}}. Superseded historical identity — direct correspondence · superseded.}\par
\label{element:6c3de7750ef1436cad6a9533731782b6}
\index{[CT3] Person Identity Administration (II.1.3.4.1.4)}
{\small Capability · Canonical · current}\par

Resolving, maintaining, and correlating identities and identifiers assigned to an individual across different authorities.



Definition — reviewed candidate



Resolve, maintain and correlate identities and identifiers assigned to a person across issuing authorities, retaining source-qualified correction history.



\textbf{Overview}



Resolve, maintain and correlate identities and identifiers assigned to a person across issuing authorities, retaining source-qualified correction history.



\textbf{Intent}



Resolve, maintain and correlate identities and identifiers assigned to a person across issuing authorities, retaining source-qualified correction history.



\textbf{Boundary}



Information administration does not confer authority over the represented person, organisation, service, place, device or clinical activity.



\medskip\hrule\medskip

{\footnotesize \href{../../../../continuum/strategy/capability/6c3de7750ef1436cad6a9533731782b6.md}{Canonical record: 6c3de7750ef1436cad6a9533731782b6}; full qualifications and provenance.}\par
\subparagraph*{II.1.3.4.1.4.1 [Feature] Cross-Authority Identifier Correlation}
\phantomsection
\label{navigation:II.1.3.4.1.4.1}
\addcontentsline{toc}{subparagraph}{II.1.3.4.1.4.1 [Feature] Cross-Authority Identifier Correlation}

\label{element:6d74d143c7454fa0bf30743ddf7ba094}
\index{[Feature] Cross-Authority Identifier Correlation (II.1.3.4.1.4.1)}
{\small Capability / Feature · Canonical · current}\par

Correlate disparate local identifiers (e.g., hospital MRN, clinic client ID, national IHI) belonging to the same individual.



Cross-authority correlation preserves the association between identifiers used for one person by different participants. It enables information to be related across regional systems without erasing the separate identifier namespaces or claiming that Harmonia originated the identity decision.


{\footnotesize \href{../../../../continuum/strategy/capability/feature/6d74d143c7454fa0bf30743ddf7ba094.md}{Canonical record: 6d74d143c7454fa0bf30743ddf7ba094}; full qualifications and provenance.}\par
\subparagraph*{II.1.3.4.1.4.2 [Feature] Governed Identity Correction}
\phantomsection
\label{navigation:II.1.3.4.1.4.2}
\addcontentsline{toc}{subparagraph}{II.1.3.4.1.4.2 [Feature] Governed Identity Correction}

\label{element:f7ab7a24e8234129b2ab02e9abd4fba2}
\index{[Feature] Governed Identity Correction (II.1.3.4.1.4.2)}
{\small Capability / Feature · Canonical · current}\par

Apply administrative identity updates, linked-record corrections, and split/merge notifications propagated from authoritative sources.



Governed correction applies externally authoritative identity updates and their consequences to Harmonia-managed associations. Administrative review concerns the legitimate application of the source outcome; it does not grant Harmonia an EMPI merge engine or authority to select a golden person record.


{\footnotesize \href{../../../../continuum/strategy/capability/feature/f7ab7a24e8234129b2ab02e9abd4fba2.md}{Canonical record: f7ab7a24e8234129b2ab02e9abd4fba2}; full qualifications and provenance.}\par
\subparagraph*{II.1.3.4.1.4.3 [Feature] Identifier Capture}
\phantomsection
\label{navigation:II.1.3.4.1.4.3}
\addcontentsline{toc}{subparagraph}{II.1.3.4.1.4.3 [Feature] Identifier Capture}

\label{element:3cae3c4810974c5f8ee5a4c942f38b07}
\index{[Feature] Identifier Capture (II.1.3.4.1.4.3)}
{\small Capability / Feature · Accepted bounded reconciliation · current}\par

The governed source has no established definition. This absence is retained for Solution Architect review; no definition is inferred from the name.


{\small Recorded unresolved matters: definition; exact scope remains in the canonical record.}\par
{\footnotesize \href{../../../../continuum/strategy/capability/feature/3cae3c4810974c5f8ee5a4c942f38b07.md}{Canonical record: 3cae3c4810974c5f8ee5a4c942f38b07}; full qualifications and provenance.}\par
\subparagraph*{II.1.3.4.1.4.4 [Feature] Identifier Resolution}
\phantomsection
\label{navigation:II.1.3.4.1.4.4}
\addcontentsline{toc}{subparagraph}{II.1.3.4.1.4.4 [Feature] Identifier Resolution}

\label{element:ae034169e1394ccdb3a2ce9651eaf305}
\index{[Feature] Identifier Resolution (II.1.3.4.1.4.4)}
{\small Capability / Feature · Canonical · current}\par

Ingest and resolve external identifiers against registered issuing authorities and namespace systems.



Identifier Resolution interprets an identifier in its issuing-authority and namespace context, so that the same local value is not assumed to identify the same person everywhere. It provides an identifiable subject reference without deciding an authoritative master-person match.


{\footnotesize \href{../../../../continuum/strategy/capability/feature/ae034169e1394ccdb3a2ce9651eaf305.md}{Canonical record: ae034169e1394ccdb3a2ce9651eaf305}; full qualifications and provenance.}\par
\subparagraph*{II.1.3.4.1.4.5 [Feature] Identity Alias Association}
\phantomsection
\label{navigation:II.1.3.4.1.4.5}
\addcontentsline{toc}{subparagraph}{II.1.3.4.1.4.5 [Feature] Identity Alias Association}

\label{element:97721dbd017e4c2daa73e1e87a8499b9}
\index{[Feature] Identity Alias Association (II.1.3.4.1.4.5)}
{\small Capability / Feature · Canonical · current}\par

Maintain explicit associations between active identifiers and historical or demographic aliases without destructive overwriting.



Alias association keeps historical and demographic identity references navigable alongside active identifiers. The non-destructive relationship preserves how earlier information identified the person rather than rewriting that history whenever the current identity representation changes.


{\footnotesize \href{../../../../continuum/strategy/capability/feature/97721dbd017e4c2daa73e1e87a8499b9.md}{Canonical record: 97721dbd017e4c2daa73e1e87a8499b9}; full qualifications and provenance.}\par
\subparagraph*{II.1.3.4.1.4.6 [Feature] Identity Source \& Correction Lineage}
\phantomsection
\label{navigation:II.1.3.4.1.4.6}
\addcontentsline{toc}{subparagraph}{II.1.3.4.1.4.6 [Feature] Identity Source \& Correction Lineage}

\label{element:d72134e75001449a9e7ffae53f61458f}
\index{[Feature] Identity Source \& Correction Lineage (II.1.3.4.1.4.6)}
{\small Capability / Feature · Proposed · current}\par

Retain source-attributed identity assertions and externally decided correction history.



\textbf{Overview}



Identifier correlations can change when an upstream authority corrects a relationship. Reviewers must be able to distinguish the source decision, its effective standing and the managed links affected.



\textbf{Boundary}



No probabilistic matching, golden-record selection or autonomous merge authority.


{\footnotesize \href{../../../../continuum/strategy/capability/feature/d72134e75001449a9e7ffae53f61458f.md}{Canonical record: d72134e75001449a9e7ffae53f61458f}; full qualifications and provenance.}\par
\subsubsection*{II.1.3.4.2 [CT2] Clinical Device Administration}
\phantomsection
\label{navigation:II.1.3.4.2}
\addcontentsline{toc}{subsubsection}{II.1.3.4.2 [CT2] Clinical Device Administration}

{\small Historical predecessor: R.2.98 \hyperref[element:3792be10541e4d68b3f740b1ba964966]{[CT2] Clinical Device Administration}; Canonical ID: \texttt{\seqsplit{3792be10541e4d68b3f740b1ba964966}}. Superseded historical identity — direct correspondence · superseded.}\par
\label{element:31bcd6fda0d14cc8b992cdca39a3cbed}
\index{[CT2] Clinical Device Administration (II.1.3.4.2)}
{\small Capability · Canonical · current}\par

Clinical Device Administration represents device types, identifiable instances, patient or location associations and communication endpoints. This lets device-originated information be interpreted in relation to the actual device and care context, rather than merely a network address. Type, instance, assignment and connectivity each answer a different question about the source of device data.



\textbf{Overview}



Clinical Device Administration represents device types, identifiable instances, patient or location associations and communication endpoints. This lets device-originated information be interpreted in relation to the actual device and care context, rather than merely a network address. Type, instance, assignment and connectivity each answer a different question about the source of device data.



\textbf{Intent}



Clinical Device Administration represents device types, identifiable instances, patient or location associations and communication endpoints. This lets device-originated information be interpreted in relation to the actual device and care context, rather than merely a network address. Type, instance, assignment and connectivity each answer a different question about the source of device data.



\textbf{Boundary}



Information administration does not confer authority over the represented person, organisation, service, place, device or clinical activity.



\medskip\hrule\medskip

{\small Recorded unresolved matters: source-qualified matters; exact scope remains in the canonical record.}\par
{\footnotesize \href{../../../../continuum/strategy/capability/31bcd6fda0d14cc8b992cdca39a3cbed.md}{Canonical record: 31bcd6fda0d14cc8b992cdca39a3cbed}; full qualifications and provenance.}\par
\paragraph*{II.1.3.4.2.1 [CT3] Clinical Device Profile Administration}
\phantomsection
\label{navigation:II.1.3.4.2.1}
\addcontentsline{toc}{paragraph}{II.1.3.4.2.1 [CT3] Clinical Device Profile Administration}

{\small Historical predecessor: R.2.94 \hyperref[element:e8c72952da144d6883caba7fd2afb88d]{[CT3] Clinical Device Profile Administration}; Canonical ID: \texttt{\seqsplit{e8c72952da144d6883caba7fd2afb88d}}. Superseded historical identity — direct correspondence · superseded.}\par
\label{element:d0a1ba80c790476ea6042290eb7ea09e}
\index{[CT3] Clinical Device Profile Administration (II.1.3.4.2.1)}
{\small Capability · Proposed · current}\par

Administer device type, instance, identifying and descriptive information needed to recognise clinical equipment and its stated capabilities.



\textbf{Overview}



Administer device type, instance, identifying and descriptive information needed to recognise clinical equipment and its stated capabilities.



\textbf{Intent}



Administer device type, instance, identifying and descriptive information needed to recognise clinical equipment and its stated capabilities.



\textbf{Boundary}



Describes device types and capabilities without taking regulatory approval or clinical-use authority.


{\footnotesize \href{../../../../continuum/strategy/capability/d0a1ba80c790476ea6042290eb7ea09e.md}{Canonical record: d0a1ba80c790476ea6042290eb7ea09e}; full qualifications and provenance.}\par
\subparagraph*{II.1.3.4.2.1.1 [Feature] Device Serial Tracking}
\phantomsection
\label{navigation:II.1.3.4.2.1.1}
\addcontentsline{toc}{subparagraph}{II.1.3.4.2.1.1 [Feature] Device Serial Tracking}

\label{element:df3dae0cf3044059875dbb888ea6a9ba}
\index{[Feature] Device Serial Tracking (II.1.3.4.2.1.1)}
{\small Capability / Feature · Canonical · current}\par

Ingest and maintain unique device identifier (UDI) records and serial numbers.



Serial tracking makes individual hardware distinguishable within the device estate. A unique device reference supports subsequent assignment and source interpretation, while remaining separate from evidence that the device is maintained, connected or fit for a particular clinical use.


{\footnotesize \href{../../../../continuum/strategy/capability/feature/df3dae0cf3044059875dbb888ea6a9ba.md}{Canonical record: df3dae0cf3044059875dbb888ea6a9ba}; full qualifications and provenance.}\par
\subparagraph*{II.1.3.4.2.1.2 [Feature] Device Type Specification}
\phantomsection
\label{navigation:II.1.3.4.2.1.2}
\addcontentsline{toc}{subparagraph}{II.1.3.4.2.1.2 [Feature] Device Type Specification}

\label{element:ca6d91b03b56415790810846649a3dca}
\index{[Feature] Device Type Specification (II.1.3.4.2.1.2)}
{\small Capability / Feature · Canonical · current}\par

Catalog medical device models, regulatory approvals, and capabilities.



Device type specification provides model-level capabilities and regulatory context against which individual devices can be understood. It describes the kind of equipment without substituting for the release, condition or deployment of a particular instance.


{\footnotesize \href{../../../../continuum/strategy/capability/feature/ca6d91b03b56415790810846649a3dca.md}{Canonical record: ca6d91b03b56415790810846649a3dca}; full qualifications and provenance.}\par
\paragraph*{II.1.3.4.2.2 [CT3] Device Association [superseded]}
\phantomsection
\label{navigation:II.1.3.4.2.2}
\addcontentsline{toc}{paragraph}{II.1.3.4.2.2 [CT3] Device Association [superseded]}

\label{element:902daf2ce9714876bbfee52eef1a7ab0}
\index{[CT3] Device Association (II.1.3.4.2.2)}
{\small Capability · Proposed · superseded}\par
\href{../../../../continuum/strategy/capability/902daf2ce9714876bbfee52eef1a7ab0.md}{Historical canonical definition and provenance}.\par
\paragraph*{II.1.3.4.2.3 [CT3] Device Connectivity [superseded]}
\phantomsection
\label{navigation:II.1.3.4.2.3}
\addcontentsline{toc}{paragraph}{II.1.3.4.2.3 [CT3] Device Connectivity [superseded]}

\label{element:658d4b6770a943bcb5801fb7ca9b47be}
\index{[CT3] Device Connectivity (II.1.3.4.2.3)}
{\small Capability · Proposed · superseded}\par
\href{../../../../continuum/strategy/capability/658d4b6770a943bcb5801fb7ca9b47be.md}{Historical canonical definition and provenance}.\par
\paragraph*{II.1.3.4.2.4 [CT3] Device Connectivity Compliance Management}
\phantomsection
\label{navigation:II.1.3.4.2.4}
\addcontentsline{toc}{paragraph}{II.1.3.4.2.4 [CT3] Device Connectivity Compliance Management}

{\small Historical predecessor: R.2.95 \hyperref[element:9225f083d3474715a75af2a827a4016d]{[CT3] Device Connectivity Compliance Management}; Canonical ID: \texttt{\seqsplit{9225f083d3474715a75af2a827a4016d}}. Superseded historical identity — direct correspondence · superseded.}\par
\label{element:b16ab359030a4583868f3aac2221074e}
\index{[CT3] Device Connectivity Compliance Management (II.1.3.4.2.4)}
{\small Capability · Proposed · current}\par

Manage information about whether device connectivity satisfies applicable standards, policies and integration constraints.



\textbf{Overview}



Manage information about whether device connectivity satisfies applicable standards, policies and integration constraints. The reviewed hierarchy assigns this enduring ability to its stated business-enabling context.



\textbf{Intent}



Manage information about whether device connectivity satisfies applicable standards, policies and integration constraints.



\textbf{Boundary}



Endpoint registration alone does not establish compliance; conformity concerns applicable standards, policies and integration constraints.


{\footnotesize \href{../../../../continuum/strategy/capability/b16ab359030a4583868f3aac2221074e.md}{Canonical record: b16ab359030a4583868f3aac2221074e}; full qualifications and provenance.}\par
\paragraph*{II.1.3.4.2.5 [CT3] Device Instance [superseded]}
\phantomsection
\label{navigation:II.1.3.4.2.5}
\addcontentsline{toc}{paragraph}{II.1.3.4.2.5 [CT3] Device Instance [superseded]}

\label{element:dd860926038f41e082c56e38cbb5ab95}
\index{[CT3] Device Instance (II.1.3.4.2.5)}
{\small Capability · Proposed · superseded}\par
\href{../../../../continuum/strategy/capability/dd860926038f41e082c56e38cbb5ab95.md}{Historical canonical definition and provenance}.\par
\paragraph*{II.1.3.4.2.6 [CT3] Device OA\&M Metadata Management}
\phantomsection
\label{navigation:II.1.3.4.2.6}
\addcontentsline{toc}{paragraph}{II.1.3.4.2.6 [CT3] Device OA\&M Metadata Management}

{\small Historical predecessor: R.2.96 \hyperref[element:a398c72bef094dee96c96d97994449fb]{[CT3] Device OA\&M Metadata Management}; Canonical ID: \texttt{\seqsplit{a398c72bef094dee96c96d97994449fb}}. Superseded historical identity — direct correspondence · superseded.}\par
\label{element:659bb194a7ef42ab9ba8a938adb5c179}
\index{[CT3] Device OA\&M Metadata Management (II.1.3.4.2.6)}
{\small Capability · Proposed · current}\par

Manage operational and administrative metadata needed to understand how a clinical device is connected, monitored and administered.



\textbf{Overview}



Manage operational and administrative metadata needed to understand how a clinical device is connected, monitored and administered. The reviewed hierarchy assigns this enduring ability to its stated business-enabling context.



\textbf{Intent}



Manage operational and administrative metadata needed to understand how a clinical device is connected, monitored and administered.



\textbf{Boundary}



Connectivity, monitoring and administration metadata do not perform device maintenance, interpret clinical telemetry or determine conformance.


{\footnotesize \href{../../../../continuum/strategy/capability/659bb194a7ef42ab9ba8a938adb5c179.md}{Canonical record: 659bb194a7ef42ab9ba8a938adb5c179}; full qualifications and provenance.}\par
\subparagraph*{II.1.3.4.2.6.1 [Feature] Device Endpoint Registration}
\phantomsection
\label{navigation:II.1.3.4.2.6.1}
\addcontentsline{toc}{subparagraph}{II.1.3.4.2.6.1 [Feature] Device Endpoint Registration}

\label{element:d32c856559c54b41941fc51fbd019195}
\index{[Feature] Device Endpoint Registration (II.1.3.4.2.6.1)}
{\small Capability / Feature · Canonical · current}\par

Register communication interfaces and security credentials for diagnostic telemetry feeds.



Endpoint registration supplies the interfaces and credential context through which a diagnostic or telemetry feed can participate. It describes provisioned connectivity rather than proving that data has arrived, been clinically interpreted or acquired authoritative standing.


{\footnotesize \href{../../../../continuum/strategy/capability/feature/d32c856559c54b41941fc51fbd019195.md}{Canonical record: d32c856559c54b41941fc51fbd019195}; full qualifications and provenance.}\par
\paragraph*{II.1.3.4.2.7 [CT3] Device Utilisation Administration}
\phantomsection
\label{navigation:II.1.3.4.2.7}
\addcontentsline{toc}{paragraph}{II.1.3.4.2.7 [CT3] Device Utilisation Administration}

{\small Historical predecessor: R.2.97 \hyperref[element:a0af0a86667d48f4ab95ed2920e878cb]{[CT3] Device Utilisation Administration}; Canonical ID: \texttt{\seqsplit{a0af0a86667d48f4ab95ed2920e878cb}}. Superseded historical identity — direct correspondence · superseded.}\par
\label{element:1a3732b5911c45b3aece38b2db2d0f15}
\index{[CT3] Device Utilisation Administration (II.1.3.4.2.7)}
{\small Capability · Proposed · current}\par

Administer qualified device-to-subject/place allocation, location, utilisation and operational readiness information, including the time at which an association applies.



\textbf{Overview}



Administer qualified device-to-subject/place allocation, location, utilisation and operational readiness information, including the time at which an association applies. The reviewed hierarchy assigns this enduring ability to its stated business-enabling context.



\textbf{Intent}



Administer qualified device-to-subject/place allocation, location, utilisation and operational readiness information, including the time at which an association applies.



\textbf{Boundary}



Information administration and coordination preserve the authority of the responsible source, clinical actor and governed domain; they confer no additional real-world management authority.


{\footnotesize \href{../../../../continuum/strategy/capability/1a3732b5911c45b3aece38b2db2d0f15.md}{Canonical record: 1a3732b5911c45b3aece38b2db2d0f15}; full qualifications and provenance.}\par
\subparagraph*{II.1.3.4.2.7.1 [Feature] Device Association Lineage}
\phantomsection
\label{navigation:II.1.3.4.2.7.1}
\addcontentsline{toc}{subparagraph}{II.1.3.4.2.7.1 [Feature] Device Association Lineage}

\label{element:849889c1a9e4495c86a6b32e49c449e3}
\index{[Feature] Device Association Lineage (II.1.3.4.2.7.1)}
{\small Capability / Feature · Proposed · current}\par

Preserve the attributed history and effective basis of device-to-subject/place associations.



\textbf{Overview}



A telemetry record needs to be interpreted against the association applicable when it arose, rather than the device current location alone.



\textbf{Boundary}



Association history does not interpret physiological observations or prove clinical correctness.


{\footnotesize \href{../../../../continuum/strategy/capability/feature/849889c1a9e4495c86a6b32e49c449e3.md}{Canonical record: 849889c1a9e4495c86a6b32e49c449e3}; full qualifications and provenance.}\par
\subparagraph*{II.1.3.4.2.7.2 [Feature] Device-Patient/Location Assignment}
\phantomsection
\label{navigation:II.1.3.4.2.7.2}
\addcontentsline{toc}{subparagraph}{II.1.3.4.2.7.2 [Feature] Device-Patient/Location Assignment}

\label{element:891ed75bdd3c490296e5f0b7cacf6c87}
\index{[Feature] Device-Patient/Location Assignment (II.1.3.4.2.7.2)}
{\small Capability / Feature · Canonical · current}\par

Track dynamic association between a telemetry monitor, patient, and physical bed.



Dynamic device association preserves which patient or physical place a device is linked to as assignments change. This is essential context for interpreting telemetry: a stable device identity does not imply that every observation concerns the same patient over time.


{\footnotesize \href{../../../../continuum/strategy/capability/feature/891ed75bdd3c490296e5f0b7cacf6c87.md}{Canonical record: 891ed75bdd3c490296e5f0b7cacf6c87}; full qualifications and provenance.}\par
\subparagraph*{II.1.3.4.2.7.3 [Feature] Mobile Clinical Equipment Tracking}
\phantomsection
\label{navigation:II.1.3.4.2.7.3}
\addcontentsline{toc}{subparagraph}{II.1.3.4.2.7.3 [Feature] Mobile Clinical Equipment Tracking}

\label{element:a013a8f41ce245f48aed8d9fe33615c5}
\index{[Feature] Mobile Clinical Equipment Tracking (II.1.3.4.2.7.3)}
{\small Capability / Feature · Canonical · current}\par

Track physical location and allocation status of shared clinical equipment (infusion pumps, ventilators, mobile telemetry).



Equipment tracking provides the dynamic location and allocation context of shared devices. It distinguishes an identifiable asset from one available for the present operational need, while leaving maintenance, ownership and clinical authorisation with their respective responsibilities.


{\footnotesize \href{../../../../continuum/strategy/capability/feature/a013a8f41ce245f48aed8d9fe33615c5.md}{Canonical record: a013a8f41ce245f48aed8d9fe33615c5}; full qualifications and provenance.}\par
\subsubsection*{II.1.3.4.3 [CT2] Health Materiel Administration}
\phantomsection
\label{navigation:II.1.3.4.3}
\addcontentsline{toc}{subsubsection}{II.1.3.4.3 [CT2] Health Materiel Administration}

{\small Historical predecessor: R.2.102 \hyperref[element:4952f86078c54a8d898ac992e6f4f577]{[CT2] Health Materiel Administration}; Canonical ID: \texttt{\seqsplit{4952f86078c54a8d898ac992e6f4f577}}. Superseded historical identity — direct correspondence · superseded.}\par
\label{element:e7081070a9f8444f81c1d8fce37ef664}
\index{[CT2] Health Materiel Administration (II.1.3.4.3)}
{\small Capability · Canonical · current}\par

Health Product Administration retains the reference context for medications, implants, blood products and consumables. Product definitions, formularies and batch or serial information help explain what is being referenced without assigning a supply-chain implementation to Harmonia. No Feature decomposition is established for this responsibility in the present baseline.



\textbf{Overview}



Health Product Administration retains the reference context for medications, implants, blood products and consumables. Product definitions, formularies and batch or serial information help explain what is being referenced without assigning a supply-chain implementation to Harmonia. No Feature decomposition is established for this responsibility in the present baseline.



\textbf{Intent}



Health Product Administration retains the reference context for medications, implants, blood products and consumables. Product definitions, formularies and batch or serial information help explain what is being referenced without assigning a supply-chain implementation to Harmonia. No Feature decomposition is established for this responsibility in the present baseline.



\textbf{Boundary}



Information administration does not confer authority over the represented person, organisation, service, place, device or clinical activity.



\medskip\hrule\medskip

{\small Recorded unresolved matters: source-qualified matters; exact scope remains in the canonical record.}\par
{\footnotesize \href{../../../../continuum/strategy/capability/e7081070a9f8444f81c1d8fce37ef664.md}{Canonical record: e7081070a9f8444f81c1d8fce37ef664}; full qualifications and provenance.}\par
\paragraph*{II.1.3.4.3.1 [CT3] Health Materiel Catalogue \& Formulary Administration}
\phantomsection
\label{navigation:II.1.3.4.3.1}
\addcontentsline{toc}{paragraph}{II.1.3.4.3.1 [CT3] Health Materiel Catalogue \& Formulary Administration}

{\small Historical predecessor: R.2.99 \hyperref[element:5c691959e2cc4a289fc25b9f3497eed9]{[CT3] Health Materiel Catalogue \& Formulary Administration}; Canonical ID: \texttt{\seqsplit{5c691959e2cc4a289fc25b9f3497eed9}}. Superseded historical identity — direct correspondence · superseded.}\par
\label{element:28b327760af94b8b986357469524d793}
\index{[CT3] Health Materiel Catalogue \& Formulary Administration (II.1.3.4.3.1)}
{\small Capability · Proposed · current}\par

Administer health-materiel catalogue and externally approved formulary information, preserving approval and version standing.



\textbf{Overview}



Administer health-materiel catalogue and externally approved formulary information, preserving approval and version standing.



\textbf{Intent}



Administer health-materiel catalogue and externally approved formulary information, preserving approval and version standing.



\textbf{Boundary}



Represents externally approved formularies; it does not acquire formulary approval or dispensing authority.


{\footnotesize \href{../../../../continuum/strategy/capability/28b327760af94b8b986357469524d793.md}{Canonical record: 28b327760af94b8b986357469524d793}; full qualifications and provenance.}\par
\paragraph*{II.1.3.4.3.2 [CT3] Health Materiel Profile Administration}
\phantomsection
\label{navigation:II.1.3.4.3.2}
\addcontentsline{toc}{paragraph}{II.1.3.4.3.2 [CT3] Health Materiel Profile Administration}

{\small Historical predecessor: R.2.100 \hyperref[element:9216f9cc905e4526b1d98cd4ab382cbf]{[CT3] Health Materiel Profile Administration}; Canonical ID: \texttt{\seqsplit{9216f9cc905e4526b1d98cd4ab382cbf}}. Superseded historical identity — direct correspondence · superseded.}\par
\label{element:9626edae517d40d2bfae619860f18e20}
\index{[CT3] Health Materiel Profile Administration (II.1.3.4.3.2)}
{\small Capability · Proposed · current}\par

Administer identifying and descriptive information about health materiel referenced in healthcare information and activity.



\textbf{Overview}



Administer identifying and descriptive information about health materiel referenced in healthcare information and activity.



\textbf{Intent}



Administer identifying and descriptive information about health materiel referenced in healthcare information and activity.



\textbf{Boundary}



Reference representation is distinct from procurement, manufacturing and prescribing.


{\footnotesize \href{../../../../continuum/strategy/capability/9626edae517d40d2bfae619860f18e20.md}{Canonical record: 9626edae517d40d2bfae619860f18e20}; full qualifications and provenance.}\par
\paragraph*{II.1.3.4.3.3 [CT3] Health Materiel Traceability Administration}
\phantomsection
\label{navigation:II.1.3.4.3.3}
\addcontentsline{toc}{paragraph}{II.1.3.4.3.3 [CT3] Health Materiel Traceability Administration}

{\small Historical predecessor: R.2.101 \hyperref[element:f1b888e8d50c4019ac1afe162d6223f3]{[CT3] Health Materiel Traceability Administration}; Canonical ID: \texttt{\seqsplit{f1b888e8d50c4019ac1afe162d6223f3}}. Superseded historical identity — direct correspondence · superseded.}\par
\label{element:2d41ed0f00a94a4b90ffe87d6fd4e9fe}
\index{[CT3] Health Materiel Traceability Administration (II.1.3.4.3.3)}
{\small Capability · Proposed · current}\par

Administer batch/lot, serial, supplier, manufacturer, receipt/delivery and relevant movement/provenance information concerning health materiel.



\textbf{Overview}



Administer batch/lot, serial, supplier, manufacturer, receipt/delivery and relevant movement/provenance information concerning health materiel.



\textbf{Intent}



Administer batch/lot, serial, supplier, manufacturer, receipt/delivery and relevant movement/provenance information concerning health materiel.



\textbf{Boundary}



Materiel traceability does not make Harmonia a procurement, stock-control or supply-chain manager; specimen movement remains Clinical Logistics Coordination.


{\footnotesize \href{../../../../continuum/strategy/capability/2d41ed0f00a94a4b90ffe87d6fd4e9fe.md}{Canonical record: 2d41ed0f00a94a4b90ffe87d6fd4e9fe}; full qualifications and provenance.}\par
\subsubsection*{II.1.3.4.4 [CT2] Health Service Administration}
\phantomsection
\label{navigation:II.1.3.4.4}
\addcontentsline{toc}{subsubsection}{II.1.3.4.4 [CT2] Health Service Administration}

{\small Historical predecessor: R.2.104 \hyperref[element:f1ab1258cf2944efb50105e7885f8027]{[CT2] Health Service Administration}; Canonical ID: \texttt{\seqsplit{f1ab1258cf2944efb50105e7885f8027}}. Superseded historical identity — direct correspondence · superseded.}\par
\label{element:ba68a52913a0499d8377553bddfa74a4}
\index{[CT2] Health Service Administration (II.1.3.4.4)}
{\small Capability · Canonical · current}\par

Health Service Administration represents identifiable services and the conditions under which organisations, locations and practitioner roles provide them. Catalogue publishing, service-location-provider association, operating schedules and eligibility definitions allow a service to be understood as more than a location or specialty label. The service context remains material wherever activity meaning or accountability depends on it.



\textbf{Overview}



Health Service Administration represents identifiable services and the conditions under which organisations, locations and practitioner roles provide them. Catalogue publishing, service-location-provider association, operating schedules and eligibility definitions allow a service to be understood as more than a location or specialty label. The service context remains material wherever activity meaning or accountability depends on it.



\textbf{Intent}



Health Service Administration represents identifiable services and the conditions under which organisations, locations and practitioner roles provide them. Catalogue publishing, service-location-provider association, operating schedules and eligibility definitions allow a service to be understood as more than a location or specialty label. The service context remains material wherever activity meaning or accountability depends on it.



\textbf{Boundary}



Information administration does not confer authority over the represented person, organisation, service, place, device or clinical activity.



\medskip\hrule\medskip

{\small Recorded unresolved matters: source-qualified matters; exact scope remains in the canonical record.}\par
{\footnotesize \href{../../../../continuum/strategy/capability/ba68a52913a0499d8377553bddfa74a4.md}{Canonical record: ba68a52913a0499d8377553bddfa74a4}; full qualifications and provenance.}\par
\paragraph*{II.1.3.4.4.1 [CT3] Healthcare Service Profile Administration}
\phantomsection
\label{navigation:II.1.3.4.4.1}
\addcontentsline{toc}{paragraph}{II.1.3.4.4.1 [CT3] Healthcare Service Profile Administration}

{\small Historical predecessor: R.2.103 \hyperref[element:7780ca2fed2346439d5f559d00cc4ee4]{[CT3] Healthcare Service Profile Administration}; Canonical ID: \texttt{\seqsplit{7780ca2fed2346439d5f559d00cc4ee4}}. Superseded historical identity — direct correspondence · superseded.}\par
\label{element:e75d5eb30936476ea7a9cbadfa43d5ef}
\index{[CT3] Healthcare Service Profile Administration (II.1.3.4.4.1)}
{\small Capability · Proposed · current}\par

Administer service definitions, provision relationships, defined availability and eligibility information describing offered healthcare services.



\textbf{Overview}



Administer service definitions, provision relationships, defined availability and eligibility information describing offered healthcare services.



\textbf{Intent}



Administer service definitions, provision relationships, defined availability and eligibility information describing offered healthcare services.



\textbf{Boundary}



Defined availability is not dynamic operational capacity, collective readiness or authority to accept clinical care.


{\footnotesize \href{../../../../continuum/strategy/capability/e75d5eb30936476ea7a9cbadfa43d5ef.md}{Canonical record: e75d5eb30936476ea7a9cbadfa43d5ef}; full qualifications and provenance.}\par
\subparagraph*{II.1.3.4.4.1.1 [Feature] Eligibility \& Catchment Specification}
\phantomsection
\label{navigation:II.1.3.4.4.1.1}
\addcontentsline{toc}{subparagraph}{II.1.3.4.4.1.1 [Feature] Eligibility \& Catchment Specification}

\label{element:c65feabade3d47c39e7d7af9bc16ebdc}
\index{[Feature] Eligibility \& Catchment Specification (II.1.3.4.4.1.1)}
{\small Capability / Feature · Canonical · current}\par

Define clinical referral thresholds and geographic catchment boundaries.



Eligibility specification captures the clinical thresholds and geographical bounds applicable to service access. It provides the criteria context needed to consider a request without making publication of those criteria an individual eligibility decision or a completed allocation.


{\footnotesize \href{../../../../continuum/strategy/capability/feature/c65feabade3d47c39e7d7af9bc16ebdc.md}{Canonical record: c65feabade3d47c39e7d7af9bc16ebdc}; full qualifications and provenance.}\par
\subparagraph*{II.1.3.4.4.1.2 [Feature] Operating Schedule Specification}
\phantomsection
\label{navigation:II.1.3.4.4.1.2}
\addcontentsline{toc}{subparagraph}{II.1.3.4.4.1.2 [Feature] Operating Schedule Specification}

\label{element:4753953665d343c5a97bc4965be3f615}
\index{[Feature] Operating Schedule Specification (II.1.3.4.4.1.2)}
{\small Capability / Feature · Canonical · current}\par

Define temporal availability, clinic operating hours, and on-call service windows.



Operating schedules make temporal availability and exceptions explicit for a service. They distinguish an existing offering from one available at the relevant time, while leaving actual resource readiness and operational capacity as separate concerns.


{\footnotesize \href{../../../../continuum/strategy/capability/feature/4753953665d343c5a97bc4965be3f615.md}{Canonical record: 4753953665d343c5a97bc4965be3f615}; full qualifications and provenance.}\par
\subparagraph*{II.1.3.4.4.1.3 [Feature] Service Catalogue Publishing}
\phantomsection
\label{navigation:II.1.3.4.4.1.3}
\addcontentsline{toc}{subparagraph}{II.1.3.4.4.1.3 [Feature] Service Catalogue Publishing}

\label{element:d397e670ae524f399e91176c64390fa7}
\index{[Feature] Service Catalogue Publishing (II.1.3.4.4.1.3)}
{\small Capability / Feature · Canonical · current}\par

Publish structured healthcare service definitions, clinical specialties, and intake criteria to regional directories.



Service catalogue publishing makes identifiable service definitions and intake context available to regional participants. A published service can be discovered as an offering rather than inferred from an organisation's name; publication alone does not establish present availability or collective readiness.


{\footnotesize \href{../../../../continuum/strategy/capability/feature/d397e670ae524f399e91176c64390fa7.md}{Canonical record: d397e670ae524f399e91176c64390fa7}; full qualifications and provenance.}\par
\subparagraph*{II.1.3.4.4.1.4 [Feature] Service-Location-Provider Binding}
\phantomsection
\label{navigation:II.1.3.4.4.1.4}
\addcontentsline{toc}{subparagraph}{II.1.3.4.4.1.4 [Feature] Service-Location-Provider Binding}

\label{element:37188bb916244f41b1b7d988ad52a2eb}
\index{[Feature] Service-Location-Provider Binding (II.1.3.4.4.1.4)}
{\small Capability / Feature · Canonical · current}\par

Correlate which organization provides a named health service at which physical location with which practitioner roles.



Service-location-provider binding explains which participant provides an identifiable service in which place and acting capacities. The mapping associations belong to Health Service Administration, while the organisations, locations and practitioner entities retain their own information responsibility.


{\footnotesize \href{../../../../continuum/strategy/capability/feature/37188bb916244f41b1b7d988ad52a2eb.md}{Canonical record: 37188bb916244f41b1b7d988ad52a2eb}; full qualifications and provenance.}\par
\paragraph*{II.1.3.4.4.2 [CT3] Service Availability Definition [superseded]}
\phantomsection
\label{navigation:II.1.3.4.4.2}
\addcontentsline{toc}{paragraph}{II.1.3.4.4.2 [CT3] Service Availability Definition [superseded]}

\label{element:9e5329a42f964199b4b0a87ab1cacc2e}
\index{[CT3] Service Availability Definition (II.1.3.4.4.2)}
{\small Capability · Proposed · superseded}\par
\href{../../../../continuum/strategy/capability/9e5329a42f964199b4b0a87ab1cacc2e.md}{Historical canonical definition and provenance}.\par
\paragraph*{II.1.3.4.4.3 [CT3] Service Eligibility Definition [superseded]}
\phantomsection
\label{navigation:II.1.3.4.4.3}
\addcontentsline{toc}{paragraph}{II.1.3.4.4.3 [CT3] Service Eligibility Definition [superseded]}

\label{element:8a3f29e5948141e2a3d0c72539759daf}
\index{[CT3] Service Eligibility Definition (II.1.3.4.4.3)}
{\small Capability · Proposed · superseded}\par
\href{../../../../continuum/strategy/capability/8a3f29e5948141e2a3d0c72539759daf.md}{Historical canonical definition and provenance}.\par
\paragraph*{II.1.3.4.4.4 [CT3] Service Provision Relationship [superseded]}
\phantomsection
\label{navigation:II.1.3.4.4.4}
\addcontentsline{toc}{paragraph}{II.1.3.4.4.4 [CT3] Service Provision Relationship [superseded]}

\label{element:718ebe1fa8d44ad9a6632a5932277a8e}
\index{[CT3] Service Provision Relationship (II.1.3.4.4.4)}
{\small Capability · Proposed · superseded}\par
\href{../../../../continuum/strategy/capability/718ebe1fa8d44ad9a6632a5932277a8e.md}{Historical canonical definition and provenance}.\par
\subsubsection*{II.1.3.4.5 [CT2] Location Administration}
\phantomsection
\label{navigation:II.1.3.4.5}
\addcontentsline{toc}{subsubsection}{II.1.3.4.5 [CT2] Location Administration}

{\small Historical predecessor: R.2.107 \hyperref[element:d412aff3a5c64e6aaba910d4c2c9356d]{[CT2] Location Administration}; Canonical ID: \texttt{\seqsplit{d412aff3a5c64e6aaba910d4c2c9356d}}. Superseded historical identity — direct correspondence · superseded.}\par
\label{element:c410005667444f99a221483bf577fad2}
\index{[CT2] Location Administration (II.1.3.4.5)}
{\small Capability · Canonical · current}\par

Location Administration represents the places and physical relationships relevant to healthcare activity, from facilities to particular care places. Address and geospatial context answer where a place is; hierarchy explains its surroundings; care-place definition records the local characteristics relevant to care. These are stable descriptive responsibilities, distinct from the changing operational state of a bed or ward.



\textbf{Overview}



Location Administration represents the places and physical relationships relevant to healthcare activity, from facilities to particular care places. Address and geospatial context answer where a place is; hierarchy explains its surroundings; care-place definition records the local characteristics relevant to care. These are stable descriptive responsibilities, distinct from the changing operational state of a bed or ward.



\textbf{Intent}



Location Administration represents the places and physical relationships relevant to healthcare activity, from facilities to particular care places. Address and geospatial context answer where a place is; hierarchy explains its surroundings; care-place definition records the local characteristics relevant to care. These are stable descriptive responsibilities, distinct from the changing operational state of a bed or ward.



\textbf{Boundary}



Information administration does not confer authority over the represented person, organisation, service, place, device or clinical activity.



\medskip\hrule\medskip

{\small Recorded unresolved matters: source-qualified matters; exact scope remains in the canonical record.}\par
{\footnotesize \href{../../../../continuum/strategy/capability/c410005667444f99a221483bf577fad2.md}{Canonical record: c410005667444f99a221483bf577fad2}; full qualifications and provenance.}\par
\paragraph*{II.1.3.4.5.1 [CT3] Care-Place Definition Administration}
\phantomsection
\label{navigation:II.1.3.4.5.1}
\addcontentsline{toc}{paragraph}{II.1.3.4.5.1 [CT3] Care-Place Definition Administration}

{\small Historical predecessor: R.2.105 \hyperref[element:a3464bd9c0dc4a01b50803acc7705b4e]{[CT3] Care-Place Definition Administration}; Canonical ID: \texttt{\seqsplit{a3464bd9c0dc4a01b50803acc7705b4e}}. Superseded historical identity — direct correspondence · superseded.}\par
\label{element:629b493c7d9844baa25c029a87ed04f6}
\index{[CT3] Care-Place Definition Administration (II.1.3.4.5.1)}
{\small Capability · Proposed · current}\par

Administer the stable description, characteristics and physical location of a bed, bay, chair or other care-place.



\textbf{Overview}



Administer the stable description, characteristics and physical location of a bed, bay, chair or other care-place.



\textbf{Intent}



Administer the stable description, characteristics and physical location of a bed, bay, chair or other care-place.



\textbf{Boundary}



Describes a place and its capabilities; operational bed state and readiness remain separately managed.


{\footnotesize \href{../../../../continuum/strategy/capability/629b493c7d9844baa25c029a87ed04f6.md}{Canonical record: 629b493c7d9844baa25c029a87ed04f6}; full qualifications and provenance.}\par
\subparagraph*{II.1.3.4.5.1.1 [Feature] Atomic Care-Place Modeling}
\phantomsection
\label{navigation:II.1.3.4.5.1.1}
\addcontentsline{toc}{subparagraph}{II.1.3.4.5.1.1 [Feature] Atomic Care-Place Modeling}

\label{element:ddf34f0844114798a4bd84ee72bbe497}
\index{[Feature] Atomic Care-Place Modeling (II.1.3.4.5.1.1)}
{\small Capability / Feature · Canonical · current}\par

Define specific beds, care-bays, and examination rooms with their physical capabilities (e.g., telemetry-capable, negative-pressure).



Care-place modelling identifies the particular physical place and characteristics relevant to care, such as isolation or telemetry capability. It supplies the descriptive frame against which changing occupancy and readiness can be considered without making those states part of the place's identity.


{\footnotesize \href{../../../../continuum/strategy/capability/feature/ddf34f0844114798a4bd84ee72bbe497.md}{Canonical record: ddf34f0844114798a4bd84ee72bbe497}; full qualifications and provenance.}\par
\paragraph*{II.1.3.4.5.2 [CT3] Healthcare Location Profile Administration}
\phantomsection
\label{navigation:II.1.3.4.5.2}
\addcontentsline{toc}{paragraph}{II.1.3.4.5.2 [CT3] Healthcare Location Profile Administration}

{\small Historical predecessor: R.2.106 \hyperref[element:8a04d1f16d0849e786d2aaec33e0ba87]{[CT3] Healthcare Location Profile Administration}; Canonical ID: \texttt{\seqsplit{8a04d1f16d0849e786d2aaec33e0ba87}}. Superseded historical identity — direct correspondence · superseded.}\par
\label{element:e579e54a08bb4f899a504aa998ba246b}
\index{[CT3] Healthcare Location Profile Administration (II.1.3.4.5.2)}
{\small Capability · Proposed · current}\par

Administer the identity, address, geospatial description and physical containment relationships of healthcare locations.



\textbf{Overview}



Administer the identity, address, geospatial description and physical containment relationships of healthcare locations.



\textbf{Intent}



Administer the identity, address, geospatial description and physical containment relationships of healthcare locations.



\textbf{Boundary}



Location description does not establish occupancy, cleanliness, service capacity or service readiness.


{\footnotesize \href{../../../../continuum/strategy/capability/e579e54a08bb4f899a504aa998ba246b.md}{Canonical record: e579e54a08bb4f899a504aa998ba246b}; full qualifications and provenance.}\par
\subparagraph*{II.1.3.4.5.2.1 [Feature] Location Geospatial \& Address Resolution}
\phantomsection
\label{navigation:II.1.3.4.5.2.1}
\addcontentsline{toc}{subparagraph}{II.1.3.4.5.2.1 [Feature] Location Geospatial \& Address Resolution}

\label{element:a27d1f97313f4ba5907c7cf7e2ec8858}
\index{[Feature] Location Geospatial \& Address Resolution (II.1.3.4.5.2.1)}
{\small Capability / Feature · Canonical · current}\par

Validate physical addresses, geospatial coordinates, and location operational status.



Address and geospatial resolution gives a recognised healthcare location an intelligible physical reference. It supports place interpretation without treating a valid address or recorded operating status as proof that a particular service or care place is currently ready.


{\footnotesize \href{../../../../continuum/strategy/capability/feature/a27d1f97313f4ba5907c7cf7e2ec8858.md}{Canonical record: a27d1f97313f4ba5907c7cf7e2ec8858}; full qualifications and provenance.}\par
\subparagraph*{II.1.3.4.5.2.2 [Feature] Physical Hierarchy Modeling}
\phantomsection
\label{navigation:II.1.3.4.5.2.2}
\addcontentsline{toc}{subparagraph}{II.1.3.4.5.2.2 [Feature] Physical Hierarchy Modeling}

\label{element:6166149dbd7b431985063d92e78ecdb7}
\index{[Feature] Physical Hierarchy Modeling (II.1.3.4.5.2.2)}
{\small Capability / Feature · Canonical · current}\par

Represent campus, building, wing, floor, and ward architectural relationships.



Physical hierarchy modelling relates a care environment to its campus, building, floor or ward context. The relationships describe where places sit physically; they do not automatically reproduce organisational hierarchy, service responsibility or operational capacity.


{\footnotesize \href{../../../../continuum/strategy/capability/feature/6166149dbd7b431985063d92e78ecdb7.md}{Canonical record: 6166149dbd7b431985063d92e78ecdb7}; full qualifications and provenance.}\par
\paragraph*{II.1.3.4.5.3 [CT3] Location Structure [superseded]}
\phantomsection
\label{navigation:II.1.3.4.5.3}
\addcontentsline{toc}{paragraph}{II.1.3.4.5.3 [CT3] Location Structure [superseded]}

\label{element:2436c87195334daf85ed4e2de4ec4bd5}
\index{[CT3] Location Structure (II.1.3.4.5.3)}
{\small Capability · Proposed · superseded}\par
\href{../../../../continuum/strategy/capability/2436c87195334daf85ed4e2de4ec4bd5.md}{Historical canonical definition and provenance}.\par
\subsubsection*{II.1.3.4.6 [CT2] Organisation Administration}
\phantomsection
\label{navigation:II.1.3.4.6}
\addcontentsline{toc}{subsubsection}{II.1.3.4.6 [CT2] Organisation Administration}

{\small Historical predecessor: R.2.109 \hyperref[element:1b0eacfbb27d4631a24f9279dab6681d]{[CT2] Organisation Administration}; Canonical ID: \texttt{\seqsplit{1b0eacfbb27d4631a24f9279dab6681d}}. Superseded historical identity — direct correspondence · superseded.}\par
\label{element:14fbdc5ce3dd4c2092c5397dab3ed743}
\index{[CT2] Organisation Administration (II.1.3.4.6)}
{\small Capability · Canonical · current}\par

Organisation Administration maintains the identity, profiles, structural relationships and contact context of participating healthcare organisations. It enables a legal or operational entity to be distinguished from a department, network relationship or communication channel. This supports consistent interpretation of providers and destinations across information flows rather than allowing each interface to create an incompatible organisational account.



\textbf{Overview}



Organisation Administration maintains the identity, profiles, structural relationships and contact context of participating healthcare organisations. It enables a legal or operational entity to be distinguished from a department, network relationship or communication channel. This supports consistent interpretation of providers and destinations across information flows rather than allowing each interface to create an incompatible organisational account.



\textbf{Intent}



Organisation Administration maintains the identity, profiles, structural relationships and contact context of participating healthcare organisations. It enables a legal or operational entity to be distinguished from a department, network relationship or communication channel. This supports consistent interpretation of providers and destinations across information flows rather than allowing each interface to create an incompatible organisational account.



\textbf{Boundary}



Information administration does not confer authority over the represented person, organisation, service, place, device or clinical activity.



\medskip\hrule\medskip

{\small Recorded unresolved matters: source-qualified matters; exact scope remains in the canonical record.}\par
{\footnotesize \href{../../../../continuum/strategy/capability/14fbdc5ce3dd4c2092c5397dab3ed743.md}{Canonical record: 14fbdc5ce3dd4c2092c5397dab3ed743}; full qualifications and provenance.}\par
\paragraph*{II.1.3.4.6.1 [CT3] Healthcare Organisation Profile Administration}
\phantomsection
\label{navigation:II.1.3.4.6.1}
\addcontentsline{toc}{paragraph}{II.1.3.4.6.1 [CT3] Healthcare Organisation Profile Administration}

{\small Historical predecessor: R.2.108 \hyperref[element:e485424b799c42aca7399d341ca02d0b]{[CT3] Healthcare Organisation Profile Administration}; Canonical ID: \texttt{\seqsplit{e485424b799c42aca7399d341ca02d0b}}. Superseded historical identity — direct correspondence · superseded.}\par
\label{element:f2de9e9394b849f2943355c7a2882a15}
\index{[CT3] Healthcare Organisation Profile Administration (II.1.3.4.6.1)}
{\small Capability · Proposed · current}\par

Administer identifying, descriptive, structural, location, contact, registered/legal and recognised endpoint information needed to represent and interact with a healthcare organisation.



\textbf{Overview}



Administer identifying, descriptive, structural, location, contact, registered/legal and recognised endpoint information needed to represent and interact with a healthcare organisation.



\textbf{Intent}



Administer identifying, descriptive, structural, location, contact, registered/legal and recognised endpoint information needed to represent and interact with a healthcare organisation.



\textbf{Boundary}



Maintains verified legal and operational identity information without corporate management authority.


{\footnotesize \href{../../../../continuum/strategy/capability/f2de9e9394b849f2943355c7a2882a15.md}{Canonical record: f2de9e9394b849f2943355c7a2882a15}; full qualifications and provenance.}\par
\subparagraph*{II.1.3.4.6.1.1 [Feature] Directory Contact Management}
\phantomsection
\label{navigation:II.1.3.4.6.1.1}
\addcontentsline{toc}{subparagraph}{II.1.3.4.6.1.1 [Feature] Directory Contact Management}

\label{element:1cdac17b7547448daa20b74ed45ec366}
\index{[Feature] Directory Contact Management (II.1.3.4.6.1.1)}
{\small Capability / Feature · Canonical · current}\par

Manage verified phone, email, and administrative contact channels for healthcare facilities.



Contact management provides verified administrative and clinical communication channels for organisations. It supports reaching the appropriate participant while distinguishing a directory contact from the role-specific destination and authority relevant to a particular clinical exchange.


{\footnotesize \href{../../../../continuum/strategy/capability/feature/1cdac17b7547448daa20b74ed45ec366.md}{Canonical record: 1cdac17b7547448daa20b74ed45ec366}; full qualifications and provenance.}\par
\subparagraph*{II.1.3.4.6.1.2 [Feature] Hierarchical Structure Navigation}
\phantomsection
\label{navigation:II.1.3.4.6.1.2}
\addcontentsline{toc}{subparagraph}{II.1.3.4.6.1.2 [Feature] Hierarchical Structure Navigation}

\label{element:7ded6f7865294fb2a696a58160d577f7}
\index{[Feature] Hierarchical Structure Navigation (II.1.3.4.6.1.2)}
{\small Capability / Feature · Canonical · current}\par

Model, traverse, and resolve nested departmental, divisional, and regional healthcare networks.



Hierarchy navigation makes departmental, divisional and network relationships usable without flattening them into unrelated organisations. Traversal reflects established structure; it does not create corporate authority or infer responsibility from the way a directory happens to display a tree.


{\footnotesize \href{../../../../continuum/strategy/capability/feature/7ded6f7865294fb2a696a58160d577f7.md}{Canonical record: 7ded6f7865294fb2a696a58160d577f7}; full qualifications and provenance.}\par
\subparagraph*{II.1.3.4.6.1.3 [Feature] National Organisation Verification}
\phantomsection
\label{navigation:II.1.3.4.6.1.3}
\addcontentsline{toc}{subparagraph}{II.1.3.4.6.1.3 [Feature] National Organisation Verification}

\label{element:0f17fe5439fe45dc8f4edbff3cca0776}
\index{[Feature] National Organisation Verification (II.1.3.4.6.1.3)}
{\small Capability / Feature · Canonical · current}\par

Validate national organization identifiers (e.g., HPI-O) and legal entity classifications.



National organisation verification grounds an organisation identifier and entity classification in their recognised context. This supports reliable interpretation of the healthcare participant without assuming that its identity establishes every service, location or authority associated with it.


{\footnotesize \href{../../../../continuum/strategy/capability/feature/0f17fe5439fe45dc8f4edbff3cca0776.md}{Canonical record: 0f17fe5439fe45dc8f4edbff3cca0776}; full qualifications and provenance.}\par
\subparagraph*{II.1.3.4.6.1.4 [Feature] Organisation Profile Management}
\phantomsection
\label{navigation:II.1.3.4.6.1.4}
\addcontentsline{toc}{subparagraph}{II.1.3.4.6.1.4 [Feature] Organisation Profile Management}

\label{element:6a4c44cfb1414f3bae5ba0fce1e58970}
\index{[Feature] Organisation Profile Management (II.1.3.4.6.1.4)}
{\small Capability / Feature · Canonical · current}\par

Govern directory profiles, operational status, and published service descriptions.



Organisation profiles provide a governed account of the participant's status and published service descriptions. They support consistent directory use across systems while leaving the organisation's real-world governance and the definition of individual offerings distinguishable.


{\footnotesize \href{../../../../continuum/strategy/capability/feature/6a4c44cfb1414f3bae5ba0fce1e58970.md}{Canonical record: 6a4c44cfb1414f3bae5ba0fce1e58970}; full qualifications and provenance.}\par
\paragraph*{II.1.3.4.6.2 [CT3] Organisation Contact [superseded]}
\phantomsection
\label{navigation:II.1.3.4.6.2}
\addcontentsline{toc}{paragraph}{II.1.3.4.6.2 [CT3] Organisation Contact [superseded]}

\label{element:b7f5108880eb429bb407c87a37b0e23f}
\index{[CT3] Organisation Contact (II.1.3.4.6.2)}
{\small Capability · Proposed · superseded}\par
\href{../../../../continuum/strategy/capability/b7f5108880eb429bb407c87a37b0e23f.md}{Historical canonical definition and provenance}.\par
\paragraph*{II.1.3.4.6.3 [CT3] Organisational Structure [superseded]}
\phantomsection
\label{navigation:II.1.3.4.6.3}
\addcontentsline{toc}{paragraph}{II.1.3.4.6.3 [CT3] Organisational Structure [superseded]}

\label{element:19d211cf601c49e5a65665330d0c5c3d}
\index{[CT3] Organisational Structure (II.1.3.4.6.3)}
{\small Capability · Proposed · superseded}\par
\href{../../../../continuum/strategy/capability/19d211cf601c49e5a65665330d0c5c3d.md}{Historical canonical definition and provenance}.\par
\subsubsection*{II.1.3.4.7 [CT2] Provider Administration}
\phantomsection
\label{navigation:II.1.3.4.7}
\addcontentsline{toc}{subsubsection}{II.1.3.4.7 [CT2] Provider Administration}

{\small Historical predecessor: R.2.113 \hyperref[element:a352fd50419d451887c85fcb98566cd5]{[CT2] Provider Administration}; Canonical ID: \texttt{\seqsplit{a352fd50419d451887c85fcb98566cd5}}. Superseded historical identity — direct correspondence · superseded.}\par
\label{element:d6312f85968a4209bafff5ed009297a6}
\index{[CT2] Provider Administration (II.1.3.4.7)}
{\small Capability · Canonical · current}\par

Provider Administration maintains the practitioner, role, affiliation and engagement context needed to interpret healthcare participation. National identifiers, professional profiles, role-specific endpoints and privilege records support different questions: who the practitioner is, in what capacity they act and how they can be reached. Directory representation makes those distinctions usable without replacing the business governance that establishes professional standing.



\textbf{Overview}



Provider Administration maintains the practitioner, role, affiliation and engagement context needed to interpret healthcare participation. National identifiers, professional profiles, role-specific endpoints and privilege records support different questions: who the practitioner is, in what capacity they act and how they can be reached. Directory representation makes those distinctions usable without replacing the business governance that establishes professional standing.



\textbf{Intent}



Provider Administration maintains the practitioner, role, affiliation and engagement context needed to interpret healthcare participation. National identifiers, professional profiles, role-specific endpoints and privilege records support different questions: who the practitioner is, in what capacity they act and how they can be reached. Directory representation makes those distinctions usable without replacing the business governance that establishes professional standing.



\textbf{Boundary}



Information administration does not confer authority over the represented person, organisation, service, place, device or clinical activity.



\medskip\hrule\medskip

{\small Recorded unresolved matters: source-qualified matters; exact scope remains in the canonical record.}\par
{\footnotesize \href{../../../../continuum/strategy/capability/d6312f85968a4209bafff5ed009297a6.md}{Canonical record: d6312f85968a4209bafff5ed009297a6}; full qualifications and provenance.}\par
\paragraph*{II.1.3.4.7.1 [CT3] Practitioner Profile Administration}
\phantomsection
\label{navigation:II.1.3.4.7.1}
\addcontentsline{toc}{paragraph}{II.1.3.4.7.1 [CT3] Practitioner Profile Administration}

{\small Historical predecessor: R.2.110 \hyperref[element:2415ddcbc73c47e1a63be310bd4342ee]{[CT3] Practitioner Profile Administration}; Canonical ID: \texttt{\seqsplit{2415ddcbc73c47e1a63be310bd4342ee}}. Superseded historical identity — direct correspondence · superseded.}\par
\label{element:40759a835f7b456e9257c38b67b0a380}
\index{[CT3] Practitioner Profile Administration (II.1.3.4.7.1)}
{\small Capability · Proposed · current}\par

Administer practitioner identity, professional profile and supplied registration information with its issuing-authority attribution.



\textbf{Overview}



Administer practitioner identity, professional profile and supplied registration information with its issuing-authority attribution.



\textbf{Intent}



Administer practitioner identity, professional profile and supplied registration information with its issuing-authority attribution.



\textbf{Boundary}



Professional registration and profile representation do not grant clinical or operational privilege.


{\footnotesize \href{../../../../continuum/strategy/capability/40759a835f7b456e9257c38b67b0a380.md}{Canonical record: 40759a835f7b456e9257c38b67b0a380}; full qualifications and provenance.}\par
\subparagraph*{II.1.3.4.7.1.1 [Feature] National Practitioner Verification}
\phantomsection
\label{navigation:II.1.3.4.7.1.1}
\addcontentsline{toc}{subparagraph}{II.1.3.4.7.1.1 [Feature] National Practitioner Verification}

\label{element:12c5a17dcc564619a6a9222dd2e15528}
\index{[Feature] National Practitioner Verification (II.1.3.4.7.1.1)}
{\small Capability / Feature · Canonical · current}\par

Ingest, validate, and verify national practitioner identifiers (e.g., HPI-I) and professional registration status.



National verification connects practitioner identifiers and registration information to their recognised source context. It helps distinguish an identifiable professional from an unverified directory assertion without implying that verification alone grants local clinical privileges.


{\footnotesize \href{../../../../continuum/strategy/capability/feature/12c5a17dcc564619a6a9222dd2e15528.md}{Canonical record: 12c5a17dcc564619a6a9222dd2e15528}; full qualifications and provenance.}\par
\subparagraph*{II.1.3.4.7.1.2 [Feature] Practitioner Profile Governance}
\phantomsection
\label{navigation:II.1.3.4.7.1.2}
\addcontentsline{toc}{subparagraph}{II.1.3.4.7.1.2 [Feature] Practitioner Profile Governance}

\label{element:d23be319c6114c28bcb9e24e1258a0fc}
\index{[Feature] Practitioner Profile Governance (II.1.3.4.7.1.2)}
{\small Capability / Feature · Canonical · current}\par

Maintain core professional demographic details, qualifications, and active practice status.



Practitioner profile governance maintains the professional context used across care information and communication. Qualifications and practice status remain attributable facts, distinct from the organisation-specific role or privilege under which the practitioner may act.


{\footnotesize \href{../../../../continuum/strategy/capability/feature/d23be319c6114c28bcb9e24e1258a0fc.md}{Canonical record: d23be319c6114c28bcb9e24e1258a0fc}; full qualifications and provenance.}\par
\subparagraph*{II.1.3.4.7.1.3 [Feature] Provider Directory Discovery \& Resolution}
\phantomsection
\label{navigation:II.1.3.4.7.1.3}
\addcontentsline{toc}{subparagraph}{II.1.3.4.7.1.3 [Feature] Provider Directory Discovery \& Resolution}

\label{element:cd3bd35f7f9e4852bbbce299bb795d8a}
\index{[Feature] Provider Directory Discovery \& Resolution (II.1.3.4.7.1.3)}
{\small Capability / Feature · Proposed · current}\par

Expose qualified practitioner discovery using identifiers, profile criteria and effective standing.



\textbf{Overview}



Routing and care participation need to find a described provider matching stated criteria, beyond merely ingesting a national identifier or updating a profile.



\textbf{Boundary}



Discovery supplies information; it does not determine privilege or availability.


{\footnotesize \href{../../../../continuum/strategy/capability/feature/cd3bd35f7f9e4852bbbce299bb795d8a.md}{Canonical record: cd3bd35f7f9e4852bbbce299bb795d8a}; full qualifications and provenance.}\par
\paragraph*{II.1.3.4.7.2 [CT3] Practitioner Role Administration}
\phantomsection
\label{navigation:II.1.3.4.7.2}
\addcontentsline{toc}{paragraph}{II.1.3.4.7.2 [CT3] Practitioner Role Administration}

{\small Historical predecessor: R.2.111 \hyperref[element:6b777f1201db474eb22dd89f2b1f68c6]{[CT3] Practitioner Role Administration}; Canonical ID: \texttt{\seqsplit{6b777f1201db474eb22dd89f2b1f68c6}}. Superseded historical identity — direct correspondence · superseded.}\par
\label{element:6651471e67ae46fb8d1c1241c8590be9}
\index{[CT3] Practitioner Role Administration (II.1.3.4.7.2)}
{\small Capability · Proposed · current}\par

Administer practitioner roles, specialties, organisational affiliations, operational capacities and recognised role endpoints.



\textbf{Overview}



Administer practitioner roles, specialties, organisational affiliations, operational capacities and recognised role endpoints.



\textbf{Intent}



Administer practitioner roles, specialties, organisational affiliations, operational capacities and recognised role endpoints.



\textbf{Boundary}



Role, specialty, capacity and endpoint information supply context; they do not constitute a privilege determination.


{\footnotesize \href{../../../../continuum/strategy/capability/6651471e67ae46fb8d1c1241c8590be9.md}{Canonical record: 6651471e67ae46fb8d1c1241c8590be9}; full qualifications and provenance.}\par
\subparagraph*{II.1.3.4.7.2.1 [Feature] Electronic Endpoint Binding}
\phantomsection
\label{navigation:II.1.3.4.7.2.1}
\addcontentsline{toc}{subparagraph}{II.1.3.4.7.2.1 [Feature] Electronic Endpoint Binding}

\label{element:5dc3a424e68d47aea64ec573dc4b1225}
\index{[Feature] Electronic Endpoint Binding (II.1.3.4.7.2.1)}
{\small Capability / Feature · Canonical · current}\par

Associate validated communication endpoints (secure messaging addresses, direct routing queues) with specific practitioner roles.



Endpoint binding makes a verified communication destination specific to the practitioner's acting role. It supports routing to the intended professional capacity rather than assuming that every address attached to the person is interchangeable across organisations and services.


{\footnotesize \href{../../../../continuum/strategy/capability/feature/5dc3a424e68d47aea64ec573dc4b1225.md}{Canonical record: 5dc3a424e68d47aea64ec573dc4b1225}; full qualifications and provenance.}\par
\subparagraph*{II.1.3.4.7.2.2 [Feature] Practitioner-Organisation Affiliation}
\phantomsection
\label{navigation:II.1.3.4.7.2.2}
\addcontentsline{toc}{subparagraph}{II.1.3.4.7.2.2 [Feature] Practitioner-Organisation Affiliation}

\label{element:9501e88ea4b7474a8aa63191687db51e}
\index{[Feature] Practitioner-Organisation Affiliation (II.1.3.4.7.2.2)}
{\small Capability / Feature · Canonical · current}\par

Maintain temporal affiliations between practitioners and healthcare organizations.



Affiliation records keep practitioner participation in an organisation temporally intelligible. They support interpretation of an actor's organisational context while preserving that affiliation, employment, qualification and permission to undertake a particular activity answer different questions.


{\footnotesize \href{../../../../continuum/strategy/capability/feature/9501e88ea4b7474a8aa63191687db51e.md}{Canonical record: 9501e88ea4b7474a8aa63191687db51e}; full qualifications and provenance.}\par
\subparagraph*{II.1.3.4.7.2.3 [Feature] Role \& Specialty Modeling}
\phantomsection
\label{navigation:II.1.3.4.7.2.3}
\addcontentsline{toc}{subparagraph}{II.1.3.4.7.2.3 [Feature] Role \& Specialty Modeling}

\label{element:ef859d4934514f4eb601c6abf9708760}
\index{[Feature] Role \& Specialty Modeling (II.1.3.4.7.2.3)}
{\small Capability / Feature · Canonical · current}\par

Define and associate clinical specialties, practice roles, and service capacities for a practitioner.



Role and specialty modelling describes the capacity in which a practitioner participates at an organisation or location. This keeps the person separate from each acting role and its service context; role resolution itself establishes neither Clinical Privilege nor Operational Privilege.


{\footnotesize \href{../../../../continuum/strategy/capability/feature/ef859d4934514f4eb601c6abf9708760.md}{Canonical record: ef859d4934514f4eb601c6abf9708760}; full qualifications and provenance.}\par
\paragraph*{II.1.3.4.7.3 [CT3] Provider Engagement [superseded]}
\phantomsection
\label{navigation:II.1.3.4.7.3}
\addcontentsline{toc}{paragraph}{II.1.3.4.7.3 [CT3] Provider Engagement [superseded]}

\label{element:a0777dbd9b0c4528817e99a7f34aa444}
\index{[CT3] Provider Engagement (II.1.3.4.7.3)}
{\small Capability · Proposed · superseded}\par
\href{../../../../continuum/strategy/capability/a0777dbd9b0c4528817e99a7f34aa444.md}{Historical canonical definition and provenance}.\par
\paragraph*{II.1.3.4.7.4 [CT3] Provider Relationship [superseded]}
\phantomsection
\label{navigation:II.1.3.4.7.4}
\addcontentsline{toc}{paragraph}{II.1.3.4.7.4 [CT3] Provider Relationship [superseded]}

\label{element:3bb8a3d5b62a4b779442290fe8d5106c}
\index{[CT3] Provider Relationship (II.1.3.4.7.4)}
{\small Capability · Proposed · superseded}\par
\href{../../../../continuum/strategy/capability/3bb8a3d5b62a4b779442290fe8d5106c.md}{Historical canonical definition and provenance}.\par
\paragraph*{II.1.3.4.7.5 [CT3] Vendor Liaison}
\phantomsection
\label{navigation:II.1.3.4.7.5}
\addcontentsline{toc}{paragraph}{II.1.3.4.7.5 [CT3] Vendor Liaison}

{\small Historical predecessor: R.2.112 \hyperref[element:29076cbb9fc5442fb25faf43947574f7]{[CT3] Vendor Liaison}; Canonical ID: \texttt{\seqsplit{29076cbb9fc5442fb25faf43947574f7}}. Superseded historical identity — direct correspondence · superseded.}\par
\label{element:66358bd6698e4050b12342522fbbb923}
\index{[CT3] Vendor Liaison (II.1.3.4.7.5)}
{\small Capability · Proposed · current}\par

Maintain the information and liaison arrangements needed to interact with vendors supporting healthcare provision.



\textbf{Overview}



Maintain the information and liaison arrangements needed to interact with vendors supporting healthcare provision. The reviewed hierarchy assigns this enduring ability to its stated business-enabling context.



\textbf{Intent}



Maintain the information and liaison arrangements needed to interact with vendors supporting healthcare provision.



\textbf{Boundary}



Information administration and coordination preserve the authority of the responsible source, clinical actor and governed domain; they confer no additional real-world management authority.


{\footnotesize \href{../../../../continuum/strategy/capability/66358bd6698e4050b12342522fbbb923.md}{Canonical record: 66358bd6698e4050b12342522fbbb923}; full qualifications and provenance.}\par
\subsection*{II.1.3.5 [CT1] Health Informatics Management}
\phantomsection
\label{navigation:II.1.3.5}
\addcontentsline{toc}{subsection}{II.1.3.5 [CT1] Health Informatics Management}

{\small Historical predecessor: R.2.118 \hyperref[element:50c56bac0dc540a2b392fc2376efc129]{[CT1] Health Informatics Management}; Canonical ID: \texttt{\seqsplit{50c56bac0dc540a2b392fc2376efc129}}. Superseded historical identity — direct correspondence · superseded.}\par
\label{element:a790b5f2c5a1497fa6ec09d7824da32e}
\index{[CT1] Health Informatics Management (II.1.3.5)}
{\small Capability · Proposed · current}\par

Establish and govern semantically meaningful health-information assets suitable for downstream research, population-health reporting and service planning in an adjacent analytical environment.



\textbf{Overview}



Establish and govern semantically meaningful health-information assets suitable for downstream research, population-health reporting and service planning in an adjacent analytical environment.



\textbf{Intent}



Establish and govern semantically meaningful health-information assets suitable for downstream research, population-health reporting and service planning in an adjacent analytical environment.



\textbf{Boundary}



Harmonia is adjacent to a Big Data/analytical environment. It establishes governed assets and reporting inputs; it does not become that analytical platform.


{\footnotesize \href{../../../../continuum/strategy/capability/a790b5f2c5a1497fa6ec09d7824da32e.md}{Canonical record: a790b5f2c5a1497fa6ec09d7824da32e}; full qualifications and provenance.}\par
\subsubsection*{II.1.3.5.1 [CT2] Population Health Information Reporting}
\phantomsection
\label{navigation:II.1.3.5.1}
\addcontentsline{toc}{subsubsection}{II.1.3.5.1 [CT2] Population Health Information Reporting}

{\small Historical predecessor: R.2.115 \hyperref[element:ced472b2d1a54016878ee541de969d9b]{[CT2] Population Health Information Reporting}; Canonical ID: \texttt{\seqsplit{ced472b2d1a54016878ee541de969d9b}}. Superseded historical identity — direct correspondence · superseded.}\par
\label{element:2135b45eab034370b0063b20e821985d}
\index{[CT2] Population Health Information Reporting (II.1.3.5.1)}
{\small Capability · Proposed · current}\par

Establish and provide governed information assets and reporting inputs for population-health purposes, including applicable notification obligations.



\textbf{Overview}



Establish and provide governed information assets and reporting inputs for population-health purposes, including applicable notification obligations.



\textbf{Intent}



Establish and provide governed information assets and reporting inputs for population-health purposes, including applicable notification obligations.



\textbf{Boundary}



Does not determine population-health findings, create statutory obligations or equate transport delivery with discharge of notification duty.


{\footnotesize \href{../../../../continuum/strategy/capability/2135b45eab034370b0063b20e821985d.md}{Canonical record: 2135b45eab034370b0063b20e821985d}; full qualifications and provenance.}\par
\subsubsection*{II.1.3.5.2 [CT2] Research Information Curation}
\phantomsection
\label{navigation:II.1.3.5.2}
\addcontentsline{toc}{subsubsection}{II.1.3.5.2 [CT2] Research Information Curation}

{\small Historical predecessor: R.2.116 \hyperref[element:11203b9c80964a9b96883eeab07f41d5]{[CT2] Research Information Curation}; Canonical ID: \texttt{\seqsplit{11203b9c80964a9b96883eeab07f41d5}}. Superseded historical identity — direct correspondence · superseded.}\par
\label{element:4f5b8cec90324c2989cdc30e80bc8cb9}
\index{[CT2] Research Information Curation (II.1.3.5.2)}
{\small Capability · Proposed · current}\par

Establish and curate governed data assets suitable for explicitly authorised research use, retaining source attribution and applicable release constraints.



\textbf{Overview}



Establish and curate governed data assets suitable for explicitly authorised research use, retaining source attribution and applicable release constraints.



\textbf{Intent}



Establish and curate governed data assets suitable for explicitly authorised research use, retaining source attribution and applicable release constraints.



\textbf{Boundary}



Does not grant ethics approval, recruit trial participants, manage trials or become the authority for derived research conclusions.


{\footnotesize \href{../../../../continuum/strategy/capability/4f5b8cec90324c2989cdc30e80bc8cb9.md}{Canonical record: 4f5b8cec90324c2989cdc30e80bc8cb9}; full qualifications and provenance.}\par
\subsubsection*{II.1.3.5.3 [CT2] Service Planning Data Management}
\phantomsection
\label{navigation:II.1.3.5.3}
\addcontentsline{toc}{subsubsection}{II.1.3.5.3 [CT2] Service Planning Data Management}

{\small Historical predecessor: R.2.117 \hyperref[element:19263e0802cc4b88af7a345f1410cffd]{[CT2] Service Planning Data Management}; Canonical ID: \texttt{\seqsplit{19263e0802cc4b88af7a345f1410cffd}}. Superseded historical identity — direct correspondence · superseded.}\par
\label{element:4873cd4b960d45d8b93560d5c8b98500}
\index{[CT2] Service Planning Data Management (II.1.3.5.3)}
{\small Capability · Proposed · current}\par

Establish and maintain governed information assets suitable for health-service planning, modelling and analysis.



\textbf{Overview}



Establish and maintain governed information assets suitable for health-service planning, modelling and analysis.



\textbf{Intent}



Establish and maintain governed information assets suitable for health-service planning, modelling and analysis.



\textbf{Boundary}



Observed utilisation, operational capacity and planned capacity are distinct; information provision does not commission services or allocate resources.


{\footnotesize \href{../../../../continuum/strategy/capability/4873cd4b960d45d8b93560d5c8b98500.md}{Canonical record: 4873cd4b960d45d8b93560d5c8b98500}; full qualifications and provenance.}\par
\subsection*{II.1.3.6 [CT1] Health Service Operations Enablement}
\phantomsection
\label{navigation:II.1.3.6}
\addcontentsline{toc}{subsection}{II.1.3.6 [CT1] Health Service Operations Enablement}

{\small Historical predecessor: R.2.160 \hyperref[element:a98424298dac46d8b2ddd109d928270a]{[CT1] Health Service Operations Enablement}; Canonical ID: \texttt{\seqsplit{a98424298dac46d8b2ddd109d928270a}}. Superseded historical identity — direct correspondence · superseded.}\par
\label{element:b13ee2c5a4ef4fd6be8a44795696dfc2}
\index{[CT1] Health Service Operations Enablement (II.1.3.6)}
{\small Capability · Proposed · current}\par

Enable health-service operations through information about flow, care-places, capacity, workforce capability, operational work, logistics and discharge progression.



\textbf{Overview}



Enable health-service operations through information about flow, care-places, capacity, workforce capability, operational work, logistics and discharge progression.



\textbf{Intent}



Enable health-service operations through information about flow, care-places, capacity, workforce capability, operational work, logistics and discharge progression.



\textbf{Boundary}



Coordination owns only explicitly assigned operational activity. It does not originate clinical clearance, determine clinical work or take workforce and resource governance from their sources.


{\footnotesize \href{../../../../continuum/strategy/capability/b13ee2c5a4ef4fd6be8a44795696dfc2.md}{Canonical record: b13ee2c5a4ef4fd6be8a44795696dfc2}; full qualifications and provenance.}\par
\subsubsection*{II.1.3.6.1 [CT2] Bed \& Care-Place Management}
\phantomsection
\label{navigation:II.1.3.6.1}
\addcontentsline{toc}{subsubsection}{II.1.3.6.1 [CT2] Bed \& Care-Place Management}

{\small Historical predecessor: R.2.121 \hyperref[element:c94d29b7d1c045dabc7814ceea795484]{[CT2] Bed \& Care-Place Management}; Canonical ID: \texttt{\seqsplit{c94d29b7d1c045dabc7814ceea795484}}. Superseded historical identity — direct correspondence · superseded.}\par
\label{element:f117ccec0bfa4f738e6ddb646b3534a2}
\index{[CT2] Bed \& Care-Place Management (II.1.3.6.1)}
{\small Capability · Canonical · current}\par

Bed \& Care-Place Management maintains operational place state and the turnover progression that makes a place usable again. Vacancy, reservation, cleaning and readiness must remain distinguishable so that an empty bed is not automatically treated as available for allocation. Notifications expose meaningful readiness changes to the responsible operational participants.



\textbf{Overview}



Bed \& Care-Place Management maintains operational place state and the turnover progression that makes a place usable again. Vacancy, reservation, cleaning and readiness must remain distinguishable so that an empty bed is not automatically treated as available for allocation. Notifications expose meaningful readiness changes to the responsible operational participants.



\textbf{Intent}



Bed \& Care-Place Management maintains operational place state and the turnover progression that makes a place usable again. Vacancy, reservation, cleaning and readiness must remain distinguishable so that an empty bed is not automatically treated as available for allocation. Notifications expose meaningful readiness changes to the responsible operational participants.



\textbf{Boundary}



Information administration does not confer authority over the represented person, organisation, service, place, device or clinical activity.



\medskip\hrule\medskip

{\footnotesize \href{../../../../continuum/strategy/capability/f117ccec0bfa4f738e6ddb646b3534a2.md}{Canonical record: f117ccec0bfa4f738e6ddb646b3534a2}; full qualifications and provenance.}\par
\paragraph*{II.1.3.6.1.1 [CT3] Bed \& Care-Place State Administration}
\phantomsection
\label{navigation:II.1.3.6.1.1}
\addcontentsline{toc}{paragraph}{II.1.3.6.1.1 [CT3] Bed \& Care-Place State Administration}

{\small Historical predecessor: R.2.119 \hyperref[element:dd2e2ee1a68845c087ba93610752181b]{[CT3] Bed \& Care-Place State Administration}; Canonical ID: \texttt{\seqsplit{dd2e2ee1a68845c087ba93610752181b}}. Superseded historical identity — direct correspondence · superseded.}\par
\label{element:26a9383d7bb146a19141558d232843de}
\index{[CT3] Bed \& Care-Place State Administration (II.1.3.6.1.1)}
{\small Capability · Proposed · current}\par

Administer the current operational occupancy, reservation, blocking and readiness states of beds and other care-places.



\textbf{Overview}



Administer the current operational occupancy, reservation, blocking and readiness states of beds and other care-places. The reviewed hierarchy assigns this enduring ability to its stated business-enabling context.



\textbf{Intent}



Administer the current operational occupancy, reservation, blocking and readiness states of beds and other care-places.



\textbf{Boundary}



Information administration and coordination preserve the authority of the responsible source, clinical actor and governed domain; they confer no additional real-world management authority.


{\footnotesize \href{../../../../continuum/strategy/capability/26a9383d7bb146a19141558d232843de.md}{Canonical record: 26a9383d7bb146a19141558d232843de}; full qualifications and provenance.}\par
\subparagraph*{II.1.3.6.1.1.1 [Feature] Bed Availability Tracking}
\phantomsection
\label{navigation:II.1.3.6.1.1.1}
\addcontentsline{toc}{subparagraph}{II.1.3.6.1.1.1 [Feature] Bed Availability Tracking}

\label{element:7c61f21e558c44ab92c267a781116d1d}
\index{[Feature] Bed Availability Tracking (II.1.3.6.1.1.1)}
{\small Capability / Feature · Canonical · current}\par

Maintain operational status of all registered beds and care-places (VACANT, OCCUPIED, BLOCKED, RESERVED, DIRTY, CLEANING, CLEAN).



Availability tracking preserves the operational distinctions among vacancy, occupancy, reservation, blockage and cleaning conditions. These states support meaningful care-place coordination without equating an unoccupied bed with a ready place or usable service capacity.


{\footnotesize \href{../../../../continuum/strategy/capability/feature/7c61f21e558c44ab92c267a781116d1d.md}{Canonical record: 7c61f21e558c44ab92c267a781116d1d}; full qualifications and provenance.}\par
\paragraph*{II.1.3.6.1.2 [CT3] Bed \& Care-Place Turnover Coordination}
\phantomsection
\label{navigation:II.1.3.6.1.2}
\addcontentsline{toc}{paragraph}{II.1.3.6.1.2 [CT3] Bed \& Care-Place Turnover Coordination}

{\small Historical predecessor: R.2.120 \hyperref[element:a50662edda954bc58c44184d84591d5b]{[CT3] Bed \& Care-Place Turnover Coordination}; Canonical ID: \texttt{\seqsplit{a50662edda954bc58c44184d84591d5b}}. Superseded historical identity — direct correspondence · superseded.}\par
\label{element:3c9464f9eae1415d8fa041bd1f3da4f2}
\index{[CT3] Bed \& Care-Place Turnover Coordination (II.1.3.6.1.2)}
{\small Capability · Proposed · current}\par

Coordinate the operational transition of a released care-place through cleaning, preparation and readiness notifications.



\textbf{Overview}



Coordinate the operational transition of a released care-place through cleaning, preparation and readiness notifications. The reviewed hierarchy assigns this enduring ability to its stated business-enabling context.



\textbf{Intent}



Coordinate the operational transition of a released care-place through cleaning, preparation and readiness notifications.



\textbf{Boundary}



Information administration and coordination preserve the authority of the responsible source, clinical actor and governed domain; they confer no additional real-world management authority.


{\footnotesize \href{../../../../continuum/strategy/capability/3c9464f9eae1415d8fa041bd1f3da4f2.md}{Canonical record: 3c9464f9eae1415d8fa041bd1f3da4f2}; full qualifications and provenance.}\par
\subparagraph*{II.1.3.6.1.2.1 [Feature] Bed Readiness Notification}
\phantomsection
\label{navigation:II.1.3.6.1.2.1}
\addcontentsline{toc}{subparagraph}{II.1.3.6.1.2.1 [Feature] Bed Readiness Notification}

\label{element:6a9cf80d6f82425cb2aaae27f4c5c349}
\index{[Feature] Bed Readiness Notification (II.1.3.6.1.2.1)}
{\small Capability / Feature · Canonical · current}\par

Notify bed allocation managers and clinical units when a cleaned bed achieves READY status.



Readiness notification exposes the established readiness of a care place to the responsible allocation participants. It makes a relevant transition actionable without deciding the clinical suitability of the next allocation or establishing whole-service readiness.


{\footnotesize \href{../../../../continuum/strategy/capability/feature/6a9cf80d6f82425cb2aaae27f4c5c349.md}{Canonical record: 6a9cf80d6f82425cb2aaae27f4c5c349}; full qualifications and provenance.}\par
\subparagraph*{II.1.3.6.1.2.2 [Feature] Bed State Progression}
\phantomsection
\label{navigation:II.1.3.6.1.2.2}
\addcontentsline{toc}{subparagraph}{II.1.3.6.1.2.2 [Feature] Bed State Progression}

\label{element:c903b83181bb456a89dc9d450063eff6}
\index{[Feature] Bed State Progression (II.1.3.6.1.2.2)}
{\small Capability / Feature · Canonical · current}\par

Progress bed states in response to clinical and environmental events (e.g., patient discharge automatically transitions bed to DIRTY).



Bed-state progression applies the consequences of relevant clinical and environmental events to the managed place state. The resulting change is an operational fact, such as a need for cleaning after departure, rather than a transfer of clinical responsibility.


{\footnotesize \href{../../../../continuum/strategy/capability/feature/c903b83181bb456a89dc9d450063eff6.md}{Canonical record: c903b83181bb456a89dc9d450063eff6}; full qualifications and provenance.}\par
\subsubsection*{II.1.3.6.2 [CT2] Care Movement \& Transition Coordination [superseded]}
\phantomsection
\label{navigation:II.1.3.6.2}
\addcontentsline{toc}{subsubsection}{II.1.3.6.2 [CT2] Care Movement \& Transition Coordination [superseded]}

\label{element:59b9331fa3c04ef09bf963ad8a5d0d71}
\index{[CT2] Care Movement \& Transition Coordination (II.1.3.6.2)}
{\small Capability · Proposed · superseded}\par
\href{../../../../continuum/strategy/capability/59b9331fa3c04ef09bf963ad8a5d0d71.md}{Historical canonical definition and provenance}.\par
\subsubsection*{II.1.3.6.3 [CT2] Care Resource \& Capacity Coordination [superseded]}
\phantomsection
\label{navigation:II.1.3.6.3}
\addcontentsline{toc}{subsubsection}{II.1.3.6.3 [CT2] Care Resource \& Capacity Coordination [superseded]}

\label{element:96f71cb8fa6844d7aa7e9f16759a3a62}
\index{[CT2] Care Resource \& Capacity Coordination (II.1.3.6.3)}
{\small Capability · Proposed · superseded}\par
\href{../../../../continuum/strategy/capability/96f71cb8fa6844d7aa7e9f16759a3a62.md}{Historical canonical definition and provenance}.\par
\paragraph*{II.1.3.6.3.1 [CT3] Clinical Resource Management [superseded]}
\phantomsection
\label{navigation:II.1.3.6.3.1}
\addcontentsline{toc}{paragraph}{II.1.3.6.3.1 [CT3] Clinical Resource Management [superseded]}

\label{element:a5faedcd97064c28b9e7e551e9106ea1}
\index{[CT3] Clinical Resource Management (II.1.3.6.3.1)}
{\small Capability · Canonical · superseded}\par
\href{../../../../continuum/strategy/capability/a5faedcd97064c28b9e7e551e9106ea1.md}{Historical canonical definition and provenance}.\par
\paragraph*{II.1.3.6.3.2 [CT3] Service Offering Readiness Coordination}
\phantomsection
\label{navigation:II.1.3.6.3.2}
\addcontentsline{toc}{paragraph}{II.1.3.6.3.2 [CT3] Service Offering Readiness Coordination}

\label{element:0caecaba5d3e4781b8671b12c13ee283}
\index{[CT3] Service Offering Readiness Coordination (II.1.3.6.3.2)}
{\small Capability · Proposed · pending\_adjudication}\par

Assemble the dependency evidence needed to determine and communicate whether a defined service offering is operationally ready in a particular context.



\textbf{Overview}



Collective readiness depends on suitable workforce, spaces, equipment, materials and supporting services. This proposal evaluates their joint availability for an offering without treating one available bed, an operating-hours declaration or aggregate capacity as service readiness.



\textbf{Intent}



Distinguish collective service deliverability from nominal availability or readiness of any single dependency.



\textbf{Boundary}



Contributing capabilities retain source readiness and resource management. Clinical go/no-go decisions and service governance are not granted by this information-based coordination.


{\small Recorded unresolved matters: source-qualified matters; exact scope remains in the canonical record.}\par
{\footnotesize \href{../../../../continuum/strategy/capability/0caecaba5d3e4781b8671b12c13ee283.md}{Canonical record: 0caecaba5d3e4781b8671b12c13ee283}; full qualifications and provenance.}\par
\subparagraph*{II.1.3.6.3.2.1 [Feature] Service Readiness Dependency Assessment}
\phantomsection
\label{navigation:II.1.3.6.3.2.1}
\addcontentsline{toc}{subparagraph}{II.1.3.6.3.2.1 [Feature] Service Readiness Dependency Assessment}

\label{element:4d917929d5cf43218049a84cf46662b0}
\index{[Feature] Service Readiness Dependency Assessment (II.1.3.6.3.2.1)}
{\small Capability / Feature · Proposed · pending\_adjudication}\par

Associate time- and context-qualified dependency evidence to establish a bounded readiness view for an offering.



\textbf{Overview}



An offering may be in operating hours but not deliverable because staffing, equipment, materials or supporting services are missing. A readiness view retains insufficiency and the owners of those dependency assertions.



\textbf{Boundary}



Readiness evidence does not originate clinical clearance or transfer resource management.


{\small Recorded unresolved matters: source-qualified matters; exact scope remains in the canonical record.}\par
{\footnotesize \href{../../../../continuum/strategy/capability/feature/4d917929d5cf43218049a84cf46662b0.md}{Canonical record: 4d917929d5cf43218049a84cf46662b0}; full qualifications and provenance.}\par
\subsubsection*{II.1.3.6.4 [CT2] Care Setting Operations Enablement [superseded]}
\phantomsection
\label{navigation:II.1.3.6.4}
\addcontentsline{toc}{subsubsection}{II.1.3.6.4 [CT2] Care Setting Operations Enablement [superseded]}

\label{element:a03741e6e98642bea44680d23fdcb434}
\index{[CT2] Care Setting Operations Enablement (II.1.3.6.4)}
{\small Capability · Proposed · superseded}\par
\href{../../../../continuum/strategy/capability/a03741e6e98642bea44680d23fdcb434.md}{Historical canonical definition and provenance}.\par
\subsubsection*{II.1.3.6.5 [CT2] Clinic \& Practice Operations}
\phantomsection
\label{navigation:II.1.3.6.5}
\addcontentsline{toc}{subsubsection}{II.1.3.6.5 [CT2] Clinic \& Practice Operations}

{\small Historical predecessor: R.2.123 \hyperref[element:0856d2daf2ba4b92b5390b38b98efb54]{[CT2] Clinic \& Practice Operations}; Canonical ID: \texttt{\seqsplit{0856d2daf2ba4b92b5390b38b98efb54}}. Superseded historical identity — direct correspondence · superseded.}\par
\label{element:3dffe63e03124efab255204ee4e9a5d9}
\index{[CT2] Clinic \& Practice Operations (II.1.3.6.5)}
{\small Capability · Canonical · current}\par

Clinic \& Practice Operations tracks operating sessions and the operational progression of people through the clinic environment. It makes delays, waiting and completion context visible alongside service-specific capacity responsibilities. This is real-world clinic operation, distinct from the technical queueing machinery that might realise the information flow.



\textbf{Overview}



Clinic \& Practice Operations tracks operating sessions and the operational progression of people through the clinic environment. It makes delays, waiting and completion context visible alongside service-specific capacity responsibilities. This is real-world clinic operation, distinct from the technical queueing machinery that might realise the information flow.



\textbf{Intent}



Clinic \& Practice Operations tracks operating sessions and the operational progression of people through the clinic environment. It makes delays, waiting and completion context visible alongside service-specific capacity responsibilities. This is real-world clinic operation, distinct from the technical queueing machinery that might realise the information flow.



\textbf{Boundary}



Information administration does not confer authority over the represented person, organisation, service, place, device or clinical activity.



\medskip\hrule\medskip

{\footnotesize \href{../../../../continuum/strategy/capability/3dffe63e03124efab255204ee4e9a5d9.md}{Canonical record: 3dffe63e03124efab255204ee4e9a5d9}; full qualifications and provenance.}\par
\paragraph*{II.1.3.6.5.1 [CT3] Session \& Queue Coordination}
\phantomsection
\label{navigation:II.1.3.6.5.1}
\addcontentsline{toc}{paragraph}{II.1.3.6.5.1 [CT3] Session \& Queue Coordination}

{\small Historical predecessor: R.2.122 \hyperref[element:5298c6170cef4f05a6d956da716071bd]{[CT3] Session \& Queue Coordination}; Canonical ID: \texttt{\seqsplit{5298c6170cef4f05a6d956da716071bd}}. Superseded historical identity — direct correspondence · superseded.}\par
\label{element:22981a95ca6e466582363706fde8fceb}
\index{[CT3] Session \& Queue Coordination (II.1.3.6.5.1)}
{\small Capability · Proposed · current}\par

Coordinate clinic/practice session state, operational queue progression and context-specific session capacity information.



\textbf{Overview}



Coordinate clinic/practice session state, operational queue progression and context-specific session capacity information. The reviewed hierarchy assigns this enduring ability to its stated business-enabling context.



\textbf{Intent}



Coordinate clinic/practice session state, operational queue progression and context-specific session capacity information.



\textbf{Boundary}



Information administration and coordination preserve the authority of the responsible source, clinical actor and governed domain; they confer no additional real-world management authority.


{\footnotesize \href{../../../../continuum/strategy/capability/22981a95ca6e466582363706fde8fceb.md}{Canonical record: 22981a95ca6e466582363706fde8fceb}; full qualifications and provenance.}\par
\subparagraph*{II.1.3.6.5.1.1 [Feature] Clinic Session Capacity Management}
\phantomsection
\label{navigation:II.1.3.6.5.1.1}
\addcontentsline{toc}{subparagraph}{II.1.3.6.5.1.1 [Feature] Clinic Session Capacity Management}

\label{element:02bdd6a6f22e47f5a29d722ab57f9915}
\index{[Feature] Clinic Session Capacity Management (II.1.3.6.5.1.1)}
{\small Capability / Feature · Proposed · current}\par

Represent available, committed and utilised session capacity under clinic-specific constraints.



\textbf{Overview}



Clinic capacity depends on session staffing and service mix as well as places; queue progression alone is insufficient.



\textbf{Boundary}



Local responsibility retains its capacity meaning; Service Capacity Management consumes it without acquiring ownership or imposing one common capacity model.


{\footnotesize \href{../../../../continuum/strategy/capability/feature/02bdd6a6f22e47f5a29d722ab57f9915.md}{Canonical record: 02bdd6a6f22e47f5a29d722ab57f9915}; full qualifications and provenance.}\par
\subparagraph*{II.1.3.6.5.1.2 [Feature] Clinic Session Tracking}
\phantomsection
\label{navigation:II.1.3.6.5.1.2}
\addcontentsline{toc}{subparagraph}{II.1.3.6.5.1.2 [Feature] Clinic Session Tracking}

\label{element:03e185a6e0d145c3975b3f4715c94207}
\index{[Feature] Clinic Session Tracking (II.1.3.6.5.1.2)}
{\small Capability / Feature · Canonical · current}\par

Ingest and track clinic operating session start, delay, and completion states.



Clinic-session tracking records the operating session's start, delay and completion as service-level context. It is distinct from an individual patient's queue state: knowing where one patient is waiting does not establish whether the clinic session is running as planned.


{\footnotesize \href{../../../../continuum/strategy/capability/feature/03e185a6e0d145c3975b3f4715c94207.md}{Canonical record: 03e185a6e0d145c3975b3f4715c94207}; full qualifications and provenance.}\par
\subparagraph*{II.1.3.6.5.1.3 [Feature] Operational Queue Progression}
\phantomsection
\label{navigation:II.1.3.6.5.1.3}
\addcontentsline{toc}{subparagraph}{II.1.3.6.5.1.3 [Feature] Operational Queue Progression}

\label{element:5852c7d4c9ac4bbf99e7df008d0ed770}
\index{[Feature] Operational Queue Progression (II.1.3.6.5.1.3)}
{\small Capability / Feature · Canonical · current}\par

Track client queue progression through reception, waiting area, consultation, and billing.



Operational queue progression makes the patient's passage through clinic stages visible. The queue concerns real-world arrival, waiting and consultation context, rather than a broker queue, and does not give Harmonia authority to determine clinical priority or own the clinician's worklist.


{\footnotesize \href{../../../../continuum/strategy/capability/feature/5852c7d4c9ac4bbf99e7df008d0ed770.md}{Canonical record: 5852c7d4c9ac4bbf99e7df008d0ed770}; full qualifications and provenance.}\par
\subsubsection*{II.1.3.6.6 [CT2] Clinical Competency Administration}
\phantomsection
\label{navigation:II.1.3.6.6}
\addcontentsline{toc}{subsubsection}{II.1.3.6.6 [CT2] Clinical Competency Administration}

{\small Historical predecessor: R.2.124 \hyperref[element:8a4b918e1c9d47bb99daf2b42008ab39]{[CT2] Clinical Competency Administration}; Canonical ID: \texttt{\seqsplit{8a4b918e1c9d47bb99daf2b42008ab39}}. Superseded historical identity — direct correspondence · superseded.}\par
\label{element:8a99f4a6e5054d6c997bacec04805c01}
\index{[CT2] Clinical Competency Administration (II.1.3.6.6)}
{\small Capability · Canonical · current}\par

Administer information about what clinical capability an organisation can provide and what a practitioner is competent to perform.



\textbf{Overview}



Administer information about what clinical capability an organisation can provide and what a practitioner is competent to perform.



\textbf{Intent}



Administer information about what clinical capability an organisation can provide and what a practitioner is competent to perform.



\textbf{Boundary}



Information administration does not confer authority over the represented person, organisation, service, place, device or clinical activity.



\medskip\hrule\medskip

{\footnotesize \href{../../../../continuum/strategy/capability/8a99f4a6e5054d6c997bacec04805c01.md}{Canonical record: 8a99f4a6e5054d6c997bacec04805c01}; full qualifications and provenance.}\par
\subsubsection*{II.1.3.6.7 [CT2] Clinical Logistics Coordination}
\phantomsection
\label{navigation:II.1.3.6.7}
\addcontentsline{toc}{subsubsection}{II.1.3.6.7 [CT2] Clinical Logistics Coordination}

{\small Historical predecessor: R.2.127 \hyperref[element:90a782ff2d004fa984ae0465528b2253]{[CT2] Clinical Logistics Coordination}; Canonical ID: \texttt{\seqsplit{90a782ff2d004fa984ae0465528b2253}}. Superseded historical identity — direct correspondence · superseded.}\par
\label{element:0c27be11b4a44cb19e992309b8146a17}
\index{[CT2] Clinical Logistics Coordination (II.1.3.6.7)}
{\small Capability · Canonical · current}\par

Coordinate specimen metadata and specimen movement traceability across responsible collection, courier and diagnostic services.



\textbf{Overview}



Coordinate specimen metadata and specimen movement traceability across responsible collection, courier and diagnostic services.



\textbf{Intent}



Coordinate specimen metadata and specimen movement traceability across responsible collection, courier and diagnostic services.



\textbf{Boundary}



Includes specimen metadata and specimen movement traceability. Supply/materiel movement belongs to Health Materiel Traceability Administration.



\medskip\hrule\medskip

{\footnotesize \href{../../../../continuum/strategy/capability/0c27be11b4a44cb19e992309b8146a17.md}{Canonical record: 0c27be11b4a44cb19e992309b8146a17}; full qualifications and provenance.}\par
\paragraph*{II.1.3.6.7.1 [CT3] Specimen Metadata Management}
\phantomsection
\label{navigation:II.1.3.6.7.1}
\addcontentsline{toc}{paragraph}{II.1.3.6.7.1 [CT3] Specimen Metadata Management}

{\small Historical predecessor: R.2.125 \hyperref[element:3b000d3712534d16a48470f79df8e33a]{[CT3] Specimen Metadata Management}; Canonical ID: \texttt{\seqsplit{3b000d3712534d16a48470f79df8e33a}}. Superseded historical identity — direct correspondence · superseded.}\par
\label{element:738736425b7446fdae47b06c91c4f69c}
\index{[CT3] Specimen Metadata Management (II.1.3.6.7.1)}
{\small Capability · Proposed · current}\par

Manage the identifying, descriptive, custody and relevant handling information needed to interpret a transported specimen.



\textbf{Overview}



Manage the identifying, descriptive, custody and relevant handling information needed to interpret a transported specimen. The reviewed hierarchy assigns this enduring ability to its stated business-enabling context.



\textbf{Intent}



Manage the identifying, descriptive, custody and relevant handling information needed to interpret a transported specimen.



\textbf{Boundary}



Information administration and coordination preserve the authority of the responsible source, clinical actor and governed domain; they confer no additional real-world management authority.


{\footnotesize \href{../../../../continuum/strategy/capability/738736425b7446fdae47b06c91c4f69c.md}{Canonical record: 738736425b7446fdae47b06c91c4f69c}; full qualifications and provenance.}\par
\paragraph*{II.1.3.6.7.2 [CT3] Specimen Movement Traceability}
\phantomsection
\label{navigation:II.1.3.6.7.2}
\addcontentsline{toc}{paragraph}{II.1.3.6.7.2 [CT3] Specimen Movement Traceability}

{\small Historical predecessor: R.2.126 \hyperref[element:129119d7b2744ffd805c58f25557fa27]{[CT3] Specimen Movement Traceability}; Canonical ID: \texttt{\seqsplit{129119d7b2744ffd805c58f25557fa27}}. Superseded historical identity — direct correspondence · superseded.}\par
\label{element:bfa2e9a84ddd4287897573dc0d917a00}
\index{[CT3] Specimen Movement Traceability (II.1.3.6.7.2)}
{\small Capability · Proposed · current}\par

Maintain attributable pickup, custody, transport-condition and delivery traceability for clinical specimens.



\textbf{Overview}



Maintain attributable pickup, custody, transport-condition and delivery traceability for clinical specimens. The reviewed hierarchy assigns this enduring ability to its stated business-enabling context.



\textbf{Intent}



Maintain attributable pickup, custody, transport-condition and delivery traceability for clinical specimens.



\textbf{Boundary}



Information administration and coordination preserve the authority of the responsible source, clinical actor and governed domain; they confer no additional real-world management authority.


{\footnotesize \href{../../../../continuum/strategy/capability/bfa2e9a84ddd4287897573dc0d917a00.md}{Canonical record: bfa2e9a84ddd4287897573dc0d917a00}; full qualifications and provenance.}\par
\subparagraph*{II.1.3.6.7.2.1 [Feature] Pathology Specimen Transport Tracking}
\phantomsection
\label{navigation:II.1.3.6.7.2.1}
\addcontentsline{toc}{subparagraph}{II.1.3.6.7.2.1 [Feature] Pathology Specimen Transport Tracking}

\label{element:34f182cfc6b945d090a880943bbb4959}
\index{[Feature] Pathology Specimen Transport Tracking (II.1.3.6.7.2.1)}
{\small Capability / Feature · Canonical · current}\par

Track pickup, custody chain, transport conditions, and laboratory delivery of clinical specimens.



Specimen transport tracking retains pickup, custody and condition context through laboratory delivery. It supports interpretation of the material's journey without treating a delivered package as proof of specimen suitability, completed analysis or an authorised diagnostic report.


{\footnotesize \href{../../../../continuum/strategy/capability/feature/34f182cfc6b945d090a880943bbb4959.md}{Canonical record: 34f182cfc6b945d090a880943bbb4959}; full qualifications and provenance.}\par
\subsubsection*{II.1.3.6.8 [CT2] Clinical Qualification Administration}
\phantomsection
\label{navigation:II.1.3.6.8}
\addcontentsline{toc}{subsubsection}{II.1.3.6.8 [CT2] Clinical Qualification Administration}

{\small Historical predecessor: R.2.130 \hyperref[element:5fea7aaac4384245b1f08a3d8a076bd9]{[CT2] Clinical Qualification Administration}; Canonical ID: \texttt{\seqsplit{5fea7aaac4384245b1f08a3d8a076bd9}}. Superseded historical identity — direct correspondence · superseded.}\par
\label{element:0511a705d74a4f48a0ce92590ae764fa}
\index{[CT2] Clinical Qualification Administration (II.1.3.6.8)}
{\small Capability · Canonical · current}\par

Administer organisation clinical certification/compliance information and provider qualification evidence with its externally established standing.



\textbf{Overview}



Administer organisation clinical certification/compliance information and provider qualification evidence with its externally established standing.



\textbf{Intent}



Administer organisation clinical certification/compliance information and provider qualification evidence with its externally established standing.



\textbf{Boundary}



Information administration does not confer authority over the represented person, organisation, service, place, device or clinical activity.



\medskip\hrule\medskip

{\footnotesize \href{../../../../continuum/strategy/capability/0511a705d74a4f48a0ce92590ae764fa.md}{Canonical record: 0511a705d74a4f48a0ce92590ae764fa}; full qualifications and provenance.}\par
\paragraph*{II.1.3.6.8.1 [CT3] Organisation Clinical Certification \& Compliance Support}
\phantomsection
\label{navigation:II.1.3.6.8.1}
\addcontentsline{toc}{paragraph}{II.1.3.6.8.1 [CT3] Organisation Clinical Certification \& Compliance Support}

{\small Historical predecessor: R.2.128 \hyperref[element:cee76a930a4e4e7d9c49357ba290c8bf]{[CT3] Organisation Clinical Certification \& Compliance Support}; Canonical ID: \texttt{\seqsplit{cee76a930a4e4e7d9c49357ba290c8bf}}. Superseded historical identity — direct correspondence · superseded.}\par
\label{element:4de624e7bf49472f91e3e5c734ee8725}
\index{[CT3] Organisation Clinical Certification \& Compliance Support (II.1.3.6.8.1)}
{\small Capability · Proposed · current}\par

Administer supplied organisation clinical certification and conformance evidence for use in service qualification and participation.



\textbf{Overview}



Administer supplied organisation clinical certification and conformance evidence for use in service qualification and participation. The reviewed hierarchy assigns this enduring ability to its stated business-enabling context.



\textbf{Intent}



Administer supplied organisation clinical certification and conformance evidence for use in service qualification and participation.



\textbf{Boundary}



Information administration and coordination preserve the authority of the responsible source, clinical actor and governed domain; they confer no additional real-world management authority.


{\footnotesize \href{../../../../continuum/strategy/capability/4de624e7bf49472f91e3e5c734ee8725.md}{Canonical record: 4de624e7bf49472f91e3e5c734ee8725}; full qualifications and provenance.}\par
\paragraph*{II.1.3.6.8.2 [CT3] Provider Clinical Qualification Tracking}
\phantomsection
\label{navigation:II.1.3.6.8.2}
\addcontentsline{toc}{paragraph}{II.1.3.6.8.2 [CT3] Provider Clinical Qualification Tracking}

{\small Historical predecessor: R.2.129 \hyperref[element:ab0415e2ea0a4da4ad3dc46997d48bc3]{[CT3] Provider Clinical Qualification Tracking}; Canonical ID: \texttt{\seqsplit{ab0415e2ea0a4da4ad3dc46997d48bc3}}. Superseded historical identity — direct correspondence · superseded.}\par
\label{element:8fd5c78eba3046fa9153fb4efb44e9c8}
\index{[CT3] Provider Clinical Qualification Tracking (II.1.3.6.8.2)}
{\small Capability · Proposed · current}\par

Track supplied provider clinical qualifications, competency certificates and privilege evidence with effective dates and issuing-authority standing.



\textbf{Overview}



Track supplied provider clinical qualifications, competency certificates and privilege evidence with effective dates and issuing-authority standing. The reviewed hierarchy assigns this enduring ability to its stated business-enabling context.



\textbf{Intent}



Track supplied provider clinical qualifications, competency certificates and privilege evidence with effective dates and issuing-authority standing.



\textbf{Boundary}



Information administration and coordination preserve the authority of the responsible source, clinical actor and governed domain; they confer no additional real-world management authority.


{\footnotesize \href{../../../../continuum/strategy/capability/8fd5c78eba3046fa9153fb4efb44e9c8.md}{Canonical record: 8fd5c78eba3046fa9153fb4efb44e9c8}; full qualifications and provenance.}\par
\subparagraph*{II.1.3.6.8.2.1 [Feature] Clinical Privileging Record}
\phantomsection
\label{navigation:II.1.3.6.8.2.1}
\addcontentsline{toc}{subparagraph}{II.1.3.6.8.2.1 [Feature] Clinical Privileging Record}

\label{element:5073df3c6bb94c218ca19b74550b19fa}
\index{[Feature] Clinical Privileging Record (II.1.3.6.8.2.1)}
{\small Capability / Feature · Canonical · current}\par

Track active clinical admitting and practice privileges within healthcare facilities.



The privileging record makes applicable admitting and practice permissions available as contextual evidence. It represents the health service's established privilege standing rather than granting new authority merely because a record exists in the directory.


{\footnotesize \href{../../../../continuum/strategy/capability/feature/5073df3c6bb94c218ca19b74550b19fa.md}{Canonical record: 5073df3c6bb94c218ca19b74550b19fa}; full qualifications and provenance.}\par
\subparagraph*{II.1.3.6.8.2.2 [Feature] Operational Credential Check}
\phantomsection
\label{navigation:II.1.3.6.8.2.2}
\addcontentsline{toc}{subparagraph}{II.1.3.6.8.2.2 [Feature] Operational Credential Check}

\label{element:1a76831985fa461c891a3428e5dea86d}
\index{[Feature] Operational Credential Check (II.1.3.6.8.2.2)}
{\small Capability / Feature · Canonical · current}\par

Validate that a practitioner possesses active credentials and required competency certs prior to dispatching specialized work.



The qualification checkpoint considers active professional credentials and competency evidence before specialised dispatch. Its historical credential terminology does not make it the operational badge or physical-access responsibility of Credential Management, nor does a qualification check alone establish all required privilege.


{\footnotesize \href{../../../../continuum/strategy/capability/feature/1a76831985fa461c891a3428e5dea86d.md}{Canonical record: 1a76831985fa461c891a3428e5dea86d}; full qualifications and provenance.}\par
\subsubsection*{II.1.3.6.9 [CT2] Clinical Rostering}
\phantomsection
\label{navigation:II.1.3.6.9}
\addcontentsline{toc}{subsubsection}{II.1.3.6.9 [CT2] Clinical Rostering}

{\small Historical predecessor: R.2.132 \hyperref[element:b6527516a69c499d8cc9c8fe881169fe]{[CT2] Clinical Rostering}; Canonical ID: \texttt{\seqsplit{b6527516a69c499d8cc9c8fe881169fe}}. Superseded historical identity — direct correspondence · superseded.}\par
\label{element:192ca88a289544dcbf6adaf643f39969}
\index{[CT2] Clinical Rostering (II.1.3.6.9)}
{\small Capability · Canonical · current}\par

Clinical Rostering provides the Adjacent context of planned workforce availability. It is distinct from consuming roster projections for operational views or resolving who is currently on call. The present architecture retains the context without establishing a Harmonia Feature decomposition for authoring or managing the authoritative roster.



\textbf{Overview}



Clinical Rostering provides the Adjacent context of planned workforce availability. It is distinct from consuming roster projections for operational views or resolving who is currently on call. The present architecture retains the context without establishing a Harmonia Feature decomposition for authoring or managing the authoritative roster.



\textbf{Intent}



Clinical Rostering provides the Adjacent context of planned workforce availability. It is distinct from consuming roster projections for operational views or resolving who is currently on call. The present architecture retains the context without establishing a Harmonia Feature decomposition for authoring or managing the authoritative roster.



\textbf{Boundary}



Information administration does not confer authority over the represented person, organisation, service, place, device or clinical activity.



\medskip\hrule\medskip

{\footnotesize \href{../../../../continuum/strategy/capability/192ca88a289544dcbf6adaf643f39969.md}{Canonical record: 192ca88a289544dcbf6adaf643f39969}; full qualifications and provenance.}\par
\paragraph*{II.1.3.6.9.1 [CT3] Clinical Roster Capture \& Communication}
\phantomsection
\label{navigation:II.1.3.6.9.1}
\addcontentsline{toc}{paragraph}{II.1.3.6.9.1 [CT3] Clinical Roster Capture \& Communication}

{\small Historical predecessor: R.2.131 \hyperref[element:2054d57bb33a4361b3ef0824545a1784]{[CT3] Clinical Roster Capture \& Communication}; Canonical ID: \texttt{\seqsplit{2054d57bb33a4361b3ef0824545a1784}}. Superseded historical identity — direct correspondence · superseded.}\par
\label{element:8243cba81fb6449781d10823ae194237}
\index{[CT3] Clinical Roster Capture \& Communication (II.1.3.6.9.1)}
{\small Capability · Proposed · current}\par

Capture and communicate supplied clinical roster information, its effective periods and subsequent coverage changes.



\textbf{Overview}



Capture and communicate supplied clinical roster information, its effective periods and subsequent coverage changes. The reviewed hierarchy assigns this enduring ability to its stated business-enabling context.



\textbf{Intent}



Capture and communicate supplied clinical roster information, its effective periods and subsequent coverage changes.



\textbf{Boundary}



Information administration and coordination preserve the authority of the responsible source, clinical actor and governed domain; they confer no additional real-world management authority.


{\footnotesize \href{../../../../continuum/strategy/capability/8243cba81fb6449781d10823ae194237.md}{Canonical record: 8243cba81fb6449781d10823ae194237}; full qualifications and provenance.}\par
\subsubsection*{II.1.3.6.10 [CT2] Clinical Supervision}
\phantomsection
\label{navigation:II.1.3.6.10}
\addcontentsline{toc}{subsubsection}{II.1.3.6.10 [CT2] Clinical Supervision}

{\small Historical predecessor: R.2.134 \hyperref[element:50c838981f3946d78e8f78d18e31c70c]{[CT2] Clinical Supervision}; Canonical ID: \texttt{\seqsplit{50c838981f3946d78e8f78d18e31c70c}}. Superseded historical identity — direct correspondence · superseded.}\par
\label{element:f697b1ac79b544749315f31819b17969}
\index{[CT2] Clinical Supervision (II.1.3.6.10)}
{\small Capability · Canonical · current}\par

Enable responsible clinical participants to identify and use the applicable clinical escalation path.



\textbf{Overview}



Enable responsible clinical participants to identify and use the applicable clinical escalation path.



\textbf{Intent}



Enable responsible clinical participants to identify and use the applicable clinical escalation path.



\textbf{Boundary}



Information administration does not confer authority over the represented person, organisation, service, place, device or clinical activity.



\medskip\hrule\medskip

{\footnotesize \href{../../../../continuum/strategy/capability/f697b1ac79b544749315f31819b17969.md}{Canonical record: f697b1ac79b544749315f31819b17969}; full qualifications and provenance.}\par
\paragraph*{II.1.3.6.10.1 [CT3] Clinical Escalation Path Management}
\phantomsection
\label{navigation:II.1.3.6.10.1}
\addcontentsline{toc}{paragraph}{II.1.3.6.10.1 [CT3] Clinical Escalation Path Management}

{\small Historical predecessor: R.2.133 \hyperref[element:759a2e52172345e8ad1bc96d7e5ddad3]{[CT3] Clinical Escalation Path Management}; Canonical ID: \texttt{\seqsplit{759a2e52172345e8ad1bc96d7e5ddad3}}. Superseded historical identity — direct correspondence · superseded.}\par
\label{element:cca86a32ba8c4dda84caeaa248e9f444}
\index{[CT3] Clinical Escalation Path Management (II.1.3.6.10.1)}
{\small Capability · Proposed · current}\par

Manage clinical escalation paths and their applicability, contacts and progression conditions for use by responsible clinical participants.



\textbf{Overview}



Manage clinical escalation paths and their applicability, contacts and progression conditions for use by responsible clinical participants. The reviewed hierarchy assigns this enduring ability to its stated business-enabling context.



\textbf{Intent}



Manage clinical escalation paths and their applicability, contacts and progression conditions for use by responsible clinical participants.



\textbf{Boundary}



Harmonia manages the escalation path; responsible clinical participants perform supervision and retain clinical judgement.


{\footnotesize \href{../../../../continuum/strategy/capability/cca86a32ba8c4dda84caeaa248e9f444.md}{Canonical record: cca86a32ba8c4dda84caeaa248e9f444}; full qualifications and provenance.}\par
\subsubsection*{II.1.3.6.11 [CT2] Clinical Training \& Education}
\phantomsection
\label{navigation:II.1.3.6.11}
\addcontentsline{toc}{subsubsection}{II.1.3.6.11 [CT2] Clinical Training \& Education}

{\small Historical predecessor: R.2.135 \hyperref[element:df8bff7df9c1469da16cac6a5eb2bf76]{[CT2] Clinical Training \& Education}; Canonical ID: \texttt{\seqsplit{df8bff7df9c1469da16cac6a5eb2bf76}}. Superseded historical identity — direct correspondence · superseded.}\par
\label{element:d016f35ffbc345d99d7281e968c14ddb}
\index{[CT2] Clinical Training \& Education (II.1.3.6.11)}
{\small Capability · Canonical · current}\par

Clinical Training \& Education remains a Reference context for developing practitioners' knowledge and skills. The catalogue recognises its enterprise significance without assigning training delivery or a Feature decomposition to Harmonia. Its separation from qualification checks prevents a represented certificate or training record from being mistaken for the education activity itself.



\textbf{Overview}



Clinical Training \& Education remains a Reference context for developing practitioners' knowledge and skills. The catalogue recognises its enterprise significance without assigning training delivery or a Feature decomposition to Harmonia. Its separation from qualification checks prevents a represented certificate or training record from being mistaken for the education activity itself.



\textbf{Intent}



Clinical Training \& Education remains a Reference context for developing practitioners' knowledge and skills. The catalogue recognises its enterprise significance without assigning training delivery or a Feature decomposition to Harmonia. Its separation from qualification checks prevents a represented certificate or training record from being mistaken for the education activity itself.



\textbf{Boundary}



Information administration does not confer authority over the represented person, organisation, service, place, device or clinical activity.



\medskip\hrule\medskip

{\footnotesize \href{../../../../continuum/strategy/capability/d016f35ffbc345d99d7281e968c14ddb.md}{Canonical record: d016f35ffbc345d99d7281e968c14ddb}; full qualifications and provenance.}\par
\subsubsection*{II.1.3.6.12 [CT2] Clinical Workforce Capability Monitoring}
\phantomsection
\label{navigation:II.1.3.6.12}
\addcontentsline{toc}{subsubsection}{II.1.3.6.12 [CT2] Clinical Workforce Capability Monitoring}

{\small Historical predecessor: R.2.136 \hyperref[element:24c015676c254e2e82e8e4287c295179]{[CT2] Clinical Workforce Capability Monitoring}; Canonical ID: \texttt{\seqsplit{24c015676c254e2e82e8e4287c295179}}. Superseded historical identity — direct correspondence · superseded.}\par
\label{element:367b3e5643f44fa785f093db8bb9a6e8}
\index{[CT2] Clinical Workforce Capability Monitoring (II.1.3.6.12)}
{\small Capability · Canonical · current}\par

Monitor supplied information about the clinical capability available from the workforce without acquiring employment or clinical allocation authority.



\textbf{Overview}



Monitor supplied information about the clinical capability available from the workforce without acquiring employment or clinical allocation authority.



\textbf{Intent}



Monitor supplied information about the clinical capability available from the workforce without acquiring employment or clinical allocation authority.



\textbf{Boundary}



Information administration does not confer authority over the represented person, organisation, service, place, device or clinical activity.



\medskip\hrule\medskip

{\footnotesize \href{../../../../continuum/strategy/capability/367b3e5643f44fa785f093db8bb9a6e8.md}{Canonical record: 367b3e5643f44fa785f093db8bb9a6e8}; full qualifications and provenance.}\par
\subsubsection*{II.1.3.6.13 [CT2] Credential Management}
\phantomsection
\label{navigation:II.1.3.6.13}
\addcontentsline{toc}{subsubsection}{II.1.3.6.13 [CT2] Credential Management}

{\small Historical predecessor: R.2.138 \hyperref[element:bab3d0ab6eb748978b962ea20fb31640]{[CT2] Credential Management}; Canonical ID: \texttt{\seqsplit{bab3d0ab6eb748978b962ea20fb31640}}. Superseded historical identity — direct correspondence · superseded.}\par
\label{element:21eed4fc8c924496b16c6ddaa32f3d01}
\index{[CT2] Credential Management (II.1.3.6.13)}
{\small Capability · Proposed · current}\par

Administer definitions of clinical-service credentials and their required applicability without granting professional privilege or managing physical badges.



\textbf{Overview}



Administer definitions of clinical-service credentials and their required applicability without granting professional privilege or managing physical badges. The reviewed hierarchy assigns this enduring ability to its stated business-enabling context.



\textbf{Intent}



Administer definitions of clinical-service credentials and their required applicability without granting professional privilege or managing physical badges.



\textbf{Boundary}



Information administration and coordination preserve the authority of the responsible source, clinical actor and governed domain; they confer no additional real-world management authority.


{\footnotesize \href{../../../../continuum/strategy/capability/21eed4fc8c924496b16c6ddaa32f3d01.md}{Canonical record: 21eed4fc8c924496b16c6ddaa32f3d01}; full qualifications and provenance.}\par
\paragraph*{II.1.3.6.13.1 [CT3] Clinical Service Credential Definition Management}
\phantomsection
\label{navigation:II.1.3.6.13.1}
\addcontentsline{toc}{paragraph}{II.1.3.6.13.1 [CT3] Clinical Service Credential Definition Management}

{\small Historical predecessor: R.2.137 \hyperref[element:0115ed50867d468a95c0fafa1c5d79e2]{[CT3] Clinical Service Credential Definition Management}; Canonical ID: \texttt{\seqsplit{0115ed50867d468a95c0fafa1c5d79e2}}. Superseded historical identity — direct correspondence · superseded.}\par
\label{element:2e5c97f51e7e48d1b865ac416dd8099d}
\index{[CT3] Clinical Service Credential Definition Management (II.1.3.6.13.1)}
{\small Capability · Proposed · current}\par

Manage the definitions of credentials required for participation in a clinical service, distinct from physical badges and authentication credentials.



\textbf{Overview}



Manage the definitions of credentials required for participation in a clinical service, distinct from physical badges and authentication credentials. The reviewed hierarchy assigns this enduring ability to its stated business-enabling context.



\textbf{Intent}



Manage the definitions of credentials required for participation in a clinical service, distinct from physical badges and authentication credentials.



\textbf{Boundary}



Information administration and coordination preserve the authority of the responsible source, clinical actor and governed domain; they confer no additional real-world management authority.


{\footnotesize \href{../../../../continuum/strategy/capability/2e5c97f51e7e48d1b865ac416dd8099d.md}{Canonical record: 2e5c97f51e7e48d1b865ac416dd8099d}; full qualifications and provenance.}\par
\subsubsection*{II.1.3.6.14 [CT2] Discharge Management}
\phantomsection
\label{navigation:II.1.3.6.14}
\addcontentsline{toc}{subsubsection}{II.1.3.6.14 [CT2] Discharge Management}

{\small Historical predecessor: R.2.140 \hyperref[element:c57c0721ab0040ccb24e88e5e92e8cf8]{[CT2] Discharge Management}; Canonical ID: \texttt{\seqsplit{c57c0721ab0040ccb24e88e5e92e8cf8}}. Superseded historical identity — direct correspondence · superseded.}\par
\label{element:ec5e8d6732764aa489c2dd7630fc8fef}
\index{[CT2] Discharge Management (II.1.3.6.14)}
{\small Capability · Canonical · current}\par

Discharge Management coordinates the reported milestones needed around discharge, including medication supply, transport and home-care confirmations. It makes readiness and impediments visible across disciplines and communicates departure information when physical exit occurs. The responsibility preserves a distinction between gathering confirmations and making the clinical decision that discharge is appropriate.



\textbf{Overview}



Discharge Management coordinates the reported milestones needed around discharge, including medication supply, transport and home-care confirmations. It makes readiness and impediments visible across disciplines and communicates departure information when physical exit occurs. The responsibility preserves a distinction between gathering confirmations and making the clinical decision that discharge is appropriate.



\textbf{Intent}



Discharge Management coordinates the reported milestones needed around discharge, including medication supply, transport and home-care confirmations. It makes readiness and impediments visible across disciplines and communicates departure information when physical exit occurs. The responsibility preserves a distinction between gathering confirmations and making the clinical decision that discharge is appropriate.



\textbf{Boundary}



Information administration does not confer authority over the represented person, organisation, service, place, device or clinical activity.



\medskip\hrule\medskip

{\footnotesize \href{../../../../continuum/strategy/capability/ec5e8d6732764aa489c2dd7630fc8fef.md}{Canonical record: ec5e8d6732764aa489c2dd7630fc8fef}; full qualifications and provenance.}\par
\paragraph*{II.1.3.6.14.1 [CT3] Discharge Progression \& Coordination}
\phantomsection
\label{navigation:II.1.3.6.14.1}
\addcontentsline{toc}{paragraph}{II.1.3.6.14.1 [CT3] Discharge Progression \& Coordination}

{\small Historical predecessor: R.2.139 \hyperref[element:68344272ca6049648a3ddae20c150f37]{[CT3] Discharge Progression \& Coordination}; Canonical ID: \texttt{\seqsplit{68344272ca6049648a3ddae20c150f37}}. Superseded historical identity — direct correspondence · superseded.}\par
\label{element:d6048993146b45639bab27860f94f34f}
\index{[CT3] Discharge Progression \& Coordination (II.1.3.6.14.1)}
{\small Capability · Proposed · current}\par

Coordinate externally established discharge milestones, multidisciplinary readiness information and departure notifications.



\textbf{Overview}



Coordinate externally established discharge milestones, multidisciplinary readiness information and departure notifications. The reviewed hierarchy assigns this enduring ability to its stated business-enabling context.



\textbf{Intent}



Coordinate externally established discharge milestones, multidisciplinary readiness information and departure notifications.



\textbf{Boundary}



Information administration and coordination preserve the authority of the responsible source, clinical actor and governed domain; they confer no additional real-world management authority.


{\footnotesize \href{../../../../continuum/strategy/capability/d6048993146b45639bab27860f94f34f.md}{Canonical record: d6048993146b45639bab27860f94f34f}; full qualifications and provenance.}\par
\subparagraph*{II.1.3.6.14.1.1 [Feature] Discharge Milestone Tracking}
\phantomsection
\label{navigation:II.1.3.6.14.1.1}
\addcontentsline{toc}{subparagraph}{II.1.3.6.14.1.1 [Feature] Discharge Milestone Tracking}

\label{element:887eec1336934b13a3f21a8e3eb5f131}
\index{[Feature] Discharge Milestone Tracking (II.1.3.6.14.1.1)}
{\small Capability / Feature · Canonical · current}\par

Track clinical criteria-led discharge milestones (medications dispensed, transportation arranged, home care confirmed).



Discharge milestone tracking records confirmations contributed by the relevant clinical and operational participants. It makes outstanding requirements visible without allowing the tracking mechanism to declare clinical readiness or assume that every local milestone defines a universal discharge pathway.


{\footnotesize \href{../../../../continuum/strategy/capability/feature/887eec1336934b13a3f21a8e3eb5f131.md}{Canonical record: 887eec1336934b13a3f21a8e3eb5f131}; full qualifications and provenance.}\par
\subparagraph*{II.1.3.6.14.1.2 [Feature] Discharge Notification Dispatch}
\phantomsection
\label{navigation:II.1.3.6.14.1.2}
\addcontentsline{toc}{subparagraph}{II.1.3.6.14.1.2 [Feature] Discharge Notification Dispatch}

\label{element:5476307ca5a34590aa7025ab59ca0c10}
\index{[Feature] Discharge Notification Dispatch (II.1.3.6.14.1.2)}
{\small Capability / Feature · Canonical · current}\par

Dispatch departure notifications, updated bed states, and finalized discharge summaries to external care providers upon physical patient exit.



Departure notification communicates the physical exit and associated record and place-state context to the applicable participants. It distinguishes departure information from earlier readiness and from the clinical decision, while retaining the separate question of responsibility actually assumed elsewhere.


{\footnotesize \href{../../../../continuum/strategy/capability/feature/5476307ca5a34590aa7025ab59ca0c10.md}{Canonical record: 5476307ca5a34590aa7025ab59ca0c10}; full qualifications and provenance.}\par
\subparagraph*{II.1.3.6.14.1.3 [Feature] Discharge Readiness \& Coordination Oversight}
\phantomsection
\label{navigation:II.1.3.6.14.1.3}
\addcontentsline{toc}{subparagraph}{II.1.3.6.14.1.3 [Feature] Discharge Readiness \& Coordination Oversight}

\label{element:5499422ed3644e0bbf7645f9ae3f20cb}
\index{[Feature] Discharge Readiness \& Coordination Oversight (II.1.3.6.14.1.3)}
{\small Capability / Feature · Canonical · current}\par

Track and aggregate multidisciplinary discharge milestones, home care confirmations, and transit readiness across care teams.



Readiness oversight brings multidisciplinary, home-care and transport confirmations into one coordination view. Aggregation supports the responsible participants' decisions but does not itself confer authority to discharge or establish acceptance of continuing care by another provider.


{\footnotesize \href{../../../../continuum/strategy/capability/feature/5499422ed3644e0bbf7645f9ae3f20cb.md}{Canonical record: 5499422ed3644e0bbf7645f9ae3f20cb}; full qualifications and provenance.}\par
\subsubsection*{II.1.3.6.15 [CT2] Emergency Department Operations}
\phantomsection
\label{navigation:II.1.3.6.15}
\addcontentsline{toc}{subsubsection}{II.1.3.6.15 [CT2] Emergency Department Operations}

{\small Historical predecessor: R.2.142 \hyperref[element:8963b571e9cc40e49c5d738e31d2f97e]{[CT2] Emergency Department Operations}; Canonical ID: \texttt{\seqsplit{8963b571e9cc40e49c5d738e31d2f97e}}. Superseded historical identity — direct correspondence · superseded.}\par
\label{element:24a4b5815faa4d369bc756a6e97d3e9e}
\index{[CT2] Emergency Department Operations (II.1.3.6.15)}
{\small Capability · Canonical · current}\par

This capability maintains the operational context of emergency places, waits and length-of-stay milestones. It supports understanding of occupancy and impediments in the department while keeping capacity's emergency-service meaning intact. Its concern differs from Emergency Care's clinical presentation context, even though both can refer to the same patient encounter.



\textbf{Overview}



This capability maintains the operational context of emergency places, waits and length-of-stay milestones. It supports understanding of occupancy and impediments in the department while keeping capacity's emergency-service meaning intact. Its concern differs from Emergency Care's clinical presentation context, even though both can refer to the same patient encounter.



\textbf{Intent}



This capability maintains the operational context of emergency places, waits and length-of-stay milestones. It supports understanding of occupancy and impediments in the department while keeping capacity's emergency-service meaning intact. Its concern differs from Emergency Care's clinical presentation context, even though both can refer to the same patient encounter.



\textbf{Boundary}



Information administration does not confer authority over the represented person, organisation, service, place, device or clinical activity.



\medskip\hrule\medskip

{\footnotesize \href{../../../../continuum/strategy/capability/24a4b5815faa4d369bc756a6e97d3e9e.md}{Canonical record: 24a4b5815faa4d369bc756a6e97d3e9e}; full qualifications and provenance.}\par
\paragraph*{II.1.3.6.15.1 [CT3] Emergency Care Progression}
\phantomsection
\label{navigation:II.1.3.6.15.1}
\addcontentsline{toc}{paragraph}{II.1.3.6.15.1 [CT3] Emergency Care Progression}

\label{element:7125e459a48946d0ba075b5ecc3571b4}
\index{[CT3] Emergency Care Progression (II.1.3.6.15.1)}
{\small Capability · Canonical · current}\par

Definition unresolved in the baseline. No definition has been synthesised from the name.



\medskip\hrule\medskip

{\footnotesize \href{../../../../continuum/strategy/capability/7125e459a48946d0ba075b5ecc3571b4.md}{Canonical record: 7125e459a48946d0ba075b5ecc3571b4}; full qualifications and provenance.}\par
\subparagraph*{II.1.3.6.15.1.1 [Feature] Emergency Care Cross-Facility Correlation}
\phantomsection
\label{navigation:II.1.3.6.15.1.1}
\addcontentsline{toc}{subparagraph}{II.1.3.6.15.1.1 [Feature] Emergency Care Cross-Facility Correlation}

\label{element:6140586e2b18424d845a1f07018e6b49}
\index{[Feature] Emergency Care Cross-Facility Correlation (II.1.3.6.15.1.1)}
{\small Capability / Feature · Canonical · current}\par

Correlate incoming emergency presentations with recent regional admissions, prior discharge summaries, and known allergy profiles.



Cross-facility correlation connects an emergency presentation with relevant recent history and allergy information. It supports situational understanding across institutional boundaries without deciding the clinical significance of the combined evidence or assigning emergency treatment.


{\footnotesize \href{../../../../continuum/strategy/capability/feature/6140586e2b18424d845a1f07018e6b49.md}{Canonical record: 6140586e2b18424d845a1f07018e6b49}; full qualifications and provenance.}\par
\subparagraph*{II.1.3.6.15.1.2 [Feature] Emergency Department Arrival \& Triage Ingestion}
\phantomsection
\label{navigation:II.1.3.6.15.1.2}
\addcontentsline{toc}{subparagraph}{II.1.3.6.15.1.2 [Feature] Emergency Department Arrival \& Triage Ingestion}

\label{element:1eb4053512164f6f976a9767da9ba7a7}
\index{[Feature] Emergency Department Arrival \& Triage Ingestion (II.1.3.6.15.1.2)}
{\small Capability / Feature · Canonical · current}\par

Ingest emergency presentation, triage category, and chief complaint records.



Emergency presentation ingestion retains the arrival, triage category and presenting complaint reported by the emergency context. It makes that event available for regional correlation while preserving that the clinical triage judgement was established by the responsible source.


{\footnotesize \href{../../../../continuum/strategy/capability/feature/1eb4053512164f6f976a9767da9ba7a7.md}{Canonical record: 1eb4053512164f6f976a9767da9ba7a7}; full qualifications and provenance.}\par
\paragraph*{II.1.3.6.15.2 [CT3] Emergency Department Operational Coordination}
\phantomsection
\label{navigation:II.1.3.6.15.2}
\addcontentsline{toc}{paragraph}{II.1.3.6.15.2 [CT3] Emergency Department Operational Coordination}

{\small Historical predecessor: R.2.141 \hyperref[element:15b6d2518e60444098b7c546ca31c49b]{[CT3] Emergency Department Operational Coordination}; Canonical ID: \texttt{\seqsplit{15b6d2518e60444098b7c546ca31c49b}}. Superseded historical identity — direct correspondence · superseded.}\par
\label{element:8f0b0f3d9e8d44a7827b0835870b4a4f}
\index{[CT3] Emergency Department Operational Coordination (II.1.3.6.15.2)}
{\small Capability · Proposed · current}\par

Coordinate emergency-department bed/bay occupancy, operational waits, length-of-stay milestones and local capacity information.



\textbf{Overview}



Coordinate emergency-department bed/bay occupancy, operational waits, length-of-stay milestones and local capacity information. The reviewed hierarchy assigns this enduring ability to its stated business-enabling context.



\textbf{Intent}



Coordinate emergency-department bed/bay occupancy, operational waits, length-of-stay milestones and local capacity information.



\textbf{Boundary}



Information administration and coordination preserve the authority of the responsible source, clinical actor and governed domain; they confer no additional real-world management authority.


{\footnotesize \href{../../../../continuum/strategy/capability/8f0b0f3d9e8d44a7827b0835870b4a4f.md}{Canonical record: 8f0b0f3d9e8d44a7827b0835870b4a4f}; full qualifications and provenance.}\par
\subparagraph*{II.1.3.6.15.2.1 [Feature] ED Bed \& Bay Tracking}
\phantomsection
\label{navigation:II.1.3.6.15.2.1}
\addcontentsline{toc}{subparagraph}{II.1.3.6.15.2.1 [Feature] ED Bed \& Bay Tracking}

\label{element:7c0251128fc942d88bf25141421fd909}
\index{[Feature] ED Bed \& Bay Tracking (II.1.3.6.15.2.1)}
{\small Capability / Feature · Canonical · current}\par

Track occupancy and status of emergency cubicles, resuscitation bays, and short-stay beds.



Emergency-place tracking distinguishes occupancy and status across cubicles, resuscitation bays and short-stay beds. It provides place context for departmental operations while leaving the service's usable capacity and clinical triage decisions separate from an inventory of spaces.


{\footnotesize \href{../../../../continuum/strategy/capability/feature/7c0251128fc942d88bf25141421fd909.md}{Canonical record: 7c0251128fc942d88bf25141421fd909}; full qualifications and provenance.}\par
\subparagraph*{II.1.3.6.15.2.2 [Feature] ED Length of Stay Tracking}
\phantomsection
\label{navigation:II.1.3.6.15.2.2}
\addcontentsline{toc}{subparagraph}{II.1.3.6.15.2.2 [Feature] ED Length of Stay Tracking}

\label{element:e544f3d76db84dfaad57880d13a873c4}
\index{[Feature] ED Length of Stay Tracking (II.1.3.6.15.2.2)}
{\small Capability / Feature · Canonical · current}\par

Monitor operational wait times and length-of-stay milestones, generating operational escalation notifications.



Length-of-stay tracking makes waits and operational timing thresholds visible for escalation. The timing information supports management of impediments without determining clinical urgency or making an elapsed-time alert an independent assurance finding.


{\footnotesize \href{../../../../continuum/strategy/capability/feature/e544f3d76db84dfaad57880d13a873c4.md}{Canonical record: e544f3d76db84dfaad57880d13a873c4}; full qualifications and provenance.}\par
\subparagraph*{II.1.3.6.15.2.3 [Feature] Emergency Treatment Capacity Management}
\phantomsection
\label{navigation:II.1.3.6.15.2.3}
\addcontentsline{toc}{subparagraph}{II.1.3.6.15.2.3 [Feature] Emergency Treatment Capacity Management}

\label{element:502876430cf747f4b3bb102fb69557a5}
\index{[Feature] Emergency Treatment Capacity Management (II.1.3.6.15.2.3)}
{\small Capability / Feature · Proposed · current}\par

Represent current emergency treatment capacity under space, staffing and acuity constraints.



\textbf{Overview}



Bay state and length of stay contribute evidence but do not alone establish the ability to accept another care demand.



\textbf{Boundary}



Local responsibility retains its capacity meaning; Service Capacity Management consumes it without acquiring ownership or imposing one common capacity model.


{\footnotesize \href{../../../../continuum/strategy/capability/feature/502876430cf747f4b3bb102fb69557a5.md}{Canonical record: 502876430cf747f4b3bb102fb69557a5}; full qualifications and provenance.}\par
\subsubsection*{II.1.3.6.16 [CT2] Mobile Staff Management}
\phantomsection
\label{navigation:II.1.3.6.16}
\addcontentsline{toc}{subsubsection}{II.1.3.6.16 [CT2] Mobile Staff Management}

{\small Historical predecessor: R.2.144 \hyperref[element:32118856636f400a95fc28daacea61c8]{[CT2] Mobile Staff Management}; Canonical ID: \texttt{\seqsplit{32118856636f400a95fc28daacea61c8}}. Superseded historical identity — direct correspondence · superseded.}\par
\label{element:79a4a48fddf74c69acf7e16404d06ea5}
\index{[CT2] Mobile Staff Management (II.1.3.6.16)}
{\small Capability · Canonical · current}\par

Mobile Staff Management integrates with external business mobile-work mechanisms and receives visit progression and safety milestones. Its established dispatch Feature exchanges assignment, dispatch and status information with the relevant worker context. This is an information integration boundary around mobile activity, rather than a platform claim to decide the clinician's work.



\textbf{Overview}



Mobile Staff Management integrates with external business mobile-work mechanisms and receives visit progression and safety milestones. Its established dispatch Feature exchanges assignment, dispatch and status information with the relevant worker context. This is an information integration boundary around mobile activity, rather than a platform claim to decide the clinician's work.



\textbf{Intent}



Mobile Staff Management integrates with external business mobile-work mechanisms and receives visit progression and safety milestones. Its established dispatch Feature exchanges assignment, dispatch and status information with the relevant worker context. This is an information integration boundary around mobile activity, rather than a platform claim to decide the clinician's work.



\textbf{Boundary}



Information administration does not confer authority over the represented person, organisation, service, place, device or clinical activity.



\medskip\hrule\medskip

{\footnotesize \href{../../../../continuum/strategy/capability/79a4a48fddf74c69acf7e16404d06ea5.md}{Canonical record: 79a4a48fddf74c69acf7e16404d06ea5}; full qualifications and provenance.}\par
\paragraph*{II.1.3.6.16.1 [CT3] Mobile Worker Dispatch \& Safety Coordination}
\phantomsection
\label{navigation:II.1.3.6.16.1}
\addcontentsline{toc}{paragraph}{II.1.3.6.16.1 [CT3] Mobile Worker Dispatch \& Safety Coordination}

{\small Historical predecessor: R.2.143 \hyperref[element:5aafd1f4bff046b2b532c20dfb191473]{[CT3] Mobile Worker Dispatch \& Safety Coordination}; Canonical ID: \texttt{\seqsplit{5aafd1f4bff046b2b532c20dfb191473}}. Superseded historical identity — direct correspondence · superseded.}\par
\label{element:d72164668dfe448bbebe74e9e1a75233}
\index{[CT3] Mobile Worker Dispatch \& Safety Coordination (II.1.3.6.16.1)}
{\small Capability · Proposed · current}\par

Exchange externally established mobile work assignments, visit progression and safety check-in information with qualified worker context.



\textbf{Overview}



Exchange externally established mobile work assignments, visit progression and safety check-in information with qualified worker context. The reviewed hierarchy assigns this enduring ability to its stated business-enabling context.



\textbf{Intent}



Exchange externally established mobile work assignments, visit progression and safety check-in information with qualified worker context.



\textbf{Boundary}



Information administration and coordination preserve the authority of the responsible source, clinical actor and governed domain; they confer no additional real-world management authority.


{\footnotesize \href{../../../../continuum/strategy/capability/d72164668dfe448bbebe74e9e1a75233.md}{Canonical record: d72164668dfe448bbebe74e9e1a75233}; full qualifications and provenance.}\par
\subparagraph*{II.1.3.6.16.1.1 [Feature] Mobile Service Capacity Management}
\phantomsection
\label{navigation:II.1.3.6.16.1.1}
\addcontentsline{toc}{subparagraph}{II.1.3.6.16.1.1 [Feature] Mobile Service Capacity Management}

\label{element:abdcb95b3725405f878f199b07779d19}
\index{[Feature] Mobile Service Capacity Management (II.1.3.6.16.1.1)}
{\small Capability / Feature · Proposed · current}\par

Represent mobile service capacity with travel, coverage, work commitment and availability context.



\textbf{Overview}



Mobile availability cannot be interpreted as a simple on-duty count because distance, assignments and suitable skills affect deliverability.



\textbf{Boundary}



Local responsibility retains its capacity meaning; Service Capacity Management consumes it without acquiring ownership or imposing one common capacity model.


{\footnotesize \href{../../../../continuum/strategy/capability/feature/abdcb95b3725405f878f199b07779d19.md}{Canonical record: abdcb95b3725405f878f199b07779d19}; full qualifications and provenance.}\par
\subparagraph*{II.1.3.6.16.1.2 [Feature] Mobile Visit Status Tracking}
\phantomsection
\label{navigation:II.1.3.6.16.1.2}
\addcontentsline{toc}{subparagraph}{II.1.3.6.16.1.2 [Feature] Mobile Visit Status Tracking}

\label{element:3cd49cb00159437ea25048c1158172c1}
\index{[Feature] Mobile Visit Status Tracking (II.1.3.6.16.1.2)}
{\small Capability / Feature · Canonical · current}\par

Ingest arrival, visit progression, and safety check-in milestones from roving staff.



Mobile visit tracking receives arrival, progression and safety checkpoints from roving workers. It makes reported activity visible in the wider operational context without originating the clinical visit, deciding its clinical content or assuming that a checkpoint proves the intended clinical outcome.


{\footnotesize \href{../../../../continuum/strategy/capability/feature/3cd49cb00159437ea25048c1158172c1.md}{Canonical record: 3cd49cb00159437ea25048c1158172c1}; full qualifications and provenance.}\par
\subparagraph*{II.1.3.6.16.1.3 [Feature] Mobile Worker Task Dispatch}
\phantomsection
\label{navigation:II.1.3.6.16.1.3}
\addcontentsline{toc}{subparagraph}{II.1.3.6.16.1.3 [Feature] Mobile Worker Task Dispatch}

\label{element:426363f2f02a4625a4feab89a2adc755}
\index{[Feature] Mobile Worker Task Dispatch (II.1.3.6.16.1.3)}
{\small Capability / Feature · Canonical · current}\par

Dispatch community care tasks and urgent visits to roving clinicians' mobile devices.



Mobile dispatch exchanges externally established assignment, dispatch and progression information with business mobile-work mechanisms. The retained name does not grant Harmonia authority to allocate clinical work, manage handover or own EMR worklists, and external task terminology does not redefine Synthetic Task.



\textbf{Boundary}



The proposed replacement name remains unconfirmed. This entry preserves the existing identity and the human-adjudicated integration boundary.


{\footnotesize \href{../../../../continuum/strategy/capability/feature/426363f2f02a4625a4feab89a2adc755.md}{Canonical record: 426363f2f02a4625a4feab89a2adc755}; full qualifications and provenance.}\par
\subsubsection*{II.1.3.6.17 [CT2] On-Call Management}
\phantomsection
\label{navigation:II.1.3.6.17}
\addcontentsline{toc}{subsubsection}{II.1.3.6.17 [CT2] On-Call Management}

{\small Historical predecessor: R.2.146 \hyperref[element:c8e6e66e8718493c9093f37d129f9ff2]{[CT2] On-Call Management}; Canonical ID: \texttt{\seqsplit{c8e6e66e8718493c9093f37d129f9ff2}}. Superseded historical identity — direct correspondence · superseded.}\par
\label{element:b16cef8f2bfd4b92948ec75cbcfa7b9a}
\index{[CT2] On-Call Management (II.1.3.6.17)}
{\small Capability · Canonical · current}\par

On-Call Management resolves the active provider for a specialty, service or hospital at the relevant time. It turns changing coverage information into usable routing context so that operational communication can reach the intended capacity. This differs from authoring the roster or deciding who is professionally qualified for the requested activity.



\textbf{Overview}



On-Call Management resolves the active provider for a specialty, service or hospital at the relevant time. It turns changing coverage information into usable routing context so that operational communication can reach the intended capacity. This differs from authoring the roster or deciding who is professionally qualified for the requested activity.



\textbf{Intent}



On-Call Management resolves the active provider for a specialty, service or hospital at the relevant time. It turns changing coverage information into usable routing context so that operational communication can reach the intended capacity. This differs from authoring the roster or deciding who is professionally qualified for the requested activity.



\textbf{Boundary}



Information administration does not confer authority over the represented person, organisation, service, place, device or clinical activity.



\medskip\hrule\medskip

{\footnotesize \href{../../../../continuum/strategy/capability/b16cef8f2bfd4b92948ec75cbcfa7b9a.md}{Canonical record: b16cef8f2bfd4b92948ec75cbcfa7b9a}; full qualifications and provenance.}\par
\paragraph*{II.1.3.6.17.1 [CT3] On-Call Availability \& Contact Resolution}
\phantomsection
\label{navigation:II.1.3.6.17.1}
\addcontentsline{toc}{paragraph}{II.1.3.6.17.1 [CT3] On-Call Availability \& Contact Resolution}

{\small Historical predecessor: R.2.145 \hyperref[element:a03075a46f75497f8c6ba95d9c9e27c4]{[CT3] On-Call Availability \& Contact Resolution}; Canonical ID: \texttt{\seqsplit{a03075a46f75497f8c6ba95d9c9e27c4}}. Superseded historical identity — direct correspondence · superseded.}\par
\label{element:677a816c42414ea694791b4ee5a14285}
\index{[CT3] On-Call Availability \& Contact Resolution (II.1.3.6.17.1)}
{\small Capability · Proposed · current}\par

Resolve currently available on-call providers and their contact paths for the applicable specialty, service and time.



\textbf{Overview}



Resolve currently available on-call providers and their contact paths for the applicable specialty, service and time. The reviewed hierarchy assigns this enduring ability to its stated business-enabling context.



\textbf{Intent}



Resolve currently available on-call providers and their contact paths for the applicable specialty, service and time.



\textbf{Boundary}



Information administration and coordination preserve the authority of the responsible source, clinical actor and governed domain; they confer no additional real-world management authority.


{\footnotesize \href{../../../../continuum/strategy/capability/677a816c42414ea694791b4ee5a14285.md}{Canonical record: 677a816c42414ea694791b4ee5a14285}; full qualifications and provenance.}\par
\subparagraph*{II.1.3.6.17.1.1 [Feature] Active On-Call Provider Resolution}
\phantomsection
\label{navigation:II.1.3.6.17.1.1}
\addcontentsline{toc}{subparagraph}{II.1.3.6.17.1.1 [Feature] Active On-Call Provider Resolution}

\label{element:b588f3cc869d460a9561657352fb189a}
\index{[Feature] Active On-Call Provider Resolution (II.1.3.6.17.1.1)}
{\small Capability / Feature · Canonical · current}\par

Resolve currently active on-call specialists, registrars, and operational managers for a given specialty and hospital at any point in time.



Active on-call resolution identifies the current provider in the required specialty and hospital context. It supports correct routing at the relevant time without authoring the roster or treating current coverage as proof of qualification, privilege or responsibility acceptance.


{\footnotesize \href{../../../../continuum/strategy/capability/feature/b588f3cc869d460a9561657352fb189a.md}{Canonical record: b588f3cc869d460a9561657352fb189a}; full qualifications and provenance.}\par
\subsubsection*{II.1.3.6.18 [CT2] Operational Workforce Coordination [superseded]}
\phantomsection
\label{navigation:II.1.3.6.18}
\addcontentsline{toc}{subsubsection}{II.1.3.6.18 [CT2] Operational Workforce Coordination [superseded]}

\label{element:66e72222f2064a23adbdb0311092918b}
\index{[CT2] Operational Workforce Coordination (II.1.3.6.18)}
{\small Capability · Proposed · superseded}\par
\href{../../../../continuum/strategy/capability/66e72222f2064a23adbdb0311092918b.md}{Historical canonical definition and provenance}.\par
\subsubsection*{II.1.3.6.19 [CT2] Outpatient Operations}
\phantomsection
\label{navigation:II.1.3.6.19}
\addcontentsline{toc}{subsubsection}{II.1.3.6.19 [CT2] Outpatient Operations}

{\small Historical predecessor: R.2.148 \hyperref[element:3db41bc06ade4f9088eec7b8aedd7448]{[CT2] Outpatient Operations}; Canonical ID: \texttt{\seqsplit{3db41bc06ade4f9088eec7b8aedd7448}}. Superseded historical identity — direct correspondence · superseded.}\par
\label{element:21399d52b9624ddfbeb786f1d72e40c3}
\index{[CT2] Outpatient Operations (II.1.3.6.19)}
{\small Capability · Canonical · current}\par

Outpatient Operations integrates arrival events and follow-up requirements generated as clinics progress. It connects a person's operational presence with the relevant service context and makes post-clinic booking or diagnostic needs available for subsequent activity. The responsibility retains outpatient-specific capacity semantics without requiring a separate arrival Function where approved reusable composition suffices.



\textbf{Overview}



Outpatient Operations integrates arrival events and follow-up requirements generated as clinics progress. It connects a person's operational presence with the relevant service context and makes post-clinic booking or diagnostic needs available for subsequent activity. The responsibility retains outpatient-specific capacity semantics without requiring a separate arrival Function where approved reusable composition suffices.



\textbf{Intent}



Outpatient Operations integrates arrival events and follow-up requirements generated as clinics progress. It connects a person's operational presence with the relevant service context and makes post-clinic booking or diagnostic needs available for subsequent activity. The responsibility retains outpatient-specific capacity semantics without requiring a separate arrival Function where approved reusable composition suffices.



\textbf{Boundary}



Information administration does not confer authority over the represented person, organisation, service, place, device or clinical activity.



\medskip\hrule\medskip

{\footnotesize \href{../../../../continuum/strategy/capability/21399d52b9624ddfbeb786f1d72e40c3.md}{Canonical record: 21399d52b9624ddfbeb786f1d72e40c3}; full qualifications and provenance.}\par
\paragraph*{II.1.3.6.19.1 [CT3] Outpatient Flow Coordination}
\phantomsection
\label{navigation:II.1.3.6.19.1}
\addcontentsline{toc}{paragraph}{II.1.3.6.19.1 [CT3] Outpatient Flow Coordination}

{\small Historical predecessor: R.2.147 \hyperref[element:e6c59c7b72fe44169a5fdfd08a5e91c4]{[CT3] Outpatient Flow Coordination}; Canonical ID: \texttt{\seqsplit{e6c59c7b72fe44169a5fdfd08a5e91c4}}. Superseded historical identity — direct correspondence · superseded.}\par
\label{element:799ea278af014168b9b534e5fdadc6d1}
\index{[CT3] Outpatient Flow Coordination (II.1.3.6.19.1)}
{\small Capability · Proposed · current}\par

Coordinate outpatient arrival, clinic flow, follow-up requirements and context-specific outpatient capacity information.



\textbf{Overview}



Coordinate outpatient arrival, clinic flow, follow-up requirements and context-specific outpatient capacity information. The reviewed hierarchy assigns this enduring ability to its stated business-enabling context.



\textbf{Intent}



Coordinate outpatient arrival, clinic flow, follow-up requirements and context-specific outpatient capacity information.



\textbf{Boundary}



Information administration and coordination preserve the authority of the responsible source, clinical actor and governed domain; they confer no additional real-world management authority.


{\footnotesize \href{../../../../continuum/strategy/capability/799ea278af014168b9b534e5fdadc6d1.md}{Canonical record: 799ea278af014168b9b534e5fdadc6d1}; full qualifications and provenance.}\par
\subparagraph*{II.1.3.6.19.1.1 [Feature] Outpatient Appointment Capacity Management}
\phantomsection
\label{navigation:II.1.3.6.19.1.1}
\addcontentsline{toc}{subparagraph}{II.1.3.6.19.1.1 [Feature] Outpatient Appointment Capacity Management}

\label{element:3c4f8b0fc23044cbbac32602e1272680}
\index{[Feature] Outpatient Appointment Capacity Management (II.1.3.6.19.1.1)}
{\small Capability / Feature · Proposed · current}\par

Represent usable outpatient session capacity from booking, staffing and local flow information.



\textbf{Overview}



The projection distinguishes scheduled slots from currently usable care capacity without taking source booking authority.



\textbf{Boundary}



Local responsibility retains its capacity meaning; Service Capacity Management consumes it without acquiring ownership or imposing one common capacity model.


{\footnotesize \href{../../../../continuum/strategy/capability/feature/3c4f8b0fc23044cbbac32602e1272680.md}{Canonical record: 3c4f8b0fc23044cbbac32602e1272680}; full qualifications and provenance.}\par
\subparagraph*{II.1.3.6.19.1.2 [Feature] Outpatient Arrival Registration}
\phantomsection
\label{navigation:II.1.3.6.19.1.2}
\addcontentsline{toc}{subparagraph}{II.1.3.6.19.1.2 [Feature] Outpatient Arrival Registration}

\label{element:4ea71fb4950848cfb4f6d916190edad0}
\index{[Feature] Outpatient Arrival Registration (II.1.3.6.19.1.2)}
{\small Capability / Feature · Canonical · current}\par

Ingest self-service check-in and clinic arrival events, notifying treating clinicians.



Outpatient arrival registration connects a check-in event with the treating context and required notification. Approved reusable composition supplies this integration responsibility; it does not require a specialised arrival Function or prove that a consultation has begun or finished.


{\footnotesize \href{../../../../continuum/strategy/capability/feature/4ea71fb4950848cfb4f6d916190edad0.md}{Canonical record: 4ea71fb4950848cfb4f6d916190edad0}; full qualifications and provenance.}\par
\subparagraph*{II.1.3.6.19.1.3 [Feature] Post-Clinic Order Coordination}
\phantomsection
\label{navigation:II.1.3.6.19.1.3}
\addcontentsline{toc}{subparagraph}{II.1.3.6.19.1.3 [Feature] Post-Clinic Order Coordination}

\label{element:439b3530b7c9443ab3e48d2cc50a3b0c}
\index{[Feature] Post-Clinic Order Coordination (II.1.3.6.19.1.3)}
{\small Capability / Feature · Canonical · current}\par

Capture follow-up booking and diagnostic order requirements generated at clinic completion.



Post-clinic coordination captures the follow-up bookings and diagnostic requirements established at clinic completion. It connects those requirements with subsequent operational activity while preserving that capture is not the clinical decision, a completed booking or diagnostic fulfilment.


{\footnotesize \href{../../../../continuum/strategy/capability/feature/439b3530b7c9443ab3e48d2cc50a3b0c.md}{Canonical record: 439b3530b7c9443ab3e48d2cc50a3b0c}; full qualifications and provenance.}\par
\subsubsection*{II.1.3.6.20 [CT2] Patient Transport}
\phantomsection
\label{navigation:II.1.3.6.20}
\addcontentsline{toc}{subsubsection}{II.1.3.6.20 [CT2] Patient Transport}

{\small Historical predecessor: R.2.150 \hyperref[element:e9c1a897884f4d758f5654b824364030]{[CT2] Patient Transport}; Canonical ID: \texttt{\seqsplit{e9c1a897884f4d758f5654b824364030}}. Superseded historical identity — direct correspondence · superseded.}\par
\label{element:6041d7c0cc5d4519aaea584bb4e6d365}
\index{[CT2] Patient Transport (II.1.3.6.20)}
{\small Capability · Canonical · current}\par

Patient Transport coordinates the operational progression of movement requests, including the patient's applicable mobility constraints. Departure, transit and arrival checkpoints make the physical journey visible to participants. The responsibility supports internal and inter-facility movement without treating arrival as acceptance of clinical responsibility at the destination.



\textbf{Overview}



Patient Transport coordinates the operational progression of movement requests, including the patient's applicable mobility constraints. Departure, transit and arrival checkpoints make the physical journey visible to participants. The responsibility supports internal and inter-facility movement without treating arrival as acceptance of clinical responsibility at the destination.



\textbf{Intent}



Patient Transport coordinates the operational progression of movement requests, including the patient's applicable mobility constraints. Departure, transit and arrival checkpoints make the physical journey visible to participants. The responsibility supports internal and inter-facility movement without treating arrival as acceptance of clinical responsibility at the destination.



\textbf{Boundary}



Information administration does not confer authority over the represented person, organisation, service, place, device or clinical activity.



\medskip\hrule\medskip

{\footnotesize \href{../../../../continuum/strategy/capability/6041d7c0cc5d4519aaea584bb4e6d365.md}{Canonical record: 6041d7c0cc5d4519aaea584bb4e6d365}; full qualifications and provenance.}\par
\paragraph*{II.1.3.6.20.1 [CT3] Patient Transfer Coordination}
\phantomsection
\label{navigation:II.1.3.6.20.1}
\addcontentsline{toc}{paragraph}{II.1.3.6.20.1 [CT3] Patient Transfer Coordination}

{\small Historical predecessor: R.2.149 \hyperref[element:b5a76d4877424b2ea9b01fb5fa22772c]{[CT3] Patient Transfer Coordination}; Canonical ID: \texttt{\seqsplit{b5a76d4877424b2ea9b01fb5fa22772c}}. Superseded historical identity — direct correspondence · superseded.}\par
\label{element:c7c7aa8e34df4568ae4b91cdde380a11}
\index{[CT3] Patient Transfer Coordination (II.1.3.6.20.1)}
{\small Capability · Proposed · current}\par

Coordinate patient transport requests, dispatch and observed departure, transit and arrival milestones.



\textbf{Overview}



Coordinate patient transport requests, dispatch and observed departure, transit and arrival milestones. The reviewed hierarchy assigns this enduring ability to its stated business-enabling context.



\textbf{Intent}



Coordinate patient transport requests, dispatch and observed departure, transit and arrival milestones.



\textbf{Boundary}



Information administration and coordination preserve the authority of the responsible source, clinical actor and governed domain; they confer no additional real-world management authority.


{\footnotesize \href{../../../../continuum/strategy/capability/c7c7aa8e34df4568ae4b91cdde380a11.md}{Canonical record: c7c7aa8e34df4568ae4b91cdde380a11}; full qualifications and provenance.}\par
\subparagraph*{II.1.3.6.20.1.1 [Feature] Transport Dispatch \& Progress Tracking}
\phantomsection
\label{navigation:II.1.3.6.20.1.1}
\addcontentsline{toc}{subparagraph}{II.1.3.6.20.1.1 [Feature] Transport Dispatch \& Progress Tracking}

\label{element:d8ef9994298f466bafce919c5e6cb701}
\index{[Feature] Transport Dispatch \& Progress Tracking (II.1.3.6.20.1.1)}
{\small Capability / Feature · Canonical · current}\par

Dispatch transport orders, tracking departure, transit, and destination arrival checkpoints.



Transport progression connects dispatch, departure, transit and arrival evidence for a physical patient journey. The checkpoints make movement intelligible while preserving that reaching a destination does not establish clinical handover or assumption of responsibility there.


{\footnotesize \href{../../../../continuum/strategy/capability/feature/d8ef9994298f466bafce919c5e6cb701.md}{Canonical record: d8ef9994298f466bafce919c5e6cb701}; full qualifications and provenance.}\par
\subparagraph*{II.1.3.6.20.1.2 [Feature] Transport Request Ingestion}
\phantomsection
\label{navigation:II.1.3.6.20.1.2}
\addcontentsline{toc}{subparagraph}{II.1.3.6.20.1.2 [Feature] Transport Request Ingestion}

\label{element:7104472d065840eb91ae4cb7a3a0cb43}
\index{[Feature] Transport Request Ingestion (II.1.3.6.20.1.2)}
{\small Capability / Feature · Canonical · current}\par

Ingest internal (portering) and external (ambulance, inter-hospital) patient transfer requests with clinical mobility requirements.



Transport ingestion establishes the requested movement and applicable clinical mobility constraints. It gives operational coordination the context needed to arrange transit without making the request itself a departure, arrival or accepted care transition.


{\footnotesize \href{../../../../continuum/strategy/capability/feature/7104472d065840eb91ae4cb7a3a0cb43.md}{Canonical record: 7104472d065840eb91ae4cb7a3a0cb43}; full qualifications and provenance.}\par
\subsubsection*{II.1.3.6.21 [CT2] Service Capacity Management}
\phantomsection
\label{navigation:II.1.3.6.21}
\addcontentsline{toc}{subsubsection}{II.1.3.6.21 [CT2] Service Capacity Management}

{\small Historical predecessor: R.2.152 \hyperref[element:ac205f374c9d41c3b6d18a6b7c356152]{[CT2] Service Capacity Management}; Canonical ID: \texttt{\seqsplit{ac205f374c9d41c3b6d18a6b7c356152}}. Superseded historical identity — direct correspondence · superseded.}\par
\label{element:d86105bf9c86445394088f65c895d37f}
\index{[CT2] Service Capacity Management (II.1.3.6.21)}
{\small Capability · Canonical · current}\par

Service Capacity Management provides a higher-order view by aggregating capacity information across services and facilities. It can consider occupancy, waits and diversion while preserving the service-specific meaning of the underlying capacity. Its purpose is not to erase differences between clinic, ward, theatre, emergency, outpatient and mobile-service operation.



\textbf{Overview}



Service Capacity Management provides a higher-order view by aggregating capacity information across services and facilities. It can consider occupancy, waits and diversion while preserving the service-specific meaning of the underlying capacity. Its purpose is not to erase differences between clinic, ward, theatre, emergency, outpatient and mobile-service operation.



\textbf{Intent}



Service Capacity Management provides a higher-order view by aggregating capacity information across services and facilities. It can consider occupancy, waits and diversion while preserving the service-specific meaning of the underlying capacity. Its purpose is not to erase differences between clinic, ward, theatre, emergency, outpatient and mobile-service operation.



\textbf{Boundary}



Information administration does not confer authority over the represented person, organisation, service, place, device or clinical activity.



\medskip\hrule\medskip

{\footnotesize \href{../../../../continuum/strategy/capability/d86105bf9c86445394088f65c895d37f.md}{Canonical record: d86105bf9c86445394088f65c895d37f}; full qualifications and provenance.}\par
\paragraph*{II.1.3.6.21.1 [CT3] Service Demand \& Capacity Coordination}
\phantomsection
\label{navigation:II.1.3.6.21.1}
\addcontentsline{toc}{paragraph}{II.1.3.6.21.1 [CT3] Service Demand \& Capacity Coordination}

{\small Historical predecessor: R.2.151 \hyperref[element:b57081e650f74ad488549fdaa5305bcb]{[CT3] Service Demand \& Capacity Coordination}; Canonical ID: \texttt{\seqsplit{b57081e650f74ad488549fdaa5305bcb}}. Superseded historical identity — direct correspondence · superseded.}\par
\label{element:93b5548ce8d04fcd999b97536c1eebf9}
\index{[CT3] Service Demand \& Capacity Coordination (II.1.3.6.21.1)}
{\small Capability · Proposed · current}\par

Coordinate demand and capacity information across services while retaining the originating responsibilities of context-specific capacity providers.



\textbf{Overview}



Coordinate demand and capacity information across services while retaining the originating responsibilities of context-specific capacity providers. The reviewed hierarchy assigns this enduring ability to its stated business-enabling context.



\textbf{Intent}



Coordinate demand and capacity information across services while retaining the originating responsibilities of context-specific capacity providers.



\textbf{Boundary}



Information administration and coordination preserve the authority of the responsible source, clinical actor and governed domain; they confer no additional real-world management authority.


{\footnotesize \href{../../../../continuum/strategy/capability/93b5548ce8d04fcd999b97536c1eebf9.md}{Canonical record: 93b5548ce8d04fcd999b97536c1eebf9}; full qualifications and provenance.}\par
\subparagraph*{II.1.3.6.21.1.1 [Feature] Operational Capacity Metric Aggregation}
\phantomsection
\label{navigation:II.1.3.6.21.1.1}
\addcontentsline{toc}{subparagraph}{II.1.3.6.21.1.1 [Feature] Operational Capacity Metric Aggregation}

\label{element:1fe9199a49744eeaa5d063d7a84544ea}
\index{[Feature] Operational Capacity Metric Aggregation (II.1.3.6.21.1.1)}
{\small Capability / Feature · Canonical · current}\par

Aggregate operational capacity telemetry (occupancy rates, wait times, diversion status) across facilities.



Capacity aggregation combines occupancy, waiting and diversion context across facilities. The higher-order view must retain the originating service-specific capacity meaning; a shared metric does not transfer clinic, ward or theatre capacity responsibility to the aggregator.


{\footnotesize \href{../../../../continuum/strategy/capability/feature/1fe9199a49744eeaa5d063d7a84544ea.md}{Canonical record: 1fe9199a49744eeaa5d063d7a84544ea}; full qualifications and provenance.}\par
\subsubsection*{II.1.3.6.22 [CT2] Service Establishment Support}
\phantomsection
\label{navigation:II.1.3.6.22}
\addcontentsline{toc}{subsubsection}{II.1.3.6.22 [CT2] Service Establishment Support}

{\small Historical predecessor: R.2.153 \hyperref[element:8a63b08f34fa42c5bda2a3e12b9ecd55]{[CT2] Service Establishment Support}; Canonical ID: \texttt{\seqsplit{8a63b08f34fa42c5bda2a3e12b9ecd55}}. Superseded historical identity — direct correspondence · superseded.}\par
\label{element:20e67a41e7f44398b0f0cf3d379d8b9c}
\index{[CT2] Service Establishment Support (II.1.3.6.22)}
{\small Capability · Proposed · current}\par

Support the information needed to establish a health service in its intended operating context.



\textbf{Overview}



Support the information needed to establish a health service in its intended operating context. The reviewed hierarchy assigns this enduring ability to its stated business-enabling context.



\textbf{Intent}



Support the information needed to establish a health service in its intended operating context.



\textbf{Boundary}



Information administration and coordination preserve the authority of the responsible source, clinical actor and governed domain; they confer no additional real-world management authority.


{\footnotesize \href{../../../../continuum/strategy/capability/20e67a41e7f44398b0f0cf3d379d8b9c.md}{Canonical record: 20e67a41e7f44398b0f0cf3d379d8b9c}; full qualifications and provenance.}\par
\subsubsection*{II.1.3.6.23 [CT2] Theatre Operations}
\phantomsection
\label{navigation:II.1.3.6.23}
\addcontentsline{toc}{subsubsection}{II.1.3.6.23 [CT2] Theatre Operations}

{\small Historical predecessor: R.2.155 \hyperref[element:9ac72fea821c45e2b18632ae9a9867dd]{[CT2] Theatre Operations}; Canonical ID: \texttt{\seqsplit{9ac72fea821c45e2b18632ae9a9867dd}}. Superseded historical identity — direct correspondence · superseded.}\par
\label{element:4bf7a8260e1444fdac44e97c3557b533}
\index{[CT2] Theatre Operations (II.1.3.6.23)}
{\small Capability · Canonical · current}\par

Theatre Operations integrates surgical case progression and the operational consequences relevant to portering, sterilisation and recovery participants. It gives reported milestones an operational coordination context without originating the surgical acts they describe. Theatre-specific capacity remains a business responsibility even when shared information machinery supports several service contexts.



\textbf{Overview}



Theatre Operations integrates surgical case progression and the operational consequences relevant to portering, sterilisation and recovery participants. It gives reported milestones an operational coordination context without originating the surgical acts they describe. Theatre-specific capacity remains a business responsibility even when shared information machinery supports several service contexts.



\textbf{Intent}



Theatre Operations integrates surgical case progression and the operational consequences relevant to portering, sterilisation and recovery participants. It gives reported milestones an operational coordination context without originating the surgical acts they describe. Theatre-specific capacity remains a business responsibility even when shared information machinery supports several service contexts.



\textbf{Boundary}



Information administration does not confer authority over the represented person, organisation, service, place, device or clinical activity.



\medskip\hrule\medskip

{\footnotesize \href{../../../../continuum/strategy/capability/4bf7a8260e1444fdac44e97c3557b533.md}{Canonical record: 4bf7a8260e1444fdac44e97c3557b533}; full qualifications and provenance.}\par
\paragraph*{II.1.3.6.23.1 [CT3] Perioperative Logistics Coordination}
\phantomsection
\label{navigation:II.1.3.6.23.1}
\addcontentsline{toc}{paragraph}{II.1.3.6.23.1 [CT3] Perioperative Logistics Coordination}

{\small Historical predecessor: R.2.154 \hyperref[element:e7813a06b9df462fa135c299261f192e]{[CT3] Perioperative Logistics Coordination}; Canonical ID: \texttt{\seqsplit{e7813a06b9df462fa135c299261f192e}}. Superseded historical identity — direct correspondence · superseded.}\par
\label{element:6b7c08b60b5b48519bfe4b6f993dd6fb}
\index{[CT3] Perioperative Logistics Coordination (II.1.3.6.23.1)}
{\small Capability · Proposed · current}\par

Coordinate theatre case progression information and notifications to supporting perioperative operational teams.



\textbf{Overview}



Coordinate theatre case progression information and notifications to supporting perioperative operational teams. The reviewed hierarchy assigns this enduring ability to its stated business-enabling context.



\textbf{Intent}



Coordinate theatre case progression information and notifications to supporting perioperative operational teams.



\textbf{Boundary}



Information administration and coordination preserve the authority of the responsible source, clinical actor and governed domain; they confer no additional real-world management authority.


{\footnotesize \href{../../../../continuum/strategy/capability/6b7c08b60b5b48519bfe4b6f993dd6fb.md}{Canonical record: 6b7c08b60b5b48519bfe4b6f993dd6fb}; full qualifications and provenance.}\par
\subparagraph*{II.1.3.6.23.1.1 [Feature] Operating Theatre Case Progression}
\phantomsection
\label{navigation:II.1.3.6.23.1.1}
\addcontentsline{toc}{subparagraph}{II.1.3.6.23.1.1 [Feature] Operating Theatre Case Progression}

\label{element:2fd65f09f7384f0486280e56ec5a9e84}
\index{[Feature] Operating Theatre Case Progression (II.1.3.6.23.1.1)}
{\small Capability / Feature · Canonical · current}\par

Track surgical case state transitions (CALLED → IN\_THEATRE → ANAESTHETISED → INCISION → CLOSED → RECOVERY).



Case progression records externally established surgical milestones in an operationally meaningful sequence. It makes the theatre situation visible without authorising or performing the clinical acts denoted by anaesthesia, incision, closure or recovery.


{\footnotesize \href{../../../../continuum/strategy/capability/feature/2fd65f09f7384f0486280e56ec5a9e84.md}{Canonical record: 2fd65f09f7384f0486280e56ec5a9e84}; full qualifications and provenance.}\par
\subparagraph*{II.1.3.6.23.1.2 [Feature] Perioperative Resource Notification}
\phantomsection
\label{navigation:II.1.3.6.23.1.2}
\addcontentsline{toc}{subparagraph}{II.1.3.6.23.1.2 [Feature] Perioperative Resource Notification}

\label{element:1f2446ccc59348a6bc54d8268ef4e6e1}
\index{[Feature] Perioperative Resource Notification (II.1.3.6.23.1.2)}
{\small Capability / Feature · Canonical · current}\par

Dispatch alerts to portering, sterilization, and recovery teams based on surgical case progression milestones.



Perioperative notification connects relevant case milestones with portering, sterilisation and recovery participants. Observation alone does not realise the Feature's dispatch responsibility; the established case-progression behaviour includes the required alerts without adding another Function or Service.


{\footnotesize \href{../../../../continuum/strategy/capability/feature/1f2446ccc59348a6bc54d8268ef4e6e1.md}{Canonical record: 1f2446ccc59348a6bc54d8268ef4e6e1}; full qualifications and provenance.}\par
\subparagraph*{II.1.3.6.23.1.3 [Feature] Theatre Session Capacity Management}
\phantomsection
\label{navigation:II.1.3.6.23.1.3}
\addcontentsline{toc}{subparagraph}{II.1.3.6.23.1.3 [Feature] Theatre Session Capacity Management}

\label{element:258c27077caf4bf88e562e8dd607bdfd}
\index{[Feature] Theatre Session Capacity Management (II.1.3.6.23.1.3)}
{\small Capability / Feature · Proposed · current}\par

Represent available and committed procedural-session capacity with its resource dependencies.



\textbf{Overview}



Theatre capacity involves session time, team capability, equipment and recovery provision rather than a count of theatres.



\textbf{Boundary}



Local responsibility retains its capacity meaning; Service Capacity Management consumes it without acquiring ownership or imposing one common capacity model.


{\footnotesize \href{../../../../continuum/strategy/capability/feature/258c27077caf4bf88e562e8dd607bdfd.md}{Canonical record: 258c27077caf4bf88e562e8dd607bdfd}; full qualifications and provenance.}\par
\subsubsection*{II.1.3.6.24 [CT2] Ward Operations}
\phantomsection
\label{navigation:II.1.3.6.24}
\addcontentsline{toc}{subsubsection}{II.1.3.6.24 [CT2] Ward Operations}

{\small Historical predecessor: R.2.157 \hyperref[element:bf3d2e3581f24beda20a6a3a3348460e]{[CT2] Ward Operations}; Canonical ID: \texttt{\seqsplit{bf3d2e3581f24beda20a6a3a3348460e}}. Superseded historical identity — direct correspondence · superseded.}\par
\label{element:f996ecca728f4942a88b034aca818a99}
\index{[CT2] Ward Operations (II.1.3.6.24)}
{\small Capability · Canonical · current}\par

Ward Operations coordinates ward occupancy and care-place context and assembles the information needed around shift changes. Its responsibility includes the ward-specific meaning of capacity and patient flow rather than a generic resource counter. It can make pending concerns and tasks visible without becoming the authority managing clinical handover.



\textbf{Overview}



Ward Operations coordinates ward occupancy and care-place context and assembles the information needed around shift changes. Its responsibility includes the ward-specific meaning of capacity and patient flow rather than a generic resource counter. It can make pending concerns and tasks visible without becoming the authority managing clinical handover.



\textbf{Intent}



Ward Operations coordinates ward occupancy and care-place context and assembles the information needed around shift changes. Its responsibility includes the ward-specific meaning of capacity and patient flow rather than a generic resource counter. It can make pending concerns and tasks visible without becoming the authority managing clinical handover.



\textbf{Boundary}



Information administration does not confer authority over the represented person, organisation, service, place, device or clinical activity.



\medskip\hrule\medskip

{\footnotesize \href{../../../../continuum/strategy/capability/f996ecca728f4942a88b034aca818a99.md}{Canonical record: f996ecca728f4942a88b034aca818a99}; full qualifications and provenance.}\par
\paragraph*{II.1.3.6.24.1 [CT3] Multidisciplinary Inpatient Care Coordination}
\phantomsection
\label{navigation:II.1.3.6.24.1}
\addcontentsline{toc}{paragraph}{II.1.3.6.24.1 [CT3] Multidisciplinary Inpatient Care Coordination}

\label{element:c798e38e12e5428dbb8b3cabc977d430}
\index{[CT3] Multidisciplinary Inpatient Care Coordination (II.1.3.6.24.1)}
{\small Capability · Canonical · current}\par

Definition unresolved in the baseline. No definition has been synthesised from the name.



\medskip\hrule\medskip

{\footnotesize \href{../../../../continuum/strategy/capability/c798e38e12e5428dbb8b3cabc977d430.md}{Canonical record: c798e38e12e5428dbb8b3cabc977d430}; full qualifications and provenance.}\par
\subparagraph*{II.1.3.6.24.1.1 [Feature] Inpatient Clinical Progression Tracking}
\phantomsection
\label{navigation:II.1.3.6.24.1.1}
\addcontentsline{toc}{subparagraph}{II.1.3.6.24.1.1 [Feature] Inpatient Clinical Progression Tracking}

\label{element:863265027ecb432595b9c6d45adf1190}
\index{[Feature] Inpatient Clinical Progression Tracking (II.1.3.6.24.1.1)}
{\small Capability / Feature · Canonical · current}\par

Capture inpatient clinical milestone completions (e.g., medical stability achieved, allied health clearance).



Inpatient progression tracking captures clinical milestones reported by the responsible participants. General clinical-information capabilities and configured workflow composition suffice for the established Feature; capture does not require inventing a specialised Function or make Harmonia the authority declaring clinical stability.


{\footnotesize \href{../../../../continuum/strategy/capability/feature/863265027ecb432595b9c6d45adf1190.md}{Canonical record: 863265027ecb432595b9c6d45adf1190}; full qualifications and provenance.}\par
\subparagraph*{II.1.3.6.24.1.2 [Feature] Inpatient Rounding Context Provision}
\phantomsection
\label{navigation:II.1.3.6.24.1.2}
\addcontentsline{toc}{subparagraph}{II.1.3.6.24.1.2 [Feature] Inpatient Rounding Context Provision}

\label{element:c0dfa2ef22354340a4f1488790359bbd}
\index{[Feature] Inpatient Rounding Context Provision (II.1.3.6.24.1.2)}
{\small Capability / Feature · Canonical · current}\par

Aggregate diagnostic results, current medication administration records, and nursing notes for multidisciplinary inpatient rounds.



Rounding context brings results, medication records and nursing notes into an accessible multidisciplinary view. Assembly serves the clinical discussion but does not perform rounds, determine pending clinical work or establish the team's professional conclusions.


{\footnotesize \href{../../../../continuum/strategy/capability/feature/c0dfa2ef22354340a4f1488790359bbd.md}{Canonical record: c0dfa2ef22354340a4f1488790359bbd}; full qualifications and provenance.}\par
\paragraph*{II.1.3.6.24.2 [CT3] Ward Logistics Coordination}
\phantomsection
\label{navigation:II.1.3.6.24.2}
\addcontentsline{toc}{paragraph}{II.1.3.6.24.2 [CT3] Ward Logistics Coordination}

{\small Historical predecessor: R.2.156 \hyperref[element:8a9dcc3b5f0143458ec71bc188a56020]{[CT3] Ward Logistics Coordination}; Canonical ID: \texttt{\seqsplit{8a9dcc3b5f0143458ec71bc188a56020}}. Superseded historical identity — direct correspondence · superseded.}\par
\label{element:3249647f7b9846c899f982ab339c81f9}
\index{[CT3] Ward Logistics Coordination (II.1.3.6.24.2)}
{\small Capability · Proposed · current}\par

Coordinate ward patient flow, handover context and ward-specific operational logistics and capacity information.



\textbf{Overview}



Coordinate ward patient flow, handover context and ward-specific operational logistics and capacity information. The reviewed hierarchy assigns this enduring ability to its stated business-enabling context.



\textbf{Intent}



Coordinate ward patient flow, handover context and ward-specific operational logistics and capacity information.



\textbf{Boundary}



Information administration and coordination preserve the authority of the responsible source, clinical actor and governed domain; they confer no additional real-world management authority.


{\footnotesize \href{../../../../continuum/strategy/capability/3249647f7b9846c899f982ab339c81f9.md}{Canonical record: 3249647f7b9846c899f982ab339c81f9}; full qualifications and provenance.}\par
\subparagraph*{II.1.3.6.24.2.1 [Feature] Clinical Handover Context Assembly}
\phantomsection
\label{navigation:II.1.3.6.24.2.1}
\addcontentsline{toc}{subparagraph}{II.1.3.6.24.2.1 [Feature] Clinical Handover Context Assembly}

\label{element:f8c748aa08444723b784887174831d51}
\index{[Feature] Clinical Handover Context Assembly (II.1.3.6.24.2.1)}
{\small Capability / Feature · Canonical · current}\par

Collate current ward patient lists, active medical concerns, and pending tasks for nursing shift handovers.



Handover context assembly brings patient lists, current concerns and pending tasks together for the participants conducting a shift handover. Assembly is the established responsibility: it does not perform clinical review, decide pending clinical work or manage the handover itself.


{\footnotesize \href{../../../../continuum/strategy/capability/feature/f8c748aa08444723b784887174831d51.md}{Canonical record: f8c748aa08444723b784887174831d51}; full qualifications and provenance.}\par
\subparagraph*{II.1.3.6.24.2.2 [Feature] Ward Care Capacity Management}
\phantomsection
\label{navigation:II.1.3.6.24.2.2}
\addcontentsline{toc}{subparagraph}{II.1.3.6.24.2.2 [Feature] Ward Care Capacity Management}

\label{element:0d45b0afc0584afba4bd0414f5f10e37}
\index{[Feature] Ward Care Capacity Management (II.1.3.6.24.2.2)}
{\small Capability / Feature · Proposed · current}\par

Represent ward care capacity in relation to staffing, care-place suitability and patient dependency.



\textbf{Overview}



A clean bed can remain unusable when staffing or patient needs constrain the ward; source clinical facts remain attributed.



\textbf{Boundary}



Local responsibility retains its capacity meaning; Service Capacity Management consumes it without acquiring ownership or imposing one common capacity model.


{\footnotesize \href{../../../../continuum/strategy/capability/feature/0d45b0afc0584afba4bd0414f5f10e37.md}{Canonical record: 0d45b0afc0584afba4bd0414f5f10e37}; full qualifications and provenance.}\par
\subparagraph*{II.1.3.6.24.2.3 [Feature] Ward Patient Flow Coordination}
\phantomsection
\label{navigation:II.1.3.6.24.2.3}
\addcontentsline{toc}{subparagraph}{II.1.3.6.24.2.3 [Feature] Ward Patient Flow Coordination}

\label{element:011d8c6481264ff58b7866a2c74a235b}
\index{[Feature] Ward Patient Flow Coordination (II.1.3.6.24.2.3)}
{\small Capability / Feature · Canonical · current}\par

Coordinate inpatient ward occupancy, care-place allocation, and patient bed status across care facilities.



Ward flow coordination connects occupancy, place allocation and bed status within the ward's operational context. Those facts contribute to ward capacity but do not exhaust its acuity, staffing and cohorting responsibilities or prove that another patient can be accepted.


{\footnotesize \href{../../../../continuum/strategy/capability/feature/011d8c6481264ff58b7866a2c74a235b.md}{Canonical record: 011d8c6481264ff58b7866a2c74a235b}; full qualifications and provenance.}\par
\subsubsection*{II.1.3.6.25 [CT2] Work Allocation \& Dispatch}
\phantomsection
\label{navigation:II.1.3.6.25}
\addcontentsline{toc}{subsubsection}{II.1.3.6.25 [CT2] Work Allocation \& Dispatch}

{\small Historical predecessor: R.2.159 \hyperref[element:f8c13b6f74b840cc96fbeada5865455a]{[CT2] Work Allocation \& Dispatch}; Canonical ID: \texttt{\seqsplit{f8c13b6f74b840cc96fbeada5865455a}}. Superseded historical identity — direct correspondence · superseded.}\par
\label{element:152a5252952c405994e9b171234efcb0}
\index{[CT2] Work Allocation \& Dispatch (II.1.3.6.25)}
{\small Capability · Canonical · current}\par

Work Allocation \& Dispatch progresses authorised operational work from a recognised requirement through matching, dispatch, acknowledgement and oversight. It distinguishes an instantiated work item from its assignment and from evidence of performance, allowing stalled activity to be identified and escalated. Examples such as cleaning or transport concern real-world support work, not technical message queues.



\textbf{Overview}



Work Allocation \& Dispatch progresses authorised operational work from a recognised requirement through matching, dispatch, acknowledgement and oversight. It distinguishes an instantiated work item from its assignment and from evidence of performance, allowing stalled activity to be identified and escalated. Examples such as cleaning or transport concern real-world support work, not technical message queues.



\textbf{Intent}



Work Allocation \& Dispatch progresses authorised operational work from a recognised requirement through matching, dispatch, acknowledgement and oversight. It distinguishes an instantiated work item from its assignment and from evidence of performance, allowing stalled activity to be identified and escalated. Examples such as cleaning or transport concern real-world support work, not technical message queues.



\textbf{Boundary}



Information administration does not confer authority over the represented person, organisation, service, place, device or clinical activity.



\medskip\hrule\medskip

{\footnotesize \href{../../../../continuum/strategy/capability/152a5252952c405994e9b171234efcb0.md}{Canonical record: 152a5252952c405994e9b171234efcb0}; full qualifications and provenance.}\par
\paragraph*{II.1.3.6.25.1 [CT3] Operational Work Assignment \& Dispatch Coordination}
\phantomsection
\label{navigation:II.1.3.6.25.1}
\addcontentsline{toc}{paragraph}{II.1.3.6.25.1 [CT3] Operational Work Assignment \& Dispatch Coordination}

{\small Historical predecessor: R.2.158 \hyperref[element:b42afda9e876463ba56fa8f73fe7a32a]{[CT3] Operational Work Assignment \& Dispatch Coordination}; Canonical ID: \texttt{\seqsplit{b42afda9e876463ba56fa8f73fe7a32a}}. Superseded historical identity — direct correspondence · superseded.}\par
\label{element:3b58935cfa974176989c3c75d76a00fe}
\index{[CT3] Operational Work Assignment \& Dispatch Coordination (II.1.3.6.25.1)}
{\small Capability · Proposed · current}\par

Coordinate qualified operational work assignment, dispatch, acknowledgement, progression and escalation within assigned authority.



\textbf{Overview}



Coordinate qualified operational work assignment, dispatch, acknowledgement, progression and escalation within assigned authority. The reviewed hierarchy assigns this enduring ability to its stated business-enabling context.



\textbf{Intent}



Coordinate qualified operational work assignment, dispatch, acknowledgement, progression and escalation within assigned authority.



\textbf{Boundary}



Information administration and coordination preserve the authority of the responsible source, clinical actor and governed domain; they confer no additional real-world management authority.


{\footnotesize \href{../../../../continuum/strategy/capability/3b58935cfa974176989c3c75d76a00fe.md}{Canonical record: 3b58935cfa974176989c3c75d76a00fe}; full qualifications and provenance.}\par
\subparagraph*{II.1.3.6.25.1.1 [Feature] Work Dispatch Delivery}
\phantomsection
\label{navigation:II.1.3.6.25.1.1}
\addcontentsline{toc}{subparagraph}{II.1.3.6.25.1.1 [Feature] Work Dispatch Delivery}

\label{element:095826f788b64784af38af31bbe77105}
\index{[Feature] Work Dispatch Delivery (II.1.3.6.25.1.1)}
{\small Capability / Feature · Canonical · current}\par

Deliver work orders to target workers, devices, or team queues with technical and operational acknowledgement capture.



Dispatch delivery brings assigned operational work to the intended worker or team and retains acknowledgement evidence. Technical receipt and operational acknowledgement must remain distinguishable; neither alone proves that the work was performed or its intended outcome achieved.


{\footnotesize \href{../../../../continuum/strategy/capability/feature/095826f788b64784af38af31bbe77105.md}{Canonical record: 095826f788b64784af38af31bbe77105}; full qualifications and provenance.}\par
\subparagraph*{II.1.3.6.25.1.2 [Feature] Work Item Instantiation}
\phantomsection
\label{navigation:II.1.3.6.25.1.2}
\addcontentsline{toc}{subparagraph}{II.1.3.6.25.1.2 [Feature] Work Item Instantiation}

\label{element:321fa656c71f4019842b5b991dbc9ff4}
\index{[Feature] Work Item Instantiation (II.1.3.6.25.1.2)}
{\small Capability / Feature · Canonical · current}\par

Instantiate operational units of work from clinical triggers (e.g., bed cleaning required, specimen transport required).



Work-item instantiation gives an authorised operational requirement an identifiable unit of work. It distinguishes recognising that support activity is needed from assigning it, dispatching it and recording its completion; a clinical trigger does not make the resulting support item clinical work authority.


{\footnotesize \href{../../../../continuum/strategy/capability/feature/321fa656c71f4019842b5b991dbc9ff4.md}{Canonical record: 321fa656c71f4019842b5b991dbc9ff4}; full qualifications and provenance.}\par
\subparagraph*{II.1.3.6.25.1.3 [Feature] Work Progress Oversight \& Escalation}
\phantomsection
\label{navigation:II.1.3.6.25.1.3}
\addcontentsline{toc}{subparagraph}{II.1.3.6.25.1.3 [Feature] Work Progress Oversight \& Escalation}

\label{element:abc12faeffd24917aad6748b0f0ca8d7}
\index{[Feature] Work Progress Oversight \& Escalation (II.1.3.6.25.1.3)}
{\small Capability / Feature · Canonical · current}\par

Monitor work progression against time-to-completion thresholds, escalating stalled or unacknowledged tasks.



Progress oversight identifies stalled or unacknowledged operational work against applicable timing expectations. Escalation serves completion management and does not itself constitute independent assurance or transfer authority over clinical decisions related to the support activity.


{\footnotesize \href{../../../../continuum/strategy/capability/feature/abc12faeffd24917aad6748b0f0ca8d7.md}{Canonical record: abc12faeffd24917aad6748b0f0ca8d7}; full qualifications and provenance.}\par
\subparagraph*{II.1.3.6.25.1.4 [Feature] Worker Matching \& Allocation}
\phantomsection
\label{navigation:II.1.3.6.25.1.4}
\addcontentsline{toc}{subparagraph}{II.1.3.6.25.1.4 [Feature] Worker Matching \& Allocation}

\label{element:caac8a83dc974cdb9c9bc1b778b9cee7}
\index{[Feature] Worker Matching \& Allocation (II.1.3.6.25.1.4)}
{\small Capability / Feature · Canonical · current}\par

Match work items against available staff roles, credentials, and locations based on allocation policies.



Worker matching relates an operational requirement to eligible roles, credentials, locations and availability under policy. It consumes those constraints rather than changing their authoritative records, and does not turn operational allocation into permission to decide clinical work.


{\footnotesize \href{../../../../continuum/strategy/capability/feature/caac8a83dc974cdb9c9bc1b778b9cee7.md}{Canonical record: caac8a83dc974cdb9c9bc1b778b9cee7}; full qualifications and provenance.}\par
\subsubsection*{II.1.3.6.26 [CT2] Workforce Qualification \& Development Enablement [superseded]}
\phantomsection
\label{navigation:II.1.3.6.26}
\addcontentsline{toc}{subsubsection}{II.1.3.6.26 [CT2] Workforce Qualification \& Development Enablement [superseded]}

\label{element:d24a1f4c4f674c45a95c317998ff4186}
\index{[CT2] Workforce Qualification \& Development Enablement (II.1.3.6.26)}
{\small Capability · Proposed · superseded}\par
\href{../../../../continuum/strategy/capability/d24a1f4c4f674c45a95c317998ff4186.md}{Historical canonical definition and provenance}.\par
\paragraph*{II.1.3.6.26.1 [CT3] Credential Management}
\phantomsection
\label{navigation:II.1.3.6.26.1}
\addcontentsline{toc}{paragraph}{II.1.3.6.26.1 [CT3] Credential Management}

\label{element:3a7a4d2d0ce8446b8a6293284fb16b83}
\index{[CT3] Credential Management (II.1.3.6.26.1)}
{\small Capability · Canonical · pending\_adjudication}\par

Definition unresolved in the baseline. No definition has been synthesised from the name.



\medskip\hrule\medskip

{\small Recorded unresolved matters: source-qualified matters; exact scope remains in the canonical record.}\par
{\footnotesize \href{../../../../continuum/strategy/capability/3a7a4d2d0ce8446b8a6293284fb16b83.md}{Canonical record: 3a7a4d2d0ce8446b8a6293284fb16b83}; full qualifications and provenance.}\par
\subsection*{II.1.3.7 [CT1] Human Collaboration \& Work Enablement [superseded]}
\phantomsection
\label{navigation:II.1.3.7}
\addcontentsline{toc}{subsection}{II.1.3.7 [CT1] Human Collaboration \& Work Enablement [superseded]}

\label{element:c9347f73b21242c5b8cce41252b74da0}
\index{[CT1] Human Collaboration \& Work Enablement (II.1.3.7)}
{\small Capability · Proposed · superseded}\par
\href{../../../../continuum/strategy/capability/c9347f73b21242c5b8cce41252b74da0.md}{Historical canonical definition and provenance}.\par
\subsubsection*{II.1.3.7.1 [CT2] Collaborative Interaction Enablement [superseded]}
\phantomsection
\label{navigation:II.1.3.7.1}
\addcontentsline{toc}{subsubsection}{II.1.3.7.1 [CT2] Collaborative Interaction Enablement [superseded]}

\label{element:6140b21c874241d092708e2411caceed}
\index{[CT2] Collaborative Interaction Enablement (II.1.3.7.1)}
{\small Capability · Proposed · superseded}\par
\href{../../../../continuum/strategy/capability/6140b21c874241d092708e2411caceed.md}{Historical canonical definition and provenance}.\par
\subsubsection*{II.1.3.7.2 [CT2] Operational Activity \& Temporal Enablement [superseded]}
\phantomsection
\label{navigation:II.1.3.7.2}
\addcontentsline{toc}{subsubsection}{II.1.3.7.2 [CT2] Operational Activity \& Temporal Enablement [superseded]}

\label{element:ca96b2b2ff434854995711e403e95f6f}
\index{[CT2] Operational Activity \& Temporal Enablement (II.1.3.7.2)}
{\small Capability · Proposed · superseded}\par
\href{../../../../continuum/strategy/capability/ca96b2b2ff434854995711e403e95f6f.md}{Historical canonical definition and provenance}.\par
\paragraph*{II.1.3.7.2.1 [CT3] Entity-Specific Operational Coordination}
\phantomsection
\label{navigation:II.1.3.7.2.1}
\addcontentsline{toc}{paragraph}{II.1.3.7.2.1 [CT3] Entity-Specific Operational Coordination}

\label{element:9f0c1429b64f4f88982fbbda6eaec233}
\index{[CT3] Entity-Specific Operational Coordination (II.1.3.7.2.1)}
{\small Capability · Proposed · pending\_adjudication}\par

Coordinate governed operational progression relative to one represented entity, including interested behaviour participation and controlled recovery.



\textbf{Overview}



Activity concerning the same patient, practitioner or bed needs a coherent entity-relative frame. Interest fan-out, serial progression, fresh information context and uncertain effects remain distinguishable from generic multi-work-unit progression.



\textbf{Intent}



Keep the progression of one represented entity coherent across interested operational behaviours.



\textbf{Boundary}



Entity coordination does not own all entity information, determine clinical work or become the source of clinical decisions. No Twin component allocation follows from this proposal.


{\small Recorded unresolved matters: source-qualified matters; exact scope remains in the canonical record.}\par
{\footnotesize \href{../../../../continuum/strategy/capability/9f0c1429b64f4f88982fbbda6eaec233.md}{Canonical record: 9f0c1429b64f4f88982fbbda6eaec233}; full qualifications and provenance.}\par
\subparagraph*{II.1.3.7.2.1.1 [Feature] Entity Behaviour Failure \& Recovery Governance}
\phantomsection
\label{navigation:II.1.3.7.2.1.1}
\addcontentsline{toc}{subparagraph}{II.1.3.7.2.1.1 [Feature] Entity Behaviour Failure \& Recovery Governance}

\label{element:bdfaeebab74e4d4d96a5e3799a3aa07b}
\index{[Feature] Entity Behaviour Failure \& Recovery Governance (II.1.3.7.2.1.1)}
{\small Capability / Feature · Proposed · pending\_adjudication}\par

Apply explicit behaviour-failure dispositions and controlled, authorised resumption for entity-relative processing.



\textbf{Overview}



Failure may stop a behaviour, an input or the entity progression and may leave effects uncertain. Resumption must preserve that distinction rather than imply rollback or automatic safe retry.



\textbf{Boundary}



Recovery cannot silently declare no effects, bypass governance or resume a halted entity automatically.


{\small Recorded unresolved matters: source-qualified matters; exact scope remains in the canonical record.}\par
{\footnotesize \href{../../../../continuum/strategy/capability/feature/bdfaeebab74e4d4d96a5e3799a3aa07b.md}{Canonical record: bdfaeebab74e4d4d96a5e3799a3aa07b}; full qualifications and provenance.}\par
\subparagraph*{II.1.3.7.2.1.2 [Feature] Entity-Relative Activity Coordination}
\phantomsection
\label{navigation:II.1.3.7.2.1.2}
\addcontentsline{toc}{subparagraph}{II.1.3.7.2.1.2 [Feature] Entity-Relative Activity Coordination}

\label{element:b9960d7de097464ea6de94929ca48c89}
\index{[Feature] Entity-Relative Activity Coordination (II.1.3.7.2.1.2)}
{\small Capability / Feature · Proposed · pending\_adjudication}\par

Sequence applicable behaviours for an identifiable entity input while preserving its operational reference frame.



\textbf{Overview}



Interest matching can select multiple behaviours while the entity progresses one input at a time. Each behaviour uses a fresh managed information context.



\textbf{Boundary}



Fan-out is not concurrent execution; clinical decisions and originating information remain with their owners.


{\small Recorded unresolved matters: source-qualified matters; exact scope remains in the canonical record.}\par
{\footnotesize \href{../../../../continuum/strategy/capability/feature/b9960d7de097464ea6de94929ca48c89.md}{Canonical record: b9960d7de097464ea6de94929ca48c89}; full qualifications and provenance.}\par
\paragraph*{II.1.3.7.2.2 [CT3] Governed Operational Activity Design}
\phantomsection
\label{navigation:II.1.3.7.2.2}
\addcontentsline{toc}{paragraph}{II.1.3.7.2.2 [CT3] Governed Operational Activity Design}

\label{element:3a96040821c14fe38830411dab4f8453}
\index{[CT3] Governed Operational Activity Design (II.1.3.7.2.2)}
{\small Capability · Proposed · pending\_adjudication}\par

Define the information context, expected outcomes, continuation conditions and housekeeping of governed operational activity.



\textbf{Overview}



An activity needs a fulfilment definition even when only one bounded execution unit is involved. This capability specifies its semantic envelope while execution and event delivery remain separate responsibilities.



\textbf{Intent}



Make the expected information and fulfilment basis of an undertaking explicit before it executes.



\textbf{Boundary}



An activity definition is neither an orchestrator nor a choreographer. It does not choose clinical care or dynamically invent execution logic.


{\small Recorded unresolved matters: source-qualified matters; exact scope remains in the canonical record.}\par
{\footnotesize \href{../../../../continuum/strategy/capability/3a96040821c14fe38830411dab4f8453.md}{Canonical record: 3a96040821c14fe38830411dab4f8453}; full qualifications and provenance.}\par
\subparagraph*{II.1.3.7.2.2.1 [Feature] Activity Outcome \& Continuation Rules}
\phantomsection
\label{navigation:II.1.3.7.2.2.1}
\addcontentsline{toc}{subparagraph}{II.1.3.7.2.2.1 [Feature] Activity Outcome \& Continuation Rules}

\label{element:34c8309244c142ec8772dcbb11ecc4ef}
\index{[Feature] Activity Outcome \& Continuation Rules (II.1.3.7.2.2.1)}
{\small Capability / Feature · Proposed · pending\_adjudication}\par

Define how output roles, expected outcomes and continuation conditions are interpreted for an undertaking.



\textbf{Overview}



An output may contribute to completion, subsequent activity or retained evidence under different rules. Those meanings need to be specified rather than inferred from a returned payload.



\textbf{Boundary}



Completion policy is not business outcome truth; uncertain effects remain explicit.


{\small Recorded unresolved matters: source-qualified matters; exact scope remains in the canonical record.}\par
{\footnotesize \href{../../../../continuum/strategy/capability/feature/34c8309244c142ec8772dcbb11ecc4ef.md}{Canonical record: 34c8309244c142ec8772dcbb11ecc4ef}; full qualifications and provenance.}\par
\subparagraph*{II.1.3.7.2.2.2 [Feature] Operational Activity Context Definition}
\phantomsection
\label{navigation:II.1.3.7.2.2.2}
\addcontentsline{toc}{subparagraph}{II.1.3.7.2.2.2 [Feature] Operational Activity Context Definition}

\label{element:58dcd59389664b7d83d14b536d7d03d0}
\index{[Feature] Operational Activity Context Definition (II.1.3.7.2.2.2)}
{\small Capability / Feature · Proposed · pending\_adjudication}\par

Define the anticipated information context and expected fulfilment basis of an operational undertaking.



\textbf{Overview}



Even a one-unit activity needs a declared information and outcome frame. The definition is separate from loading values or running the business logic.



\textbf{Boundary}



No dynamically invented resource set or orchestration engine is implied.


{\small Recorded unresolved matters: source-qualified matters; exact scope remains in the canonical record.}\par
{\footnotesize \href{../../../../continuum/strategy/capability/feature/58dcd59389664b7d83d14b536d7d03d0.md}{Canonical record: 58dcd59389664b7d83d14b536d7d03d0}; full qualifications and provenance.}\par
\paragraph*{II.1.3.7.2.3 [CT3] Workflow \& Activity Coordination}
\phantomsection
\label{navigation:II.1.3.7.2.3}
\addcontentsline{toc}{paragraph}{II.1.3.7.2.3 [CT3] Workflow \& Activity Coordination}

\label{element:94c683c7588f475e973b58cb226345db}
\index{[CT3] Workflow \& Activity Coordination (II.1.3.7.2.3)}
{\small Capability · Canonical · pending\_adjudication}\par

Definition unresolved in the baseline. No definition has been synthesised from the name.



\medskip\hrule\medskip

{\small Recorded unresolved matters: source-qualified matters; exact scope remains in the canonical record.}\par
{\footnotesize \href{../../../../continuum/strategy/capability/94c683c7588f475e973b58cb226345db.md}{Canonical record: 94c683c7588f475e973b58cb226345db}; full qualifications and provenance.}\par
\subparagraph*{II.1.3.7.2.3.1 [Feature] Controlled Administrative Recovery}
\phantomsection
\label{navigation:II.1.3.7.2.3.1}
\addcontentsline{toc}{subparagraph}{II.1.3.7.2.3.1 [Feature] Controlled Administrative Recovery}

\label{element:21d0601a852043a4b0b1d4fe841108d2}
\index{[Feature] Controlled Administrative Recovery (II.1.3.7.2.3.1)}
{\small Capability / Feature · Proposed · pending\_adjudication}\par

Perform authorised recovery as identifiable governed activity through normal execution and evidence boundaries.



\textbf{Overview}



Operators may need to resume, reconcile or repair an operational concern. That is first-class work with permission and outcome standing, rather than direct mutation through a console.



\textbf{Boundary}



Administrative interaction grants no raw cache/database or subsystem-internal bypass.


{\small Recorded unresolved matters: source-qualified matters; exact scope remains in the canonical record.}\par
{\footnotesize \href{../../../../continuum/strategy/capability/feature/21d0601a852043a4b0b1d4fe841108d2.md}{Canonical record: 21d0601a852043a4b0b1d4fe841108d2}; full qualifications and provenance.}\par
\subparagraph*{II.1.3.7.2.3.2 [Feature] Synthetic Task Orchestration}
\phantomsection
\label{navigation:II.1.3.7.2.3.2}
\addcontentsline{toc}{subparagraph}{II.1.3.7.2.3.2 [Feature] Synthetic Task Orchestration}

\label{element:7cf9dd3aeeb84a83b2445f897db9bd4e}
\index{[Feature] Synthetic Task Orchestration (II.1.3.7.2.3.2)}
{\small Capability / Feature · Canonical · pending\_adjudication}\par

Execute non-human, automated platform activity units and sequence orchestration.



\medskip\hrule\medskip

{\small Recorded unresolved matters: source-qualified matters; exact scope remains in the canonical record.}\par
{\footnotesize \href{../../../../continuum/strategy/capability/feature/7cf9dd3aeeb84a83b2445f897db9bd4e.md}{Canonical record: 7cf9dd3aeeb84a83b2445f897db9bd4e}; full qualifications and provenance.}\par
\subparagraph*{II.1.3.7.2.3.3 [Feature] To Do Decision \& Approval Coordination}
\phantomsection
\label{navigation:II.1.3.7.2.3.3}
\addcontentsline{toc}{subparagraph}{II.1.3.7.2.3.3 [Feature] To Do Decision \& Approval Coordination}

\label{element:377a40e8f2d64840a10dd55f87a9d005}
\index{[Feature] To Do Decision \& Approval Coordination (II.1.3.7.2.3.3)}
{\small Capability / Feature · Canonical · pending\_adjudication}\par

Coordinate human clinical review items, sign-offs, and approval workflows.



\medskip\hrule\medskip

{\small Recorded unresolved matters: source-qualified matters; exact scope remains in the canonical record.}\par
{\footnotesize \href{../../../../continuum/strategy/capability/feature/377a40e8f2d64840a10dd55f87a9d005.md}{Canonical record: 377a40e8f2d64840a10dd55f87a9d005}; full qualifications and provenance.}\par
\subparagraph*{II.1.3.7.2.3.4 [Feature] Uncertain Activity Outcome Reconciliation}
\phantomsection
\label{navigation:II.1.3.7.2.3.4}
\addcontentsline{toc}{subparagraph}{II.1.3.7.2.3.4 [Feature] Uncertain Activity Outcome Reconciliation}

\label{element:eb06ecc956404d0397701ceafb5fbd46}
\index{[Feature] Uncertain Activity Outcome Reconciliation (II.1.3.7.2.3.4)}
{\small Capability / Feature · Proposed · pending\_adjudication}\par

Reconcile indeterminate activity effects through authoritative observation before establishing a safe next disposition.



\textbf{Overview}



A timeout or absent response leaves a materially significant outcome unknown. Reconciliation can confirm an effect, establish absence or retain uncertainty without automatically retrying the original work.



\textbf{Boundary}



Execution uncertainty differs from insufficient assurance evidence; a retry does not prove the earlier attempt had no effect.


{\small Recorded unresolved matters: source-qualified matters; exact scope remains in the canonical record.}\par
{\footnotesize \href{../../../../continuum/strategy/capability/feature/eb06ecc956404d0397701ceafb5fbd46.md}{Canonical record: eb06ecc956404d0397701ceafb5fbd46}; full qualifications and provenance.}\par
\subparagraph*{II.1.3.7.2.3.5 [Feature] Work Order Progression}
\phantomsection
\label{navigation:II.1.3.7.2.3.5}
\addcontentsline{toc}{subparagraph}{II.1.3.7.2.3.5 [Feature] Work Order Progression}

\label{element:5ca494e85869419c867618fd9648d03a}
\index{[Feature] Work Order Progression (II.1.3.7.2.3.5)}
{\small Capability / Feature · Canonical · pending\_adjudication}\par

Manage execution states and dispatch of human operational work orders.



\medskip\hrule\medskip

{\small Recorded unresolved matters: source-qualified matters; exact scope remains in the canonical record.}\par
{\footnotesize \href{../../../../continuum/strategy/capability/feature/5ca494e85869419c867618fd9648d03a.md}{Canonical record: 5ca494e85869419c867618fd9648d03a}; full qualifications and provenance.}\par
\subparagraph*{II.1.3.7.2.3.6 [Feature] Workflow Timeout \& Escalation Oversight}
\phantomsection
\label{navigation:II.1.3.7.2.3.6}
\addcontentsline{toc}{subparagraph}{II.1.3.7.2.3.6 [Feature] Workflow Timeout \& Escalation Oversight}

\label{element:ec55d06e393e4988a84c6157b719bd40}
\index{[Feature] Workflow Timeout \& Escalation Oversight (II.1.3.7.2.3.6)}
{\small Capability / Feature · Canonical · pending\_adjudication}\par

Supervise operational deadlines, executing automated timeout notifications and supervisory escalations.



\medskip\hrule\medskip

{\small Recorded unresolved matters: source-qualified matters; exact scope remains in the canonical record.}\par
{\footnotesize \href{../../../../continuum/strategy/capability/feature/ec55d06e393e4988a84c6157b719bd40.md}{Canonical record: ec55d06e393e4988a84c6157b719bd40}; full qualifications and provenance.}\par
\subsection*{II.1.3.8 [CT1] Independent Assurance Enablement}
\phantomsection
\label{navigation:II.1.3.8}
\addcontentsline{toc}{subsection}{II.1.3.8 [CT1] Independent Assurance Enablement}

{\small Historical predecessor: R.2.166 \hyperref[element:88bf57b7a92c4c1da1b5f0a67d3ebd82]{[CT1] Independent Assurance Enablement}; Canonical ID: \texttt{\seqsplit{88bf57b7a92c4c1da1b5f0a67d3ebd82}}. Superseded historical identity — direct correspondence · superseded.}\par
\label{element:616682b612d244429e656b040e002d56}
\index{[CT1] Independent Assurance Enablement (II.1.3.8)}
{\small Capability · Proposed · current}\par

Enable assurance specification and independent review of the assurance service while preserving the distinct Canonical responsibility for assurance against a governed subject.



\textbf{Overview}



Enable assurance specification and independent review of the assurance service while preserving the distinct Canonical responsibility for assurance against a governed subject.



\textbf{Intent}



Enable assurance specification and independent review of the assurance service while preserving the distinct Canonical responsibility for assurance against a governed subject.



\textbf{Boundary}



Assurance cannot manage, remediate or approve its subject merely because it observes it. Clinical adequacy, financial governance and source-information authority remain outside this mandate.


{\footnotesize \href{../../../../continuum/strategy/capability/616682b612d244429e656b040e002d56.md}{Canonical record: 616682b612d244429e656b040e002d56}; full qualifications and provenance.}\par
\subsubsection*{II.1.3.8.1 [CT2] Assurance Service Review}
\phantomsection
\label{navigation:II.1.3.8.1}
\addcontentsline{toc}{subsubsection}{II.1.3.8.1 [CT2] Assurance Service Review}

{\small Historical predecessor: R.2.162 \hyperref[element:b345cb12b45746409302899514126170]{[CT2] Assurance Service Review}; Canonical ID: \texttt{\seqsplit{b345cb12b45746409302899514126170}}. Superseded historical identity — direct correspondence · superseded.}\par
\label{element:b1a791458d644d4794895e4fa42ec83c}
\index{[CT2] Assurance Service Review (II.1.3.8.1)}
{\small Capability · Proposed · current}\par

Independently review the operation, effectiveness, appropriateness and conformance of the assurance service itself.



\textbf{Overview}



Independently review the operation, effectiveness, appropriateness and conformance of the assurance service itself. The reviewed hierarchy assigns this enduring ability to its stated business-enabling context.



\textbf{Intent}



Independently review the operation, effectiveness, appropriateness and conformance of the assurance service itself.



\textbf{Boundary}



Reviews the assurance service itself. Canonical evidence association, subject assessment, adjudication, findings and conclusions are retained separately pending adjudication.


{\footnotesize \href{../../../../continuum/strategy/capability/b1a791458d644d4794895e4fa42ec83c.md}{Canonical record: b1a791458d644d4794895e4fa42ec83c}; full qualifications and provenance.}\par
\paragraph*{II.1.3.8.1.1 [CT3] Assurance Service Evaluation}
\phantomsection
\label{navigation:II.1.3.8.1.1}
\addcontentsline{toc}{paragraph}{II.1.3.8.1.1 [CT3] Assurance Service Evaluation}

{\small Historical predecessor: R.2.161 \hyperref[element:be34b54314c6444f8418d645f6c83b6d]{[CT3] Assurance Service Evaluation}; Canonical ID: \texttt{\seqsplit{be34b54314c6444f8418d645f6c83b6d}}. Superseded historical identity — direct correspondence · superseded.}\par
\label{element:039343518db542d48cde6afab84a3e1e}
\index{[CT3] Assurance Service Evaluation (II.1.3.8.1.1)}
{\small Capability · Proposed · current}\par

Evaluate the operation, effectiveness, appropriateness and conformance of the assurance service itself through independent review.



\textbf{Overview}



Evaluate the operation, effectiveness, appropriateness and conformance of the assurance service itself through independent review. The reviewed hierarchy assigns this enduring ability to its stated business-enabling context.



\textbf{Intent}



Evaluate the operation, effectiveness, appropriateness and conformance of the assurance service itself through independent review.



\textbf{Boundary}



Evaluation concerns operation, effectiveness, appropriateness and conformance of the assurance service; it is not execution of assurance against a subject.


{\footnotesize \href{../../../../continuum/strategy/capability/039343518db542d48cde6afab84a3e1e.md}{Canonical record: 039343518db542d48cde6afab84a3e1e}; full qualifications and provenance.}\par
\subsubsection*{II.1.3.8.2 [CT2] Assurance Specification}
\phantomsection
\label{navigation:II.1.3.8.2}
\addcontentsline{toc}{subsubsection}{II.1.3.8.2 [CT2] Assurance Specification}

{\small Historical predecessor: R.2.165 \hyperref[element:8006ee0dd95f4d7fb1f408fa09b6ce19]{[CT2] Assurance Specification}; Canonical ID: \texttt{\seqsplit{8006ee0dd95f4d7fb1f408fa09b6ce19}}. Superseded historical identity — direct correspondence · superseded.}\par
\label{element:b81775d89d774cf8ba61731d13c67a54}
\index{[CT2] Assurance Specification (II.1.3.8.2)}
{\small Capability · Proposed · current}\par

Model how satisfaction will be assured and maintain applicable, temporally qualified criteria.



\textbf{Overview}



Model how satisfaction will be assured and maintain applicable, temporally qualified criteria.



\textbf{Intent}



Model how satisfaction will be assured and maintain applicable, temporally qualified criteria.



\textbf{Boundary}



Modelling and criteria stewardship do not approve governing requirements or perform independent conclusion-making.


{\footnotesize \href{../../../../continuum/strategy/capability/b81775d89d774cf8ba61731d13c67a54.md}{Canonical record: b81775d89d774cf8ba61731d13c67a54}; full qualifications and provenance.}\par
\paragraph*{II.1.3.8.2.1 [CT3] Assurance Criteria Administration}
\phantomsection
\label{navigation:II.1.3.8.2.1}
\addcontentsline{toc}{paragraph}{II.1.3.8.2.1 [CT3] Assurance Criteria Administration}

{\small Historical predecessor: R.2.163 \hyperref[element:6dab08ab1e3249d7b314843ca9ed0149]{[CT3] Assurance Criteria Administration}; Canonical ID: \texttt{\seqsplit{6dab08ab1e3249d7b314843ca9ed0149}}. Superseded historical identity — direct correspondence · superseded.}\par
\label{element:1bfd71f24321428c8871b587d413f47e}
\index{[CT3] Assurance Criteria Administration (II.1.3.8.2.1)}
{\small Capability · Canonical · current}\par

The capability to establish and manage reusable, governed and temporally identifiable assurance criteria through which subjects may be evaluated.



Definition — reviewed candidate



Administer reusable governed assurance criteria and their version-specific and temporal applicability without changing governing requirement authority.



\textbf{Overview}



Administer reusable governed assurance criteria and their version-specific and temporal applicability without changing governing requirement authority.



\textbf{Intent}



Administer reusable governed assurance criteria and their version-specific and temporal applicability without changing governing requirement authority.



\textbf{Boundary}



Information administration does not confer authority over the represented person, organisation, service, place, device or clinical activity.



\medskip\hrule\medskip

{\small Recorded unresolved matters: source-qualified matters; exact scope remains in the canonical record.}\par
{\footnotesize \href{../../../../continuum/strategy/capability/1bfd71f24321428c8871b587d413f47e.md}{Canonical record: 1bfd71f24321428c8871b587d413f47e}; full qualifications and provenance.}\par
\subparagraph*{II.1.3.8.2.1.1 [Feature] Temporal Assurance Criteria Stewardship}
\phantomsection
\label{navigation:II.1.3.8.2.1.1}
\addcontentsline{toc}{subparagraph}{II.1.3.8.2.1.1 [Feature] Temporal Assurance Criteria Stewardship}

\label{element:e345a4b1eef846a7b77568f519e2c75c}
\index{[Feature] Temporal Assurance Criteria Stewardship (II.1.3.8.2.1.1)}
{\small Capability / Feature · Proposed · current}\par

Manage identifiable assurance criteria and their relevant temporal and version applicability.



\textbf{Overview}



Evidence must be evaluated against the criteria applicable to its purpose and time rather than an unqualified current rule.



\textbf{Boundary}



Criteria stewardship does not approve governing requirements or select a universal result taxonomy.


{\footnotesize \href{../../../../continuum/strategy/capability/feature/e345a4b1eef846a7b77568f519e2c75c.md}{Canonical record: e345a4b1eef846a7b77568f519e2c75c}; full qualifications and provenance.}\par
\paragraph*{II.1.3.8.2.2 [CT3] Assurance Design}
\phantomsection
\label{navigation:II.1.3.8.2.2}
\addcontentsline{toc}{paragraph}{II.1.3.8.2.2 [CT3] Assurance Design}

{\small Historical predecessor: R.2.164 \hyperref[element:e1b8db7c95fa4e7ea296351a03532e24]{[CT3] Assurance Design}; Canonical ID: \texttt{\seqsplit{e1b8db7c95fa4e7ea296351a03532e24}}. Superseded historical identity — direct correspondence · superseded.}\par
\label{element:71c8a612b88f4ecba049e94d506bd0c2}
\index{[CT3] Assurance Design (II.1.3.8.2.2)}
{\small Capability · Canonical · current}\par

The capability to determine and define how satisfaction of a governed requirement, constraint or expected behaviour will be assured, including its assurance disposition, applicable criteria, required evidence and evaluation expectations.



Definition — reviewed candidate



Model how satisfaction of a governed subject or requirement will be assured, preserving applicable assurance dispositions and boundaries.



\textbf{Overview}



Model how satisfaction of a governed subject or requirement will be assured, preserving applicable assurance dispositions and boundaries.



\textbf{Intent}



Model how satisfaction of a governed subject or requirement will be assured, preserving applicable assurance dispositions and boundaries.



\textbf{Boundary}



Information administration does not confer authority over the represented person, organisation, service, place, device or clinical activity.



\medskip\hrule\medskip

{\small Recorded unresolved matters: source-qualified matters; exact scope remains in the canonical record.}\par
{\footnotesize \href{../../../../continuum/strategy/capability/71c8a612b88f4ecba049e94d506bd0c2.md}{Canonical record: 71c8a612b88f4ecba049e94d506bd0c2}; full qualifications and provenance.}\par
\subparagraph*{II.1.3.8.2.2.1 [Feature] Assurance Disposition \& Design Modelling}
\phantomsection
\label{navigation:II.1.3.8.2.2.1}
\addcontentsline{toc}{subparagraph}{II.1.3.8.2.2.1 [Feature] Assurance Disposition \& Design Modelling}

\label{element:c36180ebffe44b8db587a00e78b4b358}
\index{[Feature] Assurance Disposition \& Design Modelling (II.1.3.8.2.2.1)}
{\small Capability / Feature · Proposed · current}\par

Model how satisfaction of an applicable governed concern will be independently established.



\textbf{Overview}



This proposal provides a Feature frame for the established design responsibility, including disposition, criteria and evidence expectations, without creating an additional assurance Function.



\textbf{Boundary}



Modelling is distinct from requirement approval, subject management and evaluation.


{\footnotesize \href{../../../../continuum/strategy/capability/feature/c36180ebffe44b8db587a00e78b4b358.md}{Canonical record: c36180ebffe44b8db587a00e78b4b358}; full qualifications and provenance.}\par
\subsubsection*{II.1.3.8.3 [CT2] Independent Assurance Evaluation Enablement}
\phantomsection
\label{navigation:II.1.3.8.3}
\addcontentsline{toc}{subsubsection}{II.1.3.8.3 [CT2] Independent Assurance Evaluation Enablement}

\label{element:7fa5112efb724b90a17a03c7c19c516a}
\index{[CT2] Independent Assurance Evaluation Enablement (II.1.3.8.3)}
{\small Capability · Proposed · pending\_adjudication}\par

Associate trustworthy evidence, assess it and adjudicate independently reportable findings about governed service behaviour.



\textbf{Overview}



Associate trustworthy evidence, assess it and adjudicate independently reportable findings about governed service behaviour. Its constituent abilities distinguish Governed Assurance.



\textbf{Intent}



Produce findings supported by applicable evidence and criteria while retaining inability to conclude.



\textbf{Boundary}



The subject can supply evidence but cannot control its conclusion. Evidence insufficiency remains distinct from execution uncertainty and clinical judgement.


{\small Recorded unresolved matters: source-qualified matters; exact scope remains in the canonical record.}\par
{\footnotesize \href{../../../../continuum/strategy/capability/7fa5112efb724b90a17a03c7c19c516a.md}{Canonical record: 7fa5112efb724b90a17a03c7c19c516a}; full qualifications and provenance.}\par
\paragraph*{II.1.3.8.3.1 [CT3] Governed Assurance}
\phantomsection
\label{navigation:II.1.3.8.3.1}
\addcontentsline{toc}{paragraph}{II.1.3.8.3.1 [CT3] Governed Assurance}

\label{element:da1a983d4dcf4d5bb772b8f69bc08426}
\index{[CT3] Governed Assurance (II.1.3.8.3.1)}
{\small Capability · Canonical · pending\_adjudication}\par

The capability to independently evaluate a governed subject against applicable assurance criteria through the assembly and evaluation of sufficient trustworthy evidence, establishing explicit findings, exceptions and assurance conclusions.



\medskip\hrule\medskip

{\small Recorded unresolved matters: source-qualified matters; exact scope remains in the canonical record.}\par
{\footnotesize \href{../../../../continuum/strategy/capability/da1a983d4dcf4d5bb772b8f69bc08426.md}{Canonical record: da1a983d4dcf4d5bb772b8f69bc08426}; full qualifications and provenance.}\par
\subparagraph*{II.1.3.8.3.1.1 [Feature] Independent Evidence Assessment \& Adjudication}
\phantomsection
\label{navigation:II.1.3.8.3.1.1}
\addcontentsline{toc}{subparagraph}{II.1.3.8.3.1.1 [Feature] Independent Evidence Assessment \& Adjudication}

\label{element:3aedee8f78484b81955866bcf16b8a6d}
\index{[Feature] Independent Evidence Assessment \& Adjudication (II.1.3.8.3.1.1)}
{\small Capability / Feature · Proposed · pending\_adjudication}\par

Independently establish assurance context, assess evidence and adjudicate findings under applicable criteria.



\textbf{Overview}



The Feature unifies the review frame while retaining the three distinct canonical Functions and the boundary between criterion-level assessment and reportable conclusion.



\textbf{Boundary}



No new assurance Function or clinical-adequacy mandate is introduced; insufficient evidence establishes neither satisfaction nor non-satisfaction.


{\small Recorded unresolved matters: source-qualified matters; exact scope remains in the canonical record.}\par
{\footnotesize \href{../../../../continuum/strategy/capability/feature/3aedee8f78484b81955866bcf16b8a6d.md}{Canonical record: 3aedee8f78484b81955866bcf16b8a6d}; full qualifications and provenance.}\par
\subsection*{II.1.3.9 [CT1] Shared Information \& Coordination Enablement}
\phantomsection
\label{navigation:II.1.3.9}
\addcontentsline{toc}{subsection}{II.1.3.9 [CT1] Shared Information \& Coordination Enablement}

{\small Historical predecessor: R.2.214 \hyperref[element:eeabb9c3c526404cbe4302906560fa7e]{[CT1] Shared Information \& Coordination Enablement}; Canonical ID: \texttt{\seqsplit{eeabb9c3c526404cbe4302906560fa7e}}. Superseded historical identity — direct correspondence · superseded.}\par
\label{element:0062fcacae1740988b6be2bda2aa1674}
\index{[CT1] Shared Information \& Coordination Enablement (II.1.3.9)}
{\small Capability · Proposed · current}\par

Enable governed shared health-information use, exchange, stewardship, knowledge, collaboration, external work coordination, calendars, presentation and information design.



\textbf{Overview}



Enable governed shared health-information use, exchange, stewardship, knowledge, collaboration, external work coordination, calendars, presentation and information design.



\textbf{Intent}



Enable governed shared health-information use, exchange, stewardship, knowledge, collaboration, external work coordination, calendars, presentation and information design.



\textbf{Boundary}



Preservation, custody, discoverability and presentation do not confer originating fact authority. Clinical documentation administration and clinical-context use retain their own responsibilities.


{\footnotesize \href{../../../../continuum/strategy/capability/0062fcacae1740988b6be2bda2aa1674.md}{Canonical record: 0062fcacae1740988b6be2bda2aa1674}; full qualifications and provenance.}\par
\subsubsection*{II.1.3.9.1 [CT2] Calendar Management}
\phantomsection
\label{navigation:II.1.3.9.1}
\addcontentsline{toc}{subsubsection}{II.1.3.9.1 [CT2] Calendar Management}

{\small Historical predecessor: R.2.168 \hyperref[element:8b42aac6083441caa78b59f5e62eebd9]{[CT2] Calendar Management}; Canonical ID: \texttt{\seqsplit{8b42aac6083441caa78b59f5e62eebd9}}. Superseded historical identity — direct correspondence · superseded.}\par
\label{element:7669036f7e6749219412fc64df5f7064}
\index{[CT2] Calendar Management (II.1.3.9.1)}
{\small Capability · Canonical · current}\par

Calendar Management projects schedules and correlates events with temporal commitments. It brings rosters, on-call windows, facility calendars and appointment milestones into operational views without replacing their authoritative sources. This distinguishes the temporal context used for coordination from the business decisions that establish the schedule.



\textbf{Overview}



Calendar Management projects schedules and correlates events with temporal commitments. It brings rosters, on-call windows, facility calendars and appointment milestones into operational views without replacing their authoritative sources. This distinguishes the temporal context used for coordination from the business decisions that establish the schedule.



\textbf{Intent}



Calendar Management projects schedules and correlates events with temporal commitments. It brings rosters, on-call windows, facility calendars and appointment milestones into operational views without replacing their authoritative sources. This distinguishes the temporal context used for coordination from the business decisions that establish the schedule.



\textbf{Boundary}



Information administration does not confer authority over the represented person, organisation, service, place, device or clinical activity.



\medskip\hrule\medskip

{\footnotesize \href{../../../../continuum/strategy/capability/7669036f7e6749219412fc64df5f7064.md}{Canonical record: 7669036f7e6749219412fc64df5f7064}; full qualifications and provenance.}\par
\paragraph*{II.1.3.9.1.1 [CT3] Operational \& Clinical Calendar Coordination}
\phantomsection
\label{navigation:II.1.3.9.1.1}
\addcontentsline{toc}{paragraph}{II.1.3.9.1.1 [CT3] Operational \& Clinical Calendar Coordination}

{\small Historical predecessor: R.2.167 \hyperref[element:c0f6e5f7f5a94d1580fa6e0e8b82c93a]{[CT3] Operational \& Clinical Calendar Coordination}; Canonical ID: \texttt{\seqsplit{c0f6e5f7f5a94d1580fa6e0e8b82c93a}}. Superseded historical identity — direct correspondence · superseded.}\par
\label{element:f6d8f1310ba543f69e8b57bf05603db3}
\index{[CT3] Operational \& Clinical Calendar Coordination (II.1.3.9.1.1)}
{\small Capability · Proposed · current}\par

Coordinate projections of operational and clinical calendars and correlate temporal events without taking ownership of source bookings or rosters.



\textbf{Overview}



Coordinate projections of operational and clinical calendars and correlate temporal events without taking ownership of source bookings or rosters. The reviewed hierarchy assigns this enduring ability to its stated business-enabling context.



\textbf{Intent}



Coordinate projections of operational and clinical calendars and correlate temporal events without taking ownership of source bookings or rosters.



\textbf{Boundary}



Information administration and coordination preserve the authority of the responsible source, clinical actor and governed domain; they confer no additional real-world management authority.


{\footnotesize \href{../../../../continuum/strategy/capability/f6d8f1310ba543f69e8b57bf05603db3.md}{Canonical record: f6d8f1310ba543f69e8b57bf05603db3}; full qualifications and provenance.}\par
\subparagraph*{II.1.3.9.1.1.1 [Feature] Schedule Projection Integration}
\phantomsection
\label{navigation:II.1.3.9.1.1.1}
\addcontentsline{toc}{subparagraph}{II.1.3.9.1.1.1 [Feature] Schedule Projection Integration}

\label{element:883a68ab282b462988d5ae7a6ea2e980}
\index{[Feature] Schedule Projection Integration (II.1.3.9.1.1.1)}
{\small Capability / Feature · Canonical · current}\par

Project clinician rosters, on-call schedules, and facility operating calendars into operational views.



Schedule projection makes external roster, on-call and facility calendars available in operational views. It preserves their temporal usefulness without making the projection the authoritative schedule or granting authority to alter staffing and coverage decisions.


{\footnotesize \href{../../../../continuum/strategy/capability/feature/883a68ab282b462988d5ae7a6ea2e980.md}{Canonical record: 883a68ab282b462988d5ae7a6ea2e980}; full qualifications and provenance.}\par
\subparagraph*{II.1.3.9.1.1.2 [Feature] Temporal Event Correlation}
\phantomsection
\label{navigation:II.1.3.9.1.1.2}
\addcontentsline{toc}{subparagraph}{II.1.3.9.1.1.2 [Feature] Temporal Event Correlation}

\label{element:05e1c23baf1c47fbac8457f991efa2be}
\index{[Feature] Temporal Event Correlation (II.1.3.9.1.1.2)}
{\small Capability / Feature · Canonical · current}\par

Correlate incoming clinical events with scheduled appointments and operational calendar milestones.



Temporal correlation relates incoming events to the relevant appointment or calendar milestone. Association provides context for interpretation and coordination without assuming that temporal proximity proves attendance, fulfilment or the clinical outcome expected at that time.


{\footnotesize \href{../../../../continuum/strategy/capability/feature/05e1c23baf1c47fbac8457f991efa2be.md}{Canonical record: 05e1c23baf1c47fbac8457f991efa2be}; full qualifications and provenance.}\par
\subsubsection*{II.1.3.9.2 [CT2] Clinical Collaboration}
\phantomsection
\label{navigation:II.1.3.9.2}
\addcontentsline{toc}{subsubsection}{II.1.3.9.2 [CT2] Clinical Collaboration}

{\small Historical predecessor: R.2.170 \hyperref[element:c5adc1955b6046d8b920b0176cebc9de]{[CT2] Clinical Collaboration}; Canonical ID: \texttt{\seqsplit{c5adc1955b6046d8b920b0176cebc9de}}. Superseded historical identity — direct correspondence · superseded.}\par
\label{element:e6429225463a44d39f82e490e254e5d4}
\index{[CT2] Clinical Collaboration (II.1.3.9.2)}
{\small Capability · Canonical · current}\par

Clinical Collaboration provides secure multidisciplinary spaces bound to clinical context, including governed access to longitudinal information during discussion. Its purpose is contextual conversation between care participants rather than a separate source of authoritative clinical facts. A statement in a conversation remains discourse until explicitly authored through the governed record boundary.



\textbf{Overview}



Clinical Collaboration provides secure multidisciplinary spaces bound to clinical context, including governed access to longitudinal information during discussion. Its purpose is contextual conversation between care participants rather than a separate source of authoritative clinical facts. A statement in a conversation remains discourse until explicitly authored through the governed record boundary.



\textbf{Intent}



Clinical Collaboration provides secure multidisciplinary spaces bound to clinical context, including governed access to longitudinal information during discussion. Its purpose is contextual conversation between care participants rather than a separate source of authoritative clinical facts. A statement in a conversation remains discourse until explicitly authored through the governed record boundary.



\textbf{Boundary}



Information administration does not confer authority over the represented person, organisation, service, place, device or clinical activity.



\medskip\hrule\medskip

{\footnotesize \href{../../../../continuum/strategy/capability/e6429225463a44d39f82e490e254e5d4.md}{Canonical record: e6429225463a44d39f82e490e254e5d4}; full qualifications and provenance.}\par
\paragraph*{II.1.3.9.2.1 [CT3] Clinical Context Collaboration}
\phantomsection
\label{navigation:II.1.3.9.2.1}
\addcontentsline{toc}{paragraph}{II.1.3.9.2.1 [CT3] Clinical Context Collaboration}

{\small Historical predecessor: R.2.169 \hyperref[element:750ef0d8e1504a8e9a1fde9bf13a6d0c]{[CT3] Clinical Context Collaboration}; Canonical ID: \texttt{\seqsplit{750ef0d8e1504a8e9a1fde9bf13a6d0c}}. Superseded historical identity — direct correspondence · superseded.}\par
\label{element:f11617eed6ff4d6fa237644ec95b28b2}
\index{[CT3] Clinical Context Collaboration (II.1.3.9.2.1)}
{\small Capability · Proposed · current}\par

Enable governed multidisciplinary collaboration with clinical context while preserving the distinction between discourse and authored clinical records.



\textbf{Overview}



Enable governed multidisciplinary collaboration with clinical context while preserving the distinction between discourse and authored clinical records. The reviewed hierarchy assigns this enduring ability to its stated business-enabling context.



\textbf{Intent}



Enable governed multidisciplinary collaboration with clinical context while preserving the distinction between discourse and authored clinical records.



\textbf{Boundary}



Information administration and coordination preserve the authority of the responsible source, clinical actor and governed domain; they confer no additional real-world management authority.


{\footnotesize \href{../../../../continuum/strategy/capability/f11617eed6ff4d6fa237644ec95b28b2.md}{Canonical record: f11617eed6ff4d6fa237644ec95b28b2}; full qualifications and provenance.}\par
\subparagraph*{II.1.3.9.2.1.1 [Feature] Care-Team Clinical Collaboration Coordination}
\phantomsection
\label{navigation:II.1.3.9.2.1.1}
\addcontentsline{toc}{subparagraph}{II.1.3.9.2.1.1 [Feature] Care-Team Clinical Collaboration Coordination}

\label{element:b0860d6b68f34b8cb2f93b25cdf2aa97}
\index{[Feature] Care-Team Clinical Collaboration Coordination (II.1.3.9.2.1.1)}
{\small Capability / Feature · Canonical · current}\par

Coordinate secure multidisciplinary clinical communication and consultation exchange bound to patient encounters or clinical topics.



Clinical collaboration coordination binds secure multidisciplinary exchange to the relevant encounter or clinical topic. It supports conversation among authorised participants without making discourse an authoritative observation or acquiring the clinical handover authority excluded by the metamodel.


{\footnotesize \href{../../../../continuum/strategy/capability/feature/b0860d6b68f34b8cb2f93b25cdf2aa97.md}{Canonical record: b0860d6b68f34b8cb2f93b25cdf2aa97}; full qualifications and provenance.}\par
\subparagraph*{II.1.3.9.2.1.2 [Feature] Collaboration Interaction Capture Qualification}
\phantomsection
\label{navigation:II.1.3.9.2.1.2}
\addcontentsline{toc}{subparagraph}{II.1.3.9.2.1.2 [Feature] Collaboration Interaction Capture Qualification}

\label{element:8e2a760f15bd454c8aba1526ee69740c}
\index{[Feature] Collaboration Interaction Capture Qualification (II.1.3.9.2.1.2)}
{\small Capability / Feature · Proposed · current}\par

Qualify when a defined collaboration interaction creates a new governed information or activity consequence.



\textbf{Overview}



A message may convey existing information or an explicit interaction may derive a new managed consequence. The distinction preserves collaboration-native records while routing governed consequences to their accountable responsibility.



\textbf{Boundary}



Conversation presence, display or read receipt does not itself make a clinical decision, approval or final record.


{\footnotesize \href{../../../../continuum/strategy/capability/feature/8e2a760f15bd454c8aba1526ee69740c.md}{Canonical record: 8e2a760f15bd454c8aba1526ee69740c}; full qualifications and provenance.}\par
\subparagraph*{II.1.3.9.2.1.3 [Feature] In-Conversation LHR Query Resolution}
\phantomsection
\label{navigation:II.1.3.9.2.1.3}
\addcontentsline{toc}{subparagraph}{II.1.3.9.2.1.3 [Feature] In-Conversation LHR Query Resolution}

\label{element:d4437b477aab42c79b5f45150c5d7f3f}
\index{[Feature] In-Conversation LHR Query Resolution (II.1.3.9.2.1.3)}
{\small Capability / Feature · Canonical · current}\par

Resolve clinical data queries within an authorized collaboration context, presenting verified clinical summaries directly to care team members.



In-conversation query supplies governed clinical summaries within an authorised collaboration context. The convenience of retrieving information during discussion does not weaken consent or access controls, and displaying a summary does not promote the surrounding conversation into the clinical record.


{\footnotesize \href{../../../../continuum/strategy/capability/feature/d4437b477aab42c79b5f45150c5d7f3f.md}{Canonical record: d4437b477aab42c79b5f45150c5d7f3f}; full qualifications and provenance.}\par
\subparagraph*{II.1.3.9.2.1.4 [Feature] Zero-PHI Collaboration Metadata Governance}
\phantomsection
\label{navigation:II.1.3.9.2.1.4}
\addcontentsline{toc}{subparagraph}{II.1.3.9.2.1.4 [Feature] Zero-PHI Collaboration Metadata Governance}

\label{element:87e9246c944a40d9af5c4d0a3a7d4e9a}
\index{[Feature] Zero-PHI Collaboration Metadata Governance (II.1.3.9.2.1.4)}
{\small Capability / Feature · Canonical · current}\par

Ensure collaboration channel titles, topics, and aliases use opaque identifiers free of patient names, MRNs, or sensitive details.



Metadata governance keeps patient names, identifiers and sensitive detail out of collaboration titles, topics and aliases. Opaque references preserve useful contextual association without placing PHI in broadly visible channel metadata or diagnostic surfaces.


{\footnotesize \href{../../../../continuum/strategy/capability/feature/87e9246c944a40d9af5c4d0a3a7d4e9a.md}{Canonical record: 87e9246c944a40d9af5c4d0a3a7d4e9a}; full qualifications and provenance.}\par
\subsubsection*{II.1.3.9.3 [CT2] Clinical Information Stewardship Enablement [superseded]}
\phantomsection
\label{navigation:II.1.3.9.3}
\addcontentsline{toc}{subsubsection}{II.1.3.9.3 [CT2] Clinical Information Stewardship Enablement [superseded]}

\label{element:3811c758ff2140c39a801bf89f75098a}
\index{[CT2] Clinical Information Stewardship Enablement (II.1.3.9.3)}
{\small Capability · Proposed · superseded}\par
\href{../../../../continuum/strategy/capability/3811c758ff2140c39a801bf89f75098a.md}{Historical canonical definition and provenance}.\par
\paragraph*{II.1.3.9.3.1 [CT3] Information Authority Representation}
\phantomsection
\label{navigation:II.1.3.9.3.1}
\addcontentsline{toc}{paragraph}{II.1.3.9.3.1 [CT3] Information Authority Representation}

\label{element:8c850faafb334f97a81a0b4cd8fdb65f}
\index{[CT3] Information Authority Representation (II.1.3.9.3.1)}
{\small Capability · Proposed · pending\_adjudication}\par

Represent the source, credibility and qualified authority standing of competing health-information assertions.



\textbf{Overview}



A persisted assertion may be authoritative, informational or anecdotal. This ability carries that standing into interpretation without treating transport reliability, freshness or successful commit as proof of credibility.



\textbf{Intent}



Keep an assertion technical acceptance separate from the standing of the source that asserted it.



\textbf{Boundary}



Classification does not originate clinical facts or universally adjudicate conflicts. Policy determines how qualified assertions are used.


{\small Recorded unresolved matters: source-qualified matters; exact scope remains in the canonical record.}\par
{\footnotesize \href{../../../../continuum/strategy/capability/8c850faafb334f97a81a0b4cd8fdb65f.md}{Canonical record: 8c850faafb334f97a81a0b4cd8fdb65f}; full qualifications and provenance.}\par
\subparagraph*{II.1.3.9.3.1.1 [Feature] Information Authority \& Source Standing Representation}
\phantomsection
\label{navigation:II.1.3.9.3.1.1}
\addcontentsline{toc}{subparagraph}{II.1.3.9.3.1.1 [Feature] Information Authority \& Source Standing Representation}

\label{element:70a444156cd246cb97608a782e4bdea8}
\index{[Feature] Information Authority \& Source Standing Representation (II.1.3.9.3.1.1)}
{\small Capability / Feature · Proposed · pending\_adjudication}\par

Attach and preserve qualified source-authority and credibility standing for managed assertions.



\textbf{Overview}



The user of an assertion needs to distinguish authoritative, informational and anecdotal standing even when each was technically valid and durably accepted.



\textbf{Boundary}



Authority classification is not a clinical truth judgement or a universal rule for choosing competing assertions.


{\small Recorded unresolved matters: source-qualified matters; exact scope remains in the canonical record.}\par
{\footnotesize \href{../../../../continuum/strategy/capability/feature/70a444156cd246cb97608a782e4bdea8.md}{Canonical record: 70a444156cd246cb97608a782e4bdea8}; full qualifications and provenance.}\par
\subsubsection*{II.1.3.9.4 [CT2] Clinical Knowledge Services}
\phantomsection
\label{navigation:II.1.3.9.4}
\addcontentsline{toc}{subsubsection}{II.1.3.9.4 [CT2] Clinical Knowledge Services}

{\small Historical predecessor: R.2.174 \hyperref[element:5b8478b08c3f42c0a3480edb8f0b9ea7]{[CT2] Clinical Knowledge Services}; Canonical ID: \texttt{\seqsplit{5b8478b08c3f42c0a3480edb8f0b9ea7}}. Superseded historical identity — direct correspondence · superseded.}\par
\label{element:bbd2bd21859f40e39963bdbe7a851686}
\index{[CT2] Clinical Knowledge Services (II.1.3.9.4)}
{\small Capability · Canonical · current}\par

Administer clinical knowledge requests, distribute governed knowledge artefacts and administer approved reference datasets.



\textbf{Overview}



Administer clinical knowledge requests, distribute governed knowledge artefacts and administer approved reference datasets.



\textbf{Intent}



Administer clinical knowledge requests, distribute governed knowledge artefacts and administer approved reference datasets.



\textbf{Boundary}



Information administration does not confer authority over the represented person, organisation, service, place, device or clinical activity.



\medskip\hrule\medskip

{\footnotesize \href{../../../../continuum/strategy/capability/bbd2bd21859f40e39963bdbe7a851686.md}{Canonical record: bbd2bd21859f40e39963bdbe7a851686}; full qualifications and provenance.}\par
\paragraph*{II.1.3.9.4.1 [CT3] Clinical Knowledge Artefact Distribution Management}
\phantomsection
\label{navigation:II.1.3.9.4.1}
\addcontentsline{toc}{paragraph}{II.1.3.9.4.1 [CT3] Clinical Knowledge Artefact Distribution Management}

{\small Historical predecessor: R.2.171 \hyperref[element:d81e0d6d76014066a682e3d9e03e5f8a]{[CT3] Clinical Knowledge Artefact Distribution Management}; Canonical ID: \texttt{\seqsplit{d81e0d6d76014066a682e3d9e03e5f8a}}. Superseded historical identity — direct correspondence · superseded.}\par
\label{element:cd29a9eb798949bc8bc81bb5d434cabf}
\index{[CT3] Clinical Knowledge Artefact Distribution Management (II.1.3.9.4.1)}
{\small Capability · Proposed · current}\par

Manage distribution of governed clinical knowledge artefacts to intended consumers under their applicable use contracts.



\textbf{Overview}



Manage distribution of governed clinical knowledge artefacts to intended consumers under their applicable use contracts. The reviewed hierarchy assigns this enduring ability to its stated business-enabling context.



\textbf{Intent}



Manage distribution of governed clinical knowledge artefacts to intended consumers under their applicable use contracts.



\textbf{Boundary}



Information administration and coordination preserve the authority of the responsible source, clinical actor and governed domain; they confer no additional real-world management authority.


{\footnotesize \href{../../../../continuum/strategy/capability/cd29a9eb798949bc8bc81bb5d434cabf.md}{Canonical record: cd29a9eb798949bc8bc81bb5d434cabf}; full qualifications and provenance.}\par
\paragraph*{II.1.3.9.4.2 [CT3] Clinical Knowledge Request Administration}
\phantomsection
\label{navigation:II.1.3.9.4.2}
\addcontentsline{toc}{paragraph}{II.1.3.9.4.2 [CT3] Clinical Knowledge Request Administration}

{\small Historical predecessor: R.2.172 \hyperref[element:eb370c01fbdb4d22956d97e9ecd74b57]{[CT3] Clinical Knowledge Request Administration}; Canonical ID: \texttt{\seqsplit{eb370c01fbdb4d22956d97e9ecd74b57}}. Superseded historical identity — direct correspondence · superseded.}\par
\label{element:543c53483a2249d78044c6635b7b7c48}
\index{[CT3] Clinical Knowledge Request Administration (II.1.3.9.4.2)}
{\small Capability · Proposed · current}\par

Administer clinical knowledge requests and their qualified resolution to applicable governed knowledge or reference content.



\textbf{Overview}



Administer clinical knowledge requests and their qualified resolution to applicable governed knowledge or reference content. The reviewed hierarchy assigns this enduring ability to its stated business-enabling context.



\textbf{Intent}



Administer clinical knowledge requests and their qualified resolution to applicable governed knowledge or reference content.



\textbf{Boundary}



Information administration and coordination preserve the authority of the responsible source, clinical actor and governed domain; they confer no additional real-world management authority.


{\footnotesize \href{../../../../continuum/strategy/capability/543c53483a2249d78044c6635b7b7c48.md}{Canonical record: 543c53483a2249d78044c6635b7b7c48}; full qualifications and provenance.}\par
\subparagraph*{II.1.3.9.4.2.1 [Feature] Bidirectional Terminology Mapping}
\phantomsection
\label{navigation:II.1.3.9.4.2.1}
\addcontentsline{toc}{subparagraph}{II.1.3.9.4.2.1 [Feature] Bidirectional Terminology Mapping}

\label{element:6134dc91af8d45e8b81cc28d18a5251f}
\index{[Feature] Bidirectional Terminology Mapping (II.1.3.9.4.2.1)}
{\small Capability / Feature · Canonical · current}\par

Translate between legacy local code systems and canonical national healthcare profiles.



Terminology mapping applies governed correspondence between local and canonical semantic systems. Translation preserves the limitations and authority of the map; it does not justify silently treating related source concepts as identical or clinically adjudicating the underlying information.


{\footnotesize \href{../../../../continuum/strategy/capability/feature/6134dc91af8d45e8b81cc28d18a5251f.md}{Canonical record: 6134dc91af8d45e8b81cc28d18a5251f}; full qualifications and provenance.}\par
\subparagraph*{II.1.3.9.4.2.2 [Feature] Clinical Concept Resolution}
\phantomsection
\label{navigation:II.1.3.9.4.2.2}
\addcontentsline{toc}{subparagraph}{II.1.3.9.4.2.2 [Feature] Clinical Concept Resolution}

\label{element:30cdb9b2a810467c9b56e9870de4367e}
\index{[Feature] Clinical Concept Resolution (II.1.3.9.4.2.2)}
{\small Capability / Feature · Canonical · current}\par

Resolve local clinical codes against standard terminologies (SNOMED CT, LOINC, ICD-10).



Concept resolution makes a coded clinical reference interpretable against its governed terminology. It validates the reference context without making a recognised code proof of clinical truth or giving the consuming system authority to redefine the concept.


{\footnotesize \href{../../../../continuum/strategy/capability/feature/30cdb9b2a810467c9b56e9870de4367e.md}{Canonical record: 30cdb9b2a810467c9b56e9870de4367e}; full qualifications and provenance.}\par
\paragraph*{II.1.3.9.4.3 [CT3] Governed Reference Dataset Administration}
\phantomsection
\label{navigation:II.1.3.9.4.3}
\addcontentsline{toc}{paragraph}{II.1.3.9.4.3 [CT3] Governed Reference Dataset Administration}

{\small Historical predecessor: R.2.173 \hyperref[element:15d972931ed6447797c30a5ddb71737e]{[CT3] Governed Reference Dataset Administration}; Canonical ID: \texttt{\seqsplit{15d972931ed6447797c30a5ddb71737e}}. Superseded historical identity — direct correspondence · superseded.}\par
\label{element:682b8684ad9e49d5a5c348a148081607}
\index{[CT3] Governed Reference Dataset Administration (II.1.3.9.4.3)}
{\small Capability · Proposed · current}\par

Administer approved runtime reference-dataset content and revisions separately from the reusable models defining its meaning.



\textbf{Overview}



Administer approved runtime reference-dataset content and revisions separately from the reusable models defining its meaning.



\textbf{Intent}



Administer approved runtime reference-dataset content and revisions separately from the reusable models defining its meaning.



\textbf{Boundary}



Does not own external ontology authority, execute authorisation policy or establish durable truth independently of the normal persistence boundary.


{\footnotesize \href{../../../../continuum/strategy/capability/682b8684ad9e49d5a5c348a148081607.md}{Canonical record: 682b8684ad9e49d5a5c348a148081607}; full qualifications and provenance.}\par
\subparagraph*{II.1.3.9.4.3.1 [Feature] Governed Dataset Content Revision}
\phantomsection
\label{navigation:II.1.3.9.4.3.1}
\addcontentsline{toc}{subparagraph}{II.1.3.9.4.3.1 [Feature] Governed Dataset Content Revision}

\label{element:89123c868b8a46c38800916110cd918d}
\index{[Feature] Governed Dataset Content Revision (II.1.3.9.4.3.1)}
{\small Capability / Feature · Proposed · current}\par

Review, accept and expose authorised changes to mutable reference content with effective revision standing.



\textbf{Overview}



A model definition can remain stable while approved terminology membership or mapping content changes. This Feature represents that operational governance without making every definition runtime mutable.



\textbf{Boundary}



Does not invent an approval authority or bypass normal durable establishment.


{\footnotesize \href{../../../../continuum/strategy/capability/feature/89123c868b8a46c38800916110cd918d.md}{Canonical record: 89123c868b8a46c38800916110cd918d}; full qualifications and provenance.}\par
\subsubsection*{II.1.3.9.5 [CT2] General Collaboration}
\phantomsection
\label{navigation:II.1.3.9.5}
\addcontentsline{toc}{subsubsection}{II.1.3.9.5 [CT2] General Collaboration}

{\small Historical predecessor: R.2.176 \hyperref[element:4f58c82aca7a4b54aeb6e2920ce5c12a]{[CT2] General Collaboration}; Canonical ID: \texttt{\seqsplit{4f58c82aca7a4b54aeb6e2920ce5c12a}}. Superseded historical identity — direct correspondence · superseded.}\par
\label{element:3bb740e6db284112b054905c0c6c80be}
\index{[CT2] General Collaboration (II.1.3.9.5)}
{\small Capability · Canonical · current}\par

General Collaboration provides secure operational messaging for departmental, administrative and facility-support participants. It gives non-clinical coordination an interaction context without assuming every discussion requires a patient-linked clinical workspace. Its role is communication enablement rather than ownership of the support activity being discussed.



\textbf{Overview}



General Collaboration provides secure operational messaging for departmental, administrative and facility-support participants. It gives non-clinical coordination an interaction context without assuming every discussion requires a patient-linked clinical workspace. Its role is communication enablement rather than ownership of the support activity being discussed.



\textbf{Intent}



General Collaboration provides secure operational messaging for departmental, administrative and facility-support participants. It gives non-clinical coordination an interaction context without assuming every discussion requires a patient-linked clinical workspace. Its role is communication enablement rather than ownership of the support activity being discussed.



\textbf{Boundary}



Information administration does not confer authority over the represented person, organisation, service, place, device or clinical activity.



\medskip\hrule\medskip

{\footnotesize \href{../../../../continuum/strategy/capability/3bb740e6db284112b054905c0c6c80be.md}{Canonical record: 3bb740e6db284112b054905c0c6c80be}; full qualifications and provenance.}\par
\paragraph*{II.1.3.9.5.1 [CT3] Operational Collaboration}
\phantomsection
\label{navigation:II.1.3.9.5.1}
\addcontentsline{toc}{paragraph}{II.1.3.9.5.1 [CT3] Operational Collaboration}

{\small Historical predecessor: R.2.175 \hyperref[element:99e9d9763ec94140b7d7642f4b5966f6]{[CT3] Operational Collaboration}; Canonical ID: \texttt{\seqsplit{99e9d9763ec94140b7d7642f4b5966f6}}. Superseded historical identity — direct correspondence · superseded.}\par
\label{element:631fdb8b47084008ac9b5de28b26972b}
\index{[CT3] Operational Collaboration (II.1.3.9.5.1)}
{\small Capability · Proposed · current}\par

Enable operational team discussion and collaboration without promoting conversational content to clinical record authority.



\textbf{Overview}



Enable operational team discussion and collaboration without promoting conversational content to clinical record authority. The reviewed hierarchy assigns this enduring ability to its stated business-enabling context.



\textbf{Intent}



Enable operational team discussion and collaboration without promoting conversational content to clinical record authority.



\textbf{Boundary}



Information administration and coordination preserve the authority of the responsible source, clinical actor and governed domain; they confer no additional real-world management authority.


{\footnotesize \href{../../../../continuum/strategy/capability/631fdb8b47084008ac9b5de28b26972b.md}{Canonical record: 631fdb8b47084008ac9b5de28b26972b}; full qualifications and provenance.}\par
\subparagraph*{II.1.3.9.5.1.1 [Feature] Operational Team Chat Enablement}
\phantomsection
\label{navigation:II.1.3.9.5.1.1}
\addcontentsline{toc}{subparagraph}{II.1.3.9.5.1.1 [Feature] Operational Team Chat Enablement}

\label{element:367a390e7a464003beba189428c5faca}
\index{[Feature] Operational Team Chat Enablement (II.1.3.9.5.1.1)}
{\small Capability / Feature · Canonical · current}\par

Provide secure instant messaging channels for departmental, administrative, and facility support staff.



Operational chat provides secure communication for administrative and support teams. It enables participants to discuss work while leaving assignment, completion and the authority of operational decisions with the capabilities responsible for that activity.


{\footnotesize \href{../../../../continuum/strategy/capability/feature/367a390e7a464003beba189428c5faca.md}{Canonical record: 367a390e7a464003beba189428c5faca}; full qualifications and provenance.}\par
\subsubsection*{II.1.3.9.6 [CT2] Health Information Access}
\phantomsection
\label{navigation:II.1.3.9.6}
\addcontentsline{toc}{subsubsection}{II.1.3.9.6 [CT2] Health Information Access}

{\small Historical predecessor: R.2.178 \hyperref[element:5019ae9f2cdc45a1b2de97bc0f5a313d]{[CT2] Health Information Access}; Canonical ID: \texttt{\seqsplit{5019ae9f2cdc45a1b2de97bc0f5a313d}}. Superseded historical identity — direct correspondence · superseded.}\par
\label{element:cc253aea3c5f489da0502df74dde1669}
\index{[CT2] Health Information Access (II.1.3.9.6)}
{\small Capability · Canonical · current}\par

Health Information Access enables authorised discovery and retrieval across clinical information. It combines semantic search and discrete retrieval with access-qualified filtering so that a discoverable item is not automatically disclosable. It consumes established policy and consent determinations rather than making the search interface an alternative security authority.



\textbf{Overview}



Health Information Access enables authorised discovery and retrieval across clinical information. It combines semantic search and discrete retrieval with access-qualified filtering so that a discoverable item is not automatically disclosable. It consumes established policy and consent determinations rather than making the search interface an alternative security authority.



\textbf{Intent}



Health Information Access enables authorised discovery and retrieval across clinical information. It combines semantic search and discrete retrieval with access-qualified filtering so that a discoverable item is not automatically disclosable. It consumes established policy and consent determinations rather than making the search interface an alternative security authority.



\textbf{Boundary}



Information administration does not confer authority over the represented person, organisation, service, place, device or clinical activity.



\medskip\hrule\medskip

{\footnotesize \href{../../../../continuum/strategy/capability/cc253aea3c5f489da0502df74dde1669.md}{Canonical record: cc253aea3c5f489da0502df74dde1669}; full qualifications and provenance.}\par
\paragraph*{II.1.3.9.6.1 [CT3] Governed Health Information Retrieval}
\phantomsection
\label{navigation:II.1.3.9.6.1}
\addcontentsline{toc}{paragraph}{II.1.3.9.6.1 [CT3] Governed Health Information Retrieval}

{\small Historical predecessor: R.2.177 \hyperref[element:90b838d703dd472b8b6d74961f3a748c]{[CT3] Governed Health Information Retrieval}; Canonical ID: \texttt{\seqsplit{90b838d703dd472b8b6d74961f3a748c}}. Superseded historical identity — direct correspondence · superseded.}\par
\label{element:76e74e46ca4d44a1a7bb116b358f231c}
\index{[CT3] Governed Health Information Retrieval (II.1.3.9.6.1)}
{\small Capability · Proposed · current}\par

Retrieve and search health information and apply access-qualified result filtering under established access and consent determinations.



\textbf{Overview}



Retrieve and search health information and apply access-qualified result filtering under established access and consent determinations. The reviewed hierarchy assigns this enduring ability to its stated business-enabling context.



\textbf{Intent}



Retrieve and search health information and apply access-qualified result filtering under established access and consent determinations.



\textbf{Boundary}



Information administration and coordination preserve the authority of the responsible source, clinical actor and governed domain; they confer no additional real-world management authority.


{\footnotesize \href{../../../../continuum/strategy/capability/76e74e46ca4d44a1a7bb116b358f231c.md}{Canonical record: 76e74e46ca4d44a1a7bb116b358f231c}; full qualifications and provenance.}\par
\subparagraph*{II.1.3.9.6.1.1 [Feature] Access-Controlled Information Filtering}
\phantomsection
\label{navigation:II.1.3.9.6.1.1}
\addcontentsline{toc}{subparagraph}{II.1.3.9.6.1.1 [Feature] Access-Controlled Information Filtering}

\label{element:a9f1d2100faa44beb69add62bb4f4b3f}
\index{[Feature] Access-Controlled Information Filtering (II.1.3.9.6.1.1)}
{\small Capability / Feature · Canonical · current}\par

Filter search results dynamically according to requesting user authorizations and patient consent directives.



Access-qualified filtering restricts the disclosed result to the verified caller's authority and applicable consent. The responsibility is more specific than retrieving a resource; the latter's shared owner and subject matter do not establish a Feature realisation relationship.


{\footnotesize \href{../../../../continuum/strategy/capability/feature/a9f1d2100faa44beb69add62bb4f4b3f.md}{Canonical record: a9f1d2100faa44beb69add62bb4f4b3f}; full qualifications and provenance.}\par
\subparagraph*{II.1.3.9.6.1.2 [Feature] Federated Clinical Query}
\phantomsection
\label{navigation:II.1.3.9.6.1.2}
\addcontentsline{toc}{subparagraph}{II.1.3.9.6.1.2 [Feature] Federated Clinical Query}

\label{element:cfcc4d64af434418a567ccf46d8052dd}
\index{[Feature] Federated Clinical Query (II.1.3.9.6.1.2)}
{\small Capability / Feature · Canonical · current}\par

Execute bounded, indexed queries across distributed clinical resources and directory graphs.



Federated query makes relevant clinical and directory information discoverable across distributed contexts using bounded criteria. Discovery is a governed information-access responsibility rather than authority to disclose every match or reinterpret the clinical meaning of the results.


{\footnotesize \href{../../../../continuum/strategy/capability/feature/cfcc4d64af434418a567ccf46d8052dd.md}{Canonical record: cfcc4d64af434418a567ccf46d8052dd}; full qualifications and provenance.}\par
\subsubsection*{II.1.3.9.7 [CT2] Health Information Communication}
\phantomsection
\label{navigation:II.1.3.9.7}
\addcontentsline{toc}{subsubsection}{II.1.3.9.7 [CT2] Health Information Communication}

{\small Historical predecessor: R.2.183 \hyperref[element:bcefb378424842fcb6d703208c3f8773]{[CT2] Health Information Communication}; Canonical ID: \texttt{\seqsplit{bcefb378424842fcb6d703208c3f8773}}. Superseded historical identity — direct correspondence · superseded.}\par
\label{element:7925605bf60f4b288e826020409a09d9}
\index{[CT2] Health Information Communication (II.1.3.9.7)}
{\small Capability · Proposed · current}\par

Enable patient notifications, patient task communication, patient self-service and access to a patient health library under applicable information-use controls.



\textbf{Overview}



Enable patient notifications, patient task communication, patient self-service and access to a patient health library under applicable information-use controls. The reviewed hierarchy assigns this enduring ability to its stated business-enabling context.



\textbf{Intent}



Enable patient notifications, patient task communication, patient self-service and access to a patient health library under applicable information-use controls.



\textbf{Boundary}



Information administration and coordination preserve the authority of the responsible source, clinical actor and governed domain; they confer no additional real-world management authority.


{\footnotesize \href{../../../../continuum/strategy/capability/7925605bf60f4b288e826020409a09d9.md}{Canonical record: 7925605bf60f4b288e826020409a09d9}; full qualifications and provenance.}\par
\paragraph*{II.1.3.9.7.1 [CT3] Patient Health Library Access}
\phantomsection
\label{navigation:II.1.3.9.7.1}
\addcontentsline{toc}{paragraph}{II.1.3.9.7.1 [CT3] Patient Health Library Access}

{\small Historical predecessor: R.2.179 \hyperref[element:b132224cfa6f435f8f2a66a3878d6df6]{[CT3] Patient Health Library Access}; Canonical ID: \texttt{\seqsplit{b132224cfa6f435f8f2a66a3878d6df6}}. Superseded historical identity — direct correspondence · superseded.}\par
\label{element:808e5a3c3f6b496fa329962088f0bee0}
\index{[CT3] Patient Health Library Access (II.1.3.9.7.1)}
{\small Capability · Proposed · current}\par

Enable authorised patient access to appropriately governed health knowledge and library information.



\textbf{Overview}



Enable authorised patient access to appropriately governed health knowledge and library information. The reviewed hierarchy assigns this enduring ability to its stated business-enabling context.



\textbf{Intent}



Enable authorised patient access to appropriately governed health knowledge and library information.



\textbf{Boundary}



Information administration and coordination preserve the authority of the responsible source, clinical actor and governed domain; they confer no additional real-world management authority.


{\footnotesize \href{../../../../continuum/strategy/capability/808e5a3c3f6b496fa329962088f0bee0.md}{Canonical record: 808e5a3c3f6b496fa329962088f0bee0}; full qualifications and provenance.}\par
\paragraph*{II.1.3.9.7.2 [CT3] Patient Notification Services}
\phantomsection
\label{navigation:II.1.3.9.7.2}
\addcontentsline{toc}{paragraph}{II.1.3.9.7.2 [CT3] Patient Notification Services}

{\small Historical predecessor: R.2.180 \hyperref[element:35c5be472bf2433fbbde9646d87aefb3]{[CT3] Patient Notification Services}; Canonical ID: \texttt{\seqsplit{35c5be472bf2433fbbde9646d87aefb3}}. Superseded historical identity — direct correspondence · superseded.}\par
\label{element:9d172cd511b742499c4e7c23f6630a35}
\index{[CT3] Patient Notification Services (II.1.3.9.7.2)}
{\small Capability · Proposed · current}\par

Enable governed patient notifications concerning health information and applicable care or operational events.



\textbf{Overview}



Enable governed patient notifications concerning health information and applicable care or operational events. The reviewed hierarchy assigns this enduring ability to its stated business-enabling context.



\textbf{Intent}



Enable governed patient notifications concerning health information and applicable care or operational events.



\textbf{Boundary}



Information administration and coordination preserve the authority of the responsible source, clinical actor and governed domain; they confer no additional real-world management authority.


{\footnotesize \href{../../../../continuum/strategy/capability/9d172cd511b742499c4e7c23f6630a35.md}{Canonical record: 9d172cd511b742499c4e7c23f6630a35}; full qualifications and provenance.}\par
\paragraph*{II.1.3.9.7.3 [CT3] Patient Self-Service}
\phantomsection
\label{navigation:II.1.3.9.7.3}
\addcontentsline{toc}{paragraph}{II.1.3.9.7.3 [CT3] Patient Self-Service}

{\small Historical predecessor: R.2.181 \hyperref[element:60f427e6b5764757b909624085a2af89]{[CT3] Patient Self-Service}; Canonical ID: \texttt{\seqsplit{60f427e6b5764757b909624085a2af89}}. Superseded historical identity — direct correspondence · superseded.}\par
\label{element:d4956cab8b2b46c396899949c9486a1c}
\index{[CT3] Patient Self-Service (II.1.3.9.7.3)}
{\small Capability · Proposed · current}\par

Enable patients to undertake permitted information-access and administrative self-service interactions.



\textbf{Overview}



Enable patients to undertake permitted information-access and administrative self-service interactions. The reviewed hierarchy assigns this enduring ability to its stated business-enabling context.



\textbf{Intent}



Enable patients to undertake permitted information-access and administrative self-service interactions.



\textbf{Boundary}



Information administration and coordination preserve the authority of the responsible source, clinical actor and governed domain; they confer no additional real-world management authority.


{\footnotesize \href{../../../../continuum/strategy/capability/d4956cab8b2b46c396899949c9486a1c.md}{Canonical record: d4956cab8b2b46c396899949c9486a1c}; full qualifications and provenance.}\par
\paragraph*{II.1.3.9.7.4 [CT3] Patient Task Communication Services}
\phantomsection
\label{navigation:II.1.3.9.7.4}
\addcontentsline{toc}{paragraph}{II.1.3.9.7.4 [CT3] Patient Task Communication Services}

{\small Historical predecessor: R.2.182 \hyperref[element:0914f2bb267449e282e9ae0f7407d9fe]{[CT3] Patient Task Communication Services}; Canonical ID: \texttt{\seqsplit{0914f2bb267449e282e9ae0f7407d9fe}}. Superseded historical identity — direct correspondence · superseded.}\par
\label{element:ba8ce1d54ace4bc9aa3521dd954f36cb}
\index{[CT3] Patient Task Communication Services (II.1.3.9.7.4)}
{\small Capability · Proposed · current}\par

Communicate patient-facing work requests and their reported progression under the applicable work and disclosure contracts.



\textbf{Overview}



Communicate patient-facing work requests and their reported progression under the applicable work and disclosure contracts. The reviewed hierarchy assigns this enduring ability to its stated business-enabling context.



\textbf{Intent}



Communicate patient-facing work requests and their reported progression under the applicable work and disclosure contracts.



\textbf{Boundary}



Information administration and coordination preserve the authority of the responsible source, clinical actor and governed domain; they confer no additional real-world management authority.


{\footnotesize \href{../../../../continuum/strategy/capability/ba8ce1d54ace4bc9aa3521dd954f36cb.md}{Canonical record: ba8ce1d54ace4bc9aa3521dd954f36cb}; full qualifications and provenance.}\par
\subsubsection*{II.1.3.9.8 [CT2] Health Information Connectivity Enablement [superseded]}
\phantomsection
\label{navigation:II.1.3.9.8}
\addcontentsline{toc}{subsubsection}{II.1.3.9.8 [CT2] Health Information Connectivity Enablement [superseded]}

\label{element:6a1030483dfc4ae6af6998b587129d1c}
\index{[CT2] Health Information Connectivity Enablement (II.1.3.9.8)}
{\small Capability · Proposed · superseded}\par
\href{../../../../continuum/strategy/capability/6a1030483dfc4ae6af6998b587129d1c.md}{Historical canonical definition and provenance}.\par
\paragraph*{II.1.3.9.8.1 [CT3] Health Information Communication}
\phantomsection
\label{navigation:II.1.3.9.8.1}
\addcontentsline{toc}{paragraph}{II.1.3.9.8.1 [CT3] Health Information Communication}

\label{element:9794b2aff43443b495a19aefe0efa67b}
\index{[CT3] Health Information Communication (II.1.3.9.8.1)}
{\small Capability · Canonical · pending\_adjudication}\par

Definition unresolved in the baseline. No definition has been synthesised from the name.



\medskip\hrule\medskip

{\small Recorded unresolved matters: source-qualified matters; exact scope remains in the canonical record.}\par
{\footnotesize \href{../../../../continuum/strategy/capability/9794b2aff43443b495a19aefe0efa67b.md}{Canonical record: 9794b2aff43443b495a19aefe0efa67b}; full qualifications and provenance.}\par
\subparagraph*{II.1.3.9.8.1.1 [Feature] Standards-Based Boundary Exchange}
\phantomsection
\label{navigation:II.1.3.9.8.1.1}
\addcontentsline{toc}{subparagraph}{II.1.3.9.8.1.1 [Feature] Standards-Based Boundary Exchange}

\label{element:3192ba1541ff41458d176184119c30bc}
\index{[Feature] Standards-Based Boundary Exchange (II.1.3.9.8.1.1)}
{\small Capability / Feature · Canonical · pending\_adjudication}\par

Provide standards-conformant boundary exchange for healthcare messaging and resource representations.



\medskip\hrule\medskip

{\small Recorded unresolved matters: source-qualified matters; exact scope remains in the canonical record.}\par
{\footnotesize \href{../../../../continuum/strategy/capability/feature/3192ba1541ff41458d176184119c30bc.md}{Canonical record: 3192ba1541ff41458d176184119c30bc}; full qualifications and provenance.}\par
\subsubsection*{II.1.3.9.9 [CT2] Health Information Control}
\phantomsection
\label{navigation:II.1.3.9.9}
\addcontentsline{toc}{subsubsection}{II.1.3.9.9 [CT2] Health Information Control}

{\small Historical predecessor: R.2.187 \hyperref[element:a3196aef21d443519de874c24f64a185]{[CT2] Health Information Control}; Canonical ID: \texttt{\seqsplit{a3196aef21d443519de874c24f64a185}}. Superseded historical identity — direct correspondence · superseded.}\par
\label{element:44cc0e17ad534675bf26dce6e46a68ff}
\index{[CT2] Health Information Control (II.1.3.9.9)}
{\small Capability · Canonical · current}\par

Govern patient health-information restrictions, filtering and disclosure using established information-access authority, client privacy/consent determinations and applicable security controls.



\textbf{Overview}



Govern patient health-information restrictions, filtering and disclosure using established information-access authority, client privacy/consent determinations and applicable security controls.



\textbf{Intent}



Govern patient health-information restrictions, filtering and disclosure using established information-access authority, client privacy/consent determinations and applicable security controls.



\textbf{Boundary}



Information administration does not confer authority over the represented person, organisation, service, place, device or clinical activity. Client consent directives remain owned by Client Privacy Guardrail Management (FEAT-EM-09); this capability references and respects those determinations without defining a second consent owner.



\medskip\hrule\medskip

{\footnotesize \href{../../../../continuum/strategy/capability/44cc0e17ad534675bf26dce6e46a68ff.md}{Canonical record: 44cc0e17ad534675bf26dce6e46a68ff}; full qualifications and provenance.}\par
\paragraph*{II.1.3.9.9.1 [CT3] Patient Health Information Disclosure Management}
\phantomsection
\label{navigation:II.1.3.9.9.1}
\addcontentsline{toc}{paragraph}{II.1.3.9.9.1 [CT3] Patient Health Information Disclosure Management}

{\small Historical predecessor: R.2.184 \hyperref[element:07cf507f183d4905a7ad2adbb8555df5]{[CT3] Patient Health Information Disclosure Management}; Canonical ID: \texttt{\seqsplit{07cf507f183d4905a7ad2adbb8555df5}}. Superseded historical identity — direct correspondence · superseded.}\par
\label{element:7f918639e23842e78fa92f1bf47c5a7b}
\index{[CT3] Patient Health Information Disclosure Management (II.1.3.9.9.1)}
{\small Capability · Proposed · current}\par

Manage permitted patient health-information disclosure and its applicable recipient, purpose and restriction conditions.



\textbf{Overview}



Manage permitted patient health-information disclosure and its applicable recipient, purpose and restriction conditions. The reviewed hierarchy assigns this enduring ability to its stated business-enabling context.



\textbf{Intent}



Manage permitted patient health-information disclosure and its applicable recipient, purpose and restriction conditions.



\textbf{Boundary}



Information administration and coordination preserve the authority of the responsible source, clinical actor and governed domain; they confer no additional real-world management authority. Client consent directives remain owned by Client Privacy Guardrail Management (FEAT-EM-09); this capability references and respects those determinations without defining a second consent owner.


{\footnotesize \href{../../../../continuum/strategy/capability/7f918639e23842e78fa92f1bf47c5a7b.md}{Canonical record: 7f918639e23842e78fa92f1bf47c5a7b}; full qualifications and provenance.}\par
\paragraph*{II.1.3.9.9.2 [CT3] Patient Health Information Filter Management}
\phantomsection
\label{navigation:II.1.3.9.9.2}
\addcontentsline{toc}{paragraph}{II.1.3.9.9.2 [CT3] Patient Health Information Filter Management}

{\small Historical predecessor: R.2.185 \hyperref[element:589f192797454d9e8960fb63774e5e9c]{[CT3] Patient Health Information Filter Management}; Canonical ID: \texttt{\seqsplit{589f192797454d9e8960fb63774e5e9c}}. Superseded historical identity — direct correspondence · superseded.}\par
\label{element:62f673613c76494eba972cad3e4e8d0b}
\index{[CT3] Patient Health Information Filter Management (II.1.3.9.9.2)}
{\small Capability · Proposed · current}\par

Manage patient health-information filtering rules and qualified projections while respecting established client privacy and consent determinations.



\textbf{Overview}



Manage patient health-information filtering rules and qualified projections while respecting established client privacy and consent determinations. The reviewed hierarchy assigns this enduring ability to its stated business-enabling context.



\textbf{Intent}



Manage patient health-information filtering rules and qualified projections while respecting established client privacy and consent determinations.



\textbf{Boundary}



Information administration and coordination preserve the authority of the responsible source, clinical actor and governed domain; they confer no additional real-world management authority. Client consent directives remain owned by Client Privacy Guardrail Management (FEAT-EM-09); this capability references and respects those determinations without defining a second consent owner.


{\footnotesize \href{../../../../continuum/strategy/capability/62f673613c76494eba972cad3e4e8d0b.md}{Canonical record: 62f673613c76494eba972cad3e4e8d0b}; full qualifications and provenance.}\par
\paragraph*{II.1.3.9.9.3 [CT3] Patient Health Information Restriction Management}
\phantomsection
\label{navigation:II.1.3.9.9.3}
\addcontentsline{toc}{paragraph}{II.1.3.9.9.3 [CT3] Patient Health Information Restriction Management}

{\small Historical predecessor: R.2.186 \hyperref[element:5ad91d045474435a96b560e66ff7de05]{[CT3] Patient Health Information Restriction Management}; Canonical ID: \texttt{\seqsplit{5ad91d045474435a96b560e66ff7de05}}. Superseded historical identity — direct correspondence · superseded.}\par
\label{element:262366990bda4ac4ace5429269cf6a98}
\index{[CT3] Patient Health Information Restriction Management (II.1.3.9.9.3)}
{\small Capability · Proposed · current}\par

Manage application of patient health-information restrictions using the client-associated privacy and consent responsibilities retained by Client Privacy Guardrail Management.



\textbf{Overview}



Manage application of patient health-information restrictions using the client-associated privacy and consent responsibilities retained by Client Privacy Guardrail Management. The reviewed hierarchy assigns this enduring ability to its stated business-enabling context.



\textbf{Intent}



Manage application of patient health-information restrictions using the client-associated privacy and consent responsibilities retained by Client Privacy Guardrail Management.



\textbf{Boundary}



Information administration and coordination preserve the authority of the responsible source, clinical actor and governed domain; they confer no additional real-world management authority. Client consent directives remain owned by Client Privacy Guardrail Management (FEAT-EM-09); this capability references and respects those determinations without defining a second consent owner.


{\footnotesize \href{../../../../continuum/strategy/capability/262366990bda4ac4ace5429269cf6a98.md}{Canonical record: 262366990bda4ac4ace5429269cf6a98}; full qualifications and provenance.}\par
\subparagraph*{II.1.3.9.9.3.1 [Feature] Health Information Access \& Consent Control}
\phantomsection
\label{navigation:II.1.3.9.9.3.1}
\addcontentsline{toc}{subparagraph}{II.1.3.9.9.3.1 [Feature] Health Information Access \& Consent Control}

\label{element:e2643fa59303401abb1b778f3295e536}
\index{[Feature] Health Information Access \& Consent Control (II.1.3.9.9.3.1)}
{\small Capability / Feature · Canonical · current}\par

Evaluate patient consent preferences, practitioner authority boundaries, and access control policies for managed clinical records.



Access and consent control evaluates whether an information operation is legitimate in its practitioner, patient and purpose context. It supplies policy determinations consumed by access and presentation paths rather than allowing those consumers to establish authority from caller-supplied content.


{\footnotesize \href{../../../../continuum/strategy/capability/feature/e2643fa59303401abb1b778f3295e536.md}{Canonical record: e2643fa59303401abb1b778f3295e536}; full qualifications and provenance.}\par
\paragraph*{II.1.3.9.9.4 [Feature] Governed Policy Assignment \& Revocation}
\phantomsection
\label{navigation:II.1.3.9.9.4}
\addcontentsline{toc}{paragraph}{II.1.3.9.9.4 [Feature] Governed Policy Assignment \& Revocation}

\label{element:d6c271ae95ea44449c5e22665d1d9bc2}
\index{[Feature] Governed Policy Assignment \& Revocation (II.1.3.9.9.4)}
{\small Capability / Feature · Proposed · pending\_adjudication}\par

Manage approved permission criteria and identity/role assignments, including temporal applicability and revocation.



\textbf{Overview}



An assignment becoming inapplicable by expiry differs from an explicit revocation; both must affect the effective authority used for operations.



\textbf{Boundary}



Management does not grant unconstrained administrative authority or derive permission from resource tags.


{\small Recorded unresolved matters: source-qualified matters; exact scope remains in the canonical record.}\par
{\footnotesize \href{../../../../continuum/strategy/capability/feature/d6c271ae95ea44449c5e22665d1d9bc2.md}{Canonical record: d6c271ae95ea44449c5e22665d1d9bc2}; full qualifications and provenance.}\par
\paragraph*{II.1.3.9.9.5 [Feature] Non-PHI Compliance Auditing}
\phantomsection
\label{navigation:II.1.3.9.9.5}
\addcontentsline{toc}{paragraph}{II.1.3.9.9.5 [Feature] Non-PHI Compliance Auditing}

\label{element:08e37326908a49988dfb787f68df840a}
\index{[Feature] Non-PHI Compliance Auditing (II.1.3.9.9.5)}
{\small Capability / Feature · Canonical · pending\_adjudication}\par

Generate tamper-evident, verifiable audit records of security-significant events without emitting unmasked PHI.



\medskip\hrule\medskip

{\small Recorded unresolved matters: source-qualified matters; exact scope remains in the canonical record.}\par
{\footnotesize \href{../../../../continuum/strategy/capability/feature/08e37326908a49988dfb787f68df840a.md}{Canonical record: 08e37326908a49988dfb787f68df840a}; full qualifications and provenance.}\par
\paragraph*{II.1.3.9.9.6 [Feature] Operation-Bound Authorisation Enforcement}
\phantomsection
\label{navigation:II.1.3.9.9.6}
\addcontentsline{toc}{paragraph}{II.1.3.9.9.6 [Feature] Operation-Bound Authorisation Enforcement}

\label{element:3b7dcee7d58d4d1c960dc92abd58b4c3}
\index{[Feature] Operation-Bound Authorisation Enforcement (II.1.3.9.9.6)}
{\small Capability / Feature · Proposed · pending\_adjudication}\par

Apply effective policy at the governed action boundary and reject actions lacking sufficient authority.



\textbf{Overview}



Consistent enforcement is required across mutation, access, execution and administration rather than relying on each caller to voluntarily invoke a check.



\textbf{Boundary}



Consumes trusted identity and criteria; it does not authenticate or establish clinical privilege from directory roles.


{\small Recorded unresolved matters: source-qualified matters; exact scope remains in the canonical record.}\par
{\footnotesize \href{../../../../continuum/strategy/capability/feature/3b7dcee7d58d4d1c960dc92abd58b4c3.md}{Canonical record: 3b7dcee7d58d4d1c960dc92abd58b4c3}; full qualifications and provenance.}\par
\paragraph*{II.1.3.9.9.7 [Feature] Security Context Propagation}
\phantomsection
\label{navigation:II.1.3.9.9.7}
\addcontentsline{toc}{paragraph}{II.1.3.9.9.7 [Feature] Security Context Propagation}

\label{element:2b071c4fb10c4f61afecd18d09387214}
\index{[Feature] Security Context Propagation (II.1.3.9.9.7)}
{\small Capability / Feature · Canonical · pending\_adjudication}\par

Propagate tamper-evident, attributable security context across asynchronous processing boundaries.



\medskip\hrule\medskip

{\small Recorded unresolved matters: source-qualified matters; exact scope remains in the canonical record.}\par
{\footnotesize \href{../../../../continuum/strategy/capability/feature/2b071c4fb10c4f61afecd18d09387214.md}{Canonical record: 2b071c4fb10c4f61afecd18d09387214}; full qualifications and provenance.}\par
\paragraph*{II.1.3.9.9.8 [Feature] Trusted Identity Context Intake}
\phantomsection
\label{navigation:II.1.3.9.9.8}
\addcontentsline{toc}{paragraph}{II.1.3.9.9.8 [Feature] Trusted Identity Context Intake}

\label{element:2ef8f67907c34cb5a6a37fe7181434e2}
\index{[Feature] Trusted Identity Context Intake (II.1.3.9.9.8)}
{\small Capability / Feature · Proposed · pending\_adjudication}\par

Accept and qualify externally authenticated identity context for Harmonia-governed operations.



\textbf{Overview}



The boundary between WHO and WHAT remains explicit: an external identity assertion is verified as trusted input before internal authorisation is resolved.



\textbf{Boundary}



Harmonia does not become an authentication or credential-issuance provider.


{\small Recorded unresolved matters: source-qualified matters; exact scope remains in the canonical record.}\par
{\footnotesize \href{../../../../continuum/strategy/capability/feature/2ef8f67907c34cb5a6a37fe7181434e2.md}{Canonical record: 2ef8f67907c34cb5a6a37fe7181434e2}; full qualifications and provenance.}\par
\subsubsection*{II.1.3.9.10 [CT2] Health Information Exchange Enablement}
\phantomsection
\label{navigation:II.1.3.9.10}
\addcontentsline{toc}{subsubsection}{II.1.3.9.10 [CT2] Health Information Exchange Enablement}

{\small Historical predecessor: R.2.191 \hyperref[element:330d1c43856a4322a2e9f331b1a5ed20]{[CT2] Health Information Exchange}; Canonical ID: \texttt{\seqsplit{330d1c43856a4322a2e9f331b1a5ed20}}. Superseded historical identity — evolved correspondence · superseded.}\par
\label{element:a61a15908ed544e6acf3e9cab63a0267}
\index{[CT2] Health Information Exchange Enablement (II.1.3.9.10)}
{\small Capability · Canonical · current}\par

As a BEC, Health Information Exchange provides reusable information-movement enablement through submission, retrieval mediation, addressed distribution and syndication. The models distinguish sending to a participant, requesting information, deliberately distributing to named recipients and publishing changes to matching subscribers. Their common responsibility is the exchange context, including routing and delivery state, rather than the business meaning of every payload.



\textbf{Overview}



As a BEC, Health Information Exchange provides reusable information-movement enablement through submission, retrieval mediation, addressed distribution and syndication. The models distinguish sending to a participant, requesting information, deliberately distributing to named recipients and publishing changes to matching subscribers. Their common responsibility is the exchange context, including routing and delivery state, rather than the business meaning of every payload.



\textbf{Intent}



As a BEC, Health Information Exchange provides reusable information-movement enablement through submission, retrieval mediation, addressed distribution and syndication. The models distinguish sending to a participant, requesting information, deliberately distributing to named recipients and publishing changes to matching subscribers. Their common responsibility is the exchange context, including routing and delivery state, rather than the business meaning of every payload.



\textbf{Boundary}



Information administration does not confer authority over the represented person, organisation, service, place, device or clinical activity.



\medskip\hrule\medskip

{\small Recorded unresolved matters: source-qualified matters; exact scope remains in the canonical record.}\par
{\footnotesize \href{../../../../continuum/strategy/capability/a61a15908ed544e6acf3e9cab63a0267.md}{Canonical record: a61a15908ed544e6acf3e9cab63a0267}; full qualifications and provenance.}\par
\paragraph*{II.1.3.9.10.1 [CT3] HIE Endpoint State Management}
\phantomsection
\label{navigation:II.1.3.9.10.1}
\addcontentsline{toc}{paragraph}{II.1.3.9.10.1 [CT3] HIE Endpoint State Management}

{\small Historical predecessor: R.2.188 \hyperref[element:697de08ac49a4b7f8df7a1ed6f04210e]{[CT3] HIE Endpoint State Management}; Canonical ID: \texttt{\seqsplit{697de08ac49a4b7f8df7a1ed6f04210e}}. Superseded historical identity — evolved correspondence · superseded.}\par
\label{element:1bec951a1c98450b8d4672c5c424f9bd}
\index{[CT3] HIE Endpoint State Management (II.1.3.9.10.1)}
{\small Capability · Proposed · current}\par

Manage recognised exchange endpoint registrations, operational standing and contract-qualified availability information.



\textbf{Overview}



Manage recognised exchange endpoint registrations, operational standing and contract-qualified availability information. The reviewed hierarchy assigns this enduring ability to its stated business-enabling context.



\textbf{Intent}



Manage recognised exchange endpoint registrations, operational standing and contract-qualified availability information.



\textbf{Boundary}



Information administration and coordination preserve the authority of the responsible source, clinical actor and governed domain; they confer no additional real-world management authority.


{\footnotesize \href{../../../../continuum/strategy/capability/1bec951a1c98450b8d4672c5c424f9bd.md}{Canonical record: 1bec951a1c98450b8d4672c5c424f9bd}; full qualifications and provenance.}\par
\paragraph*{II.1.3.9.10.2 [CT3] HIE Implementation Management}
\phantomsection
\label{navigation:II.1.3.9.10.2}
\addcontentsline{toc}{paragraph}{II.1.3.9.10.2 [CT3] HIE Implementation Management}

{\small Historical predecessor: R.2.189 \hyperref[element:a99b4119339545a89dc54ebb216ddb85]{[CT3] HIE Implementation Management}; Canonical ID: \texttt{\seqsplit{a99b4119339545a89dc54ebb216ddb85}}. Superseded historical identity — evolved correspondence · superseded.}\par
\label{element:8ce29be373e04252b7aa881b6a72f375}
\index{[CT3] HIE Implementation Management (II.1.3.9.10.2)}
{\small Capability · Proposed · current}\par

Manage implementation information and progression for established exchange requirements, including submission, retrieval, distribution and syndication interactions.



\textbf{Overview}



Manage implementation information and progression for established exchange requirements, including submission, retrieval, distribution and syndication interactions. The reviewed hierarchy assigns this enduring ability to its stated business-enabling context.



\textbf{Intent}



Manage implementation information and progression for established exchange requirements, including submission, retrieval, distribution and syndication interactions.



\textbf{Boundary}



Information administration and coordination preserve the authority of the responsible source, clinical actor and governed domain; they confer no additional real-world management authority.


{\footnotesize \href{../../../../continuum/strategy/capability/8ce29be373e04252b7aa881b6a72f375.md}{Canonical record: 8ce29be373e04252b7aa881b6a72f375}; full qualifications and provenance.}\par
\subparagraph*{II.1.3.9.10.2.1 [Feature] Addressed Clinical Distribution}
\phantomsection
\label{navigation:II.1.3.9.10.2.1}
\addcontentsline{toc}{subparagraph}{II.1.3.9.10.2.1 [Feature] Addressed Clinical Distribution}

\label{element:b476c86edab644c394da33d3de7bd0bd}
\index{[Feature] Addressed Clinical Distribution (II.1.3.9.10.2.1)}
{\small Capability / Feature · Canonical · current}\par

Fan out clinical events and diagnostic reports to designated external and internal recipients with per-destination delivery tracking.



Addressed distribution sends information to explicitly designated recipients and tracks each destination's delivery separately. This preserves partial outcomes in a fan-out rather than collapsing several deliveries into a single success that hides a missing recipient.


{\footnotesize \href{../../../../continuum/strategy/capability/feature/b476c86edab644c394da33d3de7bd0bd.md}{Canonical record: b476c86edab644c394da33d3de7bd0bd}; full qualifications and provenance.}\par
\subparagraph*{II.1.3.9.10.2.2 [Feature] Clinical Record Retrieval Mediation}
\phantomsection
\label{navigation:II.1.3.9.10.2.2}
\addcontentsline{toc}{subparagraph}{II.1.3.9.10.2.2 [Feature] Clinical Record Retrieval Mediation}

\label{element:6f4a9c8f7916488bb0dd838ccccb1df4}
\index{[Feature] Clinical Record Retrieval Mediation (II.1.3.9.10.2.2)}
{\small Capability / Feature · Canonical · current}\par

Mediate query-based retrieval requests across external and internal repositories.



Retrieval mediation connects a request with information held in participating repositories. It makes distributed retrieval possible without transferring repository authority, acquiring the clinical payload or treating a resolvable request as automatic permission to disclose.


{\footnotesize \href{../../../../continuum/strategy/capability/feature/6f4a9c8f7916488bb0dd838ccccb1df4.md}{Canonical record: 6f4a9c8f7916488bb0dd838ccccb1df4}; full qualifications and provenance.}\par
\subparagraph*{II.1.3.9.10.2.3 [Feature] Direct Submission Ingestion \& Delivery}
\phantomsection
\label{navigation:II.1.3.9.10.2.3}
\addcontentsline{toc}{subparagraph}{II.1.3.9.10.2.3 [Feature] Direct Submission Ingestion \& Delivery}

\label{element:e8bfad88315c4833b6772ef580ca55a5}
\index{[Feature] Direct Submission Ingestion \& Delivery (II.1.3.9.10.2.3)}
{\small Capability / Feature · Canonical · current}\par

Support push submissions of clinical documents and messages across boundaries.



Direct submission supports a sender's deliberate push of information across an exchange boundary. It preserves the distinction between ingress or delivery acceptance and the recipient's business handling, under the actual contract rather than an assumed meaning of transport success.


{\footnotesize \href{../../../../continuum/strategy/capability/feature/e8bfad88315c4833b6772ef580ca55a5.md}{Canonical record: e8bfad88315c4833b6772ef580ca55a5}; full qualifications and provenance.}\par
\subparagraph*{II.1.3.9.10.2.4 [Feature] Event-Driven Syndication}
\phantomsection
\label{navigation:II.1.3.9.10.2.4}
\addcontentsline{toc}{subparagraph}{II.1.3.9.10.2.4 [Feature] Event-Driven Syndication}

\label{element:30bb720ed4e64a47a1fcf3be811c1931}
\index{[Feature] Event-Driven Syndication (II.1.3.9.10.2.4)}
{\small Capability / Feature · Canonical · current}\par

Publish clinical and directory change events to subscribed systems matching interest filters.



Syndication connects published changes with subscribers whose registered interests match them. It differs from named-recipient distribution in how consumers are selected, while retaining governed delivery and the distinction between notification and the subscriber's eventual business action.


{\footnotesize \href{../../../../continuum/strategy/capability/feature/30bb720ed4e64a47a1fcf3be811c1931.md}{Canonical record: 30bb720ed4e64a47a1fcf3be811c1931}; full qualifications and provenance.}\par
\paragraph*{II.1.3.9.10.3 [CT3] HIE Requirement Management}
\phantomsection
\label{navigation:II.1.3.9.10.3}
\addcontentsline{toc}{paragraph}{II.1.3.9.10.3 [CT3] HIE Requirement Management}

{\small Historical predecessor: R.2.190 \hyperref[element:a32c8bff0c284f5c9d6ab1f46991c232]{[CT3] HIE Requirement Management}; Canonical ID: \texttt{\seqsplit{a32c8bff0c284f5c9d6ab1f46991c232}}. Superseded historical identity — evolved correspondence · superseded.}\par
\label{element:1c852a1799984feba50011172983c865}
\index{[CT3] HIE Requirement Management (II.1.3.9.10.3)}
{\small Capability · Proposed · current}\par

Manage the identified information-exchange requirements, interaction contracts and participant obligations applicable to a health-information exchange.



\textbf{Overview}



Manage the identified information-exchange requirements, interaction contracts and participant obligations applicable to a health-information exchange. The reviewed hierarchy assigns this enduring ability to its stated business-enabling context.



\textbf{Intent}



Manage the identified information-exchange requirements, interaction contracts and participant obligations applicable to a health-information exchange.



\textbf{Boundary}



Information administration and coordination preserve the authority of the responsible source, clinical actor and governed domain; they confer no additional real-world management authority.


{\footnotesize \href{../../../../continuum/strategy/capability/1c852a1799984feba50011172983c865.md}{Canonical record: 1c852a1799984feba50011172983c865}; full qualifications and provenance.}\par
\subsubsection*{II.1.3.9.11 [CT2] Information Asset Stewardship Management}
\phantomsection
\label{navigation:II.1.3.9.11}
\addcontentsline{toc}{subsubsection}{II.1.3.9.11 [CT2] Information Asset Stewardship Management}

{\small Historical predecessor: R.2.194 \hyperref[element:cee085c690314802851de71f4741d8db]{[CT2] Information Asset Stewardship Management}; Canonical ID: \texttt{\seqsplit{cee085c690314802851de71f4741d8db}}. Superseded historical identity — direct correspondence · superseded.}\par
\label{element:5bac15316d8d4f4db6e7a253f1acfbe3}
\index{[CT2] Information Asset Stewardship Management (II.1.3.9.11)}
{\small Capability · Proposed · current}\par

Manage definitions and tags that make governed information assets recognisable and interpretable for their permitted uses.



\textbf{Overview}



Manage definitions and tags that make governed information assets recognisable and interpretable for their permitted uses. The reviewed hierarchy assigns this enduring ability to its stated business-enabling context.



\textbf{Intent}



Manage definitions and tags that make governed information assets recognisable and interpretable for their permitted uses.



\textbf{Boundary}



Information administration and coordination preserve the authority of the responsible source, clinical actor and governed domain; they confer no additional real-world management authority.


{\footnotesize \href{../../../../continuum/strategy/capability/5bac15316d8d4f4db6e7a253f1acfbe3.md}{Canonical record: 5bac15316d8d4f4db6e7a253f1acfbe3}; full qualifications and provenance.}\par
\paragraph*{II.1.3.9.11.1 [CT3] Information Asset Definition}
\phantomsection
\label{navigation:II.1.3.9.11.1}
\addcontentsline{toc}{paragraph}{II.1.3.9.11.1 [CT3] Information Asset Definition}

{\small Historical predecessor: R.2.192 \hyperref[element:166a22463d2c4534bda7ae31273398b6]{[CT3] Information Asset Definition}; Canonical ID: \texttt{\seqsplit{166a22463d2c4534bda7ae31273398b6}}. Superseded historical identity — direct correspondence · superseded.}\par
\label{element:05c58b2ccfcc45e9a6c70d9a569c4f61}
\index{[CT3] Information Asset Definition (II.1.3.9.11.1)}
{\small Capability · Proposed · current}\par

Define the identifying, semantic and governed-use characteristics of an information asset.



\textbf{Overview}



Define the identifying, semantic and governed-use characteristics of an information asset. The reviewed hierarchy assigns this enduring ability to its stated business-enabling context.



\textbf{Intent}



Define the identifying, semantic and governed-use characteristics of an information asset.



\textbf{Boundary}



Information administration and coordination preserve the authority of the responsible source, clinical actor and governed domain; they confer no additional real-world management authority.


{\footnotesize \href{../../../../continuum/strategy/capability/05c58b2ccfcc45e9a6c70d9a569c4f61.md}{Canonical record: 05c58b2ccfcc45e9a6c70d9a569c4f61}; full qualifications and provenance.}\par
\paragraph*{II.1.3.9.11.2 [CT3] Information Asset Tag Management}
\phantomsection
\label{navigation:II.1.3.9.11.2}
\addcontentsline{toc}{paragraph}{II.1.3.9.11.2 [CT3] Information Asset Tag Management}

{\small Historical predecessor: R.2.193 \hyperref[element:ae44a4ee78414d1098814315793bd751]{[CT3] Information Asset Tag Management}; Canonical ID: \texttt{\seqsplit{ae44a4ee78414d1098814315793bd751}}. Superseded historical identity — direct correspondence · superseded.}\par
\label{element:df27c3ba6e164e56af7bff4e1eb1954a}
\index{[CT3] Information Asset Tag Management (II.1.3.9.11.2)}
{\small Capability · Proposed · current}\par

Manage attributable tags describing information assets without treating caller-controlled tags as security authority.



\textbf{Overview}



Manage attributable tags describing information assets without treating caller-controlled tags as security authority. The reviewed hierarchy assigns this enduring ability to its stated business-enabling context.



\textbf{Intent}



Manage attributable tags describing information assets without treating caller-controlled tags as security authority.



\textbf{Boundary}



Information administration and coordination preserve the authority of the responsible source, clinical actor and governed domain; they confer no additional real-world management authority.


{\footnotesize \href{../../../../continuum/strategy/capability/df27c3ba6e164e56af7bff4e1eb1954a.md}{Canonical record: df27c3ba6e164e56af7bff4e1eb1954a}; full qualifications and provenance.}\par
\subsubsection*{II.1.3.9.12 [CT2] Information Content Maintainer Management}
\phantomsection
\label{navigation:II.1.3.9.12}
\addcontentsline{toc}{subsubsection}{II.1.3.9.12 [CT2] Information Content Maintainer Management}

{\small Historical predecessor: R.2.198 \hyperref[element:acb977e879204f1491293e5452e0ddf0]{[CT2] Information Content Maintainer Management}; Canonical ID: \texttt{\seqsplit{acb977e879204f1491293e5452e0ddf0}}. Superseded historical identity — direct correspondence · superseded.}\par
\label{element:b81a1ab42af24c2e9896b8f3f78d291a}
\index{[CT2] Information Content Maintainer Management (II.1.3.9.12)}
{\small Capability · Proposed · current}\par

Administer information-content maintainer participation, source provenance and governed information products.



\textbf{Overview}



Administer information-content maintainer participation, source provenance and governed information products. The reviewed hierarchy assigns this enduring ability to its stated business-enabling context.



\textbf{Intent}



Administer information-content maintainer participation, source provenance and governed information products.



\textbf{Boundary}



Information administration and coordination preserve the authority of the responsible source, clinical actor and governed domain; they confer no additional real-world management authority.


{\footnotesize \href{../../../../continuum/strategy/capability/b81a1ab42af24c2e9896b8f3f78d291a.md}{Canonical record: b81a1ab42af24c2e9896b8f3f78d291a}; full qualifications and provenance.}\par
\paragraph*{II.1.3.9.12.1 [CT3] Data Provenance Capture}
\phantomsection
\label{navigation:II.1.3.9.12.1}
\addcontentsline{toc}{paragraph}{II.1.3.9.12.1 [CT3] Data Provenance Capture}

{\small Historical predecessor: R.2.195 \hyperref[element:61258c7fe41f4541871a67ec9d7185cf]{[CT3] Data Provenance Capture}; Canonical ID: \texttt{\seqsplit{61258c7fe41f4541871a67ec9d7185cf}}. Superseded historical identity — direct correspondence · superseded.}\par
\label{element:44a7937520a54945b06b5fd4a2bbd870}
\index{[CT3] Data Provenance Capture (II.1.3.9.12.1)}
{\small Capability · Proposed · current}\par

Capture meaningful source attribution and transformation lineage for governed information content without transferring originating authority.



\textbf{Overview}



Capture meaningful source attribution and transformation lineage for governed information content without transferring originating authority.



\textbf{Intent}



Capture meaningful source attribution and transformation lineage for governed information content without transferring originating authority.



\textbf{Boundary}



Lineage does not grant originating authority, decide clinical correctness or replace activity audit and independent assurance.


{\footnotesize \href{../../../../continuum/strategy/capability/44a7937520a54945b06b5fd4a2bbd870.md}{Canonical record: 44a7937520a54945b06b5fd4a2bbd870}; full qualifications and provenance.}\par
\subparagraph*{II.1.3.9.12.1.1 [Feature] Information Transformation Lineage}
\phantomsection
\label{navigation:II.1.3.9.12.1.1}
\addcontentsline{toc}{subparagraph}{II.1.3.9.12.1.1 [Feature] Information Transformation Lineage}

\label{element:cc47f49b85a34e87bafc49de7f0426e2}
\index{[Feature] Information Transformation Lineage (II.1.3.9.12.1.1)}
{\small Capability / Feature · Proposed · current}\par

Record meaningful source-to-derived information relationships and transformation attribution.



\textbf{Overview}



A consumer should be able to explain which source assertions contributed to a managed representation and what significant transformation occurred.



\textbf{Boundary}



Routine cache reconstruction, retries and diagnostics do not automatically become permanent lineage.


{\footnotesize \href{../../../../continuum/strategy/capability/feature/cc47f49b85a34e87bafc49de7f0426e2.md}{Canonical record: cc47f49b85a34e87bafc49de7f0426e2}; full qualifications and provenance.}\par
\paragraph*{II.1.3.9.12.2 [CT3] Data Scientist Administration}
\phantomsection
\label{navigation:II.1.3.9.12.2}
\addcontentsline{toc}{paragraph}{II.1.3.9.12.2 [CT3] Data Scientist Administration}

{\small Historical predecessor: R.2.196 \hyperref[element:1cb497aec93c4b82ba7f5b622813e54e]{[CT3] Data Scientist Administration}; Canonical ID: \texttt{\seqsplit{1cb497aec93c4b82ba7f5b622813e54e}}. Superseded historical identity — direct correspondence · superseded.}\par
\label{element:3bb1a11fba234717bf40c9e6452f401a}
\index{[CT3] Data Scientist Administration (II.1.3.9.12.2)}
{\small Capability · Proposed · current}\par

Administer profiles and qualified participation information for data scientists maintaining governed information content.



\textbf{Overview}



Administer profiles and qualified participation information for data scientists maintaining governed information content. The reviewed hierarchy assigns this enduring ability to its stated business-enabling context.



\textbf{Intent}



Administer profiles and qualified participation information for data scientists maintaining governed information content.



\textbf{Boundary}



Information administration and coordination preserve the authority of the responsible source, clinical actor and governed domain; they confer no additional real-world management authority.


{\footnotesize \href{../../../../continuum/strategy/capability/3bb1a11fba234717bf40c9e6452f401a.md}{Canonical record: 3bb1a11fba234717bf40c9e6452f401a}; full qualifications and provenance.}\par
\paragraph*{II.1.3.9.12.3 [CT3] Information Product Administration}
\phantomsection
\label{navigation:II.1.3.9.12.3}
\addcontentsline{toc}{paragraph}{II.1.3.9.12.3 [CT3] Information Product Administration}

{\small Historical predecessor: R.2.197 \hyperref[element:5d9b7f8d2d7b41759b6d4931c1d2aaa2]{[CT3] Information Product Administration}; Canonical ID: \texttt{\seqsplit{5d9b7f8d2d7b41759b6d4931c1d2aaa2}}. Superseded historical identity — direct correspondence · superseded.}\par
\label{element:b5a6a31855834ea58870ebb6c51ed4b0}
\index{[CT3] Information Product Administration (II.1.3.9.12.3)}
{\small Capability · Proposed · current}\par

Administer identifiable governed information products, their purpose, content standing and qualified release/use context.



\textbf{Overview}



Administer identifiable governed information products, their purpose, content standing and qualified release/use context. The reviewed hierarchy assigns this enduring ability to its stated business-enabling context.



\textbf{Intent}



Administer identifiable governed information products, their purpose, content standing and qualified release/use context.



\textbf{Boundary}



Information administration and coordination preserve the authority of the responsible source, clinical actor and governed domain; they confer no additional real-world management authority.


{\footnotesize \href{../../../../continuum/strategy/capability/b5a6a31855834ea58870ebb6c51ed4b0.md}{Canonical record: b5a6a31855834ea58870ebb6c51ed4b0}; full qualifications and provenance.}\par
\subparagraph*{II.1.3.9.12.3.1 [Feature] Authorised Research Release Assembly}
\phantomsection
\label{navigation:II.1.3.9.12.3.1}
\addcontentsline{toc}{subparagraph}{II.1.3.9.12.3.1 [Feature] Authorised Research Release Assembly}

\label{element:dcebc2f2b5b7494a8e80a02402664078}
\index{[Feature] Authorised Research Release Assembly (II.1.3.9.12.3.1)}
{\small Capability / Feature · Proposed · current}\par

Assemble and qualify the existing information allowed by a particular authorised research purpose.



\textbf{Overview}



The release carries purpose and applicable privacy/consent/use constraints; the requester does not gain unrestricted clinical-record access by having a research role.



\textbf{Boundary}



No ethics approval, de-identification guarantee or clinical-trial management is inferred.


{\footnotesize \href{../../../../continuum/strategy/capability/feature/dcebc2f2b5b7494a8e80a02402664078.md}{Canonical record: dcebc2f2b5b7494a8e80a02402664078}; full qualifications and provenance.}\par
\subparagraph*{II.1.3.9.12.3.2 [Feature] Demand \& Capacity Planning Projection}
\phantomsection
\label{navigation:II.1.3.9.12.3.2}
\addcontentsline{toc}{subparagraph}{II.1.3.9.12.3.2 [Feature] Demand \& Capacity Planning Projection}

\label{element:4f97baa3aed24b8385bf3ecda4eab622}
\index{[Feature] Demand \& Capacity Planning Projection (II.1.3.9.12.3.2)}
{\small Capability / Feature · Proposed · current}\par

Project qualified observed demand, capacity and utilisation for authorised planning consumers.



\textbf{Overview}



Comparisons need service, temporal and context qualifiers so that planned future capacity is not confused with present operational capacity.



\textbf{Boundary}



Projection does not establish a planning conclusion or allocate services and staff.


{\footnotesize \href{../../../../continuum/strategy/capability/feature/4f97baa3aed24b8385bf3ecda4eab622.md}{Canonical record: 4f97baa3aed24b8385bf3ecda4eab622}; full qualifications and provenance.}\par
\subparagraph*{II.1.3.9.12.3.3 [Feature] Population Notification Obligation Tracking}
\phantomsection
\label{navigation:II.1.3.9.12.3.3}
\addcontentsline{toc}{subparagraph}{II.1.3.9.12.3.3 [Feature] Population Notification Obligation Tracking}

\label{element:62f6a91a10684876b7d08b536700f12f}
\index{[Feature] Population Notification Obligation Tracking (II.1.3.9.12.3.3)}
{\small Capability / Feature · Proposed · current}\par

Track the relation between an applicable notification obligation, the supplied source evidence and observed notification disposition.



\textbf{Overview}



Sending a required notification and discharging the governing obligation are different assertions. This Feature keeps the explicit contract and indeterminate responses visible.



\textbf{Boundary}



Does not define notifiability, author statutory rules or own public-health response.


{\footnotesize \href{../../../../continuum/strategy/capability/feature/62f6a91a10684876b7d08b536700f12f.md}{Canonical record: 62f6a91a10684876b7d08b536700f12f}; full qualifications and provenance.}\par
\subsubsection*{II.1.3.9.13 [CT2] Information Design Governance}
\phantomsection
\label{navigation:II.1.3.9.13}
\addcontentsline{toc}{subsubsection}{II.1.3.9.13 [CT2] Information Design Governance}

{\small Historical predecessor: R.2.203 \hyperref[element:0bdd7f8c0cd74d9fabc49801af037fc2]{[CT2] Information Design Governance}; Canonical ID: \texttt{\seqsplit{0bdd7f8c0cd74d9fabc49801af037fc2}}. Superseded historical identity — direct correspondence · superseded.}\par
\label{element:f39ed9ebcc724857a4bf7a2b8e045c7a}
\index{[CT2] Information Design Governance (II.1.3.9.13)}
{\small Capability · Canonical · current}\par

Information Design Governance maintains canonical definitions, semantic bindings, exchange profiles and the conformance rules applied to representations. It makes meaning and constraints available across regional information flows so that an accepted technical format is not mistaken for semantic correctness. Governance and verification remain distinguishable from the clinical facts represented.



\textbf{Overview}



Information Design Governance maintains canonical definitions, semantic bindings, exchange profiles and the conformance rules applied to representations. It makes meaning and constraints available across regional information flows so that an accepted technical format is not mistaken for semantic correctness. Governance and verification remain distinguishable from the clinical facts represented.



\textbf{Intent}



Information Design Governance maintains canonical definitions, semantic bindings, exchange profiles and the conformance rules applied to representations. It makes meaning and constraints available across regional information flows so that an accepted technical format is not mistaken for semantic correctness. Governance and verification remain distinguishable from the clinical facts represented.



\textbf{Boundary}



Information administration does not confer authority over the represented person, organisation, service, place, device or clinical activity.



\medskip\hrule\medskip

{\footnotesize \href{../../../../continuum/strategy/capability/f39ed9ebcc724857a4bf7a2b8e045c7a.md}{Canonical record: f39ed9ebcc724857a4bf7a2b8e045c7a}; full qualifications and provenance.}\par
\paragraph*{II.1.3.9.13.1 [CT3] Enterprise Canonical/Reusable Model Design}
\phantomsection
\label{navigation:II.1.3.9.13.1}
\addcontentsline{toc}{paragraph}{II.1.3.9.13.1 [CT3] Enterprise Canonical/Reusable Model Design}

{\small Historical predecessor: R.2.199 \hyperref[element:dbc37445ad944a85a7c4f26f410967b2]{[CT3] Enterprise Canonical/Reusable Model Design}; Canonical ID: \texttt{\seqsplit{dbc37445ad944a85a7c4f26f410967b2}}. Superseded historical identity — direct correspondence · superseded.}\par
\label{element:90e9cd401ddf4c7690c978d8dd131413}
\index{[CT3] Enterprise Canonical/Reusable Model Design (II.1.3.9.13.1)}
{\small Capability · Proposed · current}\par

Design canonical and reusable enterprise information models and semantic bindings under applicable information-design governance.



\textbf{Overview}



Design canonical and reusable enterprise information models and semantic bindings under applicable information-design governance. The reviewed hierarchy assigns this enduring ability to its stated business-enabling context.



\textbf{Intent}



Design canonical and reusable enterprise information models and semantic bindings under applicable information-design governance.



\textbf{Boundary}



Information administration and coordination preserve the authority of the responsible source, clinical actor and governed domain; they confer no additional real-world management authority.


{\footnotesize \href{../../../../continuum/strategy/capability/90e9cd401ddf4c7690c978d8dd131413.md}{Canonical record: 90e9cd401ddf4c7690c978d8dd131413}; full qualifications and provenance.}\par
\subparagraph*{II.1.3.9.13.1.1 [Feature] Exchange Profile Governance}
\phantomsection
\label{navigation:II.1.3.9.13.1.1}
\addcontentsline{toc}{subparagraph}{II.1.3.9.13.1.1 [Feature] Exchange Profile Governance}

\label{element:5e59d8f2e23e4fb5b59a5bd1dd3f31fa}
\index{[Feature] Exchange Profile Governance (II.1.3.9.13.1.1)}
{\small Capability / Feature · Proposed · current}\par

Govern the semantic and conformance constraints of an interoperability profile.



\textbf{Overview}



An external contract must state which representations and extensions are permitted; validation applies that contract but does not define it.



\textbf{Boundary}



Profile governance does not transmit messages or publish internal operational metadata.


{\footnotesize \href{../../../../continuum/strategy/capability/feature/5e59d8f2e23e4fb5b59a5bd1dd3f31fa.md}{Canonical record: 5e59d8f2e23e4fb5b59a5bd1dd3f31fa}; full qualifications and provenance.}\par
\subparagraph*{II.1.3.9.13.1.2 [Feature] Information Definition \& Schema Governance}
\phantomsection
\label{navigation:II.1.3.9.13.1.2}
\addcontentsline{toc}{subparagraph}{II.1.3.9.13.1.2 [Feature] Information Definition \& Schema Governance}

\label{element:3317bbb7e2db4384abf44a50bf67422c}
\index{[Feature] Information Definition \& Schema Governance (II.1.3.9.13.1.2)}
{\small Capability / Feature · Proposed · current}\par

Govern structural definitions, data dictionaries and constraint meaning.



\textbf{Overview}



Existing definition-governance behaviour needs a distinct Feature frame from general semantic-model authorship and runtime validation.



\textbf{Boundary}



Governed definitions do not own runtime instances or their clinical facts.


{\footnotesize \href{../../../../continuum/strategy/capability/feature/3317bbb7e2db4384abf44a50bf67422c.md}{Canonical record: 3317bbb7e2db4384abf44a50bf67422c}; full qualifications and provenance.}\par
\subparagraph*{II.1.3.9.13.1.3 [Feature] Information Standards \& Semantic Governance}
\phantomsection
\label{navigation:II.1.3.9.13.1.3}
\addcontentsline{toc}{subparagraph}{II.1.3.9.13.1.3 [Feature] Information Standards \& Semantic Governance}

\label{element:af5d5249302f47ebb393c91afcc60392}
\index{[Feature] Information Standards \& Semantic Governance (II.1.3.9.13.1.3)}
{\small Capability / Feature · Canonical · current}\par

Govern canonical clinical data definitions, terminology bindings, and exchange profiles across regional systems.



Semantic governance provides canonical definitions, bindings and profiles through which regional representations can be understood consistently. Its purpose concerns meaning and constraints rather than a particular schema technology, and consuming a definition does not transfer its governing authority.


{\footnotesize \href{../../../../continuum/strategy/capability/feature/af5d5249302f47ebb393c91afcc60392.md}{Canonical record: af5d5249302f47ebb393c91afcc60392}; full qualifications and provenance.}\par
\subparagraph*{II.1.3.9.13.1.4 [Feature] Terminology Binding Governance}
\phantomsection
\label{navigation:II.1.3.9.13.1.4}
\addcontentsline{toc}{subparagraph}{II.1.3.9.13.1.4 [Feature] Terminology Binding Governance}

\label{element:b999df692ca64010b4adeedbcb0d3a5a}
\index{[Feature] Terminology Binding Governance (II.1.3.9.13.1.4)}
{\small Capability / Feature · Proposed · current}\par

Govern how information definitions are bound to applicable terminology and value sets.



\textbf{Overview}



Binding a data element to a vocabulary is distinct from resolving a patient code or changing the runtime membership of a dataset.



\textbf{Boundary}



Definition of a binding does not take external ontology authority or execute policy.


{\footnotesize \href{../../../../continuum/strategy/capability/feature/b999df692ca64010b4adeedbcb0d3a5a.md}{Canonical record: b999df692ca64010b4adeedbcb0d3a5a}; full qualifications and provenance.}\par
\paragraph*{II.1.3.9.13.2 [CT3] Enterprise Canonical/Reusable Model Library Management}
\phantomsection
\label{navigation:II.1.3.9.13.2}
\addcontentsline{toc}{paragraph}{II.1.3.9.13.2 [CT3] Enterprise Canonical/Reusable Model Library Management}

{\small Historical predecessor: R.2.200 \hyperref[element:f5eccd4f65df49bba98104e5abfa5e19]{[CT3] Enterprise Canonical/Reusable Model Library Management}; Canonical ID: \texttt{\seqsplit{f5eccd4f65df49bba98104e5abfa5e19}}. Superseded historical identity — direct correspondence · superseded.}\par
\label{element:76943b4b68644c4dac1368fce8bb20d9}
\index{[CT3] Enterprise Canonical/Reusable Model Library Management (II.1.3.9.13.2)}
{\small Capability · Proposed · current}\par

Manage the identifiable library of canonical and reusable enterprise model definitions and their governed revisions.



\textbf{Overview}



Manage the identifiable library of canonical and reusable enterprise model definitions and their governed revisions. The reviewed hierarchy assigns this enduring ability to its stated business-enabling context.



\textbf{Intent}



Manage the identifiable library of canonical and reusable enterprise model definitions and their governed revisions.



\textbf{Boundary}



Information administration and coordination preserve the authority of the responsible source, clinical actor and governed domain; they confer no additional real-world management authority.


{\footnotesize \href{../../../../continuum/strategy/capability/76943b4b68644c4dac1368fce8bb20d9.md}{Canonical record: 76943b4b68644c4dac1368fce8bb20d9}; full qualifications and provenance.}\par
\paragraph*{II.1.3.9.13.3 [CT3] Enterprise Canonical/Reusable Model Utilisation Support}
\phantomsection
\label{navigation:II.1.3.9.13.3}
\addcontentsline{toc}{paragraph}{II.1.3.9.13.3 [CT3] Enterprise Canonical/Reusable Model Utilisation Support}

{\small Historical predecessor: R.2.201 \hyperref[element:7ae95677a00346f18a0ba01522a20153]{[CT3] Enterprise Canonical/Reusable Model Utilisation Support}; Canonical ID: \texttt{\seqsplit{7ae95677a00346f18a0ba01522a20153}}. Superseded historical identity — direct correspondence · superseded.}\par
\label{element:5744cb9d86494064af87afb3cbf2cf97}
\index{[CT3] Enterprise Canonical/Reusable Model Utilisation Support (II.1.3.9.13.3)}
{\small Capability · Proposed · current}\par

Support correct utilisation and conformance of canonical and reusable models across their intended consumers.



\textbf{Overview}



Support correct utilisation and conformance of canonical and reusable models across their intended consumers. The reviewed hierarchy assigns this enduring ability to its stated business-enabling context.



\textbf{Intent}



Support correct utilisation and conformance of canonical and reusable models across their intended consumers.



\textbf{Boundary}



Information administration and coordination preserve the authority of the responsible source, clinical actor and governed domain; they confer no additional real-world management authority.


{\footnotesize \href{../../../../continuum/strategy/capability/5744cb9d86494064af87afb3cbf2cf97.md}{Canonical record: 5744cb9d86494064af87afb3cbf2cf97}; full qualifications and provenance.}\par
\subparagraph*{II.1.3.9.13.3.1 [Feature] Semantic Conformance Verification}
\phantomsection
\label{navigation:II.1.3.9.13.3.1}
\addcontentsline{toc}{subparagraph}{II.1.3.9.13.3.1 [Feature] Semantic Conformance Verification}

\label{element:d170a68efcae4aa59350f565ca69443a}
\index{[Feature] Semantic Conformance Verification (II.1.3.9.13.3.1)}
{\small Capability / Feature · Canonical · current}\par

Verify inbound and outbound representations against published semantic rules and conformance constraints.



Conformance verification evaluates a representation against published semantic rules and constraints. It distinguishes meaningful conformity from simple parseability while retaining that a conformant message is not necessarily clinically true, authoritative or permitted for disclosure.


{\footnotesize \href{../../../../continuum/strategy/capability/feature/d170a68efcae4aa59350f565ca69443a.md}{Canonical record: d170a68efcae4aa59350f565ca69443a}; full qualifications and provenance.}\par
\paragraph*{II.1.3.9.13.4 [CT3] Enterprise Reference Data Stewardship}
\phantomsection
\label{navigation:II.1.3.9.13.4}
\addcontentsline{toc}{paragraph}{II.1.3.9.13.4 [CT3] Enterprise Reference Data Stewardship}

{\small Historical predecessor: R.2.202 \hyperref[element:fd9f3239fdb64ea687825b66bb4d1f17]{[CT3] Enterprise Reference Data Stewardship}; Canonical ID: \texttt{\seqsplit{fd9f3239fdb64ea687825b66bb4d1f17}}. Superseded historical identity — direct correspondence · superseded.}\par
\label{element:971d4b079ad24e89b70a8ed8dec0c8bf}
\index{[CT3] Enterprise Reference Data Stewardship (II.1.3.9.13.4)}
{\small Capability · Proposed · current}\par

Steward enterprise reference-data definitions, bindings and applicable governance independently of mutable runtime dataset administration.



\textbf{Overview}



Steward enterprise reference-data definitions, bindings and applicable governance independently of mutable runtime dataset administration. The reviewed hierarchy assigns this enduring ability to its stated business-enabling context.



\textbf{Intent}



Steward enterprise reference-data definitions, bindings and applicable governance independently of mutable runtime dataset administration.



\textbf{Boundary}



Information administration and coordination preserve the authority of the responsible source, clinical actor and governed domain; they confer no additional real-world management authority.


{\footnotesize \href{../../../../continuum/strategy/capability/971d4b079ad24e89b70a8ed8dec0c8bf.md}{Canonical record: 971d4b079ad24e89b70a8ed8dec0c8bf}; full qualifications and provenance.}\par
\subsubsection*{II.1.3.9.14 [CT2] Information Protection \& Accountability Enablement [superseded]}
\phantomsection
\label{navigation:II.1.3.9.14}
\addcontentsline{toc}{subsubsection}{II.1.3.9.14 [CT2] Information Protection \& Accountability Enablement [superseded]}

\label{element:e922cba7cf1d4ebeb9abf6472a3803ce}
\index{[CT2] Information Protection \& Accountability Enablement (II.1.3.9.14)}
{\small Capability · Proposed · superseded}\par
\href{../../../../continuum/strategy/capability/e922cba7cf1d4ebeb9abf6472a3803ce.md}{Historical canonical definition and provenance}.\par
\paragraph*{II.1.3.9.14.1 [CT3] Governance Requirement Publication Support}
\phantomsection
\label{navigation:II.1.3.9.14.1}
\addcontentsline{toc}{paragraph}{II.1.3.9.14.1 [CT3] Governance Requirement Publication Support}

\label{element:7e780751d50e4ffe942f8443b7cb6567}
\index{[CT3] Governance Requirement Publication Support (II.1.3.9.14.1)}
{\small Capability · Proposed · pending\_adjudication}\par

Represent and publish versioned governing requirements for use by authorised operational and assurance consumers.



\textbf{Overview}



Consumers need an identifiable statement of which rule applies and who established it. This proposal fills BC-17 information support without giving the platform authority to make legislation, approve governance requirements or reinterpret their force.



\textbf{Intent}



Let operational and assurance consumers identify the governing statement and applicability basis they actually used.



\textbf{Boundary}



Governance approval remains externally mandated; operational policy, assurance criteria and requirement authority are not interchangeable.


{\small Recorded unresolved matters: source-qualified matters; exact scope remains in the canonical record.}\par
{\footnotesize \href{../../../../continuum/strategy/capability/7e780751d50e4ffe942f8443b7cb6567.md}{Canonical record: 7e780751d50e4ffe942f8443b7cb6567}; full qualifications and provenance.}\par
\subparagraph*{II.1.3.9.14.1.1 [Feature] Governance Requirement Version Publication}
\phantomsection
\label{navigation:II.1.3.9.14.1.1}
\addcontentsline{toc}{subparagraph}{II.1.3.9.14.1.1 [Feature] Governance Requirement Version Publication}

\label{element:001a067a4b3b46bebfe945f1bef541ce}
\index{[Feature] Governance Requirement Version Publication (II.1.3.9.14.1.1)}
{\small Capability / Feature · Proposed · pending\_adjudication}\par

Publish attributable versions and applicability of externally approved governing requirements.



\textbf{Overview}



Operational and assurance consumers need to identify the governing basis they used rather than silently substitute a current rule for a historical one.



\textbf{Boundary}



Publication does not confer requirement approval, authorisation policy or assurance-criteria authority.


{\small Recorded unresolved matters: source-qualified matters; exact scope remains in the canonical record.}\par
{\footnotesize \href{../../../../continuum/strategy/capability/feature/001a067a4b3b46bebfe945f1bef541ce.md}{Canonical record: 001a067a4b3b46bebfe945f1bef541ce}; full qualifications and provenance.}\par
\paragraph*{II.1.3.9.14.2 [CT3] Significant Activity Audit Capture}
\phantomsection
\label{navigation:II.1.3.9.14.2}
\addcontentsline{toc}{paragraph}{II.1.3.9.14.2 [CT3] Significant Activity Audit Capture}

\label{element:d903948169eb4785b0f03d916c842a8a}
\index{[CT3] Significant Activity Audit Capture (II.1.3.9.14.2)}
{\small Capability · Proposed · pending\_adjudication}\par

Capture required security- and information-significant activity as governed, attributable audit evidence.



\textbf{Overview}



Mandatory evidence capture has its own policy-governed platform authority and must remain possible when the audited identity is denied an action. The captured fact is distinguished from debug output and the identity whose activity is recorded.



\textbf{Intent}



Ensure required significant activity evidence can be captured under the authority specifically assigned to that responsibility.



\textbf{Boundary}



The specific intrinsic authority for mandatory audit capture is not generalised to other activities. Capture is not an assurance finding and cannot expose unmasked PHI in diagnostic logs.


{\small Recorded unresolved matters: source-qualified matters; exact scope remains in the canonical record.}\par
{\footnotesize \href{../../../../continuum/strategy/capability/d903948169eb4785b0f03d916c842a8a.md}{Canonical record: d903948169eb4785b0f03d916c842a8a}; full qualifications and provenance.}\par
\subparagraph*{II.1.3.9.14.2.1 [Feature] Mandatory Significant Activity Audit Capture}
\phantomsection
\label{navigation:II.1.3.9.14.2.1}
\addcontentsline{toc}{subparagraph}{II.1.3.9.14.2.1 [Feature] Mandatory Significant Activity Audit Capture}

\label{element:b69936757c4446a7a1b57a455aae7a5f}
\index{[Feature] Mandatory Significant Activity Audit Capture (II.1.3.9.14.2.1)}
{\small Capability / Feature · Proposed · pending\_adjudication}\par

Capture required significant activity evidence using the specifically governed authoritative audit-capture boundary.



\textbf{Overview}



The actor whose action is recorded need not be entitled to persist the audit. Required capture uses its own established authority and failure can block the governed operation.



\textbf{Boundary}



Specific audit authority is not borrowed from the audited identity or generalised to other operations.


{\small Recorded unresolved matters: source-qualified matters; exact scope remains in the canonical record.}\par
{\footnotesize \href{../../../../continuum/strategy/capability/feature/b69936757c4446a7a1b57a455aae7a5f.md}{Canonical record: b69936757c4446a7a1b57a455aae7a5f}; full qualifications and provenance.}\par
\subsubsection*{II.1.3.9.15 [CT2] Patient Clinical Record}
\phantomsection
\label{navigation:II.1.3.9.15}
\addcontentsline{toc}{subsubsection}{II.1.3.9.15 [CT2] Patient Clinical Record}

{\small Historical predecessor: R.2.206 \hyperref[element:2e4de521bda24ad7b0ea64ae535b7f84]{[CT2] Patient Clinical Record}; Canonical ID: \texttt{\seqsplit{2e4de521bda24ad7b0ea64ae535b7f84}}. Superseded historical identity — direct correspondence · superseded.}\par
\label{element:5478878e0e6e4d34a8d5023ac6ac2d54}
\index{[CT2] Patient Clinical Record (II.1.3.9.15)}
{\small Capability · Canonical · current}\par

Patient Clinical Record provides a governed longitudinal representation across encounters, providers and source systems. Assembly supplies the timeline, active access supports current care context, and durable preservation retains the history of what Harmonia received, knew or managed. These contributions are complementary: a responsive working view must not become the authority for durable history.



\textbf{Overview}



Patient Clinical Record provides a governed longitudinal representation across encounters, providers and source systems. Assembly supplies the timeline, active access supports current care context, and durable preservation retains the history of what Harmonia received, knew or managed. These contributions are complementary: a responsive working view must not become the authority for durable history.



\textbf{Intent}



Patient Clinical Record provides a governed longitudinal representation across encounters, providers and source systems. Assembly supplies the timeline, active access supports current care context, and durable preservation retains the history of what Harmonia received, knew or managed. These contributions are complementary: a responsive working view must not become the authority for durable history.



\textbf{Boundary}



Information administration does not confer authority over the represented person, organisation, service, place, device or clinical activity.



\medskip\hrule\medskip

{\footnotesize \href{../../../../continuum/strategy/capability/5478878e0e6e4d34a8d5023ac6ac2d54.md}{Canonical record: 5478878e0e6e4d34a8d5023ac6ac2d54}; full qualifications and provenance.}\par
\paragraph*{II.1.3.9.15.1 [CT3] Linked Information Curation}
\phantomsection
\label{navigation:II.1.3.9.15.1}
\addcontentsline{toc}{paragraph}{II.1.3.9.15.1 [CT3] Linked Information Curation}

{\small Historical predecessor: R.2.204 \hyperref[element:9d0925bb39664a5985cacf3cea7442d1]{[CT3] Linked Information Curation}; Canonical ID: \texttt{\seqsplit{9d0925bb39664a5985cacf3cea7442d1}}. Superseded historical identity — direct correspondence · superseded.}\par
\label{element:5edb2a8b5ee54ddda8104d1fab323a94}
\index{[CT3] Linked Information Curation (II.1.3.9.15.1)}
{\small Capability · Proposed · current}\par

Curate associations to information and binary content needed to interpret and use a clinical record, retaining custody and source-authority distinctions.



\textbf{Overview}



Curate associations to information and binary content needed to interpret and use a clinical record, retaining custody and source-authority distinctions.



\textbf{Intent}



Curate associations to information and binary content needed to interpret and use a clinical record, retaining custody and source-authority distinctions.



\textbf{Boundary}



Association is not physical custody. Persistence authority, source clinical meaning and external repository authority remain separate.


{\footnotesize \href{../../../../continuum/strategy/capability/5edb2a8b5ee54ddda8104d1fab323a94.md}{Canonical record: 5edb2a8b5ee54ddda8104d1fab323a94}; full qualifications and provenance.}\par
\subparagraph*{II.1.3.9.15.1.1 [Feature] Binary Content Resolution \& Assembly}
\phantomsection
\label{navigation:II.1.3.9.15.1.1}
\addcontentsline{toc}{subparagraph}{II.1.3.9.15.1.1 [Feature] Binary Content Resolution \& Assembly}

\label{element:8fb4036f7a9c47c3883b2507528fab5e}
\index{[Feature] Binary Content Resolution \& Assembly (II.1.3.9.15.1.1)}
{\small Capability / Feature · Proposed · current}\par

Resolve governed payload associations and retrieve or assemble the content required by a requesting information context.



\textbf{Overview}



A logical clinical item may reference externally held or multipart payloads. The Feature supplies a usable content result while retaining custody and missing-content standing.



\textbf{Boundary}



No assumption of local copying, shared storage, universal content retention or originating clinical authority.


{\footnotesize \href{../../../../continuum/strategy/capability/feature/8fb4036f7a9c47c3883b2507528fab5e.md}{Canonical record: 8fb4036f7a9c47c3883b2507528fab5e}; full qualifications and provenance.}\par
\paragraph*{II.1.3.9.15.2 [CT3] Longitudinal Clinical Record Administration}
\phantomsection
\label{navigation:II.1.3.9.15.2}
\addcontentsline{toc}{paragraph}{II.1.3.9.15.2 [CT3] Longitudinal Clinical Record Administration}

{\small Historical predecessor: R.2.205 \hyperref[element:bd91d848cf424c42bec7c80cd32edd28]{[CT3] Longitudinal Clinical Record Administration}; Canonical ID: \texttt{\seqsplit{bd91d848cf424c42bec7c80cd32edd28}}. Superseded historical identity — direct correspondence · superseded.}\par
\label{element:14748dc6117f4b6ba7c0fc554b76bd89}
\index{[CT3] Longitudinal Clinical Record Administration (II.1.3.9.15.2)}
{\small Capability · Proposed · current}\par

Administer longitudinal record assembly, active clinical summaries and durable received/managed history with source attribution and supersession.



\textbf{Overview}



Administer longitudinal record assembly, active clinical summaries and durable received/managed history with source attribution and supersession. The reviewed hierarchy assigns this enduring ability to its stated business-enabling context.



\textbf{Intent}



Administer longitudinal record assembly, active clinical summaries and durable received/managed history with source attribution and supersession.



\textbf{Boundary}



Information administration and coordination preserve the authority of the responsible source, clinical actor and governed domain; they confer no additional real-world management authority.


{\footnotesize \href{../../../../continuum/strategy/capability/14748dc6117f4b6ba7c0fc554b76bd89.md}{Canonical record: 14748dc6117f4b6ba7c0fc554b76bd89}; full qualifications and provenance.}\par
\subparagraph*{II.1.3.9.15.2.1 [Feature] Active Clinical Record Access}
\phantomsection
\label{navigation:II.1.3.9.15.2.1}
\addcontentsline{toc}{subparagraph}{II.1.3.9.15.2.1 [Feature] Active Clinical Record Access}

\label{element:76b35ecac42b4781bfd7911c9931f5f1}
\index{[Feature] Active Clinical Record Access (II.1.3.9.15.2.1)}
{\small Capability / Feature · Canonical · current}\par

Provide low-latency, active distributed access to consolidated patient clinical summaries.



Active access supports responsive use of consolidated clinical summaries in ongoing care. The working representation remains reconstructable active state; speed and availability do not promote it to authoritative durable truth when persistence or convergence is unavailable.


{\footnotesize \href{../../../../continuum/strategy/capability/feature/76b35ecac42b4781bfd7911c9931f5f1.md}{Canonical record: 76b35ecac42b4781bfd7911c9931f5f1}; full qualifications and provenance.}\par
\subparagraph*{II.1.3.9.15.2.2 [Feature] Clinical Information Historical Reconstruction}
\phantomsection
\label{navigation:II.1.3.9.15.2.2}
\addcontentsline{toc}{subparagraph}{II.1.3.9.15.2.2 [Feature] Clinical Information Historical Reconstruction}

\label{element:838728ef4fd04b4da0c28dc4fce0cfe9}
\index{[Feature] Clinical Information Historical Reconstruction (II.1.3.9.15.2.2)}
{\small Capability / Feature · Proposed · current}\par

Reconstruct what information Harmonia received, knew or managed at a relevant historical basis.



\textbf{Overview}



Current state and historical operational truth answer different questions. Reconstruction retains superseded and corrected knowledge without rewriting the earlier record.



\textbf{Boundary}



Does not infer that historical information was clinically true or that a present source state applied in the past.


{\footnotesize \href{../../../../continuum/strategy/capability/feature/838728ef4fd04b4da0c28dc4fce0cfe9.md}{Canonical record: 838728ef4fd04b4da0c28dc4fce0cfe9}; full qualifications and provenance.}\par
\subparagraph*{II.1.3.9.15.2.3 [Feature] Durable Clinical Record Preservation}
\phantomsection
\label{navigation:II.1.3.9.15.2.3}
\addcontentsline{toc}{subparagraph}{II.1.3.9.15.2.3 [Feature] Durable Clinical Record Preservation}

\label{element:aba48dbf729d477aa87e112725dc1511}
\index{[Feature] Durable Clinical Record Preservation (II.1.3.9.15.2.3)}
{\small Capability / Feature · Canonical · current}\par

Durably maintain governed, vendor-neutral longitudinal clinical representation and preservation with verifiable integrity controls.



Durable preservation retains the governed longitudinal representation and the history of what Harmonia received or managed. Corrections establish changed knowledge without retrospectively rewriting that history, and preservation neither originates clinical facts nor automatically establishes legal qualification of the record.


{\footnotesize \href{../../../../continuum/strategy/capability/feature/aba48dbf729d477aa87e112725dc1511.md}{Canonical record: aba48dbf729d477aa87e112725dc1511}; full qualifications and provenance.}\par
\subparagraph*{II.1.3.9.15.2.4 [Feature] Longitudinal Health Record (LHR) Assembly}
\phantomsection
\label{navigation:II.1.3.9.15.2.4}
\addcontentsline{toc}{subparagraph}{II.1.3.9.15.2.4 [Feature] Longitudinal Health Record (LHR) Assembly}

\label{element:c9c1c7a58cb5470db0141329ecc349f4}
\index{[Feature] Longitudinal Health Record (LHR) Assembly (II.1.3.9.15.2.4)}
{\small Capability / Feature · Canonical · current}\par

Assemble disparate encounters, diagnostic reports, medications, and clinical documents into a unified, vendor-neutral longitudinal timeline.



Longitudinal assembly relates contributions from disparate systems into a usable chronology. The unified view preserves source and subject context rather than turning consolidation into originating clinical authority or erasing differences in the standing of individual assertions.


{\footnotesize \href{../../../../continuum/strategy/capability/feature/c9c1c7a58cb5470db0141329ecc349f4.md}{Canonical record: c9c1c7a58cb5470db0141329ecc349f4}; full qualifications and provenance.}\par
\paragraph*{II.1.3.9.15.3 [CT3] Primary Care Coordination}
\phantomsection
\label{navigation:II.1.3.9.15.3}
\addcontentsline{toc}{paragraph}{II.1.3.9.15.3 [CT3] Primary Care Coordination}

\label{element:f80e73ee91624ce5baa11ae15480fd35}
\index{[CT3] Primary Care Coordination (II.1.3.9.15.3)}
{\small Capability · Canonical · current}\par

Definition unresolved in the baseline. No definition has been synthesised from the name.



\medskip\hrule\medskip

{\footnotesize \href{../../../../continuum/strategy/capability/f80e73ee91624ce5baa11ae15480fd35.md}{Canonical record: f80e73ee91624ce5baa11ae15480fd35}; full qualifications and provenance.}\par
\subparagraph*{II.1.3.9.15.3.1 [Feature] General Practice Shared Record Ingestion}
\phantomsection
\label{navigation:II.1.3.9.15.3.1}
\addcontentsline{toc}{subparagraph}{II.1.3.9.15.3.1 [Feature] General Practice Shared Record Ingestion}

\label{element:869b77462ece488788c5a20598e7d587}
\index{[Feature] General Practice Shared Record Ingestion (II.1.3.9.15.3.1)}
{\small Capability / Feature · Canonical · current}\par

Ingest shared health summaries, encounter notes, and immunization records from primary care practice management systems.



Shared-record ingestion brings general-practice contributions into cross-provider continuity. The summary or encounter note retains its clinical source while becoming available alongside regional information; consolidation does not transfer the GP's professional authority to the platform.


{\footnotesize \href{../../../../continuum/strategy/capability/feature/869b77462ece488788c5a20598e7d587.md}{Canonical record: 869b77462ece488788c5a20598e7d587}; full qualifications and provenance.}\par
\subparagraph*{II.1.3.9.15.3.2 [Feature] Primary Care Notification Dispatch}
\phantomsection
\label{navigation:II.1.3.9.15.3.2}
\addcontentsline{toc}{subparagraph}{II.1.3.9.15.3.2 [Feature] Primary Care Notification Dispatch}

\label{element:56739765fcfd4084b38d4550593ee74d}
\index{[Feature] Primary Care Notification Dispatch (II.1.3.9.15.3.2)}
{\small Capability / Feature · Canonical · current}\par

Dispatch emergency department arrival, acute admission, and discharge notifications to the patient's registered general practitioner.



Primary-care notifications inform the registered GP of significant acute-care events. This outward notification responsibility is distinct from querying a chronic-care timeline: the latter shares context but has no established equivalence to dispatching the required arrival, admission or discharge notice.


{\footnotesize \href{../../../../continuum/strategy/capability/feature/56739765fcfd4084b38d4550593ee74d.md}{Canonical record: 56739765fcfd4084b38d4550593ee74d}; full qualifications and provenance.}\par
\paragraph*{II.1.3.9.15.4 [CT3] Remote Patient Monitoring}
\phantomsection
\label{navigation:II.1.3.9.15.4}
\addcontentsline{toc}{paragraph}{II.1.3.9.15.4 [CT3] Remote Patient Monitoring}

\label{element:2d1addc7dd1d4821b71a11c692d35b27}
\index{[CT3] Remote Patient Monitoring (II.1.3.9.15.4)}
{\small Capability · Canonical · current}\par

Definition unresolved in the baseline. No definition has been synthesised from the name.



\medskip\hrule\medskip

{\footnotesize \href{../../../../continuum/strategy/capability/2d1addc7dd1d4821b71a11c692d35b27.md}{Canonical record: 2d1addc7dd1d4821b71a11c692d35b27}; full qualifications and provenance.}\par
\subparagraph*{II.1.3.9.15.4.1 [Feature] Remote Telemetry Ingestion}
\phantomsection
\label{navigation:II.1.3.9.15.4.1}
\addcontentsline{toc}{subparagraph}{II.1.3.9.15.4.1 [Feature] Remote Telemetry Ingestion}

\label{element:969dde193fab4c5d82d7dc210c7f5a71}
\index{[Feature] Remote Telemetry Ingestion (II.1.3.9.15.4.1)}
{\small Capability / Feature · Canonical · current}\par

Ingest home patient monitoring data, biometric telemetry, and patient-reported outcome measures (PROMs).



Remote telemetry ingestion accepts measurements and patient-reported contributions from outside the usual care site. The information can inform subsequent monitoring and interpretation, but acquisition alone does not establish clinical significance or authorise an intervention.


{\footnotesize \href{../../../../continuum/strategy/capability/feature/969dde193fab4c5d82d7dc210c7f5a71.md}{Canonical record: 969dde193fab4c5d82d7dc210c7f5a71}; full qualifications and provenance.}\par
\subparagraph*{II.1.3.9.15.4.2 [Feature] Virtual Care Session Context Binding}
\phantomsection
\label{navigation:II.1.3.9.15.4.2}
\addcontentsline{toc}{subparagraph}{II.1.3.9.15.4.2 [Feature] Virtual Care Session Context Binding}

\label{element:c25d66fd6aff4299857eadc7c0243544}
\index{[Feature] Virtual Care Session Context Binding (II.1.3.9.15.4.2)}
{\small Capability / Feature · Canonical · current}\par

Associate synchronous virtual/telehealth consultations with the longitudinal health record and active care episode.



Virtual-session binding associates a consultation with the enduring record and the relevant episode. It preserves the clinical context across delivery modes without turning the communication session itself into proof of clinical conclusions or adequacy of care.


{\footnotesize \href{../../../../continuum/strategy/capability/feature/c25d66fd6aff4299857eadc7c0243544.md}{Canonical record: c25d66fd6aff4299857eadc7c0243544}; full qualifications and provenance.}\par
\subsubsection*{II.1.3.9.16 [CT2] Presentation Services}
\phantomsection
\label{navigation:II.1.3.9.16}
\addcontentsline{toc}{subsubsection}{II.1.3.9.16 [CT2] Presentation Services}

{\small Historical predecessor: R.2.208 \hyperref[element:c89e7c8b277247d9965fdb534aad2d1f]{[CT2] Presentation Services}; Canonical ID: \texttt{\seqsplit{c89e7c8b277247d9965fdb534aad2d1f}}. Superseded historical identity — direct correspondence · superseded.}\par
\label{element:e9b6ba1917ba44cca7466444d0836c98}
\index{[CT2] Presentation Services (II.1.3.9.16)}
{\small Capability · Canonical · current}\par

Presentation Services supplies contextual human interaction through clinical timelines, directory views and operational dashboards. It makes governed information usable by people while consuming the identity, security and clinical context established elsewhere. The responsibility is the interaction and presentation experience rather than a new authority for the information shown.



\textbf{Overview}



Presentation Services supplies contextual human interaction through clinical timelines, directory views and operational dashboards. It makes governed information usable by people while consuming the identity, security and clinical context established elsewhere. The responsibility is the interaction and presentation experience rather than a new authority for the information shown.



\textbf{Intent}



Presentation Services supplies contextual human interaction through clinical timelines, directory views and operational dashboards. It makes governed information usable by people while consuming the identity, security and clinical context established elsewhere. The responsibility is the interaction and presentation experience rather than a new authority for the information shown.



\textbf{Boundary}



Information administration does not confer authority over the represented person, organisation, service, place, device or clinical activity.



\medskip\hrule\medskip

{\footnotesize \href{../../../../continuum/strategy/capability/e9b6ba1917ba44cca7466444d0836c98.md}{Canonical record: e9b6ba1917ba44cca7466444d0836c98}; full qualifications and provenance.}\par
\paragraph*{II.1.3.9.16.1 [CT3] Contextual Information Presentation}
\phantomsection
\label{navigation:II.1.3.9.16.1}
\addcontentsline{toc}{paragraph}{II.1.3.9.16.1 [CT3] Contextual Information Presentation}

{\small Historical predecessor: R.2.207 \hyperref[element:e349d5f6a8774bbe8d150f9c51af770c]{[CT3] Contextual Information Presentation}; Canonical ID: \texttt{\seqsplit{e349d5f6a8774bbe8d150f9c51af770c}}. Superseded historical identity — direct correspondence · superseded.}\par
\label{element:9f900c53ae56488cad69e53c2205e898}
\index{[CT3] Contextual Information Presentation (II.1.3.9.16.1)}
{\small Capability · Proposed · current}\par

Present governed longitudinal clinical and operational information in its applicable human interaction context.



\textbf{Overview}



Present governed longitudinal clinical and operational information in its applicable human interaction context. The reviewed hierarchy assigns this enduring ability to its stated business-enabling context.



\textbf{Intent}



Present governed longitudinal clinical and operational information in its applicable human interaction context.



\textbf{Boundary}



Information administration and coordination preserve the authority of the responsible source, clinical actor and governed domain; they confer no additional real-world management authority.


{\footnotesize \href{../../../../continuum/strategy/capability/9f900c53ae56488cad69e53c2205e898.md}{Canonical record: 9f900c53ae56488cad69e53c2205e898}; full qualifications and provenance.}\par
\subparagraph*{II.1.3.9.16.1.1 [Feature] Longitudinal Clinical \& Operational Presentation}
\phantomsection
\label{navigation:II.1.3.9.16.1.1}
\addcontentsline{toc}{subparagraph}{II.1.3.9.16.1.1 [Feature] Longitudinal Clinical \& Operational Presentation}

\label{element:e82eaafa2e4046f48d4c86e15f282a0b}
\index{[Feature] Longitudinal Clinical \& Operational Presentation (II.1.3.9.16.1.1)}
{\small Capability / Feature · Canonical · current}\par

Deliver contextual clinical timelines, directory administration views, and operational management dashboards.



Contextual presentation makes clinical timelines, directories and operational information usable for people. It consumes established identity and security context while preserving the originating authority of what is shown; a user interface does not become an alternative persistence or policy boundary.


{\footnotesize \href{../../../../continuum/strategy/capability/feature/e82eaafa2e4046f48d4c86e15f282a0b.md}{Canonical record: e82eaafa2e4046f48d4c86e15f282a0b}; full qualifications and provenance.}\par
\subsubsection*{II.1.3.9.17 [CT2] Semantic \& Reference Knowledge Enablement [superseded]}
\phantomsection
\label{navigation:II.1.3.9.17}
\addcontentsline{toc}{subsubsection}{II.1.3.9.17 [CT2] Semantic \& Reference Knowledge Enablement [superseded]}

\label{element:92a71292710c41b29775d19b98c197d6}
\index{[CT2] Semantic \& Reference Knowledge Enablement (II.1.3.9.17)}
{\small Capability · Proposed · superseded}\par
\href{../../../../continuum/strategy/capability/92a71292710c41b29775d19b98c197d6.md}{Historical canonical definition and provenance}.\par
\subsubsection*{II.1.3.9.18 [CT2] Workflow \& Activity Coordination}
\phantomsection
\label{navigation:II.1.3.9.18}
\addcontentsline{toc}{subsubsection}{II.1.3.9.18 [CT2] Workflow \& Activity Coordination}

{\small Historical predecessor: R.2.213 \hyperref[element:4422d865978c4353ad575759c6026d27]{[CT2] Workflow \& Activity Coordination}; Canonical ID: \texttt{\seqsplit{4422d865978c4353ad575759c6026d27}}. Superseded historical identity — direct correspondence · superseded.}\par
\label{element:c724f3d38bfa4567983a11ad22095d61}
\index{[CT2] Workflow \& Activity Coordination (II.1.3.9.18)}
{\small Capability · Proposed · current}\par

Coordinate tasks entering or leaving the Harmonia execution boundary and work undertaken by other systems, with responsibility/destination routing, external status/traceability and governance of the completion bridge.



\textbf{Overview}



Coordinate tasks entering or leaving the Harmonia execution boundary and work undertaken by other systems, with responsibility/destination routing, external status/traceability and governance of the completion bridge. The reviewed hierarchy assigns this enduring ability to its stated business-enabling context.



\textbf{Intent}



Coordinate tasks entering or leaving the Harmonia execution boundary and work undertaken by other systems, with responsibility/destination routing, external status/traceability and governance of the completion bridge.



\textbf{Boundary}



External means outside Harmonia execution, potentially inside the same health organisation. Work coordination includes completion responsibilities and traceability; it is not generic interoperability or message routing.


{\footnotesize \href{../../../../continuum/strategy/capability/c724f3d38bfa4567983a11ad22095d61.md}{Canonical record: c724f3d38bfa4567983a11ad22095d61}; full qualifications and provenance.}\par
\paragraph*{II.1.3.9.18.1 [CT3] Activity Routing Management}
\phantomsection
\label{navigation:II.1.3.9.18.1}
\addcontentsline{toc}{paragraph}{II.1.3.9.18.1 [CT3] Activity Routing Management}

{\small Historical predecessor: R.2.209 \hyperref[element:243ea93e7b7448fe97dc74f90bfc8946]{[CT3] Activity Routing Management}; Canonical ID: \texttt{\seqsplit{243ea93e7b7448fe97dc74f90bfc8946}}. Superseded historical identity — direct correspondence · superseded.}\par
\label{element:27411c143e6e41ffba1af139fbbcd7e2}
\index{[CT3] Activity Routing Management (II.1.3.9.18.1)}
{\small Capability · Proposed · current}\par

Manage responsibility and destination routing for activity crossing the Harmonia execution boundary, respecting the intended work and completion responsibility.



\textbf{Overview}



Manage responsibility and destination routing for activity crossing the Harmonia execution boundary, respecting the intended work and completion responsibility. The reviewed hierarchy assigns this enduring ability to its stated business-enabling context.



\textbf{Intent}



Manage responsibility and destination routing for activity crossing the Harmonia execution boundary, respecting the intended work and completion responsibility.



\textbf{Boundary}



Information administration and coordination preserve the authority of the responsible source, clinical actor and governed domain; they confer no additional real-world management authority.


{\footnotesize \href{../../../../continuum/strategy/capability/27411c143e6e41ffba1af139fbbcd7e2.md}{Canonical record: 27411c143e6e41ffba1af139fbbcd7e2}; full qualifications and provenance.}\par
\paragraph*{II.1.3.9.18.2 [CT3] External Task Coordination}
\phantomsection
\label{navigation:II.1.3.9.18.2}
\addcontentsline{toc}{paragraph}{II.1.3.9.18.2 [CT3] External Task Coordination}

{\small Historical predecessor: R.2.210 \hyperref[element:32be309d65594592bac4d7b068def0d2]{[CT3] External Task Coordination}; Canonical ID: \texttt{\seqsplit{32be309d65594592bac4d7b068def0d2}}. Superseded historical identity — direct correspondence · superseded.}\par
\label{element:ff3447c929724245a13ee864874b5561}
\index{[CT3] External Task Coordination (II.1.3.9.18.2)}
{\small Capability · Proposed · current}\par

Coordinate work entering or leaving the Harmonia execution boundary with the system or actor responsible for undertaking and completing that work.



\textbf{Overview}



Coordinate work entering or leaving the Harmonia execution boundary with the system or actor responsible for undertaking and completing that work. The reviewed hierarchy assigns this enduring ability to its stated business-enabling context.



\textbf{Intent}



Coordinate work entering or leaving the Harmonia execution boundary with the system or actor responsible for undertaking and completing that work.



\textbf{Boundary}



Information administration and coordination preserve the authority of the responsible source, clinical actor and governed domain; they confer no additional real-world management authority.


{\footnotesize \href{../../../../continuum/strategy/capability/ff3447c929724245a13ee864874b5561.md}{Canonical record: ff3447c929724245a13ee864874b5561}; full qualifications and provenance.}\par
\paragraph*{II.1.3.9.18.3 [CT3] External Task Status \& Traceability Management}
\phantomsection
\label{navigation:II.1.3.9.18.3}
\addcontentsline{toc}{paragraph}{II.1.3.9.18.3 [CT3] External Task Status \& Traceability Management}

{\small Historical predecessor: R.2.211 \hyperref[element:95985d244a234571892a9e9b4024f594]{[CT3] External Task Status \& Traceability Management}; Canonical ID: \texttt{\seqsplit{95985d244a234571892a9e9b4024f594}}. Superseded historical identity — direct correspondence · superseded.}\par
\label{element:b25bb04a8aa94dcfac105a646878ccfc}
\index{[CT3] External Task Status \& Traceability Management (II.1.3.9.18.3)}
{\small Capability · Proposed · current}\par

Manage the observable status and attributable traceability of work as it progresses through systems outside the Harmonia execution boundary.



\textbf{Overview}



Manage the observable status and attributable traceability of work as it progresses through systems outside the Harmonia execution boundary. The reviewed hierarchy assigns this enduring ability to its stated business-enabling context.



\textbf{Intent}



Manage the observable status and attributable traceability of work as it progresses through systems outside the Harmonia execution boundary.



\textbf{Boundary}



Information administration and coordination preserve the authority of the responsible source, clinical actor and governed domain; they confer no additional real-world management authority.


{\footnotesize \href{../../../../continuum/strategy/capability/b25bb04a8aa94dcfac105a646878ccfc.md}{Canonical record: b25bb04a8aa94dcfac105a646878ccfc}; full qualifications and provenance.}\par
\paragraph*{II.1.3.9.18.4 [CT3] Work Boundary Governance}
\phantomsection
\label{navigation:II.1.3.9.18.4}
\addcontentsline{toc}{paragraph}{II.1.3.9.18.4 [CT3] Work Boundary Governance}

{\small Historical predecessor: R.2.212 \hyperref[element:b4d14abf930a4230a7f5b93e9ce3df91]{[CT3] Work Boundary Governance}; Canonical ID: \texttt{\seqsplit{b4d14abf930a4230a7f5b93e9ce3df91}}. Superseded historical identity — direct correspondence · superseded.}\par
\label{element:35b66a0b38db4568849ecb842ffdcbf3}
\index{[CT3] Work Boundary Governance (II.1.3.9.18.4)}
{\small Capability · Proposed · current}\par

Govern boundary conditions and completion responsibilities of the bridge between Harmonia and another system undertaking work.



\textbf{Overview}



Govern boundary conditions and completion responsibilities of the bridge between Harmonia and another system undertaking work. The reviewed hierarchy assigns this enduring ability to its stated business-enabling context.



\textbf{Intent}



Govern boundary conditions and completion responsibilities of the bridge between Harmonia and another system undertaking work.



\textbf{Boundary}



Information administration and coordination preserve the authority of the responsible source, clinical actor and governed domain; they confer no additional real-world management authority.


{\footnotesize \href{../../../../continuum/strategy/capability/35b66a0b38db4568849ecb842ffdcbf3.md}{Canonical record: 35b66a0b38db4568849ecb842ffdcbf3}; full qualifications and provenance.}\par
\chapter*{II.2 Courses of Action}
\phantomsection
\label{navigation:II.2}
\addcontentsline{toc}{chapter}{II.2 Courses of Action}

\section*{II.2.1 Boundary Membrane Sovereignty}
\phantomsection
\label{navigation:II.2.1}
\addcontentsline{toc}{section}{II.2.1 Boundary Membrane Sovereignty}

\label{element:25d8509146fd4f2794a17f4a5253b7a3}
\index{Boundary Membrane Sovereignty (II.2.1)}
{\small CourseOfAction · Canonical · current}\par

\par\medskip\textbf{Strategic Approach}\par


\begin{quote}


\textbf{Preserve strict separation between external standards-compliant exchange representations and internal Harmonia-governed operational semantics.}



\end{quote}


Harmonia acts as an integration environment bridging heterogeneous external healthcare systems and internal integration services. This course of action establishes that:
- External exchange contracts (such as FHIR REST interactions, MLLP message streams, or directory protocols) govern only the interaction semantics at the enterprise boundary.
- External or standards-based representations MAY remain native Harmonia information representations where appropriate. Their use SHALL NOT cause an external exchange contract or representation to define Harmonia-private management, execution, governance or coordination semantics.
- Egress across the boundary terminates Harmonia's governance of the emitted representation.



\par\medskip\textbf{Motivational Grounding \& Traceability}\par


\begin{itemize}


\item \textbf{Architectural Axioms}:

\item \texttt{\seqsplit{AX-02}} (Boundary Membrane): External representations are boundary projections; internal semantics are sovereign.

\item \texttt{\seqsplit{AX-03}} (Native Standards): Standards govern interoperability at the boundary; standards are not distorted to fit private operational models.

\item \texttt{\seqsplit{AX-13}} (Explicit Management Boundary): Egress terminates platform management; internal operational metadata does not leak across the boundary.

\item \texttt{\seqsplit{AX-14}} (Preserve Distinctions): Distinctions between external exchange syntax and internal operational meaning are preserved.

\item \textbf{Strategic Drivers \& Goals}: Interoperability Mandates, Vendor Independence, Unified Ingress/Egress Governance.

\end{itemize}


\par\medskip\textbf{Capability Realisation \& Configuration}\par


\begin{itemize}


\item \textbf{Primary Capabilities}: Shapes \texttt{\seqsplit{EC-08}} (Interoperability \& Exchange) and \texttt{\seqsplit{EC-13}} (Semantic Governance \& Conformance).

\item \textbf{Component Boundary Impact}: Informs the boundary responsibility of \textbf{Pylai}, establishing it as the standards-facing membrane rather than an internal processing engine or generic transport layer.

\end{itemize}


\par\medskip\textbf{Downstream Direction (Domains 03–13)}\par


\begin{itemize}


\item Application Architecture (Domain 05) and Integration Architecture (Domain 06) must implement fail-closed boundary projections, ensuring private coordination metadata (e.g., distributed transaction IDs, internal routing tags) is never emitted externally.

\item Information Architecture (Domain 04) preserves conceptual meaning and responsibility independently of external contracts. Native standards representations remain permitted under AX-03; no parallel information model is required merely to accommodate Harmonia-private management concerns.

\end{itemize}

{\small Recorded unresolved matters: canonicalName; exact scope remains in the canonical record.}\par
{\footnotesize \href{../../../../continuum/strategy/course-of-action/25d8509146fd4f2794a17f4a5253b7a3.md}{Canonical record: 25d8509146fd4f2794a17f4a5253b7a3}; full qualifications and provenance.}\par
\section*{II.2.2 Collaborative Reusable Capability Realisation}
\phantomsection
\label{navigation:II.2.2}
\addcontentsline{toc}{section}{II.2.2 Collaborative Reusable Capability Realisation}

\label{element:503939e1fa464b78b1d980ddedfeb3cf}
\index{Collaborative Reusable Capability Realisation (II.2.2)}
{\small CourseOfAction · Canonical · current}\par

\par\medskip\textbf{Strategic Approach}\par


\begin{quote}


\textbf{Realize bounded reusable platform capabilities (context, policy, provenance, EC-12 Operational Assurance) collaboratively across participating components rather than through centralized, bottlenecked runtime services.}



\end{quote}


Harmonia requires pervasive platform capabilities—such as security policy evaluation, contextual propagation, operational assurance, and provenance tracking—across all transactions. This course of action establishes that:
- \textbf{Reusable Capability \$\textbackslash{}neq\$ Centralised Service (Guardrail G1)}: The existence of a reusable enterprise capability does not imply that a single monolithic service must execute it centrally.
- Reusable capabilities define bounded responsibilities and common semantic rules. Participating components consume or enforce them through governed contracts without acquiring shared semantic ownership. \textbf{Distributed participation does not imply distributed responsibility. Cross-cutting concern does not imply cross-cutting responsibility.}
- Centralized choke points and synchronous runtime bottlenecks are strictly avoided, ensuring horizontal scalability and fault isolation.



\par\medskip\textbf{Motivational Grounding \& Traceability}\par


\begin{itemize}


\item \textbf{Architectural Axioms}:

\item \texttt{\seqsplit{AX-04}} (Own Semantics / Use Machinery): Maintain sovereign capability rules while delegating execution across components.

\item \texttt{\seqsplit{AX-07}} (Intrinsic Security): Security and policy evaluation must be pervasive and intrinsic to all operations, not an external proxy.

\item \texttt{\seqsplit{AX-11}} (High Availability \& Responsiveness): Eliminate single points of congestion and failure across the integration fabric.

\item \textbf{Strategic Drivers \& Goals}: High Scalability, Fault Tolerance, Zero Bottleneck Architecture, Decentralised Resilience.

\end{itemize}


\par\medskip\textbf{Capability Realisation \& Configuration}\par


\begin{itemize}


\item \textbf{Primary Capabilities}: Shapes \texttt{\seqsplit{EC-02}} (Context Management), \texttt{\seqsplit{EC-06}} (Policy \& Control), \texttt{\seqsplit{EC-07}} (Provenance \& Traceability), and \texttt{\seqsplit{EC-12}} (Operational Assurance). Assurance here means the established EC-12 operational contribution; no independent Governed Assurance or EC-14 relationship is established.

\item \textbf{Component Boundary Impact}: Prevents components like \textbf{Themis} (Policy \& Control) from becoming synchronous network choke points for every micro-operation; policy contracts are distributed and evaluated collaboratively.

\end{itemize}


\par\medskip\textbf{Downstream Direction (Domains 03–13)}\par


\begin{itemize}


\item Application Architecture (Domain 05) and Technology Architecture (Domain 07) realise reusable responsibilities through governed contracts while preserving bounded ownership, scalability and fault isolation. Existing \texttt{\seqsplit{themis-api}} and \texttt{\seqsplit{petasos-api}} packages are illustrative downstream traceability only; Strategy does not mandate API packaging, local embedding, frameworks or deployment choices.

\end{itemize}

{\small Recorded unresolved matters: canonicalName; exact scope remains in the canonical record.}\par
{\footnotesize \href{../../../../continuum/strategy/course-of-action/503939e1fa464b78b1d980ddedfeb3cf.md}{Canonical record: 503939e1fa464b78b1d980ddedfeb3cf}; full qualifications and provenance.}\par
\section*{II.2.3 Distinct Management and Durable Preservation of Information and State}
\phantomsection
\label{navigation:II.2.3}
\addcontentsline{toc}{section}{II.2.3 Distinct Management and Durable Preservation of Information and State}

\label{element:4c71b89477b34733b31e478485e89b1a}
\index{Distinct Management and Durable Preservation of Information and State (II.2.3)}
{\small CourseOfAction · Canonical · current}\par

\par\medskip\textbf{Strategic Approach}\par


\begin{quote}


\textbf{Maintain a clear separation between the governed management of Harmonia information and state and its durable preservation and recovery, allowing runtime management and persistence concerns to evolve independently without compromising information integrity.}



\end{quote}


This course of action establishes that:
- Runtime management of information, active relationships, context, and state is an operational governance concern requiring high availability, concurrent access, and responsive query navigation.
- Authoritative durable state establishment, preservation and recovery atomically establish authoritative state and authoritative version progression, preserve the durable management metadata required to interpret that state, and ensure historical retention, immutability, and state reconstruction.
- The distinction is strictly one of \textbf{architectural responsibility}, not a crude division between "ephemeral data" and "durable data". Information managed at runtime frequently represents durable healthcare concepts, but runtime management and durable preservation fulfill distinct architectural purposes.



\par\medskip\textbf{Motivational Grounding \& Traceability}\par


\begin{itemize}


\item \textbf{Architectural Axioms}:

\item \texttt{\seqsplit{AX-05}} (Active vs Durable State): Application-facing active state access/coordination and authoritative durable state/version establishment remain separate architectural concerns.

\item \texttt{\seqsplit{AX-09}} (Ephemeral Operational State): Transient execution states are reconstructable and decoupled from durable health records.

\item \texttt{\seqsplit{AX-11}} (High Availability \& Responsiveness): Information access must not be blocked by backend persistence latencies or recovery locks.

\item \texttt{\seqsplit{AX-14}} (Preserve Distinctions): Preserve the boundary between runtime information governance and storage mechanics.

\item \textbf{Strategic Drivers \& Goals}: 24/7 Clinical Operational Continuity, High-Throughput Access, Authoritative Preservation.

\end{itemize}


\par\medskip\textbf{Capability Realisation \& Configuration}\par


\begin{itemize}


\item \textbf{Primary Capabilities}: Shapes \texttt{\seqsplit{EC-03}} (Managed State \& Lifecycle) and \texttt{\seqsplit{EC-04}} (Information Management).

\item \textbf{Component Boundary Impact}: Explicitly establishes the seam between \textbf{Mneme} (application-facing access and active information/state management) and \textbf{Mnemosyne} (authoritative durable state/version establishment, preservation and recovery). Neither component subsumes the other; Mnemosyne does not own operational activity progression, workflow execution, active distributed state, Digital Twin coordination, or application-facing query interfaces. Mneme may reject or coordinate a proposed state progression before persistence; only Mnemosyne establishes the new authoritative durable state. Following authoritative commit, Mneme converges its active representation toward that state. Mneme active-state generation and Mnemosyne authoritative version remain distinct concurrency domains (\texttt{\seqsplit{AX-05}}).

\end{itemize}


\par\medskip\textbf{Downstream Direction (Domains 03–13)}\par


\begin{itemize}


\item Application Architecture (Domain 05) must forbid direct presentation/application access to database/JPA layers (enforcing Invariant 8).

\item Technology Architecture (Domain 07) must allow distributed in-memory data fabrics (for runtime state) and durable relational/document stores (for preservation) to be configured and scaled independently.

\end{itemize}

{\small Recorded unresolved matters: canonicalName; exact scope remains in the canonical record.}\par
{\footnotesize \href{../../../../continuum/strategy/course-of-action/4c71b89477b34733b31e478485e89b1a.md}{Canonical record: 4c71b89477b34733b31e478485e89b1a}; full qualifications and provenance.}\par
\section*{II.2.4 Entity-Centred Operational Coordination}
\phantomsection
\label{navigation:II.2.4}
\addcontentsline{toc}{section}{II.2.4 Entity-Centred Operational Coordination}

\label{element:f601fe1f39a3447f8e9fb2cc4ffe9dee}
\index{Entity-Centred Operational Coordination (II.2.4)}
{\small CourseOfAction · Canonical · current}\par

\par\medskip\textbf{Strategic Approach}\par


\begin{quote}


\textbf{Coordinate governed information, state and operational activity around real-world healthcare entities where those entities are operationally significant in their own right and require active management, without requiring every managed entity to maintain a permanently active execution construct.}



\end{quote}


Healthcare operations naturally center around authentic real-world entities: patients, practitioners, care teams, wards, and beds. This course of action establishes that:
- Harmonia organizes the coordination of state and ongoing activities around the operational identity of the real-world entity itself.
- \textbf{The Digital Twin is the architectural construct} through which this entity-centred coordination strategy is realised. A Twin coordinates information, state, and activity across the Mneme/Ponos seam for a specific real-world entity.
- \textbf{Demand-Driven Lifecycle}: A managed entity does not require a permanently running execution thread or process. Twins are activated when entity-centred operational activity or active monitoring is required, and become quiescent when work concludes.
- \textbf{Representation \$\textbackslash{}neq\$ Twin}: Simply possessing a database record or FHIR resource does not justify or constitute a Digital Twin. Coordination is justified only when the entity is operationally active and evolving in state.



\par\medskip\textbf{Motivational Grounding \& Traceability}\par


\begin{itemize}


\item \textbf{Architectural Axioms}:

\item \texttt{\seqsplit{AX-01}} (Health-Information Centric): Architecture aligns with genuine healthcare entities and clinical relationships.

\item \texttt{\seqsplit{AX-11}} (High Availability \& Responsiveness): High concurrency is achieved by isolating entity coordination without central serialisation bottlenecks.

\item \texttt{\seqsplit{AX-16}} (Activity \& State Progress Together): Entity state and active operational workflows are coordinated dynamically.

\item \textbf{Strategic Drivers \& Goals}: Holistic Patient-Centric Care, High Concurrency without Serialization, Real-Time Healthcare Coordination.

\end{itemize}


\par\medskip\textbf{Capability Realisation \& Configuration}\par


\begin{itemize}


\item \textbf{Primary Capabilities}: Shapes \texttt{\seqsplit{EC-01}} (Managed Entity \& Relationship), \texttt{\seqsplit{EC-02}} (Context Management), and \texttt{\seqsplit{EC-10}} (Activity \& Execution).

\item \textbf{Component Boundary Impact}: Validates the \textbf{Digital Twin} as an active management construct bridging the \textbf{Mneme ↔ Ponos} boundary, while asserting that the Twin is an architectural construct rather than an independent deployable platform component.

\end{itemize}


\par\medskip\textbf{Downstream Direction (Domains 03–13)}\par


\begin{itemize}


\item Application Architecture (Domain 05) must model Twin archetypes using composite context (combining multiple resources such as Practitioner, PractitionerRole, and Location) rather than equating twins to single FHIR resources or permanent background threads.

\end{itemize}

{\small Recorded unresolved matters: canonicalName; exact scope remains in the canonical record.}\par
{\footnotesize \href{../../../../continuum/strategy/course-of-action/f601fe1f39a3447f8e9fb2cc4ffe9dee.md}{Canonical record: f601fe1f39a3447f8e9fb2cc4ffe9dee}; full qualifications and provenance.}\par
\section*{II.2.5 Governed Asynchronous Activity Progression}
\phantomsection
\label{navigation:II.2.5}
\addcontentsline{toc}{section}{II.2.5 Governed Asynchronous Activity Progression}

\label{element:a8ab09ee29454376bd45ba90b320d6ec}
\index{Governed Asynchronous Activity Progression (II.2.5)}
{\small CourseOfAction · Canonical · current}\par

\par\medskip\textbf{Strategic Approach}\par


\begin{quote}


\textbf{Progress multi-stage operational activities through explicit, observable unit-of-work transitions coordinated with entity state, separating execution from boundary exchange.}



\end{quote}


Complex healthcare integrations (e.g., closed-loop diagnostic ordering, multi-destination clinical document fan-out) span multiple distributed systems and cannot rely on brittle synchronous distributed transactions. This course of action establishes that:
- Work progresses through discrete, observable units of work whose lifecycle state transitions are explicitly tracked and auditable.
- Operational progression progresses in lockstep with governed entity state (\texttt{\seqsplit{AX-16}}), preserving explicit uncertainty (\texttt{\seqsplit{AX-15}}) when external responses are pending or failed.
- Execution progression is strictly decoupled from boundary transport and exchange mechanics: activity execution determines \emph{that} an external interaction is needed, but does not own transport protocols or wire-level retries.



\par\medskip\textbf{Motivational Grounding \& Traceability}\par


\begin{itemize}


\item \textbf{Architectural Axioms}:

\item \texttt{\seqsplit{AX-10}} (Distribution \& Normal Failure): Systems fail normally; activity progression must be resilient, idempotent, and restartable.

\item \texttt{\seqsplit{AX-15}} (Preserve Uncertainty): Never assume success or failure; pending, timed-out, or ambiguous outcomes must be explicitly modeled.

\item \texttt{\seqsplit{AX-16}} (Activity \& State Progress Together): Operational activity and managed information state advance in coordinated lockstep.

\item \textbf{Strategic Drivers \& Goals}: Resilient Regional Integration, Asynchronous Clinical Decoupling, Bounded Workflow Recovery.

\end{itemize}


\par\medskip\textbf{Capability Realisation \& Configuration}\par


\begin{itemize}


\item \textbf{Primary Capabilities}: Shapes \texttt{\seqsplit{EC-03}} (Managed State \& Lifecycle), \texttt{\seqsplit{EC-09}} (Event \& Subscription), and \texttt{\seqsplit{EC-10}} (Activity \& Execution).

\item \textbf{Component Boundary Impact}: Grounds the responsibility of \textbf{Ponos} as the execution engine for governed units of work, while strictly enforcing Guardrail G4 (Ponos progresses execution, but does not acquire transport adapters or wire protocols).

\end{itemize}


\par\medskip\textbf{Downstream Direction (Domains 03–13)}\par


\begin{itemize}


\item Application Architecture (Domain 05) and Integration Architecture (Domain 06) realise governed work context, recoverable activity handoffs and explicit per-destination progression according to the applicable acceptance and delivery contracts (preserving Invariants 4 and 5). Established downstream constructs such as Pragma envelopes and Petasos messaging illustrate traceability; their APIs, queues, packaging and runtime mechanisms are not selected by this strategic responsibility.

\end{itemize}

{\small Recorded unresolved matters: canonicalName; exact scope remains in the canonical record.}\par
{\footnotesize \href{../../../../continuum/strategy/course-of-action/a8ab09ee29454376bd45ba90b320d6ec.md}{Canonical record: a8ab09ee29454376bd45ba90b320d6ec}; full qualifications and provenance.}\par
\section*{II.2.6 Meaning-Centric Provenance and Traceability}
\phantomsection
\label{navigation:II.2.6}
\addcontentsline{toc}{section}{II.2.6 Meaning-Centric Provenance and Traceability}

\label{element:fd00bdbdc2234a208c5045aa69279db8}
\index{Meaning-Centric Provenance and Traceability (II.2.6)}
{\small CourseOfAction · Canonical · current}\par

\par\medskip\textbf{Strategic Approach}\par


\begin{quote}


\textbf{Preserve provenance, authority, attribution and traceability for information-significant and business-significant actions and state changes, while avoiding unnecessary elevation of transient operational mechanics into enduring business evidence.}



\end{quote}


Healthcare integration requires undeniable accountability without overwhelming storage or logging systems with low-level operational noise. This course of action establishes that:
- Harmonia SHALL preserve provenance, audit and operational evidence where required by governing obligations, architectural significance, accountability, reconstruction, security or assurance needs. Routine technical activity SHALL NOT automatically require durable evidentiary recording.
- Required semantic attribution records who, what, when, why and under whose authority. Telemetry, operational facts, provenance, audit and assurance evidence remain distinct; being evidence is contextual, not intrinsic to every recorded fact.
- Transient operational mechanics (e.g., thread switches, socket retries, queue polling, network pinging) represent operational telemetry, not business or clinical evidence.
- Strategy sets the architectural requirement for semantic attribution and non-repudiation without dictating specific implementation mechanisms (such as cryptographic signatures, write-once ledgers, or specific storage engines) at this layer.



\par\medskip\textbf{Motivational Grounding \& Traceability}\par


\begin{itemize}


\item \textbf{Architectural Axioms}:

\item \texttt{\seqsplit{AX-06}} (Explicit Authority): Information authority and credibility remain explicit and independent of technical transport or persistence. Required attribution preserves the applicable source or actor authority without making every access evaluation durable evidence.

\item \texttt{\seqsplit{AX-07}} (Intrinsic Security): Security, authorization, and provenance are embedded within operations, not bolted on.

\item \texttt{\seqsplit{AX-08}} (Meaning over Machinery): Record semantic and business meaning, not transient technical machinery.

\item \texttt{\seqsplit{AX-14}} (Preserve Distinctions): Distinguish enduring clinical/business evidence from transient technical logs.

\item \textbf{Strategic Drivers \& Goals}: Regulatory Accountability, Clinical Governance, Legal Auditability, Evidentiary Integrity.

\end{itemize}


\par\medskip\textbf{Capability Realisation \& Configuration}\par


\begin{itemize}


\item \textbf{Primary Capabilities}: Shapes \texttt{\seqsplit{EC-06}} (Policy \& Control) and \texttt{\seqsplit{EC-07}} (Provenance \& Traceability).

\item \textbf{Component Boundary Impact}: Governs how all logical components emit evidence, ensuring that audit trails capture governed business assertions rather than unmasked PHI or low-level machine noise (supporting Invariants 6 and 7).

\end{itemize}


\par\medskip\textbf{Downstream Direction (Domains 03–13)}\par


\begin{itemize}


\item Security Architecture (Domain 08) and Information Architecture (Domain 04) define the governed audit and provenance representations (with Kleio as established downstream evidence traceability), preserving PHI-safe capture and verifiable integrity according to applicable obligations. Strategy does not select a verification technology.

\end{itemize}

{\small Recorded unresolved matters: canonicalName; exact scope remains in the canonical record.}\par
{\footnotesize \href{../../../../continuum/strategy/course-of-action/fd00bdbdc2234a208c5045aa69279db8.md}{Canonical record: fd00bdbdc2234a208c5045aa69279db8}; full qualifications and provenance.}\par
\section*{II.2.7 Implementation Guardrails}
\phantomsection
\label{navigation:II.2.7}
\addcontentsline{toc}{section}{II.2.7 Implementation Guardrails}

\subsection*{II.2.7.1 [ImplementationGuideline] Default Development Test and POC Identity Provider}
\phantomsection
\label{navigation:II.2.7.1}
\addcontentsline{toc}{subsection}{II.2.7.1 [ImplementationGuideline] Default Development Test and POC Identity Provider}

\label{element:f8874dd3bc5c4b68b1ece4eb2859d1d9}
\index{[ImplementationGuideline] Default Development Test and POC Identity Provider (II.2.7.1)}
{\small CourseOfAction / ImplementationGuideline · Adjudicated Domain05 direction — U12/U15/U16; holistic result acceptance pending · current}\par

\begin{quote}


\textbf{Keycloak + PostgreSQL SHALL provide the default FOSS Identity Provider for Harmonia development, test and POC deployments. Keycloak remains an external, replaceable implementation component and SHALL NOT introduce Keycloak-specific dependencies into Harmonia's Identity.Authentication architecture.}



\end{quote}


This assertion is normalised from the governed DesignNote under U12/U15/U16 and the ARB release Praxis. Original candidate/release qualifications and separately deferred implementation details remain in force.



Architectural subject: \hyperref[element:e4998b0a132c41afbe610b11869c3297]{SEC-001 — Trusted Authentication}.


{\small Recorded unresolved matters: canonicalName; exact scope remains in the canonical record.}\par
{\footnotesize \href{../../../../continuum/strategy/course-of-action/implementation-guideline/f8874dd3bc5c4b68b1ece4eb2859d1d9.md}{Canonical record: f8874dd3bc5c4b68b1ece4eb2859d1d9}; full qualifications and provenance.}\par
\chapter*{II.3 Resources}
\phantomsection
\label{navigation:II.3}
\addcontentsline{toc}{chapter}{II.3 Resources}

\section*{II.3.1 Australian National Healthcare Directory \& Identifier Specifications}
\phantomsection
\label{navigation:II.3.1}
\addcontentsline{toc}{section}{II.3.1 Australian National Healthcare Directory \& Identifier Specifications}

\label{element:eac71c56624d43bda2b9d3648f14603a}
\index{Australian National Healthcare Directory \& Identifier Specifications (II.3.1)}
{\small Resource · Canonical · current}\par

\par\medskip\textbf{Strategic Definition}\par


\textbf{Australian National Healthcare Directory \& Identifier Specifications} comprise the jurisdictional normative specifications, identifier schemes, and federation protocols defining how healthcare providers, organisations, and service delivery endpoints are identified, addressed, and verified across Australian digital health.



\par\medskip\textbf{Enduring Strategic Role \& Distinct Facets}\par


Harmonia depends upon external national specifications to participate in federated healthcare discovery and directed clinical message routing. To maintain architectural clarity, Harmonia distinguishes the complementary facets of this resource:
1. \textbf{National Healthcare Directory \& Locator Specifications}:
   - Normative specifications governing federated directory lookup, provider service discovery, and electronic service delivery endpoint resolution (e.g., Australian National Directory specifications, Endpoint Directory Service, and Location Directory Service).
   - Enables Harmonia to resolve destination endpoints for clinical communications without relying on brittle point-to-point configurations.
2. \textbf{Healthcare Identifier (HI) Ecosystem Specifications}:
   - Normative identifier schemes, metadata structures, and verification service specifications for Australian healthcare entities:
     - \textbf{HPI-I}: Healthcare Provider Identifier - Individual (practitioners).
     - \textbf{HPI-O}: Healthcare Provider Identifier - Organisation (health service organisations).
     - \textbf{IHI}: Individual Healthcare Identifier (patients / healthcare recipients).
   - Establishes jurisdictional identity anchors that Harmonia must correlate across diverse regional identifiers (e.g., hospital MRNs, practice local IDs).



\par\medskip\textbf{Traceability}\par


\begin{itemize}


\item \textbf{Architectural Axioms}: \texttt{\seqsplit{AX-01}} (Health-Information Centric), \texttt{\seqsplit{AX-06}} (Explicit Authority), \texttt{\seqsplit{AX-14}} (Preserve Distinctions).

\item \textbf{Enterprise Capabilities}: \texttt{\seqsplit{EC-01}} (Managed Entity \& Relationship), \texttt{\seqsplit{EC-08}} (Interoperability \& Exchange).

\item \textbf{Courses of Action}: \texttt{\seqsplit{COA-01}} (Boundary Membrane Sovereignty), \texttt{\seqsplit{COA-05}} (Entity-Centred Operational Coordination).

\end{itemize}

{\small Recorded unresolved matters: canonicalName; exact scope remains in the canonical record.}\par
{\footnotesize \href{../../../../continuum/strategy/resource/eac71c56624d43bda2b9d3648f14603a.md}{Canonical record: eac71c56624d43bda2b9d3648f14603a}; full qualifications and provenance.}\par
\section*{II.3.2 Healthcare Interoperability Standards and Specifications}
\phantomsection
\label{navigation:II.3.2}
\addcontentsline{toc}{section}{II.3.2 Healthcare Interoperability Standards and Specifications}

\label{element:4d725c804fb34dfa90486b21118bab8e}
\index{Healthcare Interoperability Standards and Specifications (II.3.2)}
{\small Resource · Canonical · current}\par

\par\medskip\textbf{Strategic Definition}\par


\textbf{Healthcare Interoperability Standards and Specifications} comprise the external, open, normative specifications governing healthcare data structures, clinical resource models, exchange semantics, and conformance criteria across enterprise and regional boundaries.



\par\medskip\textbf{Enduring Strategic Role}\par


Harmonia is architected to facilitate vendor-neutral health data exchange across acute, primary, community, and diagnostic healthcare systems. This interoperability cannot be accomplished using proprietary vendor protocols. Harmonia depends directly on consensus-based, open standards to establish:
- Standardised structural and semantic representations for clinical entities (encounters, observations, conditions, medications, diagnostic reports).
- Conformance validation frameworks to assert that emitted and ingested payloads comply with normative specifications.
- Standards-conformant external interaction contracts across the Harmonia boundary.



\par\medskip\textbf{Current Strategic Instance}\par


\begin{itemize}


\item \textbf{HL7® FHIR® Release 5 Core \& Australian Profile Specifications}: The principal contemporary instance of this resource is HL7 FHIR Release 5, alongside Australian jurisdictional profile specifications (AU Base and AU Core).

\item \textbf{Technology Independence}: The strategic resource is the \textbf{specification and standard itself}, not any specific implementation framework or runtime library (such as HAPI FHIR). The strategic dependency endures regardless of which runtime software serialises or deserialises the specifications.

\end{itemize}


\par\medskip\textbf{Traceability}\par


\begin{itemize}


\item \textbf{Architectural Axioms}: \texttt{\seqsplit{AX-02}} (Boundary Membrane), \texttt{\seqsplit{AX-03}} (Native Standards), \texttt{\seqsplit{AX-13}} (Explicit Management Boundary).

\item \textbf{Enterprise Capabilities}: \texttt{\seqsplit{EC-08}} (Interoperability \& Exchange), \texttt{\seqsplit{EC-13}} (Semantic Governance \& Conformance).

\item \textbf{Courses of Action}: \texttt{\seqsplit{COA-01}} (Boundary Membrane Sovereignty).

\end{itemize}

{\small Recorded unresolved matters: canonicalName; exact scope remains in the canonical record.}\par
{\footnotesize \href{../../../../continuum/strategy/resource/4d725c804fb34dfa90486b21118bab8e.md}{Canonical record: 4d725c804fb34dfa90486b21118bab8e}; full qualifications and provenance.}\par
\section*{II.3.3 National Clinical Terminology Assets (SNOMED CT-AU / AMT)}
\phantomsection
\label{navigation:II.3.3}
\addcontentsline{toc}{section}{II.3.3 National Clinical Terminology Assets (SNOMED CT-AU / AMT)}

\label{element:430f641e40aa4bd4a8e85a5cbc1f93be}
\index{National Clinical Terminology Assets (SNOMED CT-AU / AMT) (II.3.3)}
{\small Resource · Canonical · current}\par

\par\medskip\textbf{Strategic Definition}\par


\textbf{National Clinical Terminology Assets} comprise the authoritative, governed clinical vocabularies, concept identifiers, descriptions, and semantic relationships used to capture clinical meaning unambiguously across healthcare systems.



\par\medskip\textbf{Enduring Strategic Role}\par


Harmonia requires authoritative clinical vocabularies to govern semantic meaning during data ingestion, correlation, and egress:
- \textbf{SNOMED CT-AU}: The Australian extension to SNOMED CT, providing the clinical concept reference hierarchy for clinical findings, diagnoses, procedures, and body structures.
- \textbf{Australian Medicines Terminology (AMT)}: The national standard terminology for medicinal products, supporting unambiguous medication reconciliation, prescribing concepts, and dispensing data.



These terminology assets enable Harmonia to evaluate clinical semantic equivalence, enforce conformance rules, and execute meaning-preserving transformations without corrupting clinical intent.



\emph{(Note on LOINC: While Logical Observation Identifiers Names and Codes (LOINC) is a recognised international standard for laboratory and diagnostic observations, its formal admission as an enduring foundational resource is governed by jurisdictional adoption and is not pre-empted ahead of explicit canonical scope).}



\par\medskip\textbf{Traceability}\par


\begin{itemize}


\item \textbf{Architectural Axioms}: \texttt{\seqsplit{AX-04}} (Own Semantics / Use Machinery), \texttt{\seqsplit{AX-08}} (Meaning over Machinery), \texttt{\seqsplit{AX-14}} (Preserve Distinctions).

\item \textbf{Enterprise Capabilities}: \texttt{\seqsplit{EC-13}} (Semantic Governance \& Conformance).

\item \textbf{Courses of Action}: \texttt{\seqsplit{COA-03}} (Meaning-Centric Provenance and Traceability).

\end{itemize}

{\small Recorded unresolved matters: canonicalName; exact scope remains in the canonical record.}\par
{\footnotesize \href{../../../../continuum/strategy/resource/430f641e40aa4bd4a8e85a5cbc1f93be.md}{Canonical record: 430f641e40aa4bd4a8e85a5cbc1f93be}; full qualifications and provenance.}\par

\part*{Part III — Business}
\label{navigation:III}
\addcontentsline{toc}{part}{Part III — Business}
The Business chapters retain the model's enabling-capability groupings, BusinessObjectMap and Products. Functions, services, interactions, roles and processes retain their own native types. Use the canonical descriptions to understand each responsibility; their co-location in a folder does not itself establish a service consumer, process sequence or execution mapping. Explicit relationships and their qualifications remain available in the Continuum Manifest.

{\small Reading sources: \hyperref[element:8882573df7ce43c6bc0f66da0a676139]{AX-17 Architectural Authority and Explicit Uncertainty}.}\par
\chapter*{III.1 Business Information Architecture}
\phantomsection
\label{navigation:III.1}
\addcontentsline{toc}{chapter}{III.1 Business Information Architecture}

\section*{III.1.1 [BusinessObject] Availability Schedule}
\phantomsection
\label{navigation:III.1.1}
\addcontentsline{toc}{section}{III.1.1 [BusinessObject] Availability Schedule}

\label{element:96638431fc39486d879d080768a68fe9}
\index{[BusinessObject] Availability Schedule (III.1.1)}
{\small BusinessObject · Established Package1 conceptual information meaning; residual qualifications in selected section remain unchanged. · current}\par

\texttt{\seqsplit{DeliverableHealthcareService}} is the canonical Information Concept representing the reified Service Provision mapping. It binds:
- \textbf{\texttt{\seqsplit{OfferedHealthcareService}}}: The governing definition of what service is offered;
- \textbf{\texttt{\seqsplit{Healthcare Organisation}}}: The providing enterprise fulfilling the Domain 03 \textbf{Business Role} \texttt{\seqsplit{Service Provider}};
- \textbf{\texttt{\seqsplit{Healthcare Location}}}: The physical care site, clinic, or ward where delivery occurs;
- \textbf{\texttt{\seqsplit{Practitioner Role Binding}}}: The clinical workforce roles qualified and assigned to deliver the service;
- \textbf{\texttt{\seqsplit{Availability Schedule}}}: The operating hours, intake periods, and appointment booking windows;
- \textbf{\texttt{\seqsplit{Eligibility Rule Set}}}: The clinical indications, catchment areas, and patient age/acuity constraints.



Source standing: Established Package1 conceptual information meaning; residual qualifications in selected section remain unchanged.


{\footnotesize \href{../../../../continuum/business/business-object/96638431fc39486d879d080768a68fe9.md}{Canonical record: 96638431fc39486d879d080768a68fe9}; full qualifications and provenance.}\par
\section*{III.1.2 [BusinessObject] Carer Relationship}
\phantomsection
\label{navigation:III.1.2}
\addcontentsline{toc}{section}{III.1.2 [BusinessObject] Carer Relationship}

\label{element:e95d07cf27c041de955665e7b2af8f52}
\index{[BusinessObject] Carer Relationship (III.1.2)}
{\small BusinessObject · Established Package1 conceptual information meaning; residual qualifications in selected section remain unchanged. · current}\par

\begin{longtable}{>{\raggedright\arraybackslash}p{\dimexpr\linewidth/5-2\tabcolsep\relax}>{\raggedright\arraybackslash}p{\dimexpr\linewidth/5-2\tabcolsep\relax}>{\raggedright\arraybackslash}p{\dimexpr\linewidth/5-2\tabcolsep\relax}>{\raggedright\arraybackslash}p{\dimexpr\linewidth/5-2\tabcolsep\relax}>{\raggedright\arraybackslash}p{\dimexpr\linewidth/5-2\tabcolsep\relax}}
Relationship & Source \& Role & Target \& Role & Type / Qualification & Governed Evidence \& Semantics \\
\textbf{Family Relationship} & \texttt{\seqsplit{Person}} (\emph{Family Member}) & \texttt{\seqsplit{Person}} (\emph{Family Member}) & \emph{Familial Association} (e.g. Parent, Child, Sibling, Spouse) & Biological, legal, or social relationship with optional effective period and verification. \\
\textbf{Carer Relationship} & \texttt{\seqsplit{Person}} (\emph{Carer / Support}) & \texttt{\seqsplit{Person}} (\emph{Care Recipient}) & \emph{Carer Association} & Nominated primary/secondary carer supporting the individual's daily living and care needs. \\
\textbf{Legal Representation} & \texttt{\seqsplit{Person}} (\emph{Authorised Representative}) & \texttt{\seqsplit{Person}} (\emph{Represented Subject}) & \emph{Legal Mandate} (e.g. Guardian, Appointed Proxy, Substitute Decision Maker) & Verified statutory or judicial mandate specifying scope of legal authority, validity period, and evidence documentation. \\
\textbf{Support Person} & \texttt{\seqsplit{Person}} (\emph{Nominated Support}) & \texttt{\seqsplit{Person}} (\emph{Supported Subject}) & \emph{Informal Support} & Nominated contact or advocate for communication and emotional support. \\
\end{longtable}


Source standing: Established Package1 conceptual information meaning; residual qualifications in selected section remain unchanged.


{\footnotesize \href{../../../../continuum/business/business-object/e95d07cf27c041de955665e7b2af8f52.md}{Canonical record: e95d07cf27c041de955665e7b2af8f52}; full qualifications and provenance.}\par
\section*{III.1.3 [BusinessObject] Eligibility Rule Set}
\phantomsection
\label{navigation:III.1.3}
\addcontentsline{toc}{section}{III.1.3 [BusinessObject] Eligibility Rule Set}

\label{element:cdb296a83e7e4a17b5e80a9c910c114f}
\index{[BusinessObject] Eligibility Rule Set (III.1.3)}
{\small BusinessObject · Established Package1 conceptual information meaning; residual qualifications in selected section remain unchanged. · current}\par

\texttt{\seqsplit{DeliverableHealthcareService}} is the canonical Information Concept representing the reified Service Provision mapping. It binds:
- \textbf{\texttt{\seqsplit{OfferedHealthcareService}}}: The governing definition of what service is offered;
- \textbf{\texttt{\seqsplit{Healthcare Organisation}}}: The providing enterprise fulfilling the Domain 03 \textbf{Business Role} \texttt{\seqsplit{Service Provider}};
- \textbf{\texttt{\seqsplit{Healthcare Location}}}: The physical care site, clinic, or ward where delivery occurs;
- \textbf{\texttt{\seqsplit{Practitioner Role Binding}}}: The clinical workforce roles qualified and assigned to deliver the service;
- \textbf{\texttt{\seqsplit{Availability Schedule}}}: The operating hours, intake periods, and appointment booking windows;
- \textbf{\texttt{\seqsplit{Eligibility Rule Set}}}: The clinical indications, catchment areas, and patient age/acuity constraints.



Source standing: Established Package1 conceptual information meaning; residual qualifications in selected section remain unchanged.


{\footnotesize \href{../../../../continuum/business/business-object/cdb296a83e7e4a17b5e80a9c910c114f.md}{Canonical record: cdb296a83e7e4a17b5e80a9c910c114f}; full qualifications and provenance.}\par
\section*{III.1.4 [BusinessObject] Family Relationship}
\phantomsection
\label{navigation:III.1.4}
\addcontentsline{toc}{section}{III.1.4 [BusinessObject] Family Relationship}

\label{element:8adbce49ffdd46a19e95cf40f2aa36fd}
\index{[BusinessObject] Family Relationship (III.1.4)}
{\small BusinessObject · Established Package1 conceptual information meaning; residual qualifications in selected section remain unchanged. · current}\par

\begin{longtable}{>{\raggedright\arraybackslash}p{\dimexpr\linewidth/5-2\tabcolsep\relax}>{\raggedright\arraybackslash}p{\dimexpr\linewidth/5-2\tabcolsep\relax}>{\raggedright\arraybackslash}p{\dimexpr\linewidth/5-2\tabcolsep\relax}>{\raggedright\arraybackslash}p{\dimexpr\linewidth/5-2\tabcolsep\relax}>{\raggedright\arraybackslash}p{\dimexpr\linewidth/5-2\tabcolsep\relax}}
Relationship & Source \& Role & Target \& Role & Type / Qualification & Governed Evidence \& Semantics \\
\textbf{Family Relationship} & \texttt{\seqsplit{Person}} (\emph{Family Member}) & \texttt{\seqsplit{Person}} (\emph{Family Member}) & \emph{Familial Association} (e.g. Parent, Child, Sibling, Spouse) & Biological, legal, or social relationship with optional effective period and verification. \\
\textbf{Carer Relationship} & \texttt{\seqsplit{Person}} (\emph{Carer / Support}) & \texttt{\seqsplit{Person}} (\emph{Care Recipient}) & \emph{Carer Association} & Nominated primary/secondary carer supporting the individual's daily living and care needs. \\
\textbf{Legal Representation} & \texttt{\seqsplit{Person}} (\emph{Authorised Representative}) & \texttt{\seqsplit{Person}} (\emph{Represented Subject}) & \emph{Legal Mandate} (e.g. Guardian, Appointed Proxy, Substitute Decision Maker) & Verified statutory or judicial mandate specifying scope of legal authority, validity period, and evidence documentation. \\
\textbf{Support Person} & \texttt{\seqsplit{Person}} (\emph{Nominated Support}) & \texttt{\seqsplit{Person}} (\emph{Supported Subject}) & \emph{Informal Support} & Nominated contact or advocate for communication and emotional support. \\
\end{longtable}


Source standing: Established Package1 conceptual information meaning; residual qualifications in selected section remain unchanged.


{\footnotesize \href{../../../../continuum/business/business-object/8adbce49ffdd46a19e95cf40f2aa36fd.md}{Canonical record: 8adbce49ffdd46a19e95cf40f2aa36fd}; full qualifications and provenance.}\par
\section*{III.1.5 [BusinessObject] Identifier}
\phantomsection
\label{navigation:III.1.5}
\addcontentsline{toc}{section}{III.1.5 [BusinessObject] Identifier}

\label{element:b3141f382ca04bf4bb214c12300ebbfb}
\index{[BusinessObject] Identifier (III.1.5)}
{\small BusinessObject · Established Package1 conceptual information meaning; residual qualifications in selected section remain unchanged. · current}\par

\begin{itemize}


\item \textbf{Semantic Classification}: \texttt{\seqsplit{Assertion}}

\item \textbf{Definition}: A distinct identifier token issued within a specific namespace authority to identify a person.

\item \textbf{Key Semantic Decomposition}:

\item \texttt{\seqsplit{Identifier Value}}: The lexical string or token representing the identity code.

\item \texttt{\seqsplit{Identifier Authority / Namespace}}: The issuing organisation, national scheme, or local domain establishing the uniqueness and validity of the value.

\item \textbf{Architectural Rule}: An identifier is NEVER modelled as an unqualified, opaque string. It is always qualified by its originating authority/namespace.

\end{itemize}


Source standing: Established Package1 conceptual information meaning; residual qualifications in selected section remain unchanged.


{\footnotesize \href{../../../../continuum/business/business-object/b3141f382ca04bf4bb214c12300ebbfb.md}{Canonical record: b3141f382ca04bf4bb214c12300ebbfb}; full qualifications and provenance.}\par
\section*{III.1.6 [BusinessObject] Identity Alias}
\phantomsection
\label{navigation:III.1.6}
\addcontentsline{toc}{section}{III.1.6 [BusinessObject] Identity Alias}

\label{element:b444496792a24c33b439f9cb59a12d98}
\index{[BusinessObject] Identity Alias (III.1.6)}
{\small BusinessObject · Established Package1 conceptual information meaning; residual qualifications in selected section remain unchanged. · current}\par

\begin{itemize}


\item \textbf{Semantic Classification}: \texttt{\seqsplit{Assertion}}

\item \textbf{Definition}: A name, title, or alias by which a person is known in a specific legal, cultural, or operational context.

\item \textbf{Key Characteristics}: Alias type (e.g. \emph{Legal Name}, \emph{Preferred Name}, \emph{Former Name}, \emph{Emergency Alias}), name components, effective period, and verification status.

\end{itemize}


Source standing: Established Package1 conceptual information meaning; residual qualifications in selected section remain unchanged.


{\footnotesize \href{../../../../continuum/business/business-object/b444496792a24c33b439f9cb59a12d98.md}{Canonical record: b444496792a24c33b439f9cb59a12d98}; full qualifications and provenance.}\par
\section*{III.1.7 [BusinessObject] Identity Correction}
\phantomsection
\label{navigation:III.1.7}
\addcontentsline{toc}{section}{III.1.7 [BusinessObject] Identity Correction}

\label{element:7c74034576cf4972bbb6e86e0f0e1e58}
\index{[BusinessObject] Identity Correction (III.1.7)}
{\small BusinessObject · Established Package1 conceptual information meaning; residual qualifications in selected section remain unchanged. · current}\par

\begin{itemize}


\item \textbf{Semantic Classification}: \texttt{\seqsplit{Activity / Governance Record}}

\item \textbf{Definition}: The audited record and managed consequences of source-established identity corrections, including merge, split/unlink outcomes and demographic rectification within established administrative scope. Harmonia does not adjudicate the originating person-merge decision.

\item \textbf{Key Characteristics}: Source-established correction type (\emph{Merge}, \emph{Split}, \emph{Rectification}), target identities, supplied rationale, originating authorising party, and timestamp. Exact correction/evidence semantics beyond the established authority boundary remain unresolved.

\item \textbf{Governance Invariant}: Identity corrections preserve full provenance and audit logs. They do not destructively overwrite historical assertions or break historical provenance chains. Source authorities remain distinct and intact.

\end{itemize}


Source standing: Established Package1 conceptual information meaning; residual qualifications in selected section remain unchanged.


{\footnotesize \href{../../../../continuum/business/business-object/7c74034576cf4972bbb6e86e0f0e1e58.md}{Canonical record: 7c74034576cf4972bbb6e86e0f0e1e58}; full qualifications and provenance.}\par
\section*{III.1.8 [BusinessObject] Identity Correlation Graph}
\phantomsection
\label{navigation:III.1.8}
\addcontentsline{toc}{section}{III.1.8 [BusinessObject] Identity Correlation Graph}

\label{element:8d0a2577d3364dbda96798afe0800200}
\index{[BusinessObject] Identity Correlation Graph (III.1.8)}
{\small BusinessObject · Established Package1 conceptual information meaning; residual qualifications in selected section remain unchanged. · current}\par

\begin{itemize}


\item \textbf{Semantic Classification}: \texttt{\seqsplit{Relationship / Correlation Graph}}

\item \textbf{Definition}: The network of qualified associations linking distinct identifiers originating from different authority namespaces to the same underlying \texttt{\seqsplit{Person}}.

\item \textbf{Key Characteristics}: Source-established linkage type, supplied confidence/verification evidence, originating decision authority, and effective period. These describe attributable received information; they do not assign matching or merge-decision authority to Harmonia. Detailed linkage/evidence semantics remain unresolved.

\item \textbf{Authority Invariant}: Cross-authority correlation links distinct identifier namespaces without merging or transferring their originating authorities, and without implying destructive consolidation into a single authoritative source identity.

\end{itemize}


Source standing: Established Package1 conceptual information meaning; residual qualifications in selected section remain unchanged.


{\footnotesize \href{../../../../continuum/business/business-object/8d0a2577d3364dbda96798afe0800200.md}{Canonical record: 8d0a2577d3364dbda96798afe0800200}; full qualifications and provenance.}\par
\section*{III.1.9 [BusinessObject] Information Assertion}
\phantomsection
\label{navigation:III.1.9}
\addcontentsline{toc}{section}{III.1.9 [BusinessObject] Information Assertion}

\label{element:f3bd4181bb394db8a505705b03972cea}
\index{[BusinessObject] Information Assertion (III.1.9)}
{\small BusinessObject · Canonical conceptual information definition · current}\par

In healthcare integration, clinical and demographic facts are rarely absolute, detached truths. Rather, they are statements made by specific practitioners, diagnostic systems, administrative staff, or patients at specific points in time.



\begin{quote}


\textbf{Information Assertion}: A formal statement about an entity, relationship, activity, event, or circumstance made by an identifiable source within a qualified context.



\end{quote}


\begin{quote}\small\ttfamily
graph TD\\
    AST["Information Assertion"]\\
    AST -->|"asserted by"| AUTH["1. Information Authority<br/>(Clinician, Lab, System, Patient)"]\\
    AST -->|"governed by"| RESP["2. Information Responsibility<br/>(Owning Domain 03 Capability)"]\\
    AST -->|"held / transported by"| CUST["3. Information Custody<br/>(Host Node, Repository, Gateway)"]\\
    AST -->|"accessed by"| CONS["4. Information Consumption<br/>(Dashboard, Workflow, Analytics)"]\\
    AST -->|"contextualised by"| QUAL["5. Qualification \& Context<br/>(Verification State, Confidence, Evidence)"]\\
    AST -->|"evidenced by"| PROV["6. Provenance \& Lineage<br/>(Time, Method, Originating Transaction)"]\\

\end{quote}


\par\medskip\textbf{Pragmatic Assertion Modelling}\par


While the assertion concept provides rigorous semantic governance, \textbf{Domain 04 does not require every single atomic attribute to be modeled as a distinct, reified Assertion object}.
- Where a fact carries its own independent legal liability, clinical attestation, or disputed confidence (e.g. an \emph{Allergy Intolerance Finding}, a \emph{Diagnostic Pathology Finding}, or a \emph{Death Notification}), it is explicitly governed as an Assertion.
- Where attributes form an indivisible, cohesive demographic or descriptive unit asserted under a single common authority (e.g. standard address elements on a profile), they are governed collectively under the parent concept.


{\footnotesize \href{../../../../continuum/business/business-object/f3bd4181bb394db8a505705b03972cea.md}{Canonical record: f3bd4181bb394db8a505705b03972cea}; full qualifications and provenance.}\par
\section*{III.1.10 [BusinessObject] Information Unit}
\phantomsection
\label{navigation:III.1.10}
\addcontentsline{toc}{section}{III.1.10 [BusinessObject] Information Unit}

\label{element:10284034c8534b639284386b064f2c54}
\index{[BusinessObject] Information Unit (III.1.10)}
{\small BusinessObject · Established conceptual governed management boundary; no downstream schema or representation mapping. · current}\par

\begin{quote}


\textbf{An Information Unit is the governed boundary within which Harmonia associates a Governed Representation with its Management Context.}



\end{quote}


An Information Unit is bounded, governed and a unit of Harmonia-managed
information. It is \textbf{not necessarily atomic}: boundedness does not imply
indivisibility. An Information Unit MAY contain structured, composite or related
information according to its defined Information Unit structure.



\textbf{Information Unit} is the canonical term for this concept. \textbf{Information
Container} is retired as its proposed name; legitimate technology or application
uses of container and the existing containment pattern retain their meaning.
No architectural mapping to \texttt{\seqsplit{PetasosParcel}} is established. Any future
implementation mapping must be derived downstream.



\par\medskip\textbf{6.1 Conceptual Partition}\par


\begin{quote}\small\ttfamily
Information Unit\\
│\\
├── Identity\\
│\\
├── Metadata\\
│\\
├── Content\\
│\\
└── Context\\

\end{quote}


\begin{longtable}{>{\raggedright\arraybackslash}p{\dimexpr\linewidth/2-2\tabcolsep\relax}>{\raggedright\arraybackslash}p{\dimexpr\linewidth/2-2\tabcolsep\relax}}
Conceptual partition & Meaning \\
\textbf{Identity} & Identifies the Information Unit. \\
\textbf{Metadata} & Describes how Harmonia recognises and manages the Information Unit. \\
\textbf{Content} & Contains the represented information and its representation characteristics. \\
\textbf{Context} & Contains the management context governing the information over time. \\
\end{longtable}


These are conceptual partitions, \textbf{not a final schema}. They SHALL NOT be
interpreted as requiring four implementation objects, database structures,
Java classes or serialised elements. The partition does not prescribe a physical
separation or a complete allocation of management concerns to implementation
representations.



All dotted notation below is illustrative conceptual notation only. It SHALL
NOT establish final attribute names, datatypes, identifier schemes, cardinalities
or implementation representations, and does not enumerate complete structures.



\par\medskip\textbf{6.2 Identity}\par


\begin{quote}


\textbf{Identity identifies the Information Unit itself.}



\end{quote}


Illustrative notation: \texttt{\seqsplit{InformationUnit.Identity.InstanceId}}.



Information Unit identity, the identity of the represented information, and the
identity of a real-world entity represented by that information are potentially
distinct architectural concerns. None SHALL be inferred to be another merely
because an implementation or external representation uses a common identifier.
The Information Concept's \href{\detokenize{../../../../../docs/markdown/04-information-architecture/metamodel/information-architecture-metamodel.md#31-identity--reference}}{Identity \& Reference characteristics}
remain distinct from this Information Unit identity. The complete Identity
structure is not established here.



\par\medskip\textbf{6.3 Metadata}\par


\begin{quote}


\textbf{Metadata describes the Information Unit and the characteristics required for Harmonia to recognise and manage that Information Unit.}



\end{quote}


Illustrative notation: \texttt{\seqsplit{InformationUnit.Metadata.StructureId}}. This MAY
conceptually identify the defined Information Unit structure/type to which the
Information Unit conforms. It concerns the \textbf{Unit structure}, not the semantic
structure/model of the information represented within it. The complete Metadata
attribute set, naming, cardinality and representation are not established here.



\par\medskip\textbf{6.4 Content}\par


\begin{quote}


\textbf{Content contains the information represented by the Information Unit together with the representation characteristics necessary to interpret that content.}



\end{quote}


Illustrative notation:



\begin{quote}\small\ttfamily
InformationUnit.Content.Data\\
InformationUnit.Content.StructureId\\
InformationUnit.Content.Format\\

\end{quote}


\textbf{Content Structure} conceptually identifies the semantic/structural definition
to which the represented content conforms. \textbf{Content Format} conceptually
identifies the concrete representation or encoding of that content.



Downstream architecture might, for example, distinguish a semantic structure,
profile or model from JSON, XML, a binary representation or a media type. These
examples are non-normative: no supported model, format, mapping or mandatory
attribute is established by them. Concrete encoding/format SHALL NOT be treated
as equivalent to semantic information structure.



\begin{longtable}{>{\raggedright\arraybackslash}p{\dimexpr\linewidth/2-2\tabcolsep\relax}>{\raggedright\arraybackslash}p{\dimexpr\linewidth/2-2\tabcolsep\relax}}
Illustrative notation & Distinct architectural concern \\
\texttt{\seqsplit{InformationUnit.Metadata.StructureId}} & The structure/type of the Information Unit. \\
\texttt{\seqsplit{InformationUnit.Content.StructureId}} & The semantic structure/model of the represented information within Content. \\
\texttt{\seqsplit{InformationUnit.Content.Format}} & The concrete representation or encoding of that Content. \\
\end{longtable}


Unit Structure, Content Structure and Content Format SHALL remain distinct.
Sharing a notation suffix or a downstream representation does not collapse
these concerns.



\par\medskip\textbf{6.5 Context}\par


\begin{quote}


\textbf{Context contains the management context through which Harmonia governs, understands and safely manages the represented information over time.}



\end{quote}


The complementary \href{\detokenize{../../../../../docs/markdown/04-information-architecture/patterns/metabole.md}}{Metabolē pattern} defines bounded
progression of one or more Information Units, with applicable Context,
provenance/derivation and accumulated outputs. Its conceptual
\hyperref[element:c34b4929f71946708d37dbe6c1da920a]{ActivityContext}
classifies the broad activity; it neither replaces this Unit Management Context
nor equates it with Mneme's loaded-value Context or Themis PEP-Context.



The following are currently identified \textbf{conceptual management concerns}. They
are not a final attribute list or schema, nor a mandatory populated field bundle.



\begin{longtable}{>{\raggedright\arraybackslash}p{\dimexpr\linewidth/2-2\tabcolsep\relax}>{\raggedright\arraybackslash}p{\dimexpr\linewidth/2-2\tabcolsep\relax}}
Concern & Conceptual meaning and boundary \\
\textbf{Provenance} & Where/how information arose or was established. \\
\textbf{Authority} & Whose assertion, relationship or determination is authoritative within the applicable context. \\
\textbf{Version} & Identifies a governed progression of an Information Unit's representation. \\
\textbf{Lifecycle} & The applicable information lifecycle and its governance; no universal lifecycle or new state set is defined. \\
\textbf{Temporal Context} & The temporal context needed to understand and manage the represented information; no new temporal model is defined. \\
\textbf{Relationships} & Applicable relationships relevant to management and interpretation; no new relationship taxonomy is defined. \\
\textbf{Active Generation} & Identifies a particular valid manifestation of an Information Unit established for active management or use. \\
\end{longtable}


\textbf{Provenance and Authority SHALL NOT be collapsed.} The source supplying
information is not necessarily the originating authority for every assertion or
relationship represented by it. The existing granular assertion/relationship
authority and provenance distinctions remain applicable; this model does not
complete the wider authority/credibility vocabulary or evidence/retention model.



\textbf{Version and Active Generation SHALL NOT be collapsed.} Under
\hyperref[element:0412a2084f0145babc4f022c71fea65f]{AX-05}, authoritative durable
progression and valid active managed-state generation remain distinct concurrency
domains. Mnemosyne establishes authoritative durable state/version; Mneme manages
active representation and coordination. Failed or degraded convergence SHALL NOT
establish or advance a valid active-state generation; untrusted active
coordination requires reconciliation with authoritative state. Neither domain
may substitute for the other or promote cached/process-local state to durable
authority. This preserves the established boundary without completing the
managed-state model or prescribing concurrency mechanisms.



Management Context does not imply that every management concern is durable
resource content. \hyperref[element:5de2ae2c3b794c2d9f3a14efd51d564f]{AX-07} preserves
security context as operational context without automatic persistence as
information content; \hyperref[element:21d6864657964ba5b03e3f6b5eb29cd5]{AX-08} and
\hyperref[element:9a5cb9e8876844e58e660f909f7df9fc]{AX-09} distinguish meaningful
durable evidence from transient operational state. No retention or security
representation mechanism is inferred here.



\par\medskip\textbf{6.5.1 Version and Active Generation}\par


\begin{quote}


\textbf{Version identifies a governed progression of an Information Unit's representation.}



\textbf{Active Generation identifies a particular valid manifestation of an Information Unit established for active management or use.}



\end{quote}


These definitions are complementary and identify different semantic dimensions:
\textbf{Version concerns governed information progression. Active Generation concerns
active manifestation.}



\textbf{Multiple Active Generations MAY derive from the same Version. A change of
Active Generation does not, by itself, indicate a change of Version.}



\begin{quote}\small\ttfamily
Information Unit\\
      │\\
      ├── Version V7\\
      │      ├── Active Generation G21\\
      │      ├── Active Generation G22\\
      │      └── Active Generation G23\\
      │\\
      └── Version V8\\
             └── Active Generation G24\\

\end{quote}


The labels are illustrative only; they establish no identifier syntax, sequence
mechanism or numbering scheme. The model also permits a Version with no current
Active Generation. It establishes no maximum number of Generations.



Version and Active Generation remain conceptual concerns within Information Unit
Context. Their inclusion in Context does not require simple scalar fields and
defines no canonical Version or Generation attributes. The architecture
establishes semantics, not schema.



\par\medskip\textbf{6.5.2 Observability and Manifestation}\par


Independent processes/functions may act upon manifestations of the same
Information Unit Version established at different observable points. For
example, Process A may act upon G21, Process B upon G22 and Process C upon G23,
with all three manifestations derived from V7:



\begin{quote}\small\ttfamily
                Information Unit\\
                   Version V7\\
                       │\\
          ┌────────────┼────────────┐\\
          │            │            │\\
          ▼            ▼            ▼\\
      Process A    Process B    Process C\\
          │            │            │\\
          ▼            ▼            ▼\\
         G21          G22          G23\\

\end{quote}


The arrows illustrate which manifestations the processes act upon; they do not
define Generation creation, assignment or ownership. The processes can therefore
share governed information progression while acting against different active
manifestations. Version alone cannot express that distinction.



\textbf{Active Generation identifies the manifestation, not the process acting upon
it.} Multiple processes/functions MAY legitimately act upon the same Active
Generation. The example requires neither a separate Generation for each process
nor any process-to-Generation cardinality or affinity.



\par\medskip\textbf{6.5.3 Generation and Temporal Context}\par


\begin{quote}


\textbf{Active Generation is not a timestamp or measure of time.}



\end{quote}


Temporal Context may help explain when a manifestation was established,
applicable or observed, but Temporal Context and Active Generation are distinct
Information Architecture concerns. A single Active Generation may remain valid
across an interval. Different Active Generations need not be distinguished by
wall-clock time. No timestamp, clock, sequence counter or temporal implementation
mechanism is established.



\begin{quote}\small\ttfamily
Version change           = governed information progression\\
Active Generation change = different valid active manifestation\\
                         ≠ necessarily information progression\\

\end{quote}


Active Generation SHALL NOT be treated as another form of Version. A difference
in active manifestation does not, by itself, establish authoritative information
progression.



\par\medskip\textbf{6.5.4 Durable Authority and Unresolved Realisation}\par


\begin{quote}


\textbf{An Active Generation SHALL NOT be treated as an authoritative durable Version merely because Harmonia has manifested or acted upon it.}



\end{quote}


Active manifestation does not manufacture durable authoritative progression.
The \hyperref[element:0412a2084f0145babc4f022c71fea65f]{AX-05 boundary} preserved above
continues to govern: authoritative durable Version and valid Active Generation
remain distinct, and failed or degraded convergence following authoritative state
progression SHALL NOT establish or advance a valid active-state generation.
Active use does not establish durable authoritative progression.



This distinction does not imply that every Active Generation must be persisted,
becomes or produces a Version, is durable, constitutes historical truth or
constitutes evidence. Applicable existing evidence and lifecycle obligations
retain their scope; those relationships are not established by manifestation
alone.



These information semantics may later inform Mneme active managed information,
cache-instance management, concurrent active manifestations, Mnemosyne
authoritative durable representation, convergence, conflict handling and
reconstruction of active state. \textbf{Domain04 SHALL establish the information
semantics before downstream architecture determines their realisation.} This
distinction prescribes no cache, persistence, concurrency or coordination
mechanism and SHALL NOT be retrofitted to existing implementation behaviour.



The remaining managed-state questions are deliberately unresolved here; the
\href{\detokenize{../../../../../docs/markdown/04-information-architecture/metamodel/information-architecture-metamodel.md#656-establishment-and-semantic-indivisibility}}{establishment and failure boundary}
below does not determine their realisation:



\begin{itemize}


\item how manifestation begins or is realised, a Generation's lifetime or retirement,
  and when it ceases to be active;

\item its ownership, visibility, or process, cache, node or cluster affinity;

\item how concurrent Generations reconcile, how active state converges, or what
  follows uncertain effect;

\item persistence or retention, association with historical truth or evidence,
  and relationships to snapshots, lifecycle state, Praxis progression or
  entity state.

\end{itemize}


The distinction does not determine whether retrieval manifests a Generation
or how returned Units or result sets are manifested. The separately adjudicated
\href{\detokenize{../../../../../docs/markdown/04-information-architecture/patterns/search-results.md}}{search semantics} establish a managed metadata
Search Result Set referencing the Information Unit Versions comprising a
managed search result; search alone creates neither copies nor new Versions
of those returned Units. External search acquisition belongs to the governed
Ergo activity that intentionally performs it. These decisions do not change
Version / Active Generation semantics. The separately established
\href{\detokenize{../../../../../docs/markdown/04-information-architecture/information-families/task-work.md}}{bounded Task / Work model} and its
retained questions also remain unchanged, as do Assurance and capacity
information questions.



\par\medskip\textbf{6.5.5 Version / Active Generation Guardrails}\par


These VG-01–VG-08 identifiers are local to this semantic distinction.



\begin{longtable}{>{\raggedright\arraybackslash}p{\dimexpr\linewidth/2-2\tabcolsep\relax}>{\raggedright\arraybackslash}p{\dimexpr\linewidth/2-2\tabcolsep\relax}}
Guardrail & Required semantic boundary \\
\textbf{VG-01 — Different Dimensions} & Version and Active Generation identify different semantic dimensions. \\
\textbf{VG-02 — Version Progression} & Version identifies governed progression of an Information Unit's representation. \\
\textbf{VG-03 — Active Manifestation} & Active Generation identifies a particular valid manifestation established for active management or use. \\
\textbf{VG-04 — Multiple Generations per Version} & Multiple Active Generations MAY derive from the same Version. \\
\textbf{VG-05 — Generation Change != Version Change} & A change in Active Generation does not by itself establish information progression or a new Version. \\
\textbf{VG-06 — Manifestation, Not Process} & Active Generation identifies the manifestation, not the process/function using it. \\
\textbf{VG-07 — Generation != Time} & Active Generation is not a timestamp or measure of time. \\
\textbf{VG-08 — No Durable Authority by Manifestation} & An Active Generation SHALL NOT become an authoritative durable Version merely because Harmonia manifested or acted upon it. \\
\end{longtable}


\par\medskip\textbf{6.5.6 Establishment and Semantic Indivisibility}\par


\begin{quote}


\textbf{Active Generation is a successful-state concept. An unsuccessful attempt to manifest an Information Unit does not create an invalid Active Generation; it fails to establish an Active Generation.}



\textbf{Active Generation establishment is semantically indivisible. Partial, degraded or incomplete manifestation SHALL NOT constitute an Active Generation.}



\end{quote}


The manifestation must be successfully established in full according to the
applicable governed architecture; otherwise no Active Generation is established
by that attempt.



\begin{quote}\small\ttfamily
Manifestation Attempt\\
      │\\
      ├── successfully established in full\\
      │           └── Active Generation established\\
      │\\
      └── otherwise\\
                  └── No Active Generation established by this attempt\\

\end{quote}


There is no partially valid, degraded, incomplete or half-established Active
Generation. Runtime work may pass through partial, intermediate, degraded or
failed conditions while attempting manifestation. Those belong to the
manifestation activity or its mechanics and SHALL NOT acquire the semantic
standing of an Active Generation. These examples define no implementation state
names, attributes or schema. This rule establishes no universal structural,
semantic, temporal, authority, security, provenance, relationship or cache
validation checklist.



\textbf{Semantic indivisibility is not transactional atomicity.} Either an Active
Generation has been successfully established or it has not. This semantic outcome
SHALL NOT imply database or distributed transactions, atomic cache operations,
distributed locks, mutexes, optimistic locking, MVCC, compare-and-swap,
synchronised execution, or single-thread, single-node execution. How downstream
architecture guarantees or implements the outcome remains to be derived.



\par\medskip\textbf{6.5.7 Operational Failure and Information State}\par


\begin{quote}


\textbf{Failure to establish an Active Generation is an operational failure within Harmonia. It describes the unsuccessful manifestation activity and SHALL NOT, by itself, alter the Version or state of the Information Unit from which manifestation was attempted.}



\textbf{Operational knowledge that Harmonia failed to manifest information is not, by itself, a change to that information.}



\end{quote}


The attempt and its outcome belong to Harmonia's operational activity. The
Information Unit and its Version remain governed by their own information
semantics. Manifestation failure does not, merely by occurring, establish that
the source Version is invalid or changed, that a new Version exists, that the
Information Unit lifecycle has changed, that the represented information is
defective or no longer authoritative, or that historical truth has changed.
Any such consequence requires separate established architecture.



Manifestation failure SHALL NOT be modelled as Information Unit state merely
because downstream realisation may need to record, audit, monitor or respond to
that failure. No failure-state attribute or lifecycle transition is defined here.
Active manifestation and its failure do not manufacture authoritative durable
progression; the existing AX-05 boundary remains applicable.



\par\medskip\textbf{6.5.8 Existing Active Generation Preservation}\par


\begin{quote}


\textbf{Failure to establish a subsequent Active Generation SHALL NOT, by itself, invalidate an already established Active Generation.}



\end{quote}


For example, successful manifestation of Version V7 establishes G31. A later
manifestation attempt from V7 fails:



\begin{quote}\small\ttfamily
Information Unit Version V7\\
      │\\
      ├── manifestation succeeds ──> G31 established\\
      │\\
      └── later manifestation fails ──> Operational Failure\\
                                       does not, by itself, invalidate G31\\

\end{quote}


The later failure does not retroactively alter G31's successful establishment.
The labels are illustrative and define no identifier or sequence mechanism.
The rule establishes only that another manifestation attempt's failure is not,
by itself, an invalidating event. It defines no indefinite lifetime, expiry,
lease, freshness, automatic fallback, selection, replacement, invalidation
trigger, precedence or current-generation semantics. AX-05's existing
failed/degraded-convergence and untrusted-coordination requirements remain
applicable; this narrow preservation rule does not settle their realisation.



\textbf{Domain04 deliberately stops at these establishment and failure semantics.}
Manifestation initiation, component responsibility, data loading, Context
assembly, completeness/coherence checks, success commit, failure detection and
operational representation, logging, audit, alerts, retry and recovery remain
downstream architecture and engineering. Generation caching, distribution,
replication, expiry, replacement, invalidation and concurrent interaction are
also undetermined. No broader managed-state reconciliation follows from this
semantic boundary.



\par\medskip\textbf{6.5.9 Active Generation Establishment / Failure Guardrails}\par


These AG-01–AG-07 identifiers are local to establishment and failure semantics
and supplement VG-01–VG-08 without replacing them.



\begin{longtable}{>{\raggedright\arraybackslash}p{\dimexpr\linewidth/2-2\tabcolsep\relax}>{\raggedright\arraybackslash}p{\dimexpr\linewidth/2-2\tabcolsep\relax}}
Guardrail & Required semantic boundary \\
\textbf{AG-01 — Successful-State Concept} & An Active Generation exists only when manifestation has been successfully established. \\
\textbf{AG-02 — Semantic Indivisibility} & Active Generation establishment is semantically indivisible. Partial, degraded or incomplete manifestation SHALL NOT constitute an Active Generation. \\
\textbf{AG-03 — Failure Establishes No Generation} & An unsuccessful manifestation attempt fails to establish an Active Generation; it does not establish an invalid or degraded Active Generation. \\
\textbf{AG-04 — Operational Failure Boundary} & Failure to establish an Active Generation is an operational failure within Harmonia and SHALL NOT, by itself, alter the Version or state of the Information Unit from which manifestation was attempted. \\
\textbf{AG-05 — Information / Operational Separation} & Operational knowledge of manifestation failure SHALL NOT, by itself, be treated as a change to the information being manifested. \\
\textbf{AG-06 — Existing Generation Preservation} & Failure to establish a subsequent Active Generation SHALL NOT, by itself, invalidate an already established Active Generation. \\
\textbf{AG-07 — No Transactional Inference} & Semantic indivisibility SHALL NOT imply any particular transactional, concurrency, persistence or cache mechanism. \\
\end{longtable}


Source standing: Established conceptual governed management boundary; no downstream schema or representation mapping.


{\footnotesize \href{../../../../continuum/business/business-object/10284034c8534b639284386b064f2c54.md}{Canonical record: 10284034c8534b639284386b064f2c54}; full qualifications and provenance.}\par
\section*{III.1.11 [BusinessObject] Legal Representation}
\phantomsection
\label{navigation:III.1.11}
\addcontentsline{toc}{section}{III.1.11 [BusinessObject] Legal Representation}

\label{element:33ddad9f1ce1424bbe0fb9a25efadb9f}
\index{[BusinessObject] Legal Representation (III.1.11)}
{\small BusinessObject · Established Package1 conceptual information meaning; residual qualifications in selected section remain unchanged. · current}\par

\begin{longtable}{>{\raggedright\arraybackslash}p{\dimexpr\linewidth/5-2\tabcolsep\relax}>{\raggedright\arraybackslash}p{\dimexpr\linewidth/5-2\tabcolsep\relax}>{\raggedright\arraybackslash}p{\dimexpr\linewidth/5-2\tabcolsep\relax}>{\raggedright\arraybackslash}p{\dimexpr\linewidth/5-2\tabcolsep\relax}>{\raggedright\arraybackslash}p{\dimexpr\linewidth/5-2\tabcolsep\relax}}
Relationship & Source \& Role & Target \& Role & Type / Qualification & Governed Evidence \& Semantics \\
\textbf{Family Relationship} & \texttt{\seqsplit{Person}} (\emph{Family Member}) & \texttt{\seqsplit{Person}} (\emph{Family Member}) & \emph{Familial Association} (e.g. Parent, Child, Sibling, Spouse) & Biological, legal, or social relationship with optional effective period and verification. \\
\textbf{Carer Relationship} & \texttt{\seqsplit{Person}} (\emph{Carer / Support}) & \texttt{\seqsplit{Person}} (\emph{Care Recipient}) & \emph{Carer Association} & Nominated primary/secondary carer supporting the individual's daily living and care needs. \\
\textbf{Legal Representation} & \texttt{\seqsplit{Person}} (\emph{Authorised Representative}) & \texttt{\seqsplit{Person}} (\emph{Represented Subject}) & \emph{Legal Mandate} (e.g. Guardian, Appointed Proxy, Substitute Decision Maker) & Verified statutory or judicial mandate specifying scope of legal authority, validity period, and evidence documentation. \\
\textbf{Support Person} & \texttt{\seqsplit{Person}} (\emph{Nominated Support}) & \texttt{\seqsplit{Person}} (\emph{Supported Subject}) & \emph{Informal Support} & Nominated contact or advocate for communication and emotional support. \\
\end{longtable}


Source standing: Established Package1 conceptual information meaning; residual qualifications in selected section remain unchanged.


{\footnotesize \href{../../../../continuum/business/business-object/33ddad9f1ce1424bbe0fb9a25efadb9f.md}{Canonical record: 33ddad9f1ce1424bbe0fb9a25efadb9f}; full qualifications and provenance.}\par
\section*{III.1.12 [BusinessObject] Metabolē}
\phantomsection
\label{navigation:III.1.12}
\addcontentsline{toc}{section}{III.1.12 [BusinessObject] Metabolē}

\label{element:c34b4929f71946708d37dbe6c1da920a}
\index{[BusinessObject] Metabolē (III.1.12)}
{\small BusinessObject · Adjudicated conceptual IA extension, 2026-10-13; no application manager. · current}\par

\begin{quote}


\textbf{Metabolē} is a conceptual information construct representing the bounded
progression of one or more Information Units through a change or activity.
A Metabolē incorporates the input information, applicable Context and
provenance/derivation associated with that progression and accumulates the
resulting output information.



\end{quote}


\begin{quote}\small\ttfamily
Metabolē\\
 ├── Input       [IU 1..*]\\
 ├── Context\\
 │    ├── ActivityContext\\
 │    └── other applicable Context\\
 ├── Provenance / Derivation\\
 └── Output      [IU 0..*]\\

\end{quote}


The \hyperref[element:10284034c8534b639284386b064f2c54]{Information Unit}
remains the governed information boundary. Metabolē describes the bounded
progression involving Units; it does not replace their identities, originating
authority, Management Context or lifecycle. Its conceptual partitions do not
prescribe nesting, copying, reference structures, physical representation or a
new universal lifecycle.



Metabolē is conceptual. It is not an application subsystem, persistence service,
execution component or runtime manager. No \texttt{\seqsplit{MetaboleManager}} or equivalent
application responsibility SHALL be inferred.



Source standing: Adjudicated conceptual IA extension, 2026-10-13; no application manager.


{\footnotesize \href{../../../../continuum/business/business-object/c34b4929f71946708d37dbe6c1da920a.md}{Canonical record: c34b4929f71946708d37dbe6c1da920a}; full qualifications and provenance.}\par
\section*{III.1.13 [BusinessObject] Search Result Set}
\phantomsection
\label{navigation:III.1.13}
\addcontentsline{toc}{section}{III.1.13 [BusinessObject] Search Result Set}

\label{element:b683ae87c72a4c16aba948a4b8cb0998}
\index{[BusinessObject] Search Result Set (III.1.13)}
{\small BusinessObject · Adjudicated conceptual managed metadata collection; Search Result Metadata-Set is the same concept, not an additional element. · current}\par

\begin{quote}


\textbf{A search over Harmonia-managed information produces a governed
metadata-based Search Result Set referencing the Information Unit Versions
comprising that result. Search does not, by itself, duplicate or create new
versions of the returned managed information.}



\end{quote}


\textbf{Search Result Set} names this managed metadata collection. It specialises
the established \href{\detokenize{../../../../../docs/markdown/04-information-architecture/metamodel/information-architecture-metamodel.md#21-semantic-classifications}}{Collection meaning}
and \href{\detokenize{../../../../../docs/markdown/04-information-architecture/patterns/containment-and-collections.md#3-collections-vs-containment-distinct-semantics}}{membership distinction}
through search-specific Version-reference membership. It introduces no new
semantic category or Information Concept-to-Information Unit equivalence.
“Search Result Metadata-Set” describes the same concept rather than an
additional concept or representation.



Management attaches to the result set's own governed search/result metadata
and membership within its Information Unit boundary. The referenced
Information Units / Versions retain their existing management boundaries,
authority and history. Creating the result artefact SHALL NOT create copies
of the returned Units, new Versions of those Units, a BLOB of duplicated
managed information, or another authority for that information.



\begin{quote}\small\ttfamily
Harmonia-Managed Information\\
        |\\
        | search\\
        v\\
Search Result Set (managed metadata collection)\\
        |\\
        +-- Search / Result Metadata\\
        |\\
        +-- Result Membership [0..*]\\
                |\\
                +-- Information Unit Version Reference\\

\end{quote}


\par\medskip\textbf{2.1 Information Unit Version Membership}\par


Membership records the applicable \textbf{Information Unit Version}, expressing:



\begin{quote}


\textbf{These were the managed Information Unit Versions comprising the result
of this search.}



\end{quote}


It SHALL NOT resolve later to whichever Active Generation or current Version
happens to exist when the result set is examined. For example, membership
identifying Unit U's Version V7 continues to identify V7 after U progresses to
V8 or has a different active manifestation. These labels establish no
identifier scheme. The \href{\detokenize{../../../../../docs/markdown/04-information-architecture/metamodel/information-architecture-metamodel.md#651-version-and-active-generation}}{closed Version / Active Generation semantics}
remain unchanged.



The set contains membership references, not copies of member content. The
conceptual \texttt{\seqsplit{0..*}} membership permits an empty result. It establishes neither
a reference encoding nor downstream retrieval, retention or manifestation
mechanics, and implies no universal atomic observation of all members.



\par\medskip\textbf{2.2 Search / Result Metadata}\par


The set SHALL contain sufficient metadata to represent the search and its
result as a governed information artefact. Semantic categories may include
search context, search criteria, execution/result time, result-set identity,
result membership and other metadata necessary to establish the governed
meaning of the result.



These categories are not a complete metadata schema, mandatory field list,
datatype or serialization. They prescribe no database, API or cache structure.
Search context does not automatically persist operational security context as
resource content; AX-07 remains applicable.



\par\medskip\textbf{2.3 Search and Separate Self-Contained Artefacts}\par


Version-specific membership preserves the meaning of a search result; it does
not make ordinary search a generic snapshot mechanism. The Search Result Set
is not a self-contained duplicated parcel of the referenced content.



Export, snapshot, document generation, archival package, transmission package
or another self-contained information artefact requires a separate semantic
act/capability. Such acts are not intrinsic to search and are not designed here.
The existing \href{\detokenize{../../../../../docs/markdown/04-information-architecture/assemblies-views/assemblies-and-views.md#3-assembly-lifecycle-vs-constituent-lifecycle}}{Assembly/View and snapshot distinctions}
remain applicable without an automatic Search Result Set-to-Assembly mapping.



Source standing: Adjudicated conceptual managed metadata collection; Search Result Metadata-Set is the same concept, not an additional element.


{\footnotesize \href{../../../../continuum/business/business-object/b683ae87c72a4c16aba948a4b8cb0998.md}{Canonical record: b683ae87c72a4c16aba948a4b8cb0998}; full qualifications and provenance.}\par
\chapter*{III.2 Products and Value Streams}
\phantomsection
\label{navigation:III.2}
\addcontentsline{toc}{chapter}{III.2 Products and Value Streams}

\section*{III.2.1 Products}
\phantomsection
\label{navigation:III.2.1}
\addcontentsline{toc}{section}{III.2.1 Products}

\subsection*{III.2.1.1 [Product] Health Information Exchange Service}
\phantomsection
\label{navigation:III.2.1.1}
\addcontentsline{toc}{subsection}{III.2.1.1 [Product] Health Information Exchange Service}

\label{element:59fb1f7c79eb4e5387fb0652649c2b9c}
\index{[Product] Health Information Exchange Service (III.2.1.1)}
{\small Product · Accepted bounded reconciliation · current}\par

The governed source has no established definition. This absence is retained for Solution Architect review; no definition is inferred from the name.


{\small Recorded unresolved matters: definition; exact scope remains in the canonical record.}\par
{\footnotesize \href{../../../../continuum/business/product/59fb1f7c79eb4e5387fb0652649c2b9c.md}{Canonical record: 59fb1f7c79eb4e5387fb0652649c2b9c}; full qualifications and provenance.}\par
\subsection*{III.2.1.2 [Product] Practitioner Collaboration Service}
\phantomsection
\label{navigation:III.2.1.2}
\addcontentsline{toc}{subsection}{III.2.1.2 [Product] Practitioner Collaboration Service}

\label{element:09a928910d3d458da3aebfb93d114654}
\index{[Product] Practitioner Collaboration Service (III.2.1.2)}
{\small Product · Accepted bounded reconciliation · current}\par

The governed source has no established definition. This absence is retained for Solution Architect review; no definition is inferred from the name.


{\small Recorded unresolved matters: definition; exact scope remains in the canonical record.}\par
{\footnotesize \href{../../../../continuum/business/product/09a928910d3d458da3aebfb93d114654.md}{Canonical record: 09a928910d3d458da3aebfb93d114654}; full qualifications and provenance.}\par
\subsection*{III.2.1.3 [Product] Provider Registry - Central}
\phantomsection
\label{navigation:III.2.1.3}
\addcontentsline{toc}{subsection}{III.2.1.3 [Product] Provider Registry - Central}

\label{element:407f93e81fd740e6b10728941a5a5172}
\index{[Product] Provider Registry - Central (III.2.1.3)}
{\small Product · Accepted bounded reconciliation · current}\par

The governed source has no established definition. This absence is retained for Solution Architect review; no definition is inferred from the name.


{\small Recorded unresolved matters: definition; exact scope remains in the canonical record.}\par
{\footnotesize \href{../../../../continuum/business/product/407f93e81fd740e6b10728941a5a5172.md}{Canonical record: 407f93e81fd740e6b10728941a5a5172}; full qualifications and provenance.}\par
\subsection*{III.2.1.4 [Product] Provider Registry - Satellite}
\phantomsection
\label{navigation:III.2.1.4}
\addcontentsline{toc}{subsection}{III.2.1.4 [Product] Provider Registry - Satellite}

\label{element:1846c512d8f5405aaa4f4e33b5177098}
\index{[Product] Provider Registry - Satellite (III.2.1.4)}
{\small Product · Accepted bounded reconciliation · current}\par

The governed source has no established definition. This absence is retained for Solution Architect review; no definition is inferred from the name.


{\small Recorded unresolved matters: definition; exact scope remains in the canonical record.}\par
{\footnotesize \href{../../../../continuum/business/product/1846c512d8f5405aaa4f4e33b5177098.md}{Canonical record: 1846c512d8f5405aaa4f4e33b5177098}; full qualifications and provenance.}\par
\section*{III.2.2 Value Streams}
\phantomsection
\label{navigation:III.2.2}
\addcontentsline{toc}{section}{III.2.2 Value Streams}

\subsection*{III.2.2.1 Distributed Clinical Knowledge → Informed Clinical Collaboration}
\phantomsection
\label{navigation:III.2.2.1}
\addcontentsline{toc}{subsection}{III.2.2.1 Distributed Clinical Knowledge → Informed Clinical Collaboration}

\label{element:3ba98c0c1170458f8c02fa76694fee67}
\index{Distributed Clinical Knowledge → Informed Clinical Collaboration (III.2.2.1)}
{\small ValueStream · Canonical · current}\par

\par\medskip\textbf{Strategic Purpose \& Value Focus}\par


Healthcare is inherently collaborative. Complex patient care requires continuous multidisciplinary interaction among medical specialists, nursing staff, allied health professionals, pharmacists, and community care providers. However, clinical communication is frequently disconnected from the longitudinal patient record, leading to ungrounded decisions, fragmented care plans, and clinical risk.



The strategic value created by \textbf{VS-04} is the \textbf{synthesis of distributed clinical knowledge into contextually relevant information supporting informed multidisciplinary collaboration}. It enables care teams to communicate, consult, and coordinate in the direct presence of verified longitudinal clinical facts without altering source clinical truth.



\begin{quote}\small\ttfamily
graph LR\\
    S1["S1: Distributed Clinical Knowledge"] --> S2["S2: Relevant Info Discovered"]\\
    S2 --> S3["S3: Access Appropriately Governed"]\\
    S3 --> S4["S4: Clinical Context Established"]\\
    S4 --> S5["S5: Info Presented in Context"]\\
    S5 --> S6["S6: Informed Clinical Collaboration"]\\

\end{quote}


\par\medskip\textbf{Value Transformation Stages}\par


\begin{enumerate}


\item \textbf{S1 — Distributed Clinical Knowledge}:

\item \emph{Stakeholder State}: Clinically relevant information, historical observations, diagnostic findings, care history, and other governed clinical knowledge exist across the healthcare ecosystem.

\item \emph{Value Generated}: Ecosystem knowledge is recognized as an active collective resource rather than isolated data fragments.

\item \textbf{S2 — Relevant Information Discovered}:

\item \emph{Stakeholder State}: Authorized queries, alerts, and subscriptions discover clinical facts pertinent to an active patient problem, clinical event, or multidisciplinary consultation.

\item \emph{Value Generated}: Practitioners are spared manual cross-system searching; relevant clinical facts are brought forward based on active clinical need.

\item \textbf{S3 — Access Appropriately Governed}:

\item \emph{Stakeholder State}: Participant identities, clinical roles, care team relationships, consent policies, and default-deny governance rules are evaluated prior to disclosing clinical information.

\item \emph{Value Generated}: Patient privacy and sensitive clinical data are rigorously safeguarded; collaboration proceeds within legitimate clinical relationships.

\item \textbf{S4 — Clinical Context Established}:

\item \emph{Stakeholder State}: Longitudinal patient timeline context, active problems, current medication regimens, allergies, and relevant diagnostic findings are bound to the collaborative space.

\item \emph{Value Generated}: Clinical discussions are anchored directly in objective patient history rather than fragmented recollection or unverified hearsay.

\item \textbf{S5 — Information Presented in Context}:

\item \emph{Stakeholder State}: Verified clinical summaries, trends, and diagnostic reports are rendered seamlessly alongside collaborative communication channels without exposing internal platform mechanics.

\item \emph{Value Generated}: Care team members share an identical, high-fidelity clinical picture, eliminating misunderstandings and information asymmetry.

\item \textbf{S6 — Informed Clinical Collaboration}:

\item \emph{Stakeholder State}: Treating clinicians and multidisciplinary care team members engage in timely, context-rich deliberation, arriving at informed, coordinated clinical care decisions.

\item \emph{Value Generated}: Care plans are coordinated rapidly, multidisciplinary consultation is streamlined, and patient safety is significantly enhanced.

\end{enumerate}


\par\medskip\textbf{Boundary Guardrail: Clinical Authority \& Expertise in Collaboration}\par


Harmonia establishes the governed information context in which clinical collaboration occurs. Harmonia explicitly maintains the following boundaries:
1. \textbf{Harmonia is not an Electronic Medical Record (EMR)}: It does not replace source departmental or enterprise EMRs.
2. \textbf{Harmonia is not the clinical decision-maker}: Clinical responsibility and decision-making authority rest entirely with the attending clinicians and care teams.
3. \textbf{Harmonia does not originate clinical assertions during collaboration}: Collaborative dialogue (e.g., in Matrix/Agora channels) does not automatically become an authoritative clinical record. Clinical practitioners bring their expertise to the collaboration; Harmonia provides the governed clinical context in which that expertise is exercised.
4. \textbf{Clinical statements require source attribution}: Any clinical order or diagnostic finding resulting from collaboration must be explicitly committed through governed clinical workflows with source practitioner attribution.



\par\medskip\textbf{Representative Strategic Traceability}\par


\begin{longtable}{>{\raggedright\arraybackslash}p{\dimexpr\linewidth/2-2\tabcolsep\relax}>{\raggedright\arraybackslash}p{\dimexpr\linewidth/2-2\tabcolsep\relax}}
Traceability level & Current architectural standing \\
Motivation & Established existing constraint: \hyperref[element:5de2ae2c3b794c2d9f3a14efd51d564f]{AX-07 — Security Is Intrinsic to Managed Operations}. Current named Driver/Goal relationships and a replacement for the retired axiom meanings remain unresolved. \\
Current Business Capability contribution & \textbf{Unresolved / not presently established.} The superseded BC relationships are retired; no relationship to BC-01 through BC-18 is supplied by identifier position or similarity. \\
Current Business Enabling contribution and contextual view & \textbf{Unresolved / not presently established.} Old view ordinals and meanings are retired; no enabling-capability or Feature relationship is inferred. \\
Existing current Enterprise Capability contributions & \textbf{Established, partial:} \href{\detokenize{../../../../../docs/markdown/02-strategy/capabilities/enterprise-capabilities.md#ec-02-context-management}}{EC-02 — Context Management}; \href{\detokenize{../../../../../docs/markdown/02-strategy/capabilities/enterprise-capabilities.md#ec-06-policy--control}}{EC-06 — Policy \& Control}. These previously documented contributions retain their canonical meanings. Other current contributions are unresolved; the retired EC labels do not establish replacements. \\
\end{longtable}


\begin{itemize}


\item \textbf{Strategic Courses of Action}: \texttt{\seqsplit{COA-01}} (Boundary Membrane Sovereignty), \texttt{\seqsplit{COA-03}} (Meaning-Centric Provenance and Traceability), \texttt{\seqsplit{COA-06}} (Collaborative Reusable Capability Realisation).

\end{itemize}

{\small Recorded unresolved matters: canonicalName; exact scope remains in the canonical record.}\par
{\footnotesize \href{../../../../continuum/strategy/value-stream/3ba98c0c1170458f8c02fa76694fee67.md}{Canonical record: 3ba98c0c1170458f8c02fa76694fee67}; full qualifications and provenance.}\par
\subsection*{III.2.2.2 Multi-Source Clinical Information → Coherent Clinical Picture}
\phantomsection
\label{navigation:III.2.2.2}
\addcontentsline{toc}{subsection}{III.2.2.2 Multi-Source Clinical Information → Coherent Clinical Picture}

\label{element:01f655eae3d8485699449c3dcf4f9b56}
\index{Multi-Source Clinical Information → Coherent Clinical Picture (III.2.2.2)}
{\small ValueStream · Canonical · current}\par

\par\medskip\textbf{Strategic Purpose \& Value Focus}\par


Modern healthcare delivery is chronically fragmented across acute facilities, primary care practices, diagnostic laboratories, imaging centres, and community services. Clinical information arrives across disjointed standards (HL7 v2, CDA, FHIR), disparate identifier namespaces, and isolated episodes of care.



The strategic value created by \textbf{VS-01} is \textbf{not} merely the technical transit of electronic messages between systems. The authentic clinical and enterprise value is that \textbf{clinically relevant information from heterogeneous sources is unified into a coherent, longitudinal, patient-centred clinical picture}, enabling clinicians to make timely, well-informed, and safer care decisions.



\begin{quote}\small\ttfamily
graph LR\\
    S1["S1: Clinically Relevant Info Available"] --> S2["S2: Identity \& Context Established"]\\
    S2 --> S3["S3: Meaning Established \& Preserved"]\\
    S3 --> S4["S4: Information Integrated Longitudinally"]\\
    S4 --> S5["S5: Coherent Understanding Available"]\\

\end{quote}


\par\medskip\textbf{Value Transformation Stages}\par


\begin{enumerate}


\item \textbf{S1 — Clinically Relevant Information Available}:

\item \emph{Stakeholder State}: Clinical observations, pathology reports, diagnostic imaging results, medication orders, discharge summaries, or encounter notes occur across distributed care settings and become accessible for integration.

\item \emph{Value Generated}: Clinical events are no longer trapped within isolated source silos; they enter the integration sphere where their clinical significance can be recognized.

\item \textbf{S2 — Identity \& Clinical Context Established}:

\item \emph{Stakeholder State}: Subject-of-care identity is established to the level required for safe association of the information, preserving source identifiers, provenance, confidence, and ambiguity where applicable; relevant encounter, episode, service, and care-delivery context is associated.

\item \emph{Value Generated}: Clinical facts are safely bound to the correct individual and operational care context without masking identity ambiguity or making premature, unsafe identity conflations.

\item \textbf{S3 — Meaning Established \& Preserved}:

\item \emph{Stakeholder State}: Clinical concepts, classifications, and terminologies are mapped to canonical standards (e.g., SNOMED CT-AU, AMT) while strictly preserving original source clinical intent, qualifiers, and provenance.

\item \emph{Value Generated}: Clinical data is semantically unambiguous across organisational boundaries while ensuring that the originating practitioner's original meaning and expression remain unaltered.

\item \textbf{S4 — Information Integrated Longitudinally}:

\item \emph{Stakeholder State}: Disparate episodic facts, temporal observations, diagnostic series, and encounter histories are synthesized into a cumulative, longitudinal patient-centred record.

\item \emph{Value Generated}: Clinical interactions are connected across time and specialty, moving beyond episodic snapshots to reveal patient trajectories, trends, and clinical history.

\item \textbf{S5 — Coherent Patient-Centred Understanding Available}:

\item \emph{Stakeholder State}: The longitudinal clinical picture is made accessible in a contextual, consumable representation, structured for clinical review, clinical decision support, or standards-compliant query.

\item \emph{Value Generated}: Treating clinicians and multidisciplinary care teams gain holistic, trusted situational awareness, reducing diagnostic errors, preventing duplicate testing, and improving patient outcomes.

\end{enumerate}


\par\medskip\textbf{Clinical Authority Boundary Guardrail}\par


Harmonia provides \textbf{governed, vendor-neutral longitudinal clinical representation and preservation}. Harmonia does \textbf{not} claim originating clinical authority for source clinical facts. The originating EMR, laboratory information system, or practitioner retains source clinical authority; Harmonia is authoritative for its governed longitudinal aggregation, contextual relationships, and verifiable provenance trail.



\par\medskip\textbf{Representative Strategic Traceability}\par


\begin{longtable}{>{\raggedright\arraybackslash}p{\dimexpr\linewidth/2-2\tabcolsep\relax}>{\raggedright\arraybackslash}p{\dimexpr\linewidth/2-2\tabcolsep\relax}}
Traceability level & Current architectural standing \\
Motivation & Unresolved: the retired block does not establish a replacement stream-specific axiom set or current named Driver/Goal relationship. \\
Current Business Capability contribution & \textbf{Unresolved / not presently established.} The superseded BC relationships are retired; no relationship to BC-01 through BC-18 is supplied by identifier position or similarity. \\
Current Business Enabling contribution and contextual view & \textbf{Unresolved / not presently established.} Old view ordinals and meanings are retired; no enabling-capability or Feature relationship is inferred. \\
Existing current Enterprise Capability contributions & \textbf{Established, partial:} \href{\detokenize{../../../../../docs/markdown/02-strategy/capabilities/enterprise-capabilities.md#ec-02-context-management}}{EC-02 — Context Management}; \href{\detokenize{../../../../../docs/markdown/02-strategy/capabilities/enterprise-capabilities.md#ec-07-provenance--traceability}}{EC-07 — Provenance \& Traceability}. These previously documented contributions retain their canonical meanings. Other current contributions are unresolved; the retired EC labels do not establish replacements. \\
\end{longtable}


\begin{itemize}


\item \textbf{Strategic Courses of Action}: \texttt{\seqsplit{COA-01}} (Boundary Membrane Sovereignty), \texttt{\seqsplit{COA-02}} (Distinct Management and Durable Preservation), \texttt{\seqsplit{COA-03}} (Meaning-Centric Provenance and Traceability).

\end{itemize}

{\small Recorded unresolved matters: canonicalName; exact scope remains in the canonical record.}\par
{\footnotesize \href{../../../../continuum/strategy/value-stream/01f655eae3d8485699449c3dcf4f9b56.md}{Canonical record: 01f655eae3d8485699449c3dcf4f9b56}; full qualifications and provenance.}\par
\subsection*{III.2.2.3 Operational Need → Coordinated Activity → Resolved Outcome}
\phantomsection
\label{navigation:III.2.2.3}
\addcontentsline{toc}{subsection}{III.2.2.3 Operational Need → Coordinated Activity → Resolved Outcome}

\label{element:beea615ca7b34bbcbd0ad3a2a6ab2cff}
\index{Operational Need → Coordinated Activity → Resolved Outcome (III.2.2.3)}
{\small ValueStream · Canonical · current}\par

\par\medskip\textbf{Strategic Purpose \& Value Focus}\par


Healthcare delivery involves dynamic, distributed operational coordination—from processing an urgent laboratory order, to coordinating an inpatient bed transfer, managing clinical task escalation, or fulfilling an outpatient referral. When operational activity is disconnected from the current state of affected patients, beds, and practitioners, clinical delays and administrative failures occur.



The strategic value created by \textbf{VS-03} is the \textbf{governed, coordinated progression of operational activity together with the corresponding state of affected real-world healthcare entities} (\texttt{\seqsplit{AX-16}}). It bridges operational demand and clinical execution without premature coupling to software execution threads or queueing systems.



\begin{quote}\small\ttfamily
graph LR\\
    S1["S1: Operational Need Identified"] --> S2["S2: Required Activity Established"]\\
    S2 --> S3["S3: Activity Coordinated"]\\
    S3 --> S4["S4: Work Progressed"]\\
    S4 --> S5["S5: Outcome Established"]\\
    S5 --> S6["S6: Affected Entity State Progressed"]\\

\end{quote}


\par\medskip\textbf{Value Transformation Stages}\par


\begin{enumerate}


\item \textbf{S1 — Operational Need Identified}:

\item \emph{Stakeholder State}: A healthcare operational trigger occurs (e.g., patient admission, clinical order placed, diagnostic test requested, bed turnover required, specialist referral initiated).

\item \emph{Value Generated}: Healthcare demands are captured systematically rather than relying on informal or uncoordinated communication.

\item \textbf{S2 — Required Activity Established}:

\item \emph{Stakeholder State}: The activity required to address the operational need is established together with its relevant context, expected outcome, and completion conditions.

\item \emph{Value Generated}: Vague operational requirements are translated into clear, governed, and accountable objectives with explicit success criteria.

\item \textbf{S3 — Activity Coordinated}:

\item \emph{Stakeholder State}: Required activity is coordinated with the appropriate people, services, resources, and affected real-world entities.

\item \emph{Value Generated}: Work is matched to capable, authorized actors and resources in harmony with real-time operational availability and clinical priorities.

\item \textbf{S4 — Work Progressed}:

\item \emph{Stakeholder State}: Execution of the activity advances through observable, auditable lifecycle states with active oversight, milestone tracking, and escalation management.

\item \emph{Value Generated}: Healthcare operations gain end-to-end transparency; bottlenecks, delays, and stalled activities are immediately visible and actionable.

\item \textbf{S5 — Outcome Established}:

\item \emph{Stakeholder State}: Clinical or operational results are achieved, verified, and recorded with full attribution and provenance.

\item \emph{Value Generated}: Activity resolution is verified against the original intent, establishing reliable evidence of care delivery.

\item \textbf{S6 — Affected Entity State Progressed}:

\item \emph{Stakeholder State}: The state of affected real-world entities (e.g., patient clinical journey, bed occupancy status, practitioner assignment, care team schedule) advances to its next governed operational state in alignment with the activity outcome (\texttt{\seqsplit{AX-16}}).

\item \emph{Value Generated}: The physical and organizational reality of the healthcare facility stays synchronized with system records, eliminating state drift and coordination failures.

\end{enumerate}


\par\medskip\textbf{Decoupling Note: Execution Constructs Below the Value Stream}\par


Digital Twin activity is demand-driven and coordinated through the established activity-execution architecture; Digital Twin identity does not imply a dedicated or permanently active execution thread, process or runtime engine. The Twin remains the \href{\detokenize{../../../../../docs/markdown/02-strategy/strategic-views/logical-component-responsibilities.md#component-7-digital-twin-entity-centred-operational-coordination-construct}}{entity-centred coordination construct}.



Activity execution constructs and downstream mechanisms (illustratively, worker pools, messaging queues and FHIR Tasks) remain \textbf{below} the Value Stream. Their names are not stage names or technology choices made by this Strategy view.



\par\medskip\textbf{Representative Strategic Traceability}\par


\begin{longtable}{>{\raggedright\arraybackslash}p{\dimexpr\linewidth/2-2\tabcolsep\relax}>{\raggedright\arraybackslash}p{\dimexpr\linewidth/2-2\tabcolsep\relax}}
Traceability level & Current architectural standing \\
Motivation & Established: \hyperref[element:ef8abbaec37e4890ad78b9f263e67906]{AX-15 — Uncertainty Is Preserved Until Resolved} and \hyperref[element:de7e90f46127455a9ec9ad3b9b48ba54]{AX-16 — Operational Activity \& Entity State Progress Together}. AX-16 is explicitly applied by this stream’s purpose and entity-state progression stage. Current named Driver/Goal relationships and replacement grounding for the other retired axiom meanings remain unresolved. \\
Current Business Capability contribution & \textbf{Unresolved / not presently established.} The superseded BC relationships are retired; no relationship to BC-01 through BC-18 is supplied by identifier position or similarity. \\
Current Business Enabling contribution and contextual view & \textbf{Unresolved / not presently established.} Old view ordinals and meanings are retired; no enabling-capability or Feature relationship is inferred. \\
Existing current Enterprise Capability contributions & \textbf{Established, partial:} \href{\detokenize{../../../../../docs/markdown/02-strategy/capabilities/enterprise-capabilities.md#ec-02-context-management}}{EC-02 — Context Management}; \href{\detokenize{../../../../../docs/markdown/02-strategy/capabilities/enterprise-capabilities.md#ec-12-operational-assurance}}{EC-12 — Operational Assurance}. These previously documented contributions retain their canonical meanings. Other current contributions are unresolved; the retired EC labels do not establish replacements. \\
\end{longtable}


\begin{itemize}


\item \textbf{Strategic Courses of Action}: \texttt{\seqsplit{COA-04}} (Governed Asynchronous Activity Progression), \texttt{\seqsplit{COA-05}} (Entity-Centred Operational Coordination), \texttt{\seqsplit{COA-06}} (Collaborative Reusable Capability Realisation).

\end{itemize}

{\small Recorded unresolved matters: canonicalName; exact scope remains in the canonical record.}\par
{\footnotesize \href{../../../../continuum/strategy/value-stream/beea615ca7b34bbcbd0ad3a2a6ab2cff.md}{Canonical record: beea615ca7b34bbcbd0ad3a2a6ab2cff}; full qualifications and provenance.}\par
\subsection*{III.2.2.4 Unqualified Information → Governed, Consumable Information}
\phantomsection
\label{navigation:III.2.2.4}
\addcontentsline{toc}{subsection}{III.2.2.4 Unqualified Information → Governed, Consumable Information}

\label{element:02e036f18bab400bac261007e39457ee}
\index{Unqualified Information → Governed, Consumable Information (III.2.2.4)}
{\small ValueStream · Canonical · current}\par

\par\medskip\textbf{Strategic Purpose \& Value Focus}\par


Information arriving at the healthcare enterprise boundary is frequently incomplete, syntactically inconsistent, ambiguous in authority, or lacking explicit privacy controls. Consuming or publishing unverified information risks clinical misinterpretation, privacy violations, and regulatory breach.



The strategic value created by \textbf{VS-02} is the progressive \textbf{qualification, governance, policy attachment, and durable integrity assurance of information}, transforming raw external data into trusted, policy-governed assets that can be safely published across enterprise boundaries.



\begin{quote}\small\ttfamily
graph LR\\
    S1["S1: Unqualified Information"] --> S2["S2: Qualified / Structured Info"]\\
    S2 --> S3["S3: Governed / Controlled Info"]\\
    S3 --> S4["S4: Durably Available Info"]\\
    S4 --> S5["S5: Published / Consumable Info"]\\

\end{quote}


\par\medskip\textbf{Value Transformation Stages}\par


\begin{enumerate}


\item \textbf{S1 — Unqualified Information}:

\item \emph{Stakeholder State}: Raw, unvalidated data streams, documents, or transactions arrive at the enterprise membrane with unverified structural conformance, incomplete provenance, and unverified authority.

\item \emph{Value Generated}: Ingress boundaries accept external input safely without allowing malformed or unauthorized data to compromise internal environments.

\item \textbf{S2 — Qualified / Structured Information}:

\item \emph{Stakeholder State}: Data structures are validated against canonical schemas, syntactically parsed, and reconciled into well-defined, verifiable information models.

\item \emph{Value Generated}: Syntactic and structural ambiguity is eliminated; information adheres to predictable, conformant architectural structures.

\item \textbf{S3 — Governed / Controlled Information}:

\item \emph{Stakeholder State}: Information is evaluated against security policies, organizational authority boundaries, consent directives, and confidentiality classifications under default-deny governance (\texttt{\seqsplit{Themis}}).

\item \emph{Value Generated}: Information carries explicit, enforceable access policies and privacy guardrails, ensuring that sensitive healthcare data cannot be disclosed without legitimate authority.

\item \textbf{S4 — Durably Available Information}:

\item \emph{Stakeholder State}: Information is preserved with verifiable integrity, non-destructive versioning, and comprehensive provenance, remaining resiliently queryable and recoverable.

\item \emph{Value Generated}: Stakeholders are assured that information is protected against unauthorised alteration or loss, establishing an enduring, verifiable and auditable history of care information.

\item \textbf{S5 — Published / Consumable Information}:

\item \emph{Stakeholder State}: Governed information is projected into standards-compliant external representations (such as FHIR R5 Core and Australian profiles) across explicit egress boundaries.

\item \emph{Value Generated}: External consumers and presentation layers receive trusted, standards-conformant information tailored to their authorised scope without exposing internal platform mechanics.

\end{enumerate}


\par\medskip\textbf{Enabling Behavioural Pattern (Subordinate to Value Stream)}\par


The progression through VS-02 is supported internally by an established architectural behavioural pattern:



\$\$\textbackslash{}text\{Aggregation\} \textbackslash{}longrightarrow \textbackslash{}text\{Processing\} \textbackslash{}longrightarrow \textbackslash{}text\{Persistence\} \textbackslash{}longrightarrow \textbackslash{}text\{Publishing\}\$\$



\textbf{Crucial Distinction}: This behavioural pattern describes the internal activity contributing to the transformation. It must \textbf{not} be substituted for the value stream stage names. The value stream expresses the \emph{evolving state and value of the information}; the behavioural pattern describes the underlying operational mechanics that support it.



\par\medskip\textbf{Representative Strategic Traceability}\par


\begin{longtable}{>{\raggedright\arraybackslash}p{\dimexpr\linewidth/2-2\tabcolsep\relax}>{\raggedright\arraybackslash}p{\dimexpr\linewidth/2-2\tabcolsep\relax}}
Traceability level & Current architectural standing \\
Motivation & Established existing constraints: \hyperref[element:0412a2084f0145babc4f022c71fea65f]{AX-05 — Active State and Authoritative Durable State Are Distinct} and \hyperref[element:5de2ae2c3b794c2d9f3a14efd51d564f]{AX-07 — Security Is Intrinsic to Managed Operations}. Current named Driver/Goal relationships and a replacement for the retired axiom meanings remain unresolved. \\
Current Business Capability contribution & \textbf{Unresolved / not presently established.} The superseded BC relationships are retired; no relationship to BC-01 through BC-18 is supplied by identifier position or similarity. \\
Current Business Enabling contribution and contextual view & \textbf{Unresolved / not presently established.} Old view ordinals and meanings are retired; no enabling-capability or Feature relationship is inferred. \\
Existing current Enterprise Capability contributions & \textbf{Established, partial:} \href{\detokenize{../../../../../docs/markdown/02-strategy/capabilities/enterprise-capabilities.md#ec-02-context-management}}{EC-02 — Context Management}; \href{\detokenize{../../../../../docs/markdown/02-strategy/capabilities/enterprise-capabilities.md#ec-06-policy--control}}{EC-06 — Policy \& Control}; \href{\detokenize{../../../../../docs/markdown/02-strategy/capabilities/enterprise-capabilities.md#ec-07-provenance--traceability}}{EC-07 — Provenance \& Traceability}. These previously documented contributions retain their canonical meanings. Other current contributions are unresolved; the retired EC labels do not establish replacements. \\
\end{longtable}


\begin{itemize}


\item \textbf{Strategic Courses of Action}: \texttt{\seqsplit{COA-01}} (Boundary Membrane Sovereignty), \texttt{\seqsplit{COA-02}} (Distinct Management and Durable Preservation), \texttt{\seqsplit{COA-03}} (Meaning-Centric Provenance and Traceability), \texttt{\seqsplit{COA-06}} (Collaborative Reusable Capability Realisation).

\end{itemize}

{\small Recorded unresolved matters: canonicalName; exact scope remains in the canonical record.}\par
{\footnotesize \href{../../../../continuum/strategy/value-stream/02e036f18bab400bac261007e39457ee.md}{Canonical record: 02e036f18bab400bac261007e39457ee}; full qualifications and provenance.}\par
\chapter*{III.3 Organisation Architecture}
\phantomsection
\label{navigation:III.3}
\addcontentsline{toc}{chapter}{III.3 Organisation Architecture}

\section*{III.3.1 Roles}
\phantomsection
\label{navigation:III.3.1}
\addcontentsline{toc}{section}{III.3.1 Roles}

\subsection*{III.3.1.1 Advocate}
\phantomsection
\label{navigation:III.3.1.1}
\addcontentsline{toc}{subsection}{III.3.1.1 Advocate}

\label{element:8845af697832409b90c51a65d7b06669}
\index{Advocate (III.3.1.1)}
{\small BusinessRole · Canonical · current}\par

\begin{itemize}


\item \textbf{Advocate}: An individual or agency acting on behalf of the patient to promote their rights, preferences, and interests within the healthcare system.

\end{itemize}

{\small Recorded unresolved matters: canonicalName; exact scope remains in the canonical record.}\par
{\footnotesize \href{../../../../continuum/business/business-role/8845af697832409b90c51a65d7b06669.md}{Canonical record: 8845af697832409b90c51a65d7b06669}; full qualifications and provenance.}\par
\subsection*{III.3.1.2 Bed Manager}
\phantomsection
\label{navigation:III.3.1.2}
\addcontentsline{toc}{subsection}{III.3.1.2 Bed Manager}

\label{element:4f44e0f4c2ae4774a652c44874449487}
\index{Bed Manager (III.3.1.2)}
{\small BusinessRole · Canonical · current}\par

\begin{itemize}


\item \textbf{Bed Manager}: An operational agent who tracks physical bed availability, assigns care-places, and coordinates bed turnover logistics.

\end{itemize}

{\small Recorded unresolved matters: canonicalName; exact scope remains in the canonical record.}\par
{\footnotesize \href{../../../../continuum/business/business-role/4f44e0f4c2ae4774a652c44874449487.md}{Canonical record: 4f44e0f4c2ae4774a652c44874449487}; full qualifications and provenance.}\par
\subsection*{III.3.1.3 Care Coordinator}
\phantomsection
\label{navigation:III.3.1.3}
\addcontentsline{toc}{subsection}{III.3.1.3 Care Coordinator}

\label{element:84cf0b0868b94a1ea5ce5791fcfa66f9}
\index{Care Coordinator (III.3.1.3)}
{\small BusinessRole · Canonical · current}\par

\begin{itemize}


\item \textbf{Care Coordinator}: A practitioner or case manager responsible for planning, synchronising, and monitoring cross-specialty or cross-setting care delivery for a patient.

\end{itemize}

{\small Recorded unresolved matters: canonicalName; exact scope remains in the canonical record.}\par
{\footnotesize \href{../../../../continuum/business/business-role/84cf0b0868b94a1ea5ce5791fcfa66f9.md}{Canonical record: 84cf0b0868b94a1ea5ce5791fcfa66f9}; full qualifications and provenance.}\par
\subsection*{III.3.1.4 Carer}
\phantomsection
\label{navigation:III.3.1.4}
\addcontentsline{toc}{subsection}{III.3.1.4 Carer}

\label{element:3c2784f8f3404e978c618311148e3422}
\index{Carer (III.3.1.4)}
{\small BusinessRole · Canonical · current}\par

\begin{itemize}


\item \textbf{Carer}: An individual who provides unpaid, ongoing personal care, support, or assistance to a patient.

\end{itemize}

{\small Recorded unresolved matters: canonicalName; exact scope remains in the canonical record.}\par
{\footnotesize \href{../../../../continuum/business/business-role/3c2784f8f3404e978c618311148e3422.md}{Canonical record: 3c2784f8f3404e978c618311148e3422}; full qualifications and provenance.}\par
\subsection*{III.3.1.5 Client}
\phantomsection
\label{navigation:III.3.1.5}
\addcontentsline{toc}{subsection}{III.3.1.5 Client}

\label{element:3b90ebcc858e4b0685c7950ceeb8f3fd}
\index{Client (III.3.1.5)}
{\small BusinessRole · Canonical · current}\par

\begin{itemize}


\item \textbf{Client}: A person or party who enters into a service relationship with a healthcare organisation or practitioner.

\end{itemize}

{\small Recorded unresolved matters: canonicalName; exact scope remains in the canonical record.}\par
{\footnotesize \href{../../../../continuum/business/business-role/3b90ebcc858e4b0685c7950ceeb8f3fd.md}{Canonical record: 3b90ebcc858e4b0685c7950ceeb8f3fd}; full qualifications and provenance.}\par
\subsection*{III.3.1.6 Clinician}
\phantomsection
\label{navigation:III.3.1.6}
\addcontentsline{toc}{subsection}{III.3.1.6 Clinician}

\label{element:94fbd35dadc543dd8c5efaca1951e005}
\index{Clinician (III.3.1.6)}
{\small BusinessRole · Canonical · current}\par

\begin{itemize}


\item \textbf{Clinician}: A Practitioner actively engaged in direct clinical assessment, diagnosis, treatment, or medical decision-making.

\end{itemize}

{\small Recorded unresolved matters: canonicalName; exact scope remains in the canonical record.}\par
{\footnotesize \href{../../../../continuum/business/business-role/94fbd35dadc543dd8c5efaca1951e005.md}{Canonical record: 94fbd35dadc543dd8c5efaca1951e005}; full qualifications and provenance.}\par
\subsection*{III.3.1.7 Delegate / Proxy}
\phantomsection
\label{navigation:III.3.1.7}
\addcontentsline{toc}{subsection}{III.3.1.7 Delegate / Proxy}

\label{element:da5d23f9d93a49f8ad11232c12a23e3f}
\index{Delegate / Proxy (III.3.1.7)}
{\small BusinessRole · Canonical · current}\par

\begin{itemize}


\item \textbf{Delegate / Proxy}: An individual nominated by the patient or authorised by law to access health information, receive communications, or perform specific administrative acts.

\end{itemize}

{\small Recorded unresolved matters: canonicalName; exact scope remains in the canonical record.}\par
{\footnotesize \href{../../../../continuum/business/business-role/da5d23f9d93a49f8ad11232c12a23e3f.md}{Canonical record: da5d23f9d93a49f8ad11232c12a23e3f}; full qualifications and provenance.}\par
\subsection*{III.3.1.8 Guardian}
\phantomsection
\label{navigation:III.3.1.8}
\addcontentsline{toc}{subsection}{III.3.1.8 Guardian}

\label{element:81942689c7b74960a23ce0d6da94859c}
\index{Guardian (III.3.1.8)}
{\small BusinessRole · Canonical · current}\par

\begin{itemize}


\item \textbf{Guardian}: A legally appointed guardian exercising general legal custody and personal welfare decision-making.

\end{itemize}

{\small Recorded unresolved matters: canonicalName; exact scope remains in the canonical record.}\par
{\footnotesize \href{../../../../continuum/business/business-role/81942689c7b74960a23ce0d6da94859c.md}{Canonical record: 81942689c7b74960a23ce0d6da94859c}; full qualifications and provenance.}\par
\subsection*{III.3.1.9 Information Client}
\phantomsection
\label{navigation:III.3.1.9}
\addcontentsline{toc}{subsection}{III.3.1.9 Information Client}

\label{element:3ac6ca6e8c9a42b8a3b48e9631073235}
\index{Information Client (III.3.1.9)}
{\small BusinessRole · Canonical · current}\par

\begin{itemize}


\item \textbf{Information Client}: A participant (person, organisation, or system) that queries, retrieves, or consumes healthcare information for clinical review or operational support.

\end{itemize}

{\small Recorded unresolved matters: canonicalName; exact scope remains in the canonical record.}\par
{\footnotesize \href{../../../../continuum/business/business-role/3ac6ca6e8c9a42b8a3b48e9631073235.md}{Canonical record: 3ac6ca6e8c9a42b8a3b48e9631073235}; full qualifications and provenance.}\par
\subsection*{III.3.1.10 Information Custodian}
\phantomsection
\label{navigation:III.3.1.10}
\addcontentsline{toc}{subsection}{III.3.1.10 Information Custodian}

\label{element:603a0fada51746ef9d26dd6fe20b2e78}
\index{Information Custodian (III.3.1.10)}
{\small BusinessRole · Canonical · current}\par

\begin{itemize}


\item \textbf{Information Custodian}: An entity or designated role holding operational custody and legal responsibility for safeguarding, preserving, and managing access to health records.

\end{itemize}

{\small Recorded unresolved matters: canonicalName; exact scope remains in the canonical record.}\par
{\footnotesize \href{../../../../continuum/business/business-role/603a0fada51746ef9d26dd6fe20b2e78.md}{Canonical record: 603a0fada51746ef9d26dd6fe20b2e78}; full qualifications and provenance.}\par
\subsection*{III.3.1.11 Information Steward}
\phantomsection
\label{navigation:III.3.1.11}
\addcontentsline{toc}{subsection}{III.3.1.11 Information Steward}

\label{element:4110a988149e4049844d4b683ff5d8d4}
\index{Information Steward (III.3.1.11)}
{\small BusinessRole · Canonical · current}\par

\begin{itemize}


\item \textbf{Information Steward}: An authority responsible for defining the governance policies, classification rules, retention periods, and privacy standards for health information assets.

\end{itemize}

{\small Recorded unresolved matters: canonicalName; exact scope remains in the canonical record.}\par
{\footnotesize \href{../../../../continuum/business/business-role/4110a988149e4049844d4b683ff5d8d4.md}{Canonical record: 4110a988149e4049844d4b683ff5d8d4}; full qualifications and provenance.}\par
\subsection*{III.3.1.12 Information Supplier}
\phantomsection
\label{navigation:III.3.1.12}
\addcontentsline{toc}{subsection}{III.3.1.12 Information Supplier}

\label{element:545713793b154ed189f4229256933053}
\index{Information Supplier (III.3.1.12)}
{\small BusinessRole · Canonical · current}\par

\begin{itemize}


\item \textbf{Information Supplier}: A participant (person, organisation, or system) that routinely produces, transmits, or submits healthcare information, diagnostic results, or clinical documents into the integration environment.

\end{itemize}

{\small Recorded unresolved matters: canonicalName; exact scope remains in the canonical record.}\par
{\footnotesize \href{../../../../continuum/business/business-role/545713793b154ed189f4229256933053.md}{Canonical record: 545713793b154ed189f4229256933053}; full qualifications and provenance.}\par
\subsection*{III.3.1.13 Patient}
\phantomsection
\label{navigation:III.3.1.13}
\addcontentsline{toc}{subsection}{III.3.1.13 Patient}

\label{element:0d4ad18a7472447ebcd30619b37dce39}
\index{Patient (III.3.1.13)}
{\small BusinessRole · Canonical · current}\par

\begin{itemize}


\item \textbf{Patient}: A Client actively receiving clinical assessment, diagnosis, medical treatment, nursing care, or direct clinical management.

\end{itemize}

{\small Recorded unresolved matters: canonicalName; exact scope remains in the canonical record.}\par
{\footnotesize \href{../../../../continuum/business/business-role/0d4ad18a7472447ebcd30619b37dce39.md}{Canonical record: 0d4ad18a7472447ebcd30619b37dce39}; full qualifications and provenance.}\par
\subsection*{III.3.1.14 Patient Flow Coordinator}
\phantomsection
\label{navigation:III.3.1.14}
\addcontentsline{toc}{subsection}{III.3.1.14 Patient Flow Coordinator}

\label{element:e8f5b7caaee84a67afda4844b1580cee}
\index{Patient Flow Coordinator (III.3.1.14)}
{\small BusinessRole · Canonical · current}\par

\begin{itemize}


\item \textbf{Patient Flow Coordinator}: A facility-wide coordinator managing admissions, emergency department egress, transfers, and bed allocations across hospital directorates.

\end{itemize}

{\small Recorded unresolved matters: canonicalName; exact scope remains in the canonical record.}\par
{\footnotesize \href{../../../../continuum/business/business-role/e8f5b7caaee84a67afda4844b1580cee.md}{Canonical record: e8f5b7caaee84a67afda4844b1580cee}; full qualifications and provenance.}\par
\subsection*{III.3.1.15 Performer}
\phantomsection
\label{navigation:III.3.1.15}
\addcontentsline{toc}{subsection}{III.3.1.15 Performer}

\label{element:c87e99ff36224b9da5a05bbd3002da34}
\index{Performer (III.3.1.15)}
{\small BusinessRole · Canonical · current}\par

\begin{itemize}


\item \textbf{Performer}: A clinician, proceduralist, technician, or care agent who executes a diagnostic test, surgical procedure, clinical observation, or therapeutic intervention.

\end{itemize}

{\small Recorded unresolved matters: canonicalName; exact scope remains in the canonical record.}\par
{\footnotesize \href{../../../../continuum/business/business-role/c87e99ff36224b9da5a05bbd3002da34.md}{Canonical record: c87e99ff36224b9da5a05bbd3002da34}; full qualifications and provenance.}\par
\subsection*{III.3.1.16 Policy Authority}
\phantomsection
\label{navigation:III.3.1.16}
\addcontentsline{toc}{subsection}{III.3.1.16 Policy Authority}

\label{element:e0ff9ce2cf4448568d1caf7a44023f48}
\index{Policy Authority (III.3.1.16)}
{\small BusinessRole · Canonical · current}\par

\begin{itemize}


\item \textbf{Policy Authority}: An entity responsible for defining enterprise or jurisdictional clinical, security, and operational policies.

\end{itemize}

{\small Recorded unresolved matters: canonicalName; exact scope remains in the canonical record.}\par
{\footnotesize \href{../../../../continuum/business/business-role/e0ff9ce2cf4448568d1caf7a44023f48.md}{Canonical record: e0ff9ce2cf4448568d1caf7a44023f48}; full qualifications and provenance.}\par
\subsection*{III.3.1.17 Practitioner}
\phantomsection
\label{navigation:III.3.1.17}
\addcontentsline{toc}{subsection}{III.3.1.17 Practitioner}

\label{element:0b5544e3d9d741b39607a5e61d79049f}
\index{Practitioner (III.3.1.17)}
{\small BusinessRole · Canonical · current}\par

\begin{itemize}


\item \textbf{Practitioner}: An individual registered, certified, or qualified to provide healthcare services.

\end{itemize}

{\small Recorded unresolved matters: canonicalName; exact scope remains in the canonical record.}\par
{\footnotesize \href{../../../../continuum/business/business-role/0b5544e3d9d741b39607a5e61d79049f.md}{Canonical record: 0b5544e3d9d741b39607a5e61d79049f}; full qualifications and provenance.}\par
\subsection*{III.3.1.18 Prescriber}
\phantomsection
\label{navigation:III.3.1.18}
\addcontentsline{toc}{subsection}{III.3.1.18 Prescriber}

\label{element:4a6e12d4f48d426a98136ea50ad91ef9}
\index{Prescriber (III.3.1.18)}
{\small BusinessRole · Canonical · current}\par

\begin{itemize}


\item \textbf{Prescriber}: A practitioner legally authorised and clinically responsible for authorising the dispensing and administration of scheduled medications.

\end{itemize}

{\small Recorded unresolved matters: canonicalName; exact scope remains in the canonical record.}\par
{\footnotesize \href{../../../../continuum/business/business-role/4a6e12d4f48d426a98136ea50ad91ef9.md}{Canonical record: 4a6e12d4f48d426a98136ea50ad91ef9}; full qualifications and provenance.}\par
\subsection*{III.3.1.19 Referrer}
\phantomsection
\label{navigation:III.3.1.19}
\addcontentsline{toc}{subsection}{III.3.1.19 Referrer}

\label{element:197b811fc8594cc3a68a372b74008ba0}
\index{Referrer (III.3.1.19)}
{\small BusinessRole · Canonical · current}\par

\begin{itemize}


\item \textbf{Referrer}: A practitioner or service agent who formally requests another service provider to assess, manage, or assume partial/total clinical responsibility for a patient.

\end{itemize}

{\small Recorded unresolved matters: canonicalName; exact scope remains in the canonical record.}\par
{\footnotesize \href{../../../../continuum/business/business-role/197b811fc8594cc3a68a372b74008ba0.md}{Canonical record: 197b811fc8594cc3a68a372b74008ba0}; full qualifications and provenance.}\par
\subsection*{III.3.1.20 Regulator}
\phantomsection
\label{navigation:III.3.1.20}
\addcontentsline{toc}{subsection}{III.3.1.20 Regulator}

\label{element:1160219638ef42b292a5612e5ac7c539}
\index{Regulator (III.3.1.20)}
{\small BusinessRole · Canonical · current}\par

\begin{itemize}


\item \textbf{Regulator}: A statutory authority responsible for professional registration, healthcare standards enforcement, and accreditation (e.g., AHPRA).

\end{itemize}

{\small Recorded unresolved matters: canonicalName; exact scope remains in the canonical record.}\par
{\footnotesize \href{../../../../continuum/business/business-role/1160219638ef42b292a5612e5ac7c539.md}{Canonical record: 1160219638ef42b292a5612e5ac7c539}; full qualifications and provenance.}\par
\subsection*{III.3.1.21 Representative}
\phantomsection
\label{navigation:III.3.1.21}
\addcontentsline{toc}{subsection}{III.3.1.21 Representative}

\label{element:ca4ae4c178eb488d98118db7773b2127}
\index{Representative (III.3.1.21)}
{\small BusinessRole · Canonical · current}\par

\begin{itemize}


\item \textbf{Representative}: An individual or legal entity empowered to act or make decisions on behalf of the patient.

\end{itemize}

{\small Recorded unresolved matters: canonicalName; exact scope remains in the canonical record.}\par
{\footnotesize \href{../../../../continuum/business/business-role/ca4ae4c178eb488d98118db7773b2127.md}{Canonical record: ca4ae4c178eb488d98118db7773b2127}; full qualifications and provenance.}\par
\subsection*{III.3.1.22 Requester}
\phantomsection
\label{navigation:III.3.1.22}
\addcontentsline{toc}{subsection}{III.3.1.22 Requester}

\label{element:dd262c6bed9a44829e7a4f6796c13b3b}
\index{Requester (III.3.1.22)}
{\small BusinessRole · Canonical · current}\par

\begin{itemize}


\item \textbf{Requester}: A practitioner or clinical agent who formally places an order for a diagnostic investigation, medication supply, or procedural intervention.

\end{itemize}

{\small Recorded unresolved matters: canonicalName; exact scope remains in the canonical record.}\par
{\footnotesize \href{../../../../continuum/business/business-role/dd262c6bed9a44829e7a4f6796c13b3b.md}{Canonical record: dd262c6bed9a44829e7a4f6796c13b3b}; full qualifications and provenance.}\par
\subsection*{III.3.1.23 Service Assurance Modeller}
\phantomsection
\label{navigation:III.3.1.23}
\addcontentsline{toc}{subsection}{III.3.1.23 Service Assurance Modeller}

\label{element:962735ea17a840e5a772f813d3ad7d45}
\index{Service Assurance Modeller (III.3.1.23)}
{\small BusinessRole · Canonical · current}\par

\begin{itemize}


\item \textbf{Canonical Name}: \texttt{\seqsplit{Service Assurance Modeller}}. \textbf{Element Type}: Business Role. Canonical ID remains unresolved; the metamodel does not establish Role identifier tokens or structural context.

\item \textbf{Fundamental Responsibility}: Determine and model how satisfaction of a governed requirement, constraint or expected behaviour is to be assured.

\item \textbf{Responsibility Contexts}: \href{\detokenize{../../../../../docs/markdown/02-strategy/capabilities/business-enabling-capabilities.md#assurance-design}}{Assurance Design} and \href{\detokenize{../../../../../docs/markdown/02-strategy/capabilities/business-enabling-capabilities.md#assurance-criteria-management}}{Assurance Criteria Management}.

\item \textbf{Performed Business Functions}: \href{\detokenize{../../../../../docs/markdown/03-business-architecture/behaviours/health-service-assurance.md#design-service-assurance}}{Design Service Assurance}, owned by Assurance Design; \href{\detokenize{../../../../../docs/markdown/03-business-architecture/behaviours/health-service-assurance.md#manage-assurance-criteria}}{Manage Assurance Criteria}, owned by Assurance Criteria Management. Role performance does not replace capability ownership.

\item \textbf{Principal Process Performance}: \href{\detokenize{../../../../../docs/markdown/03-business-architecture/processes/business-processes.md#assurance-process-design}}{Assurance Process Design}, owned by Assurance Design, with both established Functions participating without transferring Assurance Criteria Management ownership. Assurance Definition remains a conceptual output/boundary; its use by Service Guardian does not establish a Modeller-to-Guardian Interaction.

\item \textbf{Authority Boundary}: Modelling does not inherently confer authority to approve the governing requirement, Assurance Definition or Assurance Criteria. Governance Authority is exercised through existing mandated authority Roles as described in \href{\detokenize{../../../../../docs/markdown/03-business-architecture/actors-roles/roles.md#37-governance-authority-and-assurance-modelling}}{§3.7}.

\item \textbf{Clinical and Operational Boundary}: The Role confers no clinical authority, clinical correctness or adequacy determination, professional clinical judgement, peer-review or clinical-governance authority, or Clinical Services Delivery Assurance. It does not manage, perform, remediate, assign, delegate or operationally escalate the subject activity merely because it models assurance concerning that activity. Explicitly required assurance concerning clinical information or behaviour does not transfer clinical responsibility.

\item \textbf{Allocation Boundary}: The Role models how satisfaction is assured; it does not perform the resulting independent assurance merely by fulfilling this Role. Detailed Actor eligibility, mandates, approval workflows, further Interactions and Collaborations remain unresolved. No application, runtime or implementation allocation is established.

\end{itemize}





{\small Recorded unresolved matters: canonicalName; exact scope remains in the canonical record.}\par
{\footnotesize \href{../../../../continuum/business/business-role/962735ea17a840e5a772f813d3ad7d45.md}{Canonical record: 962735ea17a840e5a772f813d3ad7d45}; full qualifications and provenance.}\par
\subsection*{III.3.1.24 Service Coordinator}
\phantomsection
\label{navigation:III.3.1.24}
\addcontentsline{toc}{subsection}{III.3.1.24 Service Coordinator}

\label{element:0dbbce90dc3e48b190890c832f52cb6f}
\index{Service Coordinator (III.3.1.24)}
{\small BusinessRole · Canonical · current}\par

\begin{itemize}


\item \textbf{Service Coordinator}: An operational administrator who manages service schedules, clinic queues, and service capacity.

\end{itemize}

{\small Recorded unresolved matters: canonicalName; exact scope remains in the canonical record.}\par
{\footnotesize \href{../../../../continuum/business/business-role/0dbbce90dc3e48b190890c832f52cb6f.md}{Canonical record: 0dbbce90dc3e48b190890c832f52cb6f}; full qualifications and provenance.}\par
\subsection*{III.3.1.25 Service Guardian}
\phantomsection
\label{navigation:III.3.1.25}
\addcontentsline{toc}{subsection}{III.3.1.25 Service Guardian}

\label{element:3ed945eeff23404d8dae43382ea746d1}
\index{Service Guardian (III.3.1.25)}
{\small BusinessRole · Canonical · current}\par

\begin{itemize}


\item \textbf{Canonical Name}: \texttt{\seqsplit{Service Guardian}}. \textbf{Element Type}: Business Role. Canonical ID remains unresolved. The previous assurance-role name collision is \href{\detokenize{../../../../../docs/markdown/03-business-architecture/actors-roles/roles.md#36-assurance-role-naming--resolved}}{resolved in §3.6}; the care-support Guardian under Representative is unchanged.

\item \textbf{Definition}: A Business Role responsible for independently evaluating governed activity, information, state or outcomes against applicable assurance criteria and establishing assurance findings and conclusions from sufficient trustworthy evidence.

\item \textbf{Derivation}: \hyperref[element:bbb7438e78604308abba93afc14f4e2e]{REQ-FND-005}, \href{\detokenize{../../../../../docs/markdown/02-strategy/capabilities/business-capabilities.md#18-health-service-assurance}}{BC-18 Health Service Assurance} and \href{\detokenize{../../../../../docs/markdown/02-strategy/capabilities/business-enabling-capabilities.md#governed-assurance}}{Governed Assurance} establish the independent evaluation/conclusion responsibility documented in the \href{\detokenize{../../../../../docs/markdown/03-business-architecture/behaviours/health-service-assurance.md}}{bounded Domain03 assurance derivation}. The shared name is intentional: \textbf{EC-14 Service Guardian is an Enterprise Capability; Service Guardian here is a Business Role.} Neither is a software component.

\item \textbf{Performed Business Functions}: \href{\detokenize{../../../../../docs/markdown/03-business-architecture/behaviours/health-service-assurance.md#establish-assurance-context}}{Establish Assurance Context}, \href{\detokenize{../../../../../docs/markdown/03-business-architecture/behaviours/health-service-assurance.md#assess-assurance-evidence}}{Assess Assurance Evidence} and \href{\detokenize{../../../../../docs/markdown/03-business-architecture/behaviours/health-service-assurance.md#adjudicate-assurance-assessment}}{Adjudicate Assurance Assessment}, all owned by Governed Assurance. Context establishment includes evidence assembly; adjudication includes establishment of findings/conclusions. No separate evidence-assembly, conclusion-establishment or reporting Function is added.

\item \textbf{Approved Process / Interaction Participation}: Principal performing Role for \href{\detokenize{../../../../../docs/markdown/03-business-architecture/processes/business-processes.md#assurance-process-execution}}{Assurance Process Execution}; communicates established outcomes through \href{\detokenize{../../../../../docs/markdown/03-business-architecture/processes/business-processes.md#assurance-process-reporting-communication}}{Assurance Process Reporting / Communication}. Participates in \href{\detokenize{../../../../../docs/markdown/03-business-architecture/collaborations-interactions/interactions.md#5-approved-service-assurance-interactions}}{Service Assurance Request, Status and Outcome Communication} with the approved requesting/status/recipient Roles. This adds no Function or Collaboration and does not make status an assurance conclusion.

\item \textbf{Management Boundary}: Service Guardian may manage progression of its own assurance activity but does not thereby manage, perform, remediate, assign, delegate or operationally escalate the subject being assured. Findings may cause another responsible party to initiate operational response; responsibility for that response remains with its applicable operational capability and authority.

\item \textbf{Clinical Boundary}: Service Guardian confers no clinical authority and does not establish clinical adequacy, clinical correctness or professional clinical judgement, replace clinical peer review or clinical governance, or assume Clinical Services Delivery Assurance. Clinical review and clinical assurance remain with their applicable clinical processes and authorities. Information and behaviour concerning clinical activity may be subject to Guardianship where an explicit applicable governed requirement establishes that concern; this does not transfer clinical responsibility.

\item \textbf{Conclusion Boundary}: Evidence assessment remains distinct from adjudication. Adjudication determines the finding or conclusion from assessed evidence and applicable criteria. Where evidence is insufficient to conclude with the required confidence, insufficiency remains explicit and implies neither satisfaction nor non-satisfaction. This is distinct from execution failure or uncertainty about whether assurance activity occurred.

\item \textbf{Allocation Boundary}: The Role performs independent assurance within Governed Assurance. It does not inherently approve governing requirements, Assurance Definitions or Assurance Criteria, or acquire modelling, clinical review or operational-response responsibility. Detailed Actor eligibility/mandates beyond the approved Role relationships, further Interactions and Collaborations remain unresolved. It has no application, runtime or implementation allocation.

\end{itemize}

{\small Recorded unresolved matters: canonicalName; exact scope remains in the canonical record.}\par
{\footnotesize \href{../../../../continuum/business/business-role/3ed945eeff23404d8dae43382ea746d1.md}{Canonical record: 3ed945eeff23404d8dae43382ea746d1}; full qualifications and provenance.}\par
\subsection*{III.3.1.26 Service Provider}
\phantomsection
\label{navigation:III.3.1.26}
\addcontentsline{toc}{subsection}{III.3.1.26 Service Provider}

\label{element:7e85b6b2c65a4e038e11711a4c7f7c60}
\index{Service Provider (III.3.1.26)}
{\small BusinessRole · Canonical · current}\par

\begin{itemize}


\item \textbf{Service Provider}: An organisation, facility, or practitioner acting in the capacity of delivering a defined healthcare or diagnostic service.

\end{itemize}

{\small Recorded unresolved matters: canonicalName; exact scope remains in the canonical record.}\par
{\footnotesize \href{../../../../continuum/business/business-role/7e85b6b2c65a4e038e11711a4c7f7c60.md}{Canonical record: 7e85b6b2c65a4e038e11711a4c7f7c60}; full qualifications and provenance.}\par
\subsection*{III.3.1.27 Subject of Care}
\phantomsection
\label{navigation:III.3.1.27}
\addcontentsline{toc}{subsection}{III.3.1.27 Subject of Care}

\label{element:ada564bbff80416aa661a1d2cde705f4}
\index{Subject of Care (III.3.1.27)}
{\small BusinessRole · Canonical · current}\par

\begin{itemize}


\item \textbf{Subject of Care}: The identifiable individual whose health, clinical observations, biological samples, or clinical documents form the subject matter of an encounter or health record (often synonymous with Patient, but distinct when specimens, genomic data, or population screening cohorts are processed).

\end{itemize}

{\small Recorded unresolved matters: canonicalName; exact scope remains in the canonical record.}\par
{\footnotesize \href{../../../../continuum/business/business-role/ada564bbff80416aa661a1d2cde705f4.md}{Canonical record: ada564bbff80416aa661a1d2cde705f4}; full qualifications and provenance.}\par
\subsection*{III.3.1.28 Substitute Decision-Maker}
\phantomsection
\label{navigation:III.3.1.28}
\addcontentsline{toc}{subsection}{III.3.1.28 Substitute Decision-Maker}

\label{element:c30a4e1e36b44fd3a953395740d946ef}
\index{Substitute Decision-Maker (III.3.1.28)}
{\small BusinessRole · Canonical · current}\par

\begin{itemize}


\item \textbf{Substitute Decision-Maker}: An individual legally authorised to make specific medical treatment decisions when the patient lacks decision-making capacity.

\end{itemize}

{\small Recorded unresolved matters: canonicalName; exact scope remains in the canonical record.}\par
{\footnotesize \href{../../../../continuum/business/business-role/c30a4e1e36b44fd3a953395740d946ef.md}{Canonical record: c30a4e1e36b44fd3a953395740d946ef}; full qualifications and provenance.}\par
\subsection*{III.3.1.29 Support Person}
\phantomsection
\label{navigation:III.3.1.29}
\addcontentsline{toc}{subsection}{III.3.1.29 Support Person}

\label{element:c2334e0e5c3b40789c672471ea82c465}
\index{Support Person (III.3.1.29)}
{\small BusinessRole · Canonical · current}\par

\begin{itemize}


\item \textbf{Support Person}: An individual nominated by the patient to accompany them, provide moral/emotional support, or assist in communication during care episodes.

\end{itemize}

{\small Recorded unresolved matters: canonicalName; exact scope remains in the canonical record.}\par
{\footnotesize \href{../../../../continuum/business/business-role/c2334e0e5c3b40789c672471ea82c465.md}{Canonical record: c2334e0e5c3b40789c672471ea82c465}; full qualifications and provenance.}\par
\subsection*{III.3.1.30 System Steward}
\phantomsection
\label{navigation:III.3.1.30}
\addcontentsline{toc}{subsection}{III.3.1.30 System Steward}

\label{element:27d2797d2f9e4a9f9131300467d212bc}
\index{System Steward (III.3.1.30)}
{\small BusinessRole · Canonical · current}\par

\begin{itemize}


\item \textbf{System Steward}: An organisational authority responsible for the governance, availability, and lifecycle integrity of digital health systems.

\end{itemize}

{\small Recorded unresolved matters: canonicalName; exact scope remains in the canonical record.}\par
{\footnotesize \href{../../../../continuum/business/business-role/27d2797d2f9e4a9f9131300467d212bc.md}{Canonical record: 27d2797d2f9e4a9f9131300467d212bc}; full qualifications and provenance.}\par
\subsection*{III.3.1.31 Terminology Steward}
\phantomsection
\label{navigation:III.3.1.31}
\addcontentsline{toc}{subsection}{III.3.1.31 Terminology Steward}

\label{element:c3cdd7a91a584942abac167af15f642c}
\index{Terminology Steward (III.3.1.31)}
{\small BusinessRole · Canonical · current}\par

\begin{itemize}


\item \textbf{Terminology Steward}: An authority responsible for authoring, validating, and governing canonical value sets, clinical ontologies, and semantic mapping tables.

\end{itemize}

{\small Recorded unresolved matters: canonicalName; exact scope remains in the canonical record.}\par
{\footnotesize \href{../../../../continuum/business/business-role/c3cdd7a91a584942abac167af15f642c.md}{Canonical record: c3cdd7a91a584942abac167af15f642c}; full qualifications and provenance.}\par
\subsection*{III.3.1.32 Ward Coordinator}
\phantomsection
\label{navigation:III.3.1.32}
\addcontentsline{toc}{subsection}{III.3.1.32 Ward Coordinator}

\label{element:9b4ea5f137c9474db1c9dc5ac57a1207}
\index{Ward Coordinator (III.3.1.32)}
{\small BusinessRole · Canonical · current}\par

\begin{itemize}


\item \textbf{Ward Coordinator}: A senior nurse or ward administrator responsible for shift-level operational oversight, staffing balance, and bed usage within a clinical unit.

\end{itemize}

{\small Recorded unresolved matters: canonicalName; exact scope remains in the canonical record.}\par
{\footnotesize \href{../../../../continuum/business/business-role/9b4ea5f137c9474db1c9dc5ac57a1207.md}{Canonical record: 9b4ea5f137c9474db1c9dc5ac57a1207}; full qualifications and provenance.}\par
\subsection*{III.3.1.33 Wardsperson}
\phantomsection
\label{navigation:III.3.1.33}
\addcontentsline{toc}{subsection}{III.3.1.33 Wardsperson}

\label{element:ebc0fbc6ad7645439a32e1f1a7944209}
\index{Wardsperson (III.3.1.33)}
{\small BusinessRole · Canonical · current}\par

\begin{itemize}


\item \textbf{Wardsperson}: An operational staff member (orderly/porter) responsible for the physical transport of patients, equipment, specimens, and clinical supplies.

\end{itemize}

{\small Recorded unresolved matters: canonicalName; exact scope remains in the canonical record.}\par
{\footnotesize \href{../../../../continuum/business/business-role/ebc0fbc6ad7645439a32e1f1a7944209.md}{Canonical record: ebc0fbc6ad7645439a32e1f1a7944209}; full qualifications and provenance.}\par
\subsection*{III.3.1.34 Work Coordinator}
\phantomsection
\label{navigation:III.3.1.34}
\addcontentsline{toc}{subsection}{III.3.1.34 Work Coordinator}

\label{element:5c52eb31e8ba45a7a4eea4fa8d66d0c5}
\index{Work Coordinator (III.3.1.34)}
{\small BusinessRole · Canonical · current}\par

\begin{itemize}


\item \textbf{Work Coordinator}: An operational manager or dispatch supervisor who creates, assigns, reallocates, and tracks operational work tasks (e.g., portering, cleaning, specimen transport).

\end{itemize}

{\small Recorded unresolved matters: canonicalName; exact scope remains in the canonical record.}\par
{\footnotesize \href{../../../../continuum/business/business-role/5c52eb31e8ba45a7a4eea4fa8d66d0c5.md}{Canonical record: 5c52eb31e8ba45a7a4eea4fa8d66d0c5}; full qualifications and provenance.}\par
\section*{III.3.2 Organisational Structures}
\phantomsection
\label{navigation:III.3.2}
\addcontentsline{toc}{section}{III.3.2 Organisational Structures}

\subsection*{III.3.2.1 Care-Team Collaboration}
\phantomsection
\label{navigation:III.3.2.1}
\addcontentsline{toc}{subsection}{III.3.2.1 Care-Team Collaboration}

\label{element:f748b6d3b83c4591a721f773a3f372be}
\index{Care-Team Collaboration (III.3.2.1)}
{\small BusinessCollaboration · Canonical · current}\par

\begin{itemize}


\item \textbf{Purpose}: Governs the active, multidisciplinary clinical team assembled around a specific patient for an acute episode, chronic disease management programme, or inpatient stay.

\item \textbf{Participating Roles}: \texttt{\seqsplit{Clinician}} (Lead \& Consulting), \texttt{\seqsplit{Care Coordinator}}, \texttt{\seqsplit{Patient}}, \texttt{\seqsplit{Carer}}, \texttt{\seqsplit{Performer}}.

\item \textbf{Governed Activities}: Shared care planning, real-time longitudinal health record (LHR) review, multidisciplinary ward rounds, clinical messaging, and transition-of-care handovers.

\end{itemize}

{\small Recorded unresolved matters: canonicalName; exact scope remains in the canonical record.}\par
{\footnotesize \href{../../../../continuum/business/business-collaboration/f748b6d3b83c4591a721f773a3f372be.md}{Canonical record: f748b6d3b83c4591a721f773a3f372be}; full qualifications and provenance.}\par
\subsection*{III.3.2.2 Information-Sharing Collaboration}
\phantomsection
\label{navigation:III.3.2.2}
\addcontentsline{toc}{subsection}{III.3.2.2 Information-Sharing Collaboration}

\label{element:232a2bc1e9604266b042ebf448e2297c}
\index{Information-Sharing Collaboration (III.3.2.2)}
{\small BusinessCollaboration · Canonical · current}\par

\begin{itemize}


\item \textbf{Purpose}: Governs formal, trusted regional health information exchange, cross-custodian clinical data syndication, and statutory reporting networks.

\item \textbf{Participating Roles}: \texttt{\seqsplit{Information Supplier}}, \texttt{\seqsplit{Information Client}}, \texttt{\seqsplit{Information Steward}}, \texttt{\seqsplit{Information Custodian}}, \texttt{\seqsplit{Policy Authority}}, \texttt{\seqsplit{Regulator}}.

\item \textbf{Governed Activities}: Federated health information exchange, public health surveillance reporting, registry submission, and privacy policy compliance monitoring.

\end{itemize}

{\small Recorded unresolved matters: canonicalName; exact scope remains in the canonical record.}\par
{\footnotesize \href{../../../../continuum/business/business-collaboration/232a2bc1e9604266b042ebf448e2297c.md}{Canonical record: 232a2bc1e9604266b042ebf448e2297c}; full qualifications and provenance.}\par
\subsection*{III.3.2.3 Operational Collaboration}
\phantomsection
\label{navigation:III.3.2.3}
\addcontentsline{toc}{subsection}{III.3.2.3 Operational Collaboration}

\label{element:20f31ce141f947af8700fe7f41d695f6}
\index{Operational Collaboration (III.3.2.3)}
{\small BusinessCollaboration · Canonical · current}\par

\begin{itemize}


\item \textbf{Purpose}: Governs facility-wide operational logistics, resource balancing, bed management, patient flow coordination, and dispatch between clinical units and non-clinical support services.

\item \textbf{Participating Roles}: \texttt{\seqsplit{Patient Flow Coordinator}}, \texttt{\seqsplit{Bed Manager}}, \texttt{\seqsplit{Ward Coordinator}}, \texttt{\seqsplit{Work Coordinator}}, \texttt{\seqsplit{Wardsperson}}.

\item \textbf{Governed Activities}: Daily operational huddles, emergency department bed pull, ward turnover coordination, porter dispatch, and discharge transport scheduling.

\end{itemize}

{\small Recorded unresolved matters: canonicalName; exact scope remains in the canonical record.}\par
{\footnotesize \href{../../../../continuum/business/business-collaboration/20f31ce141f947af8700fe7f41d695f6.md}{Canonical record: 20f31ce141f947af8700fe7f41d695f6}; full qualifications and provenance.}\par
\subsection*{III.3.2.4 Patient Collaboration}
\phantomsection
\label{navigation:III.3.2.4}
\addcontentsline{toc}{subsection}{III.3.2.4 Patient Collaboration}

\label{element:05cb06fbfc3d4b119f067f5699f67efa}
\index{Patient Collaboration (III.3.2.4)}
{\small BusinessCollaboration · Canonical · current}\par

\begin{itemize}


\item \textbf{Purpose}: Governs the structured, longitudinal relationship between a patient (or their authorised carer/representative) and their primary care network, health navigators, and participating health services.

\item \textbf{Participating Roles}: \texttt{\seqsplit{Patient}}, \texttt{\seqsplit{Carer}}, \texttt{\seqsplit{Support Person}}, \texttt{\seqsplit{Advocate}}, \texttt{\seqsplit{Representative}}, \texttt{\seqsplit{Clinician}}, \texttt{\seqsplit{Care Coordinator}}.

\item \textbf{Governed Activities}: Shared goal setting, patient-reported outcome tracking, consent directive establishment, and preference communication.

\end{itemize}

{\small Recorded unresolved matters: canonicalName; exact scope remains in the canonical record.}\par
{\footnotesize \href{../../../../continuum/business/business-collaboration/05cb06fbfc3d4b119f067f5699f67efa.md}{Canonical record: 05cb06fbfc3d4b119f067f5699f67efa}; full qualifications and provenance.}\par
\subsection*{III.3.2.5 Practitioner Collaboration}
\phantomsection
\label{navigation:III.3.2.5}
\addcontentsline{toc}{subsection}{III.3.2.5 Practitioner Collaboration}

\label{element:aeca28004a834fb1acbbf6946ab84a33}
\index{Practitioner Collaboration (III.3.2.5)}
{\small BusinessCollaboration · Canonical · current}\par

\begin{itemize}


\item \textbf{Purpose}: Governs professional peer-to-peer clinical engagement, informal secondary consultations, multidisciplinary review panels, and clinical knowledge sharing across practitioner networks.

\item \textbf{Participating Roles}: \texttt{\seqsplit{Practitioner}}, \texttt{\seqsplit{Clinician}}, \texttt{\seqsplit{Care Coordinator}}, \texttt{\seqsplit{Terminology Steward}}.

\item \textbf{Governed Activities}: Peer review, diagnostic consultation, case conferencing, and shared clinical governance discussions.

\end{itemize}

{\small Recorded unresolved matters: canonicalName; exact scope remains in the canonical record.}\par
{\footnotesize \href{../../../../continuum/business/business-collaboration/aeca28004a834fb1acbbf6946ab84a33.md}{Canonical record: aeca28004a834fb1acbbf6946ab84a33}; full qualifications and provenance.}\par
\subsection*{III.3.2.6 Service Provider Collaboration}
\phantomsection
\label{navigation:III.3.2.6}
\addcontentsline{toc}{subsection}{III.3.2.6 Service Provider Collaboration}

\label{element:f558b3e0948a409a87b26989156063dc}
\index{Service Provider Collaboration (III.3.2.6)}
{\small BusinessCollaboration · Canonical · current}\par

\begin{itemize}


\item \textbf{Purpose}: Governs formal inter-organisational healthcare service partnerships, diagnostic service level agreements, shared service networks, and regional commissioning alliances.

\item \textbf{Participants}: Actors in the \texttt{\seqsplit{Organisation}} and \texttt{\seqsplit{Organisational Unit}} categories; participation through the \texttt{\seqsplit{Service Provider}} and \texttt{\seqsplit{Service Coordinator}} Business Roles. Actor categories and Roles remain distinct; no universal Actor-to-Role assignment is established.

\item \textbf{Governed Activities}: Service directory coordination, regional capacity sharing, diagnostic panel contracting, and cross-organisational service eligibility alignment.

\end{itemize}

{\small Recorded unresolved matters: canonicalName; exact scope remains in the canonical record.}\par
{\footnotesize \href{../../../../continuum/business/business-collaboration/f558b3e0948a409a87b26989156063dc.md}{Canonical record: f558b3e0948a409a87b26989156063dc}; full qualifications and provenance.}\par
\subsection*{III.3.2.7 Service-Delivery Collaboration}
\phantomsection
\label{navigation:III.3.2.7}
\addcontentsline{toc}{subsection}{III.3.2.7 Service-Delivery Collaboration}

\label{element:db39b636a6904a50a9225c98b30bc5a9}
\index{Service-Delivery Collaboration (III.3.2.7)}
{\small BusinessCollaboration · Canonical · current}\par

\begin{itemize}


\item \textbf{Purpose}: Governs the operational and clinical partnership between requesting clinicians and service-fulfilment teams during the active execution of a complex clinical procedure, diagnostic pathway, or surgical episode.

\item \textbf{Participating Roles}: \texttt{\seqsplit{Requester}}, \texttt{\seqsplit{Performer}}, \texttt{\seqsplit{Clinician}}, \texttt{\seqsplit{Service Coordinator}}, \texttt{\seqsplit{Ward Coordinator}}.

\item \textbf{Governed Activities}: Theatre list progression, procedural preparation, perioperative care coordination, and diagnostic order closed-loop fulfilment.

\end{itemize}

{\small Recorded unresolved matters: canonicalName; exact scope remains in the canonical record.}\par
{\footnotesize \href{../../../../continuum/business/business-collaboration/db39b636a6904a50a9225c98b30bc5a9.md}{Canonical record: db39b636a6904a50a9225c98b30bc5a9}; full qualifications and provenance.}\par
\chapter*{III.4 Business Behaviour}
\phantomsection
\label{navigation:III.4}
\addcontentsline{toc}{chapter}{III.4 Business Behaviour}

\section*{III.4.1 [CT1] Care Administration Enablement — Capability context}
\phantomsection
\label{navigation:III.4.1}
\addcontentsline{toc}{section}{III.4.1 [CT1] Care Administration Enablement — Capability context}

{\small Canonical Strategy Capability: II.1.3.1 \hyperref[element:28df27d55c514a37995145ab46945f78]{[CT1] Care Administration Enablement}; Canonical ID: \texttt{\seqsplit{28df27d55c514a37995145ab46945f78}}. Proposed · current.}\par
Contextual reference only; no additional architectural instance or semantic relationship.\par
\subsection*{III.4.1.1 [CT2] Care Plan Administration — Capability context}
\phantomsection
\label{navigation:III.4.1.1}
\addcontentsline{toc}{subsection}{III.4.1.1 [CT2] Care Plan Administration — Capability context}

{\small Canonical Strategy Capability: II.1.3.1.5 \hyperref[element:cf06eec5699f406fa30a00d2f76e3098]{[CT2] Care Plan Administration}; Canonical ID: \texttt{\seqsplit{cf06eec5699f406fa30a00d2f76e3098}}. Canonical · current.}\par
Contextual reference only; no additional architectural instance or semantic relationship.\par
\subsubsection*{III.4.1.1.1 [CT3] Community Care Coordination — Capability context}
\phantomsection
\label{navigation:III.4.1.1.1}
\addcontentsline{toc}{subsubsection}{III.4.1.1.1 [CT3] Community Care Coordination — Capability context}

{\small Canonical Strategy Capability: II.1.3.1.5.2 \hyperref[element:8c8c6266d9ed43e09dbe44fcdd669132]{[CT3] Community Care Coordination}; Canonical ID: \texttt{\seqsplit{8c8c6266d9ed43e09dbe44fcdd669132}}. Canonical · current.}\par
Contextual reference only; no additional architectural instance or semantic relationship.\par
\paragraph*{III.4.1.1.1.1 Business Services}
\phantomsection
\label{navigation:III.4.1.1.1.1}
\addcontentsline{toc}{paragraph}{III.4.1.1.1.1 Business Services}

\subparagraph*{III.4.1.1.1.1.1 Community Care Encounter Ingress}
\phantomsection
\label{navigation:III.4.1.1.1.1.1}
\addcontentsline{toc}{subparagraph}{III.4.1.1.1.1.1 Community Care Encounter Ingress}

\label{element:5806fe40300044b280a25fb86667e789}
\index{Community Care Encounter Ingress (III.4.1.1.1.1.1)}
{\small BusinessService · Canonical · current}\par

Ingests community care documentation.


{\footnotesize \href{../../../../continuum/business/business-service/5806fe40300044b280a25fb86667e789.md}{Canonical record: 5806fe40300044b280a25fb86667e789}; full qualifications and provenance.}\par
\subparagraph*{III.4.1.1.1.1.2 Community Care Team Query}
\phantomsection
\label{navigation:III.4.1.1.1.1.2}
\addcontentsline{toc}{subparagraph}{III.4.1.1.1.1.2 Community Care Team Query}

\label{element:5292120098cc4b7e8280b22f46b60b95}
\index{Community Care Team Query (III.4.1.1.1.1.2)}
{\small BusinessService · Canonical · current}\par

Exposes community provider contact rosters.



Feature association not established: Strategy FEAT-SD-16 — Community Care Plan Synchronization remains valid, but semantic evidence does not establish that the following Function or Service belongs to, realises, specialises or is equivalent to it. No replacement Feature is assigned.


{\small Recorded unresolved matters: source-qualified matters; exact scope remains in the canonical record.}\par
{\footnotesize \href{../../../../continuum/business/business-service/5292120098cc4b7e8280b22f46b60b95.md}{Canonical record: 5292120098cc4b7e8280b22f46b60b95}; full qualifications and provenance.}\par
\paragraph*{III.4.1.1.1.2 Business Functions}
\phantomsection
\label{navigation:III.4.1.1.1.2}
\addcontentsline{toc}{paragraph}{III.4.1.1.1.2 Business Functions}

\subparagraph*{III.4.1.1.1.2.1 Bind Community Support Network Context}
\phantomsection
\label{navigation:III.4.1.1.1.2.1}
\addcontentsline{toc}{subparagraph}{III.4.1.1.1.2.1 Bind Community Support Network Context}

\label{element:321fdf544ef044b8be100b656e16d7dc}
\index{Bind Community Support Network Context (III.4.1.1.1.2.1)}
{\small BusinessFunction · Canonical · current}\par

Links community nursing teams with local non-clinical social care and home care packages.



Feature association not established: Strategy FEAT-SD-16 — Community Care Plan Synchronization remains valid, but semantic evidence does not establish that the following Function or Service belongs to, realises, specialises or is equivalent to it. No replacement Feature is assigned.


{\small Recorded unresolved matters: source-qualified matters; exact scope remains in the canonical record.}\par
{\footnotesize \href{../../../../continuum/business/business-function/321fdf544ef044b8be100b656e16d7dc.md}{Canonical record: 321fdf544ef044b8be100b656e16d7dc}; full qualifications and provenance.}\par
\subparagraph*{III.4.1.1.1.2.2 Synchronize Community Visit Record}
\phantomsection
\label{navigation:III.4.1.1.1.2.2}
\addcontentsline{toc}{subparagraph}{III.4.1.1.1.2.2 Synchronize Community Visit Record}

\label{element:cbd4817eb757449a8f0a0e8ce02c5f07}
\index{Synchronize Community Visit Record (III.4.1.1.1.2.2)}
{\small BusinessFunction · Canonical · current}\par

Ingests clinical notes, wound assessments, and vital signs recorded in community and home settings.


{\footnotesize \href{../../../../continuum/business/business-function/cbd4817eb757449a8f0a0e8ce02c5f07.md}{Canonical record: cbd4817eb757449a8f0a0e8ce02c5f07}; full qualifications and provenance.}\par
\subsubsection*{III.4.1.1.2 [CT3] Outreach Care Continuity — Capability context}
\phantomsection
\label{navigation:III.4.1.1.2}
\addcontentsline{toc}{subsubsection}{III.4.1.1.2 [CT3] Outreach Care Continuity — Capability context}

{\small Canonical Strategy Capability: II.1.3.1.5.3 \hyperref[element:8526d971b07f40de84ef05b0d192089c]{[CT3] Outreach Care Continuity}; Canonical ID: \texttt{\seqsplit{8526d971b07f40de84ef05b0d192089c}}. Canonical · current.}\par
Contextual reference only; no additional architectural instance or semantic relationship.\par
\paragraph*{III.4.1.1.2.1 Business Services}
\phantomsection
\label{navigation:III.4.1.1.2.1}
\addcontentsline{toc}{paragraph}{III.4.1.1.2.1 Business Services}

\subparagraph*{III.4.1.1.2.1.1 Outreach Batch Reconciliation Service}
\phantomsection
\label{navigation:III.4.1.1.2.1.1}
\addcontentsline{toc}{subparagraph}{III.4.1.1.2.1.1 Outreach Batch Reconciliation Service}

\label{element:9488822da1764a1898303a9c656a6062}
\index{Outreach Batch Reconciliation Service (III.4.1.1.2.1.1)}
{\small BusinessService · Canonical · current}\par

Reconciles batch clinical encounters upon network reconnection.


{\footnotesize \href{../../../../continuum/business/business-service/9488822da1764a1898303a9c656a6062.md}{Canonical record: 9488822da1764a1898303a9c656a6062}; full qualifications and provenance.}\par
\paragraph*{III.4.1.1.2.2 Business Functions}
\phantomsection
\label{navigation:III.4.1.1.2.2}
\addcontentsline{toc}{paragraph}{III.4.1.1.2.2 Business Functions}

\subparagraph*{III.4.1.1.2.2.1 Provide Outreach Visit Context}
\phantomsection
\label{navigation:III.4.1.1.2.2.1}
\addcontentsline{toc}{subparagraph}{III.4.1.1.2.2.1 Provide Outreach Visit Context}

\label{element:4984d2f87f5a46869400c60cd5c5e0ab}
\index{Provide Outreach Visit Context (III.4.1.1.2.2.1)}
{\small BusinessFunction · Canonical · current}\par

Packages and delivers portable clinical summaries and offline-capable care summaries for remote outreach clinics.



Derivation: FEAT-SD-17 in Strategy's Outreach Care Continuity (\url{https://github.com/fhirfactory/harmonia/blob/98410db60155198a75f74c79d98060070209faa1/docs/markdown/02-strategy/capabilities/business-enabling-capabilities.md#15-outreach-care-harmonia-relevant}) establishes this preparation/provision responsibility. Outcome reconciliation below does not supply pre-visit context. This Function remains owned by Outreach Care, uses governed source information without acquiring its originating authority and establishes no copying, storage or offline implementation mechanism.


{\footnotesize \href{../../../../continuum/business/business-function/4984d2f87f5a46869400c60cd5c5e0ab.md}{Canonical record: 4984d2f87f5a46869400c60cd5c5e0ab}; full qualifications and provenance.}\par
\subparagraph*{III.4.1.1.2.2.2 Reconcile Outreach Clinical Batch}
\phantomsection
\label{navigation:III.4.1.1.2.2.2}
\addcontentsline{toc}{subparagraph}{III.4.1.1.2.2.2 Reconcile Outreach Clinical Batch}

\label{element:53a7392c0ceb413ba8a00505b29f0f28}
\index{Reconcile Outreach Clinical Batch (III.4.1.1.2.2.2)}
{\small BusinessFunction · Canonical · current}\par

Reconciles and merges clinical documentation recorded during disconnected remote clinic operations.


{\footnotesize \href{../../../../continuum/business/business-function/53a7392c0ceb413ba8a00505b29f0f28.md}{Canonical record: 53a7392c0ceb413ba8a00505b29f0f28}; full qualifications and provenance.}\par
\subsubsection*{III.4.1.1.3 Business Services}
\phantomsection
\label{navigation:III.4.1.1.3}
\addcontentsline{toc}{subsubsection}{III.4.1.1.3 Business Services}

\paragraph*{III.4.1.1.3.1 Care Plan Change Notification}
\phantomsection
\label{navigation:III.4.1.1.3.1}
\addcontentsline{toc}{paragraph}{III.4.1.1.3.1 Care Plan Change Notification}

\label{element:08a8309372734287b2c22a06b95080a4}
\index{Care Plan Change Notification (III.4.1.1.3.1)}
{\small BusinessService · Canonical · current}\par

Emits care plan revision notices.


{\small Recorded unresolved matters: source-qualified matters; exact scope remains in the canonical record.}\par
{\footnotesize \href{../../../../continuum/business/business-service/08a8309372734287b2c22a06b95080a4.md}{Canonical record: 08a8309372734287b2c22a06b95080a4}; full qualifications and provenance.}\par
\paragraph*{III.4.1.1.3.2 Care Plan Context Query}
\phantomsection
\label{navigation:III.4.1.1.3.2}
\addcontentsline{toc}{paragraph}{III.4.1.1.3.2 Care Plan Context Query}

\label{element:0bb4cceb0f1b4311930635d8b4c320e0}
\index{Care Plan Context Query (III.4.1.1.3.2)}
{\small BusinessService · Canonical · current}\par

Discloses active care plans to treating clinicians.


{\small Recorded unresolved matters: source-qualified matters; exact scope remains in the canonical record.}\par
{\footnotesize \href{../../../../continuum/business/business-service/0bb4cceb0f1b4311930635d8b4c320e0.md}{Canonical record: 0bb4cceb0f1b4311930635d8b4c320e0}; full qualifications and provenance.}\par
\paragraph*{III.4.1.1.3.3 Care Plan Ingress}
\phantomsection
\label{navigation:III.4.1.1.3.3}
\addcontentsline{toc}{paragraph}{III.4.1.1.3.3 Care Plan Ingress}

\label{element:180d3f29c2324905ab19337a4dcbce7b}
\index{Care Plan Ingress (III.4.1.1.3.3)}
{\small BusinessService · Canonical · current}\par

Ingests shared care plans.


{\footnotesize \href{../../../../continuum/business/business-service/180d3f29c2324905ab19337a4dcbce7b.md}{Canonical record: 180d3f29c2324905ab19337a4dcbce7b}; full qualifications and provenance.}\par
\subsubsection*{III.4.1.1.4 Business Functions}
\phantomsection
\label{navigation:III.4.1.1.4}
\addcontentsline{toc}{subsubsection}{III.4.1.1.4 Business Functions}

\paragraph*{III.4.1.1.4.1 Associate Care Plan Context}
\phantomsection
\label{navigation:III.4.1.1.4.1}
\addcontentsline{toc}{paragraph}{III.4.1.1.4.1 Associate Care Plan Context}

\label{element:98cd313fdaa74634a8e1ae56a5bb06dc}
\index{Associate Care Plan Context (III.4.1.1.4.1)}
{\small BusinessFunction · Canonical · current}\par

Links care plans to active diagnoses, health problems, and care team members.


{\small Recorded unresolved matters: source-qualified matters; exact scope remains in the canonical record.}\par
{\footnotesize \href{../../../../continuum/business/business-function/98cd313fdaa74634a8e1ae56a5bb06dc.md}{Canonical record: 98cd313fdaa74634a8e1ae56a5bb06dc}; full qualifications and provenance.}\par
\paragraph*{III.4.1.1.4.2 Distribute Care Plan Change}
\phantomsection
\label{navigation:III.4.1.1.4.2}
\addcontentsline{toc}{paragraph}{III.4.1.1.4.2 Distribute Care Plan Change}

\label{element:cbfa51308103418a9dffe23a2d8cfd87}
\index{Distribute Care Plan Change (III.4.1.1.4.2)}
{\small BusinessFunction · Canonical · current}\par

Broadcasts care plan updates to all enrolled care team participants.


{\small Recorded unresolved matters: source-qualified matters; exact scope remains in the canonical record.}\par
{\footnotesize \href{../../../../continuum/business/business-function/cbfa51308103418a9dffe23a2d8cfd87.md}{Canonical record: cbfa51308103418a9dffe23a2d8cfd87}; full qualifications and provenance.}\par
\paragraph*{III.4.1.1.4.3 Receive Care Plan}
\phantomsection
\label{navigation:III.4.1.1.4.3}
\addcontentsline{toc}{paragraph}{III.4.1.1.4.3 Receive Care Plan}

\label{element:abb5a2ac72424923ae2da671b94b3b48}
\index{Receive Care Plan (III.4.1.1.4.3)}
{\small BusinessFunction · Canonical · current}\par

Ingests multidisciplinary care plans, clinical goals, and patient action steps.


{\footnotesize \href{../../../../continuum/business/business-function/abb5a2ac72424923ae2da671b94b3b48.md}{Canonical record: abb5a2ac72424923ae2da671b94b3b48}; full qualifications and provenance.}\par
\subsection*{III.4.1.2 [CT2] Client-Centred Administration — Capability context}
\phantomsection
\label{navigation:III.4.1.2}
\addcontentsline{toc}{subsection}{III.4.1.2 [CT2] Client-Centred Administration — Capability context}

{\small Canonical Strategy Capability: II.1.3.1.6 \hyperref[element:8752b7a937754b2faa296300c6a48911]{[CT2] Client-Centred Administration}; Canonical ID: \texttt{\seqsplit{8752b7a937754b2faa296300c6a48911}}. Canonical · current.}\par
Contextual reference only; no additional architectural instance or semantic relationship.\par
\subsubsection*{III.4.1.2.1 [CT3] Screening and Immunisation Enablement — Capability context}
\phantomsection
\label{navigation:III.4.1.2.1}
\addcontentsline{toc}{subsubsection}{III.4.1.2.1 [CT3] Screening and Immunisation Enablement — Capability context}

{\small Canonical Strategy Capability: II.1.3.1.6.2 \hyperref[element:3e9016eaebae4ed0b67a5cacba2652f5]{[CT3] Screening and Immunisation Enablement}; Canonical ID: \texttt{\seqsplit{3e9016eaebae4ed0b67a5cacba2652f5}}. Canonical · current.}\par
Contextual reference only; no additional architectural instance or semantic relationship.\par
\paragraph*{III.4.1.2.1.1 Business Services}
\phantomsection
\label{navigation:III.4.1.2.1.1}
\addcontentsline{toc}{paragraph}{III.4.1.2.1.1 Business Services}

\subparagraph*{III.4.1.2.1.1.1 Immunisation History Query}
\phantomsection
\label{navigation:III.4.1.2.1.1.1}
\addcontentsline{toc}{subparagraph}{III.4.1.2.1.1.1 Immunisation History Query}

\label{element:f2a0107d73664aa0be9d6695fc455d74}
\index{Immunisation History Query (III.4.1.2.1.1.1)}
{\small BusinessService · Canonical · current}\par

Discloses complete vaccination history.


{\footnotesize \href{../../../../continuum/business/business-service/f2a0107d73664aa0be9d6695fc455d74.md}{Canonical record: f2a0107d73664aa0be9d6695fc455d74}; full qualifications and provenance.}\par
\subparagraph*{III.4.1.2.1.1.2 Preventive Screening Status Query}
\phantomsection
\label{navigation:III.4.1.2.1.1.2}
\addcontentsline{toc}{subparagraph}{III.4.1.2.1.1.2 Preventive Screening Status Query}

\label{element:bc679a6f23de44b79a452fb5ef10c5ee}
\index{Preventive Screening Status Query (III.4.1.2.1.1.2)}
{\small BusinessService · Canonical · current}\par

Discloses screening due dates to primary care providers.



Feature association not established: FEAT-SD-14 — Screening Recall Notification Distribution is adjudicated sufficient through reusable composition/workflow above. Projecting/querying eligibility dates remains distinct from distribution; no dedicated recall Function or replacement association is assigned.


{\small Recorded unresolved matters: source-qualified matters; exact scope remains in the canonical record.}\par
{\footnotesize \href{../../../../continuum/business/business-service/bc679a6f23de44b79a452fb5ef10c5ee.md}{Canonical record: bc679a6f23de44b79a452fb5ef10c5ee}; full qualifications and provenance.}\par
\paragraph*{III.4.1.2.1.2 Business Functions}
\phantomsection
\label{navigation:III.4.1.2.1.2}
\addcontentsline{toc}{paragraph}{III.4.1.2.1.2 Business Functions}

\subparagraph*{III.4.1.2.1.2.1 Project Screening Eligibility Timeline}
\phantomsection
\label{navigation:III.4.1.2.1.2.1}
\addcontentsline{toc}{subparagraph}{III.4.1.2.1.2.1 Project Screening Eligibility Timeline}

\label{element:afb08bdb99cd41d3a86ef7bf7ad5e92d}
\index{Project Screening Eligibility Timeline (III.4.1.2.1.2.1)}
{\small BusinessFunction · Canonical · current}\par

Tracks age-based and risk-based screening schedules across primary care.



Feature association not established: FEAT-SD-14 — Screening Recall Notification Distribution is adjudicated sufficient through reusable composition/workflow above. Projecting/querying eligibility dates remains distinct from distribution; no dedicated recall Function or replacement association is assigned.


{\small Recorded unresolved matters: source-qualified matters; exact scope remains in the canonical record.}\par
{\footnotesize \href{../../../../continuum/business/business-function/afb08bdb99cd41d3a86ef7bf7ad5e92d.md}{Canonical record: afb08bdb99cd41d3a86ef7bf7ad5e92d}; full qualifications and provenance.}\par
\subparagraph*{III.4.1.2.1.2.2 Syndicate Immunisation Record}
\phantomsection
\label{navigation:III.4.1.2.1.2.2}
\addcontentsline{toc}{subparagraph}{III.4.1.2.1.2.2 Syndicate Immunisation Record}

\label{element:5a2c53df2e92444ab9b5018b224d98f9}
\index{Syndicate Immunisation Record (III.4.1.2.1.2.2)}
{\small BusinessFunction · Canonical · current}\par

Records vaccine encounters and syndicates records to national immunisation registers (e.g., AIR).


{\footnotesize \href{../../../../continuum/business/business-function/5a2c53df2e92444ab9b5018b224d98f9.md}{Canonical record: 5a2c53df2e92444ab9b5018b224d98f9}; full qualifications and provenance.}\par
\subsubsection*{III.4.1.2.2 Business Services}
\phantomsection
\label{navigation:III.4.1.2.2}
\addcontentsline{toc}{subsubsection}{III.4.1.2.2 Business Services}

\paragraph*{III.4.1.2.2.1 Demographic Intake Service}
\phantomsection
\label{navigation:III.4.1.2.2.1}
\addcontentsline{toc}{paragraph}{III.4.1.2.2.1 Demographic Intake Service}

\label{element:b6d8aab087944a6e8cc246b85d32843e}
\index{Demographic Intake Service (III.4.1.2.2.1)}
{\small BusinessService · Canonical · current}\par

Exposed to departmental registration systems.


{\footnotesize \href{../../../../continuum/business/business-service/b6d8aab087944a6e8cc246b85d32843e.md}{Canonical record: b6d8aab087944a6e8cc246b85d32843e}; full qualifications and provenance.}\par
\paragraph*{III.4.1.2.2.2 Demographic Update Notification}
\phantomsection
\label{navigation:III.4.1.2.2.2}
\addcontentsline{toc}{paragraph}{III.4.1.2.2.2 Demographic Update Notification}

\label{element:10b54e9ef0ed47e491d18731c58247f5}
\index{Demographic Update Notification (III.4.1.2.2.2)}
{\small BusinessService · Canonical · current}\par

Emits change notifications across the HIE.


{\footnotesize \href{../../../../continuum/business/business-service/10b54e9ef0ed47e491d18731c58247f5.md}{Canonical record: 10b54e9ef0ed47e491d18731c58247f5}; full qualifications and provenance.}\par
\subsubsection*{III.4.1.2.3 Business Functions}
\phantomsection
\label{navigation:III.4.1.2.3}
\addcontentsline{toc}{subsubsection}{III.4.1.2.3 Business Functions}

\paragraph*{III.4.1.2.3.1 Distribute Governed Demographic Change}
\phantomsection
\label{navigation:III.4.1.2.3.1}
\addcontentsline{toc}{paragraph}{III.4.1.2.3.1 Distribute Governed Demographic Change}

\label{element:61a4ae56602c4e2d9428169188e458d8}
\index{Distribute Governed Demographic Change (III.4.1.2.3.1)}
{\small BusinessFunction · Canonical · current}\par

Propagates verified demographic modifications to subscribed clinical subsystems.


{\footnotesize \href{../../../../continuum/business/business-function/61a4ae56602c4e2d9428169188e458d8.md}{Canonical record: 61a4ae56602c4e2d9428169188e458d8}; full qualifications and provenance.}\par
\paragraph*{III.4.1.2.3.2 Receive Client Demographic Change}
\phantomsection
\label{navigation:III.4.1.2.3.2}
\addcontentsline{toc}{paragraph}{III.4.1.2.3.2 Receive Client Demographic Change}

\label{element:459db1af34724863a57961ed921763d6}
\index{Receive Client Demographic Change (III.4.1.2.3.2)}
{\small BusinessFunction · Canonical · current}\par

Ingests administrative demographic updates from PAS/registration feeds.


{\footnotesize \href{../../../../continuum/business/business-function/459db1af34724863a57961ed921763d6.md}{Canonical record: 459db1af34724863a57961ed921763d6}; full qualifications and provenance.}\par
\subsection*{III.4.1.3 [CT2] Clinical Communication Administration — Capability context}
\phantomsection
\label{navigation:III.4.1.3}
\addcontentsline{toc}{subsection}{III.4.1.3 [CT2] Clinical Communication Administration — Capability context}

{\small Canonical Strategy Capability: II.1.3.1.8 \hyperref[element:b38f2c2811724734b5887b6dbaaabbfc]{[CT2] Clinical Communication Administration}; Canonical ID: \texttt{\seqsplit{b38f2c2811724734b5887b6dbaaabbfc}}. Canonical · current.}\par
Contextual reference only; no additional architectural instance or semantic relationship.\par
\subsubsection*{III.4.1.3.1 Business Services}
\phantomsection
\label{navigation:III.4.1.3.1}
\addcontentsline{toc}{subsubsection}{III.4.1.3.1 Business Services}

\paragraph*{III.4.1.3.1.1 Communication Audit Query}
\phantomsection
\label{navigation:III.4.1.3.1.1}
\addcontentsline{toc}{paragraph}{III.4.1.3.1.1 Communication Audit Query}

\label{element:7c871898fd634a16902048da75586fd9}
\index{Communication Audit Query (III.4.1.3.1.1)}
{\small BusinessService · Canonical · current}\par

Discloses transmission history for medico-legal verification.


{\footnotesize \href{../../../../continuum/business/business-service/7c871898fd634a16902048da75586fd9.md}{Canonical record: 7c871898fd634a16902048da75586fd9}; full qualifications and provenance.}\par
\paragraph*{III.4.1.3.1.2 Communication Delivery Status Query}
\phantomsection
\label{navigation:III.4.1.3.1.2}
\addcontentsline{toc}{paragraph}{III.4.1.3.1.2 Communication Delivery Status Query}

\label{element:1b3447db9891406b951f6a6ccf20763e}
\index{Communication Delivery Status Query (III.4.1.3.1.2)}
{\small BusinessService · Canonical · current}\par

Discloses delivery confirmation to message senders.


{\footnotesize \href{../../../../continuum/business/business-service/1b3447db9891406b951f6a6ccf20763e.md}{Canonical record: 1b3447db9891406b951f6a6ccf20763e}; full qualifications and provenance.}\par
\paragraph*{III.4.1.3.1.3 Secure Clinical Message Dispatch}
\phantomsection
\label{navigation:III.4.1.3.1.3}
\addcontentsline{toc}{paragraph}{III.4.1.3.1.3 Secure Clinical Message Dispatch}

\label{element:d28fb3b4123b4e569a95a823fd83f837}
\index{Secure Clinical Message Dispatch (III.4.1.3.1.3)}
{\small BusinessService · Canonical · current}\par

Exposes message dispatch to clinical authors.


{\footnotesize \href{../../../../continuum/business/business-service/d28fb3b4123b4e569a95a823fd83f837.md}{Canonical record: d28fb3b4123b4e569a95a823fd83f837}; full qualifications and provenance.}\par
\subsubsection*{III.4.1.3.2 Business Functions}
\phantomsection
\label{navigation:III.4.1.3.2}
\addcontentsline{toc}{subsubsection}{III.4.1.3.2 Business Functions}

\paragraph*{III.4.1.3.2.1 Distribute Clinical Communication}
\phantomsection
\label{navigation:III.4.1.3.2.1}
\addcontentsline{toc}{paragraph}{III.4.1.3.2.1 Distribute Clinical Communication}

\label{element:5b8d4d4a3f7147fc9ee5f114155c2db8}
\index{Distribute Clinical Communication (III.4.1.3.2.1)}
{\small BusinessFunction · Canonical · current}\par

Routes secure messages, referral letters, and critical alerts to practitioner endpoints.


{\footnotesize \href{../../../../continuum/business/business-function/5b8d4d4a3f7147fc9ee5f114155c2db8.md}{Canonical record: 5b8d4d4a3f7147fc9ee5f114155c2db8}; full qualifications and provenance.}\par
\paragraph*{III.4.1.3.2.2 Record Communication Transaction}
\phantomsection
\label{navigation:III.4.1.3.2.2}
\addcontentsline{toc}{paragraph}{III.4.1.3.2.2 Record Communication Transaction}

\label{element:d64ced698052400fa7892e4d768b3f8f}
\index{Record Communication Transaction (III.4.1.3.2.2)}
{\small BusinessFunction · Canonical · current}\par

Maintains an immutable audit trail of clinical message dispatch and receipt.


{\footnotesize \href{../../../../continuum/business/business-function/d64ced698052400fa7892e4d768b3f8f.md}{Canonical record: d64ced698052400fa7892e4d768b3f8f}; full qualifications and provenance.}\par
\paragraph*{III.4.1.3.2.3 Track Business Delivery Acknowledgement}
\phantomsection
\label{navigation:III.4.1.3.2.3}
\addcontentsline{toc}{paragraph}{III.4.1.3.2.3 Track Business Delivery Acknowledgement}

\label{element:75fb22db4c2346b99fd7c788264abb5d}
\index{Track Business Delivery Acknowledgement (III.4.1.3.2.3)}
{\small BusinessFunction · Canonical · current}\par

Tracks end-to-end delivery state and records business-level delivery/rejection acknowledgements.


{\footnotesize \href{../../../../continuum/business/business-function/75fb22db4c2346b99fd7c788264abb5d.md}{Canonical record: 75fb22db4c2346b99fd7c788264abb5d}; full qualifications and provenance.}\par
\subsection*{III.4.1.4 [CT2] Clinical Record Administration — Capability context}
\phantomsection
\label{navigation:III.4.1.4}
\addcontentsline{toc}{subsection}{III.4.1.4 [CT2] Clinical Record Administration — Capability context}

{\small Canonical Strategy Capability: II.1.3.1.9 \hyperref[element:ce9e7e7b22054cefb707a650d5be7219]{[CT2] Clinical Record Administration}; Canonical ID: \texttt{\seqsplit{ce9e7e7b22054cefb707a650d5be7219}}. Canonical · current.}\par
Contextual reference only; no additional architectural instance or semantic relationship.\par
\subsubsection*{III.4.1.4.1 Business Services}
\phantomsection
\label{navigation:III.4.1.4.1}
\addcontentsline{toc}{subsubsection}{III.4.1.4.1 Business Services}

\paragraph*{III.4.1.4.1.1 Clinical Document Ingress}
\phantomsection
\label{navigation:III.4.1.4.1.1}
\addcontentsline{toc}{paragraph}{III.4.1.4.1.1 Clinical Document Ingress}

\label{element:2a69a33c93954a868249270d919dabfd}
\index{Clinical Document Ingress (III.4.1.4.1.1)}
{\small BusinessService · Canonical · current}\par

Exposes document submission endpoints across the HIE.


{\footnotesize \href{../../../../continuum/business/business-service/2a69a33c93954a868249270d919dabfd.md}{Canonical record: 2a69a33c93954a868249270d919dabfd}; full qualifications and provenance.}\par
\paragraph*{III.4.1.4.1.2 Clinical Document Lifecycle Service}
\phantomsection
\label{navigation:III.4.1.4.1.2}
\addcontentsline{toc}{paragraph}{III.4.1.4.1.2 Clinical Document Lifecycle Service}

\label{element:fbeffe86f67347e3b39ad4544b13cbe7}
\index{Clinical Document Lifecycle Service (III.4.1.4.1.2)}
{\small BusinessService · Canonical · current}\par

Allows authoring clinicians to publish addenda or corrections.


{\footnotesize \href{../../../../continuum/business/business-service/fbeffe86f67347e3b39ad4544b13cbe7.md}{Canonical record: fbeffe86f67347e3b39ad4544b13cbe7}; full qualifications and provenance.}\par
\paragraph*{III.4.1.4.1.3 Document Metadata Registry Query}
\phantomsection
\label{navigation:III.4.1.4.1.3}
\addcontentsline{toc}{paragraph}{III.4.1.4.1.3 Document Metadata Registry Query}

\label{element:63254ecfe5554eeeb63d0c9f1720388b}
\index{Document Metadata Registry Query (III.4.1.4.1.3)}
{\small BusinessService · Canonical · current}\par

Discloses document registry metadata to clinical viewers.


{\footnotesize \href{../../../../continuum/business/business-service/63254ecfe5554eeeb63d0c9f1720388b.md}{Canonical record: 63254ecfe5554eeeb63d0c9f1720388b}; full qualifications and provenance.}\par
\subsubsection*{III.4.1.4.2 Business Functions}
\phantomsection
\label{navigation:III.4.1.4.2}
\addcontentsline{toc}{subsubsection}{III.4.1.4.2 Business Functions}

\paragraph*{III.4.1.4.2.1 Govern Clinical Document Lifecycle}
\phantomsection
\label{navigation:III.4.1.4.2.1}
\addcontentsline{toc}{paragraph}{III.4.1.4.2.1 Govern Clinical Document Lifecycle}

\label{element:e6d00e00eba64ef4bc41d2af8bc3b944}
\index{Govern Clinical Document Lifecycle (III.4.1.4.2.1)}
{\small BusinessFunction · Canonical · current}\par

Manages document states (scoped illustration, not a universal transition model): Draft → Preliminary → Final → Amended → Superseded → Entered-in-Error.


{\footnotesize \href{../../../../continuum/business/business-function/e6d00e00eba64ef4bc41d2af8bc3b944.md}{Canonical record: e6d00e00eba64ef4bc41d2af8bc3b944}; full qualifications and provenance.}\par
\paragraph*{III.4.1.4.2.2 Index Clinical Document Metadata}
\phantomsection
\label{navigation:III.4.1.4.2.2}
\addcontentsline{toc}{paragraph}{III.4.1.4.2.2 Index Clinical Document Metadata}

\label{element:0ee35bd069924319a7a9fe115d3c52b9}
\index{Index Clinical Document Metadata (III.4.1.4.2.2)}
{\small BusinessFunction · Canonical · current}\par

Indexes author, specialty, encounter, date, and document type (LOINC/SNOMED).


{\footnotesize \href{../../../../continuum/business/business-function/0ee35bd069924319a7a9fe115d3c52b9.md}{Canonical record: 0ee35bd069924319a7a9fe115d3c52b9}; full qualifications and provenance.}\par
\paragraph*{III.4.1.4.2.3 Receive Clinical Document}
\phantomsection
\label{navigation:III.4.1.4.2.3}
\addcontentsline{toc}{paragraph}{III.4.1.4.2.3 Receive Clinical Document}

\label{element:d8f80ba3a5a9401b84e6476cbec8cace}
\index{Receive Clinical Document (III.4.1.4.2.3)}
{\small BusinessFunction · Canonical · current}\par

Receives clinical documents from information suppliers.


{\footnotesize \href{../../../../continuum/business/business-function/d8f80ba3a5a9401b84e6476cbec8cace.md}{Canonical record: d8f80ba3a5a9401b84e6476cbec8cace}; full qualifications and provenance.}\par
\subsubsection*{III.4.1.4.3 Business Processes}
\phantomsection
\label{navigation:III.4.1.4.3}
\addcontentsline{toc}{subsubsection}{III.4.1.4.3 Business Processes}

\paragraph*{III.4.1.4.3.1 Clinical Document Lifecycle Process}
\phantomsection
\label{navigation:III.4.1.4.3.1}
\addcontentsline{toc}{paragraph}{III.4.1.4.3.1 Clinical Document Lifecycle Process}

\label{element:eacdec16b96844b4963b8a960d1a8666}
\index{Clinical Document Lifecycle Process (III.4.1.4.3.1)}
{\small BusinessProcess · Canonical · current}\par

•
Owning Capability: Clinical Record Administration



•
Process Purpose: Governs the versioning, clinical sign-off, addenda, superseding, and legal status of clinical documents (discharge summaries, specialist letters, advance care directives).



•
State Progression Lifecycle: Draft → Preliminary → Final Signed → Amended / Addended → Superseded → Entered-in-Error   (Scoped illustration; no universal document lifecycle)



•
Key State Dispositions:



•
Draft: Incomplete clinical document saved during consultation.



•
Preliminary: Document authored pending senior registrar or consultant countersignature.



•
Final Signed: Authoritative clinical document legally signed and published.



•
Amended / Addended: Governed supplementary clinical addendum appended to signed document.



•
Superseded: Entire document replaced by a newer revision; prior version retained for audit.



•
Entered-in-Error: Document formally retracted due to wrong patient or invalid content; content struck through with retraction notice.



•
Scope qualification: Final ≠ automatically Final Signed. The described signing/countersigning and legal qualifications remain local requirements, not a universal document lifecycle. Authorship, attestation, approval, authentication, verification, signature, finalisation, authority and legal qualification remain distinct. Amendment, supersession and entered-in-error are not mandatory stages for every document, and consumer processing does not establish originating authority or legal status.



This Process makes authorship and changing document standing intelligible through the locally described review, amendment and withdrawal circumstances. It preserves prior versions rather than treating correction as erasure. The illustration does not require every document to pass through every stage or make finalisation, signing, legal qualification and originating authority equivalent.


{\small Recorded unresolved matters: source-qualified matters; exact scope remains in the canonical record.}\par
{\footnotesize \href{../../../../continuum/business/business-process/eacdec16b96844b4963b8a960d1a8666.md}{Canonical record: eacdec16b96844b4963b8a960d1a8666}; full qualifications and provenance.}\par
\subsection*{III.4.1.5 [CT2] Diagnostic Administration — Capability context}
\phantomsection
\label{navigation:III.4.1.5}
\addcontentsline{toc}{subsection}{III.4.1.5 [CT2] Diagnostic Administration — Capability context}

{\small Canonical Strategy Capability: II.1.3.1.10 \hyperref[element:3cc720ba2e8b4d7a94bb349391d61895]{[CT2] Diagnostic Administration}; Canonical ID: \texttt{\seqsplit{3cc720ba2e8b4d7a94bb349391d61895}}. Canonical · current.}\par
Contextual reference only; no additional architectural instance or semantic relationship.\par
\subsubsection*{III.4.1.5.1 [CT3] Diagnostic Service Integration — Capability context}
\phantomsection
\label{navigation:III.4.1.5.1}
\addcontentsline{toc}{subsubsection}{III.4.1.5.1 [CT3] Diagnostic Service Integration — Capability context}

{\small Canonical Strategy Capability: II.1.3.1.10.1 \hyperref[element:9ded2f3a28d24017bb921a0cdcf65a2b]{[CT3] Diagnostic Service Integration}; Canonical ID: \texttt{\seqsplit{9ded2f3a28d24017bb921a0cdcf65a2b}}. Canonical · current.}\par
Contextual reference only; no additional architectural instance or semantic relationship.\par
\paragraph*{III.4.1.5.1.1 Business Services}
\phantomsection
\label{navigation:III.4.1.5.1.1}
\addcontentsline{toc}{paragraph}{III.4.1.5.1.1 Business Services}

\subparagraph*{III.4.1.5.1.1.1 Diagnostic Publication Service}
\phantomsection
\label{navigation:III.4.1.5.1.1.1}
\addcontentsline{toc}{subparagraph}{III.4.1.5.1.1.1 Diagnostic Publication Service}

\label{element:600ca5384b4643b0b35579a02c307dcf}
\index{Diagnostic Publication Service (III.4.1.5.1.1.1)}
{\small BusinessService · Canonical · current}\par

Publishes verified diagnostic reports across the enterprise.



Feature association not established: Strategy FEAT-SD-10 — Diagnostic History Consolidation remains valid, but semantic evidence does not establish that the following Function or Service belongs to, realises, specialises or is equivalent to it. No replacement Feature is assigned.


{\small Recorded unresolved matters: source-qualified matters; exact scope remains in the canonical record.}\par
{\footnotesize \href{../../../../continuum/business/business-service/600ca5384b4643b0b35579a02c307dcf.md}{Canonical record: 600ca5384b4643b0b35579a02c307dcf}; full qualifications and provenance.}\par
\subparagraph*{III.4.1.5.1.1.2 Modality Worklist Distribution}
\phantomsection
\label{navigation:III.4.1.5.1.1.2}
\addcontentsline{toc}{subparagraph}{III.4.1.5.1.1.2 Modality Worklist Distribution}

\label{element:e254b6a140c94b33a206e74db1aecfbb}
\index{Modality Worklist Distribution (III.4.1.5.1.1.2)}
{\small BusinessService · Canonical · current}\par

Supplies order worklists to diagnostic laboratory/imaging systems.



Feature association not established: Strategy FEAT-SD-09 — Laboratory \& Imaging Result Distribution remains valid, but semantic evidence does not establish that the following Function or Service belongs to, realises, specialises or is equivalent to it. No replacement Feature is assigned.


{\small Recorded unresolved matters: source-qualified matters; exact scope remains in the canonical record.}\par
{\footnotesize \href{../../../../continuum/business/business-service/e254b6a140c94b33a206e74db1aecfbb.md}{Canonical record: e254b6a140c94b33a206e74db1aecfbb}; full qualifications and provenance.}\par
\paragraph*{III.4.1.5.1.2 Business Functions}
\phantomsection
\label{navigation:III.4.1.5.1.2}
\addcontentsline{toc}{paragraph}{III.4.1.5.1.2 Business Functions}

\subparagraph*{III.4.1.5.1.2.1 Publish Governed Diagnostic Report}
\phantomsection
\label{navigation:III.4.1.5.1.2.1}
\addcontentsline{toc}{subparagraph}{III.4.1.5.1.2.1 Publish Governed Diagnostic Report}

\label{element:f9fcdc3162044865ac6f26015fdd5b17}
\index{Publish Governed Diagnostic Report (III.4.1.5.1.2.1)}
{\small BusinessFunction · Canonical · current}\par

Normalises and distributes structured diagnostic observations, reference ranges, and diagnostic image study links into the LHR.



Feature association not established: Strategy FEAT-SD-10 — Diagnostic History Consolidation remains valid, but semantic evidence does not establish that the following Function or Service belongs to, realises, specialises or is equivalent to it. No replacement Feature is assigned.


{\small Recorded unresolved matters: source-qualified matters; exact scope remains in the canonical record.}\par
{\footnotesize \href{../../../../continuum/business/business-function/f9fcdc3162044865ac6f26015fdd5b17.md}{Canonical record: f9fcdc3162044865ac6f26015fdd5b17}; full qualifications and provenance.}\par
\subparagraph*{III.4.1.5.1.2.2 Receive Diagnostic Requisition Feed}
\phantomsection
\label{navigation:III.4.1.5.1.2.2}
\addcontentsline{toc}{subparagraph}{III.4.1.5.1.2.2 Receive Diagnostic Requisition Feed}

\label{element:af3d47ada486425abc5ab5d054db7004}
\index{Receive Diagnostic Requisition Feed (III.4.1.5.1.2.2)}
{\small BusinessFunction · Canonical · current}\par

Normalises order requisitions from across the health network into modality worklists.



Feature association not established: Strategy FEAT-SD-09 — Laboratory \& Imaging Result Distribution remains valid, but semantic evidence does not establish that the following Function or Service belongs to, realises, specialises or is equivalent to it. No replacement Feature is assigned.


{\small Recorded unresolved matters: source-qualified matters; exact scope remains in the canonical record.}\par
{\footnotesize \href{../../../../continuum/business/business-function/af3d47ada486425abc5ab5d054db7004.md}{Canonical record: af3d47ada486425abc5ab5d054db7004}; full qualifications and provenance.}\par
\subsubsection*{III.4.1.5.2 Business Services}
\phantomsection
\label{navigation:III.4.1.5.2}
\addcontentsline{toc}{subsubsection}{III.4.1.5.2 Business Services}

\paragraph*{III.4.1.5.2.1 Diagnostic Correlation Resolution}
\phantomsection
\label{navigation:III.4.1.5.2.1}
\addcontentsline{toc}{paragraph}{III.4.1.5.2.1 Diagnostic Correlation Resolution}

\label{element:15f3582e02ae47e1b7f9117573a8ddf7}
\index{Diagnostic Correlation Resolution (III.4.1.5.2.1)}
{\small BusinessService · Canonical · current}\par

Discloses correlated order/result linkages.


{\footnotesize \href{../../../../continuum/business/business-service/15f3582e02ae47e1b7f9117573a8ddf7.md}{Canonical record: 15f3582e02ae47e1b7f9117573a8ddf7}; full qualifications and provenance.}\par
\paragraph*{III.4.1.5.2.2 Diagnostic Report Ingress}
\phantomsection
\label{navigation:III.4.1.5.2.2}
\addcontentsline{toc}{paragraph}{III.4.1.5.2.2 Diagnostic Report Ingress}

\label{element:b151ad42b9194d9e84b1dd9beecdf6dd}
\index{Diagnostic Report Ingress (III.4.1.5.2.2)}
{\small BusinessService · Canonical · current}\par

Ingests results from external LIS/RIS systems.


{\footnotesize \href{../../../../continuum/business/business-service/b151ad42b9194d9e84b1dd9beecdf6dd.md}{Canonical record: b151ad42b9194d9e84b1dd9beecdf6dd}; full qualifications and provenance.}\par
\paragraph*{III.4.1.5.2.3 Diagnostic Result Delivery Service}
\phantomsection
\label{navigation:III.4.1.5.2.3}
\addcontentsline{toc}{paragraph}{III.4.1.5.2.3 Diagnostic Result Delivery Service}

\label{element:86a9d1868c8c45788791a9436794531c}
\index{Diagnostic Result Delivery Service (III.4.1.5.2.3)}
{\small BusinessService · Canonical · current}\par

Emits secure result notifications.


{\footnotesize \href{../../../../continuum/business/business-service/86a9d1868c8c45788791a9436794531c.md}{Canonical record: 86a9d1868c8c45788791a9436794531c}; full qualifications and provenance.}\par
\subsubsection*{III.4.1.5.3 Business Functions}
\phantomsection
\label{navigation:III.4.1.5.3}
\addcontentsline{toc}{subsubsection}{III.4.1.5.3 Business Functions}

\paragraph*{III.4.1.5.3.1 Bind Diagnostic Report}
\phantomsection
\label{navigation:III.4.1.5.3.1}
\addcontentsline{toc}{paragraph}{III.4.1.5.3.1 Bind Diagnostic Report}

\label{element:ee4bd6218f3a4029815bfc16912b88af}
\index{Bind Diagnostic Report (III.4.1.5.3.1)}
{\small BusinessFunction · Canonical · current}\par

Associates structured observations, reference ranges, and abnormal flags with patient health records.


{\footnotesize \href{../../../../continuum/business/business-function/ee4bd6218f3a4029815bfc16912b88af.md}{Canonical record: ee4bd6218f3a4029815bfc16912b88af}; full qualifications and provenance.}\par
\paragraph*{III.4.1.5.3.2 Correlate Diagnostic Request}
\phantomsection
\label{navigation:III.4.1.5.3.2}
\addcontentsline{toc}{paragraph}{III.4.1.5.3.2 Correlate Diagnostic Request}

\label{element:6797819b02e0405c8169fa57544d9606}
\index{Correlate Diagnostic Request (III.4.1.5.3.2)}
{\small BusinessFunction · Canonical · current}\par

Matches incoming lab/radiology results to active patient encounters and order requisitions.


{\footnotesize \href{../../../../continuum/business/business-function/6797819b02e0405c8169fa57544d9606.md}{Canonical record: 6797819b02e0405c8169fa57544d9606}; full qualifications and provenance.}\par
\paragraph*{III.4.1.5.3.3 Distribute Diagnostic Report}
\phantomsection
\label{navigation:III.4.1.5.3.3}
\addcontentsline{toc}{paragraph}{III.4.1.5.3.3 Distribute Diagnostic Report}

\label{element:ab6d94a9e06742aabe087f11d93f4334}
\index{Distribute Diagnostic Report (III.4.1.5.3.3)}
{\small BusinessFunction · Canonical · current}\par

Delivers preliminary, final, and corrected reports to ordering and copied practitioners.


{\footnotesize \href{../../../../continuum/business/business-function/ab6d94a9e06742aabe087f11d93f4334.md}{Canonical record: ab6d94a9e06742aabe087f11d93f4334}; full qualifications and provenance.}\par
\subsection*{III.4.1.6 [CT2] Episode \& Encounter Administration — Capability context}
\phantomsection
\label{navigation:III.4.1.6}
\addcontentsline{toc}{subsection}{III.4.1.6 [CT2] Episode \& Encounter Administration — Capability context}

{\small Canonical Strategy Capability: II.1.3.1.11 \hyperref[element:61d7067aee0540259b3daca4287d66d2]{[CT2] Episode \& Encounter Administration}; Canonical ID: \texttt{\seqsplit{61d7067aee0540259b3daca4287d66d2}}. Canonical · current.}\par
Contextual reference only; no additional architectural instance or semantic relationship.\par
\subsubsection*{III.4.1.6.1 [CT3] Acute Care Monitoring \& Event Coordination — Capability context}
\phantomsection
\label{navigation:III.4.1.6.1}
\addcontentsline{toc}{subsubsection}{III.4.1.6.1 [CT3] Acute Care Monitoring \& Event Coordination — Capability context}

{\small Canonical Strategy Capability: II.1.3.1.11.1 \hyperref[element:ee17b26b8c874cc3af50d07a787ce488]{[CT3] Acute Care Monitoring \& Event Coordination}; Canonical ID: \texttt{\seqsplit{ee17b26b8c874cc3af50d07a787ce488}}. Canonical · current.}\par
Contextual reference only; no additional architectural instance or semantic relationship.\par
\paragraph*{III.4.1.6.1.1 Business Services}
\phantomsection
\label{navigation:III.4.1.6.1.1}
\addcontentsline{toc}{paragraph}{III.4.1.6.1.1 Business Services}

\subparagraph*{III.4.1.6.1.1.1 Acute Admission Dossier Query}
\phantomsection
\label{navigation:III.4.1.6.1.1.1}
\addcontentsline{toc}{subparagraph}{III.4.1.6.1.1.1 Acute Admission Dossier Query}

\label{element:7276d5e4d484440f9466ec24c183bdab}
\index{Acute Admission Dossier Query (III.4.1.6.1.1.1)}
{\small BusinessService · Canonical · current}\par

Delivers a high-priority clinical briefing to admitting teams.


{\footnotesize \href{../../../../continuum/business/business-service/7276d5e4d484440f9466ec24c183bdab.md}{Canonical record: 7276d5e4d484440f9466ec24c183bdab}; full qualifications and provenance.}\par
\subparagraph*{III.4.1.6.1.1.2 Acute Clinical Alert Dispatch}
\phantomsection
\label{navigation:III.4.1.6.1.1.2}
\addcontentsline{toc}{subparagraph}{III.4.1.6.1.1.2 Acute Clinical Alert Dispatch}

\label{element:34d7f0bfff8f470bb97c00582f5e8a1d}
\index{Acute Clinical Alert Dispatch (III.4.1.6.1.1.2)}
{\small BusinessService · Canonical · current}\par

Dispatches urgent interruptive alerts to on-duty ward teams.



Feature association not established: FEAT-SD-03 — Acute Clinical State Event Capture is adjudicated sufficient through general clinical-information responsibilities above. Routing a critical-result alert remains distinct from capture; no specialised capture Function or replacement association is assigned.


{\small Recorded unresolved matters: source-qualified matters; exact scope remains in the canonical record.}\par
{\footnotesize \href{../../../../continuum/business/business-service/34d7f0bfff8f470bb97c00582f5e8a1d.md}{Canonical record: 34d7f0bfff8f470bb97c00582f5e8a1d}; full qualifications and provenance.}\par
\paragraph*{III.4.1.6.1.2 Business Functions}
\phantomsection
\label{navigation:III.4.1.6.1.2}
\addcontentsline{toc}{paragraph}{III.4.1.6.1.2 Business Functions}

\subparagraph*{III.4.1.6.1.2.1 Assemble Acute Clinical Baseline}
\phantomsection
\label{navigation:III.4.1.6.1.2.1}
\addcontentsline{toc}{subparagraph}{III.4.1.6.1.2.1 Assemble Acute Clinical Baseline}

\label{element:be64da1302bf4c4bae357295612e8e6e}
\index{Assemble Acute Clinical Baseline (III.4.1.6.1.2.1)}
{\small BusinessFunction · Canonical · current}\par

Rapidly aggregates previous discharge summaries, active medications, allergies, and diagnostic alerts upon acute hospital presentation.


{\footnotesize \href{../../../../continuum/business/business-function/be64da1302bf4c4bae357295612e8e6e.md}{Canonical record: be64da1302bf4c4bae357295612e8e6e}; full qualifications and provenance.}\par
\subparagraph*{III.4.1.6.1.2.2 Route Acute Critical Result Alert}
\phantomsection
\label{navigation:III.4.1.6.1.2.2}
\addcontentsline{toc}{subparagraph}{III.4.1.6.1.2.2 Route Acute Critical Result Alert}

\label{element:66d0295d5a6c4eedba6406d6f672906c}
\index{Route Acute Critical Result Alert (III.4.1.6.1.2.2)}
{\small BusinessFunction · Canonical · current}\par

Immediately directs urgent lab/radiology alerts (e.g., critical troponin, intracranial bleed) to active acute attending clinicians.



Feature association not established: FEAT-SD-03 — Acute Clinical State Event Capture is adjudicated sufficient through general clinical-information responsibilities above. Routing a critical-result alert remains distinct from capture; no specialised capture Function or replacement association is assigned.


{\small Recorded unresolved matters: source-qualified matters; exact scope remains in the canonical record.}\par
{\footnotesize \href{../../../../continuum/business/business-function/66d0295d5a6c4eedba6406d6f672906c.md}{Canonical record: 66d0295d5a6c4eedba6406d6f672906c}; full qualifications and provenance.}\par
\subsubsection*{III.4.1.6.2 Business Services}
\phantomsection
\label{navigation:III.4.1.6.2}
\addcontentsline{toc}{subsubsection}{III.4.1.6.2 Business Services}

\paragraph*{III.4.1.6.2.1 Encounter Context Resolution}
\phantomsection
\label{navigation:III.4.1.6.2.1}
\addcontentsline{toc}{paragraph}{III.4.1.6.2.1 Encounter Context Resolution}

\label{element:ca22b0a4e51b4f5eb6b528bc36fa5d9c}
\index{Encounter Context Resolution (III.4.1.6.2.1)}
{\small BusinessService · Canonical · current}\par

Discloses active encounter details to clinical applications.


{\footnotesize \href{../../../../continuum/business/business-service/ca22b0a4e51b4f5eb6b528bc36fa5d9c.md}{Canonical record: ca22b0a4e51b4f5eb6b528bc36fa5d9c}; full qualifications and provenance.}\par
\paragraph*{III.4.1.6.2.2 Encounter Lifecycle Event Notification}
\phantomsection
\label{navigation:III.4.1.6.2.2}
\addcontentsline{toc}{paragraph}{III.4.1.6.2.2 Encounter Lifecycle Event Notification}

\label{element:e9c6b81172984785a47d222ac0c28d45}
\index{Encounter Lifecycle Event Notification (III.4.1.6.2.2)}
{\small BusinessService · Canonical · current}\par

Emits encounter transition events across the HIE.


{\footnotesize \href{../../../../continuum/business/business-service/e9c6b81172984785a47d222ac0c28d45.md}{Canonical record: e9c6b81172984785a47d222ac0c28d45}; full qualifications and provenance.}\par
\paragraph*{III.4.1.6.2.3 Encounter Location History Query}
\phantomsection
\label{navigation:III.4.1.6.2.3}
\addcontentsline{toc}{paragraph}{III.4.1.6.2.3 Encounter Location History Query}

\label{element:5f0f9fec1d16458aaf914b0957793953}
\index{Encounter Location History Query (III.4.1.6.2.3)}
{\small BusinessService · Canonical · current}\par

Discloses patient location tracking history.


{\footnotesize \href{../../../../continuum/business/business-service/5f0f9fec1d16458aaf914b0957793953.md}{Canonical record: 5f0f9fec1d16458aaf914b0957793953}; full qualifications and provenance.}\par
\subsubsection*{III.4.1.6.3 Business Functions}
\phantomsection
\label{navigation:III.4.1.6.3}
\addcontentsline{toc}{subsubsection}{III.4.1.6.3 Business Functions}

\paragraph*{III.4.1.6.3.1 Establish Encounter Context}
\phantomsection
\label{navigation:III.4.1.6.3.1}
\addcontentsline{toc}{paragraph}{III.4.1.6.3.1 Establish Encounter Context}

\label{element:df00c187d509483a8408304a066abec4}
\index{Establish Encounter Context (III.4.1.6.3.1)}
{\small BusinessFunction · Canonical · current}\par

Binds patient, attending practitioner, clinical class (e.g., Inpatient, Emergency, Ambulatory), and admitting diagnosis to a unique encounter identifier.


{\footnotesize \href{../../../../continuum/business/business-function/df00c187d509483a8408304a066abec4.md}{Canonical record: df00c187d509483a8408304a066abec4}; full qualifications and provenance.}\par
\paragraph*{III.4.1.6.3.2 Maintain Encounter Care-Place Association}
\phantomsection
\label{navigation:III.4.1.6.3.2}
\addcontentsline{toc}{paragraph}{III.4.1.6.3.2 Maintain Encounter Care-Place Association}

\label{element:3ff076243fe141759e7d4401389b3527}
\index{Maintain Encounter Care-Place Association (III.4.1.6.3.2)}
{\small BusinessFunction · Canonical · current}\par

Binds the encounter to physical wards, rooms, and beds over time.


{\footnotesize \href{../../../../continuum/business/business-function/3ff076243fe141759e7d4401389b3527.md}{Canonical record: 3ff076243fe141759e7d4401389b3527}; full qualifications and provenance.}\par
\paragraph*{III.4.1.6.3.3 Progress Encounter State}
\phantomsection
\label{navigation:III.4.1.6.3.3}
\addcontentsline{toc}{paragraph}{III.4.1.6.3.3 Progress Encounter State}

\label{element:223b1140628b444482bfb0a2c870d751}
\index{Progress Encounter State (III.4.1.6.3.3)}
{\small BusinessFunction · Canonical · current}\par

Coordinates encounter progression (scoped illustration; ON\_LEAVE relationship unresolved): Planned → Arrived → In-Progress → Discharged → Completed.


{\footnotesize \href{../../../../continuum/business/business-function/223b1140628b444482bfb0a2c870d751.md}{Canonical record: 223b1140628b444482bfb0a2c870d751}; full qualifications and provenance.}\par
\subsubsection*{III.4.1.6.4 Business Processes}
\phantomsection
\label{navigation:III.4.1.6.4}
\addcontentsline{toc}{subsubsection}{III.4.1.6.4 Business Processes}

\paragraph*{III.4.1.6.4.1 Encounter Lifecycle Process}
\phantomsection
\label{navigation:III.4.1.6.4.1}
\addcontentsline{toc}{paragraph}{III.4.1.6.4.1 Encounter Lifecycle Process}

\label{element:df49e2e8bb3748f3abbdd1c28ae1be8f}
\index{Encounter Lifecycle Process (III.4.1.6.4.1)}
{\small BusinessProcess · Canonical · current}\par

•
Owning Capability: Episode \& Encounter Administration



•
Process Purpose: Governs the clinical and administrative lifecycle of an acute, emergency, inpatient, or outpatient encounter between a patient and healthcare services.



•
State Progression Lifecycle: Planned / Booked → Arrived → Triaged / Ingested → Active In-Progress → Discharged → Completed / Encoded   (Scoped illustration; ON\_LEAVE relationship unresolved)



•
Key State Dispositions:



•
Planned / Booked: Elective admission or clinic appointment scheduled.



•
Arrived: Patient presents at facility or emergency desk.



•
Triaged / Ingested: Clinical triage category assigned; encounter record activated.



•
Active In-Progress: Inpatient care, bedside monitoring, and clinical orders actively underway.



•
Discharged: Patient clinically discharged; care-place vacated.



•
Completed / Encoded: Clinical documentation finalized, ICD/DRG coding completed, and encounter legally closed.



•
Scope qualification: These completion criteria retain their locally described applicability; they do not establish universal signing, legal closure or originating authority over clinical information. Finalisation, signing and legal qualification remain distinct. The relationship with Strategy ON\_LEAVE, including applicability and transitions, remains unresolved.



The encounter lifecycle distinguishes the planned opportunity, presence, active encounter and locally applicable completion context. It supplies an administrative progression frame for care information rather than evidence that every clinical objective has been met. The relation to Strategy's ON\_LEAVE state remains unresolved, and local signing or legal-closure descriptions do not become universal completion rules.


{\small Recorded unresolved matters: source-qualified matters; exact scope remains in the canonical record.}\par
{\footnotesize \href{../../../../continuum/business/business-process/df49e2e8bb3748f3abbdd1c28ae1be8f.md}{Canonical record: df49e2e8bb3748f3abbdd1c28ae1be8f}; full qualifications and provenance.}\par
\subsection*{III.4.1.7 [CT2] Medication Administration — Capability context}
\phantomsection
\label{navigation:III.4.1.7}
\addcontentsline{toc}{subsection}{III.4.1.7 [CT2] Medication Administration — Capability context}

{\small Canonical Strategy Capability: II.1.3.1.13 \hyperref[element:c907e2a95629435b9951941ad4e2ebe9]{[CT2] Medication Administration}; Canonical ID: \texttt{\seqsplit{c907e2a95629435b9951941ad4e2ebe9}}. Canonical · current.}\par
Contextual reference only; no additional architectural instance or semantic relationship.\par
\subsubsection*{III.4.1.7.1 [CT3] Medication Therapy Provision and Monitoring — Capability context}
\phantomsection
\label{navigation:III.4.1.7.1}
\addcontentsline{toc}{subsubsection}{III.4.1.7.1 [CT3] Medication Therapy Provision and Monitoring — Capability context}

{\small Canonical Strategy Capability: II.1.3.1.13.2 \hyperref[element:d1821b4f6d5746f9a697dfe3b792a186]{[CT3] Medication Therapy Provision and Monitoring}; Canonical ID: \texttt{\seqsplit{d1821b4f6d5746f9a697dfe3b792a186}}. Canonical · current.}\par
Contextual reference only; no additional architectural instance or semantic relationship.\par
\paragraph*{III.4.1.7.1.1 Business Services}
\phantomsection
\label{navigation:III.4.1.7.1.1}
\addcontentsline{toc}{paragraph}{III.4.1.7.1.1 Business Services}

\subparagraph*{III.4.1.7.1.1.1 Allergy \& Adverse Reaction Query}
\phantomsection
\label{navigation:III.4.1.7.1.1.1}
\addcontentsline{toc}{subparagraph}{III.4.1.7.1.1.1 Allergy \& Adverse Reaction Query}

\label{element:7820fc4fb98c452aba29f35d8c7e74fb}
\index{Allergy \& Adverse Reaction Query (III.4.1.7.1.1.1)}
{\small BusinessService · Canonical · current}\par

Exposes active allergy profiles to clinical decision support systems.



Feature association not established: Strategy FEAT-SD-12 — Medication Administration Tracking remains valid, but semantic evidence does not establish that the following Function or Service belongs to, realises, specialises or is equivalent to it. No replacement Feature is assigned.


{\small Recorded unresolved matters: source-qualified matters; exact scope remains in the canonical record.}\par
{\footnotesize \href{../../../../continuum/business/business-service/7820fc4fb98c452aba29f35d8c7e74fb.md}{Canonical record: 7820fc4fb98c452aba29f35d8c7e74fb}; full qualifications and provenance.}\par
\subparagraph*{III.4.1.7.1.1.2 Longitudinal Medication History Query}
\phantomsection
\label{navigation:III.4.1.7.1.1.2}
\addcontentsline{toc}{subparagraph}{III.4.1.7.1.1.2 Longitudinal Medication History Query}

\label{element:6a42771bbe5449ab87e8ece94519d03d}
\index{Longitudinal Medication History Query (III.4.1.7.1.1.2)}
{\small BusinessService · Canonical · current}\par

Discloses full medication timelines for medication reconciliation.


{\footnotesize \href{../../../../continuum/business/business-service/6a42771bbe5449ab87e8ece94519d03d.md}{Canonical record: 6a42771bbe5449ab87e8ece94519d03d}; full qualifications and provenance.}\par
\paragraph*{III.4.1.7.1.2 Business Functions}
\phantomsection
\label{navigation:III.4.1.7.1.2}
\addcontentsline{toc}{paragraph}{III.4.1.7.1.2 Business Functions}

\subparagraph*{III.4.1.7.1.2.1 Correlate Adverse Reaction \& Allergy Profile}
\phantomsection
\label{navigation:III.4.1.7.1.2.1}
\addcontentsline{toc}{subparagraph}{III.4.1.7.1.2.1 Correlate Adverse Reaction \& Allergy Profile}

\label{element:834d18fe21f14af7ad0c9773d2aaad1e}
\index{Correlate Adverse Reaction \& Allergy Profile (III.4.1.7.1.2.1)}
{\small BusinessFunction · Canonical · current}\par

Binds confirmed substance allergies, adverse drug reactions, and severity grades to the patient identity.



Feature association not established: Strategy FEAT-SD-12 — Medication Administration Tracking remains valid, but semantic evidence does not establish that the following Function or Service belongs to, realises, specialises or is equivalent to it. No replacement Feature is assigned.


{\small Recorded unresolved matters: source-qualified matters; exact scope remains in the canonical record.}\par
{\footnotesize \href{../../../../continuum/business/business-function/834d18fe21f14af7ad0c9773d2aaad1e.md}{Canonical record: 834d18fe21f14af7ad0c9773d2aaad1e}; full qualifications and provenance.}\par
\subparagraph*{III.4.1.7.1.2.2 Project Longitudinal Medication History}
\phantomsection
\label{navigation:III.4.1.7.1.2.2}
\addcontentsline{toc}{subparagraph}{III.4.1.7.1.2.2 Project Longitudinal Medication History}

\label{element:494f529ada064f77836d40cfe9c7ab48}
\index{Project Longitudinal Medication History (III.4.1.7.1.2.2)}
{\small BusinessFunction · Canonical · current}\par

Aggregates community dispensing, hospital e-MAR administrations, and specialist prescriptions into a unified chronological medication profile.


{\footnotesize \href{../../../../continuum/business/business-function/494f529ada064f77836d40cfe9c7ab48.md}{Canonical record: 494f529ada064f77836d40cfe9c7ab48}; full qualifications and provenance.}\par
\subsubsection*{III.4.1.7.2 Business Services}
\phantomsection
\label{navigation:III.4.1.7.2}
\addcontentsline{toc}{subsubsection}{III.4.1.7.2 Business Services}

\paragraph*{III.4.1.7.2.1 Dispense Event Ingress}
\phantomsection
\label{navigation:III.4.1.7.2.1}
\addcontentsline{toc}{paragraph}{III.4.1.7.2.1 Dispense Event Ingress}

\label{element:b81512ca8e36434e9fa29493e645d31d}
\index{Dispense Event Ingress (III.4.1.7.2.1)}
{\small BusinessService · Canonical · current}\par

Ingests pharmacy supply notifications.


{\footnotesize \href{../../../../continuum/business/business-service/b81512ca8e36434e9fa29493e645d31d.md}{Canonical record: b81512ca8e36434e9fa29493e645d31d}; full qualifications and provenance.}\par
\paragraph*{III.4.1.7.2.2 Medication Administration Ingress}
\phantomsection
\label{navigation:III.4.1.7.2.2}
\addcontentsline{toc}{paragraph}{III.4.1.7.2.2 Medication Administration Ingress}

\label{element:8d224a68f01648d39096d39ebac893df}
\index{Medication Administration Ingress (III.4.1.7.2.2)}
{\small BusinessService · Canonical · current}\par

Ingests administration records.


{\footnotesize \href{../../../../continuum/business/business-service/8d224a68f01648d39096d39ebac893df.md}{Canonical record: 8d224a68f01648d39096d39ebac893df}; full qualifications and provenance.}\par
\paragraph*{III.4.1.7.2.3 Medication Order Ingress}
\phantomsection
\label{navigation:III.4.1.7.2.3}
\addcontentsline{toc}{paragraph}{III.4.1.7.2.3 Medication Order Ingress}

\label{element:3bcfd2c922aa429ba77618ac94982a56}
\index{Medication Order Ingress (III.4.1.7.2.3)}
{\small BusinessService · Canonical · current}\par

Ingests prescription events.


{\footnotesize \href{../../../../continuum/business/business-service/3bcfd2c922aa429ba77618ac94982a56.md}{Canonical record: 3bcfd2c922aa429ba77618ac94982a56}; full qualifications and provenance.}\par
\subsubsection*{III.4.1.7.3 Business Functions}
\phantomsection
\label{navigation:III.4.1.7.3}
\addcontentsline{toc}{subsubsection}{III.4.1.7.3 Business Functions}

\paragraph*{III.4.1.7.3.1 Receive Medication Administration Event}
\phantomsection
\label{navigation:III.4.1.7.3.1}
\addcontentsline{toc}{paragraph}{III.4.1.7.3.1 Receive Medication Administration Event}

\label{element:7437a97ccc3947b38f70d7da47bd0714}
\index{Receive Medication Administration Event (III.4.1.7.3.1)}
{\small BusinessFunction · Canonical · current}\par

Records medication administration-event information, including nurse-administered doses, e-MAR documentation and infusion telemetry as non-exhaustive examples; scope is not limited to nursing, bedside or inpatient settings.


{\footnotesize \href{../../../../continuum/business/business-function/7437a97ccc3947b38f70d7da47bd0714.md}{Canonical record: 7437a97ccc3947b38f70d7da47bd0714}; full qualifications and provenance.}\par
\paragraph*{III.4.1.7.3.2 Receive Medication Order}
\phantomsection
\label{navigation:III.4.1.7.3.2}
\addcontentsline{toc}{paragraph}{III.4.1.7.3.2 Receive Medication Order}

\label{element:96e2ae3ccb684635acf78037e75a1322}
\index{Receive Medication Order (III.4.1.7.3.2)}
{\small BusinessFunction · Canonical · current}\par

Ingests inpatient and outpatient drug prescriptions.


{\footnotesize \href{../../../../continuum/business/business-function/96e2ae3ccb684635acf78037e75a1322.md}{Canonical record: 96e2ae3ccb684635acf78037e75a1322}; full qualifications and provenance.}\par
\paragraph*{III.4.1.7.3.3 Track Dispense Event}
\phantomsection
\label{navigation:III.4.1.7.3.3}
\addcontentsline{toc}{paragraph}{III.4.1.7.3.3 Track Dispense Event}

\label{element:e4c3b7e257b3448aa79d66992e9e9bbc}
\index{Track Dispense Event (III.4.1.7.3.3)}
{\small BusinessFunction · Canonical · current}\par

Ingests community and hospital pharmacy dispensing records.


{\footnotesize \href{../../../../continuum/business/business-function/e4c3b7e257b3448aa79d66992e9e9bbc.md}{Canonical record: e4c3b7e257b3448aa79d66992e9e9bbc}; full qualifications and provenance.}\par
\subsection*{III.4.1.8 [CT2] Order Administration — Capability context}
\phantomsection
\label{navigation:III.4.1.8}
\addcontentsline{toc}{subsection}{III.4.1.8 [CT2] Order Administration — Capability context}

{\small Canonical Strategy Capability: II.1.3.1.14 \hyperref[element:c3c5abd7987f4433a46cbfed8167914c]{[CT2] Order Administration}; Canonical ID: \texttt{\seqsplit{c3c5abd7987f4433a46cbfed8167914c}}. Canonical · current.}\par
Contextual reference only; no additional architectural instance or semantic relationship.\par
\subsubsection*{III.4.1.8.1 Business Services}
\phantomsection
\label{navigation:III.4.1.8.1}
\addcontentsline{toc}{subsubsection}{III.4.1.8.1 Business Services}

\paragraph*{III.4.1.8.1.1 Order Cancellation Service}
\phantomsection
\label{navigation:III.4.1.8.1.1}
\addcontentsline{toc}{paragraph}{III.4.1.8.1.1 Order Cancellation Service}

\label{element:1c8467ef0b6047329bdb899a52c23efc}
\index{Order Cancellation Service (III.4.1.8.1.1)}
{\small BusinessService · Canonical · current}\par

Dispatches cancel directives to performing systems.


{\footnotesize \href{../../../../continuum/business/business-service/1c8467ef0b6047329bdb899a52c23efc.md}{Canonical record: 1c8467ef0b6047329bdb899a52c23efc}; full qualifications and provenance.}\par
\paragraph*{III.4.1.8.1.2 Order Dispatch Service}
\phantomsection
\label{navigation:III.4.1.8.1.2}
\addcontentsline{toc}{paragraph}{III.4.1.8.1.2 Order Dispatch Service}

\label{element:a8b0f7284a854fdfb43a9af446aa1270}
\index{Order Dispatch Service (III.4.1.8.1.2)}
{\small BusinessService · Canonical · current}\par

Transmits orders to performing diagnostic systems.


{\footnotesize \href{../../../../continuum/business/business-service/a8b0f7284a854fdfb43a9af446aa1270.md}{Canonical record: a8b0f7284a854fdfb43a9af446aa1270}; full qualifications and provenance.}\par
\paragraph*{III.4.1.8.1.3 Order Outcome Notification}
\phantomsection
\label{navigation:III.4.1.8.1.3}
\addcontentsline{toc}{paragraph}{III.4.1.8.1.3 Order Outcome Notification}

\label{element:1d961373ab484e5f887ecf4c7b166130}
\index{Order Outcome Notification (III.4.1.8.1.3)}
{\small BusinessService · Canonical · current}\par

Emits completion notices to ordering practitioners.


{\small Recorded unresolved matters: source-qualified matters; exact scope remains in the canonical record.}\par
{\footnotesize \href{../../../../continuum/business/business-service/1d961373ab484e5f887ecf4c7b166130.md}{Canonical record: 1d961373ab484e5f887ecf4c7b166130}; full qualifications and provenance.}\par
\paragraph*{III.4.1.8.1.4 Order Requisition Ingress}
\phantomsection
\label{navigation:III.4.1.8.1.4}
\addcontentsline{toc}{paragraph}{III.4.1.8.1.4 Order Requisition Ingress}

\label{element:ba824a1ec2c64d9598c409803bcf9d86}
\index{Order Requisition Ingress (III.4.1.8.1.4)}
{\small BusinessService · Canonical · current}\par

Ingests electronic orders from clinical ordering systems.


{\footnotesize \href{../../../../continuum/business/business-service/ba824a1ec2c64d9598c409803bcf9d86.md}{Canonical record: ba824a1ec2c64d9598c409803bcf9d86}; full qualifications and provenance.}\par
\paragraph*{III.4.1.8.1.5 Order Status \& Tracking Query}
\phantomsection
\label{navigation:III.4.1.8.1.5}
\addcontentsline{toc}{paragraph}{III.4.1.8.1.5 Order Status \& Tracking Query}

\label{element:a4c049391dd64cd7af83bb5b4c7af0c5}
\index{Order Status \& Tracking Query (III.4.1.8.1.5)}
{\small BusinessService · Canonical · current}\par

Provides real-time order tracking to ordering clinicians.


{\footnotesize \href{../../../../continuum/business/business-service/a4c049391dd64cd7af83bb5b4c7af0c5.md}{Canonical record: a4c049391dd64cd7af83bb5b4c7af0c5}; full qualifications and provenance.}\par
\subsubsection*{III.4.1.8.2 Business Functions}
\phantomsection
\label{navigation:III.4.1.8.2}
\addcontentsline{toc}{subsubsection}{III.4.1.8.2 Business Functions}

\paragraph*{III.4.1.8.2.1 Associate Order Outcome}
\phantomsection
\label{navigation:III.4.1.8.2.1}
\addcontentsline{toc}{paragraph}{III.4.1.8.2.1 Associate Order Outcome}

\label{element:8bc2b528ca354c8ba95ddb44facdcd05}
\index{Associate Order Outcome (III.4.1.8.2.1)}
{\small BusinessFunction · Canonical · current}\par

Correlates returning diagnostic reports and observations back to the originating order requisition.


{\small Recorded unresolved matters: source-qualified matters; exact scope remains in the canonical record.}\par
{\footnotesize \href{../../../../continuum/business/business-function/8bc2b528ca354c8ba95ddb44facdcd05.md}{Canonical record: 8bc2b528ca354c8ba95ddb44facdcd05}; full qualifications and provenance.}\par
\paragraph*{III.4.1.8.2.2 Coordinate Order Modification / Cancellation}
\phantomsection
\label{navigation:III.4.1.8.2.2}
\addcontentsline{toc}{paragraph}{III.4.1.8.2.2 Coordinate Order Modification / Cancellation}

\label{element:6e445f8bc64e427fbe76b6dda20ed80f}
\index{Coordinate Order Modification / Cancellation (III.4.1.8.2.2)}
{\small BusinessFunction · Canonical · current}\par

Handles requests to amend clinical details or cancel unexecuted orders.


{\footnotesize \href{../../../../continuum/business/business-function/6e445f8bc64e427fbe76b6dda20ed80f.md}{Canonical record: 6e445f8bc64e427fbe76b6dda20ed80f}; full qualifications and provenance.}\par
\paragraph*{III.4.1.8.2.3 Manage Order Progression}
\phantomsection
\label{navigation:III.4.1.8.2.3}
\addcontentsline{toc}{paragraph}{III.4.1.8.2.3 Manage Order Progression}

\label{element:52f4b2dfb0c64e5aa4b378a9e63e0867}
\index{Manage Order Progression (III.4.1.8.2.3)}
{\small BusinessFunction · Canonical · current}\par

Tracks order state (scoped diagnostic/procedural illustration): Placed → Received → Specimen Collected → In-Progress → Preliminary Result → Final Result → Closed.


{\footnotesize \href{../../../../continuum/business/business-function/52f4b2dfb0c64e5aa4b378a9e63e0867.md}{Canonical record: 52f4b2dfb0c64e5aa4b378a9e63e0867}; full qualifications and provenance.}\par
\paragraph*{III.4.1.8.2.4 Receive Order Request}
\phantomsection
\label{navigation:III.4.1.8.2.4}
\addcontentsline{toc}{paragraph}{III.4.1.8.2.4 Receive Order Request}

\label{element:e5387fca59324407b8ed5d6db66ef93b}
\index{Receive Order Request (III.4.1.8.2.4)}
{\small BusinessFunction · Canonical · current}\par

Ingests structured orders containing clinical indication, requesting clinician, and test specifications.


{\footnotesize \href{../../../../continuum/business/business-function/e5387fca59324407b8ed5d6db66ef93b.md}{Canonical record: e5387fca59324407b8ed5d6db66ef93b}; full qualifications and provenance.}\par
\paragraph*{III.4.1.8.2.5 Resolve Order Destination}
\phantomsection
\label{navigation:III.4.1.8.2.5}
\addcontentsline{toc}{paragraph}{III.4.1.8.2.5 Resolve Order Destination}

\label{element:5e21ed48205446a9b05bc8017d1eccfc}
\index{Resolve Order Destination (III.4.1.8.2.5)}
{\small BusinessFunction · Canonical · current}\par

Evaluates order routing rules to determine target laboratory, imaging centre, or procedural unit.


{\footnotesize \href{../../../../continuum/business/business-function/5e21ed48205446a9b05bc8017d1eccfc.md}{Canonical record: 5e21ed48205446a9b05bc8017d1eccfc}; full qualifications and provenance.}\par
\subsubsection*{III.4.1.8.3 Business Processes}
\phantomsection
\label{navigation:III.4.1.8.3}
\addcontentsline{toc}{subsubsection}{III.4.1.8.3 Business Processes}

\paragraph*{III.4.1.8.3.1 Closed-Loop Order Progression Process}
\phantomsection
\label{navigation:III.4.1.8.3.1}
\addcontentsline{toc}{paragraph}{III.4.1.8.3.1 Closed-Loop Order Progression Process}

\label{element:e19a7c9594a4401e92450e4c31598982}
\index{Closed-Loop Order Progression Process (III.4.1.8.3.1)}
{\small BusinessProcess · Canonical · current}\par

•
Owning Capability: Order Administration



•
Process Purpose: Governs the rigorous, closed-loop tracking of diagnostic pathology, radiology, and procedural orders from requisition to result correlation, preventing dropped or unfulfilled investigations.



•
Scope qualification: These examples describe a diagnostic/procedural progression and do not exclude medication from the broader Order Administration responsibility. Specimen and result checkpoints are scoped. The signed requisition/report requirements remain within this illustrated scope; signing is not universally equivalent to finalisation or authority. Result consumption/binding does not define originating validity, authority, authorship or legal status.



•
State Progression Lifecycle: Requisition Placed → Order Dispatched → Specimen Collected / Scheduled → In-Execution → Preliminary Result Bound → Final Result Bound → Closed / Verified   (Scoped diagnostic/procedural illustration)



•
Key State Dispositions:



•
Requisition Placed: Electronic order created and signed by requesting clinician.



•
Order Dispatched: Order routed and accepted by performing diagnostic service. This local compound checkpoint retains routing and recipient acceptance as distinct facts; dispatch alone establishes neither receipt nor acceptance. Behaviour Received is not equated with this checkpoint without evidence.



•
Specimen Collected / Scheduled: Bio-specimen collected or imaging appointment booked.



•
In-Execution: Laboratory analysis or diagnostic scanning underway.



•
Preliminary Result Bound: Critical or interim observations published and linked to order.



•
Final Result Bound: Authoritative diagnostic report signed and bound to order.



•
Closed / Verified: Requesting clinician acknowledges result; order closed. This Order-closure checkpoint does not establish diagnostic-content verification, approval or clinical incorporation.



Closed-loop progression preserves the connection between the requisition, performing context and returned result so that a request does not disappear after dispatch. The diagnostic/procedural illustration includes specimen and report checkpoints without excluding medication from the broader owning capability. Dispatch, recipient acceptance, result binding and requester closure remain distinct facts even where a source checkpoint combines them.


{\small Recorded unresolved matters: source-qualified matters; exact scope remains in the canonical record.}\par
{\footnotesize \href{../../../../continuum/business/business-process/e19a7c9594a4401e92450e4c31598982.md}{Canonical record: e19a7c9594a4401e92450e4c31598982}; full qualifications and provenance.}\par
\subsection*{III.4.1.9 [CT2] Procedure Administration — Capability context}
\phantomsection
\label{navigation:III.4.1.9}
\addcontentsline{toc}{subsection}{III.4.1.9 [CT2] Procedure Administration — Capability context}

{\small Canonical Strategy Capability: II.1.3.1.15 \hyperref[element:649dd7cc26094e539386a050086dfb91]{[CT2] Procedure Administration}; Canonical ID: \texttt{\seqsplit{649dd7cc26094e539386a050086dfb91}}. Canonical · current.}\par
Contextual reference only; no additional architectural instance or semantic relationship.\par
\subsubsection*{III.4.1.9.1 Business Services}
\phantomsection
\label{navigation:III.4.1.9.1}
\addcontentsline{toc}{subsubsection}{III.4.1.9.1 Business Services}

\paragraph*{III.4.1.9.1.1 Procedural Documentation Ingress}
\phantomsection
\label{navigation:III.4.1.9.1.1}
\addcontentsline{toc}{paragraph}{III.4.1.9.1.1 Procedural Documentation Ingress}

\label{element:b58405366b7a43dc909ca3074503e11a}
\index{Procedural Documentation Ingress (III.4.1.9.1.1)}
{\small BusinessService · Canonical · current}\par

Ingests procedural documentation, including surgical notes, into patient record streams.


{\footnotesize \href{../../../../continuum/business/business-service/b58405366b7a43dc909ca3074503e11a.md}{Canonical record: b58405366b7a43dc909ca3074503e11a}; full qualifications and provenance.}\par
\paragraph*{III.4.1.9.1.2 Procedure Booking Ingress}
\phantomsection
\label{navigation:III.4.1.9.1.2}
\addcontentsline{toc}{paragraph}{III.4.1.9.1.2 Procedure Booking Ingress}

\label{element:64c52e8e6f304dee8d86a5da96f927e3}
\index{Procedure Booking Ingress (III.4.1.9.1.2)}
{\small BusinessService · Canonical · current}\par

Receives surgical/procedural requests.


{\footnotesize \href{../../../../continuum/business/business-service/64c52e8e6f304dee8d86a5da96f927e3.md}{Canonical record: 64c52e8e6f304dee8d86a5da96f927e3}; full qualifications and provenance.}\par
\subsubsection*{III.4.1.9.2 Business Functions}
\phantomsection
\label{navigation:III.4.1.9.2}
\addcontentsline{toc}{subsubsection}{III.4.1.9.2 Business Functions}

\paragraph*{III.4.1.9.2.1 Receive Procedural Documentation}
\phantomsection
\label{navigation:III.4.1.9.2.1}
\addcontentsline{toc}{paragraph}{III.4.1.9.2.1 Receive Procedural Documentation}

\label{element:62a47d22a07e4fcc8b3f99ccdffff390}
\index{Receive Procedural Documentation (III.4.1.9.2.1)}
{\small BusinessFunction · Canonical · current}\par

Ingests operation reports, anaesthetic records, and post-procedure notes.


{\footnotesize \href{../../../../continuum/business/business-function/62a47d22a07e4fcc8b3f99ccdffff390.md}{Canonical record: 62a47d22a07e4fcc8b3f99ccdffff390}; full qualifications and provenance.}\par
\paragraph*{III.4.1.9.2.2 Receive Procedure Booking / Request}
\phantomsection
\label{navigation:III.4.1.9.2.2}
\addcontentsline{toc}{paragraph}{III.4.1.9.2.2 Receive Procedure Booking / Request}

\label{element:6d29d7e7422443b4acf3cbe85c4074e2}
\index{Receive Procedure Booking / Request (III.4.1.9.2.2)}
{\small BusinessFunction · Canonical · current}\par

Ingests procedural and surgical requisitions.


{\footnotesize \href{../../../../continuum/business/business-function/6d29d7e7422443b4acf3cbe85c4074e2.md}{Canonical record: 6d29d7e7422443b4acf3cbe85c4074e2}; full qualifications and provenance.}\par
\subsection*{III.4.1.10 [CT2] Referral Administration — Capability context}
\phantomsection
\label{navigation:III.4.1.10}
\addcontentsline{toc}{subsection}{III.4.1.10 [CT2] Referral Administration — Capability context}

{\small Canonical Strategy Capability: II.1.3.1.16 \hyperref[element:ae0de25600ee416aa24d2a167b64f230]{[CT2] Referral Administration}; Canonical ID: \texttt{\seqsplit{ae0de25600ee416aa24d2a167b64f230}}. Canonical · current.}\par
Contextual reference only; no additional architectural instance or semantic relationship.\par
\subsubsection*{III.4.1.10.1 Business Services}
\phantomsection
\label{navigation:III.4.1.10.1}
\addcontentsline{toc}{subsubsection}{III.4.1.10.1 Business Services}

\paragraph*{III.4.1.10.1.1 Referral Context Query}
\phantomsection
\label{navigation:III.4.1.10.1.1}
\addcontentsline{toc}{paragraph}{III.4.1.10.1.1 Referral Context Query}

\label{element:d688ee73f68a4dda9380404e22b3539b}
\index{Referral Context Query (III.4.1.10.1.1)}
{\small BusinessService · Canonical · current}\par

Provides comprehensive referral dossiers to intake triage clinicians.


{\footnotesize \href{../../../../continuum/business/business-service/d688ee73f68a4dda9380404e22b3539b.md}{Canonical record: d688ee73f68a4dda9380404e22b3539b}; full qualifications and provenance.}\par
\paragraph*{III.4.1.10.1.2 Referral Status \& Outcome Service}
\phantomsection
\label{navigation:III.4.1.10.1.2}
\addcontentsline{toc}{paragraph}{III.4.1.10.1.2 Referral Status \& Outcome Service}

\label{element:dfd0481be6124da0a791db41a5e9c238}
\index{Referral Status \& Outcome Service (III.4.1.10.1.2)}
{\small BusinessService · Canonical · current}\par

Exposes progression milestones to referrers and patients.


{\footnotesize \href{../../../../continuum/business/business-service/dfd0481be6124da0a791db41a5e9c238.md}{Canonical record: dfd0481be6124da0a791db41a5e9c238}; full qualifications and provenance.}\par
\paragraph*{III.4.1.10.1.3 Referral Submission Service}
\phantomsection
\label{navigation:III.4.1.10.1.3}
\addcontentsline{toc}{paragraph}{III.4.1.10.1.3 Referral Submission Service}

\label{element:42b4e9da22764480a0bea7019927860c}
\index{Referral Submission Service (III.4.1.10.1.3)}
{\small BusinessService · Canonical · current}\par

Exposed to external GP practices and referring health facilities.


{\footnotesize \href{../../../../continuum/business/business-service/42b4e9da22764480a0bea7019927860c.md}{Canonical record: 42b4e9da22764480a0bea7019927860c}; full qualifications and provenance.}\par
\subsubsection*{III.4.1.10.2 Business Functions}
\phantomsection
\label{navigation:III.4.1.10.2}
\addcontentsline{toc}{subsubsection}{III.4.1.10.2 Business Functions}

\paragraph*{III.4.1.10.2.1 Assemble Referral Context}
\phantomsection
\label{navigation:III.4.1.10.2.1}
\addcontentsline{toc}{paragraph}{III.4.1.10.2.1 Assemble Referral Context}

\label{element:b4a05bf97e844c24b04d27a6b325627a}
\index{Assemble Referral Context (III.4.1.10.2.1)}
{\small BusinessFunction · Canonical · current}\par

Gathers relevant clinical history, medications, and diagnostic investigations into the referral package.


{\footnotesize \href{../../../../continuum/business/business-function/b4a05bf97e844c24b04d27a6b325627a.md}{Canonical record: b4a05bf97e844c24b04d27a6b325627a}; full qualifications and provenance.}\par
\paragraph*{III.4.1.10.2.2 Manage Referral Progression}
\phantomsection
\label{navigation:III.4.1.10.2.2}
\addcontentsline{toc}{paragraph}{III.4.1.10.2.2 Manage Referral Progression}

\label{element:7f9f01be4fd14facb943de5f63f1d8fd}
\index{Manage Referral Progression (III.4.1.10.2.2)}
{\small BusinessFunction · Canonical · current}\par

Tracks referral disposition in an illustrative specialist pathway; dispositions vary: Submitted → Triaged → Accepted/Waitlisted → Scheduled → Discharged.


{\footnotesize \href{../../../../continuum/business/business-function/7f9f01be4fd14facb943de5f63f1d8fd.md}{Canonical record: 7f9f01be4fd14facb943de5f63f1d8fd}; full qualifications and provenance.}\par
\paragraph*{III.4.1.10.2.3 Receive Referral}
\phantomsection
\label{navigation:III.4.1.10.2.3}
\addcontentsline{toc}{paragraph}{III.4.1.10.2.3 Receive Referral}

\label{element:28775f5af15646369539a4fc0a7ff8a9}
\index{Receive Referral (III.4.1.10.2.3)}
{\small BusinessFunction · Canonical · current}\par

Validates and ingests structured electronic referrals, clinical indications, and priority categories.


{\footnotesize \href{../../../../continuum/business/business-function/28775f5af15646369539a4fc0a7ff8a9.md}{Canonical record: 28775f5af15646369539a4fc0a7ff8a9}; full qualifications and provenance.}\par
\subsubsection*{III.4.1.10.3 Business Processes}
\phantomsection
\label{navigation:III.4.1.10.3}
\addcontentsline{toc}{subsubsection}{III.4.1.10.3 Business Processes}

\paragraph*{III.4.1.10.3.1 Referral Progression Process}
\phantomsection
\label{navigation:III.4.1.10.3.1}
\addcontentsline{toc}{paragraph}{III.4.1.10.3.1 Referral Progression Process}

\label{element:cfb714e9320647ffa1dffc80dddd01d9}
\index{Referral Progression Process (III.4.1.10.3.1)}
{\small BusinessProcess · Canonical · current}\par

•
Owning Capability: Referral Administration



•
Process Purpose: Governs the end-to-end operational progression of an incoming clinical referral from receipt through specialist clinical triage, booking, and final service acceptance.



•
Scope qualification: This illustrates a specialist booking/attendance pathway; Referral does not universally require acceptance, scheduling, attendance or discharge, or establish delivery, responsibility actually assumed or Transfer of Care. Decline, rejection and redirection remain valid dispositions without invented transition paths. Accepted / Waitlisted is a local compound checkpoint, not universal equivalence between acceptance and waitlisting.



•
State Progression Lifecycle: Submitted → Intake Validated → Clinically Triaged → Accepted / Waitlisted → Scheduled → Consultation Attended → Discharged / Rejected   (Illustrative specialist pathway; dispositions vary)



•
Key State Dispositions:



•
Submitted: Electronic referral received from GP or external facility.



•
Intake Validated: Administrative validation of patient details, mandatory fields, and tests.



•
Clinically Triaged: Senior clinician assigns urgency category (e.g., Cat 1 urgent within 30 days).



•
Accepted / Waitlisted: Referral accepted and placed onto the specialty waitlist.



•
Scheduled: Outpatient appointment allocated and patient notified.



•
Consultation Attended: Patient seen; initial specialist assessment completed.



•
Discharged / Rejected: Referral completed and discharged back to primary care, or formally rejected with clinical rationale.



This Process illustrates one specialist pathway through validation, triage and booking. It makes that pathway's progression visible without imposing it on every referral: rejection, decline or redirection can be legitimate dispositions, and a request need not universally lead to attendance. Neither a transmitted referral nor the local compound acceptance/waitlist checkpoint establishes actual responsibility transfer.


{\small Recorded unresolved matters: source-qualified matters; exact scope remains in the canonical record.}\par
{\footnotesize \href{../../../../continuum/business/business-process/cfb714e9320647ffa1dffc80dddd01d9.md}{Canonical record: cfb714e9320647ffa1dffc80dddd01d9}; full qualifications and provenance.}\par
\subsection*{III.4.1.11 [CT2] Scheduling Administration — Capability context}
\phantomsection
\label{navigation:III.4.1.11}
\addcontentsline{toc}{subsection}{III.4.1.11 [CT2] Scheduling Administration — Capability context}

{\small Canonical Strategy Capability: II.1.3.1.17 \hyperref[element:35fb3744c900470f9895056246a13fb1]{[CT2] Scheduling Administration}; Canonical ID: \texttt{\seqsplit{35fb3744c900470f9895056246a13fb1}}. Canonical · current.}\par
Contextual reference only; no additional architectural instance or semantic relationship.\par
\subsubsection*{III.4.1.11.1 Business Services}
\phantomsection
\label{navigation:III.4.1.11.1}
\addcontentsline{toc}{subsubsection}{III.4.1.11.1 Business Services}

\paragraph*{III.4.1.11.1.1 Appointment Notification Ingress}
\phantomsection
\label{navigation:III.4.1.11.1.1}
\addcontentsline{toc}{paragraph}{III.4.1.11.1.1 Appointment Notification Ingress}

\label{element:67045c71e2194f6ba19002edcea94ee0}
\index{Appointment Notification Ingress (III.4.1.11.1.1)}
{\small BusinessService · Canonical · current}\par

Ingests booking events from departmental schedulers.


{\footnotesize \href{../../../../continuum/business/business-service/67045c71e2194f6ba19002edcea94ee0.md}{Canonical record: 67045c71e2194f6ba19002edcea94ee0}; full qualifications and provenance.}\par
\paragraph*{III.4.1.11.1.2 Appointment Schedule Query}
\phantomsection
\label{navigation:III.4.1.11.1.2}
\addcontentsline{toc}{paragraph}{III.4.1.11.1.2 Appointment Schedule Query}

\label{element:2a7eb5b22dc04256aa049c2fcf84e62c}
\index{Appointment Schedule Query (III.4.1.11.1.2)}
{\small BusinessService · Canonical · current}\par

Exposes consolidated patient appointment schedules across facilities.


{\footnotesize \href{../../../../continuum/business/business-service/2a7eb5b22dc04256aa049c2fcf84e62c.md}{Canonical record: 2a7eb5b22dc04256aa049c2fcf84e62c}; full qualifications and provenance.}\par
\subsubsection*{III.4.1.11.2 Business Functions}
\phantomsection
\label{navigation:III.4.1.11.2}
\addcontentsline{toc}{subsubsection}{III.4.1.11.2 Business Functions}

\paragraph*{III.4.1.11.2.1 Receive Appointment Notification}
\phantomsection
\label{navigation:III.4.1.11.2.1}
\addcontentsline{toc}{paragraph}{III.4.1.11.2.1 Receive Appointment Notification}

\label{element:8aa2e7b197ae4a0c99b654b9f1a312e8}
\index{Receive Appointment Notification (III.4.1.11.2.1)}
{\small BusinessFunction · Canonical · current}\par

Ingests appointment creation, rescheduling, cancellation, and check-in events from host PAS engines.


{\footnotesize \href{../../../../continuum/business/business-function/8aa2e7b197ae4a0c99b654b9f1a312e8.md}{Canonical record: 8aa2e7b197ae4a0c99b654b9f1a312e8}; full qualifications and provenance.}\par
\paragraph*{III.4.1.11.2.2 Synchronise Appointment Status}
\phantomsection
\label{navigation:III.4.1.11.2.2}
\addcontentsline{toc}{paragraph}{III.4.1.11.2.2 Synchronise Appointment Status}

\label{element:c60fc8ce05104185b5240a910bfd5f67}
\index{Synchronise Appointment Status (III.4.1.11.2.2)}
{\small BusinessFunction · Canonical · current}\par

Projects synchronised appointment timelines to clinical dashboards and care coordinators.


{\footnotesize \href{../../../../continuum/business/business-function/c60fc8ce05104185b5240a910bfd5f67.md}{Canonical record: c60fc8ce05104185b5240a910bfd5f67}; full qualifications and provenance.}\par
\section*{III.4.2 [CT1] Entity Information Enablement — Capability context}
\phantomsection
\label{navigation:III.4.2}
\addcontentsline{toc}{section}{III.4.2 [CT1] Entity Information Enablement — Capability context}

{\small Canonical Strategy Capability: II.1.3.4 \hyperref[element:b5a92874f8ac4509a30af5253590edc1]{[CT1] Entity Information Enablement}; Canonical ID: \texttt{\seqsplit{b5a92874f8ac4509a30af5253590edc1}}. Proposed · current.}\par
Contextual reference only; no additional architectural instance or semantic relationship.\par
\subsection*{III.4.2.1 [CT2] Client Administration — Capability context}
\phantomsection
\label{navigation:III.4.2.1}
\addcontentsline{toc}{subsection}{III.4.2.1 [CT2] Client Administration — Capability context}

{\small Canonical Strategy Capability: II.1.3.4.1 \hyperref[element:f0885472cced4bccbcce1568bf0360dc]{[CT2] Client Administration}; Canonical ID: \texttt{\seqsplit{f0885472cced4bccbcce1568bf0360dc}}. Canonical · current.}\par
Contextual reference only; no additional architectural instance or semantic relationship.\par
\subsubsection*{III.4.2.1.1 [CT3] Client Privacy Guardrail Management — Capability context}
\phantomsection
\label{navigation:III.4.2.1.1}
\addcontentsline{toc}{subsubsection}{III.4.2.1.1 [CT3] Client Privacy Guardrail Management — Capability context}

{\small Canonical Strategy Capability: II.1.3.4.1.1 \hyperref[element:7577a6e1815243648dd5356c21cca151]{[CT3] Client Privacy Guardrail Management}; Canonical ID: \texttt{\seqsplit{7577a6e1815243648dd5356c21cca151}}. Canonical · current.}\par
Contextual reference only; no additional architectural instance or semantic relationship.\par
\paragraph*{III.4.2.1.1.1 Business Services}
\phantomsection
\label{navigation:III.4.2.1.1.1}
\addcontentsline{toc}{paragraph}{III.4.2.1.1.1 Business Services}

\subparagraph*{III.4.2.1.1.1.1 Confidentiality Restriction Evaluation}
\phantomsection
\label{navigation:III.4.2.1.1.1.1}
\addcontentsline{toc}{subparagraph}{III.4.2.1.1.1.1 Confidentiality Restriction Evaluation}

\label{element:87cfcfb20b7a4570b502b4a9a53803ad}
\index{Confidentiality Restriction Evaluation (III.4.2.1.1.1.1)}
{\small BusinessService · Canonical · current}\par

Enforces record-level masking directives across the HIE.


{\footnotesize \href{../../../../continuum/business/business-service/87cfcfb20b7a4570b502b4a9a53803ad.md}{Canonical record: 87cfcfb20b7a4570b502b4a9a53803ad}; full qualifications and provenance.}\par
\subparagraph*{III.4.2.1.1.1.2 Consent Decision Service}
\phantomsection
\label{navigation:III.4.2.1.1.1.2}
\addcontentsline{toc}{subparagraph}{III.4.2.1.1.1.2 Consent Decision Service}

\label{element:e1bebefbdbef4a1981ec31341eb71dfd}
\index{Consent Decision Service (III.4.2.1.1.1.2)}
{\small BusinessService · Canonical · current}\par

Provides real-time consent evaluation decisions to Health Information Control.


{\footnotesize \href{../../../../continuum/business/business-service/e1bebefbdbef4a1981ec31341eb71dfd.md}{Canonical record: e1bebefbdbef4a1981ec31341eb71dfd}; full qualifications and provenance.}\par
\paragraph*{III.4.2.1.1.2 Business Functions}
\phantomsection
\label{navigation:III.4.2.1.1.2}
\addcontentsline{toc}{paragraph}{III.4.2.1.1.2 Business Functions}

\subparagraph*{III.4.2.1.1.2.1 Evaluate Client Confidentiality Restrictions}
\phantomsection
\label{navigation:III.4.2.1.1.2.1}
\addcontentsline{toc}{subparagraph}{III.4.2.1.1.2.1 Evaluate Client Confidentiality Restrictions}

\label{element:a61324c3ef734910af5886f048cd3bf5}
\index{Evaluate Client Confidentiality Restrictions (III.4.2.1.1.2.1)}
{\small BusinessFunction · Canonical · current}\par

Evaluates sensitive record flags (e.g., VIP masking, restricted department exclusions).


{\footnotesize \href{../../../../continuum/business/business-function/a61324c3ef734910af5886f048cd3bf5.md}{Canonical record: a61324c3ef734910af5886f048cd3bf5}; full qualifications and provenance.}\par
\subparagraph*{III.4.2.1.1.2.2 Evaluate Client Consent}
\phantomsection
\label{navigation:III.4.2.1.1.2.2}
\addcontentsline{toc}{subparagraph}{III.4.2.1.1.2.2 Evaluate Client Consent}

\label{element:b9e978a651cf4ff3bb690af50b8528cd}
\index{Evaluate Client Consent (III.4.2.1.1.2.2)}
{\small BusinessFunction · Canonical · current}\par

Evaluates patient opt-in, opt-out, and general information sharing preferences against proposed exchanges.


{\footnotesize \href{../../../../continuum/business/business-function/b9e978a651cf4ff3bb690af50b8528cd.md}{Canonical record: b9e978a651cf4ff3bb690af50b8528cd}; full qualifications and provenance.}\par
\subsubsection*{III.4.2.1.2 [CT3] Client Relationship Administration — Capability context}
\phantomsection
\label{navigation:III.4.2.1.2}
\addcontentsline{toc}{subsubsection}{III.4.2.1.2 [CT3] Client Relationship Administration — Capability context}

{\small Canonical Strategy Capability: II.1.3.4.1.2 \hyperref[element:d682c11d3eb04770b1f20be478fd475e]{[CT3] Client Relationship Administration}; Canonical ID: \texttt{\seqsplit{d682c11d3eb04770b1f20be478fd475e}}. Canonical · current.}\par
Contextual reference only; no additional architectural instance or semantic relationship.\par
\paragraph*{III.4.2.1.2.1 Business Services}
\phantomsection
\label{navigation:III.4.2.1.2.1}
\addcontentsline{toc}{paragraph}{III.4.2.1.2.1 Business Services}

\subparagraph*{III.4.2.1.2.1.1 Client Relationship Query}
\phantomsection
\label{navigation:III.4.2.1.2.1.1}
\addcontentsline{toc}{subparagraph}{III.4.2.1.2.1.1 Client Relationship Query}

\label{element:7b3eb2e91baa4e9c887927433979a8f6}
\index{Client Relationship Query (III.4.2.1.2.1.1)}
{\small BusinessService · Canonical · current}\par

Discloses verified support network members and contact details to clinical teams.


{\footnotesize \href{../../../../continuum/business/business-service/7b3eb2e91baa4e9c887927433979a8f6.md}{Canonical record: 7b3eb2e91baa4e9c887927433979a8f6}; full qualifications and provenance.}\par
\subparagraph*{III.4.2.1.2.1.2 Representative Authority Verification}
\phantomsection
\label{navigation:III.4.2.1.2.1.2}
\addcontentsline{toc}{subparagraph}{III.4.2.1.2.1.2 Representative Authority Verification}

\label{element:2d3205c785ba4db5af5868ebb6ab1495}
\index{Representative Authority Verification (III.4.2.1.2.1.2)}
{\small BusinessService · Canonical · current}\par

Confirms legal decision-making authority for substitute consent.


{\footnotesize \href{../../../../continuum/business/business-service/2d3205c785ba4db5af5868ebb6ab1495.md}{Canonical record: 2d3205c785ba4db5af5868ebb6ab1495}; full qualifications and provenance.}\par
\paragraph*{III.4.2.1.2.2 Business Functions}
\phantomsection
\label{navigation:III.4.2.1.2.2}
\addcontentsline{toc}{paragraph}{III.4.2.1.2.2 Business Functions}

\subparagraph*{III.4.2.1.2.2.1 Govern Client Relationship Validity}
\phantomsection
\label{navigation:III.4.2.1.2.2.1}
\addcontentsline{toc}{subparagraph}{III.4.2.1.2.2.1 Govern Client Relationship Validity}

\label{element:dc87537df6254fc5bdeeca0ccc011679}
\index{Govern Client Relationship Validity (III.4.2.1.2.2.1)}
{\small BusinessFunction · Canonical · current}\par

Verifies power-of-attorney documents, guardianship orders, and proxy validity periods.


{\footnotesize \href{../../../../continuum/business/business-function/dc87537df6254fc5bdeeca0ccc011679.md}{Canonical record: dc87537df6254fc5bdeeca0ccc011679}; full qualifications and provenance.}\par
\subparagraph*{III.4.2.1.2.2.2 Maintain Client Relationships}
\phantomsection
\label{navigation:III.4.2.1.2.2.2}
\addcontentsline{toc}{subparagraph}{III.4.2.1.2.2.2 Maintain Client Relationships}

\label{element:1b2ce4a240384672b78976904b83a011}
\index{Maintain Client Relationships (III.4.2.1.2.2.2)}
{\small BusinessFunction · Canonical · current}\par

Records links between patients and carers, next of kin, nominated representatives, and legal guardians.


{\footnotesize \href{../../../../continuum/business/business-function/1b2ce4a240384672b78976904b83a011.md}{Canonical record: 1b2ce4a240384672b78976904b83a011}; full qualifications and provenance.}\par
\subsubsection*{III.4.2.1.3 [CT3] Healthcare Subject Profile Administration — Capability context}
\phantomsection
\label{navigation:III.4.2.1.3}
\addcontentsline{toc}{subsubsection}{III.4.2.1.3 [CT3] Healthcare Subject Profile Administration — Capability context}

{\small Canonical Strategy Capability: II.1.3.4.1.3 \hyperref[element:df214b5499694b768eb4a7a175a9d471]{[CT3] Healthcare Subject Profile Administration}; Canonical ID: \texttt{\seqsplit{df214b5499694b768eb4a7a175a9d471}}. Canonical · current.}\par
Contextual reference only; no additional architectural instance or semantic relationship.\par
\paragraph*{III.4.2.1.3.1 Business Services}
\phantomsection
\label{navigation:III.4.2.1.3.1}
\addcontentsline{toc}{paragraph}{III.4.2.1.3.1 Business Services}

\subparagraph*{III.4.2.1.3.1.1 Demographic Change Distribution}
\phantomsection
\label{navigation:III.4.2.1.3.1.1}
\addcontentsline{toc}{subparagraph}{III.4.2.1.3.1.1 Demographic Change Distribution}

\label{element:2c4a8c1251294885922f8ccdc3d7b069}
\index{Demographic Change Distribution (III.4.2.1.3.1.1)}
{\small BusinessService · Canonical · current}\par

Emits governed subject demographic updates to subscribed clinical systems.


{\footnotesize \href{../../../../continuum/business/business-service/2c4a8c1251294885922f8ccdc3d7b069.md}{Canonical record: 2c4a8c1251294885922f8ccdc3d7b069}; full qualifications and provenance.}\par
\subparagraph*{III.4.2.1.3.1.2 Healthcare Subject Context Resolution}
\phantomsection
\label{navigation:III.4.2.1.3.1.2}
\addcontentsline{toc}{subparagraph}{III.4.2.1.3.1.2 Healthcare Subject Context Resolution}

\label{element:99a739efc7be4d7b92d90181ecd69d15}
\index{Healthcare Subject Context Resolution (III.4.2.1.3.1.2)}
{\small BusinessService · Canonical · current}\par

Provides comprehensive demographic, language, and indigenous status context to clinical applications.


{\footnotesize \href{../../../../continuum/business/business-service/99a739efc7be4d7b92d90181ecd69d15.md}{Canonical record: 99a739efc7be4d7b92d90181ecd69d15}; full qualifications and provenance.}\par
\paragraph*{III.4.2.1.3.2 Business Functions}
\phantomsection
\label{navigation:III.4.2.1.3.2}
\addcontentsline{toc}{paragraph}{III.4.2.1.3.2 Business Functions}

\subparagraph*{III.4.2.1.3.2.1 Establish Healthcare Subject Context}
\phantomsection
\label{navigation:III.4.2.1.3.2.1}
\addcontentsline{toc}{subparagraph}{III.4.2.1.3.2.1 Establish Healthcare Subject Context}

\label{element:f445f86f9de24461aabee527e5b41287}
\index{Establish Healthcare Subject Context (III.4.2.1.3.2.1)}
{\small BusinessFunction · Canonical · current}\par

Binds demographic, cultural, and communication requirements to a patient record.


{\footnotesize \href{../../../../continuum/business/business-function/f445f86f9de24461aabee527e5b41287.md}{Canonical record: f445f86f9de24461aabee527e5b41287}; full qualifications and provenance.}\par
\subparagraph*{III.4.2.1.3.2.2 Govern Subject Demographic Context}
\phantomsection
\label{navigation:III.4.2.1.3.2.2}
\addcontentsline{toc}{subparagraph}{III.4.2.1.3.2.2 Govern Subject Demographic Context}

\label{element:76288dacb9034e0bbf75b43d8afe59d2}
\index{Govern Subject Demographic Context (III.4.2.1.3.2.2)}
{\small BusinessFunction · Canonical · current}\par

Applies validation rules to address changes, date of birth verification, and death notifications.


{\footnotesize \href{../../../../continuum/business/business-function/76288dacb9034e0bbf75b43d8afe59d2.md}{Canonical record: 76288dacb9034e0bbf75b43d8afe59d2}; full qualifications and provenance.}\par
\subsubsection*{III.4.2.1.4 [CT3] Person Identity Administration — Capability context}
\phantomsection
\label{navigation:III.4.2.1.4}
\addcontentsline{toc}{subsubsection}{III.4.2.1.4 [CT3] Person Identity Administration — Capability context}

{\small Canonical Strategy Capability: II.1.3.4.1.4 \hyperref[element:6c3de7750ef1436cad6a9533731782b6]{[CT3] Person Identity Administration}; Canonical ID: \texttt{\seqsplit{6c3de7750ef1436cad6a9533731782b6}}. Canonical · current.}\par
Contextual reference only; no additional architectural instance or semantic relationship.\par
\paragraph*{III.4.2.1.4.1 [Feature] Identifier Capture — Capability context}
\phantomsection
\label{navigation:III.4.2.1.4.1}
\addcontentsline{toc}{paragraph}{III.4.2.1.4.1 [Feature] Identifier Capture — Capability context}

{\small Canonical Strategy Capability: II.1.3.4.1.4.3 \hyperref[element:3cae3c4810974c5f8ee5a4c942f38b07]{[Feature] Identifier Capture}; Canonical ID: \texttt{\seqsplit{3cae3c4810974c5f8ee5a4c942f38b07}}. Accepted bounded reconciliation · current.}\par
Contextual reference only; no additional architectural instance or semantic relationship.\par
\subparagraph*{III.4.2.1.4.1.1 Business Services}
\phantomsection
\label{navigation:III.4.2.1.4.1.1}
\addcontentsline{toc}{subparagraph}{III.4.2.1.4.1.1 Business Services}

\subsubparagraph*{III.4.2.1.4.1.1.1 [Business Service] Identifier Create}
\phantomsection
\label{navigation:III.4.2.1.4.1.1.1}
\addcontentsline{toc}{subsubparagraph}{III.4.2.1.4.1.1.1 [Business Service] Identifier Create}

\label{element:063b7a6b4ad549428a7566a4ce2c5de5}
\index{[Business Service] Identifier Create (III.4.2.1.4.1.1.1)}
{\small BusinessService · Accepted bounded reconciliation · current}\par

The governed source has no established definition. This absence is retained for Solution Architect review; no definition is inferred from the name.


{\small Recorded unresolved matters: definition; exact scope remains in the canonical record.}\par
{\footnotesize \href{../../../../continuum/business/business-service/063b7a6b4ad549428a7566a4ce2c5de5.md}{Canonical record: 063b7a6b4ad549428a7566a4ce2c5de5}; full qualifications and provenance.}\par
\subsubparagraph*{III.4.2.1.4.1.1.2 [Business Service] Identifier Update}
\phantomsection
\label{navigation:III.4.2.1.4.1.1.2}
\addcontentsline{toc}{subsubparagraph}{III.4.2.1.4.1.1.2 [Business Service] Identifier Update}

\label{element:281155a0820942b884a13fa08d2e0b04}
\index{[Business Service] Identifier Update (III.4.2.1.4.1.1.2)}
{\small BusinessService · Accepted bounded reconciliation · current}\par

The governed source has no established definition. This absence is retained for Solution Architect review; no definition is inferred from the name.


{\small Recorded unresolved matters: definition; exact scope remains in the canonical record.}\par
{\footnotesize \href{../../../../continuum/business/business-service/281155a0820942b884a13fa08d2e0b04.md}{Canonical record: 281155a0820942b884a13fa08d2e0b04}; full qualifications and provenance.}\par
\subparagraph*{III.4.2.1.4.1.2 Business Processes}
\phantomsection
\label{navigation:III.4.2.1.4.1.2}
\addcontentsline{toc}{subparagraph}{III.4.2.1.4.1.2 Business Processes}

\subsubparagraph*{III.4.2.1.4.1.2.1 [Business Process] Identifier Verification}
\phantomsection
\label{navigation:III.4.2.1.4.1.2.1}
\addcontentsline{toc}{subsubparagraph}{III.4.2.1.4.1.2.1 [Business Process] Identifier Verification}

\label{element:6ea7f8afe17d4319ab6fa9691db4328d}
\index{[Business Process] Identifier Verification (III.4.2.1.4.1.2.1)}
{\small BusinessProcess · Accepted bounded reconciliation · current}\par

The governed source has no established definition. This absence is retained for Solution Architect review; no definition is inferred from the name.


{\small Recorded unresolved matters: definition; exact scope remains in the canonical record.}\par
{\footnotesize \href{../../../../continuum/business/business-process/6ea7f8afe17d4319ab6fa9691db4328d.md}{Canonical record: 6ea7f8afe17d4319ab6fa9691db4328d}; full qualifications and provenance.}\par
\paragraph*{III.4.2.1.4.2 [Feature] Identity Source \& Correction Lineage — Capability context}
\phantomsection
\label{navigation:III.4.2.1.4.2}
\addcontentsline{toc}{paragraph}{III.4.2.1.4.2 [Feature] Identity Source \& Correction Lineage — Capability context}

{\small Canonical Strategy Capability: II.1.3.4.1.4.6 \hyperref[element:d72134e75001449a9e7ffae53f61458f]{[Feature] Identity Source \& Correction Lineage}; Canonical ID: \texttt{\seqsplit{d72134e75001449a9e7ffae53f61458f}}. Proposed · current.}\par
Contextual reference only; no additional architectural instance or semantic relationship.\par
\subparagraph*{III.4.2.1.4.2.1 Business Services}
\phantomsection
\label{navigation:III.4.2.1.4.2.1}
\addcontentsline{toc}{subparagraph}{III.4.2.1.4.2.1 Business Services}

\subsubparagraph*{III.4.2.1.4.2.1.1 [Business Service] Verify Identity Source}
\phantomsection
\label{navigation:III.4.2.1.4.2.1.1}
\addcontentsline{toc}{subsubparagraph}{III.4.2.1.4.2.1.1 [Business Service] Verify Identity Source}

\label{element:b95f652107f4420b86cc10fb893c9b01}
\index{[Business Service] Verify Identity Source (III.4.2.1.4.2.1.1)}
{\small BusinessService · Accepted bounded reconciliation · current}\par

The governed source has no established definition. This absence is retained for Solution Architect review; no definition is inferred from the name.


{\small Recorded unresolved matters: definition; exact scope remains in the canonical record.}\par
{\footnotesize \href{../../../../continuum/business/business-service/b95f652107f4420b86cc10fb893c9b01.md}{Canonical record: b95f652107f4420b86cc10fb893c9b01}; full qualifications and provenance.}\par
\paragraph*{III.4.2.1.4.3 Business Services}
\phantomsection
\label{navigation:III.4.2.1.4.3}
\addcontentsline{toc}{paragraph}{III.4.2.1.4.3 Business Services}

\subparagraph*{III.4.2.1.4.3.1 Person Identifier Correlation}
\phantomsection
\label{navigation:III.4.2.1.4.3.1}
\addcontentsline{toc}{subparagraph}{III.4.2.1.4.3.1 Person Identifier Correlation}

\label{element:71220901361e42f686faf7d852f57e45}
\index{Person Identifier Correlation (III.4.2.1.4.3.1)}
{\small BusinessService · Canonical · current}\par

Exposed across the HIE to establish patient cross-reference mappings.


{\footnotesize \href{../../../../continuum/business/business-service/71220901361e42f686faf7d852f57e45.md}{Canonical record: 71220901361e42f686faf7d852f57e45}; full qualifications and provenance.}\par
\subparagraph*{III.4.2.1.4.3.2 Person Identifier Resolution}
\phantomsection
\label{navigation:III.4.2.1.4.3.2}
\addcontentsline{toc}{subparagraph}{III.4.2.1.4.3.2 Person Identifier Resolution}

\label{element:a533cf0d23364f398f2a71843cce98c4}
\index{Person Identifier Resolution (III.4.2.1.4.3.2)}
{\small BusinessService · Canonical · current}\par

Exposed to Service Administration and Service Delivery to disambiguate client identities.


{\footnotesize \href{../../../../continuum/business/business-service/a533cf0d23364f398f2a71843cce98c4.md}{Canonical record: a533cf0d23364f398f2a71843cce98c4}; full qualifications and provenance.}\par
\subparagraph*{III.4.2.1.4.3.3 Person Identity Correction}
\phantomsection
\label{navigation:III.4.2.1.4.3.3}
\addcontentsline{toc}{subparagraph}{III.4.2.1.4.3.3 Person Identity Correction}

\label{element:dd5ed5ba58f94ecbbdac913c88d2b7b6}
\index{Person Identity Correction (III.4.2.1.4.3.3)}
{\small BusinessService · Canonical · current}\par

Exposed to authorized HIM/Medical Records administrators.


{\footnotesize \href{../../../../continuum/business/business-service/dd5ed5ba58f94ecbbdac913c88d2b7b6.md}{Canonical record: dd5ed5ba58f94ecbbdac913c88d2b7b6}; full qualifications and provenance.}\par
\paragraph*{III.4.2.1.4.4 Business Functions}
\phantomsection
\label{navigation:III.4.2.1.4.4}
\addcontentsline{toc}{paragraph}{III.4.2.1.4.4 Business Functions}

\subparagraph*{III.4.2.1.4.4.1 Apply Governed Person Identity Correction}
\phantomsection
\label{navigation:III.4.2.1.4.4.1}
\addcontentsline{toc}{subparagraph}{III.4.2.1.4.4.1 Apply Governed Person Identity Correction}

\label{element:844b9abe470d4f4684f35adddf1b880d}
\index{Apply Governed Person Identity Correction (III.4.2.1.4.4.1)}
{\small BusinessFunction · Canonical · current}\par

Processes externally authoritative corrections and merge/unlink outcomes, updates governed identity associations and propagates resulting managed-state consequences with audit history. Harmonia does not adjudicate the originating person-merge decision.


{\footnotesize \href{../../../../continuum/business/business-function/844b9abe470d4f4684f35adddf1b880d.md}{Canonical record: 844b9abe470d4f4684f35adddf1b880d}; full qualifications and provenance.}\par
\subparagraph*{III.4.2.1.4.4.2 Correlate Person Identifiers}
\phantomsection
\label{navigation:III.4.2.1.4.4.2}
\addcontentsline{toc}{subparagraph}{III.4.2.1.4.4.2 Correlate Person Identifiers}

\label{element:3050366e2e1d47eeaa4cfd6db92c0a56}
\index{Correlate Person Identifiers (III.4.2.1.4.4.2)}
{\small BusinessFunction · Canonical · current}\par

Correlates local identifiers (e.g., Hospital MRN, National IHI, Pathology ID) using relationships established by authoritative information, retaining their source and namespace.


{\footnotesize \href{../../../../continuum/business/business-function/3050366e2e1d47eeaa4cfd6db92c0a56.md}{Canonical record: 3050366e2e1d47eeaa4cfd6db92c0a56}; full qualifications and provenance.}\par
\subparagraph*{III.4.2.1.4.4.3 Maintain Person Identity Aliases}
\phantomsection
\label{navigation:III.4.2.1.4.4.3}
\addcontentsline{toc}{subparagraph}{III.4.2.1.4.4.3 Maintain Person Identity Aliases}

\label{element:9db2528e180a418eacbad3acf86b3956}
\index{Maintain Person Identity Aliases (III.4.2.1.4.4.3)}
{\small BusinessFunction · Canonical · current}\par

Records and qualifies alias names, titles, and emergency aliases for a governed person identity.


{\footnotesize \href{../../../../continuum/business/business-function/9db2528e180a418eacbad3acf86b3956.md}{Canonical record: 9db2528e180a418eacbad3acf86b3956}; full qualifications and provenance.}\par
\subparagraph*{III.4.2.1.4.4.4 Resolve Person Identifier}
\phantomsection
\label{navigation:III.4.2.1.4.4.4}
\addcontentsline{toc}{subparagraph}{III.4.2.1.4.4.4 Resolve Person Identifier}

\label{element:2f128e1d8de642478c18b7a0ccb8e083}
\index{Resolve Person Identifier (III.4.2.1.4.4.4)}
{\small BusinessFunction · Canonical · current}\par

Validates and retrieves the authoritative demographic record bound to a specific identifier.


{\footnotesize \href{../../../../continuum/business/business-function/2f128e1d8de642478c18b7a0ccb8e083.md}{Canonical record: 2f128e1d8de642478c18b7a0ccb8e083}; full qualifications and provenance.}\par
\paragraph*{III.4.2.1.4.5 Business Processes}
\phantomsection
\label{navigation:III.4.2.1.4.5}
\addcontentsline{toc}{paragraph}{III.4.2.1.4.5 Business Processes}

\subparagraph*{III.4.2.1.4.5.1 Governed Person Identity Correction Process}
\phantomsection
\label{navigation:III.4.2.1.4.5.1}
\addcontentsline{toc}{subparagraph}{III.4.2.1.4.5.1 Governed Person Identity Correction Process}

\label{element:72f3c63da45d44868579a3bb15d1568c}
\index{Governed Person Identity Correction Process (III.4.2.1.4.5.1)}
{\small BusinessProcess · Canonical · current}\par

•
Owning Capability: Person Identity (under Client Administration)



•
Process Purpose: Processes externally authoritative identity corrections and merge/unlink outcomes, applies their consequences to Harmonia-managed identity information and identifier associations, and communicates resulting changes with an audit history.



•
Authority boundary: Person correction/merge decisions originate externally. Internal triage, evidence verification and approval concern applicability and governed processing of that authoritative outcome, not person matching, golden-record determination, master-person selection or originating merge adjudication.



•
State Progression Lifecycle: Correction Requested → HIM Triage → Evidence Verified → Correction Approved → Merge/Unlink Executed → Change Broadcasted → Closed



•
Key State Dispositions:



•
Correction Requested: Externally authoritative correction or merge/unlink outcome received for governed application.



•
HIM Triage: Health Information Manager reviews the received decision/outcome and its applicability.



•
Evidence Verified: Authoritative source decision and supporting provenance/documentation verified.



•
Correction Approved: Governed application of the externally authoritative outcome approved; this does not originate a person-merge decision.



•
Merge/Unlink Executed: External merge/unlink consequences applied to Harmonia's managed identifier associations and audit history.



•
Change Broadcasted: Downstream clinical systems notified of merged or rectified identity.



This Process governs the application of an externally authoritative correction rather than originating a person-merge judgement. Triage, evidence verification and approval establish that the received outcome should be applied to Harmonia's associations. The progression and change communication preserve why the managed identity account changed without making internal approval a master-person decision.


{\small Recorded unresolved matters: source-qualified matters; exact scope remains in the canonical record.}\par
{\footnotesize \href{../../../../continuum/business/business-process/72f3c63da45d44868579a3bb15d1568c.md}{Canonical record: 72f3c63da45d44868579a3bb15d1568c}; full qualifications and provenance.}\par
\subsection*{III.4.2.2 [CT2] Clinical Device Administration — Capability context}
\phantomsection
\label{navigation:III.4.2.2}
\addcontentsline{toc}{subsection}{III.4.2.2 [CT2] Clinical Device Administration — Capability context}

{\small Canonical Strategy Capability: II.1.3.4.2 \hyperref[element:31bcd6fda0d14cc8b992cdca39a3cbed]{[CT2] Clinical Device Administration}; Canonical ID: \texttt{\seqsplit{31bcd6fda0d14cc8b992cdca39a3cbed}}. Canonical · current.}\par
Contextual reference only; no additional architectural instance or semantic relationship.\par
\subsubsection*{III.4.2.2.1 [CT3] Device Utilisation Administration — Capability context}
\phantomsection
\label{navigation:III.4.2.2.1}
\addcontentsline{toc}{subsubsection}{III.4.2.2.1 [CT3] Device Utilisation Administration — Capability context}

{\small Canonical Strategy Capability: II.1.3.4.2.7 \hyperref[element:1a3732b5911c45b3aece38b2db2d0f15]{[CT3] Device Utilisation Administration}; Canonical ID: \texttt{\seqsplit{1a3732b5911c45b3aece38b2db2d0f15}}. Proposed · current.}\par
Contextual reference only; no additional architectural instance or semantic relationship.\par
\paragraph*{III.4.2.2.1.1 Business Services}
\phantomsection
\label{navigation:III.4.2.2.1.1}
\addcontentsline{toc}{paragraph}{III.4.2.2.1.1 Business Services}

\subparagraph*{III.4.2.2.1.1.1 Clinical Resource Availability Query}
\phantomsection
\label{navigation:III.4.2.2.1.1.1}
\addcontentsline{toc}{subparagraph}{III.4.2.2.1.1.1 Clinical Resource Availability Query}

\label{element:0e5ed162a288425aad2833ed5e07ed64}
\index{Clinical Resource Availability Query (III.4.2.2.1.1.1)}
{\small BusinessService · Canonical · current}\par

Exposes mobile equipment availability.


{\footnotesize \href{../../../../continuum/business/business-service/0e5ed162a288425aad2833ed5e07ed64.md}{Canonical record: 0e5ed162a288425aad2833ed5e07ed64}; full qualifications and provenance.}\par
\paragraph*{III.4.2.2.1.2 Business Functions}
\phantomsection
\label{navigation:III.4.2.2.1.2}
\addcontentsline{toc}{paragraph}{III.4.2.2.1.2 Business Functions}

\subparagraph*{III.4.2.2.1.2.1 Track Mobile Clinical Asset Operational State}
\phantomsection
\label{navigation:III.4.2.2.1.2.1}
\addcontentsline{toc}{subparagraph}{III.4.2.2.1.2.1 Track Mobile Clinical Asset Operational State}

\label{element:0a31e03a6c7847c0b7c04310ca0c29e9}
\index{Track Mobile Clinical Asset Operational State (III.4.2.2.1.2.1)}
{\small BusinessFunction · Canonical · current}\par

Manages the real-time location and operational readiness of mobile ventilators, infusion pumps, and telemetry transmitters.


{\footnotesize \href{../../../../continuum/business/business-function/0a31e03a6c7847c0b7c04310ca0c29e9.md}{Canonical record: 0a31e03a6c7847c0b7c04310ca0c29e9}; full qualifications and provenance.}\par
\subsubsection*{III.4.2.2.2 Business Services}
\phantomsection
\label{navigation:III.4.2.2.2}
\addcontentsline{toc}{subsubsection}{III.4.2.2.2 Business Services}

\paragraph*{III.4.2.2.2.1 Device Association Query}
\phantomsection
\label{navigation:III.4.2.2.2.1}
\addcontentsline{toc}{paragraph}{III.4.2.2.2.1 Device Association Query}

\label{element:01637c40914945d68d261c2094b86904}
\index{Device Association Query (III.4.2.2.2.1)}
{\small BusinessService · Canonical · current}\par

Discloses active device-to-patient and device-to-location bindings.


{\footnotesize \href{../../../../continuum/business/business-service/01637c40914945d68d261c2094b86904.md}{Canonical record: 01637c40914945d68d261c2094b86904}; full qualifications and provenance.}\par
\paragraph*{III.4.2.2.2.2 Device Definition Query}
\phantomsection
\label{navigation:III.4.2.2.2.2}
\addcontentsline{toc}{paragraph}{III.4.2.2.2.2 Device Definition Query}

\label{element:49325c789401465aa565fd7fa7109070}
\index{Device Definition Query (III.4.2.2.2.2)}
{\small BusinessService · Canonical · current}\par

Discloses approved medical device models and technical parameters.


{\footnotesize \href{../../../../continuum/business/business-service/49325c789401465aa565fd7fa7109070.md}{Canonical record: 49325c789401465aa565fd7fa7109070}; full qualifications and provenance.}\par
\paragraph*{III.4.2.2.2.3 Device Endpoint Resolution}
\phantomsection
\label{navigation:III.4.2.2.2.3}
\addcontentsline{toc}{paragraph}{III.4.2.2.2.3 Device Endpoint Resolution}

\label{element:5806d6e030b543d0867a862d5b6d8fad}
\index{Device Endpoint Resolution (III.4.2.2.2.3)}
{\small BusinessService · Canonical · current}\par

Exposes network addresses for telemetry routing.


{\footnotesize \href{../../../../continuum/business/business-service/5806d6e030b543d0867a862d5b6d8fad.md}{Canonical record: 5806d6e030b543d0867a862d5b6d8fad}; full qualifications and provenance.}\par
\paragraph*{III.4.2.2.2.4 Device Instance Resolution}
\phantomsection
\label{navigation:III.4.2.2.2.4}
\addcontentsline{toc}{paragraph}{III.4.2.2.2.4 Device Instance Resolution}

\label{element:14fb1e3f17dd40d88087a00c9707578d}
\index{Device Instance Resolution (III.4.2.2.2.4)}
{\small BusinessService · Canonical · current}\par

Resolves physical device instances from scanned barcodes or UDIs.


{\footnotesize \href{../../../../continuum/business/business-service/14fb1e3f17dd40d88087a00c9707578d.md}{Canonical record: 14fb1e3f17dd40d88087a00c9707578d}; full qualifications and provenance.}\par
\subsubsection*{III.4.2.2.3 Business Functions}
\phantomsection
\label{navigation:III.4.2.2.3}
\addcontentsline{toc}{subsubsection}{III.4.2.2.3 Business Functions}

\paragraph*{III.4.2.2.3.1 Maintain Device Association}
\phantomsection
\label{navigation:III.4.2.2.3.1}
\addcontentsline{toc}{paragraph}{III.4.2.2.3.1 Maintain Device Association}

\label{element:fc0f4d9c781241359d783d26f038b7ff}
\index{Maintain Device Association (III.4.2.2.3.1)}
{\small BusinessFunction · Canonical · current}\par

Records active associations between clinical devices, patient encounters, and specific care-places (e.g., Bed 4 Vital Signs Monitor).


{\footnotesize \href{../../../../continuum/business/business-function/fc0f4d9c781241359d783d26f038b7ff.md}{Canonical record: fc0f4d9c781241359d783d26f038b7ff}; full qualifications and provenance.}\par
\paragraph*{III.4.2.2.3.2 Maintain Device Communication Endpoints}
\phantomsection
\label{navigation:III.4.2.2.3.2}
\addcontentsline{toc}{paragraph}{III.4.2.2.3.2 Maintain Device Communication Endpoints}

\label{element:ab5a7538fbd449ae8214d584b7d53a43}
\index{Maintain Device Communication Endpoints (III.4.2.2.3.2)}
{\small BusinessFunction · Canonical · current}\par

Manages network telemetry addresses, gateway ports, and protocol bindings for medical devices.


{\footnotesize \href{../../../../continuum/business/business-function/ab5a7538fbd449ae8214d584b7d53a43.md}{Canonical record: ab5a7538fbd449ae8214d584b7d53a43}; full qualifications and provenance.}\par
\paragraph*{III.4.2.2.3.3 Maintain Device Definition}
\phantomsection
\label{navigation:III.4.2.2.3.3}
\addcontentsline{toc}{paragraph}{III.4.2.2.3.3 Maintain Device Definition}

\label{element:d44083e67fb540579737337cabf2c227}
\index{Maintain Device Definition (III.4.2.2.3.3)}
{\small BusinessFunction · Canonical · current}\par

Models medical device models, manufacturers, regulatory classifications (e.g., TGA/FDA classes), and calibration specifications.


{\footnotesize \href{../../../../continuum/business/business-function/d44083e67fb540579737337cabf2c227.md}{Canonical record: d44083e67fb540579737337cabf2c227}; full qualifications and provenance.}\par
\paragraph*{III.4.2.2.3.4 Maintain Device Instance Identity}
\phantomsection
\label{navigation:III.4.2.2.3.4}
\addcontentsline{toc}{paragraph}{III.4.2.2.3.4 Maintain Device Instance Identity}

\label{element:927796d77f5e4f61835bd2a7bcceb681}
\index{Maintain Device Instance Identity (III.4.2.2.3.4)}
{\small BusinessFunction · Canonical · current}\par

Tracks physical serial numbers, Unique Device Identifiers (UDI), asset tags, and firmware versions.


{\footnotesize \href{../../../../continuum/business/business-function/927796d77f5e4f61835bd2a7bcceb681.md}{Canonical record: 927796d77f5e4f61835bd2a7bcceb681}; full qualifications and provenance.}\par
\subsection*{III.4.2.3 [CT2] Health Service Administration — Capability context}
\phantomsection
\label{navigation:III.4.2.3}
\addcontentsline{toc}{subsection}{III.4.2.3 [CT2] Health Service Administration — Capability context}

{\small Canonical Strategy Capability: II.1.3.4.4 \hyperref[element:ba68a52913a0499d8377553bddfa74a4]{[CT2] Health Service Administration}; Canonical ID: \texttt{\seqsplit{ba68a52913a0499d8377553bddfa74a4}}. Canonical · current.}\par
Contextual reference only; no additional architectural instance or semantic relationship.\par
\subsubsection*{III.4.2.3.1 Business Services}
\phantomsection
\label{navigation:III.4.2.3.1}
\addcontentsline{toc}{subsubsection}{III.4.2.3.1 Business Services}

\paragraph*{III.4.2.3.1.1 Healthcare Service Directory Query}
\phantomsection
\label{navigation:III.4.2.3.1.1}
\addcontentsline{toc}{paragraph}{III.4.2.3.1.1 Healthcare Service Directory Query}

\label{element:c7bf105e1217441a9423706b1481d153}
\index{Healthcare Service Directory Query (III.4.2.3.1.1)}
{\small BusinessService · Canonical · current}\par

Exposes published clinical service offerings.


{\footnotesize \href{../../../../continuum/business/business-service/c7bf105e1217441a9423706b1481d153.md}{Canonical record: c7bf105e1217441a9423706b1481d153}; full qualifications and provenance.}\par
\paragraph*{III.4.2.3.1.2 Service Availability Query}
\phantomsection
\label{navigation:III.4.2.3.1.2}
\addcontentsline{toc}{paragraph}{III.4.2.3.1.2 Service Availability Query}

\label{element:a3846cb2960a4000a5db9ef18942ee20}
\index{Service Availability Query (III.4.2.3.1.2)}
{\small BusinessService · Canonical · current}\par

Discloses operational intake hours and open status.


{\footnotesize \href{../../../../continuum/business/business-service/a3846cb2960a4000a5db9ef18942ee20.md}{Canonical record: a3846cb2960a4000a5db9ef18942ee20}; full qualifications and provenance.}\par
\paragraph*{III.4.2.3.1.3 Service Eligibility Evaluation}
\phantomsection
\label{navigation:III.4.2.3.1.3}
\addcontentsline{toc}{paragraph}{III.4.2.3.1.3 Service Eligibility Evaluation}

\label{element:7e005ce374f347088ae854442fc3edef}
\index{Service Eligibility Evaluation (III.4.2.3.1.3)}
{\small BusinessService · Canonical · current}\par

Evaluates referral eligibility for intake coordinators.


{\footnotesize \href{../../../../continuum/business/business-service/7e005ce374f347088ae854442fc3edef.md}{Canonical record: 7e005ce374f347088ae854442fc3edef}; full qualifications and provenance.}\par
\paragraph*{III.4.2.3.1.4 Service Provision Resolution}
\phantomsection
\label{navigation:III.4.2.3.1.4}
\addcontentsline{toc}{paragraph}{III.4.2.3.1.4 Service Provision Resolution}

\label{element:e54866d88b1e442f8ea170f06f98028b}
\index{Service Provision Resolution (III.4.2.3.1.4)}
{\small BusinessService · Canonical · current}\par

Exposes multi-dimensional resolution queries answering who, where, and when healthcare services are delivered.


{\footnotesize \href{../../../../continuum/business/business-service/e54866d88b1e442f8ea170f06f98028b.md}{Canonical record: e54866d88b1e442f8ea170f06f98028b}; full qualifications and provenance.}\par
\subsubsection*{III.4.2.3.2 Business Functions}
\phantomsection
\label{navigation:III.4.2.3.2}
\addcontentsline{toc}{subsubsection}{III.4.2.3.2 Business Functions}

\paragraph*{III.4.2.3.2.1 Maintain Healthcare Service Definition}
\phantomsection
\label{navigation:III.4.2.3.2.1}
\addcontentsline{toc}{paragraph}{III.4.2.3.2.1 Maintain Healthcare Service Definition}

\label{element:833a98c58dae416da783ef6e77290576}
\index{Maintain Healthcare Service Definition (III.4.2.3.2.1)}
{\small BusinessFunction · Canonical · current}\par

Defines clinical service offerings (e.g., Cardiology Outpatient Clinic, MRI Imaging, Acute Stroke Care), specialty classifications, and clinical scope.


{\footnotesize \href{../../../../continuum/business/business-function/833a98c58dae416da783ef6e77290576.md}{Canonical record: 833a98c58dae416da783ef6e77290576}; full qualifications and provenance.}\par
\paragraph*{III.4.2.3.2.2 Maintain Service / Location / Provider Map}
\phantomsection
\label{navigation:III.4.2.3.2.2}
\addcontentsline{toc}{paragraph}{III.4.2.3.2.2 Maintain Service / Location / Provider Map}

\label{element:ea06c77ea3c841cc94cb1686295322b4}
\index{Maintain Service / Location / Provider Map (III.4.2.3.2.2)}
{\small BusinessFunction · Canonical · current}\par

Binds health service definitions to delivering organisations, physical clinic locations, and participating practitioner roles.


{\footnotesize \href{../../../../continuum/business/business-function/ea06c77ea3c841cc94cb1686295322b4.md}{Canonical record: ea06c77ea3c841cc94cb1686295322b4}; full qualifications and provenance.}\par
\paragraph*{III.4.2.3.2.3 Maintain Service Availability}
\phantomsection
\label{navigation:III.4.2.3.2.3}
\addcontentsline{toc}{paragraph}{III.4.2.3.2.3 Maintain Service Availability}

\label{element:38321dd2581b46a182dcdf52865d44be}
\index{Maintain Service Availability (III.4.2.3.2.3)}
{\small BusinessFunction · Canonical · current}\par

Governs regular operating hours, emergency intake availability, and seasonal clinic schedules.


{\footnotesize \href{../../../../continuum/business/business-function/38321dd2581b46a182dcdf52865d44be.md}{Canonical record: 38321dd2581b46a182dcdf52865d44be}; full qualifications and provenance.}\par
\paragraph*{III.4.2.3.2.4 Maintain Service Eligibility Rules}
\phantomsection
\label{navigation:III.4.2.3.2.4}
\addcontentsline{toc}{paragraph}{III.4.2.3.2.4 Maintain Service Eligibility Rules}

\label{element:e1e1cf2aa30840368d107ab9dfd0257f}
\index{Maintain Service Eligibility Rules (III.4.2.3.2.4)}
{\small BusinessFunction · Canonical · current}\par

Configures catchment criteria, age thresholds, referral requirements, and clinical indications for service access.


{\footnotesize \href{../../../../continuum/business/business-function/e1e1cf2aa30840368d107ab9dfd0257f.md}{Canonical record: e1e1cf2aa30840368d107ab9dfd0257f}; full qualifications and provenance.}\par
\subsection*{III.4.2.4 [CT2] Location Administration — Capability context}
\phantomsection
\label{navigation:III.4.2.4}
\addcontentsline{toc}{subsection}{III.4.2.4 [CT2] Location Administration — Capability context}

{\small Canonical Strategy Capability: II.1.3.4.5 \hyperref[element:c410005667444f99a221483bf577fad2]{[CT2] Location Administration}; Canonical ID: \texttt{\seqsplit{c410005667444f99a221483bf577fad2}}. Canonical · current.}\par
Contextual reference only; no additional architectural instance or semantic relationship.\par
\subsubsection*{III.4.2.4.1 Business Services}
\phantomsection
\label{navigation:III.4.2.4.1}
\addcontentsline{toc}{subsubsection}{III.4.2.4.1 Business Services}

\paragraph*{III.4.2.4.1.1 Care-Place Specification Lookup}
\phantomsection
\label{navigation:III.4.2.4.1.1}
\addcontentsline{toc}{paragraph}{III.4.2.4.1.1 Care-Place Specification Lookup}

\label{element:43b011b0c233408eaf6e555ecbf05d9e}
\index{Care-Place Specification Lookup (III.4.2.4.1.1)}
{\small BusinessService · Canonical · current}\par

Discloses care-place physical capabilities to Bed \& Care-Place Management.


{\footnotesize \href{../../../../continuum/business/business-service/43b011b0c233408eaf6e555ecbf05d9e.md}{Canonical record: 43b011b0c233408eaf6e555ecbf05d9e}; full qualifications and provenance.}\par
\paragraph*{III.4.2.4.1.2 Location Directory Query}
\phantomsection
\label{navigation:III.4.2.4.1.2}
\addcontentsline{toc}{paragraph}{III.4.2.4.1.2 Location Directory Query}

\label{element:5ef66362eb674794b335c435227b0c95}
\index{Location Directory Query (III.4.2.4.1.2)}
{\small BusinessService · Canonical · current}\par

Exposes location profiles and geospatial addresses.


{\footnotesize \href{../../../../continuum/business/business-service/5ef66362eb674794b335c435227b0c95.md}{Canonical record: 5ef66362eb674794b335c435227b0c95}; full qualifications and provenance.}\par
\paragraph*{III.4.2.4.1.3 Location Hierarchy Resolution}
\phantomsection
\label{navigation:III.4.2.4.1.3}
\addcontentsline{toc}{paragraph}{III.4.2.4.1.3 Location Hierarchy Resolution}

\label{element:389af7aaa9d74ab8b42e55f0fa72ef99}
\index{Location Hierarchy Resolution (III.4.2.4.1.3)}
{\small BusinessService · Canonical · current}\par

Discloses physical containment and navigation paths for clinical facilities.


{\footnotesize \href{../../../../continuum/business/business-service/389af7aaa9d74ab8b42e55f0fa72ef99.md}{Canonical record: 389af7aaa9d74ab8b42e55f0fa72ef99}; full qualifications and provenance.}\par
\subsubsection*{III.4.2.4.2 Business Functions}
\phantomsection
\label{navigation:III.4.2.4.2}
\addcontentsline{toc}{subsubsection}{III.4.2.4.2 Business Functions}

\paragraph*{III.4.2.4.2.1 Maintain Care-Place Definition}
\phantomsection
\label{navigation:III.4.2.4.2.1}
\addcontentsline{toc}{paragraph}{III.4.2.4.2.1 Maintain Care-Place Definition}

\label{element:ad9aa4eeecf84eb08460234ed0dc4ecc}
\index{Maintain Care-Place Definition (III.4.2.4.2.1)}
{\small BusinessFunction · Canonical · current}\par

Defines static bed bays, treatment chairs, surgical tables, and emergency resuscitation cubicles, including equipment capabilities and oxygen/telemetry infrastructure.


{\footnotesize \href{../../../../continuum/business/business-function/ad9aa4eeecf84eb08460234ed0dc4ecc.md}{Canonical record: ad9aa4eeecf84eb08460234ed0dc4ecc}; full qualifications and provenance.}\par
\paragraph*{III.4.2.4.2.2 Maintain Healthcare Location}
\phantomsection
\label{navigation:III.4.2.4.2.2}
\addcontentsline{toc}{paragraph}{III.4.2.4.2.2 Maintain Healthcare Location}

\label{element:448f1c9317914498ab246f74d4ce2bd6}
\index{Maintain Healthcare Location (III.4.2.4.2.2)}
{\small BusinessFunction · Canonical · current}\par

Records physical sites, campuses, outpatient clinics, and community health centres.


{\footnotesize \href{../../../../continuum/business/business-function/448f1c9317914498ab246f74d4ce2bd6.md}{Canonical record: 448f1c9317914498ab246f74d4ce2bd6}; full qualifications and provenance.}\par
\paragraph*{III.4.2.4.2.3 Maintain Physical Location Hierarchy}
\phantomsection
\label{navigation:III.4.2.4.2.3}
\addcontentsline{toc}{paragraph}{III.4.2.4.2.3 Maintain Physical Location Hierarchy}

\label{element:d217a5cd3a3d40e9a89068d1680ec45f}
\index{Maintain Physical Location Hierarchy (III.4.2.4.2.3)}
{\small BusinessFunction · Canonical · current}\par

Models the nested structural containment tree: Campus → Building → Floor → Ward/Suite → Room/Bay → Bed Care-Place.


{\footnotesize \href{../../../../continuum/business/business-function/d217a5cd3a3d40e9a89068d1680ec45f.md}{Canonical record: d217a5cd3a3d40e9a89068d1680ec45f}; full qualifications and provenance.}\par
\subsection*{III.4.2.5 [CT2] Organisation Administration — Capability context}
\phantomsection
\label{navigation:III.4.2.5}
\addcontentsline{toc}{subsection}{III.4.2.5 [CT2] Organisation Administration — Capability context}

{\small Canonical Strategy Capability: II.1.3.4.6 \hyperref[element:14fbdc5ce3dd4c2092c5397dab3ed743]{[CT2] Organisation Administration}; Canonical ID: \texttt{\seqsplit{14fbdc5ce3dd4c2092c5397dab3ed743}}. Canonical · current.}\par
Contextual reference only; no additional architectural instance or semantic relationship.\par
\subsubsection*{III.4.2.5.1 Business Services}
\phantomsection
\label{navigation:III.4.2.5.1}
\addcontentsline{toc}{subsubsection}{III.4.2.5.1 Business Services}

\paragraph*{III.4.2.5.1.1 Organisation Contact Lookup}
\phantomsection
\label{navigation:III.4.2.5.1.1}
\addcontentsline{toc}{paragraph}{III.4.2.5.1.1 Organisation Contact Lookup}

\label{element:551f1e73fb73449a83f3155cb8158f30}
\index{Organisation Contact Lookup (III.4.2.5.1.1)}
{\small BusinessService · Canonical · current}\par

Provides official operational contact details.


{\footnotesize \href{../../../../continuum/business/business-service/551f1e73fb73449a83f3155cb8158f30.md}{Canonical record: 551f1e73fb73449a83f3155cb8158f30}; full qualifications and provenance.}\par
\paragraph*{III.4.2.5.1.2 Organisation Directory Query}
\phantomsection
\label{navigation:III.4.2.5.1.2}
\addcontentsline{toc}{paragraph}{III.4.2.5.1.2 Organisation Directory Query}

\label{element:c5acf1dc97874d0a8264f9a5c1b735fa}
\index{Organisation Directory Query (III.4.2.5.1.2)}
{\small BusinessService · Canonical · current}\par

Discloses healthcare organisation profiles across the region.


{\footnotesize \href{../../../../continuum/business/business-service/c5acf1dc97874d0a8264f9a5c1b735fa.md}{Canonical record: c5acf1dc97874d0a8264f9a5c1b735fa}; full qualifications and provenance.}\par
\paragraph*{III.4.2.5.1.3 Organisation Verification Service}
\phantomsection
\label{navigation:III.4.2.5.1.3}
\addcontentsline{toc}{paragraph}{III.4.2.5.1.3 Organisation Verification Service}

\label{element:223b2cbe762b40efb9f05bb5cdd66bcb}
\index{Organisation Verification Service (III.4.2.5.1.3)}
{\small BusinessService · Canonical · current}\par

Provides verified organisation credentials to external and internal partners.


{\footnotesize \href{../../../../continuum/business/business-service/223b2cbe762b40efb9f05bb5cdd66bcb.md}{Canonical record: 223b2cbe762b40efb9f05bb5cdd66bcb}; full qualifications and provenance.}\par
\paragraph*{III.4.2.5.1.4 Organisational Hierarchy Resolution}
\phantomsection
\label{navigation:III.4.2.5.1.4}
\addcontentsline{toc}{paragraph}{III.4.2.5.1.4 Organisational Hierarchy Resolution}

\label{element:c9187704dd7f49f69f9e674c16e9181b}
\index{Organisational Hierarchy Resolution (III.4.2.5.1.4)}
{\small BusinessService · Canonical · current}\par

Resolves department nesting and reporting hierarchies.


{\footnotesize \href{../../../../continuum/business/business-service/c9187704dd7f49f69f9e674c16e9181b.md}{Canonical record: c9187704dd7f49f69f9e674c16e9181b}; full qualifications and provenance.}\par
\subsubsection*{III.4.2.5.2 Business Functions}
\phantomsection
\label{navigation:III.4.2.5.2}
\addcontentsline{toc}{subsubsection}{III.4.2.5.2 Business Functions}

\paragraph*{III.4.2.5.2.1 Govern Organisation Profile}
\phantomsection
\label{navigation:III.4.2.5.2.1}
\addcontentsline{toc}{paragraph}{III.4.2.5.2.1 Govern Organisation Profile}

\label{element:08990a2aa346403f8dd7d716acf188bf}
\index{Govern Organisation Profile (III.4.2.5.2.1)}
{\small BusinessFunction · Canonical · current}\par

Maintains official trading names, enterprise registration numbers, and governing bodies.


{\footnotesize \href{../../../../continuum/business/business-function/08990a2aa346403f8dd7d716acf188bf.md}{Canonical record: 08990a2aa346403f8dd7d716acf188bf}; full qualifications and provenance.}\par
\paragraph*{III.4.2.5.2.2 Maintain Organisation Contacts}
\phantomsection
\label{navigation:III.4.2.5.2.2}
\addcontentsline{toc}{paragraph}{III.4.2.5.2.2 Maintain Organisation Contacts}

\label{element:4584f74d08834bc2bc595463ad51bd37}
\index{Maintain Organisation Contacts (III.4.2.5.2.2)}
{\small BusinessFunction · Canonical · current}\par

Manages administrative, clinical, and emergency contact details for facilities.


{\footnotesize \href{../../../../continuum/business/business-function/4584f74d08834bc2bc595463ad51bd37.md}{Canonical record: 4584f74d08834bc2bc595463ad51bd37}; full qualifications and provenance.}\par
\paragraph*{III.4.2.5.2.3 Maintain Organisational Structure}
\phantomsection
\label{navigation:III.4.2.5.2.3}
\addcontentsline{toc}{paragraph}{III.4.2.5.2.3 Maintain Organisational Structure}

\label{element:5ae4edf0813f4a738fdb3160f3e77b0d}
\index{Maintain Organisational Structure (III.4.2.5.2.3)}
{\small BusinessFunction · Canonical · current}\par

Models parent-subsidiary relationships, clinical directorates, departments, and operational divisions.


{\footnotesize \href{../../../../continuum/business/business-function/5ae4edf0813f4a738fdb3160f3e77b0d.md}{Canonical record: 5ae4edf0813f4a738fdb3160f3e77b0d}; full qualifications and provenance.}\par
\paragraph*{III.4.2.5.2.4 Verify Healthcare Organisation}
\phantomsection
\label{navigation:III.4.2.5.2.4}
\addcontentsline{toc}{paragraph}{III.4.2.5.2.4 Verify Healthcare Organisation}

\label{element:c39236ca89864f089684242df5647c9d}
\index{Verify Healthcare Organisation (III.4.2.5.2.4)}
{\small BusinessFunction · Canonical · current}\par

Validates enterprise legal identity, national facility identifiers (e.g., HPI-O), and accreditation status.


{\footnotesize \href{../../../../continuum/business/business-function/c39236ca89864f089684242df5647c9d.md}{Canonical record: c39236ca89864f089684242df5647c9d}; full qualifications and provenance.}\par
\subsection*{III.4.2.6 [CT2] Provider Administration — Capability context}
\phantomsection
\label{navigation:III.4.2.6}
\addcontentsline{toc}{subsection}{III.4.2.6 [CT2] Provider Administration — Capability context}

{\small Canonical Strategy Capability: II.1.3.4.7 \hyperref[element:d6312f85968a4209bafff5ed009297a6]{[CT2] Provider Administration}; Canonical ID: \texttt{\seqsplit{d6312f85968a4209bafff5ed009297a6}}. Canonical · current.}\par
Contextual reference only; no additional architectural instance or semantic relationship.\par
\subsubsection*{III.4.2.6.1 Business Services}
\phantomsection
\label{navigation:III.4.2.6.1}
\addcontentsline{toc}{subsubsection}{III.4.2.6.1 Business Services}

\paragraph*{III.4.2.6.1.1 Clinical Privilege Verification}
\phantomsection
\label{navigation:III.4.2.6.1.1}
\addcontentsline{toc}{paragraph}{III.4.2.6.1.1 Clinical Privilege Verification}

\label{element:f6599b4283194577a977bfc86ea854c2}
\index{Clinical Privilege Verification (III.4.2.6.1.1)}
{\small BusinessService · Canonical · current}\par

Verifies procedural and prescribing authority for clinical ordering systems.


{\footnotesize \href{../../../../continuum/business/business-service/f6599b4283194577a977bfc86ea854c2.md}{Canonical record: f6599b4283194577a977bfc86ea854c2}; full qualifications and provenance.}\par
\paragraph*{III.4.2.6.1.2 Practitioner Affiliation Resolution}
\phantomsection
\label{navigation:III.4.2.6.1.2}
\addcontentsline{toc}{paragraph}{III.4.2.6.1.2 Practitioner Affiliation Resolution}

\label{element:65baca2a21f040d49be74fab1727f96d}
\index{Practitioner Affiliation Resolution (III.4.2.6.1.2)}
{\small BusinessService · Canonical · current}\par

Discloses organizational ties and practice locations.


{\footnotesize \href{../../../../continuum/business/business-service/65baca2a21f040d49be74fab1727f96d.md}{Canonical record: 65baca2a21f040d49be74fab1727f96d}; full qualifications and provenance.}\par
\paragraph*{III.4.2.6.1.3 Practitioner Directory Query}
\phantomsection
\label{navigation:III.4.2.6.1.3}
\addcontentsline{toc}{paragraph}{III.4.2.6.1.3 Practitioner Directory Query}

\label{element:ed8e6a3d891340fe92fc3201abd15c46}
\index{Practitioner Directory Query (III.4.2.6.1.3)}
{\small BusinessService · Canonical · current}\par

Provides provider directory lookups across the healthcare network.


{\footnotesize \href{../../../../continuum/business/business-service/ed8e6a3d891340fe92fc3201abd15c46.md}{Canonical record: ed8e6a3d891340fe92fc3201abd15c46}; full qualifications and provenance.}\par
\paragraph*{III.4.2.6.1.4 Practitioner Endpoint Lookup}
\phantomsection
\label{navigation:III.4.2.6.1.4}
\addcontentsline{toc}{paragraph}{III.4.2.6.1.4 Practitioner Endpoint Lookup}

\label{element:2efb21629acf41e9b4f5c59db00721d0}
\index{Practitioner Endpoint Lookup (III.4.2.6.1.4)}
{\small BusinessService · Canonical · current}\par

Provides secure message routing destinations for clinical communication.


{\footnotesize \href{../../../../continuum/business/business-service/2efb21629acf41e9b4f5c59db00721d0.md}{Canonical record: 2efb21629acf41e9b4f5c59db00721d0}; full qualifications and provenance.}\par
\paragraph*{III.4.2.6.1.5 Practitioner Role Resolution}
\phantomsection
\label{navigation:III.4.2.6.1.5}
\addcontentsline{toc}{paragraph}{III.4.2.6.1.5 Practitioner Role Resolution}

\label{element:c3e5c04e2a194b4a8d49b8776fce2c51}
\index{Practitioner Role Resolution (III.4.2.6.1.5)}
{\small BusinessService · Canonical · current}\par

Resolves the specific clinical capacities of a practitioner.


{\footnotesize \href{../../../../continuum/business/business-service/c3e5c04e2a194b4a8d49b8776fce2c51.md}{Canonical record: c3e5c04e2a194b4a8d49b8776fce2c51}; full qualifications and provenance.}\par
\paragraph*{III.4.2.6.1.6 Practitioner Verification Service}
\phantomsection
\label{navigation:III.4.2.6.1.6}
\addcontentsline{toc}{paragraph}{III.4.2.6.1.6 Practitioner Verification Service}

\label{element:adf039a87e9243a1b82612a4919f6f14}
\index{Practitioner Verification Service (III.4.2.6.1.6)}
{\small BusinessService · Canonical · current}\par

Discloses verified registration and status to credentialing and clinical systems.


{\footnotesize \href{../../../../continuum/business/business-service/adf039a87e9243a1b82612a4919f6f14.md}{Canonical record: adf039a87e9243a1b82612a4919f6f14}; full qualifications and provenance.}\par
\subsubsection*{III.4.2.6.2 Business Functions}
\phantomsection
\label{navigation:III.4.2.6.2}
\addcontentsline{toc}{subsubsection}{III.4.2.6.2 Business Functions}

\paragraph*{III.4.2.6.2.1 Bind Practitioner Electronic Endpoints}
\phantomsection
\label{navigation:III.4.2.6.2.1}
\addcontentsline{toc}{paragraph}{III.4.2.6.2.1 Bind Practitioner Electronic Endpoints}

\label{element:11962189d1734cca840842209613c8d1}
\index{Bind Practitioner Electronic Endpoints (III.4.2.6.2.1)}
{\small BusinessFunction · Canonical · current}\par

Associates secure messaging addresses, direct communication endpoints, and notification preferences to practitioner roles.


{\footnotesize \href{../../../../continuum/business/business-function/11962189d1734cca840842209613c8d1.md}{Canonical record: 11962189d1734cca840842209613c8d1}; full qualifications and provenance.}\par
\paragraph*{III.4.2.6.2.2 Govern Practitioner Profile}
\phantomsection
\label{navigation:III.4.2.6.2.2}
\addcontentsline{toc}{paragraph}{III.4.2.6.2.2 Govern Practitioner Profile}

\label{element:a9a1a2c5ab544c029ba0c362a305eaa5}
\index{Govern Practitioner Profile (III.4.2.6.2.2)}
{\small BusinessFunction · Canonical · current}\par

Maintains official names, professional identifiers (e.g., HPI-I, Prescriber Number), and clinical qualifications.


{\footnotesize \href{../../../../continuum/business/business-function/a9a1a2c5ab544c029ba0c362a305eaa5.md}{Canonical record: a9a1a2c5ab544c029ba0c362a305eaa5}; full qualifications and provenance.}\par
\paragraph*{III.4.2.6.2.3 Maintain Clinical Privileges}
\phantomsection
\label{navigation:III.4.2.6.2.3}
\addcontentsline{toc}{paragraph}{III.4.2.6.2.3 Maintain Clinical Privileges}

\label{element:4965ae4538cc42cd8af1bd9a1905b108}
\index{Maintain Clinical Privileges (III.4.2.6.2.3)}
{\small BusinessFunction · Canonical · current}\par

Records approved clinical scopes of practice, procedural authorizations, and prescribing privileges.


{\footnotesize \href{../../../../continuum/business/business-function/4965ae4538cc42cd8af1bd9a1905b108.md}{Canonical record: 4965ae4538cc42cd8af1bd9a1905b108}; full qualifications and provenance.}\par
\paragraph*{III.4.2.6.2.4 Maintain Practitioner Affiliations}
\phantomsection
\label{navigation:III.4.2.6.2.4}
\addcontentsline{toc}{paragraph}{III.4.2.6.2.4 Maintain Practitioner Affiliations}

\label{element:7dc829546e99424dbbf8f19b4565d3fc}
\index{Maintain Practitioner Affiliations (III.4.2.6.2.4)}
{\small BusinessFunction · Canonical · current}\par

Tracks employment, admitting rights, and consulting appointments at specific healthcare facilities.


{\footnotesize \href{../../../../continuum/business/business-function/7dc829546e99424dbbf8f19b4565d3fc.md}{Canonical record: 7dc829546e99424dbbf8f19b4565d3fc}; full qualifications and provenance.}\par
\paragraph*{III.4.2.6.2.5 Maintain Practitioner Roles}
\phantomsection
\label{navigation:III.4.2.6.2.5}
\addcontentsline{toc}{paragraph}{III.4.2.6.2.5 Maintain Practitioner Roles}

\label{element:08a269d84dad46a58d6b9181b465ddc4}
\index{Maintain Practitioner Roles (III.4.2.6.2.5)}
{\small BusinessFunction · Canonical · current}\par

Assigns and maintains clinical roles (e.g., Specialist Physician, Registrar, General Practitioner) fulfilled by a practitioner.


{\footnotesize \href{../../../../continuum/business/business-function/08a269d84dad46a58d6b9181b465ddc4.md}{Canonical record: 08a269d84dad46a58d6b9181b465ddc4}; full qualifications and provenance.}\par
\paragraph*{III.4.2.6.2.6 Verify Practitioner Registration}
\phantomsection
\label{navigation:III.4.2.6.2.6}
\addcontentsline{toc}{paragraph}{III.4.2.6.2.6 Verify Practitioner Registration}

\label{element:a9f4fe6090bf4a4b946ae7064213945c}
\index{Verify Practitioner Registration (III.4.2.6.2.6)}
{\small BusinessFunction · Canonical · current}\par

Validates national registration, specialty endorsements, and disciplinary sanctions against regulatory registries (e.g., AHPRA).


{\footnotesize \href{../../../../continuum/business/business-function/a9f4fe6090bf4a4b946ae7064213945c.md}{Canonical record: a9f4fe6090bf4a4b946ae7064213945c}; full qualifications and provenance.}\par
\subsubsection*{III.4.2.6.3 Business Processes}
\phantomsection
\label{navigation:III.4.2.6.3}
\addcontentsline{toc}{subsubsection}{III.4.2.6.3 Business Processes}

\paragraph*{III.4.2.6.3.1 Practitioner Verification Process}
\phantomsection
\label{navigation:III.4.2.6.3.1}
\addcontentsline{toc}{paragraph}{III.4.2.6.3.1 Practitioner Verification Process}

\label{element:32a3e1f948474d3aa77ca71bfa515f58}
\index{Practitioner Verification Process (III.4.2.6.3.1)}
{\small BusinessProcess · Canonical · current}\par

•
Owning Capability: Provider Administration



•
Process Purpose: Governs the lifecycle of validating, certifying, and periodically re-verifying a healthcare practitioner's professional credentials, registration status, and clinical scope of practice.



•
State Progression Lifecycle: Verification Initiated → Regulatory Registry Queried → Credentials Validated → Privileges Endorsed → Active Verified → Expired / Suspended



•
Key State Dispositions:



•
Verification Initiated: New practitioner registration or periodic re-credentialing triggered.



•
Regulatory Registry Queried: Real-time query executed against national registry (e.g., AHPRA).



•
Credentials Validated: Primary qualifications, specialist endorsements, and sanctions checked.



•
Privileges Endorsed: Health service medical administration endorses admitting/procedural scope.



•
Active Verified: Practitioner verified for clinical ordering and electronic communication.



•
Expired / Suspended: Verification lapses or regulatory sanction triggers privilege suspension.



Verification progresses from checking the recognised professional source to the locally applicable endorsement and continuing verification status. The distinction matters because a nationally identifiable practitioner is not automatically privileged for every service activity. The illustrated progression retains the health service's endorsement responsibility rather than letting directory ingestion grant clinical authority.


{\small Recorded unresolved matters: source-qualified matters; exact scope remains in the canonical record.}\par
{\footnotesize \href{../../../../continuum/business/business-process/32a3e1f948474d3aa77ca71bfa515f58.md}{Canonical record: 32a3e1f948474d3aa77ca71bfa515f58}; full qualifications and provenance.}\par
\section*{III.4.3 [CT1] Health Informatics Management — Capability context}
\phantomsection
\label{navigation:III.4.3}
\addcontentsline{toc}{section}{III.4.3 [CT1] Health Informatics Management — Capability context}

{\small Canonical Strategy Capability: II.1.3.5 \hyperref[element:a790b5f2c5a1497fa6ec09d7824da32e]{[CT1] Health Informatics Management}; Canonical ID: \texttt{\seqsplit{a790b5f2c5a1497fa6ec09d7824da32e}}. Proposed · current.}\par
Contextual reference only; no additional architectural instance or semantic relationship.\par
\subsection*{III.4.3.1 [CT2] Research Information Curation — Capability context}
\phantomsection
\label{navigation:III.4.3.1}
\addcontentsline{toc}{subsection}{III.4.3.1 [CT2] Research Information Curation — Capability context}

{\small Canonical Strategy Capability: II.1.3.5.2 \hyperref[element:4f5b8cec90324c2989cdc30e80bc8cb9]{[CT2] Research Information Curation}; Canonical ID: \texttt{\seqsplit{4f5b8cec90324c2989cdc30e80bc8cb9}}. Proposed · current.}\par
Contextual reference only; no additional architectural instance or semantic relationship.\par
\subsubsection*{III.4.3.1.1 Business Services}
\phantomsection
\label{navigation:III.4.3.1.1}
\addcontentsline{toc}{subsubsection}{III.4.3.1.1 Business Services}

\paragraph*{III.4.3.1.1.1 Authorised Research Information Export}
\phantomsection
\label{navigation:III.4.3.1.1.1}
\addcontentsline{toc}{paragraph}{III.4.3.1.1.1 Authorised Research Information Export}

\label{element:134e1ff0331040b3a8069c9dc5e331d4}
\index{Authorised Research Information Export (III.4.3.1.1.1)}
{\small BusinessService · Proposed · current}\par

Provide an approved research consumer with purpose-qualified existing health information under its release contract.



\textbf{Overview}



The service exposes the permitted release and its applicable use basis rather than unrestricted clinical repository access.



\textbf{Boundary}



No new research-governance mandate or universal release format is allocated.


{\footnotesize \href{../../../../continuum/business/business-service/134e1ff0331040b3a8069c9dc5e331d4.md}{Canonical record: 134e1ff0331040b3a8069c9dc5e331d4}; full qualifications and provenance.}\par
\subsubsection*{III.4.3.1.2 Business Functions}
\phantomsection
\label{navigation:III.4.3.1.2}
\addcontentsline{toc}{subsubsection}{III.4.3.1.2 Business Functions}

\paragraph*{III.4.3.1.2.1 Assemble Purpose-Qualified Research Release}
\phantomsection
\label{navigation:III.4.3.1.2.1}
\addcontentsline{toc}{paragraph}{III.4.3.1.2.1 Assemble Purpose-Qualified Research Release}

\label{element:bf4354a421174a9ba56c281f4dcd11f6}
\index{Assemble Purpose-Qualified Research Release (III.4.3.1.2.1)}
{\small BusinessFunction · Proposed · current}\par

Construct a release containing only existing information permitted for the explicitly authorised research purpose.



\textbf{Overview}



Moves a qualified request to a bounded release with source attribution and applicable use constraints.



\textbf{Boundary}



Does not author ethics approvals, promise de-identification or acquire ownership of research conclusions.


{\footnotesize \href{../../../../continuum/business/business-function/bf4354a421174a9ba56c281f4dcd11f6.md}{Canonical record: bf4354a421174a9ba56c281f4dcd11f6}; full qualifications and provenance.}\par
\section*{III.4.4 [CT1] Health Service Operations Enablement — Capability context}
\phantomsection
\label{navigation:III.4.4}
\addcontentsline{toc}{section}{III.4.4 [CT1] Health Service Operations Enablement — Capability context}

{\small Canonical Strategy Capability: II.1.3.6 \hyperref[element:b13ee2c5a4ef4fd6be8a44795696dfc2]{[CT1] Health Service Operations Enablement}; Canonical ID: \texttt{\seqsplit{b13ee2c5a4ef4fd6be8a44795696dfc2}}. Proposed · current.}\par
Contextual reference only; no additional architectural instance or semantic relationship.\par
\subsection*{III.4.4.1 [CT2] Bed \& Care-Place Management — Capability context}
\phantomsection
\label{navigation:III.4.4.1}
\addcontentsline{toc}{subsection}{III.4.4.1 [CT2] Bed \& Care-Place Management — Capability context}

{\small Canonical Strategy Capability: II.1.3.6.1 \hyperref[element:f117ccec0bfa4f738e6ddb646b3534a2]{[CT2] Bed \& Care-Place Management}; Canonical ID: \texttt{\seqsplit{f117ccec0bfa4f738e6ddb646b3534a2}}. Canonical · current.}\par
Contextual reference only; no additional architectural instance or semantic relationship.\par
\subsubsection*{III.4.4.1.1 Business Services}
\phantomsection
\label{navigation:III.4.4.1.1}
\addcontentsline{toc}{subsubsection}{III.4.4.1.1 Business Services}

\paragraph*{III.4.4.1.1.1 Bed Availability Query}
\phantomsection
\label{navigation:III.4.4.1.1.1}
\addcontentsline{toc}{paragraph}{III.4.4.1.1.1 Bed Availability Query}

\label{element:c51894eccd7b45019c2f084993f4cfc6}
\index{Bed Availability Query (III.4.4.1.1.1)}
{\small BusinessService · Canonical · current}\par

Discloses real-time bed availability across hospital directorates.


{\footnotesize \href{../../../../continuum/business/business-service/c51894eccd7b45019c2f084993f4cfc6.md}{Canonical record: c51894eccd7b45019c2f084993f4cfc6}; full qualifications and provenance.}\par
\paragraph*{III.4.4.1.1.2 Bed Turnover Lifecycle Service}
\phantomsection
\label{navigation:III.4.4.1.1.2}
\addcontentsline{toc}{paragraph}{III.4.4.1.1.2 Bed Turnover Lifecycle Service}

\label{element:f0ab3847b1bf41869a85b1a5ef217429}
\index{Bed Turnover Lifecycle Service (III.4.4.1.1.2)}
{\small BusinessService · Canonical · current}\par

Coordinates turnover requests with environmental services.



Additional Feature responsibility: FEAT-HSO-13 — Bed Readiness Notification requires the existing turnover behaviour/exposure to notify bed allocation managers and clinical units when a cleaned bed is ready. Readiness notification remains distinct from availability query and from physical placement; no new lifecycle state or Service is introduced.


{\footnotesize \href{../../../../continuum/business/business-service/f0ab3847b1bf41869a85b1a5ef217429.md}{Canonical record: f0ab3847b1bf41869a85b1a5ef217429}; full qualifications and provenance.}\par
\subsubsection*{III.4.4.1.2 Business Functions}
\phantomsection
\label{navigation:III.4.4.1.2}
\addcontentsline{toc}{subsubsection}{III.4.4.1.2 Business Functions}

\paragraph*{III.4.4.1.2.1 Maintain Operational Bed State}
\phantomsection
\label{navigation:III.4.4.1.2.1}
\addcontentsline{toc}{paragraph}{III.4.4.1.2.1 Maintain Operational Bed State}

\label{element:22f3016539cf4c2db68baa9921acf31b}
\index{Maintain Operational Bed State (III.4.4.1.2.1)}
{\small BusinessFunction · Canonical · current}\par

Tracks real-time bed status and environmental locks.


{\footnotesize \href{../../../../continuum/business/business-function/22f3016539cf4c2db68baa9921acf31b.md}{Canonical record: 22f3016539cf4c2db68baa9921acf31b}; full qualifications and provenance.}\par
\paragraph*{III.4.4.1.2.2 Manage Bed Turnover Progression}
\phantomsection
\label{navigation:III.4.4.1.2.2}
\addcontentsline{toc}{paragraph}{III.4.4.1.2.2 Manage Bed Turnover Progression}

\label{element:7637b9bcb38d4e3abbe01a59dc23f485}
\index{Manage Bed Turnover Progression (III.4.4.1.2.2)}
{\small BusinessFunction · Canonical · current}\par

Coordinates cleaning, disinfection, and preparation: Vacated → Cleaning Requested → Cleaning In-Progress → Inspected/Ready → Occupied.



Additional Feature responsibility: FEAT-HSO-13 — Bed Readiness Notification requires the existing turnover behaviour/exposure to notify bed allocation managers and clinical units when a cleaned bed is ready. Readiness notification remains distinct from availability query and from physical placement; no new lifecycle state or Service is introduced.


{\footnotesize \href{../../../../continuum/business/business-function/7637b9bcb38d4e3abbe01a59dc23f485.md}{Canonical record: 7637b9bcb38d4e3abbe01a59dc23f485}; full qualifications and provenance.}\par
\subsubsection*{III.4.4.1.3 Business Processes}
\phantomsection
\label{navigation:III.4.4.1.3}
\addcontentsline{toc}{subsubsection}{III.4.4.1.3 Business Processes}

\paragraph*{III.4.4.1.3.1 Bed Turnover Process}
\phantomsection
\label{navigation:III.4.4.1.3.1}
\addcontentsline{toc}{paragraph}{III.4.4.1.3.1 Bed Turnover Process}

\label{element:69f74289926942f4baed4acf4c654a3b}
\index{Bed Turnover Process (III.4.4.1.3.1)}
{\small BusinessProcess · Canonical · current}\par

•
Owning Capability: Bed \& Care-Place Management



•
Process Purpose: Governs the rapid physical turnover, cleaning, sanitisation, and re-allocation of inpatient and emergency care-places.



•
State Progression Lifecycle: Bed Vacated → Cleaning Dispatched → Sanitisation In-Progress → Terminal Cleaning Verified → Bed Ready / Available → Bed Allocated / Occupied



•
Key State Dispositions:



•
Bed Vacated: Prior patient discharged or transferred.



•
Cleaning Dispatched: Cleaning job allocated to environmental services team.



•
Sanitisation In-Progress: Ward cleaning or isolation decontamination underway.



•
Terminal Cleaning Verified: Ward coordinator or supervisor verifies care-place readiness.



•
Bed Ready / Available: Bed marked available for incoming patient allocation.



•
Bed Allocated / Occupied: Incoming patient allocated and physically placed in bed.



Bed turnover connects the release of a care place with the operational work needed to return it to use. Its progression preserves the difference between vacancy, cleaning and readiness. A place can be empty while work remains outstanding, and a completed turnover does not by itself prove that the service has capacity for the next patient.


{\small Recorded unresolved matters: source-qualified matters; exact scope remains in the canonical record.}\par
{\footnotesize \href{../../../../continuum/business/business-process/69f74289926942f4baed4acf4c654a3b.md}{Canonical record: 69f74289926942f4baed4acf4c654a3b}; full qualifications and provenance.}\par
\subsection*{III.4.4.2 [CT2] Care Resource \& Capacity Coordination — Capability context [superseded]}
\phantomsection
\label{navigation:III.4.4.2}
\addcontentsline{toc}{subsection}{III.4.4.2 [CT2] Care Resource \& Capacity Coordination — Capability context [superseded]}

{\small Canonical Strategy Capability: II.1.3.6.3 \hyperref[element:96f71cb8fa6844d7aa7e9f16759a3a62]{[CT2] Care Resource \& Capacity Coordination}; Canonical ID: \texttt{\seqsplit{96f71cb8fa6844d7aa7e9f16759a3a62}}. Proposed · superseded.}\par
Contextual reference only; no additional architectural instance or semantic relationship.\par
\subsubsection*{III.4.4.2.1 [CT3] Service Offering Readiness Coordination — Capability context}
\phantomsection
\label{navigation:III.4.4.2.1}
\addcontentsline{toc}{subsubsection}{III.4.4.2.1 [CT3] Service Offering Readiness Coordination — Capability context}

{\small Canonical Strategy Capability: II.1.3.6.3.2 \hyperref[element:0caecaba5d3e4781b8671b12c13ee283]{[CT3] Service Offering Readiness Coordination}; Canonical ID: \texttt{\seqsplit{0caecaba5d3e4781b8671b12c13ee283}}. Proposed · pending\_adjudication.}\par
Contextual reference only; no additional architectural instance or semantic relationship.\par
\paragraph*{III.4.4.2.1.1 Business Services}
\phantomsection
\label{navigation:III.4.4.2.1.1}
\addcontentsline{toc}{paragraph}{III.4.4.2.1.1 Business Services}

\subparagraph*{III.4.4.2.1.1.1 Service Offering Readiness Query}
\phantomsection
\label{navigation:III.4.4.2.1.1.1}
\addcontentsline{toc}{subparagraph}{III.4.4.2.1.1.1 Service Offering Readiness Query}

\label{element:3b0bdc7bc18940208511f39e43d5357a}
\index{Service Offering Readiness Query (III.4.4.2.1.1.1)}
{\small BusinessService · Proposed · pending\_adjudication}\par

Expose the qualified readiness view of a defined offering and operating context to operational planning consumers.



\textbf{Overview}



The response identifies relevant evidence basis and insufficiency rather than an unconditional ready flag.



\textbf{Boundary}



A readiness answer grants no resource allocation or clinical go/no-go authority.


{\small Recorded unresolved matters: source-qualified matters; exact scope remains in the canonical record.}\par
{\footnotesize \href{../../../../continuum/business/business-service/3b0bdc7bc18940208511f39e43d5357a.md}{Canonical record: 3b0bdc7bc18940208511f39e43d5357a}; full qualifications and provenance.}\par
\paragraph*{III.4.4.2.1.2 Business Functions}
\phantomsection
\label{navigation:III.4.4.2.1.2}
\addcontentsline{toc}{paragraph}{III.4.4.2.1.2 Business Functions}

\subparagraph*{III.4.4.2.1.2.1 Assess Service Offering Readiness}
\phantomsection
\label{navigation:III.4.4.2.1.2.1}
\addcontentsline{toc}{subparagraph}{III.4.4.2.1.2.1 Assess Service Offering Readiness}

\label{element:a7eb76c1ff5b450b9835ac19662c8148}
\index{Assess Service Offering Readiness (III.4.4.2.1.2.1)}
{\small BusinessFunction · Proposed · pending\_adjudication}\par

Evaluate applicable dependency evidence to establish a context- and time-qualified service readiness view.



\textbf{Overview}



Combines qualified dependency standing without converting missing evidence into ready or assuming that an available resource establishes deliverability.



\textbf{Boundary}



Resource ownership and clinical clearance remain external to this coordination.


{\small Recorded unresolved matters: source-qualified matters; exact scope remains in the canonical record.}\par
{\footnotesize \href{../../../../continuum/business/business-function/a7eb76c1ff5b450b9835ac19662c8148.md}{Canonical record: a7eb76c1ff5b450b9835ac19662c8148}; full qualifications and provenance.}\par
\subsection*{III.4.4.3 [CT2] Clinic \& Practice Operations — Capability context}
\phantomsection
\label{navigation:III.4.4.3}
\addcontentsline{toc}{subsection}{III.4.4.3 [CT2] Clinic \& Practice Operations — Capability context}

{\small Canonical Strategy Capability: II.1.3.6.5 \hyperref[element:3dffe63e03124efab255204ee4e9a5d9]{[CT2] Clinic \& Practice Operations}; Canonical ID: \texttt{\seqsplit{3dffe63e03124efab255204ee4e9a5d9}}. Canonical · current.}\par
Contextual reference only; no additional architectural instance or semantic relationship.\par
\subsubsection*{III.4.4.3.1 Business Services}
\phantomsection
\label{navigation:III.4.4.3.1}
\addcontentsline{toc}{subsubsection}{III.4.4.3.1 Business Services}

\paragraph*{III.4.4.3.1.1 Clinic Queue Telemetry Service}
\phantomsection
\label{navigation:III.4.4.3.1.1}
\addcontentsline{toc}{paragraph}{III.4.4.3.1.1 Clinic Queue Telemetry Service}

\label{element:c720457349e644179c98d13be59bd581}
\index{Clinic Queue Telemetry Service (III.4.4.3.1.1)}
{\small BusinessService · Canonical · current}\par

Discloses real-time clinic waiting times and queue states.


{\footnotesize \href{../../../../continuum/business/business-service/c720457349e644179c98d13be59bd581.md}{Canonical record: c720457349e644179c98d13be59bd581}; full qualifications and provenance.}\par
\subsubsection*{III.4.4.3.2 Business Functions}
\phantomsection
\label{navigation:III.4.4.3.2}
\addcontentsline{toc}{subsubsection}{III.4.4.3.2 Business Functions}

\paragraph*{III.4.4.3.2.1 Manage Clinic \& Practice Capacity}
\phantomsection
\label{navigation:III.4.4.3.2.1}
\addcontentsline{toc}{paragraph}{III.4.4.3.2.1 Manage Clinic \& Practice Capacity}

\label{element:389fb97b2bcd4527959a61f2a288bc39}
\index{Manage Clinic \& Practice Capacity (III.4.4.3.2.1)}
{\small BusinessFunction · Canonical · current}\par

Understands and manages available, committed and utilised clinic/practice capacity in relation to practitioner availability, sessions, rooms, appointments and patient demand.



Responsibility boundary: Operating-session tracking and individual queue progression contribute to this responsibility; neither alone expresses overall service capacity. Capacity coordination does not own source bookings, clinical decisions or practitioner rostering.


{\small Recorded unresolved matters: source-qualified matters; exact scope remains in the canonical record.}\par
{\footnotesize \href{../../../../continuum/business/business-function/389fb97b2bcd4527959a61f2a288bc39.md}{Canonical record: 389fb97b2bcd4527959a61f2a288bc39}; full qualifications and provenance.}\par
\paragraph*{III.4.4.3.2.2 Manage Clinic Queue Progression}
\phantomsection
\label{navigation:III.4.4.3.2.2}
\addcontentsline{toc}{paragraph}{III.4.4.3.2.2 Manage Clinic Queue Progression}

\label{element:cb3f8c48623b4479988c45de86d5a59f}
\index{Manage Clinic Queue Progression (III.4.4.3.2.2)}
{\small BusinessFunction · Canonical · current}\par

Tracks patient arrival, rooming, consultation in-progress, and check-out states.


{\footnotesize \href{../../../../continuum/business/business-function/cb3f8c48623b4479988c45de86d5a59f.md}{Canonical record: cb3f8c48623b4479988c45de86d5a59f}; full qualifications and provenance.}\par
\paragraph*{III.4.4.3.2.3 Track Clinic Operating Session}
\phantomsection
\label{navigation:III.4.4.3.2.3}
\addcontentsline{toc}{paragraph}{III.4.4.3.2.3 Track Clinic Operating Session}

\label{element:83a85a08803c4c00b8de4a38e8c9d77e}
\index{Track Clinic Operating Session (III.4.4.3.2.3)}
{\small BusinessFunction · Canonical · current}\par

Ingests and tracks clinic operating-session start, delay and completion, distinct from an individual patient's queue progression.



Derivation: FEAT-HSO-01 directly establishes this session-level behaviour; Manage Clinic Queue Progression and its Process concern individual patients and do not supply it.


{\footnotesize \href{../../../../continuum/business/business-function/83a85a08803c4c00b8de4a38e8c9d77e.md}{Canonical record: 83a85a08803c4c00b8de4a38e8c9d77e}; full qualifications and provenance.}\par
\subsubsection*{III.4.4.3.3 Business Processes}
\phantomsection
\label{navigation:III.4.4.3.3}
\addcontentsline{toc}{subsubsection}{III.4.4.3.3 Business Processes}

\paragraph*{III.4.4.3.3.1 Clinic / Operational Queue Progression Process}
\phantomsection
\label{navigation:III.4.4.3.3.1}
\addcontentsline{toc}{paragraph}{III.4.4.3.3.1 Clinic / Operational Queue Progression Process}

\label{element:319d6956276a4cdeb43e350e1b01d078}
\index{Clinic / Operational Queue Progression Process (III.4.4.3.3.1)}
{\small BusinessProcess · Canonical · current}\par

•
Owning Capability: Clinic \& Practice Operations



•
Process Purpose: Governs the movement and state tracking of outpatients through an ambulatory specialty clinic session.



•
State Progression Lifecycle: Scheduled → Patient Arrived → Checked-In → Roomed / Pre-Consult → Consultation In-Progress → Consultation Concluded → Departed



•
Key State Dispositions:



•
Scheduled: Appointment booked on session list.



•
Patient Arrived: Patient enters waiting area.



•
Checked-In: Reception verifies demographics and marks patient ready.



•
Roomed / Pre-Consult: Nursing staff record baseline vitals and room patient.



•
Consultation In-Progress: Clinician actively conducting consultation.



•
Consultation Concluded: Clinical discussion finished; follow-up orders placed.



•
Departed: Patient checks out and departs clinic facility.



This Process follows an individual's operational passage through the clinic. Waiting, rooming, consultation and departure have different implications for flow and communication, while the clinic's overall session has its own state. The progression is therefore patient-specific operational context rather than a complete account of clinic capacity or clinical priority.


{\small Recorded unresolved matters: source-qualified matters; exact scope remains in the canonical record.}\par
{\footnotesize \href{../../../../continuum/business/business-process/319d6956276a4cdeb43e350e1b01d078.md}{Canonical record: 319d6956276a4cdeb43e350e1b01d078}; full qualifications and provenance.}\par
\subsection*{III.4.4.4 [CT2] Clinical Logistics Coordination — Capability context}
\phantomsection
\label{navigation:III.4.4.4}
\addcontentsline{toc}{subsection}{III.4.4.4 [CT2] Clinical Logistics Coordination — Capability context}

{\small Canonical Strategy Capability: II.1.3.6.7 \hyperref[element:0c27be11b4a44cb19e992309b8146a17]{[CT2] Clinical Logistics Coordination}; Canonical ID: \texttt{\seqsplit{0c27be11b4a44cb19e992309b8146a17}}. Canonical · current.}\par
Contextual reference only; no additional architectural instance or semantic relationship.\par
\subsubsection*{III.4.4.4.1 Business Services}
\phantomsection
\label{navigation:III.4.4.4.1}
\addcontentsline{toc}{subsubsection}{III.4.4.4.1 Business Services}

\paragraph*{III.4.4.4.1.1 Specimen Transit Tracking Service}
\phantomsection
\label{navigation:III.4.4.4.1.1}
\addcontentsline{toc}{paragraph}{III.4.4.4.1.1 Specimen Transit Tracking Service}

\label{element:100e740770bc4ab28cfe3c9af4b541cd}
\index{Specimen Transit Tracking Service (III.4.4.4.1.1)}
{\small BusinessService · Canonical · current}\par

Exposes chain-of-custody tracking.


{\footnotesize \href{../../../../continuum/business/business-service/100e740770bc4ab28cfe3c9af4b541cd.md}{Canonical record: 100e740770bc4ab28cfe3c9af4b541cd}; full qualifications and provenance.}\par
\subsubsection*{III.4.4.4.2 Business Functions}
\phantomsection
\label{navigation:III.4.4.4.2}
\addcontentsline{toc}{subsubsection}{III.4.4.4.2 Business Functions}

\paragraph*{III.4.4.4.2.1 Manage Specimen Transport Progression}
\phantomsection
\label{navigation:III.4.4.4.2.1}
\addcontentsline{toc}{paragraph}{III.4.4.4.2.1 Manage Specimen Transport Progression}

\label{element:4c9d50efc9e640fcb3abb02b8dd05486}
\index{Manage Specimen Transport Progression (III.4.4.4.2.1)}
{\small BusinessFunction · Canonical · pending\_adjudication}\par

Tracks the physical transit of pathology bio-specimens, blood products, and surgical trays: Collected → Courier Picked Up → In-Transit → Laboratory Ingress Received.


{\small Recorded unresolved matters: source-qualified matters; exact scope remains in the canonical record.}\par
{\footnotesize \href{../../../../continuum/business/business-function/4c9d50efc9e640fcb3abb02b8dd05486.md}{Canonical record: 4c9d50efc9e640fcb3abb02b8dd05486}; full qualifications and provenance.}\par
\subsubsection*{III.4.4.4.3 Business Processes}
\phantomsection
\label{navigation:III.4.4.4.3}
\addcontentsline{toc}{subsubsection}{III.4.4.4.3 Business Processes}

\paragraph*{III.4.4.4.3.1 Specimen Transport Process}
\phantomsection
\label{navigation:III.4.4.4.3.1}
\addcontentsline{toc}{paragraph}{III.4.4.4.3.1 Specimen Transport Process}

\label{element:11bb000294fc44a8af295f4682e2fe6c}
\index{Specimen Transport Process (III.4.4.4.3.1)}
{\small BusinessProcess · Canonical · pending\_adjudication}\par

•
Owning Capability: Clinical Logistics Coordination



•
Process Purpose: Governs the physical courier chain-of-custody, transport, and delivery of pathology bio-specimens, blood products, and surgical biopsies.



•
State Progression Lifecycle: Specimen Packaged → Courier Dispatched → Chain-of-Custody Signed → In-Transit (Monitored) → Lab Ingress Received → Specimen Accepted



•
Key State Dispositions:



•
Specimen Packaged: Clinical unit packages bio-sample with barcode manifest.



•
Courier Dispatched: Courier route assigned for pickup.



•
Chain-of-Custody Signed: Courier scans barcode and accepts legal custody.



•
In-Transit (Monitored): Specimen in transit under temperature/time constraints.



•
Lab Ingress Received: Central pathology specimen reception logs physical receipt.



•
Specimen Accepted: Specimen integrity verified and passed to analytical track.



\medskip\hrule\medskip

{\small Recorded unresolved matters: source-qualified matters; exact scope remains in the canonical record.}\par
{\footnotesize \href{../../../../continuum/business/business-process/11bb000294fc44a8af295f4682e2fe6c.md}{Canonical record: 11bb000294fc44a8af295f4682e2fe6c}; full qualifications and provenance.}\par
\subsection*{III.4.4.5 [CT2] Clinical Qualification Administration — Capability context}
\phantomsection
\label{navigation:III.4.4.5}
\addcontentsline{toc}{subsection}{III.4.4.5 [CT2] Clinical Qualification Administration — Capability context}

{\small Canonical Strategy Capability: II.1.3.6.8 \hyperref[element:0511a705d74a4f48a0ce92590ae764fa]{[CT2] Clinical Qualification Administration}; Canonical ID: \texttt{\seqsplit{0511a705d74a4f48a0ce92590ae764fa}}. Canonical · current.}\par
Contextual reference only; no additional architectural instance or semantic relationship.\par
\subsubsection*{III.4.4.5.1 Business Functions}
\phantomsection
\label{navigation:III.4.4.5.1}
\addcontentsline{toc}{subsubsection}{III.4.4.5.1 Business Functions}

\paragraph*{III.4.4.5.1.1 Verify Clinical Qualification for Work Dispatch}
\phantomsection
\label{navigation:III.4.4.5.1.1}
\addcontentsline{toc}{paragraph}{III.4.4.5.1.1 Verify Clinical Qualification for Work Dispatch}

\label{element:1dfebc3aa5f949e5a20faa504eadbee1}
\index{Verify Clinical Qualification for Work Dispatch (III.4.4.5.1.1)}
{\small BusinessFunction · Canonical · current}\par

Validates that a practitioner possesses active clinical credentials and required competency certificates before dispatch of specialised work.



Derivation: FEAT-HSO-23 belongs to Strategy Clinical Qualification Management. Provider Administration establishes professional qualification/registration and privilege information, but the dispatch-specific qualification/competency checkpoint was not expressed by that information management alone.


{\footnotesize \href{../../../../continuum/business/business-function/1dfebc3aa5f949e5a20faa504eadbee1.md}{Canonical record: 1dfebc3aa5f949e5a20faa504eadbee1}; full qualifications and provenance.}\par
\subsection*{III.4.4.6 [CT2] Discharge Management — Capability context}
\phantomsection
\label{navigation:III.4.4.6}
\addcontentsline{toc}{subsection}{III.4.4.6 [CT2] Discharge Management — Capability context}

{\small Canonical Strategy Capability: II.1.3.6.14 \hyperref[element:ec5e8d6732764aa489c2dd7630fc8fef]{[CT2] Discharge Management}; Canonical ID: \texttt{\seqsplit{ec5e8d6732764aa489c2dd7630fc8fef}}. Canonical · current.}\par
Contextual reference only; no additional architectural instance or semantic relationship.\par
\subsubsection*{III.4.4.6.1 Business Services}
\phantomsection
\label{navigation:III.4.4.6.1}
\addcontentsline{toc}{subsubsection}{III.4.4.6.1 Business Services}

\paragraph*{III.4.4.6.1.1 Discharge Readiness Telemetry}
\phantomsection
\label{navigation:III.4.4.6.1.1}
\addcontentsline{toc}{paragraph}{III.4.4.6.1.1 Discharge Readiness Telemetry}

\label{element:03d01a9b024c4497badf3fe011a8dd20}
\index{Discharge Readiness Telemetry (III.4.4.6.1.1)}
{\small BusinessService · Canonical · current}\par

Discloses discharge barrier checklists and planned departure times.


{\footnotesize \href{../../../../continuum/business/business-service/03d01a9b024c4497badf3fe011a8dd20.md}{Canonical record: 03d01a9b024c4497badf3fe011a8dd20}; full qualifications and provenance.}\par
\subsubsection*{III.4.4.6.2 Business Functions}
\phantomsection
\label{navigation:III.4.4.6.2}
\addcontentsline{toc}{subsubsection}{III.4.4.6.2 Business Functions}

\paragraph*{III.4.4.6.2.1 Manage Discharge Coordination Progression}
\phantomsection
\label{navigation:III.4.4.6.2.1}
\addcontentsline{toc}{paragraph}{III.4.4.6.2.1 Manage Discharge Coordination Progression}

\label{element:ecda6d08166d4df78c15c87387fbe108}
\index{Manage Discharge Coordination Progression (III.4.4.6.2.1)}
{\small BusinessFunction · Canonical · current}\par

Synchronises multi-agency discharge readiness (illustrative preparation; publication, authorisation and confirmed exit are distinct): Discharge Planning Initiated → Medications Reconciled → Transport Booked → Discharge Summary Finalised → Physically Departed.


{\footnotesize \href{../../../../continuum/business/business-function/ecda6d08166d4df78c15c87387fbe108.md}{Canonical record: ecda6d08166d4df78c15c87387fbe108}; full qualifications and provenance.}\par
\subsubsection*{III.4.4.6.3 Business Processes}
\phantomsection
\label{navigation:III.4.4.6.3}
\addcontentsline{toc}{subsubsection}{III.4.4.6.3 Business Processes}

\paragraph*{III.4.4.6.3.1 Discharge Progression Process}
\phantomsection
\label{navigation:III.4.4.6.3.1}
\addcontentsline{toc}{paragraph}{III.4.4.6.3.1 Discharge Progression Process}

\label{element:7165955d9c454175a04b9ad00293f011}
\index{Discharge Progression Process (III.4.4.6.3.1)}
{\small BusinessProcess · Canonical · current}\par

•
Owning Capability: Discharge Management



•
Process Purpose: Governs the multidisciplinary coordination of patient discharge planning, pharmacy reconciliation, transport, and community handover.



•
State Progression Lifecycle: Discharge Planning Initiated → Clinical Readiness Confirmed → Medications Reconciled → Transport \& Services Booked → Discharge Summary Signed → Physically Departed   (Illustrative preparation; publication, authorisation and exit distinct)



•
Key State Dispositions:



•
Discharge Planning Initiated: Estimated Date of Discharge (EDD) set on admission.



•
Clinical Readiness Confirmed: Attending medical team declares patient fit for discharge.



•
Medications Reconciled: Hospital pharmacy reconciles and dispenses discharge medications.



•
Transport \& Services Booked: Community nursing, home equipment, and patient transport booked.



•
Discharge Summary Signed: Applicable signed summary is available/published and may be communicated before departure where the Business behaviour requires it; signing, finalisation, publication and discharge authorisation remain distinct.



•
Physically Departed: Confirmed physical exit triggers departure notifications, updated bed-state communication and finalised-summary dispatch to external care providers under FEAT-HSO-29; care-place turnover follows its own responsibility.



•
Timing reconciliation: Preparation, information availability/publication, applicable discharge authorisation and physical exit are distinct. Earlier preparation/communication is permitted; it does not realise a Feature triggered by confirmed exit. This illustrated local pathway establishes neither universal publication timing nor a universal discharge state machine, and adds no inferred publication event or authorisation state.



•
Preserved uncertainty: Summary signing and finalisation are not universal equivalents. Detailed applicability, discharge-authority workflows and timing beyond the established preparation/confirmed-exit distinction remain unestablished.



Discharge progression coordinates the confirmations and practical requirements around leaving care. Clinical readiness, medications, destination arrangements and transport are related but distinct contributions, and the source progression retains its stated local applicability. Operational coordination does not originate the clinical discharge decision or prove that another provider assumed responsibility.


{\small Recorded unresolved matters: source-qualified matters; exact scope remains in the canonical record.}\par
{\footnotesize \href{../../../../continuum/business/business-process/7165955d9c454175a04b9ad00293f011.md}{Canonical record: 7165955d9c454175a04b9ad00293f011}; full qualifications and provenance.}\par
\subsection*{III.4.4.7 [CT2] Emergency Department Operations — Capability context}
\phantomsection
\label{navigation:III.4.4.7}
\addcontentsline{toc}{subsection}{III.4.4.7 [CT2] Emergency Department Operations — Capability context}

{\small Canonical Strategy Capability: II.1.3.6.15 \hyperref[element:24a4b5815faa4d369bc756a6e97d3e9e]{[CT2] Emergency Department Operations}; Canonical ID: \texttt{\seqsplit{24a4b5815faa4d369bc756a6e97d3e9e}}. Canonical · current.}\par
Contextual reference only; no additional architectural instance or semantic relationship.\par
\subsubsection*{III.4.4.7.1 [CT3] Emergency Care Progression — Capability context}
\phantomsection
\label{navigation:III.4.4.7.1}
\addcontentsline{toc}{subsubsection}{III.4.4.7.1 [CT3] Emergency Care Progression — Capability context}

{\small Canonical Strategy Capability: II.1.3.6.15.1 \hyperref[element:7125e459a48946d0ba075b5ecc3571b4]{[CT3] Emergency Care Progression}; Canonical ID: \texttt{\seqsplit{7125e459a48946d0ba075b5ecc3571b4}}. Canonical · current.}\par
Contextual reference only; no additional architectural instance or semantic relationship.\par
\paragraph*{III.4.4.7.1.1 Business Services}
\phantomsection
\label{navigation:III.4.4.7.1.1}
\addcontentsline{toc}{paragraph}{III.4.4.7.1.1 Business Services}

\subparagraph*{III.4.4.7.1.1.1 ED Triage Stream Ingress}
\phantomsection
\label{navigation:III.4.4.7.1.1.1}
\addcontentsline{toc}{subparagraph}{III.4.4.7.1.1.1 ED Triage Stream Ingress}

\label{element:677947e357b541f58f36e9b4e36157e8}
\index{ED Triage Stream Ingress (III.4.4.7.1.1.1)}
{\small BusinessService · Canonical · current}\par

Ingests emergency triage presentations.


{\footnotesize \href{../../../../continuum/business/business-service/677947e357b541f58f36e9b4e36157e8.md}{Canonical record: 677947e357b541f58f36e9b4e36157e8}; full qualifications and provenance.}\par
\subparagraph*{III.4.4.7.1.1.2 Emergency Fast-Path LHR Query}
\phantomsection
\label{navigation:III.4.4.7.1.1.2}
\addcontentsline{toc}{subparagraph}{III.4.4.7.1.1.2 Emergency Fast-Path LHR Query}

\label{element:0588ef78691143cdb2d884960811db1e}
\index{Emergency Fast-Path LHR Query (III.4.4.7.1.1.2)}
{\small BusinessService · Canonical · current}\par

Provides low-latency emergency record summaries.


{\footnotesize \href{../../../../continuum/business/business-service/0588ef78691143cdb2d884960811db1e.md}{Canonical record: 0588ef78691143cdb2d884960811db1e}; full qualifications and provenance.}\par
\paragraph*{III.4.4.7.1.2 Business Functions}
\phantomsection
\label{navigation:III.4.4.7.1.2}
\addcontentsline{toc}{paragraph}{III.4.4.7.1.2 Business Functions}

\subparagraph*{III.4.4.7.1.2.1 Bind ED Triage Assessment}
\phantomsection
\label{navigation:III.4.4.7.1.2.1}
\addcontentsline{toc}{subparagraph}{III.4.4.7.1.2.1 Bind ED Triage Assessment}

\label{element:fa51fad9b9344546953ff8ec5b5f39ca}
\index{Bind ED Triage Assessment (III.4.4.7.1.2.1)}
{\small BusinessFunction · Canonical · current}\par

Captures ATS (Australasian Triage Scale) category, presenting complaint, and ambulance paramedic handover data.


{\footnotesize \href{../../../../continuum/business/business-function/fa51fad9b9344546953ff8ec5b5f39ca.md}{Canonical record: fa51fad9b9344546953ff8ec5b5f39ca}; full qualifications and provenance.}\par
\subparagraph*{III.4.4.7.1.2.2 Provide Emergency LHR Summary}
\phantomsection
\label{navigation:III.4.4.7.1.2.2}
\addcontentsline{toc}{subparagraph}{III.4.4.7.1.2.2 Provide Emergency LHR Summary}

\label{element:a7235e8b249a4914b951e0d795132316}
\index{Provide Emergency LHR Summary (III.4.4.7.1.2.2)}
{\small BusinessFunction · Canonical · current}\par

Delivers rapid, override-capable access to Advance Care Directives, allergy records, and recent hospital admissions.


{\footnotesize \href{../../../../continuum/business/business-function/a7235e8b249a4914b951e0d795132316.md}{Canonical record: a7235e8b249a4914b951e0d795132316}; full qualifications and provenance.}\par
\subsubsection*{III.4.4.7.2 Business Services}
\phantomsection
\label{navigation:III.4.4.7.2}
\addcontentsline{toc}{subsubsection}{III.4.4.7.2 Business Services}

\paragraph*{III.4.4.7.2.1 ED Operational Status Query}
\phantomsection
\label{navigation:III.4.4.7.2.1}
\addcontentsline{toc}{paragraph}{III.4.4.7.2.1 ED Operational Status Query}

\label{element:f07828653c06425f9effd92b1d9886c9}
\index{ED Operational Status Query (III.4.4.7.2.1)}
{\small BusinessService · Canonical · current}\par

Discloses facility ED occupancy and surge level.



Capacity boundary: Existing departmental capacity responsibility is clarified, not duplicated. Care-place state contributes to service capacity; capacity coordination neither performs clinical triage nor establishes clinical acuity assertions.


{\small Recorded unresolved matters: source-qualified matters; exact scope remains in the canonical record.}\par
{\footnotesize \href{../../../../continuum/business/business-service/f07828653c06425f9effd92b1d9886c9.md}{Canonical record: f07828653c06425f9effd92b1d9886c9}; full qualifications and provenance.}\par
\subsubsection*{III.4.4.7.3 Business Functions}
\phantomsection
\label{navigation:III.4.4.7.3}
\addcontentsline{toc}{subsubsection}{III.4.4.7.3 Business Functions}

\paragraph*{III.4.4.7.3.1 Coordinate ED Departmental Capacity}
\phantomsection
\label{navigation:III.4.4.7.3.1}
\addcontentsline{toc}{paragraph}{III.4.4.7.3.1 Coordinate ED Departmental Capacity}

\label{element:a1559e7ed12c4f249128fba2ff74d536}
\index{Coordinate ED Departmental Capacity (III.4.4.7.3.1)}
{\small BusinessFunction · Canonical · current}\par

Understands and manages available, committed and utilised ED capacity in relation to care-place availability, demand, patient presence/progression, acuity and operational constraints; retains waiting room occupancy, resuscitation bay availability and time-to-triage metrics.



Capacity boundary: Existing departmental capacity responsibility is clarified, not duplicated. Care-place state contributes to service capacity; capacity coordination neither performs clinical triage nor establishes clinical acuity assertions.


{\footnotesize \href{../../../../continuum/business/business-function/a1559e7ed12c4f249128fba2ff74d536.md}{Canonical record: a1559e7ed12c4f249128fba2ff74d536}; full qualifications and provenance.}\par
\subsection*{III.4.4.8 [CT2] Mobile Staff Management — Capability context}
\phantomsection
\label{navigation:III.4.4.8}
\addcontentsline{toc}{subsection}{III.4.4.8 [CT2] Mobile Staff Management — Capability context}

{\small Canonical Strategy Capability: II.1.3.6.16 \hyperref[element:79a4a48fddf74c69acf7e16404d06ea5]{[CT2] Mobile Staff Management}; Canonical ID: \texttt{\seqsplit{79a4a48fddf74c69acf7e16404d06ea5}}. Canonical · current.}\par
Contextual reference only; no additional architectural instance or semantic relationship.\par
\subsubsection*{III.4.4.8.1 Business Services}
\phantomsection
\label{navigation:III.4.4.8.1}
\addcontentsline{toc}{subsubsection}{III.4.4.8.1 Business Services}

\paragraph*{III.4.4.8.1.1 Staff Presence Telemetry Query}
\phantomsection
\label{navigation:III.4.4.8.1.1}
\addcontentsline{toc}{paragraph}{III.4.4.8.1.1 Staff Presence Telemetry Query}

\label{element:b3cd0d8e40f44b9fa2b64367237345fa}
\index{Staff Presence Telemetry Query (III.4.4.8.1.1)}
{\small BusinessService · Canonical · current}\par

Discloses on-duty staff availability by hospital zone.



Feature association not established: FEAT-HSO-16 — Mobile Worker Task Dispatch retains its Strategy identity with the integration boundary below. Presence/zone tracking remains distinct from mobile work-management integration; no replacement association to presence is assigned.


{\small Recorded unresolved matters: source-qualified matters; exact scope remains in the canonical record.}\par
{\footnotesize \href{../../../../continuum/business/business-service/b3cd0d8e40f44b9fa2b64367237345fa.md}{Canonical record: b3cd0d8e40f44b9fa2b64367237345fa}; full qualifications and provenance.}\par
\subsubsection*{III.4.4.8.2 Business Functions}
\phantomsection
\label{navigation:III.4.4.8.2}
\addcontentsline{toc}{subsubsection}{III.4.4.8.2 Business Functions}

\paragraph*{III.4.4.8.2.1 Integrate Mobile Work Management}
\phantomsection
\label{navigation:III.4.4.8.2.1}
\addcontentsline{toc}{paragraph}{III.4.4.8.2.1 Integrate Mobile Work Management}

\label{element:81f97ad819da48f09f26029066550e08}
\index{Integrate Mobile Work Management (III.4.4.8.2.1)}
{\small BusinessFunction · Canonical · current}\par

Exchanges externally established work/assignment, dispatch and operational progression/status information, with applicable worker/context information, between mobile work-management mechanisms and Harmonia operational coordination.



Derivation and boundary: The human adjudication of FEAT-HSO-16 establishes integration with mobile work/task management. This responsibility is absent from presence observation and visit-specific milestone ingestion. It does not determine, allocate or manage the underlying clinical work, clinical handover or EMR worklists. Work Allocation \& Dispatch retains non-clinical operational Work Order responsibility.


{\footnotesize \href{../../../../continuum/business/business-function/81f97ad819da48f09f26029066550e08.md}{Canonical record: 81f97ad819da48f09f26029066550e08}; full qualifications and provenance.}\par
\paragraph*{III.4.4.8.2.2 Manage Mobile Service Capacity}
\phantomsection
\label{navigation:III.4.4.8.2.2}
\addcontentsline{toc}{paragraph}{III.4.4.8.2.2 Manage Mobile Service Capacity}

\label{element:ad5cb5451ef645acb55ceb0c942ea013}
\index{Manage Mobile Service Capacity (III.4.4.8.2.2)}
{\small BusinessFunction · Canonical · current}\par

Understands and manages available, committed and utilised mobile-service capacity in relation to workforce availability, geography, time, workload and applicable service demand.



Responsibility boundary: Staff presence and externally established assignments/visit progression contribute to capacity; they do not alone express it. Capacity coordination does not determine clinical work, own the source workforce roster or allocate underlying clinical duties.


{\small Recorded unresolved matters: source-qualified matters; exact scope remains in the canonical record.}\par
{\footnotesize \href{../../../../continuum/business/business-function/ad5cb5451ef645acb55ceb0c942ea013.md}{Canonical record: ad5cb5451ef645acb55ceb0c942ea013}; full qualifications and provenance.}\par
\paragraph*{III.4.4.8.2.3 Receive Mobile Visit Progression}
\phantomsection
\label{navigation:III.4.4.8.2.3}
\addcontentsline{toc}{paragraph}{III.4.4.8.2.3 Receive Mobile Visit Progression}

\label{element:71700f211ef4403f98a2b9b11d8d7848}
\index{Receive Mobile Visit Progression (III.4.4.8.2.3)}
{\small BusinessFunction · Canonical · current}\par

Ingests arrival, visit-progression and safety check-in milestones from roving staff.



Derivation: FEAT-HSO-17 directly requires visit-specific milestone ingestion. It remains distinct from presence, capacity assessment and work-management integration. Received progression can inform mobile capacity without determining clinical work or clinical completion.


{\footnotesize \href{../../../../continuum/business/business-function/71700f211ef4403f98a2b9b11d8d7848.md}{Canonical record: 71700f211ef4403f98a2b9b11d8d7848}; full qualifications and provenance.}\par
\paragraph*{III.4.4.8.2.4 Track Operational Staff Presence}
\phantomsection
\label{navigation:III.4.4.8.2.4}
\addcontentsline{toc}{paragraph}{III.4.4.8.2.4 Track Operational Staff Presence}

\label{element:9ccdb235a69946e0ad5a99001a1d16e0}
\index{Track Operational Staff Presence (III.4.4.8.2.4)}
{\small BusinessFunction · Canonical · current}\par

Tracks active on-duty presence and physical zone assignment for porters, phlebotomists, and roving nurses.



Feature association not established: FEAT-HSO-16 — Mobile Worker Task Dispatch retains its Strategy identity with the integration boundary below. Presence/zone tracking remains distinct from mobile work-management integration; no replacement association to presence is assigned.


{\small Recorded unresolved matters: source-qualified matters; exact scope remains in the canonical record.}\par
{\footnotesize \href{../../../../continuum/business/business-function/9ccdb235a69946e0ad5a99001a1d16e0.md}{Canonical record: 9ccdb235a69946e0ad5a99001a1d16e0}; full qualifications and provenance.}\par
\subsection*{III.4.4.9 [CT2] On-Call Management — Capability context}
\phantomsection
\label{navigation:III.4.4.9}
\addcontentsline{toc}{subsection}{III.4.4.9 [CT2] On-Call Management — Capability context}

{\small Canonical Strategy Capability: II.1.3.6.17 \hyperref[element:b16cef8f2bfd4b92948ec75cbcfa7b9a]{[CT2] On-Call Management}; Canonical ID: \texttt{\seqsplit{b16cef8f2bfd4b92948ec75cbcfa7b9a}}. Canonical · current.}\par
Contextual reference only; no additional architectural instance or semantic relationship.\par
\subsubsection*{III.4.4.9.1 Business Services}
\phantomsection
\label{navigation:III.4.4.9.1}
\addcontentsline{toc}{subsubsection}{III.4.4.9.1 Business Services}

\paragraph*{III.4.4.9.1.1 On-Call Specialist Lookup}
\phantomsection
\label{navigation:III.4.4.9.1.1}
\addcontentsline{toc}{paragraph}{III.4.4.9.1.1 On-Call Specialist Lookup}

\label{element:e0ba307dc51e47a395c1eac9fe728a98}
\index{On-Call Specialist Lookup (III.4.4.9.1.1)}
{\small BusinessService · Canonical · current}\par

Discloses active on-call contacts for switchboard and emergency teams.


{\footnotesize \href{../../../../continuum/business/business-service/e0ba307dc51e47a395c1eac9fe728a98.md}{Canonical record: e0ba307dc51e47a395c1eac9fe728a98}; full qualifications and provenance.}\par
\subsubsection*{III.4.4.9.2 Business Functions}
\phantomsection
\label{navigation:III.4.4.9.2}
\addcontentsline{toc}{subsubsection}{III.4.4.9.2 Business Functions}

\paragraph*{III.4.4.9.2.1 Resolve Active On-Call Coverage}
\phantomsection
\label{navigation:III.4.4.9.2.1}
\addcontentsline{toc}{paragraph}{III.4.4.9.2.1 Resolve Active On-Call Coverage}

\label{element:a88d3779091140c6b90f98d550b5387f}
\index{Resolve Active On-Call Coverage (III.4.4.9.2.1)}
{\small BusinessFunction · Canonical · current}\par

Resolves primary, secondary, and tertiary on-call clinical specialists for emergency escalation.


{\footnotesize \href{../../../../continuum/business/business-function/a88d3779091140c6b90f98d550b5387f.md}{Canonical record: a88d3779091140c6b90f98d550b5387f}; full qualifications and provenance.}\par
\subsection*{III.4.4.10 [CT2] Outpatient Operations — Capability context}
\phantomsection
\label{navigation:III.4.4.10}
\addcontentsline{toc}{subsection}{III.4.4.10 [CT2] Outpatient Operations — Capability context}

{\small Canonical Strategy Capability: II.1.3.6.19 \hyperref[element:21399d52b9624ddfbeb786f1d72e40c3]{[CT2] Outpatient Operations}; Canonical ID: \texttt{\seqsplit{21399d52b9624ddfbeb786f1d72e40c3}}. Canonical · current.}\par
Contextual reference only; no additional architectural instance or semantic relationship.\par
\subsubsection*{III.4.4.10.1 Business Services}
\phantomsection
\label{navigation:III.4.4.10.1}
\addcontentsline{toc}{subsubsection}{III.4.4.10.1 Business Services}

\paragraph*{III.4.4.10.1.1 Ambulatory Session Status Query}
\phantomsection
\label{navigation:III.4.4.10.1.1}
\addcontentsline{toc}{paragraph}{III.4.4.10.1.1 Ambulatory Session Status Query}

\label{element:a790a785f2e148d7a1c99655d53db4e9}
\index{Ambulatory Session Status Query (III.4.4.10.1.1)}
{\small BusinessService · Canonical · current}\par

Exposes clinic room occupancy.



Exposure qualification: Query behaviour and Outpatient Operations ownership remain established; a meaningful outside consumer/mandate is not identified by the current description. That exposure relationship remains unestablished rather than populated with an inferred Actor, Role or Capability.



Capacity boundary: Operational coordination does not acquire authoritative booking/rostering or clinical-work management. Room-occupancy query remains narrower than the capacity responsibility; no enlarged exposure or new named Service is inferred.


{\small Recorded unresolved matters: source-qualified matters; exact scope remains in the canonical record.}\par
{\footnotesize \href{../../../../continuum/business/business-service/a790a785f2e148d7a1c99655d53db4e9.md}{Canonical record: a790a785f2e148d7a1c99655d53db4e9}; full qualifications and provenance.}\par
\subsubsection*{III.4.4.10.2 Business Functions}
\phantomsection
\label{navigation:III.4.4.10.2}
\addcontentsline{toc}{subsubsection}{III.4.4.10.2 Business Functions}

\paragraph*{III.4.4.10.2.1 Capture Post-Clinic Follow-up Requirements}
\phantomsection
\label{navigation:III.4.4.10.2.1}
\addcontentsline{toc}{paragraph}{III.4.4.10.2.1 Capture Post-Clinic Follow-up Requirements}

\label{element:9136f75069bf4770bcc8dc395204d475}
\index{Capture Post-Clinic Follow-up Requirements (III.4.4.10.2.1)}
{\small BusinessFunction · Canonical · current}\par

Captures follow-up booking and diagnostic-order requirements generated at clinic completion.



Derivation: FEAT-HSO-10 directly establishes this capture responsibility. The Clinic queue Process's consultation-concluded checkpoint describes orders placed, not Outpatient Operations' capture of follow-up requirements. Capture does not authorise clinical orders, own authoritative bookings or establish a new Service dependency.


{\footnotesize \href{../../../../continuum/business/business-function/9136f75069bf4770bcc8dc395204d475.md}{Canonical record: 9136f75069bf4770bcc8dc395204d475}; full qualifications and provenance.}\par
\paragraph*{III.4.4.10.2.2 Manage Outpatient Capacity}
\phantomsection
\label{navigation:III.4.4.10.2.2}
\addcontentsline{toc}{paragraph}{III.4.4.10.2.2 Manage Outpatient Capacity}

\label{element:45fd6190b7334f838575881725f6a66b}
\index{Manage Outpatient Capacity (III.4.4.10.2.2)}
{\small BusinessFunction · Canonical · current}\par

Understands and manages available, committed and utilised outpatient capacity through clinic/session, practitioner and room availability, appointment demand, patient arrival/waiting and operational progression. Retains specialty room allocation and clinician arrival tracking.



Exposure qualification: Query behaviour and Outpatient Operations ownership remain established; a meaningful outside consumer/mandate is not identified by the current description. That exposure relationship remains unestablished rather than populated with an inferred Actor, Role or Capability.



Capacity boundary: Operational coordination does not acquire authoritative booking/rostering or clinical-work management. Room-occupancy query remains narrower than the capacity responsibility; no enlarged exposure or new named Service is inferred.


{\footnotesize \href{../../../../continuum/business/business-function/45fd6190b7334f838575881725f6a66b.md}{Canonical record: 45fd6190b7334f838575881725f6a66b}; full qualifications and provenance.}\par
\subsection*{III.4.4.11 [CT2] Patient Transport — Capability context}
\phantomsection
\label{navigation:III.4.4.11}
\addcontentsline{toc}{subsection}{III.4.4.11 [CT2] Patient Transport — Capability context}

{\small Canonical Strategy Capability: II.1.3.6.20 \hyperref[element:6041d7c0cc5d4519aaea584bb4e6d365]{[CT2] Patient Transport}; Canonical ID: \texttt{\seqsplit{6041d7c0cc5d4519aaea584bb4e6d365}}. Canonical · current.}\par
Contextual reference only; no additional architectural instance or semantic relationship.\par
\subsubsection*{III.4.4.11.1 Business Services}
\phantomsection
\label{navigation:III.4.4.11.1}
\addcontentsline{toc}{subsubsection}{III.4.4.11.1 Business Services}

\paragraph*{III.4.4.11.1.1 Patient Transport Tracking Service}
\phantomsection
\label{navigation:III.4.4.11.1.1}
\addcontentsline{toc}{paragraph}{III.4.4.11.1.1 Patient Transport Tracking Service}

\label{element:8d71168ac7da48e990757b2a0ae3670f}
\index{Patient Transport Tracking Service (III.4.4.11.1.1)}
{\small BusinessService · Canonical · current}\par

Provides real-time transit telemetry.



Additional Feature responsibility: FEAT-HSO-24 — Transport Request Ingestion is directly supported by the request stage of this same Function/Process. It ingests internal portering and external ambulance/inter-hospital requests with clinical mobility requirements; request ingestion is distinct from dispatch, transit and arrival.


{\footnotesize \href{../../../../continuum/business/business-service/8d71168ac7da48e990757b2a0ae3670f.md}{Canonical record: 8d71168ac7da48e990757b2a0ae3670f}; full qualifications and provenance.}\par
\subsubsection*{III.4.4.11.2 Business Functions}
\phantomsection
\label{navigation:III.4.4.11.2}
\addcontentsline{toc}{subsubsection}{III.4.4.11.2 Business Functions}

\paragraph*{III.4.4.11.2.1 Manage Patient Transport Progression}
\phantomsection
\label{navigation:III.4.4.11.2.1}
\addcontentsline{toc}{paragraph}{III.4.4.11.2.1 Manage Patient Transport Progression}

\label{element:8a2136544646458fb60413afed2b66e1}
\index{Manage Patient Transport Progression (III.4.4.11.2.1)}
{\small BusinessFunction · Canonical · current}\par

Coordinates intra-facility and inter-facility patient movements: Transport Requested → Porter Dispatched → Patient Collected → In-Transit → Delivered at Destination → Handover Completed.



Additional Feature responsibility: FEAT-HSO-24 — Transport Request Ingestion is directly supported by the request stage of this same Function/Process. It ingests internal portering and external ambulance/inter-hospital requests with clinical mobility requirements; request ingestion is distinct from dispatch, transit and arrival.


{\footnotesize \href{../../../../continuum/business/business-function/8a2136544646458fb60413afed2b66e1.md}{Canonical record: 8a2136544646458fb60413afed2b66e1}; full qualifications and provenance.}\par
\subsubsection*{III.4.4.11.3 Business Processes}
\phantomsection
\label{navigation:III.4.4.11.3}
\addcontentsline{toc}{subsubsection}{III.4.4.11.3 Business Processes}

\paragraph*{III.4.4.11.3.1 Patient Transport Process}
\phantomsection
\label{navigation:III.4.4.11.3.1}
\addcontentsline{toc}{paragraph}{III.4.4.11.3.1 Patient Transport Process}

\label{element:f2e58fd18cd94b0faaa364d3890a7c9e}
\index{Patient Transport Process (III.4.4.11.3.1)}
{\small BusinessProcess · Canonical · current}\par

•
Owning Capability: Patient Transport



•
Process Purpose: Governs the physical portering and transport of patients between hospital departments, wards, diagnostic suites, and external facilities.



•
State Progression Lifecycle: Transport Requested → Porter Dispatched → Patient Collected → In-Transit → Delivered at Destination → Clinical Handover Complete



•
Key State Dispositions:



•
Transport Requested: Clinical unit books wheelchair, bed, or ambulance transport.



•
Porter Dispatched: Wardsperson assigned and en route to pickup location.



•
Patient Collected: Patient identity verified and transfer initiated.



•
In-Transit: Patient actively moving through facility corridors or ambulance transit.



•
Delivered at Destination: Patient arrives at target diagnostic suite or ward.



•
Clinical Handover Complete: Safe custody and clinical briefing transferred to receiving staff.



Patient transport progression records the physical journey from a movement request through departure, transit and destination arrival. Its checkpoints allow operational participants to follow the journey with relevant mobility context intact. Arrival closes a movement concern only under its applicable scope; it does not establish clinical handover or destination acceptance of care responsibility.


{\small Recorded unresolved matters: source-qualified matters; exact scope remains in the canonical record.}\par
{\footnotesize \href{../../../../continuum/business/business-process/f2e58fd18cd94b0faaa364d3890a7c9e.md}{Canonical record: f2e58fd18cd94b0faaa364d3890a7c9e}; full qualifications and provenance.}\par
\subsection*{III.4.4.12 [CT2] Service Capacity Management — Capability context}
\phantomsection
\label{navigation:III.4.4.12}
\addcontentsline{toc}{subsection}{III.4.4.12 [CT2] Service Capacity Management — Capability context}

{\small Canonical Strategy Capability: II.1.3.6.21 \hyperref[element:d86105bf9c86445394088f65c895d37f]{[CT2] Service Capacity Management}; Canonical ID: \texttt{\seqsplit{d86105bf9c86445394088f65c895d37f}}. Canonical · current.}\par
Contextual reference only; no additional architectural instance or semantic relationship.\par
\subsubsection*{III.4.4.12.1 Business Services}
\phantomsection
\label{navigation:III.4.4.12.1}
\addcontentsline{toc}{subsubsection}{III.4.4.12.1 Business Services}

\paragraph*{III.4.4.12.1.1 Regional Health System Capacity Query}
\phantomsection
\label{navigation:III.4.4.12.1.1}
\addcontentsline{toc}{paragraph}{III.4.4.12.1.1 Regional Health System Capacity Query}

\label{element:ecbb1944572d44dbba55aaf565559d13}
\index{Regional Health System Capacity Query (III.4.4.12.1.1)}
{\small BusinessService · Canonical · current}\par

Provides executive operational situational awareness.


{\footnotesize \href{../../../../continuum/business/business-service/ecbb1944572d44dbba55aaf565559d13.md}{Canonical record: ecbb1944572d44dbba55aaf565559d13}; full qualifications and provenance.}\par
\subsubsection*{III.4.4.12.2 Business Functions}
\phantomsection
\label{navigation:III.4.4.12.2}
\addcontentsline{toc}{subsubsection}{III.4.4.12.2 Business Functions}

\paragraph*{III.4.4.12.2.1 Aggregate Service Capacity Telemetry}
\phantomsection
\label{navigation:III.4.4.12.2.1}
\addcontentsline{toc}{paragraph}{III.4.4.12.2.1 Aggregate Service Capacity Telemetry}

\label{element:ac08378c955048ef81ed5d9ff9c5cc7e}
\index{Aggregate Service Capacity Telemetry (III.4.4.12.2.1)}
{\small BusinessFunction · Canonical · current}\par

Consumes and coordinates context-specific service-capacity information for broader situational awareness, alongside bed occupancy, ICU surge capacity and ventilator availability across regional hospitals.


{\footnotesize \href{../../../../continuum/business/business-function/ac08378c955048ef81ed5d9ff9c5cc7e.md}{Canonical record: ac08378c955048ef81ed5d9ff9c5cc7e}; full qualifications and provenance.}\par
\subsection*{III.4.4.13 [CT2] Theatre Operations — Capability context}
\phantomsection
\label{navigation:III.4.4.13}
\addcontentsline{toc}{subsection}{III.4.4.13 [CT2] Theatre Operations — Capability context}

{\small Canonical Strategy Capability: II.1.3.6.23 \hyperref[element:4bf7a8260e1444fdac44e97c3557b533]{[CT2] Theatre Operations}; Canonical ID: \texttt{\seqsplit{4bf7a8260e1444fdac44e97c3557b533}}. Canonical · current.}\par
Contextual reference only; no additional architectural instance or semantic relationship.\par
\subsubsection*{III.4.4.13.1 Business Services}
\phantomsection
\label{navigation:III.4.4.13.1}
\addcontentsline{toc}{subsubsection}{III.4.4.13.1 Business Services}

\paragraph*{III.4.4.13.1.1 Theatre Case Progression Telemetry}
\phantomsection
\label{navigation:III.4.4.13.1.1}
\addcontentsline{toc}{paragraph}{III.4.4.13.1.1 Theatre Case Progression Telemetry}

\label{element:15fca4d52810448785ca3c70b899a294}
\index{Theatre Case Progression Telemetry (III.4.4.13.1.1)}
{\small BusinessService · Canonical · current}\par

Emits real-time surgical milestone updates.



Additional Feature responsibility: FEAT-HSO-06 — Perioperative Resource Notification is realised by dispatching progression-triggered alerts to portering, sterilisation and recovery teams within this same case-progression behaviour/exposure. Milestone observation alone does not establish notification; no new Function or Service is created.


{\footnotesize \href{../../../../continuum/business/business-service/15fca4d52810448785ca3c70b899a294.md}{Canonical record: 15fca4d52810448785ca3c70b899a294}; full qualifications and provenance.}\par
\subsubsection*{III.4.4.13.2 Business Functions}
\phantomsection
\label{navigation:III.4.4.13.2}
\addcontentsline{toc}{subsubsection}{III.4.4.13.2 Business Functions}

\paragraph*{III.4.4.13.2.1 Manage Theatre Capacity}
\phantomsection
\label{navigation:III.4.4.13.2.1}
\addcontentsline{toc}{paragraph}{III.4.4.13.2.1 Manage Theatre Capacity}

\label{element:3566d549770f490dba823fe3b0421446}
\index{Manage Theatre Capacity (III.4.4.13.2.1)}
{\small BusinessFunction · Canonical · current}\par

Understands and manages available, committed and utilised theatre capacity in relation to theatre, team and equipment availability, scheduled cases and case progression.



Responsibility boundary: Case milestones and targeted resource notifications contribute; they do not by themselves establish available or committed theatre service capacity. Clinical case determination and originating team/equipment information remain separately governed.


{\small Recorded unresolved matters: source-qualified matters; exact scope remains in the canonical record.}\par
{\footnotesize \href{../../../../continuum/business/business-function/3566d549770f490dba823fe3b0421446.md}{Canonical record: 3566d549770f490dba823fe3b0421446}; full qualifications and provenance.}\par
\paragraph*{III.4.4.13.2.2 Manage Theatre Case Progression}
\phantomsection
\label{navigation:III.4.4.13.2.2}
\addcontentsline{toc}{paragraph}{III.4.4.13.2.2 Manage Theatre Case Progression}

\label{element:88d2d17cccfb4a578b835be588043bab}
\index{Manage Theatre Case Progression (III.4.4.13.2.2)}
{\small BusinessFunction · Canonical · current}\par

Coordinates surgical case milestones: Case Called → Patient in Anaesthetic Bay → Anaesthesia Commenced → Knife to Skin → Procedure Finished → In PACU/Recovery → Ward Handover Complete.



Additional Feature responsibility: FEAT-HSO-06 — Perioperative Resource Notification is realised by dispatching progression-triggered alerts to portering, sterilisation and recovery teams within this same case-progression behaviour/exposure. Milestone observation alone does not establish notification; no new Function or Service is created.


{\footnotesize \href{../../../../continuum/business/business-function/88d2d17cccfb4a578b835be588043bab.md}{Canonical record: 88d2d17cccfb4a578b835be588043bab}; full qualifications and provenance.}\par
\subsubsection*{III.4.4.13.3 Business Processes}
\phantomsection
\label{navigation:III.4.4.13.3}
\addcontentsline{toc}{subsubsection}{III.4.4.13.3 Business Processes}

\paragraph*{III.4.4.13.3.1 Theatre Case Progression Process}
\phantomsection
\label{navigation:III.4.4.13.3.1}
\addcontentsline{toc}{paragraph}{III.4.4.13.3.1 Theatre Case Progression Process}

\label{element:8ad97120b5b54a31a64a4522e6bc9fd1}
\index{Theatre Case Progression Process (III.4.4.13.3.1)}
{\small BusinessProcess · Canonical · current}\par

•
Owning Capability: Theatre Operations



•
Process Purpose: Governs the perioperative operational milestones of a surgical procedure within the operating theatre suite.



•
State Progression Lifecycle: Case Called → In Anaesthetic Bay → Anaesthesia Commenced → Patient in Theatre → Knife-to-Skin → Procedure Finished → In PACU Recovery → Ward Handover Complete



•
Key State Dispositions:



•
Case Called: Theatre team requests ward to send patient for surgery.



•
In Anaesthetic Bay: Patient arrives in theatre holding and identity/consent verified.



•
Anaesthesia Commenced: Anaesthetic induction underway.



•
Patient in Theatre: Patient transferred onto operating table.



•
Knife-to-Skin: Surgical incision initiated (operative start).



•
Procedure Finished: Wound closure completed and dressings applied.



•
In PACU Recovery: Patient monitored in post-anaesthesia care unit.



•
Ward Handover Complete: Post-operative clinical handover to inpatient ward completed.



Theatre case progression follows reported milestones of the surgical encounter so that supporting participants can coordinate around the actual situation. The state account concerns what the clinical service reports has occurred. It does not make the integration mechanism the actor administering anaesthesia, performing surgery or managing clinical handover.


{\small Recorded unresolved matters: source-qualified matters; exact scope remains in the canonical record.}\par
{\footnotesize \href{../../../../continuum/business/business-process/8ad97120b5b54a31a64a4522e6bc9fd1.md}{Canonical record: 8ad97120b5b54a31a64a4522e6bc9fd1}; full qualifications and provenance.}\par
\subsection*{III.4.4.14 [CT2] Ward Operations — Capability context}
\phantomsection
\label{navigation:III.4.4.14}
\addcontentsline{toc}{subsection}{III.4.4.14 [CT2] Ward Operations — Capability context}

{\small Canonical Strategy Capability: II.1.3.6.24 \hyperref[element:f996ecca728f4942a88b034aca818a99]{[CT2] Ward Operations}; Canonical ID: \texttt{\seqsplit{f996ecca728f4942a88b034aca818a99}}. Canonical · current.}\par
Contextual reference only; no additional architectural instance or semantic relationship.\par
\subsubsection*{III.4.4.14.1 [CT3] Multidisciplinary Inpatient Care Coordination — Capability context}
\phantomsection
\label{navigation:III.4.4.14.1}
\addcontentsline{toc}{subsubsection}{III.4.4.14.1 [CT3] Multidisciplinary Inpatient Care Coordination — Capability context}

{\small Canonical Strategy Capability: II.1.3.6.24.1 \hyperref[element:c798e38e12e5428dbb8b3cabc977d430]{[CT3] Multidisciplinary Inpatient Care Coordination}; Canonical ID: \texttt{\seqsplit{c798e38e12e5428dbb8b3cabc977d430}}. Canonical · current.}\par
Contextual reference only; no additional architectural instance or semantic relationship.\par
\paragraph*{III.4.4.14.1.1 Business Services}
\phantomsection
\label{navigation:III.4.4.14.1.1}
\addcontentsline{toc}{paragraph}{III.4.4.14.1.1 Business Services}

\subparagraph*{III.4.4.14.1.1.1 Discharge Dossier Packaging Service}
\phantomsection
\label{navigation:III.4.4.14.1.1.1}
\addcontentsline{toc}{subparagraph}{III.4.4.14.1.1.1 Discharge Dossier Packaging Service}

\label{element:02a17f4f9b0a425299141261130a1768}
\index{Discharge Dossier Packaging Service (III.4.4.14.1.1.1)}
{\small BusinessService · Canonical · current}\par

Emits canonical discharge summaries to GPs and community providers.



Feature association not established: FEAT-SD-08 — Inpatient Clinical Progression Tracking is adjudicated sufficient through general clinical-information responsibilities above. Discharge-dossier assembly/publication remains distinct; no separate clinical-tracking Function or replacement association is assigned.


{\small Recorded unresolved matters: source-qualified matters; exact scope remains in the canonical record.}\par
{\footnotesize \href{../../../../continuum/business/business-service/02a17f4f9b0a425299141261130a1768.md}{Canonical record: 02a17f4f9b0a425299141261130a1768}; full qualifications and provenance.}\par
\subparagraph*{III.4.4.14.1.1.2 Ward Round Clinical Briefing Query}
\phantomsection
\label{navigation:III.4.4.14.1.1.2}
\addcontentsline{toc}{subparagraph}{III.4.4.14.1.1.2 Ward Round Clinical Briefing Query}

\label{element:2b0f17a8a2da4c0f92a5ad4a17ee526b}
\index{Ward Round Clinical Briefing Query (III.4.4.14.1.1.2)}
{\small BusinessService · Canonical · current}\par

Discloses consolidated 24-hour clinical briefing views.


{\footnotesize \href{../../../../continuum/business/business-service/2b0f17a8a2da4c0f92a5ad4a17ee526b.md}{Canonical record: 2b0f17a8a2da4c0f92a5ad4a17ee526b}; full qualifications and provenance.}\par
\paragraph*{III.4.4.14.1.2 Business Functions}
\phantomsection
\label{navigation:III.4.4.14.1.2}
\addcontentsline{toc}{paragraph}{III.4.4.14.1.2 Business Functions}

\subparagraph*{III.4.4.14.1.2.1 Assemble Inpatient Discharge Dossier}
\phantomsection
\label{navigation:III.4.4.14.1.2.1}
\addcontentsline{toc}{subparagraph}{III.4.4.14.1.2.1 Assemble Inpatient Discharge Dossier}

\label{element:eeefe77a935a46aface61aca4e2bcfb6}
\index{Assemble Inpatient Discharge Dossier (III.4.4.14.1.2.1)}
{\small BusinessFunction · Canonical · current}\par

Synthesises inpatient diagnoses, procedural summaries, reconciled discharge medications, and follow-up instructions into a canonical discharge summary package.



Feature association not established: FEAT-SD-08 — Inpatient Clinical Progression Tracking is adjudicated sufficient through general clinical-information responsibilities above. Discharge-dossier assembly/publication remains distinct; no separate clinical-tracking Function or replacement association is assigned.


{\small Recorded unresolved matters: source-qualified matters; exact scope remains in the canonical record.}\par
{\footnotesize \href{../../../../continuum/business/business-function/eeefe77a935a46aface61aca4e2bcfb6.md}{Canonical record: eeefe77a935a46aface61aca4e2bcfb6}; full qualifications and provenance.}\par
\subparagraph*{III.4.4.14.1.2.2 Consolidate Inpatient Clinical Round Feed}
\phantomsection
\label{navigation:III.4.4.14.1.2.2}
\addcontentsline{toc}{subparagraph}{III.4.4.14.1.2.2 Consolidate Inpatient Clinical Round Feed}

\label{element:76eb51b9cc6a4417bbc8491830e43ad1}
\index{Consolidate Inpatient Clinical Round Feed (III.4.4.14.1.2.2)}
{\small BusinessFunction · Canonical · current}\par

Aggregates daily vital signs, overnight pathology, allied health notes, and medication administration records.


{\footnotesize \href{../../../../continuum/business/business-function/76eb51b9cc6a4417bbc8491830e43ad1.md}{Canonical record: 76eb51b9cc6a4417bbc8491830e43ad1}; full qualifications and provenance.}\par
\subsubsection*{III.4.4.14.2 Business Services}
\phantomsection
\label{navigation:III.4.4.14.2}
\addcontentsline{toc}{subsubsection}{III.4.4.14.2 Business Services}

\paragraph*{III.4.4.14.2.1 Ward Operational Summary Query}
\phantomsection
\label{navigation:III.4.4.14.2.1}
\addcontentsline{toc}{paragraph}{III.4.4.14.2.1 Ward Operational Summary Query}

\label{element:55a7f84dfb4142e2aa96c5253b4b6394}
\index{Ward Operational Summary Query (III.4.4.14.2.1)}
{\small BusinessService · Canonical · current}\par

Discloses ward-level operational telemetry.



Capacity boundary: Bed \& Care-Place Management contributes resource state. Ward capacity is assessed in its service context; an available bed alone does not establish capacity to accept a patient. Source staffing and clinical assertions retain their authority.


{\small Recorded unresolved matters: source-qualified matters; exact scope remains in the canonical record.}\par
{\footnotesize \href{../../../../continuum/business/business-service/55a7f84dfb4142e2aa96c5253b4b6394.md}{Canonical record: 55a7f84dfb4142e2aa96c5253b4b6394}; full qualifications and provenance.}\par
\subsubsection*{III.4.4.14.3 Business Functions}
\phantomsection
\label{navigation:III.4.4.14.3}
\addcontentsline{toc}{subsubsection}{III.4.4.14.3 Business Functions}

\paragraph*{III.4.4.14.3.1 Assemble Ward Handover Context}
\phantomsection
\label{navigation:III.4.4.14.3.1}
\addcontentsline{toc}{paragraph}{III.4.4.14.3.1 Assemble Ward Handover Context}

\label{element:6006d70ed3794c57ab2bc57119292e4b}
\index{Assemble Ward Handover Context (III.4.4.14.3.1)}
{\small BusinessFunction · Canonical · current}\par

Collates current ward patient lists, active medical concerns and pending tasks for nursing shift handover.



Derivation: FEAT-HSO-04 directly requires this assembly. Ward acuity/cohorting telemetry alone does not provide the stated handover content. Source clinical information retains its originating authority; assembly does not perform clinical review, manage clinical handover, determine pending clinical work or acquire clinical judgement.


{\footnotesize \href{../../../../continuum/business/business-function/6006d70ed3794c57ab2bc57119292e4b.md}{Canonical record: 6006d70ed3794c57ab2bc57119292e4b}; full qualifications and provenance.}\par
\paragraph*{III.4.4.14.3.2 Manage Ward Capacity}
\phantomsection
\label{navigation:III.4.4.14.3.2}
\addcontentsline{toc}{paragraph}{III.4.4.14.3.2 Manage Ward Capacity}

\label{element:eca5f4fbf92743f987a01930dc4ee763}
\index{Manage Ward Capacity (III.4.4.14.3.2)}
{\small BusinessFunction · Canonical · current}\par

Coordinates available, committed and utilised ward capacity and its patient-flow implications, retaining acuity distribution, nurse-to-patient ratios and isolation cohorting alongside bed/care-place state, patient demand and operational constraints.



Capacity boundary: Bed \& Care-Place Management contributes resource state. Ward capacity is assessed in its service context; an available bed alone does not establish capacity to accept a patient. Source staffing and clinical assertions retain their authority.


{\footnotesize \href{../../../../continuum/business/business-function/eca5f4fbf92743f987a01930dc4ee763.md}{Canonical record: eca5f4fbf92743f987a01930dc4ee763}; full qualifications and provenance.}\par
\subsection*{III.4.4.15 [CT2] Work Allocation \& Dispatch — Capability context}
\phantomsection
\label{navigation:III.4.4.15}
\addcontentsline{toc}{subsection}{III.4.4.15 [CT2] Work Allocation \& Dispatch — Capability context}

{\small Canonical Strategy Capability: II.1.3.6.25 \hyperref[element:152a5252952c405994e9b171234efcb0]{[CT2] Work Allocation \& Dispatch}; Canonical ID: \texttt{\seqsplit{152a5252952c405994e9b171234efcb0}}. Canonical · current.}\par
Contextual reference only; no additional architectural instance or semantic relationship.\par
\subsubsection*{III.4.4.15.1 Business Services}
\phantomsection
\label{navigation:III.4.4.15.1}
\addcontentsline{toc}{subsubsection}{III.4.4.15.1 Business Services}

\paragraph*{III.4.4.15.1.1 Operational Work Status Query}
\phantomsection
\label{navigation:III.4.4.15.1.1}
\addcontentsline{toc}{paragraph}{III.4.4.15.1.1 Operational Work Status Query}

\label{element:8659408923e641b99727ef218e2b8dfc}
\index{Operational Work Status Query (III.4.4.15.1.1)}
{\small BusinessService · Canonical · current}\par

Discloses task progression to requesting wards.


{\footnotesize \href{../../../../continuum/business/business-service/8659408923e641b99727ef218e2b8dfc.md}{Canonical record: 8659408923e641b99727ef218e2b8dfc}; full qualifications and provenance.}\par
\paragraph*{III.4.4.15.1.2 Work Dispatch Service}
\phantomsection
\label{navigation:III.4.4.15.1.2}
\addcontentsline{toc}{paragraph}{III.4.4.15.1.2 Work Dispatch Service}

\label{element:7410926f15d3480eaa856d998d67d83f}
\index{Work Dispatch Service (III.4.4.15.1.2)}
{\small BusinessService · Canonical · current}\par

Delivers allocated work and captures distinct technical and operational acknowledgements; acknowledgement alone does not establish completion.


{\footnotesize \href{../../../../continuum/business/business-service/7410926f15d3480eaa856d998d67d83f.md}{Canonical record: 7410926f15d3480eaa856d998d67d83f}; full qualifications and provenance.}\par
\subsubsection*{III.4.4.15.2 Business Functions}
\phantomsection
\label{navigation:III.4.4.15.2}
\addcontentsline{toc}{subsubsection}{III.4.4.15.2 Business Functions}

\paragraph*{III.4.4.15.2.1 Dispatch Operational Work Order}
\phantomsection
\label{navigation:III.4.4.15.2.1}
\addcontentsline{toc}{paragraph}{III.4.4.15.2.1 Dispatch Operational Work Order}

\label{element:05a14489d3e943499581620356130007}
\index{Dispatch Operational Work Order (III.4.4.15.2.1)}
{\small BusinessFunction · Canonical · current}\par

Allocates work requests to available wardspersons or support teams using applicable allocation policies, staff roles, credentials, location and skill, and delivers allocated work to the target worker, device or team queue.


{\footnotesize \href{../../../../continuum/business/business-function/05a14489d3e943499581620356130007.md}{Canonical record: 05a14489d3e943499581620356130007}; full qualifications and provenance.}\par
\paragraph*{III.4.4.15.2.2 Manage Operational Work Progression}
\phantomsection
\label{navigation:III.4.4.15.2.2}
\addcontentsline{toc}{paragraph}{III.4.4.15.2.2 Manage Operational Work Progression}

\label{element:ebc35951617a4f92863922b74dd1c1f4}
\index{Manage Operational Work Progression (III.4.4.15.2.2)}
{\small BusinessFunction · Canonical · current}\par

Tracks execution states: Queued → Assigned → Accepted → In-Progress → Completed / Aborted.


{\footnotesize \href{../../../../continuum/business/business-function/ebc35951617a4f92863922b74dd1c1f4.md}{Canonical record: ebc35951617a4f92863922b74dd1c1f4}; full qualifications and provenance.}\par
\subsubsection*{III.4.4.15.3 Business Processes}
\phantomsection
\label{navigation:III.4.4.15.3}
\addcontentsline{toc}{subsubsection}{III.4.4.15.3 Business Processes}

\paragraph*{III.4.4.15.3.1 Operational Work Progression Process}
\phantomsection
\label{navigation:III.4.4.15.3.1}
\addcontentsline{toc}{paragraph}{III.4.4.15.3.1 Operational Work Progression Process}

\label{element:5ebba522194b4de5a5b3963e855d712d}
\index{Operational Work Progression Process (III.4.4.15.3.1)}
{\small BusinessProcess · Canonical · current}\par

•
Owning Capability: Work Allocation \& Dispatch



•
Process Purpose: Governs the dispatch, assignment, execution, and completion of facility operational jobs (e.g., equipment moves, waste removal, linen supply).



•
State Progression Lifecycle: Work Requested → Queued in Dispatch → Assigned to Worker → Job Accepted → Work In-Progress → Work Completed / Aborted



•
Key State Dispositions:



•
Work Requested: Facility support task submitted by clinical unit.



•
Queued in Dispatch: Job prioritized in operational queue.



•
Assigned to Worker: Job allocated to specific staff member or team.



•
Job Accepted: Worker acknowledges dispatch on mobile device.



•
Work In-Progress: Task actively being executed.



•
Work Completed / Aborted: Job marked finished or formally cancelled with reason.



Operational work progression makes the path from a recognised support requirement to assigned and reported work explicit. Acknowledgement and activity progress help distinguish a dispatch from actual performance and expose impediments requiring operational attention. This responsibility remains bounded to authorised support work rather than ownership of the clinician's worklist.


{\small Recorded unresolved matters: source-qualified matters; exact scope remains in the canonical record.}\par
{\footnotesize \href{../../../../continuum/business/business-process/5ebba522194b4de5a5b3963e855d712d.md}{Canonical record: 5ebba522194b4de5a5b3963e855d712d}; full qualifications and provenance.}\par
\subsection*{III.4.4.16 [CT2] Workforce Qualification \& Development Enablement — Capability context [superseded]}
\phantomsection
\label{navigation:III.4.4.16}
\addcontentsline{toc}{subsection}{III.4.4.16 [CT2] Workforce Qualification \& Development Enablement — Capability context [superseded]}

{\small Canonical Strategy Capability: II.1.3.6.26 \hyperref[element:d24a1f4c4f674c45a95c317998ff4186]{[CT2] Workforce Qualification \& Development Enablement}; Canonical ID: \texttt{\seqsplit{d24a1f4c4f674c45a95c317998ff4186}}. Proposed · superseded.}\par
Contextual reference only; no additional architectural instance or semantic relationship.\par
\subsubsection*{III.4.4.16.1 [CT3] Credential Management — Capability context}
\phantomsection
\label{navigation:III.4.4.16.1}
\addcontentsline{toc}{subsubsection}{III.4.4.16.1 [CT3] Credential Management — Capability context}

{\small Canonical Strategy Capability: II.1.3.6.26.1 \hyperref[element:3a7a4d2d0ce8446b8a6293284fb16b83]{[CT3] Credential Management}; Canonical ID: \texttt{\seqsplit{3a7a4d2d0ce8446b8a6293284fb16b83}}. Canonical · pending\_adjudication.}\par
Contextual reference only; no additional architectural instance or semantic relationship.\par
\paragraph*{III.4.4.16.1.1 Business Services}
\phantomsection
\label{navigation:III.4.4.16.1.1}
\addcontentsline{toc}{paragraph}{III.4.4.16.1.1 Business Services}

\subparagraph*{III.4.4.16.1.1.1 Operational Credential Checkpoint Service}
\phantomsection
\label{navigation:III.4.4.16.1.1.1}
\addcontentsline{toc}{subparagraph}{III.4.4.16.1.1.1 Operational Credential Checkpoint Service}

\label{element:7ecc60e34e07491588b106b9075ab09a}
\index{Operational Credential Checkpoint Service (III.4.4.16.1.1.1)}
{\small BusinessService · Canonical · pending\_adjudication}\par

Evaluates physical/operational access rights.


{\small Recorded unresolved matters: source-qualified matters; exact scope remains in the canonical record.}\par
{\footnotesize \href{../../../../continuum/business/business-service/7ecc60e34e07491588b106b9075ab09a.md}{Canonical record: 7ecc60e34e07491588b106b9075ab09a}; full qualifications and provenance.}\par
\paragraph*{III.4.4.16.1.2 Business Functions}
\phantomsection
\label{navigation:III.4.4.16.1.2}
\addcontentsline{toc}{paragraph}{III.4.4.16.1.2 Business Functions}

\subparagraph*{III.4.4.16.1.2.1 Verify Operational Access Credential}
\phantomsection
\label{navigation:III.4.4.16.1.2.1}
\addcontentsline{toc}{subparagraph}{III.4.4.16.1.2.1 Verify Operational Access Credential}

\label{element:a0eb34a7cb7046fd880248e799632810}
\index{Verify Operational Access Credential (III.4.4.16.1.2.1)}
{\small BusinessFunction · Canonical · pending\_adjudication}\par

Verifies physical access badges, specialty ward entry authorizations, and restricted area permissions.


{\small Recorded unresolved matters: source-qualified matters; exact scope remains in the canonical record.}\par
{\footnotesize \href{../../../../continuum/business/business-function/a0eb34a7cb7046fd880248e799632810.md}{Canonical record: a0eb34a7cb7046fd880248e799632810}; full qualifications and provenance.}\par
\section*{III.4.5 [CT1] Human Collaboration \& Work Enablement — Capability context [superseded]}
\phantomsection
\label{navigation:III.4.5}
\addcontentsline{toc}{section}{III.4.5 [CT1] Human Collaboration \& Work Enablement — Capability context [superseded]}

{\small Canonical Strategy Capability: II.1.3.7 \hyperref[element:c9347f73b21242c5b8cce41252b74da0]{[CT1] Human Collaboration \& Work Enablement}; Canonical ID: \texttt{\seqsplit{c9347f73b21242c5b8cce41252b74da0}}. Proposed · superseded.}\par
Contextual reference only; no additional architectural instance or semantic relationship.\par
\subsection*{III.4.5.1 [CT2] Operational Activity \& Temporal Enablement — Capability context [superseded]}
\phantomsection
\label{navigation:III.4.5.1}
\addcontentsline{toc}{subsection}{III.4.5.1 [CT2] Operational Activity \& Temporal Enablement — Capability context [superseded]}

{\small Canonical Strategy Capability: II.1.3.7.2 \hyperref[element:ca96b2b2ff434854995711e403e95f6f]{[CT2] Operational Activity \& Temporal Enablement}; Canonical ID: \texttt{\seqsplit{ca96b2b2ff434854995711e403e95f6f}}. Proposed · superseded.}\par
Contextual reference only; no additional architectural instance or semantic relationship.\par
\subsubsection*{III.4.5.1.1 [CT3] Entity-Specific Operational Coordination — Capability context}
\phantomsection
\label{navigation:III.4.5.1.1}
\addcontentsline{toc}{subsubsection}{III.4.5.1.1 [CT3] Entity-Specific Operational Coordination — Capability context}

{\small Canonical Strategy Capability: II.1.3.7.2.1 \hyperref[element:9f0c1429b64f4f88982fbbda6eaec233]{[CT3] Entity-Specific Operational Coordination}; Canonical ID: \texttt{\seqsplit{9f0c1429b64f4f88982fbbda6eaec233}}. Proposed · pending\_adjudication.}\par
Contextual reference only; no additional architectural instance or semantic relationship.\par
\paragraph*{III.4.5.1.1.1 Business Services}
\phantomsection
\label{navigation:III.4.5.1.1.1}
\addcontentsline{toc}{paragraph}{III.4.5.1.1.1 Business Services}

\subparagraph*{III.4.5.1.1.1.1 Entity Operational Input Submission}
\phantomsection
\label{navigation:III.4.5.1.1.1.1}
\addcontentsline{toc}{subparagraph}{III.4.5.1.1.1.1 Entity Operational Input Submission}

\label{element:6a27b36b624444d4b09e68f286045d46}
\index{Entity Operational Input Submission (III.4.5.1.1.1.1)}
{\small BusinessService · Proposed · pending\_adjudication}\par

Allow authorised domain coordinators to submit an entity-relative stimulus or requested operational effect.



\textbf{Overview}



Acceptance concerns the stated handoff contract; the caller observes completion/effect through separately qualified progression information.



\textbf{Boundary}



Input acceptance does not establish successful completion or authority over source clinical facts.


{\small Recorded unresolved matters: source-qualified matters; exact scope remains in the canonical record.}\par
{\footnotesize \href{../../../../continuum/business/business-service/6a27b36b624444d4b09e68f286045d46.md}{Canonical record: 6a27b36b624444d4b09e68f286045d46}; full qualifications and provenance.}\par
\paragraph*{III.4.5.1.1.2 Business Functions}
\phantomsection
\label{navigation:III.4.5.1.1.2}
\addcontentsline{toc}{paragraph}{III.4.5.1.1.2 Business Functions}

\subparagraph*{III.4.5.1.1.2.1 Coordinate Entity-Relative Input}
\phantomsection
\label{navigation:III.4.5.1.1.2.1}
\addcontentsline{toc}{subparagraph}{III.4.5.1.1.2.1 Coordinate Entity-Relative Input}

\label{element:78417d5c46cc406f95e6a41d22e3f143}
\index{Coordinate Entity-Relative Input (III.4.5.1.1.2.1)}
{\small BusinessFunction · Proposed · pending\_adjudication}\par

Select interested behaviours and progress one entity-relative input with explicit participation and effect standing.



\textbf{Overview}



Moves an accepted input through the applicable stable behaviour sequence using fresh managed contexts.



\textbf{Boundary}



No concurrent intra-entity progression or implicit clinical authority is introduced.


{\small Recorded unresolved matters: source-qualified matters; exact scope remains in the canonical record.}\par
{\footnotesize \href{../../../../continuum/business/business-function/78417d5c46cc406f95e6a41d22e3f143.md}{Canonical record: 78417d5c46cc406f95e6a41d22e3f143}; full qualifications and provenance.}\par
\subsubsection*{III.4.5.1.2 [CT3] Governed Operational Activity Design — Capability context}
\phantomsection
\label{navigation:III.4.5.1.2}
\addcontentsline{toc}{subsubsection}{III.4.5.1.2 [CT3] Governed Operational Activity Design — Capability context}

{\small Canonical Strategy Capability: II.1.3.7.2.2 \hyperref[element:3a96040821c14fe38830411dab4f8453]{[CT3] Governed Operational Activity Design}; Canonical ID: \texttt{\seqsplit{3a96040821c14fe38830411dab4f8453}}. Proposed · pending\_adjudication.}\par
Contextual reference only; no additional architectural instance or semantic relationship.\par
\paragraph*{III.4.5.1.2.1 Business Functions}
\phantomsection
\label{navigation:III.4.5.1.2.1}
\addcontentsline{toc}{paragraph}{III.4.5.1.2.1 Business Functions}

\subparagraph*{III.4.5.1.2.1.1 Maintain Operational Fulfilment Definition}
\phantomsection
\label{navigation:III.4.5.1.2.1.1}
\addcontentsline{toc}{subparagraph}{III.4.5.1.2.1.1 Maintain Operational Fulfilment Definition}

\label{element:9f7085926221436bb672e52e59bca512}
\index{Maintain Operational Fulfilment Definition (III.4.5.1.2.1.1)}
{\small BusinessFunction · Proposed · pending\_adjudication}\par

Establish and revise the declared context needs, expected outcome basis and housekeeping of an operational undertaking.



\textbf{Overview}



Changes governed activity design while preserving the applicability of the definition used by an existing undertaking.



\textbf{Boundary}



A definition revision is not a command to run work or a dynamic clinical workflow generator.


{\small Recorded unresolved matters: source-qualified matters; exact scope remains in the canonical record.}\par
{\footnotesize \href{../../../../continuum/business/business-function/9f7085926221436bb672e52e59bca512.md}{Canonical record: 9f7085926221436bb672e52e59bca512}; full qualifications and provenance.}\par
\subsubsection*{III.4.5.1.3 [CT3] Workflow \& Activity Coordination — Capability context}
\phantomsection
\label{navigation:III.4.5.1.3}
\addcontentsline{toc}{subsubsection}{III.4.5.1.3 [CT3] Workflow \& Activity Coordination — Capability context}

{\small Canonical Strategy Capability: II.1.3.7.2.3 \hyperref[element:94c683c7588f475e973b58cb226345db]{[CT3] Workflow \& Activity Coordination}; Canonical ID: \texttt{\seqsplit{94c683c7588f475e973b58cb226345db}}. Canonical · pending\_adjudication.}\par
Contextual reference only; no additional architectural instance or semantic relationship.\par
\paragraph*{III.4.5.1.3.1 Business Services}
\phantomsection
\label{navigation:III.4.5.1.3.1}
\addcontentsline{toc}{paragraph}{III.4.5.1.3.1 Business Services}

\subparagraph*{III.4.5.1.3.1.1 Activity Escalation Telemetry}
\phantomsection
\label{navigation:III.4.5.1.3.1.1}
\addcontentsline{toc}{subparagraph}{III.4.5.1.3.1.1 Activity Escalation Telemetry}

\label{element:f6c353c9349d47daadd59e57821ad03b}
\index{Activity Escalation Telemetry (III.4.5.1.3.1.1)}
{\small BusinessService · Canonical · pending\_adjudication}\par

Discloses SLA breaches.


{\small Recorded unresolved matters: source-qualified matters; exact scope remains in the canonical record.}\par
{\footnotesize \href{../../../../continuum/business/business-service/f6c353c9349d47daadd59e57821ad03b.md}{Canonical record: f6c353c9349d47daadd59e57821ad03b}; full qualifications and provenance.}\par
\subparagraph*{III.4.5.1.3.1.2 Synthetic Task Execution Service}
\phantomsection
\label{navigation:III.4.5.1.3.1.2}
\addcontentsline{toc}{subparagraph}{III.4.5.1.3.1.2 Synthetic Task Execution Service}

\label{element:57317da133ba4e29bde3997921f92e2e}
\index{Synthetic Task Execution Service (III.4.5.1.3.1.2)}
{\small BusinessService · Canonical · pending\_adjudication}\par

Coordinates automated tasks.


{\small Recorded unresolved matters: source-qualified matters; exact scope remains in the canonical record.}\par
{\footnotesize \href{../../../../continuum/business/business-service/57317da133ba4e29bde3997921f92e2e.md}{Canonical record: 57317da133ba4e29bde3997921f92e2e}; full qualifications and provenance.}\par
\subparagraph*{III.4.5.1.3.1.3 To Do Management Service}
\phantomsection
\label{navigation:III.4.5.1.3.1.3}
\addcontentsline{toc}{subparagraph}{III.4.5.1.3.1.3 To Do Management Service}

\label{element:6bc69c5517024c43b1ab701aac5e39b7}
\index{To Do Management Service (III.4.5.1.3.1.3)}
{\small BusinessService · Canonical · pending\_adjudication}\par

Provides To Do inbox and sign-off coordination without acquiring underlying clinical-work management.


{\small Recorded unresolved matters: source-qualified matters; exact scope remains in the canonical record.}\par
{\footnotesize \href{../../../../continuum/business/business-service/6bc69c5517024c43b1ab701aac5e39b7.md}{Canonical record: 6bc69c5517024c43b1ab701aac5e39b7}; full qualifications and provenance.}\par
\subparagraph*{III.4.5.1.3.1.4 Work Order Coordination Service}
\phantomsection
\label{navigation:III.4.5.1.3.1.4}
\addcontentsline{toc}{subparagraph}{III.4.5.1.3.1.4 Work Order Coordination Service}

\label{element:8a1987ed11d44598b325359287e503f0}
\index{Work Order Coordination Service (III.4.5.1.3.1.4)}
{\small BusinessService · Canonical · pending\_adjudication}\par

Manages operational work progression.


{\small Recorded unresolved matters: source-qualified matters; exact scope remains in the canonical record.}\par
{\footnotesize \href{../../../../continuum/business/business-service/8a1987ed11d44598b325359287e503f0.md}{Canonical record: 8a1987ed11d44598b325359287e503f0}; full qualifications and provenance.}\par
\paragraph*{III.4.5.1.3.2 Business Functions}
\phantomsection
\label{navigation:III.4.5.1.3.2}
\addcontentsline{toc}{paragraph}{III.4.5.1.3.2 Business Functions}

\subparagraph*{III.4.5.1.3.2.1 Coordinate Synthetic Task}
\phantomsection
\label{navigation:III.4.5.1.3.2.1}
\addcontentsline{toc}{subparagraph}{III.4.5.1.3.2.1 Coordinate Synthetic Task}

\label{element:6cb88c7e689648b1bd0a4d1e26b12f9b}
\index{Coordinate Synthetic Task (III.4.5.1.3.2.1)}
{\small BusinessFunction · Canonical · pending\_adjudication}\par

Orchestrates automated system workflows and operational progression.


{\small Recorded unresolved matters: source-qualified matters; exact scope remains in the canonical record.}\par
{\footnotesize \href{../../../../continuum/business/business-function/6cb88c7e689648b1bd0a4d1e26b12f9b.md}{Canonical record: 6cb88c7e689648b1bd0a4d1e26b12f9b}; full qualifications and provenance.}\par
\subparagraph*{III.4.5.1.3.2.2 Coordinate To Do}
\phantomsection
\label{navigation:III.4.5.1.3.2.2}
\addcontentsline{toc}{subparagraph}{III.4.5.1.3.2.2 Coordinate To Do}

\label{element:a4c9a281578d493eb876eda9fb73cac1}
\index{Coordinate To Do (III.4.5.1.3.2.2)}
{\small BusinessFunction · Canonical · pending\_adjudication}\par

Coordinates human review/update/approval items, reminders and sign-off queues within Harmonia's operational responsibility; clinical requirements and decisions remain with accountable clinical authority.


{\small Recorded unresolved matters: source-qualified matters; exact scope remains in the canonical record.}\par
{\footnotesize \href{../../../../continuum/business/business-function/a4c9a281578d493eb876eda9fb73cac1.md}{Canonical record: a4c9a281578d493eb876eda9fb73cac1}; full qualifications and provenance.}\par
\subparagraph*{III.4.5.1.3.2.3 Coordinate Work Order}
\phantomsection
\label{navigation:III.4.5.1.3.2.3}
\addcontentsline{toc}{subparagraph}{III.4.5.1.3.2.3 Coordinate Work Order}

\label{element:0161308bad6945b5bcfdeea5519341cd}
\index{Coordinate Work Order (III.4.5.1.3.2.3)}
{\small BusinessFunction · Canonical · pending\_adjudication}\par

Coordinates operational task lifecycles and dispatches.


{\small Recorded unresolved matters: source-qualified matters; exact scope remains in the canonical record.}\par
{\footnotesize \href{../../../../continuum/business/business-function/0161308bad6945b5bcfdeea5519341cd.md}{Canonical record: 0161308bad6945b5bcfdeea5519341cd}; full qualifications and provenance.}\par
\subparagraph*{III.4.5.1.3.2.4 Determine Reconciled Activity Continuation}
\phantomsection
\label{navigation:III.4.5.1.3.2.4}
\addcontentsline{toc}{subparagraph}{III.4.5.1.3.2.4 Determine Reconciled Activity Continuation}

\label{element:2e326eb5a06641048ad8a919723c4c51}
\index{Determine Reconciled Activity Continuation (III.4.5.1.3.2.4)}
{\small BusinessFunction · Proposed · pending\_adjudication}\par

Establish the permitted operational continuation from reconciled effect standing and applicable authorisation.



\textbf{Overview}



Moves a reconciled undertaking to justified continuation, authorised retry or explicitly retained uncertainty.



\textbf{Boundary}



Does not confer clinical decision authority or infer absence of effects from silence; policy evaluation remains separately owned.


{\small Recorded unresolved matters: source-qualified matters; exact scope remains in the canonical record.}\par
{\footnotesize \href{../../../../continuum/business/business-function/2e326eb5a06641048ad8a919723c4c51.md}{Canonical record: 2e326eb5a06641048ad8a919723c4c51}; full qualifications and provenance.}\par
\subparagraph*{III.4.5.1.3.2.5 Establish Activity Effect Observation Basis}
\phantomsection
\label{navigation:III.4.5.1.3.2.5}
\addcontentsline{toc}{subparagraph}{III.4.5.1.3.2.5 Establish Activity Effect Observation Basis}

\label{element:4038dd6a1baa4856afc76d7adfa6a457}
\index{Establish Activity Effect Observation Basis (III.4.5.1.3.2.5)}
{\small BusinessFunction · Proposed · pending\_adjudication}\par

Associate the applicable handoff contract and authoritative observations with an indeterminate operational undertaking.



\textbf{Overview}



Changes an unqualified uncertainty concern into a time- and attempt-specific observation basis for reconciliation.



\textbf{Boundary}



Source observations retain their authority; this is an operational observation basis rather than an independent assurance assessment.


{\small Recorded unresolved matters: source-qualified matters; exact scope remains in the canonical record.}\par
{\footnotesize \href{../../../../continuum/business/business-function/4038dd6a1baa4856afc76d7adfa6a457.md}{Canonical record: 4038dd6a1baa4856afc76d7adfa6a457}; full qualifications and provenance.}\par
\subparagraph*{III.4.5.1.3.2.6 Reconcile Indeterminate Activity Effects}
\phantomsection
\label{navigation:III.4.5.1.3.2.6}
\addcontentsline{toc}{subparagraph}{III.4.5.1.3.2.6 Reconcile Indeterminate Activity Effects}

\label{element:eace49fe7b304ad786d460ccef622378}
\index{Reconcile Indeterminate Activity Effects (III.4.5.1.3.2.6)}
{\small BusinessFunction · Proposed · pending\_adjudication}\par

Inspect authoritative observations to establish, qualify or retain uncertainty about a previously indeterminate activity effect.



\textbf{Overview}



Moves an uncertainty concern to a justified disposition while preserving the possibility that the original attempt produced an effect.



\textbf{Boundary}



Missing evidence does not establish absence of effect and cannot automatically justify a repeat mutation.


{\small Recorded unresolved matters: source-qualified matters; exact scope remains in the canonical record.}\par
{\footnotesize \href{../../../../continuum/business/business-function/eace49fe7b304ad786d460ccef622378.md}{Canonical record: eace49fe7b304ad786d460ccef622378}; full qualifications and provenance.}\par
\subparagraph*{III.4.5.1.3.2.7 Supervise Activity Timeout / Escalation}
\phantomsection
\label{navigation:III.4.5.1.3.2.7}
\addcontentsline{toc}{subparagraph}{III.4.5.1.3.2.7 Supervise Activity Timeout / Escalation}

\label{element:cbf6bfd71d5d4ea9923f10f288ff24dd}
\index{Supervise Activity Timeout / Escalation (III.4.5.1.3.2.7)}
{\small BusinessFunction · Canonical · pending\_adjudication}\par

Monitors SLAs and triggers escalations when clinical or operational deadlines elapse.


{\small Recorded unresolved matters: source-qualified matters; exact scope remains in the canonical record.}\par
{\footnotesize \href{../../../../continuum/business/business-function/cbf6bfd71d5d4ea9923f10f288ff24dd.md}{Canonical record: cbf6bfd71d5d4ea9923f10f288ff24dd}; full qualifications and provenance.}\par
\paragraph*{III.4.5.1.3.3 Business Processes}
\phantomsection
\label{navigation:III.4.5.1.3.3}
\addcontentsline{toc}{paragraph}{III.4.5.1.3.3 Business Processes}

\subparagraph*{III.4.5.1.3.3.1 Indeterminate Activity Outcome Reconciliation Process}
\phantomsection
\label{navigation:III.4.5.1.3.3.1}
\addcontentsline{toc}{subparagraph}{III.4.5.1.3.3.1 Indeterminate Activity Outcome Reconciliation Process}

\label{element:96baf3721ce94397a11486184d3df7cf}
\index{Indeterminate Activity Outcome Reconciliation Process (III.4.5.1.3.3.1)}
{\small BusinessProcess · Proposed · pending\_adjudication}\par

Order the observation, effect qualification and authorised disposition of a materially indeterminate operational outcome.



\textbf{Overview}



First inspect the governing handoff and available authoritative state; then reconcile observed effects; finally authorise a safe continuation, retry or retained uncertainty. Each participating Function keeps its responsibility owner.



\textbf{Boundary}



Not a universal lifecycle or automatic retry loop. A safe disposition may remain unresolved.


{\small Recorded unresolved matters: source-qualified matters; exact scope remains in the canonical record.}\par
{\footnotesize \href{../../../../continuum/business/business-process/96baf3721ce94397a11486184d3df7cf.md}{Canonical record: 96baf3721ce94397a11486184d3df7cf}; full qualifications and provenance.}\par
\subparagraph*{III.4.5.1.3.3.2 Synthetic Task Progression Process}
\phantomsection
\label{navigation:III.4.5.1.3.3.2}
\addcontentsline{toc}{subparagraph}{III.4.5.1.3.3.2 Synthetic Task Progression Process}

\label{element:618acdb38fc74952ab360fdde4f7d865}
\index{Synthetic Task Progression Process (III.4.5.1.3.3.2)}
{\small BusinessProcess · Canonical · pending\_adjudication}\par

•
Owning Capability: Workflow \& Activity Coordination



•
Activity Archetype: Non-Human Executable Work — automated system workflows, batch syndications, or policy evaluation tasks.



•
State Progression Lifecycle:



RAD diagram: textual source notation.
\textbackslash{}begin\{Verbatim\}[breaklines,fontsize=
]
graph LR
    D["Draft / Scheduled"] --> E["Executing"]
    E --> W["Awaiting-Dependency"]
    W --> C["Completed"]
    F["Progression Stalled / Failed\newline{}Outcome/position unresolved"]
    NOTE["Illustrative progression, not a normative transition model.\newline{}No Completed-to-Stalled/Failed transition is established."]
    classDef note fill:\#fff8dc,stroke:\#8a6d3b,stroke-dasharray:5 5;
    class NOTE,F note;
\textbackslash{}end\{Verbatim\}



•
Key State Dispositions:



•
Draft / Scheduled: Automated task instantiated and scheduled for processing.



•
Executing: System actively processing task.



•
Awaiting-Dependency: Paused awaiting asynchronous external reply or event.



•
Completed: Automated execution succeeded and output artifact produced.



•
Progression Stalled / Failed: Transient fault triggers scheduled recovery attempt; permanent failure records failure evidence.



•
Preserved uncertainty: Whether stalled/failed is an alternative outcome, reopening/post-completion behaviour or a presentation defect remains unresolved. Its disconnected diagram node does not make success and failure equivalent or establish post-completion transitions.



\medskip\hrule\medskip

{\small Recorded unresolved matters: source-qualified matters; exact scope remains in the canonical record.}\par
{\footnotesize \href{../../../../continuum/business/business-process/618acdb38fc74952ab360fdde4f7d865.md}{Canonical record: 618acdb38fc74952ab360fdde4f7d865}; full qualifications and provenance.}\par
\subparagraph*{III.4.5.1.3.3.3 To Do Progression Process}
\phantomsection
\label{navigation:III.4.5.1.3.3.3}
\addcontentsline{toc}{subparagraph}{III.4.5.1.3.3.3 To Do Progression Process}

\label{element:c5e4c3076b094c98a628179b91ae03d0}
\index{To Do Progression Process (III.4.5.1.3.3.3)}
{\small BusinessProcess · Canonical · pending\_adjudication}\par

•
Owning Capability: Workflow \& Activity Coordination



•
Activity Archetype: Human Reviewing / Updating / Deciding — clinical review, document countersignature, or administrative authorization task.



•
State Progression Lifecycle:



RAD diagram: textual source notation.
\textbackslash{}begin\{Verbatim\}[breaklines,fontsize=
]
graph LR
    I["Issued"] --> B["In-Inbox"]
    B --> R["Opened / Under-Review"]
    R --> A["Actioned / Decided"]
    D["Dismissed / Delegated\newline{}Disposition/position unresolved"]
    NOTE["Illustrative progression, not a normative transition model.\newline{}Relationship of dismissal/delegation to action/decision is unestablished."]
    classDef note fill:\#fff8dc,stroke:\#8a6d3b,stroke-dasharray:5 5;
    class NOTE,D note;
\textbackslash{}end\{Verbatim\}



•
Key State Dispositions:



•
Issued: To Do generated (e.g., review abnormal potassium result).



•
In-Inbox: Displayed in practitioner's actionable task list.



•
Opened / Under-Review: Clinician inspecting relevant clinical context.



•
Actioned / Decided: Decision executed (e.g., signed off, order placed, acknowledged).



•
Dismissed / Delegated: Task reassigned to registrar or dismissed with comment.



•
Preserved uncertainty: Whether dismissal/delegation is an alternative disposition, a subsequent action or a scoped variant remains unresolved; its disconnected diagram node establishes no transition or position in the illustrative progression.



\medskip\hrule\medskip

{\small Recorded unresolved matters: source-qualified matters; exact scope remains in the canonical record.}\par
{\footnotesize \href{../../../../continuum/business/business-process/c5e4c3076b094c98a628179b91ae03d0.md}{Canonical record: c5e4c3076b094c98a628179b91ae03d0}; full qualifications and provenance.}\par
\subparagraph*{III.4.5.1.3.3.4 Work Order Progression Process}
\phantomsection
\label{navigation:III.4.5.1.3.3.4}
\addcontentsline{toc}{subparagraph}{III.4.5.1.3.3.4 Work Order Progression Process}

\label{element:73fbbe4e027b4805809fad430e791eba}
\index{Work Order Progression Process (III.4.5.1.3.3.4)}
{\small BusinessProcess · Canonical · pending\_adjudication}\par

•
Owning Capability: Workflow \& Activity Coordination



•
Activity Archetype: Human Doing — physical or operational task executed by healthcare staff.



•
State Progression Lifecycle: Created → Ready → Dispatched → Accepted → In-Progress → Completed / Failed / Cancelled



•
Key State Dispositions:



•
Created: Work order instantiated with scope, target, and required skills.



•
Ready: Preconditions satisfied; available for assignment.



•
Dispatched: Sent to target individual or worker group queue.



•
Accepted: Worker claims task.



•
In-Progress: Execution underway.



•
Completed: Task finished with recorded completion telemetry.



\medskip\hrule\medskip

{\small Recorded unresolved matters: source-qualified matters; exact scope remains in the canonical record.}\par
{\footnotesize \href{../../../../continuum/business/business-process/73fbbe4e027b4805809fad430e791eba.md}{Canonical record: 73fbbe4e027b4805809fad430e791eba}; full qualifications and provenance.}\par
\section*{III.4.6 [CT1] Independent Assurance Enablement — Capability context}
\phantomsection
\label{navigation:III.4.6}
\addcontentsline{toc}{section}{III.4.6 [CT1] Independent Assurance Enablement — Capability context}

{\small Canonical Strategy Capability: II.1.3.8 \hyperref[element:616682b612d244429e656b040e002d56]{[CT1] Independent Assurance Enablement}; Canonical ID: \texttt{\seqsplit{616682b612d244429e656b040e002d56}}. Proposed · current.}\par
Contextual reference only; no additional architectural instance or semantic relationship.\par
\subsection*{III.4.6.1 [CT2] Assurance Specification — Capability context}
\phantomsection
\label{navigation:III.4.6.1}
\addcontentsline{toc}{subsection}{III.4.6.1 [CT2] Assurance Specification — Capability context}

{\small Canonical Strategy Capability: II.1.3.8.2 \hyperref[element:b81775d89d774cf8ba61731d13c67a54]{[CT2] Assurance Specification}; Canonical ID: \texttt{\seqsplit{b81775d89d774cf8ba61731d13c67a54}}. Proposed · current.}\par
Contextual reference only; no additional architectural instance or semantic relationship.\par
\subsubsection*{III.4.6.1.1 [CT3] Assurance Criteria Administration — Capability context}
\phantomsection
\label{navigation:III.4.6.1.1}
\addcontentsline{toc}{subsubsection}{III.4.6.1.1 [CT3] Assurance Criteria Administration — Capability context}

{\small Canonical Strategy Capability: II.1.3.8.2.1 \hyperref[element:1bfd71f24321428c8871b587d413f47e]{[CT3] Assurance Criteria Administration}; Canonical ID: \texttt{\seqsplit{1bfd71f24321428c8871b587d413f47e}}. Canonical · current.}\par
Contextual reference only; no additional architectural instance or semantic relationship.\par
\paragraph*{III.4.6.1.1.1 Business Functions}
\phantomsection
\label{navigation:III.4.6.1.1.1}
\addcontentsline{toc}{paragraph}{III.4.6.1.1.1 Business Functions}

\subparagraph*{III.4.6.1.1.1.1 Manage Assurance Criteria}
\phantomsection
\label{navigation:III.4.6.1.1.1.1}
\addcontentsline{toc}{subparagraph}{III.4.6.1.1.1.1 Manage Assurance Criteria}

\label{element:d7e9ba5addd94db2bce164c7c9809e4e}
\index{Manage Assurance Criteria (III.4.6.1.1.1.1)}
{\small BusinessFunction · Canonical · current}\par

•
Canonical Name: Manage Assurance Criteria.



•
Definition: Establishes and manages reusable, governed and temporally identifiable assurance criteria through which governed subjects may be evaluated.



•
Scope: Criteria may express governance rule boundaries, quality constraints, performance expectations, compliance criteria, temporal criteria, evidence requirements or completeness expectations. These examples are not an exhaustive taxonomy. Criteria may have lifecycle, version and temporal applicability.



•
Boundary: Managing a criterion does not transfer authority for the governing requirement from which it may derive. Criteria may operationalise a requirement without silently altering, weakening, strengthening or replacing it; not every criterion must derive from formal Governance. Modelling/managing criteria does not inherently confer approval authority. Detailed criterion lifecycle states and approval workflows remain unresolved.



Criteria management maintains the reusable evaluative basis and its temporal standing. It permits an assessment to be understood against the criterion applicable at the relevant time, rather than today's unqualified version. Management of that basis does not inherently grant approval authority or permit alteration of the governing obligation.


{\small Recorded unresolved matters: source-qualified matters; exact scope remains in the canonical record.}\par
{\footnotesize \href{../../../../continuum/business/business-function/d7e9ba5addd94db2bce164c7c9809e4e.md}{Canonical record: d7e9ba5addd94db2bce164c7c9809e4e}; full qualifications and provenance.}\par
\subsubsection*{III.4.6.1.2 [CT3] Assurance Design — Capability context}
\phantomsection
\label{navigation:III.4.6.1.2}
\addcontentsline{toc}{subsubsection}{III.4.6.1.2 [CT3] Assurance Design — Capability context}

{\small Canonical Strategy Capability: II.1.3.8.2.2 \hyperref[element:71c8a612b88f4ecba049e94d506bd0c2]{[CT3] Assurance Design}; Canonical ID: \texttt{\seqsplit{71c8a612b88f4ecba049e94d506bd0c2}}. Canonical · current.}\par
Contextual reference only; no additional architectural instance or semantic relationship.\par
\paragraph*{III.4.6.1.2.1 Business Functions}
\phantomsection
\label{navigation:III.4.6.1.2.1}
\addcontentsline{toc}{paragraph}{III.4.6.1.2.1 Business Functions}

\subparagraph*{III.4.6.1.2.1.1 Design Service Assurance}
\phantomsection
\label{navigation:III.4.6.1.2.1.1}
\addcontentsline{toc}{subparagraph}{III.4.6.1.2.1.1 Design Service Assurance}

\label{element:6668aa8daf964fec93fce42d48b37ad1}
\index{Design Service Assurance (III.4.6.1.2.1.1)}
{\small BusinessFunction · Canonical · current}\par

•
Canonical Name: Design Service Assurance.



•
Definition: Defines how satisfaction of a governed requirement, constraint or expected behaviour is to be independently established, including the applicable assurance disposition, assurance criteria, required evidence and evaluation expectations.



•
Boundary: Concerns assurance design. It does not establish the underlying governance requirement, grant authority to it, perform the resulting assurance, manage the subject activity or prescribe application or implementation mechanisms. Modelling the resulting Assurance Definition does not grant approval authority. The Strategy peer relationship to behaviour design and failure/recovery design creates no additional Capability or Function.



Designing service assurance turns a requirement's assurance question into an explicit disposition and evaluation basis. The modeller determines how satisfaction is to be assessed, not what the authoritative requirement should mean or who may approve it. Design completion is not evidence that the implemented behaviour satisfies that requirement.


{\small Recorded unresolved matters: source-qualified matters; exact scope remains in the canonical record.}\par
{\footnotesize \href{../../../../continuum/business/business-function/6668aa8daf964fec93fce42d48b37ad1.md}{Canonical record: 6668aa8daf964fec93fce42d48b37ad1}; full qualifications and provenance.}\par
\paragraph*{III.4.6.1.2.2 Business Processes}
\phantomsection
\label{navigation:III.4.6.1.2.2}
\addcontentsline{toc}{paragraph}{III.4.6.1.2.2 Business Processes}

\subparagraph*{III.4.6.1.2.2.1 Assurance Process Design}
\phantomsection
\label{navigation:III.4.6.1.2.2.1}
\addcontentsline{toc}{subparagraph}{III.4.6.1.2.2.1 Assurance Process Design}

\label{element:aec3177dcdd84961814f2ba4f688ad33}
\index{Assurance Process Design (III.4.6.1.2.2.1)}
{\small BusinessProcess · Canonical · current}\par

•
Canonical Name: Assurance Process Design.



•
Owning Capability: Assurance Design (\url{https://github.com/fhirfactory/harmonia/blob/98410db60155198a75f74c79d98060070209faa1/docs/markdown/02-strategy/capabilities/business-enabling-capabilities.md#assurance-design}).



•
Process Purpose: Establishes and maintains how satisfaction of a governed requirement, constraint or expected behaviour is to be assured.



•
Principal Performing Role: Service Assurance Modeller (\url{https://github.com/fhirfactory/harmonia/blob/98410db60155198a75f74c79d98060070209faa1/docs/markdown/03-business-architecture/actors-roles/roles.md#service-assurance-modeller}).



•
Participating Functions: Design Service Assurance (\url{https://github.com/fhirfactory/harmonia/blob/98410db60155198a75f74c79d98060070209faa1/docs/markdown/03-business-architecture/behaviours/health-service-assurance.md#design-service-assurance}) and Manage Assurance Criteria (\url{https://github.com/fhirfactory/harmonia/blob/98410db60155198a75f74c79d98060070209faa1/docs/markdown/03-business-architecture/behaviours/health-service-assurance.md#manage-assurance-criteria}). Manage Assurance Criteria retains ownership within Assurance Criteria Management; its participation does not subsume that Capability's criteria lifecycle or transfer its information responsibility to Assurance Design.



•
Conceptual Output: Assurance Definition (\url{https://github.com/fhirfactory/harmonia/blob/98410db60155198a75f74c79d98060070209faa1/docs/markdown/03-business-architecture/behaviours/health-service-assurance.md#42-assurance-definition--conceptual-modelling-output}), retaining its conceptual Business Architecture output/boundary treatment.



•
Boundary: Governance Authority establishes what is authoritative. Design does not grant Service Assurance Modeller approval authority. Detailed Assurance Definition lifecycle, approval/publication workflows, criteria lifecycle and information modelling remain unresolved. Participation of the two Functions does not establish a cross-capability lifecycle or a new exchange/Service between their owners.



Assurance design establishes a reusable basis for how satisfaction will be evaluated, rather than running the evaluation itself. Criteria-management participation retains its separate owner and does not create an exchange Service or a shared criteria lifecycle. An Assurance Definition can be reused, so each execution need not imply a freshly performed Design Process.


{\small Recorded unresolved matters: source-qualified matters; exact scope remains in the canonical record.}\par
{\footnotesize \href{../../../../continuum/business/business-process/aec3177dcdd84961814f2ba4f688ad33.md}{Canonical record: aec3177dcdd84961814f2ba4f688ad33}; full qualifications and provenance.}\par
\subsection*{III.4.6.2 [CT2] Independent Assurance Evaluation Enablement — Capability context}
\phantomsection
\label{navigation:III.4.6.2}
\addcontentsline{toc}{subsection}{III.4.6.2 [CT2] Independent Assurance Evaluation Enablement — Capability context}

{\small Canonical Strategy Capability: II.1.3.8.3 \hyperref[element:7fa5112efb724b90a17a03c7c19c516a]{[CT2] Independent Assurance Evaluation Enablement}; Canonical ID: \texttt{\seqsplit{7fa5112efb724b90a17a03c7c19c516a}}. Proposed · pending\_adjudication.}\par
Contextual reference only; no additional architectural instance or semantic relationship.\par
\subsubsection*{III.4.6.2.1 [CT3] Governed Assurance — Capability context}
\phantomsection
\label{navigation:III.4.6.2.1}
\addcontentsline{toc}{subsubsection}{III.4.6.2.1 [CT3] Governed Assurance — Capability context}

{\small Canonical Strategy Capability: II.1.3.8.3.1 \hyperref[element:da1a983d4dcf4d5bb772b8f69bc08426]{[CT3] Governed Assurance}; Canonical ID: \texttt{\seqsplit{da1a983d4dcf4d5bb772b8f69bc08426}}. Canonical · pending\_adjudication.}\par
Contextual reference only; no additional architectural instance or semantic relationship.\par
\paragraph*{III.4.6.2.1.1 Business Services}
\phantomsection
\label{navigation:III.4.6.2.1.1}
\addcontentsline{toc}{paragraph}{III.4.6.2.1.1 Business Services}

\subparagraph*{III.4.6.2.1.1.1 Communicate Service Assurance Outcome}
\phantomsection
\label{navigation:III.4.6.2.1.1.1}
\addcontentsline{toc}{subparagraph}{III.4.6.2.1.1.1 Communicate Service Assurance Outcome}

\label{element:371b04d6bbf84b53a3c5667ae6df2809}
\index{Communicate Service Assurance Outcome (III.4.6.2.1.1.1)}
{\small BusinessService · Canonical · pending\_adjudication}\par

•
Canonical Name: Communicate Service Assurance Outcome.



•
Definition: Communicates applicable established Assurance Findings and Conclusions to authorised recipients according to their responsibilities, authority and information requirements.



•
Approved Potential Recipients: Care Coordinator, Service Coordinator, Regulator, Policy Authority, System Steward.



•
Associated Interaction: Service Assurance Outcome Communication (\url{https://github.com/fhirfactory/harmonia/blob/98410db60155198a75f74c79d98060070209faa1/docs/markdown/03-business-architecture/collaborations-interactions/interactions.md#service-assurance-outcome-communication}), with Service Guardian and one or more authorised recipient Roles.



•
Process Relationship: Exposes the communication behaviour of Assurance Process Reporting / Communication (\url{https://github.com/fhirfactory/harmonia/blob/98410db60155198a75f74c79d98060070209faa1/docs/markdown/03-business-architecture/processes/business-processes.md#assurance-process-reporting-communication}), beginning from an already established finding/conclusion. It does not perform assessment or adjudication or alter/reinterpret the established outcome. A documentary report is not required.



•
Recipient Boundary: No mandatory broadcast is established. A particular outcome may be communicated to one, some or all legitimate recipients according to applicability, authority, responsibility and information requirements. Requester, status consumer and outcome recipient need not be the same Role. Communication does not transfer responsibility for recipients' interpretation within their authority, operational/clinical management, remediation, escalation, policy/compliance response, process improvement or other subsequent action to Service Guardian.



\medskip\hrule\medskip

{\small Recorded unresolved matters: source-qualified matters; exact scope remains in the canonical record.}\par
{\footnotesize \href{../../../../continuum/business/business-service/371b04d6bbf84b53a3c5667ae6df2809.md}{Canonical record: 371b04d6bbf84b53a3c5667ae6df2809}; full qualifications and provenance.}\par
\subparagraph*{III.4.6.2.1.1.2 Request Service Assurance}
\phantomsection
\label{navigation:III.4.6.2.1.1.2}
\addcontentsline{toc}{subparagraph}{III.4.6.2.1.1.2 Request Service Assurance}

\label{element:a76c0be369744fcfb5eb91b90c9b59e7}
\index{Request Service Assurance (III.4.6.2.1.1.2)}
{\small BusinessService · Canonical · pending\_adjudication}\par

•
Canonical Name: Request Service Assurance.



•
Definition: Enables an authorised consumer to request that Service Assurance be performed for an applicable governed subject in accordance with an established Assurance Definition.



•
Approved Consumers: Care Coordinator, Service Coordinator, Regulator, Policy Authority, System Steward.



•
Associated Interaction: Service Assurance Request (\url{https://github.com/fhirfactory/harmonia/blob/98410db60155198a75f74c79d98060070209faa1/docs/markdown/03-business-architecture/collaborations-interactions/interactions.md#service-assurance-request}), with Service Guardian and one authorised requesting Role.



•
Process Relationship: Requests Assurance Process Execution (\url{https://github.com/fhirfactory/harmonia/blob/98410db60155198a75f74c79d98060070209faa1/docs/markdown/03-business-architecture/processes/business-processes.md#assurance-process-execution}) under an applicable established Assurance Definition. It does not require a new Assurance Process Design instance for each request.



•
Independence Boundary: The requesting Role identifies the assurance need within its responsibility and authority. Requesting does not grant authority to define or alter the Assurance Definition, alter Assurance Criteria, direct evidence assessment, determine adjudication or control the resulting Assurance Finding / Conclusion. REQ-FND-005 control independence remains applicable, including where a requester performs or manages the subject. Detailed payloads, contracts, initiation mechanics and states remain unresolved.



\medskip\hrule\medskip

{\small Recorded unresolved matters: source-qualified matters; exact scope remains in the canonical record.}\par
{\footnotesize \href{../../../../continuum/business/business-service/a76c0be369744fcfb5eb91b90c9b59e7.md}{Canonical record: a76c0be369744fcfb5eb91b90c9b59e7}; full qualifications and provenance.}\par
\subparagraph*{III.4.6.2.1.1.3 Request Service Assurance Status}
\phantomsection
\label{navigation:III.4.6.2.1.1.3}
\addcontentsline{toc}{subparagraph}{III.4.6.2.1.1.3 Request Service Assurance Status}

\label{element:60f8cbc814614f669f31d4e78734c6d6}
\index{Request Service Assurance Status (III.4.6.2.1.1.3)}
{\small BusinessService · Canonical · pending\_adjudication}\par

•
Canonical Name: Request Service Assurance Status.



•
Definition: Enables the System Steward to obtain the current business progression state of an identifiable Service Assurance activity for stewardship purposes.



•
Only Approved Consumer: System Steward.



•
Associated Interaction: Service Assurance Status (\url{https://github.com/fhirfactory/harmonia/blob/98410db60155198a75f74c79d98060070209faa1/docs/markdown/03-business-architecture/collaborations-interactions/interactions.md#service-assurance-status}), between System Steward and Service Guardian.



•
Status Boundary: Concerns activity progression state. It does not communicate provisional findings, incomplete assessment results as conclusions, predicted outcomes or unadjudicated conclusions. Other requesting Roles do not gain status access merely because they requested assurance. No detailed activity-status taxonomy is established.



Service Assurance Status answers "where is the assurance activity in its progression?" It does not answer "what has the assurance established?" The latter requires an adjudicated Assurance Finding / Conclusion and subsequent outcome communication; neither status nor execution completion establishes that conclusion.



Management Monitoring evaluates activity state in order to progress or manage activity. Assurance evaluates activity against governing criteria in order to establish an assurance conclusion. System Steward's use of this Service is stewardship / management observation of assurance progression, not itself assurance or recursive assurance. The non-recursion boundary in §7.2 applies.



\medskip\hrule\medskip

{\small Recorded unresolved matters: source-qualified matters; exact scope remains in the canonical record.}\par
{\footnotesize \href{../../../../continuum/business/business-service/60f8cbc814614f669f31d4e78734c6d6.md}{Canonical record: 60f8cbc814614f669f31d4e78734c6d6}; full qualifications and provenance.}\par
\paragraph*{III.4.6.2.1.2 Business Functions}
\phantomsection
\label{navigation:III.4.6.2.1.2}
\addcontentsline{toc}{paragraph}{III.4.6.2.1.2 Business Functions}

\subparagraph*{III.4.6.2.1.2.1 Adjudicate Assurance Assessment}
\phantomsection
\label{navigation:III.4.6.2.1.2.1}
\addcontentsline{toc}{subparagraph}{III.4.6.2.1.2.1 Adjudicate Assurance Assessment}

\label{element:7d7cd216415d4b839ea3edfbf1538110}
\index{Adjudicate Assurance Assessment (III.4.6.2.1.2.1)}
{\small BusinessFunction · Canonical · pending\_adjudication}\par

•
Canonical Name: Adjudicate Assurance Assessment.



•
Definition: Determines the assurance finding or conclusion that follows from the assessed evidence and applicable assurance criteria.



•
Boundary: Establishment of the finding/conclusion is part of adjudication. It remains distinct from evidence assessment, operational response, remediation, escalation, workflow progression and clinical judgement. No separate conclusion-establishment or reporting Function is introduced merely because a conclusion is established or may be communicated.



\medskip\hrule\medskip

{\small Recorded unresolved matters: source-qualified matters; exact scope remains in the canonical record.}\par
{\footnotesize \href{../../../../continuum/business/business-function/7d7cd216415d4b839ea3edfbf1538110.md}{Canonical record: 7d7cd216415d4b839ea3edfbf1538110}; full qualifications and provenance.}\par
\subparagraph*{III.4.6.2.1.2.2 Assess Assurance Evidence}
\phantomsection
\label{navigation:III.4.6.2.1.2.2}
\addcontentsline{toc}{subparagraph}{III.4.6.2.1.2.2 Assess Assurance Evidence}

\label{element:e1f6402d47c0455b95b45a1a4d393484}
\index{Assess Assurance Evidence (III.4.6.2.1.2.2)}
{\small BusinessFunction · Canonical · pending\_adjudication}\par

•
Canonical Name: Assess Assurance Evidence.



•
Definition: Evaluates the available trustworthy evidence against the applicable assurance criteria and establishes what that evidence demonstrates for each relevant criterion or assessment concern.



•
Boundary: Assessment remains distinct from adjudication. It may identify satisfaction, exception, insufficiency or other criterion-level results where applicable; no universal result taxonomy is established. Quantitative, qualitative, compound, temporal and confidence-bearing assessment remain possible. Insufficient evidence is not equivalent to either satisfaction or non-satisfaction.



\medskip\hrule\medskip

{\small Recorded unresolved matters: source-qualified matters; exact scope remains in the canonical record.}\par
{\footnotesize \href{../../../../continuum/business/business-function/e1f6402d47c0455b95b45a1a4d393484.md}{Canonical record: e1f6402d47c0455b95b45a1a4d393484}; full qualifications and provenance.}\par
\subparagraph*{III.4.6.2.1.2.3 Establish Assurance Context}
\phantomsection
\label{navigation:III.4.6.2.1.2.3}
\addcontentsline{toc}{subparagraph}{III.4.6.2.1.2.3 Establish Assurance Context}

\label{element:ffffe352e28c4e878d986b0e15ddff15}
\index{Establish Assurance Context (III.4.6.2.1.2.3)}
{\small BusinessFunction · Canonical · pending\_adjudication}\par

•
Canonical Name: Establish Assurance Context.



•
Definition: Establishes the governed subject, applicable assurance purpose and criteria, relevant temporal context, and contextual association of sufficient trustworthy source information as evidence required for an assurance activity.



•
Scope: Includes the evidence-assembly responsibility previously described as Assurance Evidence Assembly; no separate evidence-assembly Function is established. Sufficiency is required for conclusion, not assumed merely because context is established; unavailable or insufficient evidence remains explicit.



•
Information Boundary: Source information remains owned by its responsible source Capability. Its use as evidence does not transfer ownership or originating authority. Governed Assurance owns the contextual association that particular information constitutes evidence for a particular assurance purpose. Evidence association and applicable criteria may each be temporal and version-specific; association distinguishes the information state relevant to assurance from its current state.



•
Realisation Boundary: Uses generic EC-02 context-management capability together with EC-14 assurance-specific semantics as documented in §4.1 (\url{https://github.com/fhirfactory/harmonia/blob/98410db60155198a75f74c79d98060070209faa1/docs/markdown/03-business-architecture/behaviours/health-service-assurance.md#41-context-management-contribution}). Establishment of Assurance Context prescribes no copying, persistence, aggregation, caching or other technical mechanism.



\medskip\hrule\medskip

{\small Recorded unresolved matters: source-qualified matters; exact scope remains in the canonical record.}\par
{\footnotesize \href{../../../../continuum/business/business-function/ffffe352e28c4e878d986b0e15ddff15.md}{Canonical record: ffffe352e28c4e878d986b0e15ddff15}; full qualifications and provenance.}\par
\paragraph*{III.4.6.2.1.3 Business Processes}
\phantomsection
\label{navigation:III.4.6.2.1.3}
\addcontentsline{toc}{paragraph}{III.4.6.2.1.3 Business Processes}

\subparagraph*{III.4.6.2.1.3.1 Assurance Process Execution}
\phantomsection
\label{navigation:III.4.6.2.1.3.1}
\addcontentsline{toc}{subparagraph}{III.4.6.2.1.3.1 Assurance Process Execution}

\label{element:3c67e57403b041bb8ccdaf3980b6ff94}
\index{Assurance Process Execution (III.4.6.2.1.3.1)}
{\small BusinessProcess · Canonical · pending\_adjudication}\par

•
Canonical Name: Assurance Process Execution.



•
Owning Capability: Governed Assurance (\url{https://github.com/fhirfactory/harmonia/blob/98410db60155198a75f74c79d98060070209faa1/docs/markdown/02-strategy/capabilities/business-enabling-capabilities.md#governed-assurance}).



•
Process Purpose: Performs defined assurance for an identifiable governed subject in accordance with an applicable Assurance Definition and establishes an Assurance Finding / Conclusion from applicable trustworthy evidence.



•
Principal Performing Role: Service Guardian (\url{https://github.com/fhirfactory/harmonia/blob/98410db60155198a75f74c79d98060070209faa1/docs/markdown/03-business-architecture/actors-roles/roles.md#service-guardian}).



•
Participating Functions: Establish Assurance Context (\url{https://github.com/fhirfactory/harmonia/blob/98410db60155198a75f74c79d98060070209faa1/docs/markdown/03-business-architecture/behaviours/health-service-assurance.md#establish-assurance-context}), Assess Assurance Evidence (\url{https://github.com/fhirfactory/harmonia/blob/98410db60155198a75f74c79d98060070209faa1/docs/markdown/03-business-architecture/behaviours/health-service-assurance.md#assess-assurance-evidence}) and Adjudicate Assurance Assessment (\url{https://github.com/fhirfactory/harmonia/blob/98410db60155198a75f74c79d98060070209faa1/docs/markdown/03-business-architecture/behaviours/health-service-assurance.md#adjudicate-assurance-assessment}), all retaining Governed Assurance ownership.



•
Conceptual Basis and Output: Applicable Assurance Definition plus identifiable governed subject; Assurance Finding / Conclusion established through adjudication where sufficient trustworthy evidence permits. Correct execution with insufficient evidence does not imply satisfaction or non-satisfaction or force a conclusion.



•
Boundary: Function completion does not define Process lifecycle states. No Context Established, Evidence Assessed, Adjudicated or Concluded states are inferred. Detailed lifecycle states, transitions and initiation mechanics remain unresolved. Requesting this Process does not grant control of evidence assessment, adjudication or the resulting finding/conclusion.



RAD diagram: textual source notation.
\textbackslash{}begin\{Verbatim\}[breaklines,fontsize=
]
graph TD
    DEF["Applicable Assurance Definition\newline{}Conceptual basis"] -->|Assurance basis| EXEC["PR: Assurance Process Execution\newline{}Owner: Governed Assurance"]
    SUBJECT["Identifiable Governed Subject"] -->|Subject of assurance| EXEC
    EXEC -->|Participating Function| CONTEXT["FN: Establish Assurance Context"]
    EXEC -->|Participating Function| ASSESS["FN: Assess Assurance Evidence"]
    EXEC -->|Participating Function| ADJ["FN: Adjudicate Assurance Assessment"]
    EXEC -->|Establishes through adjudication| FINDING["Assurance Finding / Conclusion"]
\textbackslash{}end\{Verbatim\}



\medskip\hrule\medskip

{\small Recorded unresolved matters: source-qualified matters; exact scope remains in the canonical record.}\par
{\footnotesize \href{../../../../continuum/business/business-process/3c67e57403b041bb8ccdaf3980b6ff94.md}{Canonical record: 3c67e57403b041bb8ccdaf3980b6ff94}; full qualifications and provenance.}\par
\subparagraph*{III.4.6.2.1.3.2 Assurance Process Reporting / Communication}
\phantomsection
\label{navigation:III.4.6.2.1.3.2}
\addcontentsline{toc}{subparagraph}{III.4.6.2.1.3.2 Assurance Process Reporting / Communication}

\label{element:2fc0b626577b43feb050bcd91ac2b0a9}
\index{Assurance Process Reporting / Communication (III.4.6.2.1.3.2)}
{\small BusinessProcess · Canonical · pending\_adjudication}\par

•
Canonical Name: Assurance Process Reporting / Communication.



•
Owning Capability: Governed Assurance (\url{https://github.com/fhirfactory/harmonia/blob/98410db60155198a75f74c79d98060070209faa1/docs/markdown/02-strategy/capabilities/business-enabling-capabilities.md#governed-assurance}).



•
Process Purpose: Communicates applicable established Assurance Findings and Conclusions to legitimate recipients according to their responsibilities, authority and information requirements, without altering the established assurance outcome or assuming responsibility for the recipient's subsequent action.



•
Starting Basis: An already established Assurance Finding / Conclusion. This Process does not perform assessment or adjudication, alter or reinterpret the established outcome.



•
Communication Participation: Service Guardian participates through Service Assurance Outcome Communication (\url{https://github.com/fhirfactory/harmonia/blob/98410db60155198a75f74c79d98060070209faa1/docs/markdown/03-business-architecture/collaborations-interactions/interactions.md#service-assurance-outcome-communication}), associated with Communicate Service Assurance Outcome (\url{https://github.com/fhirfactory/harmonia/blob/98410db60155198a75f74c79d98060070209faa1/docs/markdown/03-business-architecture/behaviours/health-service-assurance.md#communicate-service-assurance-outcome}). No additional Function is introduced to populate this Process.



•
Boundary: Reporting / Communication is not limited to documentary report production; an appropriate Business Interaction may communicate the outcome without a report document. It assumes no operational response, remediation, escalation, policy/compliance response, process improvement, clinical management, clinical judgement or other recipient action. Detailed communication progression and completion semantics remain unresolved.



\medskip\hrule\medskip

{\small Recorded unresolved matters: source-qualified matters; exact scope remains in the canonical record.}\par
{\footnotesize \href{../../../../continuum/business/business-process/2fc0b626577b43feb050bcd91ac2b0a9.md}{Canonical record: 2fc0b626577b43feb050bcd91ac2b0a9}; full qualifications and provenance.}\par
\section*{III.4.7 [CT1] Shared Information \& Coordination Enablement — Capability context}
\phantomsection
\label{navigation:III.4.7}
\addcontentsline{toc}{section}{III.4.7 [CT1] Shared Information \& Coordination Enablement — Capability context}

{\small Canonical Strategy Capability: II.1.3.9 \hyperref[element:0062fcacae1740988b6be2bda2aa1674]{[CT1] Shared Information \& Coordination Enablement}; Canonical ID: \texttt{\seqsplit{0062fcacae1740988b6be2bda2aa1674}}. Proposed · current.}\par
Contextual reference only; no additional architectural instance or semantic relationship.\par
\subsection*{III.4.7.1 [CT2] Calendar Management — Capability context}
\phantomsection
\label{navigation:III.4.7.1}
\addcontentsline{toc}{subsection}{III.4.7.1 [CT2] Calendar Management — Capability context}

{\small Canonical Strategy Capability: II.1.3.9.1 \hyperref[element:7669036f7e6749219412fc64df5f7064]{[CT2] Calendar Management}; Canonical ID: \texttt{\seqsplit{7669036f7e6749219412fc64df5f7064}}. Canonical · current.}\par
Contextual reference only; no additional architectural instance or semantic relationship.\par
\subsubsection*{III.4.7.1.1 Business Services}
\phantomsection
\label{navigation:III.4.7.1.1}
\addcontentsline{toc}{subsubsection}{III.4.7.1.1 Business Services}

\paragraph*{III.4.7.1.1.1 Temporal Context Resolution}
\phantomsection
\label{navigation:III.4.7.1.1.1}
\addcontentsline{toc}{paragraph}{III.4.7.1.1.1 Temporal Context Resolution}

\label{element:47f6491361794c2393cd3b64caec2d1a}
\index{Temporal Context Resolution (III.4.7.1.1.1)}
{\small BusinessService · Canonical · current}\par

Resolves operational shifts.


{\footnotesize \href{../../../../continuum/business/business-service/47f6491361794c2393cd3b64caec2d1a.md}{Canonical record: 47f6491361794c2393cd3b64caec2d1a}; full qualifications and provenance.}\par
\paragraph*{III.4.7.1.1.2 Unified Calendar Projection Query}
\phantomsection
\label{navigation:III.4.7.1.1.2}
\addcontentsline{toc}{paragraph}{III.4.7.1.1.2 Unified Calendar Projection Query}

\label{element:c78bdb7ffc7d42fda450c024ca252322}
\index{Unified Calendar Projection Query (III.4.7.1.1.2)}
{\small BusinessService · Canonical · current}\par

Discloses consolidated temporal schedules.


{\footnotesize \href{../../../../continuum/business/business-service/c78bdb7ffc7d42fda450c024ca252322.md}{Canonical record: c78bdb7ffc7d42fda450c024ca252322}; full qualifications and provenance.}\par
\subsubsection*{III.4.7.1.2 Business Functions}
\phantomsection
\label{navigation:III.4.7.1.2}
\addcontentsline{toc}{subsubsection}{III.4.7.1.2 Business Functions}

\paragraph*{III.4.7.1.2.1 Correlate Activity with Temporal Context}
\phantomsection
\label{navigation:III.4.7.1.2.1}
\addcontentsline{toc}{paragraph}{III.4.7.1.2.1 Correlate Activity with Temporal Context}

\label{element:b45fa964568843398661cc119b4fc693}
\index{Correlate Activity with Temporal Context (III.4.7.1.2.1)}
{\small BusinessFunction · Canonical · current}\par

Associates clinical events with shift boundaries and appointment slots.


{\footnotesize \href{../../../../continuum/business/business-function/b45fa964568843398661cc119b4fc693.md}{Canonical record: b45fa964568843398661cc119b4fc693}; full qualifications and provenance.}\par
\paragraph*{III.4.7.1.2.2 Project Operational Schedule}
\phantomsection
\label{navigation:III.4.7.1.2.2}
\addcontentsline{toc}{paragraph}{III.4.7.1.2.2 Project Operational Schedule}

\label{element:3ca2e6a5ca614df991a2e1a060f0f1d8}
\index{Project Operational Schedule (III.4.7.1.2.2)}
{\small BusinessFunction · Canonical · current}\par

Aggregates and overlays multidisciplinary schedules onto a unified temporal view.


{\footnotesize \href{../../../../continuum/business/business-function/3ca2e6a5ca614df991a2e1a060f0f1d8.md}{Canonical record: 3ca2e6a5ca614df991a2e1a060f0f1d8}; full qualifications and provenance.}\par
\subsection*{III.4.7.2 [CT2] Clinical Collaboration — Capability context}
\phantomsection
\label{navigation:III.4.7.2}
\addcontentsline{toc}{subsection}{III.4.7.2 [CT2] Clinical Collaboration — Capability context}

{\small Canonical Strategy Capability: II.1.3.9.2 \hyperref[element:e6429225463a44d39f82e490e254e5d4]{[CT2] Clinical Collaboration}; Canonical ID: \texttt{\seqsplit{e6429225463a44d39f82e490e254e5d4}}. Canonical · current.}\par
Contextual reference only; no additional architectural instance or semantic relationship.\par
\subsubsection*{III.4.7.2.1 Business Services}
\phantomsection
\label{navigation:III.4.7.2.1}
\addcontentsline{toc}{subsubsection}{III.4.7.2.1 Business Services}

\paragraph*{III.4.7.2.1.1 Clinical Collaboration Space Service}
\phantomsection
\label{navigation:III.4.7.2.1.1}
\addcontentsline{toc}{paragraph}{III.4.7.2.1.1 Clinical Collaboration Space Service}

\label{element:17579c34d3014825a878a9ff121e30a2}
\index{Clinical Collaboration Space Service (III.4.7.2.1.1)}
{\small BusinessService · Canonical · current}\par

Manages team collaboration spaces and discussions.


{\footnotesize \href{../../../../continuum/business/business-service/17579c34d3014825a878a9ff121e30a2.md}{Canonical record: 17579c34d3014825a878a9ff121e30a2}; full qualifications and provenance.}\par
\paragraph*{III.4.7.2.1.2 Collaboration Clinical Summary Resolution}
\phantomsection
\label{navigation:III.4.7.2.1.2}
\addcontentsline{toc}{paragraph}{III.4.7.2.1.2 Collaboration Clinical Summary Resolution}

\label{element:6a25de08f1b5461b86bd410ea21b80a6}
\index{Collaboration Clinical Summary Resolution (III.4.7.2.1.2)}
{\small BusinessService · Canonical · current}\par

Embeds clinical summaries in collaboration feeds.


{\footnotesize \href{../../../../continuum/business/business-service/6a25de08f1b5461b86bd410ea21b80a6.md}{Canonical record: 6a25de08f1b5461b86bd410ea21b80a6}; full qualifications and provenance.}\par
\paragraph*{III.4.7.2.1.3 Collaboration Space Governance Service}
\phantomsection
\label{navigation:III.4.7.2.1.3}
\addcontentsline{toc}{paragraph}{III.4.7.2.1.3 Collaboration Space Governance Service}

\label{element:f94eef008b214ed88ff92bc64640ccb0}
\index{Collaboration Space Governance Service (III.4.7.2.1.3)}
{\small BusinessService · Canonical · current}\par

Governs space lifecycles and membership.


{\footnotesize \href{../../../../continuum/business/business-service/f94eef008b214ed88ff92bc64640ccb0.md}{Canonical record: f94eef008b214ed88ff92bc64640ccb0}; full qualifications and provenance.}\par
\subsubsection*{III.4.7.2.2 Business Functions}
\phantomsection
\label{navigation:III.4.7.2.2}
\addcontentsline{toc}{subsubsection}{III.4.7.2.2 Business Functions}

\paragraph*{III.4.7.2.2.1 Coordinate Care-Team Collaboration}
\phantomsection
\label{navigation:III.4.7.2.2.1}
\addcontentsline{toc}{paragraph}{III.4.7.2.2.1 Coordinate Care-Team Collaboration}

\label{element:5443b253f0e4425f984702d396e77a3d}
\index{Coordinate Care-Team Collaboration (III.4.7.2.2.1)}
{\small BusinessFunction · Canonical · current}\par

Manages care-team communication spaces, participant memberships, and clinical handovers.



This Function maintains clinical collaboration spaces and participant context, enabling team discourse around care. The governed source also says it manages clinical handovers, which conflicts with the metamodel's explicit exclusion of that authority. The conflict is preserved for adjudication; this explanatory catalogue does not resolve it by quietly deleting either statement.


{\small Recorded unresolved matters: source-qualified matters; exact scope remains in the canonical record.}\par
{\footnotesize \href{../../../../continuum/business/business-function/5443b253f0e4425f984702d396e77a3d.md}{Canonical record: 5443b253f0e4425f984702d396e77a3d}; full qualifications and provenance.}\par
\paragraph*{III.4.7.2.2.2 Govern Collaboration Metadata}
\phantomsection
\label{navigation:III.4.7.2.2.2}
\addcontentsline{toc}{paragraph}{III.4.7.2.2.2 Govern Collaboration Metadata}

\label{element:94a92326ef2141efb414c4602d878caa}
\index{Govern Collaboration Metadata (III.4.7.2.2.2)}
{\small BusinessFunction · Canonical · current}\par

Enforces privacy policies, participant access gates, and zero-PHI space metadata rules.


{\footnotesize \href{../../../../continuum/business/business-function/94a92326ef2141efb414c4602d878caa.md}{Canonical record: 94a92326ef2141efb414c4602d878caa}; full qualifications and provenance.}\par
\paragraph*{III.4.7.2.2.3 Resolve LHR Query within Collaboration}
\phantomsection
\label{navigation:III.4.7.2.2.3}
\addcontentsline{toc}{paragraph}{III.4.7.2.2.3 Resolve LHR Query within Collaboration}

\label{element:1d4ac561d963484480e6b06a09416616}
\index{Resolve LHR Query within Collaboration (III.4.7.2.2.3)}
{\small BusinessFunction · Canonical · current}\par

Contextually projects active patient summaries directly into collaborative clinical discussion channels.


{\footnotesize \href{../../../../continuum/business/business-function/1d4ac561d963484480e6b06a09416616.md}{Canonical record: 1d4ac561d963484480e6b06a09416616}; full qualifications and provenance.}\par
\subsection*{III.4.7.3 [CT2] Clinical Information Stewardship Enablement — Capability context [superseded]}
\phantomsection
\label{navigation:III.4.7.3}
\addcontentsline{toc}{subsection}{III.4.7.3 [CT2] Clinical Information Stewardship Enablement — Capability context [superseded]}

{\small Canonical Strategy Capability: II.1.3.9.3 \hyperref[element:3811c758ff2140c39a801bf89f75098a]{[CT2] Clinical Information Stewardship Enablement}; Canonical ID: \texttt{\seqsplit{3811c758ff2140c39a801bf89f75098a}}. Proposed · superseded.}\par
Contextual reference only; no additional architectural instance or semantic relationship.\par
\subsubsection*{III.4.7.3.1 [CT3] Information Authority Representation — Capability context}
\phantomsection
\label{navigation:III.4.7.3.1}
\addcontentsline{toc}{subsubsection}{III.4.7.3.1 [CT3] Information Authority Representation — Capability context}

{\small Canonical Strategy Capability: II.1.3.9.3.1 \hyperref[element:8c850faafb334f97a81a0b4cd8fdb65f]{[CT3] Information Authority Representation}; Canonical ID: \texttt{\seqsplit{8c850faafb334f97a81a0b4cd8fdb65f}}. Proposed · pending\_adjudication.}\par
Contextual reference only; no additional architectural instance or semantic relationship.\par
\paragraph*{III.4.7.3.1.1 Business Functions}
\phantomsection
\label{navigation:III.4.7.3.1.1}
\addcontentsline{toc}{paragraph}{III.4.7.3.1.1 Business Functions}

\subparagraph*{III.4.7.3.1.1.1 Qualify Information Authority Standing}
\phantomsection
\label{navigation:III.4.7.3.1.1.1}
\addcontentsline{toc}{subparagraph}{III.4.7.3.1.1.1 Qualify Information Authority Standing}

\label{element:f5b57bbba47a4e86bd2f79f700d6977a}
\index{Qualify Information Authority Standing (III.4.7.3.1.1.1)}
{\small BusinessFunction · Proposed · pending\_adjudication}\par

Associate source and credibility qualifications with an information assertion under applicable governed semantics.



\textbf{Overview}



Moves an unqualified received assertion to explicitly attributed authority standing without equating persistence success to credibility.



\textbf{Boundary}



Does not adjudicate all competing claims or originate clinical truth.


{\small Recorded unresolved matters: source-qualified matters; exact scope remains in the canonical record.}\par
{\footnotesize \href{../../../../continuum/business/business-function/f5b57bbba47a4e86bd2f79f700d6977a.md}{Canonical record: f5b57bbba47a4e86bd2f79f700d6977a}; full qualifications and provenance.}\par
\subsection*{III.4.7.4 [CT2] Clinical Knowledge Services — Capability context}
\phantomsection
\label{navigation:III.4.7.4}
\addcontentsline{toc}{subsection}{III.4.7.4 [CT2] Clinical Knowledge Services — Capability context}

{\small Canonical Strategy Capability: II.1.3.9.4 \hyperref[element:bbd2bd21859f40e39963bdbe7a851686]{[CT2] Clinical Knowledge Services}; Canonical ID: \texttt{\seqsplit{bbd2bd21859f40e39963bdbe7a851686}}. Canonical · current.}\par
Contextual reference only; no additional architectural instance or semantic relationship.\par
\subsubsection*{III.4.7.4.1 [CT3] Governed Reference Dataset Administration — Capability context}
\phantomsection
\label{navigation:III.4.7.4.1}
\addcontentsline{toc}{subsubsection}{III.4.7.4.1 [CT3] Governed Reference Dataset Administration — Capability context}

{\small Canonical Strategy Capability: II.1.3.9.4.3 \hyperref[element:682b8684ad9e49d5a5c348a148081607]{[CT3] Governed Reference Dataset Administration}; Canonical ID: \texttt{\seqsplit{682b8684ad9e49d5a5c348a148081607}}. Proposed · current.}\par
Contextual reference only; no additional architectural instance or semantic relationship.\par
\paragraph*{III.4.7.4.1.1 Business Services}
\phantomsection
\label{navigation:III.4.7.4.1.1}
\addcontentsline{toc}{paragraph}{III.4.7.4.1.1 Business Services}

\subparagraph*{III.4.7.4.1.1.1 Governed Reference Dataset Content Query}
\phantomsection
\label{navigation:III.4.7.4.1.1.1}
\addcontentsline{toc}{subparagraph}{III.4.7.4.1.1.1 Governed Reference Dataset Content Query}

\label{element:625ef1036a7442b2b8aed06f72b8afff}
\index{Governed Reference Dataset Content Query (III.4.7.4.1.1.1)}
{\small BusinessService · Proposed · current}\par

Provide authorised consumers with approved content at an identified revision and applicability basis.



\textbf{Overview}



Consumers can resolve the membership or mapping content they actually used without treating build-time model definitions as runtime data.



\textbf{Boundary}



Does not grant unrestricted dataset mutation or external terminology authority.


{\footnotesize \href{../../../../continuum/business/business-service/625ef1036a7442b2b8aed06f72b8afff.md}{Canonical record: 625ef1036a7442b2b8aed06f72b8afff}; full qualifications and provenance.}\par
\paragraph*{III.4.7.4.1.2 Business Functions}
\phantomsection
\label{navigation:III.4.7.4.1.2}
\addcontentsline{toc}{paragraph}{III.4.7.4.1.2 Business Functions}

\subparagraph*{III.4.7.4.1.2.1 Accept Governed Dataset Revision}
\phantomsection
\label{navigation:III.4.7.4.1.2.1}
\addcontentsline{toc}{subparagraph}{III.4.7.4.1.2.1 Accept Governed Dataset Revision}

\label{element:96c1c089e7f2443dbd3ad5beab6a3011}
\index{Accept Governed Dataset Revision (III.4.7.4.1.2.1)}
{\small BusinessFunction · Proposed · current}\par

Apply an explicitly authorised content revision and its effective standing through governed information-state progression.



\textbf{Overview}



Moves a reviewed change to the applicable approved content revision, retaining attribution and uncertainty if durable establishment cannot be confirmed.



\textbf{Boundary}



The Function does not invent who approves or become a persistence manager.


{\footnotesize \href{../../../../continuum/business/business-function/96c1c089e7f2443dbd3ad5beab6a3011.md}{Canonical record: 96c1c089e7f2443dbd3ad5beab6a3011}; full qualifications and provenance.}\par
\subsubsection*{III.4.7.4.2 Business Services}
\phantomsection
\label{navigation:III.4.7.4.2}
\addcontentsline{toc}{subsubsection}{III.4.7.4.2 Business Services}

\paragraph*{III.4.7.4.2.1 Terminology Concept Resolution Service}
\phantomsection
\label{navigation:III.4.7.4.2.1}
\addcontentsline{toc}{paragraph}{III.4.7.4.2.1 Terminology Concept Resolution Service}

\label{element:3be22abcc49045c7a609bdf6428fe15a}
\index{Terminology Concept Resolution Service (III.4.7.4.2.1)}
{\small BusinessService · Canonical · current}\par

Exposes concept validation lookups.


{\footnotesize \href{../../../../continuum/business/business-service/3be22abcc49045c7a609bdf6428fe15a.md}{Canonical record: 3be22abcc49045c7a609bdf6428fe15a}; full qualifications and provenance.}\par
\paragraph*{III.4.7.4.2.2 Terminology Mapping Service}
\phantomsection
\label{navigation:III.4.7.4.2.2}
\addcontentsline{toc}{paragraph}{III.4.7.4.2.2 Terminology Mapping Service}

\label{element:27dfb51817344e17ab78aa6f33f959ba}
\index{Terminology Mapping Service (III.4.7.4.2.2)}
{\small BusinessService · Canonical · current}\par

Translates codes across terminology systems.


{\footnotesize \href{../../../../continuum/business/business-service/27dfb51817344e17ab78aa6f33f959ba.md}{Canonical record: 27dfb51817344e17ab78aa6f33f959ba}; full qualifications and provenance.}\par
\subsubsection*{III.4.7.4.3 Business Functions}
\phantomsection
\label{navigation:III.4.7.4.3}
\addcontentsline{toc}{subsubsection}{III.4.7.4.3 Business Functions}

\paragraph*{III.4.7.4.3.1 Map Clinical Concept}
\phantomsection
\label{navigation:III.4.7.4.3.1}
\addcontentsline{toc}{paragraph}{III.4.7.4.3.1 Map Clinical Concept}

\label{element:9549208d56cd47d8a9ba4e586323ef3b}
\index{Map Clinical Concept (III.4.7.4.3.1)}
{\small BusinessFunction · Canonical · current}\par

Translates local lab codes and legacy terms to canonical SNOMED/LOINC codes using governed translation maps.


{\footnotesize \href{../../../../continuum/business/business-function/9549208d56cd47d8a9ba4e586323ef3b.md}{Canonical record: 9549208d56cd47d8a9ba4e586323ef3b}; full qualifications and provenance.}\par
\paragraph*{III.4.7.4.3.2 Resolve Clinical Concept}
\phantomsection
\label{navigation:III.4.7.4.3.2}
\addcontentsline{toc}{paragraph}{III.4.7.4.3.2 Resolve Clinical Concept}

\label{element:8f89eacbc53b4697b5a747db2e143510}
\index{Resolve Clinical Concept (III.4.7.4.3.2)}
{\small BusinessFunction · Canonical · current}\par

Validates concept identifiers and retrieves preferred clinical rubrics and hierarchy paths.


{\footnotesize \href{../../../../continuum/business/business-function/8f89eacbc53b4697b5a747db2e143510.md}{Canonical record: 8f89eacbc53b4697b5a747db2e143510}; full qualifications and provenance.}\par
\subsection*{III.4.7.5 [CT2] General Collaboration — Capability context}
\phantomsection
\label{navigation:III.4.7.5}
\addcontentsline{toc}{subsection}{III.4.7.5 [CT2] General Collaboration — Capability context}

{\small Canonical Strategy Capability: II.1.3.9.5 \hyperref[element:3bb740e6db284112b054905c0c6c80be]{[CT2] General Collaboration}; Canonical ID: \texttt{\seqsplit{3bb740e6db284112b054905c0c6c80be}}. Canonical · current.}\par
Contextual reference only; no additional architectural instance or semantic relationship.\par
\subsubsection*{III.4.7.5.1 Business Services}
\phantomsection
\label{navigation:III.4.7.5.1}
\addcontentsline{toc}{subsubsection}{III.4.7.5.1 Business Services}

\paragraph*{III.4.7.5.1.1 Operational Collaboration Service}
\phantomsection
\label{navigation:III.4.7.5.1.1}
\addcontentsline{toc}{paragraph}{III.4.7.5.1.1 Operational Collaboration Service}

\label{element:4096602ad3084bc2a9ad9a06276cbcdc}
\index{Operational Collaboration Service (III.4.7.5.1.1)}
{\small BusinessService · Canonical · current}\par

Provides non-clinical operational messaging channels.


{\footnotesize \href{../../../../continuum/business/business-service/4096602ad3084bc2a9ad9a06276cbcdc.md}{Canonical record: 4096602ad3084bc2a9ad9a06276cbcdc}; full qualifications and provenance.}\par
\subsubsection*{III.4.7.5.2 Business Functions}
\phantomsection
\label{navigation:III.4.7.5.2}
\addcontentsline{toc}{subsubsection}{III.4.7.5.2 Business Functions}

\paragraph*{III.4.7.5.2.1 Coordinate Operational Collaboration}
\phantomsection
\label{navigation:III.4.7.5.2.1}
\addcontentsline{toc}{paragraph}{III.4.7.5.2.1 Coordinate Operational Collaboration}

\label{element:3b5604d966504661af2cb6654ae0de5c}
\index{Coordinate Operational Collaboration (III.4.7.5.2.1)}
{\small BusinessFunction · Canonical · current}\par

Provides communication spaces for facility logistics, disaster management teams, and administrative committees.


{\footnotesize \href{../../../../continuum/business/business-function/3b5604d966504661af2cb6654ae0de5c.md}{Canonical record: 3b5604d966504661af2cb6654ae0de5c}; full qualifications and provenance.}\par
\subsection*{III.4.7.6 [CT2] Health Information Access — Capability context}
\phantomsection
\label{navigation:III.4.7.6}
\addcontentsline{toc}{subsection}{III.4.7.6 [CT2] Health Information Access — Capability context}

{\small Canonical Strategy Capability: II.1.3.9.6 \hyperref[element:cc253aea3c5f489da0502df74dde1669]{[CT2] Health Information Access}; Canonical ID: \texttt{\seqsplit{cc253aea3c5f489da0502df74dde1669}}. Canonical · current.}\par
Contextual reference only; no additional architectural instance or semantic relationship.\par
\subsubsection*{III.4.7.6.1 Business Services}
\phantomsection
\label{navigation:III.4.7.6.1}
\addcontentsline{toc}{subsubsection}{III.4.7.6.1 Business Services}

\paragraph*{III.4.7.6.1.1 Clinical Information Search Service}
\phantomsection
\label{navigation:III.4.7.6.1.1}
\addcontentsline{toc}{paragraph}{III.4.7.6.1.1 Clinical Information Search Service}

\label{element:ef87ad2b419443a9989849bf18dfb385}
\index{Clinical Information Search Service (III.4.7.6.1.1)}
{\small BusinessService · Canonical · current}\par

Discloses search capabilities to clinical user interfaces.


{\footnotesize \href{../../../../continuum/business/business-service/ef87ad2b419443a9989849bf18dfb385.md}{Canonical record: ef87ad2b419443a9989849bf18dfb385}; full qualifications and provenance.}\par
\paragraph*{III.4.7.6.1.2 Clinical Resource Retrieval Service}
\phantomsection
\label{navigation:III.4.7.6.1.2}
\addcontentsline{toc}{paragraph}{III.4.7.6.1.2 Clinical Resource Retrieval Service}

\label{element:3cdde3bb6a204c46b08468e080f2545a}
\index{Clinical Resource Retrieval Service (III.4.7.6.1.2)}
{\small BusinessService · Canonical · current}\par

Delivers structured clinical information to requesting applications.



Feature association not established: FEAT-ISE-09 — Access-Controlled Information Filtering requires access-controlled disclosure/filtering; retrieval of discrete information alone does not materially realise that responsibility. No replacement Feature is assigned.


{\small Recorded unresolved matters: source-qualified matters; exact scope remains in the canonical record.}\par
{\footnotesize \href{../../../../continuum/business/business-service/3cdde3bb6a204c46b08468e080f2545a.md}{Canonical record: 3cdde3bb6a204c46b08468e080f2545a}; full qualifications and provenance.}\par
\subsubsection*{III.4.7.6.2 Business Functions}
\phantomsection
\label{navigation:III.4.7.6.2}
\addcontentsline{toc}{subsubsection}{III.4.7.6.2 Business Functions}

\paragraph*{III.4.7.6.2.1 Apply Access-Qualified Result Filtering}
\phantomsection
\label{navigation:III.4.7.6.2.1}
\addcontentsline{toc}{paragraph}{III.4.7.6.2.1 Apply Access-Qualified Result Filtering}

\label{element:598e46ab52644e70809b3039b7a4c21e}
\index{Apply Access-Qualified Result Filtering (III.4.7.6.2.1)}
{\small BusinessFunction · Canonical · current}\par

Redacts or masks search results and clinical content based on the caller's verified security context and patient consent directives.



Result filtering changes what can be disclosed according to the verified context and applicable consent determinations. It applies established control outcomes without owning their policy. Its capability-scoped standing is retained instead of promoting the behaviour to a Feature simply to obtain traceability symmetry.


{\small Recorded unresolved matters: source-qualified matters; exact scope remains in the canonical record.}\par
{\footnotesize \href{../../../../continuum/business/business-function/598e46ab52644e70809b3039b7a4c21e.md}{Canonical record: 598e46ab52644e70809b3039b7a4c21e}; full qualifications and provenance.}\par
\paragraph*{III.4.7.6.2.2 Retrieve Health Information}
\phantomsection
\label{navigation:III.4.7.6.2.2}
\addcontentsline{toc}{paragraph}{III.4.7.6.2.2 Retrieve Health Information}

\label{element:0afb265973744f6f9e95ee54e6ea5b77}
\index{Retrieve Health Information (III.4.7.6.2.2)}
{\small BusinessFunction · Canonical · current}\par

Fetches discrete clinical observations and original attachments.



Feature association not established: FEAT-ISE-09 — Access-Controlled Information Filtering requires access-controlled disclosure/filtering; retrieval of discrete information alone does not materially realise that responsibility. No replacement Feature is assigned.



Retrieval obtains discrete clinical information or attachments under the owning access responsibility. It is distinct from applying the access-qualified filtering required by FEAT-ISE-09. A shared owner and related data do not justify making that unresolved realisation association authoritative.


{\small Recorded unresolved matters: source-qualified matters; exact scope remains in the canonical record.}\par
{\footnotesize \href{../../../../continuum/business/business-function/0afb265973744f6f9e95ee54e6ea5b77.md}{Canonical record: 0afb265973744f6f9e95ee54e6ea5b77}; full qualifications and provenance.}\par
\paragraph*{III.4.7.6.2.3 Search Health Information}
\phantomsection
\label{navigation:III.4.7.6.2.3}
\addcontentsline{toc}{paragraph}{III.4.7.6.2.3 Search Health Information}

\label{element:2e75f678b9814305b1cb2e810f4f0e9e}
\index{Search Health Information (III.4.7.6.2.3)}
{\small BusinessFunction · Canonical · current}\par

Executes structured semantic searches across diagnostic reports, clinical notes, and medication histories.


{\footnotesize \href{../../../../continuum/business/business-function/2e75f678b9814305b1cb2e810f4f0e9e.md}{Canonical record: 2e75f678b9814305b1cb2e810f4f0e9e}; full qualifications and provenance.}\par
\subsection*{III.4.7.7 [CT2] Health Information Connectivity Enablement — Capability context [superseded]}
\phantomsection
\label{navigation:III.4.7.7}
\addcontentsline{toc}{subsection}{III.4.7.7 [CT2] Health Information Connectivity Enablement — Capability context [superseded]}

{\small Canonical Strategy Capability: II.1.3.9.8 \hyperref[element:6a1030483dfc4ae6af6998b587129d1c]{[CT2] Health Information Connectivity Enablement}; Canonical ID: \texttt{\seqsplit{6a1030483dfc4ae6af6998b587129d1c}}. Proposed · superseded.}\par
Contextual reference only; no additional architectural instance or semantic relationship.\par
\subsubsection*{III.4.7.7.1 [CT3] Health Information Communication — Capability context}
\phantomsection
\label{navigation:III.4.7.7.1}
\addcontentsline{toc}{subsubsection}{III.4.7.7.1 [CT3] Health Information Communication — Capability context}

{\small Canonical Strategy Capability: II.1.3.9.8.1 \hyperref[element:9794b2aff43443b495a19aefe0efa67b]{[CT3] Health Information Communication}; Canonical ID: \texttt{\seqsplit{9794b2aff43443b495a19aefe0efa67b}}. Canonical · pending\_adjudication.}\par
Contextual reference only; no additional architectural instance or semantic relationship.\par
\paragraph*{III.4.7.7.1.1 Business Services}
\phantomsection
\label{navigation:III.4.7.7.1.1}
\addcontentsline{toc}{paragraph}{III.4.7.7.1.1 Business Services}

\subparagraph*{III.4.7.7.1.1.1 Standards Communication Gateway Service}
\phantomsection
\label{navigation:III.4.7.7.1.1.1}
\addcontentsline{toc}{subparagraph}{III.4.7.7.1.1.1 Standards Communication Gateway Service}

\label{element:80e239faf1bd476e973eb602b3730380}
\index{Standards Communication Gateway Service (III.4.7.7.1.1.1)}
{\small BusinessService · Canonical · pending\_adjudication}\par

Protocol interface for external client integrations.


{\small Recorded unresolved matters: source-qualified matters; exact scope remains in the canonical record.}\par
{\footnotesize \href{../../../../continuum/business/business-service/80e239faf1bd476e973eb602b3730380.md}{Canonical record: 80e239faf1bd476e973eb602b3730380}; full qualifications and provenance.}\par
\paragraph*{III.4.7.7.1.2 Business Functions}
\phantomsection
\label{navigation:III.4.7.7.1.2}
\addcontentsline{toc}{paragraph}{III.4.7.7.1.2 Business Functions}

\subparagraph*{III.4.7.7.1.2.1 Mediate Standards-Based Health Information Communication}
\phantomsection
\label{navigation:III.4.7.7.1.2.1}
\addcontentsline{toc}{subparagraph}{III.4.7.7.1.2.1 Mediate Standards-Based Health Information Communication}

\label{element:6c14589215f744a2b76605a2023b04fb}
\index{Mediate Standards-Based Health Information Communication (III.4.7.7.1.2.1)}
{\small BusinessFunction · Canonical · pending\_adjudication}\par

Handles standards-based communication mediation, transaction coordination, and technical error handling.


{\small Recorded unresolved matters: source-qualified matters; exact scope remains in the canonical record.}\par
{\footnotesize \href{../../../../continuum/business/business-function/6c14589215f744a2b76605a2023b04fb.md}{Canonical record: 6c14589215f744a2b76605a2023b04fb}; full qualifications and provenance.}\par
\subsection*{III.4.7.8 [CT2] Health Information Control — Capability context}
\phantomsection
\label{navigation:III.4.7.8}
\addcontentsline{toc}{subsection}{III.4.7.8 [CT2] Health Information Control — Capability context}

{\small Canonical Strategy Capability: II.1.3.9.9 \hyperref[element:44cc0e17ad534675bf26dce6e46a68ff]{[CT2] Health Information Control}; Canonical ID: \texttt{\seqsplit{44cc0e17ad534675bf26dce6e46a68ff}}. Canonical · current.}\par
Contextual reference only; no additional architectural instance or semantic relationship.\par
\subsubsection*{III.4.7.8.1 Business Services}
\phantomsection
\label{navigation:III.4.7.8.1}
\addcontentsline{toc}{subsubsection}{III.4.7.8.1 Business Services}

\paragraph*{III.4.7.8.1.1 Consent Enforcement Service}
\phantomsection
\label{navigation:III.4.7.8.1.1}
\addcontentsline{toc}{paragraph}{III.4.7.8.1.1 Consent Enforcement Service}

\label{element:67bf392d2dd6464292619e259cc37f8e}
\index{Consent Enforcement Service (III.4.7.8.1.1)}
{\small BusinessService · Canonical · current}\par

Returns masking and redaction instructions.



Feature association not established: FEAT-ISE-12 — Security Context Propagation propagates attributable security context; evaluating consent constraints does not materially realise that responsibility. No replacement Feature is assigned.


{\small Recorded unresolved matters: source-qualified matters; exact scope remains in the canonical record.}\par
{\footnotesize \href{../../../../continuum/business/business-service/67bf392d2dd6464292619e259cc37f8e.md}{Canonical record: 67bf392d2dd6464292619e259cc37f8e}; full qualifications and provenance.}\par
\paragraph*{III.4.7.8.1.2 Policy Evaluation \& Authorisation Service}
\phantomsection
\label{navigation:III.4.7.8.1.2}
\addcontentsline{toc}{paragraph}{III.4.7.8.1.2 Policy Evaluation \& Authorisation Service}

\label{element:a025ac002b6a461abaac07fcb87f411f}
\index{Policy Evaluation \& Authorisation Service (III.4.7.8.1.2)}
{\small BusinessService · Canonical · current}\par

Evaluates access requests across the platform.


{\footnotesize \href{../../../../continuum/business/business-service/a025ac002b6a461abaac07fcb87f411f.md}{Canonical record: a025ac002b6a461abaac07fcb87f411f}; full qualifications and provenance.}\par
\paragraph*{III.4.7.8.1.3 Security Audit Ingress}
\phantomsection
\label{navigation:III.4.7.8.1.3}
\addcontentsline{toc}{paragraph}{III.4.7.8.1.3 Security Audit Ingress}

\label{element:43c4a2aa7b4548ec8fdc2844e3d666b6}
\index{Security Audit Ingress (III.4.7.8.1.3)}
{\small BusinessService · Canonical · current}\par

Ingests audit events across all capabilities.


{\footnotesize \href{../../../../continuum/business/business-service/43c4a2aa7b4548ec8fdc2844e3d666b6.md}{Canonical record: 43c4a2aa7b4548ec8fdc2844e3d666b6}; full qualifications and provenance.}\par
\paragraph*{III.4.7.8.1.4 Security Context Validation Service}
\phantomsection
\label{navigation:III.4.7.8.1.4}
\addcontentsline{toc}{paragraph}{III.4.7.8.1.4 Security Context Validation Service}

\label{element:f1ff9e6542004a7584a385062cb0ed21}
\index{Security Context Validation Service (III.4.7.8.1.4)}
{\small BusinessService · Canonical · current}\par

Verifies transaction security credentials.



Exposure qualification: Validation is the declared Service behaviour. Its juxtaposition with Propagate Security Context does not establish a one-to-one Function/Service relationship; the precise internal mapping remains unestablished.



This Service lets a client obtain validation of transaction security credentials. The exposure is established even though its precise internal Function mapping is not. Validation must not be rewritten as context propagation merely to manufacture a symmetrical Function-to-Service pair.


{\small Recorded unresolved matters: source-qualified matters; exact scope remains in the canonical record.}\par
{\footnotesize \href{../../../../continuum/business/business-service/f1ff9e6542004a7584a385062cb0ed21.md}{Canonical record: f1ff9e6542004a7584a385062cb0ed21}; full qualifications and provenance.}\par
\subsubsection*{III.4.7.8.2 Business Functions}
\phantomsection
\label{navigation:III.4.7.8.2}
\addcontentsline{toc}{subsubsection}{III.4.7.8.2 Business Functions}

\paragraph*{III.4.7.8.2.1 Evaluate Consent Constraint}
\phantomsection
\label{navigation:III.4.7.8.2.1}
\addcontentsline{toc}{paragraph}{III.4.7.8.2.1 Evaluate Consent Constraint}

\label{element:4a881488296d43d09112cb046fa836d3}
\index{Evaluate Consent Constraint (III.4.7.8.2.1)}
{\small BusinessFunction · Canonical · current}\par

Cross-references access requests against client consent directives and sensitive health category restrictions.



Feature association not established: FEAT-ISE-12 — Security Context Propagation propagates attributable security context; evaluating consent constraints does not materially realise that responsibility. No replacement Feature is assigned.


{\small Recorded unresolved matters: source-qualified matters; exact scope remains in the canonical record.}\par
{\footnotesize \href{../../../../continuum/business/business-function/4a881488296d43d09112cb046fa836d3.md}{Canonical record: 4a881488296d43d09112cb046fa836d3}; full qualifications and provenance.}\par
\paragraph*{III.4.7.8.2.2 Evaluate Information Access Authority}
\phantomsection
\label{navigation:III.4.7.8.2.2}
\addcontentsline{toc}{paragraph}{III.4.7.8.2.2 Evaluate Information Access Authority}

\label{element:4707b7fc77254ea68c30fc89f5119a77}
\index{Evaluate Information Access Authority (III.4.7.8.2.2)}
{\small BusinessFunction · Canonical · current}\par

Evaluates user role, organisation, patient relationship, purpose of use, and break-glass overrides against default-deny policies.


{\footnotesize \href{../../../../continuum/business/business-function/4707b7fc77254ea68c30fc89f5119a77.md}{Canonical record: 4707b7fc77254ea68c30fc89f5119a77}; full qualifications and provenance.}\par
\paragraph*{III.4.7.8.2.3 Propagate Security Context}
\phantomsection
\label{navigation:III.4.7.8.2.3}
\addcontentsline{toc}{paragraph}{III.4.7.8.2.3 Propagate Security Context}

\label{element:ee79679c6b5041659a2efca15190454f}
\index{Propagate Security Context (III.4.7.8.2.3)}
{\small BusinessFunction · Canonical · current}\par

Generates and binds canonical, immutable security contexts to all intra-platform transactions.



Exposure qualification: Validation is the declared Service behaviour. Its juxtaposition with Propagate Security Context does not establish a one-to-one Function/Service relationship; the precise internal mapping remains unestablished.



Propagation preserves attributable operational security context across managed transactions. It provides continuity of context rather than independently deciding permission for every downstream action. The nearby validation Service has its own declared exposure; juxtaposition does not establish an exact one-to-one implementation relationship.


{\small Recorded unresolved matters: source-qualified matters; exact scope remains in the canonical record.}\par
{\footnotesize \href{../../../../continuum/business/business-function/ee79679c6b5041659a2efca15190454f.md}{Canonical record: ee79679c6b5041659a2efca15190454f}; full qualifications and provenance.}\par
\paragraph*{III.4.7.8.2.4 Record Security-Significant Activity}
\phantomsection
\label{navigation:III.4.7.8.2.4}
\addcontentsline{toc}{paragraph}{III.4.7.8.2.4 Record Security-Significant Activity}

\label{element:6177ed0d53e5407ab1bc4b721619ed44}
\index{Record Security-Significant Activity (III.4.7.8.2.4)}
{\small BusinessFunction · Canonical · current}\par

Records tamper-evident audit evidence (who, what, when, why, patient MRN, policy outcome) without logging unmasked PHI.


{\footnotesize \href{../../../../continuum/business/business-function/6177ed0d53e5407ab1bc4b721619ed44.md}{Canonical record: 6177ed0d53e5407ab1bc4b721619ed44}; full qualifications and provenance.}\par
\subsection*{III.4.7.9 [CT2] Health Information Exchange Enablement — Capability context}
\phantomsection
\label{navigation:III.4.7.9}
\addcontentsline{toc}{subsection}{III.4.7.9 [CT2] Health Information Exchange Enablement — Capability context}

{\small Canonical Strategy Capability: II.1.3.9.10 \hyperref[element:a61a15908ed544e6acf3e9cab63a0267]{[CT2] Health Information Exchange Enablement}; Canonical ID: \texttt{\seqsplit{a61a15908ed544e6acf3e9cab63a0267}}. Canonical · current.}\par
Contextual reference only; no additional architectural instance or semantic relationship.\par
\subsubsection*{III.4.7.9.1 Business Services}
\phantomsection
\label{navigation:III.4.7.9.1}
\addcontentsline{toc}{subsubsection}{III.4.7.9.1 Business Services}

\paragraph*{III.4.7.9.1.1 Addressed Document Delivery Service}
\phantomsection
\label{navigation:III.4.7.9.1.1}
\addcontentsline{toc}{paragraph}{III.4.7.9.1.1 Addressed Document Delivery Service}

\label{element:97080460bacf40ef8b6db2091cc3bb6e}
\index{Addressed Document Delivery Service (III.4.7.9.1.1)}
{\small BusinessService · Canonical · current}\par

Securely delivers documents to designated recipient mailboxes.


{\footnotesize \href{../../../../continuum/business/business-service/97080460bacf40ef8b6db2091cc3bb6e.md}{Canonical record: 97080460bacf40ef8b6db2091cc3bb6e}; full qualifications and provenance.}\par
\paragraph*{III.4.7.9.1.2 Clinical Event Syndication Stream}
\phantomsection
\label{navigation:III.4.7.9.1.2}
\addcontentsline{toc}{paragraph}{III.4.7.9.1.2 Clinical Event Syndication Stream}

\label{element:337f39bb62d9491bb6e41c4af32b220c}
\index{Clinical Event Syndication Stream (III.4.7.9.1.2)}
{\small BusinessService · Canonical · current}\par

Emits publish-subscribe clinical event feeds.


{\footnotesize \href{../../../../continuum/business/business-service/337f39bb62d9491bb6e41c4af32b220c.md}{Canonical record: 337f39bb62d9491bb6e41c4af32b220c}; full qualifications and provenance.}\par
\paragraph*{III.4.7.9.1.3 Federated Health Information Query}
\phantomsection
\label{navigation:III.4.7.9.1.3}
\addcontentsline{toc}{paragraph}{III.4.7.9.1.3 Federated Health Information Query}

\label{element:b4db5850f7114c968901966b4f6e36f5}
\index{Federated Health Information Query (III.4.7.9.1.3)}
{\small BusinessService · Canonical · current}\par

Resolves distributed record queries.


{\footnotesize \href{../../../../continuum/business/business-service/b4db5850f7114c968901966b4f6e36f5.md}{Canonical record: b4db5850f7114c968901966b4f6e36f5}; full qualifications and provenance.}\par
\paragraph*{III.4.7.9.1.4 HIE Submission Gateway Service}
\phantomsection
\label{navigation:III.4.7.9.1.4}
\addcontentsline{toc}{paragraph}{III.4.7.9.1.4 HIE Submission Gateway Service}

\label{element:816c1013b609440fb772060396fd2dac}
\index{HIE Submission Gateway Service (III.4.7.9.1.4)}
{\small BusinessService · Canonical · current}\par

Ingress endpoint for regional data contributions.


{\footnotesize \href{../../../../continuum/business/business-service/816c1013b609440fb772060396fd2dac.md}{Canonical record: 816c1013b609440fb772060396fd2dac}; full qualifications and provenance.}\par
\subsubsection*{III.4.7.9.2 Business Functions}
\phantomsection
\label{navigation:III.4.7.9.2}
\addcontentsline{toc}{subsubsection}{III.4.7.9.2 Business Functions}

\paragraph*{III.4.7.9.2.1 Distribute Addressed Information}
\phantomsection
\label{navigation:III.4.7.9.2.1}
\addcontentsline{toc}{paragraph}{III.4.7.9.2.1 Distribute Addressed Information}

\label{element:6c9090a3d0a74a968475e05eb617e84f}
\index{Distribute Addressed Information (III.4.7.9.2.1)}
{\small BusinessFunction · Canonical · current}\par

Directs point-to-point clinical documents to specific recipient practitioner endpoints.


{\footnotesize \href{../../../../continuum/business/business-function/6c9090a3d0a74a968475e05eb617e84f.md}{Canonical record: 6c9090a3d0a74a968475e05eb617e84f}; full qualifications and provenance.}\par
\paragraph*{III.4.7.9.2.2 Mediate Information Retrieval}
\phantomsection
\label{navigation:III.4.7.9.2.2}
\addcontentsline{toc}{paragraph}{III.4.7.9.2.2 Mediate Information Retrieval}

\label{element:7667abb329724b7d894fca020acb8014}
\index{Mediate Information Retrieval (III.4.7.9.2.2)}
{\small BusinessFunction · Canonical · current}\par

Federates queries across participating hospital and diagnostic node repositories.


{\footnotesize \href{../../../../continuum/business/business-function/7667abb329724b7d894fca020acb8014.md}{Canonical record: 7667abb329724b7d894fca020acb8014}; full qualifications and provenance.}\par
\paragraph*{III.4.7.9.2.3 Receive Information Submission}
\phantomsection
\label{navigation:III.4.7.9.2.3}
\addcontentsline{toc}{paragraph}{III.4.7.9.2.3 Receive Information Submission}

\label{element:26abb23ed9f54dcf9002eb894f1d02ec}
\index{Receive Information Submission (III.4.7.9.2.3)}
{\small BusinessFunction · Canonical · current}\par

Validates and ingests standards-compliant clinical messages from external health providers.


{\footnotesize \href{../../../../continuum/business/business-function/26abb23ed9f54dcf9002eb894f1d02ec.md}{Canonical record: 26abb23ed9f54dcf9002eb894f1d02ec}; full qualifications and provenance.}\par
\paragraph*{III.4.7.9.2.4 Syndicate Information Change}
\phantomsection
\label{navigation:III.4.7.9.2.4}
\addcontentsline{toc}{paragraph}{III.4.7.9.2.4 Syndicate Information Change}

\label{element:86241540124a40458f7d1dc2f5df84ca}
\index{Syndicate Information Change (III.4.7.9.2.4)}
{\small BusinessFunction · Canonical · current}\par

Broadcasts event notifications and record updates to subscribed regional care network systems.


{\footnotesize \href{../../../../continuum/business/business-function/86241540124a40458f7d1dc2f5df84ca.md}{Canonical record: 86241540124a40458f7d1dc2f5df84ca}; full qualifications and provenance.}\par
\subsection*{III.4.7.10 [CT2] Information Content Maintainer Management — Capability context}
\phantomsection
\label{navigation:III.4.7.10}
\addcontentsline{toc}{subsection}{III.4.7.10 [CT2] Information Content Maintainer Management — Capability context}

{\small Canonical Strategy Capability: II.1.3.9.12 \hyperref[element:b81a1ab42af24c2e9896b8f3f78d291a]{[CT2] Information Content Maintainer Management}; Canonical ID: \texttt{\seqsplit{b81a1ab42af24c2e9896b8f3f78d291a}}. Proposed · current.}\par
Contextual reference only; no additional architectural instance or semantic relationship.\par
\subsubsection*{III.4.7.10.1 [CT3] Data Provenance Capture — Capability context}
\phantomsection
\label{navigation:III.4.7.10.1}
\addcontentsline{toc}{subsubsection}{III.4.7.10.1 [CT3] Data Provenance Capture — Capability context}

{\small Canonical Strategy Capability: II.1.3.9.12.1 \hyperref[element:44a7937520a54945b06b5fd4a2bbd870]{[CT3] Data Provenance Capture}; Canonical ID: \texttt{\seqsplit{44a7937520a54945b06b5fd4a2bbd870}}. Proposed · current.}\par
Contextual reference only; no additional architectural instance or semantic relationship.\par
\paragraph*{III.4.7.10.1.1 Business Functions}
\phantomsection
\label{navigation:III.4.7.10.1.1}
\addcontentsline{toc}{paragraph}{III.4.7.10.1.1 Business Functions}

\subparagraph*{III.4.7.10.1.1.1 Record Significant Information Lineage}
\phantomsection
\label{navigation:III.4.7.10.1.1.1}
\addcontentsline{toc}{subparagraph}{III.4.7.10.1.1.1 Record Significant Information Lineage}

\label{element:afdb0358393f4560b926330a14c1aedd}
\index{Record Significant Information Lineage (III.4.7.10.1.1.1)}
{\small BusinessFunction · Proposed · current}\par

Associate a managed assertion with its attributed sources and information-significant transformation.



\textbf{Overview}



Changes an unexplained derived representation into an attributable relationship to the source information and transformation basis.



\textbf{Boundary}



Routine retries, cache mechanics and debug traces are excluded unless information-significant.


{\footnotesize \href{../../../../continuum/business/business-function/afdb0358393f4560b926330a14c1aedd.md}{Canonical record: afdb0358393f4560b926330a14c1aedd}; full qualifications and provenance.}\par
\subsection*{III.4.7.11 [CT2] Information Design Governance — Capability context}
\phantomsection
\label{navigation:III.4.7.11}
\addcontentsline{toc}{subsection}{III.4.7.11 [CT2] Information Design Governance — Capability context}

{\small Canonical Strategy Capability: II.1.3.9.13 \hyperref[element:f39ed9ebcc724857a4bf7a2b8e045c7a]{[CT2] Information Design Governance}; Canonical ID: \texttt{\seqsplit{f39ed9ebcc724857a4bf7a2b8e045c7a}}. Canonical · current.}\par
Contextual reference only; no additional architectural instance or semantic relationship.\par
\subsubsection*{III.4.7.11.1 Business Services}
\phantomsection
\label{navigation:III.4.7.11.1}
\addcontentsline{toc}{subsubsection}{III.4.7.11.1 Business Services}

\paragraph*{III.4.7.11.1.1 Exchange Profile Specification Query}
\phantomsection
\label{navigation:III.4.7.11.1.1}
\addcontentsline{toc}{paragraph}{III.4.7.11.1.1 Exchange Profile Specification Query}

\label{element:98bc4e8f57f5493ebeb88f8bfc8104a8}
\index{Exchange Profile Specification Query (III.4.7.11.1.1)}
{\small BusinessService · Canonical · current}\par

Discloses interoperability schemas.


{\small Recorded unresolved matters: source-qualified matters; exact scope remains in the canonical record.}\par
{\footnotesize \href{../../../../continuum/business/business-service/98bc4e8f57f5493ebeb88f8bfc8104a8.md}{Canonical record: 98bc4e8f57f5493ebeb88f8bfc8104a8}; full qualifications and provenance.}\par
\paragraph*{III.4.7.11.1.2 Information Dictionary Query}
\phantomsection
\label{navigation:III.4.7.11.1.2}
\addcontentsline{toc}{paragraph}{III.4.7.11.1.2 Information Dictionary Query}

\label{element:ea491740e6cd40c69f65ad7d9d4f81fe}
\index{Information Dictionary Query (III.4.7.11.1.2)}
{\small BusinessService · Canonical · current}\par

Provides data dictionary specifications.


{\small Recorded unresolved matters: source-qualified matters; exact scope remains in the canonical record.}\par
{\footnotesize \href{../../../../continuum/business/business-service/ea491740e6cd40c69f65ad7d9d4f81fe.md}{Canonical record: ea491740e6cd40c69f65ad7d9d4f81fe}; full qualifications and provenance.}\par
\paragraph*{III.4.7.11.1.3 Semantic Conformance Verification Service}
\phantomsection
\label{navigation:III.4.7.11.1.3}
\addcontentsline{toc}{paragraph}{III.4.7.11.1.3 Semantic Conformance Verification Service}

\label{element:3d492f97eb3f4d50a1c448477c61bf2b}
\index{Semantic Conformance Verification Service (III.4.7.11.1.3)}
{\small BusinessService · Canonical · current}\par

Validates representations for compliance.


{\footnotesize \href{../../../../continuum/business/business-service/3d492f97eb3f4d50a1c448477c61bf2b.md}{Canonical record: 3d492f97eb3f4d50a1c448477c61bf2b}; full qualifications and provenance.}\par
\paragraph*{III.4.7.11.1.4 Semantic Model Definition Service}
\phantomsection
\label{navigation:III.4.7.11.1.4}
\addcontentsline{toc}{paragraph}{III.4.7.11.1.4 Semantic Model Definition Service}

\label{element:0f6b41152f294a87b84970c74ab5a014}
\index{Semantic Model Definition Service (III.4.7.11.1.4)}
{\small BusinessService · Canonical · current}\par

Discloses canonical model definitions.


{\footnotesize \href{../../../../continuum/business/business-service/0f6b41152f294a87b84970c74ab5a014.md}{Canonical record: 0f6b41152f294a87b84970c74ab5a014}; full qualifications and provenance.}\par
\paragraph*{III.4.7.11.1.5 Terminology Binding Specification Lookup}
\phantomsection
\label{navigation:III.4.7.11.1.5}
\addcontentsline{toc}{paragraph}{III.4.7.11.1.5 Terminology Binding Specification Lookup}

\label{element:1b9e7c961a054a6bbdb7876d08a82ef3}
\index{Terminology Binding Specification Lookup (III.4.7.11.1.5)}
{\small BusinessService · Canonical · current}\par

Discloses value set bindings.


{\small Recorded unresolved matters: source-qualified matters; exact scope remains in the canonical record.}\par
{\footnotesize \href{../../../../continuum/business/business-service/1b9e7c961a054a6bbdb7876d08a82ef3.md}{Canonical record: 1b9e7c961a054a6bbdb7876d08a82ef3}; full qualifications and provenance.}\par
\subsubsection*{III.4.7.11.2 Business Functions}
\phantomsection
\label{navigation:III.4.7.11.2}
\addcontentsline{toc}{subsubsection}{III.4.7.11.2 Business Functions}

\paragraph*{III.4.7.11.2.1 Govern Exchange Profile}
\phantomsection
\label{navigation:III.4.7.11.2.1}
\addcontentsline{toc}{paragraph}{III.4.7.11.2.1 Govern Exchange Profile}

\label{element:935f9a62438944599b913f2614c9e595}
\index{Govern Exchange Profile (III.4.7.11.2.1)}
{\small BusinessFunction · Canonical · current}\par

Defines standards-based interchange profiles and message constraints.


{\small Recorded unresolved matters: source-qualified matters; exact scope remains in the canonical record.}\par
{\footnotesize \href{../../../../continuum/business/business-function/935f9a62438944599b913f2614c9e595.md}{Canonical record: 935f9a62438944599b913f2614c9e595}; full qualifications and provenance.}\par
\paragraph*{III.4.7.11.2.2 Govern Information Definition}
\phantomsection
\label{navigation:III.4.7.11.2.2}
\addcontentsline{toc}{paragraph}{III.4.7.11.2.2 Govern Information Definition}

\label{element:820b4d41490a4389bb864f1c4a04e881}
\index{Govern Information Definition (III.4.7.11.2.2)}
{\small BusinessFunction · Canonical · current}\par

Maintains governed data element dictionaries, data types, and structural rules.


{\small Recorded unresolved matters: source-qualified matters; exact scope remains in the canonical record.}\par
{\footnotesize \href{../../../../continuum/business/business-function/820b4d41490a4389bb864f1c4a04e881.md}{Canonical record: 820b4d41490a4389bb864f1c4a04e881}; full qualifications and provenance.}\par
\paragraph*{III.4.7.11.2.3 Govern Information Semantics}
\phantomsection
\label{navigation:III.4.7.11.2.3}
\addcontentsline{toc}{paragraph}{III.4.7.11.2.3 Govern Information Semantics}

\label{element:39a297b6c40c44ecba64857f6c3f3e0f}
\index{Govern Information Semantics (III.4.7.11.2.3)}
{\small BusinessFunction · Canonical · current}\par

Authors and validates canonical data concepts and entity relationships.


{\footnotesize \href{../../../../continuum/business/business-function/39a297b6c40c44ecba64857f6c3f3e0f.md}{Canonical record: 39a297b6c40c44ecba64857f6c3f3e0f}; full qualifications and provenance.}\par
\paragraph*{III.4.7.11.2.4 Govern Terminology Binding}
\phantomsection
\label{navigation:III.4.7.11.2.4}
\addcontentsline{toc}{paragraph}{III.4.7.11.2.4 Govern Terminology Binding}

\label{element:9e96ce12b7494594a6129f79f7491223}
\index{Govern Terminology Binding (III.4.7.11.2.4)}
{\small BusinessFunction · Canonical · current}\par

Binds canonical data elements to authoritative terminology value sets.


{\small Recorded unresolved matters: source-qualified matters; exact scope remains in the canonical record.}\par
{\footnotesize \href{../../../../continuum/business/business-function/9e96ce12b7494594a6129f79f7491223.md}{Canonical record: 9e96ce12b7494594a6129f79f7491223}; full qualifications and provenance.}\par
\paragraph*{III.4.7.11.2.5 Verify Semantic Conformance}
\phantomsection
\label{navigation:III.4.7.11.2.5}
\addcontentsline{toc}{paragraph}{III.4.7.11.2.5 Verify Semantic Conformance}

\label{element:5445184252c94c05874c52350ff13d2b}
\index{Verify Semantic Conformance (III.4.7.11.2.5)}
{\small BusinessFunction · Canonical · current}\par

Validates candidate messages against canonical governance schemas.


{\footnotesize \href{../../../../continuum/business/business-function/5445184252c94c05874c52350ff13d2b.md}{Canonical record: 5445184252c94c05874c52350ff13d2b}; full qualifications and provenance.}\par
\subsection*{III.4.7.12 [CT2] Information Protection \& Accountability Enablement — Capability context [superseded]}
\phantomsection
\label{navigation:III.4.7.12}
\addcontentsline{toc}{subsection}{III.4.7.12 [CT2] Information Protection \& Accountability Enablement — Capability context [superseded]}

{\small Canonical Strategy Capability: II.1.3.9.14 \hyperref[element:e922cba7cf1d4ebeb9abf6472a3803ce]{[CT2] Information Protection \& Accountability Enablement}; Canonical ID: \texttt{\seqsplit{e922cba7cf1d4ebeb9abf6472a3803ce}}. Proposed · superseded.}\par
Contextual reference only; no additional architectural instance or semantic relationship.\par
\subsubsection*{III.4.7.12.1 [CT3] Significant Activity Audit Capture — Capability context}
\phantomsection
\label{navigation:III.4.7.12.1}
\addcontentsline{toc}{subsubsection}{III.4.7.12.1 [CT3] Significant Activity Audit Capture — Capability context}

{\small Canonical Strategy Capability: II.1.3.9.14.2 \hyperref[element:d903948169eb4785b0f03d916c842a8a]{[CT3] Significant Activity Audit Capture}; Canonical ID: \texttt{\seqsplit{d903948169eb4785b0f03d916c842a8a}}. Proposed · pending\_adjudication.}\par
Contextual reference only; no additional architectural instance or semantic relationship.\par
\paragraph*{III.4.7.12.1.1 Business Services}
\phantomsection
\label{navigation:III.4.7.12.1.1}
\addcontentsline{toc}{paragraph}{III.4.7.12.1.1 Business Services}

\subparagraph*{III.4.7.12.1.1.1 Mandatory Audit Evidence Capture}
\phantomsection
\label{navigation:III.4.7.12.1.1.1}
\addcontentsline{toc}{subparagraph}{III.4.7.12.1.1.1 Mandatory Audit Evidence Capture}

\label{element:6ab3f8628ba14f17bddbf6092cf56ba7}
\index{Mandatory Audit Evidence Capture (III.4.7.12.1.1.1)}
{\small BusinessService · Proposed · pending\_adjudication}\par

Allow governed-operation boundaries to require authoritative capture of significant activity evidence.



\textbf{Overview}



The client contract distinguishes confirmed capture from inability to capture, allowing required operations to fail closed.



\textbf{Boundary}



The exposed boundary does not permit callers to assign themselves audit platform authority.


{\small Recorded unresolved matters: source-qualified matters; exact scope remains in the canonical record.}\par
{\footnotesize \href{../../../../continuum/business/business-service/6ab3f8628ba14f17bddbf6092cf56ba7.md}{Canonical record: 6ab3f8628ba14f17bddbf6092cf56ba7}; full qualifications and provenance.}\par
\paragraph*{III.4.7.12.1.2 Business Functions}
\phantomsection
\label{navigation:III.4.7.12.1.2}
\addcontentsline{toc}{paragraph}{III.4.7.12.1.2 Business Functions}

\subparagraph*{III.4.7.12.1.2.1 Capture Required Significant Audit Evidence}
\phantomsection
\label{navigation:III.4.7.12.1.2.1}
\addcontentsline{toc}{subparagraph}{III.4.7.12.1.2.1 Capture Required Significant Audit Evidence}

\label{element:c8e614257cf14bd1878df894455f21fb}
\index{Capture Required Significant Audit Evidence (III.4.7.12.1.2.1)}
{\small BusinessFunction · Proposed · pending\_adjudication}\par

Establish required significant activity evidence under the governed authoritative audit-capture responsibility.



\textbf{Overview}



Moves an auditable fact to confirmed captured standing or an explicit capture failure affecting the protected operation.



\textbf{Boundary}



Specific capture authority is not borrowed from the originating identity and is not generalised.


{\small Recorded unresolved matters: source-qualified matters; exact scope remains in the canonical record.}\par
{\footnotesize \href{../../../../continuum/business/business-function/c8e614257cf14bd1878df894455f21fb.md}{Canonical record: c8e614257cf14bd1878df894455f21fb}; full qualifications and provenance.}\par
\subsection*{III.4.7.13 [CT2] Patient Clinical Record — Capability context}
\phantomsection
\label{navigation:III.4.7.13}
\addcontentsline{toc}{subsection}{III.4.7.13 [CT2] Patient Clinical Record — Capability context}

{\small Canonical Strategy Capability: II.1.3.9.15 \hyperref[element:5478878e0e6e4d34a8d5023ac6ac2d54]{[CT2] Patient Clinical Record}; Canonical ID: \texttt{\seqsplit{5478878e0e6e4d34a8d5023ac6ac2d54}}. Canonical · current.}\par
Contextual reference only; no additional architectural instance or semantic relationship.\par
\subsubsection*{III.4.7.13.1 [CT3] Linked Information Curation — Capability context}
\phantomsection
\label{navigation:III.4.7.13.1}
\addcontentsline{toc}{subsubsection}{III.4.7.13.1 [CT3] Linked Information Curation — Capability context}

{\small Canonical Strategy Capability: II.1.3.9.15.1 \hyperref[element:5edb2a8b5ee54ddda8104d1fab323a94]{[CT3] Linked Information Curation}; Canonical ID: \texttt{\seqsplit{5edb2a8b5ee54ddda8104d1fab323a94}}. Proposed · current.}\par
Contextual reference only; no additional architectural instance or semantic relationship.\par
\paragraph*{III.4.7.13.1.1 Business Services}
\phantomsection
\label{navigation:III.4.7.13.1.1}
\addcontentsline{toc}{paragraph}{III.4.7.13.1.1 Business Services}

\subparagraph*{III.4.7.13.1.1.1 Associated Binary Content Retrieval}
\phantomsection
\label{navigation:III.4.7.13.1.1.1}
\addcontentsline{toc}{subparagraph}{III.4.7.13.1.1.1 Associated Binary Content Retrieval}

\label{element:f0e0c03e407e4534b9e422df103fc35f}
\index{Associated Binary Content Retrieval (III.4.7.13.1.1.1)}
{\small BusinessService · Proposed · current}\par

Allow information-consuming capabilities to request resolution and retrieval of associated content.



\textbf{Overview}



The client contract returns authorised content or explicit incomplete/unavailable standing while retaining custody attribution.



\textbf{Boundary}



The client acquires neither physical repository access nor originating clinical authority.


{\footnotesize \href{../../../../continuum/business/business-service/f0e0c03e407e4534b9e422df103fc35f.md}{Canonical record: f0e0c03e407e4534b9e422df103fc35f}; full qualifications and provenance.}\par
\paragraph*{III.4.7.13.1.2 Business Functions}
\phantomsection
\label{navigation:III.4.7.13.1.2}
\addcontentsline{toc}{paragraph}{III.4.7.13.1.2 Business Functions}

\subparagraph*{III.4.7.13.1.2.1 Assemble Associated Binary Content}
\phantomsection
\label{navigation:III.4.7.13.1.2.1}
\addcontentsline{toc}{subparagraph}{III.4.7.13.1.2.1 Assemble Associated Binary Content}

\label{element:0a0a5f3e40d241e38a5b155ca66eadd5}
\index{Assemble Associated Binary Content (III.4.7.13.1.2.1)}
{\small BusinessFunction · Proposed · current}\par

Retrieve and assemble the requested logical content from its permitted constituent payloads.



\textbf{Overview}



Moves a qualified retrieval request to an assembled result or explicit unavailable/partial standing.



\textbf{Boundary}



No storage topology or universal local custody is selected.


{\footnotesize \href{../../../../continuum/business/business-function/0a0a5f3e40d241e38a5b155ca66eadd5.md}{Canonical record: 0a0a5f3e40d241e38a5b155ca66eadd5}; full qualifications and provenance.}\par
\subparagraph*{III.4.7.13.1.2.2 Resolve Associated Binary Content}
\phantomsection
\label{navigation:III.4.7.13.1.2.2}
\addcontentsline{toc}{subparagraph}{III.4.7.13.1.2.2 Resolve Associated Binary Content}

\label{element:118867dc136f4f689265b173abdf626b}
\index{Resolve Associated Binary Content (III.4.7.13.1.2.2)}
{\small BusinessFunction · Proposed · current}\par

Resolve logical content associations to an authorised, custody-qualified retrieval basis.



\textbf{Overview}



Changes an unresolved content request into an identified retrieval basis or an explicit inability to resolve it.



\textbf{Boundary}



Does not assert payload existence, copy all content or establish clinical meaning.


{\footnotesize \href{../../../../continuum/business/business-function/118867dc136f4f689265b173abdf626b.md}{Canonical record: 118867dc136f4f689265b173abdf626b}; full qualifications and provenance.}\par
\subsubsection*{III.4.7.13.2 [CT3] Primary Care Coordination — Capability context}
\phantomsection
\label{navigation:III.4.7.13.2}
\addcontentsline{toc}{subsubsection}{III.4.7.13.2 [CT3] Primary Care Coordination — Capability context}

{\small Canonical Strategy Capability: II.1.3.9.15.3 \hyperref[element:f80e73ee91624ce5baa11ae15480fd35]{[CT3] Primary Care Coordination}; Canonical ID: \texttt{\seqsplit{f80e73ee91624ce5baa11ae15480fd35}}. Canonical · current.}\par
Contextual reference only; no additional architectural instance or semantic relationship.\par
\paragraph*{III.4.7.13.2.1 Business Services}
\phantomsection
\label{navigation:III.4.7.13.2.1}
\addcontentsline{toc}{paragraph}{III.4.7.13.2.1 Business Services}

\subparagraph*{III.4.7.13.2.1.1 Chronic Care Timeline Query}
\phantomsection
\label{navigation:III.4.7.13.2.1.1}
\addcontentsline{toc}{subparagraph}{III.4.7.13.2.1.1 Chronic Care Timeline Query}

\label{element:88243738a3c941efb369f5dd6da6252b}
\index{Chronic Care Timeline Query (III.4.7.13.2.1.1)}
{\small BusinessService · Canonical · current}\par

Exposes aggregated multidisciplinary care timelines.



Feature association not established: Strategy FEAT-SD-02 — Primary Care Notification Dispatch remains valid, but semantic evidence does not establish that the following Function or Service belongs to, realises, specialises or is equivalent to it. No replacement Feature is assigned.


{\small Recorded unresolved matters: source-qualified matters; exact scope remains in the canonical record.}\par
{\footnotesize \href{../../../../continuum/business/business-service/88243738a3c941efb369f5dd6da6252b.md}{Canonical record: 88243738a3c941efb369f5dd6da6252b}; full qualifications and provenance.}\par
\subparagraph*{III.4.7.13.2.1.2 Primary Care Summary Ingress}
\phantomsection
\label{navigation:III.4.7.13.2.1.2}
\addcontentsline{toc}{subparagraph}{III.4.7.13.2.1.2 Primary Care Summary Ingress}

\label{element:7f4ca66b911e48bc8bf859e239f5ba00}
\index{Primary Care Summary Ingress (III.4.7.13.2.1.2)}
{\small BusinessService · Canonical · current}\par

Ingests GP event summaries into the HIE.


{\footnotesize \href{../../../../continuum/business/business-service/7f4ca66b911e48bc8bf859e239f5ba00.md}{Canonical record: 7f4ca66b911e48bc8bf859e239f5ba00}; full qualifications and provenance.}\par
\paragraph*{III.4.7.13.2.2 Business Functions}
\phantomsection
\label{navigation:III.4.7.13.2.2}
\addcontentsline{toc}{paragraph}{III.4.7.13.2.2 Business Functions}

\subparagraph*{III.4.7.13.2.2.1 Coordinate Multi-Agency Chronic Care Plan}
\phantomsection
\label{navigation:III.4.7.13.2.2.1}
\addcontentsline{toc}{subparagraph}{III.4.7.13.2.2.1 Coordinate Multi-Agency Chronic Care Plan}

\label{element:fc49b716c2a845408198591e9a4dc515}
\index{Coordinate Multi-Agency Chronic Care Plan (III.4.7.13.2.2.1)}
{\small BusinessFunction · Canonical · current}\par

Enables GPs to track multi-provider disease management milestones (e.g., Diabetes or Cardiovascular Care Plans).



Feature association not established: Strategy FEAT-SD-02 — Primary Care Notification Dispatch remains valid, but semantic evidence does not establish that the following Function or Service belongs to, realises, specialises or is equivalent to it. No replacement Feature is assigned.


{\small Recorded unresolved matters: source-qualified matters; exact scope remains in the canonical record.}\par
{\footnotesize \href{../../../../continuum/business/business-function/fc49b716c2a845408198591e9a4dc515.md}{Canonical record: fc49b716c2a845408198591e9a4dc515}; full qualifications and provenance.}\par
\subparagraph*{III.4.7.13.2.2.2 Correlate Primary Care Consultation Context}
\phantomsection
\label{navigation:III.4.7.13.2.2.2}
\addcontentsline{toc}{subparagraph}{III.4.7.13.2.2.2 Correlate Primary Care Consultation Context}

\label{element:e6c9088f502742afb538ee654ac0f962}
\index{Correlate Primary Care Consultation Context (III.4.7.13.2.2.2)}
{\small BusinessFunction · Canonical · current}\par

Binds GP consultation notes, vital signs, and encounter episodes to the regional Longitudinal Health Record (LHR).


{\footnotesize \href{../../../../continuum/business/business-function/e6c9088f502742afb538ee654ac0f962.md}{Canonical record: e6c9088f502742afb538ee654ac0f962}; full qualifications and provenance.}\par
\subsubsection*{III.4.7.13.3 [CT3] Remote Patient Monitoring — Capability context}
\phantomsection
\label{navigation:III.4.7.13.3}
\addcontentsline{toc}{subsubsection}{III.4.7.13.3 [CT3] Remote Patient Monitoring — Capability context}

{\small Canonical Strategy Capability: II.1.3.9.15.4 \hyperref[element:2d1addc7dd1d4821b71a11c692d35b27]{[CT3] Remote Patient Monitoring}; Canonical ID: \texttt{\seqsplit{2d1addc7dd1d4821b71a11c692d35b27}}. Canonical · current.}\par
Contextual reference only; no additional architectural instance or semantic relationship.\par
\paragraph*{III.4.7.13.3.1 Business Services}
\phantomsection
\label{navigation:III.4.7.13.3.1}
\addcontentsline{toc}{paragraph}{III.4.7.13.3.1 Business Services}

\subparagraph*{III.4.7.13.3.1.1 Remote Telemetry Ingress}
\phantomsection
\label{navigation:III.4.7.13.3.1.1}
\addcontentsline{toc}{subparagraph}{III.4.7.13.3.1.1 Remote Telemetry Ingress}

\label{element:89014429b6454ced82c26cbe777fa33a}
\index{Remote Telemetry Ingress (III.4.7.13.3.1.1)}
{\small BusinessService · Canonical · current}\par

Ingests remote patient data streams into the LHR.


{\footnotesize \href{../../../../continuum/business/business-service/89014429b6454ced82c26cbe777fa33a.md}{Canonical record: 89014429b6454ced82c26cbe777fa33a}; full qualifications and provenance.}\par
\subparagraph*{III.4.7.13.3.1.2 Virtual Encounter Context Resolution}
\phantomsection
\label{navigation:III.4.7.13.3.1.2}
\addcontentsline{toc}{subparagraph}{III.4.7.13.3.1.2 Virtual Encounter Context Resolution}

\label{element:db1dfa6db11342128e82df186877c477}
\index{Virtual Encounter Context Resolution (III.4.7.13.3.1.2)}
{\small BusinessService · Canonical · current}\par

Binds virtual session identifiers to patient encounters.


{\footnotesize \href{../../../../continuum/business/business-service/db1dfa6db11342128e82df186877c477.md}{Canonical record: db1dfa6db11342128e82df186877c477}; full qualifications and provenance.}\par
\paragraph*{III.4.7.13.3.2 Business Functions}
\phantomsection
\label{navigation:III.4.7.13.3.2}
\addcontentsline{toc}{paragraph}{III.4.7.13.3.2 Business Functions}

\subparagraph*{III.4.7.13.3.2.1 Correlate Virtual Care Consultation Context}
\phantomsection
\label{navigation:III.4.7.13.3.2.1}
\addcontentsline{toc}{subparagraph}{III.4.7.13.3.2.1 Correlate Virtual Care Consultation Context}

\label{element:d315a7eab4d9453f83ecd44ed8b57b56}
\index{Correlate Virtual Care Consultation Context (III.4.7.13.3.2.1)}
{\small BusinessFunction · Canonical · current}\par

Binds video consultation metadata, remote practitioner identity, and shared digital whiteboard artifacts to the encounter.


{\footnotesize \href{../../../../continuum/business/business-function/d315a7eab4d9453f83ecd44ed8b57b56.md}{Canonical record: d315a7eab4d9453f83ecd44ed8b57b56}; full qualifications and provenance.}\par
\subparagraph*{III.4.7.13.3.2.2 Ingest Remote Patient Telemetry}
\phantomsection
\label{navigation:III.4.7.13.3.2.2}
\addcontentsline{toc}{subparagraph}{III.4.7.13.3.2.2 Ingest Remote Patient Telemetry}

\label{element:cb03c3040e1c4a4a9fdadb2502113014}
\index{Ingest Remote Patient Telemetry (III.4.7.13.3.2.2)}
{\small BusinessFunction · Canonical · current}\par

Ingests continuous or periodic physiological measurements (e.g., blood pressure, pulse oximetry, continuous glucose) from home monitoring devices.


{\footnotesize \href{../../../../continuum/business/business-function/cb03c3040e1c4a4a9fdadb2502113014.md}{Canonical record: cb03c3040e1c4a4a9fdadb2502113014}; full qualifications and provenance.}\par
\subsubsection*{III.4.7.13.4 Business Services}
\phantomsection
\label{navigation:III.4.7.13.4}
\addcontentsline{toc}{subsubsection}{III.4.7.13.4 Business Services}

\paragraph*{III.4.7.13.4.1 Active Problem \& Allergy Summary Query}
\phantomsection
\label{navigation:III.4.7.13.4.1}
\addcontentsline{toc}{paragraph}{III.4.7.13.4.1 Active Problem \& Allergy Summary Query}

\label{element:a382bc84c1d043849b8cd8ebe8a482c1}
\index{Active Problem \& Allergy Summary Query (III.4.7.13.4.1)}
{\small BusinessService · Canonical · current}\par

Exposes immediate active clinical summaries.


{\footnotesize \href{../../../../continuum/business/business-service/a382bc84c1d043849b8cd8ebe8a482c1.md}{Canonical record: a382bc84c1d043849b8cd8ebe8a482c1}; full qualifications and provenance.}\par
\paragraph*{III.4.7.13.4.2 Historical Clinical Record Access}
\phantomsection
\label{navigation:III.4.7.13.4.2}
\addcontentsline{toc}{paragraph}{III.4.7.13.4.2 Historical Clinical Record Access}

\label{element:be4938981b43403282694631b40f0bb2}
\index{Historical Clinical Record Access (III.4.7.13.4.2)}
{\small BusinessService · Canonical · current}\par

Provides governed historical record retrieval for authorised clinical and legal review; preservation does not itself establish legal qualification or originating authority.


{\footnotesize \href{../../../../continuum/business/business-service/be4938981b43403282694631b40f0bb2.md}{Canonical record: be4938981b43403282694631b40f0bb2}; full qualifications and provenance.}\par
\paragraph*{III.4.7.13.4.3 Longitudinal Clinical Record Query}
\phantomsection
\label{navigation:III.4.7.13.4.3}
\addcontentsline{toc}{paragraph}{III.4.7.13.4.3 Longitudinal Clinical Record Query}

\label{element:df10431b93bf47dfae2b8ec7b3e6552b}
\index{Longitudinal Clinical Record Query (III.4.7.13.4.3)}
{\small BusinessService · Canonical · current}\par

Discloses the comprehensive clinical timeline to authorised clinicians.


{\footnotesize \href{../../../../continuum/business/business-service/df10431b93bf47dfae2b8ec7b3e6552b.md}{Canonical record: df10431b93bf47dfae2b8ec7b3e6552b}; full qualifications and provenance.}\par
\subsubsection*{III.4.7.13.5 Business Functions}
\phantomsection
\label{navigation:III.4.7.13.5}
\addcontentsline{toc}{subsubsection}{III.4.7.13.5 Business Functions}

\paragraph*{III.4.7.13.5.1 Assemble Longitudinal Clinical Record}
\phantomsection
\label{navigation:III.4.7.13.5.1}
\addcontentsline{toc}{paragraph}{III.4.7.13.5.1 Assemble Longitudinal Clinical Record}

\label{element:65ea2d2427dd474d84f153671de79d63}
\index{Assemble Longitudinal Clinical Record (III.4.7.13.5.1)}
{\small BusinessFunction · Canonical · current}\par

Synthesises clinical documents, encounters, diagnostic reports, and medication histories across disparate source systems into a unified chronological clinical view.


{\footnotesize \href{../../../../continuum/business/business-function/65ea2d2427dd474d84f153671de79d63.md}{Canonical record: 65ea2d2427dd474d84f153671de79d63}; full qualifications and provenance.}\par
\paragraph*{III.4.7.13.5.2 Maintain Active Clinical Record}
\phantomsection
\label{navigation:III.4.7.13.5.2}
\addcontentsline{toc}{paragraph}{III.4.7.13.5.2 Maintain Active Clinical Record}

\label{element:dbccf500cf4a44e0987f972311bab03a}
\index{Maintain Active Clinical Record (III.4.7.13.5.2)}
{\small BusinessFunction · Canonical · current}\par

Maintains active problem lists, current medication regimens, allergies, and open clinical alerts for active care coordination.


{\footnotesize \href{../../../../continuum/business/business-function/dbccf500cf4a44e0987f972311bab03a.md}{Canonical record: dbccf500cf4a44e0987f972311bab03a}; full qualifications and provenance.}\par
\paragraph*{III.4.7.13.5.3 Preserve Durable Clinical Record}
\phantomsection
\label{navigation:III.4.7.13.5.3}
\addcontentsline{toc}{paragraph}{III.4.7.13.5.3 Preserve Durable Clinical Record}

\label{element:ace0f4afad104482b7970534a45d74d7}
\index{Preserve Durable Clinical Record (III.4.7.13.5.3)}
{\small BusinessFunction · Canonical · current}\par

Preserves the governed longitudinal representation and durable, immutable, append-only history of what Harmonia received, knew, asserted or managed, with provenance and versioning. Corrections/supersession establish changed knowledge without retrospectively altering that history or acquiring originating clinical authority.


{\footnotesize \href{../../../../continuum/business/business-function/ace0f4afad104482b7970534a45d74d7.md}{Canonical record: ace0f4afad104482b7970534a45d74d7}; full qualifications and provenance.}\par
\subsection*{III.4.7.14 [CT2] Presentation Services — Capability context}
\phantomsection
\label{navigation:III.4.7.14}
\addcontentsline{toc}{subsection}{III.4.7.14 [CT2] Presentation Services — Capability context}

{\small Canonical Strategy Capability: II.1.3.9.16 \hyperref[element:e9b6ba1917ba44cca7466444d0836c98]{[CT2] Presentation Services}; Canonical ID: \texttt{\seqsplit{e9b6ba1917ba44cca7466444d0836c98}}. Canonical · current.}\par
Contextual reference only; no additional architectural instance or semantic relationship.\par
\subsubsection*{III.4.7.14.1 Business Services}
\phantomsection
\label{navigation:III.4.7.14.1}
\addcontentsline{toc}{subsubsection}{III.4.7.14.1 Business Services}

\paragraph*{III.4.7.14.1.1 Contextual Presentation Service}
\phantomsection
\label{navigation:III.4.7.14.1.1}
\addcontentsline{toc}{paragraph}{III.4.7.14.1.1 Contextual Presentation Service}

\label{element:905d44b44f434fa4974e09dcafac016a}
\index{Contextual Presentation Service (III.4.7.14.1.1)}
{\small BusinessService · Canonical · current}\par

Delivers contextual health and operational information for presentation.


{\footnotesize \href{../../../../continuum/business/business-service/905d44b44f434fa4974e09dcafac016a.md}{Canonical record: 905d44b44f434fa4974e09dcafac016a}; full qualifications and provenance.}\par
\subsubsection*{III.4.7.14.2 Business Functions}
\phantomsection
\label{navigation:III.4.7.14.2}
\addcontentsline{toc}{subsubsection}{III.4.7.14.2 Business Functions}

\paragraph*{III.4.7.14.2.1 Present Contextual Health and Operational Information}
\phantomsection
\label{navigation:III.4.7.14.2.1}
\addcontentsline{toc}{paragraph}{III.4.7.14.2.1 Present Contextual Health and Operational Information}

\label{element:28ad1d4dc0444c7e8f98906a70e9319e}
\index{Present Contextual Health and Operational Information (III.4.7.14.2.1)}
{\small BusinessFunction · Canonical · current}\par

Composes clinical summaries, work queues, and patient records into responsive role-based user interfaces.


{\footnotesize \href{../../../../continuum/business/business-function/28ad1d4dc0444c7e8f98906a70e9319e.md}{Canonical record: 28ad1d4dc0444c7e8f98906a70e9319e}; full qualifications and provenance.}\par
\chapter*{III.5 Unclassified Business}
\phantomsection
\label{navigation:III.5}
\addcontentsline{toc}{chapter}{III.5 Unclassified Business}

\section*{III.5.1 Authorisation — Consent}
\phantomsection
\label{navigation:III.5.1}
\addcontentsline{toc}{section}{III.5.1 Authorisation — Consent}

\label{element:8c8a371f6a324d908656c3c935c7656f}
\index{Authorisation — Consent (III.5.1)}
{\small BusinessInteraction · Canonical · current}\par

\begin{itemize}


\item \textbf{Authorisation — Consent}: The explicit granting or revocation of authority by a patient or substitute decision-maker to collect, use, or disclose health data.

\end{itemize}

{\small Recorded unresolved matters: canonicalName; exact scope remains in the canonical record.}\par
{\small Navigation classification: Native BusinessInteraction definition is determinate, but no governed capability context is represented for this interaction. No Application Interaction, integration flow or capability allocation is inferred.}\par
{\footnotesize \href{../../../../continuum/business/business-interaction/8c8a371f6a324d908656c3c935c7656f.md}{Canonical record: 8c8a371f6a324d908656c3c935c7656f}; full qualifications and provenance.}\par
\section*{III.5.2 Authorisation — Governance}
\phantomsection
\label{navigation:III.5.2}
\addcontentsline{toc}{section}{III.5.2 Authorisation — Governance}

\label{element:a1a6a8eeb4d54fa9ac184198a506d318}
\index{Authorisation — Governance (III.5.2)}
{\small BusinessInteraction · Canonical · current}\par

\begin{itemize}


\item \textbf{Authorisation — Governance}: The formal granting of administrative, system, or operational access privileges based on organizational authority.

\end{itemize}

{\small Recorded unresolved matters: canonicalName; exact scope remains in the canonical record.}\par
{\small Navigation classification: Native BusinessInteraction definition is determinate, but no governed capability context is represented for this interaction. No Application Interaction, integration flow or capability allocation is inferred.}\par
{\footnotesize \href{../../../../continuum/business/business-interaction/a1a6a8eeb4d54fa9ac184198a506d318.md}{Canonical record: a1a6a8eeb4d54fa9ac184198a506d318}; full qualifications and provenance.}\par
\section*{III.5.3 Clinical Alert}
\phantomsection
\label{navigation:III.5.3}
\addcontentsline{toc}{section}{III.5.3 Clinical Alert}

\label{element:8666d4168fe545fd9ce222217887bd61}
\index{Clinical Alert (III.5.3)}
{\small BusinessInteraction · Canonical · current}\par

\begin{itemize}


\item \textbf{Clinical Alert}: An urgent, interruptive warning requiring prompt clinical attention (e.g., critical lab value, severe drug interaction).

\end{itemize}

{\small Recorded unresolved matters: canonicalName; exact scope remains in the canonical record.}\par
{\small Navigation classification: Native BusinessInteraction definition is determinate, but no governed capability context is represented for this interaction. No Application Interaction, integration flow or capability allocation is inferred.}\par
{\footnotesize \href{../../../../continuum/business/business-interaction/8666d4168fe545fd9ce222217887bd61.md}{Canonical record: 8666d4168fe545fd9ce222217887bd61}; full qualifications and provenance.}\par
\section*{III.5.4 Clinical Collaboration}
\phantomsection
\label{navigation:III.5.4}
\addcontentsline{toc}{section}{III.5.4 Clinical Collaboration}

\label{element:f54492b30639431e80cd0723217d6be8}
\index{Clinical Collaboration (III.5.4)}
{\small BusinessInteraction · Canonical · current}\par

\begin{itemize}


\item \textbf{Clinical Collaboration}: Structured, multi-party dialogue, handover messaging, or clinical commentary among care team members.

\end{itemize}

{\small Recorded unresolved matters: canonicalName; exact scope remains in the canonical record.}\par
{\small Navigation classification: Native BusinessInteraction definition is determinate, but no governed capability context is represented for this interaction. No Application Interaction, integration flow or capability allocation is inferred.}\par
{\footnotesize \href{../../../../continuum/business/business-interaction/f54492b30639431e80cd0723217d6be8.md}{Canonical record: f54492b30639431e80cd0723217d6be8}; full qualifications and provenance.}\par
\section*{III.5.5 Clinical Event Notification}
\phantomsection
\label{navigation:III.5.5}
\addcontentsline{toc}{section}{III.5.5 Clinical Event Notification}

\label{element:d2153d1057d740fea1665d71ff66dd1c}
\index{Clinical Event Notification (III.5.5)}
{\small BusinessInteraction · Canonical · current}\par

\begin{itemize}


\item \textbf{Clinical Event Notification}: An automated notice indicating that a clinically significant event has occurred (e.g., patient admission, vital sign threshold breached).

\end{itemize}

{\small Recorded unresolved matters: canonicalName; exact scope remains in the canonical record.}\par
{\small Navigation classification: Native BusinessInteraction definition is determinate, but no governed capability context is represented for this interaction. No Application Interaction, integration flow or capability allocation is inferred.}\par
{\footnotesize \href{../../../../continuum/business/business-interaction/d2153d1057d740fea1665d71ff66dd1c.md}{Canonical record: d2153d1057d740fea1665d71ff66dd1c}; full qualifications and provenance.}\par
\section*{III.5.6 Clinical Information Notification}
\phantomsection
\label{navigation:III.5.6}
\addcontentsline{toc}{section}{III.5.6 Clinical Information Notification}

\label{element:31e03beac32946a09985ad52e0f9cce4}
\index{Clinical Information Notification (III.5.6)}
{\small BusinessInteraction · Canonical · current}\par

\begin{itemize}


\item \textbf{Clinical Information Notification}: An automated or operational notice informing participants that new or updated clinical information is now available.

\end{itemize}

{\small Recorded unresolved matters: canonicalName; exact scope remains in the canonical record.}\par
{\small Navigation classification: Native BusinessInteraction definition is determinate, but no governed capability context is represented for this interaction. No Application Interaction, integration flow or capability allocation is inferred.}\par
{\footnotesize \href{../../../../continuum/business/business-interaction/31e03beac32946a09985ad52e0f9cce4.md}{Canonical record: 31e03beac32946a09985ad52e0f9cce4}; full qualifications and provenance.}\par
\section*{III.5.7 Clinical Information Provision}
\phantomsection
\label{navigation:III.5.7}
\addcontentsline{toc}{section}{III.5.7 Clinical Information Provision}

\label{element:4d34b3705cb0405a910ddcacfdd6b6d3}
\index{Clinical Information Provision (III.5.7)}
{\small BusinessInteraction · Canonical · current}\par

\begin{itemize}


\item \textbf{Clinical Information Provision}: The delivery of requested clinical record content to an authorised clinician or system.

\end{itemize}

{\small Recorded unresolved matters: canonicalName; exact scope remains in the canonical record.}\par
{\small Navigation classification: Native BusinessInteraction definition is determinate, but no governed capability context is represented for this interaction. No Application Interaction, integration flow or capability allocation is inferred.}\par
{\footnotesize \href{../../../../continuum/business/business-interaction/4d34b3705cb0405a910ddcacfdd6b6d3.md}{Canonical record: 4d34b3705cb0405a910ddcacfdd6b6d3}; full qualifications and provenance.}\par
\section*{III.5.8 Clinical Information Query}
\phantomsection
\label{navigation:III.5.8}
\addcontentsline{toc}{section}{III.5.8 Clinical Information Query}

\label{element:8aabfdf9c83c4de486cf5f50c41c0d5e}
\index{Clinical Information Query (III.5.8)}
{\small BusinessInteraction · Canonical · current}\par

\begin{itemize}


\item \textbf{Clinical Information Query}: A targeted request to discover or search for clinical records, observations, or history for a specific healthcare subject.

\end{itemize}

{\small Recorded unresolved matters: canonicalName; exact scope remains in the canonical record.}\par
{\small Navigation classification: Native BusinessInteraction definition is determinate, but no governed capability context is represented for this interaction. No Application Interaction, integration flow or capability allocation is inferred.}\par
{\footnotesize \href{../../../../continuum/business/business-interaction/8aabfdf9c83c4de486cf5f50c41c0d5e.md}{Canonical record: 8aabfdf9c83c4de486cf5f50c41c0d5e}; full qualifications and provenance.}\par
\section*{III.5.9 Clinical Information Review}
\phantomsection
\label{navigation:III.5.9}
\addcontentsline{toc}{section}{III.5.9 Clinical Information Review}

\label{element:06dbfdb6b8a84bd29f67ca955ba23006}
\index{Clinical Information Review (III.5.9)}
{\small BusinessInteraction · Canonical · current}\par

\begin{itemize}


\item \textbf{Clinical Information Review}: The structured examination, validation, or clinical evaluation of clinical record content by a practitioner.

\end{itemize}

{\small Recorded unresolved matters: canonicalName; exact scope remains in the canonical record.}\par
{\small Navigation classification: Native BusinessInteraction definition is determinate, but no governed capability context is represented for this interaction. No Application Interaction, integration flow or capability allocation is inferred.}\par
{\footnotesize \href{../../../../continuum/business/business-interaction/06dbfdb6b8a84bd29f67ca955ba23006.md}{Canonical record: 06dbfdb6b8a84bd29f67ca955ba23006}; full qualifications and provenance.}\par
\section*{III.5.10 Clinical Information Supply}
\phantomsection
\label{navigation:III.5.10}
\addcontentsline{toc}{section}{III.5.10 Clinical Information Supply}

\label{element:a6e2525ea3294ee3ab9a4ec3dd4dfe9f}
\index{Clinical Information Supply (III.5.10)}
{\small BusinessInteraction · Canonical · current}\par

\begin{itemize}


\item \textbf{Clinical Information Supply}: Routine, authorised transmission or contribution of clinical information, discharge summaries, or diagnostic results into the integration environment.

\end{itemize}

{\small Recorded unresolved matters: canonicalName; exact scope remains in the canonical record.}\par
{\small Navigation classification: Native BusinessInteraction definition is determinate, but no governed capability context is represented for this interaction. No Application Interaction, integration flow or capability allocation is inferred.}\par
{\footnotesize \href{../../../../continuum/business/business-interaction/a6e2525ea3294ee3ab9a4ec3dd4dfe9f.md}{Canonical record: a6e2525ea3294ee3ab9a4ec3dd4dfe9f}; full qualifications and provenance.}\par
\section*{III.5.11 Complaint}
\phantomsection
\label{navigation:III.5.11}
\addcontentsline{toc}{section}{III.5.11 Complaint}

\label{element:29f2a74c8234411c9ac15458e66ca4bc}
\index{Complaint (III.5.11)}
{\small BusinessInteraction · Canonical · current}\par

\begin{itemize}


\item \textbf{Complaint}: A formal grievance lodged regarding clinical care, privacy breaches, or operational conduct.

\end{itemize}

{\small Recorded unresolved matters: canonicalName; exact scope remains in the canonical record.}\par
{\small Navigation classification: Native BusinessInteraction definition is determinate, but no governed capability context is represented for this interaction. No Application Interaction, integration flow or capability allocation is inferred.}\par
{\footnotesize \href{../../../../continuum/business/business-interaction/29f2a74c8234411c9ac15458e66ca4bc.md}{Canonical record: 29f2a74c8234411c9ac15458e66ca4bc}; full qualifications and provenance.}\par
\section*{III.5.12 Coordination Outcome}
\phantomsection
\label{navigation:III.5.12}
\addcontentsline{toc}{section}{III.5.12 Coordination Outcome}

\label{element:abb4947b64944790b20fd52524f0babc}
\index{Coordination Outcome (III.5.12)}
{\small BusinessInteraction · Canonical · current}\par

\begin{itemize}


\item \textbf{Coordination Outcome}: A recorded summary of coordinated care decisions, agreed milestones, or integrated care plans.

\end{itemize}

{\small Recorded unresolved matters: canonicalName; exact scope remains in the canonical record.}\par
{\small Navigation classification: Native BusinessInteraction definition is determinate, but no governed capability context is represented for this interaction. No Application Interaction, integration flow or capability allocation is inferred.}\par
{\footnotesize \href{../../../../continuum/business/business-interaction/abb4947b64944790b20fd52524f0babc.md}{Canonical record: abb4947b64944790b20fd52524f0babc}; full qualifications and provenance.}\par
\section*{III.5.13 FOI Request}
\phantomsection
\label{navigation:III.5.13}
\addcontentsline{toc}{section}{III.5.13 FOI Request}

\label{element:2c7ca48559fb4a72b4d9c250df4015ce}
\index{FOI Request (III.5.13)}
{\small BusinessInteraction · Canonical · current}\par

\begin{itemize}


\item \textbf{FOI Request}: A formal statutory request for access to records under Freedom of Information or privacy legislation.

\end{itemize}

{\small Recorded unresolved matters: canonicalName; exact scope remains in the canonical record.}\par
{\small Navigation classification: Native BusinessInteraction definition is determinate, but no governed capability context is represented for this interaction. No Application Interaction, integration flow or capability allocation is inferred.}\par
{\footnotesize \href{../../../../continuum/business/business-interaction/2c7ca48559fb4a72b4d9c250df4015ce.md}{Canonical record: 2c7ca48559fb4a72b4d9c250df4015ce}; full qualifications and provenance.}\par
\section*{III.5.14 FOI Response}
\phantomsection
\label{navigation:III.5.14}
\addcontentsline{toc}{section}{III.5.14 FOI Response}

\label{element:909faede10b5452cbee2d7d052635f47}
\index{FOI Response (III.5.14)}
{\small BusinessInteraction · Canonical · current}\par

\begin{itemize}


\item \textbf{FOI Response}: The governed disclosure, redaction, or refusal response to a statutory access request.

\end{itemize}

{\small Recorded unresolved matters: canonicalName; exact scope remains in the canonical record.}\par
{\small Navigation classification: Native BusinessInteraction definition is determinate, but no governed capability context is represented for this interaction. No Application Interaction, integration flow or capability allocation is inferred.}\par
{\footnotesize \href{../../../../continuum/business/business-interaction/909faede10b5452cbee2d7d052635f47.md}{Canonical record: 909faede10b5452cbee2d7d052635f47}; full qualifications and provenance.}\par
\section*{III.5.15 Identity Assertion}
\phantomsection
\label{navigation:III.5.15}
\addcontentsline{toc}{section}{III.5.15 Identity Assertion}

\label{element:47b9abea0964451eb0fa419b464b4505}
\index{Identity Assertion (III.5.15)}
{\small BusinessInteraction · Canonical · current}\par

\begin{itemize}


\item \textbf{Identity Assertion}: A claim presented by an actor regarding their identity or an identifier representing them.

\end{itemize}

{\small Recorded unresolved matters: canonicalName; exact scope remains in the canonical record.}\par
{\small Navigation classification: Native BusinessInteraction definition is determinate, but no governed capability context is represented for this interaction. No Application Interaction, integration flow or capability allocation is inferred.}\par
{\footnotesize \href{../../../../continuum/business/business-interaction/47b9abea0964451eb0fa419b464b4505.md}{Canonical record: 47b9abea0964451eb0fa419b464b4505}; full qualifications and provenance.}\par
\section*{III.5.16 Identity Establishment}
\phantomsection
\label{navigation:III.5.16}
\addcontentsline{toc}{section}{III.5.16 Identity Establishment}

\label{element:3d3459edfa884ef0abdcfd3772b538e4}
\index{Identity Establishment (III.5.16)}
{\small BusinessInteraction · Canonical · current}\par

\begin{itemize}


\item \textbf{Identity Establishment}: The governed recording and official recognition of a newly identified entity in an authoritative registry.

\end{itemize}

{\small Recorded unresolved matters: canonicalName; exact scope remains in the canonical record.}\par
{\small Navigation classification: Native BusinessInteraction definition is determinate, but no governed capability context is represented for this interaction. No Application Interaction, integration flow or capability allocation is inferred.}\par
{\footnotesize \href{../../../../continuum/business/business-interaction/3d3459edfa884ef0abdcfd3772b538e4.md}{Canonical record: 3d3459edfa884ef0abdcfd3772b538e4}; full qualifications and provenance.}\par
\section*{III.5.17 Identity Query}
\phantomsection
\label{navigation:III.5.17}
\addcontentsline{toc}{section}{III.5.17 Identity Query}

\label{element:d119aa62e95f4d7aa01822401ca4203c}
\index{Identity Query (III.5.17)}
{\small BusinessInteraction · Canonical · current}\par

\begin{itemize}


\item \textbf{Identity Query}: A search or lookup request to discover or disambiguate identity records matching supplied demographic traits.

\end{itemize}

{\small Recorded unresolved matters: canonicalName; exact scope remains in the canonical record.}\par
{\small Navigation classification: Native BusinessInteraction definition is determinate, but no governed capability context is represented for this interaction. No Application Interaction, integration flow or capability allocation is inferred.}\par
{\footnotesize \href{../../../../continuum/business/business-interaction/d119aa62e95f4d7aa01822401ca4203c.md}{Canonical record: d119aa62e95f4d7aa01822401ca4203c}; full qualifications and provenance.}\par
\section*{III.5.18 Identity Reconciliation}
\phantomsection
\label{navigation:III.5.18}
\addcontentsline{toc}{section}{III.5.18 Identity Reconciliation}

\label{element:332cdc30ae0346c6b86233db71caf683}
\index{Identity Reconciliation (III.5.18)}
{\small BusinessInteraction · Canonical · current}\par

\begin{itemize}


\item \textbf{Identity Reconciliation}: The governed merging, unlinking, or cross-authority correlation of duplicate or fragmented identity records.

\end{itemize}

{\small Recorded unresolved matters: canonicalName; exact scope remains in the canonical record.}\par
{\small Navigation classification: Native BusinessInteraction definition is determinate, but no governed capability context is represented for this interaction. No Application Interaction, integration flow or capability allocation is inferred.}\par
{\footnotesize \href{../../../../continuum/business/business-interaction/332cdc30ae0346c6b86233db71caf683.md}{Canonical record: 332cdc30ae0346c6b86233db71caf683}; full qualifications and provenance.}\par
\section*{III.5.19 Identity Verification}
\phantomsection
\label{navigation:III.5.19}
\addcontentsline{toc}{section}{III.5.19 Identity Verification}

\label{element:1a56f778873a4094b5a0f84cae842d35}
\index{Identity Verification (III.5.19)}
{\small BusinessInteraction · Canonical · current}\par

\begin{itemize}


\item \textbf{Identity Verification}: The formal validation of an asserted identity against trusted credentials, identity documents, or biometric proofs.

\end{itemize}

{\small Recorded unresolved matters: canonicalName; exact scope remains in the canonical record.}\par
{\small Navigation classification: Native BusinessInteraction definition is determinate, but no governed capability context is represented for this interaction. No Application Interaction, integration flow or capability allocation is inferred.}\par
{\footnotesize \href{../../../../continuum/business/business-interaction/1a56f778873a4094b5a0f84cae842d35.md}{Canonical record: 1a56f778873a4094b5a0f84cae842d35}; full qualifications and provenance.}\par
\section*{III.5.20 Incident Registration}
\phantomsection
\label{navigation:III.5.20}
\addcontentsline{toc}{section}{III.5.20 Incident Registration}

\label{element:126276873bde4105a08e21bb3187f5b3}
\index{Incident Registration (III.5.20)}
{\small BusinessInteraction · Canonical · current}\par

\begin{itemize}


\item \textbf{Incident Registration}: The formal recording of a clinical, operational, or safety event into a governed incident management framework.

\end{itemize}

{\small Recorded unresolved matters: canonicalName; exact scope remains in the canonical record.}\par
{\small Navigation classification: Native BusinessInteraction definition is determinate, but no governed capability context is represented for this interaction. No Application Interaction, integration flow or capability allocation is inferred.}\par
{\footnotesize \href{../../../../continuum/business/business-interaction/126276873bde4105a08e21bb3187f5b3.md}{Canonical record: 126276873bde4105a08e21bb3187f5b3}; full qualifications and provenance.}\par
\section*{III.5.21 Incident Report}
\phantomsection
\label{navigation:III.5.21}
\addcontentsline{toc}{section}{III.5.21 Incident Report}

\label{element:7cb1eb055be44eefb7092356510acc77}
\index{Incident Report (III.5.21)}
{\small BusinessInteraction · Canonical · current}\par

\begin{itemize}


\item \textbf{Incident Report}: The formal communication and investigation findings regarding an incident.

\end{itemize}

{\small Recorded unresolved matters: canonicalName; exact scope remains in the canonical record.}\par
{\small Navigation classification: Native BusinessInteraction definition is determinate, but no governed capability context is represented for this interaction. No Application Interaction, integration flow or capability allocation is inferred.}\par
{\footnotesize \href{../../../../continuum/business/business-interaction/7cb1eb055be44eefb7092356510acc77.md}{Canonical record: 7cb1eb055be44eefb7092356510acc77}; full qualifications and provenance.}\par
\section*{III.5.22 Information Access}
\phantomsection
\label{navigation:III.5.22}
\addcontentsline{toc}{section}{III.5.22 Information Access}

\label{element:b2624626d06044fb9767a96045cd0302}
\index{Information Access (III.5.22)}
{\small BusinessInteraction · Canonical · current}\par

\begin{itemize}


\item \textbf{Information Access}: An authorised consumer exercising legitimate authority to retrieve or inspect governed health information.

\end{itemize}

{\small Recorded unresolved matters: canonicalName; exact scope remains in the canonical record.}\par
{\small Navigation classification: Native BusinessInteraction definition is determinate, but no governed capability context is represented for this interaction. No Application Interaction, integration flow or capability allocation is inferred.}\par
{\footnotesize \href{../../../../continuum/business/business-interaction/b2624626d06044fb9767a96045cd0302.md}{Canonical record: b2624626d06044fb9767a96045cd0302}; full qualifications and provenance.}\par
\section*{III.5.23 Information Challenge}
\phantomsection
\label{navigation:III.5.23}
\addcontentsline{toc}{section}{III.5.23 Information Challenge}

\label{element:44d71c390a8f435cbb596601e627bdfa}
\index{Information Challenge (III.5.23)}
{\small BusinessInteraction · Canonical · current}\par

\begin{itemize}


\item \textbf{Information Challenge}: A formal dispute or query raised regarding the accuracy, completeness, or attribution of a health record.

\end{itemize}

{\small Recorded unresolved matters: canonicalName; exact scope remains in the canonical record.}\par
{\small Navigation classification: Native BusinessInteraction definition is determinate, but no governed capability context is represented for this interaction. No Application Interaction, integration flow or capability allocation is inferred.}\par
{\footnotesize \href{../../../../continuum/business/business-interaction/44d71c390a8f435cbb596601e627bdfa.md}{Canonical record: 44d71c390a8f435cbb596601e627bdfa}; full qualifications and provenance.}\par
\section*{III.5.24 Information Control}
\phantomsection
\label{navigation:III.5.24}
\addcontentsline{toc}{section}{III.5.24 Information Control}

\label{element:773133f23f7540eb9b90d39ca872aebd}
\index{Information Control (III.5.24)}
{\small BusinessInteraction · Canonical · current}\par

\begin{itemize}


\item \textbf{Information Control}: The assertion or enforcement of confidentiality restrictions, consent directives, or masking rules.

\end{itemize}

{\small Recorded unresolved matters: canonicalName; exact scope remains in the canonical record.}\par
{\small Navigation classification: Native BusinessInteraction definition is determinate, but no governed capability context is represented for this interaction. No Application Interaction, integration flow or capability allocation is inferred.}\par
{\footnotesize \href{../../../../continuum/business/business-interaction/773133f23f7540eb9b90d39ca872aebd.md}{Canonical record: 773133f23f7540eb9b90d39ca872aebd}; full qualifications and provenance.}\par
\section*{III.5.25 Information Correction}
\phantomsection
\label{navigation:III.5.25}
\addcontentsline{toc}{section}{III.5.25 Information Correction}

\label{element:3f1db661c97b4fa692e3d70278a40226}
\index{Information Correction (III.5.25)}
{\small BusinessInteraction · Canonical · current}\par

\begin{itemize}


\item \textbf{Information Correction}: The governed rectification, amendment, or retraction of erroneous or misattributed health information.

\end{itemize}

{\small Recorded unresolved matters: canonicalName; exact scope remains in the canonical record.}\par
{\small Navigation classification: Native BusinessInteraction definition is determinate, but no governed capability context is represented for this interaction. No Application Interaction, integration flow or capability allocation is inferred.}\par
{\footnotesize \href{../../../../continuum/business/business-interaction/3f1db661c97b4fa692e3d70278a40226.md}{Canonical record: 3f1db661c97b4fa692e3d70278a40226}; full qualifications and provenance.}\par
\section*{III.5.26 Information Disclosure}
\phantomsection
\label{navigation:III.5.26}
\addcontentsline{toc}{section}{III.5.26 Information Disclosure}

\label{element:47536a5d81aa4d8eb2d13054b864d1bf}
\index{Information Disclosure (III.5.26)}
{\small BusinessInteraction · Canonical · current}\par

\begin{itemize}


\item \textbf{Information Disclosure}: An information custodian or system exercising authority or legal obligation to release records to an external entity.

\end{itemize}

{\small Recorded unresolved matters: canonicalName; exact scope remains in the canonical record.}\par
{\small Navigation classification: Native BusinessInteraction definition is determinate, but no governed capability context is represented for this interaction. No Application Interaction, integration flow or capability allocation is inferred.}\par
{\footnotesize \href{../../../../continuum/business/business-interaction/47536a5d81aa4d8eb2d13054b864d1bf.md}{Canonical record: 47536a5d81aa4d8eb2d13054b864d1bf}; full qualifications and provenance.}\par
\section*{III.5.27 Information Publishing}
\phantomsection
\label{navigation:III.5.27}
\addcontentsline{toc}{section}{III.5.27 Information Publishing}

\label{element:6fb6f66531874c2d983071a04b5934c2}
\index{Information Publishing (III.5.27)}
{\small BusinessInteraction · Canonical · current}\par

\begin{itemize}


\item \textbf{Information Publishing}: The distribution of governed, approved health information into public or shared access registries.

\end{itemize}

{\small Recorded unresolved matters: canonicalName; exact scope remains in the canonical record.}\par
{\small Navigation classification: Native BusinessInteraction definition is determinate, but no governed capability context is represented for this interaction. No Application Interaction, integration flow or capability allocation is inferred.}\par
{\footnotesize \href{../../../../continuum/business/business-interaction/6fb6f66531874c2d983071a04b5934c2.md}{Canonical record: 6fb6f66531874c2d983071a04b5934c2}; full qualifications and provenance.}\par
\section*{III.5.28 Information Qualification}
\phantomsection
\label{navigation:III.5.28}
\addcontentsline{toc}{section}{III.5.28 Information Qualification}

\label{element:6b606df4675f48b3baba03f4b2ca7602}
\index{Information Qualification (III.5.28)}
{\small BusinessInteraction · Canonical · current}\par

\begin{itemize}


\item \textbf{Information Qualification}: The formal annotation, verification stamp, or clinical validation applied to an information asset.

\end{itemize}

{\small Recorded unresolved matters: canonicalName; exact scope remains in the canonical record.}\par
{\small Navigation classification: Native BusinessInteraction definition is determinate, but no governed capability context is represented for this interaction. No Application Interaction, integration flow or capability allocation is inferred.}\par
{\footnotesize \href{../../../../continuum/business/business-interaction/6b606df4675f48b3baba03f4b2ca7602.md}{Canonical record: 6b606df4675f48b3baba03f4b2ca7602}; full qualifications and provenance.}\par
\section*{III.5.29 Information Submission}
\phantomsection
\label{navigation:III.5.29}
\addcontentsline{toc}{section}{III.5.29 Information Submission}

\label{element:d3186ad3c8a74e92a80cb4c1f8f32fb6}
\index{Information Submission (III.5.29)}
{\small BusinessInteraction · Canonical · current}\par

\begin{itemize}


\item \textbf{Information Submission}: The formal submission of clinical documents or datasets for governance validation, indexing, or archival.

\end{itemize}

{\small Recorded unresolved matters: canonicalName; exact scope remains in the canonical record.}\par
{\small Navigation classification: Native BusinessInteraction definition is determinate, but no governed capability context is represented for this interaction. No Application Interaction, integration flow or capability allocation is inferred.}\par
{\footnotesize \href{../../../../continuum/business/business-interaction/d3186ad3c8a74e92a80cb4c1f8f32fb6.md}{Canonical record: d3186ad3c8a74e92a80cb4c1f8f32fb6}; full qualifications and provenance.}\par
\section*{III.5.30 Operational Alert}
\phantomsection
\label{navigation:III.5.30}
\addcontentsline{toc}{section}{III.5.30 Operational Alert}

\label{element:3697066c0d634076be5a8d745e3fdd91}
\index{Operational Alert (III.5.30)}
{\small BusinessInteraction · Canonical · current}\par

\begin{itemize}


\item \textbf{Operational Alert}: An urgent operational warning indicating capacity exhaustion, system degradation, or logistics bottlenecks.

\end{itemize}

{\small Recorded unresolved matters: canonicalName; exact scope remains in the canonical record.}\par
{\small Navigation classification: Native BusinessInteraction definition is determinate, but no governed capability context is represented for this interaction. No Application Interaction, integration flow or capability allocation is inferred.}\par
{\footnotesize \href{../../../../continuum/business/business-interaction/3697066c0d634076be5a8d745e3fdd91.md}{Canonical record: 3697066c0d634076be5a8d745e3fdd91}; full qualifications and provenance.}\par
\section*{III.5.31 Operational Coordination}
\phantomsection
\label{navigation:III.5.31}
\addcontentsline{toc}{section}{III.5.31 Operational Coordination}

\label{element:cdb5d645d3984eaa8ad5e56a81023c59}
\index{Operational Coordination (III.5.31)}
{\small BusinessInteraction · Canonical · current}\par

\begin{itemize}


\item \textbf{Operational Coordination}: Tactical alignment regarding resource allocation, bed distribution, queue balancing, and facility throughput.

\end{itemize}

{\small Recorded unresolved matters: canonicalName; exact scope remains in the canonical record.}\par
{\small Navigation classification: Native BusinessInteraction definition is determinate, but no governed capability context is represented for this interaction. No Application Interaction, integration flow or capability allocation is inferred.}\par
{\footnotesize \href{../../../../continuum/business/business-interaction/cdb5d645d3984eaa8ad5e56a81023c59.md}{Canonical record: cdb5d645d3984eaa8ad5e56a81023c59}; full qualifications and provenance.}\par
\section*{III.5.32 Operational Coordination Outcome}
\phantomsection
\label{navigation:III.5.32}
\addcontentsline{toc}{section}{III.5.32 Operational Coordination Outcome}

\label{element:d40c5ee0dbaa450dab073103f25ded6d}
\index{Operational Coordination Outcome (III.5.32)}
{\small BusinessInteraction · Canonical · current}\par

\begin{itemize}


\item \textbf{Operational Coordination Outcome}: Recorded agreements from operational huddles or bed-meeting decisions.

\end{itemize}

{\small Recorded unresolved matters: canonicalName; exact scope remains in the canonical record.}\par
{\small Navigation classification: Native BusinessInteraction definition is determinate, but no governed capability context is represented for this interaction. No Application Interaction, integration flow or capability allocation is inferred.}\par
{\footnotesize \href{../../../../continuum/business/business-interaction/d40c5ee0dbaa450dab073103f25ded6d.md}{Canonical record: d40c5ee0dbaa450dab073103f25ded6d}; full qualifications and provenance.}\par
\section*{III.5.33 Operational Event Notification}
\phantomsection
\label{navigation:III.5.33}
\addcontentsline{toc}{section}{III.5.33 Operational Event Notification}

\label{element:7a2b6e8802e24e018af89f265271f5ca}
\index{Operational Event Notification (III.5.33)}
{\small BusinessInteraction · Canonical · current}\par

\begin{itemize}


\item \textbf{Operational Event Notification}: A notice indicating a facility or logistics state change (e.g., bed vacated, cleaning completed).

\end{itemize}

{\small Recorded unresolved matters: canonicalName; exact scope remains in the canonical record.}\par
{\small Navigation classification: Native BusinessInteraction definition is determinate, but no governed capability context is represented for this interaction. No Application Interaction, integration flow or capability allocation is inferred.}\par
{\footnotesize \href{../../../../continuum/business/business-interaction/7a2b6e8802e24e018af89f265271f5ca.md}{Canonical record: 7a2b6e8802e24e018af89f265271f5ca}; full qualifications and provenance.}\par
\section*{III.5.34 Participation Query}
\phantomsection
\label{navigation:III.5.34}
\addcontentsline{toc}{section}{III.5.34 Participation Query}

\label{element:1ce0ba25408c47f9bd82953fba605c40}
\index{Participation Query (III.5.34)}
{\small BusinessInteraction · Canonical · current}\par

\begin{itemize}


\item \textbf{Participation Query}: A lookup to determine whether an actor is actively enrolled in a specific care programme or team.

\end{itemize}

{\small Recorded unresolved matters: canonicalName; exact scope remains in the canonical record.}\par
{\small Navigation classification: Native BusinessInteraction definition is determinate, but no governed capability context is represented for this interaction. No Application Interaction, integration flow or capability allocation is inferred.}\par
{\footnotesize \href{../../../../continuum/business/business-interaction/1ce0ba25408c47f9bd82953fba605c40.md}{Canonical record: 1ce0ba25408c47f9bd82953fba605c40}; full qualifications and provenance.}\par
\section*{III.5.35 Participation Registration}
\phantomsection
\label{navigation:III.5.35}
\addcontentsline{toc}{section}{III.5.35 Participation Registration}

\label{element:906c61eb392f4b538f6d0a79917d21cc}
\index{Participation Registration (III.5.35)}
{\small BusinessInteraction · Canonical · current}\par

\begin{itemize}


\item \textbf{Participation Registration}: The formal enrolment or addition of an actor into a multidisciplinary programme, registry, or care collective.

\end{itemize}

{\small Recorded unresolved matters: canonicalName; exact scope remains in the canonical record.}\par
{\small Navigation classification: Native BusinessInteraction definition is determinate, but no governed capability context is represented for this interaction. No Application Interaction, integration flow or capability allocation is inferred.}\par
{\footnotesize \href{../../../../continuum/business/business-interaction/906c61eb392f4b538f6d0a79917d21cc.md}{Canonical record: 906c61eb392f4b538f6d0a79917d21cc}; full qualifications and provenance.}\par
\section*{III.5.36 Participation Withdrawal}
\phantomsection
\label{navigation:III.5.36}
\addcontentsline{toc}{section}{III.5.36 Participation Withdrawal}

\label{element:ce1901f036364b7291451bfbd4f25c00}
\index{Participation Withdrawal (III.5.36)}
{\small BusinessInteraction · Canonical · current}\par

\begin{itemize}


\item \textbf{Participation Withdrawal}: The formal removal or resignation of an actor from a participation context.

\end{itemize}

{\small Recorded unresolved matters: canonicalName; exact scope remains in the canonical record.}\par
{\small Navigation classification: Native BusinessInteraction definition is determinate, but no governed capability context is represented for this interaction. No Application Interaction, integration flow or capability allocation is inferred.}\par
{\footnotesize \href{../../../../continuum/business/business-interaction/ce1901f036364b7291451bfbd4f25c00.md}{Canonical record: ce1901f036364b7291451bfbd4f25c00}; full qualifications and provenance.}\par
\section*{III.5.37 Policy Direction}
\phantomsection
\label{navigation:III.5.37}
\addcontentsline{toc}{section}{III.5.37 Policy Direction}

\label{element:b2c573b855ad4c7abc4f02b2ae484a65}
\index{Policy Direction (III.5.37)}
{\small BusinessInteraction · Canonical · current}\par

\begin{itemize}


\item \textbf{Policy Direction}: An enterprise or jurisdictional directive establishing mandatory clinical or operational policy rules.

\end{itemize}

{\small Recorded unresolved matters: canonicalName; exact scope remains in the canonical record.}\par
{\small Navigation classification: Native BusinessInteraction definition is determinate, but no governed capability context is represented for this interaction. No Application Interaction, integration flow or capability allocation is inferred.}\par
{\footnotesize \href{../../../../continuum/business/business-interaction/b2c573b855ad4c7abc4f02b2ae484a65.md}{Canonical record: b2c573b855ad4c7abc4f02b2ae484a65}; full qualifications and provenance.}\par
\section*{III.5.38 Reassignment}
\phantomsection
\label{navigation:III.5.38}
\addcontentsline{toc}{section}{III.5.38 Reassignment}

\label{element:d4fb961b6148436686632668c98a6139}
\index{Reassignment (III.5.38)}
{\small BusinessInteraction · Canonical · current}\par

\begin{itemize}


\item \textbf{Reassignment}: The redirection of an assigned operational task to an alternative worker or unit due to operational constraints.

\end{itemize}

{\small Recorded unresolved matters: canonicalName; exact scope remains in the canonical record.}\par
{\small Navigation classification: Native BusinessInteraction definition is determinate, but no governed capability context is represented for this interaction. No Application Interaction, integration flow or capability allocation is inferred.}\par
{\footnotesize \href{../../../../continuum/business/business-interaction/d4fb961b6148436686632668c98a6139.md}{Canonical record: d4fb961b6148436686632668c98a6139}; full qualifications and provenance.}\par
\section*{III.5.39 Regulatory Direction}
\phantomsection
\label{navigation:III.5.39}
\addcontentsline{toc}{section}{III.5.39 Regulatory Direction}

\label{element:baf4be48ac11472aa130c71e6474d9a4}
\index{Regulatory Direction (III.5.39)}
{\small BusinessInteraction · Canonical · current}\par

\begin{itemize}


\item \textbf{Regulatory Direction}: A legally binding order or compliance requirement issued by a statutory regulatory body.

\end{itemize}

{\small Recorded unresolved matters: canonicalName; exact scope remains in the canonical record.}\par
{\small Navigation classification: Native BusinessInteraction definition is determinate, but no governed capability context is represented for this interaction. No Application Interaction, integration flow or capability allocation is inferred.}\par
{\footnotesize \href{../../../../continuum/business/business-interaction/baf4be48ac11472aa130c71e6474d9a4.md}{Canonical record: baf4be48ac11472aa130c71e6474d9a4}; full qualifications and provenance.}\par
\section*{III.5.40 Relationship Assertion}
\phantomsection
\label{navigation:III.5.40}
\addcontentsline{toc}{section}{III.5.40 Relationship Assertion}

\label{element:9a53f10eac7748a2b4a9d88da75f8814}
\index{Relationship Assertion (III.5.40)}
{\small BusinessInteraction · Canonical · current}\par

\begin{itemize}


\item \textbf{Relationship Assertion}: A claim that a specific relationship exists between two entities (e.g., claiming to be a patient's carer).

\end{itemize}

{\small Recorded unresolved matters: canonicalName; exact scope remains in the canonical record.}\par
{\small Navigation classification: Native BusinessInteraction definition is determinate, but no governed capability context is represented for this interaction. No Application Interaction, integration flow or capability allocation is inferred.}\par
{\footnotesize \href{../../../../continuum/business/business-interaction/9a53f10eac7748a2b4a9d88da75f8814.md}{Canonical record: 9a53f10eac7748a2b4a9d88da75f8814}; full qualifications and provenance.}\par
\section*{III.5.41 Relationship Establishment}
\phantomsection
\label{navigation:III.5.41}
\addcontentsline{toc}{section}{III.5.41 Relationship Establishment}

\label{element:f451d56104944498909549ab7f864003}
\index{Relationship Establishment (III.5.41)}
{\small BusinessInteraction · Canonical · current}\par

\begin{itemize}


\item \textbf{Relationship Establishment}: The governed recording and formal activation of an acknowledged relationship in an authoritative registry.

\end{itemize}

{\small Recorded unresolved matters: canonicalName; exact scope remains in the canonical record.}\par
{\small Navigation classification: Native BusinessInteraction definition is determinate, but no governed capability context is represented for this interaction. No Application Interaction, integration flow or capability allocation is inferred.}\par
{\footnotesize \href{../../../../continuum/business/business-interaction/f451d56104944498909549ab7f864003.md}{Canonical record: f451d56104944498909549ab7f864003}; full qualifications and provenance.}\par
\section*{III.5.42 Relationship Query}
\phantomsection
\label{navigation:III.5.42}
\addcontentsline{toc}{section}{III.5.42 Relationship Query}

\label{element:b30441d9b44a4bc29771e8f6e94f6b5a}
\index{Relationship Query (III.5.42)}
{\small BusinessInteraction · Canonical · current}\par

\begin{itemize}


\item \textbf{Relationship Query}: An inquiry to discover existing care relationships, nominated representatives, or practitioner affiliations.

\end{itemize}

{\small Recorded unresolved matters: canonicalName; exact scope remains in the canonical record.}\par
{\small Navigation classification: Native BusinessInteraction definition is determinate, but no governed capability context is represented for this interaction. No Application Interaction, integration flow or capability allocation is inferred.}\par
{\footnotesize \href{../../../../continuum/business/business-interaction/b30441d9b44a4bc29771e8f6e94f6b5a.md}{Canonical record: b30441d9b44a4bc29771e8f6e94f6b5a}; full qualifications and provenance.}\par
\section*{III.5.43 Relationship Termination}
\phantomsection
\label{navigation:III.5.43}
\addcontentsline{toc}{section}{III.5.43 Relationship Termination}

\label{element:e536ab3a1bbe4a36a99635458ba208f6}
\index{Relationship Termination (III.5.43)}
{\small BusinessInteraction · Canonical · current}\par

\begin{itemize}


\item \textbf{Relationship Termination}: The formal dissolution, expiry, or revocation of an established relationship.

\end{itemize}

{\small Recorded unresolved matters: canonicalName; exact scope remains in the canonical record.}\par
{\small Navigation classification: Native BusinessInteraction definition is determinate, but no governed capability context is represented for this interaction. No Application Interaction, integration flow or capability allocation is inferred.}\par
{\footnotesize \href{../../../../continuum/business/business-interaction/e536ab3a1bbe4a36a99635458ba208f6.md}{Canonical record: e536ab3a1bbe4a36a99635458ba208f6}; full qualifications and provenance.}\par
\section*{III.5.44 Relationship Verification}
\phantomsection
\label{navigation:III.5.44}
\addcontentsline{toc}{section}{III.5.44 Relationship Verification}

\label{element:c71453e1b0ff46a7b91e27a1745d842b}
\index{Relationship Verification (III.5.44)}
{\small BusinessInteraction · Canonical · current}\par

\begin{itemize}


\item \textbf{Relationship Verification}: The validation of legal or clinical authority supporting an asserted relationship (e.g., verifying power of attorney).

\end{itemize}

{\small Recorded unresolved matters: canonicalName; exact scope remains in the canonical record.}\par
{\small Navigation classification: Native BusinessInteraction definition is determinate, but no governed capability context is represented for this interaction. No Application Interaction, integration flow or capability allocation is inferred.}\par
{\footnotesize \href{../../../../continuum/business/business-interaction/c71453e1b0ff46a7b91e27a1745d842b.md}{Canonical record: c71453e1b0ff46a7b91e27a1745d842b}; full qualifications and provenance.}\par
\section*{III.5.45 Resource Allocation}
\phantomsection
\label{navigation:III.5.45}
\addcontentsline{toc}{section}{III.5.45 Resource Allocation}

\label{element:8479c1d4b2304d9dabc3716c6cb20e83}
\index{Resource Allocation (III.5.45)}
{\small BusinessInteraction · Canonical · current}\par

\begin{itemize}


\item \textbf{Resource Allocation}: The formal assignment or reservation of a physical care-place, bed, theatre slot, or device to a patient.

\end{itemize}

{\small Recorded unresolved matters: canonicalName; exact scope remains in the canonical record.}\par
{\small Navigation classification: Native BusinessInteraction definition is determinate, but no governed capability context is represented for this interaction. No Application Interaction, integration flow or capability allocation is inferred.}\par
{\footnotesize \href{../../../../continuum/business/business-interaction/8479c1d4b2304d9dabc3716c6cb20e83.md}{Canonical record: 8479c1d4b2304d9dabc3716c6cb20e83}; full qualifications and provenance.}\par
\section*{III.5.46 Resource Reporting}
\phantomsection
\label{navigation:III.5.46}
\addcontentsline{toc}{section}{III.5.46 Resource Reporting}

\label{element:05bf8e557ce14f0eb1c869d07c5a93ca}
\index{Resource Reporting (III.5.46)}
{\small BusinessInteraction · Canonical · current}\par

\begin{itemize}


\item \textbf{Resource Reporting}: Telemetry or status reporting on the availability, capacity, or status of physical beds, theatres, or clinical equipment.

\end{itemize}

{\small Recorded unresolved matters: canonicalName; exact scope remains in the canonical record.}\par
{\small Navigation classification: Native BusinessInteraction definition is determinate, but no governed capability context is represented for this interaction. No Application Interaction, integration flow or capability allocation is inferred.}\par
{\footnotesize \href{../../../../continuum/business/business-interaction/05bf8e557ce14f0eb1c869d07c5a93ca.md}{Canonical record: 05bf8e557ce14f0eb1c869d07c5a93ca}; full qualifications and provenance.}\par
\section*{III.5.47 Review Outcome}
\phantomsection
\label{navigation:III.5.47}
\addcontentsline{toc}{section}{III.5.47 Review Outcome}

\label{element:45bf28eb44e64554b568db48a30ac43c}
\index{Review Outcome (III.5.47)}
{\small BusinessInteraction · Canonical · current}\par

\begin{itemize}


\item \textbf{Review Outcome}: The documented findings, consensus decisions, and quality recommendations resulting from a clinical review.

\end{itemize}

{\small Recorded unresolved matters: canonicalName; exact scope remains in the canonical record.}\par
{\small Navigation classification: Native BusinessInteraction definition is determinate, but no governed capability context is represented for this interaction. No Application Interaction, integration flow or capability allocation is inferred.}\par
{\footnotesize \href{../../../../continuum/business/business-interaction/45bf28eb44e64554b568db48a30ac43c.md}{Canonical record: 45bf28eb44e64554b568db48a30ac43c}; full qualifications and provenance.}\par
\section*{III.5.48 Review Request}
\phantomsection
\label{navigation:III.5.48}
\addcontentsline{toc}{section}{III.5.48 Review Request}

\label{element:480556d544ad4dc18f6945f213044e9b}
\index{Review Request (III.5.48)}
{\small BusinessInteraction · Canonical · current}\par

\begin{itemize}


\item \textbf{Review Request}: A formal submission requesting clinical peer review, adverse event evaluation, or diagnostic second opinion.

\end{itemize}

{\small Recorded unresolved matters: canonicalName; exact scope remains in the canonical record.}\par
{\small Navigation classification: Native BusinessInteraction definition is determinate, but no governed capability context is represented for this interaction. No Application Interaction, integration flow or capability allocation is inferred.}\par
{\footnotesize \href{../../../../continuum/business/business-interaction/480556d544ad4dc18f6945f213044e9b.md}{Canonical record: 480556d544ad4dc18f6945f213044e9b}; full qualifications and provenance.}\par
\section*{III.5.49 Sanction}
\phantomsection
\label{navigation:III.5.49}
\addcontentsline{toc}{section}{III.5.49 Sanction}

\label{element:f4dc8ed1c034473784a5c8ec4227c7d6}
\index{Sanction (III.5.49)}
{\small BusinessInteraction · Canonical · current}\par

\begin{itemize}


\item \textbf{Sanction}: A formal disciplinary, regulatory, or administrative restriction imposed upon a practitioner, organisation, or system.

\end{itemize}

{\small Recorded unresolved matters: canonicalName; exact scope remains in the canonical record.}\par
{\small Navigation classification: Native BusinessInteraction definition is determinate, but no governed capability context is represented for this interaction. No Application Interaction, integration flow or capability allocation is inferred.}\par
{\footnotesize \href{../../../../continuum/business/business-interaction/f4dc8ed1c034473784a5c8ec4227c7d6.md}{Canonical record: f4dc8ed1c034473784a5c8ec4227c7d6}; full qualifications and provenance.}\par
\section*{III.5.50 Service Assurance Outcome Communication}
\phantomsection
\label{navigation:III.5.50}
\addcontentsline{toc}{section}{III.5.50 Service Assurance Outcome Communication}

\label{element:73b6d1ab8ecc48c4bdea98f013288c9c}
\index{Service Assurance Outcome Communication (III.5.50)}
{\small BusinessInteraction · Canonical · current}\par

\begin{itemize}


\item \textbf{Canonical Name}: \texttt{\seqsplit{Service Assurance Outcome Communication}}.

\item \textbf{Participants}: \texttt{\seqsplit{Service Guardian}} and \textbf{one or more authorised recipient Roles} from \texttt{\seqsplit{Care Coordinator}}, \texttt{\seqsplit{Service Coordinator}}, \texttt{\seqsplit{Regulator}}, \texttt{\seqsplit{Policy Authority}}, \texttt{\seqsplit{System Steward}}.

\item \textbf{Purpose}: Communicates an established Assurance Finding or Conclusion to authorised Business Roles for use within their respective responsibilities.

\item \textbf{Associated Business Service}: \href{\detokenize{../../../../../docs/markdown/03-business-architecture/behaviours/health-service-assurance.md#communicate-service-assurance-outcome}}{Communicate Service Assurance Outcome}.

\item \textbf{Boundary}: The outcome is already established; communication does not assess/adjudicate, reinterpret or alter it. This Interaction is not necessarily a broadcast: one, some or all legitimate recipients may receive the applicable outcome according to authority, responsibility and information requirements. The requesting Role, status consumer and outcome recipient need not coincide. Communication need not require a documentary report.

\end{itemize}


Service Guardian does not become responsible for the recipient's interpretation within its own authority, management action, remediation, escalation, policy change, process improvement, clinical action or other subsequent response. Financial relevance does not establish Harmonia financial-governance responsibility; the \href{\detokenize{../../../../../docs/markdown/03-business-architecture/behaviours/health-service-assurance.md#financial-governance-exclusion}}{deliberate Financial Governance exclusion} applies, with exact external recipients still unresolved.


{\small Recorded unresolved matters: canonicalName; exact scope remains in the canonical record.}\par
{\small Navigation classification: Native BusinessInteraction definition is determinate, but no governed capability context is represented for this interaction. No Application Interaction, integration flow or capability allocation is inferred.}\par
{\footnotesize \href{../../../../continuum/business/business-interaction/73b6d1ab8ecc48c4bdea98f013288c9c.md}{Canonical record: 73b6d1ab8ecc48c4bdea98f013288c9c}; full qualifications and provenance.}\par
\section*{III.5.51 Service Assurance Request}
\phantomsection
\label{navigation:III.5.51}
\addcontentsline{toc}{section}{III.5.51 Service Assurance Request}

\label{element:dd80af764f88480ea8b8b5898bb15ce7}
\index{Service Assurance Request (III.5.51)}
{\small BusinessInteraction · Canonical · current}\par

\begin{itemize}


\item \textbf{Canonical Name}: \texttt{\seqsplit{Service Assurance Request}}.

\item \textbf{Participants}: \texttt{\seqsplit{Service Guardian}} and \textbf{one authorised requesting Role} from \texttt{\seqsplit{Care Coordinator}}, \texttt{\seqsplit{Service Coordinator}}, \texttt{\seqsplit{Regulator}}, \texttt{\seqsplit{Policy Authority}}, \texttt{\seqsplit{System Steward}}.

\item \textbf{Purpose}: Establishes a request by an authorised Business Role for Service Assurance to be performed against an identifiable governed subject under an applicable assurance basis.

\item \textbf{Associated Business Service}: \href{\detokenize{../../../../../docs/markdown/03-business-architecture/behaviours/health-service-assurance.md#request-service-assurance}}{Request Service Assurance}.

\item \textbf{Boundary}: The requesting Role identifies the need within its responsibility/authority; it does not control assurance evidence assessment, adjudication or the resulting finding/conclusion, nor acquire authority to define/alter the Assurance Definition or Assurance Criteria. REQ-FND-005 independence remains controlling. Participant variation does not create five separate Interactions. Detailed request information and initiation mechanics remain unresolved.

\end{itemize}

{\small Recorded unresolved matters: canonicalName; exact scope remains in the canonical record.}\par
{\small Navigation classification: Native BusinessInteraction definition is determinate, but no governed capability context is represented for this interaction. No Application Interaction, integration flow or capability allocation is inferred.}\par
{\footnotesize \href{../../../../continuum/business/business-interaction/dd80af764f88480ea8b8b5898bb15ce7.md}{Canonical record: dd80af764f88480ea8b8b5898bb15ce7}; full qualifications and provenance.}\par
\section*{III.5.52 Service Assurance Status}
\phantomsection
\label{navigation:III.5.52}
\addcontentsline{toc}{section}{III.5.52 Service Assurance Status}

\label{element:6bad04918445465cb350c3252af73377}
\index{Service Assurance Status (III.5.52)}
{\small BusinessInteraction · Canonical · current}\par

\begin{itemize}


\item \textbf{Canonical Name}: \texttt{\seqsplit{Service Assurance Status}}.

\item \textbf{Participants}: \texttt{\seqsplit{System Steward}}, \texttt{\seqsplit{Service Guardian}}.

\item \textbf{Purpose}: Enables the System Steward to obtain the current business progression state of an identifiable Service Assurance activity for stewardship purposes.

\item \textbf{Associated Business Service}: \href{\detokenize{../../../../../docs/markdown/03-business-architecture/behaviours/health-service-assurance.md#request-service-assurance-status}}{Request Service Assurance Status}.

\item \textbf{Boundary}: \textbf{Service Assurance Status concerns progression of the assurance activity and SHALL NOT constitute, imply or expose an unadjudicated Assurance Finding or Conclusion.} It supplies no provisional/predicted finding or incomplete assessment presented as a conclusion. Only System Steward consumes the associated Service; requesting assurance does not grant another Role access to status. Status taxonomy remains unresolved.

\end{itemize}


This Interaction is stewardship / management observation of assurance activity progression, not recursive assurance. Management Monitoring observes activity state to progress/manage activity; assurance evaluates a governed subject against governing criteria to establish a conclusion. Obtaining status is not itself assurance. The \href{\detokenize{../../../../../docs/markdown/03-business-architecture/behaviours/health-service-assurance.md#assurance-non-recursion-boundary}}{non-recursion boundary} applies.


{\small Recorded unresolved matters: canonicalName; exact scope remains in the canonical record.}\par
{\small Navigation classification: Native BusinessInteraction definition is determinate, but no governed capability context is represented for this interaction. No Application Interaction, integration flow or capability allocation is inferred.}\par
{\footnotesize \href{../../../../continuum/business/business-interaction/6bad04918445465cb350c3252af73377.md}{Canonical record: 6bad04918445465cb350c3252af73377}; full qualifications and provenance.}\par
\section*{III.5.53 Service Coordination}
\phantomsection
\label{navigation:III.5.53}
\addcontentsline{toc}{section}{III.5.53 Service Coordination}

\label{element:c00bba65740947fea75542c13563553c}
\index{Service Coordination (III.5.53)}
{\small BusinessInteraction · Canonical · current}\par

\begin{itemize}


\item \textbf{Service Coordination}: Dialogue aimed at synchronising multiple care delivery activities across clinical specialties or care settings.

\end{itemize}

{\small Recorded unresolved matters: canonicalName; exact scope remains in the canonical record.}\par
{\small Navigation classification: Native BusinessInteraction definition is determinate, but no governed capability context is represented for this interaction. No Application Interaction, integration flow or capability allocation is inferred.}\par
{\footnotesize \href{../../../../continuum/business/business-interaction/c00bba65740947fea75542c13563553c.md}{Canonical record: c00bba65740947fea75542c13563553c}; full qualifications and provenance.}\par
\section*{III.5.54 Service Outcome}
\phantomsection
\label{navigation:III.5.54}
\addcontentsline{toc}{section}{III.5.54 Service Outcome}

\label{element:e52b7159e28441769fcf3214cc6be6a6}
\index{Service Outcome (III.5.54)}
{\small BusinessInteraction · Canonical · current}\par

\begin{itemize}


\item \textbf{Service Outcome}: The formal clinical finding, procedural report, or discharge disposition resulting from completed service delivery.

\end{itemize}

{\small Recorded unresolved matters: canonicalName; exact scope remains in the canonical record.}\par
{\small Navigation classification: Native BusinessInteraction definition is determinate, but no governed capability context is represented for this interaction. No Application Interaction, integration flow or capability allocation is inferred.}\par
{\footnotesize \href{../../../../continuum/business/business-interaction/e52b7159e28441769fcf3214cc6be6a6.md}{Canonical record: e52b7159e28441769fcf3214cc6be6a6}; full qualifications and provenance.}\par
\section*{III.5.55 Service Referral}
\phantomsection
\label{navigation:III.5.55}
\addcontentsline{toc}{section}{III.5.55 Service Referral}

\label{element:b7ea2864b6c044c1a9b7763ee6428d7c}
\index{Service Referral (III.5.55)}
{\small BusinessInteraction · Canonical · current}\par

\begin{itemize}


\item \textbf{Service Referral}: A formal request for a specialist, service provider, or clinical department to assess, manage, or assume care responsibility for a patient.

\end{itemize}

{\small Recorded unresolved matters: canonicalName; exact scope remains in the canonical record.}\par
{\small Navigation classification: Native BusinessInteraction definition is determinate, but no governed capability context is represented for this interaction. No Application Interaction, integration flow or capability allocation is inferred.}\par
{\footnotesize \href{../../../../continuum/business/business-interaction/b7ea2864b6c044c1a9b7763ee6428d7c.md}{Canonical record: b7ea2864b6c044c1a9b7763ee6428d7c}; full qualifications and provenance.}\par
\section*{III.5.56 Service Request}
\phantomsection
\label{navigation:III.5.56}
\addcontentsline{toc}{section}{III.5.56 Service Request}

\label{element:be3915edc56d4020969b7e1f4837f112}
\index{Service Request (III.5.56)}
{\small BusinessInteraction · Canonical · current}\par

\begin{itemize}


\item \textbf{Service Request}: A formal order placed for a specific diagnostic test, therapeutic procedure, or clinical service.

\end{itemize}

{\small Recorded unresolved matters: canonicalName; exact scope remains in the canonical record.}\par
{\small Navigation classification: Native BusinessInteraction definition is determinate, but no governed capability context is represented for this interaction. No Application Interaction, integration flow or capability allocation is inferred.}\par
{\footnotesize \href{../../../../continuum/business/business-interaction/be3915edc56d4020969b7e1f4837f112.md}{Canonical record: be3915edc56d4020969b7e1f4837f112}; full qualifications and provenance.}\par
\section*{III.5.57 Service Response}
\phantomsection
\label{navigation:III.5.57}
\addcontentsline{toc}{section}{III.5.57 Service Response}

\label{element:6bc0e6df5ba3405a9037e34f9810aefc}
\index{Service Response (III.5.57)}
{\small BusinessInteraction · Canonical · current}\par

\begin{itemize}


\item \textbf{Service Response}: A communication indicating the status, triage decision, scheduled timing, or acceptance disposition of a referral or request.

\end{itemize}

{\small Recorded unresolved matters: canonicalName; exact scope remains in the canonical record.}\par
{\small Navigation classification: Native BusinessInteraction definition is determinate, but no governed capability context is represented for this interaction. No Application Interaction, integration flow or capability allocation is inferred.}\par
{\footnotesize \href{../../../../continuum/business/business-interaction/6bc0e6df5ba3405a9037e34f9810aefc.md}{Canonical record: 6bc0e6df5ba3405a9037e34f9810aefc}; full qualifications and provenance.}\par
\section*{III.5.58 ServiceCatalogue Publish}
\phantomsection
\label{navigation:III.5.58}
\addcontentsline{toc}{section}{III.5.58 ServiceCatalogue Publish}

\label{element:397a4200bb894479a54a03dc84451efb}
\index{ServiceCatalogue Publish (III.5.58)}
{\small BusinessInteraction · Canonical · current}\par

\begin{itemize}


\item \textbf{ServiceCatalogue Publish}: Releasing updated service directory definitions to consumers.

\end{itemize}

{\small Recorded unresolved matters: canonicalName; exact scope remains in the canonical record.}\par
{\small Navigation classification: Native BusinessInteraction definition is determinate, but no governed capability context is represented for this interaction. No Application Interaction, integration flow or capability allocation is inferred.}\par
{\footnotesize \href{../../../../continuum/business/business-interaction/397a4200bb894479a54a03dc84451efb.md}{Canonical record: 397a4200bb894479a54a03dc84451efb}; full qualifications and provenance.}\par
\section*{III.5.59 ServiceCatalogue Request}
\phantomsection
\label{navigation:III.5.59}
\addcontentsline{toc}{section}{III.5.59 ServiceCatalogue Request}

\label{element:d6c56d2bb1274fc68cb0ed6c3b8a84b9}
\index{ServiceCatalogue Request (III.5.59)}
{\small BusinessInteraction · Canonical · current}\par

\begin{itemize}


\item \textbf{ServiceCatalogue Request}: Requesting published service catalogue specifications.

\end{itemize}

{\small Recorded unresolved matters: canonicalName; exact scope remains in the canonical record.}\par
{\small Navigation classification: Native BusinessInteraction definition is determinate, but no governed capability context is represented for this interaction. No Application Interaction, integration flow or capability allocation is inferred.}\par
{\footnotesize \href{../../../../continuum/business/business-interaction/d6c56d2bb1274fc68cb0ed6c3b8a84b9.md}{Canonical record: d6c56d2bb1274fc68cb0ed6c3b8a84b9}; full qualifications and provenance.}\par
\section*{III.5.60 ServiceCatalogue Update}
\phantomsection
\label{navigation:III.5.60}
\addcontentsline{toc}{section}{III.5.60 ServiceCatalogue Update}

\label{element:c15051a06c1a46b1acb75d18a5e728dd}
\index{ServiceCatalogue Update (III.5.60)}
{\small BusinessInteraction · Canonical · current}\par

\begin{itemize}


\item \textbf{ServiceCatalogue Update}: Modifying service definitions, operating hours, or delivery modalities.

\end{itemize}

{\small Recorded unresolved matters: canonicalName; exact scope remains in the canonical record.}\par
{\small Navigation classification: Native BusinessInteraction definition is determinate, but no governed capability context is represented for this interaction. No Application Interaction, integration flow or capability allocation is inferred.}\par
{\footnotesize \href{../../../../continuum/business/business-interaction/c15051a06c1a46b1acb75d18a5e728dd.md}{Canonical record: c15051a06c1a46b1acb75d18a5e728dd}; full qualifications and provenance.}\par
\section*{III.5.61 ServiceEligibility Assessment Report}
\phantomsection
\label{navigation:III.5.61}
\addcontentsline{toc}{section}{III.5.61 ServiceEligibility Assessment Report}

\label{element:db44999c87d545618969319d3f63725a}
\index{ServiceEligibility Assessment Report (III.5.61)}
{\small BusinessInteraction · Canonical · current}\par

\begin{itemize}


\item \textbf{ServiceEligibility Assessment Report}: The formal determination of patient eligibility for a service.

\end{itemize}

{\small Recorded unresolved matters: canonicalName; exact scope remains in the canonical record.}\par
{\small Navigation classification: Native BusinessInteraction definition is determinate, but no governed capability context is represented for this interaction. No Application Interaction, integration flow or capability allocation is inferred.}\par
{\footnotesize \href{../../../../continuum/business/business-interaction/db44999c87d545618969319d3f63725a.md}{Canonical record: db44999c87d545618969319d3f63725a}; full qualifications and provenance.}\par
\section*{III.5.62 ServiceEligibility Assessment Request}
\phantomsection
\label{navigation:III.5.62}
\addcontentsline{toc}{section}{III.5.62 ServiceEligibility Assessment Request}

\label{element:1508d1e33b264534bbcde367eb3dc6a3}
\index{ServiceEligibility Assessment Request (III.5.62)}
{\small BusinessInteraction · Canonical · current}\par

\begin{itemize}


\item \textbf{ServiceEligibility Assessment Request}: Evaluating whether a patient meets clinical or funding eligibility criteria.

\end{itemize}

{\small Recorded unresolved matters: canonicalName; exact scope remains in the canonical record.}\par
{\small Navigation classification: Native BusinessInteraction definition is determinate, but no governed capability context is represented for this interaction. No Application Interaction, integration flow or capability allocation is inferred.}\par
{\footnotesize \href{../../../../continuum/business/business-interaction/1508d1e33b264534bbcde367eb3dc6a3.md}{Canonical record: 1508d1e33b264534bbcde367eb3dc6a3}; full qualifications and provenance.}\par
\section*{III.5.63 ServiceOffering Query}
\phantomsection
\label{navigation:III.5.63}
\addcontentsline{toc}{section}{III.5.63 ServiceOffering Query}

\label{element:40105b22c930413692cc64826d3ad32f}
\index{ServiceOffering Query (III.5.63)}
{\small BusinessInteraction · Canonical · current}\par

\begin{itemize}


\item \textbf{ServiceOffering Query}: Discovering services offered by a facility or provider.

\end{itemize}

{\small Recorded unresolved matters: canonicalName; exact scope remains in the canonical record.}\par
{\small Navigation classification: Native BusinessInteraction definition is determinate, but no governed capability context is represented for this interaction. No Application Interaction, integration flow or capability allocation is inferred.}\par
{\footnotesize \href{../../../../continuum/business/business-interaction/40105b22c930413692cc64826d3ad32f.md}{Canonical record: 40105b22c930413692cc64826d3ad32f}; full qualifications and provenance.}\par
\section*{III.5.64 ServiceOffering Selection}
\phantomsection
\label{navigation:III.5.64}
\addcontentsline{toc}{section}{III.5.64 ServiceOffering Selection}

\label{element:dbb4668e39994da6bd7ec0251210d4bc}
\index{ServiceOffering Selection (III.5.64)}
{\small BusinessInteraction · Canonical · current}\par

\begin{itemize}


\item \textbf{ServiceOffering Selection}: Choosing an appropriate service offering for a referral or order.

\end{itemize}

{\small Recorded unresolved matters: canonicalName; exact scope remains in the canonical record.}\par
{\small Navigation classification: Native BusinessInteraction definition is determinate, but no governed capability context is represented for this interaction. No Application Interaction, integration flow or capability allocation is inferred.}\par
{\footnotesize \href{../../../../continuum/business/business-interaction/dbb4668e39994da6bd7ec0251210d4bc.md}{Canonical record: dbb4668e39994da6bd7ec0251210d4bc}; full qualifications and provenance.}\par
\section*{III.5.65 Transfer of Care}
\phantomsection
\label{navigation:III.5.65}
\addcontentsline{toc}{section}{III.5.65 Transfer of Care}

\label{element:8bb02a34807f4bb183bb42a722f4509b}
\index{Transfer of Care (III.5.65)}
{\small BusinessInteraction · Canonical · current}\par

\begin{itemize}


\item \textbf{Transfer of Care}: The formal transition of clinical and medico-legal accountability for a patient between care teams, wards, or organisations.

\end{itemize}

{\small Recorded unresolved matters: canonicalName; exact scope remains in the canonical record.}\par
{\small Navigation classification: Native BusinessInteraction definition is determinate, but no governed capability context is represented for this interaction. No Application Interaction, integration flow or capability allocation is inferred.}\par
{\footnotesize \href{../../../../continuum/business/business-interaction/8bb02a34807f4bb183bb42a722f4509b.md}{Canonical record: 8bb02a34807f4bb183bb42a722f4509b}; full qualifications and provenance.}\par
\section*{III.5.66 Work Acceptance}
\phantomsection
\label{navigation:III.5.66}
\addcontentsline{toc}{section}{III.5.66 Work Acceptance}

\label{element:64d7dfac2f5f4b61998df127b44c0521}
\index{Work Acceptance (III.5.66)}
{\small BusinessInteraction · Canonical · current}\par

\begin{itemize}


\item \textbf{Work Acceptance}: Acknowledgement and commitment by an operational worker to execute an assigned task.

\end{itemize}

{\small Recorded unresolved matters: canonicalName; exact scope remains in the canonical record.}\par
{\small Navigation classification: Native BusinessInteraction definition is determinate, but no governed capability context is represented for this interaction. No Application Interaction, integration flow or capability allocation is inferred.}\par
{\footnotesize \href{../../../../continuum/business/business-interaction/64d7dfac2f5f4b61998df127b44c0521.md}{Canonical record: 64d7dfac2f5f4b61998df127b44c0521}; full qualifications and provenance.}\par
\section*{III.5.67 Work Assignment}
\phantomsection
\label{navigation:III.5.67}
\addcontentsline{toc}{section}{III.5.67 Work Assignment}

\label{element:c5611769a52d44aebeef55bd48e1f433}
\index{Work Assignment (III.5.67)}
{\small BusinessInteraction · Canonical · current}\par

\begin{itemize}


\item \textbf{Work Assignment}: The formal dispatch or allocation of an operational task (e.g., patient transfer, cleaning, specimen delivery) to a staff member or role.

\end{itemize}

{\small Recorded unresolved matters: canonicalName; exact scope remains in the canonical record.}\par
{\small Navigation classification: Native BusinessInteraction definition is determinate, but no governed capability context is represented for this interaction. No Application Interaction, integration flow or capability allocation is inferred.}\par
{\footnotesize \href{../../../../continuum/business/business-interaction/c5611769a52d44aebeef55bd48e1f433.md}{Canonical record: c5611769a52d44aebeef55bd48e1f433}; full qualifications and provenance.}\par
\section*{III.5.68 Work Outcome}
\phantomsection
\label{navigation:III.5.68}
\addcontentsline{toc}{section}{III.5.68 Work Outcome}

\label{element:a950de4051d04cfa8bb180f8193f7495}
\index{Work Outcome (III.5.68)}
{\small BusinessInteraction · Canonical · current}\par

\begin{itemize}


\item \textbf{Work Outcome}: The formal completion report and outcome disposition of an operational task.

\end{itemize}

{\small Recorded unresolved matters: canonicalName; exact scope remains in the canonical record.}\par
{\small Navigation classification: Native BusinessInteraction definition is determinate, but no governed capability context is represented for this interaction. No Application Interaction, integration flow or capability allocation is inferred.}\par
{\footnotesize \href{../../../../continuum/business/business-interaction/a950de4051d04cfa8bb180f8193f7495.md}{Canonical record: a950de4051d04cfa8bb180f8193f7495}; full qualifications and provenance.}\par
\section*{III.5.69 Work Progress}
\phantomsection
\label{navigation:III.5.69}
\addcontentsline{toc}{section}{III.5.69 Work Progress}

\label{element:a687efed1f7e456bba53e356b157f0d8}
\index{Work Progress (III.5.69)}
{\small BusinessInteraction · Canonical · current}\par

\begin{itemize}


\item \textbf{Work Progress}: Milestone updates regarding the active execution status of an in-flight operational task.

\end{itemize}

{\small Recorded unresolved matters: canonicalName; exact scope remains in the canonical record.}\par
{\small Navigation classification: Native BusinessInteraction definition is determinate, but no governed capability context is represented for this interaction. No Application Interaction, integration flow or capability allocation is inferred.}\par
{\footnotesize \href{../../../../continuum/business/business-interaction/a687efed1f7e456bba53e356b157f0d8.md}{Canonical record: a687efed1f7e456bba53e356b157f0d8}; full qualifications and provenance.}\par

\part*{Part IV — Application}
\label{navigation:IV}
\addcontentsline{toc}{part}{Part IV — Application}
Begin with the distinction between managed active information and authoritative persistence. Mneme manages Context and Pragma and provides Context Load, Context Save, Adhoc Resource Access and Adhoc Search. Mnemosyne establishes authoritative durable state and persistence outcomes. Cache retention, completed execution and publication of outputs do not independently establish durable authority. The component definitions retain the candidate standing and detailed-contract reservations of the accepted baseline.

{\small Reading sources: \hyperref[element:1af8a56358f7429fabfb57be66542a99]{Mneme}; \hyperref[element:7985e647bf9e4ca3ba70faf9e3e4ae02]{Mnemosyne}.}\par
Execution, operational management and messaging are separate responsibilities. Ponos provides governed execution; Energeia owns Praxis and Journey and observes operational progression. Energeia's R1.x/R2.x boundary is Observe to Report, without autonomous changes to execution or workload. Petasos manages durable queues and PetasosParcel handoff; referenced information retains its own authority. Read these responsibilities together without turning their collaboration into an inferred component hierarchy.

{\small Reading sources: \hyperref[element:b521762c1ea444b48cd5f66be3752d30]{Ponos}; \hyperref[element:fb3f37a6ac614415ac486a6de9d35994]{Energeia}; \hyperref[element:195e5f10db8e419391354197b5d7a8cc]{Petasos}.}\par
At the interaction boundaries, Pylai supplies external interoperability while Iris supplies contextual human interaction. Themis owns authorisation semantics and consumes trusted authenticated Identity rather than authenticating it. Agora controls collaboration semantics and managed Room/Space intent; Omilia performs collaboration operations. These responsibilities do not confer ownership of information merely conveyed through them. Detailed representations and deferred mechanisms remain as qualified in their definitions.

{\small Reading sources: \hyperref[element:e3e4147fea3640198dc251ed41ee854d]{Pylai}; \hyperref[element:e927c1af96c64fcd979be659bee2cbe3]{Iris}; \hyperref[element:899d34ec491948b68aad09e507dea6d1]{Themis}; \hyperref[element:d1c4053e68384080b8eb79d2008632ba]{Agora}.}\par
Calliope-Models and Calliope-DataSets distinguish build-time representation from runtime content management. Models supplies structural definitions and library contracts. DataSets owns mutable governed content and uses the established information-management and persistence responsibilities. This distinction does not select a new deployment or storage allocation. The information map that follows remains a separate view of information concepts, not an inferred physical schema.

{\small Reading sources: \hyperref[element:307b3440a507403f8bc214d290bffb6a]{Calliope-Models}; \hyperref[element:53bf6815c7a54ed0abde9b1aad311843]{Calliope-DataSets}.}\par
\begin{center}\begin{tikzpicture}[>=Stealth,every node/.style={draw=Ink!40,rounded corners,align=left,font=\small,text width=5.2cm,inner sep=6pt}]
\node (source) at (0,0) {\hyperref[element:4af969b8c23a4dcab7569dc51b263383]{[Topic] Individual}};
\node (target0) at (8,4.25) {\hyperref[element:99dc695355384ea29ac3515b9178dfd0]{[Topic] Authority}};
\draw[->] (source.east) -- (3,0) |- (target0.west);
\node[draw=none,text width=1.8cm,font=\tiny,align=center,inner sep=0pt] at (4.15,4.63) {Composition\newline{}source-qualified / unclassified proposal};
\node (target1) at (8,2.55) {\hyperref[element:5c33b01e23964ef5935be79364e7c0f9]{[Topic] Contact}};
\draw[->] (source.east) -- (3,0) |- (target1.west);
\node[draw=none,text width=1.8cm,font=\tiny,align=center,inner sep=0pt] at (4.15,2.93) {Composition\newline{}source-qualified / unclassified proposal};
\node (target2) at (8,0.85) {\hyperref[element:95a13e957f5c471797869b901085266c]{[Topic] Demographics}};
\draw[->] (source.east) -- (3,0) |- (target2.west);
\node[draw=none,text width=1.8cm,font=\tiny,align=center,inner sep=0pt] at (4.15,1.23) {Composition\newline{}source-qualified / unclassified proposal};
\node (target3) at (8,-0.85) {\hyperref[element:3b72456b5686449f95ad26bfbbae07f3]{[Topic] Identity}};
\draw[->] (source.east) -- (3,0) |- (target3.west);
\node[draw=none,text width=1.8cm,font=\tiny,align=center,inner sep=0pt] at (4.15,-0.47) {Composition\newline{}source-qualified / unclassified proposal};
\node (target4) at (8,-2.55) {\hyperref[element:e92ea15616ea43c9a1d698e34ed1376b]{[Topic] Patient}};
\draw[->] (source.east) -- (3,0) |- (target4.west);
\node[draw=none,text width=1.8cm,font=\tiny,align=center,inner sep=0pt] at (4.15,-2.17) {Composition\newline{}source-qualified / unclassified proposal};
\node (target5) at (8,-4.25) {\hyperref[element:1163ffac87814f319578ef5cb95dcd4d]{[Topic] Person}};
\draw[->] (source.east) -- (3,0) |- (target5.west);
\node[draw=none,text width=1.8cm,font=\tiny,align=center,inner sep=0pt] at (4.15,-3.87) {Composition\newline{}source-qualified / unclassified proposal};
\end{tikzpicture}\end{center}
{\small Selected outgoing Composition assertions; arrows retain source-to-target direction. Labels retain assertion standing and record state; element standing accompanies each description. \href{structure.json}{Complete relationship qualifications and provenance} remain source-defined. No omitted edge or semantic effect is inferred.}\par
\chapter*{IV.1 Information Architecture}
\phantomsection
\label{navigation:IV.1}
\addcontentsline{toc}{chapter}{IV.1 Information Architecture}

\section*{IV.1.1 Topics}
\phantomsection
\label{navigation:IV.1.1}
\addcontentsline{toc}{section}{IV.1.1 Topics}

\subsection*{IV.1.1.1 [Topic] Care}
\phantomsection
\label{navigation:IV.1.1.1}
\addcontentsline{toc}{subsection}{IV.1.1.1 [Topic] Care}

\label{element:f702d357183847dd8f0fd09545679ffd}
\index{[Topic] Care (IV.1.1.1)}
{\small DataObject / Topic · Accepted bounded reconciliation · current}\par

\textbf{Description}: Information establishing the contexts within which healthcare occurs, including encounters, episodes, care-place associations, progression events and longitudinal relationships that provide continuity and context across care interactions.


{\footnotesize \href{../../../../continuum/application/data-object/topic/f702d357183847dd8f0fd09545679ffd.md}{Canonical record: f702d357183847dd8f0fd09545679ffd}; full qualifications and provenance.}\par
\subsubsection*{IV.1.1.1.1 [Topic] CareLocation}
\phantomsection
\label{navigation:IV.1.1.1.1}
\addcontentsline{toc}{subsubsection}{IV.1.1.1.1 [Topic] CareLocation}

\label{element:cfc241a7b6904bd19ccc1b9adbac0393}
\index{[Topic] CareLocation (IV.1.1.1.1)}
{\small DataObject / Topic · Accepted bounded reconciliation · current}\par

\textbf{Description}
Information describing the association between an Encounter and the place or places in which care occurs, including changes in location during the Encounter.


{\footnotesize \href{../../../../continuum/application/data-object/topic/cfc241a7b6904bd19ccc1b9adbac0393.md}{Canonical record: cfc241a7b6904bd19ccc1b9adbac0393}; full qualifications and provenance.}\par
\subsubsection*{IV.1.1.1.2 [Topic] Encounter}
\phantomsection
\label{navigation:IV.1.1.1.2}
\addcontentsline{toc}{subsubsection}{IV.1.1.1.2 [Topic] Encounter}

\label{element:fb464349f13b41c4ae9920a743fc6c12}
\index{[Topic] Encounter (IV.1.1.1.2)}
{\small DataObject / Topic · Explicitly APPROVED G2 decision only; G2 remains candidate/open and all residual qualifications remain unresolved. · current}\par

\textbf{Status: RESOLVED — APPROVED, with explicit unresolved further authority/evidence details.}



\textbf{Description}: Information describing a bounded interaction in which healthcare is delivered or healthcare-related activity occurs for an individual within a particular context.



\textbf{Definition}
An Encounter may provide the care-interaction context within which zero or more Healthcare Service Deliveries occur. A Healthcare Service Delivery may occur without an associated Encounter.



Encounter and HealthcareServiceDelivery are distinct Information Concepts. A single Encounter may contextualise multiple HealthcareServiceDelivery instances [0..*], and delivery may exist independently of Encounter context.



\begin{itemize}


\item Encounter ≠ HealthcareServiceDelivery. Encounter existence does not establish that any service was delivered.

\item The relationship is optional contextual association, not mandatory lifecycle progression, identity, containment or a component/state of Encounter.

\item The approved [0..*] multiplicity concerns deliveries contextualised by an Encounter; no further reverse cardinality is inferred.

\item Package1's DeliverableHealthcareService → optional Order → HealthcareServiceDelivery → ServiceOutcome semantics remain unchanged. Encounter is not made mandatory anywhere in that chain.

\item Further source/correlation qualification and universal delivery/outcome ownership are not established by this decision.

\end{itemize}


\textbf{Review consequence:} G2-Q05's Encounter/delivery existence, basic meaning and multiplicity questions are resolved. Further qualifications and the Referral portion remain unreviewed; no Referral Disposition review is commenced.



Source standing: Explicitly APPROVED G2 decision only; G2 remains candidate/open and all residual qualifications remain unresolved.


{\footnotesize \href{../../../../continuum/application/data-object/topic/fb464349f13b41c4ae9920a743fc6c12.md}{Canonical record: fb464349f13b41c4ae9920a743fc6c12}; full qualifications and provenance.}\par
\subsubsection*{IV.1.1.1.3 [Topic] EncounterContext}
\phantomsection
\label{navigation:IV.1.1.1.3}
\addcontentsline{toc}{subsubsection}{IV.1.1.1.3 [Topic] EncounterContext}

\label{element:a8be2b6474cd43cd9845b93258bd2fe3}
\index{[Topic] EncounterContext (IV.1.1.1.3)}
{\small DataObject / Topic · Accepted bounded reconciliation · current}\par

\textbf{Description}
Information bringing together the relevant subject, participants, services, places and other contextual information required to understand an Encounter.


{\footnotesize \href{../../../../continuum/application/data-object/topic/a8be2b6474cd43cd9845b93258bd2fe3.md}{Canonical record: a8be2b6474cd43cd9845b93258bd2fe3}; full qualifications and provenance.}\par
\subsubsection*{IV.1.1.1.4 [Topic] EncounterProgression}
\phantomsection
\label{navigation:IV.1.1.1.4}
\addcontentsline{toc}{subsubsection}{IV.1.1.1.4 [Topic] EncounterProgression}

\label{element:a3badcf73eaa4331b75d7b3eb6e67fc8}
\index{[Topic] EncounterProgression (IV.1.1.1.4)}
{\small DataObject / Topic · Accepted bounded reconciliation · current}\par

\textbf{Description}
Information describing changes in the administrative or operational progression of an Encounter over time, preserving the significant states and transitions associated with that interaction.


{\footnotesize \href{../../../../continuum/application/data-object/topic/a3badcf73eaa4331b75d7b3eb6e67fc8.md}{Canonical record: a3badcf73eaa4331b75d7b3eb6e67fc8}; full qualifications and provenance.}\par
\subsubsection*{IV.1.1.1.5 [Topic] Episode}
\phantomsection
\label{navigation:IV.1.1.1.5}
\addcontentsline{toc}{subsubsection}{IV.1.1.1.5 [Topic] Episode}

\label{element:3927d161556d48f4b2cc57bd46ab785d}
\index{[Topic] Episode (IV.1.1.1.5)}
{\small DataObject / Topic · Explicitly APPROVED G2 decision only; G2 remains candidate/open and all residual qualifications remain unresolved. · current}\par

\textbf{Status: RESOLVED — APPROVED, with explicit unresolved association/responsibility details.}



\textbf{Description}
Information describing a longitudinal context that relates healthcare activity over time for one or more health problems or conditions, providing continuity across relevant Encounters and other care interactions.



\textbf{Definition}
Episode is a longitudinal information context representing the treatment of one or more Problems or Conditions across a related series of Encounters, providing a coherent perspective for assurance, workflow and financial purposes.



Episode is now an independently meaningful Information Concept. The former non-authoritative Episode Context suggestion is promoted to Episode with this definition; its original rationale and standing remain historical in §7.



\begin{itemize}


\item Episode is not merely a Collection of Encounters. Encounter association contributes to longitudinal treatment context but does not define its complete meaning.

\item Episode relates one or more Problems/Conditions [1..\emph{] to a longitudinal course of treatment and may relate to multiple Encounters [0..}]. Not every Encounter for a healthcare subject necessarily belongs to a particular Episode.

\item Encounter-association mechanisms/rules, detailed identity, originating authority, Harmonia information-responsibility allocation and temporal/lifecycle semantics remain unresolved. The decision allocates no reverse cardinality, upstream owner, Feature or Service.

\item No Episode ownership of related Encounters is inferred; the association is not implementation containment. Assurance/workflow/financial perspective is not an allocation of authority for those activities.

\item This architecturally approved meaning stands independently of representation. No FHIR EpisodeOfCare derivation, mapping or equivalence is established.

\end{itemize}


\textbf{Review consequence:} G2-Q04 no longer asks whether Episode exists or what its basic meaning is; only genuinely unresolved details remain. No Problems/Conditions family is derived.



Source standing: Explicitly APPROVED G2 decision only; G2 remains candidate/open and all residual qualifications remain unresolved.


{\footnotesize \href{../../../../continuum/application/data-object/topic/3927d161556d48f4b2cc57bd46ab785d.md}{Canonical record: 3927d161556d48f4b2cc57bd46ab785d}; full qualifications and provenance.}\par
\subsection*{IV.1.1.2 [Topic] Coordination}
\phantomsection
\label{navigation:IV.1.1.2}
\addcontentsline{toc}{subsection}{IV.1.1.2 [Topic] Coordination}

\label{element:882234c291ca4c6492ad6dfbe1e23fea}
\index{[Topic] Coordination (IV.1.1.2)}
{\small DataObject / Topic · Accepted bounded reconciliation · current}\par

Information supporting the coordination of care between participants and providers, including referrals, their purpose and context, progression, supporting information and outcomes, without presuming that coordination necessarily transfers care responsibility.


{\footnotesize \href{../../../../continuum/application/data-object/topic/882234c291ca4c6492ad6dfbe1e23fea.md}{Canonical record: 882234c291ca4c6492ad6dfbe1e23fea}; full qualifications and provenance.}\par
\subsubsection*{IV.1.1.2.1 [Topic] Referral}
\phantomsection
\label{navigation:IV.1.1.2.1}
\addcontentsline{toc}{subsubsection}{IV.1.1.2.1 [Topic] Referral}

\label{element:d79d14c20bf14ecdab0434d4aa7db9b8}
\index{[Topic] Referral (IV.1.1.2.1)}
{\small DataObject / Topic · Accepted bounded reconciliation · current}\par

\textbf{Description}
Information describing a request for another healthcare participant or provider to assess, manage, provide care for, or assume responsibility concerning an individual.


{\footnotesize \href{../../../../continuum/application/data-object/topic/d79d14c20bf14ecdab0434d4aa7db9b8.md}{Canonical record: d79d14c20bf14ecdab0434d4aa7db9b8}; full qualifications and provenance.}\par
\subsubsection*{IV.1.1.2.2 [Topic] ReferralContext}
\phantomsection
\label{navigation:IV.1.1.2.2}
\addcontentsline{toc}{subsubsection}{IV.1.1.2.2 [Topic] ReferralContext}

\label{element:9eb5906e93ca4625bdb675a3397c4c8d}
\index{[Topic] ReferralContext (IV.1.1.2.2)}
{\small DataObject / Topic · Accepted bounded reconciliation · current}\par

\textbf{Description}
Information bringing together the subject, requester, recipient, purpose and other contextual information necessary to understand and act upon a Referral.


{\footnotesize \href{../../../../continuum/application/data-object/topic/9eb5906e93ca4625bdb675a3397c4c8d.md}{Canonical record: 9eb5906e93ca4625bdb675a3397c4c8d}; full qualifications and provenance.}\par
\subsubsection*{IV.1.1.2.3 [Topic] ReferralOutcome}
\phantomsection
\label{navigation:IV.1.1.2.3}
\addcontentsline{toc}{subsubsection}{IV.1.1.2.3 [Topic] ReferralOutcome}

\label{element:18ad2e58913b4a938e172df7b03f366a}
\index{[Topic] ReferralOutcome (IV.1.1.2.3)}
{\small DataObject / Topic · Explicitly APPROVED G2-D04; no promotion of Referral family/package or independent Referral Disposition. · current}\par

\textbf{Description}
Information describing the outcome or resolution of a Referral, including the resulting disposition and relevant information returned to participants involved in the referral.



\textbf{Status: RESOLVED — APPROVED. G2-Q01 is resolved; G2 remains open.}



Referral Disposition is \textbf{not retained as an independent Information Concept}. The assurance and governance requirement previously represented by that B candidate is instead represented by \textbf{Referral Outcome}.



\begin{quote}


\textbf{Referral Outcome is the governed information concerning what ultimately became of a Referral, sufficient to support assurance, governance and traceability over the Referral and its progression.}



\end{quote}


\textbf{Semantic classification:} Outcome; provenance-bearing information with Assertion characteristics where a source reports what became of the Referral. This is not ServiceOutcome equivalence or a downstream representation category.



Referral Outcome exists principally to support the ability to account for Referrals within the applicable governance/assurance scope. Where established and appropriate, it may include the resulting outcome, basis/reason, source/originator, authority, temporal context and relationships to relevant resulting activity or information. No canonical Referral Outcome value set is established.



\begin{quote}\small\ttfamily
Referral ≠ Referral Progression Event ≠ Referral Outcome\\
Referral Outcome ≠ Responsibility Assumption ≠ Transfer of Care ≠ HealthcareServiceDelivery ≠ ServiceOutcome\\
Referral information semantics ≠ Referral Administration workflow progression\\

\end{quote}


Referral Outcome may reference or relate to those Concepts where appropriate; its existence establishes none of them. Recording that a Referral was accepted, progressed or otherwise reached an outcome does not itself prove that care responsibility was actually assumed or transferred. G1 K1's intent/disposition/responsibility/Transfer of Care distinctions remain intact. Outcome is not equated with acceptance.



\textbf{Referral Progression Event:} Retained as historical information concerning significant occurrences in the administrative/business progression of a Referral. Not every Event represents a Referral Outcome, and the historical progression record is not collapsed into the final/current Outcome.



\textbf{Assurance / governance coverage — RF-IR5:} Harmonia must support assurance that Referrals within the applicable governance scope can be accounted for through sufficient progression and outcome information to establish what became of them. This requirement is applied to established Referral Administration / Manage Referral Progression / Referral Status \& Outcome Service responsibility without allocating a new owner, Feature, Service or Collection.



\begin{quote}\small\ttfamily
Referral population within governance scope → progression/outcome evidence → assurance coverage\\

\end{quote}


An inability to establish what became of a Referral is itself an assurance gap. No Referral Outcome is manufactured to make coverage appear complete. Where the outcome is not known, its absence is preserved explicitly.



\textbf{Workflow closure / finalisation:} This belongs to the progression of the Business Process / Referral Administration activity, not intrinsic Referral disposition information. The workflow may reach a state in which no further Referral Administration activity is required. This does not require a Referral Finalisation or Referral Closure Information Concept, a terminal state intrinsic to Referral, or Referral Outcome to double as workflow state. No workflow terminal states are defined and the Business Process state machine is not reproduced as the Information Concept lifecycle.



\textbf{Authority, temporal meaning and provenance:} Harmonia governs the outcome information for assurance and traceability; it does not infer substantive authority from recording or possessing it. Outcome may require provenance sufficient to establish what outcome was understood, source/originator, authority where relevant, when established/effective, and rationale/qualification where relevant. These are not all universally mandatory attributes. Detailed governance-scope criteria, evidence sufficiency, authority eligibility and contextual provenance/temporal qualification remain unestablished where existing architecture supplies no detail.



\textbf{Review history and consequence:} Q01 retains its original B question. The former Referral Disposition candidate proposed an independent Assertion for an attributable decision/reported disposition, with acceptance, decline, redirection and context-specific waitlisting examples. Architectural review rejects that separate Concept and reframes its assurance/governance purpose as Referral Outcome. Recipient-decision meanings and the canonical upstream Referral Disposition Log label are not erased. Q01 is resolved, without answering acceptance's clinical responsibility effects or the precise Referral–delivery/outcome relationships under Q05. Responsibility Assumption Assertion remains a non-authoritative D suggestion.



\textbf{Persistence boundary:} Only the Referral family, this review record and the necessary Referral navigation text are reconciled. Q02–Q07 and the complete G2-D01–D03 record remain unchanged. Frozen upstream architecture, Foundation, Package1 and initial historical report/diff are not modified. No other G2 question, G3, later Package2 family, FHIR mapping or downstream representation is commenced. G1 remains CLOSED; G2 is not CLOSED.



Source standing: Explicitly APPROVED G2-D04; no promotion of Referral family/package or independent Referral Disposition.


{\footnotesize \href{../../../../continuum/application/data-object/topic/18ad2e58913b4a938e172df7b03f366a.md}{Canonical record: 18ad2e58913b4a938e172df7b03f366a}; full qualifications and provenance.}\par
\subsubsection*{IV.1.1.2.4 [Topic] ReferralProgression}
\phantomsection
\label{navigation:IV.1.1.2.4}
\addcontentsline{toc}{subsubsection}{IV.1.1.2.4 [Topic] ReferralProgression}

\label{element:bc70a1055ca34eed8143856816587f67}
\index{[Topic] ReferralProgression (IV.1.1.2.4)}
{\small DataObject / Topic · G2-D04 explicitly retains significant historical progression information; historical describes represented occurrences, not historical architectural standing. · current}\par

\textbf{Description} 
Information describing the progression of a Referral over time, including significant changes in its state, handling or disposition.



\textbf{Status: RESOLVED — APPROVED. G2-Q01 is resolved; G2 remains open.}



Referral Disposition is \textbf{not retained as an independent Information Concept}. The assurance and governance requirement previously represented by that B candidate is instead represented by \textbf{Referral Outcome}.



\begin{quote}


\textbf{Referral Outcome is the governed information concerning what ultimately became of a Referral, sufficient to support assurance, governance and traceability over the Referral and its progression.}



\end{quote}


\textbf{Semantic classification:} Outcome; provenance-bearing information with Assertion characteristics where a source reports what became of the Referral. This is not ServiceOutcome equivalence or a downstream representation category.



Referral Outcome exists principally to support the ability to account for Referrals within the applicable governance/assurance scope. Where established and appropriate, it may include the resulting outcome, basis/reason, source/originator, authority, temporal context and relationships to relevant resulting activity or information. No canonical Referral Outcome value set is established.



\begin{quote}\small\ttfamily
Referral ≠ Referral Progression Event ≠ Referral Outcome\\
Referral Outcome ≠ Responsibility Assumption ≠ Transfer of Care ≠ HealthcareServiceDelivery ≠ ServiceOutcome\\
Referral information semantics ≠ Referral Administration workflow progression\\

\end{quote}


Referral Outcome may reference or relate to those Concepts where appropriate; its existence establishes none of them. Recording that a Referral was accepted, progressed or otherwise reached an outcome does not itself prove that care responsibility was actually assumed or transferred. G1 K1's intent/disposition/responsibility/Transfer of Care distinctions remain intact. Outcome is not equated with acceptance.



\textbf{Referral Progression Event:} Retained as historical information concerning significant occurrences in the administrative/business progression of a Referral. Not every Event represents a Referral Outcome, and the historical progression record is not collapsed into the final/current Outcome.



\textbf{Assurance / governance coverage — RF-IR5:} Harmonia must support assurance that Referrals within the applicable governance scope can be accounted for through sufficient progression and outcome information to establish what became of them. This requirement is applied to established Referral Administration / Manage Referral Progression / Referral Status \& Outcome Service responsibility without allocating a new owner, Feature, Service or Collection.



\begin{quote}\small\ttfamily
Referral population within governance scope → progression/outcome evidence → assurance coverage\\

\end{quote}


An inability to establish what became of a Referral is itself an assurance gap. No Referral Outcome is manufactured to make coverage appear complete. Where the outcome is not known, its absence is preserved explicitly.



\textbf{Workflow closure / finalisation:} This belongs to the progression of the Business Process / Referral Administration activity, not intrinsic Referral disposition information. The workflow may reach a state in which no further Referral Administration activity is required. This does not require a Referral Finalisation or Referral Closure Information Concept, a terminal state intrinsic to Referral, or Referral Outcome to double as workflow state. No workflow terminal states are defined and the Business Process state machine is not reproduced as the Information Concept lifecycle.



\textbf{Authority, temporal meaning and provenance:} Harmonia governs the outcome information for assurance and traceability; it does not infer substantive authority from recording or possessing it. Outcome may require provenance sufficient to establish what outcome was understood, source/originator, authority where relevant, when established/effective, and rationale/qualification where relevant. These are not all universally mandatory attributes. Detailed governance-scope criteria, evidence sufficiency, authority eligibility and contextual provenance/temporal qualification remain unestablished where existing architecture supplies no detail.



\textbf{Review history and consequence:} Q01 retains its original B question. The former Referral Disposition candidate proposed an independent Assertion for an attributable decision/reported disposition, with acceptance, decline, redirection and context-specific waitlisting examples. Architectural review rejects that separate Concept and reframes its assurance/governance purpose as Referral Outcome. Recipient-decision meanings and the canonical upstream Referral Disposition Log label are not erased. Q01 is resolved, without answering acceptance's clinical responsibility effects or the precise Referral–delivery/outcome relationships under Q05. Responsibility Assumption Assertion remains a non-authoritative D suggestion.



\textbf{Persistence boundary:} Only the Referral family, this review record and the necessary Referral navigation text are reconciled. Q02–Q07 and the complete G2-D01–D03 record remain unchanged. Frozen upstream architecture, Foundation, Package1 and initial historical report/diff are not modified. No other G2 question, G3, later Package2 family, FHIR mapping or downstream representation is commenced. G1 remains CLOSED; G2 is not CLOSED.



Source standing: G2-D04 explicitly retains significant historical progression information; historical describes represented occurrences, not historical architectural standing.


{\footnotesize \href{../../../../continuum/application/data-object/topic/bc70a1055ca34eed8143856816587f67.md}{Canonical record: bc70a1055ca34eed8143856816587f67}; full qualifications and provenance.}\par
\subsection*{IV.1.1.3 [Topic] Enterprise}
\phantomsection
\label{navigation:IV.1.1.3}
\addcontentsline{toc}{subsection}{IV.1.1.3 [Topic] Enterprise}

\label{element:5c5da6a06dac413aa2809a1397f8a3b9}
\index{[Topic] Enterprise (IV.1.1.3)}
{\small DataObject / Topic · Accepted bounded reconciliation · current}\par

Information concerning healthcare organisations and organisational units, including their identity, classification, structure, profile, relationships, and contact information.


{\footnotesize \href{../../../../continuum/application/data-object/topic/5c5da6a06dac413aa2809a1397f8a3b9.md}{Canonical record: 5c5da6a06dac413aa2809a1397f8a3b9}; full qualifications and provenance.}\par
\subsubsection*{IV.1.1.3.1 [Topic] Organisation}
\phantomsection
\label{navigation:IV.1.1.3.1}
\addcontentsline{toc}{subsubsection}{IV.1.1.3.1 [Topic] Organisation}

\label{element:dd3ba70128874caeb2d5da419d006c49}
\index{[Topic] Organisation (IV.1.1.3.1)}
{\small DataObject / Topic · Established Package1 conceptual information meaning; residual qualifications in selected section remain unchanged. · current}\par

\textbf{Description}
Information describing an organisation participating in healthcare, including its identity, characteristics and role within the healthcare environment.



\begin{itemize}


\item \textbf{Semantic Classification}: \texttt{\seqsplit{Entity}}

\item \textbf{Definition}: An identifiable legal, corporate, or administrative organisation operating within the healthcare domain (e.g. Health Service District, Hospital Network, Private Healthcare Provider, Diagnostic Laboratory Enterprise).

\item \textbf{Architectural Scope}: Represents the overarching enterprise entity possessing legal identity, operational mandate, and administrative responsibility.

\end{itemize}


Source standing: Established Package1 conceptual information meaning; residual qualifications in selected section remain unchanged.


{\footnotesize \href{../../../../continuum/application/data-object/topic/dd3ba70128874caeb2d5da419d006c49.md}{Canonical record: dd3ba70128874caeb2d5da419d006c49}; full qualifications and provenance.}\par
\subsubsection*{IV.1.1.3.2 [Topic] OrganisationalUnit}
\phantomsection
\label{navigation:IV.1.1.3.2}
\addcontentsline{toc}{subsubsection}{IV.1.1.3.2 [Topic] OrganisationalUnit}

\label{element:2c884cdd6a4b41a184818e3a7e540897}
\index{[Topic] OrganisationalUnit (IV.1.1.3.2)}
{\small DataObject / Topic · Established Package1 conceptual information meaning; residual qualifications in selected section remain unchanged. · current}\par

\textbf{Description}
Information describing a defined subdivision of an organisation and its position within the organisation's structural context.



\begin{itemize}


\item \textbf{Semantic Classification}: \texttt{\seqsplit{Entity / Administrative Subdivision}}

\item \textbf{Definition}: An identifiable internal division, clinical directorate, operational department, or specialty unit within a healthcare organisation (e.g. \emph{Division of Surgery}, \emph{Department of Cardiology}, \emph{Pharmacy Department}).

\item \textbf{Architectural Scope}: Represents structured administrative subdivisions that manage resources, staff, and services under the umbrella of the parent organisation.

\end{itemize}


Source standing: Established Package1 conceptual information meaning; residual qualifications in selected section remain unchanged.


{\footnotesize \href{../../../../continuum/application/data-object/topic/2c884cdd6a4b41a184818e3a7e540897.md}{Canonical record: 2c884cdd6a4b41a184818e3a7e540897}; full qualifications and provenance.}\par
\subsubsection*{IV.1.1.3.3 [Topic] OrganisationClassification}
\phantomsection
\label{navigation:IV.1.1.3.3}
\addcontentsline{toc}{subsubsection}{IV.1.1.3.3 [Topic] OrganisationClassification}

\label{element:227974cf26054bdea17038ecd6c0aa3a}
\index{[Topic] OrganisationClassification (IV.1.1.3.3)}
{\small DataObject / Topic · Established Package1 conceptual information meaning; residual qualifications in selected section remain unchanged. · current}\par

\textbf{Description}
Information describing the classifications applicable to an organisation, including its type, function or other governed categorisation.



\begin{itemize}


\item \textbf{Semantic Classification}: \texttt{\seqsplit{Assertion / Classification}}

\item \textbf{Definition}: The categorisation of the organisation according to its structural, operational, and clinical mandate (e.g. \emph{Tertiary Hospital}, \emph{Regional Health District}, \emph{Ambulatory Primary Care Clinic}, \emph{Specialist Pathology Service}).

\end{itemize}


Source standing: Established Package1 conceptual information meaning; residual qualifications in selected section remain unchanged.


{\footnotesize \href{../../../../continuum/application/data-object/topic/227974cf26054bdea17038ecd6c0aa3a.md}{Canonical record: 227974cf26054bdea17038ecd6c0aa3a}; full qualifications and provenance.}\par
\subsubsection*{IV.1.1.3.4 [Topic] OrganisationContact}
\phantomsection
\label{navigation:IV.1.1.3.4}
\addcontentsline{toc}{subsubsection}{IV.1.1.3.4 [Topic] OrganisationContact}

\label{element:cad820b906fa40f28824278fb18b7497}
\index{[Topic] OrganisationContact (IV.1.1.3.4)}
{\small DataObject / Topic · Established Package1 conceptual information meaning; residual qualifications in selected section remain unchanged. · current}\par

\textbf{Description}
Information describing how an organisation may be contacted or communicated with, including applicable addresses and communication channels.



\begin{itemize}


\item \textbf{Semantic Classification}: \texttt{\seqsplit{Assertion / Contact Bundle}}

\item \textbf{Definition}: The official administrative communication details, registered physical and postal addresses, corporate telephone/electronic points of contact, and nominated official representatives.

\end{itemize}


Source standing: Established Package1 conceptual information meaning; residual qualifications in selected section remain unchanged.


{\footnotesize \href{../../../../continuum/application/data-object/topic/cad820b906fa40f28824278fb18b7497.md}{Canonical record: cad820b906fa40f28824278fb18b7497}; full qualifications and provenance.}\par
\subsubsection*{IV.1.1.3.5 [Topic] OrganisationIdentity}
\phantomsection
\label{navigation:IV.1.1.3.5}
\addcontentsline{toc}{subsubsection}{IV.1.1.3.5 [Topic] OrganisationIdentity}

\label{element:c240be16de76415682237a238b3d3fc8}
\index{[Topic] OrganisationIdentity (IV.1.1.3.5)}
{\small DataObject / Topic · Established Package1 conceptual information meaning; residual qualifications in selected section remain unchanged. · current}\par

\textbf{Description}
Information used to identify, distinguish and correlate an organisation across authorities, systems and healthcare contexts.



\begin{itemize}


\item \textbf{Semantic Classification}: \texttt{\seqsplit{Assertion / Identity Bundle}}

\item \textbf{Definition}: The governed set of official enterprise identifiers, registry bindings, and external directory identifiers establishing the identity of the organisation.

\item \textbf{Key Semantic Decomposition}:

\item \texttt{\seqsplit{Identifier Value}}: The token or registration number identifying the organisation.

\item \texttt{\seqsplit{Identifier Authority / Namespace}}: The statutory authority, national healthcare directory, or enterprise namespace issuing the identifier.

\end{itemize}


Source standing: Established Package1 conceptual information meaning; residual qualifications in selected section remain unchanged.


{\footnotesize \href{../../../../continuum/application/data-object/topic/c240be16de76415682237a238b3d3fc8.md}{Canonical record: c240be16de76415682237a238b3d3fc8}; full qualifications and provenance.}\par
\subsubsection*{IV.1.1.3.6 [Topic] OrganisationProfile}
\phantomsection
\label{navigation:IV.1.1.3.6}
\addcontentsline{toc}{subsubsection}{IV.1.1.3.6 [Topic] OrganisationProfile}

\label{element:04adbfc449bb4b9e9ae6e7fdf8a66bf9}
\index{[Topic] OrganisationProfile (IV.1.1.3.6)}
{\small DataObject / Topic · Established Package1 conceptual information meaning; residual qualifications in selected section remain unchanged. · current}\par

\textbf{Description}
Information describing the characteristics of an organisation relevant to its participation in healthcare, including descriptive and operational attributes.



\begin{itemize}


\item \textbf{Semantic Classification}: \texttt{\seqsplit{Assertion / Profile Bundle}}

\item \textbf{Definition}: The verified set of administrative and corporate attributes describing the organisation, including registered legal name, public trading names, corporate registration numbers, and formal accreditation status.

\end{itemize}


Source standing: Established Package1 conceptual information meaning; residual qualifications in selected section remain unchanged.


{\footnotesize \href{../../../../continuum/application/data-object/topic/04adbfc449bb4b9e9ae6e7fdf8a66bf9.md}{Canonical record: 04adbfc449bb4b9e9ae6e7fdf8a66bf9}; full qualifications and provenance.}\par
\subsection*{IV.1.1.4 [Topic] Equipment}
\phantomsection
\label{navigation:IV.1.1.4}
\addcontentsline{toc}{subsection}{IV.1.1.4 [Topic] Equipment}

\label{element:8e658834311d41e68a8ecf71d8ccd5db}
\index{[Topic] Equipment (IV.1.1.4)}
{\small DataObject / Topic · Accepted bounded reconciliation · current}\par

\textbf{Description}
Information concerning clinical devices and equipment, including device definitions, physical instances, identity, operational associations, location, subject and service relationships, and communication endpoints.


{\footnotesize \href{../../../../continuum/application/data-object/topic/8e658834311d41e68a8ecf71d8ccd5db.md}{Canonical record: 8e658834311d41e68a8ecf71d8ccd5db}; full qualifications and provenance.}\par
\subsubsection*{IV.1.1.4.1 [Topic] Device}
\phantomsection
\label{navigation:IV.1.1.4.1}
\addcontentsline{toc}{subsubsection}{IV.1.1.4.1 [Topic] Device}

\label{element:25c1b751956446cea17065b6254419d5}
\index{[Topic] Device (IV.1.1.4.1)}
{\small DataObject / Topic · Established Package1 conceptual information meaning; residual qualifications in selected section remain unchanged. · current}\par

\textbf{Description}
Information describing an identifiable physical device or item of equipment, including its identity, characteristics and relationship to its defining device type.



\begin{itemize}


\item \textbf{Semantic Classification}: \texttt{\seqsplit{Entity}}

\item \textbf{Definition}: An identifiable, physical hardware asset manufactured as an instance of a \texttt{\seqsplit{Device Definition}}.

\item \textbf{Key Characteristics}: Device Identifier (such as physical serial number or asset identification), manufacturing date, device version / configuration assertion, and asset operational status (with schemes such as Unique Device Identifier [UDI] and enterprise asset tags serving as illustrative examples).

\end{itemize}


Source standing: Established Package1 conceptual information meaning; residual qualifications in selected section remain unchanged.


{\footnotesize \href{../../../../continuum/application/data-object/topic/25c1b751956446cea17065b6254419d5.md}{Canonical record: 25c1b751956446cea17065b6254419d5}; full qualifications and provenance.}\par
\subsubsection*{IV.1.1.4.2 [Topic] DeviceAssociation}
\phantomsection
\label{navigation:IV.1.1.4.2}
\addcontentsline{toc}{subsubsection}{IV.1.1.4.2 [Topic] DeviceAssociation}

\label{element:303f8586953d450d8d82d3d5f76e488a}
\index{[Topic] DeviceAssociation (IV.1.1.4.2)}
{\small DataObject / Topic · Established Package1 conceptual information meaning; residual qualifications in selected section remain unchanged. · current}\par

\textbf{Description}
Information describing the contextual association of a device with a subject, location, service or other entity, including the period for which that association applies.



\begin{itemize}


\item \textbf{Semantic Classification}: \texttt{\seqsplit{Contextual / Temporal Relationship}}

\item \textbf{Definition}: A dynamic, time-bounded association binding a \texttt{\seqsplit{Device Instance}} to a patient, physical care-place, or clinical service.

\item \textbf{Key Characteristics}: Association type, effective start and end timestamps, commissioning authority, and clinical context.

\end{itemize}


Source standing: Established Package1 conceptual information meaning; residual qualifications in selected section remain unchanged.


{\footnotesize \href{../../../../continuum/application/data-object/topic/303f8586953d450d8d82d3d5f76e488a.md}{Canonical record: 303f8586953d450d8d82d3d5f76e488a}; full qualifications and provenance.}\par
\subsubsection*{IV.1.1.4.3 [Topic] DeviceDefinition}
\phantomsection
\label{navigation:IV.1.1.4.3}
\addcontentsline{toc}{subsubsection}{IV.1.1.4.3 [Topic] DeviceDefinition}

\label{element:2a879d06b17143be8bbfc35cdaab8bdf}
\index{[Topic] DeviceDefinition (IV.1.1.4.3)}
{\small DataObject / Topic · Established Package1 conceptual information meaning; residual qualifications in selected section remain unchanged. · current}\par

\textbf{Description}
Information describing a type or model of device, including the characteristics and specifications shared by device instances of that definition.



\begin{itemize}


\item \textbf{Semantic Classification}: \texttt{\seqsplit{Definition}}

\item \textbf{Definition}: The governed type, manufacturer catalogue model, regulatory classification, and technical specification of a medical device (e.g. \emph{Infusion Pump Model XYZ}, \emph{12-Lead ECG Monitor Model ABC}).

\item \textbf{Key Characteristics}: Manufacturer name, model identifier, commercial name, regulatory classification (e.g. medical device risk class), calibration requirements, and operational capabilities.

\end{itemize}


Source standing: Established Package1 conceptual information meaning; residual qualifications in selected section remain unchanged.


{\footnotesize \href{../../../../continuum/application/data-object/topic/2a879d06b17143be8bbfc35cdaab8bdf.md}{Canonical record: 2a879d06b17143be8bbfc35cdaab8bdf}; full qualifications and provenance.}\par
\subsection*{IV.1.1.5 [Topic] Individual}
\phantomsection
\label{navigation:IV.1.1.5}
\addcontentsline{toc}{subsection}{IV.1.1.5 [Topic] Individual}

\label{element:4af969b8c23a4dcab7569dc51b263383}
\index{[Topic] Individual (IV.1.1.5)}
{\small DataObject / Topic · Accepted bounded reconciliation · current}\par

\textbf{Description}: Information concerning an individual, their identity, identifiers, demographic characteristics, relationships, and their contextual participation as a subject of healthcare.


{\footnotesize \href{../../../../continuum/application/data-object/topic/4af969b8c23a4dcab7569dc51b263383.md}{Canonical record: 4af969b8c23a4dcab7569dc51b263383}; full qualifications and provenance.}\par
\subsubsection*{IV.1.1.5.1 [Topic] Authority}
\phantomsection
\label{navigation:IV.1.1.5.1}
\addcontentsline{toc}{subsubsection}{IV.1.1.5.1 [Topic] Authority}

\label{element:99dc695355384ea29ac3515b9178dfd0}
\index{[Topic] Authority (IV.1.1.5.1)}
{\small DataObject / Topic · Accepted bounded reconciliation · current}\par

\textbf{Description}: Information describing another party's authority to represent, consent for, make decisions for, or otherwise act on behalf of an individual.


{\footnotesize \href{../../../../continuum/application/data-object/topic/99dc695355384ea29ac3515b9178dfd0.md}{Canonical record: 99dc695355384ea29ac3515b9178dfd0}; full qualifications and provenance.}\par
\subsubsection*{IV.1.1.5.2 [Topic] Contact}
\phantomsection
\label{navigation:IV.1.1.5.2}
\addcontentsline{toc}{subsubsection}{IV.1.1.5.2 [Topic] Contact}

\label{element:5c33b01e23964ef5935be79364e7c0f9}
\index{[Topic] Contact (IV.1.1.5.2)}
{\small DataObject / Topic · Accepted bounded reconciliation · current}\par

\textbf{Description}: Information describing how an individual may be contacted or communicated with, including applicable addresses and communication channels.


{\footnotesize \href{../../../../continuum/application/data-object/topic/5c33b01e23964ef5935be79364e7c0f9.md}{Canonical record: 5c33b01e23964ef5935be79364e7c0f9}; full qualifications and provenance.}\par
\subsubsection*{IV.1.1.5.3 [Topic] Demographics}
\phantomsection
\label{navigation:IV.1.1.5.3}
\addcontentsline{toc}{subsubsection}{IV.1.1.5.3 [Topic] Demographics}

\label{element:95a13e957f5c471797869b901085266c}
\index{[Topic] Demographics (IV.1.1.5.3)}
{\small DataObject / Topic · Accepted bounded reconciliation · current}\par

\textbf{Description}: Information describing the demographic characteristics of an individual relevant to their identification, administration, healthcare or communication.


{\footnotesize \href{../../../../continuum/application/data-object/topic/95a13e957f5c471797869b901085266c.md}{Canonical record: 95a13e957f5c471797869b901085266c}; full qualifications and provenance.}\par
\subsubsection*{IV.1.1.5.4 [Topic] Identity}
\phantomsection
\label{navigation:IV.1.1.5.4}
\addcontentsline{toc}{subsubsection}{IV.1.1.5.4 [Topic] Identity}

\label{element:3b72456b5686449f95ad26bfbbae07f3}
\index{[Topic] Identity (IV.1.1.5.4)}
{\small DataObject / Topic · Established Package1 conceptual information meaning; residual qualifications in selected section remain unchanged. · current}\par

\textbf{Definition}: The governed set of identity traits, identifiers, names, and demographic assertions that establish the recognised identity of a \texttt{\seqsplit{Person}}.



\textbf{Description}: Information used to establish, distinguish and correlate an individual across identities, identifiers, aliases, sources and healthcare contexts.



\textbf{Semantic Classification}: \texttt{\seqsplit{Assertion / Identity Bundle}}



\textbf{Architectural Scope}: Governed by the \texttt{\seqsplit{Person Identity}} capability. Acts as the governed identity anchor linking identifiers, aliases, and correlation graphs.



Source standing: Established Package1 conceptual information meaning; residual qualifications in selected section remain unchanged.


{\footnotesize \href{../../../../continuum/application/data-object/topic/3b72456b5686449f95ad26bfbbae07f3.md}{Canonical record: 3b72456b5686449f95ad26bfbbae07f3}; full qualifications and provenance.}\par
\subsubsection*{IV.1.1.5.5 [Topic] Patient}
\phantomsection
\label{navigation:IV.1.1.5.5}
\addcontentsline{toc}{subsubsection}{IV.1.1.5.5 [Topic] Patient}

\label{element:e92ea15616ea43c9a1d698e34ed1376b}
\index{[Topic] Patient (IV.1.1.5.5)}
{\small DataObject / Topic · Established Package1 conceptual information meaning; residual qualifications in selected section remain unchanged. · current}\par

\textbf{Definition} 
The subject-oriented context in which health information, clinical activity, or care delivery concerns a \texttt{\seqsplit{Person}}.



\textbf{Description} 
Information describing an individual in their context as a recipient or subject of healthcare, including the healthcare-specific identity and administrative context required to support care.



\textbf{Semantic Classification}: \texttt{\seqsplit{Contextual Binding}}



\textbf{Key Characteristics}: Demographic history, language/interpreter requirements, cultural and spiritual preferences, communication preferences, and vital status assertions (e.g., verified birth date, death notification).



\textbf{Architectural Scope}: Captures the individual's profile as a consumer of health services without assuming active patient status in a specific facility.



Source standing: Established Package1 conceptual information meaning; residual qualifications in selected section remain unchanged.


{\footnotesize \href{../../../../continuum/application/data-object/topic/e92ea15616ea43c9a1d698e34ed1376b.md}{Canonical record: e92ea15616ea43c9a1d698e34ed1376b}; full qualifications and provenance.}\par
\subsubsection*{IV.1.1.5.6 [Topic] Person}
\phantomsection
\label{navigation:IV.1.1.5.6}
\addcontentsline{toc}{subsubsection}{IV.1.1.5.6 [Topic] Person}

\label{element:1163ffac87814f319578ef5cb95dcd4d}
\index{[Topic] Person (IV.1.1.5.6)}
{\small DataObject / Topic · Established Package1 conceptual information meaning; residual qualifications in selected section remain unchanged. · current}\par

\textbf{Definition}: The identifiable biological and legal human entity.



\textbf{Description}: Information describing an individual as an enduring human entity, independent of any particular healthcare interaction, role or episode of care.



\textbf{Semantic Classification}: \texttt{\seqsplit{Entity}}



\textbf{Architectural Scope}: Represents the enduring human being independently of any specific clinical encounter, healthcare facility, or operational role. A \texttt{\seqsplit{Person}} exists prior to and outside of any interaction with the healthcare system.



Source standing: Established Package1 conceptual information meaning; residual qualifications in selected section remain unchanged.


{\footnotesize \href{../../../../continuum/application/data-object/topic/1163ffac87814f319578ef5cb95dcd4d.md}{Canonical record: 1163ffac87814f319578ef5cb95dcd4d}; full qualifications and provenance.}\par
\subsubsection*{IV.1.1.5.7 [Topic] Relationship}
\phantomsection
\label{navigation:IV.1.1.5.7}
\addcontentsline{toc}{subsubsection}{IV.1.1.5.7 [Topic] Relationship}

\label{element:c0ec3c75d9c54b3097e12c240ee816fe}
\index{[Topic] Relationship (IV.1.1.5.7)}
{\small DataObject / Topic · Accepted bounded reconciliation · current}\par

\textbf{Description}: Information describing relevant personal, family, carer, representative or other relationships between an individual and other people.


{\footnotesize \href{../../../../continuum/application/data-object/topic/c0ec3c75d9c54b3097e12c240ee816fe.md}{Canonical record: c0ec3c75d9c54b3097e12c240ee816fe}; full qualifications and provenance.}\par
\subsection*{IV.1.1.6 [Topic] Ordering}
\phantomsection
\label{navigation:IV.1.1.6}
\addcontentsline{toc}{subsection}{IV.1.1.6 [Topic] Ordering}

\label{element:1c4e2d5ce4574856a7d3605838e8afcb}
\index{[Topic] Ordering (IV.1.1.6)}
{\small DataObject / Topic · Accepted bounded reconciliation · current}\par

\textbf{Description}
Information concerning formal requests for services or activities, including their intent, routing, progression, modification, contextual information and correlation with resulting information or fulfilment.


{\footnotesize \href{../../../../continuum/application/data-object/topic/1c4e2d5ce4574856a7d3605838e8afcb.md}{Canonical record: 1c4e2d5ce4574856a7d3605838e8afcb}; full qualifications and provenance.}\par
\subsubsection*{IV.1.1.6.1 [Topic] Order}
\phantomsection
\label{navigation:IV.1.1.6.1}
\addcontentsline{toc}{subsubsection}{IV.1.1.6.1 [Topic] Order}

\label{element:2a791e0e1f6b45c0a57fb8233f73faa5}
\index{[Topic] Order (IV.1.1.6.1)}
{\small DataObject / Topic · Established Package1 conceptual information meaning; residual qualifications in selected section remain unchanged. · current}\par

\textbf{Description}
Information describing a formal request for a service, activity or other fulfilment to be performed, including its subject, requester, intent and requested outcome.



Harmonia enforces an essential architectural rule regarding clinical requisitions and service delivery:



\$\$\textbackslash{}text\{Order\} \textbackslash{}neq \textbackslash{}text\{Healthcare Service\}\$\$



\begin{quote}\small\ttfamily
graph TD\\
    ORDER["Order<br/>(Authoritative Requisition / Directive)"]\\
    DELIV["DeliverableHealthcareService<br/>(Bound Service Context)"]\\
    DELIVERY["HealthcareServiceDelivery<br/>(Information Referring to Service Execution)"]\\
\\
    DELIV -.->|"may be requested / directed via"| ORDER\\
    ORDER -.->|"requests / directs"| DELIVERY\\
    DELIV -->|"fulfilled by"| DELIVERY\\

\end{quote}


\par\medskip\textbf{Key Semantics \& Invariants}\par


\begin{enumerate}


\item \textbf{Associated Request Mechanism}: An \texttt{\seqsplit{Order}} is an authoritative clinical requisition, directive, or prescription requesting the delivery of a service. It is \textbf{NOT} a mandatory lifecycle stage of \texttt{\seqsplit{Healthcare Service}}.

\item \textbf{Delivery Without an Order}: Service delivery may legitimately occur without an \texttt{\seqsplit{Order}} where business or clinical semantics permit (e.g. \emph{Emergency Resuscitation}, \emph{Triage Assessment}, \emph{Direct Routine Nursing Care}, \emph{Scheduled Facility Maintenance}).

\item \textbf{Independent Lifecycles}: An \texttt{\seqsplit{Order}} possesses its own lifecycle states (\emph{Requested}, \emph{Accepted}, \emph{In-Progress}, \emph{Completed}, \emph{Discontinued}), owned by \texttt{\seqsplit{Order Administration}}, distinct from the lifecycle of the service delivery itself.

\end{enumerate}


Source standing: Established Package1 conceptual information meaning; residual qualifications in selected section remain unchanged.


{\footnotesize \href{../../../../continuum/application/data-object/topic/2a791e0e1f6b45c0a57fb8233f73faa5.md}{Canonical record: 2a791e0e1f6b45c0a57fb8233f73faa5}; full qualifications and provenance.}\par
\subsubsection*{IV.1.1.6.2 [Topic] OrderChange}
\phantomsection
\label{navigation:IV.1.1.6.2}
\addcontentsline{toc}{subsubsection}{IV.1.1.6.2 [Topic] OrderChange}

\label{element:73abd9430cfa4395ac38b41acfb169ec}
\index{[Topic] OrderChange (IV.1.1.6.2)}
{\small DataObject / Topic · Accepted bounded reconciliation · current}\par

Information describing a requested modification to an existing Order, including the nature, context and disposition of the requested change.


{\footnotesize \href{../../../../continuum/application/data-object/topic/73abd9430cfa4395ac38b41acfb169ec.md}{Canonical record: 73abd9430cfa4395ac38b41acfb169ec}; full qualifications and provenance.}\par
\subsubsection*{IV.1.1.6.3 [Topic] OrderContext}
\phantomsection
\label{navigation:IV.1.1.6.3}
\addcontentsline{toc}{subsubsection}{IV.1.1.6.3 [Topic] OrderContext}

\label{element:614c48fd853643578788b738f8017b53}
\index{[Topic] OrderContext (IV.1.1.6.3)}
{\small DataObject / Topic · Accepted bounded reconciliation · current}\par

Information bringing together the subject, requester, requested activity, intended performer and other context necessary to understand and act upon an Order.


{\footnotesize \href{../../../../continuum/application/data-object/topic/614c48fd853643578788b738f8017b53.md}{Canonical record: 614c48fd853643578788b738f8017b53}; full qualifications and provenance.}\par
\subsubsection*{IV.1.1.6.4 [Topic] OrderProgressionEvent}
\phantomsection
\label{navigation:IV.1.1.6.4}
\addcontentsline{toc}{subsubsection}{IV.1.1.6.4 [Topic] OrderProgressionEvent}

\label{element:8fb3a767d7a04d638ef76f98d61496b1}
\index{[Topic] OrderProgressionEvent (IV.1.1.6.4)}
{\small DataObject / Topic · Accepted bounded reconciliation · current}\par

\textbf{Description}
Information describing significant changes in the state and handling of an Order as it progresses from request through fulfilment or other resolution.


{\footnotesize \href{../../../../continuum/application/data-object/topic/8fb3a767d7a04d638ef76f98d61496b1.md}{Canonical record: 8fb3a767d7a04d638ef76f98d61496b1}; full qualifications and provenance.}\par
\subsubsection*{IV.1.1.6.5 [Topic] OrderResult}
\phantomsection
\label{navigation:IV.1.1.6.5}
\addcontentsline{toc}{subsubsection}{IV.1.1.6.5 [Topic] OrderResult}

\label{element:24187967f53446479c87d7676c66ba12}
\index{[Topic] OrderResult (IV.1.1.6.5)}
{\small DataObject / Topic · Accepted bounded reconciliation · current}\par

\textbf{Description}
Information establishing the relationship between an Order and information produced as a consequence of its fulfilment.


{\footnotesize \href{../../../../continuum/application/data-object/topic/24187967f53446479c87d7676c66ba12.md}{Canonical record: 24187967f53446479c87d7676c66ba12}; full qualifications and provenance.}\par
\subsubsection*{IV.1.1.6.6 [Topic] OrderRouting}
\phantomsection
\label{navigation:IV.1.1.6.6}
\addcontentsline{toc}{subsubsection}{IV.1.1.6.6 [Topic] OrderRouting}

\label{element:06bde79727884b27b3599e373b3cee5f}
\index{[Topic] OrderRouting (IV.1.1.6.6)}
{\small DataObject / Topic · Accepted bounded reconciliation · current}\par

\textbf{Description}
Information describing where and how an Order is to be directed for acceptance, processing or fulfilment.


{\footnotesize \href{../../../../continuum/application/data-object/topic/06bde79727884b27b3599e373b3cee5f.md}{Canonical record: 06bde79727884b27b3599e373b3cee5f}; full qualifications and provenance.}\par
\subsection*{IV.1.1.7 [Topic] Place}
\phantomsection
\label{navigation:IV.1.1.7}
\addcontentsline{toc}{subsection}{IV.1.1.7 [Topic] Place}

\label{element:36539e3eee7041e585aa3c5a06424bcc}
\index{[Topic] Place (IV.1.1.7)}
{\small DataObject / Topic · Accepted bounded reconciliation · current}\par

\textbf{Description}
Information concerning physical, geospatial, functional and care locations, including their identity, classification, spatial relationships, containment, characteristics, and care-place definitions.


{\footnotesize \href{../../../../continuum/application/data-object/topic/36539e3eee7041e585aa3c5a06424bcc.md}{Canonical record: 36539e3eee7041e585aa3c5a06424bcc}; full qualifications and provenance.}\par
\subsubsection*{IV.1.1.7.1 [Topic] CarePlace}
\phantomsection
\label{navigation:IV.1.1.7.1}
\addcontentsline{toc}{subsubsection}{IV.1.1.7.1 [Topic] CarePlace}

\label{element:779f2295b0b648dda9dbd21c4ade3ba0}
\index{[Topic] CarePlace (IV.1.1.7.1)}
{\small DataObject / Topic · Established Package1 conceptual information meaning; residual qualifications in selected section remain unchanged. · current}\par

\textbf{Description}
Information describing the characteristics and intended healthcare function of a place in which care may be provided, independently of its transient operational state.



\begin{itemize}


\item \textbf{Semantic Classification}: \texttt{\seqsplit{Definition}}

\item \textbf{Definition}: The static engineering, physical infrastructure, and architectural specification of a bed bay or clinical treatment point.

\item \textbf{Key Characteristics}: Physical capacity, isolation classification (e.g. \emph{Negative Pressure}, \emph{Positive Pressure}, \emph{Standard}), medical gas outlet provisions (e.g. \emph{Oxygen}, \emph{Suction}, \emph{Medical Air}), emergency power backup, and fixed telemetry wiring.

\item \textbf{Architectural Boundary}: Captures what the physical care-place \textbf{is} built to support, independent of its transient operational occupancy.

\end{itemize}


Source standing: Established Package1 conceptual information meaning; residual qualifications in selected section remain unchanged.


{\footnotesize \href{../../../../continuum/application/data-object/topic/779f2295b0b648dda9dbd21c4ade3ba0.md}{Canonical record: 779f2295b0b648dda9dbd21c4ade3ba0}; full qualifications and provenance.}\par
\subsubsection*{IV.1.1.7.2 [Topic] Geolocation}
\phantomsection
\label{navigation:IV.1.1.7.2}
\addcontentsline{toc}{subsubsection}{IV.1.1.7.2 [Topic] Geolocation}

\label{element:771b2e814946401bbe6dd6b4fdd3fb42}
\index{[Topic] Geolocation (IV.1.1.7.2)}
{\small DataObject / Topic · Established Package1 conceptual information meaning; residual qualifications in selected section remain unchanged. · current}\par

\textbf{Description}
Information describing the geographical position or address of a location, supporting its physical identification and location.



\begin{itemize}


\item \textbf{Semantic Classification}: \texttt{\seqsplit{Assertion / Spatial Coordinate}}

\item \textbf{Definition}: The physical address, geospatial coordinates (latitude/longitude), building elevation, and GIS geofence defining the physical location on Earth.

\end{itemize}


Source standing: Established Package1 conceptual information meaning; residual qualifications in selected section remain unchanged.


{\footnotesize \href{../../../../continuum/application/data-object/topic/771b2e814946401bbe6dd6b4fdd3fb42.md}{Canonical record: 771b2e814946401bbe6dd6b4fdd3fb42}; full qualifications and provenance.}\par
\subsubsection*{IV.1.1.7.3 [Topic] Location}
\phantomsection
\label{navigation:IV.1.1.7.3}
\addcontentsline{toc}{subsubsection}{IV.1.1.7.3 [Topic] Location}

\label{element:9ee4df4d76864bb1ac90fe1077c56c6f}
\index{[Topic] Location (IV.1.1.7.3)}
{\small DataObject / Topic · Established Package1 conceptual information meaning; residual qualifications in selected section remain unchanged. · current}\par

\textbf{Description}
Information describing an identifiable physical, geographical or functional location relevant to healthcare delivery or administration.



\begin{itemize}


\item \textbf{Semantic Classification}: \texttt{\seqsplit{Entity}}

\item \textbf{Definition}: An identifiable physical, geospatial, or functional place where healthcare is delivered, managed, or supported.

\item \textbf{Architectural Scope}: Represents physical sites (e.g. \emph{Main Hospital Campus}), buildings, wards, rooms, and specific bed bays, as well as mobile care units (e.g. \emph{Mobile Outreach Clinic}).

\end{itemize}


Source standing: Established Package1 conceptual information meaning; residual qualifications in selected section remain unchanged.


{\footnotesize \href{../../../../continuum/application/data-object/topic/9ee4df4d76864bb1ac90fe1077c56c6f.md}{Canonical record: 9ee4df4d76864bb1ac90fe1077c56c6f}; full qualifications and provenance.}\par
\subsubsection*{IV.1.1.7.4 [Topic] LocationClassification}
\phantomsection
\label{navigation:IV.1.1.7.4}
\addcontentsline{toc}{subsubsection}{IV.1.1.7.4 [Topic] LocationClassification}

\label{element:3896c910c28646098f0707efe88169ee}
\index{[Topic] LocationClassification (IV.1.1.7.4)}
{\small DataObject / Topic · Established Package1 conceptual information meaning; residual qualifications in selected section remain unchanged. · current}\par

\textbf{Description}
Information describing the type, function or other governed classification applicable to a location.



\begin{itemize}


\item \textbf{Semantic Classification}: \texttt{\seqsplit{Assertion / Classification}}

\item \textbf{Definition}: The structural and functional categorisation of the location (e.g. \emph{Campus}, \emph{Building}, \emph{Floor}, \emph{Ward / Wing}, \emph{Patient Room}, \emph{Care-Place / Bed Bay}, \emph{Operating Theatre}, \emph{Consulting Suite}, \emph{Mobile Unit}).

\end{itemize}


Source standing: Established Package1 conceptual information meaning; residual qualifications in selected section remain unchanged.


{\footnotesize \href{../../../../continuum/application/data-object/topic/3896c910c28646098f0707efe88169ee.md}{Canonical record: 3896c910c28646098f0707efe88169ee}; full qualifications and provenance.}\par
\subsubsection*{IV.1.1.7.5 [Topic] LocationIdentity}
\phantomsection
\label{navigation:IV.1.1.7.5}
\addcontentsline{toc}{subsubsection}{IV.1.1.7.5 [Topic] LocationIdentity}

\label{element:5f8299e31ee5410583007473c75d578f}
\index{[Topic] LocationIdentity (IV.1.1.7.5)}
{\small DataObject / Topic · Established Package1 conceptual information meaning; residual qualifications in selected section remain unchanged. · current}\par

\textbf{Description}
Information used to identify, distinguish and correlate a location across organisations, authorities, systems and healthcare contexts.



\begin{itemize}


\item \textbf{Semantic Classification}: \texttt{\seqsplit{Assertion / Identity Bundle}}

\item \textbf{Definition}: The governed set of enterprise and local identifiers that uniquely distinguish the location within an enterprise or regional jurisdiction.

\item \textbf{Key Semantic Decomposition}:

\item \texttt{\seqsplit{Identifier Value}}: The location code or name (e.g. \emph{CAMPUS-NORTH}, \emph{WARD-3B}, \emph{BED-304-A}).

\item \texttt{\seqsplit{Identifier Authority / Namespace}}: The health enterprise, government department, or facilities management namespace.

\end{itemize}


Source standing: Established Package1 conceptual information meaning; residual qualifications in selected section remain unchanged.


{\footnotesize \href{../../../../continuum/application/data-object/topic/5f8299e31ee5410583007473c75d578f.md}{Canonical record: 5f8299e31ee5410583007473c75d578f}; full qualifications and provenance.}\par
\subsection*{IV.1.1.8 [Topic] Scheduling}
\phantomsection
\label{navigation:IV.1.1.8}
\addcontentsline{toc}{subsection}{IV.1.1.8 [Topic] Scheduling}

\label{element:1225bfab0a6c4312875091a027f24275}
\index{[Topic] Scheduling (IV.1.1.8)}
{\small DataObject / Topic · Accepted bounded reconciliation · current}\par

\textbf{Description}
Information concerning the scheduling and administrative context of healthcare interactions, including appointments, scheduling changes, synchronisation state, attendance-related information and consolidated schedule perspectives.


{\footnotesize \href{../../../../continuum/application/data-object/topic/1225bfab0a6c4312875091a027f24275.md}{Canonical record: 1225bfab0a6c4312875091a027f24275}; full qualifications and provenance.}\par
\subsubsection*{IV.1.1.8.1 [Topic] Appointment}
\phantomsection
\label{navigation:IV.1.1.8.1}
\addcontentsline{toc}{subsubsection}{IV.1.1.8.1 [Topic] Appointment}

\label{element:9badf0d360b64215bc60e74eec954018}
\index{[Topic] Appointment (IV.1.1.8.1)}
{\small DataObject / Topic · Explicitly APPROVED G2 decision only; G2 remains candidate/open and all residual qualifications remain unresolved. · current}\par

""Description**
Information describing a planned healthcare interaction, including its subject, participants, intended service, time, place and scheduling context.



\textbf{Status: RESOLVED — APPROVED, with explicit unresolved cardinality/qualification details.}



\begin{quote}


\textbf{An Appointment may provide the prospective scheduling context for an Encounter. Where an Encounter occurs without a pre-existing Appointment, such as an emergency or walk-in interaction, an Appointment may be created retrospectively to represent the scheduling/administrative context associated with that Encounter.}



\end{quote}


Appointment and Encounter are distinct Information Concepts. Their relationship is principally scheduling context, not mandatory lifecycle progression.



\begin{quote}\small\ttfamily
Appointment ≠ Encounter\\
Appointment existence does not imply Encounter occurred\\
Encounter does not prove prospective Appointment existed\\
Retrospective Appointment establishment is not evidence of prospective scheduling\\

\end{quote}


\begin{itemize}


\item No Appointment → Encounter cardinality is established. Cancellation, rescheduling, recurring appointments, group appointments and other scheduling contexts remain unreviewed for multiplicity.

\item Prospective and retrospective context do not create separate Appointment Concepts. Scheduling-origin/context qualification is identified as a property/qualification for later review, without a defined representation or mandatory property structure.

\item The decision does not assign the originator/authority for retrospective creation or transfer external booking authority to Harmonia.

\end{itemize}


\textbf{Review consequence:} G2-Q02 reflects approved prospective/retrospective meaning and retains only unreviewed qualification/classification/authority/cardinality details. G2-Q03 preserves actual care-interaction meaning while retaining administrative tracking/Assembly questions.



Source standing: Explicitly APPROVED G2 decision only; G2 remains candidate/open and all residual qualifications remain unresolved.


{\footnotesize \href{../../../../continuum/application/data-object/topic/9badf0d360b64215bc60e74eec954018.md}{Canonical record: 9badf0d360b64215bc60e74eec954018}; full qualifications and provenance.}\par
\subsubsection*{IV.1.1.8.2 [Topic] PatientSchedule}
\phantomsection
\label{navigation:IV.1.1.8.2}
\addcontentsline{toc}{subsubsection}{IV.1.1.8.2 [Topic] PatientSchedule}

\label{element:c109f726f6ae4ee1b7fde1c012cfae74}
\index{[Topic] PatientSchedule (IV.1.1.8.2)}
{\small DataObject / Topic · Accepted bounded reconciliation · current}\par

\textbf{Schedule}
Information presenting the scheduled healthcare interactions relevant to an individual, providing a consolidated patient-oriented perspective across applicable appointments and scheduling sources.


{\footnotesize \href{../../../../continuum/application/data-object/topic/c109f726f6ae4ee1b7fde1c012cfae74.md}{Canonical record: c109f726f6ae4ee1b7fde1c012cfae74}; full qualifications and provenance.}\par
\subsubsection*{IV.1.1.8.3 [Topic] ScheduleChange}
\phantomsection
\label{navigation:IV.1.1.8.3}
\addcontentsline{toc}{subsubsection}{IV.1.1.8.3 [Topic] ScheduleChange}

\label{element:0e98b52f18274a81ad3e3b9ff9cca747}
\index{[Topic] ScheduleChange (IV.1.1.8.3)}
{\small DataObject / Topic · Accepted bounded reconciliation · current}\par

\textbf{Description}
Information describing a significant change to scheduled healthcare activity, including creation, rescheduling, cancellation or other alteration affecting the planned interaction.


{\footnotesize \href{../../../../continuum/application/data-object/topic/0e98b52f18274a81ad3e3b9ff9cca747.md}{Canonical record: 0e98b52f18274a81ad3e3b9ff9cca747}; full qualifications and provenance.}\par
\subsubsection*{IV.1.1.8.4 [Topic] ScheduleState}
\phantomsection
\label{navigation:IV.1.1.8.4}
\addcontentsline{toc}{subsubsection}{IV.1.1.8.4 [Topic] ScheduleState}

\label{element:6c4e9a91baf7460387d298366c5c6d0a}
\index{[Topic] ScheduleState (IV.1.1.8.4)}
{\small DataObject / Topic · Accepted bounded reconciliation · current}\par

\textbf{Description}
Information describing the synchronisation and current known state of scheduling information across participating sources or contexts.


{\footnotesize \href{../../../../continuum/application/data-object/topic/6c4e9a91baf7460387d298366c5c6d0a.md}{Canonical record: 6c4e9a91baf7460387d298366c5c6d0a}; full qualifications and provenance.}\par
\subsection*{IV.1.1.9 [Topic] Service}
\phantomsection
\label{navigation:IV.1.1.9}
\addcontentsline{toc}{subsection}{IV.1.1.9 [Topic] Service}

\label{element:0b990e0c64684bb8a45643e48eb4470a}
\index{[Topic] Service (IV.1.1.9)}
{\small DataObject / Topic · Accepted bounded reconciliation · current}\par

\textbf{Description}
Information describing healthcare services from definition and contextualised availability through delivery, outcome and assurance, including the relationships between services, providers, locations and the contexts in which services are provided.


{\footnotesize \href{../../../../continuum/application/data-object/topic/0b990e0c64684bb8a45643e48eb4470a.md}{Canonical record: 0b990e0c64684bb8a45643e48eb4470a}; full qualifications and provenance.}\par
\subsubsection*{IV.1.1.9.1 [Topic] ServiceAssurance}
\phantomsection
\label{navigation:IV.1.1.9.1}
\addcontentsline{toc}{subsubsection}{IV.1.1.9.1 [Topic] ServiceAssurance}

\label{element:07a1f83c9b4e462eaef1595c2a50d560}
\index{[Topic] ServiceAssurance (IV.1.1.9.1)}
{\small DataObject / Topic · Established Package1 conceptual information meaning; residual qualifications in selected section remain unchanged. · current}\par

\textbf{Description}
Information describing the assurance context associated with a healthcare service, including the evidence or determination that its provision satisfies applicable expectations, requirements or obligations.



\begin{itemize}


\item \textbf{Semantic Classification}: \texttt{\seqsplit{Accountability / Audit Record}}

\item \textbf{Definition}: The formal retrospective quality evaluation, clinical audit, accreditation review, or funding/billing reconciliation assessed against delivered services and outcomes.

\item \textbf{Architectural Scope \& Ownership Assessment}: Evaluates single service deliveries or aggregated cohorts across an episode of care. Because Harmonia's core operational capabilities do not claim R1 architectural ownership for statutory health insurance billing reconciliation or external accreditation governance, \texttt{\seqsplit{AssuredHealthcareService}} is designated as a \textbf{Reference / Contextual Concept} where Harmonia consumes or supports assurance interfaces without owning the external statutory authority.

\end{itemize}


Source standing: Established Package1 conceptual information meaning; residual qualifications in selected section remain unchanged.


{\footnotesize \href{../../../../continuum/application/data-object/topic/07a1f83c9b4e462eaef1595c2a50d560.md}{Canonical record: 07a1f83c9b4e462eaef1595c2a50d560}; full qualifications and provenance.}\par
\subsubsection*{IV.1.1.9.2 [Topic] ServiceDelivery}
\phantomsection
\label{navigation:IV.1.1.9.2}
\addcontentsline{toc}{subsubsection}{IV.1.1.9.2 [Topic] ServiceDelivery}

\label{element:4627ee33845744f2ae2e10c84e313755}
\index{[Topic] ServiceDelivery (IV.1.1.9.2)}
{\small DataObject / Topic · Established Package1 conceptual information meaning; residual qualifications in selected section remain unchanged. · current}\par

\textbf{Description}
Information describing the actual provision of a healthcare service within a particular subject, provider, place and delivery context.



\begin{itemize}


\item \textbf{Semantic Classification}: \texttt{\seqsplit{Activity / Fulfilment}}

\item \textbf{Definition}: An Information Concept referring to the real-world execution or performance of a healthcare service for a specific target subject (or facility target) at a concrete point in time. It represents information about the activity rather than becoming the activity or implying that Harmonia performs it.

\item \textbf{Key Characteristics}: Actual start/end timestamps, performing practitioners, delivery location, executing devices/equipment, and execution status progression.

\item \textbf{Healthcare Subject Invariant}: While clinical service deliveries typically associate with a \texttt{\seqsplit{Healthcare Subject Context}}, broader service semantics (e.g., \emph{Facility Sanitisation}, \emph{Water Quality Testing}, \emph{Equipment Calibration}) do not universally require a Healthcare Subject.

\end{itemize}


Source standing: Established Package1 conceptual information meaning; residual qualifications in selected section remain unchanged.


{\footnotesize \href{../../../../continuum/application/data-object/topic/4627ee33845744f2ae2e10c84e313755.md}{Canonical record: 4627ee33845744f2ae2e10c84e313755}; full qualifications and provenance.}\par
\subsubsection*{IV.1.1.9.3 [Topic] ServiceOffering}
\phantomsection
\label{navigation:IV.1.1.9.3}
\addcontentsline{toc}{subsubsection}{IV.1.1.9.3 [Topic] ServiceOffering}

\label{element:441b0891822f4add834ea930c3fe58b7}
\index{[Topic] ServiceOffering (IV.1.1.9.3)}
{\small DataObject / Topic · Established Package1 conceptual information meaning; residual qualifications in selected section remain unchanged. · current}\par

\textbf{Description}
Information describing a healthcare service that an organisation or provider makes available, including its nature, characteristics and conditions of availability.



\begin{itemize}


\item \textbf{Semantic Classification}: \texttt{\seqsplit{Definition}}

\item \textbf{Definition}: The canonical service offering defined within a healthcare enterprise or jurisdictional catalogue (e.g. \emph{12-Lead Electrocardiography}, \emph{Acute Inpatient Haemodialysis}, \emph{Wound Debridement Protocol}, \emph{Ward Sanitisation Protocol}).

\item \textbf{Key Characteristics}: Service classification, specialty taxonomy, clinical indications, contraindications, required workforce competencies, supported delivery modalities, and permitted composition structures.

\item \textbf{Recursive Composition}: An \texttt{\seqsplit{OfferedHealthcareService}} may contain other \texttt{\seqsplit{OfferedHealthcareService}} definitions (e.g. \emph{Cardiology Service} contains \emph{Echocardiography} and \emph{Stress ECG}).

\end{itemize}


Source standing: Established Package1 conceptual information meaning; residual qualifications in selected section remain unchanged.


{\footnotesize \href{../../../../continuum/application/data-object/topic/441b0891822f4add834ea930c3fe58b7.md}{Canonical record: 441b0891822f4add834ea930c3fe58b7}; full qualifications and provenance.}\par
\subsubsection*{IV.1.1.9.4 [Topic] ServiceOutcome}
\phantomsection
\label{navigation:IV.1.1.9.4}
\addcontentsline{toc}{subsubsection}{IV.1.1.9.4 [Topic] ServiceOutcome}

\label{element:4fb589629e7f405db44d6a87a2b5e36b}
\index{[Topic] ServiceOutcome (IV.1.1.9.4)}
{\small DataObject / Topic · Established Package1 conceptual information meaning; residual qualifications in selected section remain unchanged. · current}\par

\textbf{Description}
Information describing the outcome resulting from the delivery of a healthcare service, including the relevant disposition or consequences of that delivery.



\begin{itemize}


\item \textbf{Semantic Classification}: \texttt{\seqsplit{Outcome}}

\item \textbf{Definition}: An Information Concept concerning a result, finding, change, consequence or administrative artefact associated with a service delivery. These information kinds remain distinct; representing them does not confer responsibility for their production, clinical interpretation or originating authority.

\item \textbf{Key Scope}: Spans across:

\item \textbf{Clinical Outcomes}: Diagnoses, physiological improvements, symptom resolutions;

\item \textbf{Diagnostic Outcomes}: Pathology findings, radiology reports, ECG interpretations;

\item \textbf{Therapeutic Outcomes}: Medication administration confirmations, surgical revisions;

\item \textbf{Operational Outcomes}: Sanitised bed bay certification, repaired asset state;

\item \textbf{Administrative Outcomes}: Completed intake assessments, eligibility confirmations.

\item \textbf{Preserved uncertainty}: Exact information-kind semantics and a universal Delivery/Outcome ownership model are not established. No single representation model, lifecycle or performing Capability is inferred.

\end{itemize}


Source standing: Established Package1 conceptual information meaning; residual qualifications in selected section remain unchanged.


{\footnotesize \href{../../../../continuum/application/data-object/topic/4fb589629e7f405db44d6a87a2b5e36b.md}{Canonical record: 4fb589629e7f405db44d6a87a2b5e36b}; full qualifications and provenance.}\par
\subsubsection*{IV.1.1.9.5 [Topic] ServiceProvision}
\phantomsection
\label{navigation:IV.1.1.9.5}
\addcontentsline{toc}{subsubsection}{IV.1.1.9.5 [Topic] ServiceProvision}

\label{element:8d41e9fb6a164deb93afb5fc4c671e10}
\index{[Topic] ServiceProvision (IV.1.1.9.5)}
{\small DataObject / Topic · Established Package1 conceptual information meaning; residual qualifications in selected section remain unchanged. · current}\par

\textbf{Description}
Information describing an offered service contextualised for actual provision, including the provider, location, applicable conditions and other bindings necessary for the service to be delivered.



\begin{itemize}


\item \textbf{Semantic Classification}: \texttt{\seqsplit{Contextual Binding}}

\item \textbf{Definition}: The contextual binding of an \texttt{\seqsplit{OfferedHealthcareService}} to a specific providing organisation, delivery location, participating practitioner roles, operational availability schedules, and eligibility rule sets.

\item \textbf{Key Characteristics}: Bound service definition, performing organisation, delivery site, participating practitioner role bindings, intake schedule, and catchment criteria.

\end{itemize}


Source standing: Established Package1 conceptual information meaning; residual qualifications in selected section remain unchanged.


{\footnotesize \href{../../../../continuum/application/data-object/topic/8d41e9fb6a164deb93afb5fc4c671e10.md}{Canonical record: 8d41e9fb6a164deb93afb5fc4c671e10}; full qualifications and provenance.}\par
\subsection*{IV.1.1.10 [Topic] TaskCompletion}
\phantomsection
\label{navigation:IV.1.1.10}
\addcontentsline{toc}{subsection}{IV.1.1.10 [Topic] TaskCompletion}

\label{element:8a9351cd170940d0b36f8d56584f15ec}
\index{[Topic] TaskCompletion (IV.1.1.10)}
{\small DataObject / Topic · Accepted bounded reconciliation · current}\par

\textbf{Description} 
Information describing actionable work from reusable definition and instantiated task through undertaking and governed resolution, including execution participation, completion information, outcomes and resulting outputs.


{\footnotesize \href{../../../../continuum/application/data-object/topic/8a9351cd170940d0b36f8d56584f15ec.md}{Canonical record: 8a9351cd170940d0b36f8d56584f15ec}; full qualifications and provenance.}\par
\subsubsection*{IV.1.1.10.1 [Topic] ExecutionConcurrency}
\phantomsection
\label{navigation:IV.1.1.10.1}
\addcontentsline{toc}{subsubsection}{IV.1.1.10.1 [Topic] ExecutionConcurrency}

\label{element:0781dd44748d42338c18bf01993af656}
\index{[Topic] ExecutionConcurrency (IV.1.1.10.1)}
{\small DataObject / Topic · Accepted bounded reconciliation · current}\par

\textbf{Description}
Information governing or describing concurrency applicable to the execution of work associated with a Task.


{\footnotesize \href{../../../../continuum/application/data-object/topic/0781dd44748d42338c18bf01993af656.md}{Canonical record: 0781dd44748d42338c18bf01993af656}; full qualifications and provenance.}\par
\subsubsection*{IV.1.1.10.2 [Topic] OutcomeConcurrency}
\phantomsection
\label{navigation:IV.1.1.10.2}
\addcontentsline{toc}{subsubsection}{IV.1.1.10.2 [Topic] OutcomeConcurrency}

\label{element:72c505782347448d8f89d74ef76359af}
\index{[Topic] OutcomeConcurrency (IV.1.1.10.2)}
{\small DataObject / Topic · Accepted bounded reconciliation · current}\par

\textbf{Description}
Information describing the concurrency rules applicable when determining or establishing the outcome of a Task.


{\footnotesize \href{../../../../continuum/application/data-object/topic/72c505782347448d8f89d74ef76359af.md}{Canonical record: 72c505782347448d8f89d74ef76359af}; full qualifications and provenance.}\par
\subsubsection*{IV.1.1.10.3 [Topic] OutputRules}
\phantomsection
\label{navigation:IV.1.1.10.3}
\addcontentsline{toc}{subsubsection}{IV.1.1.10.3 [Topic] OutputRules}

\label{element:a035fe1913db4d2589a2e4edde8db737}
\index{[Topic] OutputRules (IV.1.1.10.3)}
{\small DataObject / Topic · Accepted bounded reconciliation · current}\par

\textbf{Description}
Information describing the rules governing the outputs established as part of Task resolution.


{\footnotesize \href{../../../../continuum/application/data-object/topic/a035fe1913db4d2589a2e4edde8db737.md}{Canonical record: a035fe1913db4d2589a2e4edde8db737}; full qualifications and provenance.}\par
\subsubsection*{IV.1.1.10.4 [Topic] Task}
\phantomsection
\label{navigation:IV.1.1.10.4}
\addcontentsline{toc}{subsubsection}{IV.1.1.10.4 [Topic] Task}

\label{element:c43139289cb649b4836521a96a51604c}
\index{[Topic] Task (IV.1.1.10.4)}
{\small DataObject / Topic · Bounded adjudicated conceptual Task meaning; section 5 reservations, including runtime TaskOutcome identity/lifecycle, remain unresolved. · current}\par

\textbf{Description}
Information describing an identifiable instance of work to be completed, including its purpose, context, state and applicable completion expectations.



For example, the archetype \textbf{"Climb a ladder"} describes a kind of work;
\textbf{"ladder-climb-0023"} identifies a particular \texttt{\seqsplit{ActionableTask}} of that kind.
These are conceptual examples, not an identifier scheme.



An \texttt{\seqsplit{ActionableTask}} SHALL NOT mean the reusable definition or the execution /
undertaking of the work. Its contextual work meaning remains distinct from
information about particular undertakings.



An \texttt{\seqsplit{ActionableTask}} explicitly governs execution concurrency, which
undertakings participate in resolving its outcome, and how their outputs
establish its \texttt{\seqsplit{Task.Output(s)}}. The following three controls are independent
semantic properties, not an implementation structure or execution mechanism.



\par\medskip\textbf{2.2.1 ActionableTask.ExecutionConcurrency}\par


\begin{quote}\small\ttfamily
ActionableTask.ExecutionConcurrency = [1..x]\\

\end{quote}


\texttt{\seqsplit{ExecutionConcurrency}} establishes how many \texttt{\seqsplit{FulfillmentTask}} undertakings of
the \texttt{\seqsplit{ActionableTask}} may execute concurrently. It concerns execution
concurrency only; it does not determine which undertakings establish the
\texttt{\seqsplit{TaskOutcome}}. This permitted concurrency range does not require an
undertaking to exist or execute: an ActionableTask may still have zero
undertakings. No universal value for \texttt{\seqsplit{x}} or concurrency execution mechanism
is established here.



\par\medskip\textbf{2.2.2 ActionableTask.OutcomeConcurrency.Mode}\par


\begin{quote}\small\ttfamily
ActionableTask.OutcomeConcurrency.Mode =\\
    FirstToFinish | AssignedToFinish | Aggregate\\

\end{quote}


This property determines \textbf{which FulfillmentTask undertaking or undertakings
participate in resolving the ActionableTask outcome}.



\begin{longtable}{>{\raggedright\arraybackslash}p{\dimexpr\linewidth/2-2\tabcolsep\relax}>{\raggedright\arraybackslash}p{\dimexpr\linewidth/2-2\tabcolsep\relax}}
Mode & Participation semantics \\
\textbf{FirstToFinish} & The applicable first finishing/qualifying \texttt{\seqsplit{FulfillmentTask}} establishes participation in resolution of the outcome. This defines neither detailed terminal-state qualification nor cancellation or execution mechanics for other undertakings. \\
\textbf{AssignedToFinish} & A particular assigned \texttt{\seqsplit{FulfillmentTask}} has responsibility for establishing the outcome. Other undertakings may exist or execute concurrently; their existence or completion does not independently establish the ActionableTask outcome. The assignment mechanism is not defined here. \\
\textbf{Aggregate} & Multiple applicable \texttt{\seqsplit{FulfillmentTask}} undertakings participate in resolving the ActionableTask outcome. This selects the participating undertakings; it does not require recursive TaskOutcome composition or determine how their outputs form the resulting output. \\
\end{longtable}


\par\medskip\textbf{2.2.3 ActionableTask.Outcome.OutputRules.Mode}\par


\begin{quote}\small\ttfamily
ActionableTask.Outcome.OutputRules.Mode = Direct | Collection | Derived\\

\end{quote}


This independent property determines \textbf{how the outputs of the participating
FulfillmentTasks establish the ActionableTask's Task.Output(s)}.



\begin{longtable}{>{\raggedright\arraybackslash}p{\dimexpr\linewidth/3-2\tabcolsep\relax}>{\raggedright\arraybackslash}p{\dimexpr\linewidth/3-2\tabcolsep\relax}>{\raggedright\arraybackslash}p{\dimexpr\linewidth/3-2\tabcolsep\relax}}
Mode & Output semantics & Conceptual path \\
\textbf{Direct} & The applicable participating \texttt{\seqsplit{FulfillmentTask}} output is used directly as the ActionableTask output. & \texttt{\seqsplit{FulfillmentTask.Output}} → \texttt{\seqsplit{TaskOutcome}} → \texttt{\seqsplit{Task.Output}}. \\
\textbf{Collection} & The applicable participating \texttt{\seqsplit{FulfillmentTask}} outputs collectively form the ActionableTask output, preserving them as the relevant collection/set. Collection does not itself imply derivation of new domain meaning. & \texttt{\seqsplit{FulfillmentTask.Output [0..*]}} → Collection → \texttt{\seqsplit{TaskOutcome}} → \texttt{\seqsplit{Task.Output(s)}}. \\
\textbf{Derived} & A new ActionableTask output is derived from the applicable participating \texttt{\seqsplit{FulfillmentTask}} outputs according to explicitly defined rules. & \texttt{\seqsplit{FulfillmentTask.Output(s)}} → Explicit Rules → \texttt{\seqsplit{TaskOutcome}} → \texttt{\seqsplit{Task.Output(s)}}. \\
\end{longtable}


\texttt{\seqsplit{Derived}} is an information semantic. It prescribes no implementation
mechanism for executing the explicit rules. Harmonia need not infer the
encapsulated domain meaning of participating outputs in order to apply those
rules; the
\hyperref[element:10284034c8534b639284386b064f2c54]{observable-information/domain-meaning boundary}
continues to apply. Developer-defined Ergo logic and optional runtime AI
retain their existing boundaries; neither is prescribed as the means of
derivation.



\par\medskip\textbf{2.2.4 Orthogonality of the Three Controls}\par


\begin{longtable}{>{\raggedright\arraybackslash}p{\dimexpr\linewidth/2-2\tabcolsep\relax}>{\raggedright\arraybackslash}p{\dimexpr\linewidth/2-2\tabcolsep\relax}}
ActionableTask concern & Question answered \\
\textbf{ExecutionConcurrency} & How many FulfillmentTasks may execute concurrently? \\
\textbf{OutcomeConcurrency.Mode} & Which FulfillmentTask undertaking(s) participate in resolving the ActionableTask outcome? \\
\textbf{Outcome.OutputRules.Mode} & How are the participating FulfillmentTask outputs used to establish the ActionableTask Task.Output(s)? \\
\end{longtable}


These controls SHALL remain orthogonal and SHALL NOT be collapsed into one
lifecycle, state or concurrency mechanism. Multiple undertakings may execute
concurrently while only one participates in the TaskOutcome. Multiple may
participate while their outputs are collected or used to derive another
output. No universal mode combination, default or execution algorithm is
inferred from these examples.



Source standing: Bounded adjudicated conceptual Task meaning; section 5 reservations, including runtime TaskOutcome identity/lifecycle, remain unresolved.


{\footnotesize \href{../../../../continuum/application/data-object/topic/c43139289cb649b4836521a96a51604c.md}{Canonical record: c43139289cb649b4836521a96a51604c}; full qualifications and provenance.}\par
\subsubsection*{IV.1.1.10.5 [Topic] TaskDefinition}
\phantomsection
\label{navigation:IV.1.1.10.5}
\addcontentsline{toc}{subsubsection}{IV.1.1.10.5 [Topic] TaskDefinition}

\label{element:1d6d4c3a2df440be94262855fa6564eb}
\index{[Topic] TaskDefinition (IV.1.1.10.5)}
{\small DataObject / Topic · Bounded adjudicated conceptual Task meaning; section 5 reservations, including runtime TaskOutcome identity/lifecycle, remain unresolved. · current}\par

\textbf{Description}
Information defining a reusable form of actionable work, including the characteristics, expectations and constraints applicable when instances of that work are created.



\textbf{ActionableTaskArchetype} is the reusable definition/archetype for a kind of
actionable work. It describes the meaning and constraints of work that may
subsequently be instantiated. It SHALL NOT be treated as an instance of work.



Its semantics may include, where applicable:



\begin{itemize}


\item the expected nature of the work;

\item permissible parameters and constraints;

\item prerequisites and performer requirements;

\item completion or satisfaction expectations; and

\item expectations concerning information/content consumed, manipulated or
  produced by the work.

\end{itemize}


An archetype may establish expectations concerning encapsulated content
without defining that content's structural representation or requiring
Harmonia to comprehend its domain semantics. The
\hyperref[element:10284034c8534b639284386b064f2c54]{observable-information boundary}
applies.



Source standing: Bounded adjudicated conceptual Task meaning; section 5 reservations, including runtime TaskOutcome identity/lifecycle, remain unresolved.


{\footnotesize \href{../../../../continuum/application/data-object/topic/1d6d4c3a2df440be94262855fa6564eb.md}{Canonical record: 1d6d4c3a2df440be94262855fa6564eb}; full qualifications and provenance.}\par
\subsubsection*{IV.1.1.10.6 [Topic] TaskFulfillment}
\phantomsection
\label{navigation:IV.1.1.10.6}
\addcontentsline{toc}{subsubsection}{IV.1.1.10.6 [Topic] TaskFulfillment}

\label{element:30f4deffcd694f07bf54b2b192695d4f}
\index{[Topic] TaskFulfillment (IV.1.1.10.6)}
{\small DataObject / Topic · Bounded adjudicated conceptual Task meaning; section 5 reservations, including runtime TaskOutcome identity/lifecycle, remain unresolved. · current}\par

\textbf{Description}
Information describing an identifiable undertaking performed toward completion of a Task, allowing the work performed and its relationship to the originating Task to be represented.



\textbf{FulfillmentTask} represents an identifiable undertaking of an
\texttt{\seqsplit{ActionableTask}}: doing or attempting to do that particular work. It carries
the execution-related information necessary to represent the undertaking,
including, where applicable:



\begin{itemize}


\item performer / executing actor and execution context;

\item commencement and timing;

\item progression and execution state;

\item completion, failure, suspension or other applicable execution information;
  and

\item other execution metadata/scaffolding needed to represent the undertaking.

\end{itemize}


It has its own execution state and may produce output. It does \textbf{not
independently own TaskOutcome}. One or more undertakings may contribute to
resolution of the parent ActionableTask according to that ActionableTask's
explicitly defined outcome policy.



An \texttt{\seqsplit{ActionableTask}} may have \textbf{zero, one or multiple \texttt{\seqsplit{FulfillmentTasks}}}.
Multiple undertakings MAY execute concurrently where the ActionableTask's
\texttt{\seqsplit{ExecutionConcurrency}} permits them. Execution state SHALL NOT be collapsed into
\texttt{\seqsplit{ActionableTask}} on the assumption that each work instance has exactly one
execution.



Completion of one \texttt{\seqsplit{FulfillmentTask}} does \textbf{not inherently establish
satisfaction or completion of its parent \texttt{\seqsplit{ActionableTask}}}. Satisfaction
depends on the definition and context of the work, with outcome resolution
governed by the ActionableTask's explicit outcome policy. No universal
execution state machine or satisfaction rule is established here. Information
representing an undertaking remains distinct from the undertaking itself.



Source standing: Bounded adjudicated conceptual Task meaning; section 5 reservations, including runtime TaskOutcome identity/lifecycle, remain unresolved.


{\footnotesize \href{../../../../continuum/application/data-object/topic/30f4deffcd694f07bf54b2b192695d4f.md}{Canonical record: 30f4deffcd694f07bf54b2b192695d4f}; full qualifications and provenance.}\par
\subsubsection*{IV.1.1.10.7 [Topic] TaskOutcome}
\phantomsection
\label{navigation:IV.1.1.10.7}
\addcontentsline{toc}{subsubsection}{IV.1.1.10.7 [Topic] TaskOutcome}

\label{element:08d8ee19a30f418f9d5bc497ced0d1ef}
\index{[Topic] TaskOutcome (IV.1.1.10.7)}
{\small DataObject / Topic · Bounded adjudicated conceptual Task meaning; section 5 reservations, including runtime TaskOutcome identity/lifecycle, remain unresolved. · current}\par

\textbf{Description}
Information describing the governed resolution of a Task, including its completion disposition, resulting information and applicable outputs.



\begin{quote}


\textbf{TaskOutcome is the governed resolution of an ActionableTask according to
that ActionableTask's explicitly defined outcome policy.}



\end{quote}


TaskOutcome belongs to and resolves the ActionableTask. It does not
independently belong to a FulfillmentTask. Conceptually it:



\begin{itemize}


\item **references FulfillmentTask [0..*]**: identifies the undertaking or
  undertakings that participated in establishing the outcome, as determined
  by \texttt{\seqsplit{OutcomeConcurrency.Mode}};

\item \textbf{contains Task Completion Metadata}: Harmonia-governed information
  concerning resolution/completion of the ActionableTask; no complete
  metadata schema is defined here; and

\item \textbf{contains / represents Task.Output(s)}: the resulting ActionableTask
  output established from the participating undertakings' outputs according
  to \texttt{\seqsplit{Outcome.OutputRules.Mode}}.

\end{itemize}


The encapsulated domain content of a \texttt{\seqsplit{Task.Output}} may remain semantically
opaque to Harmonia. Neither output information nor TaskOutcome becomes the
physical, clinical or other phenomenon represented. No detailed Task.Output
payload structure is established here.



\textbf{Recursive TaskOutcome containment is excluded from this model.} Multiple
execution results are represented through references from the single
governed TaskOutcome to the relevant participating FulfillmentTask
undertakings, rather than through recursive outcome aggregation hierarchies.
Aggregate participation does not imply recursive composition. The remaining
questions in §5 remain unresolved.



Source standing: Bounded adjudicated conceptual Task meaning; section 5 reservations, including runtime TaskOutcome identity/lifecycle, remain unresolved.



Under \hyperref[element:8882573df7ce43c6bc0f66da0a676139]{AX-17}, the following
remaining Task / Work matters require subsequent architectural adjudication
or separately governed detailed definition:



\begin{itemize}


\item TaskOutcome identity rules;

\item detailed cardinalities beyond the established relationships, including
  whether an outcome is mandatory;

\item detailed TaskOutcome lifecycle/state semantics;

\item detailed Task Completion Metadata and Task.Output payload structures, and
  their representation mechanisms;

\item detailed finishing/qualification rules, assignment mechanisms and
  concurrency execution mechanics; and

\item task-specific explicit output rules and their downstream execution
  mechanisms.

\end{itemize}


Outcome ownership, participation modes, output-rule modes, their orthogonality
and exclusion of recursive TaskOutcome composition are adjudicated above;
they SHALL NOT be presented as still provisional or unresolved. This
refinement allocates no previously unestablished archetype-authoring or
outcome-authority owner.



\textbf{ReportedTask remains unresolved.} Its existing intersection with
Definition → Contextualisation → Fulfilment → Outcome → Accountability does
not reaffirm its previous definition or require a separate accountable report.
Its necessity and semantics may overlap accountability, assurance, audit,
historical representation and evidence, and require separate adjudication.
It is neither removed nor normalised into one of those concerns here. No
ReportedTask relationships, identity, immutability or lifecycle are settled.



The metamodel, reusable pattern and general lifecycle/containment guidance
SHALL NOT be used to fill these reserved questions by inference. The bounded
semantics above are established; the remaining Task / Work family and Domain04
are not declared complete.


{\footnotesize \href{../../../../continuum/application/data-object/topic/08d8ee19a30f418f9d5bc497ced0d1ef.md}{Canonical record: 08d8ee19a30f418f9d5bc497ced0d1ef}; full qualifications and provenance.}\par
\subsection*{IV.1.1.11 [Topic] Workforce}
\phantomsection
\label{navigation:IV.1.1.11}
\addcontentsline{toc}{subsection}{IV.1.1.11 [Topic] Workforce}

\label{element:9f689e261140436eb94295d63e2b3370}
\index{[Topic] Workforce (IV.1.1.11)}
{\small DataObject / Topic · Accepted bounded reconciliation · current}\par

\textbf{Description}: Information concerning healthcare practitioners and workforce participants, including professional identity, registration, qualifications, privileges, organisational affiliations, roles, and communication endpoints.


{\footnotesize \href{../../../../continuum/application/data-object/topic/9f689e261140436eb94295d63e2b3370.md}{Canonical record: 9f689e261140436eb94295d63e2b3370}; full qualifications and provenance.}\par
\subsubsection*{IV.1.1.11.1 [Topic] CommunicationAddress}
\phantomsection
\label{navigation:IV.1.1.11.1}
\addcontentsline{toc}{subsubsection}{IV.1.1.11.1 [Topic] CommunicationAddress}

\label{element:410be08eea0b4653a88f4e4235354cea}
\index{[Topic] CommunicationAddress (IV.1.1.11.1)}
{\small DataObject / Topic · Accepted bounded reconciliation · current}\par

\textbf{Description}
Information identifying an electronic communication destination associated with a practitioner or professional context, including the addressing and connection information required to communicate with that destination.


{\footnotesize \href{../../../../continuum/application/data-object/topic/410be08eea0b4653a88f4e4235354cea.md}{Canonical record: 410be08eea0b4653a88f4e4235354cea}; full qualifications and provenance.}\par
\subsubsection*{IV.1.1.11.2 [Topic] Practitioner}
\phantomsection
\label{navigation:IV.1.1.11.2}
\addcontentsline{toc}{subsubsection}{IV.1.1.11.2 [Topic] Practitioner}

\label{element:4e8f2bd581a5418889813cf7ee414750}
\index{[Topic] Practitioner (IV.1.1.11.2)}
{\small DataObject / Topic · Established Package1 conceptual information meaning; residual qualifications in selected section remain unchanged. · current}\par

\textbf{Definition}: An identifiable individual healthcare workforce or professional entity.



\textbf{Description}: Information describing an individual participating professionally in healthcare, independent of a particular organisational position, assignment or service context.



\textbf{Semantic Classification}: \texttt{\seqsplit{Entity}}



\textbf{Architectural Scope}: Represents the individual professional independently of statutory registration status. A \texttt{\seqsplit{Practitioner}} may be known to Harmonia prior to verification, when registration is pending or not applicable (e.g. unregistered healthcare workers, support staff), or when registration is externally authoritative.



\textbf{Boundary Rule}: \$\textbackslash{}text\{Practitioner\} \textbackslash{}neq \textbackslash{}text\{Professional Registration\}\$. Registration is a governed assertion concerning a Practitioner, not the definition of the Practitioner entity itself.



Source standing: Established Package1 conceptual information meaning; residual qualifications in selected section remain unchanged.


{\footnotesize \href{../../../../continuum/application/data-object/topic/4e8f2bd581a5418889813cf7ee414750.md}{Canonical record: 4e8f2bd581a5418889813cf7ee414750}; full qualifications and provenance.}\par
\subsubsection*{IV.1.1.11.3 [Topic] PractitionerRole}
\phantomsection
\label{navigation:IV.1.1.11.3}
\addcontentsline{toc}{subsubsection}{IV.1.1.11.3 [Topic] PractitionerRole}

\label{element:cb7ac98f62aa4599ba57263338cb8b3c}
\index{[Topic] PractitionerRole (IV.1.1.11.3)}
{\small DataObject / Topic · Established Package1 conceptual information meaning; residual qualifications in selected section remain unchanged. · current}\par

\textbf{Definition}: The contextual binding of a \texttt{\seqsplit{Practitioner}} to a specific clinical capacity, service, facility, or department within an organisation (e.g. \emph{Specialist Physician in Cardiology Department}, \emph{On-Call Registrar}).



\textbf{Description}: Information describing the contextual relationship between a practitioner and the organisation, service, location or role through which they participate in healthcare delivery.



\textbf{Semantic Classification}: \texttt{\seqsplit{Contextual Binding}}



\textbf{Key Characteristics}: Contextual title/role, participating organisation, participating healthcare service, physical practice locations, active schedule, and linked communication endpoints.



Source standing: Established Package1 conceptual information meaning; residual qualifications in selected section remain unchanged.


{\footnotesize \href{../../../../continuum/application/data-object/topic/cb7ac98f62aa4599ba57263338cb8b3c.md}{Canonical record: cb7ac98f62aa4599ba57263338cb8b3c}; full qualifications and provenance.}\par
\subsubsection*{IV.1.1.11.4 [Topic] Privilege}
\phantomsection
\label{navigation:IV.1.1.11.4}
\addcontentsline{toc}{subsubsection}{IV.1.1.11.4 [Topic] Privilege}

\label{element:75e194f6168945588ed30cc6c78d04cc}
\index{[Topic] Privilege (IV.1.1.11.4)}
{\small DataObject / Topic · Established Package1 conceptual information meaning; residual qualifications in selected section remain unchanged. · current}\par

\textbf{Definition}: An enterprise-approved scope of clinical practice, procedural authorization, admission entitlement, or prescribing privilege granted by a healthcare organisation.



\textbf{Description}: Information describing the activities, services or clinical practices a practitioner is authorised to undertake within a particular organisational or service context.



\textbf{Semantic Classification}: \texttt{\seqsplit{Assertion / Entitlement Directive}}



\textbf{Key Characteristics}: Granting organisation, authorized clinical scope / procedure codes, supervision requirements, effective validity period, and approving clinical governance authority.



Source standing: Established Package1 conceptual information meaning; residual qualifications in selected section remain unchanged.


{\footnotesize \href{../../../../continuum/application/data-object/topic/75e194f6168945588ed30cc6c78d04cc.md}{Canonical record: 75e194f6168945588ed30cc6c78d04cc}; full qualifications and provenance.}\par
\subsubsection*{IV.1.1.11.5 [Topic] ProfessionalIdentity}
\phantomsection
\label{navigation:IV.1.1.11.5}
\addcontentsline{toc}{subsubsection}{IV.1.1.11.5 [Topic] ProfessionalIdentity}

\label{element:2f2868f077fc47afbb541925a9801b2c}
\index{[Topic] ProfessionalIdentity (IV.1.1.11.5)}
{\small DataObject / Topic · Established Package1 conceptual information meaning; residual qualifications in selected section remain unchanged. · current}\par

\textbf{Definition}: The governed set of professional identifiers, credential assertions, and demographic traits establishing the professional identity of a \texttt{\seqsplit{Practitioner}}.



\textbf{Description}: Information establishing and correlating the professional identity of a practitioner across organisations, authorities and healthcare contexts.



\textbf{Semantic Classification}: \texttt{\seqsplit{Assertion / Identity Bundle}}
\textbf{Key Semantic Decomposition}:
  - \texttt{\seqsplit{Identifier Value}}: The token or registration number identifying the practitioner.
  - \texttt{\seqsplit{Identifier Authority / Namespace}}: The statutory body, national authority, or enterprise issuing the identifier.



Source standing: Established Package1 conceptual information meaning; residual qualifications in selected section remain unchanged.


{\footnotesize \href{../../../../continuum/application/data-object/topic/2f2868f077fc47afbb541925a9801b2c.md}{Canonical record: 2f2868f077fc47afbb541925a9801b2c}; full qualifications and provenance.}\par
\subsubsection*{IV.1.1.11.6 [Topic] ProfessionalProfile}
\phantomsection
\label{navigation:IV.1.1.11.6}
\addcontentsline{toc}{subsubsection}{IV.1.1.11.6 [Topic] ProfessionalProfile}

\label{element:703bb80c062743eeb7ed8dea3ffeaaa0}
\index{[Topic] ProfessionalProfile (IV.1.1.11.6)}
{\small DataObject / Topic · Established Package1 conceptual information meaning; residual qualifications in selected section remain unchanged. · current}\par

\textbf{Definition}: The verified set of educational qualifications, clinical competencies, certified subspecialties, and academic credentials possessed by a practitioner.



\textbf{Description}: Information describing the professional characteristics of a practitioner, including qualifications, competencies, specialties and other information relevant to professional practice.



\textbf{Semantic Classification}: \texttt{\seqsplit{Assertion / Profile Bundle}}



Source standing: Established Package1 conceptual information meaning; residual qualifications in selected section remain unchanged.


{\footnotesize \href{../../../../continuum/application/data-object/topic/703bb80c062743eeb7ed8dea3ffeaaa0.md}{Canonical record: 703bb80c062743eeb7ed8dea3ffeaaa0}; full qualifications and provenance.}\par
\subsubsection*{IV.1.1.11.7 [Topic] Registration}
\phantomsection
\label{navigation:IV.1.1.11.7}
\addcontentsline{toc}{subsubsection}{IV.1.1.11.7 [Topic] Registration}

\label{element:26c72794993044af98b5329b8b14ae89}
\index{[Topic] Registration (IV.1.1.11.7)}
{\small DataObject / Topic · Established Package1 conceptual information meaning; residual qualifications in selected section remain unchanged. · current}\par

\textbf{Definition}: A formal, externally governed assertion attesting that a practitioner is registered with a statutory registration board or regulatory authority.



\textbf{Description}: Information describing a practitioner's professional registration, including the registering authority, profession, registration status, scope and applicable period.



\textbf{Semantic Classification}: \texttt{\seqsplit{Assertion / Governed Relationship}}



\textbf{Key Characteristics}: Registration authority, registration category, endorsed specialties, condition/restriction status, and validity period.



\textbf{Authority Invariant}: Originates from an external regulatory authority. Harmonia records verification provenance without claiming originating registration authority.



Source standing: Established Package1 conceptual information meaning; residual qualifications in selected section remain unchanged.


{\footnotesize \href{../../../../continuum/application/data-object/topic/26c72794993044af98b5329b8b14ae89.md}{Canonical record: 26c72794993044af98b5329b8b14ae89}; full qualifications and provenance.}\par
\section*{IV.1.2 Reusable Information Concepts}
\phantomsection
\label{navigation:IV.1.2}
\addcontentsline{toc}{section}{IV.1.2 Reusable Information Concepts}

\subsection*{IV.1.2.1 PractitionerCapacity}
\phantomsection
\label{navigation:IV.1.2.1}
\addcontentsline{toc}{subsection}{IV.1.2.1 PractitionerCapacity}

\label{element:008f622bf9ce4d40b8a85024748bd485}
\index{PractitionerCapacity (IV.1.2.1)}
{\small DataObject · Candidate — adjudicated Domain05 Step-B baseline; final Solution Architecture unapproved · current}\par

For a Practitioner, the available PractitionerCapacity set is constrained by
the practitioner's governed profile and applicable operational context. The
Practitioner selects \textbf{one active PractitionerCapacity} for the applicable
session/operational or roster context. It supplies operational context used by
Themis when deriving effective authorisation; \textbf{it does not itself confer authority}.
Selection UI and the exact PractitionerCapacity ↔ FHIR PractitionerRole mapping
remain deferred.



\begin{quote}


\textbf{Authentication establishes who the Individual is. PractitionerCapacity
establishes the capacity in which a Practitioner is currently acting. Themis
establishes what that Identity, acting in that capacity, may do.}



\end{quote}


\begin{quote}\small\ttfamily
graph TD\\
    Auth["Authenticated Practitioner"] --> Eligible["Eligible PractitionerCapacity set"]\\
    Profile["Governed practitioner profile"] --> Eligible\\
    Context["Applicable operational / roster context"] --> Eligible\\
    Eligible -->|"Practitioner selects ONE"| Capacity["Active PractitionerCapacity: context, no authority by itself"]\\
    Capacity -->|"authorisation context"| PDP["Themis: resolve governed Criteria / assignments / policy state"]\\
    PDP --> Criteria["Identity-specific precomputed PEP-Criteria through PEP-Context"]\\
    Criteria --> PEP["Constrained local Themis PEP"]\\
    Capacity -.->|"operational responsibility context"| Responsibility["Agora: membership intent; detailed context use deferred"]\\

\end{quote}


Themis answers \textbf{“May this Identity perform this action on this information?”}
Future Agora collaboration detail concerns \textbf{“Is this Practitioner currently
responsible for receiving/responding to this operational activity?”} These are
separate responsibilities even where they consume the same PractitionerCapacity.
Receipt of an escalation or approval notification grants no information access.
Following it into Iris or another application surface requires normal Themis PEP
enforcement using current valid PEP-Criteria. U16 establishes
\hyperref[element:d1c4053e68384080b8eb79d2008632ba]{Agora control and Omilia operations};
detailed membership, escalation and approval-routing mechanics remain deferred.


{\small Recorded unresolved matters: canonicalName; exact scope remains in the canonical record.}\par
{\footnotesize \href{../../../../continuum/application/data-object/008f622bf9ce4d40b8a85024748bd485.md}{Canonical record: 008f622bf9ce4d40b8a85024748bd485}; full qualifications and provenance.}\par
\section*{IV.1.3 Logical Packages}
\phantomsection
\label{navigation:IV.1.3}
\addcontentsline{toc}{section}{IV.1.3 Logical Packages}

\subsection*{IV.1.3.1 [Package] model}
\phantomsection
\label{navigation:IV.1.3.1}
\addcontentsline{toc}{subsection}{IV.1.3.1 [Package] model}

\label{element:6bb218f2128e463482f505585e624889}
\index{[Package] model (IV.1.3.1)}
{\small Grouping / Package · Accepted bounded architecture — HCRL 1.1.15-001 DR-1 · current}\par

Top-level Calliope model package grouping representing the architectural model namespace root.


{\footnotesize \href{../../../../continuum/grouping/package/6bb218f2128e463482f505585e624889.md}{Canonical record: 6bb218f2128e463482f505585e624889}; full qualifications and provenance.}\par
\subsubsection*{IV.1.3.1.1 [Package] core}
\phantomsection
\label{navigation:IV.1.3.1.1}
\addcontentsline{toc}{subsubsection}{IV.1.3.1.1 [Package] core}

\label{element:2d587cb7547e4e6a8c20aa64da0da78a}
\index{[Package] core (IV.1.3.1.1)}
{\small Grouping / Package · Accepted bounded architecture — HCRL 1.1.15-001 DR-1 · current}\par

Calliope core package grouping nested under model.


{\footnotesize \href{../../../../continuum/grouping/package/2d587cb7547e4e6a8c20aa64da0da78a.md}{Canonical record: 2d587cb7547e4e6a8c20aa64da0da78a}; full qualifications and provenance.}\par
\paragraph*{IV.1.3.1.1.1 [Package] information}
\phantomsection
\label{navigation:IV.1.3.1.1.1}
\addcontentsline{toc}{paragraph}{IV.1.3.1.1.1 [Package] information}

\label{element:13c028886fde4682812b45772e21522f}
\index{[Package] information (IV.1.3.1.1.1)}
{\small Grouping / Package · Accepted bounded architecture — HCRL 1.1.15-001 DR-1 · current}\par

Calliope information package grouping nested under model.core containing logical model classes.


{\footnotesize \href{../../../../continuum/grouping/package/13c028886fde4682812b45772e21522f.md}{Canonical record: 13c028886fde4682812b45772e21522f}; full qualifications and provenance.}\par
\subparagraph*{IV.1.3.1.1.1.1 [Class] InformationUnit}
\phantomsection
\label{navigation:IV.1.3.1.1.1.1}
\addcontentsline{toc}{subparagraph}{IV.1.3.1.1.1.1 [Class] InformationUnit}

\label{element:74ec609de1ae498fbf9d13e4ef8c6cbe}
\index{[Class] InformationUnit (IV.1.3.1.1.1.1)}
{\small DataObject / Class · Accepted bounded architecture — HCRL 1.1.15-001 DR-1 · current}\par

The governed boundary associating a Governed Representation with its Management Context.


{\footnotesize \href{../../../../continuum/application/data-object/class/74ec609de1ae498fbf9d13e4ef8c6cbe.md}{Canonical record: 74ec609de1ae498fbf9d13e4ef8c6cbe}; full qualifications and provenance.}\par
\chapter*{IV.2 System-of-Systems Architecture}
\phantomsection
\label{navigation:IV.2}
\addcontentsline{toc}{chapter}{IV.2 System-of-Systems Architecture}

\section*{IV.2.1 System Map}
\phantomsection
\label{navigation:IV.2.1}
\addcontentsline{toc}{section}{IV.2.1 System Map}

The existing application responsibilities form the system map. These links identify independently defined systems and components; no new enclosing system or relationship is asserted.\par
\begin{longtable}{p{.22\linewidth}p{.73\linewidth}}
IV.3.1.1.1 & \hyperref[element:d1c4053e68384080b8eb79d2008632ba]{Agora} \\
IV.3.2.1.1 & \hyperref[element:466ad8f75a954c0581b5bba090443849]{Calliope} \\
IV.3.3.1.1 & \hyperref[element:3f64d339648246e99e8866c8089f8960]{DigitalTwin} \\
IV.3.4.1.1 & \hyperref[element:d739c75dffe2453081735710b285f03b]{Dokimasia} \\
IV.3.5.1.1 & \hyperref[element:fb3f37a6ac614415ac486a6de9d35994]{Energeia} \\
IV.3.6.1.1 & \hyperref[element:33b8421ffc444b78b2d059201a58a688]{Harmonia.Omilia.Client} \\
IV.3.7.1.1 & \hyperref[element:369907025d06492bb414a72cea834a77]{Harmonia.Omilia.Server} \\
IV.3.8.1.1 & \hyperref[element:feea742463b84021b66853d287d8916f]{Harmonia.Sylloge} \\
IV.3.9.1.1 & \hyperref[element:e927c1af96c64fcd979be659bee2cbe3]{Iris} \\
IV.3.10.1.1 & \hyperref[element:fefda110e2274884ad58906733606efd]{Kleio} \\
IV.3.11.1.1 & \hyperref[element:1af8a56358f7429fabfb57be66542a99]{Mneme} \\
IV.3.12.1.1 & \hyperref[element:7985e647bf9e4ca3ba70faf9e3e4ae02]{Mnemosyne} \\
IV.3.13.1.1 & \hyperref[element:b521762c1ea444b48cd5f66be3752d30]{Ponos} \\
IV.3.14.1.1 & \hyperref[element:e3e4147fea3640198dc251ed41ee854d]{Pylai} \\
IV.3.15.1.1 & \hyperref[element:899d34ec491948b68aad09e507dea6d1]{Themis} \\
IV.2.3.1.1.1 & \hyperref[element:195e5f10db8e419391354197b5d7a8cc]{Petasos} \\
\end{longtable}
\section*{IV.2.2 System Responsibility Map}
\phantomsection
\label{navigation:IV.2.2}
\addcontentsline{toc}{section}{IV.2.2 System Responsibility Map}

\subsection*{IV.2.2.1 Provider Registry solution collaboration}
\phantomsection
\label{navigation:IV.2.2.1}
\addcontentsline{toc}{subsection}{IV.2.2.1 Provider Registry solution collaboration}

\label{element:9305dc8d55b34b70950190e28b5464cf}
\index{Provider Registry solution collaboration (IV.2.2.1)}
{\small ApplicationCollaboration · Candidate — accepted Step-B baseline; final Solution Architecture unapproved · current}\par

\begin{itemize}


\item \textbf{Role/responsibility:} documented application solution for directory resources
  and governed changes across existing components, rather than a new enclosing
  platform component. \href{\detokenize{../../../../../docs/markdown/05-information-systems-architecture/sources.md#l15}}{L15/L18}.

\item \textbf{Features/contained constructs:} synchronous read/search, request-for-change,
  per-resource validation/identifier uniqueness/reference checks, candidate automated
  approval, durable commit and observable task progress. Seven per-resource Ergon
  activities and a registry change sequence are documented behavioural specialisations;
  Erga supplies the distinct collection/assembly/library responsibility.

\item \textbf{Information:} Practitioner, PractitionerRole, Organization, Location,
  HealthcareService, Endpoint and Group; change Pragma/input/output, errors and
  resulting versions. Endpoint and Group remain independently referenced resources.

\item \textbf{Persistence:} clinical/FHIR Mnemosyne slice; current-row versus history
  structures unclear. Iris is presentation-only; server-side governance remains
  with Mnemosyne Clinical, Ergon activities within Erga, and Themis under
  AGENTS Invariant 3.

\item \textbf{Collaboration/exclusions:} Pylai interaction, Themis, Petasos, Ponos/Praxis/Erga,
  Mneme/Mnemosyne and Iris Administration. Does not automatically own all provider,
  care-team or national-source facts, or imply an EMPI.

\item \textbf{Standing:} documented candidate domain solution, CDG-03/11/12/21/32/33.
  Direct persistence search and published async response claims are not approved
  external contracts by inclusion here.

\end{itemize}

{\small Recorded unresolved matters: canonicalName; exact scope remains in the canonical record.}\par
{\footnotesize \href{../../../../continuum/application/application-collaboration/9305dc8d55b34b70950190e28b5464cf.md}{Canonical record: 9305dc8d55b34b70950190e28b5464cf}; full qualifications and provenance.}\par
\section*{IV.2.3 Middleware Architecture}
\phantomsection
\label{navigation:IV.2.3}
\addcontentsline{toc}{section}{IV.2.3 Middleware Architecture}

\subsection*{IV.2.3.1 Petasos}
\phantomsection
\label{navigation:IV.2.3.1}
\addcontentsline{toc}{subsection}{IV.2.3.1 Petasos}

\subsubsection*{IV.2.3.1.1 Component Architecture}
\phantomsection
\label{navigation:IV.2.3.1.1}
\addcontentsline{toc}{subsubsection}{IV.2.3.1.1 Component Architecture}

\paragraph*{IV.2.3.1.1.1 Petasos}
\phantomsection
\label{navigation:IV.2.3.1.1.1}
\addcontentsline{toc}{paragraph}{IV.2.3.1.1.1 Petasos}

\label{element:195e5f10db8e419391354197b5d7a8cc}
\index{Petasos (IV.2.3.1.1.1)}
{\small ApplicationComponent · Candidate — accepted Step-B baseline; final Solution Architecture unapproved · current}\par

\begin{itemize}


\item \textbf{Role/responsibility:} \textbf{Harmonia's durable message/event queue-management
  component}, a first-class shared messaging/event capability. Petasos manages
  creation, operation and operational state/status of its durable queues.
  \href{\detokenize{../../../../../docs/markdown/05-information-systems-architecture/sources.md#u09}}{U09} extends the U05 handoff boundary, preserving AX-10.

\item \textbf{Purpose-defined queues:} purposes include inter-Ergon information flow,
  defined communication/event sharing, notification/event sharing and other
  architecturally established Harmonia event-sharing services. Queue purpose is
  architecturally defined; queues are not arbitrary generic queues. Final taxonomy,
  names, registration/configuration and topology remain unresolved.

\item \textbf{Features:} queue management and durable Parcel communication/event handling;
  persistent handoff obligations, bounded backlog, destination mechanics and
  handoff redelivery/recovery/dead-letter and designated replay responsibilities
  remain. Detailed delivery/ordering/duplicate/replay/recovery contracts are not
  newly selected; CDG-16/17/19/34/37/44 retain uncertainty.

\end{itemize}






\par\medskip\textbf{PetasosParcel and information ownership}\par


Every Petasos queue carries the single common datatype \textbf{PetasosParcel}.
It represents the queued Harmonia communication/event and is conceptually
analogous to \textbf{FHIR Communication}, without established literal equivalence.
No schema or mapping to historical PetasosMessage/ErgonEvent is inferred.



\begin{quote}\small\ttfamily
PetasosParcel          pass-by-value; the queued communication/event\\
Referenced Information predominantly pass-by-reference;\\
                       remains managed by its authoritative component\\

\end{quote}


\textbf{PetasosParcel is passed by value.} Petasos does not transport only references
and does not duplicate all associated information into the Parcel. Petasos owns
Parcel messaging/event semantics and queue operations; references confer no
ownership of the referenced information.



For the \textbf{ActionableTask handoff special case}, ActionableTask/Pragma and payload
remain managed through Mneme. A Parcel may reference that ActionableTask and is
itself queued by value. Petasos does not transport another consumer's
FulfillmentTask; each consumer undertaking the ActionableTask establishes its own
distinct FulfillmentTask. The established no-payload-copy rule for this case
does not make all Parcels reference-only. \href{\detokenize{../../../../../docs/markdown/05-information-systems-architecture/pathways-journeys-and-handoff.md#persistent-handoff}}{Handoff view}.



\begin{itemize}


\item \textbf{Persistence:} durable queues/Parcels and handoff recovery state are Petasos
  responsibility; referenced managed-information durable establishment remains
  Mnemosyne. Queue persistence mechanisms are not selected and queues do not
  become automatically permanent audit evidence.

\item \textbf{Operational observation:} Petasos owns queue operational state/status.
  Energeia may \textbf{observe/report} relevant state, without acquiring queue management
  or autonomous queue control. Detailed metrics/state structures remain unresolved.

\item \textbf{Themis collaboration:} for a governed FORWARD action, Petasos is the host/
  enforcement location because it controls the queue operation. Applicable
  enforcement SHALL use \textbf{Themis-owned PEP capability} and relevant \textbf{PEP-Criteria through PEP-Context}. Themis owns authorisation/enforcement semantics; Petasos owns the
  operation. Detailed protocol remains unresolved.

\item \textbf{Other collaborations/exclusions:} serves Ergon, ingress, execution,
  publication, Agora and architecturally established communication/notification
  services. It is not subordinate to Ponos, ActionableTask-only, a reference
  transport helper or merely a generic enterprise broker abstraction. Mneme owns
  Publishing/Subscription/Pattern semantic selection; Ponos owns execution.

\item \textbf{Standing:} U09 resolves queue responsibility and common datatype/pass semantics;
  \href{\detokenize{../../../../../docs/markdown/05-information-systems-architecture/gaps.md#petasos-detailed-contracts}}{CDG-44} retains schema, queue and interaction
  detail. Internal middleware remains Domain05; external protocols remain Domain06.

\end{itemize}

{\small Recorded unresolved matters: canonicalName; exact scope remains in the canonical record.}\par
{\footnotesize \href{../../../../continuum/application/application-component/195e5f10db8e419391354197b5d7a8cc.md}{Canonical record: 195e5f10db8e419391354197b5d7a8cc}; full qualifications and provenance.}\par
\subsubsection*{IV.2.3.1.2 Information Architecture}
\phantomsection
\label{navigation:IV.2.3.1.2}
\addcontentsline{toc}{subsubsection}{IV.2.3.1.2 Information Architecture}

\paragraph*{IV.2.3.1.2.1 PetasosParcel Topic}
\phantomsection
\label{navigation:IV.2.3.1.2.1}
\addcontentsline{toc}{paragraph}{IV.2.3.1.2.1 PetasosParcel Topic}

\label{element:f0a88f05a0cb438cb42572ec07fb8513}
\index{PetasosParcel Topic (IV.2.3.1.2.1)}
{\small DataObject · Bounded adjudicated application information meaning; Domain05 remains the accepted candidate Step-B baseline, not final Solution Architecture. Detailed representation is deferred; only this defined logical object is projected. · current}\par

\begin{quote}


\textbf{Mneme determines semantic eligibility and eligible recipient/consumer
relationships under Publishing, Subscription and Pattern semantics. Petasos
establishes and manages the concrete persistent handoff obligation to the
destination and reliably effects that handoff. Ponos manages execution.}



\end{quote}


Eligible consumer relationships express semantic selection. They do not assign
queue binding, destination mechanics or concrete handoff state to Mneme.



Petasos is Harmonia's durable message/event queue-management component.
It creates, operates and owns the operational state/status of \textbf{purpose-defined
queues} for inter-Ergon flow and other architecturally established communication,
notification and event sharing. Every queue carries \textbf{PetasosParcel by value};
associated information is predominantly by reference to governed information
managed elsewhere. Parcel is conceptually analogous to FHIR Communication;
literal equivalence and detailed schema are not established.



For the ActionableTask handoff special case, governed ActionableTask/Pragma and
Information Unit payload remain managed through Mneme. A by-value Parcel may
carry references and communication/handoff metadata required for resolution and
recovery. Petasos acquires no management authority over the referenced information
and does not copy the managed ActionableTask payload in this pattern.



\begin{quote}\small\ttfamily
Mneme manages ActionableTask / Pragma and payload\\
  --> semantic selection / eligibility\\
Petasos persistently manages handoff obligation\\
  --> queues PetasosParcel by value, referring to ActionableTask\\
Consumer resolves governed information through Mneme\\
  --> establishes its own distinct FulfillmentTask undertaking\\
Ponos manages execution / progression\\
  --> Ergon supplies bounded behaviour within Ergo's contract\\

\end{quote}


\textbf{Petasos distributes the opportunity to undertake work, not another consumer's
FulfillmentTask}, in this pattern. Each consumer undertaking the ActionableTask
establishes a distinct FulfillmentTask; no undertaking is transported on its behalf.
The Ergo framework works with \textbf{Mneme Context Load} to materialise anticipated
values before \textbf{Context Loader} exposes them to business logic. \textbf{Context Unloader}
returns changed working values and new Content as \textbf{Pragma Output} across that
boundary; \textbf{Mneme Context Save} manages subsequent consequences. The Unloader is
not persistence. Bounded Adhoc Resource Access/Search remains supported; Mnemosyne
establishes/reports durable state. \hyperref[element:49c9eb2bef0c47b8b5328621d5fc5479]{U10/U11 lifecycle}. Handoff acceptance, information commit, undertaking
completion, work resolution and external acknowledgement remain distinct.



ActionableTask handoff is one Petasos use case. Communication/Notification Events
and other established event-sharing purposes use the common PetasosParcel datatype
without becoming ActionableTasks or losing their own semantics. The no-payload-copy
rule above SHALL NOT be generalised to reference-only Petasos transport or to
copying all referenced information into every Parcel. Final purpose taxonomy,
names, Parcel schema and detailed delivery contracts remain \href{\detokenize{../../../../../docs/markdown/05-information-systems-architecture/gaps.md#petasos-detailed-contracts}}{CDG-44}.



Petasos owns queue operational state/status; Energeia may observe/report it,
without acquiring queue management or autonomous control. For governed FORWARD,
Petasos hosts enforcement because it controls the operation, using \textbf{Themis-owned
PEP capability and relevant PEP-Criteria through PEP-Context}. Themis owns semantics; Petasos
owns the queue operation. Detailed protocols/state metrics remain unresolved.



Source standing: Bounded adjudicated application information meaning; Domain05 remains the accepted candidate Step-B baseline, not final Solution Architecture. Detailed representation is deferred; only this defined logical object is projected.


{\footnotesize \href{../../../../continuum/application/data-object/f0a88f05a0cb438cb42572ec07fb8513.md}{Canonical record: f0a88f05a0cb438cb42572ec07fb8513}; full qualifications and provenance.}\par
\chapter*{IV.3 Applications}
\phantomsection
\label{navigation:IV.3}
\addcontentsline{toc}{chapter}{IV.3 Applications}

\section*{IV.3.1 Agora}
\phantomsection
\label{navigation:IV.3.1}
\addcontentsline{toc}{section}{IV.3.1 Agora}

\subsection*{IV.3.1.1 Application Overview}
\phantomsection
\label{navigation:IV.3.1.1}
\addcontentsline{toc}{subsection}{IV.3.1.1 Application Overview}

\subsubsection*{IV.3.1.1.1 Agora}
\phantomsection
\label{navigation:IV.3.1.1.1}
\addcontentsline{toc}{subsubsection}{IV.3.1.1.1 Agora}

\label{element:d1c4053e68384080b8eb79d2008632ba}
\index{Agora (IV.3.1.1.1)}
{\small ApplicationComponent · Candidate — accepted Step-B baseline; final Solution Architecture unapproved · current}\par

\begin{itemize}


\item \textbf{Role/responsibility:} \textbf{Harmonia.Agora}, collaboration control and semantics;
  native Teams/Slack-like collaboration. \textbf{Agora controls; Omilia operates.}
  \hyperref[element:d1c4053e68384080b8eb79d2008632ba]{U16 canonical view}
  replaces the earlier projection-only account.

\item \textbf{Features:} controls designated managed Room/Space lifecycle and membership
  intent and Harmonia publication; listens to dialogue/activity in \textbf{all} managed
  Rooms/Spaces through Omilia.Server's Application Services API. Defined meaningful
  interactions may be captured as Communication Information Units and/or Pragma
  for governed Harmonia processing. Not every message becomes work.

\item \textbf{Information:} owns/manages native collaboration record/state: dialogue/actions,
  Room/Space and membership state, references/activity and native collaboration
  media. Information merely conveyed through Agora retains its originating owner.
  Capture derives new governed information/activity; original dialogue remains Agora's.

\item \textbf{Membership:} Agora controls intent; Omilia.Server performs membership
  operations. Roster Management is SEP; only applicable Shift Start/End periods
  are required for R1.x/R2.x operational context. No roster/workforce authority or
  inference that PractitionerCapacity/Provider Registry is a roster.

\item \textbf{Persistence/audit:} application ownership is established, not a storage
  allocation. AX-05 remains controlling; collaboration persistence/media mechanics,
  mapping/transaction structures and retention are deferred. Dialogue retention
  and logging are distinct from audit; applicable occurrences go to Kleio, which
  retains authoritative Audit Information Management and existing capture rules.

\item \textbf{Collaborations/exclusions:} Omilia.Server performs underlying operations;
  Agora acquires no ownership of those mechanics. Petasos mediates execution
  collaboration, with no direct Ponos dependency. AGENTS Invariant 10 preserves
  Themis default-deny, Matrix DTO encapsulation and zero-PHI Room metadata;
  diagnostic logging remains zero-PHI. Membership/notification confers no access
  grant, clinical-work authority or originating-information ownership.

\item \textbf{Standing:} U16 adjudicated candidate; exact taxonomy, membership derivation,
  API, representations, persistence, security/privacy and human-experience
  integration remain \hyperref[element:d1c4053e68384080b8eb79d2008632ba]{deferred};
  no direct Iris-to-Agora/Omilia interaction is permitted (U19).
  Partial transaction effects remain CDG-19. \hyperref[element:feea742463b84021b66853d287d8916f]{Sylloge C49} owns F31/D15
  Binary Content Management under U17; Agora retains collaboration references.
  Historical class/table names are implementation evidence.

\end{itemize}

{\small Recorded unresolved matters: canonicalName; exact scope remains in the canonical record.}\par
{\footnotesize \href{../../../../continuum/application/application-component/d1c4053e68384080b8eb79d2008632ba.md}{Canonical record: d1c4053e68384080b8eb79d2008632ba}; full qualifications and provenance.}\par
\section*{IV.3.2 Calliope}
\phantomsection
\label{navigation:IV.3.2}
\addcontentsline{toc}{section}{IV.3.2 Calliope}

\subsection*{IV.3.2.1 Application Overview}
\phantomsection
\label{navigation:IV.3.2.1}
\addcontentsline{toc}{subsection}{IV.3.2.1 Application Overview}

\subsubsection*{IV.3.2.1.1 Calliope}
\phantomsection
\label{navigation:IV.3.2.1.1}
\addcontentsline{toc}{subsubsection}{IV.3.2.1.1 Calliope}

\label{element:466ad8f75a954c0581b5bba090443849}
\index{Calliope (IV.3.2.1.1)}
{\small ApplicationComponent · Candidate — accepted Step-B baseline; final Solution Architecture unapproved · current}\par

\begin{itemize}


\item \textbf{Role/responsibility:} management/governance of foundational models, definitions
  and mutable datasets, including first-class Security and Assurance Nomenclature
  and authoritative review of ontology sets for Harmonia use. \href{\detokenize{../../../../../docs/markdown/05-information-systems-architecture/sources.md#u13}}{U13},
  \href{\detokenize{../../../../../docs/markdown/05-information-systems-architecture/calliope-models-and-datasets.md}}{canonical Calliope view}.

\item \textbf{Distinct submodules:} \textbf{Calliope-Models C41} supplies design/build-time
  Models, Metadata Types, Criteria Definitions, Topics, Filter Types, Filter Elements,
  core Harmonia element libraries and security/assurance definitions.
  \textbf{Calliope-DataSets C42} manages runtime mutable ValueSets, CodeSets, Structured
  Metadata and applicable governed ontology content. No common lifecycle or
  undifferentiated metadata repository is inferred.

\item \textbf{Authority:} development representations support type-safe human/IDE/AI-assisted
  programming; current membership/content uses runtime authority. Models/build
  definitions establish structural access contracts; DataSets do not bootstrap them.
  Definition ownership confers neither runtime execution nor originating clinical authority.

\item \textbf{Persistence/collaboration:} DataSets uses established Mneme/Mnemosyne
  responsibilities; Calliope is not a Persistence Manager. Governed workflow may
  reuse Ponos/Ergo/Praxis, without another engine or detailed management UI.
  Themis decides/enforces, Dokimasia assures and Kleio audits.

\item \textbf{Standing:} U13 settles the two responsibilities, not their packaging, runtime
  ontology boundary or representation/workflow mechanisms (CDG-24/49). Existing
  pure-library dependency constraints apply to Models; U15 resolves the
  \href{\detokenize{../../../../../docs/markdown/05-information-systems-architecture/calliope-models-and-datasets.md#calliope-realisation}}{realisation boundary}.
  General converter-support
  versus activity allocation remains CDG-05; no terminology product is selected.

\end{itemize}






\par\medskip\textbf{Controlled authorisation nomenclature}\par


\href{\detokenize{../../../../../docs/markdown/05-information-systems-architecture/sources.md#u12}}{U12} establishes these \textbf{release-bounded ValueSets}:



\begin{longtable}{>{\raggedright\arraybackslash}p{\dimexpr\linewidth/2-2\tabcolsep\relax}>{\raggedright\arraybackslash}p{\dimexpr\linewidth/2-2\tabcolsep\relax}}
ValueSet & Architectural shorthand / meaning \\
\texttt{\seqsplit{Application.Function.AuthorisationPoint.Names}} & \textbf{Function.Name}: architecturally significant application AuthorisationPoints at which enforcement is meaningful. An AuthorisationPoint identifies a governed function; it is not itself a permission. \\
\texttt{\seqsplit{Application.Function.AllowedAction.Names}} & \textbf{Function.Scope}: the bounded set of actions applicable at AuthorisationPoints. Existing concepts include \textbf{ACCESS}, \textbf{CHANGE} and \textbf{FORWARD}; these are not asserted to be the complete final ValueSet. \\
\texttt{\seqsplit{Information.Compartments}} & Governed information compartments; an applicable dimension is expressed as \textbf{Information.Compartment}. \\
\texttt{\seqsplit{Information.Classification}} & Governed information classifications. \\
\end{longtable}


\texttt{\seqsplit{Application.Function.AuthorisationPoint.Names}} SHALL be comprehensive enough
to satisfy business security requirements for the release while constrained
enough to avoid unnecessary management, development and support burden.
Complete release-specific contents remain \href{\detokenize{../../../../../docs/markdown/05-information-systems-architecture/gaps.md#identity-deferred-design}}{deferred}.



\textbf{Calliope defines/governs the nomenclature. Themis owns policy association,
resolution and enforcement.} ValueSet membership alone grants no authority
and does not establish every Cartesian combination as meaningful or permitted.



The controlled vocabularies fall within \hyperref[element:466ad8f75a954c0581b5bba090443849]{Calliope-DataSets};
their structural definitions/contracts belong to Calliope-Models. This allocation
does not reopen U12 or permit unrestricted runtime changes to release-bounded
content; detailed release/change rules remain deferred.


{\small Recorded unresolved matters: canonicalName; exact scope remains in the canonical record.}\par
{\footnotesize \href{../../../../continuum/application/application-component/466ad8f75a954c0581b5bba090443849.md}{Canonical record: 466ad8f75a954c0581b5bba090443849}; full qualifications and provenance.}\par
\subsection*{IV.3.2.2 Component Architecture}
\phantomsection
\label{navigation:IV.3.2.2}
\addcontentsline{toc}{subsection}{IV.3.2.2 Component Architecture}

\subsubsection*{IV.3.2.2.1 Calliope-Models}
\phantomsection
\label{navigation:IV.3.2.2.1}
\addcontentsline{toc}{subsubsection}{IV.3.2.2.1 Calliope-Models}

\label{element:307b3440a507403f8bc214d290bffb6a}
\index{Calliope-Models (IV.3.2.2.1)}
{\small ApplicationComponent · Candidate — accepted Step-B baseline; final Solution Architecture unapproved · current}\par

Distinct Calliope design/build-time submodule: Models, Metadata Types, Criteria
Definitions, Topics, Filter Types/Elements, core element libraries and security/
assurance nomenclature definitions. Design → Code → Build → Test → Release → Deploy.
Provides structural access contracts and useful development representations of
known runtime DataSets as normal library dependencies; those representations do
not establish current membership. Models' pure-library dependency semantics remain.
Definition ownership is not runtime operational ownership. \hyperref[element:307b3440a507403f8bc214d290bffb6a]{U13 lifecycle and authority}.


{\small Recorded unresolved matters: canonicalName; exact scope remains in the canonical record.}\par
{\footnotesize \href{../../../../continuum/application/application-component/307b3440a507403f8bc214d290bffb6a.md}{Canonical record: 307b3440a507403f8bc214d290bffb6a}; full qualifications and provenance.}\par
\subsubsection*{IV.3.2.2.2 Calliope-DataSets}
\phantomsection
\label{navigation:IV.3.2.2.2}
\addcontentsline{toc}{subsubsection}{IV.3.2.2.2 Calliope-DataSets}

\label{element:53bf6815c7a54ed0abde9b1aad311843}
\index{Calliope-DataSets (IV.3.2.2.2)}
{\small ApplicationComponent · Candidate — accepted Step-B baseline; final Solution Architecture unapproved · current}\par

Distinct Calliope runtime management/repository responsibility for mutable governed
ValueSets, CodeSets, Structured Metadata and applicable ontology content. Owns
governed content, not persistence machinery; uses established Mneme/Mnemosyne and,
where workflow is required, existing Ponos/Ergo/Praxis. Models/build definitions
establish its structural contracts. Detailed ontology boundary, workflow, UI and
physical realisation remain \hyperref[element:466ad8f75a954c0581b5bba090443849]{deferred}.


{\small Recorded unresolved matters: canonicalName; exact scope remains in the canonical record.}\par
{\footnotesize \href{../../../../continuum/application/application-component/53bf6815c7a54ed0abde9b1aad311843.md}{Canonical record: 53bf6815c7a54ed0abde9b1aad311843}; full qualifications and provenance.}\par
\section*{IV.3.3 DigitalTwin}
\phantomsection
\label{navigation:IV.3.3}
\addcontentsline{toc}{section}{IV.3.3 DigitalTwin}

\subsection*{IV.3.3.1 Application Overview}
\phantomsection
\label{navigation:IV.3.3.1}
\addcontentsline{toc}{subsection}{IV.3.3.1 Application Overview}

\subsubsection*{IV.3.3.1.1 DigitalTwin}
\phantomsection
\label{navigation:IV.3.3.1.1}
\addcontentsline{toc}{subsubsection}{IV.3.3.1.1 DigitalTwin}

\label{element:3f64d339648246e99e8866c8089f8960}
\index{DigitalTwin (IV.3.3.1.1)}
{\small ApplicationComponent · Candidate — accepted Step-B baseline; final Solution Architecture unapproved · current}\par

\begin{itemize}


\item \textbf{Role/features:} specialised Ponos-based execution component for entity-specific
  operational coordination associated with Energeia. \href{\detokenize{../../../../../docs/markdown/05-information-systems-architecture/sources.md#u02}}{U02} establishes
  specialisation/reuse, not containment by a Ponos runtime. The Digital Twin
  responsibility in \href{\detokenize{../../../../../docs/markdown/05-information-systems-architecture/sources.md#a02}}{A02} remains controlling: entity-related
  execution SHOULD use the Twin unless explicitly justified otherwise.

\item \textbf{Twin-specific execution:} one input queue containing Stimulus and Effector.
  At most one queued input is processed at a time; interested Behaviours execute
  \textbf{sequentially}, subject to the failure dispositions below. Each receives its
  own BehaviourStimulus/BehaviourEffector and fresh Praxis-defined Context after
  its predecessor's completion/save on the normal path. No shared immutable
  input-level Context.
  \href{\detokenize{../../../../../docs/markdown/05-information-systems-architecture/design-notes/DesignNote-Digital-Twin-Behaviour-and-Execution-Semantics.md}}{U19 serialisation and lifecycle}.
  Bounded resources service runnable pending work; stopped Twins retain pending
  inputs. Concurrency is across independent Twins, without a thread per Twin.

\item \textbf{Sequence/failure/recovery:} R1.x/R2.x framework loads/exposes Behaviours and
  provides/maintains the sequencing API; developer configuration supplies policy.
  Without it, order is unspecified. The applicable sequence remains stable for
  the current input; changes apply to subsequently processed inputs.
  \texttt{\seqsplit{ContinueBehaviourOnFailure}} contains failure to the current execution while
  keeping the Behaviour active and other interested Behaviours and queue processing
  unaffected. \texttt{\seqsplit{StopBehaviourOnFailure}} stops the failing Behaviour immediately,
  lets remaining interested Behaviours finish the current input, then halts the
  Twin. \texttt{\seqsplit{StopDigitalTwinOnFailure}} terminates the failing Behaviour and halts the
  Twin immediately, preventing remaining current-input and subsequent queued
  processing. Stopped Twins do not
  automatically resume: authorised administrative input through Iris/Iris-BackEnd
  or Omilia/Agora leads to Administrative Function → Administrative Ergon →
  governed Twin operation. Neither channel directly manipulates the Twin.
  Detailed recovery of a stopped Behaviour remains deferred (CDG-54).

\item \textbf{Information:} uses context drawn from governed entity information and
  relationships; entity/resource/Twin identity are not aliases. Historical patient,
  practitioner, wardsperson, organisation, care-team, ward and bed examples are
  illustrative archetypes, not an approved exhaustive component catalogue.

\item \textbf{Persistence/significant information:} definition/configuration is a C\&G
  classification candidate; operational state/activity is an O\&A candidate under
  U06. Detailed coordination-state Data Objects, activation/deactivation,
  preservation and transient/durable qualification remain CDG-07. A Twin is not
  inherently a permanent object, table or FHIR resource.

\item \textbf{Collaboration/exclusions:} executes Behaviour C25, which specialises Ergon
  and retains all common Ergo/Ergon facilities. Uses Mneme Pragma Management and
  candidate Stimulus/Effector specialisations; Ponos supplies common execution,
  Petasos supplies messaging/recoverable handoff. Follows the
  \href{\detokenize{../../../../../docs/markdown/05-information-systems-architecture/README.md#specialisation-principle}}{specialisation principle}; adds genuinely
  Twin-specific semantics instead of parallel platform mechanisms.
  Not a second engine, unmanaged persistent actor, separate
  policy/search service or authority over originating clinical facts/clinical work.

\item \textbf{Standing:} U02 specialisation and U19 Behaviour execution semantics established;
  representation/runtime details remain CDG-07/54. DigitalTwin and Dokimasia are sibling
  Ponos specialisations; neither contains the other or manages Pragma by execution.

\end{itemize}

{\small Recorded unresolved matters: canonicalName; exact scope remains in the canonical record.}\par
{\footnotesize \href{../../../../continuum/application/application-component/3f64d339648246e99e8866c8089f8960.md}{Canonical record: 3f64d339648246e99e8866c8089f8960}; full qualifications and provenance.}\par
\subsection*{IV.3.3.2 Component Architecture}
\phantomsection
\label{navigation:IV.3.3.2}
\addcontentsline{toc}{subsection}{IV.3.3.2 Component Architecture}

\subsubsection*{IV.3.3.2.1 Behaviour}
\phantomsection
\label{navigation:IV.3.3.2.1}
\addcontentsline{toc}{subsubsection}{IV.3.3.2.1 Behaviour}

\label{element:a26f06d66bec4981b266e202c46a9afb}
\index{Behaviour (IV.3.3.2.1)}
{\small ApplicationComponent · Candidate — accepted Step-B baseline; final Solution Architecture unapproved · current}\par

\begin{itemize}


\item \textbf{Realisation:} each DigitalTwin has **Behaviour [1..*]**. Each Behaviour is
  realised through \textbf{Praxis → Pathway → Ergo [exactly 1]}. A one-Ergo Praxis
  retains required Context, expected criteria/outcomes, fulfilment semantics and
  housekeeping; it is neither an orchestrator nor a choreographer.
  \href{\detokenize{../../../../../docs/markdown/05-information-systems-architecture/design-notes/DesignNote-Digital-Twin-Behaviour-and-Execution-Semantics.md}}{U19 DesignNote}.

\item \textbf{Role/features:} Ergon specialisation for DigitalTwin operation. It inherits/
  reuses Context Loading and Context Unloading, process-management mechanisms,
  Petasos/platform integration, assurance integration, logging/observability,
  reporting, checkpoint/platform integration and other common Ergo/Ergon facilities.
  \href{\detokenize{../../../../../docs/markdown/05-information-systems-architecture/sources.md#u02}}{U02}, \href{\detokenize{../../../../../docs/markdown/05-information-systems-architecture/README.md#specialisation-principle}}{common specialisation principle}.

\item \textbf{Information/persistence/collaboration:} receives its own identifiable
  \textbf{BehaviourStimulus / BehaviourEffector} derived from an input in which it has
  registered interest. Operates relative to the Twin's singular represented
  entity. Its Praxis-defined Context is loaded fresh after the preceding
  Behaviour completes and saves its consequences; Mneme manages Load/Save and
  Pragma, Mnemosyne establishes durable state and Petasos handles handoff.

\item \textbf{Execution/output/failure:} interested Behaviours execute sequentially for
  each input. Outputs may enter normal governed processing, resolving Twin-relative
  references only where required. Same-Twin Stimulus/Effector outputs join its
  normal queue without recursive interruption. The developer's
  \texttt{\seqsplit{ContinueBehaviourOnFailure}} contains failure to the current execution; the
  Behaviour remains active, other interested Behaviours and normal queue processing
  continue. \texttt{\seqsplit{StopBehaviourOnFailure}} stops the failing Behaviour immediately,
  lets remaining interested Behaviours complete the current input, then halts
  the Twin. \texttt{\seqsplit{StopDigitalTwinOnFailure}} terminates the failing Behaviour and halts
  the Twin immediately; remaining current-input and subsequent queued processing
  stop. Sequencing policy is developer-owned through the framework API; the
  applicable sequence stays stable throughout an input, and configuration changes
  apply to subsequent inputs. No priority metadata or Pathway orchestration (U19).

\item \textbf{Standing/exclusions:} discrete candidate activity specialisation, not an
  independent activity framework parallel to Ergo/Ergon. DigitalTwin adds only
  genuinely Twin-specific behaviour instead of duplicating platform mechanisms;
  Behaviour acquires no standalone Information Resource merely by being activity.

\end{itemize}

{\small Recorded unresolved matters: canonicalName; exact scope remains in the canonical record.}\par
{\footnotesize \href{../../../../continuum/application/application-component/a26f06d66bec4981b266e202c46a9afb.md}{Canonical record: a26f06d66bec4981b266e202c46a9afb}; full qualifications and provenance.}\par
\section*{IV.3.4 Dokimasia}
\phantomsection
\label{navigation:IV.3.4}
\addcontentsline{toc}{section}{IV.3.4 Dokimasia}

\subsection*{IV.3.4.1 Application Overview}
\phantomsection
\label{navigation:IV.3.4.1}
\addcontentsline{toc}{subsection}{IV.3.4.1 Application Overview}

\subsubsection*{IV.3.4.1.1 Dokimasia}
\phantomsection
\label{navigation:IV.3.4.1.1}
\addcontentsline{toc}{subsubsection}{IV.3.4.1.1 Dokimasia}

\label{element:d739c75dffe2453081735710b285f03b}
\index{Dokimasia (IV.3.4.1.1)}
{\small ApplicationComponent · Candidate — accepted Step-B baseline; final Solution Architecture unapproved · current}\par

\begin{itemize}


\item \textbf{Role/features:} specialised Ponos-based execution component associated with Energeia
  coordinating independent governed assurance activity over a subject. Journey
  C35 is the common observed-progression evidence subject/substrate for Pathway
  and Pub/Sub
  assurance under U04/U05; this does not limit all assurance subjects to Journey.
  \href{\detokenize{../../../../../docs/markdown/05-information-systems-architecture/sources.md#u02}}{U02} settles the candidate specialisation relationship; actual
  Assurance Praxis evaluates available
  evidence against applicable basis and produces governed assurance information.
  \href{\detokenize{../../../../../docs/markdown/05-information-systems-architecture/sources.md#a04}}{A04 §2}, \href{\detokenize{../../../../../docs/markdown/05-information-systems-architecture/sources.md#l31}}{L31}. \textbf{Praxis is the assurance basis;
  Journey is the observed evidence subject} under U06. Pathway/Route expectations
  contribute through applicable Praxis rather than constituting the general basis.

\item \textbf{Information:} evidence/context/basis and assurance assessment/finding/outcome
  areas are documented. Journey observations remain distinct from inferred
  interpretation and assurance findings. Dokimasia activity/findings/results are
  O\&A classification candidates under U06; detailed Data Objects, identity and
  preservation remain unresolved. Candidate vocabulary is not promoted into an
  Information Family.

\item \textbf{Activity/collaboration:} executes discrete Dokimasia Ergon C26 specialisations,
  including candidate AI Assurance Ergon C27. Reuses Ponos execution, Ergo/Ergon
  activity, Mneme Pragma Management, process-management, platform integration,
  reporting and observability. This is specialisation/reuse, not runtime containment
  or a parallel assurance execution platform.

\item \textbf{Information/persistence:} corresponding Pragma specialisation semantics are
  required; assurance request/input and assurance result/finding are candidate
  roles whose names, schemas, cardinalities and lifecycle remain unresolved.
  Evidence-access/control independence, detailed workflow composition and durable
  structures remain CDG-08/09. Common framework reuse does not let the assured
  activity determine, suppress, manufacture or alter its own assurance outcome.

\item \textbf{Exclusions:} evidence, workflow completion and operational status are not
  assurance; custody does not make Kleio the evaluator. Dokimasia does not manage
  operational subject escalation/remediation and is not an established Twin subtype.
  AI invocation confers no authority on the model/service. Dokimasia can manage
  its own assurance work while the subject retains its operational management.

\item \textbf{Standing:} conceptual independence remains established by A04; execution
  and discrete activity responsibilities are adjudicated for the candidate by U02.
  Detailed AI architecture requires CDG-35 discovery. No upstream responsibility
  is reopened; detailed Praxis/Journey and assurance contracts remain unresolved.

\end{itemize}

{\small Recorded unresolved matters: canonicalName; exact scope remains in the canonical record.}\par
{\footnotesize \href{../../../../continuum/application/application-component/d739c75dffe2453081735710b285f03b.md}{Canonical record: d739c75dffe2453081735710b285f03b}; full qualifications and provenance.}\par
\subsection*{IV.3.4.2 Component Architecture}
\phantomsection
\label{navigation:IV.3.4.2}
\addcontentsline{toc}{subsection}{IV.3.4.2 Component Architecture}

\subsubsection*{IV.3.4.2.1 Dokimasia Ergon}
\phantomsection
\label{navigation:IV.3.4.2.1}
\addcontentsline{toc}{subsubsection}{IV.3.4.2.1 Dokimasia Ergon}

\label{element:34d7c844da3344a09e96ad821ff9fcea}
\index{Dokimasia Ergon (IV.3.4.2.1)}
{\small ApplicationComponent · Candidate — accepted Step-B baseline; final Solution Architecture unapproved · current}\par

\begin{itemize}


\item \textbf{Role/features:} discrete Ergon specialisation representing a bounded assurance
  activity executed within a Dokimasia workflow. Reuses normal Ergo/Ergon
  execution, Context Loading/Unloading, process-management, integration,
  checkpoint, logging/observability and reporting contracts. \href{\detokenize{../../../../../docs/markdown/05-information-systems-architecture/sources.md#u02}}{U02}.

\item \textbf{Information/persistence/collaboration:} uses governed assurance input,
  applicable basis and evidence, and may produce assurance information through
  Mneme Pragma Management. No detailed schema, evidence-access API, storage or
  additional Information Resource is inferred. Executed through Dokimasia's
  specialised Ponos model.

\item \textbf{Standing/exclusions:} discrete candidate component responsibility established;
  subject to \href{\detokenize{../../../../../docs/markdown/05-information-systems-architecture/component-register.md#dokimasia-terminal-boundaries}}{Dokimasia terminal constraints}.
  AI Assurance Ergon C27 specialises this component. Other potential assurance
  classifications, including deterministic/human subclasses, require adjudication.

\end{itemize}

{\small Recorded unresolved matters: canonicalName; exact scope remains in the canonical record.}\par
{\footnotesize \href{../../../../continuum/application/application-component/34d7c844da3344a09e96ad821ff9fcea.md}{Canonical record: 34d7c844da3344a09e96ad821ff9fcea}; full qualifications and provenance.}\par
\subsubsection*{IV.3.4.2.2 AI Assurance Ergon}
\phantomsection
\label{navigation:IV.3.4.2.2}
\addcontentsline{toc}{subsubsection}{IV.3.4.2.2 AI Assurance Ergon}

\label{element:8fe50975372b4d399a4bcb9f3d9348f7}
\index{AI Assurance Ergon (IV.3.4.2.2)}
{\small ApplicationComponent · Candidate — accepted Step-B baseline; final Solution Architecture unapproved · current}\par

\begin{itemize}


\item \textbf{Role/features:} candidate Dokimasia Ergon specialisation where AI-based
  processing is itself an explicit governed participant in Dokimasia workflow.
  Its existence and discrete responsibility are established by \href{\detokenize{../../../../../docs/markdown/05-information-systems-architecture/sources.md#u02}}{U02};
  its detailed Information Systems architecture is not finalised.

\item \textbf{Distinction:} ordinary Ergon may call AI internally as an implementation
  choice. AI Assurance Ergon makes that AI activity an explicit governed assurance
  workflow activity. This allocates neither every use of AI to this component
  nor a mandatory AI step to every Dokimasia workflow.

\item \textbf{Information/persistence/collaboration:} inherits C26/C23/C05 contracts and
  common Pragma Management, process-management, platform integration and reporting.
  AI invocation confers no originating domain or assurance-determination authority
  on the service/model. Context, disclosure, results, provenance/evidence and
  detailed workflow interactions require \href{\detokenize{../../../../../docs/markdown/05-information-systems-architecture/gaps.md#dokimasia-ai-discovery}}{CDG-35 discovery}.

\item \textbf{Standing/exclusions:} established component boundary with candidate detailed
  design; no LLM, agent framework, SDK, model, provider or service technology
  selected. Discovery precedes final implementation modelling. Inherits both
  Dokimasia terminal rules; unresolved design does not invalidate its boundary.

\end{itemize}

{\small Recorded unresolved matters: canonicalName; exact scope remains in the canonical record.}\par
{\footnotesize \href{../../../../continuum/application/application-component/8fe50975372b4d399a4bcb9f3d9348f7.md}{Canonical record: 8fe50975372b4d399a4bcb9f3d9348f7}; full qualifications and provenance.}\par
\section*{IV.3.5 Energeia}
\phantomsection
\label{navigation:IV.3.5}
\addcontentsline{toc}{section}{IV.3.5 Energeia}

\subsection*{IV.3.5.1 Application Overview}
\phantomsection
\label{navigation:IV.3.5.1}
\addcontentsline{toc}{subsection}{IV.3.5.1 Application Overview}

\subsubsection*{IV.3.5.1.1 Energeia}
\phantomsection
\label{navigation:IV.3.5.1.1}
\addcontentsline{toc}{subsubsection}{IV.3.5.1.1 Energeia}

\label{element:fb3f37a6ac614415ac486a6de9d35994}
\index{Energeia (IV.3.5.1.1)}
{\small ApplicationComponent · Candidate — accepted Step-B baseline; final Solution Architecture unapproved · current}\par

\begin{itemize}


\item \textbf{Role/responsibility:} discrete Software Component owning Harmonia operational
  management: “Is Harmonia itself operating as expected?” It observes work-to-be-done,
  work-being-done and work-done, outstanding ActionableTasks, FulfillmentTask
  activity, Journey progression, stalled/incomplete activity, Praxis context,
  abnormal behaviour and operational information for reporting/assurance.

\item \textbf{Candidate services/features:} \textbf{Praxis Management Services} (F11) and
  \textbf{Journey Management Services} (F38), plus operational observation
  and reporting (F39). Detailed Feature/Service/internal-construct/subordinate-
  Component decomposition remains unresolved (CDG-39); service-set names do not
  manufacture independently deployable components.

\item \textbf{Information:} owns Praxis C06 and Journey C35 at the application level.
  Mneme manages their governed information and retains ActionableTask,
  FulfillmentTask and Pragma management. Observing operational significance does
  not transfer management of that information or referenced clinical facts.

\item \textbf{Persistence:} Mnemosyne C\&G provides Praxis/configuration durability;
  Mnemosyne O\&A provides Journey/operational-assurance durability, through Mneme
  governed management. Energeia application ownership is distinct from durable
  persistence responsibility; physical structures remain unresolved.

\item \textbf{Collaboration/relationships:} associated execution/activity responsibilities
  retain U02 specialisation/reuse. Praxis aggregates Pathway [0..1] and has a
  working Aggregation relationship to Journey with cardinalities unresolved.
  Dokimasia uses Praxis as assurance basis and Journey as observed evidence
  subject, preserving independence and terminal rules. Iris presents operational
  views; presentation does not own Harmonia operational management.

\item \textbf{Release boundary:} for \textbf{R1.x/R2.x, Observe → Report}. Energeia
  SHALL NOT autonomously change execution, redistribute workload, modify Routes,
  Pathways, Pub/Sub or Praxis, intervene in resilience or adjust capacity merely
  from workload, failure, Journey, capacity, assurance or anomaly observations.
  Existing explicitly governed execution/recovery behaviour retains its own
  authority; observation grants no new adaptive authority.

\item \textbf{Separate management plane:} external IT/Ops manages infrastructure, network,
  containers/Kubernetes, JVM/runtime, databases, technical health/capacity and
  platform resilience. Energeia manages/observes Harmonia as Harmonia. This
  distinction selects no deployment topology or infrastructure integration.

\item \textbf{Deferred direction:} \textbf{R3 Adaptive Load and Resilience Service} (F40):
  Observe → Analyse → Decide → Adapt, with feedback to observation. R1.x/R2.x
  observability should permit later capability without embedding it now.
  Policy/control loops, algorithms, mechanisms, infrastructure integration and
  AI decision architecture are deliberately undesigned (CDG-41).

\item \textbf{Standing:} U06 adjudicated these candidate responsibilities; it did not
  itself establish whole-candidate acceptance or Step B. Current A4/B1 standing
  is recorded in the \href{\detokenize{../../../../../docs/markdown/05-information-systems-architecture/README.md}}{orientation}. Possible upstream omissions require
  later review, without deriving Business Architecture or capability demand here.

\end{itemize}

{\small Recorded unresolved matters: canonicalName; exact scope remains in the canonical record.}\par
{\footnotesize \href{../../../../continuum/application/application-component/fb3f37a6ac614415ac486a6de9d35994.md}{Canonical record: fb3f37a6ac614415ac486a6de9d35994}; full qualifications and provenance.}\par
\subsection*{IV.3.5.2 Information Architecture}
\phantomsection
\label{navigation:IV.3.5.2}
\addcontentsline{toc}{subsection}{IV.3.5.2 Information Architecture}

\subsubsection*{IV.3.5.2.1 Praxis}
\phantomsection
\label{navigation:IV.3.5.2.1}
\addcontentsline{toc}{subsubsection}{IV.3.5.2.1 Praxis}

\label{element:0f961462340348b290fd844357d54d67}
\index{Praxis (IV.3.5.2.1)}
{\small DataObject · Candidate — accepted Step-B baseline; final Solution Architecture unapproved · current}\par

\begin{itemize}


\item \textbf{Role/ownership:} higher-order operational construct defining expected/permitted
  context and the operational/assurance envelope. \textbf{Energeia C39 owns Praxis}
  and provides \textbf{Praxis Management Services F11} under \href{\detokenize{../../../../../docs/markdown/05-information-systems-architecture/sources.md#u06}}{U06}.
  It is not another execution engine or necessarily an independent component.

\item \textbf{Fulfilment semantics:} Pathway defines the conceptual Ergo sequence; Praxis
  defines how it is fulfilled: applicable Context definitions required by the
  Ergo framework, segment criteria/outcomes, continuation derivation semantics
  and other Praxis-level housekeeping. Ponos/Ergo performs output processing and
  routing. Context requirements are concrete in R1.x/R2.x; criteria-driven dynamic
  Context and dynamically configured reusable/generic Ergo remain deferred R3.x
  direction, without an AI-driven selection/configuration mechanism.

\item \textbf{Information/relationships:} candidate interface configuration, Topic exclusions
  and potentially inclusions, operational constraints/configuration and assurance
  rules. \textbf{Aggregation: Pathway [0..1]}; working \textbf{Aggregation: Journey} with
  cardinalities unresolved. No Composition or shared lifecycle is inferred.
  Praxis may exist without a Pathway for Publish/Subscribe activity.
  The specialised \hyperref[element:a26f06d66bec4981b266e202c46a9afb]{Digital Twin Behaviour pattern} uses a Pathway with
  exactly one Ergo; it retains full Praxis semantics and does not change the
  general Aggregation cardinality.

\item \textbf{Management/persistence:} Mneme manages governed representations; Mnemosyne
  \textbf{C\&G} provides durable persistence responsibility. Ponos consumes Praxis for
  execution and Dokimasia uses it as the primary normative assurance basis;
  neither use nor persistence transfers application ownership.

\item \textbf{Retained detail:} schema, lifecycle/versioning, activation/deactivation,
  replacement/retirement and in-flight effects remain CDG-01/06. Historical
  TaskSequence/PraxisDefinition/PraxisImplementation and A04 actual-assurance-
  activity descriptions do not establish equivalence or a definition-to-instance
  model. U06 supersedes the earlier wholly unresolved broader-Praxis standing.

\item \textbf{Logical boundary:} \href{\detokenize{../../../../../docs/markdown/05-information-systems-architecture/component-register.md#c46}}{Praxis.Horos C46}, with IngressHoros/EgressHoros,
  names the Praxis boundary. EgressHoros hands interaction state/outcome back
  before Ergo.ContextUnloader completes; it is not a technical endpoint.

\item \textbf{Exclusions:} Praxis is not Pathway, Route, a sequence of Ergo/Ergon, an
  orchestrator, choreographer, runtime engine or Pub/Sub graph. Mneme semantic
  selection, Petasos concrete persistent handoff and Ponos execution remain distinct.
  No detailed Praxis schema is invented.

\end{itemize}

{\small Recorded unresolved matters: canonicalName; exact scope remains in the canonical record.}\par
{\footnotesize \href{../../../../continuum/application/data-object/0f961462340348b290fd844357d54d67.md}{Canonical record: 0f961462340348b290fd844357d54d67}; full qualifications and provenance.}\par
\subsubsection*{IV.3.5.2.2 Journey}
\phantomsection
\label{navigation:IV.3.5.2.2}
\addcontentsline{toc}{subsubsection}{IV.3.5.2.2 Journey}

\label{element:359b901df27f4b43b5a2ff14a8484e34}
\index{Journey (IV.3.5.2.2)}
{\small DataObject · Candidate — accepted Step-B baseline; final Solution Architecture unapproved · current}\par

First-class observational construct / candidate Application Data Object: governed
observed progression reconstructed from available execution evidence, whether
Pathway-controlled, Pub/Sub or mixed. It may be graph-shaped and have no Pathway.
Primary reporting, analysis and assurance subject; Dokimasia may
interpret it under applicable basis without replacing observations with inference.
\textbf{Energeia owns Journey and Journey Management Services F38} under U06. Mneme
manages governed information; Mnemosyne \textbf{O\&A} provides durable persistence
responsibility. Journey aggregates/references observations/information from Mneme,
Petasos, Ponos, Ergon, DigitalTwin and other components/evidence mechanisms where
appropriate, retaining their ownership. Comparatively small discrete Journey
information may describe substantial activity/evidence: aggregate semantically,
without indiscriminate duplication. Detailed representation/construction, lifecycle,
retention, graph/evidence/causal/correlation/reference schema and inference contracts
remain unresolved (\href{\detokenize{../../../../../docs/markdown/05-information-systems-architecture/gaps.md#journey-contracts}}{CDG-38}); Journey is not a shard key.
Journey is neither an engine, AI narrative nor automatic permanent AuditEvent.
For Pathway-based Praxis fulfilment, continuation-Pragma lineage provides an
important observable progression thread. \textbf{Journey lineage is not complete
Pragma provenance}: independent outputs may establish their own activity,
Pub/Sub interactions, Praxes and Journeys. Mneme may manage richer relationships.
Journey is not a complete provenance graph, execution log, audit trail or Petasos
queue history. \hyperref[element:359b901df27f4b43b5a2ff14a8484e34]{U18 reconciliation}.


{\small Recorded unresolved matters: canonicalName; exact scope remains in the canonical record.}\par
{\footnotesize \href{../../../../continuum/application/data-object/359b901df27f4b43b5a2ff14a8484e34.md}{Canonical record: 359b901df27f4b43b5a2ff14a8484e34}; full qualifications and provenance.}\par
\section*{IV.3.6 Harmonia.Omilia.Client}
\phantomsection
\label{navigation:IV.3.6}
\addcontentsline{toc}{section}{IV.3.6 Harmonia.Omilia.Client}

\subsection*{IV.3.6.1 Application Overview}
\phantomsection
\label{navigation:IV.3.6.1}
\addcontentsline{toc}{subsection}{IV.3.6.1 Application Overview}

\subsubsection*{IV.3.6.1.1 Harmonia.Omilia.Client}
\phantomsection
\label{navigation:IV.3.6.1.1}
\addcontentsline{toc}{subsubsection}{IV.3.6.1.1 Harmonia.Omilia.Client}

\label{element:33b8421ffc444b78b2d059201a58a688}
\index{Harmonia.Omilia.Client (IV.3.6.1.1)}
{\small ApplicationComponent · Candidate — accepted Step-B baseline; final Solution Architecture unapproved · current}\par

\begin{itemize}


\item \textbf{Role/responsibility:} Harmonia collaboration client providing the associated
  collaboration user experience; communicates with Omilia.Server using Matrix
  protocol. Presentation does not confer information or collaboration-control
  authority.

\item \textbf{Provenance/standing:} derived from Element; U16 adjudicated candidate.
  Detailed modifications and human-experience integration remain deferred;
  Iris has no direct Agora/Omilia interaction (U19).

\end{itemize}

{\small Recorded unresolved matters: canonicalName; exact scope remains in the canonical record.}\par
{\footnotesize \href{../../../../continuum/application/application-component/33b8421ffc444b78b2d059201a58a688.md}{Canonical record: 33b8421ffc444b78b2d059201a58a688}; full qualifications and provenance.}\par
\section*{IV.3.7 Harmonia.Omilia.Server}
\phantomsection
\label{navigation:IV.3.7}
\addcontentsline{toc}{section}{IV.3.7 Harmonia.Omilia.Server}

\subsection*{IV.3.7.1 Application Overview}
\phantomsection
\label{navigation:IV.3.7.1}
\addcontentsline{toc}{subsection}{IV.3.7.1 Application Overview}

\subsubsection*{IV.3.7.1.1 Harmonia.Omilia.Server}
\phantomsection
\label{navigation:IV.3.7.1.1}
\addcontentsline{toc}{subsubsection}{IV.3.7.1.1 Harmonia.Omilia.Server}

\label{element:369907025d06492bb414a72cea834a77}
\index{Harmonia.Omilia.Server (IV.3.7.1.1)}
{\small ApplicationComponent · Candidate — accepted Step-B baseline; final Solution Architecture unapproved · current}\par

\begin{itemize}


\item \textbf{Role/responsibility:} owns/performs underlying Room/Space, user, membership,
  dialogue/message, collaboration-event and applicable collaboration-media
  operations. Exposes the \textbf{Application Services API} for Agora operation requests
  and dialogue/activity delivery from Agora-managed Rooms/Spaces.

\item \textbf{Boundary:} operational worker, not owner of Harmonia collaboration semantics,
  membership intent or originating clinical/operational information. Agora owns
  native collaboration information; AX-05 durable-state responsibility remains.

\item \textbf{Provenance/standing:} derived from Matrix Synapse; U16 adjudicated candidate.
  API technology, modifications and persistence mechanics are deferred, not
  inferred from historical Synapse deployment/configuration.

\end{itemize}

{\small Recorded unresolved matters: canonicalName; exact scope remains in the canonical record.}\par
{\footnotesize \href{../../../../continuum/application/application-component/369907025d06492bb414a72cea834a77.md}{Canonical record: 369907025d06492bb414a72cea834a77}; full qualifications and provenance.}\par
\subsection*{IV.3.7.2 Interfaces}
\phantomsection
\label{navigation:IV.3.7.2}
\addcontentsline{toc}{subsection}{IV.3.7.2 Interfaces}

\subsubsection*{IV.3.7.2.1 Omilia Application Services API}
\phantomsection
\label{navigation:IV.3.7.2.1}
\addcontentsline{toc}{subsubsection}{IV.3.7.2.1 Omilia Application Services API}

\label{element:49c52c0466bc4760a11e0ca6fe8142c0}
\index{Omilia Application Services API (IV.3.7.2.1)}
{\small ApplicationInterface · Candidate — adjudicated Domain05 Step-B baseline; final Solution Architecture unapproved · current}\par

Omilia.Server exposes an \textbf{Application Services API} used by Agora both to
request collaboration operations and to receive managed-Room dialogue/activity.



\begin{quote}\small\ttfamily
flowchart LR\\
    Agora["Harmonia.Agora: collaboration control"] -->|"Application Services API: operation requests"| Server["Harmonia.Omilia.Server: collaboration operations"]\\
    Server -->|"Application Services API: managed-Room dialogue / activity"| Agora\\
    Client["Harmonia.Omilia.Client: collaboration experience"] <-->|"Matrix protocol"| Server\\

\end{quote}


Agora decides to create/change a managed Room/Space, add/remove a participant or
publish information; Omilia.Server performs the operation. Room activity returns
through the same architectural boundary. The exact API contract, technology and
mapping to upstream Matrix interfaces remain deferred. Historical Admin,
Client-Server and AS-transaction endpoints do not define this new contract.



Workflow/execution collaboration continues through Petasos with \textbf{no direct
Agora-to-Ponos dependency}. Matrix DTO encapsulation, Themis default-deny
governance and zero-PHI Room aliases/names/topics remain repository constraints.
Membership or receipt of a notification does not confer information-access rights.



\href{\detokenize{../../../../../docs/markdown/05-information-systems-architecture/sources.md#u19}}{U19} establishes the administrative interaction boundary:
Administrator → Omilia → Agora → Administrative Function → specialised
Administrative Ergon → governed subsystem operation. Stopped-Twin resumption
uses this pattern, with normal authorisation, Context, execution and Kleio audit
controls; Agora and Omilia do not directly resume or manipulate the Twin.
The arrows describe semantic progression and retain Petasos execution handoff
and Agora's no-direct-Ponos-dependency rule. Iris provides a separate path through
Iris-BackEnd and does \textbf{not} directly interact with Agora or Omilia. Broader
Administrative Function reconciliation is \href{\detokenize{../../../../../docs/markdown/05-information-systems-architecture/gaps.md#administrative-function-finding}}{deferred}.



Identity mapping/provisioning is collaboration integration, not authentication
authority. Existing isolation constraints are not relaxed; detailed security,
privacy and external/federated interactions remain deferred.


{\small Recorded unresolved matters: canonicalName; exact scope remains in the canonical record.}\par
{\footnotesize \href{../../../../continuum/application/application-interface/49c52c0466bc4760a11e0ca6fe8142c0.md}{Canonical record: 49c52c0466bc4760a11e0ca6fe8142c0}; full qualifications and provenance.}\par
\section*{IV.3.8 Harmonia.Sylloge}
\phantomsection
\label{navigation:IV.3.8}
\addcontentsline{toc}{section}{IV.3.8 Harmonia.Sylloge}

\subsection*{IV.3.8.1 Application Overview}
\phantomsection
\label{navigation:IV.3.8.1}
\addcontentsline{toc}{subsection}{IV.3.8.1 Application Overview}

\subsubsection*{IV.3.8.1.1 Harmonia.Sylloge}
\phantomsection
\label{navigation:IV.3.8.1.1}
\addcontentsline{toc}{subsubsection}{IV.3.8.1.1 Harmonia.Sylloge}

\label{element:feea742463b84021b66853d287d8916f}
\index{Harmonia.Sylloge (IV.3.8.1.1)}
{\small ApplicationComponent · Candidate — accepted Step-B baseline; final Solution Architecture unapproved · current}\par

\begin{itemize}


\item \textbf{Role/responsibility:} \textbf{Binary Content Management subsystem} responsible
  for governed association, resolution, retrieval and assembly of binary content
  associated with Harmonia-managed information. \hyperref[element:feea742463b84021b66853d287d8916f]{U17 canonical view}.

\item \textbf{Feature/information:} F31/D15; binary-content metadata, Information Unit
  associations, location/resolution knowledge, retrieval, on-demand assembly from
  one or more payloads and knowledge of Harmonia-managed versus externally
  managed content. Logical information and binary payload remain distinct.

\item \textbf{Custody/persistence:} physical custody may remain external; management does
  not require copying every payload into Harmonia. Where applicable, Mnemosyne
  provides Harmonia persistence services without acquiring Binary Content
  Management. Sylloge is not a generic persistence manager.

\item \textbf{Collaborations/exclusions:} Mneme remains Context and Pragma Manager with its
  established service boundary; applicable processing components may obtain/resolve
  content through Sylloge. Agora retains native collaboration media and the
  record/reference/activity concerning formal binary content; Sylloge retains
  that content's Binary Content Management. Omilia acquires no such responsibility.
  Source information/domain authority and external physical custody are preserved.

\item \textbf{Standing:} U17 adjudicated Candidate Architecture; former A2 responsibility
  question CDG-25 closed. \hyperref[element:feea742463b84021b66853d287d8916f]{Storage, representations, APIs, algorithms and lifecycle
  mechanisms} remain
  non-blocking later design. No module, dependency or deployment allocation.

\end{itemize}

{\small Recorded unresolved matters: canonicalName; exact scope remains in the canonical record.}\par
{\footnotesize \href{../../../../continuum/application/application-component/feea742463b84021b66853d287d8916f.md}{Canonical record: feea742463b84021b66853d287d8916f}; full qualifications and provenance.}\par
\section*{IV.3.9 Iris}
\phantomsection
\label{navigation:IV.3.9}
\addcontentsline{toc}{section}{IV.3.9 Iris}

\subsection*{IV.3.9.1 Application Overview}
\phantomsection
\label{navigation:IV.3.9.1}
\addcontentsline{toc}{subsection}{IV.3.9.1 Application Overview}

\subsubsection*{IV.3.9.1.1 Iris}
\phantomsection
\label{navigation:IV.3.9.1.1}
\addcontentsline{toc}{subsubsection}{IV.3.9.1.1 Iris}

\label{element:e927c1af96c64fcd979be659bee2cbe3}
\index{Iris (IV.3.9.1.1)}
{\small ApplicationComponent · Candidate — accepted Step-B baseline; final Solution Architecture unapproved · current}\par

\begin{itemize}


\item \textbf{Role/responsibility:} contextual human interaction, presentation, input and
  governed requests. \href{\detokenize{../../../../../docs/markdown/05-information-systems-architecture/sources.md#a02}}{A02}, \href{\detokenize{../../../../../docs/markdown/05-information-systems-architecture/sources.md#l11}}{L11/L21/L23}.

\item \textbf{Features/contained constructs:} clinical, operations and administration
  applications (C16), mediation (C17), shared shell/design system and role-tailored
  workspaces. These are documented presentation responsibilities, not product modules
  promoted solely by their directory names.

\item \textbf{Information/persistence:} displays contextual resource and operational views,
  captures user intent/assertions, transient browser selections/filters/state.
  No authoritative clinical/registry or persistence ownership; authority of human
  input comes from the actor/process/source, not Iris.

\item \textbf{Collaboration/exclusions:} uses defined governed information/activity interfaces
  and Themis. No raw-cache/database alternative, JPA/SQL, backend registry validation,
  autonomous workflow engine or clinical decision authority.
  Iris does not directly interact with Agora or Omilia. Administrative intent
  passes through Iris-BackEnd to an Administrative Function and specialised
  Administrative Ergon, not direct subsystem manipulation (U19).

\item \textbf{Standing:} established high-level responsibility; candidate detailed views,
  CDG-11/24/26/28. Diagnostic zero-PHI controls do not turn clinical presentation
  into a technical-identifiers-only application.

\end{itemize}

{\small Recorded unresolved matters: canonicalName; exact scope remains in the canonical record.}\par
{\footnotesize \href{../../../../continuum/application/application-component/e927c1af96c64fcd979be659bee2cbe3.md}{Canonical record: e927c1af96c64fcd979be659bee2cbe3}; full qualifications and provenance.}\par
\subsection*{IV.3.9.2 Component Architecture}
\phantomsection
\label{navigation:IV.3.9.2}
\addcontentsline{toc}{subsection}{IV.3.9.2 Component Architecture}

\subsubsection*{IV.3.9.2.1 Iris backend-for-frontend mediation}
\phantomsection
\label{navigation:IV.3.9.2.1}
\addcontentsline{toc}{subsubsection}{IV.3.9.2.1 Iris backend-for-frontend mediation}

\label{element:37b82d161c5640298e951acc361b338a}
\index{Iris backend-for-frontend mediation (IV.3.9.2.1)}
{\small ApplicationComponent · Candidate — accepted Step-B baseline; final Solution Architecture unapproved · current}\par

\begin{itemize}


\item \textbf{Role/features:} decoupled presentation service mediation and clinical versus
  operations request separation; telemetry aggregation and provider contracts.
  \href{\detokenize{../../../../../docs/markdown/05-information-systems-architecture/sources.md#l11}}{L11/L21/L23}.

\item \textbf{Information/contained constructs:} projected resources, task/sequence views,
  subsystem status, dependency/queue/gateway observations and alert diagnostics;
  operations aggregator and health-provider support, not a new information authority.

\item \textbf{Persistence/collaboration:} uses Mneme-facing governed access, Themis and
  approved execution requests; backend/presentation state is not durable truth.
  Documented raw Hot Rod/cache paths remain conflicting (CDG-11).
  \textbf{Iris-BackEnd} mediates administrative intent into governed Administrative
  Function/Administrative Ergon execution. It does not directly resume a stopped
  Twin or provide a privileged subsystem backdoor (U19).

\item \textbf{Exclusions/standing:} no direct JPA/PostgreSQL or server-side registry governance.
  Candidate mediation component; ports/server frameworks/provider class catalogues
  are support details, and telemetry persistence/provenance remains CDG-28.

\end{itemize}

{\small Recorded unresolved matters: canonicalName; exact scope remains in the canonical record.}\par
{\footnotesize \href{../../../../continuum/application/application-component/37b82d161c5640298e951acc361b338a.md}{Canonical record: 37b82d161c5640298e951acc361b338a}; full qualifications and provenance.}\par
\section*{IV.3.10 Kleio}
\phantomsection
\label{navigation:IV.3.10}
\addcontentsline{toc}{section}{IV.3.10 Kleio}

\subsection*{IV.3.10.1 Application Overview}
\phantomsection
\label{navigation:IV.3.10.1}
\addcontentsline{toc}{subsection}{IV.3.10.1 Application Overview}

\subsubsection*{IV.3.10.1.1 Kleio}
\phantomsection
\label{navigation:IV.3.10.1.1}
\addcontentsline{toc}{subsubsection}{IV.3.10.1.1 Kleio}

\label{element:fefda110e2274884ad58906733606efd}
\index{Kleio (IV.3.10.1.1)}
{\small ApplicationComponent · Candidate — accepted Step-B baseline; final Solution Architecture unapproved · current}\par

\begin{itemize}


\item \textbf{Role/responsibility:} \textbf{Harmonia's authoritative Audit Information Management
  component.} Kleio owns authoritative Audit Capture, creation/management of
  authoritative audit information, canonical FHIR \textbf{AuditEvent} representation,
  preservation/retrieval of audit history, semantically significant user/system
  audit activity and provision/preparation of historical audit context for
  operational and assurance use. \href{\detokenize{../../../../../docs/markdown/05-information-systems-architecture/sources.md#u07}}{U07}, \hyperref[element:21d6864657964ba5b03e3f6b5eb29cd5]{AX-08}.

\item \textbf{Audit Capture:} clients report semantically significant auditable occurrences
  and necessary context through a \textbf{Kleio-provided contract/library/API}.
  \textbf{Clients SHALL NOT be required to construct FHIR AuditEvent resources.} Kleio
  constructs and manages the authoritative representation; client API and mapping
  details remain unresolved.

\item \textbf{Information:} FHIR AuditEvent is the authoritative audit-event representation,
  without requiring every internal Kleio structure to be an AuditEvent. Indexes,
  correlation structures and prepared historical context may support management.
  \textbf{Authoritative audit evidence SHALL remain distinguishable from derived,
  correlated, contextualised or prepared historical views.} Existing meaningful
  provenance responsibilities remain; AuditEvent does not redefine all evidence.

\item \textbf{User activity:} semantically significant access/view, governed search,
  create/change/delete, submit/approve/reject, disclosure/export, operational action
  and relevant denied activity may require capture. These are examples, not a
  universal event taxonomy; “delete” grants no physical managed-information DELETE.
  Generic UI telemetry/clickstream is outside Kleio; UI security assurance and
  governance remain \textbf{Somebody Else's Problem (SEP)}.

\item \textbf{Collaboration/exclusions:} \textbf{Audit is not Governance. Audit is not Assurance.}
  Kleio may consume Journey information and Dokimasia outputs and provide
  historical audit context back to Dokimasia, without evaluating assurance or
  governing the subject. Dokimasia's governed record/result remains authoritative
  evidence of its assurance activity; its \href{\detokenize{../../../../../docs/markdown/05-information-systems-architecture/component-register.md#dokimasia-terminal-boundaries}}{terminal rules}
  prohibit a separate own-execution AuditEvent and recursive assurance. Consumption
  of its outputs does not circumvent those rules.

\end{itemize}






\par\medskip\textbf{Mandatory capture authority and the audited Identity}\par


\begin{quote}


\textbf{Subject of audit evidence ≠ authority performing audit capture.}



\end{quote}


Kleio performs mandatory authoritative audit capture under \textbf{intrinsic platform
authority} associated with its established architectural audit responsibility.
The originating/acting Identity remains evidence recorded by Kleio; it is not
the authority under which the mandatory AuditEvent is persisted. Kleio does not
borrow that Identity's PEP-Criteria merely to record its activity, including
activity denied to that Identity.



Identity-derived authority granted to human/system Identities remains resolved
and enforced through Themis. Kleio capture executes within established security
context and remains governed by Harmonia security policy, with Themis retaining
policy evaluation ownership. Mandatory capture is neither a security bypass nor
a policy exception. \href{\detokenize{../../../../governance/architectural-axioms.md#kleio-mandatory-capture-authority}}{AX-07}
incorporates this adjudication.



Where authoritative audit capture is required for a governed operation and
Kleio cannot exercise that platform authority sufficiently to capture the
required evidence, the governed operation \textbf{SHALL fail closed}. Durable target
availability remains necessary; authority alone establishes no successful capture.



This authority applies specifically to Kleio mandatory authoritative audit
capture. It grants no equivalent authority to other components or other Kleio
activities. Any proposed equivalent responsibility requires separate architecture
adjudication; representation, policy/enforcement mechanics and transaction
coordination remain \href{\detokenize{../../../../../docs/markdown/05-information-systems-architecture/gaps.md#kleio-detailed-contracts}}{CDG-43}.







\par\medskip\textbf{Immutable history and longitudinal evidence}\par


Authoritative AuditEvents are immutable. An existing AuditEvent SHALL NOT be
corrected, amended, overwritten or silently replaced. Correction or qualification
SHALL create a \textbf{new AuditEvent}, preserve the earlier event unchanged and refer
back to it as appropriate. The earlier assertion remains part of historical record.



Journey is the governed representation of observed Harmonia activity/progression;
AuditEvent is the authoritative audit representation of a significant auditable
occurrence. Journey is an important audit source, \textbf{not the audit boundary}.
A Journey may relate to many AuditEvents; an AuditEvent may exist without Journey.
\textbf{Journey participation SHALL NOT be a prerequisite for Audit Capture.}



\begin{quote}


\textbf{Reference by default; replicate only where necessary to preserve longitudinal
audit evidence.}



\end{quote}


Kleio SHALL NOT duplicate authoritative source information merely because it is
audit-relevant. Where historical state can be reliably resolved longitudinally,
Kleio references it with sufficient \textbf{identity, version, provenance and temporal
detail}. Version identity alone is not assumed to supply all temporal detail.
Only where the authoritative source cannot guarantee retrieval of the information
as it existed at the audited point in time may Kleio clone/replicate the necessary
source information within or in association with the authoritative AuditEvent to
preserve required longitudinal evidence. Detailed FHIR attachment/reference
mechanisms remain unresolved.







\par\medskip\textbf{Auditability and R1.x/R2.x availability}\par


\begin{quote}


\textbf{Harmonia SHALL NOT perform an auditable activity unless the required audit
information can be authoritatively and durably captured.}



\textbf{If Kleio cannot authoritatively capture required audit information, Harmonia
is DOWN.}



\end{quote}


Audit Capture is not best-effort and SHALL NOT be assumed reconstructable later
from logs, Journey information or incidental operational state. Transaction,
retry/failover and health/readiness mechanisms are deliberately unresolved.



Kleio retains audit application ownership; Mneme provides governed Context/Pragma access
and Mnemosyne \textbf{Harmonia Operations \& Assurance} establishes durable audit state.
For R1.x/R2.x, multiple \textbf{independent authoritative audit persistence targets}
through Mneme/Mnemosyne are permitted. At least \textbf{one configured target capable
of durably establishing the required audit information} is sufficient for
continued Harmonia operation. If none can establish it, Harmonia is unavailable
and fails closed. Target count, selection and whether one/some/all receive an
event are unresolved; targets SHALL NOT be called replicas without established
replication semantics.



\begin{quote}


\textbf{For R1.x/R2.x, audit data harmonisation between independent Mnemosyne audit
persistence targets is outside Harmonia scope — Somebody Else's Problem (SEP).}



\end{quote}


R1.x/R2.x guarantees neither identical histories, complete copies of all
AuditEvents, automatic convergence, reconciliation after disconnection nor unified
cross-target longitudinal history. Different targets may expose different histories;
an immutable AuditEvent remains authoritative without its target holding the whole
corpus. No future harmonisation/convergence mechanism is selected.



\begin{itemize}


\item \textbf{Standing:} U07 resolves responsibility, representation, immutability,
  longitudinal evidence, availability and release scope; \href{\detokenize{../../../../../docs/markdown/05-information-systems-architecture/gaps.md#kleio-detailed-contracts}}{CDG-43}
  retains detailed capture/query/correlation/retention/processing and persistence
  mechanisms. Historical FHIR/operational storage structures remain provenance,
  without competing with the established O\&A semantic alignment.

\end{itemize}

{\small Recorded unresolved matters: canonicalName; exact scope remains in the canonical record.}\par
{\footnotesize \href{../../../../continuum/application/application-component/fefda110e2274884ad58906733606efd.md}{Canonical record: fefda110e2274884ad58906733606efd}; full qualifications and provenance.}\par
\section*{IV.3.11 Mneme}
\phantomsection
\label{navigation:IV.3.11}
\addcontentsline{toc}{section}{IV.3.11 Mneme}

\subsection*{IV.3.11.1 Application Overview}
\phantomsection
\label{navigation:IV.3.11.1}
\addcontentsline{toc}{subsection}{IV.3.11.1 Application Overview}

\subsubsection*{IV.3.11.1.1 Mneme}
\phantomsection
\label{navigation:IV.3.11.1.1}
\addcontentsline{toc}{subsubsection}{IV.3.11.1.1 Mneme}

\label{element:1af8a56358f7429fabfb57be66542a99}
\index{Mneme (IV.3.11.1.1)}
{\small ApplicationComponent · Candidate — accepted Step-B baseline; final Solution Architecture unapproved · current}\par

\begin{itemize}


\item \textbf{Role/responsibility:} \textbf{Harmonia's Context and Pragma Manager}. Mneme manages
  the information context of activity through Pragma/Context, providing controlled,
  concurrent and seamless access to existing values and management of new Content.
  \href{\detokenize{../../../../../docs/markdown/05-information-systems-architecture/sources.md#u10}}{U10} refines AX-05/11 active representation/coordination.

\item \textbf{Principal exposed services:} \textbf{Context Load, Context Save, Adhoc Resource
  Access, Adhoc Search} — four services. \hyperref[element:49c9eb2bef0c47b8b5328621d5fc5479]{Lifecycle}
  distinguishes Mneme materialisation/save from Ergo Loader/Unloader boundaries.
  Likely resource sets are non-exhaustive; bounded runtime acquisition is supported.

\item \textbf{Internal responsibilities:} Pragma Management C28 retains PragmaFactory and
  Publishing/Subscription/Pattern semantic eligibility/selection. Context management
  composes authoritative resource values into a versioned temporal snapshot;
  Resource Harmonisation Manager C40 manages exceptional delivery conditions.
  C11 supporting contracts do not add exposed services.

\item \textbf{Information:} Context D26 is the working value environment; Pragma C07 is
  the input/output envelope. Changed Context values and new Content expressed as
  Pragma Output D28 remain distinct. Mneme manages their subsequent consequences.

\item \textbf{Persistence:} active state is reconstructable; Mnemosyne alone establishes
  durable state and authoritatively reports persistence outcomes. Working changes,
  cache retention/replication and output production confer no durable authority.
  Context/resource versions and active generation remain distinct.

\item \textbf{Collaboration/exclusions:} Ergo framework anticipates/loads; Ergon primarily
  acts on Context values and may acquire Adhoc information. Resource harmonisation
  uses new governed Harmonisation Ergon work through common Ponos execution, not
  originating-Ergon conflict resolution. Adhoc Search uses Mnemosyne Search Management.
  Petasos retains handoff; Themis owns enforcement. Mneme is not a generic repository,
  execution framework, durable version authority or raw database/cache access path.

\item \textbf{Standing:} U10 adjudicates Context/Pragma responsibility; detailed representation,
  concurrency/access/Output and harmonisation contracts remain CDG-45/46; older
  cache/write-behind and access variants remain CDG-10/11/16/27/31.

\end{itemize}

{\small Recorded unresolved matters: canonicalName; exact scope remains in the canonical record.}\par
{\footnotesize \href{../../../../continuum/application/application-component/1af8a56358f7429fabfb57be66542a99.md}{Canonical record: 1af8a56358f7429fabfb57be66542a99}; full qualifications and provenance.}\par
\subsection*{IV.3.11.2 Component Architecture}
\phantomsection
\label{navigation:IV.3.11.2}
\addcontentsline{toc}{subsection}{IV.3.11.2 Component Architecture}

\subsubsection*{IV.3.11.2.1 PragmaFactory}
\phantomsection
\label{navigation:IV.3.11.2.1}
\addcontentsline{toc}{subsubsection}{IV.3.11.2.1 PragmaFactory}

\label{element:75e77d8dfd1d4870aa6151485d28d1dd}
\index{PragmaFactory (IV.3.11.2.1)}
{\small ApplicationComponent · Candidate — accepted Step-B baseline; final Solution Architecture unapproved · current}\par

\begin{itemize}


\item \textbf{Role:} candidate logical application embodiment associated with the
  ActionableTaskArchetype meta-concept, within Mneme Pragma Management C28.
  \href{\detokenize{../../../../../docs/markdown/05-information-systems-architecture/sources.md#u02}}{U01/U02}. The meta-concept is not itself an instantiated
  Information Asset; the factory is not a new upstream Information Resource.

\item \textbf{Responsibilities/features:} creates/initialises appropriate Pragma
  representations. Detailed instantiation/selection/validation, parameters and
  policy contracts remain \textbf{unknown}; the basic responsibility does not finalise
  a factory feature contract.

\item \textbf{Contained constructs/information/persistence:} subcomponents, realised
  definition Data Objects, inputs/outputs and persistence responsibility unknown.

\item \textbf{Collaborations/exclusions:} management allocation is established by U02;
  detailed callers and contracts remain unresolved. It must not be equated to a
  TaskSequence seeder or PraxisDefinition by inference.

\item \textbf{Standing:} explicitly instructed candidate solution meaning, CDG-04; not
  invented from a matching class and not rejected for missing implementation.

\end{itemize}

{\small Recorded unresolved matters: canonicalName; exact scope remains in the canonical record.}\par
{\footnotesize \href{../../../../continuum/application/application-component/75e77d8dfd1d4870aa6151485d28d1dd.md}{Canonical record: 75e77d8dfd1d4870aa6151485d28d1dd}; full qualifications and provenance.}\par
\subsection*{IV.3.11.3 Information Architecture}
\phantomsection
\label{navigation:IV.3.11.3}
\addcontentsline{toc}{subsection}{IV.3.11.3 Information Architecture}

\subsubsection*{IV.3.11.3.1 Pragma Topic}
\phantomsection
\label{navigation:IV.3.11.3.1}
\addcontentsline{toc}{subsubsection}{IV.3.11.3.1 Pragma Topic}

\label{element:27ad143e7c784e378556af8a57c438f7}
\index{Pragma Topic (IV.3.11.3.1)}
{\small DataObject · Bounded adjudicated application information meaning; Domain05 remains the accepted candidate Step-B baseline, not final Solution Architecture. Detailed representation is deferred; only this defined logical object is projected. · current}\par

\textbf{Mneme manages Pragma.} Context is the value-based working environment;
Pragma is the task/activity information envelope through which inputs and outputs
are expressed. They are related and distinct.



\textbf{Pragma realises \hyperref[element:c34b4929f71946708d37dbe6c1da920a]{Metabolē}}
for governed Harmonia information/activity progression. This is architectural
Realisation, not synonymy or inheritance from Diadromí. Mneme retains Pragma
management; the conceptual construct introduces no common manager, persistence
requirement or shared lifecycle with Pylai.Diadromí.



The conceptual
\hyperref[element:c34b4929f71946708d37dbe6c1da920a]{ActivityContext}
classifies the broad activity and is distinct from request action, this loaded
Context and Themis PEP-Context. Their detailed physical mapping remains deferred.



An Ergon may read existing Context values, act upon/change working Context
values, and generate \textbf{new Content}. New Content is expressed as \textbf{Pragma
Output}, rather than merely an incidental Context mutation:



\begin{quote}\small\ttfamily
Existing Context values → Ergon → changed working Context values\\
Context + bounded business logic → Ergon → Pragma Output → new Content\\
                  Context Unloader returns both, distinctly\\
                                    ↓\\
                 Mneme manages subsequent information consequences\\

\end{quote}


Changed Context values and newly generated Pragma Outputs SHALL remain
semantically distinguishable. A working change or produced output does not
itself establish authoritative durable state. Task/undertaking/outcome ownership
and output rules remain unchanged.



Ergon execution may produce one or more outputs, whose processing may result in
multiple Pragma. \textbf{Only a Pragma playing the Pathway-continuation role identifies
information for continuing the current Praxis along its Pathway}; other resulting
Pragma are independently available for governed Pub/Sub, another Praxis, other
processing or no further processing. Praxis defines continuation semantics;
the execution framework performs applicable output processing and routing,
without making Praxis an orchestrator/choreographer or transferring Mneme's
information management. See the \hyperref[element:27ad143e7c784e378556af8a57c438f7]{shared semantic view}
and \href{\detokenize{../../../../../docs/markdown/05-information-systems-architecture/design-notes/DesignNote-Pragma-Outputs-and-Pathway-Continuation.md}}{DesignNote}.
The continuation role creates no subtype. Whether a new Pragma is constructed,
an output Pragma is designated, or another governed representation is used remains
\href{\detokenize{../../../../../docs/markdown/05-information-systems-architecture/gaps.md#continuation-representation}}{explicitly deferred}. Exact Output
representation, routing mechanisms and Pragma/Information Unit mappings remain open.



Mneme may manage richer resulting Pragma relationships and provenance than the
originating Praxis's Journey represents. Continuation lineage supports meaningful
Journey progression without making Journey a complete provenance graph, execution
log, audit trail or queue history; \hyperref[element:359b901df27f4b43b5a2ff14a8484e34]{Journey's boundary}
and Energeia ownership remain intact.



Source standing: Bounded adjudicated application information meaning; Domain05 remains the accepted candidate Step-B baseline, not final Solution Architecture. Detailed representation is deferred; only this defined logical object is projected.


{\footnotesize \href{../../../../continuum/application/data-object/27ad143e7c784e378556af8a57c438f7.md}{Canonical record: 27ad143e7c784e378556af8a57c438f7}; full qualifications and provenance.}\par
\subsubsection*{IV.3.11.3.2 Context Topic}
\phantomsection
\label{navigation:IV.3.11.3.2}
\addcontentsline{toc}{subsubsection}{IV.3.11.3.2 Context Topic}

\label{element:76ded4866bf443dabc0849d68b215690}
\index{Context Topic (IV.3.11.3.2)}
{\small DataObject · Bounded adjudicated application information meaning; Domain05 remains the accepted candidate Step-B baseline, not final Solution Architecture. Detailed representation is deferred; only this defined logical object is projected. · current}\par

\begin{quote}


\textbf{Context is a versioned, temporal load-snapshot composed from authoritative
resource values materialised for a Harmonia activity's anticipated information
requirements.}



\end{quote}


Context contains \textbf{values, not merely references}. Each value represents the
defined authoritative resource state/version actually materialised when it is
incorporated. Loaded values are composed into the Context bucket:



\begin{quote}\small\ttfamily
Context — identity / Context version / temporal load basis\\
  └── Resource Bucket                         [Composition]\\
        ├── Resource Value A @ authoritative state/version\\
        ├── Resource Value B @ authoritative state/version\\
        └── Resource Value C @ authoritative state/version\\

\end{quote}


\textbf{Composition} concerns loaded values belonging to the snapshot. The Context
does not existentially own the underlying authoritative resources. This is
distinct from the separately established Aggregation relationships for Praxis,
Journey and PragmaFactory.



Context version and constituent resource authoritative state/version SHALL
remain distinct. Neither is an alias for a valid active-state generation or
security PEP-Context. Materialisation does \textbf{not} imply a globally
transactionally consistent distributed snapshot or a universal frozen instant.
Context is the Ergon's managed working information environment, not a closed
universe of all information it may access or produce.



Authoritative constituent resource state may progress while an Ergon operates.
Mneme needs sufficient state/version information to recognise delivery that cannot
safely proceed normally; generic conflict handling remains outside the originating
Ergon. No concurrency or conflict-detection algorithm is established here.



This activity Context does not redefine Domain04's Information Unit
\href{\detokenize{../../../../../docs/markdown/04-information-architecture/metamodel/information-architecture-metamodel.md#65-context}}{Management Context}.
Their detailed mapping, physical bucket representation, Context identity/version
mechanisms and concurrency algorithms remain unresolved under AX-17.



Source standing: Bounded adjudicated application information meaning; Domain05 remains the accepted candidate Step-B baseline, not final Solution Architecture. Detailed representation is deferred; only this defined logical object is projected.


{\footnotesize \href{../../../../continuum/application/data-object/76ded4866bf443dabc0849d68b215690.md}{Canonical record: 76ded4866bf443dabc0849d68b215690}; full qualifications and provenance.}\par
\subsection*{IV.3.11.4 Services}
\phantomsection
\label{navigation:IV.3.11.4}
\addcontentsline{toc}{subsection}{IV.3.11.4 Services}

\subsubsection*{IV.3.11.4.1 Context Load}
\phantomsection
\label{navigation:IV.3.11.4.1}
\addcontentsline{toc}{subsubsection}{IV.3.11.4.1 Context Load}

\label{element:1ccfe420541945bb9c0756a949799d95}
\index{Context Load (IV.3.11.4.1)}
{\small ApplicationService · Candidate — adjudicated Domain05 Step-B baseline; final Solution Architecture unapproved · current}\par

Materialise a managed Context from authoritative resource values for anticipated activity requirements.



Mneme exposes this principal service. Resource Harmonisation Manager remains internal, and supporting contracts are not additional exposed services. The Context/Pragma and authoritative-state boundaries apply.


{\small Recorded unresolved matters: canonicalName; exact scope remains in the canonical record.}\par
{\footnotesize \href{../../../../continuum/application/application-service/1ccfe420541945bb9c0756a949799d95.md}{Canonical record: 1ccfe420541945bb9c0756a949799d95}; full qualifications and provenance.}\par
\subsubsection*{IV.3.11.4.2 Context Save}
\phantomsection
\label{navigation:IV.3.11.4.2}
\addcontentsline{toc}{subsubsection}{IV.3.11.4.2 Context Save}

\label{element:24d4fc42d9524d53a747f0550833f1a4}
\index{Context Save (IV.3.11.4.2)}
{\small ApplicationService · Candidate — adjudicated Domain05 Step-B baseline; final Solution Architecture unapproved · current}\par

Return resulting Context state to Mneme management for appropriate subsequent information handling, alongside distinguishable Pragma Outputs.



Mneme exposes this principal service. Resource Harmonisation Manager remains internal, and supporting contracts are not additional exposed services. The Context/Pragma and authoritative-state boundaries apply.


{\small Recorded unresolved matters: canonicalName; exact scope remains in the canonical record.}\par
{\footnotesize \href{../../../../continuum/application/application-service/24d4fc42d9524d53a747f0550833f1a4.md}{Canonical record: 24d4fc42d9524d53a747f0550833f1a4}; full qualifications and provenance.}\par
\subsubsection*{IV.3.11.4.3 Adhoc Resource Access}
\phantomsection
\label{navigation:IV.3.11.4.3}
\addcontentsline{toc}{subsubsection}{IV.3.11.4.3 Adhoc Resource Access}

\label{element:410f2f55bbda481b973c6c67b924f0be}
\index{Adhoc Resource Access (IV.3.11.4.3)}
{\small ApplicationService · Candidate — adjudicated Domain05 Step-B baseline; final Solution Architecture unapproved · current}\par

Governed runtime acquisition of resource information outside the anticipated loaded Context where the bounded activity requires it.



Mneme exposes this principal service. Resource Harmonisation Manager remains internal, and supporting contracts are not additional exposed services. The Context/Pragma and authoritative-state boundaries apply.


{\small Recorded unresolved matters: canonicalName; exact scope remains in the canonical record.}\par
{\footnotesize \href{../../../../continuum/application/application-service/410f2f55bbda481b973c6c67b924f0be.md}{Canonical record: 410f2f55bbda481b973c6c67b924f0be}; full qualifications and provenance.}\par
\subsubsection*{IV.3.11.4.4 Adhoc Search}
\phantomsection
\label{navigation:IV.3.11.4.4}
\addcontentsline{toc}{subsubsection}{IV.3.11.4.4 Adhoc Search}

\label{element:02bf54d699b44ed0b26fa5b8125f0f01}
\index{Adhoc Search (IV.3.11.4.4)}
{\small ApplicationService · Candidate — adjudicated Domain05 Step-B baseline; final Solution Architecture unapproved · current}\par

Activity/context-facing governed search, using Mnemosyne's underlying authoritative persistence Search Management.



Mneme exposes this principal service. Resource Harmonisation Manager remains internal, and supporting contracts are not additional exposed services. The Context/Pragma and authoritative-state boundaries apply.


{\small Recorded unresolved matters: canonicalName; exact scope remains in the canonical record.}\par
{\footnotesize \href{../../../../continuum/application/application-service/02bf54d699b44ed0b26fa5b8125f0f01.md}{Canonical record: 02bf54d699b44ed0b26fa5b8125f0f01}; full qualifications and provenance.}\par
\subsection*{IV.3.11.5 Internal Functions}
\phantomsection
\label{navigation:IV.3.11.5}
\addcontentsline{toc}{subsection}{IV.3.11.5 Internal Functions}

\subsubsection*{IV.3.11.5.1 Pragma Management}
\phantomsection
\label{navigation:IV.3.11.5.1}
\addcontentsline{toc}{subsubsection}{IV.3.11.5.1 Pragma Management}

\label{element:ab0cdaa7ed4149ee856216efa5a265c4}
\index{Pragma Management (IV.3.11.5.1)}
{\small ApplicationFunction · Candidate — accepted Step-B baseline; final Solution Architecture unapproved · current}\par

\begin{itemize}


\item \textbf{Role/responsibility:} Mneme-associated information-management responsibility
  area with logical Aggregation of PragmaFactory C08 and Publishing C36,
  Subscription C37, Pattern C38 responsibilities, managing Pragma C07.
  \href{\detokenize{../../../../../docs/markdown/05-information-systems-architecture/sources.md#u05}}{U02/U03/U05}. Aggregation/management identifies responsibility,
  not Composition or a required process/service.

\item \textbf{Features/information:} creation/initialisation through PragmaFactory and
  governed management of Pragma representations and common process-management,
  assurance, reporting and platform integration associated with that information.
  Publishing/Subscription definitions, Pattern semantics/evaluation, matching,
  selection, semantic eligibility and eligible recipient/consumer relationships
  belong here; concrete persistent destination/handoff mechanics remain Petasos.
  Pragma expresses activity inputs/outputs; new Content is \textbf{Pragma Output}, distinct
  from changed Context working values. Mneme manages its subsequent consequences
  under \hyperref[element:27ad143e7c784e378556af8a57c438f7]{U10}.
  U18 distinguishes independently resulting Pragma from the Pathway-continuation
  role without transferring Praxis semantics or execution/routing to Mneme.
  Stimulus/Effector and candidate Dokimasia Pragma semantic roles
  reuse these facilities instead of requiring specialist equivalents.

\item \textbf{Persistence/collaboration:} active use/coordination through Mneme; authoritative
  durable establishment through Mnemosyne. Ponos executes, Petasos provides
  messaging/recoverable handoff and Ergon supplies bounded activity behaviour.

\item \textbf{Standing/exclusions:} candidate responsibility allocations established by
  U02/U03/U05;
  detailed management contracts and persistence structures remain unresolved.
  Execution of, behaviour upon, or handoff of managed information does not
  confer information-management responsibility for that information.

\end{itemize}

{\small Recorded unresolved matters: canonicalName; exact scope remains in the canonical record.}\par
{\footnotesize \href{../../../../continuum/application/application-function/ab0cdaa7ed4149ee856216efa5a265c4.md}{Canonical record: ab0cdaa7ed4149ee856216efa5a265c4}; full qualifications and provenance.}\par
\subsubsection*{IV.3.11.5.2 Resource Harmonisation Manager}
\phantomsection
\label{navigation:IV.3.11.5.2}
\addcontentsline{toc}{subsubsection}{IV.3.11.5.2 Resource Harmonisation Manager}

\label{element:a6580e0b9f0f40c9a33d731d20110057}
\index{Resource Harmonisation Manager (IV.3.11.5.2)}
{\small ApplicationFunction · Candidate — accepted Step-B baseline; final Solution Architecture unapproved · current}\par

\textbf{Internal Mneme responsibility/capability, not a fifth exposed service or an
independent execution framework.} \href{\detokenize{../../../../../docs/markdown/05-information-systems-architecture/sources.md#u10}}{U10} establishes the
\hyperref[element:a6580e0b9f0f40c9a33d731d20110057]{exceptional delivery path}.
It identifies/manages conditions that cannot proceed normally and causes appropriate
bounded Harmonisation Ergon work; it does not perform domain harmonisation itself.



\begin{itemize}


\item \textbf{Expected use:} approximately 5\% or less at worst of resource-delivery activity,
  a design expectation, not a hard SLA. Normal handling remains efficient.

\item \textbf{Execution:} new governed work through the established Harmonia work/execution
  contract reuses Ponos; no resumed originating Ergon, duplicate/embedded execution
  framework or originating-Ergon generic conflict resolution.

\item \textbf{Information/reuse:} relevant condition/state/Context is preserved for governed
  handling; a reusable Ergo Library remains a candidate reuse mechanism only.

\item \textbf{Unresolved:} categories, detection/concurrency/harmonisation algorithms,
  Harmonisation Ergon catalogue, detailed work contracts and repeated-harmonisation
  terminal/escalation handling (CDG-46). No UI/forms architecture is selected.
  Kleio cross-target audit harmonisation remains R1.x/R2.x SEP.

\end{itemize}

{\small Recorded unresolved matters: canonicalName; exact scope remains in the canonical record.}\par
{\footnotesize \href{../../../../continuum/application/application-function/a6580e0b9f0f40c9a33d731d20110057.md}{Canonical record: a6580e0b9f0f40c9a33d731d20110057}; full qualifications and provenance.}\par
\section*{IV.3.12 Mnemosyne}
\phantomsection
\label{navigation:IV.3.12}
\addcontentsline{toc}{section}{IV.3.12 Mnemosyne}

\subsection*{IV.3.12.1 Application Overview}
\phantomsection
\label{navigation:IV.3.12.1}
\addcontentsline{toc}{subsection}{IV.3.12.1 Application Overview}

\subsubsection*{IV.3.12.1.1 Mnemosyne}
\phantomsection
\label{navigation:IV.3.12.1.1}
\addcontentsline{toc}{subsubsection}{IV.3.12.1.1 Mnemosyne}

\label{element:7985e647bf9e4ca3ba70faf9e3e4ae02}
\index{Mnemosyne (IV.3.12.1.1)}
{\small ApplicationComponent · Candidate — accepted Step-B baseline; final Solution Architecture unapproved · current}\par

\begin{itemize}


\item \textbf{Role/responsibility:} \textbf{Harmonia's Persistence Manager} under
  \href{\detokenize{../../../../../docs/markdown/05-information-systems-architecture/sources.md#u11}}{U11}, establishing authoritative durable state/version and the
  information/metadata required to sustain activity Context. Persistence conveys
  no originating information/application ownership.

\item \textbf{First-class responsibilities:} transaction oversight; persistent-service
  oversight/reporting; shard management; authoritative persistence reporting;
  persistence logging; compartment validation/reapplication; hosting Themis-owned
  PEP enforcement; \textbf{Search Management}; \textbf{Research Metadata Management}.
  \hyperref[element:7985e647bf9e4ca3ba70faf9e3e4ae02]{The management view}
  establishes each boundary. Underlying services may retain their own transactions;
  Mnemosyne does not become infrastructure management.

\item \textbf{Information/outcomes:} durable representations/version/metadata, persistence
  observations and outcomes, search capability metadata and distinct Research Metadata.
  API/engine success is not automatically the authoritative Harmonia outcome;
  uncertain effects remain explicit. Persistence logging is distinct from Kleio audit.

\item \textbf{Search:} Mneme exposes Adhoc Search; Mnemosyne performs underlying authoritative
  managed-persistence search, without treating a transient cache as complete.
  Persistence-domain/shard/service/search-index capability metadata may be managed.

\item \textbf{Research Metadata:} manages capture, persistence, management and search utilisation
  of AI-derived meaning/relationships across content sets based on proximity. Source
  Content remains unchanged and derived metadata distinguishable. Deriving AI
  processing executes \textbf{outside Mnemosyne} through appropriate Harmonia behaviour/
  execution. Detailed provenance/confidence/schema remain unresolved.

\item \textbf{Persistence domains:} Clinical Information, C\&G and O\&A retain U06 guardrails:
  coherent C\&G authority; O\&A horizontal partitionability preserving identity,
  provenance, relationships and meaning. Domains select no database/service topology
  or shard key. Mnemosyne manages shards within semantic constraints, without
  independently deciding sharding safety.

\item \textbf{Governance:} compartments are validated/reapplied without loss or weakening;
  semantic applicability is not independently determined here. Themis owns PEP
  semantics and Criteria; Mnemosyne controls/prevents the persistence operation.
  Kleio mandatory AuditEvent persistence uses its
  \href{\detokenize{../../../../../docs/markdown/05-information-systems-architecture/component-register.md#kleio-capture-authority}}{governed intrinsic platform authority}, distinct
  from the audited Identity; Mnemosyne gains no equivalent intrinsic authority.

\item \textbf{Existing registry support:} server-side Provider Registry durable integrity,
  preservation/history and recovery responsibilities remain as described in C13
  and the data view; this correction does not redesign their semantic validation.

\item \textbf{Collaboration/exclusions:} Mneme manages Context/Pragma information consequences;
  Mnemosyne establishes/reports persistence. Energeia observes/reports persistence
  operational state, without acquiring management. No application-facing raw storage,
  AI derivation inside Mnemosyne, technology selection or workflow engine is established.

\item \textbf{Standing:} U11 and explicit task decisions settle responsibilities; detailed
  oversight/transaction/compartment/outcome/search/Research Metadata contracts remain
  CDG-47/48, alongside historical structures CDG-12/13/40.

\end{itemize}

{\small Recorded unresolved matters: canonicalName; exact scope remains in the canonical record.}\par
{\footnotesize \href{../../../../continuum/application/application-component/7985e647bf9e4ca3ba70faf9e3e4ae02.md}{Canonical record: 7985e647bf9e4ca3ba70faf9e3e4ae02}; full qualifications and provenance.}\par
\section*{IV.3.13 Ponos}
\phantomsection
\label{navigation:IV.3.13}
\addcontentsline{toc}{section}{IV.3.13 Ponos}

\subsection*{IV.3.13.1 Application Overview}
\phantomsection
\label{navigation:IV.3.13.1}
\addcontentsline{toc}{subsection}{IV.3.13.1 Application Overview}

\subsubsection*{IV.3.13.1.1 Ponos}
\phantomsection
\label{navigation:IV.3.13.1.1}
\addcontentsline{toc}{subsubsection}{IV.3.13.1.1 Ponos}

\label{element:b521762c1ea444b48cd5f66be3752d30}
\index{Ponos (IV.3.13.1.1)}
{\small ApplicationComponent · Candidate — accepted Step-B baseline; final Solution Architecture unapproved · current}\par

\begin{itemize}


\item \textbf{Role/responsibility:} reusable execution framework/component pattern, peer
  to Energeia C39, for governed execution and observable activity progression.
  \href{\detokenize{../../../../../docs/markdown/05-information-systems-architecture/sources.md#u02}}{U02} adjudicates this role; \href{\detokenize{../../../../../docs/markdown/05-information-systems-architecture/sources.md#a02}}{A02} and
  \href{\detokenize{../../../../../docs/markdown/05-information-systems-architecture/sources.md#l02}}{L02/L01/L23} retain responsibility context and historical detail.

\item \textbf{Features/contained constructs:} work ingress and dispatcher, workflow/sequence
  resolution where governed control is established, activity invocation, execution
  concurrency/timeouts, execution retry/re-attempts, checkpoint coordination and
  administrative inspection. Applies Praxis-defined output-processing/continuation
  semantics and routing; not every output-derived Pragma continues that Praxis.
  \href{\detokenize{../../../../../docs/markdown/05-information-systems-architecture/sources.md#u18}}{U18} establishes the distinction, leaving representation open.
  Subscription matching belongs to Mneme; handoff
  redelivery/recovery and dead-letter semantics belong to Petasos.
  Consumer loops/thread pools are documented mechanisms, not logical children.
  Ordinary semantic interest fan-out to different Ergon types is orthogonal to
  ActionableTask.ExecutionConcurrency permitting concurrent undertakings by
  instances of the same type, each through its own FulfillmentTask. DigitalTwin
  specialises execution with serial Behaviour processing within a Twin; common
  concurrency facilities do not permit intra-Twin Behaviour concurrency (U19).

\item \textbf{Information:} executes consumer undertakings using governed information/context
  resolved through Mneme; consults established control definitions and produces
  execution/progression observations contributing to Journey. Petasos handoff
  does not make an opaque payload into a Ponos-owned Pragma.

\item \textbf{Persistence:} coordinates recoverable handoff through Petasos and governed
  information change through Mneme/Mnemosyne. Does not establish durable state
  by completing an activity or writing a checkpoint to a cache.

\item \textbf{Collaboration/exclusions:} tightly coupled to peer Ergo/Ergon responsibilities,
  executing/managing bounded Ergon activity conforming to the Ergo contract;
  does not structurally contain Ergo or Ergon. Uses Praxis, Erga, Themis, Mneme
  Pragma Management and Petasos. DigitalTwin C19 and Dokimasia C20 specialise/
  reuse this pattern; they are not children inside a Ponos runtime. Conceptual
  supertype language does not mandate literal implementation inheritance.
  Does not own Pragma information management, Publishing/Subscription/Pattern
  semantic selection, workflow definitions, concrete domain logic, broker lifecycle,
  external standards/transport or universal clinical-work determination.

\item \textbf{Standing:} specialisation relationships adjudicated for the candidate by U02;
  detailed model remains candidate. CDG-03/08/16–18 retain incomplete control
  and continuation semantics. Numerical worker/retry defaults are not baselined.

\end{itemize}

{\small Recorded unresolved matters: canonicalName; exact scope remains in the canonical record.}\par
{\footnotesize \href{../../../../continuum/application/application-component/b521762c1ea444b48cd5f66be3752d30.md}{Canonical record: b521762c1ea444b48cd5f66be3752d30}; full qualifications and provenance.}\par
\subsection*{IV.3.13.2 Component Architecture}
\phantomsection
\label{navigation:IV.3.13.2}
\addcontentsline{toc}{subsection}{IV.3.13.2 Component Architecture}

\subsubsection*{IV.3.13.2.1 Ergon}
\phantomsection
\label{navigation:IV.3.13.2.1}
\addcontentsline{toc}{subsubsection}{IV.3.13.2.1 Ergon}

\label{element:49c9eb2bef0c47b8b5328621d5fc5479}
\index{Ergon (IV.3.13.2.1)}
{\small ApplicationComponent · Candidate — accepted Step-B baseline; final Solution Architecture unapproved · current}\par

\begin{itemize}


\item \textbf{Role/responsibility:} a single well-defined software component encapsulating
  a bounded fine-grained task-completion activity set within the Ergo contract,
  executed/managed through Ponos. \href{\detokenize{../../../../../docs/markdown/05-information-systems-architecture/sources.md#u02}}{U01/U02}, \href{\detokenize{../../../../../docs/markdown/05-information-systems-architecture/sources.md#l03}}{L03/L19}.

\item \textbf{Features:} specialised transformation/mapping, validation, identifier/reference
  checking, enrichment, routing, result processing, fan-out or registry-change
  behaviour. Reuses C05 common facilities; \texttt{\seqsplit{ErgonBase}} is historical contract
  support, not an architectural inheritance requirement.

\item \textbf{Information:} acts upon or generates Pragma, payload/resource information,
  outputs, checkpoints and observations through governed contracts. It primarily
  acts on loaded Context values, may change working values and may generate new
  Content as \textbf{Pragma Output}, distinct from Context mutation. Bounded \textbf{Adhoc
  Resource Access / Adhoc Search} remains permitted for emerging needs. \textbf{No discrete
  Information Architecture representation or Ergon Information Resource exists
  merely because the component is an Ergon.} Activity is not information about it.

\item \textbf{Output progression:} one or more execution outputs may yield multiple Pragma.
  The \hyperref[element:27ad143e7c784e378556af8a57c438f7]{Pathway-continuation role}
  identifies information continuing the current Praxis; independently resulting
  Pragma remain available for governed use, without automatically advancing it.

\item \textbf{Persistence/collaboration:} requests governed information use/change through
  Mneme; does not own Pragma management or durable state. Erga collects components;
  Ponos executes them; Petasos effects persistent handoff. Every Ergon exposes
  IngressEndpoint and EgressEndpoint C29; traversal does not expose internal logic.
  Praxis relationship detail remains open. The originating Ergon owns no persistence
  mechanics, shard/database transactions or generic conflict resolution; exceptional
  harmonisation is new governed work under U10.

\item \textbf{Specialisations:} Behaviour C25, Dokimasia Ergon C26 and, through the latter,
  AI Assurance Ergon C27. They retain normal Ergo/Ergon facilities, constrained
  by the specialised role and Dokimasia's terminal rules where applicable.

\item \textbf{Standing/exclusions:} bounded behaviour established for the candidate;
  competing atomic/pure/idempotent/branching claims remain CDG-05. An ordinary
  Ergon may intentionally invoke an AI service internally as an implementation
  choice without becoming AI Assurance Ergon or transferring authority.

\end{itemize}

{\small Recorded unresolved matters: canonicalName; exact scope remains in the canonical record.}\par
{\footnotesize \href{../../../../continuum/application/application-component/49c9eb2bef0c47b8b5328621d5fc5479.md}{Canonical record: 49c9eb2bef0c47b8b5328621d5fc5479}; full qualifications and provenance.}\par
\subsection*{IV.3.13.3 Interfaces}
\phantomsection
\label{navigation:IV.3.13.3}
\addcontentsline{toc}{subsection}{IV.3.13.3 Interfaces}

\subsubsection*{IV.3.13.3.1 PathwayEndpoint}
\phantomsection
\label{navigation:IV.3.13.3.1}
\addcontentsline{toc}{subsubsection}{IV.3.13.3.1 PathwayEndpoint}

\label{element:800789f54ba24f53a354600d5cd43f6f}
\index{PathwayEndpoint (IV.3.13.3.1)}
{\small ApplicationInterface · Candidate — accepted Step-B baseline; final Solution Architecture unapproved · current}\par

Application Interface conceptually equivalent to ArchiMate ApplicationInterface.
Every Ergon exposes IngressEndpoint and EgressEndpoint. These surfaces do not
expose internal processing; endpoint identity and technical interfaces remain open.


{\small Recorded unresolved matters: canonicalName; exact scope remains in the canonical record.}\par
{\footnotesize \href{../../../../continuum/application/application-interface/800789f54ba24f53a354600d5cd43f6f.md}{Canonical record: 800789f54ba24f53a354600d5cd43f6f}; full qualifications and provenance.}\par
\section*{IV.3.14 Pylai}
\phantomsection
\label{navigation:IV.3.14}
\addcontentsline{toc}{section}{IV.3.14 Pylai}

\subsection*{IV.3.14.1 Application Overview}
\phantomsection
\label{navigation:IV.3.14.1}
\addcontentsline{toc}{subsection}{IV.3.14.1 Application Overview}

\subsubsection*{IV.3.14.1.1 Pylai}
\phantomsection
\label{navigation:IV.3.14.1.1}
\addcontentsline{toc}{subsubsection}{IV.3.14.1.1 Pylai}

\label{element:e3e4147fea3640198dc251ed41ee854d}
\index{Pylai (IV.3.14.1.1)}
{\small ApplicationComponent · Candidate — accepted Step-B baseline; final Solution Architecture unapproved · current}\par

\begin{itemize}


\item \textbf{Role/responsibility:} \textbf{Harmonia's External Interoperability Service},
  enabling/managing system-to-system communication irrespective of style, protocol,
  standard or technology. Messaging, synchronous/asynchronous RMI, APIs, files and
  polling are supported architectural possibilities; RMI does not mean Java RMI.
  \href{\detokenize{../../../../../docs/markdown/05-information-systems-architecture/sources.md#u14}}{U14}, \href{\detokenize{../../../../../docs/markdown/05-information-systems-architecture/pylai-external-interoperability.md}}{canonical Pylai/Horos view}.

\item \textbf{Constructs/patterns:} Terma C43 acts at the logical external interface;
  Diadromí C44 carries/represents transient interaction. PylaiPathway connects
  governed Praxis/Ergo work; PylaiBroker serves executing BusinessLogic;
  PylaiServer changes use Pathway and READ/SEARCH uses Diadromí-Controller C45.
  These patterns do not mandate classes, services or deployment boundaries.

\item \textbf{Information/persistence:} Communication Information Unit preserves raw payload,
  communication metadata and applicable external identity; exact schema unresolved.
  Ingress Loader establishes Task.input/Context, with semantic mapping in paired
  BusinessLogic. Durable acceptance requires authoritative Pragma/Task persistence
  through Mneme/Mnemosyne. Protocol ACK and queue/cache success are distinct.
  Diadromí is transient and non-recoverable; Petasos owns durable Parcel handoff.

\item \textbf{Security/evidence:} carries external identity; Themis establishes Harmonia
  authority/enforcement. Kleio constructs authoritative AuditEvent from reported
  significant transactions; Dokimasia assures; Journey records applicable observed
  business progression or interoperability failure. None of these are synonyms.

\item \textbf{Exclusions/standing:} no execution framework, persistence manager, authority
  inference, fabricated READ/SEARCH work or mandatory Controller for all interactions.
  AX-02/13 non-destructive fail-closed publication remains. Technology/handover
  details remain CDG-32/34/50; CDG-51 records the resolved Strategy boundary.
  Kleio mandatory capture uses its own governed platform authority, not the
  external/acting Identity's grants; Pylai gains no such authority.

\end{itemize}

{\small Recorded unresolved matters: canonicalName; exact scope remains in the canonical record.}\par
{\footnotesize \href{../../../../continuum/application/application-component/e3e4147fea3640198dc251ed41ee854d.md}{Canonical record: e3e4147fea3640198dc251ed41ee854d}; full qualifications and provenance.}\par
\subsection*{IV.3.14.2 Information Architecture}
\phantomsection
\label{navigation:IV.3.14.2}
\addcontentsline{toc}{subsection}{IV.3.14.2 Information Architecture}

\subsubsection*{IV.3.14.2.1 Pylai.Diadromí Topic}
\phantomsection
\label{navigation:IV.3.14.2.1}
\addcontentsline{toc}{subsubsection}{IV.3.14.2.1 Pylai.Diadromí Topic}

\label{element:978f5c2bd2d343ffaf3a13dff700c258}
\index{Pylai.Diadromí Topic (IV.3.14.2.1)}
{\small DataObject · Bounded adjudicated application information meaning; Domain05 remains the accepted candidate Step-B baseline, not final Solution Architecture. Detailed representation is deferred; only this defined logical object is projected. · current}\par

\textbf{Pylai.Terma} is the logical external interoperability interface/boundary.
\texttt{\seqsplit{Pylai.IngressTerma}} and \texttt{\seqsplit{Pylai.EgressTerma}} specialise it where useful. Terma owns
external communication mechanics: connection/session handling, listening, polling,
transport/protocol interaction, framing, segmentation/reassembly, protocol
acknowledgements, API/RMI/file interaction, interaction-local communication retry
and external identity carriage. One Terma may use one or more technical endpoints/connectors;
it is not reducible to a socket, URL, Camel endpoint or deployed adapter.



\begin{quote}


\textbf{Pylai.Terma manages the communication stack; the applicable Harmonia
execution framework manages the application meaning.}



\end{quote}


Terma owns the mechanics necessary to establish, conduct and complete an external
communication interaction. It does not own the Harmonia application semantics
of the information communicated. Application interpretation and processing use
the applicable Harmonia execution framework, principally Praxis/Ergo/Ergon,
except where an explicitly defined interaction pattern provides otherwise.
The \href{\detokenize{../../../../../docs/markdown/05-information-systems-architecture/pylai-external-interoperability.md#pylai-server-controller}}{Server READ/SEARCH pattern} is such a defined
pattern; it grants Terma no application/business interpretation responsibility.



\begin{quote}\small\ttfamily
External System → Pylai.Terma → Pylai.Diadromí → Praxis / Ergo / Ergon\\
                  communication mechanics       application interpretation,\\
                                                transformation / processing\\

\end{quote}


Ingress/Egress describe \textbf{information direction}, irrespective of client/server
role, who opened the connection or push/pull behaviour. An IngressTerma may poll;
an EgressTerma may expose an external file interaction.



\begin{quote}


\textbf{Pylai.Diadromí is the transient carriage and interaction context through which
a Pylai.Terma exchanges communication with its Harmonia-side participant for the
duration of an external interoperability interaction.}



\textbf{Terma acts; Diadromí carries/represents transient interaction.}



\end{quote}


A Diadromí may span multiple protocol exchanges, acknowledgements, retries within
the active interaction and segmented communication. An intermediate protocol
acknowledgement need not establish logical handover or authoritative Harmonia
acceptance. Diadromí is neither an Ergon, Task/Pragma alias nor a durable queue.
Its exact representation is unresolved.



\textbf{Pylai.Diadromí realises
\hyperref[element:c34b4929f71946708d37dbe6c1da920a]{Metabolē}}
for transient information progression through Pylai interoperability activity.
This is genuine architectural Realisation, not analogy, synonymy with Metabolē
or inheritance from Pragma. Pylai retains responsibility; Diadromí's transient,
non-recoverable lifecycle does not acquire Pragma persistence or Mneme management
requirements. Conceptual ActivityContext does not replace a specific interaction
action or grant authorisation.



Source standing: Bounded adjudicated application information meaning; Domain05 remains the accepted candidate Step-B baseline, not final Solution Architecture. Detailed representation is deferred; only this defined logical object is projected.


{\footnotesize \href{../../../../continuum/application/data-object/978f5c2bd2d343ffaf3a13dff700c258.md}{Canonical record: 978f5c2bd2d343ffaf3a13dff700c258}; full qualifications and provenance.}\par
\subsection*{IV.3.14.3 Services}
\phantomsection
\label{navigation:IV.3.14.3}
\addcontentsline{toc}{subsection}{IV.3.14.3 Services}

\subsubsection*{IV.3.14.3.1 External Interoperability Service}
\phantomsection
\label{navigation:IV.3.14.3.1}
\addcontentsline{toc}{subsubsection}{IV.3.14.3.1 External Interoperability Service}

\label{element:43b79574792443faa57daa76b9a47a19}
\index{External Interoperability Service (IV.3.14.3.1)}
{\small ApplicationService · Candidate — adjudicated Domain05 Step-B baseline; final Solution Architecture unapproved · current}\par

\begin{quote}


\textbf{Pylai is Harmonia's External Interoperability Service, responsible for enabling
and managing system-to-system communication between Harmonia and external systems
and clients, irrespective of communication style, protocol, standard or technology.}



\end{quote}


Pylai may support asynchronous messaging, synchronous or asynchronous remote
method invocation (RMI), APIs, files, polling and other external interoperability
mechanisms. RMI is an architectural interaction concept; it does not mean Java RMI.
HL7/FHIR gateway slices are examples, not the complete Pylai architecture.
Apache Camel or another integration framework may realise external connectors/
technical endpoints; it is a candidate implementation mechanism, not Pylai itself.



External representations conform to their applicable interoperability contracts.
Publication SHALL be non-destructive and fail-closed: only permitted information/
extensions may leave. Harmonia-private authority, security-context, cache,
persistence, governance and concurrency metadata are not automatically exportable.
Egress terminates management of the emitted representation (AX-02/13); retained
provenance does not extend control over the external copy.



The \href{\detokenize{../../../../../docs/markdown/05-information-systems-architecture/gaps.md#pylai-strategy-conflict}}{Strategy/Terma reconciliation} distinguishes
external communication mechanics from Harmonia application meaning and generic
network/infrastructure management.


{\small Recorded unresolved matters: canonicalName; exact scope remains in the canonical record.}\par
{\footnotesize \href{../../../../continuum/application/application-service/43b79574792443faa57daa76b9a47a19.md}{Canonical record: 43b79574792443faa57daa76b9a47a19}; full qualifications and provenance.}\par
\section*{IV.3.15 Themis}
\phantomsection
\label{navigation:IV.3.15}
\addcontentsline{toc}{section}{IV.3.15 Themis}

\subsection*{IV.3.15.1 Application Overview}
\phantomsection
\label{navigation:IV.3.15.1}
\addcontentsline{toc}{subsection}{IV.3.15.1 Application Overview}

\subsubsection*{IV.3.15.1.1 Themis}
\phantomsection
\label{navigation:IV.3.15.1.1}
\addcontentsline{toc}{subsubsection}{IV.3.15.1.1 Themis}

\label{element:899d34ec491948b68aad09e507dea6d1}
\index{Themis (IV.3.15.1.1)}
{\small ApplicationComponent · Candidate — accepted Step-B baseline; final Solution Architecture unapproved · current}\par

\begin{itemize}


\item \textbf{Role/responsibility:} \textbf{Themis owns and manages Harmonia's Authorisation
  Model}, including PAP, PRP, PDP, Criteria/collections, AuthorisationRoles,
  Identity assignments, role/group resolution, Criteria compilation and the
  distributed PEP framework. Themis consumes trusted authenticated Identity;
  it does not authenticate that Identity. \href{\detokenize{../../../../../docs/markdown/05-information-systems-architecture/sources.md#u12}}{U12} refines
  \href{\detokenize{../../../../../docs/markdown/05-information-systems-architecture/sources.md#u08}}{U08} within \hyperref[element:5de2ae2c3b794c2d9f3a14efd51d564f]{AX-07}.

\item \textbf{PAP — Policy Administration Point:} administration and lifecycle management
  of the Themis-owned model. Management/review/approval activity uses
  \href{\detokenize{../../../../../docs/markdown/05-information-systems-architecture/component-register.md#authorisation-administration}}{Ergo/Praxis and Iris}; detailed Praxes/Erga,
  UI and versioning mechanisms remain deferred. No general-purpose policy
  language is implied by policy administration.

\item \textbf{PRP — Policy Repository Point:} authoritative repository semantics for the
  Authorisation Model. Durable persistence may be delegated through \textbf{Themis PRP
  → Mneme → Mnemosyne Harmonia Configuration \& Governance}. This requires no
  dedicated Themis physical database and transfers no model ownership to persistence.

\item \textbf{PDP — Policy Decision Point:} resolves authoritative policy, authenticated
  Identity and applicable context into precomputed Identity-specific
  \textbf{PEP-Criteria}, supplied through \textbf{PEP-Context} for local enforcement.
  An immediate decision remains possible; it uses the same governed semantics.

\item \textbf{Governed actions:} \textbf{Can I ACCESS XYZ? Can I CHANGE ABC? Can I FORWARD DEF?}
  ACCESS, CHANGE and FORWARD are distinct; they SHALL NOT collapse into CRUD.
  FORWARD may depend on destination/context independently of ACCESS.
  Unauthenticated or unauthorised operations default to \textbf{DENY}; unavailable
  authoritative policy determination remains distinct from that decision.

\end{itemize}






\par\medskip\textbf{Identity.Authentication — separately provided trusted Identity}\par


\textbf{Identity.Authentication establishes WHO an Individual or System is. It is a
Separately Provided Capability (SEP).} The primary pattern is standards-based
OIDC/OAuth integration with an externally managed Identity Provider (IdP).
Harmonia does not own or provide an identity-management, credential-management,
password-management or MFA framework.



\begin{quote}


\textbf{Keycloak + PostgreSQL SHALL provide the default FOSS Identity Provider for
development, test and POC environments.}



\end{quote}


Keycloak remains an external, replaceable implementation component. Harmonia's
Identity.Authentication architecture SHALL NOT acquire Keycloak-specific
dependencies. This target does not prescribe a production IdP or an internal
Harmonia credential store. Authentication establishes trusted Identity only;
it SHALL NOT determine Harmonia Function permissions, Function.Scope or PEP-Criteria.



\textbf{Identity.Authorisation} establishes what that authenticated Identity may do
within Harmonia and is Themis-owned. The same authorisation model applies to
human and system Identities; differing allowable assignments do not establish
separate architectures.







\par\medskip\textbf{Identity and role terminology}\par


\begin{longtable}{>{\raggedright\arraybackslash}p{\dimexpr\linewidth/2-2\tabcolsep\relax}>{\raggedright\arraybackslash}p{\dimexpr\linewidth/2-2\tabcolsep\relax}}
Term & Candidate meaning and boundary \\
\textbf{Identity} & The authenticated Individual or System consumed by Themis. This use does not redefine Information Unit Identity or the Domain04 professional identifier assertions. \\
\textbf{Practitioner} & The individual practitioner, retaining the \href{\detokenize{../../../../../docs/markdown/04-information-architecture/information-families/practitioner.md#31-practitioner}}{Domain04 workforce-entity meaning}; registration status is a separate assertion. \\
\textbf{FHIR PractitionerRole} & The formal practitioner role/service relationship represented through FHIR. \\
\textbf{PractitionerCapacity} & The selected operational capacity in which a Practitioner is currently acting within Harmonia; it does not itself confer authority. \\
\textbf{AuthorisationRole} & A named, governed, assignable collection of additive Themis Criteria and/or Criteria collections. It is distinct from PractitionerCapacity, FHIR PractitionerRole and a Domain03 Business Role. \\
\textbf{Criterion / Criteria} & One governed permitted combination / multiple such combinations in the Themis Authorisation Model. \\
\textbf{PEP-Criteria} & The precomputed, Identity-specific set of permitted combinations derived by Themis for constrained local enforcement. \\
\textbf{PEP-Context} & The Themis-derived operational enforcement context for an authenticated Identity, through which applicable PEP-Criteria is supplied. It is distinct from the general Mneme working Context. \\
\end{longtable}


The \href{\detokenize{../../../../../docs/markdown/04-information-architecture/information-families/practitioner.md#36-practitioner-role-binding}}{Domain04 Practitioner Role Binding}
and \href{\detokenize{../../../../../docs/markdown/03-business-architecture/actors-roles/roles.md}}{Domain03 Business Roles}
retain their meanings. Neither lexical similarity nor FHIR convention establishes
their mapping to PractitionerCapacity or AuthorisationRole.







\par\medskip\textbf{Criterion conjunction and additive grants}\par


A Criterion represents a \textbf{governed permitted combination} across applicable
dimensions, using \href{\detokenize{../../../../../docs/markdown/05-information-systems-architecture/component-register.md#authorisation-nomenclature}}{Calliope-controlled nomenclature}:



\begin{quote}\small\ttfamily
Criterion =\\
    Information.Compartment\\
    AND Information.Classification\\
    AND Function.Name\\
    AND Function.Scope\\
    AND applicable Temporal constraints\\

\end{quote}


The four principal dimensions describe \textbf{what} is authorised; applicable
temporal constraints describe \textbf{when}. Within an individual Criterion,
applicable dimensions are \textbf{conjunctive}. Between Criteria, grants are
\textbf{additive / UNION}: each Criterion independently contributes its permitted
combination. Only explicitly governed combinations have authorisation meaning.
Absence of a matching effective Criterion results in \textbf{DENY}.



\textbf{Explicit negative/DENY Criteria are not part of this Candidate Architecture.}
Neither arbitrary Boolean policy expressions, deny-overrides semantics nor a
general-purpose policy language is introduced by this model.







\par\medskip\textbf{Criteria containment, AuthorisationRoles and Identity assignments}\par


Criteria support containment, collection and nesting for management and reuse:



\begin{quote}\small\ttfamily
CriteriaCollection                 (conceptual label; final name unresolved)\\
  ├── Criterion\\
  ├── Criterion\\
  └── CriteriaCollection\\
       ├── Criterion\\
       └── Criterion\\

\end{quote}


Flattening nested collections produces the same additive set of effective
Criteria. Nesting creates no arbitrary Boolean semantics and requires no Role
hierarchy. The exact implementation name/class hierarchy and representation
remain deferred; this diagram establishes no lifecycle/ownership semantics
beyond the stated management/reuse requirement.



An \textbf{AuthorisationRole} is named, governed and assignable and contains Criteria
and/or Criteria collections. Its contents are additive:



\begin{quote}\small\ttfamily
EffectiveCriteria(AuthorisationRole)\\
    = UNION(all contained Criteria)\\
\\
IdentityEffectiveCriteria =\\
    DirectCriteria\\
    UNION AuthorisationRole1.Criteria\\
    UNION AuthorisationRole2.Criteria\\
    UNION ...\\

\end{quote}


An Identity may have directly assigned Criteria, one or more AuthorisationRoles,
or both. This applies to human and system Identities. The UNION is of applicable
grants, subject to governed assignments, temporal validity, revocation and
applicable operational context; it does not remove a Criterion's conjunctions
or require every contained Criterion to match simultaneously.







\par\medskip\textbf{Predetermined temporal validity and explicit revocation}\par


\textbf{Temporal validity} establishes predetermined applicability/validity bounds.
Temporal constraints may apply to a Criterion, a direct Criterion assignment,
an AuthorisationRole assignment and applicable operational context such as
PractitionerCapacity. The exact temporal expression model remains deferred.



\textbf{Revocation} withdraws an otherwise valid authorisation. It SHALL NOT be
modelled merely as temporal expiry. Both concepts are required. Propagation of
revocation and invalidation of already-issued PEP-Criteria/PEP-Context remain
unresolved; no mechanism or continued validity after revocation is inferred.







\par\medskip\textbf{PractitionerCapacity and operational responsibility}\par


For a Practitioner, the available PractitionerCapacity set is constrained by
the practitioner's governed profile and applicable operational context. The
Practitioner selects \textbf{one active PractitionerCapacity} for the applicable
session/operational or roster context. It supplies operational context used by
Themis when deriving effective authorisation; \textbf{it does not itself confer authority}.
Selection UI and the exact PractitionerCapacity ↔ FHIR PractitionerRole mapping
remain deferred.



\begin{quote}


\textbf{Authentication establishes who the Individual is. PractitionerCapacity
establishes the capacity in which a Practitioner is currently acting. Themis
establishes what that Identity, acting in that capacity, may do.}



\end{quote}


\begin{quote}\small\ttfamily
graph TD\\
    Auth["Authenticated Practitioner"] --> Eligible["Eligible PractitionerCapacity set"]\\
    Profile["Governed practitioner profile"] --> Eligible\\
    Context["Applicable operational / roster context"] --> Eligible\\
    Eligible -->|"Practitioner selects ONE"| Capacity["Active PractitionerCapacity: context, no authority by itself"]\\
    Capacity -->|"authorisation context"| PDP["Themis: resolve governed Criteria / assignments / policy state"]\\
    PDP --> Criteria["Identity-specific precomputed PEP-Criteria through PEP-Context"]\\
    Criteria --> PEP["Constrained local Themis PEP"]\\
    Capacity -.->|"operational responsibility context"| Responsibility["Agora: membership intent; detailed context use deferred"]\\

\end{quote}


Themis answers \textbf{“May this Identity perform this action on this information?”}
Future Agora collaboration detail concerns \textbf{“Is this Practitioner currently
responsible for receiving/responding to this operational activity?”} These are
separate responsibilities even where they consume the same PractitionerCapacity.
Receipt of an escalation or approval notification grants no information access.
Following it into Iris or another application surface requires normal Themis PEP
enforcement using current valid PEP-Criteria. U16 establishes
\hyperref[element:d1c4053e68384080b8eb79d2008632ba]{Agora control and Omilia operations};
detailed membership, escalation and approval-routing mechanics remain deferred.







\par\medskip\textbf{PEP-Context and distributed enforcement}\par


\begin{quote}


\textbf{PEP-Context is the Themis-derived policy enforcement context established for
an authenticated Identity. It supplies the precomputed PEP-Criteria applicable
to that Identity that a Themis Policy Enforcement Point requires for local
enforcement decisions.}



\end{quote}


\textbf{PEP-Criteria is a precomputed, Identity-specific set of permitted combinations},
not a general-purpose executable policy language distributed to PEPs. Themis
resolves direct Criteria, AuthorisationRole-derived Criteria, nested collections,
applicable temporal constraints, PractitionerCapacity context where applicable
and authoritative policy state. PEPs SHALL NOT receive the complete model merely
because they need enforcement information.



\begin{quote}\small\ttfamily
Policy                Authorisation Model managed by Themis\\
Criteria              governed permitted combinations in the Authorisation Model\\
PEP-Criteria          precomputed Identity-specific permitted combinations\\
PEP-Context           operational context supplying applicable PEP-Criteria\\
Enforcement Decision  local PEP determination for a particular proposed action\\

\end{quote}


Enforcement occurs at the component capable of preventing the action. A PEP may
be a Themis-provided library/framework element embedded within a host module.
\textbf{The distributed PEP capability, implementation and enforcement semantics are
Themis-owned/derived/built.} The host controls the actual operation; Themis
supplies enforcement mechanism and semantics. Hosts SHALL NOT independently
reinterpret the Authorisation Model in place of the Themis PEP. No separate
central Policy Enforcement Software Component is created.



A local PEP SHALL perform deliberately \textbf{constrained, deterministic enforcement}
against authenticated Identity context, effective Information.Compartment and
Information.Classification, Function.Name, Function.Scope, applicable temporal
state and the Identity's valid PEP-Criteria. It SHALL NOT become a general-purpose
policy interpreter.



Themis SHALL retain sufficient provenance to explain effective PEP-Criteria
derivation: authoritative Criteria, collections, direct/AuthorisationRole
assignments and applicable policy state. Runtime PEP-Criteria need not carry
the full derivation history; stable references may support it. No final wire/data
structure, cache, refresh or distribution mechanism is prescribed.



PEP-Context and existing security-context contracts are operational context;
they do not automatically become persisted resource content. Their detailed
relationship, schema and distribution mechanisms remain unresolved;
PEP-Context is not silently equated with a historical class or transport carrier.







\par\medskip\textbf{Untagged information and effective security attributes}\par


\begin{quote}


\textbf{Absence of explicit information security metadata does not imply absence
of security control. Untagged information is not unrestricted information.}



\end{quote}


Where required Information.Compartment and/or Information.Classification
attributes are absent, effective values SHALL be established through governed,
deterministic rules. If required effective attributes cannot be established,
enforcement SHALL fail closed.



Deriving effective attributes for an authorisation decision is distinct from
persistently assigning/changing security metadata on an Information Unit.
Themis enforces the resulting attributes; it is not assumed to own all semantic
derivation rules. Exact derivation ownership and precedence, including possible
context/containment relationships, remain \href{\detokenize{../../../../../docs/markdown/05-information-systems-architecture/gaps.md#identity-deferred-design}}{deferred}.
Mnemosyne's established compartment validation/reapplication responsibility does
not settle those semantic rules.







\par\medskip\textbf{Availability and delegation}\par


An in-flight activity with a \textbf{valid PEP-Context sufficient for its required
enforcement decisions} may continue during temporary PDP unavailability.
PDP availability is required for a \textbf{new, renewed or replacement PEP-Context}.
If the PDP cannot establish that required context, the activity SHALL NOT proceed
beyond the point requiring it and must wait. \textbf{PDP unavailability is not policy
denial}: policy permits, policy denies, and authoritative policy determination
cannot presently be established remain distinct. This grants no unauthenticated
or unauthorised action permission and weakens no required audit-capture boundary.
Detailed validity, timeout/retry, revocation and renewal mechanisms are unresolved.



At this stage delegation is \textbf{Identity / PEP-Criteria pass-off (forwarding)},
subject to applicable \textbf{FORWARD} authority. Provenance SHALL preserve both the
originating Identity and receiving/acting Identity; the receiver does not become
the originator. No separate general-purpose delegation subsystem is introduced.
Detailed forwarding lifecycle/representation, PDP interaction and governance
mechanisms remain deferred; forwarding establishes no unrestricted impersonation
or authority-broadening rule.







\par\medskip\textbf{Governed authorisation administration and target allocation}\par


Authorisation administration is governed Harmonia work. Criteria creation/change,
collection management, AuthorisationRole management, direct Identity Criteria
assignment, AuthorisationRole assignment, temporal changes and revocation SHALL
be subject to appropriate governed management/review/approval processes through
\textbf{Ergo/Praxis}, with \textbf{Iris} providing human interaction. Detailed PAP management
Praxes/Erga and approval rules remain deferred; this does not make every change
require an identical approval sequence.



\begin{longtable}{>{\raggedright\arraybackslash}p{\dimexpr\linewidth/2-2\tabcolsep\relax}>{\raggedright\arraybackslash}p{\dimexpr\linewidth/2-2\tabcolsep\relax}}
Responsibility & Candidate target \\
Identity authentication & Keycloak / external OIDC IdP; SEP, Keycloak + PostgreSQL default for development/test/POC \\
Controlled authorisation nomenclature & Calliope \\
Authorisation model, Criteria/collections, AuthorisationRoles and Identity assignments & Themis-owned \\
Role/group resolution & Themis-owned; no separate human/system authorisation architecture \\
Criteria compilation & Themis-owned \\
Policy repository & Themis PRP, persisted through Harmonia via Mneme / Mnemosyne C\&G \\
Policy administration & Themis PAP \\
Governed management / review / approval workflows & Ergo/Praxis \\
Administration UI & Iris; no authorisation-model ownership or detailed reusable-form design \\
Runtime enforcement & Themis PEP library at action-controlling hosts \\
Untagged information handling & Governed effective-attribute resolution; exact ownership/precedence deferred \\
\end{longtable}


Harmonia owns its authorisation semantics. \textbf{jCasbin, OPA, Cedar or another
general-purpose policy framework is not an architectural dependency.} FOSS
implementations may later supply implementation references where appropriate;
their conceptual models SHALL NOT redefine Harmonia semantics.



\begin{itemize}


\item \textbf{Collaboration/exclusions:} governed access, activity, publication and
  collaboration hosts use Themis-owned PEP capability. Petasos hosts governed
  FORWARD enforcement because it controls the queue operation. Praxis may supply
  relevant context but is not the Authorisation Model; authorisation does not
  establish Praxis conformance. Themis is neither Audit nor Assurance. Significant
  policy/enforcement occurrences are reported through Kleio Audit Capture under
  applicable policy; clients need not construct FHIR AuditEvent.
  \textbf{Calliope defines security/assurance nomenclature; Themis decides/enforces;
  Dokimasia assures; Kleio audits.} Governance, Authorisation, Assurance and Audit
  remain distinct responsibilities; Praxis conformance and assurance findings
  are not authorisation grants. Dokimasia terminal constraints remain unchanged.

\item \textbf{Standing:} U12 consolidates the adjudicated Identity Candidate Architecture
  and refines U08's Criteria, policy-language deferral and delegation wording.
  It does not accept all Domain05 candidate material or begin Step B.
  \href{\detokenize{../../../../../docs/markdown/05-information-systems-architecture/gaps.md#themis-detailed-contracts}}{CDG-42} retains detailed design.
  U15 resolves the former U08 §15 conflict through Kleio-specific
  \href{\detokenize{../../../../../docs/markdown/05-information-systems-architecture/component-register.md#kleio-capture-authority}}{mandatory capture authority}, governed under AX-07.
  Historical fixed gate counts and automatic context persistence remain
  CDG-15/30 provenance, without reopening established security boundaries.

\end{itemize}

{\small Recorded unresolved matters: canonicalName; exact scope remains in the canonical record.}\par
{\footnotesize \href{../../../../continuum/application/application-component/899d34ec491948b68aad09e507dea6d1.md}{Canonical record: 899d34ec491948b68aad09e507dea6d1}; full qualifications and provenance.}\par
\subsection*{IV.3.15.2 Information Architecture}
\phantomsection
\label{navigation:IV.3.15.2}
\addcontentsline{toc}{subsection}{IV.3.15.2 Information Architecture}

\subsubsection*{IV.3.15.2.1 PEP-Context}
\phantomsection
\label{navigation:IV.3.15.2.1}
\addcontentsline{toc}{subsubsection}{IV.3.15.2.1 PEP-Context}

\label{element:90c3f7e03448452baec5c425c34fd77a}
\index{PEP-Context (IV.3.15.2.1)}
{\small DataObject · Candidate — adjudicated Domain05 Step-B baseline; final Solution Architecture unapproved · current}\par

\begin{quote}


\textbf{PEP-Context is the Themis-derived policy enforcement context established for
an authenticated Identity. It supplies the precomputed PEP-Criteria applicable
to that Identity that a Themis Policy Enforcement Point requires for local
enforcement decisions.}



\end{quote}


\textbf{PEP-Criteria is a precomputed, Identity-specific set of permitted combinations},
not a general-purpose executable policy language distributed to PEPs. Themis
resolves direct Criteria, AuthorisationRole-derived Criteria, nested collections,
applicable temporal constraints, PractitionerCapacity context where applicable
and authoritative policy state. PEPs SHALL NOT receive the complete model merely
because they need enforcement information.



\begin{quote}\small\ttfamily
Policy                Authorisation Model managed by Themis\\
Criteria              governed permitted combinations in the Authorisation Model\\
PEP-Criteria          precomputed Identity-specific permitted combinations\\
PEP-Context           operational context supplying applicable PEP-Criteria\\
Enforcement Decision  local PEP determination for a particular proposed action\\

\end{quote}


Enforcement occurs at the component capable of preventing the action. A PEP may
be a Themis-provided library/framework element embedded within a host module.
\textbf{The distributed PEP capability, implementation and enforcement semantics are
Themis-owned/derived/built.} The host controls the actual operation; Themis
supplies enforcement mechanism and semantics. Hosts SHALL NOT independently
reinterpret the Authorisation Model in place of the Themis PEP. No separate
central Policy Enforcement Software Component is created.



A local PEP SHALL perform deliberately \textbf{constrained, deterministic enforcement}
against authenticated Identity context, effective Information.Compartment and
Information.Classification, Function.Name, Function.Scope, applicable temporal
state and the Identity's valid PEP-Criteria. It SHALL NOT become a general-purpose
policy interpreter.



Themis SHALL retain sufficient provenance to explain effective PEP-Criteria
derivation: authoritative Criteria, collections, direct/AuthorisationRole
assignments and applicable policy state. Runtime PEP-Criteria need not carry
the full derivation history; stable references may support it. No final wire/data
structure, cache, refresh or distribution mechanism is prescribed.



PEP-Context and existing security-context contracts are operational context;
they do not automatically become persisted resource content. Their detailed
relationship, schema and distribution mechanisms remain unresolved;
PEP-Context is not silently equated with a historical class or transport carrier.


{\small Recorded unresolved matters: canonicalName; exact scope remains in the canonical record.}\par
{\footnotesize \href{../../../../continuum/application/data-object/90c3f7e03448452baec5c425c34fd77a.md}{Canonical record: 90c3f7e03448452baec5c425c34fd77a}; full qualifications and provenance.}\par
\subsubsection*{IV.3.15.2.2 PEP-Criteria}
\phantomsection
\label{navigation:IV.3.15.2.2}
\addcontentsline{toc}{subsubsection}{IV.3.15.2.2 PEP-Criteria}

\label{element:7e10287b7e634bce8aa2481eb481bb76}
\index{PEP-Criteria (IV.3.15.2.2)}
{\small DataObject · Candidate — adjudicated Domain05 Step-B baseline; final Solution Architecture unapproved · current}\par

\textbf{PEP-Criteria is a precomputed, Identity-specific set of permitted combinations},
not a general-purpose executable policy language distributed to PEPs. Themis
resolves direct Criteria, AuthorisationRole-derived Criteria, nested collections,
applicable temporal constraints, PractitionerCapacity context where applicable
and authoritative policy state. PEPs SHALL NOT receive the complete model merely
because they need enforcement information.



Themis derives this information; PEP-Context supplies it to constrained local enforcement. It is not a general-purpose policy language or a persisted resource by default.


{\small Recorded unresolved matters: canonicalName; exact scope remains in the canonical record.}\par
{\footnotesize \href{../../../../continuum/application/data-object/7e10287b7e634bce8aa2481eb481bb76.md}{Canonical record: 7e10287b7e634bce8aa2481eb481bb76}; full qualifications and provenance.}\par
\subsubsection*{IV.3.15.2.3 Criterion}
\phantomsection
\label{navigation:IV.3.15.2.3}
\addcontentsline{toc}{subsubsection}{IV.3.15.2.3 Criterion}

\label{element:b5be4886ec77466e860209443a03771b}
\index{Criterion (IV.3.15.2.3)}
{\small DataObject · Candidate — adjudicated Domain05 Step-B baseline; final Solution Architecture unapproved · current}\par

A Criterion represents a \textbf{governed permitted combination} across applicable
dimensions, using \href{\detokenize{../../../../../docs/markdown/05-information-systems-architecture/component-register.md#authorisation-nomenclature}}{Calliope-controlled nomenclature}:



\begin{quote}\small\ttfamily
Criterion =\\
    Information.Compartment\\
    AND Information.Classification\\
    AND Function.Name\\
    AND Function.Scope\\
    AND applicable Temporal constraints\\

\end{quote}


The four principal dimensions describe \textbf{what} is authorised; applicable
temporal constraints describe \textbf{when}. Within an individual Criterion,
applicable dimensions are \textbf{conjunctive}. Between Criteria, grants are
\textbf{additive / UNION}: each Criterion independently contributes its permitted
combination. Only explicitly governed combinations have authorisation meaning.
Absence of a matching effective Criterion results in \textbf{DENY}.



\textbf{Explicit negative/DENY Criteria are not part of this Candidate Architecture.}
Neither arbitrary Boolean policy expressions, deny-overrides semantics nor a
general-purpose policy language is introduced by this model.


{\small Recorded unresolved matters: canonicalName; exact scope remains in the canonical record.}\par
{\footnotesize \href{../../../../continuum/application/data-object/b5be4886ec77466e860209443a03771b.md}{Canonical record: b5be4886ec77466e860209443a03771b}; full qualifications and provenance.}\par
\subsubsection*{IV.3.15.2.4 AuthorisationRole}
\phantomsection
\label{navigation:IV.3.15.2.4}
\addcontentsline{toc}{subsubsection}{IV.3.15.2.4 AuthorisationRole}

\label{element:bad6431f87bd4bf59645c0962cc219a7}
\index{AuthorisationRole (IV.3.15.2.4)}
{\small DataObject · Candidate — adjudicated Domain05 Step-B baseline; final Solution Architecture unapproved · current}\par

Criteria support containment, collection and nesting for management and reuse:



\begin{quote}\small\ttfamily
CriteriaCollection                 (conceptual label; final name unresolved)\\
  ├── Criterion\\
  ├── Criterion\\
  └── CriteriaCollection\\
       ├── Criterion\\
       └── Criterion\\

\end{quote}


Flattening nested collections produces the same additive set of effective
Criteria. Nesting creates no arbitrary Boolean semantics and requires no Role
hierarchy. The exact implementation name/class hierarchy and representation
remain deferred; this diagram establishes no lifecycle/ownership semantics
beyond the stated management/reuse requirement.



An \textbf{AuthorisationRole} is named, governed and assignable and contains Criteria
and/or Criteria collections. Its contents are additive:



\begin{quote}\small\ttfamily
EffectiveCriteria(AuthorisationRole)\\
    = UNION(all contained Criteria)\\
\\
IdentityEffectiveCriteria =\\
    DirectCriteria\\
    UNION AuthorisationRole1.Criteria\\
    UNION AuthorisationRole2.Criteria\\
    UNION ...\\

\end{quote}


An Identity may have directly assigned Criteria, one or more AuthorisationRoles,
or both. This applies to human and system Identities. The UNION is of applicable
grants, subject to governed assignments, temporal validity, revocation and
applicable operational context; it does not remove a Criterion's conjunctions
or require every contained Criterion to match simultaneously.


{\small Recorded unresolved matters: canonicalName; exact scope remains in the canonical record.}\par
{\footnotesize \href{../../../../continuum/application/data-object/bad6431f87bd4bf59645c0962cc219a7.md}{Canonical record: bad6431f87bd4bf59645c0962cc219a7}; full qualifications and provenance.}\par

\part*{Part V — Technology}
\label{navigation:V}
\addcontentsline{toc}{part}{Part V — Technology}
No individually represented Technology architecture is present in the current Continuum. No data, platform or infrastructure subsections are inferred.\par

\part*{Part VI — Implementation \& Migration}
\label{navigation:VI}
\addcontentsline{toc}{part}{Part VI — Implementation \& Migration}
No individually represented architectural release-family transition is present in the current Continuum. Governance records and execution tasks remain in their existing governance repository.\par

\part*{Reference navigation — R Cross-Cutting References}
\label{navigation:R}
\addcontentsline{toc}{part}{Reference navigation — R Cross-Cutting References}
Navigation and discovery only, outside the architectural decomposition of Parts I–VI. Canonical relationships establish associations; index placement does not.\par
\chapter*{R.1 Stakeholder Index — Reference navigation}
\phantomsection
\label{navigation:R.1}
\addcontentsline{toc}{chapter}{R.1 Stakeholder Index — Reference navigation}

All canonical Stakeholders, including those without represented Motivation associations. This index is cross-cutting reference material, not a Motivation architectural subject.\par
\section*{R.1.1 Care Team — Stakeholder index entry [historical\_reference]}
\phantomsection
\label{navigation:R.1.1}
\addcontentsline{toc}{section}{R.1.1 Care Team — Stakeholder index entry [historical\_reference]}

\label{element:6f5a203d81e24a81b384a80781987daa}
\index{Care Team (R.1.1)}
{\small Stakeholder · Historical — current applicability unresolved · historical\_reference}\par
{\small Existing associated Motivation material: \hyperref[element:babb9f5d63d8435890c1db3579ec78d9]{Clinical Interoperability} (I.1.3); canonical associations: id-457a3dd70bb358ff8989bc3c9c725c70.}\par
\href{../../../../continuum/motivation/stakeholder/6f5a203d81e24a81b384a80781987daa.md}{Historical canonical definition and provenance}.\par
\section*{R.1.2 Clinicians \& Care Teams — Stakeholder index entry}
\phantomsection
\label{navigation:R.1.2}
\addcontentsline{toc}{section}{R.1.2 Clinicians \& Care Teams — Stakeholder index entry}

\label{element:37264af5adb64b37b7537e3d48b6747b}
\index{Clinicians \& Care Teams (R.1.2)}
{\small Stakeholder · Agreed · current}\par
No canonical Motivation association is represented. Index placement asserts none.\par

\par\medskip\textbf{Profile \& Role}\par


Clinicians and Care Teams represent the front-line healthcare workforce: medical practitioners, nurses, clinical specialists, allied health professionals, pharmacists, and multi-disciplinary care coordinators operating at the direct point of care.



\par\medskip\textbf{Core Responsibilities}\par


\begin{itemize}


\item Diagnosing clinical conditions, formulating treatment plans, and administering patient care.

\item Ordering diagnostic investigations, prescribing medications, and documenting clinical findings.

\item Reviewing historical health records, previous test results, and discharge summaries to inform clinical decision-making.

\end{itemize}


\par\medskip\textbf{Principal Architectural Concerns}\par


\begin{itemize}


\item \textbf{Timely Information Availability}: Point-of-care access to complete, up-to-date clinical information without lag or missing historical encounters.

\item \textbf{Semantic Fidelity \& Clinical Context}: Clinical observations, units of measure, reference ranges, and diagnostic notes must be preserved without lossy truncation, confusion of normal/abnormal flags, or distortion of clinical nuance.

\item \textbf{Elimination of Data Fragmentation}: Synthesized, longitudinal patient views that reduce duplicate diagnostic testing and eliminate dangerous blind spots.

\item \textbf{Explicit Information Authority}: Clear visual and semantic distinction between authoritative hospital-validated records and preliminary, informational, or anecdotal patient-reported data.

\end{itemize}

{\footnotesize \href{../../../../continuum/motivation/stakeholder/37264af5adb64b37b7537e3d48b6747b.md}{Canonical record: 37264af5adb64b37b7537e3d48b6747b}; full qualifications and provenance.}\par
\section*{R.1.3 Health Network — Stakeholder index entry [historical\_reference]}
\phantomsection
\label{navigation:R.1.3}
\addcontentsline{toc}{section}{R.1.3 Health Network — Stakeholder index entry [historical\_reference]}

\label{element:55e57c01dfe146a2a225228b31196862}
\index{Health Network (R.1.3)}
{\small Stakeholder · Historical — current applicability unresolved · historical\_reference}\par
{\small Existing associated Motivation material: \hyperref[element:babb9f5d63d8435890c1db3579ec78d9]{Clinical Interoperability} (I.1.3); canonical associations: id-d5fa8c64975852f88132df579c806836.}\par
\href{../../../../continuum/motivation/stakeholder/55e57c01dfe146a2a225228b31196862.md}{Historical canonical definition and provenance}.\par
\section*{R.1.4 Healthcare Delivery Organizations — Stakeholder index entry}
\phantomsection
\label{navigation:R.1.4}
\addcontentsline{toc}{section}{R.1.4 Healthcare Delivery Organizations — Stakeholder index entry}

\label{element:31a8e1054c02483a9d5e314a91a91364}
\index{Healthcare Delivery Organizations (R.1.4)}
{\small Stakeholder · Agreed · current}\par
No canonical Motivation association is represented. Index placement asserts none.\par

\par\medskip\textbf{Profile \& Role}\par


Healthcare Delivery Organizations (HDOs) encompass the institutional entities providing clinical care across the region: acute public hospitals, private hospital networks, primary health clinics, specialist outpatient centers, and diagnostic pathology and medical imaging facilities.



\par\medskip\textbf{Core Responsibilities}\par


\begin{itemize}


\item Providing direct patient care across emergency, inpatient, ambulatory, and diagnostic settings.

\item Generating, maintaining, and ingesting operational clinical transactions (admissions, transfers, discharges, orders, results, clinical notes) through local hospital information systems (HIS), patient administration systems (PAS), and laboratory information systems (LIS).

\item Maintaining local boundary integration gateways and adhering to regional messaging protocols.

\end{itemize}


\par\medskip\textbf{Principal Architectural Concerns}\par


\begin{itemize}


\item \textbf{Reliable Ingress \& Egress Integration}: Guaranteeing that clinical messages dispatched to Harmonia are reliably accepted and preserved without silent loss or unhandled dropouts.

\item \textbf{Low Latency \& High Throughput}: Fast turnaround times for critical pathology results, radiology reports, and patient admission feeds without queuing delays.

\item \textbf{Standards-Compliant Boundary Interoperability}: Seamless support for legacy messaging standards (HL7 v2.x via MLLP) alongside modern FHIR RESTful interactions, avoiding disruptive upgrades to legacy hospital software.

\item \textbf{Clear Demarcation of Responsibility}: Explicit integration boundaries defining where the delivery organization's operational responsibility ends and Harmonia's management begins.

\end{itemize}

{\footnotesize \href{../../../../continuum/motivation/stakeholder/31a8e1054c02483a9d5e314a91a91364.md}{Canonical record: 31a8e1054c02483a9d5e314a91a91364}; full qualifications and provenance.}\par
\section*{R.1.5 Healthcare Directory Stewards \& Registrars — Stakeholder index entry}
\phantomsection
\label{navigation:R.1.5}
\addcontentsline{toc}{section}{R.1.5 Healthcare Directory Stewards \& Registrars — Stakeholder index entry}

\label{element:1c8e2c898a28436693cf2ae69776e125}
\index{Healthcare Directory Stewards \& Registrars (R.1.5)}
{\small Stakeholder · Agreed · current}\par
No canonical Motivation association is represented. Index placement asserts none.\par

\par\medskip\textbf{Profile \& Role}\par


Healthcare Directory Stewards and Registrars are the specialized operational custodians responsible for maintaining authoritative regional registries of healthcare providers, individual clinicians, clinical organizations, and electronic service delivery endpoints.



\par\medskip\textbf{Core Responsibilities}\par


\begin{itemize}


\item Curating authoritative master directories of healthcare practitioners, qualifications, active credentials, and organizational affiliations.

\item Verifying and maintaining electronic service delivery endpoints (secure messaging addresses, FHIR endpoints, MLLP port mappings).

\item Managing provider registry updates, deprecations, practitioner relocations, and retirement events.

\end{itemize}


\par\medskip\textbf{Principal Architectural Concerns}\par


\begin{itemize}


\item \textbf{Secure Endpoint Governance}: Ensuring clinical messages and notifications are routed exclusively to authenticated, currently authorized practitioner endpoints.

\item \textbf{Accurate Practitioner Role Representation}: Distinguishing between an individual clinician, their legal practitioner role, their accredited specialty, and their specific institutional practice location.

\item \textbf{Reliable Directory Change Propagation}: Ensuring directory updates (e.g., changes in clinic address or revocation of practitioner credentials) propagate through integration workflows rapidly and deterministically.

\end{itemize}

{\footnotesize \href{../../../../continuum/motivation/stakeholder/1c8e2c898a28436693cf2ae69776e125.md}{Canonical record: 1c8e2c898a28436693cf2ae69776e125}; full qualifications and provenance.}\par
\section*{R.1.6 Healthcare Provider — Stakeholder index entry [historical\_reference]}
\phantomsection
\label{navigation:R.1.6}
\addcontentsline{toc}{section}{R.1.6 Healthcare Provider — Stakeholder index entry [historical\_reference]}

\label{element:2018200a493d458096b72031cc2ae779}
\index{Healthcare Provider (R.1.6)}
{\small Stakeholder · Historical — current applicability unresolved · historical\_reference}\par
{\small Existing associated Motivation material: \hyperref[element:782f53e2190f4a8a8d0dc0a33cd6fbb3]{24/7 Availability} (I.1.1); canonical associations: id-78a34270ced656e7901d367bfd405999.}\par
\href{../../../../continuum/motivation/stakeholder/2018200a493d458096b72031cc2ae779.md}{Historical canonical definition and provenance}.\par
\section*{R.1.7 Patient — Stakeholder index entry [historical\_reference]}
\phantomsection
\label{navigation:R.1.7}
\addcontentsline{toc}{section}{R.1.7 Patient — Stakeholder index entry [historical\_reference]}

\label{element:5f6d333e94634420b8a8bbf8b9cb6e2a}
\index{Patient (R.1.7)}
{\small Stakeholder · Historical — current applicability unresolved · historical\_reference}\par
{\small Existing associated Motivation material: \hyperref[element:499313a0a5d84520a3061120ae9b6dac]{Audit \& PHI Privacy} (I.1.2); canonical associations: id-daee593684c55e55bf1ce04c223cf585.}\par
\href{../../../../continuum/motivation/stakeholder/5f6d333e94634420b8a8bbf8b9cb6e2a.md}{Historical canonical definition and provenance}.\par
\section*{R.1.8 Patients \& Care Recipients — Stakeholder index entry}
\phantomsection
\label{navigation:R.1.8}
\addcontentsline{toc}{section}{R.1.8 Patients \& Care Recipients — Stakeholder index entry}

\label{element:ea5a315450ce4cfb87417263560c2b6d}
\index{Patients \& Care Recipients (R.1.8)}
{\small Stakeholder · Agreed · current}\par
No canonical Motivation association is represented. Index placement asserts none.\par

\par\medskip\textbf{Profile \& Role}\par


Patients and Care Recipients are the individual citizens, families, and communities receiving healthcare services across the regional network. They are the ultimate beneficiaries and primary subjects of all health data managed by Harmonia.



\par\medskip\textbf{Core Responsibilities}\par


\begin{itemize}


\item Engaging with clinicians and healthcare providers throughout episodes of care.

\item Providing accurate demographic, personal, and medical history.

\item Exercising personal privacy choices and statutory rights regarding information sharing where applicable.

\end{itemize}


\par\medskip\textbf{Principal Architectural Concerns}\par


\begin{itemize}


\item \textbf{Patient Safety \& Identity Protection}: Guaranteeing that personal clinical events are never cross-associated with another individual, preventing severe medication errors, misdiagnoses, and incorrect surgical interventions.

\item \textbf{Confidentiality \& Privacy}: Ensuring sensitive health records (mental health, sexual health, genetic testing) are strictly protected against unauthorized access, eavesdropping, and inadvertent leakage into operational log streams.

\item \textbf{Continuity of Care}: Seamless information transfer between disparate healthcare facilities so that clinical history follows the patient across the continuum of care.

\end{itemize}

{\footnotesize \href{../../../../continuum/motivation/stakeholder/ea5a315450ce4cfb87417263560c2b6d.md}{Canonical record: ea5a315450ce4cfb87417263560c2b6d}; full qualifications and provenance.}\par
\section*{R.1.9 Platform Operations \& Integration Engineers — Stakeholder index entry}
\phantomsection
\label{navigation:R.1.9}
\addcontentsline{toc}{section}{R.1.9 Platform Operations \& Integration Engineers — Stakeholder index entry}

\label{element:c20aa17ce4aa4a58be5eb3cff2ecdd1d}
\index{Platform Operations \& Integration Engineers (R.1.9)}
{\small Stakeholder · Agreed · current}\par
No canonical Motivation association is represented. Index placement asserts none.\par

\par\medskip\textbf{Profile \& Role}\par


Platform Operations and Integration Engineers are the systems reliability engineers (SREs), DevOps professionals, integration middleware developers, and technical support teams tasked with running, monitoring, and maintaining the Harmonia platform.



\par\medskip\textbf{Core Responsibilities}\par


\begin{itemize}


\item Deploying, scaling, and configuring Harmonia runtime services across regional data centers and cloud infrastructure.

\item Monitoring system health, message queues, throughput, network latency, and integration error rates.

\item Diagnosing integration failures, investigating communication anomalies, and conducting incident retrospectives.

\end{itemize}


\par\medskip\textbf{Principal Architectural Concerns}\par


\begin{itemize}


\item \textbf{Deep Diagnosability Without PHI Leakage}: Rich operational metrics, traces, and diagnostic logs that provide full visibility into message processing without emitting Protected Health Information (PHI) into non-clinical tooling.

\item \textbf{Zero Unmonitored Message Loss}: Deterministic verification regarding the progression state of incoming events, ensuring that failed downstream deliveries are immediately flagged rather than lost in silent dead-letter sinks.

\item \textbf{Observable Multi-Stage Activity State}: Transparent visibility into complex, distributed workflows (such as fan-out distributions) where tasks progress across multiple hops and destination endpoints.

\item \textbf{Deterministic Failure Semantics}: Clean separation of active operational coordination from authoritative durable state, preventing data corruption during unexpected node crashes or network partitions.

\end{itemize}

{\footnotesize \href{../../../../continuum/motivation/stakeholder/c20aa17ce4aa4a58be5eb3cff2ecdd1d.md}{Canonical record: c20aa17ce4aa4a58be5eb3cff2ecdd1d}; full qualifications and provenance.}\par
\section*{R.1.10 Regional Health Network Operator — Stakeholder index entry}
\phantomsection
\label{navigation:R.1.10}
\addcontentsline{toc}{section}{R.1.10 Regional Health Network Operator — Stakeholder index entry}

\label{element:e76216c494fa4500bb80eda6f83513c1}
\index{Regional Health Network Operator (R.1.10)}
{\small Stakeholder · Agreed · current}\par
No canonical Motivation association is represented. Index placement asserts none.\par

\par\medskip\textbf{Profile \& Role}\par


The Regional Health Network Operator is the statutory or corporate governing body responsible for orchestrating, funding, and overseeing healthcare delivery across a defined regional jurisdiction (e.g., a state health department, local hospital network, or integrated care board).



\par\medskip\textbf{Core Responsibilities}\par


\begin{itemize}


\item Strategic governance and regional healthcare delivery policy.

\item Establishing and monitoring cross-facility service level agreements (SLAs) and clinical performance indicators.

\item Ensuring statutory compliance with jurisdictional health privacy, clinical governance, and data protection legislation.

\item Funding, procuring, and overseeing regional digital health infrastructure and integration platforms.

\end{itemize}


\par\medskip\textbf{Principal Architectural Concerns}\par


\begin{itemize}


\item \textbf{Network-Wide Availability \& Resilience}: The platform must operate continuously without single points of failure across regional public and private hospital networks.

\item \textbf{Vendor-Independent Longevity}: Avoiding commercial lock-in to proprietary Electronic Medical Record (EMR) vendors, ensuring regional health data remains an enduring asset under public governance.

\item \textbf{Auditability \& Clinical Governance}: Comprehensive, tamper-evident audit trails that prove compliance with regulatory directives and enable forensic review of clinical incidents.

\item \textbf{Systemic Risk Mitigation}: Minimizing the probability of severe systemic failures, data corruption, or unrecoverable message loss during regional health emergencies.

\end{itemize}

{\footnotesize \href{../../../../continuum/motivation/stakeholder/e76216c494fa4500bb80eda6f83513c1.md}{Canonical record: e76216c494fa4500bb80eda6f83513c1}; full qualifications and provenance.}\par
\chapter*{R.2 Historical Capability Identities — Reference navigation}
\phantomsection
\label{navigation:R.2}
\addcontentsline{toc}{chapter}{R.2 Historical Capability Identities — Reference navigation}

Reserved historical identities with governed successor correspondence. These are architectural history and provenance, not additional current Capability subjects. Definitions, standing, relationships and authored Views remain preserved.\par
\section*{R.2.1 [CT1] Care Access \& Coordination — Historical identity [superseded]}
\phantomsection
\label{navigation:R.2.1}
\addcontentsline{toc}{section}{R.2.1 [CT1] Care Access \& Coordination — Historical identity [superseded]}

{\small Governed successor (direct correspondence): II.1.2.1 \hyperref[element:ed342cb8656b4f80b55ad519ecbd2d3a]{[CT1] Care Access \& Coordination}; Canonical ID: \texttt{\seqsplit{ed342cb8656b4f80b55ad519ecbd2d3a}}. Canonical · current.}\par
\label{element:0dec07b6ddbe48bda93a2295dd3e2d7d}
\index{[CT1] Care Access \& Coordination (R.2.1)}
{\small Capability · Superseded historical identity — direct correspondence · superseded}\par
\href{../../../../continuum/strategy/capability/0dec07b6ddbe48bda93a2295dd3e2d7d.md}{Historical canonical definition and provenance}.\par
\section*{R.2.2 [CT3] Clinical Quality Criteria Management — Historical identity [superseded]}
\phantomsection
\label{navigation:R.2.2}
\addcontentsline{toc}{section}{R.2.2 [CT3] Clinical Quality Criteria Management — Historical identity [superseded]}

{\small Governed successor (direct correspondence): II.1.2.2.1.1 \hyperref[element:1f48527a9a7c4382a7f8a660ccb9f29f]{[CT3] Clinical Quality Criteria Management}; Canonical ID: \texttt{\seqsplit{1f48527a9a7c4382a7f8a660ccb9f29f}}. Canonical · current.}\par
\label{element:7a51321b00464016920976b269981edd}
\index{[CT3] Clinical Quality Criteria Management (R.2.2)}
{\small Capability · Superseded historical identity — direct correspondence · superseded}\par
\href{../../../../continuum/strategy/capability/7a51321b00464016920976b269981edd.md}{Historical canonical definition and provenance}.\par
\section*{R.2.3 [CT3] Clinical Quality Evidence Management — Historical identity [superseded]}
\phantomsection
\label{navigation:R.2.3}
\addcontentsline{toc}{section}{R.2.3 [CT3] Clinical Quality Evidence Management — Historical identity [superseded]}

{\small Governed successor (direct correspondence): II.1.2.2.1.2 \hyperref[element:1b0a1c6945ed4fadba045ff4753c104c]{[CT3] Clinical Quality Evidence Management}; Canonical ID: \texttt{\seqsplit{1b0a1c6945ed4fadba045ff4753c104c}}. Canonical · current.}\par
\label{element:f31902d474894670aeef3e0abbedccf0}
\index{[CT3] Clinical Quality Evidence Management (R.2.3)}
{\small Capability · Superseded historical identity — direct correspondence · superseded}\par
\href{../../../../continuum/strategy/capability/f31902d474894670aeef3e0abbedccf0.md}{Historical canonical definition and provenance}.\par
\section*{R.2.4 [CT3] Clinical Quality Reporting — Historical identity [superseded]}
\phantomsection
\label{navigation:R.2.4}
\addcontentsline{toc}{section}{R.2.4 [CT3] Clinical Quality Reporting — Historical identity [superseded]}

{\small Governed successor (direct correspondence): II.1.2.2.1.3 \hyperref[element:70f34839c23042d59469e6ef96944057]{[CT3] Clinical Quality Reporting}; Canonical ID: \texttt{\seqsplit{70f34839c23042d59469e6ef96944057}}. Canonical · current.}\par
\label{element:e5812fcdf3f5465b9b1d8ba536075c5f}
\index{[CT3] Clinical Quality Reporting (R.2.4)}
{\small Capability · Superseded historical identity — direct correspondence · superseded}\par
\href{../../../../continuum/strategy/capability/e5812fcdf3f5465b9b1d8ba536075c5f.md}{Historical canonical definition and provenance}.\par
\section*{R.2.5 [CT3] Clinical Quality Review \& Assessment Management — Historical identity [superseded]}
\phantomsection
\label{navigation:R.2.5}
\addcontentsline{toc}{section}{R.2.5 [CT3] Clinical Quality Review \& Assessment Management — Historical identity [superseded]}

{\small Governed successor (direct correspondence): II.1.2.2.1.4 \hyperref[element:8318c35b302b4dcaa0e795fbfdc93e3b]{[CT3] Clinical Quality Review \& Assessment Management}; Canonical ID: \texttt{\seqsplit{8318c35b302b4dcaa0e795fbfdc93e3b}}. Canonical · current.}\par
\label{element:4b17ec92f0af4c4b9556912dbdd0f29f}
\index{[CT3] Clinical Quality Review \& Assessment Management (R.2.5)}
{\small Capability · Superseded historical identity — direct correspondence · superseded}\par
\href{../../../../continuum/strategy/capability/4b17ec92f0af4c4b9556912dbdd0f29f.md}{Historical canonical definition and provenance}.\par
\section*{R.2.6 [CT2] Clinical Quality Management — Historical identity [superseded]}
\phantomsection
\label{navigation:R.2.6}
\addcontentsline{toc}{section}{R.2.6 [CT2] Clinical Quality Management — Historical identity [superseded]}

{\small Governed successor (direct correspondence): II.1.2.2.1 \hyperref[element:89761f2d38f846c39df1df0975964288]{[CT2] Clinical Quality Management}; Canonical ID: \texttt{\seqsplit{89761f2d38f846c39df1df0975964288}}. Canonical · current.}\par
\label{element:dc1bece79eb7452e951d03a2ee6dfaa6}
\index{[CT2] Clinical Quality Management (R.2.6)}
{\small Capability · Superseded historical identity — direct correspondence · superseded}\par
\href{../../../../continuum/strategy/capability/dc1bece79eb7452e951d03a2ee6dfaa6.md}{Historical canonical definition and provenance}.\par
\section*{R.2.7 [CT3] Clinical Risk Identification — Historical identity [superseded]}
\phantomsection
\label{navigation:R.2.7}
\addcontentsline{toc}{section}{R.2.7 [CT3] Clinical Risk Identification — Historical identity [superseded]}

{\small Governed successor (direct correspondence): II.1.2.2.2.1 \hyperref[element:d54c19c6efaf42af82ad71a4142149d4]{[CT3] Clinical Risk Identification}; Canonical ID: \texttt{\seqsplit{d54c19c6efaf42af82ad71a4142149d4}}. Canonical · current.}\par
\label{element:690da55f1f5b4c44ba58cf67a717a564}
\index{[CT3] Clinical Risk Identification (R.2.7)}
{\small Capability · Superseded historical identity — direct correspondence · superseded}\par
\href{../../../../continuum/strategy/capability/690da55f1f5b4c44ba58cf67a717a564.md}{Historical canonical definition and provenance}.\par
\section*{R.2.8 [CT3] Clinical Risk Mitigation Management — Historical identity [superseded]}
\phantomsection
\label{navigation:R.2.8}
\addcontentsline{toc}{section}{R.2.8 [CT3] Clinical Risk Mitigation Management — Historical identity [superseded]}

{\small Governed successor (direct correspondence): II.1.2.2.2.2 \hyperref[element:2c232ade21ec45619c3ee4e928bd97ee]{[CT3] Clinical Risk Mitigation Management}; Canonical ID: \texttt{\seqsplit{2c232ade21ec45619c3ee4e928bd97ee}}. Canonical · current.}\par
\label{element:0bd3311d89af45bb935fde20e719369c}
\index{[CT3] Clinical Risk Mitigation Management (R.2.8)}
{\small Capability · Superseded historical identity — direct correspondence · superseded}\par
\href{../../../../continuum/strategy/capability/0bd3311d89af45bb935fde20e719369c.md}{Historical canonical definition and provenance}.\par
\section*{R.2.9 [CT3] Clinical Risk Response Management — Historical identity [superseded]}
\phantomsection
\label{navigation:R.2.9}
\addcontentsline{toc}{section}{R.2.9 [CT3] Clinical Risk Response Management — Historical identity [superseded]}

{\small Governed successor (direct correspondence): II.1.2.2.2.3 \hyperref[element:5e36cb137e4e4782b2ddbf78cdf0da7b]{[CT3] Clinical Risk Response Management}; Canonical ID: \texttt{\seqsplit{5e36cb137e4e4782b2ddbf78cdf0da7b}}. Canonical · current.}\par
\label{element:20f719ab4b6741b393ef570b4d1b05df}
\index{[CT3] Clinical Risk Response Management (R.2.9)}
{\small Capability · Superseded historical identity — direct correspondence · superseded}\par
\href{../../../../continuum/strategy/capability/20f719ab4b6741b393ef570b4d1b05df.md}{Historical canonical definition and provenance}.\par
\section*{R.2.10 [CT2] Clinical Safety Management — Historical identity [superseded]}
\phantomsection
\label{navigation:R.2.10}
\addcontentsline{toc}{section}{R.2.10 [CT2] Clinical Safety Management — Historical identity [superseded]}

{\small Governed successor (direct correspondence): II.1.2.2.2 \hyperref[element:28d1bf3e514e4c29837ebedaf9e4917b]{[CT2] Clinical Safety Management}; Canonical ID: \texttt{\seqsplit{28d1bf3e514e4c29837ebedaf9e4917b}}. Canonical · current.}\par
\label{element:650ccec030fa46cc8c9d6b508a841778}
\index{[CT2] Clinical Safety Management (R.2.10)}
{\small Capability · Superseded historical identity — direct correspondence · superseded}\par
\href{../../../../continuum/strategy/capability/650ccec030fa46cc8c9d6b508a841778.md}{Historical canonical definition and provenance}.\par
\section*{R.2.11 [CT1] Clinical Quality, Safety \& Improvement — Historical identity [superseded]}
\phantomsection
\label{navigation:R.2.11}
\addcontentsline{toc}{section}{R.2.11 [CT1] Clinical Quality, Safety \& Improvement — Historical identity [superseded]}

{\small Governed successor (direct correspondence): II.1.2.2 \hyperref[element:63837bf848774572a06b935af95581a2]{[CT1] Clinical Quality, Safety \& Improvement}; Canonical ID: \texttt{\seqsplit{63837bf848774572a06b935af95581a2}}. Canonical · current.}\par
\label{element:b967c383dbcd4c24abe6a83ea5e79e00}
\index{[CT1] Clinical Quality, Safety \& Improvement (R.2.11)}
{\small Capability · Superseded historical identity — direct correspondence · superseded}\par
\href{../../../../continuum/strategy/capability/b967c383dbcd4c24abe6a83ea5e79e00.md}{Historical canonical definition and provenance}.\par
\section*{R.2.12 [CT3] Preventative Health Cohort Management — Historical identity [superseded]}
\phantomsection
\label{navigation:R.2.12}
\addcontentsline{toc}{section}{R.2.12 [CT3] Preventative Health Cohort Management — Historical identity [superseded]}

{\small Governed successor (direct correspondence): II.1.2.3.1.1 \hyperref[element:97625bb28e5745d5994e1bffd4037188]{[CT3] Preventative Health Cohort Management}; Canonical ID: \texttt{\seqsplit{97625bb28e5745d5994e1bffd4037188}}. Canonical · current.}\par
\label{element:fef479f0fd5048b19b2aec770f0233fe}
\index{[CT3] Preventative Health Cohort Management (R.2.12)}
{\small Capability · Superseded historical identity — direct correspondence · superseded}\par
\href{../../../../continuum/strategy/capability/fef479f0fd5048b19b2aec770f0233fe.md}{Historical canonical definition and provenance}.\par
\section*{R.2.13 [CT3] Preventative Health Education Services — Historical identity [superseded]}
\phantomsection
\label{navigation:R.2.13}
\addcontentsline{toc}{section}{R.2.13 [CT3] Preventative Health Education Services — Historical identity [superseded]}

{\small Governed successor (direct correspondence): II.1.2.3.1.2 \hyperref[element:7541d9b3b65d47d286dabdc093cb0eed]{[CT3] Preventative Health Education Services}; Canonical ID: \texttt{\seqsplit{7541d9b3b65d47d286dabdc093cb0eed}}. Canonical · current.}\par
\label{element:dbe10959237f405caeadef2a5748303a}
\index{[CT3] Preventative Health Education Services (R.2.13)}
{\small Capability · Superseded historical identity — direct correspondence · superseded}\par
\href{../../../../continuum/strategy/capability/dbe10959237f405caeadef2a5748303a.md}{Historical canonical definition and provenance}.\par
\section*{R.2.14 [CT3] Preventative Health Problem Identification — Historical identity [superseded]}
\phantomsection
\label{navigation:R.2.14}
\addcontentsline{toc}{section}{R.2.14 [CT3] Preventative Health Problem Identification — Historical identity [superseded]}

{\small Governed successor (direct correspondence): II.1.2.3.1.3 \hyperref[element:c1567022671f433fa1eeb37a9d156d92]{[CT3] Preventative Health Problem Identification}; Canonical ID: \texttt{\seqsplit{c1567022671f433fa1eeb37a9d156d92}}. Canonical · current.}\par
\label{element:b35d3521141c419fb676a79b2ede40bd}
\index{[CT3] Preventative Health Problem Identification (R.2.14)}
{\small Capability · Superseded historical identity — direct correspondence · superseded}\par
\href{../../../../continuum/strategy/capability/b35d3521141c419fb676a79b2ede40bd.md}{Historical canonical definition and provenance}.\par
\section*{R.2.15 [CT3] Preventative Health Programme Management — Historical identity [superseded]}
\phantomsection
\label{navigation:R.2.15}
\addcontentsline{toc}{section}{R.2.15 [CT3] Preventative Health Programme Management — Historical identity [superseded]}

{\small Governed successor (direct correspondence): II.1.2.3.1.4 \hyperref[element:563d36b179d34cb6913253d06e8d77f6]{[CT3] Preventative Health Programme Management}; Canonical ID: \texttt{\seqsplit{563d36b179d34cb6913253d06e8d77f6}}. Canonical · current.}\par
\label{element:d71def9e76a544ff9761c672bb24f3be}
\index{[CT3] Preventative Health Programme Management (R.2.15)}
{\small Capability · Superseded historical identity — direct correspondence · superseded}\par
\href{../../../../continuum/strategy/capability/d71def9e76a544ff9761c672bb24f3be.md}{Historical canonical definition and provenance}.\par
\section*{R.2.16 [CT3] Preventative Health Programme Monitoring — Historical identity [superseded]}
\phantomsection
\label{navigation:R.2.16}
\addcontentsline{toc}{section}{R.2.16 [CT3] Preventative Health Programme Monitoring — Historical identity [superseded]}

{\small Governed successor (direct correspondence): II.1.2.3.1.5 \hyperref[element:945b40ced8004e1cbf7c98c731d17daf]{[CT3] Preventative Health Programme Monitoring}; Canonical ID: \texttt{\seqsplit{945b40ced8004e1cbf7c98c731d17daf}}. Canonical · current.}\par
\label{element:e05f6310f6d544efaca5568604f27dd8}
\index{[CT3] Preventative Health Programme Monitoring (R.2.16)}
{\small Capability · Superseded historical identity — direct correspondence · superseded}\par
\href{../../../../continuum/strategy/capability/e05f6310f6d544efaca5568604f27dd8.md}{Historical canonical definition and provenance}.\par
\section*{R.2.17 [CT3] Preventative Health Reporting — Historical identity [superseded]}
\phantomsection
\label{navigation:R.2.17}
\addcontentsline{toc}{section}{R.2.17 [CT3] Preventative Health Reporting — Historical identity [superseded]}

{\small Governed successor (direct correspondence): II.1.2.3.1.6 \hyperref[element:d2ef01c4cab348df822e802ab084faae]{[CT3] Preventative Health Reporting}; Canonical ID: \texttt{\seqsplit{d2ef01c4cab348df822e802ab084faae}}. Canonical · current.}\par
\label{element:4faef45b4ebe4fa69e57c3d81bfea05f}
\index{[CT3] Preventative Health Reporting (R.2.17)}
{\small Capability · Superseded historical identity — direct correspondence · superseded}\par
\href{../../../../continuum/strategy/capability/4faef45b4ebe4fa69e57c3d81bfea05f.md}{Historical canonical definition and provenance}.\par
\section*{R.2.18 [CT2] Preventative Health — Historical identity [superseded]}
\phantomsection
\label{navigation:R.2.18}
\addcontentsline{toc}{section}{R.2.18 [CT2] Preventative Health — Historical identity [superseded]}

{\small Governed successor (direct correspondence): II.1.2.3.1 \hyperref[element:81be04a2d58548a4aea94f9aae90feb6]{[CT2] Preventative Health}; Canonical ID: \texttt{\seqsplit{81be04a2d58548a4aea94f9aae90feb6}}. Canonical · current.}\par
\label{element:80affa7bcd284ecbbcbc5092f8497ddf}
\index{[CT2] Preventative Health (R.2.18)}
{\small Capability · Superseded historical identity — direct correspondence · superseded}\par
\href{../../../../continuum/strategy/capability/80affa7bcd284ecbbcbc5092f8497ddf.md}{Historical canonical definition and provenance}.\par
\section*{R.2.19 [CT1] Community Health \& Wellbeing — Historical identity [superseded]}
\phantomsection
\label{navigation:R.2.19}
\addcontentsline{toc}{section}{R.2.19 [CT1] Community Health \& Wellbeing — Historical identity [superseded]}

{\small Governed successor (direct correspondence): II.1.2.3 \hyperref[element:b326b4aaf02a41a986ed22d184cbc152]{[CT1] Community Health \& Wellbeing}; Canonical ID: \texttt{\seqsplit{b326b4aaf02a41a986ed22d184cbc152}}. Canonical · current.}\par
\label{element:175174c5cce549d2b154a4f6704721dc}
\index{[CT1] Community Health \& Wellbeing (R.2.19)}
{\small Capability · Superseded historical identity — direct correspondence · superseded}\par
\href{../../../../continuum/strategy/capability/175174c5cce549d2b154a4f6704721dc.md}{Historical canonical definition and provenance}.\par
\section*{R.2.20 [CT1] Connected Health Services — Historical identity [superseded]}
\phantomsection
\label{navigation:R.2.20}
\addcontentsline{toc}{section}{R.2.20 [CT1] Connected Health Services — Historical identity [superseded]}

{\small Governed successor (direct correspondence): II.1.2.4 \hyperref[element:3ef397463e584737b367e93f0bf1d85d]{[CT1] Connected Health Services}; Canonical ID: \texttt{\seqsplit{3ef397463e584737b367e93f0bf1d85d}}. Canonical · current.}\par
\label{element:48f01ef8b0e6431c85cf4f40cf8f94c3}
\index{[CT1] Connected Health Services (R.2.20)}
{\small Capability · Superseded historical identity — direct correspondence · superseded}\par
\href{../../../../continuum/strategy/capability/48f01ef8b0e6431c85cf4f40cf8f94c3.md}{Historical canonical definition and provenance}.\par
\section*{R.2.21 [CT1] Corporate Resources \& Enterprise Services — Historical identity [superseded]}
\phantomsection
\label{navigation:R.2.21}
\addcontentsline{toc}{section}{R.2.21 [CT1] Corporate Resources \& Enterprise Services — Historical identity [superseded]}

{\small Governed successor (direct correspondence): II.1.2.5 \hyperref[element:b79628f156e14a8cbb82a0dea2fc57ff]{[CT1] Corporate Resources \& Enterprise Services}; Canonical ID: \texttt{\seqsplit{b79628f156e14a8cbb82a0dea2fc57ff}}. Canonical · current.}\par
\label{element:a2dbd74847a44f4c879983229f5c9d08}
\index{[CT1] Corporate Resources \& Enterprise Services (R.2.21)}
{\small Capability · Superseded historical identity — direct correspondence · superseded}\par
\href{../../../../continuum/strategy/capability/a2dbd74847a44f4c879983229f5c9d08.md}{Historical canonical definition and provenance}.\par
\section*{R.2.22 [CT1] Diagnostic, Therapeutic \& Clinical Support Services — Historical identity [superseded]}
\phantomsection
\label{navigation:R.2.22}
\addcontentsline{toc}{section}{R.2.22 [CT1] Diagnostic, Therapeutic \& Clinical Support Services — Historical identity [superseded]}

{\small Governed successor (direct correspondence): II.1.2.6 \hyperref[element:4d98c3503a9f4890a677f09b0c3b6f82]{[CT1] Diagnostic, Therapeutic \& Clinical Support Services}; Canonical ID: \texttt{\seqsplit{4d98c3503a9f4890a677f09b0c3b6f82}}. Canonical · current.}\par
\label{element:e003b9cd0730459a84775d2837d1d173}
\index{[CT1] Diagnostic, Therapeutic \& Clinical Support Services (R.2.22)}
{\small Capability · Superseded historical identity — direct correspondence · superseded}\par
\href{../../../../continuum/strategy/capability/e003b9cd0730459a84775d2837d1d173.md}{Historical canonical definition and provenance}.\par
\section*{R.2.23 [CT1] Health Information \& Knowledge Management — Historical identity [superseded]}
\phantomsection
\label{navigation:R.2.23}
\addcontentsline{toc}{section}{R.2.23 [CT1] Health Information \& Knowledge Management — Historical identity [superseded]}

{\small Governed successor (direct correspondence): II.1.2.7 \hyperref[element:cca5d5abec204b339eb5b9bd924ed64e]{[CT1] Health Information \& Knowledge Management}; Canonical ID: \texttt{\seqsplit{cca5d5abec204b339eb5b9bd924ed64e}}. Canonical · current.}\par
\label{element:a1c557bd263047d0b3ff566cf50b98c1}
\index{[CT1] Health Information \& Knowledge Management (R.2.23)}
{\small Capability · Superseded historical identity — direct correspondence · superseded}\par
\href{../../../../continuum/strategy/capability/a1c557bd263047d0b3ff566cf50b98c1.md}{Historical canonical definition and provenance}.\par
\section*{R.2.24 [CT1] Health Products \& Clinical Technology — Historical identity [superseded]}
\phantomsection
\label{navigation:R.2.24}
\addcontentsline{toc}{section}{R.2.24 [CT1] Health Products \& Clinical Technology — Historical identity [superseded]}

{\small Governed successor (direct correspondence): II.1.2.8 \hyperref[element:54388e3202c2403089b5c476c78c4567]{[CT1] Health Products \& Clinical Technology}; Canonical ID: \texttt{\seqsplit{54388e3202c2403089b5c476c78c4567}}. Canonical · current.}\par
\label{element:a631a8600b8b49df9ef452c48c91349a}
\index{[CT1] Health Products \& Clinical Technology (R.2.24)}
{\small Capability · Superseded historical identity — direct correspondence · superseded}\par
\href{../../../../continuum/strategy/capability/a631a8600b8b49df9ef452c48c91349a.md}{Historical canonical definition and provenance}.\par
\section*{R.2.25 [CT1] Health Research \& Innovation — Historical identity [superseded]}
\phantomsection
\label{navigation:R.2.25}
\addcontentsline{toc}{section}{R.2.25 [CT1] Health Research \& Innovation — Historical identity [superseded]}

{\small Governed successor (direct correspondence): II.1.2.9 \hyperref[element:582bf62145fd42118e5d57bb714841e8]{[CT1] Health Research \& Innovation}; Canonical ID: \texttt{\seqsplit{582bf62145fd42118e5d57bb714841e8}}. Canonical · current.}\par
\label{element:2cc31384c8864c73989025cbe96dcfa7}
\index{[CT1] Health Research \& Innovation (R.2.25)}
{\small Capability · Superseded historical identity — direct correspondence · superseded}\par
\href{../../../../continuum/strategy/capability/2cc31384c8864c73989025cbe96dcfa7.md}{Historical canonical definition and provenance}.\par
\section*{R.2.26 [CT1] Health Rights, Advocacy \& Participation — Historical identity [superseded]}
\phantomsection
\label{navigation:R.2.26}
\addcontentsline{toc}{section}{R.2.26 [CT1] Health Rights, Advocacy \& Participation — Historical identity [superseded]}

{\small Governed successor (direct correspondence): II.1.2.10 \hyperref[element:7ab5c82a500b45f39267883a6e78c482]{[CT1] Health Rights, Advocacy \& Participation}; Canonical ID: \texttt{\seqsplit{7ab5c82a500b45f39267883a6e78c482}}. Canonical · current.}\par
\label{element:58376e5c53bf4af0a54f99f21ac52666}
\index{[CT1] Health Rights, Advocacy \& Participation (R.2.26)}
{\small Capability · Superseded historical identity — direct correspondence · superseded}\par
\href{../../../../continuum/strategy/capability/58376e5c53bf4af0a54f99f21ac52666.md}{Historical canonical definition and provenance}.\par
\section*{R.2.27 [CT1] Health Service Assurance — Historical identity [superseded]}
\phantomsection
\label{navigation:R.2.27}
\addcontentsline{toc}{section}{R.2.27 [CT1] Health Service Assurance — Historical identity [superseded]}

{\small Governed successor (direct correspondence): II.1.2.11 \hyperref[element:14860102964c441681d41ec3da423fe0]{[CT1] Health Service Assurance}; Canonical ID: \texttt{\seqsplit{14860102964c441681d41ec3da423fe0}}. Canonical · current.}\par
\label{element:6a82dc1423374e809dbf542eeaa69482}
\index{[CT1] Health Service Assurance (R.2.27)}
{\small Capability · Superseded historical identity — direct correspondence · superseded}\par
\href{../../../../continuum/strategy/capability/6a82dc1423374e809dbf542eeaa69482.md}{Historical canonical definition and provenance}.\par
\section*{R.2.28 [CT1] Health Service Direction \& Stewardship — Historical identity [superseded]}
\phantomsection
\label{navigation:R.2.28}
\addcontentsline{toc}{section}{R.2.28 [CT1] Health Service Direction \& Stewardship — Historical identity [superseded]}

{\small Governed successor (direct correspondence): II.1.2.12 \hyperref[element:d51e7125bf9c4d2c9c59b60346b4d86f]{[CT1] Health Service Direction \& Stewardship}; Canonical ID: \texttt{\seqsplit{d51e7125bf9c4d2c9c59b60346b4d86f}}. Canonical · current.}\par
\label{element:360bf50f92484448a01cc57f2dbce23c}
\index{[CT1] Health Service Direction \& Stewardship (R.2.28)}
{\small Capability · Superseded historical identity — direct correspondence · superseded}\par
\href{../../../../continuum/strategy/capability/360bf50f92484448a01cc57f2dbce23c.md}{Historical canonical definition and provenance}.\par
\section*{R.2.29 [CT1] Health Service Governance — Historical identity [superseded]}
\phantomsection
\label{navigation:R.2.29}
\addcontentsline{toc}{section}{R.2.29 [CT1] Health Service Governance — Historical identity [superseded]}

{\small Governed successor (direct correspondence): II.1.2.13 \hyperref[element:209aada781e844e28202a2535a264c21]{[CT1] Health Service Governance}; Canonical ID: \texttt{\seqsplit{209aada781e844e28202a2535a264c21}}. Canonical · current.}\par
\label{element:24f3dd30223142efa35a5826552d53ab}
\index{[CT1] Health Service Governance (R.2.29)}
{\small Capability · Superseded historical identity — direct correspondence · superseded}\par
\href{../../../../continuum/strategy/capability/24f3dd30223142efa35a5826552d53ab.md}{Historical canonical definition and provenance}.\par
\section*{R.2.30 [CT3] Clinical Finding Interpretation — Historical identity [superseded]}
\phantomsection
\label{navigation:R.2.30}
\addcontentsline{toc}{section}{R.2.30 [CT3] Clinical Finding Interpretation — Historical identity [superseded]}

{\small Governed successor (direct correspondence): II.1.2.14.1.1 \hyperref[element:deb78ba3b1514074a8ca5bc5da0673f5]{[CT3] Clinical Finding Interpretation}; Canonical ID: \texttt{\seqsplit{deb78ba3b1514074a8ca5bc5da0673f5}}. Proposed · current.}\par
\label{element:cf8ddfca7a8e4ed893ec39b795200d77}
\index{[CT3] Clinical Finding Interpretation (R.2.30)}
{\small Capability · Superseded historical identity — direct correspondence · superseded}\par
\href{../../../../continuum/strategy/capability/cf8ddfca7a8e4ed893ec39b795200d77.md}{Historical canonical definition and provenance}.\par
\section*{R.2.31 [CT3] Diagnosis \& Problem Determination — Historical identity [superseded]}
\phantomsection
\label{navigation:R.2.31}
\addcontentsline{toc}{section}{R.2.31 [CT3] Diagnosis \& Problem Determination — Historical identity [superseded]}

{\small Governed successor (direct correspondence): II.1.2.14.1.2 \hyperref[element:d965032637e64275b7a61c0413834590]{[CT3] Diagnosis \& Problem Determination}; Canonical ID: \texttt{\seqsplit{d965032637e64275b7a61c0413834590}}. Proposed · current.}\par
\label{element:e575c583783546ef85cf898523d18797}
\index{[CT3] Diagnosis \& Problem Determination (R.2.31)}
{\small Capability · Superseded historical identity — direct correspondence · superseded}\par
\href{../../../../continuum/strategy/capability/e575c583783546ef85cf898523d18797.md}{Historical canonical definition and provenance}.\par
\section*{R.2.32 [CT3] Health Risk \& Prognostic Determination — Historical identity [superseded]}
\phantomsection
\label{navigation:R.2.32}
\addcontentsline{toc}{section}{R.2.32 [CT3] Health Risk \& Prognostic Determination — Historical identity [superseded]}

{\small Governed successor (direct correspondence): II.1.2.14.1.3 \hyperref[element:5bba0ea969814b5fa7ad11b67ab40a0c]{[CT3] Health Risk \& Prognostic Determination}; Canonical ID: \texttt{\seqsplit{5bba0ea969814b5fa7ad11b67ab40a0c}}. Proposed · current.}\par
\label{element:3215f0b0e1ce45e8ab16dae272ea6fe7}
\index{[CT3] Health Risk \& Prognostic Determination (R.2.32)}
{\small Capability · Superseded historical identity — direct correspondence · superseded}\par
\href{../../../../continuum/strategy/capability/3215f0b0e1ce45e8ab16dae272ea6fe7.md}{Historical canonical definition and provenance}.\par
\section*{R.2.33 [CT2] Clinical Determination — Historical identity [superseded]}
\phantomsection
\label{navigation:R.2.33}
\addcontentsline{toc}{section}{R.2.33 [CT2] Clinical Determination — Historical identity [superseded]}

{\small Governed successor (direct correspondence): II.1.2.14.1 \hyperref[element:39219464b6304b59bd7c3de4e237066e]{[CT2] Clinical Determination}; Canonical ID: \texttt{\seqsplit{39219464b6304b59bd7c3de4e237066e}}. Canonical · current.}\par
\label{element:018ccd818ec142fba1ffb086d900acc6}
\index{[CT2] Clinical Determination (R.2.33)}
{\small Capability · Superseded historical identity — direct correspondence · superseded}\par
\href{../../../../continuum/strategy/capability/018ccd818ec142fba1ffb086d900acc6.md}{Historical canonical definition and provenance}.\par
\section*{R.2.34 [CT2] Clinical Intervention — Historical identity [superseded]}
\phantomsection
\label{navigation:R.2.34}
\addcontentsline{toc}{section}{R.2.34 [CT2] Clinical Intervention — Historical identity [superseded]}

{\small Governed successor (direct correspondence): II.1.2.14.2 \hyperref[element:5359d1b0ff814a0e871384469f5ebcd7]{[CT2] Clinical Intervention}; Canonical ID: \texttt{\seqsplit{5359d1b0ff814a0e871384469f5ebcd7}}. Canonical · current.}\par
\label{element:958345ee65f446d3b1259b8895ab100f}
\index{[CT2] Clinical Intervention (R.2.34)}
{\small Capability · Superseded historical identity — direct correspondence · superseded}\par
\href{../../../../continuum/strategy/capability/958345ee65f446d3b1259b8895ab100f.md}{Historical canonical definition and provenance}.\par
\section*{R.2.35 [CT3] Care Goal Definition — Historical identity [superseded]}
\phantomsection
\label{navigation:R.2.35}
\addcontentsline{toc}{section}{R.2.35 [CT3] Care Goal Definition — Historical identity [superseded]}

{\small Governed successor (direct correspondence): II.1.2.14.3.1 \hyperref[element:609e1bef516b44bb976b5dbebe7f684b]{[CT3] Care Goal Definition}; Canonical ID: \texttt{\seqsplit{609e1bef516b44bb976b5dbebe7f684b}}. Proposed · current.}\par
\label{element:8bfbcc0a455d423c90ae2f1cfb732ca8}
\index{[CT3] Care Goal Definition (R.2.35)}
{\small Capability · Superseded historical identity — direct correspondence · superseded}\par
\href{../../../../continuum/strategy/capability/8bfbcc0a455d423c90ae2f1cfb732ca8.md}{Historical canonical definition and provenance}.\par
\section*{R.2.36 [CT3] Care Response Planning — Historical identity [superseded]}
\phantomsection
\label{navigation:R.2.36}
\addcontentsline{toc}{section}{R.2.36 [CT3] Care Response Planning — Historical identity [superseded]}

{\small Governed successor (direct correspondence): II.1.2.14.3.2 \hyperref[element:2636e43903e44a97ad5b8db32692ce8b]{[CT3] Care Response Planning}; Canonical ID: \texttt{\seqsplit{2636e43903e44a97ad5b8db32692ce8b}}. Proposed · current.}\par
\label{element:d64569acaf2c44f3b2a5332f74fff7c7}
\index{[CT3] Care Response Planning (R.2.36)}
{\small Capability · Superseded historical identity — direct correspondence · superseded}\par
\href{../../../../continuum/strategy/capability/d64569acaf2c44f3b2a5332f74fff7c7.md}{Historical canonical definition and provenance}.\par
\section*{R.2.37 [CT2] Individual Care Planning — Historical identity [superseded]}
\phantomsection
\label{navigation:R.2.37}
\addcontentsline{toc}{section}{R.2.37 [CT2] Individual Care Planning — Historical identity [superseded]}

{\small Governed successor (direct correspondence): II.1.2.14.3 \hyperref[element:f036a9866459429a8fb5f5e10727b5e8]{[CT2] Individual Care Planning}; Canonical ID: \texttt{\seqsplit{f036a9866459429a8fb5f5e10727b5e8}}. Canonical · current.}\par
\label{element:574e61e230e643f78c58649adb4db313}
\index{[CT2] Individual Care Planning (R.2.37)}
{\small Capability · Superseded historical identity — direct correspondence · superseded}\par
\href{../../../../continuum/strategy/capability/574e61e230e643f78c58649adb4db313.md}{Historical canonical definition and provenance}.\par
\section*{R.2.38 [CT3] Care Effectiveness Review — Historical identity [superseded]}
\phantomsection
\label{navigation:R.2.38}
\addcontentsline{toc}{section}{R.2.38 [CT3] Care Effectiveness Review — Historical identity [superseded]}

{\small Governed successor (direct correspondence): II.1.2.14.4.2 \hyperref[element:79cfeaefeb174a54ba43343d37f035c1]{[CT3] Care Effectiveness Review}; Canonical ID: \texttt{\seqsplit{79cfeaefeb174a54ba43343d37f035c1}}. Canonical · current.}\par
\label{element:31f49cdea9f44f188459503580c2b86e}
\index{[CT3] Care Effectiveness Review (R.2.38)}
{\small Capability · Superseded historical identity — direct correspondence · superseded}\par
\href{../../../../continuum/strategy/capability/31f49cdea9f44f188459503580c2b86e.md}{Historical canonical definition and provenance}.\par
\section*{R.2.39 [CT3] Care Intent Appropriateness Review — Historical identity [superseded]}
\phantomsection
\label{navigation:R.2.39}
\addcontentsline{toc}{section}{R.2.39 [CT3] Care Intent Appropriateness Review — Historical identity [superseded]}

{\small Governed successor (evolved correspondence): II.1.2.14.4.1 \hyperref[element:5eac5fb73e31416288d8c992805b33cb]{[CT3] Care Appropriateness Review}; Canonical ID: \texttt{\seqsplit{5eac5fb73e31416288d8c992805b33cb}}. Canonical · current.}\par
\label{element:905eb2de1ac341ac96331e3d7f76079b}
\index{[CT3] Care Intent Appropriateness Review (R.2.39)}
{\small Capability · Superseded historical identity — evolved correspondence · superseded}\par
\href{../../../../continuum/strategy/capability/905eb2de1ac341ac96331e3d7f76079b.md}{Historical canonical definition and provenance}.\par
\section*{R.2.40 [CT2] Individual Care Review — Historical identity [superseded]}
\phantomsection
\label{navigation:R.2.40}
\addcontentsline{toc}{section}{R.2.40 [CT2] Individual Care Review — Historical identity [superseded]}

{\small Governed successor (direct correspondence): II.1.2.14.4 \hyperref[element:625a93f132f242f5a8e2ae1dc19979fb]{[CT2] Individual Care Review}; Canonical ID: \texttt{\seqsplit{625a93f132f242f5a8e2ae1dc19979fb}}. Canonical · current.}\par
\label{element:e1048f6c635d4e16b5384cc2cb45a71b}
\index{[CT2] Individual Care Review (R.2.40)}
{\small Capability · Superseded historical identity — direct correspondence · superseded}\par
\href{../../../../continuum/strategy/capability/e1048f6c635d4e16b5384cc2cb45a71b.md}{Historical canonical definition and provenance}.\par
\section*{R.2.41 [CT3] Clinical Examination — Historical identity [superseded]}
\phantomsection
\label{navigation:R.2.41}
\addcontentsline{toc}{section}{R.2.41 [CT3] Clinical Examination — Historical identity [superseded]}

{\small Governed successor (direct correspondence): II.1.2.14.5.1 \hyperref[element:8855d4d9e2954e0ab66ab65a556fc61a]{[CT3] Clinical Examination}; Canonical ID: \texttt{\seqsplit{8855d4d9e2954e0ab66ab65a556fc61a}}. Canonical · current.}\par
\label{element:8bc15e2b8f114deebcbbe3f2013a8147}
\index{[CT3] Clinical Examination (R.2.41)}
{\small Capability · Superseded historical identity — direct correspondence · superseded}\par
\href{../../../../continuum/strategy/capability/8bc15e2b8f114deebcbbe3f2013a8147.md}{Historical canonical definition and provenance}.\par
\section*{R.2.42 [CT3] Health History Assessment — Historical identity [superseded]}
\phantomsection
\label{navigation:R.2.42}
\addcontentsline{toc}{section}{R.2.42 [CT3] Health History Assessment — Historical identity [superseded]}

{\small Governed successor (direct correspondence): II.1.2.14.5.2 \hyperref[element:7ea7e3f74e2940a7934b588a1a038834]{[CT3] Health History Assessment}; Canonical ID: \texttt{\seqsplit{7ea7e3f74e2940a7934b588a1a038834}}. Canonical · current.}\par
\label{element:bcaacf446f45443eafd8e21a1fd74ee1}
\index{[CT3] Health History Assessment (R.2.42)}
{\small Capability · Superseded historical identity — direct correspondence · superseded}\par
\href{../../../../continuum/strategy/capability/bcaacf446f45443eafd8e21a1fd74ee1.md}{Historical canonical definition and provenance}.\par
\section*{R.2.43 [CT3] Health Status Measurement — Historical identity [superseded]}
\phantomsection
\label{navigation:R.2.43}
\addcontentsline{toc}{section}{R.2.43 [CT3] Health Status Measurement — Historical identity [superseded]}

{\small Governed successor (direct correspondence): II.1.2.14.5.3 \hyperref[element:a796bfb3335d44f6aeb9970628ed8554]{[CT3] Health Status Measurement}; Canonical ID: \texttt{\seqsplit{a796bfb3335d44f6aeb9970628ed8554}}. Canonical · current.}\par
\label{element:3717906f7f8840d3a8a3e3877543dd18}
\index{[CT3] Health Status Measurement (R.2.43)}
{\small Capability · Superseded historical identity — direct correspondence · superseded}\par
\href{../../../../continuum/strategy/capability/3717906f7f8840d3a8a3e3877543dd18.md}{Historical canonical definition and provenance}.\par
\section*{R.2.44 [CT2] Individual Health Assessment — Historical identity [superseded]}
\phantomsection
\label{navigation:R.2.44}
\addcontentsline{toc}{section}{R.2.44 [CT2] Individual Health Assessment — Historical identity [superseded]}

{\small Governed successor (direct correspondence): II.1.2.14.5 \hyperref[element:045f819bb3694d3c9db7f20886e2993f]{[CT2] Individual Health Assessment}; Canonical ID: \texttt{\seqsplit{045f819bb3694d3c9db7f20886e2993f}}. Canonical · current.}\par
\label{element:9111446aa2d045bd9acd98542b0eef1f}
\index{[CT2] Individual Health Assessment (R.2.44)}
{\small Capability · Superseded historical identity — direct correspondence · superseded}\par
\href{../../../../continuum/strategy/capability/9111446aa2d045bd9acd98542b0eef1f.md}{Historical canonical definition and provenance}.\par
\section*{R.2.45 [CT3] Care Response Monitoring — Historical identity [superseded]}
\phantomsection
\label{navigation:R.2.45}
\addcontentsline{toc}{section}{R.2.45 [CT3] Care Response Monitoring — Historical identity [superseded]}

{\small Governed successor (direct correspondence): II.1.2.14.6.1 \hyperref[element:d6951c9357fa40a4943cde581d95d5b1]{[CT3] Care Response Monitoring}; Canonical ID: \texttt{\seqsplit{d6951c9357fa40a4943cde581d95d5b1}}. Canonical · current.}\par
\label{element:615f93d05fbb4a8481f07778bd4f4803}
\index{[CT3] Care Response Monitoring (R.2.45)}
{\small Capability · Superseded historical identity — direct correspondence · superseded}\par
\href{../../../../continuum/strategy/capability/615f93d05fbb4a8481f07778bd4f4803.md}{Historical canonical definition and provenance}.\par
\section*{R.2.46 [CT3] Health State Monitoring — Historical identity [superseded]}
\phantomsection
\label{navigation:R.2.46}
\addcontentsline{toc}{section}{R.2.46 [CT3] Health State Monitoring — Historical identity [superseded]}

{\small Governed successor (direct correspondence): II.1.2.14.6.2 \hyperref[element:d77f22909f8b47dba4cbd95982e576df]{[CT3] Health State Monitoring}; Canonical ID: \texttt{\seqsplit{d77f22909f8b47dba4cbd95982e576df}}. Canonical · current.}\par
\label{element:346f024f74b447a8b6056b9d521837ff}
\index{[CT3] Health State Monitoring (R.2.46)}
{\small Capability · Superseded historical identity — direct correspondence · superseded}\par
\href{../../../../continuum/strategy/capability/346f024f74b447a8b6056b9d521837ff.md}{Historical canonical definition and provenance}.\par
\section*{R.2.47 [CT2] Individual Health Monitoring — Historical identity [superseded]}
\phantomsection
\label{navigation:R.2.47}
\addcontentsline{toc}{section}{R.2.47 [CT2] Individual Health Monitoring — Historical identity [superseded]}

{\small Governed successor (direct correspondence): II.1.2.14.6 \hyperref[element:c7d0e498e1ad44138df940af2f91cd2d]{[CT2] Individual Health Monitoring}; Canonical ID: \texttt{\seqsplit{c7d0e498e1ad44138df940af2f91cd2d}}. Canonical · current.}\par
\label{element:31e7c27556314cafb74d040441d2dac6}
\index{[CT2] Individual Health Monitoring (R.2.47)}
{\small Capability · Superseded historical identity — direct correspondence · superseded}\par
\href{../../../../continuum/strategy/capability/31e7c27556314cafb74d040441d2dac6.md}{Historical canonical definition and provenance}.\par
\section*{R.2.48 [CT1] Individual Care Delivery — Historical identity [superseded]}
\phantomsection
\label{navigation:R.2.48}
\addcontentsline{toc}{section}{R.2.48 [CT1] Individual Care Delivery — Historical identity [superseded]}

{\small Governed successor (direct correspondence): II.1.2.14 \hyperref[element:b964a11cde594318aa14ff935414158a]{[CT1] Individual Care Delivery}; Canonical ID: \texttt{\seqsplit{b964a11cde594318aa14ff935414158a}}. Canonical · current.}\par
\label{element:162deb7e292042f78f677a14476a3a1d}
\index{[CT1] Individual Care Delivery (R.2.48)}
{\small Capability · Superseded historical identity — direct correspondence · superseded}\par
\href{../../../../continuum/strategy/capability/162deb7e292042f78f677a14476a3a1d.md}{Historical canonical definition and provenance}.\par
\section*{R.2.49 [CT1] Population Health \& Health-System Planning — Historical identity [superseded]}
\phantomsection
\label{navigation:R.2.49}
\addcontentsline{toc}{section}{R.2.49 [CT1] Population Health \& Health-System Planning — Historical identity [superseded]}

{\small Governed successor (direct correspondence): II.1.2.15 \hyperref[element:c63b38b37e9b4e9b97dc87ed54604def]{[CT1] Population Health \& Health-System Planning}; Canonical ID: \texttt{\seqsplit{c63b38b37e9b4e9b97dc87ed54604def}}. Canonical · current.}\par
\label{element:a5003436a7f348b1bff60ba593adfa27}
\index{[CT1] Population Health \& Health-System Planning (R.2.49)}
{\small Capability · Superseded historical identity — direct correspondence · superseded}\par
\href{../../../../continuum/strategy/capability/a5003436a7f348b1bff60ba593adfa27.md}{Historical canonical definition and provenance}.\par
\section*{R.2.50 [CT1] Security, Privacy \& Digital Trust — Historical identity [superseded]}
\phantomsection
\label{navigation:R.2.50}
\addcontentsline{toc}{section}{R.2.50 [CT1] Security, Privacy \& Digital Trust — Historical identity [superseded]}

{\small Governed successor (direct correspondence): II.1.2.16 \hyperref[element:71f0431fdfcd4f779e16d9e5ea609833]{[CT1] Security, Privacy \& Digital Trust}; Canonical ID: \texttt{\seqsplit{71f0431fdfcd4f779e16d9e5ea609833}}. Canonical · current.}\par
\label{element:b01441b2f5584fc2b9c3988643a7c93b}
\index{[CT1] Security, Privacy \& Digital Trust (R.2.50)}
{\small Capability · Superseded historical identity — direct correspondence · superseded}\par
\href{../../../../continuum/strategy/capability/b01441b2f5584fc2b9c3988643a7c93b.md}{Historical canonical definition and provenance}.\par
\section*{R.2.51 [CT1] Standards, Semantics \& Reference Governance — Historical identity [superseded]}
\phantomsection
\label{navigation:R.2.51}
\addcontentsline{toc}{section}{R.2.51 [CT1] Standards, Semantics \& Reference Governance — Historical identity [superseded]}

{\small Governed successor (direct correspondence): II.1.2.17 \hyperref[element:5ffeb7401d314f06b05d79c3aff98242]{[CT1] Standards, Semantics \& Reference Governance}; Canonical ID: \texttt{\seqsplit{5ffeb7401d314f06b05d79c3aff98242}}. Canonical · current.}\par
\label{element:3b65a3d007f54aa9b875e655475e1786}
\index{[CT1] Standards, Semantics \& Reference Governance (R.2.51)}
{\small Capability · Superseded historical identity — direct correspondence · superseded}\par
\href{../../../../continuum/strategy/capability/3b65a3d007f54aa9b875e655475e1786.md}{Historical canonical definition and provenance}.\par
\section*{R.2.52 [CT1] Workforce \& Organisational Capability — Historical identity [superseded]}
\phantomsection
\label{navigation:R.2.52}
\addcontentsline{toc}{section}{R.2.52 [CT1] Workforce \& Organisational Capability — Historical identity [superseded]}

{\small Governed successor (direct correspondence): II.1.2.18 \hyperref[element:631b43664bfb4c50ac8cf4958b2c347d]{[CT1] Workforce \& Organisational Capability}; Canonical ID: \texttt{\seqsplit{631b43664bfb4c50ac8cf4958b2c347d}}. Canonical · current.}\par
\label{element:b1161fe1a315420fbc7d048a1ea58126}
\index{[CT1] Workforce \& Organisational Capability (R.2.52)}
{\small Capability · Superseded historical identity — direct correspondence · superseded}\par
\href{../../../../continuum/strategy/capability/b1161fe1a315420fbc7d048a1ea58126.md}{Historical canonical definition and provenance}.\par
\section*{R.2.53 [CT3] Care Plan Lifecycle \& Coordination — Historical identity [superseded]}
\phantomsection
\label{navigation:R.2.53}
\addcontentsline{toc}{section}{R.2.53 [CT3] Care Plan Lifecycle \& Coordination — Historical identity [superseded]}

{\small Governed successor (direct correspondence): II.1.3.1.5.1 \hyperref[element:4cb7db65392f4fab801d7794af4d4b29]{[CT3] Care Plan Lifecycle \& Coordination}; Canonical ID: \texttt{\seqsplit{4cb7db65392f4fab801d7794af4d4b29}}. Proposed · current.}\par
\label{element:689bb3f8a24249b79f787e1f48eb268d}
\index{[CT3] Care Plan Lifecycle \& Coordination (R.2.53)}
{\small Capability · Superseded historical identity — direct correspondence · superseded}\par
\href{../../../../continuum/strategy/capability/689bb3f8a24249b79f787e1f48eb268d.md}{Historical canonical definition and provenance}.\par
\section*{R.2.54 [CT2] Care Plan Administration — Historical identity [superseded]}
\phantomsection
\label{navigation:R.2.54}
\addcontentsline{toc}{section}{R.2.54 [CT2] Care Plan Administration — Historical identity [superseded]}

{\small Governed successor (direct correspondence): II.1.3.1.5 \hyperref[element:cf06eec5699f406fa30a00d2f76e3098]{[CT2] Care Plan Administration}; Canonical ID: \texttt{\seqsplit{cf06eec5699f406fa30a00d2f76e3098}}. Canonical · current.}\par
\label{element:e31407cea7624632a8ed2dc53b09d726}
\index{[CT2] Care Plan Administration (R.2.54)}
{\small Capability · Superseded historical identity — direct correspondence · superseded}\par
\href{../../../../continuum/strategy/capability/e31407cea7624632a8ed2dc53b09d726.md}{Historical canonical definition and provenance}.\par
\section*{R.2.55 [CT2] Client-Centred Administration — Historical identity [superseded]}
\phantomsection
\label{navigation:R.2.55}
\addcontentsline{toc}{section}{R.2.55 [CT2] Client-Centred Administration — Historical identity [superseded]}

{\small Governed successor (direct correspondence): II.1.3.1.6 \hyperref[element:8752b7a937754b2faa296300c6a48911]{[CT2] Client-Centred Administration}; Canonical ID: \texttt{\seqsplit{8752b7a937754b2faa296300c6a48911}}. Canonical · current.}\par
\label{element:232b3fb2ff204c9b9ff7a60ff6a7a32d}
\index{[CT2] Client-Centred Administration (R.2.55)}
{\small Capability · Superseded historical identity — direct correspondence · superseded}\par
\href{../../../../continuum/strategy/capability/232b3fb2ff204c9b9ff7a60ff6a7a32d.md}{Historical canonical definition and provenance}.\par
\section*{R.2.56 [CT2] Clinical Charging \& Claims Administration — Historical identity [superseded]}
\phantomsection
\label{navigation:R.2.56}
\addcontentsline{toc}{section}{R.2.56 [CT2] Clinical Charging \& Claims Administration — Historical identity [superseded]}

{\small Governed successor (direct correspondence): II.1.3.1.7 \hyperref[element:af78ab7c89454a608cd4d34ec7703331]{[CT2] Clinical Charging \& Claims Administration}; Canonical ID: \texttt{\seqsplit{af78ab7c89454a608cd4d34ec7703331}}. Canonical · current.}\par
\label{element:44433108cc8f4f5f9d7581f9e5b29eb5}
\index{[CT2] Clinical Charging \& Claims Administration (R.2.56)}
{\small Capability · Superseded historical identity — direct correspondence · superseded}\par
\href{../../../../continuum/strategy/capability/44433108cc8f4f5f9d7581f9e5b29eb5.md}{Historical canonical definition and provenance}.\par
\section*{R.2.57 [CT2] Clinical Communication Administration — Historical identity [superseded]}
\phantomsection
\label{navigation:R.2.57}
\addcontentsline{toc}{section}{R.2.57 [CT2] Clinical Communication Administration — Historical identity [superseded]}

{\small Governed successor (direct correspondence): II.1.3.1.8 \hyperref[element:b38f2c2811724734b5887b6dbaaabbfc]{[CT2] Clinical Communication Administration}; Canonical ID: \texttt{\seqsplit{b38f2c2811724734b5887b6dbaaabbfc}}. Canonical · current.}\par
\label{element:bd4d6f4a7f03402ba7857383c7be7a4c}
\index{[CT2] Clinical Communication Administration (R.2.57)}
{\small Capability · Superseded historical identity — direct correspondence · superseded}\par
\href{../../../../continuum/strategy/capability/bd4d6f4a7f03402ba7857383c7be7a4c.md}{Historical canonical definition and provenance}.\par
\section*{R.2.58 [CT3] Clinical Documentation Lifecycle Administration — Historical identity [superseded]}
\phantomsection
\label{navigation:R.2.58}
\addcontentsline{toc}{section}{R.2.58 [CT3] Clinical Documentation Lifecycle Administration — Historical identity [superseded]}

{\small Governed successor (direct correspondence): II.1.3.1.9.1 \hyperref[element:df683100067b47e785037c27cfd17a3c]{[CT3] Clinical Documentation Lifecycle Administration}; Canonical ID: \texttt{\seqsplit{df683100067b47e785037c27cfd17a3c}}. Proposed · current.}\par
\label{element:5b1636373a8d4e788b3ddbc1659a47ef}
\index{[CT3] Clinical Documentation Lifecycle Administration (R.2.58)}
{\small Capability · Superseded historical identity — direct correspondence · superseded}\par
\href{../../../../continuum/strategy/capability/5b1636373a8d4e788b3ddbc1659a47ef.md}{Historical canonical definition and provenance}.\par
\section*{R.2.59 [CT2] Clinical Record Administration — Historical identity [superseded]}
\phantomsection
\label{navigation:R.2.59}
\addcontentsline{toc}{section}{R.2.59 [CT2] Clinical Record Administration — Historical identity [superseded]}

{\small Governed successor (direct correspondence): II.1.3.1.9 \hyperref[element:ce9e7e7b22054cefb707a650d5be7219]{[CT2] Clinical Record Administration}; Canonical ID: \texttt{\seqsplit{ce9e7e7b22054cefb707a650d5be7219}}. Canonical · current.}\par
\label{element:be3f2e695eda403fbce8118daf83d386}
\index{[CT2] Clinical Record Administration (R.2.59)}
{\small Capability · Superseded historical identity — direct correspondence · superseded}\par
\href{../../../../continuum/strategy/capability/be3f2e695eda403fbce8118daf83d386.md}{Historical canonical definition and provenance}.\par
\section*{R.2.60 [CT2] Diagnostic Administration — Historical identity [superseded]}
\phantomsection
\label{navigation:R.2.60}
\addcontentsline{toc}{section}{R.2.60 [CT2] Diagnostic Administration — Historical identity [superseded]}

{\small Governed successor (direct correspondence): II.1.3.1.10 \hyperref[element:3cc720ba2e8b4d7a94bb349391d61895]{[CT2] Diagnostic Administration}; Canonical ID: \texttt{\seqsplit{3cc720ba2e8b4d7a94bb349391d61895}}. Canonical · current.}\par
\label{element:2b1a27d5d9944e05b4cf0fcaad4ca433}
\index{[CT2] Diagnostic Administration (R.2.60)}
{\small Capability · Superseded historical identity — direct correspondence · superseded}\par
\href{../../../../continuum/strategy/capability/2b1a27d5d9944e05b4cf0fcaad4ca433.md}{Historical canonical definition and provenance}.\par
\section*{R.2.61 [CT2] Episode \& Encounter Administration — Historical identity [superseded]}
\phantomsection
\label{navigation:R.2.61}
\addcontentsline{toc}{section}{R.2.61 [CT2] Episode \& Encounter Administration — Historical identity [superseded]}

{\small Governed successor (direct correspondence): II.1.3.1.11 \hyperref[element:61d7067aee0540259b3daca4287d66d2]{[CT2] Episode \& Encounter Administration}; Canonical ID: \texttt{\seqsplit{61d7067aee0540259b3daca4287d66d2}}. Canonical · current.}\par
\label{element:2778f988e8d44e3091749542c39d38a4}
\index{[CT2] Episode \& Encounter Administration (R.2.61)}
{\small Capability · Superseded historical identity — direct correspondence · superseded}\par
\href{../../../../continuum/strategy/capability/2778f988e8d44e3091749542c39d38a4.md}{Historical canonical definition and provenance}.\par
\section*{R.2.62 [CT2] Funding \& Coverage Administration — Historical identity [superseded]}
\phantomsection
\label{navigation:R.2.62}
\addcontentsline{toc}{section}{R.2.62 [CT2] Funding \& Coverage Administration — Historical identity [superseded]}

{\small Governed successor (direct correspondence): II.1.3.1.12 \hyperref[element:640432d413cf47928ea2b3d221f92313]{[CT2] Funding \& Coverage Administration}; Canonical ID: \texttt{\seqsplit{640432d413cf47928ea2b3d221f92313}}. Canonical · current.}\par
\label{element:58beafcbcd084331be710eb00d27f6cf}
\index{[CT2] Funding \& Coverage Administration (R.2.62)}
{\small Capability · Superseded historical identity — direct correspondence · superseded}\par
\href{../../../../continuum/strategy/capability/58beafcbcd084331be710eb00d27f6cf.md}{Historical canonical definition and provenance}.\par
\section*{R.2.63 [CT2] Medication Administration — Historical identity [superseded]}
\phantomsection
\label{navigation:R.2.63}
\addcontentsline{toc}{section}{R.2.63 [CT2] Medication Administration — Historical identity [superseded]}

{\small Governed successor (direct correspondence): II.1.3.1.13 \hyperref[element:c907e2a95629435b9951941ad4e2ebe9]{[CT2] Medication Administration}; Canonical ID: \texttt{\seqsplit{c907e2a95629435b9951941ad4e2ebe9}}. Canonical · current.}\par
\label{element:5a38bf2527a8404bac714121f2dc76a9}
\index{[CT2] Medication Administration (R.2.63)}
{\small Capability · Superseded historical identity — direct correspondence · superseded}\par
\href{../../../../continuum/strategy/capability/5a38bf2527a8404bac714121f2dc76a9.md}{Historical canonical definition and provenance}.\par
\section*{R.2.64 [CT3] Closed-Loop Order Progression — Historical identity [superseded]}
\phantomsection
\label{navigation:R.2.64}
\addcontentsline{toc}{section}{R.2.64 [CT3] Closed-Loop Order Progression — Historical identity [superseded]}

{\small Governed successor (direct correspondence): II.1.3.1.14.1 \hyperref[element:2f79ebd7392c46ef92def6f25bc9c7ce]{[CT3] Closed-Loop Order Progression}; Canonical ID: \texttt{\seqsplit{2f79ebd7392c46ef92def6f25bc9c7ce}}. Proposed · current.}\par
\label{element:2ad29cebc1e94053a300d869e7797232}
\index{[CT3] Closed-Loop Order Progression (R.2.64)}
{\small Capability · Superseded historical identity — direct correspondence · superseded}\par
\href{../../../../continuum/strategy/capability/2ad29cebc1e94053a300d869e7797232.md}{Historical canonical definition and provenance}.\par
\section*{R.2.65 [CT2] Order Administration — Historical identity [superseded]}
\phantomsection
\label{navigation:R.2.65}
\addcontentsline{toc}{section}{R.2.65 [CT2] Order Administration — Historical identity [superseded]}

{\small Governed successor (direct correspondence): II.1.3.1.14 \hyperref[element:c3c5abd7987f4433a46cbfed8167914c]{[CT2] Order Administration}; Canonical ID: \texttt{\seqsplit{c3c5abd7987f4433a46cbfed8167914c}}. Canonical · current.}\par
\label{element:ebb735a2915744819d4001d9b8425cc6}
\index{[CT2] Order Administration (R.2.65)}
{\small Capability · Superseded historical identity — direct correspondence · superseded}\par
\href{../../../../continuum/strategy/capability/ebb735a2915744819d4001d9b8425cc6.md}{Historical canonical definition and provenance}.\par
\section*{R.2.66 [CT2] Procedure Administration — Historical identity [superseded]}
\phantomsection
\label{navigation:R.2.66}
\addcontentsline{toc}{section}{R.2.66 [CT2] Procedure Administration — Historical identity [superseded]}

{\small Governed successor (direct correspondence): II.1.3.1.15 \hyperref[element:649dd7cc26094e539386a050086dfb91]{[CT2] Procedure Administration}; Canonical ID: \texttt{\seqsplit{649dd7cc26094e539386a050086dfb91}}. Canonical · current.}\par
\label{element:466a4057001846bb8ae0481453d3e410}
\index{[CT2] Procedure Administration (R.2.66)}
{\small Capability · Superseded historical identity — direct correspondence · superseded}\par
\href{../../../../continuum/strategy/capability/466a4057001846bb8ae0481453d3e410.md}{Historical canonical definition and provenance}.\par
\section*{R.2.67 [CT3] Referral Intake \& Progression — Historical identity [superseded]}
\phantomsection
\label{navigation:R.2.67}
\addcontentsline{toc}{section}{R.2.67 [CT3] Referral Intake \& Progression — Historical identity [superseded]}

{\small Governed successor (direct correspondence): II.1.3.1.16.1 \hyperref[element:d11bdb9464b2476894aa372ab9b929a7]{[CT3] Referral Intake \& Progression}; Canonical ID: \texttt{\seqsplit{d11bdb9464b2476894aa372ab9b929a7}}. Proposed · current.}\par
\label{element:7180f6265d5c4687a001220c6126e301}
\index{[CT3] Referral Intake \& Progression (R.2.67)}
{\small Capability · Superseded historical identity — direct correspondence · superseded}\par
\href{../../../../continuum/strategy/capability/7180f6265d5c4687a001220c6126e301.md}{Historical canonical definition and provenance}.\par
\section*{R.2.68 [CT2] Referral Administration — Historical identity [superseded]}
\phantomsection
\label{navigation:R.2.68}
\addcontentsline{toc}{section}{R.2.68 [CT2] Referral Administration — Historical identity [superseded]}

{\small Governed successor (direct correspondence): II.1.3.1.16 \hyperref[element:ae0de25600ee416aa24d2a167b64f230]{[CT2] Referral Administration}; Canonical ID: \texttt{\seqsplit{ae0de25600ee416aa24d2a167b64f230}}. Canonical · current.}\par
\label{element:5b0ccf4b30054145952c298b348ed266}
\index{[CT2] Referral Administration (R.2.68)}
{\small Capability · Superseded historical identity — direct correspondence · superseded}\par
\href{../../../../continuum/strategy/capability/5b0ccf4b30054145952c298b348ed266.md}{Historical canonical definition and provenance}.\par
\section*{R.2.69 [CT3] Appointment \& Booking Administration — Historical identity [superseded]}
\phantomsection
\label{navigation:R.2.69}
\addcontentsline{toc}{section}{R.2.69 [CT3] Appointment \& Booking Administration — Historical identity [superseded]}

{\small Governed successor (direct correspondence): II.1.3.1.17.1 \hyperref[element:7b3ba9bcac564dd4a7a08631581db2ca]{[CT3] Appointment \& Booking Administration}; Canonical ID: \texttt{\seqsplit{7b3ba9bcac564dd4a7a08631581db2ca}}. Proposed · current.}\par
\label{element:cdc16dd0f84049f8af3f1ec76f1e1428}
\index{[CT3] Appointment \& Booking Administration (R.2.69)}
{\small Capability · Superseded historical identity — direct correspondence · superseded}\par
\href{../../../../continuum/strategy/capability/cdc16dd0f84049f8af3f1ec76f1e1428.md}{Historical canonical definition and provenance}.\par
\section*{R.2.70 [CT2] Scheduling Administration — Historical identity [superseded]}
\phantomsection
\label{navigation:R.2.70}
\addcontentsline{toc}{section}{R.2.70 [CT2] Scheduling Administration — Historical identity [superseded]}

{\small Governed successor (direct correspondence): II.1.3.1.17 \hyperref[element:35fb3744c900470f9895056246a13fb1]{[CT2] Scheduling Administration}; Canonical ID: \texttt{\seqsplit{35fb3744c900470f9895056246a13fb1}}. Canonical · current.}\par
\label{element:1a6630940ffb44fcad4543b93a367b67}
\index{[CT2] Scheduling Administration (R.2.70)}
{\small Capability · Superseded historical identity — direct correspondence · superseded}\par
\href{../../../../continuum/strategy/capability/1a6630940ffb44fcad4543b93a367b67.md}{Historical canonical definition and provenance}.\par
\section*{R.2.71 [CT1] Care Administration Enablement — Historical identity [superseded]}
\phantomsection
\label{navigation:R.2.71}
\addcontentsline{toc}{section}{R.2.71 [CT1] Care Administration Enablement — Historical identity [superseded]}

{\small Governed successor (direct correspondence): II.1.3.1 \hyperref[element:28df27d55c514a37995145ab46945f78]{[CT1] Care Administration Enablement}; Canonical ID: \texttt{\seqsplit{28df27d55c514a37995145ab46945f78}}. Proposed · current.}\par
\label{element:b611a1fcfadb4b30af9ba74a4c993d8a}
\index{[CT1] Care Administration Enablement (R.2.71)}
{\small Capability · Superseded historical identity — direct correspondence · superseded}\par
\href{../../../../continuum/strategy/capability/b611a1fcfadb4b30af9ba74a4c993d8a.md}{Historical canonical definition and provenance}.\par
\section*{R.2.72 [CT2] Acute Care — Historical identity [superseded]}
\phantomsection
\label{navigation:R.2.72}
\addcontentsline{toc}{section}{R.2.72 [CT2] Acute Care — Historical identity [superseded]}

{\small Governed successor (direct correspondence): II.1.3.3.1 \hyperref[element:9bb988c526404522b83aeed5069e0eb5]{[CT2] Acute Care}; Canonical ID: \texttt{\seqsplit{9bb988c526404522b83aeed5069e0eb5}}. Proposed · current.}\par
\label{element:7b591957f7424424abc80f699e9f1960}
\index{[CT2] Acute Care (R.2.72)}
{\small Capability · Superseded historical identity — direct correspondence · superseded}\par
\href{../../../../continuum/strategy/capability/7b591957f7424424abc80f699e9f1960.md}{Historical canonical definition and provenance}.\par
\section*{R.2.73 [CT2] Allied Health — Historical identity [superseded]}
\phantomsection
\label{navigation:R.2.73}
\addcontentsline{toc}{section}{R.2.73 [CT2] Allied Health — Historical identity [superseded]}

{\small Governed successor (direct correspondence): II.1.3.3.2 \hyperref[element:63cd714a3e2a4354aa5383b189fea566]{[CT2] Allied Health}; Canonical ID: \texttt{\seqsplit{63cd714a3e2a4354aa5383b189fea566}}. Canonical · current.}\par
\label{element:6cd45c27d0c344b9807174dbd812d572}
\index{[CT2] Allied Health (R.2.73)}
{\small Capability · Superseded historical identity — direct correspondence · superseded}\par
\href{../../../../continuum/strategy/capability/6cd45c27d0c344b9807174dbd812d572.md}{Historical canonical definition and provenance}.\par
\section*{R.2.74 [CT2] Ambulatory Care — Historical identity [superseded]}
\phantomsection
\label{navigation:R.2.74}
\addcontentsline{toc}{section}{R.2.74 [CT2] Ambulatory Care — Historical identity [superseded]}

{\small Governed successor (direct correspondence): II.1.3.3.3 \hyperref[element:f74d7c3a39c540c4b462576e8f09fc58]{[CT2] Ambulatory Care}; Canonical ID: \texttt{\seqsplit{f74d7c3a39c540c4b462576e8f09fc58}}. Canonical · current.}\par
\label{element:e996b61b6e7b4619913f427aed2b55b1}
\index{[CT2] Ambulatory Care (R.2.74)}
{\small Capability · Superseded historical identity — direct correspondence · superseded}\par
\href{../../../../continuum/strategy/capability/e996b61b6e7b4619913f427aed2b55b1.md}{Historical canonical definition and provenance}.\par
\section*{R.2.75 [CT2] Clinical Assessment — Historical identity [superseded]}
\phantomsection
\label{navigation:R.2.75}
\addcontentsline{toc}{section}{R.2.75 [CT2] Clinical Assessment — Historical identity [superseded]}

{\small Governed successor (direct correspondence): II.1.3.3.4 \hyperref[element:a7d84c95e48c453797419813ca80a66c]{[CT2] Clinical Assessment}; Canonical ID: \texttt{\seqsplit{a7d84c95e48c453797419813ca80a66c}}. Canonical · current.}\par
\label{element:1c3b553c87704e27ba211fbbd02d95cb}
\index{[CT2] Clinical Assessment (R.2.75)}
{\small Capability · Superseded historical identity — direct correspondence · superseded}\par
\href{../../../../continuum/strategy/capability/1c3b553c87704e27ba211fbbd02d95cb.md}{Historical canonical definition and provenance}.\par
\section*{R.2.76 [CT2] Community Care — Historical identity [superseded]}
\phantomsection
\label{navigation:R.2.76}
\addcontentsline{toc}{section}{R.2.76 [CT2] Community Care — Historical identity [superseded]}

{\small Governed successor (direct correspondence): II.1.3.3.5 \hyperref[element:942aa78cad004ef99de6b3aae8ddd014]{[CT2] Community Care}; Canonical ID: \texttt{\seqsplit{942aa78cad004ef99de6b3aae8ddd014}}. Proposed · current.}\par
\label{element:4b7177f3f40449518fba381acae6ba1e}
\index{[CT2] Community Care (R.2.76)}
{\small Capability · Superseded historical identity — direct correspondence · superseded}\par
\href{../../../../continuum/strategy/capability/4b7177f3f40449518fba381acae6ba1e.md}{Historical canonical definition and provenance}.\par
\section*{R.2.77 [CT2] Diagnostic Services — Historical identity [superseded]}
\phantomsection
\label{navigation:R.2.77}
\addcontentsline{toc}{section}{R.2.77 [CT2] Diagnostic Services — Historical identity [superseded]}

{\small Governed successor (direct correspondence): II.1.3.3.6 \hyperref[element:b825ebab42f444d8b738b4c1f489f785]{[CT2] Diagnostic Services}; Canonical ID: \texttt{\seqsplit{b825ebab42f444d8b738b4c1f489f785}}. Proposed · current.}\par
\label{element:03a57e0c9ea34f7fb7b78e5f5c85379f}
\index{[CT2] Diagnostic Services (R.2.77)}
{\small Capability · Superseded historical identity — direct correspondence · superseded}\par
\href{../../../../continuum/strategy/capability/03a57e0c9ea34f7fb7b78e5f5c85379f.md}{Historical canonical definition and provenance}.\par
\section*{R.2.78 [CT2] Emergency Care — Historical identity [superseded]}
\phantomsection
\label{navigation:R.2.78}
\addcontentsline{toc}{section}{R.2.78 [CT2] Emergency Care — Historical identity [superseded]}

{\small Governed successor (direct correspondence): II.1.3.3.7 \hyperref[element:aa088f0647ad4b78966fe28fa80af674]{[CT2] Emergency Care}; Canonical ID: \texttt{\seqsplit{aa088f0647ad4b78966fe28fa80af674}}. Proposed · current.}\par
\label{element:e5edbb54258c490aaba0816d980dc624}
\index{[CT2] Emergency Care (R.2.78)}
{\small Capability · Superseded historical identity — direct correspondence · superseded}\par
\href{../../../../continuum/strategy/capability/e5edbb54258c490aaba0816d980dc624.md}{Historical canonical definition and provenance}.\par
\section*{R.2.79 [CT2] Inpatient Care — Historical identity [superseded]}
\phantomsection
\label{navigation:R.2.79}
\addcontentsline{toc}{section}{R.2.79 [CT2] Inpatient Care — Historical identity [superseded]}

{\small Governed successor (direct correspondence): II.1.3.3.8 \hyperref[element:b39fb49c01b545a5b51d3313531efb1b]{[CT2] Inpatient Care}; Canonical ID: \texttt{\seqsplit{b39fb49c01b545a5b51d3313531efb1b}}. Proposed · current.}\par
\label{element:bcc40a14465e48aca42c9505a6853081}
\index{[CT2] Inpatient Care (R.2.79)}
{\small Capability · Superseded historical identity — direct correspondence · superseded}\par
\href{../../../../continuum/strategy/capability/bcc40a14465e48aca42c9505a6853081.md}{Historical canonical definition and provenance}.\par
\section*{R.2.80 [CT2] Medication Therapy — Historical identity [superseded]}
\phantomsection
\label{navigation:R.2.80}
\addcontentsline{toc}{section}{R.2.80 [CT2] Medication Therapy — Historical identity [superseded]}

{\small Governed successor (direct correspondence): II.1.3.3.9 \hyperref[element:bb04996164c94e0c827a020432f4e498]{[CT2] Medication Therapy}; Canonical ID: \texttt{\seqsplit{bb04996164c94e0c827a020432f4e498}}. Proposed · current.}\par
\label{element:db064e0cd4794479be1374e97a3d7e63}
\index{[CT2] Medication Therapy (R.2.80)}
{\small Capability · Superseded historical identity — direct correspondence · superseded}\par
\href{../../../../continuum/strategy/capability/db064e0cd4794479be1374e97a3d7e63.md}{Historical canonical definition and provenance}.\par
\section*{R.2.81 [CT2] Outreach Care — Historical identity [superseded]}
\phantomsection
\label{navigation:R.2.81}
\addcontentsline{toc}{section}{R.2.81 [CT2] Outreach Care — Historical identity [superseded]}

{\small Governed successor (direct correspondence): II.1.3.3.10 \hyperref[element:b9ddbe2dd6a7485ba9cd3482215ea09e]{[CT2] Outreach Care}; Canonical ID: \texttt{\seqsplit{b9ddbe2dd6a7485ba9cd3482215ea09e}}. Proposed · current.}\par
\label{element:75f8b32a3d5048ccbdfb30d61ae73dad}
\index{[CT2] Outreach Care (R.2.81)}
{\small Capability · Superseded historical identity — direct correspondence · superseded}\par
\href{../../../../continuum/strategy/capability/75f8b32a3d5048ccbdfb30d61ae73dad.md}{Historical canonical definition and provenance}.\par
\section*{R.2.82 [CT2] Preventive Care — Historical identity [superseded]}
\phantomsection
\label{navigation:R.2.82}
\addcontentsline{toc}{section}{R.2.82 [CT2] Preventive Care — Historical identity [superseded]}

{\small Governed successor (direct correspondence): II.1.3.3.11 \hyperref[element:dfdc4f0273ba4910848b3b65c9ef6ab2]{[CT2] Preventive Care}; Canonical ID: \texttt{\seqsplit{dfdc4f0273ba4910848b3b65c9ef6ab2}}. Proposed · current.}\par
\label{element:97f0d1acc37b4ca994c6c1ca006f2daa}
\index{[CT2] Preventive Care (R.2.82)}
{\small Capability · Superseded historical identity — direct correspondence · superseded}\par
\href{../../../../continuum/strategy/capability/97f0d1acc37b4ca994c6c1ca006f2daa.md}{Historical canonical definition and provenance}.\par
\section*{R.2.83 [CT2] Primary Care — Historical identity [superseded]}
\phantomsection
\label{navigation:R.2.83}
\addcontentsline{toc}{section}{R.2.83 [CT2] Primary Care — Historical identity [superseded]}

{\small Governed successor (direct correspondence): II.1.3.3.12 \hyperref[element:d118673d77d74c8cb8cec24f03adddd8]{[CT2] Primary Care}; Canonical ID: \texttt{\seqsplit{d118673d77d74c8cb8cec24f03adddd8}}. Proposed · current.}\par
\label{element:fe9d58b5cfe14e8cb7a83a1ac6686d65}
\index{[CT2] Primary Care (R.2.83)}
{\small Capability · Superseded historical identity — direct correspondence · superseded}\par
\href{../../../../continuum/strategy/capability/fe9d58b5cfe14e8cb7a83a1ac6686d65.md}{Historical canonical definition and provenance}.\par
\section*{R.2.84 [CT2] Rehabilitation — Historical identity [superseded]}
\phantomsection
\label{navigation:R.2.84}
\addcontentsline{toc}{section}{R.2.84 [CT2] Rehabilitation — Historical identity [superseded]}

{\small Governed successor (direct correspondence): II.1.3.3.13 \hyperref[element:99ed51abcc5e4e8ba3b3d0161b17544a]{[CT2] Rehabilitation}; Canonical ID: \texttt{\seqsplit{99ed51abcc5e4e8ba3b3d0161b17544a}}. Canonical · current.}\par
\label{element:28ec21d5918a49aaac12740c69f42d77}
\index{[CT2] Rehabilitation (R.2.84)}
{\small Capability · Superseded historical identity — direct correspondence · superseded}\par
\href{../../../../continuum/strategy/capability/28ec21d5918a49aaac12740c69f42d77.md}{Historical canonical definition and provenance}.\par
\section*{R.2.85 [CT2] Remote \& Virtual Care — Historical identity [superseded]}
\phantomsection
\label{navigation:R.2.85}
\addcontentsline{toc}{section}{R.2.85 [CT2] Remote \& Virtual Care — Historical identity [superseded]}

{\small Governed successor (direct correspondence): II.1.3.3.14 \hyperref[element:30dff22d40b64da29ede7c41673db085]{[CT2] Remote \& Virtual Care}; Canonical ID: \texttt{\seqsplit{30dff22d40b64da29ede7c41673db085}}. Proposed · current.}\par
\label{element:75e2b447659c4fdf94107d94758337ed}
\index{[CT2] Remote \& Virtual Care (R.2.85)}
{\small Capability · Superseded historical identity — direct correspondence · superseded}\par
\href{../../../../continuum/strategy/capability/75e2b447659c4fdf94107d94758337ed.md}{Historical canonical definition and provenance}.\par
\section*{R.2.86 [CT2] Surgical \& Procedural Care — Historical identity [superseded]}
\phantomsection
\label{navigation:R.2.86}
\addcontentsline{toc}{section}{R.2.86 [CT2] Surgical \& Procedural Care — Historical identity [superseded]}

{\small Governed successor (direct correspondence): II.1.3.3.15 \hyperref[element:c9a26d1aa1014545b9ae759b663c807b]{[CT2] Surgical \& Procedural Care}; Canonical ID: \texttt{\seqsplit{c9a26d1aa1014545b9ae759b663c807b}}. Canonical · current.}\par
\label{element:8bc2a106165a4702b5f72fe46a206b69}
\index{[CT2] Surgical \& Procedural Care (R.2.86)}
{\small Capability · Superseded historical identity — direct correspondence · superseded}\par
\href{../../../../continuum/strategy/capability/8bc2a106165a4702b5f72fe46a206b69.md}{Historical canonical definition and provenance}.\par
\section*{R.2.87 [CT2] Therapeutic Services — Historical identity [superseded]}
\phantomsection
\label{navigation:R.2.87}
\addcontentsline{toc}{section}{R.2.87 [CT2] Therapeutic Services — Historical identity [superseded]}

{\small Governed successor (direct correspondence): II.1.3.3.16 \hyperref[element:b512b6d326004cb7b91a627fe6dbd8be]{[CT2] Therapeutic Services}; Canonical ID: \texttt{\seqsplit{b512b6d326004cb7b91a627fe6dbd8be}}. Canonical · current.}\par
\label{element:54032b52ee444bfe8b611fb9e8e2ba10}
\index{[CT2] Therapeutic Services (R.2.87)}
{\small Capability · Superseded historical identity — direct correspondence · superseded}\par
\href{../../../../continuum/strategy/capability/54032b52ee444bfe8b611fb9e8e2ba10.md}{Historical canonical definition and provenance}.\par
\section*{R.2.88 [CT1] Clinical Service Dependency Management — Historical identity [superseded]}
\phantomsection
\label{navigation:R.2.88}
\addcontentsline{toc}{section}{R.2.88 [CT1] Clinical Service Dependency Management — Historical identity [superseded]}

{\small Governed successor (direct correspondence): II.1.3.3 \hyperref[element:232db9b8b47c450cb0aa9acaf539c09d]{[CT1] Clinical Service Dependency Management}; Canonical ID: \texttt{\seqsplit{232db9b8b47c450cb0aa9acaf539c09d}}. Proposed · current.}\par
\label{element:984ee17123eb4f92b3e55e55d24459b8}
\index{[CT1] Clinical Service Dependency Management (R.2.88)}
{\small Capability · Superseded historical identity — direct correspondence · superseded}\par
\href{../../../../continuum/strategy/capability/984ee17123eb4f92b3e55e55d24459b8.md}{Historical canonical definition and provenance}.\par
\section*{R.2.89 [CT3] Client Privacy Guardrail Management — Historical identity [superseded]}
\phantomsection
\label{navigation:R.2.89}
\addcontentsline{toc}{section}{R.2.89 [CT3] Client Privacy Guardrail Management — Historical identity [superseded]}

{\small Governed successor (direct correspondence): II.1.3.4.1.1 \hyperref[element:7577a6e1815243648dd5356c21cca151]{[CT3] Client Privacy Guardrail Management}; Canonical ID: \texttt{\seqsplit{7577a6e1815243648dd5356c21cca151}}. Canonical · current.}\par
\label{element:d2e6078b5db041cb904e0cadcebce6eb}
\index{[CT3] Client Privacy Guardrail Management (R.2.89)}
{\small Capability · Superseded historical identity — direct correspondence · superseded}\par
\href{../../../../continuum/strategy/capability/d2e6078b5db041cb904e0cadcebce6eb.md}{Historical canonical definition and provenance}.\par
\section*{R.2.90 [CT3] Client Relationship Administration — Historical identity [superseded]}
\phantomsection
\label{navigation:R.2.90}
\addcontentsline{toc}{section}{R.2.90 [CT3] Client Relationship Administration — Historical identity [superseded]}

{\small Governed successor (direct correspondence): II.1.3.4.1.2 \hyperref[element:d682c11d3eb04770b1f20be478fd475e]{[CT3] Client Relationship Administration}; Canonical ID: \texttt{\seqsplit{d682c11d3eb04770b1f20be478fd475e}}. Canonical · current.}\par
\label{element:2c799ec0c4c74c97a75c12e949887b95}
\index{[CT3] Client Relationship Administration (R.2.90)}
{\small Capability · Superseded historical identity — direct correspondence · superseded}\par
\href{../../../../continuum/strategy/capability/2c799ec0c4c74c97a75c12e949887b95.md}{Historical canonical definition and provenance}.\par
\section*{R.2.91 [CT3] Healthcare Subject Profile Administration — Historical identity [superseded]}
\phantomsection
\label{navigation:R.2.91}
\addcontentsline{toc}{section}{R.2.91 [CT3] Healthcare Subject Profile Administration — Historical identity [superseded]}

{\small Governed successor (direct correspondence): II.1.3.4.1.3 \hyperref[element:df214b5499694b768eb4a7a175a9d471]{[CT3] Healthcare Subject Profile Administration}; Canonical ID: \texttt{\seqsplit{df214b5499694b768eb4a7a175a9d471}}. Canonical · current.}\par
\label{element:d653c3247ddd49479f985f6b435933a0}
\index{[CT3] Healthcare Subject Profile Administration (R.2.91)}
{\small Capability · Superseded historical identity — direct correspondence · superseded}\par
\href{../../../../continuum/strategy/capability/d653c3247ddd49479f985f6b435933a0.md}{Historical canonical definition and provenance}.\par
\section*{R.2.92 [CT3] Person Identity Administration — Historical identity [superseded]}
\phantomsection
\label{navigation:R.2.92}
\addcontentsline{toc}{section}{R.2.92 [CT3] Person Identity Administration — Historical identity [superseded]}

{\small Governed successor (direct correspondence): II.1.3.4.1.4 \hyperref[element:6c3de7750ef1436cad6a9533731782b6]{[CT3] Person Identity Administration}; Canonical ID: \texttt{\seqsplit{6c3de7750ef1436cad6a9533731782b6}}. Canonical · current.}\par
\label{element:7550fd5b17e643118fa76d2616e1ba20}
\index{[CT3] Person Identity Administration (R.2.92)}
{\small Capability · Superseded historical identity — direct correspondence · superseded}\par
\href{../../../../continuum/strategy/capability/7550fd5b17e643118fa76d2616e1ba20.md}{Historical canonical definition and provenance}.\par
\section*{R.2.93 [CT2] Client Administration — Historical identity [superseded]}
\phantomsection
\label{navigation:R.2.93}
\addcontentsline{toc}{section}{R.2.93 [CT2] Client Administration — Historical identity [superseded]}

{\small Governed successor (direct correspondence): II.1.3.4.1 \hyperref[element:f0885472cced4bccbcce1568bf0360dc]{[CT2] Client Administration}; Canonical ID: \texttt{\seqsplit{f0885472cced4bccbcce1568bf0360dc}}. Canonical · current.}\par
\label{element:70d2a44a0442441f8b306b82b297bb8a}
\index{[CT2] Client Administration (R.2.93)}
{\small Capability · Superseded historical identity — direct correspondence · superseded}\par
\href{../../../../continuum/strategy/capability/70d2a44a0442441f8b306b82b297bb8a.md}{Historical canonical definition and provenance}.\par
\section*{R.2.94 [CT3] Clinical Device Profile Administration — Historical identity [superseded]}
\phantomsection
\label{navigation:R.2.94}
\addcontentsline{toc}{section}{R.2.94 [CT3] Clinical Device Profile Administration — Historical identity [superseded]}

{\small Governed successor (direct correspondence): II.1.3.4.2.1 \hyperref[element:d0a1ba80c790476ea6042290eb7ea09e]{[CT3] Clinical Device Profile Administration}; Canonical ID: \texttt{\seqsplit{d0a1ba80c790476ea6042290eb7ea09e}}. Proposed · current.}\par
\label{element:e8c72952da144d6883caba7fd2afb88d}
\index{[CT3] Clinical Device Profile Administration (R.2.94)}
{\small Capability · Superseded historical identity — direct correspondence · superseded}\par
\href{../../../../continuum/strategy/capability/e8c72952da144d6883caba7fd2afb88d.md}{Historical canonical definition and provenance}.\par
\section*{R.2.95 [CT3] Device Connectivity Compliance Management — Historical identity [superseded]}
\phantomsection
\label{navigation:R.2.95}
\addcontentsline{toc}{section}{R.2.95 [CT3] Device Connectivity Compliance Management — Historical identity [superseded]}

{\small Governed successor (direct correspondence): II.1.3.4.2.4 \hyperref[element:b16ab359030a4583868f3aac2221074e]{[CT3] Device Connectivity Compliance Management}; Canonical ID: \texttt{\seqsplit{b16ab359030a4583868f3aac2221074e}}. Proposed · current.}\par
\label{element:9225f083d3474715a75af2a827a4016d}
\index{[CT3] Device Connectivity Compliance Management (R.2.95)}
{\small Capability · Superseded historical identity — direct correspondence · superseded}\par
\href{../../../../continuum/strategy/capability/9225f083d3474715a75af2a827a4016d.md}{Historical canonical definition and provenance}.\par
\section*{R.2.96 [CT3] Device OA\&M Metadata Management — Historical identity [superseded]}
\phantomsection
\label{navigation:R.2.96}
\addcontentsline{toc}{section}{R.2.96 [CT3] Device OA\&M Metadata Management — Historical identity [superseded]}

{\small Governed successor (direct correspondence): II.1.3.4.2.6 \hyperref[element:659bb194a7ef42ab9ba8a938adb5c179]{[CT3] Device OA\&M Metadata Management}; Canonical ID: \texttt{\seqsplit{659bb194a7ef42ab9ba8a938adb5c179}}. Proposed · current.}\par
\label{element:a398c72bef094dee96c96d97994449fb}
\index{[CT3] Device OA\&M Metadata Management (R.2.96)}
{\small Capability · Superseded historical identity — direct correspondence · superseded}\par
\href{../../../../continuum/strategy/capability/a398c72bef094dee96c96d97994449fb.md}{Historical canonical definition and provenance}.\par
\section*{R.2.97 [CT3] Device Utilisation Administration — Historical identity [superseded]}
\phantomsection
\label{navigation:R.2.97}
\addcontentsline{toc}{section}{R.2.97 [CT3] Device Utilisation Administration — Historical identity [superseded]}

{\small Governed successor (direct correspondence): II.1.3.4.2.7 \hyperref[element:1a3732b5911c45b3aece38b2db2d0f15]{[CT3] Device Utilisation Administration}; Canonical ID: \texttt{\seqsplit{1a3732b5911c45b3aece38b2db2d0f15}}. Proposed · current.}\par
\label{element:a0af0a86667d48f4ab95ed2920e878cb}
\index{[CT3] Device Utilisation Administration (R.2.97)}
{\small Capability · Superseded historical identity — direct correspondence · superseded}\par
\href{../../../../continuum/strategy/capability/a0af0a86667d48f4ab95ed2920e878cb.md}{Historical canonical definition and provenance}.\par
\section*{R.2.98 [CT2] Clinical Device Administration — Historical identity [superseded]}
\phantomsection
\label{navigation:R.2.98}
\addcontentsline{toc}{section}{R.2.98 [CT2] Clinical Device Administration — Historical identity [superseded]}

{\small Governed successor (direct correspondence): II.1.3.4.2 \hyperref[element:31bcd6fda0d14cc8b992cdca39a3cbed]{[CT2] Clinical Device Administration}; Canonical ID: \texttt{\seqsplit{31bcd6fda0d14cc8b992cdca39a3cbed}}. Canonical · current.}\par
\label{element:3792be10541e4d68b3f740b1ba964966}
\index{[CT2] Clinical Device Administration (R.2.98)}
{\small Capability · Superseded historical identity — direct correspondence · superseded}\par
\href{../../../../continuum/strategy/capability/3792be10541e4d68b3f740b1ba964966.md}{Historical canonical definition and provenance}.\par
\section*{R.2.99 [CT3] Health Materiel Catalogue \& Formulary Administration — Historical identity [superseded]}
\phantomsection
\label{navigation:R.2.99}
\addcontentsline{toc}{section}{R.2.99 [CT3] Health Materiel Catalogue \& Formulary Administration — Historical identity [superseded]}

{\small Governed successor (direct correspondence): II.1.3.4.3.1 \hyperref[element:28b327760af94b8b986357469524d793]{[CT3] Health Materiel Catalogue \& Formulary Administration}; Canonical ID: \texttt{\seqsplit{28b327760af94b8b986357469524d793}}. Proposed · current.}\par
\label{element:5c691959e2cc4a289fc25b9f3497eed9}
\index{[CT3] Health Materiel Catalogue \& Formulary Administration (R.2.99)}
{\small Capability · Superseded historical identity — direct correspondence · superseded}\par
\href{../../../../continuum/strategy/capability/5c691959e2cc4a289fc25b9f3497eed9.md}{Historical canonical definition and provenance}.\par
\section*{R.2.100 [CT3] Health Materiel Profile Administration — Historical identity [superseded]}
\phantomsection
\label{navigation:R.2.100}
\addcontentsline{toc}{section}{R.2.100 [CT3] Health Materiel Profile Administration — Historical identity [superseded]}

{\small Governed successor (direct correspondence): II.1.3.4.3.2 \hyperref[element:9626edae517d40d2bfae619860f18e20]{[CT3] Health Materiel Profile Administration}; Canonical ID: \texttt{\seqsplit{9626edae517d40d2bfae619860f18e20}}. Proposed · current.}\par
\label{element:9216f9cc905e4526b1d98cd4ab382cbf}
\index{[CT3] Health Materiel Profile Administration (R.2.100)}
{\small Capability · Superseded historical identity — direct correspondence · superseded}\par
\href{../../../../continuum/strategy/capability/9216f9cc905e4526b1d98cd4ab382cbf.md}{Historical canonical definition and provenance}.\par
\section*{R.2.101 [CT3] Health Materiel Traceability Administration — Historical identity [superseded]}
\phantomsection
\label{navigation:R.2.101}
\addcontentsline{toc}{section}{R.2.101 [CT3] Health Materiel Traceability Administration — Historical identity [superseded]}

{\small Governed successor (direct correspondence): II.1.3.4.3.3 \hyperref[element:2d41ed0f00a94a4b90ffe87d6fd4e9fe]{[CT3] Health Materiel Traceability Administration}; Canonical ID: \texttt{\seqsplit{2d41ed0f00a94a4b90ffe87d6fd4e9fe}}. Proposed · current.}\par
\label{element:f1b888e8d50c4019ac1afe162d6223f3}
\index{[CT3] Health Materiel Traceability Administration (R.2.101)}
{\small Capability · Superseded historical identity — direct correspondence · superseded}\par
\href{../../../../continuum/strategy/capability/f1b888e8d50c4019ac1afe162d6223f3.md}{Historical canonical definition and provenance}.\par
\section*{R.2.102 [CT2] Health Materiel Administration — Historical identity [superseded]}
\phantomsection
\label{navigation:R.2.102}
\addcontentsline{toc}{section}{R.2.102 [CT2] Health Materiel Administration — Historical identity [superseded]}

{\small Governed successor (direct correspondence): II.1.3.4.3 \hyperref[element:e7081070a9f8444f81c1d8fce37ef664]{[CT2] Health Materiel Administration}; Canonical ID: \texttt{\seqsplit{e7081070a9f8444f81c1d8fce37ef664}}. Canonical · current.}\par
\label{element:4952f86078c54a8d898ac992e6f4f577}
\index{[CT2] Health Materiel Administration (R.2.102)}
{\small Capability · Superseded historical identity — direct correspondence · superseded}\par
\href{../../../../continuum/strategy/capability/4952f86078c54a8d898ac992e6f4f577.md}{Historical canonical definition and provenance}.\par
\section*{R.2.103 [CT3] Healthcare Service Profile Administration — Historical identity [superseded]}
\phantomsection
\label{navigation:R.2.103}
\addcontentsline{toc}{section}{R.2.103 [CT3] Healthcare Service Profile Administration — Historical identity [superseded]}

{\small Governed successor (direct correspondence): II.1.3.4.4.1 \hyperref[element:e75d5eb30936476ea7a9cbadfa43d5ef]{[CT3] Healthcare Service Profile Administration}; Canonical ID: \texttt{\seqsplit{e75d5eb30936476ea7a9cbadfa43d5ef}}. Proposed · current.}\par
\label{element:7780ca2fed2346439d5f559d00cc4ee4}
\index{[CT3] Healthcare Service Profile Administration (R.2.103)}
{\small Capability · Superseded historical identity — direct correspondence · superseded}\par
\href{../../../../continuum/strategy/capability/7780ca2fed2346439d5f559d00cc4ee4.md}{Historical canonical definition and provenance}.\par
\section*{R.2.104 [CT2] Health Service Administration — Historical identity [superseded]}
\phantomsection
\label{navigation:R.2.104}
\addcontentsline{toc}{section}{R.2.104 [CT2] Health Service Administration — Historical identity [superseded]}

{\small Governed successor (direct correspondence): II.1.3.4.4 \hyperref[element:ba68a52913a0499d8377553bddfa74a4]{[CT2] Health Service Administration}; Canonical ID: \texttt{\seqsplit{ba68a52913a0499d8377553bddfa74a4}}. Canonical · current.}\par
\label{element:f1ab1258cf2944efb50105e7885f8027}
\index{[CT2] Health Service Administration (R.2.104)}
{\small Capability · Superseded historical identity — direct correspondence · superseded}\par
\href{../../../../continuum/strategy/capability/f1ab1258cf2944efb50105e7885f8027.md}{Historical canonical definition and provenance}.\par
\section*{R.2.105 [CT3] Care-Place Definition Administration — Historical identity [superseded]}
\phantomsection
\label{navigation:R.2.105}
\addcontentsline{toc}{section}{R.2.105 [CT3] Care-Place Definition Administration — Historical identity [superseded]}

{\small Governed successor (direct correspondence): II.1.3.4.5.1 \hyperref[element:629b493c7d9844baa25c029a87ed04f6]{[CT3] Care-Place Definition Administration}; Canonical ID: \texttt{\seqsplit{629b493c7d9844baa25c029a87ed04f6}}. Proposed · current.}\par
\label{element:a3464bd9c0dc4a01b50803acc7705b4e}
\index{[CT3] Care-Place Definition Administration (R.2.105)}
{\small Capability · Superseded historical identity — direct correspondence · superseded}\par
\href{../../../../continuum/strategy/capability/a3464bd9c0dc4a01b50803acc7705b4e.md}{Historical canonical definition and provenance}.\par
\section*{R.2.106 [CT3] Healthcare Location Profile Administration — Historical identity [superseded]}
\phantomsection
\label{navigation:R.2.106}
\addcontentsline{toc}{section}{R.2.106 [CT3] Healthcare Location Profile Administration — Historical identity [superseded]}

{\small Governed successor (direct correspondence): II.1.3.4.5.2 \hyperref[element:e579e54a08bb4f899a504aa998ba246b]{[CT3] Healthcare Location Profile Administration}; Canonical ID: \texttt{\seqsplit{e579e54a08bb4f899a504aa998ba246b}}. Proposed · current.}\par
\label{element:8a04d1f16d0849e786d2aaec33e0ba87}
\index{[CT3] Healthcare Location Profile Administration (R.2.106)}
{\small Capability · Superseded historical identity — direct correspondence · superseded}\par
\href{../../../../continuum/strategy/capability/8a04d1f16d0849e786d2aaec33e0ba87.md}{Historical canonical definition and provenance}.\par
\section*{R.2.107 [CT2] Location Administration — Historical identity [superseded]}
\phantomsection
\label{navigation:R.2.107}
\addcontentsline{toc}{section}{R.2.107 [CT2] Location Administration — Historical identity [superseded]}

{\small Governed successor (direct correspondence): II.1.3.4.5 \hyperref[element:c410005667444f99a221483bf577fad2]{[CT2] Location Administration}; Canonical ID: \texttt{\seqsplit{c410005667444f99a221483bf577fad2}}. Canonical · current.}\par
\label{element:d412aff3a5c64e6aaba910d4c2c9356d}
\index{[CT2] Location Administration (R.2.107)}
{\small Capability · Superseded historical identity — direct correspondence · superseded}\par
\href{../../../../continuum/strategy/capability/d412aff3a5c64e6aaba910d4c2c9356d.md}{Historical canonical definition and provenance}.\par
\section*{R.2.108 [CT3] Healthcare Organisation Profile Administration — Historical identity [superseded]}
\phantomsection
\label{navigation:R.2.108}
\addcontentsline{toc}{section}{R.2.108 [CT3] Healthcare Organisation Profile Administration — Historical identity [superseded]}

{\small Governed successor (direct correspondence): II.1.3.4.6.1 \hyperref[element:f2de9e9394b849f2943355c7a2882a15]{[CT3] Healthcare Organisation Profile Administration}; Canonical ID: \texttt{\seqsplit{f2de9e9394b849f2943355c7a2882a15}}. Proposed · current.}\par
\label{element:e485424b799c42aca7399d341ca02d0b}
\index{[CT3] Healthcare Organisation Profile Administration (R.2.108)}
{\small Capability · Superseded historical identity — direct correspondence · superseded}\par
\href{../../../../continuum/strategy/capability/e485424b799c42aca7399d341ca02d0b.md}{Historical canonical definition and provenance}.\par
\section*{R.2.109 [CT2] Organisation Administration — Historical identity [superseded]}
\phantomsection
\label{navigation:R.2.109}
\addcontentsline{toc}{section}{R.2.109 [CT2] Organisation Administration — Historical identity [superseded]}

{\small Governed successor (direct correspondence): II.1.3.4.6 \hyperref[element:14fbdc5ce3dd4c2092c5397dab3ed743]{[CT2] Organisation Administration}; Canonical ID: \texttt{\seqsplit{14fbdc5ce3dd4c2092c5397dab3ed743}}. Canonical · current.}\par
\label{element:1b0eacfbb27d4631a24f9279dab6681d}
\index{[CT2] Organisation Administration (R.2.109)}
{\small Capability · Superseded historical identity — direct correspondence · superseded}\par
\href{../../../../continuum/strategy/capability/1b0eacfbb27d4631a24f9279dab6681d.md}{Historical canonical definition and provenance}.\par
\section*{R.2.110 [CT3] Practitioner Profile Administration — Historical identity [superseded]}
\phantomsection
\label{navigation:R.2.110}
\addcontentsline{toc}{section}{R.2.110 [CT3] Practitioner Profile Administration — Historical identity [superseded]}

{\small Governed successor (direct correspondence): II.1.3.4.7.1 \hyperref[element:40759a835f7b456e9257c38b67b0a380]{[CT3] Practitioner Profile Administration}; Canonical ID: \texttt{\seqsplit{40759a835f7b456e9257c38b67b0a380}}. Proposed · current.}\par
\label{element:2415ddcbc73c47e1a63be310bd4342ee}
\index{[CT3] Practitioner Profile Administration (R.2.110)}
{\small Capability · Superseded historical identity — direct correspondence · superseded}\par
\href{../../../../continuum/strategy/capability/2415ddcbc73c47e1a63be310bd4342ee.md}{Historical canonical definition and provenance}.\par
\section*{R.2.111 [CT3] Practitioner Role Administration — Historical identity [superseded]}
\phantomsection
\label{navigation:R.2.111}
\addcontentsline{toc}{section}{R.2.111 [CT3] Practitioner Role Administration — Historical identity [superseded]}

{\small Governed successor (direct correspondence): II.1.3.4.7.2 \hyperref[element:6651471e67ae46fb8d1c1241c8590be9]{[CT3] Practitioner Role Administration}; Canonical ID: \texttt{\seqsplit{6651471e67ae46fb8d1c1241c8590be9}}. Proposed · current.}\par
\label{element:6b777f1201db474eb22dd89f2b1f68c6}
\index{[CT3] Practitioner Role Administration (R.2.111)}
{\small Capability · Superseded historical identity — direct correspondence · superseded}\par
\href{../../../../continuum/strategy/capability/6b777f1201db474eb22dd89f2b1f68c6.md}{Historical canonical definition and provenance}.\par
\section*{R.2.112 [CT3] Vendor Liaison — Historical identity [superseded]}
\phantomsection
\label{navigation:R.2.112}
\addcontentsline{toc}{section}{R.2.112 [CT3] Vendor Liaison — Historical identity [superseded]}

{\small Governed successor (direct correspondence): II.1.3.4.7.5 \hyperref[element:66358bd6698e4050b12342522fbbb923]{[CT3] Vendor Liaison}; Canonical ID: \texttt{\seqsplit{66358bd6698e4050b12342522fbbb923}}. Proposed · current.}\par
\label{element:29076cbb9fc5442fb25faf43947574f7}
\index{[CT3] Vendor Liaison (R.2.112)}
{\small Capability · Superseded historical identity — direct correspondence · superseded}\par
\href{../../../../continuum/strategy/capability/29076cbb9fc5442fb25faf43947574f7.md}{Historical canonical definition and provenance}.\par
\section*{R.2.113 [CT2] Provider Administration — Historical identity [superseded]}
\phantomsection
\label{navigation:R.2.113}
\addcontentsline{toc}{section}{R.2.113 [CT2] Provider Administration — Historical identity [superseded]}

{\small Governed successor (direct correspondence): II.1.3.4.7 \hyperref[element:d6312f85968a4209bafff5ed009297a6]{[CT2] Provider Administration}; Canonical ID: \texttt{\seqsplit{d6312f85968a4209bafff5ed009297a6}}. Canonical · current.}\par
\label{element:a352fd50419d451887c85fcb98566cd5}
\index{[CT2] Provider Administration (R.2.113)}
{\small Capability · Superseded historical identity — direct correspondence · superseded}\par
\href{../../../../continuum/strategy/capability/a352fd50419d451887c85fcb98566cd5.md}{Historical canonical definition and provenance}.\par
\section*{R.2.114 [CT1] Entity Information Enablement — Historical identity [superseded]}
\phantomsection
\label{navigation:R.2.114}
\addcontentsline{toc}{section}{R.2.114 [CT1] Entity Information Enablement — Historical identity [superseded]}

{\small Governed successor (direct correspondence): II.1.3.4 \hyperref[element:b5a92874f8ac4509a30af5253590edc1]{[CT1] Entity Information Enablement}; Canonical ID: \texttt{\seqsplit{b5a92874f8ac4509a30af5253590edc1}}. Proposed · current.}\par
\label{element:e4d959d6c71d498c9ce00d2be4c048e3}
\index{[CT1] Entity Information Enablement (R.2.114)}
{\small Capability · Superseded historical identity — direct correspondence · superseded}\par
\href{../../../../continuum/strategy/capability/e4d959d6c71d498c9ce00d2be4c048e3.md}{Historical canonical definition and provenance}.\par
\section*{R.2.115 [CT2] Population Health Information Reporting — Historical identity [superseded]}
\phantomsection
\label{navigation:R.2.115}
\addcontentsline{toc}{section}{R.2.115 [CT2] Population Health Information Reporting — Historical identity [superseded]}

{\small Governed successor (direct correspondence): II.1.3.5.1 \hyperref[element:2135b45eab034370b0063b20e821985d]{[CT2] Population Health Information Reporting}; Canonical ID: \texttt{\seqsplit{2135b45eab034370b0063b20e821985d}}. Proposed · current.}\par
\label{element:ced472b2d1a54016878ee541de969d9b}
\index{[CT2] Population Health Information Reporting (R.2.115)}
{\small Capability · Superseded historical identity — direct correspondence · superseded}\par
\href{../../../../continuum/strategy/capability/ced472b2d1a54016878ee541de969d9b.md}{Historical canonical definition and provenance}.\par
\section*{R.2.116 [CT2] Research Information Curation — Historical identity [superseded]}
\phantomsection
\label{navigation:R.2.116}
\addcontentsline{toc}{section}{R.2.116 [CT2] Research Information Curation — Historical identity [superseded]}

{\small Governed successor (direct correspondence): II.1.3.5.2 \hyperref[element:4f5b8cec90324c2989cdc30e80bc8cb9]{[CT2] Research Information Curation}; Canonical ID: \texttt{\seqsplit{4f5b8cec90324c2989cdc30e80bc8cb9}}. Proposed · current.}\par
\label{element:11203b9c80964a9b96883eeab07f41d5}
\index{[CT2] Research Information Curation (R.2.116)}
{\small Capability · Superseded historical identity — direct correspondence · superseded}\par
\href{../../../../continuum/strategy/capability/11203b9c80964a9b96883eeab07f41d5.md}{Historical canonical definition and provenance}.\par
\section*{R.2.117 [CT2] Service Planning Data Management — Historical identity [superseded]}
\phantomsection
\label{navigation:R.2.117}
\addcontentsline{toc}{section}{R.2.117 [CT2] Service Planning Data Management — Historical identity [superseded]}

{\small Governed successor (direct correspondence): II.1.3.5.3 \hyperref[element:4873cd4b960d45d8b93560d5c8b98500]{[CT2] Service Planning Data Management}; Canonical ID: \texttt{\seqsplit{4873cd4b960d45d8b93560d5c8b98500}}. Proposed · current.}\par
\label{element:19263e0802cc4b88af7a345f1410cffd}
\index{[CT2] Service Planning Data Management (R.2.117)}
{\small Capability · Superseded historical identity — direct correspondence · superseded}\par
\href{../../../../continuum/strategy/capability/19263e0802cc4b88af7a345f1410cffd.md}{Historical canonical definition and provenance}.\par
\section*{R.2.118 [CT1] Health Informatics Management — Historical identity [superseded]}
\phantomsection
\label{navigation:R.2.118}
\addcontentsline{toc}{section}{R.2.118 [CT1] Health Informatics Management — Historical identity [superseded]}

{\small Governed successor (direct correspondence): II.1.3.5 \hyperref[element:a790b5f2c5a1497fa6ec09d7824da32e]{[CT1] Health Informatics Management}; Canonical ID: \texttt{\seqsplit{a790b5f2c5a1497fa6ec09d7824da32e}}. Proposed · current.}\par
\label{element:50c56bac0dc540a2b392fc2376efc129}
\index{[CT1] Health Informatics Management (R.2.118)}
{\small Capability · Superseded historical identity — direct correspondence · superseded}\par
\href{../../../../continuum/strategy/capability/50c56bac0dc540a2b392fc2376efc129.md}{Historical canonical definition and provenance}.\par
\section*{R.2.119 [CT3] Bed \& Care-Place State Administration — Historical identity [superseded]}
\phantomsection
\label{navigation:R.2.119}
\addcontentsline{toc}{section}{R.2.119 [CT3] Bed \& Care-Place State Administration — Historical identity [superseded]}

{\small Governed successor (direct correspondence): II.1.3.6.1.1 \hyperref[element:26a9383d7bb146a19141558d232843de]{[CT3] Bed \& Care-Place State Administration}; Canonical ID: \texttt{\seqsplit{26a9383d7bb146a19141558d232843de}}. Proposed · current.}\par
\label{element:dd2e2ee1a68845c087ba93610752181b}
\index{[CT3] Bed \& Care-Place State Administration (R.2.119)}
{\small Capability · Superseded historical identity — direct correspondence · superseded}\par
\href{../../../../continuum/strategy/capability/dd2e2ee1a68845c087ba93610752181b.md}{Historical canonical definition and provenance}.\par
\section*{R.2.120 [CT3] Bed \& Care-Place Turnover Coordination — Historical identity [superseded]}
\phantomsection
\label{navigation:R.2.120}
\addcontentsline{toc}{section}{R.2.120 [CT3] Bed \& Care-Place Turnover Coordination — Historical identity [superseded]}

{\small Governed successor (direct correspondence): II.1.3.6.1.2 \hyperref[element:3c9464f9eae1415d8fa041bd1f3da4f2]{[CT3] Bed \& Care-Place Turnover Coordination}; Canonical ID: \texttt{\seqsplit{3c9464f9eae1415d8fa041bd1f3da4f2}}. Proposed · current.}\par
\label{element:a50662edda954bc58c44184d84591d5b}
\index{[CT3] Bed \& Care-Place Turnover Coordination (R.2.120)}
{\small Capability · Superseded historical identity — direct correspondence · superseded}\par
\href{../../../../continuum/strategy/capability/a50662edda954bc58c44184d84591d5b.md}{Historical canonical definition and provenance}.\par
\section*{R.2.121 [CT2] Bed \& Care-Place Management — Historical identity [superseded]}
\phantomsection
\label{navigation:R.2.121}
\addcontentsline{toc}{section}{R.2.121 [CT2] Bed \& Care-Place Management — Historical identity [superseded]}

{\small Governed successor (direct correspondence): II.1.3.6.1 \hyperref[element:f117ccec0bfa4f738e6ddb646b3534a2]{[CT2] Bed \& Care-Place Management}; Canonical ID: \texttt{\seqsplit{f117ccec0bfa4f738e6ddb646b3534a2}}. Canonical · current.}\par
\label{element:c94d29b7d1c045dabc7814ceea795484}
\index{[CT2] Bed \& Care-Place Management (R.2.121)}
{\small Capability · Superseded historical identity — direct correspondence · superseded}\par
\href{../../../../continuum/strategy/capability/c94d29b7d1c045dabc7814ceea795484.md}{Historical canonical definition and provenance}.\par
\section*{R.2.122 [CT3] Session \& Queue Coordination — Historical identity [superseded]}
\phantomsection
\label{navigation:R.2.122}
\addcontentsline{toc}{section}{R.2.122 [CT3] Session \& Queue Coordination — Historical identity [superseded]}

{\small Governed successor (direct correspondence): II.1.3.6.5.1 \hyperref[element:22981a95ca6e466582363706fde8fceb]{[CT3] Session \& Queue Coordination}; Canonical ID: \texttt{\seqsplit{22981a95ca6e466582363706fde8fceb}}. Proposed · current.}\par
\label{element:5298c6170cef4f05a6d956da716071bd}
\index{[CT3] Session \& Queue Coordination (R.2.122)}
{\small Capability · Superseded historical identity — direct correspondence · superseded}\par
\href{../../../../continuum/strategy/capability/5298c6170cef4f05a6d956da716071bd.md}{Historical canonical definition and provenance}.\par
\section*{R.2.123 [CT2] Clinic \& Practice Operations — Historical identity [superseded]}
\phantomsection
\label{navigation:R.2.123}
\addcontentsline{toc}{section}{R.2.123 [CT2] Clinic \& Practice Operations — Historical identity [superseded]}

{\small Governed successor (direct correspondence): II.1.3.6.5 \hyperref[element:3dffe63e03124efab255204ee4e9a5d9]{[CT2] Clinic \& Practice Operations}; Canonical ID: \texttt{\seqsplit{3dffe63e03124efab255204ee4e9a5d9}}. Canonical · current.}\par
\label{element:0856d2daf2ba4b92b5390b38b98efb54}
\index{[CT2] Clinic \& Practice Operations (R.2.123)}
{\small Capability · Superseded historical identity — direct correspondence · superseded}\par
\href{../../../../continuum/strategy/capability/0856d2daf2ba4b92b5390b38b98efb54.md}{Historical canonical definition and provenance}.\par
\section*{R.2.124 [CT2] Clinical Competency Administration — Historical identity [superseded]}
\phantomsection
\label{navigation:R.2.124}
\addcontentsline{toc}{section}{R.2.124 [CT2] Clinical Competency Administration — Historical identity [superseded]}

{\small Governed successor (direct correspondence): II.1.3.6.6 \hyperref[element:8a99f4a6e5054d6c997bacec04805c01]{[CT2] Clinical Competency Administration}; Canonical ID: \texttt{\seqsplit{8a99f4a6e5054d6c997bacec04805c01}}. Canonical · current.}\par
\label{element:8a4b918e1c9d47bb99daf2b42008ab39}
\index{[CT2] Clinical Competency Administration (R.2.124)}
{\small Capability · Superseded historical identity — direct correspondence · superseded}\par
\href{../../../../continuum/strategy/capability/8a4b918e1c9d47bb99daf2b42008ab39.md}{Historical canonical definition and provenance}.\par
\section*{R.2.125 [CT3] Specimen Metadata Management — Historical identity [superseded]}
\phantomsection
\label{navigation:R.2.125}
\addcontentsline{toc}{section}{R.2.125 [CT3] Specimen Metadata Management — Historical identity [superseded]}

{\small Governed successor (direct correspondence): II.1.3.6.7.1 \hyperref[element:738736425b7446fdae47b06c91c4f69c]{[CT3] Specimen Metadata Management}; Canonical ID: \texttt{\seqsplit{738736425b7446fdae47b06c91c4f69c}}. Proposed · current.}\par
\label{element:3b000d3712534d16a48470f79df8e33a}
\index{[CT3] Specimen Metadata Management (R.2.125)}
{\small Capability · Superseded historical identity — direct correspondence · superseded}\par
\href{../../../../continuum/strategy/capability/3b000d3712534d16a48470f79df8e33a.md}{Historical canonical definition and provenance}.\par
\section*{R.2.126 [CT3] Specimen Movement Traceability — Historical identity [superseded]}
\phantomsection
\label{navigation:R.2.126}
\addcontentsline{toc}{section}{R.2.126 [CT3] Specimen Movement Traceability — Historical identity [superseded]}

{\small Governed successor (direct correspondence): II.1.3.6.7.2 \hyperref[element:bfa2e9a84ddd4287897573dc0d917a00]{[CT3] Specimen Movement Traceability}; Canonical ID: \texttt{\seqsplit{bfa2e9a84ddd4287897573dc0d917a00}}. Proposed · current.}\par
\label{element:129119d7b2744ffd805c58f25557fa27}
\index{[CT3] Specimen Movement Traceability (R.2.126)}
{\small Capability · Superseded historical identity — direct correspondence · superseded}\par
\href{../../../../continuum/strategy/capability/129119d7b2744ffd805c58f25557fa27.md}{Historical canonical definition and provenance}.\par
\section*{R.2.127 [CT2] Clinical Logistics Coordination — Historical identity [superseded]}
\phantomsection
\label{navigation:R.2.127}
\addcontentsline{toc}{section}{R.2.127 [CT2] Clinical Logistics Coordination — Historical identity [superseded]}

{\small Governed successor (direct correspondence): II.1.3.6.7 \hyperref[element:0c27be11b4a44cb19e992309b8146a17]{[CT2] Clinical Logistics Coordination}; Canonical ID: \texttt{\seqsplit{0c27be11b4a44cb19e992309b8146a17}}. Canonical · current.}\par
\label{element:90a782ff2d004fa984ae0465528b2253}
\index{[CT2] Clinical Logistics Coordination (R.2.127)}
{\small Capability · Superseded historical identity — direct correspondence · superseded}\par
\href{../../../../continuum/strategy/capability/90a782ff2d004fa984ae0465528b2253.md}{Historical canonical definition and provenance}.\par
\section*{R.2.128 [CT3] Organisation Clinical Certification \& Compliance Support — Historical identity [superseded]}
\phantomsection
\label{navigation:R.2.128}
\addcontentsline{toc}{section}{R.2.128 [CT3] Organisation Clinical Certification \& Compliance Support — Historical identity [superseded]}

{\small Governed successor (direct correspondence): II.1.3.6.8.1 \hyperref[element:4de624e7bf49472f91e3e5c734ee8725]{[CT3] Organisation Clinical Certification \& Compliance Support}; Canonical ID: \texttt{\seqsplit{4de624e7bf49472f91e3e5c734ee8725}}. Proposed · current.}\par
\label{element:cee76a930a4e4e7d9c49357ba290c8bf}
\index{[CT3] Organisation Clinical Certification \& Compliance Support (R.2.128)}
{\small Capability · Superseded historical identity — direct correspondence · superseded}\par
\href{../../../../continuum/strategy/capability/cee76a930a4e4e7d9c49357ba290c8bf.md}{Historical canonical definition and provenance}.\par
\section*{R.2.129 [CT3] Provider Clinical Qualification Tracking — Historical identity [superseded]}
\phantomsection
\label{navigation:R.2.129}
\addcontentsline{toc}{section}{R.2.129 [CT3] Provider Clinical Qualification Tracking — Historical identity [superseded]}

{\small Governed successor (direct correspondence): II.1.3.6.8.2 \hyperref[element:8fd5c78eba3046fa9153fb4efb44e9c8]{[CT3] Provider Clinical Qualification Tracking}; Canonical ID: \texttt{\seqsplit{8fd5c78eba3046fa9153fb4efb44e9c8}}. Proposed · current.}\par
\label{element:ab0415e2ea0a4da4ad3dc46997d48bc3}
\index{[CT3] Provider Clinical Qualification Tracking (R.2.129)}
{\small Capability · Superseded historical identity — direct correspondence · superseded}\par
\href{../../../../continuum/strategy/capability/ab0415e2ea0a4da4ad3dc46997d48bc3.md}{Historical canonical definition and provenance}.\par
\section*{R.2.130 [CT2] Clinical Qualification Administration — Historical identity [superseded]}
\phantomsection
\label{navigation:R.2.130}
\addcontentsline{toc}{section}{R.2.130 [CT2] Clinical Qualification Administration — Historical identity [superseded]}

{\small Governed successor (direct correspondence): II.1.3.6.8 \hyperref[element:0511a705d74a4f48a0ce92590ae764fa]{[CT2] Clinical Qualification Administration}; Canonical ID: \texttt{\seqsplit{0511a705d74a4f48a0ce92590ae764fa}}. Canonical · current.}\par
\label{element:5fea7aaac4384245b1f08a3d8a076bd9}
\index{[CT2] Clinical Qualification Administration (R.2.130)}
{\small Capability · Superseded historical identity — direct correspondence · superseded}\par
\href{../../../../continuum/strategy/capability/5fea7aaac4384245b1f08a3d8a076bd9.md}{Historical canonical definition and provenance}.\par
\section*{R.2.131 [CT3] Clinical Roster Capture \& Communication — Historical identity [superseded]}
\phantomsection
\label{navigation:R.2.131}
\addcontentsline{toc}{section}{R.2.131 [CT3] Clinical Roster Capture \& Communication — Historical identity [superseded]}

{\small Governed successor (direct correspondence): II.1.3.6.9.1 \hyperref[element:8243cba81fb6449781d10823ae194237]{[CT3] Clinical Roster Capture \& Communication}; Canonical ID: \texttt{\seqsplit{8243cba81fb6449781d10823ae194237}}. Proposed · current.}\par
\label{element:2054d57bb33a4361b3ef0824545a1784}
\index{[CT3] Clinical Roster Capture \& Communication (R.2.131)}
{\small Capability · Superseded historical identity — direct correspondence · superseded}\par
\href{../../../../continuum/strategy/capability/2054d57bb33a4361b3ef0824545a1784.md}{Historical canonical definition and provenance}.\par
\section*{R.2.132 [CT2] Clinical Rostering — Historical identity [superseded]}
\phantomsection
\label{navigation:R.2.132}
\addcontentsline{toc}{section}{R.2.132 [CT2] Clinical Rostering — Historical identity [superseded]}

{\small Governed successor (direct correspondence): II.1.3.6.9 \hyperref[element:192ca88a289544dcbf6adaf643f39969]{[CT2] Clinical Rostering}; Canonical ID: \texttt{\seqsplit{192ca88a289544dcbf6adaf643f39969}}. Canonical · current.}\par
\label{element:b6527516a69c499d8cc9c8fe881169fe}
\index{[CT2] Clinical Rostering (R.2.132)}
{\small Capability · Superseded historical identity — direct correspondence · superseded}\par
\href{../../../../continuum/strategy/capability/b6527516a69c499d8cc9c8fe881169fe.md}{Historical canonical definition and provenance}.\par
\section*{R.2.133 [CT3] Clinical Escalation Path Management — Historical identity [superseded]}
\phantomsection
\label{navigation:R.2.133}
\addcontentsline{toc}{section}{R.2.133 [CT3] Clinical Escalation Path Management — Historical identity [superseded]}

{\small Governed successor (direct correspondence): II.1.3.6.10.1 \hyperref[element:cca86a32ba8c4dda84caeaa248e9f444]{[CT3] Clinical Escalation Path Management}; Canonical ID: \texttt{\seqsplit{cca86a32ba8c4dda84caeaa248e9f444}}. Proposed · current.}\par
\label{element:759a2e52172345e8ad1bc96d7e5ddad3}
\index{[CT3] Clinical Escalation Path Management (R.2.133)}
{\small Capability · Superseded historical identity — direct correspondence · superseded}\par
\href{../../../../continuum/strategy/capability/759a2e52172345e8ad1bc96d7e5ddad3.md}{Historical canonical definition and provenance}.\par
\section*{R.2.134 [CT2] Clinical Supervision — Historical identity [superseded]}
\phantomsection
\label{navigation:R.2.134}
\addcontentsline{toc}{section}{R.2.134 [CT2] Clinical Supervision — Historical identity [superseded]}

{\small Governed successor (direct correspondence): II.1.3.6.10 \hyperref[element:f697b1ac79b544749315f31819b17969]{[CT2] Clinical Supervision}; Canonical ID: \texttt{\seqsplit{f697b1ac79b544749315f31819b17969}}. Canonical · current.}\par
\label{element:50c838981f3946d78e8f78d18e31c70c}
\index{[CT2] Clinical Supervision (R.2.134)}
{\small Capability · Superseded historical identity — direct correspondence · superseded}\par
\href{../../../../continuum/strategy/capability/50c838981f3946d78e8f78d18e31c70c.md}{Historical canonical definition and provenance}.\par
\section*{R.2.135 [CT2] Clinical Training \& Education — Historical identity [superseded]}
\phantomsection
\label{navigation:R.2.135}
\addcontentsline{toc}{section}{R.2.135 [CT2] Clinical Training \& Education — Historical identity [superseded]}

{\small Governed successor (direct correspondence): II.1.3.6.11 \hyperref[element:d016f35ffbc345d99d7281e968c14ddb]{[CT2] Clinical Training \& Education}; Canonical ID: \texttt{\seqsplit{d016f35ffbc345d99d7281e968c14ddb}}. Canonical · current.}\par
\label{element:df8bff7df9c1469da16cac6a5eb2bf76}
\index{[CT2] Clinical Training \& Education (R.2.135)}
{\small Capability · Superseded historical identity — direct correspondence · superseded}\par
\href{../../../../continuum/strategy/capability/df8bff7df9c1469da16cac6a5eb2bf76.md}{Historical canonical definition and provenance}.\par
\section*{R.2.136 [CT2] Clinical Workforce Capability Monitoring — Historical identity [superseded]}
\phantomsection
\label{navigation:R.2.136}
\addcontentsline{toc}{section}{R.2.136 [CT2] Clinical Workforce Capability Monitoring — Historical identity [superseded]}

{\small Governed successor (direct correspondence): II.1.3.6.12 \hyperref[element:367b3e5643f44fa785f093db8bb9a6e8]{[CT2] Clinical Workforce Capability Monitoring}; Canonical ID: \texttt{\seqsplit{367b3e5643f44fa785f093db8bb9a6e8}}. Canonical · current.}\par
\label{element:24c015676c254e2e82e8e4287c295179}
\index{[CT2] Clinical Workforce Capability Monitoring (R.2.136)}
{\small Capability · Superseded historical identity — direct correspondence · superseded}\par
\href{../../../../continuum/strategy/capability/24c015676c254e2e82e8e4287c295179.md}{Historical canonical definition and provenance}.\par
\section*{R.2.137 [CT3] Clinical Service Credential Definition Management — Historical identity [superseded]}
\phantomsection
\label{navigation:R.2.137}
\addcontentsline{toc}{section}{R.2.137 [CT3] Clinical Service Credential Definition Management — Historical identity [superseded]}

{\small Governed successor (direct correspondence): II.1.3.6.13.1 \hyperref[element:2e5c97f51e7e48d1b865ac416dd8099d]{[CT3] Clinical Service Credential Definition Management}; Canonical ID: \texttt{\seqsplit{2e5c97f51e7e48d1b865ac416dd8099d}}. Proposed · current.}\par
\label{element:0115ed50867d468a95c0fafa1c5d79e2}
\index{[CT3] Clinical Service Credential Definition Management (R.2.137)}
{\small Capability · Superseded historical identity — direct correspondence · superseded}\par
\href{../../../../continuum/strategy/capability/0115ed50867d468a95c0fafa1c5d79e2.md}{Historical canonical definition and provenance}.\par
\section*{R.2.138 [CT2] Credential Management — Historical identity [superseded]}
\phantomsection
\label{navigation:R.2.138}
\addcontentsline{toc}{section}{R.2.138 [CT2] Credential Management — Historical identity [superseded]}

{\small Governed successor (direct correspondence): II.1.3.6.13 \hyperref[element:21eed4fc8c924496b16c6ddaa32f3d01]{[CT2] Credential Management}; Canonical ID: \texttt{\seqsplit{21eed4fc8c924496b16c6ddaa32f3d01}}. Proposed · current.}\par
\label{element:bab3d0ab6eb748978b962ea20fb31640}
\index{[CT2] Credential Management (R.2.138)}
{\small Capability · Superseded historical identity — direct correspondence · superseded}\par
\href{../../../../continuum/strategy/capability/bab3d0ab6eb748978b962ea20fb31640.md}{Historical canonical definition and provenance}.\par
\section*{R.2.139 [CT3] Discharge Progression \& Coordination — Historical identity [superseded]}
\phantomsection
\label{navigation:R.2.139}
\addcontentsline{toc}{section}{R.2.139 [CT3] Discharge Progression \& Coordination — Historical identity [superseded]}

{\small Governed successor (direct correspondence): II.1.3.6.14.1 \hyperref[element:d6048993146b45639bab27860f94f34f]{[CT3] Discharge Progression \& Coordination}; Canonical ID: \texttt{\seqsplit{d6048993146b45639bab27860f94f34f}}. Proposed · current.}\par
\label{element:68344272ca6049648a3ddae20c150f37}
\index{[CT3] Discharge Progression \& Coordination (R.2.139)}
{\small Capability · Superseded historical identity — direct correspondence · superseded}\par
\href{../../../../continuum/strategy/capability/68344272ca6049648a3ddae20c150f37.md}{Historical canonical definition and provenance}.\par
\section*{R.2.140 [CT2] Discharge Management — Historical identity [superseded]}
\phantomsection
\label{navigation:R.2.140}
\addcontentsline{toc}{section}{R.2.140 [CT2] Discharge Management — Historical identity [superseded]}

{\small Governed successor (direct correspondence): II.1.3.6.14 \hyperref[element:ec5e8d6732764aa489c2dd7630fc8fef]{[CT2] Discharge Management}; Canonical ID: \texttt{\seqsplit{ec5e8d6732764aa489c2dd7630fc8fef}}. Canonical · current.}\par
\label{element:c57c0721ab0040ccb24e88e5e92e8cf8}
\index{[CT2] Discharge Management (R.2.140)}
{\small Capability · Superseded historical identity — direct correspondence · superseded}\par
\href{../../../../continuum/strategy/capability/c57c0721ab0040ccb24e88e5e92e8cf8.md}{Historical canonical definition and provenance}.\par
\section*{R.2.141 [CT3] Emergency Department Operational Coordination — Historical identity [superseded]}
\phantomsection
\label{navigation:R.2.141}
\addcontentsline{toc}{section}{R.2.141 [CT3] Emergency Department Operational Coordination — Historical identity [superseded]}

{\small Governed successor (direct correspondence): II.1.3.6.15.2 \hyperref[element:8f0b0f3d9e8d44a7827b0835870b4a4f]{[CT3] Emergency Department Operational Coordination}; Canonical ID: \texttt{\seqsplit{8f0b0f3d9e8d44a7827b0835870b4a4f}}. Proposed · current.}\par
\label{element:15b6d2518e60444098b7c546ca31c49b}
\index{[CT3] Emergency Department Operational Coordination (R.2.141)}
{\small Capability · Superseded historical identity — direct correspondence · superseded}\par
\href{../../../../continuum/strategy/capability/15b6d2518e60444098b7c546ca31c49b.md}{Historical canonical definition and provenance}.\par
\section*{R.2.142 [CT2] Emergency Department Operations — Historical identity [superseded]}
\phantomsection
\label{navigation:R.2.142}
\addcontentsline{toc}{section}{R.2.142 [CT2] Emergency Department Operations — Historical identity [superseded]}

{\small Governed successor (direct correspondence): II.1.3.6.15 \hyperref[element:24a4b5815faa4d369bc756a6e97d3e9e]{[CT2] Emergency Department Operations}; Canonical ID: \texttt{\seqsplit{24a4b5815faa4d369bc756a6e97d3e9e}}. Canonical · current.}\par
\label{element:8963b571e9cc40e49c5d738e31d2f97e}
\index{[CT2] Emergency Department Operations (R.2.142)}
{\small Capability · Superseded historical identity — direct correspondence · superseded}\par
\href{../../../../continuum/strategy/capability/8963b571e9cc40e49c5d738e31d2f97e.md}{Historical canonical definition and provenance}.\par
\section*{R.2.143 [CT3] Mobile Worker Dispatch \& Safety Coordination — Historical identity [superseded]}
\phantomsection
\label{navigation:R.2.143}
\addcontentsline{toc}{section}{R.2.143 [CT3] Mobile Worker Dispatch \& Safety Coordination — Historical identity [superseded]}

{\small Governed successor (direct correspondence): II.1.3.6.16.1 \hyperref[element:d72164668dfe448bbebe74e9e1a75233]{[CT3] Mobile Worker Dispatch \& Safety Coordination}; Canonical ID: \texttt{\seqsplit{d72164668dfe448bbebe74e9e1a75233}}. Proposed · current.}\par
\label{element:5aafd1f4bff046b2b532c20dfb191473}
\index{[CT3] Mobile Worker Dispatch \& Safety Coordination (R.2.143)}
{\small Capability · Superseded historical identity — direct correspondence · superseded}\par
\href{../../../../continuum/strategy/capability/5aafd1f4bff046b2b532c20dfb191473.md}{Historical canonical definition and provenance}.\par
\section*{R.2.144 [CT2] Mobile Staff Management — Historical identity [superseded]}
\phantomsection
\label{navigation:R.2.144}
\addcontentsline{toc}{section}{R.2.144 [CT2] Mobile Staff Management — Historical identity [superseded]}

{\small Governed successor (direct correspondence): II.1.3.6.16 \hyperref[element:79a4a48fddf74c69acf7e16404d06ea5]{[CT2] Mobile Staff Management}; Canonical ID: \texttt{\seqsplit{79a4a48fddf74c69acf7e16404d06ea5}}. Canonical · current.}\par
\label{element:32118856636f400a95fc28daacea61c8}
\index{[CT2] Mobile Staff Management (R.2.144)}
{\small Capability · Superseded historical identity — direct correspondence · superseded}\par
\href{../../../../continuum/strategy/capability/32118856636f400a95fc28daacea61c8.md}{Historical canonical definition and provenance}.\par
\section*{R.2.145 [CT3] On-Call Availability \& Contact Resolution — Historical identity [superseded]}
\phantomsection
\label{navigation:R.2.145}
\addcontentsline{toc}{section}{R.2.145 [CT3] On-Call Availability \& Contact Resolution — Historical identity [superseded]}

{\small Governed successor (direct correspondence): II.1.3.6.17.1 \hyperref[element:677a816c42414ea694791b4ee5a14285]{[CT3] On-Call Availability \& Contact Resolution}; Canonical ID: \texttt{\seqsplit{677a816c42414ea694791b4ee5a14285}}. Proposed · current.}\par
\label{element:a03075a46f75497f8c6ba95d9c9e27c4}
\index{[CT3] On-Call Availability \& Contact Resolution (R.2.145)}
{\small Capability · Superseded historical identity — direct correspondence · superseded}\par
\href{../../../../continuum/strategy/capability/a03075a46f75497f8c6ba95d9c9e27c4.md}{Historical canonical definition and provenance}.\par
\section*{R.2.146 [CT2] On-Call Management — Historical identity [superseded]}
\phantomsection
\label{navigation:R.2.146}
\addcontentsline{toc}{section}{R.2.146 [CT2] On-Call Management — Historical identity [superseded]}

{\small Governed successor (direct correspondence): II.1.3.6.17 \hyperref[element:b16cef8f2bfd4b92948ec75cbcfa7b9a]{[CT2] On-Call Management}; Canonical ID: \texttt{\seqsplit{b16cef8f2bfd4b92948ec75cbcfa7b9a}}. Canonical · current.}\par
\label{element:c8e6e66e8718493c9093f37d129f9ff2}
\index{[CT2] On-Call Management (R.2.146)}
{\small Capability · Superseded historical identity — direct correspondence · superseded}\par
\href{../../../../continuum/strategy/capability/c8e6e66e8718493c9093f37d129f9ff2.md}{Historical canonical definition and provenance}.\par
\section*{R.2.147 [CT3] Outpatient Flow Coordination — Historical identity [superseded]}
\phantomsection
\label{navigation:R.2.147}
\addcontentsline{toc}{section}{R.2.147 [CT3] Outpatient Flow Coordination — Historical identity [superseded]}

{\small Governed successor (direct correspondence): II.1.3.6.19.1 \hyperref[element:799ea278af014168b9b534e5fdadc6d1]{[CT3] Outpatient Flow Coordination}; Canonical ID: \texttt{\seqsplit{799ea278af014168b9b534e5fdadc6d1}}. Proposed · current.}\par
\label{element:e6c59c7b72fe44169a5fdfd08a5e91c4}
\index{[CT3] Outpatient Flow Coordination (R.2.147)}
{\small Capability · Superseded historical identity — direct correspondence · superseded}\par
\href{../../../../continuum/strategy/capability/e6c59c7b72fe44169a5fdfd08a5e91c4.md}{Historical canonical definition and provenance}.\par
\section*{R.2.148 [CT2] Outpatient Operations — Historical identity [superseded]}
\phantomsection
\label{navigation:R.2.148}
\addcontentsline{toc}{section}{R.2.148 [CT2] Outpatient Operations — Historical identity [superseded]}

{\small Governed successor (direct correspondence): II.1.3.6.19 \hyperref[element:21399d52b9624ddfbeb786f1d72e40c3]{[CT2] Outpatient Operations}; Canonical ID: \texttt{\seqsplit{21399d52b9624ddfbeb786f1d72e40c3}}. Canonical · current.}\par
\label{element:3db41bc06ade4f9088eec7b8aedd7448}
\index{[CT2] Outpatient Operations (R.2.148)}
{\small Capability · Superseded historical identity — direct correspondence · superseded}\par
\href{../../../../continuum/strategy/capability/3db41bc06ade4f9088eec7b8aedd7448.md}{Historical canonical definition and provenance}.\par
\section*{R.2.149 [CT3] Patient Transfer Coordination — Historical identity [superseded]}
\phantomsection
\label{navigation:R.2.149}
\addcontentsline{toc}{section}{R.2.149 [CT3] Patient Transfer Coordination — Historical identity [superseded]}

{\small Governed successor (direct correspondence): II.1.3.6.20.1 \hyperref[element:c7c7aa8e34df4568ae4b91cdde380a11]{[CT3] Patient Transfer Coordination}; Canonical ID: \texttt{\seqsplit{c7c7aa8e34df4568ae4b91cdde380a11}}. Proposed · current.}\par
\label{element:b5a76d4877424b2ea9b01fb5fa22772c}
\index{[CT3] Patient Transfer Coordination (R.2.149)}
{\small Capability · Superseded historical identity — direct correspondence · superseded}\par
\href{../../../../continuum/strategy/capability/b5a76d4877424b2ea9b01fb5fa22772c.md}{Historical canonical definition and provenance}.\par
\section*{R.2.150 [CT2] Patient Transport — Historical identity [superseded]}
\phantomsection
\label{navigation:R.2.150}
\addcontentsline{toc}{section}{R.2.150 [CT2] Patient Transport — Historical identity [superseded]}

{\small Governed successor (direct correspondence): II.1.3.6.20 \hyperref[element:6041d7c0cc5d4519aaea584bb4e6d365]{[CT2] Patient Transport}; Canonical ID: \texttt{\seqsplit{6041d7c0cc5d4519aaea584bb4e6d365}}. Canonical · current.}\par
\label{element:e9c1a897884f4d758f5654b824364030}
\index{[CT2] Patient Transport (R.2.150)}
{\small Capability · Superseded historical identity — direct correspondence · superseded}\par
\href{../../../../continuum/strategy/capability/e9c1a897884f4d758f5654b824364030.md}{Historical canonical definition and provenance}.\par
\section*{R.2.151 [CT3] Service Demand \& Capacity Coordination — Historical identity [superseded]}
\phantomsection
\label{navigation:R.2.151}
\addcontentsline{toc}{section}{R.2.151 [CT3] Service Demand \& Capacity Coordination — Historical identity [superseded]}

{\small Governed successor (direct correspondence): II.1.3.6.21.1 \hyperref[element:93b5548ce8d04fcd999b97536c1eebf9]{[CT3] Service Demand \& Capacity Coordination}; Canonical ID: \texttt{\seqsplit{93b5548ce8d04fcd999b97536c1eebf9}}. Proposed · current.}\par
\label{element:b57081e650f74ad488549fdaa5305bcb}
\index{[CT3] Service Demand \& Capacity Coordination (R.2.151)}
{\small Capability · Superseded historical identity — direct correspondence · superseded}\par
\href{../../../../continuum/strategy/capability/b57081e650f74ad488549fdaa5305bcb.md}{Historical canonical definition and provenance}.\par
\section*{R.2.152 [CT2] Service Capacity Management — Historical identity [superseded]}
\phantomsection
\label{navigation:R.2.152}
\addcontentsline{toc}{section}{R.2.152 [CT2] Service Capacity Management — Historical identity [superseded]}

{\small Governed successor (direct correspondence): II.1.3.6.21 \hyperref[element:d86105bf9c86445394088f65c895d37f]{[CT2] Service Capacity Management}; Canonical ID: \texttt{\seqsplit{d86105bf9c86445394088f65c895d37f}}. Canonical · current.}\par
\label{element:ac205f374c9d41c3b6d18a6b7c356152}
\index{[CT2] Service Capacity Management (R.2.152)}
{\small Capability · Superseded historical identity — direct correspondence · superseded}\par
\href{../../../../continuum/strategy/capability/ac205f374c9d41c3b6d18a6b7c356152.md}{Historical canonical definition and provenance}.\par
\section*{R.2.153 [CT2] Service Establishment Support — Historical identity [superseded]}
\phantomsection
\label{navigation:R.2.153}
\addcontentsline{toc}{section}{R.2.153 [CT2] Service Establishment Support — Historical identity [superseded]}

{\small Governed successor (direct correspondence): II.1.3.6.22 \hyperref[element:20e67a41e7f44398b0f0cf3d379d8b9c]{[CT2] Service Establishment Support}; Canonical ID: \texttt{\seqsplit{20e67a41e7f44398b0f0cf3d379d8b9c}}. Proposed · current.}\par
\label{element:8a63b08f34fa42c5bda2a3e12b9ecd55}
\index{[CT2] Service Establishment Support (R.2.153)}
{\small Capability · Superseded historical identity — direct correspondence · superseded}\par
\href{../../../../continuum/strategy/capability/8a63b08f34fa42c5bda2a3e12b9ecd55.md}{Historical canonical definition and provenance}.\par
\section*{R.2.154 [CT3] Perioperative Logistics Coordination — Historical identity [superseded]}
\phantomsection
\label{navigation:R.2.154}
\addcontentsline{toc}{section}{R.2.154 [CT3] Perioperative Logistics Coordination — Historical identity [superseded]}

{\small Governed successor (direct correspondence): II.1.3.6.23.1 \hyperref[element:6b7c08b60b5b48519bfe4b6f993dd6fb]{[CT3] Perioperative Logistics Coordination}; Canonical ID: \texttt{\seqsplit{6b7c08b60b5b48519bfe4b6f993dd6fb}}. Proposed · current.}\par
\label{element:e7813a06b9df462fa135c299261f192e}
\index{[CT3] Perioperative Logistics Coordination (R.2.154)}
{\small Capability · Superseded historical identity — direct correspondence · superseded}\par
\href{../../../../continuum/strategy/capability/e7813a06b9df462fa135c299261f192e.md}{Historical canonical definition and provenance}.\par
\section*{R.2.155 [CT2] Theatre Operations — Historical identity [superseded]}
\phantomsection
\label{navigation:R.2.155}
\addcontentsline{toc}{section}{R.2.155 [CT2] Theatre Operations — Historical identity [superseded]}

{\small Governed successor (direct correspondence): II.1.3.6.23 \hyperref[element:4bf7a8260e1444fdac44e97c3557b533]{[CT2] Theatre Operations}; Canonical ID: \texttt{\seqsplit{4bf7a8260e1444fdac44e97c3557b533}}. Canonical · current.}\par
\label{element:9ac72fea821c45e2b18632ae9a9867dd}
\index{[CT2] Theatre Operations (R.2.155)}
{\small Capability · Superseded historical identity — direct correspondence · superseded}\par
\href{../../../../continuum/strategy/capability/9ac72fea821c45e2b18632ae9a9867dd.md}{Historical canonical definition and provenance}.\par
\section*{R.2.156 [CT3] Ward Logistics Coordination — Historical identity [superseded]}
\phantomsection
\label{navigation:R.2.156}
\addcontentsline{toc}{section}{R.2.156 [CT3] Ward Logistics Coordination — Historical identity [superseded]}

{\small Governed successor (direct correspondence): II.1.3.6.24.2 \hyperref[element:3249647f7b9846c899f982ab339c81f9]{[CT3] Ward Logistics Coordination}; Canonical ID: \texttt{\seqsplit{3249647f7b9846c899f982ab339c81f9}}. Proposed · current.}\par
\label{element:8a9dcc3b5f0143458ec71bc188a56020}
\index{[CT3] Ward Logistics Coordination (R.2.156)}
{\small Capability · Superseded historical identity — direct correspondence · superseded}\par
\href{../../../../continuum/strategy/capability/8a9dcc3b5f0143458ec71bc188a56020.md}{Historical canonical definition and provenance}.\par
\section*{R.2.157 [CT2] Ward Operations — Historical identity [superseded]}
\phantomsection
\label{navigation:R.2.157}
\addcontentsline{toc}{section}{R.2.157 [CT2] Ward Operations — Historical identity [superseded]}

{\small Governed successor (direct correspondence): II.1.3.6.24 \hyperref[element:f996ecca728f4942a88b034aca818a99]{[CT2] Ward Operations}; Canonical ID: \texttt{\seqsplit{f996ecca728f4942a88b034aca818a99}}. Canonical · current.}\par
\label{element:bf3d2e3581f24beda20a6a3a3348460e}
\index{[CT2] Ward Operations (R.2.157)}
{\small Capability · Superseded historical identity — direct correspondence · superseded}\par
\href{../../../../continuum/strategy/capability/bf3d2e3581f24beda20a6a3a3348460e.md}{Historical canonical definition and provenance}.\par
\section*{R.2.158 [CT3] Operational Work Assignment \& Dispatch Coordination — Historical identity [superseded]}
\phantomsection
\label{navigation:R.2.158}
\addcontentsline{toc}{section}{R.2.158 [CT3] Operational Work Assignment \& Dispatch Coordination — Historical identity [superseded]}

{\small Governed successor (direct correspondence): II.1.3.6.25.1 \hyperref[element:3b58935cfa974176989c3c75d76a00fe]{[CT3] Operational Work Assignment \& Dispatch Coordination}; Canonical ID: \texttt{\seqsplit{3b58935cfa974176989c3c75d76a00fe}}. Proposed · current.}\par
\label{element:b42afda9e876463ba56fa8f73fe7a32a}
\index{[CT3] Operational Work Assignment \& Dispatch Coordination (R.2.158)}
{\small Capability · Superseded historical identity — direct correspondence · superseded}\par
\href{../../../../continuum/strategy/capability/b42afda9e876463ba56fa8f73fe7a32a.md}{Historical canonical definition and provenance}.\par
\section*{R.2.159 [CT2] Work Allocation \& Dispatch — Historical identity [superseded]}
\phantomsection
\label{navigation:R.2.159}
\addcontentsline{toc}{section}{R.2.159 [CT2] Work Allocation \& Dispatch — Historical identity [superseded]}

{\small Governed successor (direct correspondence): II.1.3.6.25 \hyperref[element:152a5252952c405994e9b171234efcb0]{[CT2] Work Allocation \& Dispatch}; Canonical ID: \texttt{\seqsplit{152a5252952c405994e9b171234efcb0}}. Canonical · current.}\par
\label{element:f8c13b6f74b840cc96fbeada5865455a}
\index{[CT2] Work Allocation \& Dispatch (R.2.159)}
{\small Capability · Superseded historical identity — direct correspondence · superseded}\par
\href{../../../../continuum/strategy/capability/f8c13b6f74b840cc96fbeada5865455a.md}{Historical canonical definition and provenance}.\par
\section*{R.2.160 [CT1] Health Service Operations Enablement — Historical identity [superseded]}
\phantomsection
\label{navigation:R.2.160}
\addcontentsline{toc}{section}{R.2.160 [CT1] Health Service Operations Enablement — Historical identity [superseded]}

{\small Governed successor (direct correspondence): II.1.3.6 \hyperref[element:b13ee2c5a4ef4fd6be8a44795696dfc2]{[CT1] Health Service Operations Enablement}; Canonical ID: \texttt{\seqsplit{b13ee2c5a4ef4fd6be8a44795696dfc2}}. Proposed · current.}\par
\label{element:a98424298dac46d8b2ddd109d928270a}
\index{[CT1] Health Service Operations Enablement (R.2.160)}
{\small Capability · Superseded historical identity — direct correspondence · superseded}\par
\href{../../../../continuum/strategy/capability/a98424298dac46d8b2ddd109d928270a.md}{Historical canonical definition and provenance}.\par
\section*{R.2.161 [CT3] Assurance Service Evaluation — Historical identity [superseded]}
\phantomsection
\label{navigation:R.2.161}
\addcontentsline{toc}{section}{R.2.161 [CT3] Assurance Service Evaluation — Historical identity [superseded]}

{\small Governed successor (direct correspondence): II.1.3.8.1.1 \hyperref[element:039343518db542d48cde6afab84a3e1e]{[CT3] Assurance Service Evaluation}; Canonical ID: \texttt{\seqsplit{039343518db542d48cde6afab84a3e1e}}. Proposed · current.}\par
\label{element:be34b54314c6444f8418d645f6c83b6d}
\index{[CT3] Assurance Service Evaluation (R.2.161)}
{\small Capability · Superseded historical identity — direct correspondence · superseded}\par
\href{../../../../continuum/strategy/capability/be34b54314c6444f8418d645f6c83b6d.md}{Historical canonical definition and provenance}.\par
\section*{R.2.162 [CT2] Assurance Service Review — Historical identity [superseded]}
\phantomsection
\label{navigation:R.2.162}
\addcontentsline{toc}{section}{R.2.162 [CT2] Assurance Service Review — Historical identity [superseded]}

{\small Governed successor (direct correspondence): II.1.3.8.1 \hyperref[element:b1a791458d644d4794895e4fa42ec83c]{[CT2] Assurance Service Review}; Canonical ID: \texttt{\seqsplit{b1a791458d644d4794895e4fa42ec83c}}. Proposed · current.}\par
\label{element:b345cb12b45746409302899514126170}
\index{[CT2] Assurance Service Review (R.2.162)}
{\small Capability · Superseded historical identity — direct correspondence · superseded}\par
\href{../../../../continuum/strategy/capability/b345cb12b45746409302899514126170.md}{Historical canonical definition and provenance}.\par
\section*{R.2.163 [CT3] Assurance Criteria Administration — Historical identity [superseded]}
\phantomsection
\label{navigation:R.2.163}
\addcontentsline{toc}{section}{R.2.163 [CT3] Assurance Criteria Administration — Historical identity [superseded]}

{\small Governed successor (direct correspondence): II.1.3.8.2.1 \hyperref[element:1bfd71f24321428c8871b587d413f47e]{[CT3] Assurance Criteria Administration}; Canonical ID: \texttt{\seqsplit{1bfd71f24321428c8871b587d413f47e}}. Canonical · current.}\par
\label{element:6dab08ab1e3249d7b314843ca9ed0149}
\index{[CT3] Assurance Criteria Administration (R.2.163)}
{\small Capability · Superseded historical identity — direct correspondence · superseded}\par
\href{../../../../continuum/strategy/capability/6dab08ab1e3249d7b314843ca9ed0149.md}{Historical canonical definition and provenance}.\par
\section*{R.2.164 [CT3] Assurance Design — Historical identity [superseded]}
\phantomsection
\label{navigation:R.2.164}
\addcontentsline{toc}{section}{R.2.164 [CT3] Assurance Design — Historical identity [superseded]}

{\small Governed successor (direct correspondence): II.1.3.8.2.2 \hyperref[element:71c8a612b88f4ecba049e94d506bd0c2]{[CT3] Assurance Design}; Canonical ID: \texttt{\seqsplit{71c8a612b88f4ecba049e94d506bd0c2}}. Canonical · current.}\par
\label{element:e1b8db7c95fa4e7ea296351a03532e24}
\index{[CT3] Assurance Design (R.2.164)}
{\small Capability · Superseded historical identity — direct correspondence · superseded}\par
\href{../../../../continuum/strategy/capability/e1b8db7c95fa4e7ea296351a03532e24.md}{Historical canonical definition and provenance}.\par
\section*{R.2.165 [CT2] Assurance Specification — Historical identity [superseded]}
\phantomsection
\label{navigation:R.2.165}
\addcontentsline{toc}{section}{R.2.165 [CT2] Assurance Specification — Historical identity [superseded]}

{\small Governed successor (direct correspondence): II.1.3.8.2 \hyperref[element:b81775d89d774cf8ba61731d13c67a54]{[CT2] Assurance Specification}; Canonical ID: \texttt{\seqsplit{b81775d89d774cf8ba61731d13c67a54}}. Proposed · current.}\par
\label{element:8006ee0dd95f4d7fb1f408fa09b6ce19}
\index{[CT2] Assurance Specification (R.2.165)}
{\small Capability · Superseded historical identity — direct correspondence · superseded}\par
\href{../../../../continuum/strategy/capability/8006ee0dd95f4d7fb1f408fa09b6ce19.md}{Historical canonical definition and provenance}.\par
\section*{R.2.166 [CT1] Independent Assurance Enablement — Historical identity [superseded]}
\phantomsection
\label{navigation:R.2.166}
\addcontentsline{toc}{section}{R.2.166 [CT1] Independent Assurance Enablement — Historical identity [superseded]}

{\small Governed successor (direct correspondence): II.1.3.8 \hyperref[element:616682b612d244429e656b040e002d56]{[CT1] Independent Assurance Enablement}; Canonical ID: \texttt{\seqsplit{616682b612d244429e656b040e002d56}}. Proposed · current.}\par
\label{element:88bf57b7a92c4c1da1b5f0a67d3ebd82}
\index{[CT1] Independent Assurance Enablement (R.2.166)}
{\small Capability · Superseded historical identity — direct correspondence · superseded}\par
\href{../../../../continuum/strategy/capability/88bf57b7a92c4c1da1b5f0a67d3ebd82.md}{Historical canonical definition and provenance}.\par
\section*{R.2.167 [CT3] Operational \& Clinical Calendar Coordination — Historical identity [superseded]}
\phantomsection
\label{navigation:R.2.167}
\addcontentsline{toc}{section}{R.2.167 [CT3] Operational \& Clinical Calendar Coordination — Historical identity [superseded]}

{\small Governed successor (direct correspondence): II.1.3.9.1.1 \hyperref[element:f6d8f1310ba543f69e8b57bf05603db3]{[CT3] Operational \& Clinical Calendar Coordination}; Canonical ID: \texttt{\seqsplit{f6d8f1310ba543f69e8b57bf05603db3}}. Proposed · current.}\par
\label{element:c0f6e5f7f5a94d1580fa6e0e8b82c93a}
\index{[CT3] Operational \& Clinical Calendar Coordination (R.2.167)}
{\small Capability · Superseded historical identity — direct correspondence · superseded}\par
\href{../../../../continuum/strategy/capability/c0f6e5f7f5a94d1580fa6e0e8b82c93a.md}{Historical canonical definition and provenance}.\par
\section*{R.2.168 [CT2] Calendar Management — Historical identity [superseded]}
\phantomsection
\label{navigation:R.2.168}
\addcontentsline{toc}{section}{R.2.168 [CT2] Calendar Management — Historical identity [superseded]}

{\small Governed successor (direct correspondence): II.1.3.9.1 \hyperref[element:7669036f7e6749219412fc64df5f7064]{[CT2] Calendar Management}; Canonical ID: \texttt{\seqsplit{7669036f7e6749219412fc64df5f7064}}. Canonical · current.}\par
\label{element:8b42aac6083441caa78b59f5e62eebd9}
\index{[CT2] Calendar Management (R.2.168)}
{\small Capability · Superseded historical identity — direct correspondence · superseded}\par
\href{../../../../continuum/strategy/capability/8b42aac6083441caa78b59f5e62eebd9.md}{Historical canonical definition and provenance}.\par
\section*{R.2.169 [CT3] Clinical Context Collaboration — Historical identity [superseded]}
\phantomsection
\label{navigation:R.2.169}
\addcontentsline{toc}{section}{R.2.169 [CT3] Clinical Context Collaboration — Historical identity [superseded]}

{\small Governed successor (direct correspondence): II.1.3.9.2.1 \hyperref[element:f11617eed6ff4d6fa237644ec95b28b2]{[CT3] Clinical Context Collaboration}; Canonical ID: \texttt{\seqsplit{f11617eed6ff4d6fa237644ec95b28b2}}. Proposed · current.}\par
\label{element:750ef0d8e1504a8e9a1fde9bf13a6d0c}
\index{[CT3] Clinical Context Collaboration (R.2.169)}
{\small Capability · Superseded historical identity — direct correspondence · superseded}\par
\href{../../../../continuum/strategy/capability/750ef0d8e1504a8e9a1fde9bf13a6d0c.md}{Historical canonical definition and provenance}.\par
\section*{R.2.170 [CT2] Clinical Collaboration — Historical identity [superseded]}
\phantomsection
\label{navigation:R.2.170}
\addcontentsline{toc}{section}{R.2.170 [CT2] Clinical Collaboration — Historical identity [superseded]}

{\small Governed successor (direct correspondence): II.1.3.9.2 \hyperref[element:e6429225463a44d39f82e490e254e5d4]{[CT2] Clinical Collaboration}; Canonical ID: \texttt{\seqsplit{e6429225463a44d39f82e490e254e5d4}}. Canonical · current.}\par
\label{element:c5adc1955b6046d8b920b0176cebc9de}
\index{[CT2] Clinical Collaboration (R.2.170)}
{\small Capability · Superseded historical identity — direct correspondence · superseded}\par
\href{../../../../continuum/strategy/capability/c5adc1955b6046d8b920b0176cebc9de.md}{Historical canonical definition and provenance}.\par
\section*{R.2.171 [CT3] Clinical Knowledge Artefact Distribution Management — Historical identity [superseded]}
\phantomsection
\label{navigation:R.2.171}
\addcontentsline{toc}{section}{R.2.171 [CT3] Clinical Knowledge Artefact Distribution Management — Historical identity [superseded]}

{\small Governed successor (direct correspondence): II.1.3.9.4.1 \hyperref[element:cd29a9eb798949bc8bc81bb5d434cabf]{[CT3] Clinical Knowledge Artefact Distribution Management}; Canonical ID: \texttt{\seqsplit{cd29a9eb798949bc8bc81bb5d434cabf}}. Proposed · current.}\par
\label{element:d81e0d6d76014066a682e3d9e03e5f8a}
\index{[CT3] Clinical Knowledge Artefact Distribution Management (R.2.171)}
{\small Capability · Superseded historical identity — direct correspondence · superseded}\par
\href{../../../../continuum/strategy/capability/d81e0d6d76014066a682e3d9e03e5f8a.md}{Historical canonical definition and provenance}.\par
\section*{R.2.172 [CT3] Clinical Knowledge Request Administration — Historical identity [superseded]}
\phantomsection
\label{navigation:R.2.172}
\addcontentsline{toc}{section}{R.2.172 [CT3] Clinical Knowledge Request Administration — Historical identity [superseded]}

{\small Governed successor (direct correspondence): II.1.3.9.4.2 \hyperref[element:543c53483a2249d78044c6635b7b7c48]{[CT3] Clinical Knowledge Request Administration}; Canonical ID: \texttt{\seqsplit{543c53483a2249d78044c6635b7b7c48}}. Proposed · current.}\par
\label{element:eb370c01fbdb4d22956d97e9ecd74b57}
\index{[CT3] Clinical Knowledge Request Administration (R.2.172)}
{\small Capability · Superseded historical identity — direct correspondence · superseded}\par
\href{../../../../continuum/strategy/capability/eb370c01fbdb4d22956d97e9ecd74b57.md}{Historical canonical definition and provenance}.\par
\section*{R.2.173 [CT3] Governed Reference Dataset Administration — Historical identity [superseded]}
\phantomsection
\label{navigation:R.2.173}
\addcontentsline{toc}{section}{R.2.173 [CT3] Governed Reference Dataset Administration — Historical identity [superseded]}

{\small Governed successor (direct correspondence): II.1.3.9.4.3 \hyperref[element:682b8684ad9e49d5a5c348a148081607]{[CT3] Governed Reference Dataset Administration}; Canonical ID: \texttt{\seqsplit{682b8684ad9e49d5a5c348a148081607}}. Proposed · current.}\par
\label{element:15d972931ed6447797c30a5ddb71737e}
\index{[CT3] Governed Reference Dataset Administration (R.2.173)}
{\small Capability · Superseded historical identity — direct correspondence · superseded}\par
\href{../../../../continuum/strategy/capability/15d972931ed6447797c30a5ddb71737e.md}{Historical canonical definition and provenance}.\par
\section*{R.2.174 [CT2] Clinical Knowledge Services — Historical identity [superseded]}
\phantomsection
\label{navigation:R.2.174}
\addcontentsline{toc}{section}{R.2.174 [CT2] Clinical Knowledge Services — Historical identity [superseded]}

{\small Governed successor (direct correspondence): II.1.3.9.4 \hyperref[element:bbd2bd21859f40e39963bdbe7a851686]{[CT2] Clinical Knowledge Services}; Canonical ID: \texttt{\seqsplit{bbd2bd21859f40e39963bdbe7a851686}}. Canonical · current.}\par
\label{element:5b8478b08c3f42c0a3480edb8f0b9ea7}
\index{[CT2] Clinical Knowledge Services (R.2.174)}
{\small Capability · Superseded historical identity — direct correspondence · superseded}\par
\href{../../../../continuum/strategy/capability/5b8478b08c3f42c0a3480edb8f0b9ea7.md}{Historical canonical definition and provenance}.\par
\section*{R.2.175 [CT3] Operational Collaboration — Historical identity [superseded]}
\phantomsection
\label{navigation:R.2.175}
\addcontentsline{toc}{section}{R.2.175 [CT3] Operational Collaboration — Historical identity [superseded]}

{\small Governed successor (direct correspondence): II.1.3.9.5.1 \hyperref[element:631fdb8b47084008ac9b5de28b26972b]{[CT3] Operational Collaboration}; Canonical ID: \texttt{\seqsplit{631fdb8b47084008ac9b5de28b26972b}}. Proposed · current.}\par
\label{element:99e9d9763ec94140b7d7642f4b5966f6}
\index{[CT3] Operational Collaboration (R.2.175)}
{\small Capability · Superseded historical identity — direct correspondence · superseded}\par
\href{../../../../continuum/strategy/capability/99e9d9763ec94140b7d7642f4b5966f6.md}{Historical canonical definition and provenance}.\par
\section*{R.2.176 [CT2] General Collaboration — Historical identity [superseded]}
\phantomsection
\label{navigation:R.2.176}
\addcontentsline{toc}{section}{R.2.176 [CT2] General Collaboration — Historical identity [superseded]}

{\small Governed successor (direct correspondence): II.1.3.9.5 \hyperref[element:3bb740e6db284112b054905c0c6c80be]{[CT2] General Collaboration}; Canonical ID: \texttt{\seqsplit{3bb740e6db284112b054905c0c6c80be}}. Canonical · current.}\par
\label{element:4f58c82aca7a4b54aeb6e2920ce5c12a}
\index{[CT2] General Collaboration (R.2.176)}
{\small Capability · Superseded historical identity — direct correspondence · superseded}\par
\href{../../../../continuum/strategy/capability/4f58c82aca7a4b54aeb6e2920ce5c12a.md}{Historical canonical definition and provenance}.\par
\section*{R.2.177 [CT3] Governed Health Information Retrieval — Historical identity [superseded]}
\phantomsection
\label{navigation:R.2.177}
\addcontentsline{toc}{section}{R.2.177 [CT3] Governed Health Information Retrieval — Historical identity [superseded]}

{\small Governed successor (direct correspondence): II.1.3.9.6.1 \hyperref[element:76e74e46ca4d44a1a7bb116b358f231c]{[CT3] Governed Health Information Retrieval}; Canonical ID: \texttt{\seqsplit{76e74e46ca4d44a1a7bb116b358f231c}}. Proposed · current.}\par
\label{element:90b838d703dd472b8b6d74961f3a748c}
\index{[CT3] Governed Health Information Retrieval (R.2.177)}
{\small Capability · Superseded historical identity — direct correspondence · superseded}\par
\href{../../../../continuum/strategy/capability/90b838d703dd472b8b6d74961f3a748c.md}{Historical canonical definition and provenance}.\par
\section*{R.2.178 [CT2] Health Information Access — Historical identity [superseded]}
\phantomsection
\label{navigation:R.2.178}
\addcontentsline{toc}{section}{R.2.178 [CT2] Health Information Access — Historical identity [superseded]}

{\small Governed successor (direct correspondence): II.1.3.9.6 \hyperref[element:cc253aea3c5f489da0502df74dde1669]{[CT2] Health Information Access}; Canonical ID: \texttt{\seqsplit{cc253aea3c5f489da0502df74dde1669}}. Canonical · current.}\par
\label{element:5019ae9f2cdc45a1b2de97bc0f5a313d}
\index{[CT2] Health Information Access (R.2.178)}
{\small Capability · Superseded historical identity — direct correspondence · superseded}\par
\href{../../../../continuum/strategy/capability/5019ae9f2cdc45a1b2de97bc0f5a313d.md}{Historical canonical definition and provenance}.\par
\section*{R.2.179 [CT3] Patient Health Library Access — Historical identity [superseded]}
\phantomsection
\label{navigation:R.2.179}
\addcontentsline{toc}{section}{R.2.179 [CT3] Patient Health Library Access — Historical identity [superseded]}

{\small Governed successor (direct correspondence): II.1.3.9.7.1 \hyperref[element:808e5a3c3f6b496fa329962088f0bee0]{[CT3] Patient Health Library Access}; Canonical ID: \texttt{\seqsplit{808e5a3c3f6b496fa329962088f0bee0}}. Proposed · current.}\par
\label{element:b132224cfa6f435f8f2a66a3878d6df6}
\index{[CT3] Patient Health Library Access (R.2.179)}
{\small Capability · Superseded historical identity — direct correspondence · superseded}\par
\href{../../../../continuum/strategy/capability/b132224cfa6f435f8f2a66a3878d6df6.md}{Historical canonical definition and provenance}.\par
\section*{R.2.180 [CT3] Patient Notification Services — Historical identity [superseded]}
\phantomsection
\label{navigation:R.2.180}
\addcontentsline{toc}{section}{R.2.180 [CT3] Patient Notification Services — Historical identity [superseded]}

{\small Governed successor (direct correspondence): II.1.3.9.7.2 \hyperref[element:9d172cd511b742499c4e7c23f6630a35]{[CT3] Patient Notification Services}; Canonical ID: \texttt{\seqsplit{9d172cd511b742499c4e7c23f6630a35}}. Proposed · current.}\par
\label{element:35c5be472bf2433fbbde9646d87aefb3}
\index{[CT3] Patient Notification Services (R.2.180)}
{\small Capability · Superseded historical identity — direct correspondence · superseded}\par
\href{../../../../continuum/strategy/capability/35c5be472bf2433fbbde9646d87aefb3.md}{Historical canonical definition and provenance}.\par
\section*{R.2.181 [CT3] Patient Self-Service — Historical identity [superseded]}
\phantomsection
\label{navigation:R.2.181}
\addcontentsline{toc}{section}{R.2.181 [CT3] Patient Self-Service — Historical identity [superseded]}

{\small Governed successor (direct correspondence): II.1.3.9.7.3 \hyperref[element:d4956cab8b2b46c396899949c9486a1c]{[CT3] Patient Self-Service}; Canonical ID: \texttt{\seqsplit{d4956cab8b2b46c396899949c9486a1c}}. Proposed · current.}\par
\label{element:60f427e6b5764757b909624085a2af89}
\index{[CT3] Patient Self-Service (R.2.181)}
{\small Capability · Superseded historical identity — direct correspondence · superseded}\par
\href{../../../../continuum/strategy/capability/60f427e6b5764757b909624085a2af89.md}{Historical canonical definition and provenance}.\par
\section*{R.2.182 [CT3] Patient Task Communication Services — Historical identity [superseded]}
\phantomsection
\label{navigation:R.2.182}
\addcontentsline{toc}{section}{R.2.182 [CT3] Patient Task Communication Services — Historical identity [superseded]}

{\small Governed successor (direct correspondence): II.1.3.9.7.4 \hyperref[element:ba8ce1d54ace4bc9aa3521dd954f36cb]{[CT3] Patient Task Communication Services}; Canonical ID: \texttt{\seqsplit{ba8ce1d54ace4bc9aa3521dd954f36cb}}. Proposed · current.}\par
\label{element:0914f2bb267449e282e9ae0f7407d9fe}
\index{[CT3] Patient Task Communication Services (R.2.182)}
{\small Capability · Superseded historical identity — direct correspondence · superseded}\par
\href{../../../../continuum/strategy/capability/0914f2bb267449e282e9ae0f7407d9fe.md}{Historical canonical definition and provenance}.\par
\section*{R.2.183 [CT2] Health Information Communication — Historical identity [superseded]}
\phantomsection
\label{navigation:R.2.183}
\addcontentsline{toc}{section}{R.2.183 [CT2] Health Information Communication — Historical identity [superseded]}

{\small Governed successor (direct correspondence): II.1.3.9.7 \hyperref[element:7925605bf60f4b288e826020409a09d9]{[CT2] Health Information Communication}; Canonical ID: \texttt{\seqsplit{7925605bf60f4b288e826020409a09d9}}. Proposed · current.}\par
\label{element:bcefb378424842fcb6d703208c3f8773}
\index{[CT2] Health Information Communication (R.2.183)}
{\small Capability · Superseded historical identity — direct correspondence · superseded}\par
\href{../../../../continuum/strategy/capability/bcefb378424842fcb6d703208c3f8773.md}{Historical canonical definition and provenance}.\par
\section*{R.2.184 [CT3] Patient Health Information Disclosure Management — Historical identity [superseded]}
\phantomsection
\label{navigation:R.2.184}
\addcontentsline{toc}{section}{R.2.184 [CT3] Patient Health Information Disclosure Management — Historical identity [superseded]}

{\small Governed successor (direct correspondence): II.1.3.9.9.1 \hyperref[element:7f918639e23842e78fa92f1bf47c5a7b]{[CT3] Patient Health Information Disclosure Management}; Canonical ID: \texttt{\seqsplit{7f918639e23842e78fa92f1bf47c5a7b}}. Proposed · current.}\par
\label{element:07cf507f183d4905a7ad2adbb8555df5}
\index{[CT3] Patient Health Information Disclosure Management (R.2.184)}
{\small Capability · Superseded historical identity — direct correspondence · superseded}\par
\href{../../../../continuum/strategy/capability/07cf507f183d4905a7ad2adbb8555df5.md}{Historical canonical definition and provenance}.\par
\section*{R.2.185 [CT3] Patient Health Information Filter Management — Historical identity [superseded]}
\phantomsection
\label{navigation:R.2.185}
\addcontentsline{toc}{section}{R.2.185 [CT3] Patient Health Information Filter Management — Historical identity [superseded]}

{\small Governed successor (direct correspondence): II.1.3.9.9.2 \hyperref[element:62f673613c76494eba972cad3e4e8d0b]{[CT3] Patient Health Information Filter Management}; Canonical ID: \texttt{\seqsplit{62f673613c76494eba972cad3e4e8d0b}}. Proposed · current.}\par
\label{element:589f192797454d9e8960fb63774e5e9c}
\index{[CT3] Patient Health Information Filter Management (R.2.185)}
{\small Capability · Superseded historical identity — direct correspondence · superseded}\par
\href{../../../../continuum/strategy/capability/589f192797454d9e8960fb63774e5e9c.md}{Historical canonical definition and provenance}.\par
\section*{R.2.186 [CT3] Patient Health Information Restriction Management — Historical identity [superseded]}
\phantomsection
\label{navigation:R.2.186}
\addcontentsline{toc}{section}{R.2.186 [CT3] Patient Health Information Restriction Management — Historical identity [superseded]}

{\small Governed successor (direct correspondence): II.1.3.9.9.3 \hyperref[element:262366990bda4ac4ace5429269cf6a98]{[CT3] Patient Health Information Restriction Management}; Canonical ID: \texttt{\seqsplit{262366990bda4ac4ace5429269cf6a98}}. Proposed · current.}\par
\label{element:5ad91d045474435a96b560e66ff7de05}
\index{[CT3] Patient Health Information Restriction Management (R.2.186)}
{\small Capability · Superseded historical identity — direct correspondence · superseded}\par
\href{../../../../continuum/strategy/capability/5ad91d045474435a96b560e66ff7de05.md}{Historical canonical definition and provenance}.\par
\section*{R.2.187 [CT2] Health Information Control — Historical identity [superseded]}
\phantomsection
\label{navigation:R.2.187}
\addcontentsline{toc}{section}{R.2.187 [CT2] Health Information Control — Historical identity [superseded]}

{\small Governed successor (direct correspondence): II.1.3.9.9 \hyperref[element:44cc0e17ad534675bf26dce6e46a68ff]{[CT2] Health Information Control}; Canonical ID: \texttt{\seqsplit{44cc0e17ad534675bf26dce6e46a68ff}}. Canonical · current.}\par
\label{element:a3196aef21d443519de874c24f64a185}
\index{[CT2] Health Information Control (R.2.187)}
{\small Capability · Superseded historical identity — direct correspondence · superseded}\par
\href{../../../../continuum/strategy/capability/a3196aef21d443519de874c24f64a185.md}{Historical canonical definition and provenance}.\par
\section*{R.2.188 [CT3] HIE Endpoint State Management — Historical identity [superseded]}
\phantomsection
\label{navigation:R.2.188}
\addcontentsline{toc}{section}{R.2.188 [CT3] HIE Endpoint State Management — Historical identity [superseded]}

{\small Governed successor (evolved correspondence): II.1.3.9.10.1 \hyperref[element:1bec951a1c98450b8d4672c5c424f9bd]{[CT3] HIE Endpoint State Management}; Canonical ID: \texttt{\seqsplit{1bec951a1c98450b8d4672c5c424f9bd}}. Proposed · current.}\par
\label{element:697de08ac49a4b7f8df7a1ed6f04210e}
\index{[CT3] HIE Endpoint State Management (R.2.188)}
{\small Capability · Superseded historical identity — evolved correspondence · superseded}\par
\href{../../../../continuum/strategy/capability/697de08ac49a4b7f8df7a1ed6f04210e.md}{Historical canonical definition and provenance}.\par
\section*{R.2.189 [CT3] HIE Implementation Management — Historical identity [superseded]}
\phantomsection
\label{navigation:R.2.189}
\addcontentsline{toc}{section}{R.2.189 [CT3] HIE Implementation Management — Historical identity [superseded]}

{\small Governed successor (evolved correspondence): II.1.3.9.10.2 \hyperref[element:8ce29be373e04252b7aa881b6a72f375]{[CT3] HIE Implementation Management}; Canonical ID: \texttt{\seqsplit{8ce29be373e04252b7aa881b6a72f375}}. Proposed · current.}\par
\label{element:a99b4119339545a89dc54ebb216ddb85}
\index{[CT3] HIE Implementation Management (R.2.189)}
{\small Capability · Superseded historical identity — evolved correspondence · superseded}\par
\href{../../../../continuum/strategy/capability/a99b4119339545a89dc54ebb216ddb85.md}{Historical canonical definition and provenance}.\par
\section*{R.2.190 [CT3] HIE Requirement Management — Historical identity [superseded]}
\phantomsection
\label{navigation:R.2.190}
\addcontentsline{toc}{section}{R.2.190 [CT3] HIE Requirement Management — Historical identity [superseded]}

{\small Governed successor (evolved correspondence): II.1.3.9.10.3 \hyperref[element:1c852a1799984feba50011172983c865]{[CT3] HIE Requirement Management}; Canonical ID: \texttt{\seqsplit{1c852a1799984feba50011172983c865}}. Proposed · current.}\par
\label{element:a32c8bff0c284f5c9d6ab1f46991c232}
\index{[CT3] HIE Requirement Management (R.2.190)}
{\small Capability · Superseded historical identity — evolved correspondence · superseded}\par
\href{../../../../continuum/strategy/capability/a32c8bff0c284f5c9d6ab1f46991c232.md}{Historical canonical definition and provenance}.\par
\section*{R.2.191 [CT2] Health Information Exchange — Historical identity [superseded]}
\phantomsection
\label{navigation:R.2.191}
\addcontentsline{toc}{section}{R.2.191 [CT2] Health Information Exchange — Historical identity [superseded]}

{\small Governed successor (evolved correspondence): II.1.3.9.10 \hyperref[element:a61a15908ed544e6acf3e9cab63a0267]{[CT2] Health Information Exchange Enablement}; Canonical ID: \texttt{\seqsplit{a61a15908ed544e6acf3e9cab63a0267}}. Canonical · current.}\par
\label{element:330d1c43856a4322a2e9f331b1a5ed20}
\index{[CT2] Health Information Exchange (R.2.191)}
{\small Capability · Superseded historical identity — evolved correspondence · superseded}\par
\href{../../../../continuum/strategy/capability/330d1c43856a4322a2e9f331b1a5ed20.md}{Historical canonical definition and provenance}.\par
\section*{R.2.192 [CT3] Information Asset Definition — Historical identity [superseded]}
\phantomsection
\label{navigation:R.2.192}
\addcontentsline{toc}{section}{R.2.192 [CT3] Information Asset Definition — Historical identity [superseded]}

{\small Governed successor (direct correspondence): II.1.3.9.11.1 \hyperref[element:05c58b2ccfcc45e9a6c70d9a569c4f61]{[CT3] Information Asset Definition}; Canonical ID: \texttt{\seqsplit{05c58b2ccfcc45e9a6c70d9a569c4f61}}. Proposed · current.}\par
\label{element:166a22463d2c4534bda7ae31273398b6}
\index{[CT3] Information Asset Definition (R.2.192)}
{\small Capability · Superseded historical identity — direct correspondence · superseded}\par
\href{../../../../continuum/strategy/capability/166a22463d2c4534bda7ae31273398b6.md}{Historical canonical definition and provenance}.\par
\section*{R.2.193 [CT3] Information Asset Tag Management — Historical identity [superseded]}
\phantomsection
\label{navigation:R.2.193}
\addcontentsline{toc}{section}{R.2.193 [CT3] Information Asset Tag Management — Historical identity [superseded]}

{\small Governed successor (direct correspondence): II.1.3.9.11.2 \hyperref[element:df27c3ba6e164e56af7bff4e1eb1954a]{[CT3] Information Asset Tag Management}; Canonical ID: \texttt{\seqsplit{df27c3ba6e164e56af7bff4e1eb1954a}}. Proposed · current.}\par
\label{element:ae44a4ee78414d1098814315793bd751}
\index{[CT3] Information Asset Tag Management (R.2.193)}
{\small Capability · Superseded historical identity — direct correspondence · superseded}\par
\href{../../../../continuum/strategy/capability/ae44a4ee78414d1098814315793bd751.md}{Historical canonical definition and provenance}.\par
\section*{R.2.194 [CT2] Information Asset Stewardship Management — Historical identity [superseded]}
\phantomsection
\label{navigation:R.2.194}
\addcontentsline{toc}{section}{R.2.194 [CT2] Information Asset Stewardship Management — Historical identity [superseded]}

{\small Governed successor (direct correspondence): II.1.3.9.11 \hyperref[element:5bac15316d8d4f4db6e7a253f1acfbe3]{[CT2] Information Asset Stewardship Management}; Canonical ID: \texttt{\seqsplit{5bac15316d8d4f4db6e7a253f1acfbe3}}. Proposed · current.}\par
\label{element:cee085c690314802851de71f4741d8db}
\index{[CT2] Information Asset Stewardship Management (R.2.194)}
{\small Capability · Superseded historical identity — direct correspondence · superseded}\par
\href{../../../../continuum/strategy/capability/cee085c690314802851de71f4741d8db.md}{Historical canonical definition and provenance}.\par
\section*{R.2.195 [CT3] Data Provenance Capture — Historical identity [superseded]}
\phantomsection
\label{navigation:R.2.195}
\addcontentsline{toc}{section}{R.2.195 [CT3] Data Provenance Capture — Historical identity [superseded]}

{\small Governed successor (direct correspondence): II.1.3.9.12.1 \hyperref[element:44a7937520a54945b06b5fd4a2bbd870]{[CT3] Data Provenance Capture}; Canonical ID: \texttt{\seqsplit{44a7937520a54945b06b5fd4a2bbd870}}. Proposed · current.}\par
\label{element:61258c7fe41f4541871a67ec9d7185cf}
\index{[CT3] Data Provenance Capture (R.2.195)}
{\small Capability · Superseded historical identity — direct correspondence · superseded}\par
\href{../../../../continuum/strategy/capability/61258c7fe41f4541871a67ec9d7185cf.md}{Historical canonical definition and provenance}.\par
\section*{R.2.196 [CT3] Data Scientist Administration — Historical identity [superseded]}
\phantomsection
\label{navigation:R.2.196}
\addcontentsline{toc}{section}{R.2.196 [CT3] Data Scientist Administration — Historical identity [superseded]}

{\small Governed successor (direct correspondence): II.1.3.9.12.2 \hyperref[element:3bb1a11fba234717bf40c9e6452f401a]{[CT3] Data Scientist Administration}; Canonical ID: \texttt{\seqsplit{3bb1a11fba234717bf40c9e6452f401a}}. Proposed · current.}\par
\label{element:1cb497aec93c4b82ba7f5b622813e54e}
\index{[CT3] Data Scientist Administration (R.2.196)}
{\small Capability · Superseded historical identity — direct correspondence · superseded}\par
\href{../../../../continuum/strategy/capability/1cb497aec93c4b82ba7f5b622813e54e.md}{Historical canonical definition and provenance}.\par
\section*{R.2.197 [CT3] Information Product Administration — Historical identity [superseded]}
\phantomsection
\label{navigation:R.2.197}
\addcontentsline{toc}{section}{R.2.197 [CT3] Information Product Administration — Historical identity [superseded]}

{\small Governed successor (direct correspondence): II.1.3.9.12.3 \hyperref[element:b5a6a31855834ea58870ebb6c51ed4b0]{[CT3] Information Product Administration}; Canonical ID: \texttt{\seqsplit{b5a6a31855834ea58870ebb6c51ed4b0}}. Proposed · current.}\par
\label{element:5d9b7f8d2d7b41759b6d4931c1d2aaa2}
\index{[CT3] Information Product Administration (R.2.197)}
{\small Capability · Superseded historical identity — direct correspondence · superseded}\par
\href{../../../../continuum/strategy/capability/5d9b7f8d2d7b41759b6d4931c1d2aaa2.md}{Historical canonical definition and provenance}.\par
\section*{R.2.198 [CT2] Information Content Maintainer Management — Historical identity [superseded]}
\phantomsection
\label{navigation:R.2.198}
\addcontentsline{toc}{section}{R.2.198 [CT2] Information Content Maintainer Management — Historical identity [superseded]}

{\small Governed successor (direct correspondence): II.1.3.9.12 \hyperref[element:b81a1ab42af24c2e9896b8f3f78d291a]{[CT2] Information Content Maintainer Management}; Canonical ID: \texttt{\seqsplit{b81a1ab42af24c2e9896b8f3f78d291a}}. Proposed · current.}\par
\label{element:acb977e879204f1491293e5452e0ddf0}
\index{[CT2] Information Content Maintainer Management (R.2.198)}
{\small Capability · Superseded historical identity — direct correspondence · superseded}\par
\href{../../../../continuum/strategy/capability/acb977e879204f1491293e5452e0ddf0.md}{Historical canonical definition and provenance}.\par
\section*{R.2.199 [CT3] Enterprise Canonical/Reusable Model Design — Historical identity [superseded]}
\phantomsection
\label{navigation:R.2.199}
\addcontentsline{toc}{section}{R.2.199 [CT3] Enterprise Canonical/Reusable Model Design — Historical identity [superseded]}

{\small Governed successor (direct correspondence): II.1.3.9.13.1 \hyperref[element:90e9cd401ddf4c7690c978d8dd131413]{[CT3] Enterprise Canonical/Reusable Model Design}; Canonical ID: \texttt{\seqsplit{90e9cd401ddf4c7690c978d8dd131413}}. Proposed · current.}\par
\label{element:dbc37445ad944a85a7c4f26f410967b2}
\index{[CT3] Enterprise Canonical/Reusable Model Design (R.2.199)}
{\small Capability · Superseded historical identity — direct correspondence · superseded}\par
\href{../../../../continuum/strategy/capability/dbc37445ad944a85a7c4f26f410967b2.md}{Historical canonical definition and provenance}.\par
\section*{R.2.200 [CT3] Enterprise Canonical/Reusable Model Library Management — Historical identity [superseded]}
\phantomsection
\label{navigation:R.2.200}
\addcontentsline{toc}{section}{R.2.200 [CT3] Enterprise Canonical/Reusable Model Library Management — Historical identity [superseded]}

{\small Governed successor (direct correspondence): II.1.3.9.13.2 \hyperref[element:76943b4b68644c4dac1368fce8bb20d9]{[CT3] Enterprise Canonical/Reusable Model Library Management}; Canonical ID: \texttt{\seqsplit{76943b4b68644c4dac1368fce8bb20d9}}. Proposed · current.}\par
\label{element:f5eccd4f65df49bba98104e5abfa5e19}
\index{[CT3] Enterprise Canonical/Reusable Model Library Management (R.2.200)}
{\small Capability · Superseded historical identity — direct correspondence · superseded}\par
\href{../../../../continuum/strategy/capability/f5eccd4f65df49bba98104e5abfa5e19.md}{Historical canonical definition and provenance}.\par
\section*{R.2.201 [CT3] Enterprise Canonical/Reusable Model Utilisation Support — Historical identity [superseded]}
\phantomsection
\label{navigation:R.2.201}
\addcontentsline{toc}{section}{R.2.201 [CT3] Enterprise Canonical/Reusable Model Utilisation Support — Historical identity [superseded]}

{\small Governed successor (direct correspondence): II.1.3.9.13.3 \hyperref[element:5744cb9d86494064af87afb3cbf2cf97]{[CT3] Enterprise Canonical/Reusable Model Utilisation Support}; Canonical ID: \texttt{\seqsplit{5744cb9d86494064af87afb3cbf2cf97}}. Proposed · current.}\par
\label{element:7ae95677a00346f18a0ba01522a20153}
\index{[CT3] Enterprise Canonical/Reusable Model Utilisation Support (R.2.201)}
{\small Capability · Superseded historical identity — direct correspondence · superseded}\par
\href{../../../../continuum/strategy/capability/7ae95677a00346f18a0ba01522a20153.md}{Historical canonical definition and provenance}.\par
\section*{R.2.202 [CT3] Enterprise Reference Data Stewardship — Historical identity [superseded]}
\phantomsection
\label{navigation:R.2.202}
\addcontentsline{toc}{section}{R.2.202 [CT3] Enterprise Reference Data Stewardship — Historical identity [superseded]}

{\small Governed successor (direct correspondence): II.1.3.9.13.4 \hyperref[element:971d4b079ad24e89b70a8ed8dec0c8bf]{[CT3] Enterprise Reference Data Stewardship}; Canonical ID: \texttt{\seqsplit{971d4b079ad24e89b70a8ed8dec0c8bf}}. Proposed · current.}\par
\label{element:fd9f3239fdb64ea687825b66bb4d1f17}
\index{[CT3] Enterprise Reference Data Stewardship (R.2.202)}
{\small Capability · Superseded historical identity — direct correspondence · superseded}\par
\href{../../../../continuum/strategy/capability/fd9f3239fdb64ea687825b66bb4d1f17.md}{Historical canonical definition and provenance}.\par
\section*{R.2.203 [CT2] Information Design Governance — Historical identity [superseded]}
\phantomsection
\label{navigation:R.2.203}
\addcontentsline{toc}{section}{R.2.203 [CT2] Information Design Governance — Historical identity [superseded]}

{\small Governed successor (direct correspondence): II.1.3.9.13 \hyperref[element:f39ed9ebcc724857a4bf7a2b8e045c7a]{[CT2] Information Design Governance}; Canonical ID: \texttt{\seqsplit{f39ed9ebcc724857a4bf7a2b8e045c7a}}. Canonical · current.}\par
\label{element:0bdd7f8c0cd74d9fabc49801af037fc2}
\index{[CT2] Information Design Governance (R.2.203)}
{\small Capability · Superseded historical identity — direct correspondence · superseded}\par
\href{../../../../continuum/strategy/capability/0bdd7f8c0cd74d9fabc49801af037fc2.md}{Historical canonical definition and provenance}.\par
\section*{R.2.204 [CT3] Linked Information Curation — Historical identity [superseded]}
\phantomsection
\label{navigation:R.2.204}
\addcontentsline{toc}{section}{R.2.204 [CT3] Linked Information Curation — Historical identity [superseded]}

{\small Governed successor (direct correspondence): II.1.3.9.15.1 \hyperref[element:5edb2a8b5ee54ddda8104d1fab323a94]{[CT3] Linked Information Curation}; Canonical ID: \texttt{\seqsplit{5edb2a8b5ee54ddda8104d1fab323a94}}. Proposed · current.}\par
\label{element:9d0925bb39664a5985cacf3cea7442d1}
\index{[CT3] Linked Information Curation (R.2.204)}
{\small Capability · Superseded historical identity — direct correspondence · superseded}\par
\href{../../../../continuum/strategy/capability/9d0925bb39664a5985cacf3cea7442d1.md}{Historical canonical definition and provenance}.\par
\section*{R.2.205 [CT3] Longitudinal Clinical Record Administration — Historical identity [superseded]}
\phantomsection
\label{navigation:R.2.205}
\addcontentsline{toc}{section}{R.2.205 [CT3] Longitudinal Clinical Record Administration — Historical identity [superseded]}

{\small Governed successor (direct correspondence): II.1.3.9.15.2 \hyperref[element:14748dc6117f4b6ba7c0fc554b76bd89]{[CT3] Longitudinal Clinical Record Administration}; Canonical ID: \texttt{\seqsplit{14748dc6117f4b6ba7c0fc554b76bd89}}. Proposed · current.}\par
\label{element:bd91d848cf424c42bec7c80cd32edd28}
\index{[CT3] Longitudinal Clinical Record Administration (R.2.205)}
{\small Capability · Superseded historical identity — direct correspondence · superseded}\par
\href{../../../../continuum/strategy/capability/bd91d848cf424c42bec7c80cd32edd28.md}{Historical canonical definition and provenance}.\par
\section*{R.2.206 [CT2] Patient Clinical Record — Historical identity [superseded]}
\phantomsection
\label{navigation:R.2.206}
\addcontentsline{toc}{section}{R.2.206 [CT2] Patient Clinical Record — Historical identity [superseded]}

{\small Governed successor (direct correspondence): II.1.3.9.15 \hyperref[element:5478878e0e6e4d34a8d5023ac6ac2d54]{[CT2] Patient Clinical Record}; Canonical ID: \texttt{\seqsplit{5478878e0e6e4d34a8d5023ac6ac2d54}}. Canonical · current.}\par
\label{element:2e4de521bda24ad7b0ea64ae535b7f84}
\index{[CT2] Patient Clinical Record (R.2.206)}
{\small Capability · Superseded historical identity — direct correspondence · superseded}\par
\href{../../../../continuum/strategy/capability/2e4de521bda24ad7b0ea64ae535b7f84.md}{Historical canonical definition and provenance}.\par
\section*{R.2.207 [CT3] Contextual Information Presentation — Historical identity [superseded]}
\phantomsection
\label{navigation:R.2.207}
\addcontentsline{toc}{section}{R.2.207 [CT3] Contextual Information Presentation — Historical identity [superseded]}

{\small Governed successor (direct correspondence): II.1.3.9.16.1 \hyperref[element:9f900c53ae56488cad69e53c2205e898]{[CT3] Contextual Information Presentation}; Canonical ID: \texttt{\seqsplit{9f900c53ae56488cad69e53c2205e898}}. Proposed · current.}\par
\label{element:e349d5f6a8774bbe8d150f9c51af770c}
\index{[CT3] Contextual Information Presentation (R.2.207)}
{\small Capability · Superseded historical identity — direct correspondence · superseded}\par
\href{../../../../continuum/strategy/capability/e349d5f6a8774bbe8d150f9c51af770c.md}{Historical canonical definition and provenance}.\par
\section*{R.2.208 [CT2] Presentation Services — Historical identity [superseded]}
\phantomsection
\label{navigation:R.2.208}
\addcontentsline{toc}{section}{R.2.208 [CT2] Presentation Services — Historical identity [superseded]}

{\small Governed successor (direct correspondence): II.1.3.9.16 \hyperref[element:e9b6ba1917ba44cca7466444d0836c98]{[CT2] Presentation Services}; Canonical ID: \texttt{\seqsplit{e9b6ba1917ba44cca7466444d0836c98}}. Canonical · current.}\par
\label{element:c89e7c8b277247d9965fdb534aad2d1f}
\index{[CT2] Presentation Services (R.2.208)}
{\small Capability · Superseded historical identity — direct correspondence · superseded}\par
\href{../../../../continuum/strategy/capability/c89e7c8b277247d9965fdb534aad2d1f.md}{Historical canonical definition and provenance}.\par
\section*{R.2.209 [CT3] Activity Routing Management — Historical identity [superseded]}
\phantomsection
\label{navigation:R.2.209}
\addcontentsline{toc}{section}{R.2.209 [CT3] Activity Routing Management — Historical identity [superseded]}

{\small Governed successor (direct correspondence): II.1.3.9.18.1 \hyperref[element:27411c143e6e41ffba1af139fbbcd7e2]{[CT3] Activity Routing Management}; Canonical ID: \texttt{\seqsplit{27411c143e6e41ffba1af139fbbcd7e2}}. Proposed · current.}\par
\label{element:243ea93e7b7448fe97dc74f90bfc8946}
\index{[CT3] Activity Routing Management (R.2.209)}
{\small Capability · Superseded historical identity — direct correspondence · superseded}\par
\href{../../../../continuum/strategy/capability/243ea93e7b7448fe97dc74f90bfc8946.md}{Historical canonical definition and provenance}.\par
\section*{R.2.210 [CT3] External Task Coordination — Historical identity [superseded]}
\phantomsection
\label{navigation:R.2.210}
\addcontentsline{toc}{section}{R.2.210 [CT3] External Task Coordination — Historical identity [superseded]}

{\small Governed successor (direct correspondence): II.1.3.9.18.2 \hyperref[element:ff3447c929724245a13ee864874b5561]{[CT3] External Task Coordination}; Canonical ID: \texttt{\seqsplit{ff3447c929724245a13ee864874b5561}}. Proposed · current.}\par
\label{element:32be309d65594592bac4d7b068def0d2}
\index{[CT3] External Task Coordination (R.2.210)}
{\small Capability · Superseded historical identity — direct correspondence · superseded}\par
\href{../../../../continuum/strategy/capability/32be309d65594592bac4d7b068def0d2.md}{Historical canonical definition and provenance}.\par
\section*{R.2.211 [CT3] External Task Status \& Traceability Management — Historical identity [superseded]}
\phantomsection
\label{navigation:R.2.211}
\addcontentsline{toc}{section}{R.2.211 [CT3] External Task Status \& Traceability Management — Historical identity [superseded]}

{\small Governed successor (direct correspondence): II.1.3.9.18.3 \hyperref[element:b25bb04a8aa94dcfac105a646878ccfc]{[CT3] External Task Status \& Traceability Management}; Canonical ID: \texttt{\seqsplit{b25bb04a8aa94dcfac105a646878ccfc}}. Proposed · current.}\par
\label{element:95985d244a234571892a9e9b4024f594}
\index{[CT3] External Task Status \& Traceability Management (R.2.211)}
{\small Capability · Superseded historical identity — direct correspondence · superseded}\par
\href{../../../../continuum/strategy/capability/95985d244a234571892a9e9b4024f594.md}{Historical canonical definition and provenance}.\par
\section*{R.2.212 [CT3] Work Boundary Governance — Historical identity [superseded]}
\phantomsection
\label{navigation:R.2.212}
\addcontentsline{toc}{section}{R.2.212 [CT3] Work Boundary Governance — Historical identity [superseded]}

{\small Governed successor (direct correspondence): II.1.3.9.18.4 \hyperref[element:35b66a0b38db4568849ecb842ffdcbf3]{[CT3] Work Boundary Governance}; Canonical ID: \texttt{\seqsplit{35b66a0b38db4568849ecb842ffdcbf3}}. Proposed · current.}\par
\label{element:b4d14abf930a4230a7f5b93e9ce3df91}
\index{[CT3] Work Boundary Governance (R.2.212)}
{\small Capability · Superseded historical identity — direct correspondence · superseded}\par
\href{../../../../continuum/strategy/capability/b4d14abf930a4230a7f5b93e9ce3df91.md}{Historical canonical definition and provenance}.\par
\section*{R.2.213 [CT2] Workflow \& Activity Coordination — Historical identity [superseded]}
\phantomsection
\label{navigation:R.2.213}
\addcontentsline{toc}{section}{R.2.213 [CT2] Workflow \& Activity Coordination — Historical identity [superseded]}

{\small Governed successor (direct correspondence): II.1.3.9.18 \hyperref[element:c724f3d38bfa4567983a11ad22095d61]{[CT2] Workflow \& Activity Coordination}; Canonical ID: \texttt{\seqsplit{c724f3d38bfa4567983a11ad22095d61}}. Proposed · current.}\par
\label{element:4422d865978c4353ad575759c6026d27}
\index{[CT2] Workflow \& Activity Coordination (R.2.213)}
{\small Capability · Superseded historical identity — direct correspondence · superseded}\par
\href{../../../../continuum/strategy/capability/4422d865978c4353ad575759c6026d27.md}{Historical canonical definition and provenance}.\par
\section*{R.2.214 [CT1] Shared Information \& Coordination Enablement — Historical identity [superseded]}
\phantomsection
\label{navigation:R.2.214}
\addcontentsline{toc}{section}{R.2.214 [CT1] Shared Information \& Coordination Enablement — Historical identity [superseded]}

{\small Governed successor (direct correspondence): II.1.3.9 \hyperref[element:0062fcacae1740988b6be2bda2aa1674]{[CT1] Shared Information \& Coordination Enablement}; Canonical ID: \texttt{\seqsplit{0062fcacae1740988b6be2bda2aa1674}}. Proposed · current.}\par
\label{element:eeabb9c3c526404cbe4302906560fa7e}
\index{[CT1] Shared Information \& Coordination Enablement (R.2.214)}
{\small Capability · Superseded historical identity — direct correspondence · superseded}\par
\href{../../../../continuum/strategy/capability/eeabb9c3c526404cbe4302906560fa7e.md}{Historical canonical definition and provenance}.\par

\clearpage
\phantomsection
\addcontentsline{toc}{chapter}{Architectural Name Index}
\renewcommand{\indexname}{Architectural Name Index}
\printindex
\end{document}
